\documentclass[11pt]{article}
\usepackage[T1]{fontenc}
\usepackage{lmodern,amsmath,amssymb,amsthm,mathtools}
\usepackage[margin=1in]{geometry}
\usepackage{booktabs,longtable,array,microtype}
\usepackage[hidelinks]{hyperref}
\usepackage{xurl}
\newtheorem{theorem}{Theorem}[section]
\newtheorem{proposition}[theorem]{Proposition}
\newtheorem{lemma}[theorem]{Lemma}
\newtheorem{corollary}[theorem]{Corollary}
\theoremstyle{remark}
\newtheorem{remark}[theorem]{Remark}
\DeclareMathOperator{\Tr}{Tr}
\DeclareMathOperator{\Var}{Var}
\DeclareMathOperator{\Cov}{Cov}
\DeclareMathOperator{\Ad}{Ad}
\newcommand{\E}{\mathbb E}
\newcommand{\SU}{\mathrm{SU}}
\newcommand{\R}{\mathbb R}
\newcommand{\dd}{\,d}
\title{Renewal Yang--Mills Mass-Gap Research Notes\\
Consolidated Restart after f16.1 and f16.2\\
Version V100: Resonances and Local Averaging}
\author{Research continuation and mathematical audit}
\date{19 September 2026}
\begin{document}
\maketitle
\begin{abstract}
This compact restart consolidates the f15 V100 archive and the two
subsequent branches, f16.1 V85 and f16.2 V89.  The interacting
four-dimensional continuum Yang--Mills theory and its mass gap remain
unconstructed in these notes.  We do not promote numerical tests,
reference-field calculations, or fixed-volume gaps to that conclusion.

The new calculation goes upstream of the two branches' current
obstructions: integrate the temporal links of the original Wilson
kernel before organizing its nonlinear corrections.  For a fixed
connected spatial graph and bounded rescaled endpoints, we prove a
complete two-order expansion, with the exact Gaussian mean, covariance,
Haar factor, cubic return, and residual root integral.  At flat
endpoints its first correction is evaluated explicitly in the grounded
graph Green function.  Trees and cycles give independent exact checks;
a noncommuting one-edge check detects an omitted cubic return.
We also correct the finite-time thermodynamic identity and record a
published two-loop benchmark.  None of the new remainder bounds is
uniform in growing spatial volume or in an unrestricted spatial history.
\end{abstract}
\tableofcontents

\section{Context and the authoritative result ledger}
\label{restart:context}

The objective is an interacting four-dimensional continuum Yang--Mills
construction with a strictly positive spectral gap in its physical
Hamiltonian.  The finite-lattice calculations below use \(\SU(2)\);
they do not prove the analogous construction for arbitrary compact
simple gauge groups.  Spatial flow time, lattice transfer time and
physical Euclidean time remain distinct.

The review began from the three requested archives, none edited by this work:
\begin{itemize}
\item \(A=\) \path{Renewal_Yang_Mills_Mass_Gap_Research_Notes_V100_f15.tex};
\item \(B=\) \path{Renewal_Yang_Mills_Mass_Gap_Research_Notes_V85_f16.1.tex};
\item \(C=\) \path{Renewal_Yang_Mills_Mass_Gap_Research_Notes_V89_f16.2.tex}.
\end{itemize}
References such as \(B\), V84 and \(C\), V84 designate different
sections.  The two branches have an identical mathematical body from
the opening context through V48, verified by direct text comparison
up to their V49 append markers (with surrounding whitespace ignored).
They diverge at V49; \(C\)'s higher version number does not supersede
\(B\)'s later perturbative work.

During verification a parallel update replaced the V85 filename by
\path{Renewal_Yang_Mills_Mass_Gap_Research_Notes_V86_f16.1.tex}.
Its new terminal section was also read.  It evaluates electric-sector
and partial orbit-measure corrections but explicitly leaves the
two-loop accounting unresolved.  Its numerical additions are not
imported as certified coefficients here.  Removing its V86 appendix
and reversing its title changes reconstructs the exact original V85
SHA--256 hash.  We preserve that external update rather than restore
its superseded filename; references to \(B\), V85 still mean the
verified predecessor body.

The review compared the section maps, the common body, the terminal
proofs and status ledgers, and the calculations feeding the present
bottleneck.  It is not an independent line-by-line recertification of
every historical proof.  Inherited claims below keep their original
hypotheses; the new results used in this restart have self-contained
proofs in Sections~\ref{restart:trace}--\ref{restart:tower}.

\subsection{What is retained, and with what scope}

\begin{longtable}{@{}p{.18\textwidth}p{.76\textwidth}@{}}
\toprule
Location & Retained result or interface, with its limitation\\
\midrule
\endhead
\(A\), V92--V100 &
Complete spatial assembly, genuine compact Gram projection, magnetic
carrier response, and the excitation-count calculation.  In particular
V100's normalized Gaussian image requires an extensive old-oscillator
degree budget; a subextensive budget loses its mass.  Those oscillator
degrees are not physical energies.  Fixed-regulator comparison does not
supply volume-uniform transfer control.\\[1ex]
Common V1--V48 &
Conditional-affinity kinetic compression, coherent fast-history
estimates, compact normalization, gauge-root bookkeeping, and
fixed-volume soft/hard spectral reductions.  The temporal-memory,
boundary, and small-volume hypotheses cannot be dropped when exporting
these results.  No bare finite-volume soft gap is a continuum mass gap.\\[1ex]
\(B\), V49--V73 &
Fixed-regulator spectral refinements, followed by the volume-uniform
reduction called N2 and perturbative running-coupling calculations.
The unresolved connected remainder and Perron-vector control are
substantive analytic inputs, not consequences of the fitted
one-loop coefficient.\\[1ex]
\(B\), V74--V85 &
A proposed two-loop chain calculation, its one-loop contraction
correction, and an independent Monte Carlo discrepancy.
V85 explicitly marks the two-loop vertex list as incomplete.
Its values \(p_2(8)\simeq0.0934\) and the extrapolation \(0.0951\)
are not retained as Wilson-theory results.\\[1ex]
\(C\), V49--V76 &
Volume/clock comparisons and the local-source route, including
obstructions to microscopic-ribbon and rough-connection shortcuts.
A source innovation at the lattice scale need not survive a fixed
physical time separation.\\[1ex]
\(C\), V77--V87 &
Derivative-free reconstruction of the \emph{original} simultaneous-flow
source; exact native finite-regulator source estimates; forward-source
completion; a missing parity sector and a separating invariant algebra.
Finite-regulator density and contact floors do not prove a limiting
positive-time Gram or uniform contraction on a dense continuum core.\\[1ex]
\(C\), V88--V89 &
A nonlinear Gaussian-reference jet calculation, common-log cancellation
in first normalized coefficients, and exact reflection-positive input
examples with identical one-slice laws and linear time covariances but
different composite correlations.  The latter examples are not Wilson
histories.  The full nonlinear native source and its joint input
correction remain unmatched.\\
\bottomrule
\end{longtable}

The strongest common lesson is that the complete joint law, source,
normalization and physical clock must be controlled together.
We therefore do not append another Gaussian-source calculation while
silently retaining the incomplete electric sector of \(B\).

\subsection{Two audit corrections and a literature benchmark}
\label{restart:benchmark}

First, \(B\), V85's concluding verifier description says that agreement
with the chain is among the checked assertions.  The actual saved
diagnostic says otherwise:
\[
 \texttt{all\_checks=true},\qquad
 \texttt{chain\_within\_5sigma=false}.
\]
Its three asserted checks concern fit quality, precision and the leading
term; the chain comparison is recorded, not asserted.  The reported fit
is \(0.1471728\pm0.0041347\), against \(0.09342554\), with a reported
difference of \(12.999\) standard errors.  We have inspected the script
and its existing report, not rerun the Monte Carlo.  Its sampling and
finite-volume systematics are not rigorous error enclosures.  The
discrepancy invalidates export of the old prediction; by itself it does
not prove which missing diagram is responsible, nor validate all
earlier off-shell one-loop calculations.

Second, the literature comparison left open in \(B\) is available.
Athenodorou, Panagopoulos and Tsapalis
\cite[Eq.~(14), with \(N_f=0\)]{APT} report, for
\(P=\langle1-N_c^{-1}\Tr U_p\rangle\),
\[
 P=c_1^Gg_0^2+c_2^Gg_0^4+\cdots,\quad
 c_1^G=\frac{N_c^2-1}{8N_c},\quad
 c_2^G=(N_c^2-1)
       \left(0.0051069297-\frac1{128N_c^2}\right).
\]
For the present Wilson convention \(\gamma=4/g_0^2\), \(N_c=2\), this gives
\begin{equation}
 P=\frac3{4\gamma}+\frac{p_2^{\rm lit}}{\gamma^2}+\cdots,\qquad
 p_2^{\rm lit}=48(0.0051069297-1/512)
                \simeq0.1513826256 .
 \label{restart:p2}
\end{equation}
This is a published infinite-volume perturbative benchmark, not a
new proof here, an exact decimal, or a certified finite-\(8^3\times16\)
coefficient.  It is numerically close to the archived Monte Carlo
intercept, but that comparison does not eliminate either finite-size or
higher-order effects.  The plaquette-scheme shift would accordingly be
\(\gamma_P-\gamma=-4p_2^{\rm lit}/3+O(\gamma^{-1})
\simeq-0.2018435+O(\gamma^{-1})\), not the old \(-0.127\).
A single plaquette expectation is still not enough to match a
two-time composite-source law.

\section{Finite Euclidean time is not the vacuum limit}
\label{restart:trace}

Let the spatial periodic lattice have \(V\) sites and let \(N\) denote
its temporal period.  Take nondegenerate periods so the usual Wilson
plaquettes have their ordinary incidence.  Use normalized Haar measure
and the unnormalized Wilson transfer \(\mathcal T_\gamma\) on the
physical, gauge-invariant Hilbert space.  Its normalization is fixed by
\[
 Z_{N,V}(\gamma)=\Tr\mathcal T_\gamma^N
 =\int \exp\{-\gamma S(U)\}\prod_\ell dH(U_\ell),\qquad
 S(U)=\sum_p(1-\tfrac12\Tr U_p).
\]
No \(\gamma\)-dependent scalar is removed from this definition.
Let \(\overline P_{N,V}\) be the average over all \(6NV\) plaquettes.
At finite anisotropic aspect ratio, a specified spatial plaquette
need not equal this orientation-averaged observable.

\begin{proposition}[Correct finite-time identity]
\label{restart:finite-identity}
Exactly,
\begin{equation}
 6V\,\overline P_{N,V}
 =-\frac1N\partial_\gamma\log Z_{N,V}
 =\partial_\gamma E_{N,V},\qquad
 E_{N,V}=-\frac1N\log\Tr\mathcal T_\gamma^N.
 \label{restart:finite-identity-eq}
\end{equation}
Writing \(\lambda_0>\lambda_1\ge\lambda_2\ge\cdots\ge0\) for the
fixed-regulator transfer spectrum, counted with multiplicity, gives
\begin{equation}
 E_{N,V}=E_{0,V}-\frac1N\log\left(1+
                 \sum_{j\ge1}(\lambda_j/\lambda_0)^N\right),
 \qquad E_{0,V}=-\log\lambda_0.
 \label{restart:thermal}
\end{equation}
Consequently \(\partial_\gamma E_{0,V}\) replaces the right side of
\eqref{restart:finite-identity-eq} only after a justified \(N\to\infty\)
limit.  At fixed \(V,\gamma>0\) that limit holds.  No uniform rate along a
continuum or thermodynamic trajectory is supplied here.
\end{proposition}
\begin{proof}
Compactness and boundedness of \(S\) justify differentiation of the
finite integral, yielding \(\partial_\gamma\log Z=-\E S\).
The spectral expansion gives \eqref{restart:thermal}.
For completeness, the fixed-regulator kernel is smooth and strictly
positive, and is a positive semidefinite compact transfer.
The top eigenvalue is simple and isolated; these properties also follow
from its positive character expansion and positivity-improving kernel.
Its derivative \(\mathcal T_\gamma'\) is bounded.  Cyclic differentiation
of the trace gives, for sufficiently large integer \(N\),
\[
 \partial_\gamma E_{N,V}
 =-\frac{\Tr(\mathcal T_\gamma'\mathcal T_\gamma^{N-1})}
         {\Tr\mathcal T_\gamma^N}.
\]
Let \(r=\lambda_1/\lambda_0<1\) and
\(S_2=\sum_{j\ge1}(\lambda_j/\lambda_0)^2<\infty\), using the
Hilbert--Schmidt property.  The contributions of the nonleading
eigenvectors in the numerator and denominator are bounded, after their
leading powers of \(\lambda_0\) are removed, by
\(\|\mathcal T_\gamma'\|S_2r^{N-3}\) and \(S_2r^{N-2}\).
The leading numerator is
\(\langle\Omega,\mathcal T_\gamma'\Omega\rangle=\lambda_0'\).
Thus the ratio tends to \(-\lambda_0'/\lambda_0\).
All constants, including \(1-r\), may depend on \(V,\gamma\).
\end{proof}

This corrects the literal finite-\(N\) statement in \(B\), V84.
It does not identify the missing electric coefficient.  Likewise,
differentiating an asymptotic energy expansion requires a bound on its
\emph{differentiated} remainder: an \(O(\gamma^{-2})\) energy error alone
does not imply an \(O(\gamma^{-3})\) plaquette error.

\section{The exact temporal-link integral and the residual root}
\label{restart:kernel}

Let \(G=(\mathcal V,\mathcal E)\) be a fixed finite connected loopless
graph, with one orientation assigned to each edge, allowing
multiplicities.  Put \(n=|\mathcal V|\), \(m=|\mathcal E|\), fix a root
\(o\), and write \(s(e),t(e)\) for the endpoints.  All Haar measures have
total mass one.  For a unit quaternion \(q\), let
\(w(q)=\frac12\Tr q\).  For two spatial configurations \(U,V\), define
\begin{equation}
 K_\gamma(U,V)=
 \int_{\SU(2)^{\mathcal V}}
 \exp\left[-\gamma\sum_{e\in\mathcal E}
 \{1-w(U_ea_{t(e)}V_e^{-1}a_{s(e)}^{-1})\}\right]
 \prod_x dH(a_x).
 \label{restart:K}
\end{equation}
Each \(a_x\) is an actual temporal link.  The integrand consists exactly
of the temporal Wilson plaquettes of one slab, with all spatial links
on its ends held fixed.  The magnetic spatial factors belong outside
this integral.

Let \(K_\gamma^o(U,V)\) be the same integral with \(a_o=1\), integrating
only the other vertices.  Fixing that root is not permission to remove
its physical integration.

\begin{lemma}[Exact root completion]
\label{restart:root}
For every pair \(U,V\),
\begin{equation}
 K_\gamma(U,V)=\int_{\SU(2)}
                 K_\gamma^o(U,\Ad_R V)\,dH(R),
 \qquad (\Ad_R V)_e=RV_eR^{-1}.
 \label{restart:root-eq}
\end{equation}
On a finite periodic history of spatial slices \(U_0,\ldots,U_{N-1}\),
the original Wilson density, after integrating all temporal links, is
proportional to
\begin{equation}
 \prod_{t=0}^{N-1}
 e^{-\gamma S_{\rm sp}(U_t)}K_\gamma(U_t,U_{t+1})
 \prod_t dH(U_t),\qquad U_N=U_0.
 \label{restart:history}
\end{equation}
Every bounded observable depending only on these spatial slices,
including the simultaneous spatial-flow source at finite regulator,
has unchanged expectations under this marginal.
\end{lemma}
\begin{proof}
Set \(R=a_o\) and \(a_x=b_xR\).  Then \(b_o=1\), product Haar measure
factorizes, and
\[
 U_ea_tV_e^{-1}a_s^{-1}
       =U_eb_t(RV_eR^{-1})^{-1}b_s^{-1}.
\]
This proves \eqref{restart:root-eq}.  Conditional on all spatial slices,
the temporal variables in different slabs occur in disjoint sets of
temporal plaquettes.  Fubini's theorem therefore yields
\eqref{restart:history}, with the same partition function.  Integrating
a spatial-slice observable uses the identical identity.
\end{proof}

Thus the kernel
\(e^{-\gamma S_{\rm sp}(U)/2}K_\gamma(U,V)
 e^{-\gamma S_{\rm sp}(V)/2}\) retains both the electric sector and the
Gauss projection.  No additional orbit determinant is guessed in these
compact variables.  A later change of spatial coordinates must of course
carry its own exact measure.

\section{Complete fixed-graph expansion at bounded endpoints}
\label{restart:expansion}

Use the quaternion convention
\(\exp x=\cos|x|+\sin|x|\,x/|x|\) for \(x\in\R^3\).
Set \(g=\gamma^{-1/2}\).  This \(g\) equals \(g_0/2\) in
\eqref{restart:p2}; it is not the identically named amplitude parameter
in \(C\), V88--V89.
Write
\[
 U_e=\exp(gX_e),\quad V_e=\exp(gY_e),\quad a_x=\exp(gA_x),\quad A_o=0.
\]
The endpoint vectors range over a fixed compact subset of
\((\R^3)^{2m}\).

Let \(D:\R^{n-1}\to\R^m\) be the rooted incidence matrix:
\((DA)_e=A_{t(e)}-A_{s(e)}\).
Set
\begin{equation}
 L=D^{\mathsf T}D,\quad C=L^{-1},\quad
 \Pi=I-DCD^{\mathsf T},\quad b=X-Y,\quad
 \mu=-CD^{\mathsf T}b,\quad d=3(n-1).
 \label{restart:gaussian-data}
\end{equation}
Matrices act separately on each of the three color components.
Connectedness gives \(L>0\); \(\Pi\) is the orthogonal projection onto
the cycle space.  Let \(\E_{X,Y}\) refer to the Gaussian law of \(A\)
with mean \(\mu\) and covariance \(C\otimes I_3\).

For each edge form the ordered list
\[
 (u_1,u_2,u_3,u_4)=(X_e,A_{t(e)},-Y_e,-A_{s(e)}).
\]
Define \(S_3,S_4\) as the coefficients specified by
\begin{equation}
 \sum_e\{1-w(e^{gu_1}e^{gu_2}e^{gu_3}e^{gu_4})\}
 =\frac{g^2}{2}\|b+DA\|^2+g^3S_3+g^4S_4+O(g^5).
 \label{restart:vertices}
\end{equation}
These are explicit finite polynomials, not fitted tensors.
For example the edge contribution to \(S_3\) is
\begin{equation}
 s_3(u)=\sum_{i<j<k}\det(u_i,u_j,u_k).
 \label{restart:cubic}
\end{equation}
For clarity the edge contribution to \(S_4\) is \(-w_4(u)\), where
\begin{align}
 w_4(u)={}&\sum_i\frac{|u_i|^4}{24}
       +\sum_{i<j}\frac{|u_i|^2|u_j|^2}{4}\notag\\
 &+\sum_{i<j}(u_i\cdot u_j)
   \left\{\frac{|u_i|^2+|u_j|^2}{6}
                   +\frac12\sum_{k\ne i,j}|u_k|^2\right\}\notag\\
 &+(u_1\cdot u_2)(u_3\cdot u_4)
  -(u_1\cdot u_3)(u_2\cdot u_4)
  +(u_1\cdot u_4)(u_2\cdot u_3).
 \label{restart:quartic}
\end{align}
The signs follow from quaternion multiplication, in particular
\(w(xyz)=-(x\times y)\cdot z\) for imaginary quaternions.
Finally put \(H(A)=\frac13\sum_{x\ne o}|A_x|^2\) and
\[
 \kappa_G=\frac{(2\pi)^{d/2}}{(2\pi^2)^{n-1}(\det L)^{3/2}}.
\]

\begin{theorem}[Complete rooted electric coefficient]
\label{restart:main}
For every fixed \(G\), uniformly on any fixed compact endpoint set,
\begin{align}
 \log K_{g^{-2}}^o(e^{gX},e^{gY})
 ={}&d\log g+\log\kappa_G-\tfrac12\|\Pi(X-Y)\|^2\notag\\
 &+g\,k_1(X,Y)+g^2k_2(X,Y)+O_{G,M}(g^3),
 \label{restart:logK}\\
 k_1={}&-\E_{X,Y}S_3,\notag\\
 k_2={}&-\E_{X,Y}S_4-\E_{X,Y}H
                         +\tfrac12\Var_{X,Y}(S_3).
 \label{restart:k2}
\end{align}
Here \(M\) bounds the endpoints; there is no restriction of the original
compact integral to an imposed small-field event.  Its complement is
estimated in the proof.  The constants are not asserted uniform in \(G\).
\end{theorem}
\begin{proof}
At \(U=V=1\) the compact rooted action vanishes precisely when
\(a_t=a_s\) on every edge; connectedness and \(a_o=1\) give the unique
zero \(a_x=1\).  Its Hessian in root coordinates is \(L\otimes I_3>0\).
Outside any sufficiently small fixed product coordinate neighborhood,
compactness gives a positive action minimum.  Since
\(e^{gX},e^{gY}\to1\) uniformly on the chosen endpoint set, that complement
contributes \(O_{G,M}(e^{-c_G/g^2})\).

Inside this neighborhood, normalized Haar measure is exactly
\[
 dH(e^{gA_x})=\frac{g^3}{2\pi^2}
             \left(\frac{\sin(g|A_x|)}{g|A_x|}\right)^2 d^3A_x.
\]
The factor \(1/(2\pi^2)\) is essential.  Its logarithm beyond the
constant prefactor is \(-g^2H(A)+O(g^4\sum_x|A_x|^4)\).
Expand the actual plaquette words to get
\eqref{restart:vertices}--\eqref{restart:quartic}.  Completing the square,
\[
 \tfrac12\|b+DA\|^2
 =\tfrac12\|\Pi b\|^2+
             \tfrac12(A-\mu)^{\mathsf T}(L\otimes I_3)(A-\mu).
\]
The resulting Gaussian integral is exactly the leading factor in
\eqref{restart:logK}.

Here are remainder details to distinguish the argument from a merely
formal cumulant series.  Insert a smooth cutoff supported in the fixed
coordinate neighborhood and equal to one on a smaller neighborhood.
Uniform positive definiteness of the unscaled action Hessian there,
together with its \(O(g)\) gradient at zero, bounds the rescaled action
below by \(c_G|A|^2-C_{G,M}\).  The Taylor formula with integral remainder
for the scalar plaquette words shows that derivatives through order
three in \(g\) of the rescaled action are bounded by fixed polynomials
in \(1+|A|\), on this support.  The same holds for the Haar factors and
the cutoff derivatives.  Thus each differentiated integrand is bounded
by a fixed polynomial times \(e^{-c_G|A|^2+C_{G,M}}\), which is integrable.
The part where cutoff derivatives occur has \(|A|\) of order \(g^{-1}\)
and is exponentially small as well.  Taylor integration consequently
has an \(O_{G,M}(g^3)\) error after division by the leading \(g^d\)
factor.  This justifies integrating its first two coefficients over the
entire Gaussian space, without retaining a moving integration boundary.

The divided integral is therefore
\[
 1-g\,\E S_3+
 g^2\left\{-\E(S_4+H)+\tfrac12\E S_3^2\right\}
 +O_{G,M}(g^3).
\]
Its leading term is one uniformly on the compact endpoint set.
Taking its logarithm replaces the last raw second moment by the
variance and proves the theorem.
\end{proof}

The conditional covariance in \eqref{restart:k2} is not an optional
extra electric diagram.  Once temporal links have been integrated using
this formula, their Haar and cubic-return contributions must not also
be added as independent vertices to the same spatial marginal.

\begin{corollary}[The root has its own normalized return]
\label{restart:root-expansion}
For fixed endpoints put \(Y^R=\Ad_RY\),
\(q(R)=\frac12\|\Pi(X-Y^R)\|^2\), and
\[
 Z_{\rm root}=\int e^{-q(R)}\,dH(R),\qquad
 d\nu_{\rm root}(R)=Z_{\rm root}^{-1}e^{-q(R)}dH(R).
\]
Then the full kernel has the expansion
\begin{align}
 \log K_{g^{-2}}(e^{gX},e^{gY})
 ={}&d\log g+\log\kappa_G+\log Z_{\rm root}
       +g\,\E_{\nu_{\rm root}}k_1(X,Y^R)\notag\\
 &+g^2\left\{\E_{\nu_{\rm root}}k_2(X,Y^R)
          +\tfrac12\Var_{\nu_{\rm root}}(k_1(X,Y^R))\right\}
       +O_{G,M}(g^3).
 \label{restart:root-expansion-eq}
\end{align}
\end{corollary}
\begin{proof}
The adjoint action preserves the endpoint bound.  Apply the uniform
expansion before the exact root integral \eqref{restart:root-eq}.
Expand the exponential and then its logarithm under the positive
normalized weight \(\nu_{\rm root}\).  Since \(q\) is uniformly bounded
on the declared endpoint set, \(Z_{\rm root}\) is bounded away from zero
there.  This proves the remainder statement and the extra variance.
\end{proof}

\section{An evaluated correction at flat endpoints}
\label{restart:flat}

Extend the rooted covariance to an \(n\times n\) matrix by
\(C_{oo}=C_{ox}=C_{xo}=0\).  For an edge \(e=(s,t)\) set
\[
 r_e=C_{ss}+C_{tt}-2C_{st}.
\]
It is the unit-conductance effective resistance between the endpoints.
This interpretation is not needed for the proof.

\begin{theorem}[Explicit graph formula for the electric normalization]
\label{restart:flat-theorem}
At \(U=V=1\), the full and rooted integrals agree and
\begin{equation}
 K_\gamma(1,1)=K_\gamma^o(1,1)
 =\kappa_G\gamma^{-3(n-1)/2}
       \left\{1+\frac{a_G}{\gamma}+O_G(\gamma^{-2})\right\},
 \label{restart:flat-asymp}
\end{equation}
where
\begin{equation}
 \boxed{\displaystyle
 a_G=\sum_{e=(s,t)}
       \left\{\frac58r_e^2+C_{ss}C_{tt}-C_{st}^2\right\}
                       -\sum_{x\ne o}C_{xx}.}
 \label{restart:graph-coefficient}
\end{equation}
The displayed coefficient is independent of the choice of root.
\end{theorem}
\begin{proof}
The root identity gives equality of the compact integrals at these
endpoints.  For \(u=A_s,v=A_t\), quaternion multiplication gives
\[
 1-w(e^{gv}e^{-gu})
 =\frac{g^2}{2}|v-u|^2
   -g^4\left\{\frac{|v-u|^4}{24}
                       +\frac{|u\times v|^2}{6}\right\}+O(g^6).
\]
There is no cubic term.  Under the centered Gaussian of
Theorem~\ref{restart:main},
\[
 \E|v-u|^4=15r_e^2,\qquad
 \E|u\times v|^2=6(C_{ss}C_{tt}-C_{st}^2),\qquad
 \E H=\sum_{x\ne o}C_{xx}.
\]
For the middle identity, use
\(|u\times v|^2=|u|^2|v|^2-(u\cdot v)^2\) and the Gaussian pairings:
the two expectations are respectively
\(9C_{ss}C_{tt}+6C_{st}^2\) and
\(3C_{ss}C_{tt}+12C_{st}^2\).
Substitution in \eqref{restart:k2} gives
\eqref{restart:graph-coefficient}.  At zero endpoints the rescaled
integrand, including Haar measure, is even in \(g\).  The same dominated
Taylor argument through order four gives an \(O_G(g^4)\) divided error,
which is the remainder in \eqref{restart:flat-asymp}.

Finally the compact integral is independent of the root by global right
Haar translation.  Its leading factor and its first relative asymptotic
coefficient are unique.  Hence \(a_G\) is root independent; this also
follows by changing the pinned Green function in
\eqref{restart:graph-coefficient}.
\end{proof}

This result is a complete coefficient for a conditional temporal slab,
not the coefficient \(p_2\) of the full four-dimensional vacuum.
The background spatial fields have not been integrated.

\subsection{Trees, cycles and a noncommuting check}
\label{restart:tests}

\begin{proposition}[Exact tree normalization]
\label{restart:tree}
If \(G\) is a tree, then for every \(U,V\) and every \(\gamma>0\),
\[
 K_\gamma^o(U,V)=K_\gamma(U,V)=z_\gamma^{\,n-1},\qquad
 z_\gamma=\int e^{-\gamma(1-w(a))}\,dH(a)
          =\frac{2e^{-\gamma}I_1(\gamma)}{\gamma}.
\]
In particular \(a_G=-3(n-1)/8\).
\end{proposition}
\begin{proof}
Integrate a nonroot leaf.  Its single incident plaquette contains its
temporal variable once, or its inverse once; Haar left/right invariance
gives the constant \(z_\gamma\), independently of its neighbor and of
the two endpoint link values.  Repeat to the root.
For the coefficient, the one-variable Gaussian expansion gives
\(\E|A|^4/24-\E|A|^2/3=15/24-1=-3/8\).
The leading constant is
\((2\pi)^{3/2}/(2\pi^2)\gamma^{-3/2}\).
Products prove the claim.  The Bessel identity is the radial Haar
integral; it is also the one-group formula already present in the
archives, not a new special-function result.
\end{proof}

\begin{proposition}[Cycle coefficient]
\label{restart:cycle}
For the cycle on \(m\ge3\) vertices,
\[
 \det L=m,\qquad
 a_{C_m}=\frac{(m-1)(m-5)}{8m}.
\]
Thus a four-cycle has coefficient \(-3/32\), not the tree coefficient
\(-9/8\) of a three-edge spanning tree.
\end{proposition}
\begin{proof}
Pin vertex zero and label the others \(1,\ldots,m-1\).  Direct
multiplication by the tridiagonal rooted Laplacian verifies
\[
 C_{ij}=\min(i,j)-ij/m.
\]
Its determinant recurrence gives \(\det L=m\).  Every edge resistance
is \((m-1)/m\).  Summing the diagonal and the edge minors gives
\[
 \sum_{i=1}^{m-1}C_{ii}=\frac{m^2-1}{6},\qquad
 \sum_{e=(s,t)}(C_{ss}C_{tt}-C_{st}^2)
                         =\frac{(m-1)(m-2)}6.
\]
For the latter sum the root edges contribute zero; the edge from \(i\)
to \(i+1\) contributes \(i(m-i-1)/m\).
Insert these sums into \eqref{restart:graph-coefficient}.
\end{proof}

\begin{proposition}[A noncommuting one-edge completeness test]
\label{restart:oneedge}
Orient a single edge from its pinned root to its other vertex.  For
arbitrary \(X,Y\in\R^3\), let \(A=Y-X+Z\), where \(Z\) is a standard
three-dimensional Gaussian.  The actual plaquette word is
\(e^{gX}e^{gA}e^{-gY}\).  Its coefficients obey
\begin{align}
 \E S_3&=0,& \Var S_3&=|X\times Y|^2,\notag\\
 \E H&=1+\tfrac13|Y-X|^2,&
 \E S_4&=-\tfrac58-\tfrac13|Y-X|^2
                              +\tfrac12|X\times Y|^2.
 \label{restart:oneedge-coeff}
\end{align}
Consequently \(k_2=-3/8\), independently of both endpoints.  If the
cubic variance is omitted, the false result is
\(-3/8-|X\times Y|^2/2\).
\end{proposition}
\begin{proof}
Equation~\eqref{restart:cubic} gives
\(S_3=-\det(X,A,Y)\).  Its mean vanishes since \(Y-X\) lies in the
span of \(X,Y\); its variance is the squared norm of \(X\times Y\).
The Haar expectation is immediate.  Substituting \(u=(X,A,-Y,0)\)
in \eqref{restart:quartic} and using
\[
 \E A=\mu,\quad
 \E|A|^2=|\mu|^2+3,\quad
 \E|A|^4=|\mu|^4+10|\mu|^2+15,\quad
 \E[|A|^2A]=(|\mu|^2+5)\mu
\]
gives the last expression in \eqref{restart:oneedge-coeff};
all its terms are polynomials of degree at most four.
Inserting them into \eqref{restart:k2} gives the cancellation.
The endpoint independence also follows independently from the exact
tree integral, so this is a check against the original compact object,
not just an internal agreement of two truncated diagrams.
\end{proof}

For example \(X=(1,0,0)\), \(Y=(0,1,0)\) gives
\(\E H=5/3\), \(\E S_4=-19/24\), and \(\Var S_3=1\).
The correct sum is \(-3/8\); deleting the return gives \(-7/8\).
A commuting or zero-endpoint one-link check would miss this error.

\section{How to combine the electric return with the joint source law}
\label{restart:tower}

The preceding theorem evaluates a conditional temporal integral.  It
must be combined with the spatial law before evaluating the two-time
source insertions of \(C\), V89.  The following elementary tower identity
specifies that combination without duplicate terms.

\begin{lemma}[Two-stage second-order normalization]
\label{restart:tower-lemma}
Let \((X,A)\) have a reference probability law \(\mu_0\), and consider
the positive, unnormalized weight
\[
 W_\varepsilon=\exp\{\varepsilon L_1+\varepsilon^2L_2+o(\varepsilon^2)\}.
\]
Assume its expansion is valid in \(L^1\), also with the weights needed
for the observables under consideration.  More precisely, assume
\[
 W_\varepsilon
 =1+\varepsilon L_1+\varepsilon^2(L_2+L_1^2/2)+o_{L^1}(\varepsilon^2),
 \qquad L_1\in L^2,\quad L_2\in L^1.
\]
Then the marginal weight after eliminating \(A\) is
\[
 \E[W_\varepsilon\mid X]
 =1+\varepsilon m_1+
          \varepsilon^2(m_2+m_1^2/2)+o_{L^1}(\varepsilon^2),
\]
where
\begin{equation}
 m_1=\E[L_1\mid X],\qquad
 m_2=\E[L_2\mid X]+\tfrac12\Var(L_1\mid X).
 \label{restart:tower-scores}
\end{equation}
Moreover the coefficient of \(\varepsilon^2\) in the global log
normalization is, equivalently,
\begin{equation}
 \E L_2+\tfrac12\Var(L_1)
 =\E m_2+\tfrac12\Var(m_1).
 \label{restart:total-variance}
\end{equation}
A conditional pointwise log expansion additionally follows wherever
the conditional expansion has the required pointwise remainder;
it is not inferred solely from an \(L^1\) remainder.
\end{lemma}
\begin{proof}
Conditional expectation contracts \(L^1\).  Apply it to the stipulated
expansion, then use
\(\E[L_1^2\mid X]=\Var(L_1\mid X)+m_1^2\).
Taking the ordinary expectation and then the scalar logarithm gives
the global coefficient.  The equality of its two forms is the law of
total variance.  For a bounded \(X\)-measurable source product \(Q\),
\(\E[QW_\varepsilon]=\E[Q\,\E(W_\varepsilon\mid X)]\) exactly;
the assumed weighted version gives the same conclusion for the
declared unbounded products.
\end{proof}

For the rooted electric calculation \(L_1=-S_3\) and
\(L_2=-S_4-H\).  The inner correction in \eqref{restart:tower-scores}
is precisely \(k_2\), with the minus sign in \(k_1\).
The subsequent spatial correction contains the variance of the
\emph{conditional mean} \(k_1\), not another copy of the entire
temporal variance.  The residual root is an additional integration
already handled in Corollary~\ref{restart:root-expansion}.
In a full finite history, conditioning on all spatial slices makes
different temporal slabs independent; their conditional variances
add.  The later spatial covariance of their conditional means can
couple different times and must not be discarded.

This is a bookkeeping rule for an actual joint expansion, not an
assertion that an unrestricted native Gaussian reference or a
cutoff-uniform expansion has been constructed.  In particular the
bounded-endpoint result cannot be integrated over a growing spatial
history without source-weighted tail and remainder bounds.

There is also a useful normalization distinction.
The flat constant \(a_G/\gamma\), when independent of every spatial
endpoint, cancels from a normalized fixed-\(\gamma\) source expectation.
It does \emph{not} cancel from a coupling derivative of the partition
function.  Conversely the endpoint-dependent terms in \(k_1,k_2\) can
change separated source covariances, even when a scalar plaquette
normalization has been matched.  This explains why the two branches
need the same complete electric calculation but different output tests.

\section{Verification, preservation, and the next unpaid estimates}
\label{restart:next}

\subsection{Reproducible checks and their limits}

The companion diagnostic is
\begin{center}
\path{scripts/verify_ym_restart_v1_electric.py}.
\end{center}
It writes no files and separates exact rational identities from
floating-point diagnostics.  Its completed checks are:
\begin{itemize}
\item 88 rooted path/star cases, 52 rooted cycle cases, and all eight
root choices of an open cube, using exact rational matrix inversion.
The cube coefficient is \(5/8\).
\item Thirty exact polynomial Gaussian-moment checks of
Proposition~\ref{restart:oneedge}, all with nonzero cubic-variance
correction; deletion of that correction fails every selected case.
\item The arithmetic conversion of the published rounded coefficient
to the present convention in \eqref{restart:p2}.
\item Independent floating-point cycle partition sums
\[
 K_\gamma(1,1)=\sum_{j=0}^\infty
 (j+1)^2\left(\frac{2e^{-\gamma}I_{j+1}(\gamma)}{\gamma}\right)^m .
\]
This formula is the trace of the \(m\)-fold central convolution kernel,
by character orthogonality.  Values at \(\gamma=512,1024\), with
Richardson extrapolation, agree with the proved coefficients for
\(m=3,4,5,6,8\) to less than \(2.7\,10^{-7}\).
The truncation and floating arithmetic are diagnostics, not interval
bounds and not premises of the analytic theorem.
\end{itemize}
No Wilson Monte Carlo is claimed for these tests.
The analytic Laplace remainder is proved in the text; the finite tests
do not establish its uniformity in volume.

The pre-edit SHA--256 baseline records 46 source files.  At verification,
45 remain byte-identical at their original paths; the externally
superseded V85 is recovered byte for byte from V86 as described above.
None of the requested historical bodies is rewritten by this work.
The audit corrections live in this new note.  The folder receives only
this new \path{.tex} source; compilation and verification outputs are
kept outside it.  There is no restored V3 and no obsolete V4 draft.
The next active version should append to this V1 and replace only this
active source after verification.  At active V5 the companion programs
are to be reread, and subsequently every five versions; at V100 the
active lineage is to be archived before another compact restart.

\subsection{What this version adds}

\emph{Proved here on the stated finite objects:}
the finite-time/vacuum distinction; exact full-history temporal
marginalization with its root integration; the complete bounded-endpoint
rooted expansion with a proved fixed-graph remainder; its root
completion; the evaluated flat graph coefficient; and the
noncommuting one-edge completeness test.  Conditional expectation,
Gaussian integration, the Laplace method and character orthogonality
are standard methods, not claimed as newly invented techniques.

\emph{Imported, not reproved:}
the published infinite-volume perturbative coefficient, and the scoped
historical interfaces in Section~\ref{restart:context}.

\emph{Not established:}
the full four-dimensional \(p_2\) from the corrected native chain,
its differentiated remainder, the source-weighted unrestricted
history expansion, or a uniform large-graph version of
Theorem~\ref{restart:main}.  In particular an inverse grounded
Laplacian can have growing norm; the finite-graph domination used in
the proof cannot be called volume-free.

\subsection{Direction of the next iteration}

The next calculation should use the complete word coefficients
\eqref{restart:vertices}--\eqref{restart:quartic} in the full finite-history
joint law, with one explicit convention for spatial coordinates and
their measure.  Integrate temporal variables exactly once using the
conditional rule, keeping the root integration and all source factors.
The first target is a finite-volume coefficient that passes both the
noncommuting checks above and an independently derived four-dimensional
plaquette check.  Agreement with the numerical literature benchmark is
a diagnostic; it is not a substitute for proving the expansion.

The subsequent analytic target is a \emph{connected, source-weighted}
remainder for the positive-time even/odd packet of \(C\), not a uniform
norm approximation to an extensive global density.  Its estimates
must cover the actual simultaneous flow at the declared physical scale.
If fixed-graph constants or endpoint cutoffs again obstruct export,
the work must go upstream to locality and scale decomposition of this
complete joint expansion, rather than insert the desired source Gram
as an assumption.

Even after that, a continuum construction, noncollapse of the
physically normalized positive-time sources, compatible thermodynamic
limits, and a common upper spectral contraction on a dense limiting
source space remain necessary.  \(C\), V84 already proves why a
finite-depth compressed contraction gives no full lower mass bound;
\(C\), V86 already distinguishes finite-regulator coverage from limiting
coverage.  Those obstructions are retained, not reproved or bypassed.

\begin{center}
\fbox{\parbox{.9\textwidth}{\textbf{Status.}
This V1 repairs and advances the finite-regulator electric calculation.
It does not establish the interacting four-dimensional continuum
Yang--Mills mass gap.  That goal remains open.}}
\end{center}

\begin{thebibliography}{9}
\bibitem{APT}
A.~Athenodorou, H.~Panagopoulos and A.~Tsapalis,
\emph{The Lattice Free Energy of QCD with Clover Fermions, up to
Three-Loops}, arXiv:0710.3856, especially equations (7), (9), (11) and (14).
\url{https://arxiv.org/abs/0710.3856}.
The pure-glue term, not the fermion correction, is used here.

\bibitem{ACFP}
B.~All\'es, M.~Campostrini, A.~Feo and H.~Panagopoulos,
\emph{The Three-Loop Lattice Free Energy},
Phys.\ Lett.\ B \textbf{324} (1994), 433--436,
arXiv:hep-lat/9306001.
\url{https://arxiv.org/abs/hep-lat/9306001}.
This is the earlier free-energy calculation cited by \cite{APT};
no new three-loop coefficient is asserted in this restart.
\end{thebibliography}
% BEGIN CONSOLIDATED V2 APPEND -- 2026-09-19
\clearpage
\section{V2: a complete finite-history coefficient and a native two-time test}
\label{r2:context}

V1 computed the temporal integral conditional on bounded spatial
endpoints.  It did not justify integrating its remainder against the
whole spatial history.  The present step performs that full integration
on a precisely specified, contractible finite Euclidean box.  There is
no imposed all-good event and no frozen reference substituted for its
Wilson measure.  An exact maximal-tree change of variables removes
local gauge redundancy without a guessed orbit determinant.  All
spatial and temporal plaquettes then enter one stationary-phase
calculation.

The new results are a complete first correction to the finite-box
partition function, a remainder that can be differentiated in local
plaquette couplings, and explicitly evaluated separated plaquette
covariances.  A nonzero mixed cubic contribution shows concretely why
adding individually named electric and magnetic pieces without fixing
their common coordinates can fail.

\emph{Scope.}  The open box is a finite-regulator completeness test,
not a replacement for the programme's physical vacuum or its periodic
spatial regulator.  On a periodic box, nontrivial flat holonomies prevent
the isolated-saddle argument used below.  No equality of open- and
periodic-boundary continuum limits is assumed.  The present calculation
advances the missing joint-law input; it does not settle those limits
or the mass gap.

The previous goal turn made concrete progress through V1's complete
conditional coefficient and its corrected audit.  The f16.1 V86 and
f16.2 V89 branch files remain unchanged at the start of this iteration.
The companion-program checkpoint remains scheduled for active V5.

\subsection{The full compact measure, with one coordinate convention}
\label{r2:measure}

Let \(\Lambda\) be a finite rectangular cubical box in dimension
\(d_{\rm E}\ge2\), with vertices
\[
 x=(x_0,\ldots,x_{d_{\rm E}-1}),\qquad 0\le x_\mu\le L_\mu,
 \qquad L_\mu\ge1.
\]
Links have positive coordinate orientation.  A plaquette is the ordered
word
\[
 U_{x,\mu}\,U_{x+\hat\mu,\nu}\,
 U_{x+\hat\nu,\mu}^{-1}\,U_{x,\nu}^{-1},\qquad \mu<\nu.
\]
There are no periodic identifications and no fixed boundary links.
All links, including boundary links, are integrated with normalized Haar
measure.  Let \(v,e,f\) be the numbers of vertices, links and plaquettes,
and put \(b=e-v+1\).  With positive numbers \(c_p\), define
\begin{equation}
 P_p(U)=1-w(U_p),\qquad
 Z_\Lambda(\gamma,c)=
 \int\exp\left\{-\gamma\sum_p c_pP_p(U)\right\}\prod_\ell dH(U_\ell).
 \label{r2:Z}
\end{equation}
These \(c_p\) are local marks for differentiation; they are set to one
in the examples.  They do not change the physical time convention.

Choose a rooted spanning tree \(\mathcal T\) in the complete
\(d_{\rm E}\)-dimensional link graph.  This is a tree of the whole
Euclidean history, not one separately chosen tree per spatial slice.
Let \(W_1,\ldots,W_b\) denote the transformed non-tree links, with all
tree links set to the identity.

\begin{lemma}[Exact tree variables and the unique flat point]
\label{r2:tree}
For every integrable gauge-invariant function \(F\),
\begin{equation}
 \int F(U)\prod_\ell dH(U_\ell)
   =\int F(U^{\mathcal T}(W))\prod_{j=1}^b dH(W_j).
 \label{r2:tree-Haar}
\end{equation}
In these variables, the zero set of \(\sum_pc_pP_p\) consists only of
\(W_1=\cdots=W_b=1\).  If \(B\) is the \(f\times b\) signed
plaquette--chord incidence matrix, then
\begin{equation}
 Q(c)=B^{\mathsf T}\operatorname{diag}(c_p)B>0.
 \label{r2:Q}
\end{equation}
Residual global conjugation is retained in the chord variables; no
extra quotient measure is inserted.
\end{lemma}
\begin{proof}
Fix the gauge transformation at the root to be the identity and solve
successively along the tree for vertex variables \(h_x\) such that
\(U_\ell=h_{s(\ell)}U_\ell^{\mathcal T}h_{t(\ell)}^{-1}\).
The nonroot vertex variables and the chord variables give global
product-group coordinates for all original links.  Successive Haar
left/right translations along the tree and then on the chords show
that product Haar measure becomes
\(\prod_{x\ne o}dH(h_x)\prod_jdH(W_j)\).
Gauge invariance removes the \(h_x\)'s from \(F\); their total integral
is one.  This proves \eqref{r2:tree-Haar}.

Every plaquette loss is nonnegative and vanishes exactly when its
holonomy is the identity.  On the rectangular box, loops are generated
by plaquette boundaries and immediate backtracks: one can commute
coordinate steps across a square and then cancel inverse steps.
Flat parallel transport is therefore path independent.  After all
tree links are fixed to one, every chord is one as well.

The linear version of the same argument says that a real link
cochain with zero plaquette circulation is a vertex gradient.
A gradient vanishing on a spanning tree is zero everywhere.
Consequently \(Bz=0\) forces \(z=0\); positive \(c_p\)'s prove
\eqref{r2:Q}.  The argument applies independently to each color.
A common conjugation of all \(W_j\)'s has not been divided out, so
there is no further Jacobian to attach to the displayed product
measure.
\end{proof}

Maximal-tree gauge coordinates are standard; for modern context see
I.~M.~Burbano and C.~W.~Bauer,
\href{https://arxiv.org/abs/2409.13812}{arXiv:2409.13812},
Section III and Appendix C.3.2.  The elementary Haar argument above,
rather than any Hamiltonian truncation from that paper, is what is
used here.

\subsection{One complete coefficient and a differentiable remainder}
\label{r2:laplace}

Set \(g=\gamma^{-1/2}\), \(W_j=e^{gX_j}\), and
\(C(c)=Q(c)^{-1}\).  Under \(\E_c\), the vectors \(X_j\in\R^3\)
are jointly centered Gaussian with
\[
 \E_c[X_i^\alpha X_j^\beta]=\delta_{\alpha\beta}C_{ij}(c).
\]
In every plaquette word replace a tree link by zero at the Lie-algebra
level and a positively or negatively traversed chord by \(X_j\) or
\(-X_j\).  Using precisely the word coefficients
\eqref{restart:cubic}--\eqref{restart:quartic}, write
\begin{align}
 \sum_pc_pP_p(e^{gX})
   &=\tfrac12g^2X^{\mathsf T}(Q(c)\otimes I_3)X
       +g^3 S_3(c,X)+g^4 S_4(c,X)+O_X(g^5),\notag\\
 H_{\mathcal T}(X)&=\tfrac13\sum_{j=1}^b|X_j|^2,\qquad
 \mathcal A_\Lambda(c)=
       -\E_c S_4-\E_c H_{\mathcal T}+\tfrac12\E_c S_3^2 .
 \label{r2:coefficient}
\end{align}
Here \(S_j=\sum_pc_ps_{j,p}\).  In particular \(S_3\) includes all
plaquettes before it is squared.  Its mean is zero by Gaussian parity.

\begin{theorem}[Full finite-box weak-coupling expansion]
\label{r2:full}
For every fixed box and every compact set
\(K\subset(0,\infty)^f\),
\begin{align}
 \log Z_\Lambda(\gamma,c)
  ={}&-\frac{3b}{2}\log\gamma
      +b\log\frac{(2\pi)^{3/2}}{2\pi^2}
      -\frac32\log\det Q(c)
      +\frac{\mathcal A_\Lambda(c)}{\gamma}
      +R_\Lambda(\gamma,c).
 \label{r2:logZ}
\end{align}
Uniformly for \(c\in K\), \(\gamma\ge\gamma_{\Lambda,K}\), and
\(|\alpha|\le2\), \(j=0,1\),
\begin{equation}
 \left|\partial_\gamma^j\partial_c^\alpha
                R_\Lambda(\gamma,c)\right|
       \le C_{\Lambda,K}\gamma^{-2-j}.
 \label{r2:differentiated-remainder}
\end{equation}
The scalar function \(\mathcal A_\Lambda(c)\) is independent of the
spanning tree.  The constants in \eqref{r2:differentiated-remainder}
need not be uniform in the box.
\end{theorem}
\begin{proof}
Apply Lemma~\ref{r2:tree} to the exact compact integral.
Its unique minimum and positive Hessian permit a fixed coordinate
neighborhood of the identity on which the action is bounded below
by \(a_{\Lambda,K}|x|^2\), uniformly on \(K\).
Outside a smaller such neighborhood the action has a positive minimum.
That part of the integral, and its finitely many \(c,\gamma\)
derivatives, is exponentially small in \(\gamma\): differentiation
costs only a fixed polynomial in \(\gamma\) and bounded plaquette
functions.

In the coordinate neighborhood Haar measure is exactly
\[
 \prod_{j=1}^b dH(e^{gX_j})
 =\frac{g^{3b}}{(2\pi^2)^b}
   \prod_j\left(\frac{\sin(g|X_j|)}{g|X_j|}\right)^2dX.
\]
Thus its first logarithmic correction is
\(-g^2H_{\mathcal T}\), counted once for each chord.
The Gaussian integral gives the first three terms of
\eqref{r2:logZ}.  The relative integrand through order two is
\[
 1-gS_3+g^2\{S_3^2/2-S_4-H_{\mathcal T}\}.
\]
The odd term integrates to zero; the second coefficient is
\eqref{r2:coefficient}.

To justify the stated differentiated error, insert an even smooth
coordinate cutoff as in V1.  Taylor's integral formula for the
rescaled action \(g^{-2}S_c(e^{gX})\) bounds each required \(g\)-derivative
by a polynomial in \(1+|X|\) on the cutoff support.  Derivatives in the
finitely many \(c_p\)'s add such polynomial factors.  Uniform positive
definiteness of \(Q(c)\), \(c\in K\), bounds these integrands by
\(C_{\Lambda,K}(1+|X|)^M e^{-a_{\Lambda,K}|X|^2}\).
Cutoff derivatives are supported at \(|X|\) of order \(g^{-1}\)
and obey the same integrable bound.  Dominated Taylor integration
therefore holds jointly through at least four \(g\)-derivatives and
two \(c\)-derivatives.  The analytic Wilson functions allow further
derivatives if needed.

After division by \(g^{3b}\), the localized integral is an even
function of \(g\): change \(X\) to \(-X\) when changing the sign of \(g\).
Hence its remainder after the \(g^2\) term is \(O(g^4)\), and the
\(g\)-derivative of that remainder is \(O(g^3)\), also after the
specified \(c\)-derivatives.  The leading Gaussian integral is bounded
away from zero on \(K\), so taking its logarithm preserves these
bounds.  Finally \(\partial_\gamma=-\frac12g^3\partial_g\) gives
\eqref{r2:differentiated-remainder}.  The exponentially small
complement does not change it.

The original compact partition function is the same for every tree
at each \(\gamma,c\).  Uniqueness of the relative asymptotic
coefficients in \eqref{r2:logZ} proves tree independence of
\(\mathcal A_\Lambda\).  Individual expectations in
\eqref{r2:coefficient} are not asserted tree independent.
\end{proof}

This proves an unrestricted compact \emph{joint-history} expansion on
the stated box.  It is stronger than integrating V1's bounded-endpoint
remainder formally against an unspecified law, but weaker than a
volume-uniform or periodic-vacuum expansion.

\subsection{A finite algebraic formula for the whole coefficient}
\label{r2:algebra}

The following formulas reduce \(\mathcal A_\Lambda\) to finite rational
linear algebra when the \(c_p\)'s are rational.
For a plaquette let \(u_1,\ldots,u_4\) be its signed chord vectors,
with a zero vector for each tree edge, and set
\(k_{ij}=\E_c[u_i^1u_j^1]\).  Thus \(k\) is its \(4\times4\)
one-color covariance matrix.  Define
\begin{align}
 \mathcal W(k)={}&\frac58\sum_i k_{ii}^2
 +\frac14\sum_{i<j}(9k_{ii}k_{jj}+6k_{ij}^2)\notag\\
 &+\frac52\sum_{i<j}(k_{ii}+k_{jj})k_{ij}\notag\\
 &+\sum_{i<j}\sum_{h\ne i,j}
       \left(\frac92k_{ij}k_{hh}+3k_{ih}k_{jh}\right)\notag\\
 &+9k_{12}k_{34}-3k_{13}k_{24}+9k_{14}k_{23}.
 \label{r2:W}
\end{align}
Let \(I=(i_1<i_2<i_3)\) run through triples of chord indices, and collect
the cubic polynomial uniquely as
\[
 S_3(c,X)=\sum_I t_I(c)\det(X_{i_1},X_{i_2},X_{i_3}).
\]
The coefficients \(t_I\) are obtained by summing the signed triples
in each ordered plaquette word.  A repeated chord index contributes
zero.  Let \(C[I,J]\) be the \(3\times3\) covariance submatrix on
ordered triples \(I,J\).

\begin{proposition}[Complete finite contraction formula]
\label{r2:finite-formula}
One has exactly
\begin{equation}
 \boxed{\displaystyle
 \mathcal A_\Lambda(c)=
 \sum_p c_p\mathcal W(k^{(p)})
       -\Tr C
       +3\sum_{I,J}t_I(c)t_J(c)\det C[I,J].}
 \label{r2:finite-A}
\end{equation}
In particular no temporal--spatial cross contraction is omitted from
the last sum.  Moreover
\begin{equation}
 \mathcal A_\Lambda(sc)=s^{-1}\mathcal A_\Lambda(c),\quad
 \sum_pc_p\partial_p\mathcal A_\Lambda=-\mathcal A_\Lambda,\quad
 \sum_{p,q}c_pc_q\partial_p\partial_q\mathcal A_\Lambda
                                  =2\mathcal A_\Lambda.
 \label{r2:homogeneity}
\end{equation}
\end{proposition}
\begin{proof}
Apply Gaussian pairings to the scalar fourth-order word
\eqref{restart:quartic}.  For example
\[
 \E|u_i|^4=15k_{ii}^2,\quad
 \E(|u_i|^2|u_j|^2)=9k_{ii}k_{jj}+6k_{ij}^2,
\]
\[
 \E[(u_i\cdot u_j)|u_h|^2]
       =9k_{ij}k_{hh}+6k_{ih}k_{jh}.
\]
For \(h=i\) this last expression is \(15k_{ii}k_{ij}\).
The three four-vector products in \eqref{restart:quartic} combine to
\(9k_{12}k_{34}-3k_{13}k_{24}+9k_{14}k_{23}\).
Consequently \(\E w_{4,p}=\mathcal W(k^{(p)})\), proving the quartic
part, since \(s_{4,p}=-w_{4,p}\).
The Haar part is \(-\E H_{\mathcal T}=-\Tr C\).

For the cubic part, expanding both determinants with the
three-dimensional alternating tensor gives
\[
 \E[\det(X_{i_1},X_{i_2},X_{i_3})
       \det(X_{j_1},X_{j_2},X_{j_3})]=6\det C[I,J].
\]
A pairing internal to one triple vanishes under antisymmetrization;
the six bijections between triples give precisely the determinant
with the factor six.  Squaring the entire \(S_3\) and multiplying by
\(1/2\) proves \eqref{r2:finite-A}.
Finally \(Q(sc)=sQ(c)\), \(C(sc)=s^{-1}C(c)\), and
\(t_I(sc)=st_I(c)\).  Every term in \eqref{r2:finite-A} is homogeneous
of degree \(-1\); differentiation gives \eqref{r2:homogeneity}.
\end{proof}

\subsection{The same calculation supplies two-time source insertions}
\label{r2:sources}

For plaquettes \(p,q\), define the real matrix
\[
 M(c)=BC(c)B^{\mathsf T}.
\]
Differentiating the original compact integral, not a fitted expansion,
gives
\[
 \partial_p\log Z_\Lambda=-\gamma\E_{\gamma,c}P_p,\qquad
 \partial_p\partial_q\log Z_\Lambda
       =\gamma^2\Cov_{\gamma,c}(P_p,P_q).
\]
These identities hold also for plaquettes in different time slices.

\begin{corollary}[Native plaquette mean and joint covariance]
\label{r2:covariance}
For the fixed box, uniformly on \(K\),
\begin{align}
 \E_{\gamma,c}P_p
   &=\frac{3M_{pp}(c)}{2\gamma}
       -\frac{\partial_p\mathcal A_\Lambda(c)}{\gamma^2}
       +O_{\Lambda,K}(\gamma^{-3}),\notag\\
 \Cov_{\gamma,c}(P_p,P_q)
   &=\frac{3M_{pq}(c)^2}{2\gamma^2}
       +\frac{\partial_p\partial_q\mathcal A_\Lambda(c)}{\gamma^3}
       +O_{\Lambda,K}(\gamma^{-4}).
 \label{r2:covariance-expansion}
\end{align}
At \(c_p=1\), the orientation-averaged plaquette satisfies
\begin{equation}
 \frac1f\sum_p\E_{\gamma,1}P_p
 =\frac{3b}{2f\gamma}
             +\frac{\mathcal A_\Lambda(1)}{f\gamma^2}
             +O_\Lambda(\gamma^{-3}).
 \label{r2:average}
\end{equation}
\end{corollary}
\begin{proof}
Since \(\partial_pQ=B_p^{\mathsf T}B_p\),
\[
 \partial_p\log\det Q=M_{pp},\qquad
 \partial_p\partial_q\log\det Q=-M_{pq}^2.
\]
Apply Theorem~\ref{r2:full} and its coupling-derivative bounds to the
exact identities above.  For the average, \(M(1)\) is an orthogonal
projection of rank \(b\), and
\(\sum_p\partial_p\mathcal A_\Lambda(1)=-\mathcal A_\Lambda(1)\).
This proves \eqref{r2:average}, equivalently obtained using the paid
\(\gamma\)-derivative remainder.
\end{proof}

There is already a limited, genuinely volume-independent estimate at
leading order:
\begin{equation}
 0\preceq \frac32(M(1)\circ M(1))\preceq\frac32 I_f.
 \label{r2:leading-uniform}
\end{equation}
Here \(\circ\) denotes entrywise product.
Indeed \(M^2=M=M^{\mathsf T}\) gives
\(\sum_qM_{pq}^2=M_{pp}\le1\).
The matrix with entries \(M_{pq}^2\) is positive semidefinite, as the
Gram matrix of the tensor squares of a Gram representation of \(M\).
Its nonnegative row sums are at most one, giving the operator bound.
This controls the leading covariance of coefficient-normalized
plaquette banks without the norm of a tree-gauge propagator.
It neither bounds the nonlinear correction nor produces a positive
mass or a uniform lower Gram.

For completeness a source need not be introduced by a local
coupling derivative.  The full-law input score is now specified on
this box, rather than assumed.

\begin{proposition}[First joint correction for a smooth gauge-invariant source]
\label{r2:general-source}
Let \(F,G\) be real \(C^6\) gauge-invariant functions on the finite link
space, held fixed as \(\gamma\) varies.  Subtract their values at the flat connection.
In tree coordinates write their homogeneous Taylor polynomials
\[
 F(e^{gX})=g^2F_2+g^3F_3+g^4F_4+\cdots,\qquad
 G(e^{gX})=g^2G_2+g^3G_3+g^4G_4+\cdots.
\]
There is no linear term.  Put
\[
 W_2=S_3^2/2-S_4-H_{\mathcal T},\qquad
 F_2^\circ=F_2-\E_cF_2,\quad G_2^\circ=G_2-\E_cG_2 .
\]
Then
\[
 \Cov_{\gamma,c}(F,G)
     =\gamma^{-2}\Gamma_4+\gamma^{-3}\Gamma_6+O_{\Lambda,K,F,G}(\gamma^{-4}),
\]
where
\begin{align}
 \Gamma_4={}&\Cov_c(F_2,G_2),\notag\\
 \Gamma_6={}&\Cov_c(F_4,G_2)+\Cov_c(F_2,G_4)+\E_c(F_3G_3)\notag\\
 &-\E_c[(F_3G_2^\circ+F_2^\circ G_3)S_3]
       +\Cov_c(F_2^\circ G_2^\circ,W_2).
 \label{r2:source-coefficient}
\end{align}
All expectations use the one complete joint Gaussian on the whole box.
The statement applies entrywise to a finite source bank.
\end{proposition}
\begin{proof}
Global conjugation fixes the flat point and acts by a common
three-dimensional rotation on all tangent vectors \(X_j\).
An invariant real linear functional on a direct sum of these rotation
representations is zero, proving the absence of first derivatives.
The homogeneous Taylor terms of odd degree have odd Gaussian parity.

Use the same compact localization as in Theorem~\ref{r2:full}.
Since \(F,G\) vanish to second order at the flat point,
\(g^{-2}F(e^{gX})\) and \(g^{-2}G(e^{gX})\), with a local cutoff,
have four controlled \(g\)-derivatives bounded by polynomials in \(X\);
Taylor's integral formula and \(C^6\) regularity suffice.
Their products with the rescaled density can therefore be integrated
through order four with the same Gaussian majorant.  The compact
complement is exponentially small even after multiplication by
\(g^{-4}\).

The normalized expectation uses the divided weight
\[
 \frac{1-gS_3+g^2W_2+\cdots}{1+g^2\E_cW_2+\cdots}.
\]
Multiply this by the two source expansions and subtract the product
of their similarly expanded means.  At order \(g^2\) in the
amplitude-normalized covariance, the source-only terms give the first
three terms of \(\Gamma_6\); the first density insertion gives the
fourth term; normalization of the second insertion gives the last
covariance in \eqref{r2:source-coefficient}.
Every odd order vanishes under \(X\mapsto-X\).
The controlled fourth-order normalized remainder becomes
\(O(g^8)=O(\gamma^{-4})\) in the unscaled covariance.
\end{proof}

For any fixed finite lattice and finite spatial flow time, composing a
smooth observable with the complete simultaneous spatial Wilson flow
is still smooth and gauge invariant.  Proposition~\ref{r2:general-source}
therefore applies to that \emph{actual} source.  It supplies a native
finite-box version of the joint-insertion ledger of \(C\), V89.
Its constants may depend strongly on the flow time and on the box.
It does not justify taking lattice flow time proportional to \(a^{-2}\)
while sending the lattice spacing \(a\) to zero.

\subsection{Exact four-dimensional values and a mixed-sector countercheck}
\label{r2:values}

The finite contraction formula can be evaluated without numerical
differentiation.  For a plaquette row \(B_p\) let
\(v_p=CB_p^{\mathsf T}\).  At \(c=1\),
\begin{equation}
 \partial_pC=-v_pv_p^{\mathsf T},\qquad
 \partial_p\partial_qC
       =M_{pq}(v_pv_q^{\mathsf T}+v_qv_p^{\mathsf T}).
 \label{r2:C-jets}
\end{equation}
Together with the linear dependence of each \(t_I(c)\), this computes
\(\mathcal A,\partial_p\mathcal A,\partial_q\mathcal A,
\partial_p\partial_q\mathcal A\) by exact rational arithmetic.

Here and in the diagnostic, a box label records the number of
\emph{cells}, not vertices, in each coordinate direction; the first
coordinate is time.  For four dimensions take the spatial plaquettes
of orientation \((1,2)\) based at
\[
 p:\ (0,0,0,0),\qquad q:\ (L_0,0,0,0).
\]
They lie on opposite temporal ends.  This is a free-boundary finite
history, not a thermal trace or an infinite-time Perron state.

\begin{center}
\small
\begin{tabular}{@{}lrrrr@{}}
\toprule
Box & \(b\) & \(f\) & \(\mathcal A_\Lambda(1)\)
  & \(\mathcal A_\Lambda(1)/f\)\\
\midrule
\(1\times1\) & 1 & 1 & \(-3/8\) & \(-3/8\)\\
\(2\times2\) & 4 & 4 & \(-3/2\) & \(-3/8\)\\
\(1\times1\times1\) & 5 & 6 & \(5/48\) & \(5/288\)\\
\(1\times1\times1\times1\) & 17 & 24
 & \(2407/768\) & \(2407/18432\)\\
\(2\times1\times1\times1\) & 29 & 42
 & \(65989891/11113200\) & \(65989891/466754400\)\\
\bottomrule
\end{tabular}
\end{center}
These are exact evaluations of \eqref{r2:finite-A} for the specified
finite incidence data, not fits.  In particular the four-dimensional
orientation-averaged means have leading coefficients \(17/16\) and
\(29/28\), respectively, not the infinite-volume \(3/4\).
The boundary is substantial on these small boxes.

For the separated pair specified above, the evaluated joint covariances
are
\begin{align}
 \Cov_\gamma(P_p,P_q)\big|_{1^4}
   &=\frac{3}{128\gamma^2}
           +\frac{7777}{442368\gamma^3}+O(\gamma^{-4}),\notag\\
 \Cov_\gamma(P_p,P_q)\big|_{2\times1^3}
   &=\frac{1369}{2381400\gamma^2}
           +\frac{15702577411}{19848730860000\gamma^3}
           +O(\gamma^{-4}).
 \label{r2:evaluated-covariances}
\end{align}
The corresponding values of \(M_{pq}\) are \(1/8\) and \(37/1890\).
These insertions involve the complete native joint law of their boxes.
A plaquette mean alone would not supply their coefficients.

To make the calculation reproducible, vertices are lexicographically
ordered and each positive link is indexed first by its initial vertex,
then its coordinate direction.  A tree is obtained by scanning links
in a declared order and retaining a link exactly when it joins two
previously disconnected components.  The four orders used are:
temporal links first; spatial links first; reverse of the original
order; and increasing \((17i+11)\bmod e\), with index \(i\) breaking
ties.  Chords retain original index order.  Plaquette words are exactly
those in Section~\ref{r2:measure}.  These rules, rational Gaussian
elimination for \(Q^{-1}\), formulas \eqref{r2:W},
\eqref{r2:finite-A}, and \eqref{r2:C-jets} specify every table entry
without discretionary vertex factors.

Tree invariance is a theorem, while these evaluations provide checks
on its implementation.  For example on \(1^4\), two decompositions are
\begin{center}
\small
\begin{tabular}{@{}lrrrr@{}}
\toprule
Tree & \(-\E S_4\) & \(-\E H_{\mathcal T}\)
 & \(\E S_3^2/2\) & total\\
\midrule
Spatial-first & \(-1247/384\) & \(-179/12\)
 & \(16357/768\) & \(2407/768\)\\
Interleaved & \(437/128\) & \(-241/12\)
 & \(15209/768\) & \(2407/768\)\\
\bottomrule
\end{tabular}
\end{center}
The separate terms change substantially, although the total and its
local coupling derivatives agree.

More specifically split the cubic sum according to the original
plaquette directions:
\(S_3=S_3^{\rm sp}+S_3^{\rm tm}\), where a temporal plaquette has
one direction equal to zero.  The mixed contribution to
\(\E S_3^2/2\) is \(\E(S_3^{\rm sp}S_3^{\rm tm})\).
On \(1^4\) it equals
\begin{equation}
 \E(S_3^{\rm sp}S_3^{\rm tm})
 =\begin{cases}
       151/144,&\text{spatial-first tree},\\
       947/1536,&\text{interleaved tree},\\
       0,&\text{temporal-first tree}.
   \end{cases}
 \label{r2:mixed}
\end{equation}
The zero in the last line is structural: that tree contains every
temporal link, so a temporal plaquette word contains at most two
nonidentity factors and has no scalar cubic term.
Other terms in \eqref{r2:finite-A} adjust correspondingly.
In the spatial-first coordinates, dropping only the mixed term changes
the total coefficient by \(151/144\); it is not a small rounding effect.

Thus a temporal-link Haar term or a so-called electric sunset is not
a coordinate-independent extra scalar that can be added to any spatial
chain.  This explicit check does not identify the precise remaining
error in \(B\), V86's different horizontal coordinates.  It tells us
what a complete reconstruction in those coordinates must reproduce
after all measure and cross terms are accounted for.

\subsection{Independent controls and the exact-computation scope}
\label{r2:controls}

For a two-dimensional simply connected rectangular lattice the compact
integral factorizes into single-plaquette integrals.  One proof fixes
an axial tree and successively changes remaining row-link variables to
plaquette holonomies, using triangular Haar translations.  Therefore
\[
 Z_\Lambda(\gamma,c)=\prod_p z_{\gamma c_p},\qquad
 \mathcal A_\Lambda(c)=-\frac38\sum_p c_p^{-1}.
\]
This independently checks the two-dimensional values,
\(\partial_p^2\mathcal A=-3/4\) at unit marks, and
\(\partial_p\partial_q\mathcal A=0\) for distinct plaquettes.
The diagonal covariance is
\(3/(2\gamma^2)-3/(4\gamma^3)+O(\gamma^{-4})\).

There is a further compact non-Abelian check for the
\(1\times1\times1\) cube.  Its six-face two-complex is a sphere.
For \(\lambda_j(\beta)=2e^{-\beta}I_{j+1}(\beta)/\beta\), character
orthogonality on every edge gives exactly
\begin{equation}
 Z_{\rm cube}(\gamma,c)=
       \sum_{j=0}^\infty (j+1)^2\prod_{p=1}^6\lambda_j(\gamma c_p).
 \label{r2:cube-character}
\end{equation}
Indeed expand each face weight as
\(\sum_j(j+1)\lambda_j\chi_j\); edge integrals force a common
representation on neighboring faces, and the remaining dimension
factor is \((j+1)^{v-e+f}=(j+1)^2\).
Absolute convergence follows, for fixed positive couplings, from the
factorial decay of the Bessel coefficients.
At equal marks this is also the flat temporal-cycle integral with
six edges in V1.  Its proved coefficient
\((6-1)(6-5)/(8\cdot6)=5/48\) agrees with the full cube calculation.

Differentiating \eqref{r2:cube-character} for two distinct faces
provides a separate check of their covariance.  If \(J\) has weights
proportional to \((j+1)^2\lambda_j(\gamma)^6\), that covariance is
\[
 \Var\left(\partial_\gamma\log\lambda_J(\gamma)\right).
\]
The exact contraction formula gives
\[
 \Cov_\gamma(P_p,P_q)
       =\frac1{24\gamma^2}+\frac7{288\gamma^3}+O(\gamma^{-4})
       \quad(p\ne q).
\]
A separate floating-point character sum at
\(\gamma=256,512,1024\) checks the second coefficient; its final
Richardson value differs from \(7/288\) by less than \(4.4\,10^{-8}\).
This last numerical comparison is not an interval enclosure, a
Wilson Monte Carlo, or a premise of the asymptotic theorem.

The diagnostic is
\begin{center}
\path{scripts/verify_ym_restart_v2_joint.py}.
\end{center}
Its coefficient tables, derivatives and tree comparisons use rational
arithmetic throughout.  Twenty box/tree combinations are checked,
including full first and mixed coupling jets for the separated pair
on each of the twelve three- or four-dimensional combinations.
The character calculation is separately labelled as floating-point.
The script is an inspectable finite-arithmetic diagnostic, not a
proof-assistant certificate.  Equations \eqref{r2:finite-A} and
\eqref{r2:C-jets}, with the fully specified incidence construction,
are the mathematical specification behind the evaluations.

\subsection{What is now paid, and what must go upstream next}
\label{r2:status}

\emph{Newly established on each fixed contractible box:}
a complete original-law expansion with no missing electric sector;
its local-coupling and coupling-size differentiated remainder;
normalized two-time source coefficients; exact nontrivial
four-dimensional examples; and explicit tree-dependent redistribution
of a mixed cubic contribution.
The leading source-covariance bound
\eqref{r2:leading-uniform} is independent of box size, but is only a
Gaussian leading-coefficient bound.

\emph{Still unpaid:}
uniform nonlinear bounds on
\(\partial_p\partial_q\mathcal A_\Lambda\) and higher connected
source insertions; a source-weighted remainder along growing histories;
control of the physical simultaneous-flow scale; periodic flat
holonomy sectors; and the limits and dense-core contraction required
for a continuum mass gap.

The new finite-box numbers are not another estimate of the
infinite-volume coefficient \(p_2^{\rm lit}\) in V1.
Changing only temporal length while keeping a one-cell spatial box
does not take a four-dimensional thermodynamic limit.
Nor may the positivity of the two examples in
\eqref{r2:evaluated-covariances} be promoted to a cutoff-uniform
physical lower Gram.

The next upstream task is now more specific.
Use the complete contraction formula to reorganize its
\emph{source derivatives} into local, gauge-independent connected
quantities before estimating them.  The norm of \(Q^{-1}\) in an
arbitrary tree gauge is not the relevant final quantity:
\eqref{r2:leading-uniform} already illustrates a cancellation that
would be lost by such a bound.  A finite limit of
\(\mathcal A_\Lambda/f\), even if proved, would control a mean
plaquette coefficient, not these joint derivatives or their
nonperturbative remainders.

All prior V1 text is preserved, except the displayed version/title
line.  Removing this delimited V2 appendix and reversing that line
recovers the predecessor byte for byte.  V2 replaces only the active
V1 source in \path{latex_notes}; the prior body remains recoverable
from V2.  All other source files and parallel branch files are
preserved.  Compilation and verification outputs remain outside
\path{latex_notes}.

\begin{center}
\fbox{\parbox{.9\textwidth}{\textbf{V2 status.}
The joint calculation is complete at the stated fixed, open
finite regulator.  Its extension to the physical vacuum and to
uniform continuum estimates is not proved.  The interacting
four-dimensional Yang--Mills mass gap remains open.}}
\end{center}
% END CONSOLIDATED V2 APPEND

% BEGIN CONSOLIDATED V3 APPEND -- 2026-09-19
\clearpage
\section{V3: a literature-driven native covariance bootstrap}
\label{r3:context}

\subsection{What is imported, and what is already in the archive}

The user's new literature sweep is a set of research suggestions, not
a list of mathematical hypotheses.  We read it in full and checked the
primary records and relevant formulations of the closest proposals.
The direction changes here: supplement V2's asymptotic calculation
with a nonperturbative, finite-regulator certificate problem for the
\emph{same compact Wilson measure}.  Do not begin another lineage or
repeat a spectral hierarchy that the archive already contains.

In particular \(C\), V84 already proves the forward Hankel hierarchy,
the direction of its Ritz bound, finite-record ambiguity, a
Chebyshev soft-weight estimate with errors, and the exact-flatness
exception.  \(C\), V85 already derives native score identities and
explains why a contact variance lower bound is not a propagating
lower Gram.  Neither construction is claimed anew below.

The literature screening has the following consequences.
\begin{itemize}
\item Abbott--Jay--Oare,
\href{https://arxiv.org/abs/2508.01377}{arXiv:2508.01377},
relate matrix Euclidean correlators to moment problems and bound
smeared spectral functions.  Abbott et al.,
\href{https://arxiv.org/abs/2605.20509}{arXiv:2605.20509},
formulate positive-measure optimization and kernel dual certificates.
Import the \emph{certificate design}, not an upper spectral endpoint
from one finite feasible moment matrix.

\item Cho et al.,
\href{https://arxiv.org/abs/2511.08560}{arXiv:2511.08560}
(first submitted November 2025; revised June 2026), use reflection
positivity and dynamical constraints for quantum-mechanical
two-point functions.  Kazakov--Zheng,
\href{https://arxiv.org/abs/2404.16925}{arXiv:2404.16925},
combine finite-\(N\) lattice loop equations and positivity; the
\(\SU(2)\) trace identities close their equations linearly on
single-trace loops.  Guo et al.,
\href{https://arxiv.org/abs/2502.14421}{arXiv:2502.14421},
extend the loop bootstrap to \(\SU(3)\) with multi-trace variables.
These are the closest imports.  Here we use a larger polynomial
coordinate algebra, so that completeness can be proved without
assuming a selected loop list is closed.

\item Bach--Ballesteros--Geisler,
\href{https://arxiv.org/abs/2508.07643}{arXiv:2508.07643},
construct a smooth Feshbach--Schur spectral renormalization flow.
Its operator class, complement invertibility, norm controls and
spectral-parameter transport would still have to be verified for
our Wilson family.  A semigroup identity alone does not bound an
accumulated return.  Keep this as a subsequent uniform-estimate route.

\item The hybrid Gibbs-chain comparison of Qin--Ju--Wang,
\href{https://arxiv.org/abs/2312.12782}{arXiv:2312.12782},
concerns a sampling operator.  The frustration-free certificate
hierarchy in
\href{https://quantum-journal.org/papers/q-2026-04-13-2065/}
{Quantum 10, 2065 (2026)}
has a different Hamiltonian hypothesis.
Both suggest ways to search for certificates; neither supplies the
physical Wilson transfer gap.

\item Cao--Nissim--Sheffield,
\href{https://arxiv.org/abs/2505.16585}{arXiv:2505.16585},
control loop equations and their merger contribution in area-law
regimes.  A restoration bound is essential before using a truncated
interaction.  Chevyrev--Shen,
\href{https://arxiv.org/abs/2503.03060}{arXiv:2503.03060},
prove uniqueness of the gauge-covariant renormalization for stochastic
three-dimensional Yang--Mills--Higgs.  This is relevant to later
identification of approximations, not a four-dimensional construction.
\end{itemize}
Concentration--compactness/rigidity, renewal Fourier arguments and
approximate-subgroup flattening from the supplied sweep are retained
as possibilities, not instantiated theorems: the required compactness,
no-collapse or native flattening hypotheses have not been proved.
This is a targeted primary-source audit of the suggested directions,
not an assertion that all recent literature has been exhausted.

General moment/SOS methods are established mathematics; see also
Magron--Forets--Henrion,
\href{https://arxiv.org/abs/1807.00754}{arXiv:1807.00754},
for polynomial invariant-measure hierarchies.
The contribution of this note is an explicit complete finite-Wilson
formulation, its connected-source and verification interfaces, and
a small populated certificate.  It is not a claim to invent
semidefinite optimization or Schwinger--Dyson identities.

\paragraph{Review cadence.}
The regular wide literature checkpoints are now V10, V20, and so on,
in addition to the companion-program review at V5, V10, and so on.
This V3 intake is an extra review.  At each literature checkpoint,
search across spectral/moment methods, lattice gauge theory, RG,
probability and functional inequalities, PDE/continuum construction,
and any newly relevant area; check original sources and record the
exact hypothesis mismatch before importing a result.  Carry the
append-count schedule through future V100 archive/restarts.
The next companion checkpoint is V5; the next regular wide sweep is V10.

\subsection{The native polynomial identities fix the finite measure}
\label{r3:SD}

Fix any finite set of integrated links with Wilson plaquettes, including
the open boxes of V2.  No weak-coupling or all-good restriction is
imposed.  Write each link as a unit quaternion
\(q_e=(q_{e0},q_{e1},q_{e2},q_{e3})\), so
\[
 K=(S^3)^E,\qquad
 g_e(q)=\sum_{\alpha=0}^3q_{e\alpha}^2-1=0 .
\]
The real polynomial
\[
 V(q)=\gamma\sum_p c_p(1-w(U_p)),\qquad
 d\mu=Z^{-1}e^{-V}\,dH^{\otimes E}
 \tag{V3.1}\label{r3:mu}
\]
has degree at most four.  Fixed boundary links, if any, are parameters,
not integrated coordinates.  In this section \(\gamma,c_p\) are fixed
finite real numbers, with the positive Wilson values used in applications.

Let \(D_{e,a}\), \(a=1,2,3\), be the derivative along
\(q_e\mapsto e^{t\mathbf i_a}q_e\).
Quaternion multiplication makes its coordinate coefficients linear.
It preserves polynomial degree, is tangent to \(S^3\), and is
divergence-free for Haar measure.  Define the first-order operator
\[
 {\cal B}_{e,a}p=D_{e,a}p-pD_{e,a}V .
 \tag{V3.2}\label{r3:B}
\]
This is a Schwinger--Dyson operator, not a Markov transition or a
physical-time generator.  Its polynomial degree is at most
\(\deg p+4\).

\begin{lemma}[Finite-law identification by the full identities]
\label{r3:identification}
A Borel probability \(\nu\) on \(K\) equals \(\mu\) if
\[
       \int_K{\cal B}_{e,a}p\,d\nu=0
       \quad\hbox{for every real polynomial \(p\) and every \(e,a\)}.
 \tag{V3.3}\label{r3:SDall}
\]
Conversely \(\mu\) satisfies these identities.
\end{lemma}
\begin{proof}
For \(\mu\), integrate
\(D_{e,a}(pe^{-V})\) against product Haar measure.
For the converse, polynomial restrictions are dense in \(C^1(K)\)
in the \(C^1\) norm.  To see the required approximation, extend a
smooth function to a neighborhood of the product of spheres and
approximate it and its first derivatives on a containing box by
multivariate Bernstein polynomials.  Restrict to \(K\).
Smooth functions suffice below, and this argument already gives
their required \(C^1\) approximation.
Thus \eqref{r3:SDall} holds for every smooth \(p\).
Take \(p=e^Vh\).  With \(d\eta=e^Vd\nu\), it gives
\(\int D_{e,a}h\,d\eta=0\) for all smooth \(h\).
For the corresponding one-parameter left translation \(\tau_t\),
differentiate \(\int h\circ\tau_t\,d\eta\); its derivative is zero
by the same identity applied to \(h\circ\tau_t\).
Hence \(\eta\) is invariant under these translations.
Their products generate left translations on \(\SU(2)^E\).
Uniqueness of finite left-invariant measure on this compact group
therefore gives \(\eta=cH^{\otimes E}\).
Normalization of \(\nu\) proves \eqref{r3:mu}.
No density or mixing rate for \(\nu\) was assumed.
\end{proof}

The use of \emph{all} polynomial tests matters.  A finite list of
loop equations may have additional feasible solutions.  The lemma
does not prove uniqueness of an infinite-volume Gibbs state.

\subsection{A finite-dimensional outer hierarchy with a completeness proof}
\label{r3:hierarchy}

Let \(x\) denote all \(4E\) coordinates.  For \(r\ge2\), let
\({\cal F}_r\) consist of real linear functionals \(\ell\) on
polynomials of degree at most \(2r\) satisfying
\begin{align}
 \ell(1)&=1,\qquad
 [\ell(x^{\alpha+\beta})]_{|\alpha|,|\beta|\le r}\succeq0,
 \notag\\
 \ell(g_eh)&=0\quad(\deg h\le2r-2),\notag\\
 \ell({\cal B}_{e,a}p)&=0\quad(\deg p\le2r-4).
 \tag{V3.4}\label{r3:hierarchy-constraints}
\end{align}
Imposing the identities on monomials imposes them on all the stated
polynomials.  These are finitely many linear equalities and one
positive-semidefinite matrix constraint.  Gauge and lattice symmetries
can be added when they preserve the declared boundary conditions.
Their use is optional for this theorem.

\begin{theorem}[Complete finite-regulator Wilson expectation bounds]
\label{r3:complete}
For every fixed real polynomial \(h\), and sufficiently large \(r\), put
\[
 \ell_r^-(h)=\min_{\ell\in{\cal F}_r}\ell(h),\qquad
 \ell_r^+(h)=\max_{\ell\in{\cal F}_r}\ell(h).
\]
These extrema exist and obey
\[
 \ell_r^-(h)\ \uparrow\ \E_\mu h,\qquad
 \ell_r^+(h)\ \downarrow\ \E_\mu h .
 \tag{V3.5}\label{r3:convergence}
\]
Every finite interval is an outer bound on the original Wilson
expectation.  No rate uniform in \(E,\gamma\), temporal extent or
lattice spacing is asserted.
\end{theorem}
\begin{proof}
The moments of \(\mu\) are feasible, so the sets are nonempty.
For \(|\beta|\le r-1\), the sphere identity multiplied by
\(x^{2\beta}\), and positivity of the diagonal moment entries, give
\[
 0\le\ell(x^{2\beta}x_i^2)
       \le\sum_{\alpha=0}^3
          \ell(x^{2\beta}q_{e\alpha}^2)
       =\ell(x^{2\beta}),
\]
where \(x_i\) is a coordinate of link \(e\).
Iterating from \(1\) bounds every diagonal moment by one.
Every monomial \(x^\alpha\), \(|\alpha|\le2r\), splits into
\(x^\beta x^\delta\) with \(|\beta|,|\delta|\le r\).
The \(2\times2\) principal moment matrix and Cauchy--Schwarz then give
\(\lvert\ell(x^\alpha)\rvert\le1\).
Thus \({\cal F}_r\) is closed and bounded in finite dimension.
Restriction from order \(r+1\) to \(r\) is feasible, proving
monotonicity and existence of extrema.

Choose any sequence of feasible functionals with orders tending to
infinity.  The moment bounds and a diagonal subsequence give a
functional \(L\) on all polynomials, positive on squares, with
\(L(1)=1\), all sphere identities and all identities \eqref{r3:SDall}.
Here is a direct representation argument to avoid assuming that
a truncated positive functional already has a representing measure.
Complexify and form the Hilbert completion of polynomials with
\(\langle p,q\rangle=L(\overline p q)\), modulo its null space.
The sphere relations imply
\[
 L(x_i^2|p|^2)\le L(|p|^2).
\]
Multiplication by \(x_i\) therefore descends to a bounded symmetric
contraction \(X_i\).  These operators are self-adjoint, commute, and
satisfy the sphere relations as operator identities.
Their joint spectral theorem, evaluated at the vector represented by
\(1\), gives a probability \(\nu\) supported on \(K\) representing \(L\).
Lemma~\ref{r3:identification} identifies it with \(\mu\).
Every convergent subsequence of optimizing functionals consequently
has the same \(h\)-moment, proving \eqref{r3:convergence}.
\end{proof}

This is an outer hierarchy, not a statement that its numerical solver
output is rigorous.  A floating optimum needs a certified feasible
dual witness or an independently validated optimization enclosure.
The theorem also does not promise an economical degree or matrix size.
For \(E\) links the unreduced order-\(r\) moment matrix has
\(\binom{4E+r}{r}\) rows.

\subsection{Connected sources require a joint two-copy problem}
\label{r3:replicas}

A connected covariance contains products of unknown means, which
are not linear moment objectives.  Introduce independent complete
histories \(U,V\) and the action
\[
               V^{(2)}(U,V)=V(U)+V(V).
\]
For real polynomial source banks \(A=(A_i)\), \(B=(B_j)\), possibly
at different time slices in one history, define
\[
 \Psi_{ij}(U,V)
       =\tfrac12[A_i(U)-A_i(V)][B_j(U)-B_j(V)].
 \tag{V3.6}\label{r3:replica}
\]
Expansion gives
\[
 \E_{\mu\otimes\mu}\Psi_{ij}
       =\E_\mu(A_iB_j)-\E_\mu A_i\,\E_\mu B_j.
 \tag{V3.7}\label{r3:replica-identity}
\]
Theorem~\ref{r3:complete}, on \(2E\) links, therefore gives convergent
finite semidefinite outer bounds on each connected covariance.
The Schwinger--Dyson tests for a derivative in either copy must
be allowed to depend on \emph{both} copies.

\begin{proposition}[Why marginal constraints alone can erase the source]
If only the two marginal laws are fixed, the same purported covariance
objective can have the wrong value even at infinite polynomial degree.
\end{proposition}
\begin{proof}
The diagonal coupling \(U=V\) has marginal law \(\mu\) in each copy,
but makes \(\Psi=0\).  For \(A=B=F\) nonconstant, the correct value
is \(\Var_\mu F>0\), since \(\mu\) has a strictly positive density.
The full joint identities instead identify the product law by
Lemma~\ref{r3:identification}, excluding this diagonal coupling.
\end{proof}

Thus independence here is proved by the joint constraint system in
the limiting problem, not imposed by nonlinear moment factorization
at a truncated level.  Each finite feasible set may still contain
spurious couplings; it remains an outer relaxation.

For \(A=B\), \(\Psi=\frac12\Delta A\,\Delta A^{\mathsf T}\) is
pointwise positive semidefinite.  For separated-time or reflected
sources that pointwise claim is generally false.  Reflection positivity
may add valid linear matrix constraints only for a boundary geometry
and reflection for which it is proved; it is not inferred by replacing
one factor of a positive matrix by a translated factor.

For a smooth, genuinely simultaneous-flow source, polynomial
approximation must precede this formulation.  If
\(\|A\|_\infty\le R_A\), \(\|B\|_\infty\le R_B\) in Euclidean bank norm,
and polynomial banks \(\widetilde A,\widetilde B\) have uniform errors
\(\epsilon_A,\epsilon_B\), then
\[
 \|\Cov(A,B)-\Cov(\widetilde A,\widetilde B)\|
 \le 2(R_A\epsilon_B+R_B\epsilon_A+\epsilon_A\epsilon_B).
 \tag{V3.8}\label{r3:source-error}
\]
Indeed apply the rank-one norm bound separately to
\(\E AB^{\mathsf T}\) and \(\E A\,(\E B)^{\mathsf T}\).
This applies to the original input variables, not a substituted
probability law for flowed links.  Polynomial density supplies
existence at each fixed regulator; quantitative physical-flow
approximation still needs its own bound.

\subsection{Exact dual witnesses and the locality target}
\label{r3:dual}

For a real polynomial objective \(h\), a sufficient certificate for
\(\E_\mu h\ge a\) is a polynomial identity
\[
 h-a=\sum_\nu s_\nu^2+\sum_e g_e t_e
                 +\sum_{e,b}{\cal B}_{e,b}v_{e,b}.
 \tag{V3.9}\label{r3:dual-identity}
\]
All coefficients and degrees are explicit; at hierarchy order \(r\),
require \(\deg s_\nu\le r\), \(\deg t_e\le2r-2\) and
\(\deg v_{e,b}\le2r-4\).
Applying any feasible \(\ell\) proves the bound.  Applying this to
\(a-h\) gives an upper certificate.  No strong-duality or
finite-degree completeness claim is required for this implication.

There is a matrix version for a symmetric polynomial objective
\(\Psi\), such as a symmetrized reflected bank:
\[
 \Psi-aS=\sum_\nu R_\nu^{\mathsf T}R_\nu+
           \sum_e g_e T_e+\sum_{e,b}{\cal B}_{e,b}W_{e,b}+{\cal E}.
 \tag{V3.10}\label{r3:matrix-dual}
\]
Here \(S\) is a declared constant coefficient metric, \(T_e,W_{e,b}\)
are symmetric matrix polynomials, and the operator acts entrywise.
If \(\|{\cal E}(q)\|\le\varepsilon\) on \(K\), integration gives
\[
                 \E\Psi\succeq aS-\varepsilon I.
 \tag{V3.11}\label{r3:matrix-bound}
\]
A fully explicit sufficient residual budget is
\(\varepsilon=\sum_\alpha\|{\cal E}_\alpha\|\), since
\(|q^\alpha|\le1\) on \(K\).  Exact rational Gram factors, or certified
positive-semidefinite Gram matrices, can establish the square terms.
A small solver residual without this positivity check is insufficient.
All statements apply as written to the two-copy action.

The proof is simply positivity of the square terms, vanishing of
the sphere ideal on \(K\), and the exact integration-by-parts
identities.  It avoids inversion of a near-singular source or
moment Gram.

\paragraph{A local certificate can be uniform without a global inverse.}
Choose any subset of differentiated links.  For each such link,
\(D_{e,b}V\) involves only incident plaquettes.  Consequently a
certificate \eqref{r3:matrix-dual} that uses only these derivatives,
includes every incident plaquette, and is pointwise valid for all
exterior link values remains valid in every larger finite Wilson
system with the same local interactions.  Terms of the outside action
have zero derivative in the chosen links.  Integrate the same identity
in the larger system to prove \eqref{r3:matrix-bound}.
Thus a genuinely local, exterior-uniform witness would supply a
volume-independent inequality.  The finite-order convergence theorem
does \emph{not} construct such a witness, and no positive-time
uniform witness is exhibited in this note.  This is the specific
locality problem to attack, not an assumption of spatial mixing.

\subsection{A populated exact certificate, not a numerical fit}
\label{r3:example}

On a single plaquette with open boundary, three tree links integrate
out exactly.  Let \(x=w(U)\).  Its native marginal has density
proportional to \(e^{\gamma x}\sqrt{1-x^2}\) on \([-1,1]\).
Ordinary integration by parts, whose boundary factor is
\((1-x^2)^{3/2}\), gives for every polynomial \(p\)
\[
 \E_\gamma{\cal A}_\gamma p=0,\qquad
 {\cal A}_\gamma p=(1-x^2)p'
                     +\{\gamma(1-x^2)-3x\}p .
 \tag{V3.12}\label{r3:radial}
\]
This is also obtained by summing the three link identities with
tests \((D_a x)p(x)\).  It is a consequence of the native measure,
not a fitted recurrence.

At \(\gamma=1\), take the rational polynomial
\[
             p(x)=\frac{-127+31x-6x^2+x^3}{400}.
\]
Direct expansion gives exactly
\[
                   x-{\cal A}_1p=\frac6{25}+\frac{x^5}{400}.
 \tag{V3.13}\label{r3:certificate}
\]
Since \(|x|\le1\), integration proves
\[
          \frac{19}{80}\le m:=\E_1x\le\frac{97}{400}.
\]
Taking \(p=1\) in \eqref{r3:radial} also gives
\(\E_1x^2=1-3m\).  The function \(1-3m-m^2\) decreases on this
positive interval.  We obtain the entirely rational enclosure
\[
 \boxed{\quad
   \frac{34191}{160000}\le\Var_1(x)=\Var_1(1-x)
                  \le\frac{1479}{6400}.\quad}
 \tag{V3.14}\label{r3:variance}
\]
This verifies a nonzero \emph{ordinary one-plaquette variance}.
It is not a two-time vacuum Gram or a mass-gap certificate.
The bound does not use Bessel-function numerics or V2's asymptotic
remainder.  The nonnegative residual bounds \(1\pm x^5\ge0\)
can, if desired, be inserted directly as valid localizers.

The diagnostic
\begin{center}
\path{scripts/verify_ym_restart_v3_bootstrap.py}
\end{center}
verifies \eqref{r3:certificate} and \eqref{r3:variance} with rational
arithmetic.  It constructs twenty such polynomial residual identities
at four rational couplings and five degrees, checks thirty-one
radial Haar identities and twenty-five two-copy covariance identities.
A Bessel-ratio comparison is separately labelled floating-point.
These are exact finite algebra checks, not verification of the
analytic completeness proof or a solved four-dimensional SDP.

\subsection{The spectral interface still has two distinct gates}
\label{r3:spectral}

The hierarchy above generates finite-history correlations of the
declared measure.  It must not silently turn the covariances of
different open temporal boxes in V2 into moments of one stationary
source measure.  Here is an exact bookkeeping identity.
Let the positive spatial Wilson transfer be
\(T=e^{-\gamma S_{\rm sp}/2}Ke^{-\gamma S_{\rm sp}/2}\),
and \(b(U)=e^{-\gamma S_{\rm sp}(U)/2}\).
An open history with \(N\) time steps has
\[
 Z_N=\langle b,T^Nb\rangle,\qquad
 \E_N[F(U_s)G(U_t)]
 =\frac{\langle b,T^s M_F T^{t-s}M_G T^{N-t}b\rangle}{Z_N},
 \quad 0\le s\le t\le N .
 \tag{V3.15}\label{r3:boundary}
\]
For real \(F,G\), this follows by writing the compact integral one
slab at a time.  Endpoint half weights from \(b\) combine with the
half weights in \(T\) to give one full spatial weight at every slice.
V1 proves that \(K\) includes all temporal links.  The identity also
specifies what happens when \(s=t\).
The source vector and normalizing partition function generally
change with \(N,s,t\).  Therefore these normalized data are not,
merely from this formula, a sequence \(J^*P^nJ\) for one fixed \(J\).
One must prove the vacuum preparation and centering, or use the
appropriate boundary/thermal spectral formulation.

Once that gate is separately paid, suppose
\(C_n=J^*P^nJ=\int_0^1x^n\,d\Sigma(x)\), with the same centered
source map \(J\) and positive contraction \(P\).
The sound finite-data dual rule is the following elementary one.
For a bounded real scalar kernel \(k\) and scalar polynomials
\(p_-\le k\le p_+\) on \([0,1]\),
\[
       \sum_j(p_-)_j C_j
       \preceq\int k\,d\Sigma
       \preceq\sum_j(p_+)_j C_j .
 \tag{V3.16}\label{r3:spectral-dual}
\]
It follows by applying the scalar inequalities to every positive
scalar measure \(z^*d\Sigma z\).
If certified input errors are \(\|\widehat C_j-C_j\|\le\delta_j\),
the upper endpoint gains
\(\sum_j|(p_+)_j|\delta_j I\), and the lower endpoint loses the
analogous amount.  Error ellipsoids from statistical data give
confidence statements, not unconditional deterministic enclosures.

Matrix-valued kernel certificates require care: if
\(B(x)=\sum_jx^jB_j\succeq K(x)\), the valid scalar conclusion is
\[
 \int\Tr(K(x)\,d\Sigma(x))
                   \le\sum_j\Tr(B_j C_j).
 \tag{V3.17}\label{r3:trace-dual}
\]
To justify it, dominate the matrix measure by its trace measure;
its Radon--Nikodym density is positive semidefinite, and the trace
of its product with \(B-K\) is nonnegative.
The product \(B_jC_j\) need not even be Hermitian.
No matrix Loewner conclusion is obtained by dropping the traces.

These rules explain the useful import from spectral bootstrapping,
but are not new spectral-exhaustion results.
A failed candidate localizer can refute an upper support edge.
Finite feasibility does not prove that the \emph{actual} measure
has no weight above it; \(C\), V84 already supplies the precise
counterexamples and the necessary growing-depth test.
Bounds on a positive amount of weight away from \(1\) alone likewise
do not exclude an additional soft component arbitrarily near \(1\).

\subsection{Result ledger and adjusted next calculation}
\label{r3:ledger}

\emph{Established here:} uniqueness of the finite compact Wilson law
from the complete polynomial Schwinger--Dyson family; a finite SDP
outer hierarchy converging to every fixed polynomial expectation;
a jointly constrained two-copy formulation converging to connected
covariances; scalar and matrix dual certificates with explicit
residual budgets; a local-certificate extension criterion; and an
exact rational one-plaquette variance enclosure.

\emph{Standard methods, not claimed new:} Haar integration by parts,
compact-group invariant measure, moment positivity and GNS
representation, polynomial optimization, the independent-copy
covariance identity, and spectral duality.  The existing V84
moment hierarchy is retained, not redeveloped.

\emph{Not established:} a practical certified four-dimensional
two-time SDP value, a hierarchy convergence rate uniform in coupling
or volume, a positive-time noncontact lower Gram, growing-depth
upper contraction, vacuum-boundary error bounds uniform under
refinement, a continuum construction, or the physical mass gap.

The immediate next task is to generate source-specific dual
witnesses for V2's smallest separated plaquette pair and compare their
certified finite-\(\gamma\) bounds with the proved expansion.
Use rational identities or validated matrix bounds, with every
incident plaquette retained.  Then test whether a witness can be
rewritten with fixed support and validity for arbitrary exterior
links.  If it cannot, expose the boundary-return term and return
to its RG/locality estimate; do not merely increase a global moment
matrix forever.  V2's local connected-coefficient calculation remains
the perturbative calibration, not a substitute for this remainder
and boundary control.

The full V2 body is preserved except its displayed title/version.
Removing this delimited appendix and restoring that title recovers
V2 byte for byte.  Only the active source is replaced in
\path{latex_notes}; archives and companion sources are unchanged.
The PDF skill is used only to inspect the compilation outside that
folder; the deliverable remains this \LaTeX\ source.

\begin{center}
\fbox{\parbox{.9\textwidth}{\textbf{V3 status.}
The literature import now has an exact native finite-measure
certificate problem and a completeness proof.  It has not supplied
the uniform physical source estimates or solved the interacting
four-dimensional Yang--Mills mass gap.}}
\end{center}
% END CONSOLIDATED V3 APPEND

% BEGIN CONSOLIDATED V4 APPEND -- 2026-09-19
\clearpage
\section{V4: a certified two-boundary Wilson source Gram}
\label{r4:context}

V3 is progress: it supplies a complete native finite-measure hierarchy,
but not a computed separated-source bound.  We now populate such a
bound on the \emph{same full \(1^4\) open Wilson box} used in V2.
No temporal plaquette, spatial interaction, or mixed term is removed.
The calculation covers the six elementary spatial plaquette sources
at both ends of its one temporal slab.

A direct unreduced implementation of V3 is an unattractive first
calculation.  This box has 32 links; two copies have 256 real
quaternion coordinates.  Already the degree-eight covariance objective
requires an order-four moment matrix with
\(\binom{260}{4}=186043585\) monomials.
An exact complete-history tree gauge reduces the unreduced count to
\(\binom{140}{4}=15329615\).  These are counts for the unreduced
constructions, not lower bounds on symmetry-reduced algorithms.
We instead exploit an exact positive kinetic compression whose
coefficients can be certified from the spatial cube character formula.
This is an analytic certificate in the positivity/duality direction
of V3, not a claim to have solved its large SDP.

The numerical certificates below are rational enclosures with proved
infinite-sum tails.  They are not Monte Carlo, floating-point fits, or
the unspecified constants in V2's asymptotic remainder.
The final comparison also proves that this particular certificate
loses too many weak-coupling powers to be a physical-scale lower Gram.

\subsection{The full two-slice law and its positive representation}
\label{r4:law}

Let \(\Gamma\) be the graph of one spatial cube: eight vertices,
twelve links, and six plaquettes.  Let \(U,V\in\SU(2)^{12}\) be its
two spatial slices.  For its six faces set
\[
 w_i(U)=\tfrac12\Tr U_{p_i},\qquad
 h_\gamma(U)=\exp\{\gamma\sum_{i=1}^6 w_i(U)\}.
\]
The loss convention \(1-w_i\) would multiply \(h_\gamma\) by a
constant \(e^{-6\gamma}\); all expressions below are unchanged.
Define
\[
 z_\gamma=\int_{\SU(2)}e^{\gamma w(g)}\,dH(g)
          =\frac{2I_1(\gamma)}{\gamma},
 \qquad
 k_\gamma(g)=e^{\gamma w(g)}/z_\gamma .
\]
Let \(C_\gamma\) be product convolution by \(k_\gamma\) on the twelve
spatial links, on product-Haar \(L^2\).
For the representation of dimension \(n\ge1\), its one-link
multiplier is
\[
 r_{n-1}(\gamma)=\frac{I_n(\gamma)}{I_1(\gamma)}>0.
 \tag{V4.1}\label{r4:multipliers}
\]
These follow by integrating the \(\SU(2)\) characters
\(\chi_{n-1}(\cos\theta)=\sin(n\theta)/\sin\theta\).
Thus \(C_\gamma\) is a positive self-adjoint contraction and preserves
constants.  Positivity follows from its positive character
multipliers, not merely from positivity of the pointwise kernel.

The temporal-link integral of V1 is the product convolution followed
by gauge averaging, times a scalar.  Product convolution commutes
with gauge averaging.  All \(h_\gamma w_i\) are gauge invariant, so
gauge averaging acts as the identity in the following scalar
products.  They are therefore representations of the original
gauge-invariant Wilson integrals, not substituted observables or
a reference-field law.

Put
\[
 {\cal Z}_\gamma=\langle h_\gamma,C_\gamma h_\gamma\rangle_H,\qquad
 m_i=\frac{\langle h_\gamma w_i,C_\gamma h_\gamma\rangle_H}
               {{\cal Z}_\gamma}.
\]
For the connected two-boundary matrix under the full \(1^4\) measure,
\[
 G_\gamma(i,j)
   =\Cov_{1^4,\gamma}(w_i(U),w_j(V)),
\]
one has exactly
\[
 G_\gamma(i,j)
 =\frac{\langle h_\gamma(w_i-m_i),
                 C_\gamma h_\gamma(w_j-m_j)\rangle_H}
           {{\cal Z}_\gamma}.
 \tag{V4.2}\label{r4:native-Gram}
\]
Indeed the marginal means agree by reflection of the slab; expansion
of the right side gives the connected covariance.  In particular
\(G_\gamma\succeq0\).
Replacing \(w_i\) by \(P_i=1-w_i\) leaves \(G_\gamma\) unchanged.

This is one \emph{lattice} time step with open ends.
It is not the infinite-time vacuum, a fixed positive physical time
as \(a\downarrow0\), or the archive's fully contact-cleaned flowed
source bank.  Separation of the two spatial link sets does not
by itself justify any of those identifications.

\subsection{A finite positive compression that eliminates unknown centering}
\label{r4:compression}

Consider the single-slice probability
\[
 d\nu_\gamma=h_\gamma\,dH/Z_\Gamma(\gamma),\qquad
 Z_\Gamma(\gamma)=\int h_\gamma\,dH,
\]
and write
\[
 \alpha_i=\E_{\nu_\gamma}w_i,\qquad
 \Sigma_{ij}=\Cov_{\nu_\gamma}(w_i,w_j),\qquad
 \lambda_\gamma=\left(\frac{I_2(\gamma)}{I_1(\gamma)}\right)^4,\qquad
 \rho_\gamma=\frac{Z_\Gamma(\gamma)^2}{Z_\Gamma(2\gamma)}.
 \tag{V4.3}\label{r4:parameters}
\]
The single-slice law enters as exact coefficients of a dual test;
it does not replace the two-slice law in \(G_\gamma\).

\begin{theorem}[Certified native two-boundary matrix lower bound]
\label{r4:bound}
For every \(\gamma>0\),
\[
 \boxed{\quad
 G_\gamma\succeq
 4\rho_\gamma\lambda_\gamma
 \Sigma\bigl(I_6+4\lambda_\gamma\alpha\alpha^{\mathsf T}\bigr)^{-1}
 \Sigma .\quad}
 \tag{V4.4}\label{r4:matrix}
\]
Every quantity on the right is a one-cube spatial integral.
The right side is positive definite.
\end{theorem}
\begin{proof}
The seven functions \(1,\phi_1,\ldots,\phi_6\), with
\(\phi_i=2w_i=\chi_1(U_{p_i})\), are orthonormal under product Haar.
Each plaquette holonomy is Haar after conditioning on its other
three edges, proving its mean zero and character norm one.
Two distinct faces have an edge belonging to only one of their
loops; integrating that edge kills their cross inner product.
Convolution on each of the four edges in a face multiplies its
fundamental matrix entries by \(I_2/I_1\).  Consequently
\[
 C_\gamma1=1,\qquad C_\gamma\phi_i=\lambda_\gamma\phi_i .
\]
The span is reducing for \(C_\gamma\), which is positive on its
orthogonal complement.

Fix a real coefficient vector \(u\), put \(F=\sum_i u_iw_i\), and
temporarily allow an arbitrary scalar \(m\) in \(f=h_\gamma(F-m)\).
The preceding compression gives
\[
 \langle f,C_\gamma f\rangle_H
 \ge Z_\Gamma(\gamma)^2
 \left\{(\alpha^{\mathsf T}u-m)^2+
 4\lambda_\gamma
 \bigl|\Sigma u+\alpha(\alpha^{\mathsf T}u-m)\bigr|^2\right\}.
\]
Write \(c=\alpha^{\mathsf T}u-m\).
Minimizing \(c^2+4\lambda_\gamma|\Sigma u+c\alpha|^2\) over \(c\)
gives
\[
 4\lambda_\gamma u^{\mathsf T}\Sigma
 (I_6+4\lambda_\gamma\alpha\alpha^{\mathsf T})^{-1}\Sigma u .
\]
This lower bound therefore holds for the actual, unknown two-slice
mean \(m=\sum_i u_im_i\); no equality of \(m_i\) and \(\alpha_i\)
has been assumed.
Finally
\[
 {\cal Z}_\gamma\le\|h_\gamma\|_H^2=Z_\Gamma(2\gamma)
\]
by the contraction property.  Divide the preceding inequality by
\({\cal Z}_\gamma\) to prove \eqref{r4:matrix}.

For strict positivity, \(\nu_\gamma\) has an everywhere positive
density.  If \(u^{\mathsf T}\Sigma u=0\), then \(\sum_i u_iw_i\)
is constant Haar-almost everywhere as well.  Haar orthogonality of
\(1,\phi_1,\ldots,\phi_6\) forces \(u=0\).
Thus \(\Sigma>0\); also \(\rho_\gamma,\lambda_\gamma>0\).
\end{proof}

This is a standard positive-compression/Schur-minimization argument
applied to a specified interacting Wilson integral.  It is not a
new general principle.  In particular a finite compression is used
to bound a quadratic form \emph{below}; it is not used to claim an
upper spectral support edge.

\subsection{Exact spatial moments, including the nonlocal face covariance}
\label{r4:cube}

Let
\[
 L_n(\gamma)=\frac{2I_n(\gamma)}{\gamma},\qquad n=1,2,\ldots.
\]
V2's exact six-face character formula, with the harmless
\(e^{-6\gamma}\) factor removed, is
\[
 Z_\Gamma(\gamma_1,\ldots,\gamma_6)
   =\sum_{n=1}^{\infty}n^2\prod_{i=1}^6L_n(\gamma_i).
 \tag{V4.5}\label{r4:Z}
\]
This is symmetric in the six face couplings even though not every
face permutation is a geometric cube symmetry.  At equal couplings,
therefore,
\[
 \alpha=\alpha_\gamma\mathbf1,\qquad
 \Sigma=(d_\gamma-t_\gamma)I_6+
                    t_\gamma\mathbf1\mathbf1^{\mathsf T}.
\]
Here \(d_\gamma\) is a one-face variance and \(t_\gamma\) is the
covariance of any two distinct faces in the \emph{spatial} cube.
To avoid numerical differentiation, define absolutely convergent sums
\begin{align*}
 Z&=\sum_n n^2L_n^6,&
 A&=\sum_n n^2L_n^5L_n',\\
 B&=\sum_n n^2L_n^4(L_n')^2,&
 D&=\sum_n n^2L_n^5L_n'' .
\end{align*}
Differentiating \eqref{r4:Z}, with independent face couplings before
setting them equal, gives
\[
 \alpha_\gamma=A/Z,\qquad
 t_\gamma=B/Z-(A/Z)^2,\qquad
 d_\gamma=D/Z-(A/Z)^2 .
 \tag{V4.6}\label{r4:moments}
\]
Factorial decay of the character coefficients and their derivatives
justifies termwise differentiation on compact positive coupling
intervals.  These are exact formulas, not infinite-volume
approximations.

Let \(\Pi_\parallel=\mathbf1\mathbf1^{\mathsf T}/6\) and
\(\Pi_\perp=I_6-\Pi_\parallel\).
The lower matrix in \eqref{r4:matrix} is
\[
 b_\perp(\gamma)\Pi_\perp+b_\parallel(\gamma)\Pi_\parallel,
 \tag{V4.7}\label{r4:bands}
\]
where
\[
 b_\perp=4\rho_\gamma\lambda_\gamma(d_\gamma-t_\gamma)^2,\qquad
 b_\parallel=
 \frac{4\rho_\gamma\lambda_\gamma(d_\gamma+5t_\gamma)^2}
       {1+24\lambda_\gamma\alpha_\gamma^2}.
 \tag{V4.8}\label{r4:explicit-bounds}
\]
In particular \(G_\gamma\succeq
\min(b_\perp,b_\parallel)I_6\).
This includes the requested opposite-time plaquette pair as a diagonal
entry, while also controlling every combination of the six sources.

\subsection{Rational tail bounds and reproducible enclosures}
\label{r4:enclosures}

The special-function input can be checked without evaluating a
floating-point Bessel function.
The defining series is
\[
 I_n(\gamma)=\sum_{k=0}^{\infty}
             \frac{(\gamma/2)^{n+2k}}{k!(n+k)!}.
 \tag{V4.9}\label{r4:Bessel}
\]
See \href{https://dlmf.nist.gov/10.25.E2}{DLMF 10.25.2};
the derivative recurrence is also recorded in
\href{https://dlmf.nist.gov/10.29.E2}{DLMF 10.29.2}.
For rational \(\gamma>0\), all series terms are positive rationals.
If terms through \(k=M\) are kept, the omitted tail is at most
\[
 \frac{a_{M+1}}{1-r_M},\qquad
 a_k=\frac{(\gamma/2)^{n+2k}}{k!(n+k)!},\qquad
 r_M=\frac{\gamma^2}{4(M+2)(n+M+2)}<1.
 \tag{V4.10}\label{r4:Bessel-tail}
\]
Indeed every ratio after the first omitted term is bounded by \(r_M\).

Comparing terms at the same \(k\) in \eqref{r4:Bessel} proves
\(L_{n+1}/L_n\le\gamma/[2(n+1)]\).
Differentiating the series, or its recurrence, gives
\begin{align}
 L_n'&=\frac{n-1}{\gamma}L_n+L_{n+1},\notag\\
 L_n''&=\frac{(n-1)(n-2)}{\gamma^2}L_n
               +\frac{2n-1}{\gamma}L_{n+1}+L_{n+2}.
 \tag{V4.11}\label{r4:derivatives}
\end{align}
All coefficients in these formulas are nonnegative for integer \(n\ge1\).
With \(b_n=(n+2)/\gamma+\gamma\), they imply
\[
 0\le L_n'\le b_nL_n,\qquad
 0\le L_n''\le b_n^2L_n .
\]
For the second inequality, the two shifted terms are bounded by
\(L_n\{(2n-1)/(2(n+1))+
\gamma^2/[4(n+1)(n+2)]\}\); the displayed \(b_n^2\) is larger.

Let \(a_n=n^2L_n^6\).  Then, for \(j=0,1,2\),
\[
 \frac{a_{n+1}b_{n+1}^j}{a_nb_n^j}
 \le
 \theta_n:=
 \frac{\gamma^6}{64n^2(n+1)^4}
              \left(\frac{n+3}{n+2}\right)^2 .
\]
The right side decreases with \(n\).  If representations \(n\le N\)
are summed and \(K=N+1\) satisfies \(\theta_K<1\), the remaining
tails of \(Z,A,B,D\) are bounded respectively by
\[
 \frac{a_K}{1-\theta_K},\quad
 \frac{a_Kb_K}{1-\theta_K},\quad
 \frac{a_Kb_K^2}{1-\theta_K},\quad
 \frac{a_Kb_K^2}{1-\theta_K}.
 \tag{V4.12}\label{r4:representation-tail}
\]
An upper enclosure of \(L_K\) may be used in \(a_K\).
These bounds justify every infinite sum used in the certificate.

The verifier performs rational interval arithmetic, rounding each
operation outward to the grid \(10^{-80}\mathbb Z\).
For each argument \(\beta\), it takes
\[
 N=\max(12,\lceil\beta\rceil+8),\qquad
 M=\max(40,2\lceil\beta\rceil+20),
\]
and uses \eqref{r4:Bessel-tail}--\eqref{r4:representation-tail}.
Arguments \(\beta=\gamma\) and \(2\gamma\) are both needed.
After evaluating \eqref{r4:explicit-bounds}, the following deliberately
rounded-down rational values are certified:
\begin{center}
\begin{tabular}{@{}cr@{}}
\toprule
\(\gamma\) & \(G_\gamma\succeq b_\gamma I_6\), with certified \(b_\gamma\)\\
\midrule
\(1/4\) & \(33869/10^{10}\)\\
\(1\)   & \(17023/10^8\)\\
\(4\)   & \(17223/10^{11}\)\\
\(16\)  & \(16377/10^{18}\)\\
\bottomrule
\end{tabular}
\end{center}
These are positive lower bounds, not estimates of the actual
smallest eigenvalue.  Both eigenvalues of the lower matrix are
enclosed before taking their minimum.
The full rational endpoints are recorded outside \path{latex_notes}
in \path{artifacts/ym_restart_v4_boundary_checks.json}.
Repeating all four couplings with four additional representation and
series terms gives compatible enclosures.  The rigorous status comes
from the inequalities and outward integer arithmetic, not from that
agreement.

\subsection{The certificate's precise weak-coupling loss}
\label{r4:loss}

The positivity just proved cannot be extrapolated into the desired
physical lower Gram.  The loss can be calculated, rather than merely
called nonuniform.

\begin{proposition}[A positive certificate with the wrong normalization scale]
\label{r4:power}
As \(\gamma\to\infty\), for this fixed spatial cube,
\[
 \min(b_\perp,b_\parallel)
 =\frac{2^8}{6^{3/2}\pi^{5/2}}\,
       \gamma^{-23/2}\{1+O(\gamma^{-1})\}.
 \tag{V4.13}\label{r4:lower-power}
\]
In contrast, the full native boundary matrix satisfies
\[
             G_\gamma=\gamma^{-2}G^{(0)}+O(\gamma^{-3}),
 \qquad \lambda_{\min}(G^{(0)})=1/48 .
 \tag{V4.14}\label{r4:actual-leading}
\]
All remainders in this proposition fix the box.
\end{proposition}
\begin{proof}
Apply V2's full-box expansion to the three-dimensional spatial cube.
It has five chord variables.  For the temporal-first tree convention
used by its diagnostic, the integer Hessian is
\[
 Q=\begin{pmatrix}
 2&0&0&0&-1\\
 0&2&-1&-1&1\\
 0&-1&2&1&-1\\
 0&-1&1&2&-1\\
 -1&1&-1&-1&2
 \end{pmatrix},
 \qquad \det Q=6 .
\]
Thus, with
\(\kappa=2^{5/2}/(6^{3/2}\pi^{5/2})\),
\[
 Z_\Gamma(\gamma)=e^{6\gamma}\kappa\gamma^{-15/2}
                         \{1+O(\gamma^{-1})\},\qquad
 \rho_\gamma=\kappa\,2^{15/2}\gamma^{-15/2}
                         \{1+O(\gamma^{-1})\}.
\]
The spatial plaquette circulation covariance is the rank-five
projection whose diagonal entries are \(5/6\) and whose off-diagonal
entries have magnitude \(1/6\).  V2's differentiated expansion gives
\[
 \alpha_\gamma=1-\frac5{4\gamma}+O(\gamma^{-2}),\quad
 d_\gamma=\frac{25}{24\gamma^2}+O(\gamma^{-3}),\quad
 t_\gamma=\frac1{24\gamma^2}+O(\gamma^{-3}).
\]
The one-link Laplace expansion gives \(\lambda_\gamma=1+O(\gamma^{-1})\).
Consequently
\[
 b_\perp=4\kappa\,2^{15/2}\gamma^{-23/2}
                   \{1+O(\gamma^{-1})\},\qquad
 b_\parallel=\tfrac14\kappa\,2^{15/2}\gamma^{-23/2}
                   \{1+O(\gamma^{-1})\}.
\]
Their minimum proves \eqref{r4:lower-power}.

For completeness the full four-dimensional leading matrix can be
evaluated exactly by V2's incidence formula
\(G^{(0)}_{ij}=\frac32 M_{p_iq_j}^2\).
Order the faces on either spatial slice as
\[
 (12,0),\ (13,0),\ (23,0),\
 (12,x_3=1),\ (13,x_2=1),\ (23,x_1=1),
\]
where the first three are based at the spatial origin.
The diagonal entries of \(G^{(0)}\) are \(3/128\).
Opposite spatial faces have entries \(1/384\), and the other
off-diagonal entries are \(1/1536\).
If \(O\) swaps the three opposite pairs and
\(J=\mathbf1\mathbf1^{\mathsf T}\), this is
\[
                 G^{(0)}=(35I_6+3O+J)/1536 .
\]
The constant vector has eigenvalue \(11/384\).
The pair-symmetric, sum-zero subspace has eigenvalue \(19/768\)
and dimension two.  The pair-antisymmetric subspace has eigenvalue
\(1/48\) and dimension three.  This proves
\eqref{r4:actual-leading}.  The finite matrix remainder follows
entrywise from V2 and hence also in operator norm.
\end{proof}

The certificate-to-true-leading-scale ratio is therefore of order
\(\gamma^{-19/2}\), and tends to zero even \emph{before} a continuum
or volume limit is attempted.  Normalizing the source by \(\gamma\)
does not repair it.

Two identifiable losses appear in \eqref{r4:matrix}.
The Schur numerator squares a single-slice covariance of order
\(\gamma^{-2}\), producing order \(\gamma^{-4}\).
The global Haar normalization comparison contributes
\(\rho_\gamma\asymp\gamma^{-15/2}\).
A fixed unweighted character bank poorly resolves a measure
concentrating in its five three-color chord directions.
This is a concrete manifestation of the extensive reference-space
issue already flagged in the f15 archive; it is not a new claim
that all character methods must fail.

The normalization loss generalizes: for any fixed contractible
spatial box with \(b\) independent chords, V2 gives
\[
 \frac{Z_\Gamma(\gamma)^2}{Z_\Gamma(2\gamma)}
             \asymp_\Gamma \gamma^{-3b/2}.
\]
Here \(Z_\Gamma\) uses the exponential plaquette convention above,
so its extensive constant factors cancel in the ratio.
The exponent itself grows with the spatial graph.
Thus this \emph{global Haar comparison} cannot be advertised as a
volume-independent local certificate.

For a separate check at the opposite end of coupling space,
the positive series yields
\(\lambda_\gamma\sim\gamma^4/256\),
\(\rho_\gamma\to1\), \(\alpha_\gamma\to0\),
\(d_\gamma\to1/4\), and \(t_\gamma\to0\).
Both certificate eigenvalues are therefore asymptotic to
\(\gamma^4/1024\) as \(\gamma\downarrow0\).
This consistency check does not authorize continuation of a
strong-coupling physical clock to weak coupling.

\subsection{Verification, limits of the result, and the upstream move}
\label{r4:ledger}

The diagnostic is
\begin{center}
\path{scripts/verify_ym_restart_v4_boundary.py}.
\end{center}
It evaluates four rational-coupling enclosures and repeats them at
larger cutoffs, checks eighteen exact Schur identities, and obtains
the four-dimensional leading Gram from the literal incidence
construction.  Its three symmetry-type eigenvectors verify the
listed leading eigenvalues; the displayed decomposition gives their
multiplicities.  No floating-point special-function evaluation is
used in any certificate.
The proofs above, not the diagnostic's success flag, identify the
integrals and establish the infinite-sum bounds.

\emph{Newly paid:} an explicit positive lower matrix for all six
opposite-time plaquette combinations under the full \(1^4\) Wilson
measure at every fixed \(\gamma>0\); certified rational numerical
instances; exact handling of the unknown two-slice centering; and
a quantified explanation of why this certificate does not preserve
the actual weak-coupling covariance scale.

\emph{Not paid:} the archive's contact-cleaned, positive-physical-time
source lower bound, a volume-uniform certificate, fixed-physical-flow
control, growing-depth upper contraction, uniform vacuum preparation,
thermodynamic/continuum transport, or a physical Hamiltonian gap.
One-step positivity of \eqref{r4:native-Gram} proves none of these.

The next step must remove the global reference normalization, not
just increase a small list of character tests.
Seek a source-weighted local Schwinger--Dyson witness or a normalized
conditional boundary-message estimate whose constants are measured
in the native law.  Keep all plaquettes incident to differentiated
links.  A certificate uniform in exterior data would use V3's local
extension argument; the present character certificate does not have
that property.  If a boundary-return term prevents localization,
estimate that term directly before introducing further bootstrap
depth.  The six-source matrix and its exact \(\gamma^{-2}\) leading
scale now give a nontrivial test any proposed repair must pass.

V5 is the next companion-program checkpoint.  The next regular wide
literature sweep remains V10.  The V3 body is retained byte for byte
apart from the title/version line; removing this appendix and
reversing that line recovers it.  Only the active source is replaced.
All compilation and interval-check artifacts stay outside
\path{latex_notes}.

\begin{center}
\fbox{\parbox{.9\textwidth}{\textbf{V4 status.}
A genuine finite four-dimensional two-boundary source certificate
has been populated and its weak-coupling loss identified.
The uniform physical source estimates and the full interacting
continuum Yang--Mills mass gap remain unproved.}}
\end{center}
% END CONSOLIDATED V4 APPEND

% BEGIN CONSOLIDATED V5 APPEND -- 2026-09-19
\clearpage
\section{V5: native midplane scores and removal of the reference-volume loss}
\label{r5:context}

V4 produced a genuine positive matrix but lost the native
\(\gamma^{-2}\) covariance scale.  We now change the certificate,
not the Wilson measure.  Conditioning a reflection-symmetric history
on its actual middle slice turns its opposite-boundary covariance
into an ordinary covariance of conditional expectations.  Smooth
scores of that \emph{marginal law} give a lower matrix.  On the full
\(2\times1\times1\times1\) open box we construct six explicit vector
fields for which this certificate captures the entire leading
\(\gamma^{-2}\) matrix.

This is a two-lattice-step, fixed-box result.  It is neither V4's
one-step integral nor a stationary vacuum correlation.
The time separation is \(2a\), not a fixed positive physical time.
No continuum, uniform-volume, or mass-gap conclusion is inferred.
The result is a repair of the certificate's normalization, not a
new claim that fixed-box asymptotics alone solve the programme.

\subsection{The scheduled companion review and the nonduplication check}
\label{r5:review}

The current companion files were inspected read-only:
\begin{itemize}
\item \path{Renewal_SM_Einstein_Continuum_Research_Notes_V72.tex}.
Its V72 derives a classical covariance flow in an explicitly specified
many-copy model, at fixed finite mode cutoff and fixed times.
The empirical covariance bounds and transport argument retain their
declared scaling.  They supply no cutoff-uniform estimate for the
Wilson marginal used here.
\item \path{Renewal_NS_Stability_Research_Notes_V26.tex}.
Its latest section reduces rank-one Oseen-kernel derivative norms
to scalar angular maximizations.  These are different operators and
different measures.  No kernel constant is imported into YM; this
review is not a recertification of its numerical comparisons.
\item \path{Renewal_Yang_Mills_Shell_Mixing_Research_Notes_V16.tex}.
Its fully assembled flat-holonomy one-loop coefficient retains the
Haar and root-return terms and has a stated summable refinement
error.  Its own conclusion excludes a common nonperturbative
coupling/volume remainder and physical-time contraction.
We retain its normalization discipline, not its coefficient as
a bound on the present source matrix.
\end{itemize}
The corresponding SHA--256 records are, respectively,
\begin{center}\scriptsize
\path{F44F9745D43379E1672C5550411B58E97195A4C9278592D221DF5D3F644DC1F5}\\
\path{82C887F8C2C69E4F28FAD7C3DC8DE5B9099CB5A4BF431EBA1FFF3E91935978AD}\\
\path{BF138E63BA826BC030D3A8B87E04FDD3B1EC712D269BB064674E96509B46E710}.
\end{center}

Several ingredients are already available and are not presented as
new general principles.  The f15 V4 retained-score bound uses
conditional expectation and Cauchy--Schwarz.  The f16 common branch
contains conditional-message estimates; f16.2 V85 proves native
local score floors and explains their reflected contact nullity.
V2 of this restart already proves the full finite-box Laplace
expansion.  What is added below is a specified \emph{midplane
marginal} score bank, its exact noncontact boundary-response
identity, and its leading-order saturation on the full two-step
four-dimensional integral.

Covariance lower bounds from integration by parts and
Cauchy--Schwarz are standard.  A targeted primary-source check
found M.~Ernst, G.~Reinert and Y.~Swan,
\emph{First order covariance inequalities via Stein's method},
Bernoulli \textbf{26} (2020), 2051--2081,
\href{https://arxiv.org/abs/1906.08372}{arXiv:1906.08372}.
That paper treats univariate targets; no theorem from it is
assumed to apply to the Wilson history.  The finite-dimensional
matrix argument required here is proved directly.
This targeted check does not replace the scheduled broad V10 sweep.
The next five-version companion checkpoint is also V10.

\subsection{Exact sewing of the full two-step history}
\label{r5:sewing}

Use V2's open-box convention on
\[
 \Lambda=\{0,1,2\}\times\{0,1\}^{3}.
\]
There are 24 vertices, 52 links and 42 plaquettes.
The action is \(S=\sum_p(1-\tfrac12\Tr U_p)\), with every link
integrated and every plaquette retained.

Choose the same rooted seven-edge spatial tree on all three
slices, together with the two temporal links at the spatial root.
Their union is a complete 23-edge history tree.  The exact tree
change of variables of V2 leaves 29 Haar chord variables:
five middle spatial chords \(x\), twelve left variables \(z_-\),
and twelve right variables \(z_+\).
Each side contains five spatial chords and seven temporal chords.
After reflecting the right side and inverting its temporal link
orientations, the action has the exact form
\[
 S(x,z_-,z_+)=S_0(x)+S_1(x,z_-)+S_1(x,z_+).
 \tag{V5.1}\label{r5:split}
\]
Here \(S_0\) consists of the six middle spatial plaquettes;
\(S_1\) consists of the six outer spatial and twelve temporal
plaquettes on one side.  The final common root conjugation has
not been divided out.  Haar tree integration is an exact identity,
not a gauge-determinant approximation.

Define
\[
 Z_1(x)=\int e^{-\gamma S_1(x,z)}\,dH(z),\qquad
 d\nu_{\gamma,x}(z)=Z_1(x)^{-1}e^{-\gamma S_1(x,z)}\,dH(z),
\]
and the actual middle-slice marginal
\[
 d\mu_\gamma(x)=p_\gamma(x)\,dH(x),\qquad
 p_\gamma(x)=
 \frac{e^{-\gamma S_0(x)}Z_1(x)^2}
      {\int e^{-\gamma S_0}Z_1^2\,dH}.
 \tag{V5.2}\label{r5:marginal}
\]
These are smooth strictly positive densities on their compact groups
at every finite \(\gamma>0\).
Let \(w_i(z)\), \(1\le i\le6\), be the outer spatial face traces
\(\tfrac12\Tr U_{p_i}\), in the face order of V4.  Put
\[
 m_i(x)=\E_{\nu_{\gamma,x}}w_i,\qquad
 G_\gamma^{[2]}(i,j)=\Cov_\Lambda(w_i(z_-),w_j(z_+)).
\]

\begin{lemma}[The native conditional-expectation Gram]
\label{r5:message-Gram}
Exactly,
\[
                 G_\gamma^{[2]}=\Cov_{\mu_\gamma}(m,m).
 \tag{V5.3}\label{r5:message}
\]
\end{lemma}
\begin{proof}
The product form \eqref{r5:split}, with its original Haar measure,
makes \(z_-\) and \(z_+\) conditionally independent given \(x\),
with the same conditional law.  Thus their product expectation is
\(\E_{\mu_\gamma}m_i m_j\), and their marginal expectations are
\(\E_{\mu_\gamma}m_i\) and \(\E_{\mu_\gamma}m_j\).
Subtract the product of means.
\end{proof}

In particular this is positivity in the original four-dimensional
measure.  No global comparison to product Haar is used.
Nevertheless the normalization \(Z_1(x)\) and its derivatives
remain part of the problem; they cannot be discarded.

\subsection{A native score certificate and the contact distinction}
\label{r5:score}

For a smooth vector field \(X\) on \(\SU(2)^5\), write
\(\operatorname{div}_H X\) for its divergence relative to product
Haar, and define its marginal score
\[
 s_X=-\operatorname{div}_H X-X\log p_\gamma
     =-\operatorname{div}_H X+\gamma XS_0
                +2\gamma\E_{\nu_{\gamma,x}}XS_1 .
 \tag{V5.4}\label{r5:score-def}
\]
There is no boundary term on the compact manifold.
Hence \(\E_{\mu_\gamma}s_X=0\), and for smooth \(f\),
\[
                 \E_{\mu_\gamma}(f s_X)=\E_{\mu_\gamma}Xf.
 \tag{V5.5}\label{r5:ibp}
\]
The last term in \eqref{r5:score-def} is the actual conditional
boundary response.  Dropping it changes the score.

Choose six such fields \(X_j\), let \(s=(s_{X_j})_{j=1}^6\), and set
\[
 B_\gamma(i,j)=\E_{\mu_\gamma}X_jm_i,\qquad
 J_\gamma(i,j)=\E_{\mu_\gamma}s_{X_i}s_{X_j}.
\]
\begin{theorem}[Exact native lower matrices and their residual]
\label{r5:certificate}
For every real \(6\times6\) matrix \(W\),
\[
 G_\gamma^{[2]}\succeq
 B_\gamma W^{\mathsf T}+WB_\gamma^{\mathsf T}
                                      -WJ_\gamma W^{\mathsf T}.
 \tag{V5.6}\label{r5:dual}
\]
The difference is exactly the Gram of
\(m-\E_{\mu_\gamma}m-Ws\).
If \(J_\gamma>0\), optimizing \(W\) gives
\(G_\gamma^{[2]}\succeq B_\gamma J_\gamma^{-1}B_\gamma^{\mathsf T}\).
Without any inverse, take \(W=\gamma^{-1}I_6\) and define
\[
 L_\gamma=\gamma^{-1}(B_\gamma+B_\gamma^{\mathsf T})
                                  -\gamma^{-2}J_\gamma .
 \tag{V5.7}\label{r5:L}
\]
Then \(G_\gamma^{[2]}\succeq L_\gamma\) for every \(\gamma>0\).
Moreover the response entries have the exact formula
\[
 B_\gamma(i,j)=
 -\gamma\E_{\mu_\gamma}
       \Cov_{\nu_{\gamma,x}}(w_i,X_jS_1).
 \tag{V5.8}\label{r5:response}
\]
\end{theorem}
\begin{proof}
Equation \eqref{r5:ibp} gives
\(\E[(m-\E m)s^{\mathsf T}]=B_\gamma\).
Expand the positive Gram of \(m-\E m-Ws\).
Completing its matrix square proves the assertion about the inverse.
For \eqref{r5:response}, differentiate the numerator and denominator
of \(m_i\) under the compact conditional integral.
The observable \(w_i(z)\) does not depend on the middle chords,
so its derivative vanishes, leaving precisely the displayed
conditional covariance.
\end{proof}

The entries of \(B_\gamma\) involve half-history responses, and
those of \(J_\gamma\) involve the marginal scores.  Neither is
defined to be the target opposite-boundary covariance.
The statement is an exact certificate formula; by itself it is
not a numerical enclosure of these entries.

\paragraph{Why this is not V85's contact-score floor in disguise.}
Let the full-history middle-coordinate score be
\[
 T_X=-\operatorname{div}_H X+
       \gamma X(S_0+S_1(x,z_-)+S_1(x,z_+)).
\]
It satisfies \(\E_\Lambda[w_i(z_-)T_X]=0\), because the outer
observable has zero middle derivative.
On the other hand \(\E[T_X\mid x]=s_X\).
Conditional covariance decomposition gives
\[
 0=\E_{\mu_\gamma}[m_i s_X]+
       \gamma\E_{\mu_\gamma}
                     \Cov_{\nu_{\gamma,x}}(w_i,XS_1).
 \tag{V5.9}\label{r5:contact}
\]
The second term is generally nonzero.  Thus replacing \(T_X\) by
its marginal projection is not allowed inside a mixed expectation
without retaining this response.  Equations
\eqref{r5:response}--\eqref{r5:contact} both prove and locate the
difference from a full local divergence that pairs to zero with a
disjoint source.  We do not infer a positive-time lower bound from
the ordinary variance of that contact divergence.

\subsection{The calibrated fields and the conditional Gaussian limit}
\label{r5:fields}

We next specify the fields, with no fitted constants.
In complete-history tree coordinates let \(D\) be the signed
plaquette--chord incidence matrix and \(Q=D^{\mathsf T}D\).
It is positive definite by V2.
Write the one-color tangent coordinates as \((x,z)\), with
\(x\in\R^5\), \(z=(z_-,z_+)\in\R^{24}\), and partition \(Q\).
Define
\[
 K=-Q_{zz}^{-1}Q_{zx},\qquad
 R=Q_{xx}-Q_{xz}Q_{zz}^{-1}Q_{zx}>0,\qquad C=R^{-1}.
 \tag{V5.10}\label{r5:Schur}
\]
The left--right block of \(Q_{zz}\) is zero.
If \(d_i=(d_{i,x},d_{i,z})\) is an outer plaquette row, let
\[
 \ell_i=d_{i,x}+d_{i,z}K,\qquad
 A_i=\ell_i^{\mathsf T}\ell_i,\qquad
 V_i=-\tfrac12 C A_i .
 \tag{V5.11}\label{r5:V}
\]
The same \(\ell_i\) is obtained from the reflected outer face.
These matrices act on chord indices and identically on all three
Lie-algebra colors.

For the compact fields, use the quaternion exponential coordinates
\(U_r=\exp(x_r)\) near the identity, with \(x_r\in\R^3\).
Choose a smooth cutoff \(\eta\) equal to one on a sufficiently
small ball, supported strictly inside that chart, and invariant
under common color rotations.  Set
\[
           X_i(x)=\eta(x)(V_i\otimes I_3)x
 \tag{V5.12}\label{r5:compact-field}
\]
in this coordinate chart and extend by zero.
This defines a smooth global vector field, not a replacement of
the compact integration domain by \(\R^{15}\).
It is equivariant under the remaining common root conjugation,
so its score and the relevant functions are gauge invariant.

Let \(Y^1,Y^2,Y^3\) be independent \(N(0,C)\) vectors, and put
\[
 \phi_i(Y)=-\frac12\left\{
      \sum_{a=1}^3 (Y^a)^{\mathsf T}A_iY^a-3\Tr(CA_i)\right\},
 \qquad
 H_{ij}=\frac32(\ell_iC\ell_j^{\mathsf T})^2 .
 \tag{V5.13}\label{r5:H}
\]
Wick's formula gives \(\E\phi_i\phi_j=H_{ij}\).

\begin{lemma}[Conditional Laplace control on the declared fixed box]
\label{r5:conditional-limit}
As \(\gamma\to\infty\), under the original middle marginal,
the scaled logarithms \(Y=\sqrt\gamma x\) converge with every fixed
polynomial moment to the three-color Gaussian above.
Using these logarithms on the identity chart, the following
differences tend to zero in \(L^2(\mu_\gamma)\):
\[
 \gamma(m_i-\E_{\mu_\gamma}m_i)-\phi_i(\sqrt\gamma x),
 \qquad s_{X_i}-\phi_i(\sqrt\gamma x).
 \tag{V5.14}\label{r5:L2}
\]
A fixed smooth continuation outside the chart, or zero continuation
of the displayed Gaussian polynomials, gives the same limits.
\end{lemma}
\begin{proof}
We give the conditional justification explicitly; joint weak
convergence alone would not imply convergence of conditional means.

At \(x=0\), the half action \(S_1(0,z)\) has a unique zero at
\(z=0\) in the complete tree gauge.  Flatness on that contractible
half box forces every remaining chord to the identity.
Its Hessian in \(z\) is positive definite.
Compactness separates its complement from the minimum.
The implicit-function theorem and positive Hessian therefore give
a unique smooth conditional minimizer \(z_*(x)\) for small \(x\),
with \(z_*(x)=K_{\rm side}x+O(|x|^2)\).
All other conditional configurations are separated by a positive
action cost, uniformly in a sufficiently small closed \(x\)-ball.

Translate to this minimizer in the local integral.  Taylor
expansion, a positive quadratic lower bound, and a smooth cutoff
give the usual Laplace expansion uniformly there, also after one
\(x\)-derivative.  Its coefficient and remainder bounds follow by
integrating polynomials against a fixed Gaussian envelope;
the compact complement is exponentially small.
This is V2's localization argument with \(x\) as a smooth parameter,
not an exchange of a volume limit with a conditional expectation.

Write \(s_*(x)=S_1(x,z_*(x))\).
The effective action
\[
                       W(x)=S_0(x)+2s_*(x)
\]
has Hessian \(R\otimes I_3\).  The two conditional determinants
supply a smooth strictly positive amplitude.
The marginal has a unique global minimum: any zero of the
minimized full action would give a flat complete history and
hence \(x=0\).  Compactness again separates the complement.
Thus its normalized local density, in \(Y=\sqrt\gamma x\),
converges to the stated Gaussian, with Gaussian bounds on a fixed
small chart and exponentially small mass outside it.  In
particular all fixed scaled moments are uniformly bounded.

For an outer trace,
\[
 w_i=1-\tfrac12\left|\sum_r d_i(r)X_r\right|^2+O(|X|^3).
\]
The conditional Gaussian mean of its circulation is \(\ell_iY\).
Its conditional covariance contributes a constant
\(3d_i^{\,*}\), where
\(d_i^{\,*}=d_{i,z}Q_{zz}^{-1}d_{i,z}^{\mathsf T}\).
Uniform conditional Laplace expansion on bounded \(Y\)-sets gives
\[
 \gamma(m_i-1)\longrightarrow
       -\tfrac12\left(\sum_a|\ell_iY^a|^2+3d_i^{\,*}\right).
\]
Centering removes the constant \(3d_i^{\,*}\) as well as the
Gaussian mean.

For scores, the logarithmic derivative of the marginal expansion
gives, uniformly on bounded scaled sets,
\(\nabla_x\log p_\gamma=-\gamma(R\otimes I_3)x+O(1)\).
Terms from the conditional amplitude and the Haar density are
included in the error at these sets.
Also \(\operatorname{div}_H X_i=3\Tr V_i+o(1)\).
Consequently
\[
 s_{X_i}\longrightarrow
 -3\Tr V_i+\sum_a(Y^a)^{\mathsf T}V_i^{\mathsf T}RY^a.
\]
By \eqref{r5:V}, \(V_i^{\mathsf T}R=-A_i/2\) and
\(-3\Tr V_i=3\Tr(CA_i)/2\), proving the second limit
pointwise on bounded scaled sets.

For the claimed \(L^2\) statements, shrink the fixed chart if
necessary.  Conditional quadratic concentration bounds
\(\gamma|m_i-1|\) by \(C(1+|Y|^2)\) there.
The uniform differentiated Laplace formula and \(X_i=O(|x|)\)
similarly bound \(|s_{X_i}|\) by \(C(1+|Y|^2)\) on its inner chart.
On the cutoff annulus the score is bounded by \(C(1+\gamma)\):
this also follows directly from \eqref{r5:score-def}, since the
actions and their derivatives are bounded on compact groups.
That annulus has exponentially small marginal mass.
Outside the chart \(|\gamma(m_i-1)|\le2\gamma\), with the same
exponential mass estimate.  The Gaussian envelope and these
bounds supply uniform integrability of squares (indeed higher
fixed moments), allowing bounded-scaled-set convergence to pass
to \(L^2\).  This proves \eqref{r5:L2}.
All constants in this argument fix this box and this chart.
\end{proof}

\begin{theorem}[The native score lower matrix has the correct leading scale]
\label{r5:saturation}
For the actual compact fields \eqref{r5:compact-field},
\[
 \gamma B_\gamma\longrightarrow H,\qquad
 J_\gamma\longrightarrow H,\qquad
 \gamma^2L_\gamma\longrightarrow H,\qquad
 \gamma^2G_\gamma^{[2]}\longrightarrow H
 \tag{V5.15}\label{r5:limits}
\]
in matrix norm.  More precisely,
\[
 0\preceq G_\gamma^{[2]}-L_\gamma=o(\gamma^{-2}).
 \tag{V5.16}\label{r5:saturation-eq}
\]
\end{theorem}
\begin{proof}
The two limits in \eqref{r5:L2} have the same leading centered
quadratic \(\phi\).  Cauchy--Schwarz therefore gives convergence
of their pairings to \(H\).  Use
\(B_\gamma=\E[(m-\E m)s^{\mathsf T}]\) and
\(J_\gamma=\E ss^{\mathsf T}\) to obtain the first two limits.
Equation \eqref{r5:L} gives the third, and
\eqref{r5:message} gives the fourth.
Finally the exact residual in Theorem~\ref{r5:certificate}, with
\(W=\gamma^{-1}I\), is
\[
 \gamma^2(G_\gamma^{[2]}-L_\gamma)
 =\E\bigl[(\gamma(m-\E m)-s)(\gamma(m-\E m)-s)^{\mathsf T}\bigr],
\]
which tends to zero by the same \(L^2\) result.
\end{proof}

This is an asymptotically sharp native lower certificate, with no
factor \(Z_\Gamma(\gamma)^2/Z_\Gamma(2\gamma)\).
Its finite-\(\gamma\) entries still require bounds on the
conditional responses.  No explicit numerical onset \(\gamma_0\)
has been established by this argument.

\subsection{Exact rational calibration on the full four-dimensional box}
\label{r5:exact}

To make the construction reproducible, enumerate vertices and links
lexicographically with time first.  Build a spanning tree by taking
spatial links first in that order, then temporal links.
The five middle chord labels, written as
\((\hbox{spatial base};\hbox{direction})\), are
\[
 ((0,1,0);3),\ ((1,0,0);2),\ ((1,0,0);3),\
 ((1,0,1);2),\ ((1,1,0);3).
\]
The spatial directions are numbered \(1,2,3\).
The incidence/Schur computation gives
\[
 R=\frac1{35}
 \begin{pmatrix}
 94&2&-4&-2&-45\\
 2&94&-47&-45&47\\
 -4&-47&94&47&-45\\
 -2&-45&47&94&-47\\
 -45&47&-45&-47&94
 \end{pmatrix},\qquad
 \begin{pmatrix}\ell_1\\\ell_2\\\ell_3\\\ell_4\\\ell_5\\\ell_6\end{pmatrix}
 =\frac1{35}
 \begin{pmatrix}
 0&6&0&1&0\\
 -1&0&6&0&1\\
 6&1&-1&-1&1\\
 0&1&0&6&0\\
 -6&0&1&0&6\\
 1&6&-6&-6&6
 \end{pmatrix}.
 \tag{V5.17}\label{r5:matrices}
\]
These data specify all six \(V_i=-R^{-1}A_i/2\) exactly.
For example
\[
 V_1=\frac1{26460}
 \begin{pmatrix}
 0&42&0&7&0\\
 0&-246&0&-41&0\\
 0&-42&0&-7&0\\
 0&-78&0&-13&0\\
 0&84&0&14&0
 \end{pmatrix}.
 \tag{V5.18}\label{r5:one-field}
\]
Thus even this small source-calibrated field is not coordinate
diagonal.

Let \(O\) exchange opposite spatial faces as in V4 and let
\(J_6=\mathbf1\mathbf1^{\mathsf T}\).  Direct rational evaluation of
\eqref{r5:H} yields
\[
 H=\frac{1344I_6+264O+25J_6}{2381400}.
 \tag{V5.19}\label{r5:H-explicit}
\]
Equivalently its diagonal, opposite-face, and remaining off-diagonal
entries are \(1369/2381400\), \(289/2381400\), and \(1/95256\).
Its eigenvalues and multiplicities are
\[
 \begin{array}{c|ccc}
 \hbox{subspace}&\hbox{constant}&
 \hbox{pair-symmetric, sum zero}&\hbox{pair-antisymmetric}\\ \hline
 \hbox{eigenvalue}&293/396900&67/99225&1/2205\\
 \hbox{dimension}&1&2&3
 \end{array}
 \tag{V5.20}\label{r5:eigen}
\]
The same entries follow independently from the full 29-chord
Gaussian covariance:
if \(p_i,q_j\) are the left and right plaquette rows, then
\[
 \ell_iR^{-1}\ell_j^{\mathsf T}
       =d_{p_i}Q^{-1}d_{q_j}^{\mathsf T}.
\]
Indeed conditional left and right fluctuations are independent
given the five middle chords.  This also checks the coefficient
against V2's original full-history formula, rather than a
standalone midpoint surrogate.

In particular there exists a finite, here unquantified
\(\gamma_0\) such that
\[
 \gamma\ge\gamma_0\quad\Longrightarrow\quad
 G_\gamma^{[2]}\succeq L_\gamma
                       \succeq\frac1{4410\gamma^2}I_6 .
 \tag{V5.21}\label{r5:floor}
\]
This follows from \eqref{r5:limits} and
\(\lambda_{\min}(H)=1/2205\).
V2 already implied an eventual positive native covariance on this
fixed box.  The additional conclusion here is that an explicit
native score certificate attains its leading scale, unlike V4's
fixed Haar-character compression.  Equation \eqref{r5:floor} is
not a certified value at a specified finite \(\gamma\).

\subsection{The exact cost of changing or localizing the fields}
\label{r5:locality}

The corrected normalization does not make the fields automatically
local or uniform in volume.  There is a precise finite-box test.
Replace \(V_i\) by any real matrices \(\widetilde V_i\), with the
same compact cutoff, and define
\[
 \widetilde D_i=\tfrac12
          (R\widetilde V_i+\widetilde V_i^{\mathsf T}R),\qquad
 E_i=\widetilde D_i+\tfrac12 A_i .
\]
The corresponding leading centered score is
\[
 \widetilde s_i^{(0)}
 =\sum_a(Y^a)^{\mathsf T}\widetilde D_iY^a
                              -3\Tr(C\widetilde D_i).
\]
Using \(\Tr(C\widetilde D_i)=\Tr\widetilde V_i\), the same proof
as Lemma~\ref{r5:conditional-limit} applies.
For the certificate \(\widetilde L_\gamma\) with the unchanged
coefficient \(W=\gamma^{-1}I\), the exact residual theorem yields
\[
 \gamma^2\widetilde L_\gamma\longrightarrow H-\mathcal E,\qquad
 \mathcal E_{ij}=6\Tr(E_i C E_j C),\qquad \mathcal E\succeq0.
 \tag{V5.22}\label{r5:error}
\]
To verify the last formula, \(\phi_i-\widetilde s_i^{(0)}\) is the
negative centered three-color quadratic with matrix \(E_i\).
One-color Wick pairing contributes
\(2\Tr(E_iCE_jC)\); summing the three colors gives the factor six.
This is the Gram of the score mismatch, so it is positive
semidefinite even if individual off-diagonal entries are negative.

A sufficient scale-preserving criterion for these replacement
fields is therefore
\[
       \bigl\|H^{-1/2}\mathcal E H^{-1/2}\bigr\|<1 .
 \tag{V5.23}\label{r5:relative-error}
\]
This is a concrete relative error to estimate, not an assertion
that a spatial truncation satisfies it.

For an exact adverse check on this very box, discard every
off-diagonal entry of each \(V_i\), leaving
\(\widetilde V_i=\operatorname{diag}(V_i)\).
For \(v=(1,1,2,1,-4,-2)^{\mathsf T}\), rational arithmetic gives
\[
       v^{\mathsf T}(H-\mathcal E)v
                  =-\frac{313995881}{1543790178000}<0 .
 \tag{V5.24}\label{r5:diagonal-fails}
\]
Thus the diagonal shortcut fails to supply a positive lower
matrix with this fixed \(W\), already before any limit of the box.
This is not a proof that every local field or every reoptimized
\(W\) fails.  Coordinate diagonality in a tree chart is also not
the same as physical spatial locality.  It is a narrowly specified,
exact countercheck against discarding the coupled response.

The matrices \(R^{-1}\), \(Q_{zz}^{-1}\) and the conditional
normalizers remain the upstream dependencies.
At larger spatial size or temporal depth they contain the
response of the omitted region.  No estimate in this note makes
their source-weighted tails summable uniformly, or controls the
nonlinear conditional Laplace remainder when those sizes grow.
Replacing the missing estimates by the finite-box Gaussian
matrices would undo the native-law distinction just proved.

\subsection{Verification and the next genuinely unpaid estimate}
\label{r5:ledger}

The exact diagnostic is
\begin{center}
 \path{scripts/verify_ym_restart_v5_midplane.py}.
\end{center}
It builds the literal 52-link, 42-plaquette incidence problem,
checks the zero left--right precision block, computes both the
full inverse and the middle Schur complement, and verifies
\eqref{r5:matrices}--\eqref{r5:eigen}.
It also checks the Gaussian response and score matrices
independently by Wick traces:
\[
 \E[\phi_i s_j^{(0)}]
    =-3\Tr(A_i C D_j C),\qquad
 \E[s_i^{(0)}s_j^{(0)}]=6\Tr(D_i C D_j C),
 \quad D_j=-A_j/2.
\]
Both equal \(H\), and their optimized lower matrix equals \(H\).
The diagonal-field mismatch and
\eqref{r5:diagonal-fails} use exact fractions.
No floating-point eigenvalue or simulation is used as a
certificate.  These checks verify finite algebra, not the
analytic conditional-limit proof or a finite-coupling enclosure.

\emph{Newly established:} an exact native marginal-score lower
matrix for opposite-boundary correlations; retention of the
conditional response that distinguishes it from null contact
scores; explicit compact score fields whose lower matrix
saturates the six-source leading Gram on the full two-step box;
and an exact relative mismatch formula for proposed simpler
fields, with a failed diagonal shortcut.

\emph{Not established:} rigorous numerical enclosures of this
new score matrix at specified finite \(\gamma\); an explicit
onset for \eqref{r5:floor}; exterior-uniform or volume-uniform
bounds; fixed positive physical time or flow; vacuum preparation;
the common growing-depth spectral inequalities; or an
interacting continuum theory and its mass gap.

The next calculation should attack the \emph{conditional response}
in \eqref{r5:score-def} and \eqref{r5:response}, with all crossing
plaquettes retained.  Use \eqref{r5:error} as a quantitative
calibration for a spatially supported or RG-organized witness,
then control its nonlinear native-law error rather than assuming
Gaussian locality.  Increasing an unweighted character bank is
no longer the selected direction.  Finite-box saturation alone
must not become another loop of finite-box examples.

The complete V4 body is preserved except for its displayed title.
Removing this delimited appendix and reversing that title recovers
V4 byte for byte.  Only the immediate active source is replaced;
all archives and companion sources remain unchanged.  All
verification and compilation artifacts are outside
\path{latex_notes}.

\begin{center}
\fbox{\parbox{.9\textwidth}{\textbf{V5 status.}
The global Haar-normalization loss is removed for a specified
native finite-box score certificate, and its leading matrix is
fully calibrated.  Uniform conditional-response control and
the full interacting continuum Yang--Mills mass gap remain
unproved.}}
\end{center}
% END CONSOLIDATED V5 APPEND

% BEGIN CONSOLIDATED V6 APPEND -- 2026-09-19
\clearpage
\section{V6: local plaquette-gradient fields and a volume-independent score budget}
\label{r6:context}

V5 removed a reference-volume normalization loss, but its calibrated
fields used the inverse of the entire middle-slice Schur complement.
The diagonal shortcut failed and was not a test of physical locality.
We now use fields defined directly on the original link manifold:
the gradient of each middle spatial plaquette trace.  Each acts on
only four links and is gauge equivariant, without any tree gauge or
small-field cutoff in its definition.

There are two distinct results.  First, the Gram of their
\emph{marginal} scores has an explicit norm bound independent of
the spatial volume, temporal depth, and number of marked plaquettes.
This bound holds at every coupling under the full compact Wilson law.
Second, on V5's full two-step box these genuinely local fields retain
more than four fifths of the leading source covariance and exactly
retain its pair-antisymmetric part.  The second statement is a
fixed-box asymptotic calibration, not a uniform response theorem.
Neither result supplies the continuum mass gap.

\subsection{The original manifold and the projected scores}
\label{r6:setup}

Let \(\Lambda\) be a finite four-dimensional hypercubic box with
the Wilson action
\[
 S(U)=\sum_p(1-w_p(U)),\qquad
 w_p(U)=\tfrac12\Tr U_p,\qquad
 d\mu_\Lambda=Z^{-1}e^{-\gamma S}\,dH .
\]
All differentiated links are integrated freely.  Open boundaries
are allowed, as are nondegenerate periodic boundaries for which an
elementary plaquette has four distinct links.  Incident plaquettes
with fixed exterior links are retained in \(S\).
Let \(I\) be the set of spatial links on a chosen middle slice.
Use the product bi-invariant metric for which \(\SU(2)\) is the unit
three-sphere: a single-link linear quaternion coordinate has
Laplacian eigenvalue \(-3\), and each link factor has Ricci tensor
\(2g\).

Let \(\mathcal P_I\) be any finite collection of distinct elementary
spatial plaquettes on this slice.  Combine repeated or reversed
marks first, since \(w_{p^{-1}}=w_p\).  For real coefficients \(c\),
set
\[
 F_c=\sum_{p\in\mathcal P_I}c_pw_p,\qquad
 X_p=\nabla_I w_p,\qquad
 \mathcal L_I=\Delta_I-\gamma\langle\nabla_I S,\nabla_I\,\cdot\,\rangle .
 \tag{V6.1}\label{r6:generator}
\]
This is an auxiliary integration-by-parts operator.
It is \emph{not} the physical transfer generator or Hamiltonian.
No spectral-gap assertion about \(\mathcal L_I\) is used.

Define the full-history score and its conditional projection by
\[
 T_p=-\mathcal L_Iw_p
      =12w_p+\gamma\langle\nabla_Iw_p,\nabla_I S\rangle,\qquad
 s_p=\E_{\mu_\Lambda}[T_p\mid U_I].
 \tag{V6.2}\label{r6:scores}
\]
The first equality uses four freely varying links in the plaquette.
The field \(X_p\) and all of its nonzero components involve only
those four links.  The function \(T_p\) includes all action
plaquettes incident to them, including temporal ones.
The projected \(s_p\) need not be spatially local.

Let \(p_I\) be the exact marginal density on the middle links.
Differentiating its compact conditional integral gives
\[
 s_p=-\Delta_Iw_p-\langle\nabla_Iw_p,\nabla_I\log p_I\rangle .
 \tag{V6.3}\label{r6:marginal-score}
\]
Thus these are the actual marginal scores of V5, now generated by
specified local fields on the ungauged manifold.
They are gauge invariant scalars: local gauge transformations are
isometries, \(w_p,S,p_I\) are invariant, and gradient indices are
contracted.  Their means are zero by compact integration by parts.

For the score Grams write
\[
       \widehat J_{pq}=\E_{\mu_\Lambda}T_pT_q,\qquad
       J_{pq}=\E_{\mu_\Lambda}s_ps_q.
\]
Conditional expectation is an \(L^2\) contraction, so exactly
\[
                        0\preceq J\preceq\widehat J .
 \tag{V6.4}\label{r6:projection}
\]
This inequality does not assert that projection preserves finite
range or individual off-diagonal covariance signs.

\subsection{Local geometric estimates with explicit constants}
\label{r6:geometry}

\begin{lemma}[Uniform plaquette derivative bounds]
\label{r6:derivatives}
For a single elementary plaquette with four distinct free links,
\[
 |\nabla_e w_p|^2=1-w_p^2\quad(e\in p),\qquad
 \|\nabla_I^2w_p\|_{\rm HS}^2\le48.
 \tag{V6.5}\label{r6:single}
\]
For \(F_c\) in \eqref{r6:generator}, pointwise,
\[
 \|\nabla_I^2F_c\|_{\rm HS}^2\le192\sum_pc_p^2,\qquad
 |\nabla_IF_c|^2\le16\sum_pc_p^2(1-w_p^2).
 \tag{V6.6}\label{r6:bank}
\]
For the full action, including all incident temporal plaquettes,
\[
                 \nabla_I^2S(v,v)\le24|v|^2 .
 \tag{V6.7}\label{r6:action}
\]
All constants are independent of the box and its boundary data.
\end{lemma}
\begin{proof}
Holding the other links fixed, the plaquette trace is a unit
linear quaternion coordinate of the chosen link, possibly after
inversion and left/right multiplication.  These maps are
isometries of the unit three-sphere.  Its gradient has squared
length \(1-w_p^2\), and its same-link Hessian is \(-w_p g\).

For two different links, parallel transport of the tangent frames
reduces the mixed derivative block, up to orthogonal factors and
transpose, to \(wI-C_v\), where the plaquette quaternion is \(w+v\)
and \(C_vz=v\times z\).  The exact identity
\[
 (wI-C_v)^{\mathsf T}(wI-C_v)=I-vv^{\mathsf T}
\]
gives operator norm at most one and squared Hilbert--Schmidt norm
\(2+w^2\le3\).  This is also the mixed-block identity already used
in f16.2 V85; it is not a new quaternion identity.
There are four diagonal blocks and twelve ordered off-diagonal
blocks.  Their squared norms sum to at most
\(4\cdot3+12\cdot3=48\).

A spatial link belongs to at most four spatial plaquettes.
For each Hessian block \((e,f)\), at most four marks contribute.
Cauchy--Schwarz in that finite sum, followed by the single-face
bound, gives \(4\cdot48\sum_pc_p^2\).
Likewise at each link the gradient sum contains at most four
terms.  Sum the resulting squares over links; each plaquette is
counted four times.  This proves the factor \(4\cdot4=16\).

Finally, for one action plaquette the absolute Hessian quadratic
form in its varying links is at most
\((\sum_{e\in p}|v_e|)^2\le4\sum_{e\in p}|v_e|^2\).
A four-dimensional link is in at most six plaquettes.
Summing gives the factor \(4\cdot6=24\).
Fixed exterior links simply have zero variation; they do not
alter these counts or the unit-quaternion bounds.
\end{proof}

\subsection{An all-coupling score norm bound without volume growth}
\label{r6:norm}

Put
\[
 \epsilon_{\mathcal P}=\max_{p\in\mathcal P_I}
                \E_{\mu_\Lambda}(1-w_p^2)\in[0,1],\qquad
 K_{\gamma,\mathcal P}
     =192+16(2+24\gamma)\epsilon_{\mathcal P}.
 \tag{V6.8}\label{r6:K}
\]
For an empty bank the statements are vacuous.
\begin{theorem}[Native projected-score budget]
\label{r6:uniform}
For every \(\gamma\ge0\) and every stated finite Wilson geometry,
\[
 \boxed{\quad
 0\preceq J\preceq\widehat J
              \preceq K_{\gamma,\mathcal P}I
              \preceq(224+384\gamma)I .
 \quad}
 \tag{V6.9}\label{r6:uniform-bound}
\]
The last bound is uniform in volume, temporal depth, and bank size.
It remains valid with arbitrary fixed exterior links provided the
middle links being differentiated remain free.
No conditional mixing or positive curvature assumption is needed.
\end{theorem}
\begin{proof}
First freeze all links outside \(I\).  The conditional law on
\(\SU(2)^I\) has smooth density proportional to \(e^{-\gamma S}\).
The weighted integrated Bochner identity is
\[
 \E[(\mathcal L_IF_c)^2\mid U_{I^c}]
 =\E\left[
    \|\nabla_I^2F_c\|_{\rm HS}^2+
    2|\nabla_IF_c|^2+
    \gamma\nabla_I^2S(\nabla_IF_c,\nabla_IF_c)
          \,\middle|\, U_{I^c}\right].
 \tag{V6.10}\label{r6:bochner}
\]
For clarity, the pointwise identity is
\[
 \tfrac12\mathcal L_I|\nabla_IF_c|^2
  -\langle\nabla_IF_c,\nabla_I\mathcal L_IF_c\rangle
 =\|\nabla_I^2F_c\|_{\rm HS}^2+
   (2g+\gamma\nabla_I^2S)(\nabla_IF_c,\nabla_IF_c).
\]
Its first term integrates to zero.  Integration by parts in the
second yields \(\E(\mathcal L_IF_c)^2\).
The compact product manifold has no boundary term.
This is the standard Bochner identity with potential \(\gamma S\),
not a lower curvature estimate.

Average \eqref{r6:bochner} over the exterior.
Equations \eqref{r6:bank}--\eqref{r6:action} give
\[
 \E T_c^2
 \le192|c|^2+(2+24\gamma)\E|\nabla_IF_c|^2
 \le K_{\gamma,\mathcal P}|c|^2 .
\]
This is the asserted matrix bound for \(\widehat J\);
\eqref{r6:projection} gives it for \(J\).
Use \(\epsilon_{\mathcal P}\le1\) for the final explicit constant.
Every step can equally be applied with any further fixed
nondifferentiated exterior data.
\end{proof}

The Bochner identity and its drift transformation are standard;
for primary-source context see K.-T.~Sturm,
\emph{Ricci Tensor for Diffusion Operators and Curvature-Dimension
Inequalities under Conformal Transformations and Time Changes},
\href{https://arxiv.org/html/1401.0687v3}{arXiv:1401.0687v3},
Sections 5 and 8.
The compact identity actually used is displayed above.
The archives already use Bochner methods in other settings.
The present addition is their explicit application to the native
\emph{projected four-link field bank}, with no volume factor.
This targeted check is not the scheduled V10 literature sweep.

\paragraph{A logarithmic refinement, only where its measure matches.}
If an independently proved native plaquette bound
\(\E(1-w_p)\le\delta_\gamma\) holds for the marked law, then
\(1-w_p^2\le2(1-w_p)\) gives
\[
 K_{\gamma,\mathcal P}
 \le192+16(2+24\gamma)\min\{1,2\delta_\gamma\}.
 \tag{V6.11}\label{r6:energy-refinement}
\]
For the equal-coupling, fully periodic hypercubic law of f16.2
V64, its proved entropy/cap estimate is
\(\delta_\gamma\le B_\gamma/\gamma\), where
\(B_\gamma=8+\log\gamma-\tfrac23\log c_H\) for \(\gamma\ge1\).
Here \(c_H\) is its fixed positive Haar-cap constant.
In precisely that law \eqref{r6:energy-refinement} is \(O(\log\gamma)\)
uniformly in the periodic volume.
This reuses the archived estimate, rather than reproving it or
upgrading it to \(O(1/\gamma)\).

The periodic estimate is not asserted for V5's open box, arbitrary
exterior data, or a vacuum obtained by a different limiting
procedure.  Equation \eqref{r6:uniform-bound}, without that refinement,
does hold in all the stated geometries.
The logarithmic refinement is a score-norm statement, not a
positive-time covariance lower bound.

\subsection{What this pays in the actual boundary-source problem}
\label{r6:response}

On any reflection-symmetric finite open history with a middle
slice, keep the original left/right conditional laws, as in V5.
For a finite vector of outer observables \(A\), let
\(m(U_I)=\E[A\mid U_I]\).
Reflection and conditional independence give
\(G=\Cov(m,m)\) for the matching opposite-boundary covariance.
For the local field bank define the rectangular response matrix
\[
                  B_{ip}=\E\langle\nabla_Iw_p,\nabla_Im_i\rangle .
 \tag{V6.12}\label{r6:response-B}
\]
Compact integration by parts gives
\(B=\E[(m-\E m)s^{\mathsf T}]\).

\begin{corollary}[The denominator can be bounded without an inverse]
\label{r6:response-corollary}
For \(K=K_{\gamma,\mathcal P}\), or the larger explicit value
\(224+384\gamma\),
\[
                         G\succeq K^{-1}BB^{\mathsf T}.
 \tag{V6.13}\label{r6:lower}
\]
\end{corollary}
\begin{proof}
V5's exact residual identity applies to any rectangular coefficient
matrix \(W\).  Choose \(W=B/K\).
Since \(J\preceq KI\), its lower matrix is
\[
 2BB^{\mathsf T}/K-BJB^{\mathsf T}/K^2
                              \succeq BB^{\mathsf T}/K .
\]
No invertibility of \(J\), \(B\), or \(G\) was assumed.
\end{proof}

For an outer observable independent of middle links,
differentiation of the normalized half-history expectation gives
\[
 B_{ip}=-\gamma\,\E
   \Cov_{\rm half}\left(A_i,
           \langle\nabla_Iw_p,\nabla_IS_{\rm half}\rangle
                         \,\middle|\,U_I\right).
 \tag{V6.14}\label{r6:half-response}
\]
Every crossing plaquette remains in this formula.
The norm problem on the score side is now bounded independently of
the size of the history.  The \emph{lower singular values of the
response \(B\)} have not been bounded uniformly.
An upper bound on \(J\) cannot create such a lower bound.

In particular, locality of the generating \(X_p\) does not make
\(s_p\), \(m_i\), or \(B_{ip}\) local after the other links are
integrated.  Equation \eqref{r6:half-response} is the remaining
boundary-response problem, not its solution.
Nor may one replace these projected scores by the full \(T_p\):
the latter have the contact nullity discussed in V5.

\subsection{Correct physical gradients in the V5 benchmark}
\label{r6:metric}

We test the local choice on the same full
\(2\times1\times1\times1\) box as V5, rather than selecting a new
reference measure.  The geometry and the native conditional
Laplace result are inherited from V5.
The extra issue is the metric: deleting tree coordinates does
\emph{not} turn the physical product-link metric into the identity
matrix on chords.

At the linearized level, let \(T:\R^{12}\to\R^5\) take the spatial
link coordinates to the five tree-gauged chords.
If \(D_T,D_C\) are the rooted vertex-incidence matrices on the
seven tree edges and the five chord edges, and \(P_T,P_C\) select
those edges, then
\[
                T=P_C-D_CD_T^{-1}P_T,\qquad M=TT^{\mathsf T}.
 \tag{V6.15}\label{r6:T}
\]
This follows by subtracting the infinitesimal gauge gradient that
sets the tree edges to zero.
In the exact V5 ordering one obtains
\[
 M=\begin{pmatrix}
 4&-1&1&1&3\\
 -1&4&1&0&-2\\
 1&1&4&-1&1\\
 1&0&-1&4&2\\
 3&-2&1&2&6
 \end{pmatrix}.
 \tag{V6.16}\label{r6:M}
\]
The middle plaquette incidence rows in these chord coordinates are
\[
 \begin{pmatrix}a_1\\a_2\\a_3\\a_4\\a_5\\a_6\end{pmatrix}
 =
 \begin{pmatrix}
 0&1&0&0&0\\
 0&0&1&0&0\\
 1&0&0&0&0\\
 0&0&0&1&0\\
 -1&0&0&0&1\\
 0&1&-1&-1&1
 \end{pmatrix}.
 \tag{V6.17}\label{r6:arows}
\]
Each full-link row \(a_pT\) has squared length four.

Since
\(w_p=1-\tfrac12\sum_a(a_px^a)^2+O(|x|^3)\),
the pushforward of its physical gradient has leading linear
coefficient
\[
                 V_p^{\rm loc}=-M a_p^{\mathsf T}a_p.
 \tag{V6.18}\label{r6:Vlocal}
\]
Indeed \(\nabla_{\rm links}w_p=T^{\mathsf T}\nabla_xw_p\)
to leading order, and pushing it forward multiplies once more by
\(T\).  The matrices in \eqref{r6:Vlocal} are generally not diagonal,
even though the original fields act on only four physical links.
This explains why V5's diagonal-field failure did not rule them out.

The use of \(M\) here is a coordinate calculation for a field
already defined globally and locally on the group.
It does not redefine the gradient or insert a quotient metric
into the Wilson action.

\subsection{A fully evaluated leading lower matrix from local fields}
\label{r6:calibration}

Retain V5's \(R,C=R^{-1},\ell_i\), and
\(A_i=\ell_i^{\mathsf T}\ell_i\).
Set
\[
 D_p=\tfrac12(RV_p^{\rm loc}+(V_p^{\rm loc})^{\mathsf T}R),
\]
and form the exact one-color Wick contractions, including their
three-color multiplicities,
\[
 B^{(0)}_{ip}=-3\Tr(A_iCD_pC),\qquad
 J^{(0)}_{pq}=6\Tr(D_pCD_qC).
 \tag{V6.19}\label{r6:BJ}
\]
As a separate expression for the response,
\[
 B^{(0)}_{ip}
       =3(\ell_iM a_p^{\mathsf T})(\ell_iC a_p^{\mathsf T}).
 \tag{V6.20}\label{r6:direct-response}
\]
It follows directly by differentiating the leading conditional
outer trace along \(V_p^{\rm loc}x\).
It also follows algebraically from \eqref{r6:BJ}.

\begin{proposition}[Native asymptotics of the local bank]
\label{r6:local-limit}
On this fixed open box, for the globally defined fields
\(X_p=\nabla_Iw_p\),
\[
                  \gamma B_\gamma\longrightarrow B^{(0)},
                  \qquad J_\gamma\longrightarrow J^{(0)}.
 \tag{V6.21}\label{r6:local-convergence}
\]
\end{proposition}
\begin{proof}
The middle marginal and conditional outer means are exactly those
of V5 before the spatial tree variables are integrated.
Gauge invariance makes their integrals equal to the corresponding
tree-gauged integrals.
The gauge-equivariant field has expansion
\(V_p^{\rm loc}x+O(|x|^2)\), by \eqref{r6:T}--\eqref{r6:Vlocal}.
Its divergence at the identity is
\(3\Tr V_p^{\rm loc}=-12\), consistent with
\(\Delta_Iw_p=-12w_p\).
Derivative terms from the smooth coordinate change contribute only
higher orders.

V5's conditional Laplace proof thus gives the centered leading score
\[
 s_p^{(0)}=\sum_a(Y^a)^{\mathsf T}D_pY^a-3\Tr(CD_p),
 \qquad Y^a\sim N(0,C),
\]
and the same centered leading outer message \(\phi_i\).
The globally smooth fields and their derivatives are bounded.
Outside a small identity chart their scores are \(O(1+\gamma)\),
while the fixed-box marginal mass is exponentially small.
Inside the chart the quadratic bounds and uniform integrability
used in V5 apply unchanged.  Thus the corresponding products
converge in expectation.  Wick pairing gives
\eqref{r6:BJ}, and the mean-response identity gives
\eqref{r6:local-convergence}.
This proves a statement about the native compact law, not only a
Gaussian test model.
\end{proof}

All of \(H\) from V5, \(B^{(0)}\), and \(J^{(0)}\) are in the algebra
spanned by \(I_6\), the opposite-face swap \(O\), and
\(\mathbf1\mathbf1^{\mathsf T}\).
Their three invariant spaces are the constant vectors, the
pair-symmetric sum-zero vectors, and the pair-antisymmetric vectors.
Rational calculation gives
\[
\begin{array}{c|ccc}
 &\text{constant}&\text{pair-symmetric, sum zero}&
                       \text{pair-antisymmetric}\\ \hline
 B^{(0)}&4/15&23/105&466/2205\\
 J^{(0)}&120&84&217156/2205
\end{array}.
 \tag{V6.22}\label{r6:bands}
\]
In particular both matrices are invertible.
As entrywise checks, the diagonal/opposite/other entries are
\[
 B^{(0)}:\quad 164/735,\ 26/2205,\ 1/126;
 \qquad
 J^{(0)}:\quad 214418/2205,\ -2738/2205,\ 6 .
\]
The negative opposite-face entry of \(J^{(0)}\) is compatible with
its positive eigenvalues.

Define the constant rational matrix
\[
 W_0=B^{(0)}(J^{(0)})^{-1}
   =\tfrac1{450}\Pi_0+
          \tfrac{23}{8820}\Pi_s+\tfrac1{466}\Pi_a ,
 \tag{V6.23}\label{r6:W}
\]
where \(\Pi_0,\Pi_s,\Pi_a\) are the orthogonal projectors onto those
three spaces.  V5's exact residual inequality, now with
\(W=\gamma^{-1}W_0\), gives for every \(\gamma>0\)
\[
 G_\gamma^{[2]}\succeq L_\gamma^{\rm loc}
 :=\gamma^{-1}(B_\gamma W_0^{\mathsf T}
                     +W_0B_\gamma^{\mathsf T})
           -\gamma^{-2}W_0J_\gamma W_0^{\mathsf T}.
 \tag{V6.24}\label{r6:local-bound}
\]
This definition needs no finite-\(\gamma\) matrix inverse.
Equations \eqref{r6:local-convergence} imply
\[
 \gamma^2 L_\gamma^{\rm loc}\longrightarrow
 H_{\rm loc}=B^{(0)}(J^{(0)})^{-1}(B^{(0)})^{\mathsf T}.
\]
Its three eigenvalues and their ratios to V5's exact native
leading \(H\) are
\[
\begin{array}{c|ccc}
 &\Pi_0&\Pi_s&\Pi_a\\ \hline
 H_{\rm loc}&2/3375&529/926100&1/2205\\
 H_{\rm loc}/H&1176/1465&1587/1876&1
\end{array}.
 \tag{V6.25}\label{r6:retention}
\]
The ratio notation in this table refers to the common eigenspaces,
not entrywise division of matrices.  Hence
\[
             \tfrac45 H\preceq H_{\rm loc}\preceq H,\qquad
                    \lambda_{\min}(H_{\rm loc})=1/2205.
 \tag{V6.26}\label{r6:comparison}
\]
The first inequality is strict as a quadratic form at coefficient
\(4/5\).  Its largest possible scalar coefficient is the minimum
actual ratio \(1176/1465\), with equality on the constant subspace.
The residual eigenvalues are exactly
\(289/1984500,\ 289/2778300,\ 0\).

Thus local generating fields retain the entire leading
pair-antisymmetric source space and more than eighty percent of
the leading Gram in every direction.
No analogous retention percentage at a specified finite coupling,
larger volume, physical flow, or growing temporal depth is claimed.
An eventual lower floor
\(L_\gamma^{\rm loc}\succeq I_6/(4410\gamma^2)\) follows on this
fixed box, but its onset is still unquantified.

\subsection{What is genuinely uniform, and what must be attacked next}
\label{r6:scope}

The result that is uniform in the history is
\eqref{r6:uniform-bound}.  It controls arbitrary coefficient banks
without multiplying by their cardinality.
It does not use a global grounded Green function, a probability
comparison with Haar, a spectral gap for a sampling dynamics, or
small plaquette assumptions.

The leading retention test \eqref{r6:retention} has a different
scope: it shows that the selected physical four-link fields do
not intrinsically lose the weak-coupling power on the full
benchmark.  Its proof fixes the history before taking
\(\gamma\to\infty\).  It is not evidence that the conditional
Laplace errors are uniform as the history grows.

The still missing input is now sharply identifiable:
a lower bound on the response matrix
\eqref{r6:half-response}, with physically normalized sources and
the actual temporal/flow scale.
The norm bound permits a lower Gram once that response is paid,
but provides no lower singular value itself.
In particular a positive \(J\), bounded norm, reflection positivity,
or the finite-box ratios cannot supply an upper spectral support
edge or the full Hamiltonian gap.

The next step is to derive a local or scale-decomposed estimate
for this conditional response, retaining its normalization and
all crossing plaquettes.  If it becomes small or ill-conditioned,
resolve the responsible boundary-return modes before adding more
source fields.  A theorem about the auxiliary diffusion
\(\mathcal L_I\) would not, without a separate bridge, solve this
physical response problem.

\subsection{Verification, provenance, and remaining obligations}
\label{r6:ledger}

The exact checker is
\begin{center}
 \path{scripts/verify_ym_restart_v6_local_fields.py}.
\end{center}
It reconstructs the spatial tree projection, checks that it
annihilates vertex gauge gradients and is the identity on chord
variations, and obtains the nonidentity metric \(M\).
Using V5's literal full-history incidence, it independently evaluates
the response via both \eqref{r6:BJ} and
\eqref{r6:direct-response}, then checks the score Gram,
optimized lower matrix, all three eigenvalue bands, and the
retained fractions with exact rational arithmetic.
Twenty-five rational unit-quaternion cases check the mixed-block
identity.  They supplement its algebraic proof, not replace it.
The uniform constants are proved by local incidence counts and the
integrated identity, not extrapolated from finite numerical boxes.

The f16.2 V85 quaternion/score results and V64 entropy estimate
are explicitly reused with their original claim boundaries.
V5's conditional asymptotic proof and the general score residual
identity are inherited.  The new content is the physical local
field bank, its explicit native norm budget uniform in volume
and depth, the correct tree-coordinate metric, and its evaluated
source retention on the original two-step Wilson history.
No broad literature or companion checkpoint is due before V10.

Still open are finite-coupling enclosures for this new response,
uniform lower response estimates at physical scales, a nonzero
contact-cleaned flowed source Gram, vacuum preparation, common
growing-depth upper spectral bounds on a dense physical source
core, and the interacting continuum Yang--Mills construction
and mass gap.  The present elementary plaquette bank is not
claimed to span every physical sector.

The V5 body is preserved byte for byte apart from its displayed
version title.  Removing this appendix and reversing that title
recovers the exact predecessor.  Only the immediate active source
is replaced; all archives and companion notes remain unchanged.
Checks and compilation artifacts remain outside \path{latex_notes}.

\begin{center}
\fbox{\parbox{.9\textwidth}{\textbf{V6 status.}
Local generating fields have a proved native score norm budget
independent of history size, and pass a nontrivial full-Wilson
leading-scale calibration.  Their uniform physical response and
the full interacting continuum Yang--Mills mass gap remain
unproved.}}
\end{center}
% END CONSOLIDATED V6 APPEND

% BEGIN CONSOLIDATED V7 APPEND -- 2026-09-19
\clearpage
\section{V7: time-depth dilution and scale-matched boundary scores}
\label{r7:context}

V6 pays a native score \emph{upper} norm independently of history
size.  It does not pay the response to a remote observable.
Before seeking such a response with the same microscopic fields
at every depth, we test that choice on the already identified
transverse Gaussian transfer.  This is an upstream test of the
certificate, not a replacement of the Wilson law.

The outcome changes the next calculation.  For the zero-spatial-
momentum magnetic-energy source, the \emph{entire} elementary
plaquette-gradient bank captures a fraction asymptotic to
\(25/(14\pi^2 n^7)\) at opposite times \(-n,n\), when the spatial
box grows faster than \(n\).  This is an optimal-bank statement,
not failure of a particular coefficient choice.  A single
spatially smoothed score at flow time \(n^2/2\) instead captures,
in the same limit, strictly more than \(8/9\).  Finally, the
exactly matched reference field is a positive integral of
spatially flowed fields with an explicit Bessel weight.

\paragraph{Scope before the calculation.}
All asymptotic retention and filter statements in this section
are theorems for a specified \emph{stationary, transverse
Gaussian reference}.  They are not new uniform asymptotics of
the compact Wilson measure.  In particular, we do not exchange
V5's fixed-open-box coupling limit with a volume, vacuum, or
time-depth limit.  The native construction at the end defines
candidate fields and exact identities, and states the still
unpaid estimates.  The full interacting continuum mass gap is
not proved.

\subsection{Nonduplication, sources, and the reference law}
\label{r7:reference}

The common f16 V38--V41 material already identifies the Gaussian
transfer normalization and dispersion; f16.1 V50 records its
transverse multiplicities and shrinking smallest frequency.
Those are inherited inputs, not new dispersion theorems.
The f16.2 V69--V70 spatial-flow analysis and V85 contact-score
analysis do not estimate the marginal response studied here.
Earlier f7 V19 already supplies the weighted-divergence identity.
The new issue is the optimal \emph{capture fraction} of V6's field
bank at growing physical-transfer depth and the filter that
repairs its reference-scale loss.

As a targeted primary-source check, M.~L\"uscher,
\href{https://arxiv.org/html/1006.4518}{\emph{Properties and uses
of the Wilson flow in lattice QCD}}, Sections 1--2, explains
finite-lattice flow and its leading heat smoothing.  Its
four-dimensional flow is not silently substituted for our
time-slice flow.  The Bessel identities used below are the
standard integer-order formulas in
\href{https://dlmf.nist.gov/10.32.E3}{DLMF 10.32.3} and
\href{https://dlmf.nist.gov/10.35.E2}{10.35.2}; the needed
resolvent identity is proved here.  Neither the scheduled broad
literature sweep nor the companion review is due before V10.

Let \(L\ge3\) be the side of a cubic spatial torus and \(V=L^3\).
In the fine-link Euclidean metric let
\[
 \mathcal Y_L=\ker d_0^*\ominus(\ker d_0^*\cap\ker d_1),
 \qquad \mathsf L=d_1^*d_1\big|_{\mathcal Y_L}>0 .
 \tag{V7.1}\label{r7:space}
\]
There are three real colors.  The constant and gauge directions
are \emph{excluded from this reference}, not declared harmless
in the compact nonlinear theory.  Its real dimension is
\(6(V-1)\).  Write
\[
 \lambda(k)=4\sum_{\mu=1}^3\sin^2(k_\mu/2),\qquad
 \theta(\lambda)=\operatorname{arcosh}(1+\lambda/2),\qquad
 R(\lambda)=2\sinh\theta=\sqrt{\lambda(\lambda+4)},
 \tag{V7.2}\label{r7:symbols}
\]
where \(k\in(2\pi/L)\mathbb Z_L^3\setminus\{0\}\).
Each nonzero momentum contributes six real degrees in spectral
traces; the usual reality condition is included in this count.
Set \(\mathsf R=R(\mathsf L)\), and let \(\nu_L\) be the centered
Gaussian law with covariance \(\mathsf R^{-1}\).

For completeness its stationary time evolution follows from
the inherited one-mode transfer kernel
\[
 t_\lambda(x,y)=(2\pi)^{-1/2}
 \exp[-(x-y)^2/2-\lambda(x^2+y^2)/4].
\]
Its ground function is proportional to \(e^{-R x^2/4}\).
Completing the square in its ground transform gives exactly
\[
 y=e^{-\theta}x+\eta,\qquad
 \eta\sim N(0,e^{-\theta}),\qquad
 \frac{e^{-\theta}}{1-e^{-2\theta}}=\frac1R .
 \tag{V7.3}\label{r7:AR}
\]
The chain is reversible under \(N(0,R^{-1})\).
Taking the product defines the two-sided stationary reference.
This is the ground transform of the specified Gaussian transfer,
not an auxiliary Langevin clock.  Its finite-history endpoints
are sampled from this stationary law, not from V5's open law.

\subsection{Exact response and score formulas at every depth}
\label{r7:response}

Define the centered total magnetic-energy source, with its sign
chosen to match the plaquette trace,
\[
 A(x)=-\frac1{2\sqrt V}
       \{x^{\mathsf T}\mathsf Lx-\Tr(\mathsf L\mathsf R^{-1})\}.
\]
This is one translation- and cubic-invariant source, not a dense
physical source family.  Conditional independence across time
zero and \eqref{r7:AR} give
\[
 m_n(x)=\E[A(x_n)\mid x_0=x]
       =-\frac1{2\sqrt V}
          \{x^{\mathsf T}a_n(\mathsf L)x
                           -\Tr(a_n(\mathsf L)\mathsf R^{-1})\},
 \quad a_n(\lambda)=\lambda e^{-2n\theta(\lambda)} .
 \tag{V7.4}\label{r7:message}
\]
Consequently the opposite-time covariance is
\(H_{n,L}=\E_{\nu_L}m_n^2=\Cov(A(x_{-n}),A(x_n))\).

For a real spectral function \(q\), let
\[
 X_q(x)=-V^{-1/2}q(\mathsf L)x,\qquad
 s_q=-\operatorname{div}_{\nu_L}X_q
   =V^{-1/2}\{\Tr q(\mathsf L)
                   -x^{\mathsf T}\mathsf Rq(\mathsf L)x\}.
 \tag{V7.5}\label{r7:score}
\]
This score is generated at the middle slice.  It is not the
contact-null full-history score of V5--V6.

\begin{proposition}[Exact scalar certificate in the reference law]
\label{r7:traces}
For \(n\ge1\), put \(B_{n,L}(q)=\E m_ns_q\) and
\(J_L(q)=\E s_q^2\).  Exactly,
\[
 \begin{split}
 H_{n,L}&=\frac1{2V}\Tr[a_n(\mathsf L)^2\mathsf R^{-2}],\\
 B_{n,L}(q)&=\frac1V\Tr[a_n(\mathsf L)q(\mathsf L)\mathsf R^{-1}],\\
 J_L(q)&=\frac2V\Tr[q(\mathsf L)^2].
 \end{split}
 \tag{V7.6}\label{r7:traces-eq}
\]
For \(J_L(q)>0\), the best coefficient multiplying this score
gives the lower covariance \(B_{n,L}(q)^2/J_L(q)\).
Its fraction of the actual reference covariance is
\[
 \rho_{n,L}(q)=\frac{B_{n,L}(q)^2}{J_L(q)H_{n,L}}\in[0,1].
 \tag{V7.7}\label{r7:rho}
\]
\end{proposition}
\begin{proof}
In independent real normal coordinates,
\(\Var(x_j^2)=2/R_j^2\); different coordinates have zero
centered quadratic covariance.  Substitute
\eqref{r7:message}--\eqref{r7:score}.  This proves all three
traces, including factors two and \(V^{-1}\).
Equivalently \(B=\E X_qm_n\) by Gaussian integration by parts.
Minimizing \(\E(m_n-cs_q)^2\) over real \(c\) gives
\(c=B/J\), residual \(H-B^2/J\), and the last assertions.
\end{proof}

The reference of V6's elementary plaquette field is
\(-d_p^*d_p x\), where \(d_p x\) is the color plaquette curl.
Summing over all spatial plaquettes gives
\(-\mathsf Lx\), hence \(q_0(\lambda)=\lambda\) in
\eqref{r7:score}.  The gradient already lies in
\(\mathcal Y_L\); no identity metric on tree chords is inserted.

\begin{lemma}[The full elementary bank cannot improve the invariant source]
\label{r7:symmetry}
For this \(m_n\), the largest certificate over all real linear
combinations of all elementary plaquette-gradient scores equals
\(B_{n,L}(q_0)^2/J_L(q_0)\).
\end{lemma}
\begin{proof}
Translations and cubic rotations preserve the real link metric,
\(\mathsf L,\mathsf R,\nu_L\), and \(m_n\); they act transitively
on the elementary unoriented spatial plaquettes.  Average any
score \(z=\sum_pc_ps_p\) over this finite group.  The result
\(\bar z\) is a scalar multiple of \(\sum_ps_p\).
Invariance gives \(\E m_nz=\E m_n\bar z\), while the averaging
operator is an orthogonal projection in \(L^2(\nu_L)\), so
\(\E\bar z^2\le\E z^2\).  Thus
\((\E m_nz)^2/\E z^2\) cannot exceed the constant-bank ratio.
Zero denominators correspond to zero functions and contribute
nothing.  The aggregate score has positive variance and
attains the ratio.  No inverse of a potentially singular bank
Gram is required.
\end{proof}

\subsection{A sharp ultraviolet dilution law for the full local bank}
\label{r7:dilution}

The limit below is genuinely joint in the reference problem:
\(n\to\infty\) and \(L/n\to\infty\).  Thus the torus is large
compared with the time separation.  It does not assume a
uniform compact-Wilson-to-Gaussian comparison.

\begin{theorem}[Optimal microscopic-bank loss]
\label{r7:loss}
In this joint limit,
\[
 \begin{split}
 H_{n,L}&\sim\frac9{1024\pi^2}\,n^{-5},\\
 B_{n,L}(q_0)&\sim\frac{45}{16\pi^2}\,n^{-6},\qquad
 J_L(q_0)=504,\\
 \boxed{\quad
 \rho_{n,L}(q_0)\sim\frac{25}{14\pi^2}\,n^{-7}.
 \quad}
 \end{split}
 \tag{V7.8}\label{r7:loss-eq}
\]
The last fraction is optimal over the full elementary bank.
In particular its capture fraction is not bounded below by a
positive constant independent of depth and volume in this
reference law.
\end{theorem}
\begin{proof}
Using the multiplicity six, \eqref{r7:traces-eq} reads
\[
 \begin{split}
 H_{n,L}&=\frac3V\sum_{k\ne0}
       \frac{\lambda(k)^2}{R(\lambda(k))^2}e^{-4n\theta(\lambda(k))},\\
 B_{n,L}(q_0)&=\frac6V\sum_{k\ne0}
       \frac{\lambda(k)^2}{R(\lambda(k))}e^{-2n\theta(\lambda(k))},\\
 J_L(q_0)&=\frac{12}V\sum_{k\ne0}\lambda(k)^2 .
 \end{split}
 \tag{V7.9}\label{r7:sums}
\]
Since \(\lambda=6-2\sum_\mu\cos k_\mu\), its lattice average
square equals \(36+4(3/2)=42\) for every \(L\ge3\).
The zero mode has zero contribution.  This proves \(J=504\).

At \(k=0\),
\[
 \lambda(k)=|k|^2+O(|k|^4),\quad
 \theta(\lambda(k))=|k|+O(|k|^3),\quad
 R(\lambda(k))=2|k|(1+O(|k|^2)).
 \tag{V7.10}\label{r7:smallk}
\]
Rescale \(k=y/n\).  The normalized sums have radial limits
\[
 \begin{split}
 n^5H_{n,L}&\longrightarrow
          \frac3{8\pi^2}\int_0^\infty r^4e^{-4r}\,dr,\\
 n^6B_{n,L}(q_0)&\longrightarrow
          \frac3{2\pi^2}\int_0^\infty r^5e^{-2r}\,dr.
 \end{split}
 \tag{V7.11}\label{r7:radial}
\]
Here is a justification that includes the varying torus.
Use representatives \(k_\mu\in[-\pi,\pi)\).
On this cube, \(\lambda\) is comparable to \(|k|^2\),
\(\theta(\lambda)\ge c|k|\), and \(R\) is comparable to
\(\sqrt\lambda\), with positive absolute constants.
After rescaling, the summands, with their powers of \(n\)
removed, are bounded by
\(C(1+|y|^5)e^{-c|y|}\).
They converge uniformly on compact sets to the displayed
radial functions; defining their values at zero by continuity
is legitimate.  The rescaled mesh \(2\pi n/L\) tends to zero.
For mesh at most one, cover each lattice point by its grid
cube.  The supremum of the exponential majorant on that cube
is bounded by \(C(1+|y|^5)e^{-c|y|/2}\) at every point of the
cube, after adjusting the constant.  The summed tail outside
radius \(M\) is therefore bounded by the tail of an integrable
function, uniformly in the mesh and \(n\).  Compact Riemann
convergence and then \(M\to\infty\) prove
\eqref{r7:radial}.

The two integrals equal \(4!/4^5=3/128\) and
\(5!/2^6=15/8\).  Insert their values to obtain the first two
constants in \eqref{r7:loss-eq}.  Finally
\[
 \frac{(45/16)^2}{504(9/1024)}=\frac{25}{14},
\]
with the remaining factor \(\pi^{-2}\).
Lemma~\ref{r7:symmetry} proves optimality over the full bank.
\end{proof}

This loss does not say the remote covariance vanishes faster
than its own natural scale: it is of order \(n^{-5}\).
Rather, the response is of order \(n^{-6}\) while the
microscopic score variance stays of order one.  Its ultraviolet
part does not track the signal at momenta of order \(n^{-1}\).
Increasing the number of translates of the same microscopic
field cannot fix that mismatch for this source.

There is no conflict with V6's more-than-\(4/5\) retention.
That was a different, fixed open box and a two-step limit.
Nor does this theorem disprove a native nonperturbative
confinement mechanism.  It rules out one proposed
scale-independent use of the reference certificate.

\subsection{A single flowed score retains a fixed fraction}
\label{r7:flow}

On the reference transverse space, spatial flow is
\(x\mapsto e^{-s\mathsf L}x\).
Its flowed energy has coefficient
\[
 q_s(\lambda)=\lambda e^{-2s\lambda},\qquad
 F_s(x)=-\frac1{2\sqrt V}x^{\mathsf T}q_s(\mathsf L)x,
 \qquad X_s=\nabla F_s .
 \tag{V7.12}\label{r7:flow-field}
\]
The flow parameter \(s\) has units of squared spatial lattice
distance.  The transfer half-separation \(n\) has units of
lattice time.  They are not the same clock.

\begin{theorem}[Depth-matched smoothing]
\label{r7:flow-theorem}
Fix \(\alpha>0\) and take \(s=\alpha n^2\).
In the joint limit of Theorem~\ref{r7:loss},
\[
 \rho_{n,L}(q_{\alpha n^2})\longrightarrow
 \mathcal R(\alpha)
 =\frac{\left(\displaystyle\int_0^\infty
                 r^5e^{-2r-2\alpha r^2}\,dr\right)^2}
        {\left(\displaystyle\int_0^\infty
                 r^6e^{-4\alpha r^2}\,dr\right)
         \left(\displaystyle\int_0^\infty
                 r^4e^{-4r}\,dr\right)}
 \in(0,1).
 \tag{V7.13}\label{r7:retention-flow}
\]
In particular,
\[
                         \mathcal R(1/2)>8/9 .
 \tag{V7.14}\label{r7:certified-flow}
\]
This is a limit constant with a rigorous rational lower
certificate, not an asserted \(8/9\) bound at every finite
\(L,n\).
\end{theorem}
\begin{proof}
Let
\[
 I_5(\alpha)=\int_0^\infty r^5e^{-2r-2\alpha r^2}\,dr,\quad
 I_6(\alpha)=\int_0^\infty r^6e^{-4\alpha r^2}\,dr,\quad
 I_4=3/128 .
\]
The same rescaled Riemann proof gives
\[
 n^6 B_{n,L}(q_{\alpha n^2})\to\frac3{2\pi^2}I_5(\alpha),
 \quad
 n^7 J_L(q_{\alpha n^2})\to\frac6{\pi^2}I_6(\alpha),
 \quad
 n^5H_{n,L}\to\frac3{8\pi^2}I_4 .
 \tag{V7.15}\label{r7:flow-scales}
\]
For the score sum the Gaussian factor \(e^{-4\alpha n^2\lambda}\)
supplies the required integrable majorant.  Constants cancel
in \(B^2/(JH)\), proving \eqref{r7:retention-flow}.
Positivity is immediate.  Cauchy--Schwarz applied to
\(r^3e^{-2\alpha r^2}\) and \(r^2e^{-2r}\) gives the upper
bound one; these functions are not proportional, so it is
strict.

We give a finite rational certificate for \(\alpha=1/2\).
Then
\[
 I_6(1/2)=\frac{15\sqrt{2\pi}}{256},\qquad
 I_5(1/2)=\int_0^\infty r^5e^{-r^2-2r}\,dr .
\]
For \(x\ge0\), Taylor's theorem with its signed integral
remainder gives
\(e^{-x}\ge\sum_{k=0}^{199}(-x)^k/k!\).
Integrate that lower polynomial with \(x=r^2+2r\) on
\(0\le r\le6\).  The completely rational resulting number is
\[
 Q=\sum_{k=0}^{199}\frac{(-1)^k}{k!}
       \sum_{j=0}^k \binom{k}{j}2^{\,k-j}
                     \frac{6^{\,6+k+j}}{6+k+j}
                >\frac{5537}{100000}.
 \tag{V7.16}\label{r7:rational-integral}
\]
The last comparison is a finite rational inequality; the
diagnostic evaluates this displayed sum with integer
numerators and denominators, with no quadrature or rounding.
The positive omitted tail gives \(I_5(1/2)\ge Q\).
Using \(\pi<22/7\) gives \(\sqrt{2\pi}<251/100\), since
\(44/7<(251/100)^2\).  Hence
\[
 \mathcal R(1/2)>
 \left(\frac{5537}{100000}\right)^2
       \frac{32768}{45(251/100)}
 =\frac{3924271232}{4412109375}>\frac89 .
 \tag{V7.17}\label{r7:rational-rho}
\]
All inequalities needed for this certificate are now explicit.
\end{proof}

The scale choice is forced by the reference signal:
its time attenuation is \(e^{-2n|k|}\) at low momentum.
The spatial heat multiplier retains the same region when
\(s\asymp n^2\).  At lattice spacing \(a\), fixed half-separation
\(t=an\) gives physical spatial flow time
\(\tau=a^2s=\alpha t^2\).  This is dimensional bookkeeping,
not a calibration of a native Yang--Mills mass.

\paragraph{If the outer source is already spatially flowed.}
Take its reference flow time to be \(\beta n^2\), with fixed
\(\beta\ge0\).  Replace \(a_n\) by
\(\lambda e^{-2n\theta-2\beta n^2\lambda}\).
The same proof gives bare-bank retention of order \(n^{-7}\),
with a positive \(\beta\)-dependent constant, while the
\(\alpha n^2\)-flowed score has limiting fraction
\[
 \frac{\left(\displaystyle\int_0^\infty
       r^5e^{-2r-2(\alpha+\beta)r^2}\,dr\right)^2}
      {\left(\displaystyle\int_0^\infty r^6e^{-4\alpha r^2}\,dr\right)
       \left(\displaystyle\int_0^\infty
                           r^4e^{-4r-4\beta r^2}\,dr\right)}>0 .
 \tag{V7.18}\label{r7:outer-flow}
\]
Thus the diagnosis is not tied to omitting flow from the outer
source.  The numerical \(8/9\) certificate above is asserted
only for \(\beta=0,\alpha=1/2\).

\subsection{The exactly matched field is a positive flow mixture}
\label{r7:mixture}

The previous repair uses only one flow time.  The exact
reference filter has an equally explicit form.

\begin{theorem}[Positive heat-mixture representation]
\label{r7:perfect}
For \(\lambda>0\), define
\[
 q_n^*(\lambda)=
       \frac{\lambda e^{-2n\theta(\lambda)}}{2R(\lambda)}.
 \tag{V7.19}\label{r7:qstar}
\]
On every finite transverse reference torus its score satisfies
\(s_{q_n^*}=m_n\), and hence \(B=J=H\) and retention is exactly
one.  Moreover,
\[
 \boxed{\quad
 q_n^*(\lambda)
   =\int_0^\infty e^{-4s}I_{2n}(4s)\,
                           \lambda e^{-2s\lambda}\,ds .
 \quad}
 \tag{V7.20}\label{r7:positive-mixture}
\]
The weight is nonnegative.  It is not a probability density
in \(s\); the formula is an improper, convergent operator
integral on the strictly positive transverse spectrum.
\end{theorem}
\begin{proof}
The coefficient equality
\(R(\lambda)q_n^*(\lambda)=a_n(\lambda)/2\) gives
\(s_{q_n^*}=m_n\), including its centering constant.

For the integral put \(p_t(j)=e^{-2t}I_j(2t)\).
The Bessel generating series, or the continuous-time
nearest-neighbor walk with rate one to each neighbor, gives
\(p_t(j)\ge0\), \(\sum_jp_t(j)=1\), \(p_0(j)=\mathbf1_{j=0}\),
and
\(\partial_tp_t(j)=p_t(j-1)+p_t(j+1)-2p_t(j)\).
For \(\lambda>0\) the integrals
\(g_j=\int_0^\infty e^{-\lambda t}p_t(j)\,dt\)
are summable in \(j\).  Integrating the differential equation
gives
\[
 (\lambda+2)g_j-g_{j-1}-g_{j+1}=\mathbf1_{j=0}.
\]
The decaying characteristic root is \(e^{-\theta(\lambda)}\).
The summable solution is unique and equals
\[
 g_j=\frac{e^{-|j|\theta(\lambda)}}{R(\lambda)}.
 \tag{V7.21}\label{r7:resolvent}
\]
Indeed the recurrence holds off zero; at zero its jump is
\((\lambda+2)-2e^{-\theta}=R\).  The growing characteristic
solution is excluded on either half-line by summability.
Multiply \(g_{2n}\) by \(\lambda/2\) and put \(t=2s\).
This proves \eqref{r7:positive-mixture}.

For a fixed positive finite spectrum, the scalar integrals
converge absolutely and may be combined by the spectral
theorem.  The extra \(e^{-2s\lambda}\) is essential.  Indeed
\(g_{2n}\to\infty\) as \(\lambda\downarrow0\), by
\eqref{r7:resolvent}.  Monotone convergence and \(t=2s\)
therefore show that the weight alone has infinite integral.
We do not normalize it as a mixing probability.
\end{proof}

At fixed \(L,n\), restrict the integral to
\(\varepsilon\le s\le M\).  Its coefficients increase to
\(q_n^*\) as \(\varepsilon\downarrow0,M\uparrow\infty\).
Because there are finitely many positive eigenvalues, the
corresponding scores converge to \(m_n\) in \(L^2(\nu_L)\).
Their optimized capture fractions therefore tend to one.
This fixed-spectrum convergence is not a claimed uniform
tail bound in the native measure.  For an outer flow
\(\beta n^2\), the reference mixture simply shifts its field
times from \(s\) to \(s+\beta n^2\).

The representation is useful because it prescribes a positive
bank of spatially flowed energy gradients, rather than
requiring an inverse covariance fitted to unknown native
correlations.  It still contains arbitrarily large flow
times before truncation.  Uniformly controlling its native
tails and derivatives is a real obligation, not supplied by
positivity of the reference weight.

\subsection{A well-defined native candidate and the precise unpaid estimate}
\label{r7:native}

Return to the original finite reflection-symmetric Wilson
history and its \emph{actual} middle marginal \(p_I\,dH\).
Let \(\Phi_s\) be the spatial Wilson gradient flow on that
slice, with normalization chosen so that its flat transverse
linearization is \(e^{-s\mathsf L}\).
For a cubic periodic middle slice define the smooth function
\[
 F_s^{\rm nat}(U_I)=V^{-1/2}
                \sum_{p\ {\rm spatial}}\{w_p(\Phi_sU_I)-1\},
 \qquad X_s^{\rm nat}=\nabla_I F_s^{\rm nat}.
 \tag{V7.22}\label{r7:native-field}
\]
Finite-dimensional smooth flow on a compact group manifold
exists for every finite \(s\); the field is gauge equivariant.
It is not supported on only four original links.
The analogous definition works with the stated finite
spatial boundary conditions.  Its normalization and the
outer observable must be kept fixed when a response is
evaluated.

At finite \(0<u<v<\infty\) one may also define, without knowing
a native covariance,
\[
 F_{n;u,v}^{\rm nat}
       =\int_u^v e^{-4s}I_{2n}(4s)F_s^{\rm nat}\,ds.
 \tag{V7.23}\label{r7:native-mixture}
\]
This is a well-defined smooth, gauge-invariant function.
We do \emph{not} pass to \(u=0,v=\infty\) in the native
derivatives or score norms.  The fact that the corresponding
reference integral converges is not such a justification.

For either smooth candidate \(F\), V5--V6 give exactly
\[
 \begin{split}
 T_F&=-\mathcal L_I F,\qquad s_F=\E[T_F\mid U_I]
       =-\Delta_I F-\langle\nabla_I F,\nabla_I\log p_I\rangle,\\
 B_F&=\E\langle\nabla_I F,\nabla_I m\rangle,\qquad
 J_F=\E s_F^2,\qquad
 G\ge B_F^2/J_F\quad(J_F>0).
 \end{split}
 \tag{V7.24}\label{r7:native-interface}
\]
Here \(m\) is the actual conditional outer expectation;
every crossing plaquette and its conditional normalizer
remains present.  The inherited Bochner argument gives
\[
 J_F\le\E T_F^2
 \le \E\|\nabla_I^2F\|_{\rm HS}^2
                 +(2+24\gamma)\E|\nabla_I F|^2 .
 \tag{V7.25}\label{r7:native-budget}
\]
Equation \eqref{r7:native-budget} holds for any such smooth
\(F\); it does not evaluate its right side at \(s\asymp n^2\).
In particular, the elementary-bank constant of V6 must not be
reused after replacing its four-link functions by
\eqref{r7:native-field}.

The new reference calculation specifies an informative target:
after a justified native source normalization, retain a
response of order \(n^{-6}\) while reducing the associated
score variance to order \(n^{-7}\), rather than paying an
order-one microscopic variance.  Those exponents refer to
the massless Gaussian calibration.  They are not asserted
native asymptotics at arbitrarily large physical separation.
For a positive physical mass, long physical-time behavior
would require different, exponential information.

More quantitatively, suppose a particular common normalization,
geometry, and time range have been fixed.  If one proves
\[
 |B_F-b_{n,L}|\le\epsilon b_{n,L},\qquad
 J_F\le(1+\delta)j_{n,L},\qquad
 b_{n,L}>0,\quad 0\le\epsilon<1,\quad\delta\ge0,
 \tag{V7.26}\label{r7:import-hyp}
\]
where \(b,j\) are the explicitly evaluated reference response
and score norm for the chosen filter, then
\[
 G\ge\frac{(1-\epsilon)^2}{1+\delta}
                         \frac{b_{n,L}^2}{j_{n,L}} .
 \tag{V7.27}\label{r7:import-conclusion}
\]
This last implication is elementary substitution into
\eqref{r7:native-interface}, not a newly paid comparison
theorem.  No estimate of \(\epsilon,\delta\) uniform in the
required Wilson limits has been established.  The reference
target \(b^2/j\) is explicit; \(B_F\) is defined by a
half-history response, not by the unknown covariance \(G\).
These are measurable, noncircular tasks, but substantial ones.

The next attack is therefore the native response and derivative
budget for \eqref{r7:native-field} at flow radius comparable to
the source separation, or for a finite positive mixture in
\eqref{r7:native-mixture}.  Enlarging the microscopic translate
bank alone is now ruled out by an exact reference obstruction.
A proposed estimate must also distinguish stationary vacuum
from finite open history.  Neither reference retention nor a
source lower Gram supplies the common growing-depth
\emph{upper} spectral inequality needed for a mass gap.

\subsection{Verification and result ledger}
\label{r7:verification}

The independent diagnostic is
\begin{center}\small
\path{scripts/verify_ym_restart_v7_scale_response.py}.
\end{center}
It checks the mode covariance, response, and perfect-score
identities on twelve exact rational-frequency/depth fixtures;
the Bessel-resolvent difference equation; the coefficient
\(25/14\); the torus moment \(42\); and the explicit degree-199
rational certificate \eqref{r7:rational-integral}.
Its optional finite-torus sums use floating arithmetic and
are labeled as diagnostics, not certified native enclosures.
The all-volume/depth asymptotic proof is the dominated
rescaled-sum argument above, not an extrapolation of those
finite sums.

\emph{New here:} optimal elementary-bank dilution at growing
depth; its exact \(n^{-7}\) coefficient; a one-flow repair
with certified limiting retention above \(8/9\); a positive
Bessel-weighted flow representation of the exactly matching
reference field; and a corresponding native finite-mixture
candidate with its obligations stated.

\emph{Inherited or standard:} the transverse Gaussian
dispersion and multiplicities, Gaussian integration by parts,
symmetry averaging, spatial heat smoothing, the Bessel heat
kernel, and V5--V6's native marginal-score residual identities.
The comparison between these reference results and the
nonlinear Wilson law is \emph{not} claimed.

The V6 body is preserved byte for byte apart from the displayed
version title.  Removing this delimited append and reversing
that title recovers V6 exactly.  Only the immediate active
source is replaced; archives and companion sources are left
unchanged.  All checker output and compilation artifacts are
outside \path{latex_notes}.  The next companion review and
broad literature sweep remain V10.

\begin{center}
\fbox{\begin{minipage}{0.94\textwidth}
\textbf{V7 status.}
A fixed microscopic score bank has a proved reference-scale
failure, and spatially flowed scores have an explicit
reference-scale repair.  Native uniform response, flowed-score
norm control, vacuum preparation, the common upper spectral
estimate, and the interacting continuum Yang--Mills mass gap
remain unproved.
\end{minipage}}
\end{center}
% END CONSOLIDATED V7 APPEND

% BEGIN CONSOLIDATED V8 APPEND -- 2026-09-19
\clearpage
\section{V8: Native flow-score tails and an explicit cutoff repair}
\label{r8:context}

V7 found a genuine reference-scale improvement from simultaneous spatial
smoothing, but left its nonlinear score norm unpaid.  This append tests
an upstream issue: smoothing an observable does not necessarily smooth
its derivatives with respect to the \emph{original} links.  We prove a
quantitative obstruction and a limited repair.  The obstruction uses
the already solved isolated-plaquette subflow; its score is evaluated
under the original full Wilson law, not a substituted product law.
The repair has a full-native norm bound, but its response guarantee is
proved only for a one-plaquette probe.

These distinctions are essential.  The isolated-plaquette flow below is
\emph{not} the simultaneous spatial flow in \eqref{r7:native-field}.
We do not prove that the latter has a divergent score norm, nor that
our local cutoff fields implement V7's spatial heat filter.  We keep
the continuum mass-gap goal open.

\subsection{What is inherited and what is being added}

Archive \(C\), V74 already solves the four-link plaquette flow and
computes its pushed-forward Haar density.  Archive \(C\), V68 constructs
the full spatial flow, and its V83 evaluates a short-flow,
strong-coupling response jet.  None of those calculations estimates
the long-flow \(L^2\) norm of the \emph{pulled-back gradient score}
under the original input law.  We use their geometric solution, without
counting it as a new result.

The targeted primary-source check was
\href{https://arxiv.org/html/1006.4518}{M.~L\"uscher,
\emph{Properties and uses of the Wilson flow in lattice QCD},
arXiv:1006.4518}, especially its finite-lattice flow and change-of-variable
discussion.  Its smoothing interpretation does not supply the
original-input score estimate needed here.  All quantitative claims
in this append are proved below.  This targeted check does not replace
the scheduled V10 broad literature sweep or companion review.

\subsection{A radial score projection in the full Wilson law}
\label{r8:radial}

Work in the finite four-dimensional Wilson geometry of V6, with
\(\gamma\ge0\), at most six action plaquettes incident to each link,
and no repeated link within an elementary plaquette.  Fix a spatial
plaquette \(p\) with four distinct free middle links, and write
\(w=w_p\).  The metric on each \(\SU(2)\) factor is the unit
three-sphere metric.  In particular,
\[
 |\nabla_I w|^2=4(1-w^2),\qquad \Delta_Iw=-12w.
 \tag{V8.1}\label{r8:radial-geometry}
\]
All expectations in this subsection refer to the full Wilson history,
equivalently its actual middle marginal when the integrand is
middle-measurable.  Arbitrary fixed exterior links are allowed.

\begin{lemma}[Uniform radial density and its logarithmic derivative]
\label{r8:density-lemma}
The distribution of \(w\) has the form
\[
 d\mu_w(w)=\frac2\pi\sqrt{1-w^2}\,h(w)\,dw,\qquad -1<w<1,
 \tag{V8.2}\label{r8:density}
\]
where \(h\) is positive and smooth on the closed interval, with
\[
 e^{-12\gamma}\le h(w)\le e^{12\gamma},\qquad
 |(\log h)'(w)|\le M_\gamma:=6\gamma+12\gamma^2.
 \tag{V8.3}\label{r8:h-bounds}
\]
The bounds are independent of volume, temporal depth and exterior data.
\end{lemma}
\begin{proof}
Condition on every link except one free edge of \(p\).
Its conditional density with respect to normalized Haar measure is
\(\exp(\gamma a\cdot q)/Z(a)\), with \(q\in S^3\) and \(|a|\le6\).
Each incident plaquette contributes one unit linear quaternion
coordinate.  The conditional exponent has oscillation at most
\(12\gamma\), proving the two density bounds.

Multiplication and inversion by the other three plaquette links are
Haar isometries.  Thus one may take \(q\) to be the plaquette
holonomy, with real coordinate \(w\).  At fixed \(w\), average over
its imaginary direction on \(S^2\).  The resulting density relative
to the semicircle law is
\[
 h_a(w)=Z(a)^{-1}e^{\gamma a_0w}\frac{\sinh z}{z},
 \qquad z=\gamma|\mathbf a|\sqrt{1-w^2},
 \tag{V8.4}\label{r8:angular-density}
\]
with the quotient interpreted as one at zero.  This is smooth at the
endpoints by its power series.  For real \(z\ge0\),
\[
                 0\le z\coth z-1\le z^2/3.
\]
For the upper inequality, compare coefficients in
\(z\cosh z-\sinh z\le(z^2/3)\sinh z\).
For \(z^{2k+1}\), \(k\ge1\), the comparison is
\(2k/(2k+1)!\le1/[3(2k-1)!]\), equivalent to \(3\le2k+1\).
The lower inequality follows from the nonnegative coefficients
of \(z\cosh z-\sinh z\).  Consequently
\[
 (\log h_a)'=\gamma a_0-
       \frac{w}{1-w^2}(z\coth z-1),\qquad
 |(\log h_a)'|\le6\gamma+12\gamma^2.
\]
The endpoint values are obtained by continuity.

The unconditional density is the mixture of these conditional
densities over the other links.  The mixing law is independent of
the remaining coordinate \(w\).  Positivity preserves the density
bounds; its logarithmic derivative is a weighted average of the
conditional logarithmic derivatives.  Differentiation under the
mixture is justified by compactness and the uniformly smooth power
series at fixed \(\gamma\).  This proves the lemma.
\end{proof}

Let \(p_I\,dH\) be the actual middle marginal.  For any smooth
\(f:[-1,1]\to\R\), put
\[
 X_f=\nabla_I f(w),\qquad s_f=-\operatorname{div}_{p_I}X_f,
 \qquad J_f=\E s_f^2.
\]
These are not scores of the pushed-forward law.

\begin{lemma}[Exact radial projection]
\label{r8:projection-lemma}
With \(\ell=\log h\),
\[
 \begin{split}
 \E[s_f\mid w]&=\sigma_f(w),\\
 \sigma_f(w)&=-4\big\{(1-w^2)f''(w)
              +[(1-w^2)\ell'(w)-3w]f'(w)\big\},\\
 J_f&\ge\E\sigma_f^2.
 \end{split}
 \tag{V8.5}\label{r8:score-projection}
\]
\end{lemma}
\begin{proof}
For a smooth test function \(g\), integration by parts on the compact
middle-link manifold gives
\(\E[s_fg(w)]=4\E[(1-w^2)f'(w)g'(w)]\).
Integrating this one-dimensional expression by parts against
\eqref{r8:density} gives the displayed formula.  The boundary flux
is proportional to \((1-w^2)^{3/2}h(w)f'(w)\), and vanishes at
both endpoints.  Smooth tests are dense in \(L^2(\mu_w)\), so the
identity determines the conditional expectation.  Conditional
expectation is an \(L^2\) contraction, proving the last assertion.
\end{proof}

\subsection{The inherited subflow and its new original-input tail}
\label{r8:tail}

Flow only the four links of \(p\), using the isolated energy \(1-w\).
Its exact equation and solution, already in archive \(C\), V74, are
\[
 \dot w_s=4(1-w_s^2),\qquad
 f_s(w):=w_s=\frac{w+\tanh(4s)}{1+w\tanh(4s)}.
\]
It is convenient to set
\[
 \begin{split}
 \delta&=e^{-8s},\qquad
 D_\delta(w)=1+w+\delta(1-w),\\
 f_s(w)&=\frac{(1+\delta)w+1-\delta}{D_\delta(w)},\\
 f_s'(w)&=\frac{4\delta}{D_\delta(w)^2},\qquad
 f_s''(w)=-\frac{8\delta(1-\delta)}{D_\delta(w)^3}.
 \end{split}
 \tag{V8.6}\label{r8:flow}
\]
The generating field is \(X_s=\nabla_I f_s(w)\).
The input measure, full action and its conditional normalizers
are all unchanged.

\begin{theorem}[Rare-input score growth despite observable concentration]
\label{r8:asymptotics}
For any fixed positive smooth \(h\) as in \eqref{r8:density}, as
\(s\to\infty\),
\[
 \begin{split}
 \E\sigma_{f_s}^2&\sim24h(-1)e^{4s},\\
 \E|\nabla_I f_s(w)|^2&\sim8h(-1)e^{-4s},\\
 \Var(f_s(w))&\sim16h(-1)e^{-12s}.
 \end{split}
 \tag{V8.7}\label{r8:asymptotic-values}
\]
For the full Wilson marginal this implies
\[
 \liminf_{s\to\infty}e^{-4s}J_{f_s}\ge24h(-1).
 \tag{V8.8}\label{r8:full-liminf}
\]
The first equivalence is for the radial projection, not an asserted
equivalence for the entire full-native score.
\end{theorem}
\begin{proof}
Use \(w=-1+\delta z\), \(0\le z\le2/\delta\), and write
\(A=2+(1-\delta)z\).  The rescaled density is
\[
 \delta^{-3/2}d\mu_w
   =\frac2\pi\sqrt{z(2-\delta z)}\,h(-1+\delta z)\,dz.
\]
Writing \(\sigma^{(0)}\) for \eqref{r8:score-projection} with
\(\ell'=0\), direct substitution gives
\[
 \begin{split}
 \delta\sigma^{(0)}_{f_s}
  &=16\frac{-6+(1+5\delta)z+\delta(1-\delta)z^2}{A^3},\\
 \delta(\sigma_{f_s}-\sigma^{(0)}_{f_s})
  &=-16\delta\frac{z(2-\delta z)}{A^2}\ell'(-1+\delta z).
 \end{split}
 \tag{V8.9}\label{r8:scaled-score}
\]
Hence, at fixed \(z\),
\(\delta\sigma_{f_s}\to16(z-6)/(z+2)^3\).

For completeness, this limit may be integrated.  For
\(0<\delta\le1/2\) and \(0\le z\le2/\delta\),
\(A\ge(2+z)/2\) and \(\delta z\le2\).
The first expression in \eqref{r8:scaled-score} is bounded by
\(C(1+z)^{-2}\); the second has the same bound because
\(\delta\le C/(1+z)\) on this domain.
The constant depends only on the fixed bounds for \(h,h'\).
The rescaled density is at most \(C\sqrt z\,dz\).
Thus the squared-score integrand is dominated by an integrable
multiple of \(\sqrt z(1+z)^{-4}\).
Dominated convergence yields
\[
 \lim_{\delta\downarrow0}\delta^{1/2}\E\sigma_{f_s}^2
   =\frac{512\sqrt2}{\pi}h(-1)
       \int_0^\infty\frac{z^{1/2}(z-6)^2}{(z+2)^6}\,dz
   =24h(-1).
\]

Similarly,
\[
 |\nabla_I f_s|^2
   =64\delta^{-1}\frac{z(2-\delta z)}{A^4},\qquad
 1-f_s=\frac{2(2-\delta z)}{A}.
\]
After multiplying by the rescaled density, the dominating functions
for their respective limits are
\(Cz^{3/2}(1+z)^{-4}\) and \(C\sqrt z(1+z)^{-2}\).
The three integrals used here are
\[
 \begin{split}
 \int_0^\infty\frac{z^{1/2}(z-6)^2}{(z+2)^6}\,dz
       &=\frac{3\pi}{64\sqrt2},\\
 \int_0^\infty\frac{z^{3/2}}{(z+2)^4}\,dz
       &=\frac{\pi}{32\sqrt2},\\
 \int_0^\infty\frac{z^{1/2}}{(z+2)^2}\,dz
       &=\frac{\pi}{2\sqrt2}.
 \end{split}
 \tag{V8.10}\label{r8:beta-integrals}
\]
They follow by \(z=2t/(1-t)\), expansion of \((z-6)^2\)
in the first integral, and the elementary beta integral.
In particular,
\(\E(1-f_s)^2\sim16h(-1)\delta^{3/2}\).

Finally \(D_\delta\ge1+w\), and
\(1-f_s=2\delta(1-w)/D_\delta\).
The function \((1-w)/(1+w)\) is integrable against
\eqref{r8:density}.  Therefore \(\E(1-f_s)=O(\delta)\),
whose square is \(o(\delta^{3/2})\).
This proves the variance equivalence and all of
\eqref{r8:asymptotic-values}.  Apply the projection inequality
to obtain \eqref{r8:full-liminf}.
\end{proof}

Together with the full-native Bochner bound
\eqref{r7:native-budget}, the theorem also gives
\[
 \liminf_{s\to\infty}e^{-4s}
       \E\|\nabla_I^2 f_s(w)\|_{\rm HS}^2\ge24h(-1).
 \tag{V8.11}\label{r8:hessian-tail}
\]
Indeed the gradient term in that bound tends to zero after this
rescaling.  Decreasing energy and decreasing first-derivative
mean square do not control the second-derivative budget.

\begin{theorem}[Explicit full-native lower bound]
\label{r8:uniform-lower}
Under the full Wilson law of Lemma~\ref{r8:density-lemma}, if
\[
 \delta=e^{-8s}\le
       \min\left\{\frac1{32},\frac1{64(1+M_\gamma)}\right\},
\]
then
\[
              \boxed{\quad J_{f_s}\ge
                      \frac1{2\pi}e^{4s-12\gamma}.\quad}
 \tag{V8.12}\label{r8:uniform-score-lower}
\]
This holds uniformly in the stated finite geometries and exterior data.
\end{theorem}
\begin{proof}
On \(1\le z\le2\), the numerator in the first line of
\eqref{r8:scaled-score} is at most
\(-4+14\delta\le-7/2\), and \(A\le4\).
That line is therefore at most \(-7/8\).
The magnitude of the second line is at most
\(16\delta M_\gamma\le1/4\), so
\(|\sigma_{f_s}|\ge1/(2\delta)\) on this layer.
Also \(z(2-\delta z)\ge1\).  By \eqref{r8:h-bounds},
\[
 \Pr\{\,\delta\le1+w\le2\delta\,\}
        \ge\frac2\pi e^{-12\gamma}\delta^{3/2}.
\]
Integrate the squared-score lower bound on this layer and use
\(J_{f_s}\ge\E\sigma_{f_s}^2\).
\end{proof}

For example, along any sequence satisfying
\(0\le\gamma(a)\le C\log(1/a)\) and \(s=\tau/a^2\), \(\tau>0\),
the right side of \eqref{r8:uniform-score-lower} diverges faster
than any power of \(a^{-1}\); its small-\(\delta\) hypothesis
eventually holds.  This is a conditional scaling consequence,
not an assertion that such a trajectory has constructed a continuum
theory.  Nor does a subflow lower bound transfer to the simultaneous
spatial flow.  What is ruled out is a blanket inference from
flowed-observable concentration to a bounded original-input score.

\subsection{A one-cell response exposes the noncommuting limits}
\label{r8:cell}

This subsection alone specializes the law to the literal one-plaquette
Wilson model
\[
 d\nu_\gamma=Z_\gamma^{-1}e^{\gamma w}\,dH^{\otimes4},\qquad
 Z_\gamma=\int e^{\gamma w}\,dH^{\otimes4}.
 \tag{V8.13}\label{r8:cell-law}
\]
Thus \(h_\gamma(w)=e^{\gamma w}/Z_\gamma\) and the complete
score is the radial score.  Take the probe \(A=w-\E w\) and define
\[
 B_s=\E(X_sw)=4\E[(1-w^2)f_s'(w)],\qquad
 \rho_s(\gamma)=\frac{B_s^2}{J_{f_s}\Var(w)}.
 \tag{V8.14}\label{r8:cell-retention}
\]
Here \(\rho_s\) is the ordinary \(L^2(\nu_\gamma)\) projection
fraction of \(A\) onto the one-dimensional score span.
It is \emph{not} a physical opposite-time source Gram.

\begin{proposition}[Loss at fixed coupling; retention at fixed flow]
\label{r8:limits}
For each finite \(\gamma\ge0\), put
\(R_\gamma=\E_{\nu_\gamma}[(1-w)/(1+w)]\).
It is finite and strictly positive, and
\[
 \rho_s(\gamma)\sim
       \frac{32R_\gamma^2}{3h_\gamma(-1)\Var_{\nu_\gamma}(w)}
                     e^{-20s}\quad(s\to\infty).
 \tag{V8.15}\label{r8:cell-loss}
\]
In contrast, for every fixed finite \(s\),
\[
                  \lim_{\gamma\to\infty}\rho_s(\gamma)=1.
 \tag{V8.16}\label{r8:cell-flat}
\]
Consequently the two iterated limits, with \(s\to\infty\) taken
first or last, are zero and one respectively.
\end{proposition}
\begin{proof}
The integrability of \(R_\gamma\) follows from the endpoint
square-root density.  Since \(D_\delta\ge1+w\), dominated
convergence gives
\[
 B_s/\delta=16\E[(1-w^2)/D_\delta^2]\longrightarrow16R_\gamma.
\]
Use \(J_{f_s}\sim24h_\gamma(-1)\delta^{-1/2}\) and
\(\delta^{5/2}=e^{-20s}\) to prove \eqref{r8:cell-loss}.

For the other limit write \(x=1-w\), \(y=\gamma x\).
The density of \(y\) is proportional to
\(e^{-y}y^{1/2}\sqrt{2-y/\gamma}\) on \([0,2\gamma]\).
Dominated convergence, with polynomial moments, gives the gamma
law of shape \(3/2\) and rate one as its limit.
At fixed \(\delta>0\), all derivatives in \eqref{r8:flow}
are bounded and \(f_s'(1)=\delta\).  It follows that
\[
 \gamma^2\Var(w)\longrightarrow\frac32,\qquad
 \gamma B_s\longrightarrow12\delta,\qquad
 J_{f_s}\longrightarrow96\delta^2.
\]
For the last limit use
\(\sigma_{f_s}=12wf_s'-4(1-w^2)f_s''-4\gamma(1-w^2)f_s'\);
in the \(y\) variable it converges in \(L^2\) to
\(\delta(12-8y)\), whose second moment is \(96\delta^2\).
For domination, the fixed-\(\delta\) derivative bounds make its
absolute value at most \(C_\delta(1+y)\).
Substitution in \eqref{r8:cell-retention} proves the assertion.
\end{proof}

Thus a fixed-flow, weak-coupling expansion can be correct and still
fail completely to control a growing-flow score.  This is an explicit
obstruction to interchanging the relevant limits, not an undecidability
statement or a counterexample to the Yang--Mills mass gap.

\subsection{Cut off the generating field, not the native measure}
\label{r8:cutoff}

Define a fixed \(C^2\) function on \([-1,1]\) by
\[
 \chi(w)=
 \begin{cases}
 0,&w\le-\tfrac12,\\
 6t^5-15t^4+10t^3,\quad t=2w+1,&-\tfrac12<w<0,\\
 1,&w\ge0.
 \end{cases}
 \tag{V8.17}\label{r8:chi}
\]
It satisfies \(0\le\chi\le1\) and \(|\chi'|\le15/4\):
the derivative of its polynomial with respect to \(w\) is
\(60t^2(1-t)^2\).  Its first two derivatives match at the joins.

For \(0\le\delta\le1\), define
\[
 a_\delta(w)=
 \begin{cases}
 0,&w\le-\tfrac12,\\
 4\chi(w)/D_\delta(w)^2,&w>-\tfrac12,
 \end{cases}
 \qquad Y_{\delta,p}=a_\delta(w_p)\nabla_Iw_p.
 \tag{V8.18}\label{r8:tamed-field}
\]
The formula at \(\delta=0\) is nonsingular because of the cutoff.
For positive \(\delta=e^{-8s}\),
\[
                  \chi(w_p)X_s=\delta Y_{\delta,p}.
 \tag{V8.19}\label{r8:rescaled-cutoff}
\]
The field \(Y_{\delta,p}\) is a gradient: take any primitive
\(H_\delta'=a_\delta\) and use \(H_\delta(w_p)\).
This primitive is \(C^3\), sufficient for the integrated
Bochner identity (or its smooth-approximation form).
No law has been conditioned on a good event and no configuration
has been removed.

\begin{lemma}[Uniform derivative bounds]
\label{r8:cutoff-bounds}
For all \(\delta\in[0,1]\) and \(w\in[-1,1]\),
\[
 \begin{gathered}
 a_\delta(1)=1,\qquad 0\le a_\delta\le16,\qquad
 |a_\delta'|\le124,\qquad
 |a_\delta(w)-1|\le124(1-w),\\
 |Y_{\delta,p}|^2\le1024,\qquad
 \|\nabla_I^2H_\delta(w_p)\|_{\rm HS}\le608.
 \end{gathered}
 \tag{V8.20}\label{r8:cutoff-estimates}
\]
\end{lemma}
\begin{proof}
On the support of \(\chi\), \(D_\delta\ge1+w\ge1/2\).
The first derivative of \(a_\delta\) is
\[
 a_\delta'=\frac{4\chi'}{D_\delta^2}
             -\frac{8\chi(1-\delta)}{D_\delta^3}.
\]
The two terms are bounded by \(60\) and \(64\) respectively.
Also \(D_\delta(1)=2\), giving \(a_\delta(1)=1\).
The mean-value theorem gives the fourth bound, including across
the cutoff joins.  Now \(|\nabla_Iw_p|\le2\), and
\[
 \nabla_I^2H_\delta(w_p)
       =a_\delta\nabla_I^2w_p+
                     a_\delta'\nabla_Iw_p\otimes\nabla_Iw_p .
\]
Use V6's bound \(\|\nabla_I^2w_p\|_{\rm HS}\le\sqrt{48}<7\):
the right side has norm at most
\(16\sqrt{48}+124\cdot4<608\).
\end{proof}

\begin{theorem}[Full-native cutoff-bank score budget]
\label{r8:bank}
Let \(\mathcal P\) be any bank of distinct spatial plaquettes
with four free middle links each.  Let \(J^{Y_\delta}\) be the
Gram matrix of their actual marginal scores
\(-\operatorname{div}_{p_I}Y_{\delta,p}\).
For every \(\gamma\ge0\) and \(\delta\in[0,1]\),
\[
 \boxed{\quad
 0\preceq J^{Y_\delta}
       \preceq(1\,486\,848+98\,304\gamma)I .
 \quad}
 \tag{V8.21}\label{r8:cutoff-bank}
\]
For the original-amplitude cutoff fields
\(\chi(w_p)\nabla_I f_s(w_p)\), the right side is multiplied
by \(\delta^2=e^{-16s}\).
All constants are independent of volume, temporal depth, flow time,
bank size and fixed exterior links in the stated geometry.
\end{theorem}
\begin{proof}
For real coefficients \(c_p\), let
\(F_c=\sum_{p\in\mathcal P}c_pH_\delta(w_p)\).
As in V6, a spatial link occurs in at most four members of the
bank, and a fixed Hessian block receives at most four contributions.
Cauchy--Schwarz and \eqref{r8:cutoff-estimates} imply, pointwise,
\[
 |\nabla_IF_c|^2\le4096\sum_pc_p^2,\qquad
 \|\nabla_I^2F_c\|_{\rm HS}^2
                  \le1\,478\,656\sum_pc_p^2.
\]
Apply \eqref{r7:native-budget}, which retains every incident
temporal plaquette and the original conditional normalizer:
\[
 \E s_{F_c}^2
    \le[1\,478\,656+4096(2+24\gamma)]\sum_pc_p^2 .
\]
This is \eqref{r8:cutoff-bank}.  Score linearity and
\eqref{r8:rescaled-cutoff} prove the last assertion.
\end{proof}

Because \(\chi=1\) in a neighborhood of \(w=1\), the cutoff
leaves the uncut field exactly unchanged there, for every fixed
finite \(s\).  In particular it preserves its flat first jet.
The price is a new, explicitly defined generating field, not a
new measure or a new outer observable.

\subsection{An all-flow response repair, only in the one-cell model}
\label{r8:repair}

We can show that the preceding norm repair does not destroy the
one-cell probe at weak coupling, uniformly in flow time.
Continue with \eqref{r8:cell-law} and \(A=w-\E w\).
For \(Y_\delta=a_\delta(w)\nabla w\) put
\[
 \begin{split}
 s_\delta^\chi&=12a_\delta w-4(1-w^2)a_\delta'
                         -4\gamma(1-w^2)a_\delta,\\
 B_\delta^\chi&=\E[Y_\delta w],\qquad
 J_\delta^\chi=\E(s_\delta^\chi)^2,\qquad
 \rho_\delta^\chi
       =\frac{(B_\delta^\chi)^2}{J_\delta^\chi\Var(w)} .
 \end{split}
 \tag{V8.22}\label{r8:cut-cell-score}
\]
For positive \(\delta\), this capture fraction is identical to
the fraction for \(\delta Y_\delta=\chi X_s\); common nonzero
scalar factors cancel.

\begin{theorem}[Uniform one-cell capture after cutoff]
\label{r8:uniform-repair}
For all \(\gamma\ge2\) and \(\delta\in[0,1]\),
\[
 0\le1-\rho_\delta^\chi\le
       \frac{161\,425\,522\,688}{\gamma^2}.
 \tag{V8.23}\label{r8:uniform-cell-capture}
\]
In particular \(\inf_{s\ge0}\rho_{e^{-8s}}^\chi\to1\) as
\(\gamma\to\infty\).
The constant is deliberately coarse; the conclusion is a uniform
limit, not a useful moderate-coupling numerical threshold.
For each fixed finite \(\gamma\ge0\), the fraction also has a
strictly positive limit as \(s\to\infty\).
\end{theorem}
\begin{proof}
Write \(x=1-w\).  Define
\[
 r=8\gamma(w-1)-s_\delta^\chi+12.
\]
Expansion of \eqref{r8:cut-cell-score} gives the exact identity
\[
 r=12(1-a_\delta+xa_\delta)
      +4(2x-x^2)a_\delta'
      +\gamma[8x(a_\delta-1)-4x^2a_\delta].
\]
All coefficients \(a_\delta,a_\delta'\) in this identity are
evaluated at \(w\).  Using \eqref{r8:cutoff-estimates},
\[
 \begin{split}
 |12(1-a_\delta)+12xa_\delta|&\le1680x,\\
 |4(2x-x^2)a_\delta'|&\le992x,\\
 |8x(a_\delta-1)-4x^2a_\delta|&\le1056x^2 .
 \end{split}
\]
Thus
\[
                       |r|\le2672x+1056\gamma x^2.
 \tag{V8.24}\label{r8:residual-bound}
\]

For \(\gamma\ge1\), the unnormalized \(x\)-density is
\(e^{-\gamma x}\sqrt{x(2-x)}\) on \([0,2]\).
Its integral is at least
\((2/(3e))\gamma^{-3/2}\), by restricting to \(0\le x\le1/\gamma\).
The numerator of its \(k\)th moment is at most
\(\sqrt2\,\Gamma(k+3/2)\gamma^{-k-3/2}\).
Using \(e<3\) and \(\sqrt{2\pi}<3\) therefore gives
\[
                 \E x^2\le26\gamma^{-2},\qquad
                 \E x^4\le399\gamma^{-4}.
\]
Equation \eqref{r8:residual-bound} implies
\[
 \E r^2\le C_r\gamma^{-2},\qquad
 C_r=2(2672^2\cdot26+1056^2\cdot399)=1\,261\,136\,896 .
 \tag{V8.25}\label{r8:residual-moment}
\]
These estimates are uniform in \(\delta\).

We also need a lower variance bound; an upper moment bound alone
would not suffice.  For \(\gamma\ge2\), \(y=\gamma x\) has
unnormalized density \(e^{-y}\sqrt y\sqrt{2-y/\gamma}\).
Its normalizer is at most
\(\sqrt2\,\Gamma(3/2)<2\).
On the intervals \(I=[1/4,1/2]\) and \(K=[1,5/4]\), the last
square root is at least one.  Hence
\[
 \Pr(y\in I)\ge e^{-1/2}/16,\qquad
 \Pr(y\in K)\ge e^{-5/4}/8.
\]
The intervals are separated by \(1/2\).
Using two independent copies of \(y\) in the identity
\(\Var(y)=\frac12\E(y-y')^2\) gives
\[
 \Var(y)\ge e^{-7/4}/512>1/8192,\qquad
                \Var(w)\ge\frac1{8192\gamma^2}.
 \tag{V8.26}\label{r8:variance-lower}
\]
The strict inequality follows, for example, from \(e<3\).

The score has mean zero.  Consequently
\(8\gamma(w-\E w)-s_\delta^\chi=r-\E r\).
The squared error of the particular coefficient \(1/(8\gamma)\)
is at most \(C_r/(64\gamma^4)\).  Orthogonal projection onto the
score span can only reduce this error.  After division by
\eqref{r8:variance-lower},
\[
 1-\rho_\delta^\chi
   \le\frac{128C_r}{\gamma^2}
   =\frac{161\,425\,522\,688}{\gamma^2}.
\]
The denominator is nonzero: integration by parts gives
\(B_\delta^\chi=4\E[a_\delta(1-w^2)]>0\);
thus \(J_\delta^\chi>0\), and Cauchy--Schwarz gives
\(0<\rho_\delta^\chi\le1\).

Finally \(a_\delta\to a_0\) in \(C^1([-1,1])\) as
\(\delta\downarrow0\), by its nonsingular support and matching
derivatives at the joins.  At fixed \(\gamma\), both the response
and the score second moment converge, with positive limiting
response and score norm.  This proves the last assertion.
\end{proof}

The theorem repairs the noncommuting-limit failure for this explicit
probe without dropping the transition-derivative term of the cutoff.
It is not a uniform response theorem for the full four-dimensional
history.  Nor does it prove vacuum preparation or a physical mass gap.

\subsection{The exact cutoff accounting required for V7's actual flow}
\label{r8:next}

For any \(C^1\) middle field \(X\) and scalar cutoff \(\eta\), with
the original middle marginal fixed, the exact product rule is
\[
 \begin{split}
 s_{\eta X}&=\eta s_X-X\eta,\\
 B_{\eta X}&=\E[\eta Xm],\\
 B_X-B_{\eta X}&=\E[(1-\eta)Xm],\\
 J_{\eta X}&\le2\E[\eta^2s_X^2]+2\E[(X\eta)^2].
 \end{split}
 \tag{V8.27}\label{r8:general-cutoff}
\]
Here \(m\) is the actual conditional outer expectation of V7.
The second line follows by integration by parts, not by replacing
\(m\) with an equal-time probe.  The last line follows by squaring
the first and applying \((a-b)^2\le2a^2+2b^2\).
Thus a useful native truncation must pay both the score transition
term \(X\eta\) and the omitted, derivative-weighted response.
A small probability for the bad set, by itself, pays neither term.

For the isolated subflow, \eqref{r8:uniform-score-lower} shows why
a naive probability-only argument can fail.  The bad layer has
width \(e^{-8s}\), but the squared score there has order \(e^{16s}\).
Its density suppression does not compensate for that derivative
growth.  The cutoff succeeds here because its support is fixed away
from that singular input layer, its derivative cost is explicit,
and the one-cell response is estimated independently.

The new full-native score budget does \emph{not} remove V7's
spatial-scale bottleneck.  Each \(Y_{\delta,p}\) still depends on
only four original links.  Its flat linearization, after the
irrelevant scalar normalization, is the same local plaquette field
as before.  It has not acquired the spatial heat filter
\(\lambda e^{-2s\lambda}\), so the reference elementary-bank
loss in V7 still applies to its leading flat bank.

The next meaningful extension is a generating field that has both
a genuine spatial smoothing radius and paid nonlinear input
derivatives.  For the simultaneous flow in \eqref{r7:native-field},
this requires bounds on its actual first and second variations,
together with a weighted-tail or cutoff analysis of the form
\eqref{r8:general-cutoff}.  An alternative field is acceptable if
its native score budget and its spatially separated source response
are proved under the same unchanged Wilson law.  Repeating the
one-cell weak-coupling argument on each face independently would
not supply those correlations.

\subsection{Verification, preservation, and the result boundary}
\label{r8:verification}

The checker
\begin{center}\small
\path{scripts/verify_ym_restart_v8_flow_score_tails.py}
\end{center}
verifies the rational beta-integral coefficients, all displayed
explicit norm and residual constants, the flow and cutoff score
product identities on exact rational fixtures, and the layer
inequalities on boundary fixtures.  It also verifies exact
predecessor recovery and the protected-source hashes.
Finite fixtures are regression checks, not substitutes for the
all-parameter proofs above.  No floating fit, finite-volume gap,
or reference covariance has been used to certify a native continuum
claim.

The new results are: the original-input radial score asymptotics;
their volume-independent full-native lower bound; a full-native
cutoff-bank upper norm bound; and an all-flow, weak-coupling response
repair for an explicitly delimited one-cell probe.
The plaquette flow solution, basic integration by parts, beta
integrals, and the V6 Bochner mechanism are inherited or standard.

The V7 body is preserved byte for byte apart from the displayed
version title.  Removing this delimited append and reversing that
title recovers V7 exactly.  Only the immediate active source is
replaced; archives and companion sources are unchanged.  Build
and diagnostic artifacts remain outside \path{latex_notes}.
The next companion review and broad literature sweep both remain V10.

\begin{center}
\fbox{\begin{minipage}{0.94\textwidth}
\textbf{V8 status.}
We have an explicit rare-input obstruction and a rigorously bounded
local cutoff repair under the native Wilson law.  Its response
guarantee is only one-cell, and the obstruction concerns an isolated
plaquette subflow.  Native scale-matched spatial response, the actual
simultaneous-flow derivative budget, vacuum preparation, the common
upper spectral inequality, and the interacting continuum
Yang--Mills mass gap remain unproved.
\end{minipage}}
\end{center}
% END CONSOLIDATED V8 APPEND

% BEGIN CONSOLIDATED V9 APPEND -- 2026-09-19
\clearpage
\section{V9: Covariant heat-force scores and their native response obstruction}
\label{r9:context}

V8 separated a bounded observable from its potentially large
original-input score.  We now construct a genuinely spatially filtered
field whose score norm is controlled for \emph{every} smoothing time.
The construction smooths the plaquette force with the covariant heat
operator of the original connection; it does not flow the links.

There is an important second conclusion.  The archives already prove
native spectral suppression for that frozen connection.  Applying that
result to the new field, and estimating the actual conditional source
gradient, shows that its absolute response is smaller than every power
of the spacing on a stated class of physical-scale trajectories.
Thus this alternative pays a norm bound but does \emph{not} pay V7's
proposed Gaussian response comparison on those trajectories.
We record both conclusions, rather than use the flat linearization to
override the native suppression.

\subsection{Ownership audit and the precise change of construction}

Archive \(C\), V67 already defines the adjoint covariant spatial heat
operator, proves its path representation and scalar domination, and
uses it to smooth curvature before forming a quadratic observable.
Its V68 exhibits deterministic rough-connection screening.  Its V73
proves a native small-eigenvalue and local heat-weight bound, including
correlated insertions.  These facts and methods are inherited.
V74's estimates after preflow have an additional density cost and
are not substituted for the frozen-connection law here.

The new object below is a \emph{middle-link vector field} and its
marginal integration-by-parts score, not V67's scalar curvature-square
observable.  The new argument factors the literal plaquette force as
a covariant divergence and differentiates a self-adjoint block operator.
This gives a uniform native score-bank bound without controlling
the derivatives of a nonlinear link flow.  The new response obstruction
is an application of V73 to that force, with the conditional-response
gradient estimated explicitly.

Covariant Laplacian smearing is established background, not a new
algorithm claimed here.  A targeted primary-source check was
\href{https://arxiv.org/html/0905.2160v1}{M.~Peardon et al.,
\emph{A novel quark-field creation operator construction for hadronic
physics in lattice QCD}, arXiv:0905.2160}, Section II.
Its quark-field construction is not a Wilson score theorem.
The proofs below use the adjoint representation and the original
Wilson measure, without importing numerical spectroscopy conclusions.
The scheduled broad literature and companion reviews remain V10.

\subsection{The literal covariant divergence of the plaquette force}
\label{r9:factorization}

Take a cubic spatial torus of side \(L\ge3\), with \(V=L^3\) vertices,
as the middle slice of a finite four-dimensional Wilson history.
All middle links are free; exterior data and all incident temporal
plaquettes remain in the action.  Use dimensionless spatial lattice
units until stated otherwise.  Store links \(U_j(x)\) in the three
positive directions.  Identify a tangent vector on link \((x,i)\)
with \(\xi_i(x)U_i(x)\), \(\xi_i(x)\in\operatorname{Im}\mathbb H\).
The unit imaginary quaternions \(T_a\), \(a=1,2,3\), form an
orthonormal frame; \([T_a,v]=2T_a\times v\).

Let
\(\mathcal H=\bigoplus_{x,i}\R^3\) and
\(\mathcal E=\bigoplus_{x,j,i}\R^3\), with counting inner products.
The index \(i\) is a spectator in the covariant difference
\[
 (D_Uv)_{x,j,i}=R_j(x)v_{x+\hat j,i}-v_{x,i},
       \qquad R_j(x)=\Ad_{U_j(x)}.
 \tag{V9.1}\label{r9:D}
\]
Its adjoint and Laplacian are
\[
 \begin{split}
 (D_U^*b)_{x,i}
   &=\sum_j\{R_j(x-\hat j)^{\mathsf T}b_{x-\hat j,j,i}
                                                   -b_{x,j,i}\},\\
 \Delta_U&=D_U^*D_U,\qquad 0\preceq\Delta_U\preceq12I .
 \end{split}
 \tag{V9.2}\label{r9:lap}
\]
The last inequality follows by applying
\(|u-v|^2\le2|u|^2+2|v|^2\) to every difference.
A gauge transformation acts at the source vertex of each component
in both spaces.  Thus \(D_U\) and \(\Delta_U\) are gauge covariant.

For \(i\ne j\), write
\[
 P_{ij}(x)=U_i(x)U_j(x+\hat i)U_i(x+\hat j)^{-1}U_j(x)^{-1}.
\]
Then \(P_{ji}(x)=P_{ij}(x)^{-1}\).
Choose any bank of distinct unoriented spatial plaquettes and real
coefficients \(c_p\), extended by zero to all other plaquettes.
Set \(c_{x,ij}=c_{x,ji}\) for their common face coefficient and put
\[
 F_c=\sum_p c_pw_p,\qquad
 (B_c)_{x,j,i}=
 \begin{cases}
 c_{x,ij}\operatorname{Im}P_{ij}(x),&i\ne j,\\
 0,&i=j.
 \end{cases}
 \tag{V9.3}\label{r9:B}
\]
Here \(B_c\) is a curvature-valued array, not the source-response
scalar used elsewhere.

\begin{lemma}[Exact force factorization and derivative budget]
\label{r9:force}
In the stated right-translated tangent coordinates,
\[
 \begin{split}
 \nabla_IF_c&=D_U^*B_c,\\
 \|B_c\|^2&=2\sum_p c_p^2(1-w_p^2)\le2|c|^2,\\
 \sum_\alpha\|\partial_\alpha B_c\|^2&\le24|c|^2,\qquad
 \|\nabla_IF_c\|^2\le16|c|^2 .
 \end{split}
 \tag{V9.4}\label{r9:force-bounds}
\]
The sum over \(\alpha\) runs over every middle link and its three
orthonormal Lie directions.  Derivatives of the displayed arrays
are ordinary derivatives in these fixed algebra coordinates.
\end{lemma}
\begin{proof}
For the occurrence of \(U_i(x)\) at the start of \(P_{ij}(x)\),
differentiation along \(T_aU_i(x)\) gives
\(\partial_aw_p=-T_a\cdot\operatorname{Im}P_{ij}(x)\).
The other occurrence of an \(i\)-link, on the opposite side of
that face, contributes the corresponding imaginary part transported
backwards along \(U_j(x)\).  Consequently the total \(i\)-component
at \(x\) is
\[
 \sum_{j\ne i}\{
   R_j(x-\hat j)^{\mathsf T}
          c_{x-\hat j,ij}\operatorname{Im}P_{ij}(x-\hat j)
       -c_{x,ij}\operatorname{Im}P_{ij}(x)\},
\]
which is \(D_U^*B_c\).  Reversing the face orientation negates its
imaginary part, so the same calculation includes the \(j\)-component.

Each face occurs twice in \(\|B_c\|^2\), and a unit quaternion
has squared imaginary norm \(1-w_p^2\).
Holding three links fixed, the map from the fourth link to its
plaquette product is a Haar isometry, including a possible inversion.
Projection onto imaginary coordinates has differential norm at most
one.  Its three basis derivatives have total squared norm at most
three.  There are four free links and two copies of each face,
giving \(2\cdot4\cdot3=24\).
Different faces occupy different array entries in \(B_c\), so this
derivative estimate needs no overlap factor.
The last bound is V6's local force estimate, equivalently
Cauchy--Schwarz over the at most four incident spatial faces at each
link.
\end{proof}

\subsection{A block functional-calculus estimate with no growing-time cost}
\label{r9:block}

Define the real self-adjoint matrix on \(\mathcal H\oplus\mathcal E\)
\[
       A_U=\begin{pmatrix}0&D_U^*\\D_U&0\end{pmatrix}.
 \tag{V9.5}\label{r9:block-A}
\]
This is an auxiliary spatial matrix; it is neither a fermion operator
nor the physical Yang--Mills Hamiltonian.

\begin{lemma}[Packed derivative bounds]
\label{r9:packed}
For fixed arrays \(v\in\mathcal H\), \(b\in\mathcal E\),
\[
 \begin{split}
 \sum_\alpha\|(\partial_\alpha D_U)v\|^2&=24\|v\|^2,\\
 \sum_\alpha\|(\partial_\alpha D_U^*)b\|^2&=8\|b\|^2,\\
 \sum_\alpha\|(\partial_\alpha A_U)(v,b)\|^2
                                      &\le24(\|v\|^2+\|b\|^2).
 \end{split}
 \tag{V9.6}\label{r9:packed-bounds}
\]
The same estimates hold after complexification.
\end{lemma}
\begin{proof}
Vary only \(U_j(x)\) by left multiplication with \(e^{uT_a}\).
Then \(\partial R_j=\operatorname{ad}_{T_a}R_j\), while
\(\partial R_j^{\mathsf T}=-R_j^{\mathsf T}\operatorname{ad}_{T_a}\).
Use
\[
                    \sum_{a=1}^3|[T_a,z]|^2=8|z|^2 .
\]
In \(D_Uv\) this differentiates only the row \((x,j,i)\),
using \(v_{x+\hat j,i}\).  Every vertex is the head of three
stored edges, yielding \(3\cdot8=24\).
In \(D_U^*b\) this differentiates only the contribution of that
row to its head vertex, giving \(8\|b\|^2\).
The two block outputs of \(\partial A_U\) are orthogonal,
which proves the final line.  The identities extend to complex
vectors by applying them to real and imaginary parts.
\end{proof}

For \(t\ge0\), set \(g_t(z)=z e^{-tz^2}\).
Define \(b_0=\sqrt{12}\) and, for \(t>0\),
\[
 b_t=\min\{\sqrt{12},(2et)^{-1/2}\}.
 \tag{V9.7}\label{r9:bt}
\]

\begin{lemma}[Uniform differentiation of the filtered divergence]
\label{r9:functional}
For every fixed vector \(z\in\mathcal H\oplus\mathcal E\),
\[
 \sum_\alpha\|(\partial_\alpha g_t(A_U))z\|^2
                                      \le24\|z\|^2
 \quad(t\ge0).
 \tag{V9.8}\label{r9:frechet}
\]
Moreover the upper-right block of \(g_t(A_U)\) is
\[
 Q_t(U)=e^{-t\Delta_U}D_U^*,\qquad \|Q_t(U)\|\le b_t.
 \tag{V9.9}\label{r9:Qt}
\]
Neither statement requires a positive lower eigenvalue of \(\Delta_U\).
\end{lemma}
\begin{proof}
The block identity follows from
\(A_U^2=\operatorname{diag}(D_U^*D_U,D_UD_U^*)\), using its
convergent power series.  The spectral bound is
\(\sup_{\lambda\in[0,12]}\sqrt\lambda e^{-t\lambda}
 \le\min\{\sqrt{12},(2et)^{-1/2}\}\).

For \(t>0\), the scalar Gaussian Fourier identity gives
\[
 g_t(A)=\int_{\R} k_t(\xi)e^{i\xi A}\,d\xi,\qquad
 k_t(\xi)=-\frac{i\xi}{4\sqrt\pi\,t^{3/2}}
                          e^{-\xi^2/(4t)},\qquad
 \int_{\R}|\xi|\,|k_t(\xi)|\,d\xi=1 .
 \tag{V9.10}\label{r9:fourier}
\]
The last equality is the elementary second Gaussian moment.
In finite dimension we can differentiate under this absolutely
integrable formula.  Duhamel's identity reads
\[
 \partial_\alpha e^{i\xi A}
   =i\xi\int_0^1 e^{i(1-u)\xi A}(\partial_\alpha A)
                                           e^{iu\xi A}\,du.
\]
Both exponentials are unitary.  Apply
\eqref{r9:packed-bounds} to \(e^{iu\xi A}z\) and use Minkowski's
inequality in the direct sum over \(\alpha\).  The resulting
square-root bound is
\(\sqrt{24}\|z\|\int|\xi||k_t(\xi)|\,d\xi=\sqrt{24}\|z\|\).
This proves \eqref{r9:frechet}.
For \(t=0\), \(g_0(A)=A\), and the assertion is exactly
\eqref{r9:packed-bounds}.
\end{proof}

Factoring the force matters.  Directly differentiating
\(e^{-t\Delta_U}\) by an unweighted operator-norm Duhamel estimate
would introduce a factor \(t\).
The combination \(e^{-t\Delta_U}D_U^*\) instead has the uniform
derivative estimate above.  It is a Fourier functional-calculus
estimate, not an assumption that the nonlinear connection is close
to identity.

\subsection{The native heat-force field and a non-gradient score identity}
\label{r9:native}

Define
\[
        X_{t,c}(U)=e^{-t\Delta_U}\nabla_IF_c(U)
                         =Q_t(U)B_c(U).
 \tag{V9.11}\label{r9:field}
\]
At every finite \(t\), it is a smooth, middle-measurable,
gauge-equivariant vector field.  It is linear in \(c\).
It need not be the gradient of a scalar function.
Its heat kernel has the inherited V67 path representation and
scalar spatial-heat domination, so this is genuine spatial
filtering, not merely an average of the coefficients of the
local bank.  No link, temporal support, transfer operator, or
probability measure is changed.

For a general smooth middle field \(X\), write
\[
 T_X=-\operatorname{div}_I X+\gamma\langle\nabla_IS,X\rangle,
 \qquad s_X=\E[T_X\mid U_I]=-\operatorname{div}_{p_I}X.
 \tag{V9.12}\label{r9:scores}
\]
The projection retains the actual middle marginal \(p_I\,dH\).
For an actual conditional outer expectation \(m\), set
\[
 \mathcal B_X=\E[Xm],\qquad J_X=\E s_X^2 .
\]
Integration by parts gives
\(\mathcal B_X=\E[(m-\E m)s_X]\), and hence
\(\Var(m)\ge\mathcal B_X^2/J_X\) when \(J_X>0\).

\begin{lemma}[Native score budget for a field that is not a gradient]
\label{r9:nongradient}
For every such \(X\),
\[
 J_X\le\E T_X^2
       \le\E\|\nabla_I X\|_{\rm HS}^2
                       +(2+24\gamma)\E\|X\|^2 .
 \tag{V9.13}\label{r9:nongradient-bound}
\]
Here \(\nabla_I X\) is its Levi--Civita covariant derivative,
not just the derivative of its algebra coordinates.
\end{lemma}
\begin{proof}
First freeze the exterior and put \(V_{\rm pot}=\gamma S\).
On the compact product manifold, weighted integration by parts
and commutation of two covariant derivatives give
\[
 \int(\operatorname{div}_{e^{-V_{\rm pot}}}X)^2\,d\mu
  =\int\left[
       \sum_{i,j}(\nabla_iX^j)(\nabla_jX^i)
       +(\operatorname{Ric}+\nabla^2V_{\rm pot})(X,X)
       \right]d\mu .
 \tag{V9.14}\label{r9:div-identity}
\]
One can obtain the identity directly by expanding
\(-\int X^j\nabla_j(\nabla_iX^i-(\nabla_iV_{\rm pot})X^i)\,d\mu\).
Commute \(\nabla_j\nabla_iX^i\), integrate the remaining
\(\nabla_i\) once, and cancel the two terms containing
\((\nabla_iV_{\rm pot})X^j\nabla_jX^i\).
The curvature commutator gives the Ricci term.
The derivative contraction is at most
\(\sum_{i,j}|\nabla_iX^j|^2\) by Cauchy--Schwarz.
The product Ricci tensor is \(2g\), and V6's bound for the
\emph{full} action is \(\nabla_I^2S\preceq24g\).
Average over the exterior and apply the \(L^2\) contraction of
conditional expectation to \eqref{r9:scores}.
\end{proof}

There is a coordinate correction in applying this lemma.
For the right-invariant frame \(E_a(U)=T_aU\),
\(\nabla_{E_a}E_b=-[T_a,T_b]/2\).
Thus an algebra-coordinate field \(X\) satisfies
\[
 \|\nabla_I X\|_{\rm HS}
     \le\left(\sum_\alpha\|\partial_\alpha X\|^2\right)^{1/2}
                                              +\sqrt2\,\|X\|.
 \tag{V9.15}\label{r9:connection}
\]
Indeed the connection correction acts only on the varied link,
and its total squared norm is
\(\sum_{e,a}|[T_a,X_e]/2|^2=2\|X\|^2\).
This correction is not omitted from the following constants.

\begin{theorem}[All-time, full-native heat-force score bound]
\label{r9:uniform-score}
Let \(J_t\) be the actual marginal-score Gram of the distinct
plaquette bank \(X_{t,p}=e^{-t\Delta_U}\nabla_Iw_p\).
Put \(d_t=\min\{4,\sqrt2\,b_t\}\).  For every \(t\ge0\),
\[
 \begin{split}
 0\preceq J_t&\preceq K_\gamma(t)I,\\
 K_\gamma(t)&=
       (4\sqrt3+2\sqrt6\,b_t+\sqrt2\,d_t)^2
                                    +(2+24\gamma)d_t^2,\\
 J_t&\preceq(906+384\gamma)I .
 \end{split}
 \tag{V9.16}\label{r9:uniform-bank}
\]
All constants are independent of spatial volume, temporal depth,
bank size and smoothing time.  Exterior data may be arbitrary
provided the middle links remain free.  For \(t>0\), also
\[
 J_t\preceq
 \left[
   \left(4\sqrt3+\frac{2\sqrt3+\sqrt2}{\sqrt{et}}\right)^2
                    +\frac{2+24\gamma}{et}
 \right]I.
 \tag{V9.17}\label{r9:large-time}
\]
\end{theorem}
\begin{proof}
The heat operator is a contraction, so
\(\|X_{t,c}\|\le4|c|\) by \eqref{r9:force-bounds}.
The factorized bound gives
\(\|X_{t,c}\|\le\sqrt2\,b_t|c|\).  Together these give \(d_t|c|\).
Differentiating \eqref{r9:field}, applying \eqref{r9:frechet}
to \((0,B_c)\), and using the triangle inequality in the derivative
direct sum yield
\[
 \left(\sum_\alpha\|\partial_\alpha X_{t,c}\|^2\right)^{1/2}
       \le\sqrt{24}\|B_c\|
                +b_t\left(\sum_\alpha\|\partial_\alpha B_c\|^2\right)^{1/2}
       \le(4\sqrt3+2\sqrt6\,b_t)|c|.
\]
Use \eqref{r9:connection} and \eqref{r9:nongradient-bound}
to obtain \(K_\gamma(t)|c|^2\), proving the matrix assertion.
For the all-time constant use \(b_t\le2\sqrt3\), \(d_t\le4\):
\[
 K_\gamma(t)\le
      (4\sqrt3+16\sqrt2)^2+16(2+24\gamma)
      =592+128\sqrt6+384\gamma<906+384\gamma.
\]
For \eqref{r9:large-time} use
\(b_t\le(2et)^{-1/2}\), \(d_t\le(et)^{-1/2}\).
\end{proof}

The inequality pays a native nonlinear score budget for a genuinely
spatial field.  It does not pay the \(n^{-7}\) reference score
asymptotic or a physical source response.

There is also an exact way to charge the connection derivative to
native plaquette defects.  For a nonempty bank let
\[
           \varepsilon_{\mathcal P}
              =\max_{p\in\mathcal P}\E(1-w_p^2)\le1 .
\]
The same proof, now using
\(\E\|B_c\|^2\le2\varepsilon_{\mathcal P}|c|^2\) and
Minkowski in \(L^2(d\mu)\), gives for \(t>0\)
\[
 J_t\preceq
 \left[
  \left(4\sqrt{3\varepsilon_{\mathcal P}}
       +\frac{2\sqrt3+\sqrt{2\varepsilon_{\mathcal P}}}{\sqrt{et}}\right)^2
       +\frac{(2+24\gamma)\varepsilon_{\mathcal P}}{et}
 \right]I .
 \tag{V9.18}\label{r9:defect-budget}
\]
This statement defines its defect expectation under the actual law;
it does not substitute an unproved rate for it.

\subsection{Exact flat calibration, without a native-law comparison}
\label{r9:flat}

Take the uniform combination
\(F_0=V^{-1/2}\sum_p(w_p-1)\), with all three spatial
orientations included.  At the identity configuration its force
and \(B_0\) vanish.  Differentiating \eqref{r9:field} there gives
\[
 DX_{t,0}(\mathbf1)[v]
       =e^{-t\Delta_{\mathbf1}}\,D(\nabla_IF_0)(\mathbf1)[v].
 \tag{V9.19}\label{r9:flat-jet}
\]
The derivative of the heat operator multiplies the zero force;
it cannot contribute an extra linear term.
The flat plaquette expansion is
\[
 F_0(e^{\eta v})=-\frac{\eta^2}{2\sqrt V}
                       \|d_{\rm lat}v\|^2+O(\eta^3).
\]
On the transverse, nonconstant sector,
\(d_{\rm lat}^*d_{\rm lat}=\mathsf L\) and
\(\Delta_{\mathbf1}=\mathsf L\), with symbol
\(\lambda(k)=4\sum_j\sin^2(k_j/2)\).
Consequently the linear field is exactly
\[
 X^{\rm lin}_{t,0}(v)
     =-\frac1{\sqrt V}\mathsf L e^{-t\mathsf L}v,\qquad
                       q_t(\lambda)=\lambda e^{-t\lambda}.
 \tag{V9.20}\label{r9:flat-filter}
\]
This is V7's reference field at its flow time \(s=t/2\).
In particular, in the explicitly stationary Gaussian reference,
\(t=n^2\) has the already proved limiting capture greater than
\(8/9\) as \(n\to\infty\), \(L/n\to\infty\).
We inherit that reference calculation and its rational certificate;
we do not rebrand it as a new probability estimate.
The exact flat jet is not uniform control of the native law over
a smoothing radius growing like \(n\).

\subsection{The archived frozen-connection bound cannot be ignored}
\label{r9:screening}

We state the quantitative input from archive \(C\), V73 before
applying it.  It concerns the same dimensionless spatial operator
\(\Delta_U\) on one adjoint component.  Put
\[
 \begin{gathered}
 \kappa_\gamma=\frac{3\pi e}{4}(1+6\gamma)^{3/2},\qquad
 \ell_r=2\lceil\sqrt r\rceil+1,\qquad r\ge2,\\
 n_r=(2\ell_r-1)^3,\qquad d_r=3(2\ell_r-2),\\
 A_{r,\gamma}=64\kappa_\gamma^r
                    [\,2\sqrt{n_r}(1+2\sqrt{d_r})\,]^{2r-2},\\
 \Theta_{r,\gamma}(a,\sigma)
       =\min\{3,\,3\Gamma(r)A_{r,\gamma}(a^2/\sigma)^{r-1}\}.
 \end{gathered}
 \tag{V9.21}\label{r9:archive-constants}
\]
For translation-invariant native spatial laws on tori with all
side lengths at least \(\ell_r\), its theorem gives
\[
 \E\operatorname{tr}_{\R^3}
                  e^{-(\sigma/a^2)\Delta_U}(x,x)
                   \le\Theta_{r,\gamma}(a,\sigma).
 \tag{V9.22}\label{r9:archive-input}
\]
The theorem also applies to the corresponding stationary time limits.
This is an inherited native inequality, not a flat spectral
assumption.  Its proof conditions on plaquette-separated non-tree
links, bounds their densities by \(\kappa_\gamma\), pays a sphere
net for their common almost-invariant adjoint direction, and uses
Neumann bracketing of ordered eigenvalues.  It does not use the
false assertion that matrix exponentiation is operator monotone.

For clarity, write
\[
 \begin{split}
 I_a(R,\sigma)&=\frac Ra\operatorname{arsinh}\frac{Ra}{2\sigma}
     -\frac{2\sigma}{a^2}
            \left(\sqrt{1+(Ra/(2\sigma))^2}-1\right),\\
 h_a(R,\sigma)&=\min\{1,6e^{-I_a(R,\sigma)}\},\\
 N_R&=(2\lfloor R/a\rfloor+1)^3,\\
 \mathcal C_r(R,\sigma)&=
           N_R\Theta_{r,\gamma}(a,2\sigma)+h_a(R,\sigma)^2.
 \end{split}
 \tag{V9.23}\label{r9:local-cost}
\]
The parameter \(\sigma>0\) is a physical \emph{spatial smoothing}
scale squared.  It is not physical Euclidean transfer time.
The kernel of \(e^{-(\sigma/a^2)\Delta_U}\) is dominated by the
scalar walk of rate \(a^{-2}\), and
\[
 \sum_y\|K_\sigma^U(x,y)\|_{\rm HS}^2
                             =\operatorname{tr}K_{2\sigma}^U(x,x).
\]
Splitting its sum at distance \(R\), Cauchy--Schwarz on the near
part and the scalar walk tail on the far part show that any
possibly \(U\)-dependent insertion with \(|v(U;y)|\le M\) satisfies
\[
       \E|(K_\sigma^Uv)(x)|^2\le2M^2\mathcal C_r(R,\sigma).
 \tag{V9.24}\label{r9:insertion}
\]
This is V73's correlated-insertion estimate, recalled with its
actual constants.  Dependence of \(v\) on the same connection does
not invalidate it: the split is pointwise before expectation.

\begin{proposition}[Application to the normalized aggregate force]
\label{r9:force-suppression}
Let \(X_{\sigma/a^2,0}\) be \eqref{r9:field} for \(F_0\) above,
under the native translation-invariant laws in
\eqref{r9:archive-input}.  Then
\[
             \E\|X_{\sigma/a^2,0}\|^2
                          \le96\mathcal C_r(R,\sigma).
 \tag{V9.25}\label{r9:force-small}
\]
Suppose spatial side lengths in lattice units tend to infinity and
\[
            \log(1+\gamma_a)=o(\log(1/a)),\qquad a\downarrow0.
 \tag{V9.26}\label{r9:clock}
\]
For every fixed \(\sigma>0\) and every \(N>0\),
\[
                    \E\|X_{\sigma/a^2,0}\|^2=o(a^N).
 \tag{V9.27}\label{r9:force-superpoly}
\]
The norm is the unweighted sum over all middle links.
The conclusion is uniform in ambient volume subject to the
stated side-length condition.
\end{proposition}
\begin{proof}
Each component \((\nabla_IF_0)_i(x)\) is a sum of at most four
unit-length imaginary plaquette parts, with coefficients
\(V^{-1/2}\); its norm is at most \(4/\sqrt V\).
Apply \eqref{r9:insertion} separately for \(i=1,2,3\), with
\(M=4/\sqrt V\), then sum over \(V\) sites.
The factor is \(3V\cdot2\cdot16/V=96\).

For each desired power choose a sufficiently large fixed \(r\).
Its constant \(A_{r,\gamma_a}\) is \(a^{-o(1)}\) by
\eqref{r9:clock}.  Take
\(R_a=\sqrt{4\sigma K\log(1/a)}\) with a sufficiently large
fixed \(K\).  For small \(a\), \(aR_a/(2\sigma)\le1\).
The inequality \(I_a(R_a,\sigma)\ge R_a^2/(8\sigma)\)
follows by differentiating
\(z\operatorname{arsinh}z-\sqrt{1+z^2}+1\) twice on \([0,1]\).
It gives \(h_a^2\le36a^K\), while
\(N_{R_a}=O(a^{-3}(\log(1/a))^{3/2})\).
The first term of \(\mathcal C_r\) is
\(O(a^{2r-5-o(1)}(\log(1/a))^{3/2})\).
Choose \(2r-5>N\) and \(K>N\), with strict margins.
The fixed block fits for all sufficiently small \(a\).
This proves \eqref{r9:force-superpoly}.
No \(r=r(a)\) or volume-dependent block constant is used.
\end{proof}

\subsection{A native conditional-response upper bound}
\label{r9:response}

The preceding small field norm is not itself a score-capture
theorem.  We can, however, estimate its response to the actual
outer source rather than leave that consequence implicit.

\begin{lemma}[Conditional gradient of a separated bounded source]
\label{r9:conditional-gradient}
Let \(A\) be a real bounded observable depending only on links
outside the middle slice, with \(\|A\|_\infty\le M\).
Under the original finite Wilson history let \(m=\E[A\mid U_I]\).
Then, pointwise in \(U_I\),
\[
 \nabla_em=-\gamma\operatorname{Cov}
                    (A,\nabla_eS\mid U_I),\qquad
 \|\nabla_Im\|^2\le108\gamma^2VM^2 .
 \tag{V9.28}\label{r9:conditional-bound}
\]
The upper bound has no temporal-depth factor.
\end{lemma}
\begin{proof}
Differentiate the compact conditional integral defining \(m\),
including its normalizing denominator.  Since \(A\) has no direct
middle-link dependence, the derivative is the displayed covariance.
For a free four-dimensional link,
\(|\nabla_eS|\le6\), as at most six unit plaquette gradients contribute.
Conditional Cauchy--Schwarz gives
\[
 |\operatorname{Cov}(A,\nabla_eS\mid U_I)|^2
 \le\Var(A\mid U_I)\,
            \E[|\nabla_eS|^2\mid U_I]\le36M^2 .
\]
There are \(3V\) middle links.  Summing proves the result.
Spatial terms of \(S\) fixed by the conditioning cancel from the
covariance automatically; retaining their bound only enlarges
the constant.  No crossing factor or conditional denominator
has been dropped.
\end{proof}

At fixed spatial regulator, the bound also passes to stationary
time-cylinder limits for bounded finite-cylinder sources.  Indeed
the middle marginal density and the conditional numerator have
uniform derivative bounds on the compact middle manifold: at each
fixed derivative order, only the finitely many action terms incident
to its differentiated middle links enter.  Their density ratios
are bounded above and below at fixed \(V,\gamma\), uniformly in
temporal depth.  Compactness therefore gives convergent subsequences
of the densities and numerators together with first derivatives;
the cylinder law identifies their limits.  Passing the conditional
identity and its bound through those subsequences proves the claim.
No uniformity in a growing spatial regulator is inferred from this
fixed-regulator passage.

\begin{theorem}[The specified frozen-connection response misses the flat target]
\label{r9:response-obstruction}
Under the translation-invariant native laws of
Proposition~\ref{r9:force-suppression}, for the preceding outer
source and \(t=\sigma/a^2\),
\[
 |\mathcal B_{X_{t,0}}|
   \le72\sqrt2\,\gamma\sqrt V\,M
                           \sqrt{\mathcal C_r(R,\sigma)} .
 \tag{V9.29}\label{r9:response-upper}
\]
Assume \eqref{r9:clock}, spatial side lengths tending to infinity,
and polynomial bounds
\(V(a)\le C a^{-q}\), \(M(a)\le C a^{-d}\) for fixed finite
\(q,d\).  Then for every \(N>0\),
\[
                         |\mathcal B_{X_{\sigma/a^2,0}}|=o(a^N).
 \tag{V9.30}\label{r9:response-superpoly}
\]
The same assertion survives any additional polynomial scalar
normalization of the field or source.
In particular it cannot satisfy a relative lower comparison
\(|\mathcal B-b_a|\le\epsilon b_a\), \(\epsilon<1\), when
\(b_a\ge c a^k>0\) has a fixed polynomial lower scale.
\end{theorem}
\begin{proof}
Cauchy--Schwarz in the unchanged native measure gives
\[
 |\E[X_{t,0}m]|^2
   \le\E\|X_{t,0}\|^2\,\E\|\nabla_Im\|^2 .
\]
Insert \eqref{r9:force-small} and
\eqref{r9:conditional-bound}; \(\sqrt{96\cdot108}=72\sqrt2\).
For any requested power, \eqref{r9:force-superpoly} can be used
with an arbitrarily larger power before taking the square root.
The polynomial volume and source bounds, and the subpolynomial
coupling, are therefore absorbed.  The same reasoning pays any
fixed extra polynomial normalization.
Finally the relative comparison would imply
\(|\mathcal B|\ge(1-\epsilon)b_a\), contradicting the
superpolynomial conclusion.
\end{proof}

For example, the outer thermal magnetic energy
\(\gamma V^{-1/2}\sum_{p\ {\rm on\ an\ outer\ slice}}(w_p-1)\)
has \(M\le6\gamma\sqrt V\).
On a fixed physical-size torus \(V\) is polynomial in \(a^{-1}\);
any subpolynomial factor can be absorbed in a slightly larger
polynomial bound for \(M\).
Thus the theorem applies to the natural bounded finite-lattice
source used to seek V7's response calibration.
If \(n\) is proportional to \(a^{-1}\) and \(t=n^2\), its flat
reference response has a polynomial scale after a declared
polynomial normalization, whereas the present frozen-connection
native response does not.
For example, when \(L,n\) are both proportional to \(a^{-1}\),
one nonzero mode with \(\lambda\asymp a^2\) already contributes
a positive multiple of \(a^6\) to V7's reference response trace:
\(V^{-1}\asymp a^3\), \(\lambda^2/R(\lambda)\asymp a^3\),
and both attenuation factors are bounded below on that mode.

Two limitations must remain explicit.  First,
\eqref{r9:response-superpoly} does \emph{not} prove that
\(\mathcal B^2/(J\Var(m))\) vanishes: both denominators may also
be small, and no matching lower bound for them has been proved.
Second, a superpolynomial scalar rescaling is outside the claimed
normalization comparison.  Such rescaling cancels from a capture
fraction, but does not by itself prove that fraction is positive
uniformly or restore a particular Gaussian calibration.
The result concerns the specified frozen-connection field, not
V7's simultaneous spatial flow.

\subsection{What the next iteration must import or prove}
\label{r9:next}

The all-time score bound \eqref{r9:uniform-bank} is a native result
not available for V7's deterministic pullback field.  Its proof
isolates a useful mechanism: factor a local force as a covariant
divergence before differentiating its spatial filter.
Nevertheless the same field, on the trajectories of
Theorem~\ref{r9:response-obstruction}, cannot pay V7's prescribed
polynomial-scale response lower bound.
A good flat multiplier plus a bounded score is not a complete
source certificate.

The next broad literature sweep at V10 must therefore target the
joint issue: control of spatial transport under the native law
\emph{and} retention of a normalized physical source response.
A coarse or regularized connection would change the field and
would require its input derivatives to be paid.  The existing
preflow density estimates do not provide this for free.
Shared-path singlet or multiscale constructions must keep their
native correlations and their actual conditional normalization.
None is installed here merely because it has the right flat symbol.

We must not repeat the already solved isolated-plaquette geometry,
V67's curvature-square heat construction, or V73's spectral argument
as new advances.  Nor should we iterate different pointwise upper
bounds while leaving the normalized response untreated.
The unfinished requirements remain a noncollapsing physical source
family, vacuum preparation, a common growing-depth upper spectral
inequality on a dense physical source space, and the interacting
four-dimensional continuum theory with positive Hamiltonian mass gap.

\subsection{Verification and preservation}
\label{r9:verification}

The independent checker is
\begin{center}\small
\path{scripts/verify_ym_restart_v9_heat_force.py}.
\end{center}
It checks three exact noncommuting rational-quaternion force
factorizations on a \(3^3\) torus, the adjoint and derivative
identities, a transverse \(4^3\)-torus mode, the explicit constants,
and protected-source hashes.  Optional floating small-matrix
Fr\'echet tests are diagnostics, not rigorous enclosures.
The uniform proof is \eqref{r9:fourier}, not those fixtures.

Removing this append and reversing the displayed title recovers
V8 byte for byte.  Only that immediate active source is replaced;
archives and companions are unchanged.  All build artifacts stay
outside \path{latex_notes}.  Companion review and the broad
primary-source sweep are both due at V10.

\begin{center}
\fbox{\begin{minipage}{0.94\textwidth}
\textbf{V9 status.}
A gauge-covariant spatial heat-force bank now has a proved all-time,
volume-independent native score bound and the desired flat filter.
Its frozen-connection response nevertheless fails the specified
polynomial-scale Gaussian calibration on the stated trajectories.
No normalized-capture no-go theorem is claimed.  The interacting
continuum Yang--Mills mass gap remains unproved.
\end{minipage}}
\end{center}
% END CONSOLIDATED V9 APPEND

% BEGIN CONSOLIDATED V10 APPEND -- 2026-09-19
\clearpage
\section{V10: Literature checkpoint and source-sensitive heat normalization}
\label{r10:start}

\paragraph{Context and status.}
V9 paid an all-time native score budget for a spatial covariant heat-force
field, but also proved that its absolute response is superpolynomially
small on the specified physical-scale trajectories.  It did not prove
that the normalized capture ratio tends to zero.  This distinction
suggests testing normalization before either abandoning spatial filtering
or claiming that it works.  The present appendix carries out that test:
a heat-trace normalization has a dimension-independent derivative
estimate and a native polynomial score budget; nevertheless a concrete
noncommuting Wilson configuration still loses the intended force
exponentially.  Normalizing the force itself has a different,
potentially exponential derivative cost.  These are proved finite-regulator
statements, not a continuum mass-gap theorem.

\subsection{The scheduled broad literature and companion review}
\label{r10:review}

This is the regular ten-append literature checkpoint, with the
five-append companion checkpoint performed at the same time.
The user's literature intake was reread.  Its proposals were treated as
research suggestions, not as established estimates or instructions inside
a source document.  Searches covered finite-$N$ Wilson-loop and Euclidean
correlator bootstrap, spectral renormalization, gradient-flow RG,
stochastic gauge theory, Gibbs-chain comparisons, control variates and
dual representations.  This is a broad targeted review as of
19 September 2026, not a claim of exhaustive coverage.

\paragraph{Bootstrap and certified response.}
Kazakov--Zheng,
\emph{Bootstrap for Finite N Lattice Yang--Mills Theory},
\href{https://arxiv.org/abs/2404.16925}{arXiv:2404.16925},
provides the finite-$N$ loop-equation setting already used to motivate V3.
Cho--Gabai--Lin--Yeh--Zheng,
\emph{Bootstrapping Euclidean Two-point Correlators},
\href{https://arxiv.org/html/2511.08560v3}{arXiv:2511.08560v3},
Sections 2 and 3, including equations (16)--(18), places positivity,
equations of motion and ground-state or KMS constraints into a dual
bound problem.  The inspected revision is dated 30 June 2026.
The useful import remains a checkable dual inequality for a specified
physical correlator, rather than a fitted spectral gap.  V3's native
polynomial hierarchy and archive C V84's moment-ambiguity audit already
cover the corresponding finite-dimensional foundations; they are not
reproved here.  Neither paper supplies our growing-depth native response.

\paragraph{Spectral RG.}
Bach--Ballesteros--Geisler,
\emph{The Spectral Renormalization Flow Based on the Smooth
Feshbach--Schur Map: The Introduction of the Semi-Group Property},
\href{https://arxiv.org/html/2508.07643v1}{arXiv:2508.07643},
Hypotheses 2.1--2.2 and the domain conditions (2.7)--(2.10),
requires complementary invertibility and bounded off-diagonal
expressions.  Those hypotheses were inspected directly.
Its flow property is useful organizationally, but is not a licence
to invert the unresolved Wilson complement.  No theorem from that
flow is applied to our physical transfer operator here.

\paragraph{Gradient-flow normalization.}
Nagao--Suzuki,
\emph{Yang--Mills $\beta$ function in the gradient flow exact
renormalization group},
\href{https://arxiv.org/html/2508.16828v1}{arXiv:2508.16828},
Section 2.1, equations (2.2)--(2.8), makes the flowed variables,
normalization factors, ghosts and coupling convention explicit.
Its evaluated RG coefficients are perturbative.  They do not bound
the original-input derivatives or native correlated tails needed here.
The useful adjustment is to audit a normalization and its derivatives
together.  The finite-matrix normalization theorem below is proved
independently; it is not attributed to that paper and is not identified
with its RG transformation.

\paragraph{New exact duality lead.}
Lemoine,
\emph{Universal dualities for Wilson loops in lattice Yang--Mills},
\href{https://arxiv.org/html/2604.16252v1}{arXiv:2604.16252},
Theorem 1.1 and its proof, separates a finite-lattice character
expansion into action-dependent weights and Haar-integral coefficients.
The proof uses absolute uniform character convergence.  This is a useful
possible representation for native insertion certificates, not a
uniform weak-coupling estimate on their sign or size.
The paper's more detailed $\mathrm U(N)$ topological and channel
formulae are not silently specialized to our $\SU(2)$ problem.
A positive character weight alone does not remove cancellations in an
inserted, connected observable.  No new duality theorem is needed
for the results below.

\paragraph{Sampling and dimension-specific results.}
Qin--Wang,
\href{https://arxiv.org/abs/2208.11299}{\emph{Spectral Telescope}},
and Qin--Ju--Wang,
\href{https://arxiv.org/abs/2312.12782}{\emph{Spectral gap bounds for
reversible hybrid Gibbs chains}}, concern auxiliary sampling chains;
a comparison to physical Euclidean time remains a separate obligation.
Bhattacharya--Lawrence--Yoo,
\href{https://arxiv.org/abs/2307.14950}{\emph{Control variates for lattice
field theory}}, suggests Schwinger--Dyson variance reduction, but an
estimator's reduced variance is not a physical source lower bound.
The strong-coupling setting of Shen--Zhu--Zhu,
\href{https://arxiv.org/abs/2204.12737}{\emph{A stochastic analysis
approach to lattice Yang--Mills at strong coupling}}, and the
three-dimensional setting of Chandra--Chevyrev--Hairer--Shen,
\href{https://arxiv.org/abs/2201.03487}{\emph{Stochastic quantisation of
Yang--Mills--Higgs in 3D}}, do not supply the missing four-dimensional
weak-coupling estimate.  These items were screened for applicability;
their hypotheses have not been verified for an import into this proof.

\paragraph{Companions and nonduplication.}
The current NS V26, SM--Einstein V72 and YM Shell-Mixing V16 files
were checked against their V5 checkpoint hashes and are unchanged.
The terminal arguments were reread.  NS V26 analyzes rank-one
Oseen-kernel contractions, not a Wilson conditional law.
SM V72's many-copy covariance transport holds with the physical mode
set fixed; its final statement expressly excludes growing cutoffs.
Shell V16 evaluates a complete measured one-loop coefficient on a flat
one-holonomy family, with a summable coefficient refinement error.
It does not control the joint nonperturbative coupling/volume remainder.
None is used as a native response estimate.

Searches of the mass-gap archives distinguish the present construction
from compact-group heat normalization, conditional heat-bath derivatives,
physical heat-trace disintegration, and C V73's volume-normalized
covariant heat weight.  We reuse V6's force Hessian estimate, V9's
packed link derivatives and non-gradient score identity, and C V68's
constant noncommuting-background test idea.  The new objects are the
configuration-wise normalized \emph{spatial} heat operator, its packed
derivative bound, and its explicit source-selection failure.
The next companion checkpoint is V15; the next regular broad sweep
is V20.  These counters continue across an archive restart.

\subsection{A normalized exponential with a paid packed derivative}
\label{r10:matrix}

Let $A=A(u)$ be a smooth Hermitian $d\times d$ matrix, with $d<\infty$,
and let $\partial_\alpha$ be a finite family of real parameter
derivatives at the point under consideration.  Write
\[
 E_\alpha=\partial_\alpha A,\qquad
 \sum_\alpha E_\alpha^2\preceq K I,\qquad
 Z_t=\Tr e^{-2tA},\qquad
 R_t=\frac{e^{-tA}}{\sqrt{Z_t}},\qquad
 \rho_t=R_t^2,\qquad t\ge0.
 \tag{V10.1}\label{r10:R}
\]
The matrix trace is not an expectation over gauge configurations.
$Z_t$ is strictly positive at finite $t$ even if $A$ has zero
eigenvalues.  No lower spectral gap is assumed.

\begin{theorem}[Dimension-independent normalized-heat derivative]
\label{r10:matrix-thm}
Under \eqref{r10:R},
\[
 \|R_t\|_{\rm HS}=1,\qquad \|R_t\|_{\rm op}\le1,\qquad
 \boxed{\ \sum_\alpha\|\partial_\alpha R_t\|_{\rm HS}^2
                    \le Kt^2.\ }
 \tag{V10.2}\label{r10:packed}
\]
The assertions are invariant under replacing $A(u)$ by
$A(u)+b(u)I$, even for a variable scalar $b$.  In particular
a small $Z_t$ does not itself cause an inverse-$Z_t$ loss in
\eqref{r10:packed}.
\end{theorem}
\begin{proof}
Fix a point and diagonalize $A$ there.  Differentiable eigenvectors
are not required.  Set
$p_i=e^{-2t\lambda_i}/Z_t$, so $\sum_i p_i=1$.
The Fr\'echet derivative of the unnormalized exponential, divided
by the value of $\sqrt{Z_t}$ at this point, has entries
\[
 (B_\alpha)_{ij}
 =-t(E_\alpha)_{ij}
           \int_0^1 p_i^{(1-s)/2}p_j^{s/2}\,ds.
 \tag{V10.3}\label{r10:divdiff}
\]
This follows directly by integrating Duhamel's formula; it includes
coincident eigenvalues without a limiting denominator.
Cauchy--Schwarz in $s$, followed by weighted arithmetic--geometric
mean, gives
\[
 \left|\int_0^1p_i^{(1-s)/2}p_j^{s/2}\,ds\right|^2
 \le \int_0^1p_i^{1-s}p_j^s\,ds
 \le\frac{p_i+p_j}{2}.
\]
Consequently Hermiticity implies
\[
 \sum_\alpha\|B_\alpha\|_{\rm HS}^2
 \le t^2\sum_{\alpha,i,j}\frac{p_i+p_j}{2}|(E_\alpha)_{ij}|^2
 =t^2\Tr\!\left(\rho_t\sum_\alpha E_\alpha^2\right)
 \le Kt^2 .
 \tag{V10.4}\label{r10:B}
\]
In the real Hilbert space of Hermitian matrices, normalization
differentiates as
\[
 \partial_\alpha R_t
 =B_\alpha-R_t\langle R_t,B_\alpha\rangle_{\rm HS}.
 \tag{V10.5}\label{r10:tangent}
\]
This is an orthogonal projection because $\|R_t\|_{\rm HS}=1$.
It can only decrease each norm in \eqref{r10:B}.
The operator norm bound follows from the Hilbert--Schmidt norm.
Finally $e^{-t(A+bI)}=e^{-tb}e^{-tA}$ and its denominator acquires
the identical positive factor, proving the invariance.
All formulas are smooth through eigenvalue crossings.
\end{proof}

The theorem is about a normalized operator, not about
$e^{-tA}v/\|e^{-tA}v\|$.  The latter has no analogous conclusion
from \eqref{r10:R} alone, as proved below.

\subsection{Application to the original Wilson middle slice}
\label{r10:native}

Use V9's spatial torus of side $L\ge3$, $V=L^3$, its right-translated
tangent space $\mathcal H=\bigoplus_{x,i}\R^3$ of dimension $9V$,
and its covariant difference $D_U$ and Laplacian
$\Delta_U=D_U^*D_U$.  Derivatives $\partial_\alpha$ run over all
middle spatial-link unit Lie directions.  V9 proved, pointwise,
\[
 \|D_U\|^2\le12,\qquad
 \sum_\alpha\|(\partial_\alpha D_U)v\|^2=24\|v\|^2,\qquad
 \sum_\alpha\|(\partial_\alpha D_U^*)b\|^2=8\|b\|^2.
 \tag{V10.6}\label{r10:inherited}
\]
These constants include the spectator direction index.
It follows, for $E_\alpha=\partial_\alpha\Delta_U$, that
\[
 \begin{split}
 \sum_\alpha\|E_\alpha v\|^2
 &\le2\sum_\alpha\|(\partial_\alpha D_U^*)D_Uv\|^2
       +2\|D_U^*\|^2\sum_\alpha\|(\partial_\alpha D_U)v\|^2\\
 &\le16\|D_Uv\|^2+576\|v\|^2
 \le768\|v\|^2 .
 \end{split}
 \tag{V10.7}\label{r10:768}
\]
Since each $E_\alpha$ is self-adjoint, this is precisely
$\sum_\alpha E_\alpha^2\preceq768I$.

Define, without changing the Wilson probability law,
\[
 Z_t(U)=\Tr_{\mathcal H}e^{-2t\Delta_U},\qquad
 R_t(U)=Z_t(U)^{-1/2}e^{-t\Delta_U},\qquad
 Y_{t,c}=R_t(U)\nabla_I F_c,\qquad
 F_c=\sum_{p\ {\rm spatial}}c_p w_p.
 \tag{V10.8}\label{r10:field}
\]
The coefficients $c_p$ are real constants, independent of $U$.
All expressions are smooth at finite $t,L$.
Local gauge transformations conjugate $\Delta_U$ unitarily on
$\mathcal H$, leave its trace invariant, and transform
$\nabla_I F_c$ covariantly.  Thus $Y_{t,c}$ is an admissible
gauge-equivariant middle field.  It need not be a gradient field,
and the normalization is spatially nonlocal.

\begin{theorem}[A native score budget after heat-trace normalization]
\label{r10:score-thm}
In the original finite four-dimensional Wilson law at $\gamma\ge0$,
with all incident temporal plaquettes retained, let
\[
 s_{t,c}=\E\!\left[
 -\operatorname{div}_IY_{t,c}
       +\gamma\langle\nabla_I S,Y_{t,c}\rangle
                         \,\middle|\,U_I\right].
\]
Its actual marginal-score Gram, as a quadratic form in $c$, obeys
\[
 \boxed{\quad
 J_t^{\rm tr}\preceq
 \left[(64\sqrt3\,t+8\sqrt3+8\sqrt2)^2
                         +32+384\gamma\right]I .
 \quad}
 \tag{V10.9}\label{r10:native-budget}
\]
The constant is independent of spatial volume, temporal depth and
bank size.  The time dependence is polynomial, not bounded uniformly
as $t\to\infty$.
\end{theorem}
\begin{proof}
Put $v_c=\nabla_I F_c$.  V6's original-metric estimates give
\[
 |v_c|\le4|c|,\qquad
 \|\nabla_I v_c\|_{\rm HS}\le8\sqrt3\,|c|.
\]
For right-translated link components the packed Levi--Civita
correction has norm $\sqrt2\,|v_c|$, as in V9.  Therefore
\[
 \left(\sum_\alpha|\partial_\alpha v_c|^2\right)^{1/2}
                 \le(8\sqrt3+4\sqrt2)|c|.
\]
Theorem \ref{r10:matrix-thm} and \eqref{r10:768} imply
\[
 |Y_{t,c}|\le4|c|,\qquad
 \left(\sum_\alpha|\partial_\alpha Y_{t,c}|^2\right)^{1/2}
 \le(64\sqrt3\,t+8\sqrt3+4\sqrt2)|c|.
\]
Here $\|\partial_\alpha R_t\|_{\rm op}
\le\|\partial_\alpha R_t\|_{\rm HS}$ pays the derivative
of the \emph{configuration-dependent denominator}, not just that
of the numerator.  Adding the output connection correction gives
\[
 \|\nabla_IY_{t,c}\|_{\rm HS}
       \le(64\sqrt3\,t+8\sqrt3+8\sqrt2)|c|.
\]
V9's weighted divergence identity for arbitrary middle fields yields
\[
 \E s_{t,c}^2
 \le \E\|\nabla_IY_{t,c}\|_{\rm HS}^2
                       +(2+24\gamma)\E|Y_{t,c}|^2.
\]
This proves \eqref{r10:native-budget}.  Linearity in $c$ gives
the full bank inequality, including off-diagonal entries.
No lower curvature, mixing, spectral gap or Gaussian-law
comparison was used.
\end{proof}

For the raw field $X_{t,c}=e^{-t\Delta_U}v_c$ and
$r=Z_t^{-1/2}$ the exact marginal identity is
\[
 s_{Y}=r\,s_X-Xr
      =r\{s_X-t\Tr(\rho_t X\Delta_U)\},\qquad
 \rho_t=Z_t^{-1}e^{-2t\Delta_U}.
 \tag{V10.10}\label{r10:scorechange}
\]
Indeed $\partial\log Z_t=-2t\Tr(\rho_t\partial\Delta_U)$
by cyclicity of the Duhamel derivative.
This exhibits the extra term missed by treating $r(U)$ as a
deterministic rescaling.  For a separated outer source $A$ with
$m=\E[A\mid U_I]$, the response remains
\[
 \mathcal B_Y=\E\langle Y,\nabla_I m\rangle
             =\E[(m-\E m)s_Y].
 \tag{V10.11}\label{r10:response}
\]
It is not $\E r$ times the old response, and a full-history
contact score cannot replace $s_Y$.

\subsection{Reference calibration survives, but source selection can fail}
\label{r10:selection}

At the trivial connection,
\[
 Z_t(I)=9\sum_{k\in(2\pi/L)\{0,\ldots,L-1\}^3}
                  e^{-2t\lambda(k)},\qquad
 \lambda(k)=4\sum_{j=1}^3\sin^2(k_j/2).
 \tag{V10.12}\label{r10:flatZ}
\]
For the aggregate generator
$F_0=V^{-1/2}\sum_p(w_p-1)$, $v_0(I)=0$.
Hence the first variation of $R_tv_0$ has no term from
$\partial R_t$ at $I$.  On the transverse reference modes of V7,
\[
 Y_{t,0}^{\rm lin}
     =-\frac{\mathcal L e^{-t\mathcal L}x}
                         {\sqrt{VZ_t(I)}}.
 \tag{V10.13}\label{r10:flat}
\]
Thus its reference score is the V9 reference score multiplied by
a nonzero deterministic scalar.  The reference capture
$\mathcal B^2/(JG)$ is unchanged.  In particular the V7 limit at
$t=n^2$, $n\to\infty$, $L/n\to\infty$, retains its proved lower
bound $8/9$.  This is a statement about that reference Gaussian
law and its linear field only.

For a general connection, if $P_{\min}$ is the spectral projector
of $\Delta_U$ at its least eigenvalue and has rank $m_{\min}$,
finite-dimensional spectral decomposition gives
\[
 \lim_{t\to\infty}R_t(U)
            =m_{\min}^{-1/2}P_{\min}.
 \tag{V10.14}\label{r10:edge}
\]
Consequently the limiting field can vanish when the force has no
component in the least-eigenvalue space, despite
$\|R_t\|_{\rm HS}=1$ at every time.

\begin{proposition}[A literal noncommuting Wilson-force countercheck]
\label{r10:background}
Let $L$ be a positive multiple of four and take constant spatial links
\[
 U_1=(1+\mathbf i)/\sqrt2,\qquad
 U_2=(1+\mathbf j)/\sqrt2,\qquad U_3=1,
 \tag{V10.15}\label{r10:configuration}
\]
where $\mathbf i,\mathbf j,\mathbf k$ are quaternion units.
For the aggregate $F_0$ above the nonzero field $v_0=\nabla_I F_0$
satisfies
\[
 |v_0|^2=2,\qquad \Delta_Uv_0=2v_0,\qquad
 \lambda_{\min}(\Delta_U)\le3-\sqrt5<2.
 \tag{V10.16}\label{r10:misalignment}
\]
In particular
\[
 \boxed{\quad |Y_{t,0}(U)|
          \le\sqrt2\,e^{-(\sqrt5-1)t}.\quad}
 \tag{V10.17}\label{r10:loss}
\]
An operator heat-trace normalization does not guarantee retention
even of this actual plaquette-force seed.
\end{proposition}
\begin{proof}
For constant $U_1=\cos\theta+\mathbf i\sin\theta$ and
$U_2=\cos\phi+\mathbf j\sin\phi$, the $12$-plaquette trace is
$1-2\sin^2\theta\sin^2\phi$; the other two spatial traces are one.
The trace of a general commutator is $1-2|u_1\times u_2|^2$
for quaternion vector parts $u_1,u_2$.
At orthogonal axes its first variation under
$U_1\mapsto e^{\epsilon a}U_1$ is
$-4\sin\theta\cos\theta\sin^2\phi\,a_1$.
At $\theta=\phi=\pi/4$ this is $-a_1$.
The analogous $U_2$ variation is $-a_2$; all other
components vanish.  Translation invariance then gives the
\emph{individual-link} aggregate gradients
\[
 (v_0)_{x,1}=-\mathbf i/\sqrt V,\qquad
 (v_0)_{x,2}=-\mathbf j/\sqrt V,\qquad (v_0)_{x,3}=0.
\]
This also follows by differentiating every oriented plaquette word,
which is independently checked in the exact diagnostic.

Set $R_j=\Ad U_j$.  In the colour basis
$(\mathbf i,\mathbf j,\mathbf k)$,
\[
 R_1=\begin{pmatrix}1&0&0\\0&0&-1\\0&1&0\end{pmatrix},
 \qquad
 R_2=\begin{pmatrix}0&0&1\\0&1&0\\-1&0&0\end{pmatrix}.
\]
The Fourier symbol, separately on each spectator direction, is
$\sum_j(2I-e^{ik_j}R_j-e^{-ik_j}R_j^{\mathsf T})$.
At $k=0$ it is $\operatorname{diag}(2,2,4)$, proving the
force eigenvalue and its norm.
At the allowed momentum $k=(\pi/2,0,0)$ it is
\[
 \begin{pmatrix}4&0&0\\0&2&2i\\0&-2i&4\end{pmatrix},
\]
whose eigenvalues are $4,3-\sqrt5,3+\sqrt5$.
Complex Fourier notation does not change the real spectrum:
conjugate momenta give the corresponding real invariant subspace.
Thus $Z_t\ge e^{-2t(3-\sqrt5)}$, whereas
$e^{-t\Delta_U}v_0=e^{-2t}v_0$.  Equation \eqref{r10:loss}
follows.
\end{proof}

This is a deterministic configuration test at every finite $L$,
not a typicality statement in the Wilson law.  A point has Haar
measure zero.  Positivity of finite-regulator Wilson density in
its neighbourhood is not a uniform lower probability as the
regulator changes.  We therefore infer neither a native expected
response no-go nor a capture no-go from \eqref{r10:loss}.
The conclusion is the narrower, proved rejection of an
unconditional pointwise source-retention argument.

The normalization is also global: at $t=0$ it is
$I/\sqrt{9V}$, and its volume dependence is not a construction
of a local thermodynamic observable.  Uniformity of the score
\emph{bound} in \eqref{r10:native-budget} does not repair this
separate locality issue.

\subsection{Normalizing the seed instead: an exact derivative tradeoff}
\label{r10:seed}

One could try to bypass \eqref{r10:loss} by defining
$e^{-tA}v/\|e^{-tA}v\|$.  The following calculation shows why
the operator theorem cannot be reused for this operation.

\begin{proposition}[Exponential direction sensitivity]
\label{r10:seed-thm}
Let $d>0$, let $O(u)$ be the real two-dimensional rotation through
angle $u$, and put
\[
 A(u)=O(u)\begin{pmatrix}0&0\\0&d\end{pmatrix}O(u)^{\mathsf T},
 \qquad v=e_2,\qquad w_t(u)=e^{-tA(u)}v .
\]
Then $\|A'(0)\|=d$ but
\[
 \left\|\left.\frac{d}{du}\frac{w_t(u)}{|w_t(u)|}
                         \right|_{u=0}\right\|=e^{dt}-1.
 \tag{V10.18}\label{r10:exponential}
\]
For any fixed-in-$u$ regularization floor $\eta\ge0$, set
$W_{t,\eta}=w_t/\sqrt{\eta^2+|w_t|^2}$.
If $|W_{t,\eta}(0)|\ge c>0$, then
\[
 |W_{t,\eta}'(0)|\ge c(e^{dt}-1).
 \tag{V10.19}\label{r10:floor}
\]
The same lower bound holds for a nonnegative differentiable
floor $\eta(u)$ whenever the indicated derivative exists.
\end{proposition}
\begin{proof}
Write $h=e^{-dt}$.  Direct conjugation gives
$w_t(0)=he_2$ and $w_t'(0)=(1-h)e_1$.
The derivative is perpendicular to $w_t(0)$, so normalization
does not subtract it and \eqref{r10:exponential} follows.
For a constant floor the derivative has length
$(1-h)/\sqrt{\eta^2+h^2}$.
The retained-length hypothesis says
$h/\sqrt{\eta^2+h^2}\ge c$, giving \eqref{r10:floor}.
For a differentiable variable floor the additional derivative
is parallel to $e_2$, while the perpendicular component remains
the same, so its norm cannot be smaller.
\end{proof}

This counterexample is an elementary matrix family, not an assertion
that its rotating two-dimensional block has already been embedded as
a native Wilson source.  It disproves the general implication from
a packed generator derivative bound to a similarly controlled
\emph{seed-normalized} derivative.  It identifies a quantitative
overlap or localization input that such a repair would have to supply.

\subsection{Decision, verification, and the next native calculation}
\label{r10:decision}

The review does not license a jump from perturbative RG or an
auxiliary sampler gap to the physical transfer gap.
The new mathematical progress is more specific: the putative
exponential \emph{operator-normalization derivative} cost can be
removed, with a fully native polynomial score budget; nevertheless
the relevant force can be spectrally misaligned, and normalizing
that force may reintroduce exponential sensitivity.
These are distinct obstructions, with explicit proofs and examples.

The next calculation should not be another absolute heat upper
bound.  It should test an insertion-aware, gauge-covariant coarse
transport against the \emph{actual} response
\eqref{r10:response}, paying both its input variations and its
source overlap.  For any heat-based candidate the decisive new
quantity is the following joint spectral insertion.
For $\omega\ge0$ define the measurable window
$P_\omega(U)={\bf1}_{[\lambda_{\min},\lambda_{\min}+\omega]}(\Delta_U)$
and set
\[
 \mathcal B_{\le\omega}:=\E\left[
 \frac{\langle\nabla_I m,\,
          e^{-t\Delta_U}P_\omega(U)\nabla_I F_0\rangle}
      {\sqrt{\Tr e^{-2t\Delta_U}}}
       \right],
 \tag{V10.20}\label{r10:joint}
\]
together with a controlled complementary contribution.
There is no claimed sign or lower bound for \eqref{r10:joint}.
A bound on the eigenvalue count alone, or on the uninserted
partition function, cannot establish it.
The bootstrap route retained from the literature should target
these correlated insertions or a specified two-time physical Gram,
not replace them by independent spectral and source averages.

\begin{proposition}[A paid complementary-window error]
\label{r10:tail}
In the stated finite Wilson law, for a separated outer source
$|A|\le M$ and $m=\E[A\mid U_I]$, the aggregate field obeys
\[
 |\mathcal B_{Y_{t,0}}-\mathcal B_{\le\omega}|
       \le72\gamma\sqrt V\,M e^{-t\omega}.
 \tag{V10.21}\label{r10:tail-bound}
\]
No smoothness or derivative bound for $P_\omega$ is required.
\end{proposition}
\begin{proof}
The spectral theorem and $\sqrt{Z_t}\ge e^{-t\lambda_{\min}}$
give $\|R_t(I-P_\omega)\|\le e^{-t\omega}$ pointwise.
V6 gives $|\nabla_I F_0|\le4\sqrt3$, since its $3V$
coefficients are $1/\sqrt V$.
V9's native conditional-gradient estimate gives
$|\nabla_I m|\le6\sqrt3\,\gamma\sqrt V M$.
Insert these bounds in the exact complementary part of
\eqref{r10:response}.  The projector is used only in this
spectral decomposition, not as a differentiated generating field.
\end{proof}

For $t=\sigma/a^2$ and $\omega=q a^2\log(1/a)/\sigma$,
the error factor is exactly $a^q$ for $0<a<1$.
When $\gamma\sqrt V M$ grows at most polynomially in $1/a$,
its prefactor can therefore be paid by a fixed sufficiently large $q$.
The unpaid term is now explicitly the \emph{correlated window near
the random spectral minimum}, not the unnormalized heat trace.
This reduction does not establish that the window has the desired
physical source overlap or a continuum interpretation.

\paragraph{Executed checks and preservation.}
The diagnostic
\path{scripts/verify_ym_restart_v10_normalization.py}
checks the literal plaquette-word force at
\eqref{r10:configuration} in exact arithmetic, its covariant
eigenvalue, the two Fourier symbols and their characteristic
polynomials, the packed constant arithmetic, and the two-dimensional
normalization tradeoff.  Separate floating matrix fixtures test the
normalized Fr\'echet derivative, scalar-shift invariance and the
analytic bound; they are diagnostics, not interval enclosures or
proof premises.  The inequalities for all sizes and times are
proved above.

V9 is exactly recoverable by removing this append, restoring the
previous displayed title and adding the document terminator.
Its SHA--256 is
\path{F687D63973DE8D876CB52DF03C185455B913502E1DB6D280166DCBE81F34C702}.
All 47 baseline source records are checked, including this recoverable
predecessor; companion and archive files are unmodified.
Compilation and rendering outputs remain outside \path{latex_notes}.
Only the immediate active V9 source is replaced.

\paragraph{Status.}
Native normalized-source retention, vacuum and continuum construction,
dense physical source coverage, and a positive interacting
four-dimensional Hamiltonian mass gap remain unproved.

% END CONSOLIDATED V10 APPEND

% BEGIN CONSOLIDATED V11 APPEND -- 2026-09-19
\clearpage
\section{V11: Conditional-source fields with positive native response}
\label{r11:start}

V10 identified a signed insertion near a random spectral minimum.
An upper bound on the heat trace did not control that insertion.
We now change the generating field rather than estimate an
unrelated force overlap: use the gradient of the \emph{actual}
conditional physical-source mean.  Its self-response is a positive
Dirichlet matrix.  No division by the field norm or by a source
overlap is made.

The new quantitative input is a first-layer, conditional-variance
bound on the first two derivatives of that mean.  It yields an
explicit native score budget, independent of temporal depth but
not of spatial volume.  Together these give positive-response
certificates for the finite history and its fixed-volume vacuum
limit.  They do not give a positive response bounded below
uniformly along a continuum trajectory.  In particular, computing
a conditional mean is still a nontrivial native integral, not an
already evaluated solution of the source problem.

\subsection{Geometry, inherited inputs and nonduplication}
\label{r11:geometry}

Take a spatial cubic torus of side $L\ge4$, let $V=L^3$, and let
$N=3V$ be its number of positive spatial links.  In this appendix
$N$ counts \emph{links}, not temporal depth or colours.
The middle configuration is $x\in\mathcal Q=\SU(2)^N$ with the
unit-$S^3$ product metric.  All links remain integrated.
Consider the ordinary Wilson history at times $-n,\ldots,n$,
$n\ge1$, with every spatial plaquette on the endpoint slices
included at full weight.  Temporal links are also integrated.
Write its action exactly as
\[
 S=S_s(x)+S_+(x,z_+)+S_+(x,z_-),\qquad
 Z_+(x)=\int e^{-\gamma S_+(x,z)}\,dz .
 \tag{V11.1}\label{r11:half}
\]
The right half action includes temporal plaquettes in its $n$
slabs and spatial plaquettes at times $1,\ldots,n$.
It excludes the middle spatial action $S_s$.
Here $dz$ is product normalized Haar measure over all half-history
links other than $x$; the notation on the left half includes reflection.
The exact middle and conditional laws are
\[
 p(dx)=\mathcal Z^{-1}e^{-\gamma S_s(x)}Z_+(x)^2\,dx,\qquad
 \nu_x(dz)=Z_+(x)^{-1}e^{-\gamma S_+(x,z)}\,dz .
 \tag{V11.2}\label{r11:laws}
\]
For a bounded real gauge-invariant outer-slice observable $A$ put
\[
 m_A(x)=\E_{\nu_x}A,\qquad
 v_A(x)=\Var_{\nu_x}(A).
 \tag{V11.3}\label{r11:message}
\]
The same derivative calculations apply to a bounded observable in
the positive half which does not depend explicitly on $x$.

V5's sewing and marginal-score identities, V6's full-action Hessian
bound, and V9's non-gradient divergence identity are inherited.
Archive f14 V52, Proposition \texttt{prop:v52-fine-transfer},
already constructs the positive injective physical Wilson transfer
and its finite-spatial-volume ground law; its proof was reread.
Its injectivity is reused, not claimed as a new result.
The constrained-interface conditional-message bounds in f12 V28
were also inspected: their small-field interface restriction is
not imposed here.  The first-layer derivative calculation below
works on the full compact middle manifold at every coupling.

For historical context, the primary bibliographic record of
M.~L\"uscher,
\href{https://link.springer.com/article/10.1007/BF01614090}
{\emph{Construction of a selfadjoint, strictly positive transfer
matrix for Euclidean lattice gauge theories}},
Commun.\ Math.\ Phys.\ \textbf{54} (1977), 283--292, was checked.
The argument uses the explicit archived Wilson construction,
not an uninspected generalization of that paper.
This is a targeted check; the next companion and broad literature
reviews remain V15 and V20, respectively.

\subsection{Only the first temporal layer is differentiated}
\label{r11:derivatives}

At fixed $z$ write
\[
 g(x,z)=\nabla_xS_+(x,z),\qquad
 H(x,z)=\nabla_x^2S_+(x,z).
\]
Each middle link belongs to exactly one temporal plaquette in this
half.  That plaquette contains no other middle link.  Holding its
other three links fixed makes its normalized trace a unit linear
quaternion coordinate of the middle link.  Consequently
\[
 |g|^2\le N,\qquad \|H\|_{\rm HS}^2\le3N .
 \tag{V11.4}\label{r11:geometry-bounds}
\]
Indeed each same-link Hessian of $1-w_p$ is $w_p$ times the metric,
and all mixed middle-link blocks of $H$ vanish.
Correlations between different middle links under $\nu_x$ do
\emph{not} vanish; they enter the next formula.

\begin{theorem}[Native variance-weighted first and second derivatives]
\label{r11:derivative-thm}
For every stated half history and every $\gamma\ge0$,
\[
 \begin{split}
 \nabla m_A&=-\gamma\,\Cov_{\nu_x}(A,g),\\
 \nabla^2m_A&=-\gamma\,\Cov_{\nu_x}(A,H)
    +\gamma^2\E_{\nu_x}\big[(A-m_A)
                    (g-\bar g)\otimes(g-\bar g)\big],
 \qquad \bar g=\E_{\nu_x}g .
 \end{split}
 \tag{V11.5}\label{r11:identities}
\]
Define
\[
 a_{\gamma,N}=\gamma\sqrt N,\qquad
 b_{\gamma,N}=\gamma\sqrt{3N}+2\gamma^2N.
 \tag{V11.6}\label{r11:ab}
\]
Pointwise in $x$,
\[
 \boxed{\quad
 |\nabla m_A|\le a_{\gamma,N}\sqrt{v_A},\qquad
 \|\nabla^2m_A\|_{\rm HS}\le b_{\gamma,N}\sqrt{v_A}.
 \quad}
 \tag{V11.7}\label{r11:derivative-bounds}
\]
The constants do not depend on $n$ or on the location of the
observable within the half, provided it has no explicit middle
dependence.
\end{theorem}
\begin{proof}
Differentiate the numerator and denominator in
\eqref{r11:message}.  Compactness and boundedness of $A$
justify all derivatives, giving the first identity.
At a given $x$ use normal orthonormal product coordinates.
A second quotient differentiation gives the displayed centered
third moment and the Hessian term.  Since this is a tensor identity
at an arbitrary point, it is the covariant Hessian identity.
In particular, the derivative of the conditional normalizer has
not been dropped.

Put $\widetilde g=g-\bar g$ and $\widetilde H=H-\E H$.
Cauchy--Schwarz for a scalar and a Hilbert-space-valued variable,
together with \eqref{r11:geometry-bounds}, gives
\[
 |\Cov(A,g)|\le\sqrt{Nv_A},\qquad
 \|\Cov(A,H)\|_{\rm HS}\le\sqrt{3Nv_A}.
\]
Also $|\widetilde g|\le2\sqrt N$ and
$\E|\widetilde g|^2\le N$, so
\[
 \E|\widetilde g|^4\le4N^2,\qquad
 \left\|\E[(A-m_A)\widetilde g\otimes\widetilde g]\right\|_{\rm HS}
 \le\sqrt{v_A\,\E|\widetilde g|^4}\le2N\sqrt{v_A}.
\]
Substitution proves \eqref{r11:derivative-bounds}.
The incidence counts in \eqref{r11:geometry-bounds} do not change
when more temporal slabs are added.
\end{proof}

This improves V9's use of the full incident-action gradient for
the specified half-history geometry.  It does not change V9's
valid broader estimate for arbitrary separated outer sources.
No claim is made that the conditional covariance tensor in
\eqref{r11:identities} is spatially local.

\subsection{A positive response matrix and a variance-weighted score budget}
\label{r11:certificate}

Let $A_1,\ldots,A_k$ be a finite real source family on the outer slice,
$m=(m_{A_i})$, and set
\[
 G=\Cov_p(m),\qquad C=\Cov(A_1,\ldots,A_k),\qquad
 \Sigma=\E_p\Cov_{\nu_x}(A_1,\ldots,A_k)=C-G .
 \tag{V11.8}\label{r11:Grams}
\]
The covariance $C$ uses the actual outer marginal of the same
history.  Conditional independence of its two halves identifies
$G$ with the opposite-boundary physical source covariance.
In particular $0\preceq G\preceq C$.

Use the fields and their \emph{middle marginal} scores
\[
 X_i=\nabla m_{A_i},\qquad
 s_i=-\operatorname{div}_p X_i,\qquad
 \mathcal E_{ij}=\E_p\langle\nabla m_{A_i},\nabla m_{A_j}\rangle,
 \qquad J_{ij}=\E_p s_i s_j .
 \tag{V11.9}\label{r11:fields}
\]
They are smooth and gauge equivariant on the unreduced compact
link manifold.  Adding a constant to any source has no effect.

\begin{theorem}[Source-conditioned native certificate]
\label{r11:certificate-thm}
The response matrix is exactly the symmetric positive matrix
$\mathcal E$:
\[
 \E_p[(m_i-\E_pm_i)s_j]=\mathcal E_{ij},\qquad
 0\preceq\mathcal E\preceq a_{\gamma,N}^2\Sigma .
 \tag{V11.10}\label{r11:E}
\]
For
\[
 K_{\gamma,N}=b_{\gamma,N}^2
                     +(2+24\gamma)a_{\gamma,N}^2
\]
one has
\[
 J\preceq K_{\gamma,N}\Sigma,\qquad
 \boxed{\quad
 G\succeq W\mathcal E+\mathcal E W^{\mathsf T}
                   -K_{\gamma,N}W\Sigma W^{\mathsf T}
 \quad}
 \tag{V11.11}\label{r11:dual}
\]
for every real $k\times k$ matrix $W$.
For $\gamma>0$, the Moore--Penrose version is
\[
 G\succeq K_{\gamma,N}^{-1}
                      \mathcal E\Sigma^\dagger\mathcal E .
 \tag{V11.12}\label{r11:inverse}
\]
At $\gamma=0$ all messages are constant and the assertions without
$K^{-1}$ remain valid.
\end{theorem}
\begin{proof}
Integration by parts under $p$ gives
$\E_p(f s_j)=\E_p\langle\nabla f,\nabla m_j\rangle$.
This proves the exact response formula, its symmetry and positivity.
Apply \eqref{r11:derivative-bounds} to
$A_c=\sum_i c_i A_i$ and integrate to obtain the matrix bound
for $\mathcal E$.

The full-history middle score associated to $X_c$ has conditional
expectation $s_c$.  V9's divergence identity, the unit-$S^3$
Ricci tensor $2g$, and V6's full-action Hessian upper bound $24g$
give
\[
 \E_p s_c^2
 \le\E_p\|\nabla^2m_{A_c}\|_{\rm HS}^2
                   +(2+24\gamma)\E_p|\nabla m_{A_c}|^2
 \le K_{\gamma,N}\,c^{\mathsf T}\Sigma c.
\]
All temporal plaquettes on \emph{both} sides are retained in this
score estimate, unlike the single half used to differentiate $m$.
Expanding the positive Gram of $m-\E m-Ws$ proves the exact
V5 lower matrix with $J$; replacing $J$ by its upper bound proves
\eqref{r11:dual}.

If $c\in\ker\Sigma$, the pointwise derivative bound gives
$\nabla m_{A_c}=0$ almost everywhere, hence $\mathcal E c=0$.
Thus the inverse formula follows by taking
$W=K_{\gamma,N}^{-1}\mathcal E\Sigma^\dagger$ in
\eqref{r11:dual}; the requisite range condition has been proved.
At zero coupling, the positive half is independent of $x$.
\end{proof}

For one source, the measurable response is
$\mathcal E=\E_p|\nabla m_A|^2$, not $G=\Var_p(m_A)$.
The theorem does not evaluate $\mathcal E$.
It gives the bound $G\ge\mathcal E^2/(K\Sigma)$ when
$\Sigma>0$, not a lower bound from $K$ alone.
The constants grow with $N$, and the conditional mean can become
very small with physical separation.  These losses remain visible.

\subsection{Strict finite-regulator response and the fixed-volume vacuum}
\label{r11:vacuum}

For clarity, let $\mathsf M$ denote multiplication by
$e^{-\gamma S_s/2}$ and let $\mathsf T$ be f14 V52's literal
Wilson transfer on gauge-invariant $L^2(\mathcal Q,dU)$.
Let $b=e^{-\gamma S_s/2}$ as a function.
Integrating temporal links and multiplying the half-actions gives
\[
 m_A=\frac{\mathsf T^n(bA)}{\mathsf T^n b},\qquad
 p(dx)=\frac{(\mathsf T^n b(x))^2}
                   {\|\mathsf T^n b\|_2^2}\,dx .
 \tag{V11.13}\label{r11:transfer}
\]
The endpoint factor $b$ gives full, not half, spatial endpoint
weight.  It is not an unspecified vacuum boundary condition.

\begin{proposition}[No exact loss of a nonconstant source at finite depth]
\label{r11:strict}
If $\gamma>0$ and $A$ is not constant almost everywhere, then
$m_A$ is nonconstant and $\mathcal E>0$.
For a finite family independent modulo constants,
$\mathcal E$ is positive definite.
These conclusions also hold in the fixed-$L,\gamma$ stationary
ground law at every finite depth $n$.
They assert strict positivity, not a uniform quantitative floor.
\end{proposition}
\begin{proof}
The denominator in \eqref{r11:transfer} is everywhere positive.
If $m_A=c$, then $\mathsf T^n[b(A-c)]=0$.
The inherited physical-space injectivity of $\mathsf T$ and
the strictly positive $b$ imply $A=c$ almost everywhere.
The compact product manifold is connected, and $p$ has a
strictly positive smooth density.  Thus a smooth nonconstant
$m_A$ has $\E_p|\nabla m_A|^2>0$.  Apply this to every
nonzero source combination for the matrix assertion.

The archived injectivity applies specifically to the physical
space: the one-link Wilson character coefficients are all
positive, their product convolution is injective, and multiplication
by $\mathsf M$ is bounded and invertible at fixed regulator.
The gauge projection has a kernel on the unreduced Hilbert space;
no injectivity there is assumed.

For the vacuum statement let $h>0$ be the normalized Perron
function, $\mathsf T h=\lambda_0h$, and
$\mathsf P f=(\lambda_0 h)^{-1}\mathsf T(hf)$.
The stationary middle law is $h^2dx$ and $m_A=\mathsf P^n A$.
Injectivity of $\mathsf T$ implies injectivity of $\mathsf P$
at this fixed regulator, proving nonconstancy as before.
\end{proof}

The derivative and score bounds also pass to this ground law.
Here are details so that this passage does not implicitly
exchange spatial or continuum limits.  Put distant boundaries at
$\pm M$, with $M\ge n$, and keep the sources at $\pm n$.
Then
\[
 m_{A,M}(x)=
 \frac{\mathsf T^n[A\,\mathsf T^{M-n}b](x)}
      {\mathsf T^M b(x)}.
 \tag{V11.14}\label{r11:distant}
\]
At fixed $L,\gamma$, compactness and the simple isolated top
eigenvalue give
$\lambda_0^{-r}\mathsf T^r b\to\langle h,b\rangle h$ in $L^2$.
The smooth finite-regulator kernel maps $L^2$ continuously into
$C^j$ for every fixed finite $j$; applying one additional transfer
step upgrades this convergence to $C^j$.
The strictly positive limiting denominator is bounded below.
It follows that \eqref{r11:distant} converges in $C^2$ to
$\mathsf P^nA$, and the same argument with $A^2$ gives convergence
of the conditional variances.
The middle densities and their first derivatives converge as well.
The finite-cylinder estimates have constants independent of $M$,
so all inequalities above pass to the stationary law.
No estimate on the rate is uniform in $L$ or $\gamma$.

\subsection{Optional positive heat preconditioning and a matrix tail}
\label{r11:heat}

One may use V10's smooth normalized operator $R_t(U)$ without
returning to a signed cross-source force:
\[
 Y_i=R_t\nabla m_i,\qquad
 B^R_{ij}=\E_p\langle\nabla m_i,R_t\nabla m_j\rangle,\qquad
 J^R_{ij}=\E_p s_{Y_i}s_{Y_j}.
 \tag{V11.15}\label{r11:heatfields}
\]
Then $B^R$ is positive semidefinite and, under the hypotheses
of Proposition \ref{r11:strict}, positive definite at finite $t$.

\begin{theorem}[Positive spectral-window response, with paid norm and tail]
\label{r11:heat-thm}
Put
\[
 K^R_{\gamma,N,t}
 =\{b_{\gamma,N}+(16\sqrt3\,t+2\sqrt2)a_{\gamma,N}\}^2
                              +(2+24\gamma)a_{\gamma,N}^2 .
\]
Then $J^R\preceq K^R_{\gamma,N,t}\Sigma$.
The certificate \eqref{r11:dual} holds with
$\mathcal E,K$ replaced by $B^R,K^R$.
For V10's window
$P_\omega={\bf1}_{[\lambda_{\min},\lambda_{\min}+\omega]}(\Delta_U)$
and
\[
 B^R_{\le\omega,ij}
 =\E_p\langle\nabla m_i,R_tP_\omega\nabla m_j\rangle,
\]
one has the native matrix inequality
\[
 \boxed{\quad
 0\preceq B^R-B^R_{\le\omega}
     \preceq e^{-t\omega}\mathcal E
     \preceq\gamma^2N e^{-t\omega}\Sigma .
 \quad}
 \tag{V11.16}\label{r11:tail}
\]
No lower bound on $B^R_{\le\omega}$ is asserted.
\end{theorem}
\begin{proof}
The operator $R_t$ is positive definite and has norm at most one.
Its packed ordinary derivative norm is at most $16\sqrt3\,t$
by V10.  For $v=\nabla m_{A_c}$, convert between ordinary
right-translated and covariant derivatives twice, once for the
input and once for the output.  Each correction has packed norm
$\sqrt2$ times the corresponding vector norm.  Thus, pointwise,
\[
 \|\nabla(R_tv)\|_{\rm HS}
 \le\|\nabla v\|_{\rm HS}+
                  (16\sqrt3\,t+2\sqrt2)|v|
 \le\{b+(16\sqrt3\,t+2\sqrt2)a\}\sqrt{v_{A_c}}.
\]
Also $|R_tv|\le a\sqrt{v_{A_c}}$.
The native non-gradient score inequality now gives the stated
budget.  Integration by parts gives response $B^R$ and the
certificate.  Positivity and strict positivity follow from
those of $R_t$ and the gradients already considered.

Finally $R_t$ commutes with $P_\omega$ and
$0\preceq R_t(I-P_\omega)\preceq e^{-t\omega}I$
by V10's spectral estimate.
Sandwich this inequality between $\nabla m_{A_c}$ and integrate.
Equation \eqref{r11:E} proves the last bound.
No derivative of the sharp projector is taken.
\end{proof}

Unlike V10's cross-force insertion, both the retained and
complementary self-response windows are nonnegative matrices.
This removes cancellation from this particular response
decomposition, but not possible smallness of the whole message
or absence of its gradient from the lowest spatial eigenspaces.
$R_t$ also retains V10's global spatial normalization.
The unfiltered fields $X_i$ avoid that extra normalization entirely.

\subsection{An evaluated depth-matched reference check}
\label{r11:reference}

The new field can be tested without guessing a smoothing time.
Use exactly V7's stationary transverse Gaussian reference and its
magnetic-energy message
\[
 m_n(x)=-\frac1{2\sqrt V}
 \{x^{\mathsf T}a_n(\mathsf L)x
                      -\Tr(a_n(\mathsf L)\mathsf R^{-1})\},\qquad
 a_n(\lambda)=\lambda e^{-2n\theta(\lambda)} .
\]
Then $\nabla m_n=-a_n(\mathsf L)x/\sqrt V$.
It is V7's $q=a_n$ field, not the elementary $q=\lambda$ bank
and not the perfectly matched $q=a_n/(2R)$ field.

\begin{proposition}[Five-sixths reference capture at growing depth]
\label{r11:five-sixths}
For $n\to\infty$ and $L/n\to\infty$, let
$H=\Var(m_n)$, $B=\E|\nabla m_n|^2$ and
$J=\E(-\operatorname{div}_{\nu_L}\nabla m_n)^2$.
Then
\[
 H\sim\frac9{1024\pi^2}n^{-5},\qquad
 B\sim\frac{45}{1024\pi^2}n^{-6},\qquad
 J\sim\frac{135}{512\pi^2}n^{-7},
 \qquad
 \boxed{\ \frac{B^2}{JH}\longrightarrow\frac56.\ }
 \tag{V11.17}\label{r11:five-sixths-eq}
\]
The explicit coefficient $W=n/6$ in the scalar exact certificate
has $(2WB-W^2J)/H\to5/6$.
\end{proposition}
\begin{proof}
V7's trace identities give, without approximation,
\[
 H=\frac1{2V}\Tr(a_n^2\mathsf R^{-2}),\quad
 B=\frac1V\Tr(a_n^2\mathsf R^{-1}),\quad
 J=\frac2V\Tr(a_n^2).
\]
Near zero momentum, $\lambda(k)\sim |k|^2$,
$\theta(k)\sim|k|$, and $R(k)\sim2|k|$.
There are six real transverse degrees per nonzero momentum.
Rescale $k=y/n$.  The normalized trace has radial density
$3\pi^{-2}r^2\,dr$, giving
\[
 \begin{split}
 \lim n^5H&=\frac3{8\pi^2}\int_0^\infty r^4e^{-4r}\,dr,\\
 \lim n^6B&=\frac3{2\pi^2}\int_0^\infty r^5e^{-4r}\,dr,\\
 \lim n^7J&=\frac6{\pi^2}\int_0^\infty r^6e^{-4r}\,dr.
 \end{split}
\]
These are the same dominated Riemann-sum limits used in V7:
the rescaled mesh is $2\pi n/L\to0$; on the Brillouin zone
$\theta(k)\ge c|k|$; and after rescaling a fixed polynomial times
$e^{-c|y|}$ controls the tails.  On every bounded rescaled ball
the displayed small-momentum ratios converge uniformly, with
their continuous values at zero.  Thus the limits hold on the
stated joint sequence, not only as iterated limits.
Using $\int_0^\infty r^j e^{-4r}\,dr=j!/4^{j+1}$ evaluates
all constants.  Their ratio is $5/6$ and $B/J\sim n/6$,
proving the assertion about the explicit coefficient.
\end{proof}

This check shows that using the actual message as generator
does not have V7's elementary-bank $n^{-7}$ capture dilution.
It is still a reference calculation, not a comparison theorem
for the interacting compact law.  Its source is unbounded, so
we use exact Gaussian traces here, not the bounded-source
hypotheses of the native derivative theorem.

\subsection{What can be computed next without hiding a normalizer}
\label{r11:next}

The new native construction pays a sign and derivative problem,
but it does not turn the conditional mean into known input.
For example, let $z_1,\ldots,z_4$ be conditionally independent
with law $\nu_x$, and put
$D_{12}=(A(z_1)-A(z_2))(g(x,z_1)-g(x,z_2))$.
Exactly,
\[
 \nabla m_A=-\frac\gamma2\E_{\nu_x^{\otimes2}}D_{12},\qquad
 \mathcal E=\frac{\gamma^2}{4}
       \E_{p(dx)\nu_x^{\otimes4}}
                    \langle D_{12},D_{34}\rangle .
 \tag{V11.18}\label{r11:replica}
\]
The first equality is the elementary two-copy covariance identity;
conditional independence gives the second.
The pointwise integrand in the second formula need not be positive.
Its conditional expectation is a squared vector norm.
Most importantly, the joint density in that expectation is
\[
 \mathcal Z^{-1} e^{-\gamma S_s(x)}
                Z_+(x)^{-2}\prod_{\ell=1}^4
                           e^{-\gamma S_+(x,z_\ell)} .
 \tag{V11.19}\label{r11:replica-density}
\]
Deleting $Z_+^{-2}$ would compute a different model.
Thus a four-half-history polynomial calculation with that
denominator omitted is not a certificate for $\mathcal E$.

The next quantitative target is a lower enclosure of this native
Dirichlet matrix or a controlled approximation of the message
gradients, with the true half-history normalizers retained.
V3's polynomial certificates can be used only after this conditional
normalization problem is explicitly paid.  The fixed-volume
strict-positivity result cannot substitute for that enclosure.
Equally, a gap of the auxiliary marginal Langevin generator
$-\operatorname{div}_p\nabla$ is not the physical transfer gap:
these operators act by different dynamics, and no uniform
comparison has been established.

\paragraph{Verification and preservation.}
The checker \path{scripts/verify_ym_restart_v11_source_fields.py}
tests the first-layer incidence on finite four-dimensional histories,
the conditional quotient jets and centered third moments in exact
rational fixtures, the two- and four-copy covariance identities,
and the factorial constants giving $5/6$.
Optional floating lattice traces are labeled as diagnostics,
not proofs of joint-limit estimates.
V10's body is preserved exactly apart from the displayed title.
Its recoverable SHA--256 is
\path{3FECA6672B8DD95A178BFB14E2822D44D17FC709DDF66EBB022E73F68C35BF83}.
Only the active predecessor is replaced; all companion and archive
sources are protected.  Build and visual-QA outputs are outside
\path{latex_notes}.

\paragraph{Status.}
Positive native response and a temporal-depth-independent score
budget are established for the declared source-conditioned fields.
A uniform lower response, useful continuum-normalized capture,
the interacting continuum construction, and its physical mass
gap remain unproved.

% END CONSOLIDATED V11 APPEND

% BEGIN CONSOLIDATED V12 APPEND -- 2026-09-19
\clearpage
\section{V12: Paired native insertions and quantitative polynomial recovery}
\label{r12:start}

V11's four-conditional-replica formula contains $Z_+(x)^{-2}$.
That formula is correct, but it is not necessary to evaluate its
integrand directly.  The ordinary reflected history already contains
two independent halves \emph{conditional on the middle}.  We use
these two halves to represent linear tests of the message gradient,
then recover its squared norm by a variational problem on the
middle slice.  No conditional density is replaced.

There are two new computational consequences.  Both the field
response and the \emph{exact} marginal-score Gram become expectations
of polynomial insertions in one ordinary Wilson history.  Moreover,
the positive Dirichlet matrix of V11 admits polynomial lower
approximations with an explicit absolute error independent of
temporal depth.  This is not a solved four-dimensional moment
problem: its number of variables and its required relative precision
remain obstacles.

\subsection{Inputs, scope and the targeted literature check}
\label{r12:inputs}

Keep V11's unreduced geometry, $N=3L^3$, $L\ge4$, and
\[
 d\mu(x,z_+,z_-)=p(dx)\nu_x(dz_+)\nu_x(dz_-),\qquad
 g_\pm=\nabla_x S_+(x,z_\pm).
\]
Every link is integrated and every Wilson plaquette is retained.
A superscript $+$ or $-$ denotes the observable on the
corresponding reflected half, not a positive or negative part.
The sources $A_i$ are real, bounded, gauge invariant, independent
of $x$ as functions of a half history, and have conditional means
$m_i(x)$.  Polynomiality will only be required when we invoke V3.

V3 already supplies the complete fixed-regulator polynomial
expectation hierarchy.  V5 already supplies the native score
certificate and its contact warning.  V11 supplies the exact
message derivatives, finite-regulator strict positivity and
their fixed-spatial-volume vacuum passage.  These are inherited.
In particular, the two-history score \emph{difference} in f1 V26
compares different initial configurations of a transfer fibre;
it is not the reflected conditional-pair construction below.

The targeted primary records checked here were A.~Hyv\"arinen,
\href{https://www.jmlr.org/papers/v6/hyvarinen05a.html}
{\emph{Estimation of Non-Normalized Statistical Models by Score
Matching}}, JMLR \textbf{6} (2005), 695--709, and
T.~Bhattacharya, S.~Lawrence and J.-S.~Yoo,
\href{https://arxiv.org/abs/2307.14950}
{\emph{Control variates for lattice field theory}}.
The first motivates squared-field objectives without an explicitly
evaluated normalizer; the second constructs lattice
Schwinger--Dyson control variates and demonstrates them in
two-dimensional scalar theory.  These are methodological connections,
not theorems imported into four-dimensional Yang--Mills.
Only their stated scope and primary records were checked here;
all identities and approximation bounds used below are proved.
The next regular companion and broad literature reviews remain
V15 and V20.

\subsection{Two reflected halves represent the gradient weakly}
\label{r12:paired-gradient}

Define a random middle tangent vector, represented in the global
right-translated orthonormal frame, by
\[
 F_i(x,z_+,z_-)
   =-\frac{\gamma}{2}(A_i^+-A_i^-)(g_+-g_-).
 \tag{V12.1}\label{r12:F}
\]
This random vector is not itself the conditional message gradient.

\begin{lemma}[Native paired response]
\label{r12:paired}
Exactly,
\[
 \E[F_i\mid x]=\nabla m_i(x).
 \tag{V12.2}\label{r12:Fmean}
\]
Consequently, for any smooth middle vector field $Y$,
\[
 \E_\mu\langle F_i,Y\rangle
       =\E_p\langle\nabla m_i,Y\rangle .
 \tag{V12.3}\label{r12:weak}
\]
\end{lemma}
\begin{proof}
Conditional independence and equality of the two half laws give
\[
 \E[(A_i^+-A_i^-)(g_+-g_-)\mid x]
       =2\{\E_{\nu_x}(A_i g)-m_i\E_{\nu_x}g\}.
\]
Use V11's differentiated conditional-mean identity.
The second assertion follows by conditioning; $Y$ depends only
on $x$.
\end{proof}

There is no extra conditional partition factor in
\eqref{r12:weak}.  Indeed
$p(dx)\nu_x(dz_+)\nu_x(dz_-)$ is exactly the original
normalized Wilson history.  This does not cancel the
$Z_+^{-2}$ in V11's \emph{different, four-copy} formula.
It changes which quantity is evaluated before the variational
step.

For a declared finite bank $Y_1,\ldots,Y_r$ write
\[
 D_{i\alpha}=\E_\mu\langle F_i,Y_\alpha\rangle,\qquad
 Q_{\alpha\beta}=\E_p\langle Y_\alpha,Y_\beta\rangle .
 \tag{V12.4}\label{r12:DQ}
\]
The notation $D$ here is a response matrix, not a link derivative.

\begin{theorem}[Dirichlet recovery from ordinary-history moments]
\label{r12:recovery}
Let $\mathcal E_{ij}=\E_p\langle\nabla m_i,\nabla m_j\rangle$
as in V11.  For every real $k\times r$ matrix $W$,
\[
 \mathcal E\succeq DW^{\mathsf T}+WD^{\mathsf T}
                                      -WQW^{\mathsf T}.
 \tag{V12.5}\label{r12:Edual}
\]
Optimizing over $W$ gives
\[
 L_{\mathcal Y}=DQ^\dagger D^{\mathsf T}\preceq\mathcal E .
 \tag{V12.6}\label{r12:L}
\]
The difference is the Gram of the orthogonal-projection residuals
of the gradients in $L^2(p;T\mathcal Q)$.
\end{theorem}
\begin{proof}
Expand the Gram of
$\nabla m_i-\sum_\alpha W_{i\alpha}Y_\alpha$ and use
\eqref{r12:weak}.  If $c\in\ker Q$, the field
$\sum_\alpha c_\alpha Y_\alpha$ vanishes in $L^2(p)$, so $Dc=0$.
The required range condition therefore holds.
The choice $W=DQ^\dagger$ gives the orthogonal projections
and the displayed value.  This proof also covers linearly
dependent field lists.
\end{proof}

In particular $\mathcal E$ is the limit of these lower matrices
along an increasing dense field family.  Below we give a
quantitative version, rather than relying on density alone.

\subsection{The exact marginal-score Gram uses the same two halves}
\label{r12:score}

For a smooth middle field $Y_\alpha$ set
\[
 h_\alpha=-\operatorname{div}Y_\alpha
                       +\gamma\langle Y_\alpha,\nabla S_s\rangle ,
 \qquad
 t_\alpha^\pm=h_\alpha+2\gamma\langle Y_\alpha,g_\pm\rangle .
 \tag{V12.7}\label{r12:t}
\]
Divergence is relative to product Haar measure.
Their conditional mean is the actual marginal score
\[
 s_\alpha=\E[t_\alpha^\pm\mid x]
                         =-\operatorname{div}_pY_\alpha .
\]
Also define the full-history score and the half-force contrast
\[
 T_\alpha=\tfrac12(t_\alpha^++t_\alpha^-),\qquad
 H_\alpha=\tfrac12(t_\alpha^+-t_\alpha^-)
                       =\gamma\langle Y_\alpha,g_+-g_-\rangle .
 \tag{V12.8}\label{r12:TH}
\]

\begin{theorem}[Exact paired marginal-score moments]
\label{r12:score-thm}
For $J_{\alpha\beta}=\E_p s_\alpha s_\beta$,
\[
 \begin{split}
 J_{\alpha\beta}
   &=\E_\mu[t_\alpha^+t_\beta^-]
     =\E_\mu[T_\alpha T_\beta-H_\alpha H_\beta],\\
 D_{i\alpha}
   &=\E_\mu[A_i^+t_\alpha^-]
     =-\E_\mu[A_i^+H_\alpha].
 \end{split}
 \tag{V12.9}\label{r12:score-identities}
\]
The first expression for $J$ is symmetric after expectation.
With $G=\Cov_p(m)$, every real $k\times r$ matrix $W$ satisfies
\[
 \boxed{\quad
 G\succeq DW^{\mathsf T}+WD^{\mathsf T}-WJW^{\mathsf T}.
 \quad}
 \tag{V12.10}\label{r12:Gdual}
\]
All entries on the right have the ordinary-history representations
above; no conditional normalizer or its derivative is an
unevaluated insertion.
\end{theorem}
\begin{proof}
Conditional independence gives
$\E[t_\alpha^+t_\beta^-\mid x]=s_\alpha s_\beta$.
For the second representation of $J$, conditional variances obey
\[
 \Cov(T_\alpha,T_\beta\mid x)
   =\E[H_\alpha H_\beta\mid x]
   =2\gamma^2\Cov_{\nu_x}
       (\langle Y_\alpha,g\rangle,\langle Y_\beta,g\rangle).
\]
The conditional mean of $H$ is zero.  Subtraction proves the
identity.  This equality removes the conditional noise in
the full score; replacing $J$ by $\E TT^{\mathsf T}$ would only
give an upper bound.

Next, $\E[A_i^+t_\alpha^-\mid x]=m_i s_\alpha$.
Compact-manifold integration by parts gives
\[
 \E_p m_i s_\alpha=\E_p\langle\nabla m_i,Y_\alpha\rangle
                  =D_{i\alpha}.
\]
Since $\E_p s_\alpha=0$, no unknown
source-centering constant is needed.
The contrast identity follows either by conditional covariance
or by reflection symmetry in \eqref{r12:F}.
Finally expand the Gram of $m-\E_pm-Ws$.
\end{proof}

This does not violate V5's contact nullity:
$T_\alpha$ is the full middle divergence and
$\E_\mu[A_i^+T_\alpha]=0$.
The field $t_\alpha^-=T_\alpha-H_\alpha$ is not that divergence.
The nonzero source response is exactly the retained contrast term.
Neither $t^+(t^-)^{\mathsf T}$ nor
$TT^{\mathsf T}-HH^{\mathsf T}$ is claimed pointwise positive.
Positivity of its expectation follows from the conditional identity.

The same identities hold in the fixed-$L,\gamma$ stationary
ground law.  One can first put the sources at $\pm n$ and
boundaries at $\pm M$, then let $M\to\infty$ by V11's
fixed-regulator transfer argument.  The differentiated half
force involves only the first temporal layer, so the insertions
are fixed local functions apart from the declared sources.
This does not give a cutoff-uniform estimate for the vacuum
boundary error or an infinite-volume conclusion.

\subsection{A positive polynomial kernel with an explicit error}
\label{r12:kernel}

Let $dy$ be normalized Haar measure on $S^3$, and for an integer
$d\ge0$ define
\[
 k_d(u)=c_d\left(\frac{1+u_0}{2}\right)^d,\qquad
 c_d=\frac{(3)_d}{(3/2)_d}.
 \tag{V12.11}\label{r12:kernel-def}
\]
Here $(a)_d=a(a+1)\cdots(a+d-1)$, and an empty product is one.
This kernel is central, nonnegative, and has integral one.

\begin{lemma}[Polynomial approximate identity on the product group]
\label{r12:kernel-lemma}
For a smooth vector-valued coefficient function $f$ on $\mathcal Q$,
put
\[
 (\mathcal K_df)(x)=
       \int_{\mathcal Q}\prod_{e=1}^N k_d(x_e y_e^{-1})f(y)\,dy .
 \tag{V12.12}\label{r12:convolution}
\]
Each component of $\mathcal K_df$ is the restriction of a
polynomial of degree at most $d$ in each link quaternion.
If $\sup_x\|df(x)\|_{\rm op}\le C_f$, then
\[
 \sup_x|f(x)-\mathcal K_df(x)|^2
     \le \frac{3\pi^2N}{2(d+3)}C_f^2
     \le \frac{15N}{d+3}C_f^2 .
 \tag{V12.13}\label{r12:kernel-error}
\]
\end{lemma}
\begin{proof}
Under Haar measure, $t=(1+u_0)/2$ has the beta density with
parameters $(3/2,3/2)$.  Thus
$\int t^d\,du=(3/2)_d/(3)_d$, proving normalization, and
\[
 \int k_d(u)(1-u_0)\,du=\frac3{d+3}.
\]
For $\theta=\operatorname{dist}_{S^3}(1,u)\in[0,\pi]$,
$1-u_0=2\sin^2(\theta/2)\ge2\theta^2/\pi^2$.
The product kernel therefore has mean squared product distance
at most $3\pi^2N/[2(d+3)]$.
A shortest product geodesic gives
$|f(x)-f(y)|\le C_f\operatorname{dist}_{\mathcal Q}(x,y)$.
Jensen's inequality proves the error bound; $\pi^2<10$ gives
the last estimate.

Finally $(x_e y_e^{-1})_0$ is the Euclidean scalar product of
their four quaternion coordinates.  Expanding the integrand
in the $x$ variables gives the asserted finite polynomial.
The expansion coefficients are integrals involving $f$;
the lemma proves approximation, not that these coefficients
are already numerically known.
\end{proof}

For vector fields, apply this kernel to their coefficients in
the global right-translated frame, then rebuild the field.
Middle gauge transformations act on these coefficients by
constant orthogonal adjoint matrices at the source vertices.
Centrality of $k_d$ makes this construction commute with those
transformations.  In particular it preserves gauge equivariance.
Alternatively, a complete polynomial field space can be
gauge-averaged without increasing the approximation error.

\subsection{An explicit depth-independent recovery bound}
\label{r12:rate}

Choose a known positive-semidefinite coefficient metric $R$ such that
\[
 \left|\sum_i c_i A_i(z)\right|^2\le c^{\mathsf T}Rc
       \quad\hbox{for every }c,z .
 \tag{V12.14}\label{r12:source-R}
\]
Known constants may first be subtracted from the sources,
which does not affect their gradients or any response.
For example $R=(\sum_i\|A_i\|_\infty^2)I_k$ is valid.

Let $\mathcal Y_d$ contain all middle vector fields with
right-translated coefficients polynomial of degree at most $d$
in each link, and let $L_d$ be \eqref{r12:L} for any spanning
list of this space.  One may instead use its gauge-equivariant
subspace.  Sphere-ideal dependencies do not matter because the
pseudoinverse formula already treats dependent lists.

\begin{theorem}[Quantitative native polynomial recovery]
\label{r12:rate-thm}
Use V11's
$a=\gamma\sqrt N$, $b=\gamma\sqrt{3N}+2\gamma^2N$, and set
\[
 \delta_d=\frac{3\pi^2N}{2(d+3)}(b+\sqrt2\,a)^2,\qquad
 \overline\delta_d=
           \frac{240\gamma^2N^2(1+\gamma^2N)}{d+3}.
 \tag{V12.15}\label{r12:delta}
\]
For every temporal depth $n\ge1$,
\[
 \boxed{\quad
 0\preceq\mathcal E-L_d\preceq\delta_d R
                                  \preceq\overline\delta_d R .
 \quad}
 \tag{V12.16}\label{r12:quantitative}
\]
The matrices $L_d$ increase with $d$.  The same estimate holds
in the fixed-spatial-volume stationary ground law.
No comparison between the native density and Haar density is used.
\end{theorem}
\begin{proof}
Write $f_c$ for the right-translated coefficient vector of
$\nabla m_{A_c}$, where $A_c=\sum_i c_i A_i$.
The unit-$S^3$ connection in this frame gives
\[
 \|df_c\|_{\rm HS}
    \le\|\nabla^2m_{A_c}\|_{\rm HS}
                          +\sqrt2\,|\nabla m_{A_c}|.
\]
For completeness, on each factor the connection coefficients
are the three-dimensional cross product.  Summing
$|e_a\times v|^2$ over its three orthonormal directions gives
$2|v|^2$; summing over factors proves the correction.
V11's pointwise derivative bounds and
$\Var(A_c\mid x)\le c^{\mathsf T}Rc$ therefore imply
\[
 \sup_x\|df_c(x)\|_{\rm op}
                    \le(b+\sqrt2 a)\sqrt{c^{\mathsf T}Rc}.
\]
The field with coefficients $\mathcal K_df_c$ belongs to
$\mathcal Y_d$.  Its uniform squared error is at most
$\delta_d\,c^{\mathsf T}Rc$ by the kernel lemma.
Integrating this pointwise error against the actual $p$ does
not change the bound.  Orthogonal projection in $L^2(p)$
has no larger error than this particular approximant.
The residual characterization in Theorem~\ref{r12:recovery}
now proves the matrix inequality and monotonicity.

Finally,
\[
 (b+\sqrt2a)^2\le2b^2+4a^2
                 \le16\gamma^2N(1+\gamma^2N),\qquad \pi^2<10,
\]
which proves the rational-coefficient envelope.
V11 gives the same derivative estimates in the fixed-volume
ground law, and the same argument applies there.
\end{proof}

The error is absolute in the declared source metric $R$,
not relative to the potentially very small $\mathcal E$.
The kernel coefficients need not be evaluated to use the
theorem: $L_d$ is specified by the ordinary-history moments
$D,Q$.  Nevertheless those moments must actually be enclosed.

The bound is explicit but not economical.  For example, at
$N=192$ and $\gamma=1$, the coarse sufficient bound
$\overline\delta_d\le1/100$ requires
$d\ge170754047997$.
This is not a lower bound on the degree actually needed; it
shows that this worst-case estimate does not furnish a
practical algorithm.  Before gauge reduction, the full field
space has dimension at most
\[
 3N\left\{\frac{(d+1)(d+2)(2d+3)}6\right\}^{N},
\]
using the dimensions $(\ell+1)^2$ of degree-$\ell$ spherical
harmonics on $S^3$.
A sparse bank with a source-adapted error estimate is therefore
the next substantive requirement, not an optional implementation
detail.  Temporal-depth independence of \eqref{r12:delta}
does not remove these spatial costs.

\subsection{Certified moments, errors and the conditional noise floor}
\label{r12:certification}

Suppose the $A_i$ are polynomial and the field coefficients have
total degree at most $q$ in the middle link quaternions.
Right-invariant link derivatives preserve polynomial degree.
All half forces and spatial-action derivatives have degree at
most four.  Consequently the insertions for $D,Q,J$ have total
degrees at most, respectively,
\[
 \max_i\deg A_i+q+4,\qquad 2q,\qquad 2q+8.
\]
Their degrees do not grow when more temporal layers are added.
The number of integrated links and the difficulty of the native
expectation problem can still grow.

V3's \emph{same one-history} moment hierarchy gives convergent
outer bounds for these polynomial objectives at each fixed
finite regulator.  It is not necessary to introduce four
conditionally independent half histories or to encode
$Z_+^{-2}$.  The force-gradient approximation theorem above is
additional to V3's convergence result; it is not a claimed
rate for that hierarchy.

For a fixed chosen $W$, suppose certified entry bounds imply
operator-norm bounds
$\|D-\widehat D\|\le\epsilon_D$ and
$\|J-\widehat J\|\le\epsilon_J$, with $\widehat J$ symmetric.
Equation \eqref{r12:Gdual} then gives
\[
 G\succeq
 \widehat DW^{\mathsf T}+W\widehat D^{\mathsf T}
      -W\widehat JW^{\mathsf T}
      -(2\|W\|\epsilon_D+\|W\|^2\epsilon_J)I_k .
 \tag{V12.17}\label{r12:robust}
\]
This follows by the triangle inequality for the three error
terms.  The identical formula with $Q,\epsilon_Q$ bounds
$\mathcal E$ from below.
Thus the certified step need not invert a nearly singular
moment matrix.  A floating optimizer may suggest $W$; a
validated moment enclosure or exact V3 dual witness must
verify the resulting inequality.  No four-dimensional
numerical certificate is claimed in this appendix.

An important distinction remains even though the normalizers
have disappeared from these insertions.  For one source put
$a_0=A-m$, $u=g-\bar g$, and $v=\E_{\nu_x}a_0^2$.
Direct expansion of the two-copy square gives
\[
 \Cov(F\mid x)=\frac{\gamma^2}{2}
       \left\{\E_{\nu_x}[a_0^2u\otimes u]
                          +v\,\E_{\nu_x}[u\otimes u]\right\}.
 \tag{V12.18}\label{r12:noise}
\]
Indeed the fourth-moment expansion of
$(a_0^+-a_0^-)^2(u^+-u^-)\otimes(u^+-u^-)$ contains
twice the first term, twice the second term, and four times
$\Cov(A,g)\otimes\Cov(A,g)$; subtracting the conditional
mean square proves the identity.
In particular,
\[
 \E_\mu|F|^2
   =\E_p|\nabla m|^2+\E_p\operatorname{tr}\Cov(F\mid x).
\]
The directly measurable second moment of $F$ is generally
\emph{larger} than the desired Dirichlet energy.  Using it as
a lower bound would be erroneous.  The weak tests and
projection in \eqref{r12:Edual} remove this conditional noise
at the population level; they do not prove a favorable
sampling signal-to-noise ratio.

\paragraph{Verification and preservation.}
The checker \path{scripts/verify_ym_restart_v12_paired_recovery.py}
tests the reflected-pair identities, the score-noise subtraction,
the projection residual including dependent fields, and
\eqref{r12:noise} with exact rational finite-support fixtures.
Exact Gaussian fixtures additionally test the crossed-source
response against the contact-null full score.
These are algebra fixtures, not quadratures of the Wilson law.
It also checks exact $S^3$ power-kernel normalization and first
moments, the dimension formula, and the quoted degree threshold.
All analytic native-law claims are proved above, not inferred
from those tests.
V11 is preserved exactly apart from the displayed title;
its recoverable SHA--256 is
\path{D47C045E828FEFA8E5D1796F11167C138980D1FDF5538198444178E3273791BA}.
Only the active predecessor is replaced.
Companion and archive sources remain unchanged; build files
and diagnostics remain outside \path{latex_notes}.

\paragraph{Status and the next attack.}
The explicit conditional-normalizer obstruction in V11's
four-copy computation has been bypassed by a different exact
two-half variational formulation.  The finite-native response
and score matrices are now ordinary polynomial expectations,
and the source-gradient recovery has a proved absolute rate.
What remains is a sparse, quantitatively useful native
certificate at growing depth and volume, together with the
vacuum/continuum uniformity and physical-transfer estimates.
The interacting four-dimensional continuum construction and
its Yang--Mills mass gap are still unproved.

% END CONSOLIDATED V12 APPEND

% BEGIN CONSOLIDATED V13 APPEND -- 2026-09-19
\clearpage
\section{V13: A local tree-curvature field and native one-sided score sources}
\label{r13:start}

V12 removed an explicit conditional-normalizer insertion, but its
complete polynomial field space was too large.  We now specify
one field at each radius, with polynomial degree proportional to
that radius, and bound all of its original-link derivatives.
The field averages curvature transported along fixed local trees;
it is not the frozen covariant heat filter of V9.

There are two separate positive results.  The new field has a
volume-independent native score upper bound and a nonvanishing
Gaussian-reference capture at a matched radius.  Separately,
one-sided score insertions give explicit polynomial history
sources whose conditional projections are marginal scores.
Their stationary native response is positive by the physical
transfer operator.  We do not identify this new source with an
original flowed plaquette, infer a uniform lower bound from
strict positivity, or assert a continuum limit.

\subsection{What is inherited and what is changed}
\label{r13:scope}

The finite native measure, unit-$S^3$ metric, actual physical
transfer and V12's paired formulas are unchanged.
The f15 V6--V7 growing-depth Gaussian comparisons and f16.2 V67's
physical-clock warning were checked: repeating a mesoscopic
comparison would not pay the missing physical-scale estimate.
The f16.2 V67 covariant heat-curvature source already averages
transported plaquettes, but averages over covariant random walks.
The f15 V73 tree-averaged Ward completion already shows why
dressing-link derivatives must be counted.
Neither is claimed as a new construction here.

The new finite-radius tree prescription below has no random
connection Laplacian, no field-dependent denominator, and no
discarded dressing derivative.  It uses a compactly supported
triangular kernel rather than importing V7's heat-filter
$8/9$ estimate.  We evaluate its own reference score constant.

A targeted primary-record check of von Hippel, J\"ager, Rae and
Wittig,
\href{https://arxiv.org/abs/1306.1440}
{\emph{The Shape of Covariantly Smeared Sources in Lattice QCD}},
JHEP 09 (2013), 014, found the relevant warning: the shape of
an iteratively covariantly smeared source can depend strongly
on the gauge background.  Their quark-source construction is
not a theorem for the present Wilson score.  Only the primary
record and abstract were used; the proofs below are independent.
This targeted check does not replace the scheduled V15 companion
review or V20 broad literature review.

\subsection{A finite-radius polynomial curvature energy}
\label{r13:construction}

Let $m\ge1$ be an integer, and take a spatial torus with
$L\ge4m$, $V=L^3$ and $N=3V$.  Define a one-dimensional
probability kernel and its three-dimensional product by
\[
 b_m(j)=\frac1m{\bf1}_{\{0,\ldots,m-1\}}(j),\qquad
 w_m(r)=\prod_{\ell=1}^3(b_m*b_m)(r_\ell).
 \tag{V13.1}\label{r13:w}
\]
Thus $w_m\ge0$, $\sum_rw_m(r)=1$, and its support is
$\{0,\ldots,2m-2\}^3$.

For each anchor $x$ and offset $r$, let $\tau_{x,r}$ be the
ordered link product along the path from $x$ to $x+r$ which
first traverses direction 1, then direction 2, then direction 3.
For a fixed anchor these paths lie in a fixed rooted tree.
They are definitions of transport, not gauge fixing.
Write $P_{y,ij}$ for the positively oriented spatial plaquette
based at $y$, and set
\[
 \begin{split}
 C_{x,r,ij}(U)
   &=\operatorname{Im}\big(\tau_{x,r}P_{x+r,ij}
                                      \tau_{x,r}^{-1}\big)\in\R^3,\\
 B_{x,ij}^{(m)}&=\sum_rw_m(r)C_{x,r,ij},\\
 f_x^{(m)}&=-\frac12\sum_{i<j}|B_{x,ij}^{(m)}|^2,\qquad
 F_m=\frac1{\sqrt V}\sum_x f_x^{(m)},\qquad Y_m=\nabla F_m .
 \end{split}
 \tag{V13.2}\label{r13:F}
\]
All imaginary parts use the quaternion convention, so
$|C_{x,r,ij}|\le1$.  Each transported plaquette transforms
by conjugation at its anchor, and therefore $f_x^{(m)}$
and $F_m$ are gauge invariant.  The field $Y_m$ is gauge
equivariant.

Every word in $C$ has length at most
\[
 p_m=12m-8,\qquad \kappa_m=(2m)^3.
 \tag{V13.3}\label{r13:p}
\]
The local scalar has polynomial degree at most $2p_m$ in
quaternion coordinates, including inverse links by conjugation.
Its support lies in the anchor cube with vertices
$x+\{0,\ldots,2m-1\}^3$.
The same degree bound holds for right-translated coefficients
of $Y_m$.  No non-polynomial projection back to the group is used.

This is constituent-curvature averaging before squaring.
A spatial average of the scalar $|\operatorname{Im}P|^2$
would be a different construction and would not have the
two-constituent Fourier filter below.

\subsection{All dressing derivatives and the native score cost}
\label{r13:derivatives}

\begin{theorem}[Volume-independent score budget for the specified field]
\label{r13:budget-thm}
At every link configuration,
\[
 |\nabla F_m|^2\le27\kappa_m p_m^2,\qquad
 \|\nabla^2F_m\|_{\rm HS}^2\le676\kappa_m p_m^4.
 \tag{V13.4}\label{r13:deriv-bound}
\]
In any of V11's reflected finite Wilson histories, the actual
middle marginal score $s_m=-\operatorname{div}_pY_m$ satisfies
\[
 \boxed{\quad
 \E_p s_m^2\le
 \kappa_m\{676p_m^4+27(2+24\gamma)p_m^2\}.
 \quad}
 \tag{V13.5}\label{r13:score-bound}
\]
This holds for every $\gamma\ge0$ and temporal depth.
The same bound holds in the fixed-spatial-volume stationary
ground law.  It is uniform in total spatial volume at fixed
$m,\gamma$, not uniform in radius or coupling.
\end{theorem}
\begin{proof}
Consider a quaternion word of length at most $p=p_m$, and
let $\nu_e$ count occurrences of link $e$, with either
orientation.  Each right-invariant unit derivative replaces
one factor by a product with a unit imaginary quaternion.
Its norm is at most one, also for an inverse link.
Thus, for its imaginary part $C$,
\[
 \|D_{e,a}C\|\le\nu_e,\qquad
 \|D_{e,a}D_{f,b}C\|\le\nu_e\nu_f.
\]
The second estimate includes repeated occurrences and
derivatives hitting the same factor.  Since
$\sum_e\nu_e\le p$, summing the three frame directions gives
\[
 \|DC\|_{\rm HS}\le\sqrt3\,p,\qquad
 \|D^2C\|_{\rm HS}\le3p^2 .
\]
Here $D^2$ is the ordered frame derivative, before the
Levi-Civita correction.  The convex average defining $B$
has the same bounds, by the triangle inequality in the
corresponding tensor spaces, and $|B|\le1$.

For the three plaquette orientations, differentiating
$f_x=-\frac12\sum_{i<j}|B_{x,ij}|^2$ now gives
\[
 |\nabla f_x|\le3\sqrt3\,p,\qquad
 \|D^2 f_x\|_{\rm HS}\le18p^2.
\]
As in V12, the frame-to-covariant-Hessian correction has
norm $\sqrt2\,|\nabla f_x|$.  Therefore
\[
 \|\nabla^2 f_x\|_{\rm HS}
          \le18p^2+3\sqrt6\,p\le26p^2 .
\]
In particular transport-link derivatives have not been
dropped on the grounds that a plaquette is at another site.

A fixed oriented link is in the supporting cube of at most
$\kappa_m$ anchors.  The same upper count holds for any
fixed pair of links.  Apply Cauchy--Schwarz to the sum over
anchors in each gradient component, or each Hessian entry,
and then sum the entries:
\[
 \left|\frac1{\sqrt V}\sum_x\nabla f_x\right|^2
     \le\frac{\kappa_m}{V}\sum_x|\nabla f_x|^2,\qquad
 \left\|\frac1{\sqrt V}\sum_x\nabla^2f_x\right\|_{\rm HS}^2
     \le\frac{\kappa_m}{V}\sum_x\|\nabla^2f_x\|_{\rm HS}^2 .
\]
There are $V$ anchors, proving \eqref{r13:deriv-bound}.
V9's native projected-score inequality, with Ricci tensor
$2g$ and V6's full-action Hessian bound $24g$, proves
\eqref{r13:score-bound}.  All incident temporal plaquettes
on both sides are retained.  Its fixed-volume vacuum
passage is the one already justified in V11.
\end{proof}

For a polynomial outer source $A$, V12 therefore applies to
this one-field bank with response and \emph{exact} score
Gram given by ordinary-history polynomial insertions.
Their degrees are at most
\[
 \deg A+2p_m+4\quad\hbox{and}\quad4p_m+8,
\]
respectively.  The coarse upper budget above is optional:
a certified paired evaluation of the actual score Gram
can be much sharper.

V9's frozen-connection heat suppression theorem does not
apply to this different field.  That observation is not
a proof of native retention.  Tree transports can still
produce cancellations in the interacting law, and the
present derivative bound does not bound the response
from below.

The same counting allows spatial tests:
for $F_{m,c}=\sum_x c_x f_x^{(m)}$, the right sides of
\eqref{r13:deriv-bound} and \eqref{r13:score-bound} are
multiplied by $\sum_x|c_x|^2$.
Indeed apply the componentwise Cauchy--Schwarz step to
$c_x\nabla f_x$ and $c_x\nabla^2 f_x$.
The displayed $F_m$ is the unit-norm constant profile.
This extension does not establish source-core coverage.

\subsection{An evaluated scale-matched reference test}
\label{r13:reference}

At $U_e(\varepsilon)=\exp(\varepsilon u_e)$, the first
variation of $C_{x,r,ij}$ is the ordinary plaquette curl
$(du)_{x+r,ij}$.  The first transport variation multiplies
the zero curvature at the identity.  Hence the quadratic
part of $F_m$ is exactly
\[
 -\frac1{2\sqrt V}\sum_{x,i<j}
                         |(w_m*(du)_{ij})(x)|^2 .
\]
In V7's transverse Gaussian reference its multiplier is
\[
 q_m(k)=\lambda(k)\,\omega_m(k),\qquad
 \omega_m(k)=\prod_{\ell=1}^3
 \left|\frac{\sin(mk_\ell/2)}
                  {m\sin(k_\ell/2)}\right|^4 .
 \tag{V13.6}\label{r13:q}
\]
The factors at zero are defined by continuity.
This derives the filter from the compact observable; it
does not replace the native law by its linearization.

Keep the unflowed magnetic-energy message of V7, with
$a_n=\lambda e^{-2n\theta}$ and $R=\sqrt{\lambda(\lambda+4)}$.
Write
\[
 \mathsf H_{n,L}=\frac1{2V}\Tr(a_n^2R^{-2}),\quad
 \mathsf B_{n,L}(q)=\frac1V\Tr(qa_nR^{-1}),\quad
 \mathsf J_L(q)=\frac2V\Tr(q^2).
 \tag{V13.7}\label{r13:reference-Grams}
\]
These are reference quantities only.

\begin{theorem}[One compact-kernel field retains a fixed reference fraction]
\label{r13:capture}
Suppose $n\to\infty$, $L/n\to\infty$ and $m/n\to\alpha>0$.
Put $\operatorname{sinc}t=\sin(t)/t$, with its continuous
value at zero, and
\[
 W_\alpha(y)=\prod_{\ell=1}^3
                    \operatorname{sinc}(\alpha y_\ell/2)^4 .
\]
Then
\[
 \begin{split}
 n^5\mathsf H_{n,L}&\longrightarrow
                       h=\frac9{1024\pi^2},\\
 n^6\mathsf B_{n,L}(q_m)&\longrightarrow
    \beta_\alpha=\frac3{8\pi^3}
          \int_{\R^3}|y|^3e^{-2|y|}W_\alpha(y)\,dy>0,\\
 n^7\mathsf J_L(q_m)&\longrightarrow
    j_\alpha=\frac{1236992}{33075\,\alpha^7}.
 \end{split}
 \tag{V13.8}\label{r13:limits}
\]
Consequently the optimized reference capture tends to
$\beta_\alpha^2/(j_\alpha h)\in(0,1)$.
For the concrete choice $m=2n$, the completely specified
rational certificate coefficient
\[
 c_*=\frac{27783}{773120},\qquad W_n=c_* n
 \tag{V13.9}\label{r13:coefficient}
\]
satisfies
\[
 \lim_{n,L}
 \frac{2W_n\mathsf B_{n,L}(q_{2n})
                   -W_n^2\mathsf J_L(q_{2n})}
      {\mathsf H_{n,L}}>\frac38.
 \tag{V13.10}\label{r13:fraction}
\]
This is a limiting lower fraction, not a finite-$n,L$
inequality or a native interacting estimate.
\end{theorem}
\begin{proof}
There are six real transverse degrees per nonzero momentum.
Set $k=y/n$ in the exact traces.  Then
$\lambda\sim|y|^2/n^2$, $R\sim2|y|/n$,
$\theta\sim|y|/n$ and $\omega_m\to W_\alpha$.
This gives the stated $h$ and $\beta_\alpha$, and
\[
 j_\alpha=\frac{12}{(2\pi)^3}
          \int_{\R^3}|y|^4 W_\alpha(y)^2\,dy .
 \tag{V13.11}\label{r13:Jintegral}
\]
For the last integral and its Riemann sums, use
\[
 \left|\frac{\sin(mk/2)}{m\sin(k/2)}\right|
            \le\min\{1,\pi/(m|k|)\},\qquad |k|\le\pi .
\]
After rescaling, a fixed multiple of
\[
 |y|^4\prod_{\ell=1}^3(1+|y_\ell|)^{-8}
\]
bounds the integrand.  It is integrable: expand $|y|^4$
into the three fourth-power terms and three mixed
second-power terms.  Each resulting product has integrable
one-dimensional tails.  The same elementary power-tail
bounds control their Riemann sums uniformly once the
mesh $2\pi n/L$ is at most one.
On bounded sets convergence is uniform.
For the response, the additional factor
$e^{-2n\theta}$ supplies an exponential tail.
These observations justify the stated joint limits.

To evaluate \eqref{r13:Jintegral}, define
\[
 I_{2a}=\int_{\R}t^{2a}\operatorname{sinc}(t/2)^8\,dt,
 \qquad a=0,1,2.
\]
The sum of eight independent uniform variables on
$[-1/2,1/2]$ has density
\[
 f_8(x)=\frac1{7!}\sum_{r=0}^8(-1)^r\binom8r(x+4-r)_+^7 .
\]
This follows by inclusion--exclusion for the volume of a
truncated cube, followed by one derivative.
Fourier inversion, differentiated up to order four
(the integrals are absolutely convergent), gives
\[
 \frac{I_0}{2\pi}=\frac{151}{315},\qquad
 \frac{I_2}{2\pi}=\frac23,\qquad
 \frac{I_4}{2\pi}=\frac83 .
\]
For example these numbers equal, with the appropriate
alternating derivative sign,
\[
 \frac1{(7-2a)!}
        \sum_{r=0}^3(-1)^r\binom8r(4-r)^{7-2a}.
\]
Expanding $|y|^4$ gives
$\alpha^{-7}(3I_4I_0^2+6I_2^2I_0)$.
Substitution yields the displayed rational $j_\alpha$.
The general capture formula follows from the exact
Gaussian score certificate.  Cauchy--Schwarz bounds it
by one.  Equality would require $q$ proportional to the
exact matching multiplier $a_n/(2R)$ in the limiting
weighted space; near zero the former is quadratic in
$|y|$ and the latter linear, so equality is impossible.

Here is a rational lower certificate at $\alpha=2$.
For $|t|\le2$, the alternating sine series gives
$\operatorname{sinc}t\ge1-t^2/6$.
On the ball $|y|\le2$, all three factors on the right are
nonnegative, and
$\prod_\ell(1-y_\ell^2/6)\ge1-|y|^2/6$.
Consequently
\[
 W_2(y)\ge(1-|y|^2/6)^4,\qquad
 \beta_2\ge\frac3{2\pi^2}\int_0^2
                     r^5e^{-2r}(1-r^2/6)^4\,dr .
\]
Taylor's signed remainder gives
$e^{-2r}\ge\sum_{k=0}^{63}(-2r)^k/k!$.
Multiplication by the nonnegative factor and integration give
the rational lower bound
\[
 Q=\sum_{j=0}^4\sum_{k=0}^{63}
 \frac{(-1)^{j+k}\binom4j\,2^{6+2j+2k}}
      {6^j k!(6+2j+k)} .
\]
Using $\pi<22/7$, exact rational arithmetic verifies
\[
 \beta_2>\frac{3Q}{2(22/7)^2}
                  >b_*:=\frac{21}{2000}.
 \tag{V13.12}\label{r13:beta-lower}
\]
The displayed finite sum and inequality are checked in the
diagnostic without quadrature or rounding.
Also $\pi>3$ gives $h<1/1024$, and
$j_2=9664/33075$.
Now $c_*=b_*/j_2$, so the limiting fraction of the
specified coefficient is
\[
 \frac{2c_*\beta_2-c_*^2j_2}{h}
   >\frac{1024b_*^2}{j_2}
   =\frac{583443}{1510000}>\frac38.
\]
All final comparisons are rational.
\end{proof}

The bound is deliberately conservative; its role is to
prove nonvanishing capture for a single explicitly compact
polynomial field, not to claim an optimal filter.
At $m=2n$ the native upper score budget
\eqref{r13:score-bound} grows as $O(n^7+\gamma n^5)$,
whereas its actual Gaussian score is of order $n^{-7}$.
The large mismatch is visible.  The coarse native budget
cannot be substituted into the reference calculation to
obtain a useful native fraction.

\subsection{Native response positivity through one-sided score sources}
\label{r13:lifts}

We now change the source, not the physical transfer.
Work in the \emph{stationary} Wilson ground law at fixed
$L,\gamma>0$.  Let $P$ be its positive injective
self-adjoint Markov transfer on gauge-invariant slice
$L^2(p)$, as constructed in the archive.
For any smooth equivariant middle field $Y$, let
$s_Y=-\operatorname{div}_pY$.
Define an explicit source supported on a slice $t$ and
the temporal slab immediately to its future by
\[
 A_Y^+(t)=
 -\operatorname{div}Y(U_t)
 +\gamma\langle Y(U_t),\nabla S_s(U_t)\rangle
 +2\gamma\langle Y(U_t),g_{t,t+1}\rangle .
 \tag{V13.13}\label{r13:lift}
\]
The force $g_{t,t+1}$ differentiates the actual plaquettes
in that one slab with respect to its initial spatial slice.
Temporal links in the slab remain integrated.
Let $A_Y^-(-t)$ denote the reflected past-supported source.
For a polynomial field, these are polynomial, gauge-invariant
history observables; no Perron function appears in their
definition.

\begin{theorem}[Polynomial lifts of native marginal-score states]
\label{r13:lift-thm}
Exactly,
\[
 \E[A_Y^+(t)\mid U_t]=s_Y(U_t),\qquad
 \E[A_Y^+(n)\mid U_0]=P^n s_Y(U_0),\quad n\ge0 .
 \tag{V13.14}\label{r13:lift-message}
\]
For fields $Y_1,\ldots,Y_r$ put $s_i=s_{Y_i}$ and
\[
 J_{ij}=\langle s_i,s_j\rangle_p,\quad
 D_n(i,j)=\langle s_i,P^n s_j\rangle_p,\quad
 G_n(i,j)=\langle s_i,P^{2n}s_j\rangle_p .
 \tag{V13.15}\label{r13:lift-Grams}
\]
Here $G_n$ is the connected reflected Gram of
$A_{Y_i}^-(-n)$ and $A_{Y_j}^+(n)$.
The field response at the middle to these sources is
$D_n$, a symmetric positive-semidefinite matrix.
It is positive definite at every finite $n$ when the
scores are linearly independent.
Moreover
\[
 G_n\succeq D_n W^{\mathsf T}+W D_n-WJW^{\mathsf T}
 \tag{V13.16}\label{r13:lift-dual}
\]
for every real $r\times r$ matrix $W$.
\end{theorem}
\begin{proof}
The first identity is V12's conditional identity for
$t_Y^+$.  It can also be obtained by differentiating the
finite-cylinder half normalizer and then passing to the
fixed-volume ground law as in V11.  Only the first
temporal slab is differentiated.
The stationary Markov property gives the second identity.

Marginal scores have mean zero.  Conditional independence
at the middle plane gives the reflected Gram $G_n$.
Integration by parts identifies the response:
\[
 \E_p\langle\nabla(P^n s_i),Y_j\rangle
     =\langle P^n s_i,s_j\rangle_p=D_n(i,j).
\]
Self-adjoint positivity of $P$ makes $D_n$ positive
semidefinite.  Its injectivity implies that
$\langle s,P^ns\rangle>0$ for nonzero $s$ and finite $n$,
proving the strict statement.
Finally the Gram of the residual vector $P^ns-Ws$
is nonnegative and expands to \eqref{r13:lift-dual}.
\end{proof}

These sources are not V85's contact-null full divergences.
The full local middle score is
$T_Y=(A_Y^++A_Y^-)/2$; replacing one one-sided lift by
$T_Y$ changes its time support and its projection onto an
earlier reflection slice.  Both have conditional mean $s_Y$
on their own slice, but for $n\ge1$ integration by parts at
slice $n$ gives $\E[T_Y(n)\mid U_0]=0$: a test depending only
on $U_0$ has zero derivative in $U_n$.  The past dependence
of $T_Y(n)$ prevents use of the forward Markov identity
\eqref{r13:lift-message} for that full score.
The source at $n$ is supported at $n,n+1$ and must not be
silently extended to include $n-1$.
Nor does positivity for this deliberately aligned source
prove positivity of the response to an original flowed
plaquette or magnetic-energy source.

For the concrete $Y_m$, its score is nonzero.
Indeed a transverse mode with momentum $(2\pi/L,0,0)$
has $q_m(k)>0$ when $L\ge4m$, so the derived quadratic
variation proves that $F_m$ is nonconstant.
If $s_m$ vanished, integration by parts would give
$0=\E_p(F_m s_m)=\E_p|\nabla F_m|^2$, contradicting
nonconstancy and the strictly positive density.
Thus this specific one-field lifted source has
$J>0$, $D_n>0$ and $G_n>0$ at every finite delay for
fixed $L,\gamma>0$.  No uniform floor is inferred.

\begin{proposition}[Fixed-regulator completeness of the lifted source core]
\label{r13:dense}
As $Y$ ranges over all gauge-equivariant polynomial fields,
the span of the states $s_Y$ is dense in the centered
gauge-invariant subspace of $L^2(p)$ at fixed regulator.
\end{proposition}
\begin{proof}
Suppose a gauge-invariant $f\in L^2(p)$ is orthogonal to
all these scores.  Polynomial coefficient fields are
$C^1$ dense in smooth coefficient fields on the compact
product of spheres, by the polynomial approximation
argument already used in V3.  Gauge averaging preserves
polynomiality and converges in $C^1$; the score map
commutes with gauge averaging because $p$ is invariant.
Orthogonality therefore holds for scores of all smooth
fields, including non-equivariant ones after averaging
against the invariant $f$.

Consequently $\int f\,\operatorname{div}(pY)\,dx=0$
for every smooth vector field $Y$.
At fixed regulator $p$ is smooth and strictly positive,
so $pY$ ranges over every smooth vector field.
Thus the distributional gradient of $f$ vanishes.
On the connected product manifold, $f$ is constant.
The orthogonal complement in the centered subspace
is therefore zero.
\end{proof}

This proposition concerns the \emph{entire} polynomial
field algebra.  The single translation-averaged
tree-curvature field $Y_m$, or just its radius family,
is not claimed to cover every momentum, parity or other
physical sector.  The missing-sector warning of
f16.2 V86 remains in force.

\subsection{The remaining quantitative task}
\label{r13:next}

The new native lift removes a sign ambiguity by a specified
change of source family, and gives a finite-regulator dense
source core.  It does not bound $D_n$ below uniformly in
$n,L,\gamma$, control the physical normalization of that
core, or prove any upper physical spectral edge.
Even a positive fixed-regulator $J$ can approach zero
along a continuum trajectory.

For $Y_m$, the lift is a single explicitly defined
polynomial source of degree at most $2p_m+4$ and stated
temporal support.  Relative to each local density, its
magnetic derivative adds one spatial plaquette collar;
that enlargement is retained.  Its native $J,D_n,G_n$
can be written using V12's paired insertions in
stationary histories.  Finite open cylinders give
ordinary polynomial integrals, but certified passage
to the stationary values still requires a boundary
error bound; their different boundary weights must
not be ignored.

The next useful calculation is a native enclosure of
these paired matrices, or a spatially controlled
comparison improving the very coarse bound
\eqref{r13:score-bound}.  For original magnetic sources,
the sign of the corresponding mixed response has not
been determined by our argument, and its size remains
to be estimated.
For the aligned score-source core, positivity is now
paid, but the lower size, dense limiting coverage and
common growing-depth upper contraction are not.

\paragraph{Verification and preservation.}
The checker \path{scripts/verify_ym_restart_v13_tree_sources.py}
checks exact kernel weights, tree incidences, noncommuting
quaternion two-jets, rational spline and capture constants,
and finite positive-transfer Grams.  These are not Wilson
quadratures; optional floating reference traces are not
certified bounds.  V12's body is preserved except for the
displayed title, with recoverable SHA--256
\path{CAC7B52932579BD28B778F36EAEA2678642167B1B6BB5CB6BEC5C1FF88FCB831}.
Only the active predecessor is replaced.  Archives and
companions are unchanged; all build files remain outside
\path{latex_notes}.

\paragraph{Status.}
The native finite-regulator results and separate reference
certificate do not establish an interacting four-dimensional
continuum or its Yang--Mills mass gap.

% END CONSOLIDATED V13 APPEND

% BEGIN CONSOLIDATED V14 APPEND -- 2026-09-19
\clearpage
\section{V14: Periodic score matrices and certified vacuum errors}
\label{r14:start}

V13's positivity is not a quantitative contraction.  This iteration
works on the missing interface between finite Wilson integrals and
stationary score matrices.  We prove an exact differential
representation of the one-sided insertions, a positive periodic
matrix formula, and a one-sided upper bound for the vacuum response.
A separate signed-history estimate pays the error in the score
normalization.  Their combination is a finite-data sufficient
condition for contraction on a declared source bank.
An additional two-period inequality below replaces the coarse
pointwise boundary error by a source-weighted spectral subtraction.
This removes the explicit amplitude envelope from that analytic
boundary step, while keeping the purity and integration costs.

All these statements concern the actual finite-regulator Wilson
transfer, at arbitrary $\gamma>0$, not its Gaussian approximation.
They do not include a verified uniform numerical or analytic
contraction margin along a continuum sequence.  The construction
of the interacting continuum is also still missing.

\subsection{Inherited facts and the nonduplication audit}
\label{r14:inputs}

Fix the spatial torus, $N=3L^3$, and the unreduced product
$\mathcal X=\SU(2)^N$ with probability Haar measure $dx$.
The transfer below acts on the gauge-invariant subspace of
$L^2(dx)$.  Write
\[
 T(x,y)=\int
 \exp\{-\gamma[S_s(x)/2+S_t(x,y,h)+S_s(y)/2]\}\,dh ,
 \qquad T\psi=\Lambda\psi,\quad \|\psi\|_2=1 .
 \tag{V14.1}\label{r14:T}
\]
Every temporal link $h$ is integrated.  Its integration implements
the Gauss projection; all traces below are physical traces, not
traces of an ungauged temporal-gauge product kernel.
The previously established Wilson properties are used: $T$ is
real, positive, trace class, smoothing, and has a simple strictly
positive vacuum $\psi$.  It is injective on the physical subspace.
Set $K=T/\Lambda$, $p=\psi^2$, and let $P$ be the unitarily
equivalent ground-state Markov transfer.  Thus $K(\psi f)=\psi Pf$.

The archive f15 V58 already defines $Z_q=\Tr T^q$ and
$R_q=Z_{2q}/Z_q^2$, and proves
$R_q\le \Lambda^q/Z_q\le\sqrt{R_q}$.  That purity ratio is
\emph{inherited}, not a new proposal.  Its positive thermal
transition ledger is also inherited.  f15 V80--V91 already gives
finite-rank bounds and tests global projective contraction; these
are not repeated.  f6 V92--V93's conditional-projection
commutators concern a different operator.  The new use here is
the explicit half-density derivative of V13's polynomial score
insertions, including their adjacent-slab normalization at $n=0$,
and the resulting certified stationary matrix interval.

A targeted primary-source check read Della Morte--Giusti,
\href{https://arxiv.org/html/0806.2601}
{\emph{Symmetries and exponential error reduction in Yang--Mills
theories on the lattice}}, Sections 2--3, especially (2.12)--(2.20)
and (3.32)--(3.33).  Its physical trace and symmetry-projection
bookkeeping are useful.  Its numerical $\SU(3)$ parity-sector
study is not a certified all-sector gap estimate for our
$\SU(2)$ sequence.  No numerical or continuum theorem is
imported.  The next regular companion and broad literature
reviews remain V15 and V20.

\subsection{The one-sided insertion is a half-density derivative}
\label{r14:derivative}

For a real smooth equivariant field $Y$, define on smooth functions
\[
 \mathcal L_Y f=Y\cdot\nabla f+\tfrac12(\operatorname{div}Y)f .
 \tag{V14.2}\label{r14:L}
\]
Haar integration by parts gives
$\mathcal L_Y^*=-\mathcal L_Y$ on this core.  Equivariance
and invariance of Haar divergence make it preserve physical
functions.  Only this integration-by-parts identity and its
compositions with smoothing kernels are needed below; no
unbounded-operator trace is asserted without such smoothing.

Let $B_Y^+$ be the single-slab kernel with the insertion
$A_Y^+$ of \eqref{r13:lift} at its initial slice.  Let $B_Y^-$
be its reflected kernel, with the past-supported insertion
at the final slice of that slab.

\begin{proposition}[Exact score-insertion operators]
\label{r14:B}
As operators defined by their smooth kernels,
\[
 B_Y^+=-2\mathcal L_YT,\qquad
 B_Y^-=2T\mathcal L_Y=(B_Y^+)^*,\qquad
 \psi s_Y=-2\mathcal L_Y\psi ,
 \quad s_Y=-\operatorname{div}_pY .
 \tag{V14.3}\label{r14:B-formula}
\]
In particular, for a field bank and $n\ge0$,
\[
 D_n(i,j)
   =4\langle\mathcal L_{Y_i}\psi,
                   K^n\mathcal L_{Y_j}\psi\rangle ,
 \qquad D_0=J .
 \tag{V14.4}\label{r14:D}
\]
\end{proposition}
\begin{proof}
Differentiating \eqref{r14:T} in its initial variable gives
the multiplier $-\gamma YS_s/2-\gamma YS_t$ inside
the temporal-link integral.  Therefore the kernel of
$-2\mathcal L_YT$ is precisely
$[-\operatorname{div}Y+\gamma YS_s+2\gamma YS_t]$
times the original slab density, integrated over the same $h$.
No derivative of a separately normalized temporal law is taken.
Reflection and integration by parts in the final variable give
the second identity.  Since $p=\psi^2$,
$\psi s_Y=-\psi\operatorname{div}Y-2Y\cdot\nabla\psi$,
which is the third identity.  Conjugating $P^n$ by
multiplication by $\psi$ proves \eqref{r14:D}.
\end{proof}

The divergence term is essential.  Replacing $\mathcal L_Y$
by $Y\cdot\nabla$ would generally destroy skew symmetry and
the following positivity.  These operators are not obtained
by replacing the one-sided source with V13's contact-null
full local score.

\subsection{A positive periodic matrix that bounds the vacuum response}
\label{r14:periodic}

Take a temporal period $M\ge4$ and $0\le n\le M-2$.  In the
literal periodic Wilson law define
\[
 C_n^{(M)}(i,j)
   =\E_M[A_{Y_i}^-(0)A_{Y_j}^+(n)],\qquad
 p_{0,M}=\frac{\Lambda^M}{Z_M}.
 \tag{V14.5}\label{r14:C}
\]
The two inserted slabs are $[-1,0]$ and $[n,n+1]$.
For $n=0$ they share the middle slice but are distinct slabs.
This is not a squared insertion on a single slab.
Each individual one-sided source has periodic mean zero:
reflection equates its forward and backward force means,
and integration by parts at its slice gives the claim.
Equivalently $\Tr(\mathcal L_Y T^M)=0$ by reality and
skew symmetry.  Thus \eqref{r14:C} already has the
appropriate centering.

\begin{theorem}[Positive periodic score matrices and an upper certificate]
\label{r14:periodic-thm}
The real matrix $C_n^{(M)}$ is symmetric positive semidefinite.
Exactly,
\[
 C_n^{(M)}(i,j)=
 -\frac4{Z_M}\Tr\!\left(
        T^{M-n}\mathcal L_{Y_i}T^n\mathcal L_{Y_j}\right).
 \tag{V14.6}\label{r14:trace}
\]
Here and below this trace can equivalently be defined by
the Hilbert--Schmidt form in the proof.  In Loewner order,
\[
 C_n^{(M)}\succeq p_{0,M}(D_n+D_{M-n})
              \succeq p_{0,M}D_n,\qquad
 C_0^{(M)}-C_n^{(M)}\succeq0 .
 \tag{V14.7}\label{r14:upper}
\]
For even $M=2k$, and $Y=\sum_i c_iY_i$,
\[
 c^{\mathsf T}(C_0^{(M)}-C_k^{(M)})c
       =\frac2{Z_M}\|[T^k,\mathcal L_Y]\|_{\rm HS}^2 .
 \tag{V14.8}\label{r14:comm}
\]
\end{theorem}
\begin{proof}
The insertion product around the periodic history is
$\Tr(B_{Y_i}^-T^n B_{Y_j}^+T^{M-n-2})/Z_M$.
Substitute \eqref{r14:B-formula} and move the outside
powers of $T$ around the trace.  This gives \eqref{r14:trace}.
More intrinsically its right side is the real Gram
\[
 \frac4{Z_M}
 \left\langle T^{n/2}\mathcal L_{Y_i}T^{(M-n)/2},
              T^{n/2}\mathcal L_{Y_j}T^{(M-n)/2}
 \right\rangle_{\rm HS}.
\]
Because $(M-n)/2\ge1$, the differentiated right factor
is smoothing and Hilbert--Schmidt; the left factor is bounded.
This also justifies the formula at $n=0$.  Integration
by parts and finite spectral truncations justify the cyclic
notation without treating $\mathcal L_Y$ as bounded.

Choose a real orthonormal eigenbasis $T\phi_a=\lambda_a\phi_a$,
with $\phi_0=\psi$, and set $r_a=\lambda_a/\Lambda$.
Let $v_{ab}$ be the bank vector whose $i$th entry is
$\langle\phi_a,\mathcal L_{Y_i}\phi_b\rangle$.
Then $v_{ba}=-v_{ab}$ and $v_{aa}=0$, and the convergent
positive matrix sum is
\[
 C_n^{(M)}
   =4p_{0,M}\sum_{a,b}r_a^{M-n}r_b^n
                              v_{ab}v_{ab}^{\mathsf T}.
 \tag{V14.9}\label{r14:spectral}
\]
The terms with $a=0$ give $p_{0,M}D_n$ and those
with $b=0$ give $p_{0,M}D_{M-n}$; their intersection
vanishes.  Dropping all other positive matrices proves
the first two inequalities.  Combining the ordered
pair $(a,b)$ with $(b,a)$, its coefficient in the
difference from $n=0$ is proportional to
\[
 (r_a^n-r_b^n)(r_a^{M-n}-r_b^{M-n})\ge0 .
\]
This proves the last inequality.  At $M=2k$ the
coefficient is a square, and expanding the
Hilbert--Schmidt commutator norm gives \eqref{r14:comm}.
\end{proof}

Thus a \emph{finite periodic} integral supplies an upper,
not merely lower, enclosure of a vacuum response once its
vacuum probability is paid.  This is useful for a mass-gap
attack.  It does not make the periodic and vacuum
normalizations interchangeable: $C_0^{(M)}$ is an
upper bound for $p_{0,M}J$, not a lower bound for $J$.

\subsection{A sharp high-purity branch and a signed-history boundary bound}
\label{r14:boundary}

The following refinement of the inherited purity bracket
will be used with certified lower bounds, not fitted traces.
Suppose $1/2<\underline R_q\le R_q$ and set
\[
 \ell_q=\frac{1+\sqrt{2\underline R_q-1}}2,\qquad
 \varepsilon_q=2(1-\ell_q)=1-\sqrt{2\underline R_q-1},
 \qquad t_q=\frac{1-\ell_q}{\ell_q}<1 .
 \tag{V14.10}\label{r14:ell}
\]
These numbers do not depend on a choice of scalar
normalization of $T$.

\begin{lemma}[Refined vacuum weight]
\label{r14:purity}
One has $p_{0,q}\ge\ell_q$ and
$u_q:=\sum_{a>0}r_a^q\le t_q$.
\end{lemma}
\begin{proof}
For the thermal probabilities $w_a=\lambda_a^q/Z_q$,
the inherited bound gives $w_0\ge R_q>1/2$.
Also $\sum_{a>0}w_a^2\le(1-w_0)^2$, so
$\underline R_q\le R_q\le w_0^2+(1-w_0)^2$.
On the branch $w_0>1/2$, solving this quadratic yields
$w_0\ge\ell_q$.  Finally $u_q=w_0^{-1}-1$.
The bound is sharp when there is just one excited
thermal probability.  No gap estimate entered its proof.
\end{proof}

\begin{theorem}[A gap-free certified boundary error for a finite history]
\label{r14:history}
Let $F$ be a real gauge-invariant observable of a segment
of $r\ge1$ slabs, including any of its temporal links,
and let $|F|\le B$.  Let $\E_\infty F$ denote its
stationary vacuum expectation, and let $\E_{q+r}F$
be the same insertion in a periodic history with
$q\ge4$ additional, uninserted slabs.  Under
\eqref{r14:ell},
\[
             |\E_{q+r}F-\E_\infty F|
                          \le B\varepsilon_q .
 \tag{V14.11}\label{r14:history-bound}
\]
The period appearing in $R_q$ is the \emph{uninserted
padding length} $q$, not $q+r$.
\end{theorem}
\begin{proof}
Let $T_F$ be the segment kernel with insertion $F$.
Its literal path integral satisfies the pointwise
kernel bound $|T_F(x,y)|\le B T^r(x,y)$.
Consequently $\|\Lambda^{-r}T_F\|\le B$, by
pointwise domination followed by the $L^2$ norm.
This is an operator-norm assertion, not the generally
false Loewner inequality $-BT^r\preceq T_F\preceq BT^r$.
Put $a_b=\langle\phi_b,\Lambda^{-r}T_F\phi_b\rangle$.
Then $|a_b|\le B$, $a_0=\E_\infty F$, and
\[
 \E_{q+r}F
    =\frac{a_0+\sum_{b>0}r_b^q a_b}{1+u_{q+r}} .
\]
The trace exists because $T^q$ is trace class and $T_F$
is bounded.  It agrees with the periodic path integral.
Writing $u=u_q$ and $v=u_{q+r}$ gives
\[
 |\E_{q+r}F-a_0|
       \le B\frac{u+v}{1+v}
       \le B\frac{2u}{1+u}
       =2B(1-p_{0,q})\le B\varepsilon_q .
\]
Here $0\le v\le u\le t_q<1$; on this interval
$(u+v)/(1+v)$ increases with $v$.
This is why the stated high-purity branch matters.
\end{proof}

The estimate is finite-volume but quantitative and a posteriori:
it does not use an unknown rate of convergence to the vacuum.
For fixed regulator $R_q\to1$, so a sufficient padding
length exists.  That qualitative existence alone neither
computes its value nor bounds it along changing regulators.

\subsection{Certified intervals and a contraction test for the score bank}
\label{r14:matrix-test}

Let $\mathbf A^\pm$ be the bank of one-sided insertions.
Choose a known positive-definite matrix $H$ such that,
pointwise on every single slab,
\[
       \mathbf A^\pm(\mathbf A^\pm)^{\mathsf T}\preceq H .
 \tag{V14.12}\label{r14:envelope}
\]
For example, if $|A_i^\pm|\le b_i$ for $d$ sources,
$H=d\,\operatorname{diag}(b_i^2)$ suffices (zero sources
can be omitted).  Sharper polynomial matrix envelopes
may be substituted with their proofs.

\begin{corollary}[Stationary score matrix enclosure]
\label{r14:interval}
If $M\ge n+6$ and \eqref{r14:ell} is certified for
$q=M-n-2$, then
\[
 -\varepsilon_{M-n-2}H
       \preceq C_n^{(M)}-D_n
       \preceq\varepsilon_{M-n-2}H .
 \tag{V14.13}\label{r14:interval-eq}
\]
In particular this encloses $J=D_0$ using padding
$M-2$.  If a lower vacuum weight $\ell_M$ is also
certified, the sharper one-sided bound
$D_n\preceq C_n^{(M)}/\ell_M$ is available.
\end{corollary}
\begin{proof}
For every real $c$, use
$F=(c^{\mathsf T}\mathbf A^-(0))
       (c^{\mathsf T}\mathbf A^+(n))$.
Its segment has $r=n+2$ slabs and
$|F|\le c^{\mathsf T}Hc$.
Its stationary expectation is $c^{\mathsf T}D_nc$
by V13 and its periodic expectation is
$c^{\mathsf T}C_n^{(M)}c$.
Apply \eqref{r14:history-bound} for every $c$.
The one-sided result follows from \eqref{r14:upper}.
\end{proof}

For V13's tree field a completely explicit, though coarse,
one-source envelope is
\[
 H_m=3N\kappa_m(26p_m^2+18\gamma p_m)^2,\qquad
       \kappa_m=(2m)^3,\quad p_m=12m-8 .
 \tag{V14.14}\label{r14:tree-envelope}
\]
Indeed $|\operatorname{div}Y_m|
\le\sqrt{3N}\|\operatorname{Hess}F_m\|_{\rm HS}$,
$|\nabla S_s|\le4\sqrt N$ and $|g|\le\sqrt N$.
Combine V13's
$|Y_m|\le\sqrt{27\kappa_m}\,p_m$ and Hessian bound.
This is a pointwise history amplitude, not the
volume-independent marginal-score second moment from V13.
Its explicit factor $N$ is not suppressed.

Suppose finite periodic polynomial integrals have produced
symmetric matrices $\widehat C_0,\widehat C_n$ with
certified errors
\[
 -\delta_jH\preceq\widehat C_j-C_j^{(M)}
                  \preceq\delta_jH,\qquad j=0,n .
\]
Put
\[
 J_-=\widehat C_0-(\delta_0+\varepsilon_{M-2})H,\qquad
 U_n=\ell_M^{-1}(\widehat C_n+\delta_nH).
 \tag{V14.15}\label{r14:JU}
\]
Then $J_-\preceq J$ and $D_n\preceq U_n$.
Consequently the verifiable conditions
\[
          J_-\succ0,\qquad U_n\preceq\theta J_-,
                      \qquad 0<\theta<1
 \tag{V14.16}\label{r14:test}
\]
imply $D_n\preceq\theta J$ on this bank.  No estimate
of an inverse unknown Gram is required.
For even $n=2t$ this is
$\|P^t s\|_p^2\le\theta\|s\|_p^2$ for all bank states $s$.

The finite-data inputs themselves can be enclosed by
the fixed-history polynomial hierarchy of V3.  In
particular if $0<Z_q^-\le Z_q\le Z_q^+$ and
$0<Z_{2q}^-\le Z_{2q}$, one may use
$\underline R_q=Z_{2q}^-/(Z_q^+)^2$ when it exceeds
$1/2$.  For the unshifted deficit action on a periodic
$q$-slab lattice, $0\le S\le4Nq$, so
$Z_q\ge e^{-4\gamma Nq}>0$ is an explicit denominator
bound.  Taylor approximation of $e^{-\gamma S}$ on
this compact interval, followed by exact polynomial
Haar moments, supplies convergent interval bounds.
Its cost is not asserted to be reasonable or uniform.
If computed lower ratios exceed one, the input
enclosures are inconsistent and must be rejected.
All periods use the same spatial lattice, coupling
and Haar normalization; no selected coupling is
silently changed on doubling the time period.

This closes a \emph{certificate interface}, not its
uniform feasibility.  No Wilson matrix satisfying
\eqref{r14:test} along the required continuum sequence
has been certified in this iteration.

\subsection{Two periods recover a lower Gram without a pointwise envelope}
\label{r14:two-periods}

The preceding pointwise bound is useful for general histories,
but its $H_m$ can be much larger than the score Gram.  The
positive spectral representation of $C_0$ permits a better
bound specific to these insertions.

\begin{theorem}[Amplitude-free two-period lower Gram]
\label{r14:two-period-thm}
Fix $q\ge4$ and a certified purity bound as in
\eqref{r14:ell}.  Then
\[
 J\succeq
 J_{\rm two}:=
 \frac{C_0^{(2q)}-(t_q/\ell_q)C_0^{(q)}}{1-t_q}.
 \tag{V14.17}\label{r14:two}
\]
This is a Loewner lower bound even if its right side
is not positive.  It uses two ordinary periodic
score integrals and $R_q$; it does not contain a
pointwise amplitude or an assumed spectral gap.
For $0\le n\le2q-2$ one also has
\[
 D_n\preceq\ell_q^{-1}C_n^{(2q)} .
 \tag{V14.18}\label{r14:two-upper}
\]
\end{theorem}
\begin{proof}
For each real eigenvector define the positive matrix
\[
 W_a(i,j)=4\langle\mathcal L_{Y_i}\phi_a,
                         \mathcal L_{Y_j}\phi_a\rangle .
\]
The $n=0$ case of \eqref{r14:spectral} gives
\[
 N_s:=\frac{C_0^{(s)}}{p_{0,s}}
       =J+E_s,\qquad
 E_s=\sum_{a>0}r_a^s W_a\succeq0,\qquad W_0=J .
\]
These sums converge by the Hilbert--Schmidt proof above.
Lemma~\ref{r14:purity} gives $r_a^q\le u_q\le t_q$
for every excited eigenvalue.  Multiplying each
positive $W_a$ by the corresponding scalar proves
$E_{2q}\preceq t_q E_q$.  Hence
\[
 (1-t_q)J\succeq N_{2q}-t_qN_q
       \succeq C_0^{(2q)}-(t_q/\ell_q)C_0^{(q)}.
\]
In the second step use $p_{0,2q}\le1$,
$p_{0,q}\ge\ell_q$, and positivity of both periodic
matrices.  Dividing by $1-t_q>0$ proves
\eqref{r14:two}.  Finally $u_{2q}\le u_q$ implies
$p_{0,2q}\ge p_{0,q}\ge\ell_q$.
Combine this with \eqref{r14:upper}.
\end{proof}

All ingredients remain finite polynomial insertion integrals.
For clarity, suppose symmetric certified enclosures are
\[
 L_{2q}\preceq C_0^{(2q)},\qquad
 C_0^{(q)}\preceq U_q,\qquad C_n^{(2q)}\preceq U_{n,2q}.
\]
Then define
\[
 \widehat J_{\rm two}
      =\frac{L_{2q}-(t_q/\ell_q)U_q}{1-t_q}.
 \tag{V14.19}\label{r14:two-certified}
\]
The entirely finite conditions
\[
 \widehat J_{\rm two}\succ0,\qquad
       U_{n,2q}\preceq\theta\ell_q\widehat J_{\rm two},
                       \quad 0<\theta<1
 \tag{V14.20}\label{r14:two-test}
\]
imply $D_n\preceq\theta J$.  Certified errors in
$L_{2q},U_q,U_{n,2q}$ must still be paid; they are
not thermal boundary errors.  The subtraction in
\eqref{r14:two-certified} can require small
\emph{relative} integration errors when $J$ is small.
We do not claim monotonicity of this lower bound
as $q$ changes.  At a fixed regulator it does
converge to $J$ with exact data: $C_0^{(q)}\to J$,
$\ell_q\to1$ and $t_q\to0$ if the purity enclosure
converges to the exact ratio.  The first limit
follows by dominated convergence in $E_q$, dominated
by the summable $E_4$.  No rate uniform in the
regulator follows from that argument.

\subsection{A misleading thermal contraction, even at high vacuum weight}
\label{r14:false-test}

The lower enclosure for $J$ cannot be replaced by
$C_0^{(M)}$.  Here is an exact finite-state countercheck
of that possible inference; it is not a Wilson model.
In a real eigenbasis choose
\[
 K=\operatorname{diag}(1,9/10,1/5,1/5),\qquad
 \mathcal L_{01}=10^{-5},\quad \mathcal L_{12}=1,
 \quad \mathcal L_{ba}=-\mathcal L_{ab},
 \tag{V14.21}\label{r14:counter}
\]
with all other entries zero.  Orthogonal conjugation
by the $4$ by $4$ Hadamard matrix divided by $2$
makes $K$ a strictly entrywise-positive symmetric
Markov matrix with constant vacuum.  It leaves all
the spectral formulas unchanged.

For this one field, at $M=100$ and $n=2$,
\[
 p_{0,100}>\frac{9999}{10000},\qquad
 \frac{C_2^{(100)}}{C_0^{(100)}}<\frac1{20},
 \qquad \frac{D_2}{J}=\frac{81}{100}.
 \tag{V14.22}\label{r14:false}
\]
These are rational inequalities.  Explicitly, with
$r=9/10$, $s=1/5$ and $e=10^{-5}$,
\[
 \frac{C_2^{(100)}}{C_0^{(100)}}=
 \frac{e^2(r^2+r^{98})+r^{98}s^2+s^{98}r^2}
      {e^2(1+r^{100})+r^{100}+s^{100}},\qquad
 p_{0,100}=(1+r^{100}+2s^{100})^{-1}.
\]
The small vacuum coupling $e$ makes the thermally
excited $1\leftrightarrow2$ transition dominate this
source's periodic normalization, despite the tiny
total excited probability.  Every quantity is from
the same transfer and insertion, not an added hidden
sector.  General bounded antisymmetric matrices
suffice to disprove the proposed inference from
periodic matrix ratios alone; they are not asserted
to be continuum vector-field generators.
The amplitude-sensitive interval
\eqref{r14:interval-eq} retains precisely the cost
which that false inference omits.
Alternatively \eqref{r14:two-certified} controls
the same contamination using a second period.
The exact checker verifies that this corrected
test, with $q=120$, $n=2$ and exact moment data,
certifies $D_2\le(9/10)J$ in this example, consistent
with the actual $81/100$ ratio.  It does not
certify the false $1/20$ vacuum ratio.

\subsection{What would control the complete spectral edge?}
\label{r14:edge}

There are two distinct possible exports of the new
finite-data bounds.  First, \eqref{r14:two-test} (or
\eqref{r14:test}) would
bound the physical transfer on the complete centered
slice space if it held with the same $n,\theta$ on
\emph{every} finite bank of polynomial scores.
V13's density then extends the quadratic inequality
by continuity, and positivity of $P$ gives
$\|P|_{1^\perp}\|\le\theta^{1/n}$.
For one finite bank alone, its image under $P^n$
need not remain in that bank, so iterating its
matrix inequality is not justified.  Uniform
limiting source normalization and a constructed
continuum Hilbert space are additional obligations.

Second, the purity information alone has an
all-sector, finite-regulator consequence:
\[
 \sup_{a>0}r_a\le t_q^{1/q},\qquad
 \operatorname{gap}(-\log K)
      \ge\frac1q\log\frac{\ell_q}{1-\ell_q}.
 \tag{V14.23}\label{r14:gap}
\]
Indeed $r_a^q\le u_q\le t_q$ for every excited
eigenvalue.  This is an elementary extension of
the inherited purity bracket, not an evaluated
new Wilson mass estimate.  At a lattice time unit
$a_{\rm t}$ the right side must be divided by
$a_{\rm t}$.  A uniform physical conclusion needs
a lower bound on this whole quotient, not merely
$\underline R_q>1/2$ at each regulator.

The volume cost can be checked without conjectures.
For a tensor product of $V$ identical two-state
positive Markov transfers with eigenvalues $1,r$,
$0<r<1$, the actual gap $\Delta=-\log r$ is
independent of $V$, while
\[
 R_q=\left(\frac{1+r^{2q}}{(1+r^q)^2}\right)^V .
 \tag{V14.24}\label{r14:tensor}
\]
At fixed $q$ this tends to zero as $V\to\infty$.
Even a gapped system therefore need not pass a
fixed-period high-purity test at growing volume.
If instead $q=\lceil A\log V/\Delta\rceil$ with
$A>1$, write $z=r^q$.  Then
$\log z=-A\log V+O(1)$ and
$u_q=(1+z)^V-1=Vz(1+o(1))$.  Expansion of the
exact expression \eqref{r14:tensor} gives
$1-R_q=2Vz(1+o(1))$ and, using the exact $R_q$
in \eqref{r14:ell}, $t_q=Vz(1+o(1))$.
Thus the bound in \eqref{r14:gap} tends to
$(1-1/A)\Delta$.  This calculation proves the
time--volume bookkeeping in a specified product
model, not a Wilson factorization.

The next upstream task is now concrete: obtain
certified bounds on the periodic score matrices
and their doubled-period normalization, with
errors small relative to the actual score Gram,
or estimate the complete excitation trace
$u_q$ at a growing physical padding time.
Sector-restricted numerical decay cannot replace
either all-sector requirement.  The old
volume-dependent kernel floor and intensive
pressure comparison do not supply these bounds.

\paragraph{Verification and preservation.}
The checker \path{scripts/verify_ym_restart_v14_periodic_scores.py}
tests rational skew-derivative trace identities, matrix
positivity and vacuum domination, signed-history
boundary errors, the two-period lower Gram and
contraction test, the misleading thermal-ratio example,
purity branches, and exact tensor-product traces.
These fixtures are algebra checks, not Wilson
quadratures or substitutes for the analytic proofs.
V13's body is preserved except for the displayed title;
its recovered SHA--256 is
\path{ABAE72F7BF1EFFF8DD640135233C7D70BB97AAABE391149068AFF6C6055436C6}.
Only active V13 is replaced.  All companion and archive
sources remain protected; builds stay outside
\path{latex_notes}.

\paragraph{Status.}
Explicit finite-period-to-vacuum enclosures, including
an amplitude-free two-period lower Gram, and a score-bank
upper-contraction certificate are now proved.
Their required uniform margins, all-sector continuum
control, the interacting four-dimensional construction
and its Yang--Mills mass gap remain unproved.

% END CONSOLIDATED V14 APPEND

% BEGIN CONSOLIDATED V15 APPEND -- 2026-09-19
\clearpage
\section{V15: bounded temporal sewing and native purity defects}
\label{r15:main}

V14 reduced the finite-period score certificate to actual periodic
Wilson integrals, including the all-sector purity $R_p=Z_{2p}/Z_p^2$.
The present step attacks the latter input.  It gives an exactly
bounded observable under the original periodic law, an error estimate
which ignores arbitrary endpoint normalizations, and a native
conditional force-covariance representation of its defect.  It also
proves that a direct spatial partial-swap upper bound has a nonzero
vacuum-entanglement floor at fixed positive coupling.  This directs
the next estimates toward temporal conditional cancellation, not a
sum of spatial entanglement defects.  No uniform smallness estimate
for the resulting Wilson expectation is asserted here.

\subsection{Companion checkpoint, literature import, and scope}
\label{r15:audit}

The scheduled five-note checkpoint reread the terminal arguments of
the following sources, without changing any of them.
\begin{itemize}
\item \path{Renewal_NS_Stability_Research_Notes_V26.tex}:
rank-one Oseen derivative-kernel estimates concern a different
evolution and do not estimate the Wilson periodic law.  Numerical
claims awaiting later certification are not imported.
\item \path{Renewal_SM_Einstein_Continuum_Research_Notes_V72.tex}:
the many-copy covariance transport holds with the physical mode
cutoff fixed.  It supplies no uniform Wilson trace-purity bound.
\item \path{Renewal_Yang_Mills_Shell_Mixing_Research_Notes_V16.tex}:
the measured flat one-loop coefficient retains its root, tangent
and Haar returns, but does not control the moving-coupling,
growing-volume remainder or the missing mixed curvature coefficient.
\item \path{Renewal_YM_Parallel_Research_Notes_V5.tex}:
its margin-protection programme assumes the relevant all-colour
cofinal convergence and contraction inputs.  These are not native
periodic estimates available for import.  Its independent-component
charge warning reinforces keeping volume costs explicit; it does
not prove connected-torus charge anti-concentration.
\end{itemize}
The first three files are byte-identical to the V10 checkpoint.
The parallel note was reviewed additionally.  The next companion
checkpoint and regular broad primary-literature sweep are both V20.

A targeted primary-source check read Hastings, Gonz\'alez, Kallin
and Melko, \href{https://arxiv.org/abs/1001.2335}
{\emph{Measuring Renyi Entanglement Entropy with Quantum Monte Carlo}}
(2010), especially the swap-purity identity (3)--(5) and the
changed-weight ratio construction (9)--(11).  The swap identity is
standard and is not claimed as a discovery of this programme.
Their Heisenberg-model numerical performance supplies neither a
Wilson sampling theorem nor a four-dimensional mass gap.  We use
the identity as motivation and prove the particular native
sewing and error formulas below directly.

The nonduplication check also reread f15 V58's thermal purity ledger,
f15 V81's Gauss-projection obstruction to partial spatial swaps,
f14 V82's auxiliary-chart thermodynamic integration, and the
Perron-visibility bounds in f16.2 V38.  None is replaced by the
claims below.  In particular we do not identify a compressed
partial swap with a physical orthogonal projection, nor identify
a spatial reduced purity with the full temporal purity.

Throughout this appendix $L\ge4$, $N=3L^3$, $\gamma>0$,
$\mathcal X=\SU(2)^N$, and $T$ is exactly \eqref{r14:T}.
All Gauss integrations are retained.  Extending $T$ by zero on
the orthogonal complement of the physical subspace gives its
positive integral operator on full product Haar space; traces
of its positive powers still equal the physical traces.
The notation $Z_p$ always uses this same coupling, spatial
volume and transfer normalization when the temporal period changes.

\subsection{Balanced temporal sewing under an unchanged law}
\label{r15:sewing}

Fix an integer $q\ge2$, and write $\kappa(x,y)=T^q(x,y)$.
This is a smooth, strictly positive, symmetric kernel.  In a
periodic history of length $2q$, the pair of antipodal slices
$(x,y)$ has density $\kappa(x,y)^2/Z_{2q}$ with respect to
$dx\,dy$.  This follows by integrating each of the two
$q$-slab segments, including every temporal link.  For two
independent such histories define
\[
\begin{split}
 a&=\kappa(x_1,y_1)\kappa(x_2,y_2),\qquad
 b=\kappa(x_1,y_2)\kappa(x_2,y_1),\\
 d\mu_q&=\frac{a^2}{Z_{2q}^2}\,dx_1\,dy_1\,dx_2\,dy_2,
 \qquad W=b/a,\quad v=\log W,\\
 A&=\frac{2ab}{a^2+b^2}=\frac{1}{\cosh v},\qquad
 d=1-A=\frac{(a-b)^2}{a^2+b^2}.
\end{split}
\tag{V15.1}\label{r15:def}
\]
Here $d$ is a scalar observable, not a measure differential.
Let $S$ exchange the \emph{whole} slices $y_1,y_2$.
It is a Haar-measure-preserving involution, $a\circ S=b$,
and $A\circ S=A$.  It need not preserve $\mu_q$.

\begin{theorem}[Exact bounded sewing observable]
\label{r15:bounded}
Under the unchanged native law $\mu_q$,
\[
\begin{gathered}
 \E_{\mu_q}W=R_{2q},\qquad \E_{\mu_q}W^2=1,\\
 0<A\le1,\quad 0\le d<1,\qquad
 \E_{\mu_q}A=R_{2q},\quad \E_{\mu_q}d=1-R_{2q},\\
 \Var_{\mu_q}(A)=\Var_{\mu_q}(d)
       \le R_{2q}(1-R_{2q}).
\end{gathered}
\tag{V15.2}\label{r15:exact}
\]
Also, with $d\omega=dx_1\,dy_1\,dx_2\,dy_2$,
\[
 \int(a-b)^2\,d\omega=2(Z_{2q}^2-Z_{4q}),\qquad
 \E_{\mu_q}(W-1)^2=2(1-R_{2q}).
\tag{V15.3}\label{r15:minor}
\]
\end{theorem}
\begin{proof}
Composition of the symmetric kernel around the four-cycle gives
$\int ab\,d\omega=\Tr(T^{4q})=Z_{4q}$, while
$\int a^2\,d\omega=\int b^2\,d\omega=Z_{2q}^2$.
These equalities prove both raw moments and \eqref{r15:minor}.
Although $\mu_q$ is not $S$-invariant, Haar change of variables
and $A\circ S=A$ give
\[
 \int a^2 A\,d\omega
 =\frac12\int(a^2+b^2)A\,d\omega
 =\int ab\,d\omega .
\]
Thus the bounded observable has exactly the same expectation.
Finally $A^2\le A$ gives
$\Var A\le\E A-(\E A)^2$.
\end{proof}

The replacement of $W$ by $A$ is an exact orbit symmetrization,
not truncation of a rare-event tail.  It improves the elementary
variance bound from $1-R_{2q}^2$ to $R_{2q}(1-R_{2q})$.
It does not prove rapid mixing of a Wilson sampler or independent
sampling of histories.  Its four factors are integrated block
kernels, not four selected microscopic path weights; replacing
one by an unintegrated weight changes the observable and invalidates
the stated expectation unless a separate identity is proved.

\subsection{Certified block-kernel errors modulo endpoint factors}
\label{r15:errors}

Let $\widetilde\kappa>0$ be any approximation.  It need not be
symmetric.  Define $\widetilde v,\widetilde A,\widetilde d$
by the same four endpoint formulas, but continue to integrate
against the \emph{actual} $\mu_q$, not the law built from
$\widetilde\kappa$.

\begin{theorem}[Square-root stability of the purity defect]
\label{r15:stability}
Put $\varepsilon_2=\|\widetilde v-v\|_{L^2(\mu_q)}$.
Then
\[
 \left|\sqrt{\E_{\mu_q}\widetilde d}
             -\sqrt{1-R_{2q}}\right|
       \le\frac{\varepsilon_2}{\sqrt2}.
\tag{V15.4}\label{r15:sqrt}
\]
In particular, if for some functions $\alpha,\beta$,
\[
 \log\widetilde\kappa(x,y)-\log\kappa(x,y)
       =\alpha(x)+\beta(y)+e(x,y),\qquad |e(x,y)|\le\eta,
\tag{V15.5}\label{r15:gauge-error}
\]
then $\varepsilon_2\le4\eta$.  A certified upper bound
$\E_{\mu_q}\widetilde d\le b_*$ consequently implies
\[
 R_{2q}\ge 1-\bigl(\sqrt{b_*}+2\sqrt2\,\eta\bigr)^2.
\tag{V15.6}\label{r15:lower}
\]
More generally $2\sqrt2\eta$ is replaced by any certified
$L^2$ rectangle-error bound $\varepsilon_2/\sqrt2$.
\end{theorem}
\begin{proof}
For $h(v)=\sqrt{1-1/\cosh v}$ and $v\ne0$, set
$s=1/\cosh v$.  Direct differentiation gives
\[
 |h'(v)|^2=\frac{s^2(1+s)}4\le\frac12,
 \qquad 2-s^2-s^3=(1-s)(s^2+2s+2)\ge0.
\]
At zero the one-sided derivatives are $\pm1/\sqrt2$.
Hence $h$ is globally $1/\sqrt2$-Lipschitz.  The reverse
triangle inequality in $L^2(\mu_q)$ proves \eqref{r15:sqrt}.
In the logarithmic rectangle, every $\alpha(x_i)$ and
$\beta(y_j)$ occurs once with each sign.  Their cancellation
leaves four errors of magnitude at most $\eta$.
The upper triangle inequality followed by squaring proves
\eqref{r15:lower}.
\end{proof}

Thus no separate accuracy is required for endpoint amplitudes,
including arbitrary positive row and column multipliers.
This is not permission to alter the periodic sampling law.
If $b_*$ is very small, the kernel error produces a quadratic
$8\eta^2$ floor plus its displayed cross term, rather than an
unspecified relative error on an exponentially small partition
function.  There is still no small uniform value of $b_*$ here.

\paragraph{Finite-regulator enclosure is possible, not efficient.}
For a literal $q$-slab segment with fixed endpoints, its action
$S^{(q)}$ includes half the spatial action at each endpoint,
full spatial actions at the $q-1$ internal slices and all temporal
plaquettes.  Therefore $0\le S^{(q)}\le4Nq$ and
\[
 e^{-4\gamma Nq}\le\kappa(x,y)\le1 .
\tag{V15.7}\label{r15:floor}
\]
For a computable coupling (in particular a positive rational
one), integrate the degree-$D$ Taylor polynomial of
$e^{-\gamma S^{(q)}}$ against the internal probability Haar
measures.  The result $\kappa_D$ is a polynomial in the endpoint
quaternion coordinates, with certified coefficient enclosures
if necessary.  The integral remainder gives, uniformly,
\[
 |\kappa_D-\kappa|\le
       \frac{(4\gamma Nq)^{D+1}}{(D+1)!}=:\delta_D.
\tag{V15.8}\label{r15:taylor}
\]
For rational $\gamma$ the polynomial Haar moments are exact
rationals.  Choose $D$ large enough that
$\delta_D\le\zeta e^{-4\gamma Nq}$, $0<\zeta<1$.
Then $\kappa_D>0$ and
$|\log\kappa_D-\log\kappa|\le-\log(1-\zeta)$.
The denominator of $\widetilde d$ is a positive polynomial
bounded away from zero on the compact endpoint space.  A
uniform polynomial approximation to its reciprocal, followed
by the same Taylor/Haar enclosures for the two periodic
weights and their normalization, bounds
$\E_{\mu_q}\widetilde d$ to arbitrary fixed precision.
For example a geometric series for the reciprocal converges
uniformly after scaling by any certified upper denominator bound.

This is an application of V3's inherited finite polynomial
enclosure method to a different, bounded observable, not a
new claim of practical complexity.  Its crude positivity floor
is exponentially small in $\gamma Nq$; no volume-uniform
quadrature cost or weak-coupling estimate follows from it.

\subsection{The exact native quantity that must decay}
\label{r15:forces}

Let $\nu_{q;x,y}$ be the normalized law of the $q$-slab
interior conditioned on its endpoints $x,y$.  Endpoint spatial
half-actions are constants in this conditional law.  Let
$g_L$ and $g_R$ be the gradients of the first and last temporal
plaquette sums with respect to $x$ and $y$, respectively,
using the unit round $S^3$ metric on each link.  Each endpoint
link occurs in one such temporal plaquette, so
$|g_L|,|g_R|\le\sqrt N$ pointwise.

\begin{proposition}[Mixed endpoint derivative and bridge covariance]
\label{r15:mixed}
For endpoint tangent vectors $\xi,\zeta$ and $q\ge2$,
\[
 D_xD_y\log\kappa(x,y)[\xi,\zeta]
 =\gamma^2\Cov_{\nu_{q;x,y}}
                  (g_L\cdot\xi,g_R\cdot\zeta).
\tag{V15.9}\label{r15:covariance}
\]
Choose constant-speed minimizing product geodesics $X_s$ from
$x_1$ to $x_2$ and $Y_t$ from $y_1$ to $y_2$.  Then, writing
\[
 \mathcal C_q=\int_0^1\!\int_0^1
 \Cov_{\nu_{q;X_s,Y_t}}
       (g_L\cdot\dot X_s,g_R\cdot\dot Y_t)\,ds\,dt,
\tag{V15.10}\label{r15:C}
\]
we have $v=-\gamma^2\mathcal C_q$ and
\[
 1-R_{2q}\le\frac{\gamma^4}{2}
                   \E_{\mu_q}\mathcal C_q^2.
\tag{V15.11}\label{r15:force-bound}
\]
\end{proposition}
\begin{proof}
At fixed endpoints, differentiating a normalized finite compact
integral gives
\[
 D_xD_y\log\int e^{-\gamma S}
 =-\gamma\E(D_xD_yS)+\gamma^2\Cov(D_xS,D_yS).
\]
For $q\ge2$ no individual action term depends on both endpoints,
even when the first and last slabs share a middle slice.
Thus the first term vanishes.  The spatial half-action derivatives
are constants under the bridge law and drop out of the covariance,
leaving \eqref{r15:covariance}.  All derivatives may pass under
the integral by smoothness and compactness at fixed regulator.

The fundamental theorem of calculus on the endpoint rectangle
gives
\[
 \log\kappa(x_2,y_2)-\log\kappa(x_2,y_1)
 -\log\kappa(x_1,y_2)+\log\kappa(x_1,y_1)
       =\gamma^2\mathcal C_q=-v.
\]
The Lipschitz estimate for $h$ in the preceding proof, with
$h(0)=0$, gives $1-1/\cosh v\le v^2/2$.
Integrate and use \eqref{r15:exact}.  Geodesics can be selected
measurably; antipodal single-link endpoints form a Haar-null set,
and the positive smooth law is absolutely continuous.
\end{proof}

This target involves the native conditional interior law, not a
Gaussian reference, guessed vacuum, or sampling-time semigroup.
The elementary covariance estimate is only
\[
 |\Cov(g_L\cdot\xi,g_R\cdot\zeta)|\le N|\xi||\zeta|,
 \qquad |\mathcal C_q|\le\pi^2N^2,
\tag{V15.12}\label{r15:coarse}
\]
because each product geodesic has speed at most $\pi\sqrt N$.
This has no decay in $q$ and becomes worse with volume.  In
particular it does not close V14.  The meaningful new primitive
is a weighted mean-square, long-bridge endpoint-force covariance
or the bounded defect itself, with the conditional normalization
and the whole endpoint rectangle retained.  A typical-point
derivative bound at a flat field would not control this integral.

\subsection{Why spatial partial-swap defects have a vacuum floor}
\label{r15:spatial}

All tensor products in this subsection refer to the full
unreduced Haar space; no factorization of the physical Gauss
subspace is assumed.  Set
$h_q(x,y)=\kappa(x,y)/\sqrt{Z_{2q}}$ and
$\Phi_q=h_q\otimes h_q$.  For each spatial link $e$ let $S_e$
swap just the $y_e$ coordinates between the two copies.
These are commuting self-adjoint unitaries on the full space.
Their product is $S$, and
\[
 \langle\Phi_q,S\Phi_q\rangle=R_{2q},\qquad
 \langle\Phi_q,S_e\Phi_q\rangle=\Tr(\rho_{q,e}^2),
 \quad \rho_q=\frac{T^{2q}}{Z_{2q}}.
\tag{V15.13}\label{r15:spatial-purity}
\]
Here $\rho_{q,e}$ is the partial trace of $\rho_q$ over all
links except $e$.  To verify the second identity, integrate
the $x$ variable in $|h_q\rangle\langle h_q|$, obtaining
$T^{2q}/Z_{2q}$, and expand the two-copy swap in a product
orthonormal basis.  The resulting contraction is precisely
the squared reduced-density trace.  Thus
\[
 1-R_{2q}\le\sum_{e=1}^N[1-\Tr(\rho_{q,e}^2)].
\tag{V15.14}\label{r15:union}
\]
Indeed on every simultaneous eigenspace $S_e=\sigma_e\in\{\pm1\}$,
$1-\prod_e\sigma_e\le\sum_e(1-\sigma_e)$.
This argument does not use the false claim that the $S_e$
preserve the physical subspace.

\begin{theorem}[Strict native vacuum floor for the spatial sum]
\label{r15:floor-theorem}
At fixed $L\ge4$ and $\gamma>0$, let $\psi$ be the normalized
positive Wilson vacuum and $\rho_{\psi,e}$ its one-link
reduced density operator on full Haar space.  Then
\[
 \lim_{q\to\infty}(1-R_{2q})=0,\qquad
 \lim_{q\to\infty}\sum_e[1-\Tr(\rho_{q,e}^2)]
     =\sum_e[1-\Tr(\rho_{\psi,e}^2)]>0.
\tag{V15.15}\label{r15:strict-floor}
\]
No regulator-uniform positive lower bound for the last sum is
claimed.
\end{theorem}
\begin{proof}
First we show that $\psi$ is not constant.  Up to a positive
scalar the transfer is $M_b C M_b$, where
$b=e^{-\gamma S_s/2}$ and $C$ is normalized product Wilson
convolution followed by the Gauss projection.  The convolution
commutes with this projection.  Its nonnegative character
coefficients, already used for \eqref{r14:T}, give
$C=C^*\ge0$ and $C1=1$.  The function $b$ is positive
and nonconstant: the spatial plaquette action is zero at the
identity configuration and positive for suitable single-link
variations on this torus.

If $1$ were the vacuum eigenfunction, then $Cb=\lambda/b$
for some $\lambda>0$.  Write $\overline f=\int f\,dx$.
Integration and $C1=1$ give
$\lambda=\overline b/\overline{1/b}$.
Strict Cauchy--Schwarz gives
$\overline b\,\overline{1/b}>1$, hence
$\lambda<(\overline b)^2$.  On the other hand,
\[
 \lambda=\langle b,Cb\rangle
   =(\overline b)^2+
       \langle b-\overline b,C(b-\overline b)\rangle
   \ge(\overline b)^2,
\]
a contradiction.

Suppose now that every one-link reduction of the pure state
$|\psi\rangle\langle\psi|$ is pure.  Successive Schmidt
decompositions in the finite tensor product imply
$\psi(U)=\prod_e f_e(U_e)$.  Since $\psi$ is smooth and
strictly positive, the factors can be chosen smooth and
strictly positive: fix all but one coordinate in the product
identity, or equivalently recover each factor from its
positive one-coordinate section.  The almost-everywhere
product identity then extends everywhere by continuity.

Gauge invariance at a vertex says that the product, over its
incident links, of ratios such as $f_e(gU_e)/f_e(U_e)$
or $f_e(U_eg^{-1})/f_e(U_e)$ is one for independent link
variables.  Holding all other variables fixed proves that
each ratio is constant in its own link variable.  Haar
invariance of the $L^2$ norm and positivity force that
constant to be one.  Each link has distinct endpoints, so
applying the transformation at its initial vertex proves
left-translation invariance of $f_e$ for every $g\in\SU(2)$.
Every factor is constant.  This contradicts the established
nonconstancy of $\psi$.  Thus at least one one-link reduction
has purity strictly less than one.

Finally, the fixed-regulator simple Perron eigenvalue and
trace-class property imply
$\rho_q\to|\psi\rangle\langle\psi|$ in trace norm.
In detail, for $r_a=\lambda_a/\Lambda<1$ on the excited
space, $\sum_{a>0}r_a^{2q}\to0$ by dominated convergence,
dominated by the summable sequence $r_a^2$.
Partial trace contracts the trace norm on self-adjoint
trace-class operators.  For density operators $\rho,\sigma$,
$|\Tr\rho^2-\Tr\sigma^2|\le2\|\rho-\sigma\|_1$.
These facts prove both limits and strict positivity of
the finite sum in \eqref{r15:strict-floor}.
\end{proof}

This is a distinct obstruction from f15 V81's compressed-swap
Gauss commutator: even allowing full-Haar swaps and using
their correct expectation values, the elementary spatial
union bound does not follow the vanishing temporal defect.
It does not rule out a carefully normalized ratio method or
conditional spatial cancellation.  Such a method must prove
its own cancellation of the vacuum entanglement; setting
each local defect to zero is not valid.

\subsection{Export to V14 and the next unsolved estimate}
\label{r15:export}

Suppose a certified calculation of the actual expectation
in \eqref{r15:lower} gives $\underline R_{2q}>1/2$.
Then V14 applies with its purity period $p=2q$, not $q$:
\[
 \ell_{2q}=\frac{1+\sqrt{2\underline R_{2q}-1}}2,
 \qquad
 \operatorname{gap}(-\log K)
    \ge\frac1{2q}\log\frac{\ell_{2q}}{1-\ell_{2q}}.
\tag{V15.16}\label{r15:export-gap}
\]
The same value supplies V14's lower score-Gram subtraction
using periods $2q$ and $4q$.  This is just the already proved
export of a purity input; no such uniform native input is
numerically or analytically certified in this appendix.
The physical-time lower bound is the right side of
\eqref{r15:export-gap} divided by $a_{\rm t}$.  Keeping
$\underline R_{2q}>1/2$ while $q a_{\rm t}\to\infty$
does not on its own give a positive limiting mass.

The next substantive target is to estimate $\E_{\mu_q}d$
or the mean-square bridge covariance in \eqref{r15:force-bound}
at matched coupling, spatial volume and physical time, with
certified errors.  It requires no score-bank choice and
retains all physical sectors.  Finite-state demonstrations
or the exponentially small floor \eqref{r15:floor} cannot
replace this task.  We have not constructed an interacting
continuum theory or proved the full Yang--Mills mass gap.

\paragraph{Verification and preservation.}
The exact-rational checker
\path{scripts/verify_ym_restart_v15_temporal_sewing.py}
tests the raw and bounded sewing identities, variance and
determinant formulas, endpoint-factor cancellation, conservative
approximation bounds, a conditional mixed-log-derivative identity,
and the spatial vacuum-floor mechanism in positive finite
matrices.  These are algebra fixtures, not Wilson quadratures.
The analytic Wilson proofs above do not depend on the fixtures.
V14's complete body is retained apart from the displayed title;
reversing that title change and removing this appendix recovers
SHA--256
\path{8152CB096B6852F6C3F8B41177126DE3ED079486C46FF2DD8DD4E6C98899F0C7}.
Only active V14 is replaced by V15.  All 46 other source
files remain protected, and compilation and checking output
stay outside \path{latex_notes}.

% END CONSOLIDATED V15 APPEND

% BEGIN CONSOLIDATED V16 APPEND -- 2026-09-19
\clearpage
\section{V16: a native linear-volume bound on the temporal purity defect}
\label{r16:main}

This appendix supplies an evaluated native Wilson input for V15,
not another sufficient condition with an unspecified expectation.
Its domain is explicitly the bare strong-coupling region
$c_\gamma=18\tanh\gamma<1$.  We prove
\[
 R_{2q}\ge\exp\{-\mathcal B_{\gamma,N,q}\},\qquad
 \mathcal B_{\gamma,N,q}
 =\frac{324\pi^4\gamma^4N}{(1-c_\gamma)^4}
                         c_\gamma^{\,2q-4},\quad q\ge2.
\tag{V16.1}\label{r16:main-bound}
\]
Thus the purity defect is at most $\mathcal B_{\gamma,N,q}$.
The volume cost is linear in $N=3L^3$, not quadratic, and
the time separation is the actual Wilson temporal separation.
At $L=4$, $\gamma=1/40$, $q=6$, exact rational inequalities
give $R_{12}=Z_{24}/Z_{12}^2\ge15/16$.

This does \emph{not} solve the weak-coupling or continuum problem.
The archive f6 V19 already proves a bare Wilson Dobrushin row,
and f6 V22 already explains strong-coupling covariance decay
and its physical spectral consequence.  Those facts and the
existence of a strong-coupling mass gap are not new claims here.
The new work is the local oscillation estimate for V15's
\emph{integrated} temporal log-sewing observable and its
conversion, under the unchanged native law, into the
linear-volume partition-ratio certificate \eqref{r16:main-bound}.
This is a quantitative test of the proposed purity method
in a regime where its inputs can actually be paid.

\subsection{Inputs, literature check, and inherited limitations}
\label{r16:inputs}

We retain V15's $T$, $\kappa=T^q$, $a,b$, $v=\log(b/a)$,
and the endpoint marginal $\mu_q$ of two independent native
periodic histories of length $2q$.  Equivalently we view
$v$ as an observable of those full histories; integrating
the interiors leaves the same $\mu_q$.  The identities
\[
 \E_{\mu_q}e^v=R_{2q},\qquad
 \E_{\mu_q}e^{2v}=1
\tag{V16.2}\label{r16:moments}
\]
are exactly \eqref{r15:exact}, not additional assumptions.

The nonduplication audit read f6 V19's one-link influence
and comparison estimates, the f6 V22 strong-coupling
spectral discussion, f15 V91's global-projective limitation,
and f16.2 V20's conditional-channel versus full-transfer
distinction.  The previous bound $|\mathcal C_q|\le\pi^2N^2$
in \eqref{r15:coarse} supplies no temporal decay.  A global
kernel floor would lose exponentially in the whole spatial
volume.  Neither is used to establish \eqref{r16:main-bound}.

A targeted primary-source check consulted Rebeschini and
van Handel, \href{https://arxiv.org/abs/1308.4117}
{\emph{Comparison Theorems for Gibbs Measures}}, particularly
Corollary 2.6 and its total-variation specialization.
The comparison and bounded-difference methods below are
classical; we give finite-volume proofs to fix constants,
the direction of the influence matrix, and their validity
after arbitrary boundary freezing.  No theorem extending
our bare Wilson condition to weak coupling is imported.
The next companion review and broad literature sweep remain V20.

\subsection{Finite Gibbs estimates with all constants paid}
\label{r16:gibbs}

Consider a finite compact product space with a smooth strictly
positive Gibbs density.  A site is a whole link variable.
For real bounded $F$ let $\delta_iF$ be its oscillation when
only site $i$ changes.  Let a symmetric nonnegative matrix
$C$, with $C_{ii}=0$, dominate the single-site conditional
total-variation influences, using
$\|\nu-\nu'\|_{\rm TV}=\sup_A|\nu(A)-\nu'(A)|$.
Assume $\sup_i\sum_jC_{ij}\le c<1$ and put
$D=(I-C)^{-1}=\sum_{n\ge0}C^n$.
The same domination is assumed for the restrictions obtained
by freezing any sites; the Wilson matrix below has this property.

\begin{lemma}[Covariance, conditioning, and concentration]
\label{r16:gibbs-lemma}
For the above law $\nu$,
\[
 |\Cov_\nu(F,G)|\le\frac14\,\delta(F)^tD\delta(G).
\tag{V16.3}\label{r16:cov}
\]
If the supports of $F,G$ are at interaction-graph distance
at least $r$, where $C_{ij}=0$ off graph edges, then
\[
 |\Cov_\nu(F,G)|
 \le\frac{c^r}{4(1-c)}
          \|\delta(F)\|_1\|\delta(G)\|_\infty.
\tag{V16.4}\label{r16:distance}
\]
For every real $t$,
\[
 \log\E_\nu e^{t(F-\E_\nu F)}
 \le\frac{t^2}{8}\|D\delta(F)\|_2^2
 \le\frac{t^2}{8(1-c)^2}\|\delta(F)\|_2^2.
\tag{V16.5}\label{r16:concentration}
\]
\end{lemma}
\begin{proof}
Let $K_i$ be conditional expectation resampling site $i$,
and $P_s$ the rate-one-per-site heat-bath semigroup with
generator $\mathcal L=\sum_i(K_i-I)$.  This is only an
auxiliary proof device; its time $s$ is not Wilson transfer
time.  Coupling two conditional updates maximally gives
\[
 \delta_i(K_iF)=0,\qquad
 \delta_j(K_iF)\le\delta_jF+C_{ij}\delta_iF\quad(j\ne i).
\]
For example, in the second inequality couple the updated
site, pay $\delta_jF$ for the unchanged external discrepancy,
and pay $\delta_iF$ only on a coupling mismatch.
Expansion by random update times, or the finite Euler
approximation followed by a limit, consequently gives
\[
 \delta(P_sF)\le e^{(-I+C^t)s}\delta(F).
\]
The positive compact density ensures convergence to $\nu$;
alternatively the last inequality forces the oscillation
to zero, and invariance fixes the limiting constant.
Reversibility and differentiation now imply
\[
 \Cov_\nu(F,G)=\int_0^\infty\sum_i
       \E_\nu\Cov(F_s,G\mid U_{j},\ j\ne i)\,ds,
 \qquad F_s=P_sF.
\]
A conditional covariance is at most one quarter of the
product of its two ranges.  Integrating the oscillation
bound proves \eqref{r16:cov}, with the transpose in the
semigroup canceling on taking the scalar product.
If the supports are $r$ apart, $(C^n)_{ij}=0$ between them
for $n<r$.  The remaining row sums are bounded by
$\sum_{n\ge r}c^n$, proving \eqref{r16:distance}.

For concentration, expose the sites in a fixed order and
write $F-\E F$ as its Doob martingale differences.
When site $i$ is revealed, the remaining conditional laws
for two values of $U_i$ have a boundary discrepancy vector
bounded by $(C_{ji})$ on the unrevealed sites $J$.
Coupled single-site updates propagate a disagreement vector
by the internal influence matrix $C_J$ and this boundary
vector.  Iterating the affine inequality, or taking the
stationary limit of the coupled updates, bounds the
disagreement probabilities by
$(I-C_J)^{-1}C_{J i}$.  This also proves directly the
finite comparison estimate needed here: the difference
of the two conditional expectations of $F$ is at most
\[
 \delta_iF+\sum_{j\in J}\delta_jF
                    [(I-C_J)^{-1}C_{J i}]_j
 \le [D\delta(F)]_i.
\]
The last step follows by expanding nonnegative path sums
and using the symmetry of $C$.  It remains valid for each
value of the already exposed coordinates.

Thus the conditional range of martingale difference $i$
is at most $B_i=[D\delta(F)]_i$.  A centered random
variable of range at most $B_i$ has log moment-generating
function at most $t^2B_i^2/8$: its second derivative is
a tilted variance at most $B_i^2/4$, and its value and
first derivative vanish at zero.  Iterating this conditional
inequality proves the first bound in \eqref{r16:concentration}.
Finally symmetry and the row bound give
$\|C\|_{2\to2}\le c$ and $\|D\|_{2\to2}\le(1-c)^{-1}$.
\end{proof}

\paragraph{The actual Wilson influence matrix.}
Let $h_{ij}$ count plaquettes containing both distinct
free links $i,j$, including plaquettes with some frozen
boundary links.  In our normalization each plaquette
trace lies in $[-1,1]$.  Changing link $j$ changes the
conditional log density at $i$ by a function of oscillation
at most $4\gamma h_{ij}$.  Two exponential tilts whose
log-density difference has oscillation $u$ have TV distance
at most $\tanh(u/4)$.  One elementary proof bounds their
density ratio between two numbers with quotient $e^u$;
the largest mean absolute deviation at fixed mean one
is attained by an endpoint-valued variable, and optimization
gives $(e^{u/2}-1)/(e^{u/2}+1)$.
Consequently the symmetric dominating matrix is
\[
 C_{ij}=h_{ij}\tanh\gamma,
 \qquad\sup_i\sum_{j\ne i}C_{ij}\le18\tanh\gamma=c_\gamma.
\tag{V16.6}\label{r16:wilson-C}
\]
Here $\tanh(k\gamma)\le k\tanh\gamma$ and each link
belongs to at most six plaquettes, each with three other
link slots.  Removing or freezing links cannot increase
this count.  The same bound therefore holds for every
endpoint-conditioned bridge, each periodic history, their
independent product, and every conditional law used while
exposing those histories.  We assume $L\ge4$ and temporal
period $2q\ge4$ to avoid short-period identifications.
This is the inherited bare row, with its scope unchanged.

\subsection{A single endpoint link has a volume-free sewing oscillation}
\label{r16:local}

The forces $g_L,g_R$ in \eqref{r15:covariance} are gradients
of sums of $N$ boundary temporal plaquettes.  Their primitive
incidence, rather than a global norm of the full slice,
is essential for the next estimate.

\begin{proposition}[Local sensitivity of the integrated log sewing]
\label{r16:oscillation}
If $c_\gamma<1$, every one-link oscillation of $v$ in one
of its four endpoint slices is at most
\[
 b_{\gamma,q}
   :=\frac{18\pi^2\gamma^2}{1-c_\gamma}
                             c_\gamma^{q-2}.
\tag{V16.7}\label{r16:local-bound}
\]
As a function on the full pair of periodic histories,
$v$ has zero oscillation at every non-endpoint link.
In particular
\[
 \|\delta(v)\|_2^2\le4N b_{\gamma,q}^2.
\tag{V16.8}\label{r16:osc-square}
\]
\end{proposition}
\begin{proof}
Change only link $e$ of $x_1$ along a minimizing round
$S^3$ geodesic of speed at most $\pi$, and connect
$y_1$ to $y_2$ by product geodesics with speed at most
$\pi$ on each link.  The change of $v$ is the mixed
endpoint rectangle integral from \eqref{r15:covariance}.
At every point of that rectangle the left force
$F=g_L\cdot\dot X$ is the derivative of just one
temporal plaquette.  With its boundary spatial link
fixed, that plaquette contains exactly three free links:
two temporal links and its opposite spatial link.
Each one-link oscillation of $F$ is at most $2\pi$, so
\[
 \|\delta(F)\|_1\le6\pi.
\]
For $G=g_R\cdot\dot Y$, an opposite spatial link occurs
in one boundary plaquette and a temporal link in at
most six.  Each summand has absolute value at most
$\pi$, hence
\[
 \|\delta(G)\|_\infty\le12\pi.
\]
These estimates retain the temporal links implementing
the Gauss integration; they are not temporal-gauge
estimates with links omitted.

The free support of $F$ lies on temporal links in slab
$0$ and spatial links at time $1$.  The free support
of $G$ lies in slab $q-1$ and at spatial time $q-1$.
Assign spatial links height $t$ and temporal links
height $t+1/2$.  Links sharing a plaquette have heights
differing by at most one.  Thus the two supports have
interaction distance at least $q-2$, including the
overlapping-support case $q=2$.
By \eqref{r16:distance}, for the \emph{actual normalized}
bridge at every point of the endpoint rectangle,
\[
 |\Cov(F,G)|\le
       \frac{18\pi^2}{1-c_\gamma}c_\gamma^{q-2}.
\]
Multiply by $\gamma^2$ and integrate over the unit
rectangle to obtain \eqref{r16:local-bound}.
The same reasoning applies to the other three endpoint
slices, by reflection and exchange of copies.
There are exactly $4N$ such link sites in the two
periodic histories.  The block kernel has already
integrated its interiors, so its value depends on no
other random coordinates.  This proves \eqref{r16:osc-square}.
\end{proof}

\subsection{The fluctuation identity removes a quadratic volume loss}
\label{r16:purity}

\begin{theorem}[Native linear-volume temporal purity certificate]
\label{r16:purity-theorem}
For $L\ge4$, $q\ge2$ and $0<c_\gamma<1$,
\eqref{r16:main-bound} holds for the physical Wilson trace,
and
\[
 \E_{\mu_q}d=1-R_{2q}
       \le1-e^{-\mathcal B_{\gamma,N,q}}
       \le\mathcal B_{\gamma,N,q}.
\tag{V16.9}\label{r16:defect}
\]
All expectations are native; no Gaussian approximation,
Perron-vector approximation, or Monte Carlo error is used.
\end{theorem}
\begin{proof}
Apply \eqref{r16:concentration} to $v$ under the two full
periodic histories, with their block-diagonal Wilson
influence matrix.  Equations \eqref{r16:moments} and
\eqref{r16:osc-square}, at $t=2$, imply
\[
 1=\E e^{2v}
 \le\exp\left\{2\E v+
              \frac{2N b_{\gamma,q}^2}{(1-c_\gamma)^2}\right\}.
\]
Hence $\E v\ge-Nb_{\gamma,q}^2/(1-c_\gamma)^2$.
Jensen's inequality now gives
$R_{2q}=\E e^v\ge\exp(\E v)$.
Substitution of \eqref{r16:local-bound} proves
\eqref{r16:main-bound}.  V15 identifies $\E d$ with
the actual purity defect.  Finally $1-e^{-x}\le x$ for
$x\ge0$.
\end{proof}

This argument uses the exact fluctuation identity at
$t=2$ to control the otherwise unknown mean of $v$.
It does not assume that mean is zero.  Bounding the
whole endpoint rectangle by its worst value and then
squaring would instead give an $N^2$ cost.  The
bounded-difference argument sums squares of the $4N$
local oscillations; the internal history variables are
paid by $(1-c_\gamma)^{-2}$, not by another factor of
the number of time slices or spatial links.

\begin{corollary}[A fully specified native test point]
\label{r16:point}
At $L=4$, $\gamma=1/40$,
\[
 R_{12}\ge\frac{15}{16},\qquad
 \operatorname{gap}(-\log(T/\Lambda))
       \ge\frac{\log29}{12}>\frac14.
\tag{V16.10}\label{r16:numeric}
\]
The last gap is in lattice transfer units, not a
constructed continuum physical Hamiltonian.
\end{corollary}
\begin{proof}
Here $N=192$, $q=6$ and
$c_\gamma\le18\gamma=9/20$.  Use $\pi<22/7$ in
\eqref{r16:main-bound}.  Its upper bound is the exact rational
\[
 324\left(\frac{22}{7}\right)^4\frac{192}{40^4}
       \frac{(9/20)^8}{(11/20)^4}
 =\frac{10460353203}{240100000000}<\frac1{16}.
\tag{V16.11}\label{r16:rational}
\]
Thus $R_{12}\ge1-\mathcal B>15/16$.
For V14's high-purity branch,
$\sqrt{2(15/16)-1}=\sqrt{7/8}>14/15$,
so $\ell_{12}>29/30$ and
$(1-\ell_{12})/\ell_{12}<1/29$.
Equation \eqref{r15:export-gap} proves the gap bound.
For example $e<3$ implies $e^3<27<29$, proving
$\log29/12>1/4$ without a decimal logarithm.
\end{proof}

\paragraph{Growing spatial volumes, with the correct clock.}
Fix a coupling in this same strong-coupling domain and
put $\mu=-\log c_\gamma>0$.  For $A>1$, choose
\[
 q_N=\max\left\{2,\left\lceil
                  \frac{A\log N}{2\mu}\right\rceil\right\}.
\tag{V16.12}\label{r16:log-volume}
\]
Then $\mathcal B_{\gamma,N,q_N}=O(N^{1-A})$.
To check the spectral scale explicitly, use the certified
lower purity $1-\mathcal B$ once $\mathcal B<1/2$.
It gives
$1-\ell=(1-\sqrt{1-2\mathcal B})/2
=\mathcal B/2+O(\mathcal B^2)$.
Since $q_N=A\log N/(2\mu)+O(1)$, the resulting
V14 lower gap bound tends to
\[
                    (1-1/A)\mu.
\tag{V16.13}\label{r16:asymptotic}
\]
Indeed the logarithm of $\mathcal B$ equals
$(1-A)\log N+O(1)$, including the bounded rounding
factor.  This is a native Wilson purity calculation,
not V14's independent-product example.  Its familiar
strong-coupling gap scale is a consistency check,
not a new weak-coupling mass claim.  At a physical
time unit $a_{\rm t}$ every such lower gap must still
be divided by $a_{\rm t}$, and a continuum theory
must be constructed with that normalization.

\subsection{What remains outside the bare contraction region}
\label{r16:remaining}

For large $\gamma$, the dominating row
$18\tanh\gamma$ exceeds one.  Its geometric series
and concentration proof are then unavailable.
This is a failure of this estimate, not evidence
of a Wilson phase transition or absence of a mass gap.
The archive already gives weak-coupling obstructions
to a uniform worst-exterior one-link contraction;
we do not propose extending \eqref{r16:wilson-C}
by analyticity or by discarding exceptional exteriors.
Keeping $\gamma=1/40$ and simply declaring
$a_{\rm t}\to0$ is not a continuum Yang--Mills construction.

There is a precise way to record the remaining
normalized fluctuation cost without this bare estimate.
For every fixed regulator, define
\[
 F_q(t)=\log\E_{\mu_q}e^{tv},\qquad
 d\mu_{q,t}=e^{tv-F_q(t)}\,d\mu_q .
\tag{V16.14}\label{r16:tilt}
\]
Haar exchange of the full $y$ boundaries sends $a$
to $b$.  Thus $F_q(t)=F_q(2-t)$, $F_q(0)=F_q(2)=0$,
$F_q(1)=\log R_{2q}$ and $F_q'(1)=0$.
Differentiation on this fixed compact space gives
$F_q''(t)=\Var_{\mu_{q,t}}v$; Taylor's integral formula
about $t=1$ therefore yields the exact identity
\[
 -\log R_{2q}=\int_0^1 t\,\Var_{\mu_{q,t}}v\,dt.
\tag{V16.15}\label{r16:variance-cost}
\]
At $t=0$ the law is the two native $2q$-periodic
histories' endpoint marginal.  At $t=1$ its density
is $ab/Z_{4q}$, the endpoints of a single sewn
$4q$-periodic history.  Intermediate powers
$a^{2-t}b^t$ are normalized integrated-kernel
interpolations, not automatically local Wilson
actions with the same Dobrushin row.

Equation \eqref{r16:variance-cost} is ordinary
thermodynamic integration applied to the exact
sewing law, not a newly established weak-coupling
bound.  It identifies what any blocked or weighted
replacement must control: the fluctuations of
this normalized cross-boundary logarithm, with
vacuum endpoint factors canceled and all sectors
retained.  A full supremum comparison of spatial
partial swaps would reintroduce V15's nonzero
vacuum-entanglement floor.  The next upstream
task is to estimate this cost, or the original-law
exponential fluctuation cost in the preceding
proof, beyond the bare small-$\gamma$ domain.

\paragraph{Verification and preservation.}
The checker \path{scripts/verify_ym_restart_v16_native_purity.py}
enumerates actual four-dimensional plaquette incidence,
the three-free-link boundary-force supports and their
temporal distance.  It checks the Gibbs covariance
and martingale-range constants in exact finite rational
models, encloses their selected log moment-generating
functions with rational series bounds, and verifies
the native arithmetic certificate \eqref{r16:rational}.
The finite models test the algebra, not Wilson quadratures;
the explicit Wilson point is certified by the analytic
theorem and exact inequalities, not by sampling.
V15's algebra checks are also rerun.  Its complete
body is retained except for the displayed title;
removing this appendix and reversing that change
recovers SHA--256
\path{31BC0F74979FD7D2632306146A4B61ABEC28B24343E5859AE96A40A558889AA8}.
Only active V15 is replaced.  The other 46 sources
remain protected, and all build output stays outside
\path{latex_notes}.

\paragraph{Status.}
A nontrivial, all-sector native periodic purity
input is now evaluated in a declared strong-coupling
regime, with linear volume cost and explicit temporal
decay.  Uniform weak-coupling control, the interacting
four-dimensional continuum construction, and its
full Yang--Mills mass gap remain unproved.

% END CONSOLIDATED V16 APPEND

% BEGIN CONSOLIDATED V17 APPEND -- 2026-09-19
\clearpage
\section{V17: blocking counterterms and a native temporal swap floor}
\label{r17:main}

V16's bare comparison row fails at weak coupling.  Before
trying a blocked replacement, we must distinguish two
operations: interpolate local Wilson weights and then
integrate a temporal block, or first integrate the block
and then interpolate its normalized endpoint law.
They are not interchangeable.  This appendix identifies
their exact conditional-affinity counterterm, evaluates
it as a product of four terminal-slab bridge integrals,
and gives an all-coupling normalization bound independent
of block length.  It also proves a native obstruction:
the uncorrected single-time-slice swap has a strictly
positive limiting defect, although the desired temporal
purity defect tends to zero at each fixed regulator.

These results do not establish weak-coupling mixing.
They rule out estimating the wrong local interpolation
and silently transferring that estimate to V16's sewing
law.  The counterterm must be retained in the next
weak-coupling calculation; it cannot be discarded as a
vanishing boundary error.

\subsection{Scope and nonduplication}
\label{r17:scope}

Fix $L\ge4$, $N=3L^3$, $\gamma>0$, and $q\ge2$.
Use V15's actual Wilson transfer and its
$\kappa=T^q$, $a,b$, $v=\log(b/a)$, $\mu_q$ and $R_{2q}$.
Every limit in this appendix holds at fixed $L,\gamma$;
none exchanges that limit with growing volume or a
continuum limit.

The archive audit reread f16.2 V1's static conditional
affinity in a kinetic compression, f16.2 V35's normalized
square-root estimates on a convex chart, and f5 V19's
Wilson--heat geometric interpolation.  Conditional
affinities and normalized interpolation are therefore
inherited methods, not new concepts here.  Those results
do not identify the temporal blocking/interpolation
counterterm below or the limiting defect of a
single-time-slice swap.  V15's spatial reduced-purity
floor is also different: our swap exchanges an entire
spatial slice and does not use a partial spatial tensor
factorization of the physical Hilbert space.

A targeted primary-source check read van Erven and
Harremo\"es, \href{https://arxiv.org/abs/1206.2459}
{\emph{R\'enyi Divergence and Kullback--Leibler Divergence}},
Theorems 1 and 9 on data processing and the conditional
density argument preceding them.  Monotonicity of
Hellinger affinity under marginalization is classical.
We prove its particular conditional factorization here
so that no conditional normalizer is hidden.  The next
companion and broad literature checkpoints remain V20.

\subsection{Two orders of integration give different laws}
\label{r17:disintegration}

Let $\mathbb P_q$ be the full Wilson probability law of
two independent periodic histories, each of length $2q$.
Write their four antipodal spatial slices as
$z=(x_1,y_1,x_2,y_2)$, at times $0,q$, and put all
remaining spatial and temporal links in $h$.  All
reference measures below are products of probability
Haar measures.  Denote the density by $p_q(z,h)$ and
its endpoint marginal by
\[
 m_q(z)=\frac{a(z)^2}{Z_{2q}^2},\qquad
 \pi_z(h)=\frac{p_q(z,h)}{m_q(z)}.
\tag{V17.1}\label{r17:p}
\]
Let $S$ exchange $y_1,y_2$ while leaving every variable
in $h$ and both $x$ slices unchanged.  It is a Haar
involution on the full history space.  The swapped
density is $p_q(Sz,h)$, with marginal
$m_q(Sz)=b(z)^2/Z_{2q}^2$ and conditional density
$\pi_{Sz}(h)$.  This full-history swap keeps the
microscopic interiors fixed.  It is not V15's already
integrated endpoint observable.

For $0\le s\le1$, define
\[
\begin{split}
 H_{q,s}(z)&=\int\pi_z(h)^{1-s}\pi_{Sz}(h)^s\,dh,\\
 J_q(s)&=\int p_q(z,h)^{1-s}p_q(Sz,h)^s\,dz\,dh,\\
 d\nu_{q,s}(z)&=J_q(s)^{-1}
       \left[\int p_q(z,h)^{1-s}p_q(Sz,h)^s\,dh\right]dz.
\end{split}
\tag{V17.2}\label{r17:def}
\]
Here $\nu_{q,s}$ is what results by interpolating the
microscopic weights first.  V16's correctly blocked
sewing law is instead $\mu_{q,2s}$ from \eqref{r16:tilt}.

\begin{theorem}[Exact blocking/interpolation counterterm]
\label{r17:factor}
For all the fixed-regulator parameters above,
\[
\begin{split}
 J_q(s)&=e^{F_q(2s)}\E_{\mu_{q,2s}}H_{q,s},\\
 \frac{d\nu_{q,s}}{d\mu_{q,2s}}(z)
   &=\frac{H_{q,s}(z)}{\E_{\mu_{q,2s}}H_{q,s}},\qquad
 0<H_{q,s}\le1.
\end{split}
\tag{V17.3}\label{r17:factor-eq}
\]
For bounded $f$ the inverse change of law is
\[
 \E_{\mu_{q,2s}}f
    =\frac{\E_{\nu_{q,s}}[f/H_{q,s}]}
           {\E_{\nu_{q,s}}[1/H_{q,s}]}.
\tag{V17.4}\label{r17:inverse}
\]
For $0<s<1$, $H_{q,s}(z)=1$ if and only if
$\pi_z=\pi_{Sz}$ almost everywhere.  Moreover
$H_{q,s}(Sz)=H_{q,1-s}(z)$.
\end{theorem}
\begin{proof}
Disintegrating the interpolated density gives exactly
\[
 \int p_q(z,h)^{1-s}p_q(Sz,h)^s\,dh
       =m_q(z)^{1-s}m_q(Sz)^s H_{q,s}(z).
\]
The first factor is $m_q(z)e^{2sv(z)}$.
Integrating it alone gives $e^{F_q(2s)}$, by definition.
This proves both normalization and density formulas,
and their inversion proves \eqref{r17:inverse}.
H\"older's inequality gives $H_{q,s}\le1$; its equality
condition and normalization give the stated iff claim.
All densities are positive and smooth on fixed compact
spaces.  Exchanging $z,Sz$ proves the symmetry.
\end{proof}

In particular the microscopic affinity
$A_q^{\rm full}:=J_q(1/2)$ satisfies
\[
 A_q^{\rm full}
       =R_{2q}\E_{\mu_{q,1}}H_{q,1/2}\le R_{2q}.
\tag{V17.5}\label{r17:half}
\]
Its deficit is an upper bound for the desired defect,
but the next sections show why that upper bound can
retain a nonvanishing error.  The probability
$\nu_{q,s}$ is biased by $H_{q,s}$ relative to the
blocked law, not merely by an ignorable scalar.

For completeness, at fixed $z$ put
$\ell_z(h)=\log(\pi_{Sz}(h)/\pi_z(h))$.
Differentiating its conditional normalizer gives
\[
 \frac{d^2}{ds^2}\log H_{q,s}(z)
       =\Var_{\pi_{z,s}}\ell_z,\qquad
 \pi_{z,s}=\frac{\pi_z^{1-s}\pi_{Sz}^s}{H_{q,s}}.
\tag{V17.6}\label{r17:conditional-variance}
\]
Thus $-\log H_{q,s}$ is a nonnegative concave
interpolation correction with zero endpoint values.
It measures fluctuations of a conditional likelihood,
not fluctuations under a freely chosen reference law.

\subsection{The counterterm is four terminal-slab bridge integrals}
\label{r17:terminal}

Orient each of the four half-period segments from its
$x$ endpoint to its $y$ endpoint; reverse time on the
return half of each history.  Wilson time reflection
and inversion of temporal links preserve action and
Haar measure.  Given $z$, the four segment interiors
are independent.  Let $\beta_{q;x,y}$ denote the
normalized full $q$-slab bridge law on one such
interior, with its temporal links retained.  Set
\[
 B_{q,s}(x;y,y')
       =\int\beta_{q;x,y}^{1-s}\beta_{q;x,y'}^s.
\tag{V17.7}\label{r17:bridge-aff}
\]
The exact factorization is
\[
 H_{q,s}(z)=B_{q,s}(x_1;y_1,y_2)^2
                     B_{q,s}(x_2;y_2,y_1)^2.
\tag{V17.8}\label{r17:four}
\]
The powers two count two actual independent half-period
interiors, not repeated Radon--Nikodym factors on the
same random bridge.

There is also an explicit scalar Wilson integral for
each factor.  Write $b(y)=e^{-\gamma S_s(y)/2}$ and
remove only the final spatial half-action:
\[
 \widehat\kappa_q(x,y)=\kappa(x,y)/b(y).
\tag{V17.9}\label{r17:hat}
\]
Extend the final endpoint to a vector $Y$ of quaternions
in the closed unit ball by using the same integral and
the same bulk action, with $Y$ substituted only in the
last temporal plaquette sum.  Quaternion conjugation
is extended linearly.  The removed spatial endpoint
action is not reinserted at nonunit $Y$.
Every last temporal plaquette is affine in its single
final endpoint quaternion, so for $Y_s=(1-s)y+sy'$,
\[
 B_{q,s}(x;y,y')
     =\frac{\widehat\kappa_q(x,Y_s)}
            {\widehat\kappa_q(x,y)^{1-s}
             \widehat\kappa_q(x,y')^s}.
\tag{V17.10}\label{r17:chord}
\]
To prove this, write the two normalized bridge densities
as exponential weights divided by their respective
$\widehat\kappa$ normalizers.  Their common bulk action
occurs with total coefficient one.  Only the last
temporal action is interpolated, and its affine
dependence gives the numerator in \eqref{r17:chord}.
The initial spatial half-action is the same in all
three integrals and cancels as well.
The chord $Y_s$ is not asserted to be a group-valued
configuration or a new physical Yang--Mills measure.
Equation \eqref{r17:chord} is a representation of
the original conditional-affinity integral.

\begin{proposition}[All-coupling counterterm bound without a depth loss]
\label{r17:cost}
For every $\gamma>0$, $q\ge2$ and all endpoints,
\[
 \frac{1}{\cosh^4(\gamma N)}\le H_{q,1/2}(z)\le1.
\tag{V17.11}\label{r17:H-lower}
\]
Consequently
\[
\begin{gathered}
 A_q^{\rm full}\le R_{2q}
       \le\min\{1,A_q^{\rm full}\cosh^4(\gamma N)\},\\
 0\le\log R_{2q}-\log A_q^{\rm full}
          \le4\log\cosh(\gamma N)\le4\gamma N.
\end{gathered}
\tag{V17.12}\label{r17:cost-eq}
\]
The cost is independent of $q$, but is not small or
volume-uniform at weak coupling.
\end{proposition}
\begin{proof}
For a fixed $x$, the log likelihood ratio between
$\beta_{q;x,y'}$ and $\beta_{q;x,y}$ depends on the
interior only through the last temporal plaquettes.
Each of the two last-slab actions lies in $[0,2N]$.
Their difference lies in $[-2N,2N]$, so the oscillation
of this normalized log likelihood is at most $4\gamma N$.
The scalar normalizer changes its value, not its
oscillation.

We use an elementary sharp affinity bound.  If
$r=d\eta'/d\eta$, $\E_\eta r=1$, and
$\operatorname{osc}\log r\le R$, then
$\E_\eta\sqrt r\ge1/\cosh(R/4)$.
Indeed, put $u=\sqrt{\inf r}$, $w=\sqrt{\sup r}$.
The chord lower bound for the concave square root gives
\[
 \E\sqrt r\ge\frac{1+uw}{u+w}.
\]
If $w=cu$, then $c\le e^{R/2}$ and the last expression
is at least $2\sqrt c/(1+c)\ge1/\cosh(R/4)$.
Infimum/supremum approximations give the same statement
without attained endpoints.  Applying it to the single
bridge yields $B_{q,1/2}\ge1/\cosh(\gamma N)$.
Four factors in \eqref{r17:four} prove
\eqref{r17:H-lower}; integration in \eqref{r17:half}
proves \eqref{r17:cost-eq}.
\end{proof}

This estimate does not incur a factor proportional to
the temporal depth in its exponent.  Nevertheless,
using its absolute cost at weak coupling would still
lose the spatial volume.  A useful future estimate
must exploit the joint cancellation in
$\log R_{2q}=\log A_q^{\rm full}
-\log\E_{\mu_{q,1}}H_{q,1/2}$, not separately bound
both contributions by their whole-volume oscillations.

\subsection{A local temporal swap has a strictly positive limiting defect}
\label{r17:floor}

To separate a physical temporal effect from a poor
choice of temporal-link coordinates, first integrate
all temporal links, but retain every spatial slice.
For one periodic history of length $m=2q$ the exact
slice density is
\[
 \frac1{Z_m}\prod_{j=0}^{m-1}T(U_j,U_{j+1}),\qquad U_m=U_0.
\tag{V17.13}\label{r17:slice-law}
\]
On two independent such histories exchange only the
slice at time $q$, retaining every other spatial slice.
Denote the Hellinger affinity between these two
densities by $A_q^{\rm slice}$.  Marginalization and
\eqref{r17:half} imply
\[
                A_q^{\rm full}\le A_q^{\rm slice}\le R_{2q}.
\tag{V17.14}\label{r17:hierarchy}
\]
For the first inequality integrate temporal links; for
the second keep only the four antipodal slices.
Both integrations commute with the specified swap.
Either inequality follows directly by conditional
Cauchy--Schwarz, so no reconstruction theorem is
being assumed for an interpolating law.

Let $p(y)=\psi(y)^2$ be the native vacuum density and
let $P$ be the ground-state Markov kernel from V14,
with Haar transition density
\[
 P(y,u)=\frac{T(y,u)\psi(u)}{\Lambda\psi(y)},\qquad
 a_P(y,y')=\int\sqrt{P(y,u)P(y',u)}\,du.
\tag{V17.15}\label{r17:row-aff}
\]
The notation $a_P$ is a row affinity; it is not V15's
product $a$ of block kernels.

\begin{theorem}[Native microscopic temporal floor]
\label{r17:floor-thm}
At fixed $L\ge4$ and $\gamma>0$,
\[
\begin{split}
 \lim_{q\to\infty}A_q^{\rm slice}
    &=A_\infty^{\rm slice}
      :=\int p(y)p(y')a_P(y,y')^4\,dy\,dy'<1,\\
 \liminf_{q\to\infty}(1-A_q^{\rm full})
    &\ge1-A_\infty^{\rm slice}>0,
 \qquad \lim_{q\to\infty}(1-R_{2q})=0.
\end{split}
\tag{V17.16}\label{r17:floor-eq}
\]
For every bounded real, nonconstant smooth physical
observable $f$,
\[
 1-A_\infty^{\rm slice}
       \ge\frac{2\Var_p(Pf)}{(\operatorname{osc}f)^2}>0.
\tag{V17.17}\label{r17:floor-response}
\]
No regulator-uniform positive lower bound is claimed.
\end{theorem}
\begin{proof}
The three-slice marginal around the exchanged slice,
with neighboring slices $u,v$, is exactly
\[
 \rho_m(u,y,v)
       =\frac{T(u,y)T(y,v)T^{m-2}(v,u)}{Z_m}.
\]
At fixed regulator the simple Perron eigenvalue,
trace-class property, and smoothing give convergence
of this density to
\[
 \rho_\infty(u,y,v)=p(y)P(y,u)P(y,v).
\]
For example the spectral convergence of $K^{m-4}$,
with $K=T/\Lambda$, sandwiched between two smoothing
factors $K$ gives uniform convergence of the kernel
$K^{m-2}$ to $\psi\otimes\psi$; also
$\Tr K^m\to1$.  This justifies the density limit
without taking a spatial or continuum limit.

The likelihood ratio for the slice swap depends only
on these two triples.  It is a ratio of four positive
one-step kernels and is bounded on the fixed compact
space.  We may therefore pass to the limit in its
square-root expectation.  Given $y$, the two neighboring
slices under $\rho_\infty$ have independent law
$P(y,du)P(y,dv)$, by reversibility.  For fixed $y,y'$,
their affinity is $a_P(y,y')^2$.  The two independent
copies contribute its square, namely $a_P(y,y')^4$.
The factors $p(y)p(y')$ are unchanged by exchanging
$y,y'$.  This proves the first equality in
\eqref{r17:floor-eq}.

For probability rows $r,s$ with affinity $a$, the
Cauchy--Schwarz inequality applied to
$(\sqrt r-\sqrt s)(\sqrt r+\sqrt s)$ gives
$\|r-s\|_{\rm TV}^2\le1-a^2\le1-a^4$.
Since
$|rf-sf|\le\operatorname{osc}(f)\|r-s\|_{\rm TV}$,
integration over independent $y,y'$ of law $p$ yields
\eqref{r17:floor-response}: the integral of
$[Pf(y)-Pf(y')]^2$ is $2\Var_p(Pf)$.
This variance is strictly positive.  Indeed if $Pf$
were constant, injectivity of the native transfer on
the physical subspace, conjugated by the strictly
positive vacuum, would imply that $f$ is constant.
This is the fixed-regulator injectivity already used
in V11; it is not injectivity on the unreduced Gauss
complement.  Nonconstant bounded smooth physical
functions, for example a spatial plaquette trace,
exist, and $p$ is strictly positive.  Thus the floor
is strict.

Finally \eqref{r17:hierarchy} gives the full-history
liminf statement.  The trace-class Perron argument
already used in V15 gives $R_{2q}\to1$.
\end{proof}

This obstruction survives integrating out temporal
links and exchanging the \emph{entire} physical slice.
It is not f15 V81's obstruction to spatial partial
swaps preserving Gauss constraints, and not V15's
one-link vacuum-entanglement floor.  Here the retained
immediate past and future can distinguish two slice
values even when two widely separated endpoints
have become asymptotically independent.

Combining \eqref{r17:half} and \eqref{r17:floor-eq}
shows in particular
\[
 \limsup_{q\to\infty}
       \E_{\mu_{q,1}}H_{q,1/2}\le A_\infty^{\rm slice}<1.
\tag{V17.18}\label{r17:nonvanishing}
\]
The missing normalizer does not approach one merely
because the block becomes long.  Nor does
\eqref{r17:cost-eq} justify dropping it after dividing
by a temporal period: the mass-gap test uses the
small \emph{purity defect}, not just an intensive
free-energy error tending to zero.

\subsection{An exact finite-state illustration and the next target}
\label{r17:example}

To make the distinction checkable without Wilson
quadrature, consider the positive symmetric Markov
transfer
\[
 P=\frac1{25}\begin{pmatrix}16&9\\9&16\end{pmatrix},
 \qquad r=\frac7{25},\quad p=(1/2,1/2).
\tag{V17.19}\label{r17:matrix}
\]
Its nontrivial eigenvalue is $r$, and the affinity
between its distinct rows is $24/25$.  Thus, exactly,
\[
 R_{2q}=\frac{1+r^{4q}}{(1+r^{2q})^2}\longrightarrow1,
 \qquad
 A_\infty^{\rm slice}
     =\tfrac12\left[1+\left(\frac{24}{25}\right)^4\right]
     =\frac{722401}{781250}<1.
\tag{V17.20}\label{r17:matrix-values}
\]
For $f=(-1,1)$ the lower floor in
\eqref{r17:floor-response} is $49/1250$, below the
actual $58849/781250$.  These are finite-state
illustrations, not numerical Wilson measurements.
The strict Wilson statement is the analytic theorem
above, for arbitrary positive coupling at each
fixed regulator.

The next weak-coupling task is now more constrained.
A local interpolation of the original weights may
still be useful, but its endpoint law must be corrected
by \eqref{r17:inverse}, with the four exact terminal
affinities \eqref{r17:chord}.  Alternatively one can
work directly with the integrated block kernels and
estimate V16's normalized log-sewing fluctuations.
Both routes require a joint estimate retaining the
conditional cancellation.  Bounding the microscopic
swap defect alone cannot tend to zero with block
depth, and replacing conditional bridge integrals by
single path weights does not repair this failure.
No weak-coupling decay rate or uniform physical-time
contraction is proved in this appendix.

\paragraph{Verification and preservation.}
The checker \path{scripts/verify_ym_restart_v17_block_counterterm.py}
tests the disintegration, inverse change of law,
four-bridge affinity factorization, literal periodic
path swap and its three-slice contraction in exact
rational positive two-state transfers.  It checks
the limiting floor and the response lower bound
independently.  Selected transfer entries have
rational square-root ratios, so these are exact
arithmetic tests, not floating-point evidence for
the continuum theory.  V16's native geometry and
analytic test-point checks are rerun as well.
V16's body is preserved apart from the displayed
title; its recovered SHA--256 is
\path{F40FB6521BC0C83067847D50C4C03513C9E60C10FF7C45882DFEAC6E0EAA1E3A}.
Only active V16 is replaced by V17.  The other 46
sources are protected, and build files remain
outside \path{latex_notes}.

\paragraph{Status.}
The temporal blocking normalization is now explicit,
with an all-coupling, depth-independent comparison
cost and a rigorous native obstruction to ignoring
it.  A useful weak-coupling estimate exploiting its
cancellation, the interacting continuum construction,
and the full Yang--Mills mass gap remain unproved.

% END CONSOLIDATED V17 APPEND

% BEGIN CONSOLIDATED V18 APPEND -- 2026-09-19
\clearpage
\section{V18: exact temporal collars and two-interface swap obstructions}
\label{r18:main}

V17 shows that a microscopic single-slice swap has a
nonzero limiting defect.  We now distinguish two
proposed repairs.  Swapping a longer temporal interval
only moves the mismatch to its two interfaces; we
compute its limiting affinity and prove that its
defect also persists.  Integrating out a temporal
collar before swapping is different: it gives a
monotone family of exactly normalized observables
ending at the desired purity.  We evaluate each
member as a positive quadratic form of the original
transfer, without a Perron approximation or an
unaccounted conditional partition function.

The collar construction does not by itself prove
weak-coupling decay.  Its fixed-regulator large-collar
defect is asymptotic to the entire excited heat trace,
including multiplicity.  This identifies exactly
what a useful estimate must control.  An exactly
solved product example quantifies the collar width
required even when the underlying gap is fixed.

\subsection{Scope, source check, and inherited methods}
\label{r18:scope}

Throughout, $T$ is the same real symmetric, positive
Wilson transfer at fixed $L\ge4$, $\gamma>0$;
spatial variables lie in $X=\SU(2)^N$, $N=3L^3$,
with probability Haar reference measure.  Temporal
links have been integrated exactly, not gauge-fixed
or replaced by a reference action.  Write
$\kappa_j=T^j$, $Z_j=\Tr T^j$, $m=2q\ge4$ and
$1\le r\le q$.  All spectral limits below are at
fixed regulator.  Positivity of $T$ as an operator
and strict positivity of its kernel are both used;
they are different assertions.

The nonduplication audit reread f16.2 V23's full-block
regeneration and retained-buffer warning, f15 V81's
signed exterior-square kernel, and V15--V17's
purity and affinity identities.  We do not repeat
the abstract conditional-resampling construction,
reinterpret its sampler clock as transfer time,
or discard the antisymmetric cancellation.
The results here concern two precisely specified
operations on the actual periodic temporal histories.

A targeted primary-source check read Appendix E,
Lemma 25 of Wolfer and Kontorovich,
\href{https://proceedings.mlr.press/v99/wolfer19a.html}
{\emph{Estimating the Mixing Time of Ergodic Markov Chains}}
(2019), which recalls Kazakos's entrywise-geometric-mean
formula for trajectory Hellinger affinities.
That elementary transfer calculation motivates our
two-interface operator below; the particular Wilson
formulas and their normalization are proved directly.
No statistical mixing-time estimator is imported as
a physical mass-gap bound.  Regular companion and
broad literature checkpoints remain due at V20.

\subsection{Erasing a collar: a normalized monotone hierarchy}
\label{r18:collar}

Begin with two independent spatial-slice histories
of period $2q$, with the exact law in
\eqref{r17:slice-law}.  In each history retain the
central slice $y=U_q$, and every slice outside the
open collar $q-r<j<q+r$; the central slice itself
is not erased.  Exchange the two retained central
slices and fix all other retained variables.
Let $\mathcal A_{q,r}$ be the affinity of the
resulting two marginal densities.  Thus $r=1$
erases nothing, whereas $r=q$ retains just
$U_0$ and $U_q$ from each history.  This operation
is marginalization followed by a swap, not a
swap of the collar contents.

\begin{theorem}[Exact collar quadratic form and monotonicity]
\label{r18:collar-thm}
For $g_{r;y,y'}(u)=\sqrt{\kappa_r(u,y)\kappa_r(u,y')}$,
\[
 \mathcal A_{q,r}
 =\frac1{Z_{2q}^2}\int_{X^2}
       \langle g_{r;y,y'},T^{2q-2r}g_{r;y,y'}\rangle^2
       \,dy\,dy'.
\tag{V18.1}\label{r18:quadratic}
\]
The power zero in this formula is the identity
operator, not a constant kernel.  In particular,
\[
 A_q^{\rm slice}=\mathcal A_{q,1}
 \le\mathcal A_{q,2}\le\cdots\le\mathcal A_{q,q}
 =R_{2q}\le1.
\tag{V18.2}\label{r18:monotone}
\]
Every quantity in \eqref{r18:quadratic} is a
finite-period integral of the original transfer.
\end{theorem}
\begin{proof}
For $r<q$, let $u=U_{q-r}$ and $v=U_{q+r}$.
The three retained slices have the normalized marginal
\[
 \rho_{q,r}(u,y,v)
  =\frac{\kappa_r(u,y)\kappa_r(y,v)
                  \kappa_{2q-2r}(v,u)}{Z_{2q}}.
\tag{V18.3}\label{r18:triple}
\]
All the remaining retained exterior variables have
a conditional law independent of $y$, given $u,v$.
Consequently the likelihood ratio for swapping the
two $y$ slices depends only on the two triples.
Integrating those exterior variables does not change
the affinity.

Take the geometric mean of
$\rho_{q,r}(u,y,v)\rho_{q,r}(u',y',v')$
and its central-slice-swapped density.
The first exterior segment contributes
\[
 \int g_{r;y,y'}(u)\kappa_{2q-2r}(u,v)
                       g_{r;y,y'}(v)\,du\,dv;
\]
the second contributes the same integral.
Their product proves \eqref{r18:quadratic}.
For $r=q$ the two exterior endpoints coincide,
$u=v=U_0$.  Then the integral is instead
\[
 \langle g,g\rangle
       =\int\kappa_q(u,y)\kappa_q(u,y')\,du
       =\kappa_{2q}(y,y'),
\]
whose squared integral is $Z_{4q}$.  This proves
the claimed endpoint and the identity-operator
interpretation at zero power.

The retained variable sets decrease with $r$ and
each marginalization commutes with the same central
swap.  Conditional Cauchy--Schwarz increases
affinity under each such marginalization.
This proves the complete monotone chain, including
its first term from V17.
\end{proof}

There is a bounded observable under the actual
two-history marginal.  For $r<q$, set
\[
 v_{q,r}=\frac12\log
 \frac{\kappa_r(u,y')\kappa_r(y',v)
       \kappa_r(u',y)\kappa_r(y,v')}
      {\kappa_r(u,y)\kappa_r(y,v)
       \kappa_r(u',y')\kappa_r(y',v')}.
\tag{V18.4}\label{r18:v}
\]
At $r=q$ put $u=v=x$, $u'=v'=x'$ in this expression.
Writing $\eta_{q,r}$ for the corresponding actual
pair marginal gives the exact identities
\[
 \E_{\eta_{q,r}}e^{2v_{q,r}}=1,\qquad
 1-\mathcal A_{q,r}
   =\E_{\eta_{q,r}}\left[1-\frac1{\cosh v_{q,r}}\right].
\tag{V18.5}\label{r18:bounded}
\]
Indeed $e^{2v}$ is its normalized swapped likelihood
ratio.  Pairing each point with its swap gives
$\E(1/\cosh v)=\E e^v=\mathcal A_{q,r}$, exactly
as in V15.  The observable lies in $[0,1]$.
This is not a new symmetrization lemma: it places
the inherited bounded estimator on the new collar
law.  V17's counterterm is avoided only because
the collar has been integrated \emph{before} the
interpolation.  Computing those $\kappa_r$ integrals
has not been made cost-free.

\subsection{Swapping a larger interval leaves two interfaces}
\label{r18:interval}

In contrast, keep every spatial slice and exchange
an interval of $\ell$ consecutive slices between the
two histories, where $1\le\ell\le m-1$.
Call its affinity $\mathcal W_{m,\ell}$.
For $z=(x,x')$, $w=(y,y')$ define the symmetric
integral operator $B$ by
\[
 B(z,w)=\sqrt{T(x,y)T(x',y')T(x,y')T(x',y)},
 \qquad D=T\otimes T.
\tag{V18.6}\label{r18:B}
\]
This is an entrywise square root, not the operator
square root of a transfer.  $B$ is Hilbert--Schmidt
on the compact configuration space; no assertion
that it is positive semidefinite is needed.

\begin{theorem}[Two-interface formula and persistent interval floor]
\label{r18:window-thm}
Exactly,
\[
 \mathcal W_{m,\ell}
   =\frac{\Tr\!\left[B D^{\ell-1}B D^{m-\ell-1}\right]}{Z_m^2}.
\tag{V18.7}\label{r18:window-trace}
\]
Let $a_P$ and $p$ be V17's native row affinity and
stationary density, and put
\[
 C=\int p(y)p(y')a_P(y,y')^2\,dy\,dy'.
\tag{V18.8}\label{r18:C}
\]
At fixed regulator, if both $\ell$ and $m-\ell$
tend to infinity, then
\[
 \mathcal W_{m,\ell}\longrightarrow C^2<1,
 \qquad C^2\le A_\infty^{\rm slice}.
\tag{V18.9}\label{r18:window-floor}
\]
Thus a growing swapped interval need not repair
even the fixed-regulator defect.  This is not the
collar erasure in \eqref{r18:monotone}.
\end{theorem}
\begin{proof}
Within the swapped interval, and within its
complement, the product of the two history weights
is unchanged.  At each of the two interfaces its
geometric mean with the swapped weight is precisely
$B(z,w)$.  There are $\ell-1$ interior edges on one
side and $m-\ell-1$ on the other.  Closing the two
paths and retaining their original partition
function proves \eqref{r18:window-trace}.

Let $K=T/\Lambda$, $\Psi=\psi\otimes\psi$ and
$G=B/\Lambda^2$.  At fixed regulator,
$(K\otimes K)^j$ converges in operator norm to
$|\Psi\rangle\langle\Psi|$, and $\Tr K^m\to1$.
In the normalized trace formula, both long powers
can therefore be replaced in the limit by that
projection.  The Hilbert--Schmidt norm of $G$
bounds the trace errors by the vanishing operator
norm errors.  The result is
$\langle\Psi,G\Psi\rangle^2$.
Substitution of
$T(x,y)=\Lambda\psi(x)P(x,y)/\psi(y)$
in the four factors gives
$\langle\Psi,G\Psi\rangle=C$.

If $C=1$, the affinities of almost every two
transition rows equal one, so those rows coincide.
Stationarity then makes $P$ the rank-one kernel
$P(y,du)=p(u)du$.  This contradicts injectivity on
the nontrivial physical subspace, as in V17.
Finally Jensen's inequality gives
$(\E_{p\otimes p}a_P^2)^2\le\E_{p\otimes p}a_P^4$,
the last expression being $A_\infty^{\rm slice}$.
\end{proof}

The assertion excludes neither exact conditional
resampling nor a suitable transport of the interface
variables.  It excludes this explicitly defined
unmodified interval swap.  The old warning about
retaining an unrefreshed buffer has an actual
transfer-time counterpart here.

\subsection{The collar defect sees the whole excited heat trace}
\label{r18:trace}

For $r\ge1$, define
\[
\begin{split}
 a_r(y,y')&=\int\sqrt{P^r(y,u)P^r(y',u)}\,du,\\
 A_r&=\int p(y)p(y')a_r(y,y')^4\,dy\,dy',
 \qquad D_r=1-A_r,\\
 S_j&=\Tr K^j-1=\sum_{k\ge1}\lambda_k^j,
 \quad 1>\lambda_1\ge\lambda_2\ge\cdots>0.
\end{split}
\tag{V18.10}\label{r18:stationary}
\]
Here the $\lambda_k$ are the nonvacuum eigenvalues
of the physical normalized transfer $K$.
The Gauss-complement zero modes add nothing to
these traces.  The same three-slice argument as
in V17, now with separation $r$, proves
$\lim_{q\to\infty}\mathcal A_{q,r}=A_r$ at each
fixed $r$.  This limit is not needed to define or
evaluate the finite-period formula \eqref{r18:quadratic}.

\begin{proposition}[Explicit finite-period remainder]
\label{r18:remainder}
With $\rho=\lambda_1$, $\tau_j=1+S_j$ and
$n=2q-2r\ge0$,
\[
 \frac{A_r}{\tau_{2q}^2}\le\mathcal A_{q,r}
 \le\min\left\{1,
 \frac{\left(\sqrt{A_r}+\rho^n\sqrt{\tau_{4r}}\right)^2}
      {\tau_{2q}^2}\right\}.
\tag{V18.11}\label{r18:remainder-eq}
\]
This bounds a specified limiting replacement;
it is not a computable gap certificate unless
the spectral quantities on the right are paid.
\end{proposition}
\begin{proof}
Replacing $T$ by $K$ in \eqref{r18:quadratic}
leaves the ratio unchanged.  Set
\[
 h(u)=\sqrt{K^r(u,y)K^r(u,y')}.
\]
On the physical subspace,
\[
 0\le\langle h,K^n h\rangle-\langle\psi,h\rangle^2
          \le\rho^n\|h\|_2^2.
\]
For $n=0$ the upper factor is interpreted as one.
The same bound holds if a Gauss-complement component
is included at $n=0$; at positive $n$ it is killed.
Direct substitution of the ground-state transform gives
\[
 \langle\psi,h\rangle^2
       =\psi(y)\psi(y')a_r(y,y')^2,\qquad
 \|h\|_2^2=K^{2r}(y,y').
\]
Their squared integrals are $A_r$ and $\tau_{4r}$,
respectively.  Positivity proves the lower bound,
and the $L^2(dy\,dy')$ triangle inequality proves
the upper bound in \eqref{r18:remainder-eq}.
\end{proof}

The next result gives a quantitative interpretation
of the remaining collar estimate, without a
dimension-dependent norm conversion.  Set
\[
 k_r(y,u)=\frac{P^r(y,u)}{p(u)},\qquad
 \varepsilon_r=\sup_{y,u}|k_r(y,u)-1|.
\tag{V18.12}\label{r18:eps}
\]

\begin{theorem}[Sharp leading trace and multiplicity]
\label{r18:trace-thm}
Whenever $\varepsilon_r<1$, let
$b(\varepsilon)=1+(1-\varepsilon)+(1-\varepsilon)^2
+(1-\varepsilon)^3$.  Then
\[
 \frac{b(\varepsilon_r)}{4(1+\varepsilon_r)}S_{2r}
       \le D_r\le\frac{S_{2r}}{1-\varepsilon_r}.
\tag{V18.13}\label{r18:trace-bounds}
\]
At every fixed native regulator,
\[
 \varepsilon_r\longrightarrow0,\qquad
 D_r\sim S_{2r}\sim d_1\lambda_1^{2r},\qquad
 -\frac{\log D_r}{2r}\longrightarrow-\log\lambda_1,
\tag{V18.14}\label{r18:trace-asymptotic}
\]
where $d_1$ is the finite multiplicity of $\lambda_1$.
No uniform bound on $\varepsilon_r$, $S_{2r}$,
$d_1$, or the limiting physical gap is inferred.
\end{theorem}
\begin{proof}
Write $e_y(u)=k_r(y,u)-1$ and let
$Q(y,y')=\int p(u)(e_y(u)-e_{y'}(u))^2du$.
Normalization of both rows gives the exact relation
\[
 1-a_r(y,y')
  =\frac12\int p(u)
    \frac{(e_y(u)-e_{y'}(u))^2}
    {(\sqrt{1+e_y(u)}+\sqrt{1+e_{y'}(u)})^2}\,du.
\]
The denominator is between $4(1-\varepsilon_r)$
and $4(1+\varepsilon_r)$, and
$1-\varepsilon_r\le a_r\le1$.
Using
$1-a_r^4=(1-a_r)(1+a_r+a_r^2+a_r^3)$
therefore bounds $1-a_r^4$ between
$b(\varepsilon_r)Q/[8(1+\varepsilon_r)]$
and $Q/[2(1-\varepsilon_r)]$.

Stationarity gives $\int p(y)e_y(u)dy=0$.
Consequently the $p(y)p(y')$ average of $Q$ is
\[
 2\int p(y)p(u)e_y(u)^2\,dy\,du
      =2\|P^r-\Pi\|_{\mathrm{HS}(L^2(p))}^2
      =2S_{2r},
\]
where $\Pi f=\int pf$.  The last equality follows
from unitary equivalence of $P$ and $K$.
This proves \eqref{r18:trace-bounds}.

At fixed regulator, sandwiching the operator-norm
convergence of $K^{r-2}$ between two continuous
kernel factors gives uniform convergence of
$K^r$ to $\psi\otimes\psi$.  The strictly positive
continuous $\psi$ has a positive minimum on the
compact space, so this convergence also gives
$\varepsilon_r\to0$.  Both coefficients in
\eqref{r18:trace-bounds} tend to one.
Finally trace class, compactness and dominated
convergence imply
$\sum_{k\ge1}(\lambda_k/\lambda_1)^{2r}\to d_1$;
one may dominate by the summable series at $r=1$.
This proves all the fixed-regulator asymptotics.
\end{proof}

This theorem prevents a circular shortcut.  The
unknown heat-trace estimate cannot be supplied by
assuming the decay of $\varepsilon_r$ uniformly
along the continuum trajectory.  At fixed regulator
that decay follows from a gap already known by
compactness; it provides neither its useful value
nor a bound uniform in the spatial volume.
The constructive object available without that
assumption remains \eqref{r18:quadratic} or its
bounded same-law estimator \eqref{r18:bounded}.

\subsection{Exactly solved collar costs and a finite-period target}
\label{r18:example}

For $0<\rho<1$, consider the illustrative two-state
transfer, with stationary law $(1/2,1/2)$,
\[
 P_\rho=\frac12
 \begin{pmatrix}1+\rho&1-\rho\\1-\rho&1+\rho\end{pmatrix}.
\]
Direct evaluation of \eqref{r18:quadratic} gives
\[
 \mathcal A_{q,r}
   =\frac12\left[1+
        \left(\frac{1-\rho^{2r}}{1+\rho^{2q}}\right)^2\right],
 \quad
 A_r=1-\rho^{2r}+\tfrac12\rho^{4r},\quad
 D_r=\rho^{2r}-\tfrac12\rho^{4r}.
\tag{V18.15}\label{r18:two-state}
\]
To verify it, the two equal-endpoint vectors $g$
have quadratic form $(1+\rho^{2q})/2$; the two
unequal-endpoint vectors have quadratic form
$(1-\rho^{2r})/2$.  Squaring, summing, and dividing
by $Z_{2q}^2=(1+\rho^{2q})^2$ proves the formula,
including $r=q$.  Also $S_{2r}=\rho^{2r}$,
so the leading coefficient in \eqref{r18:trace-asymptotic}
is one, not two.  In contrast the growing-interval
affinity tends to $(1-\rho^2/2)^2$, which remains
strictly below one.

Take $M$ genuinely independent copies of this
two-state system.  Factorization of the probability
law gives exactly
\[
 \mathcal A_{q,r}^{(M)}=\mathcal A_{q,r}^{,M},\qquad
 D_r^{(M)}=1-(1-t+t^2/2)^M,\quad t=\rho^{2r}.
\tag{V18.16}\label{r18:product}
\]
The transfer gap $-\log\rho$ is independent of $M$,
but its first excited multiplicity is $M$.  Elementary
Bernoulli and exponential bounds give
\[
 1-e^{-Mt/2}\le D_r^{(M)}\le Mt,\qquad
 D_r^{(M)}\le1-\mathcal A_{q,r}^{(M)}
        \le M(t+\rho^{2q})\le2Mt.
\tag{V18.17}\label{r18:volume}
\]
For the first pair, use $t/2\le t-t^2/2\le t$.
For the finite-period bound, \eqref{r18:two-state}
gives
\[
 1-\mathcal A_{q,r}
 =\frac{(t+\rho^{2q})(2+\rho^{2q}-t)}
             {2(1+\rho^{2q})^2}\le t+\rho^{2q};
\]
then use $1-(1-x)^M\le Mx$.
Thus stationary defect at most $\delta\in(0,1)$
requires $M\rho^{2r}\le-2\log(1-\delta)$;
$M\rho^{2r}\le\delta$ suffices for that stationary
defect, and $2M\rho^{2r}\le\delta$ suffices at
every $q\ge r$.  The resulting logarithmic
dependence of collar width on $M$ is a genuine
multiplicity cost even for a gapped product system.
This is an exact model diagnostic, not a factorization
of a connected Wilson gauge theory into its links.

We can now state the next concrete task without
renaming the unknown gap.  Estimate
$1-\mathcal A_{q,r}$ from its finite Wilson
integrals, with the coupling, spatial volume,
collar width and physical time all explicit.
An upper bound on this quantity is an upper bound
on $1-R_{2q}$ by \eqref{r18:monotone}.  The fixed
positive floors of interval swaps are not an
obstacle for this \emph{different} observable,
but neither the monotone hierarchy nor its
spectral interpretation supplies the needed
weak-coupling bound.  The product calculation
shows why testing only a single local mode could
miss the required multiplicity budget.

\paragraph{Verification and preservation.}
\path{scripts/verify_ym_restart_v18_temporal_collars.py}
checks the collar integral, its bounded estimator,
monotonicity, and purity endpoint with exact rational
two-state transfers.  Independent literal path sums
check the two-interface trace, including a nonstochastic
three-state positive transfer.  Product checks and
the trace-comparison constants are verified separately;
V17's selected algebra and V16's native analytic
test point are rerun.  These diagnostics are not
Wilson quadrature or formal verification of the proofs.
The recovered V17 source has SHA--256
\path{7581CBADB4A3BA88DFD6645EC1BD4BC9A1FD5BA0572E33F6033082A893357295}.
Its body is preserved apart from the displayed title.
Only active V17 is replaced by V18; the other 46
sources are unchanged and all build outputs remain
outside \path{latex_notes}.

\paragraph{Status.}
We have an exact, normalization-preserving temporal
collar hierarchy and an explicit evaluation formula,
as well as a proof that merely expanding the swapped
interval fails.  The sharp collar asymptotics retain
the full excited heat trace rather than a selected
source correlation.  A useful native weak-coupling
estimate, the interacting continuum construction,
and the full Yang--Mills mass gap remain unproved.

% END CONSOLIDATED V18 APPEND

% BEGIN CONSOLIDATED V19 APPEND -- 2026-09-19
\clearpage
\section{V19: positive plaquette truncations and certified purity transfer}
\label{r19:main}

V18 supplies an exactly normalized collar observable
but does not evaluate its interacting Wilson integrals.
We now give a local approximation which simultaneously
preserves pointwise positivity, positive transfer
spectrum, and exact Gauss projection.  Unlike a
global alternating action Taylor expansion, it has
a relative error before any integration.  Its
periodic partition functions are finite polynomial
Haar integrals.  A square-root-defect estimate then
transfers a certified truncated purity to the
original Wilson purity, including the change of law.

We evaluate the scalar error budget at a positive
coupling outside V16's strong-coupling comparison
range.  We do \emph{not} evaluate the corresponding
interacting truncated partition functions.  The
large joint integral, rather than the local scalar
tail, remains the bottleneck.  In particular no
native weak-coupling purity margin is certified.

\paragraph{Minor notation correction.}
The exponent printed as a comma followed by $M$
in \eqref{r18:product} denotes the ordinary power
$(\mathcal A_{q,r})^M$; the comma is typographical.
The product proof and its exact checker use this
ordinary power.  The V18 source is otherwise
retained literally.

\subsection{Nonduplication and the local positive polynomial}
\label{r19:local}

The audit reread V15's global Taylor enclosure,
f15 V79--V80's positive kinetic finite-rank Gram,
f15 V82's charge-walk Poisson tail, and f15 V84's
local Bessel-tail improvement.  Finite-rank Gram
certification, fusion positivity and Poisson tails
are inherited methods.  V82's charge polynomial
need not be a positive multiplier; a character
cutoff in V79--V84 need not be a pointwise positive
kernel.  Our approximation instead changes every
plaquette weight with a uniform \emph{relative}
bound and uses its square root only for the
magnetic half-action.  We do not claim its cutoff
is smaller than the archived Bessel cutoff.

A targeted source check read the Euclidean transfer
factorization and character-diagonal kinetic
construction in Section II of Kanwar and Wagman,
\href{https://arxiv.org/abs/2103.02602}
{\emph{Real-time lattice gauge theory actions:
unitarity, convergence, and path integral contour
deformations}} (2021).  Only that standard
Euclidean mechanism is used here; no real-time
claim or continuum conclusion is imported.
The construction and constants below are proved
directly.  The regular broad literature and
companion checkpoints are still next at V20.

Write $t(U)=\tfrac12\Tr U\in[-1,1]$ and define
\[
 p_d(x)=\sum_{k=0}^d\frac{x^k}{k!},\qquad
 w_{\gamma,d}(U)=p_d\bigl(\gamma(1+t(U))\bigr).
\tag{V19.1}\label{r19:poly}
\]
The comparison weight is
$w_\gamma^+(U)=e^{\gamma(1+t(U))}$.
Its ratio to the original Wilson plaquette
weight is the constant $e^{2\gamma}$.
No unknown vacuum normalization has been used.

\begin{lemma}[Uniform relative tail, with a rational certificate]
\label{r19:tail}
Let $\lambda=2\gamma$ and
$\delta_d=\mathbb P\{\operatorname{Poisson}(\lambda)>d\}$.
Then, for every $U$,
\[
 1\le w_{\gamma,d}(U),\qquad
 (1-\delta_d)w_\gamma^+(U)
       \le w_{\gamma,d}(U)\le w_\gamma^+(U).
\tag{V19.2}\label{r19:relative}
\]
If $d+2>\lambda$, put
\[
 s_d=p_d(\lambda),\quad
 u_d=\frac{\lambda^{d+1}}{(d+1)!}
                       \frac1{1-\lambda/(d+2)},\quad
 \overline\delta_d=\frac{u_d}{s_d+u_d}.
\tag{V19.3}\label{r19:rational-tail}
\]
Then $\delta_d\le\overline\delta_d<1$.
These last quantities are exact rationals when
$\gamma$ is rational.
\end{lemma}
\begin{proof}
For $x\ge0$, $e^{-x}p_d(x)$ is the probability
that a Poisson variable of mean $x$ is at most $d$.
This probability decreases with $x$: differentiating
gives $-e^{-x}x^d/d!$.  Since
$0\le\gamma(1+t(U))\le\lambda$, this proves
\eqref{r19:relative}, including the elementary
lower bound $p_d(x)\ge1$.
The positive tail $e^\lambda-s_d$ has first term
$\lambda^{d+1}/(d+1)!$ and subsequent ratios at
most $\lambda/(d+2)$.  It is at most $u_d$.
The increasing function $u\mapsto u/(s_d+u)$
therefore gives \eqref{r19:rational-tail}.
\end{proof}

For orientation, the classical Chernoff calculation
also gives, for $a\ge0$,
\[
 \mathbb P\{X\ge\lambda+a\}
       \le\exp\left[-\frac{a^2}{2(\lambda+a/3)}\right].
\tag{V19.4}\label{r19:chernoff}
\]
Here is enough detail to fix its constants.
For $0<s<3$,
$\log\E e^{s(X-\lambda)}=\lambda(e^s-1-s)
\le\lambda s^2/[2(1-s/3)]$;
the series inequality follows from
$k!\ge2\cdot3^{k-2}$ for $k\ge2$.
Markov's inequality with
$s=a/(\lambda+a/3)$ proves \eqref{r19:chernoff},
with $a=0$ handled directly.
Thus, for $H=\log(1/\delta)>0$, it suffices to choose
\[
 d+1\ge\lambda+\sqrt{2\lambda H}+\tfrac23 H
\tag{V19.5}\label{r19:degree}
\]
to ensure $\delta_d\le\delta$.
This is a local degree, not the total degree of
the complete many-plaquette polynomial.

\subsection{Positive character coefficients and the exact physical Gram}
\label{r19:gram}

Label $\SU(2)$ characters by $n=2j$, so that
$\chi_0=1$, $\chi_1=2t$, and the dimension is
$n+1$.  Let $c_{d,n}$ be the coefficients in
\[
 w_{\gamma,d}(U)=\sum_{n=0}^d c_{d,n}\chi_n(U).
\tag{V19.6}\label{r19:char}
\]

\begin{proposition}[Strict character positivity on the retained range]
\label{r19:char-positive}
For $\gamma>0$,
\[
 c_{d,n}\ge\frac{\gamma^n}{2^n n!}>0
       \quad(0\le n\le d),\qquad c_{d,n}=0\quad(n>d).
\tag{V19.7}\label{r19:char-lower}
\]
One-link convolution by $w_{\gamma,d}$ is
positive semidefinite with eigenvalue
$c_{d,n}/(n+1)$ on the spin-$n/2$ matrix
coefficients.  Its rank is
\[
 h_d=\sum_{n=0}^d(n+1)^2
       =\frac{(d+1)(d+2)(2d+3)}6.
\tag{V19.8}\label{r19:rank-one}
\]
\end{proposition}
\begin{proof}
Each power $(1+\chi_1/2)^k$ is a nonnegative
linear combination of characters, by the tensor
product rule $\chi_1\chi_n=\chi_{n+1}+\chi_{n-1}$
with $\chi_{-1}=0$.  Its support is $n\le k$.
The term $k=n$, taking $\chi_1/2$ in every factor,
contains $\chi_n$ with coefficient $2^{-n}$.
This proves the lower bound and finite support.
Character orthogonality gives the convolution
eigenvalues and their multiplicities $(n+1)^2$.
Summing the multiplicities proves the rank formula.
\end{proof}

Fix the original spatial three-torus with $L\ge4$
and $N=3L^3$ links and spatial plaquettes.
Let $\mathcal C_d$ be the tensor product of
these one-link convolutions, $\mathsf P$ the
\emph{exact} spatial Gauss projection, and
\[
 M_d(U)=\prod_{p\ {\rm spatial}}w_{\gamma,d}(U_p),
 \qquad b_d(U)=\sqrt{M_d(U)},\qquad
 T_d=M_{b_d}\mathcal C_d\mathsf P M_{b_d}.
\tag{V19.9}\label{r19:Td}
\]
The notation $M_f$ means multiplication by $f$;
thus $M_d$ in \eqref{r19:Td} is a function, not
an operator or a count of copies.

\begin{theorem}[Simultaneous kernel positivity and finite positive spectrum]
\label{r19:transfer}
The operator $T_d$ is self-adjoint, positive
semidefinite, finite rank, and has a strictly
positive continuous kernel.  It acts on the
original physical space and is zero on the
Gauss complement.  Its rank is at most $h_d^N$.
If $\mathcal E_d$ is the full tensor Peter--Weyl
space with $n_e\le d$ on each link, its nonzero
spectrum, with multiplicity, is the spectrum
of the nonzero part of
\[
 G_d=\left.
 \mathsf P\mathcal C_d^{1/2}M_{M_d}
                 \mathcal C_d^{1/2}\mathsf P
           \right|_{\mathcal E_d}.
\tag{V19.10}\label{r19:Gram}
\]
For rational $\gamma$ the Gram entries can be
constructed by finite polynomial Haar integration
with algebraic coefficients.  In particular
$\Tr G_d^m=\Tr T_d^m$ is a positive rational
number for every integer $m\ge2$.
\end{theorem}
\begin{proof}
The central one-link convolutions commute with
left and right translations and hence with
$\mathsf P$.  The gauge-invariant $b_d$ also
commutes with $\mathsf P$.  All factors preserve
the physical space.  Put
$F=M_{b_d}\mathcal C_d^{1/2}\mathsf P$
on $\mathcal E_d$.  Then $FF^*=T_d$ and
$F^*F=G_d$, proving positivity, finite rank
and the equality of nonzero spectra.
This argument does not remove the Gauss projection.

Before temporal-link integration the kernel is
$b_d(U)b_d(V)$ times a product of $N$ strictly
positive temporal plaquette polynomials.
Its exact Haar gauge average remains strictly
positive and continuous.  Finite rank therefore
does not imply a strictly positive lower spectral
bound on the whole infinite-dimensional physical
space: the approximant has many zero modes.

Standard spin-$n/2$ matrix elements are finite
polynomials in link quaternions with algebraic
coefficients.  The finite gauge averages in
$\mathsf P$ are polynomial Haar integrals as well.
Crucially $b_d$ occurs squared in $G_d$, leaving
the polynomial $M_d$ rather than its nonpolynomial
square root.  Normalized $S^3$ Haar moments are
zero if any exponent is odd; otherwise
\[
 \int_{S^3}\prod_{j=0}^3 a_j^{2k_j}\,da
   =\frac1{(K+1)!}
       \prod_{j=0}^3\frac{(2k_j)!}{4^{k_j}k_j!},
 \qquad K=\sum_jk_j.
\tag{V19.11}\label{r19:Haar}
\]
This follows, for example, by separating radius
from direction in four independent centered
Gaussian integrals.  It proves the assertion
about algebraic Gram entries.

For the stronger rational trace assertion, use
the full periodic path integral instead.  In a
cycle every $b_d$ occurs twice, so the integrand
is exactly the product of $w_{\gamma,d}$ over
all spatial and temporal plaquettes.  At rational
$\gamma$ this is a finite polynomial with rational
quaternion coefficients.  Formula \eqref{r19:Haar}
then gives a rational trace, strictly positive
because the integrand is positive everywhere.
\end{proof}

This is not a finite subgroup replacement and not
a deletion of physical boundary sectors.
All omitted electric modes will be paid for by
the comparison with the original Wilson law.
Nonnegative character coefficients also do not
assert that an arbitrary recoupled tensor-network
expansion has no cancellations.

\subsection{Relative bounds survive blocks and periodic normalization}
\label{r19:comparison}

Let $T^+$ be the transfer obtained with
$w_\gamma^+$ and its spatial square-root weight.
It differs from the original $T$ only by
\[
                         T^+=e^{4\gamma N}T.
\tag{V19.12}\label{r19:scalar}
\]
The factor counts $N$ spatial and $N$ temporal
plaquettes per time step.  It cancels from
every normalized history and purity.
Choose any certified $\delta\ge\delta_d$,
$\delta<1$, and put $\eta=-\log(1-\delta)$.
Pointwise, \eqref{r19:relative} gives
\[
 e^{-2N\eta}T^+(U,V)\le T_d(U,V)\le T^+(U,V).
\tag{V19.13}\label{r19:kernel-order}
\]
The two magnetic half-factors together cost
$N\eta$, and the temporal factors cost $N\eta$.
These are inequalities between kernel values,
not assertions of operator order.
Integrating products of positive kernels proves
\[
 e^{-2Nm\eta}(T^+)^m(U,V)
       \le T_d^m(U,V)\le(T^+)^m(U,V),
\tag{V19.14}\label{r19:block-order}
\]
and the same bounds for their diagonal integrals.
No inverse of an exponentially small partition
function is used to convert an absolute error.

Let $\mu$ and $\mu_d$ now be the laws of two
independent full periodic histories of length
$2q$, for Wilson and for the polynomial model.
There are $8Nq$ plaquette factors in the pair.
The ratio of their \emph{unnormalized} densities
therefore lies in $[e^{-B},1]$, where
\[
                           B=8Nq\eta.
\tag{V19.15}\label{r19:B}
\]
After normalizing, its logarithmic oscillation
is still at most $B$.  The same statement holds
after taking any common marginal: its unnormalized
ratio is a conditional average of the original
ratio.  Thus every V18 collar marginal retains
this bound, including the change of its
conditional partition functions.

\begin{theorem}[Native collar and purity certification with a paid law change]
\label{r19:purity}
Let $\mathcal A_{q,r}^{(d)}$ be the exact V18
collar affinity for $T_d$.  Uniformly for
$1\le r\le q$,
\[
 \left|\sqrt{1-\mathcal A_{q,r}}
             -\sqrt{1-\mathcal A_{q,r}^{(d)}}\right|
 \le 2\sqrt{1-\frac1{\cosh(B/4)}}
 \le 2\sqrt2\,Nq\eta.
\tag{V19.16}\label{r19:sqrt-bound}
\]
In particular, set
\[
 R_{2q}^{(d)}=\frac{\Tr G_d^{4q}}{(\Tr G_d^{2q})^2}.
\tag{V19.17}\label{r19:finite-purity}
\]
If $b\ge0$ is certified to satisfy
$1-R_{2q}^{(d)}\le b^2$, then the original
Wilson purity obeys the entirely explicit bound
\[
 R_{2q}\ge
 \max\left\{0,1-
      \left(b+\frac{3Nq\delta}{1-\delta}\right)^2\right\}.
\tag{V19.18}\label{r19:rational-export}
\]
For rational $\gamma,\delta,b$, the input
\eqref{r19:finite-purity} and the exported lower
bound can both be certified by rational arithmetic.
\end{theorem}
\begin{proof}
Use the two corresponding collar marginals,
with density square roots $f$ and $g$, and let
$S$ be the unitary involution swapping their
central slices.  Each square-root affinity defect
equals $\|(I-S)f\|_2/\sqrt2$ or
$\|(I-S)g\|_2/\sqrt2$.  Their difference in
absolute value is at most $\sqrt2\|f-g\|_2$.
V17's normalized likelihood-oscillation lemma
and \eqref{r19:B} give
$\langle f,g\rangle\ge1/\cosh(B/4)$.
Since $\|f-g\|_2^2=2(1-\langle f,g\rangle)$,
the first bound follows.  V15's Lipschitz estimate
for $\sqrt{1-1/\cosh x}$ gives the second bound,
namely $B/(2\sqrt2)=2\sqrt2Nq\eta$.

At $r=q$, the two collar affinities are their
respective trace purities, by V18 and the
finite Gram identity.  The upper triangle
inequality proves the export with
$2\sqrt2Nq\eta$; finally
$2\sqrt2<3$ and $\eta\le\delta/(1-\delta)$
give \eqref{r19:rational-export}.
The proof compares the actual two laws.  It
does not evaluate a polynomial-model observable
under the Wilson law without reweighting.
\end{proof}

If a useful lower purity is obtained, V14's
all-sector spectral conversion remains available.
No conclusion about the original gap follows from
the spectrum of $G_d$ alone.  As a simple warning,
$d=0$ gives $w_{\gamma,0}=1$ and a rank-one
transfer with purity one at every period; its
tail is $1-e^{-2\gamma}$.  At large coupling
the error in \eqref{r19:rational-export} is
vacuous, correctly preventing that artificial
rank-one spectrum from proving a Wilson gap.

\subsection{An evaluated accuracy budget, not an evaluated gap}
\label{r19:budget}

Take the actual lattice size and coupling
\[
 L=4,\quad N=192,\quad \gamma=10,\quad
 q=12,\quad d=50.
\tag{V19.19}\label{r19:point}
\]
This lies outside $18\tanh\gamma<1$.
The rational tail formula with $\lambda=20$
gives the exact certified inequalities
\[
 \overline\delta_{50}<\frac1{200000000},\qquad
 \frac{3Nq\overline\delta_{50}}
                  {1-\overline\delta_{50}}<\frac1{25000}.
\tag{V19.20}\label{r19:budget-values}
\]
For scale only, the rational tail is approximately
$4.8624\times10^{-9}$ and the last error is
approximately $3.3609\times10^{-5}$; the decimal
values are not used as certificates.
Consequently a \emph{separately proved} bound
$R_{24}^{(50)}\ge99/100$ would imply
\[
 R_{24}\ge1-\left(\frac1{10}+\frac1{25000}\right)^2
           =\frac{618744999}{625000000}.
\tag{V19.21}\label{r19:conditional-point}
\]
The premise $R_{24}^{(50)}\ge99/100$ has not
been established here.  We have evaluated the
approximation error, not substituted a conjectured
partition-function ratio for measured data.

There is a serious dimension cost.  At this
cutoff $h_{50}=45526$; the ambient tensor
Peter--Weyl basis before Gauss reduction has size
\[
                  45526^{192},\qquad
        10^{894}<45526^{192}<10^{895}.
\tag{V19.22}\label{r19:dimension}
\]
This is a stated ambient construction size,
not a lower bound on every possible algorithm
or on the dimension after physical reduction.
The local degree has become manageable; the
uncontracted global Gram certainly has not.
An exact tensor contraction, a sparse algebraic
bound, or a connected estimate must exploit
additional structure before this can deliver a
useful interacting purity.

Nor should the chosen local polynomial be advertised
as a square-root-in-coupling cutoff.  If it is to
have relative error at most $\delta<1$ at $U=I$,
then $\mathbb P\{X\le d\}\ge1-\delta$ for
$X\sim\operatorname{Poisson}(2\gamma)$.
When $d<\lambda=2\gamma$, the one-sided variance
inequality gives
\[
 1-\delta\le\frac{\lambda}{\lambda+(\lambda-d)^2},
 \qquad
 d\ge\lambda-\sqrt{\frac{\lambda\delta}{1-\delta}}.
\tag{V19.23}\label{r19:degree-cost}
\]
For completeness, if $Y=X-\lambda$ and $a>0$,
Markov's inequality applied to $(Y-\lambda/a)^2$
on $\{Y\le-a\}$ yields
$\mathbb P(Y\le-a)\le\lambda/(\lambda+a^2)$.
This proves the displayed necessary condition.
It concerns this positive shifted Taylor family,
not a lower bound for all polynomial approximations.
The archived Bessel cutoff controls a different
operator error and is not contradicted.

\paragraph{What changed and what remains.}
V15's existence of finite-regulator enclosures
has been refined to a specific local, positive,
gauge-exact approximation and a paid two-law
purity transfer.  No exponentially small global
kernel floor enters the local relative-error
conversion.  However, a small certified
$1-R_{2q}^{(d)}$ is still missing.  The next
iteration should attack its connected joint
polynomial Haar integral or an equivalent
finite physical Gram estimate, and the V20
literature checkpoint should prioritize methods
that can bound that object without assuming
the unknown Wilson gap.

\paragraph{Verification and preservation.}
\path{scripts/verify_ym_restart_v19_positive_plaquettes.py}
checks character coefficients by fusion and
independent rational Haar moments, convolution
spectra on compact-group test points, normalized
law-change inequalities on finite exact fixtures,
and the rational budget \eqref{r19:budget-values}.
Selected V18 identities and the earlier native
strong-coupling arithmetic are rerun.
No full four-dimensional Gram contraction or
Wilson Monte Carlo is represented as performed.
The V18 body is preserved apart from its title;
the recovered SHA--256 is
\path{1586B7C2390D630BCB218607246689735B8E63A5E9910FAA13FBAED88870ECE0}.
Only active V18 is replaced.  The other 46
sources remain unchanged, and build artifacts
are outside \path{latex_notes}.

\paragraph{Status.}
The all-coupling approximation and its scalar
accuracy certificate are proved.  A useful
weak-coupling native purity bound, continuum
construction, physical-time transport and the
full Yang--Mills mass gap remain unproved.

% END CONSOLIDATED V19 APPEND

% BEGIN CONSOLIDATED V20 APPEND -- 2026-09-19
\clearpage
\section{V20: exact positive cubature and shared-link contraction certificates}
\label{r20:main}

V19 made the local plaquette approximation error small,
but left the interacting partition-function ratio unevaluated.
The present required literature checkpoint changes the
computational representation of that remaining problem.
We construct a positive, degree-exact cubature on the
actual compact group.  Applied link by link, it gives
an exact finite positive factor graph for the polynomial
model.  More strongly, its spatial transfer matrix has
exactly the nonzero physical spectrum of V19's $T_d$.
Thus it need not be assembled as a dense Peter--Weyl Gram.
There is no assertion that the resulting contraction is cheap.

We then prove a pointwise message-error certificate for
that graph and evaluate a six-plaquette shared-link
integral at $\gamma=10$.  This last calculation rules
out uncharged independent-factor splitting even locally.
It does not evaluate the global purity or close the gap.
All statements below use $\SU(2)$, spatial side $L\ge4$,
and, for ordinary periodic spacetime incidence counting,
time period $m\ge4$.  The purity application has $m=2q$,
$q\ge2$, and $N=3L^3$ spatial links.

\subsection{Scheduled literature, companion and nonduplication audit}
\label{r20:review}

The broad primary-source sweep considered the following
eight method families.  It is a targeted research review,
not an exhaustive survey or an endorsement of search hits
claiming a complete Yang--Mills solution.
\begin{enumerate}
\item \emph{Finite-$N$ lattice bootstrap.}
\href{https://arxiv.org/abs/2404.16925}{Kazakov--Zheng}
and \href{https://arxiv.org/html/2502.14421v2}
{Guo--Li--Yang--Zhu} combine Wilson-loop equations with
positivity.  The latter's Section 3 was inspected.
These remain possible sources of certified constraints,
not a supplied all-sector trace-purity margin.  Active
V3 already contains the polynomial-moment interface;
we do not repackage it as a new result.
\item \emph{Euclidean correlator bootstrap.}
\href{https://arxiv.org/html/2511.08560v3}{Cho et al.},
Section 2 and the ground-state constraints, were revisited.
Their correlator formulation keeps equations of motion
and state positivity together.  A bounded collection of
source moments is not automatically the complete physical
trace required here, and no new native moments are imported.
\item \emph{Certified equilibrium hierarchies.}
\href{https://arxiv.org/html/2311.18706v2}
{Fawzi--Fawzi--Scalet}, Theorems 1.1 and 1.5 and their
assumptions, distinguish convergence from rapid convergence.
The quantitative theorem has commuting and dimension/regime
restrictions.  We do not assume that $-\log T_d$ is a
finite-range Hamiltonian merely because a spacetime
factorization exists.
\item \emph{Tensor contraction.}
The \href{https://arxiv.org/abs/2206.07044}
{Gray--Chan} abstract and method description motivate
optimizing contraction order and compression.  Such
approximations are candidate generators; a discarded
singular weight alone is not our partition-ratio certificate.
The positive, observable-specific bounds below must be paid.
\item \emph{Positive cubature.}
\href{https://people.math.ethz.ch/~jteichma/tchakaloff120405.pdf}
{Bayer--Teichmann}, Theorem 2 and its moment-map proof,
give the classical positive finite-moment representation.
We use that viewpoint and prove an explicit sphere rule
instead of relying on a nonconstructive node choice.
\item \emph{Graphical-model elimination.}
\href{https://icml.cc/2011/papers/459_icmlpaper.pdf}
{Liu--Ihler}, Sections 2.2--3, give min/max and
weighted H\"older partition bounds.  Their elementary
positive-factor mechanism transfers to our exact cubature
graph.  Its quality at weak coupling must still be checked;
the shared-link test below supplies such a check.
\item \emph{Stochastic lattice Yang--Mills.}
The stated strong-coupling hypotheses in
\href{https://arxiv.org/abs/2204.12737}
{Shen--Zhu--Zhu} were rechecked.  Their functional-inequality
route does not remove the bare influence-row obstruction
at the present large $\gamma$.
\item \emph{Spectral renormalization.}
The semigroup Feshbach construction of
\href{https://arxiv.org/abs/2508.07643}
{Bach--Ballesteros--Geisler} remains relevant downstream.
The complementary inverse needed for a native Wilson
application has not been acquired; the present exact
finite representation is not such an inverse estimate.
\end{enumerate}
Only the cubature and positive-elimination mechanisms
are adopted in this iteration.  Their application is
proved below rather than inferred from a numerical
success in a different dimension or theory.
The next regular broad sweep is V30.

The compulsory companion review rehashed and reread
the terminal arguments of NS V26, SM V72, shell-mixing
V16 and parallel YM V5.  All four are unchanged since
the V15 audit.  NS's Oseen derivative norms do not
estimate this compact conditional law.  SM's many-copy
transport fixes the physical mode set.  Shell-mixing's
flat one-loop coefficient explicitly leaves the common
nonperturbative coupling/volume remainder open.  Parallel
YM's global-charge tensorization warning does not prove
connected-torus anti-concentration.  None supplies the
missing purity or justifies dropping a physical sector.
The next companion review is V25.

For nonduplication we reread f14 V21's rational
24-point angular rule, f15 V79--V80's physical finite
Gram and relative trace remainder, and active V19.
The final audit also reread f6 V28's cancelling
$\SU(3)$ staple star and f13 V11's exact $\SU(2)$
single-link conditional law.  Staple cancellation
and the one-link Bessel normalization are inherited;
their existence is not a new obstruction claimed here.
Degree-exact cubature, finite-rank approximation and
positive-message inequalities are not claimed as new
general methods.  The present additions are the
degree-$6d$ native incidence budget, the exact
Gauss-preserving spectral realization, its positive
local factor graph, and the explicit conditional
weak-coupling stress test.  In particular f14 V21's
quartic angular calculation is not being repeated.

\subsection{An explicit positive sphere rule at every required degree}
\label{r20:cubature}

Let $Q_D$ denote a cubature for normalized Haar measure
on $\SU(2)=S^3$, exact for real coordinate polynomials
of total degree at most $D$.  Complex polynomials then
follow by linearity.  Such a rule can be chosen with
positive algebraic weights and algebraic nodes.

\begin{theorem}[A constructive degree-exact positive rule]
\label{r20:Hopf}
Set $M=D+1$ and $n=\lfloor D/4\rfloor+1$.
Take the $n$-point Gaussian rule $(x_j,\omega_j)$
for Lebesgue probability measure on $[0,1]$.
For $0\le r,s<M$, use the nodes
\[
 u_{jrs}=\bigl(\sqrt{x_j}\cos\theta_r,
       \sqrt{x_j}\sin\theta_r,
       \sqrt{1-x_j}\cos\theta_s,
       \sqrt{1-x_j}\sin\theta_s\bigr),\qquad
 \theta_r=2\pi r/M,
\tag{V20.1}\label{r20:nodes}
\]
with weights $\omega_j/M^2$.
These $nM^2$ nodes define $Q_D$.
For $D=6d$ the count is
\[
                  (\lfloor3d/2\rfloor+1)(6d+1)^2.
\tag{V20.2}\label{r20:node-count}
\]
\end{theorem}
\begin{proof}
Haar measure in these coordinates is
$dx\,d\theta\,d\phi/(2\pi)^2$.
For example, normalize two independent complex
Gaussian coordinates: their phases are independent
uniform angles, their squared radii are independent
exponentials, and the ratio of one squared radius
to their sum is uniform on $[0,1]$.

The equally spaced $M$-angle rule integrates every
Fourier mode of frequency of absolute value at most
$D$ exactly, since the only multiple of $M$ in that
range is zero.  In a coordinate monomial of degree
at most $D$, each angular factor is such a trigonometric
polynomial.  Its integral vanishes if one exponent
in its sine/cosine pair is odd.  Otherwise the angular
integrals leave a multiple of
$x^{(a+b)/2}(1-x)^{(c+e)/2}$, of degree at most
$\lfloor D/2\rfloor\le2n-1$.
Gaussian quadrature integrates this exactly.

For completeness, the Gaussian assertion used here
does not require a numerical quadrature hypothesis.
Let $P_n$ be the degree-$n$ orthogonal polynomial for
$dx$ on $[0,1]$.  Its $n$ roots are distinct and
interior: otherwise the product of its interior
sign-changing linear factors would have degree less
than $n$ and its product with $P_n$ would have constant
nonzero sign, contradicting orthogonality.
With $\ell_j$ the Lagrange polynomials at these roots,
put $\omega_j=\int_0^1\ell_j(x)\,dx$.
Division by $P_n$ and orthogonality prove exactness
through degree $2n-1$.  Applying this to $\ell_j^2$
gives $\omega_j=\int\ell_j^2>0$, and applying it to
one gives $\sum\omega_j=1$.
The shifted Legendre coefficients are rational,
so their roots and weights are algebraic.
Roots of unity and the positive square roots in
\eqref{r20:nodes} are algebraic as well.
\end{proof}

There is a useful small rule for $d=1$.
Take the following 48 points, each with weight $1/48$:
\[
 \{\pm e_i\},\qquad
 \{(\pm1,\pm1,\pm1,\pm1)/2\},\qquad
 \{(\pm e_i\pm e_j)/\sqrt2:i<j\}.
\tag{V20.3}\label{r20:48}
\]
Here the signs in each set are independent.
Sign symmetry kills odd coordinate moments.
Permutation symmetry reduces the even moments
through degree six to
\[
 \begin{array}{c|rrrrrr}
 \text{monomial}&u_1^2&u_1^4&u_1^2u_2^2&
                    u_1^6&u_1^4u_2^2&u_1^2u_2^2u_3^2\\\hline
 Q&1/4&1/8&1/24&5/64&1/64&1/192 .
 \end{array}
\tag{V20.4}\label{r20:48-moments}
\]
Directly counting the three sets gives these values,
which equal \eqref{r19:Haar}.  This proves exactness
through degree seven.  It must not be used at arbitrary
degree: $Q[u_1^8]-\int u_1^8=1/256$.
The first two sets alone are the archived 24-point rule;
the present third set and weights pay the higher degree.

\subsection{The finite positive graph has the exact physical spectrum}
\label{r20:spectrum}

Keep V19's $w_{\gamma,d}$, $M_d$, $b_d$, $T_d$,
$\mathcal E_d$, and exact Gauss projector $\mathsf P$.
Write
\[
 K_d(U,V)=(\mathcal C_d\mathsf P)(U,V),\qquad
                      T_d(U,V)=b_d(U)K_d(U,V)b_d(V).
\tag{V20.5}\label{r20:K}
\]
Let $\mathcal X$ be the tensor product of $Q_{6d}$
over the $N$ spatial links, with positive product
weights $\nu_x$.  Write the one-link rule as
$Q_{6d}=\sum_a\omega_a\delta_{u_a}$.
Define the real symmetric matrix
\[
 (D_d)_{xy}=\sqrt{\nu_x}\,b_d(x)K_d(x,y)b_d(y)\sqrt{\nu_y}.
\tag{V20.6}\label{r20:D}
\]
Its entries are strictly positive.  The notation
$D_d$ is a finite matrix, not the cubature degree.

\begin{theorem}[Exact positive discrete realization]
\label{r20:exact-spectrum}
The matrix $D_d$ is positive semidefinite and its
nonzero spectrum, including multiplicities, equals
the nonzero spectrum of the original physical
polynomial transfer $T_d$ and of $G_d$.
Consequently
\[
 \Tr D_d^m=\Tr T_d^m=\Tr G_d^m,\qquad
 R_{2q}^{(d)}=\frac{\Tr D_d^{4q}}{(\Tr D_d^{2q})^2}.
\tag{V20.7}\label{r20:traces}
\]
The trace identity holds for every positive integer
$m$.  For ordinary periodic lattices with $m\ge4$
these traces are also the exact positive factor sums
\[
 Z_m^{(d)}=
   \sum_{\{a_e\}}
       \prod_e\omega_{a_e}
       \prod_p w_{\gamma,d}
        \bigl(u_{a_{e_1}}u_{a_{e_2}}
                         u_{a_{e_3}}^{-1}u_{a_{e_4}}^{-1}\bigr),
\tag{V20.8}\label{r20:graph}
\]
with the inherited plaquette orientations and one
node index for each spacetime link.  No physical
charge sector has been discarded.
\end{theorem}
\begin{proof}
A spin-$n/2$ matrix element is a polynomial of
degree $n$ in link quaternion coordinates.
Both $\mathcal C_d^{1/2}$ and $\mathsf P$ preserve
$\mathcal E_d$, because linkwise left and right
translations preserve each representation.
Thus every column $h_\alpha$ of
$\mathcal C_d^{1/2}\mathsf P|_{\mathcal E_d}$
has degree at most $d$ in each spatial link.

In three spatial dimensions each link belongs to
four spatial plaquettes.  Hence $M_d=b_d^2$ has
degree at most $4d$ in each link.  The Gram integrand
$M_d\overline h_\alpha h_\beta$ has degree at most
$6d$ in each link, irrespective of its total degree
over all links.  Therefore the matrix
\[
 A_{x\alpha}=\sqrt{\nu_x}\,b_d(x)h_\alpha(x)
\]
satisfies exactly $A^*A=G_d$ by tensor cubature.
The kernel feature identity gives $AA^*=D_d$.
Equality of nonzero spectra of these two products
proves the first assertions, including positivity.
V19 already proved $G_d=F^*F$ and $T_d=FF^*$.

One may evaluate $K_d(x,y)$ without a recoupled
signed tensor expansion: integrate its positive
temporal plaquette product over one gauge variable
at each spatial vertex.  That variable occurs in
six link factors and to degree at most $6d$.
The same positive rule evaluates this gauge average
exactly, not approximately.  Thus its strictly
positive kernel and its exact Gauss projection
are both retained in \eqref{r20:D}.

Alternatively start with the full periodic path
integral.  Every plaquette trace is multilinear in
the four distinct link quaternions; inversion is
linear on $\SU(2)$ quaternions.  A spacetime link
belongs to six plaquettes, so the full integrand
has degree at most $6d$ in each individual link.
Repeated cubature gives \eqref{r20:graph}.  Each
factor is at least one and each node weight is
positive.  The count is independent of the time
period.  This proves the positive local realization
of the same trace, without replacing a connected
lattice by independent plaquettes.
\end{proof}

\paragraph{What exactness does not say.}
The nodes are not asserted to form a group.  The
discrete measure is not invariant under arbitrary
continuous gauge transformations, and the Gauss
projector must not be replaced by a node permutation
average.  Gauge information is retained by the
exact polynomial integration and the Gram identity.
Nor are the continuous and atomic history laws close
in total variation: they are mutually singular.
One must not insert the cubature into an arbitrary
V18 collar affinity involving nonpolynomial square
roots.  The full-collar trace purity is exact by
\eqref{r20:traces}; V19 then compares the continuous
polynomial model to the continuous Wilson law.
These are two distinct steps.

At $d=50$ our explicit rule has $6{,}885{,}676$
nodes per link.  Its unassembled slice matrix would
have that number raised to $192$ as its dimension
when $L=4$.  This is not a reduction in dense
matrix size relative to V19's Gram.
The gain is the positive local factorization:
exact elimination or certified sparse compression
can be attempted without a recoupling sign problem.
No polynomial complexity, small required rank,
or uniform continuum cost has been proved.

\subsection{A contraction certificate with the denominator still present}
\label{r20:elimination}

The following applies directly to the finite
positive graph \eqref{r20:graph}.  Absorb its
positive node weights into the summations.
At any elimination step let $x$ be the variable
being removed and let $y$ denote the union of
remaining arguments of the factors involving $x$.
The exact outgoing function for the current graph is
\[
 h(y)=\sum_x\omega_x\prod_{f\in B} f(x,y_f)>0.
\tag{V20.9}\label{r20:message}
\]
This function is often called a message.  It is
not generally a product over the separate $y_f$.

\begin{proposition}[Pointwise message certificates propagate]
\label{r20:message-cert}
At step $j$, suppose a positive replacement
$\widehat h_j$ is accompanied by proved constants
$0<l_j\le u_j$ such that, at every boundary tuple,
\[
              l_j\widehat h_j\le h_j\le u_j\widehat h_j.
\tag{V20.10}\label{r20:message-error}
\]
Here $h_j$ is computed from the current graph,
including earlier replacements.  If the final
scalar contraction is $\widehat Z_m$, then
\[
 L_m:=\widehat Z_m\prod_jl_j
      \le Z_m^{(d)}\le
 U_m:=\widehat Z_m\prod_ju_j.
\tag{V20.11}\label{r20:Z-interval}
\]
In particular
\[
 \frac{L_{4q}}{U_{2q}^2}\le R_{2q}^{(d)}
       \le\min\left\{1,\frac{U_{4q}}{L_{2q}^2}\right\}.
\tag{V20.12}\label{r20:R-interval}
\]
Combining its lower endpoint with
\eqref{r19:rational-export} gives a Wilson certificate
whenever an appropriate rational $b$ has actually
been verified.  No such small $b$ is asserted here.
\end{proposition}
\begin{proof}
The product of all other factors and weights is
nonnegative.  Multiplying \eqref{r20:message-error}
by that product and summing shows that the partition
function before replacement lies between $l_j$
and $u_j$ times the partition function after it.
Successive inequalities multiply, proving
\eqref{r20:Z-interval}; all exact eliminations have
$l_j=u_j=1$.  Positivity of the denominators gives
\eqref{r20:R-interval}.  The upper bound one follows
from the positive transfer spectrum, not from an
assumption about the accuracy of the contraction.
\end{proof}

For example, writing the bucket product as
$F(x,y)\prod_i H_i(x,y)$ gives the valid bounds
\[
 \Bigl(\sum_x\omega_x F(x,y)\Bigr)
         \prod_i\min_x H_i(x,y)
 \le h(y)\le
 \Bigl(\sum_x\omega_x F(x,y)\Bigr)
         \prod_i\max_x H_i(x,y).
\tag{V20.13}\label{r20:minmax}
\]
For positive functions $f_i$ and $a_i>0$ with
$\sum_i a_i=1$, H\"older also gives
\[
 \sum_x\omega_x\prod_i f_i(x)
       \le\prod_i\left(\sum_x\omega_x f_i(x)^{1/a_i}\right)^{a_i}.
\tag{V20.14}\label{r20:Holder}
\]
The first statement is pointwise multiplication;
the second is ordinary H\"older for the same
probability measure.  These are classical bounds,
included to specify a valid interface.  Choosing
independent integrations without an inequality is
not an alternative certificate.  At rational
$\gamma$, the algebraic factor values can be enclosed with rational outward
intervals; floating point values alone do not
establish \eqref{r20:message-error}.

There is a useful quantitative warning about
free-energy accuracy.  Suppose certified one-sided
logarithmic errors relative to proposed values
$\widehat Z_{2q},\widehat Z_{4q}$ are each at most
$\epsilon$ per plaquette.  The graph has $2Nm$
plaquettes at time period $m$, so
\[
 R_{2q}^{(d)}\ge
     \frac{\widehat Z_{4q}}{\widehat Z_{2q}^2}
                          e^{-16Nq\epsilon}.
\tag{V20.15}\label{r20:volume-error}
\]
At $N=192,q=12$, the exponent coefficient is
$36864$.  The sufficient budget
\[
                   \epsilon\le\frac1{3686400}
\tag{V20.16}\label{r20:precision}
\]
limits this multiplicative loss to at most one
percent, since $e^{-1/100}\ge99/100$.
Good agreement of free energies per plaquette is
not by itself this two-period certificate.
Scalar rescalings that occur with matching total
plaquette counts cancel exactly.  Nonconstant
messages do not cancel merely because both periods
contain the same number of plaquettes per unit time.

\subsection{An evaluated native shared-link stress test}
\label{r20:star}

Freeze all links except an interior link $U$ in
a nondegenerate four-dimensional lattice.  Its
six incident plaquettes have normalized traces
$a_j\cdot u$, where $u\in S^3$ is the quaternion
of $U$ and $a_j\in S^3$ are the six oriented
staples determined by the frozen links.
Define
\[
 \Phi(z)=\int_{S^3}e^{zu_1}\,du
       =\sum_{k=0}^{\infty}\frac{(z^2/4)^k}{k!(k+1)!}.
\tag{V20.17}\label{r20:Phi}
\]
The second equality follows by expanding the
exponential and applying \eqref{r19:Haar}.
It is the familiar one-link Bessel integral;
the integral itself is not claimed as new.

For the shifted Wilson weights, the true
shared-link integration and independent surrogate are
\[
 I(a_1,\ldots,a_6)
    =e^{6\gamma}\Phi\left(\gamma\left|\sum_ja_j\right|\right),
 \qquad I_{\rm ind}=e^{6\gamma}\Phi(\gamma)^6.
\tag{V20.18}\label{r20:star-integrals}
\]
Rotation invariance proves the first formula;
performing six separate one-link integrations proves
the second.  The scalar shift relative to the original
Wilson weights cancels from their ratio.

Consider the aligned choice $a_1=\cdots=a_6=e_1$
and the balanced choice with three $e_1$ and three
$-e_1$.  Both are attainable frozen boundaries:
choose identity links on the ends of the six
length-three staple paths, and choose their distinct
middle links to be the required center elements.
Side and time lengths at least four ensure these
middle links are distinct.  These choices concern
a conditional star; no probability for such
boundaries in the full Gibbs law is assumed.

\begin{proposition}[Certified failure of independent plaquette integration]
\label{r20:star-values}
At the native Wilson coupling $\gamma=10$,
\[
 8{,}000{,}000<\frac{I_{\rm aligned}}{I_{\rm ind}}
       <9{,}000{,}000,
 \qquad
 0<\frac{I_{\rm balanced}}{I_{\rm ind}}<10^{-16}.
\tag{V20.19}\label{r20:star-bounds}
\]
The equal-weight H\"older upper bound, splitting
all six plaquettes, exceeds $I_{\rm balanced}$
by a factor strictly between $10^{23}$ and
$2\cdot10^{23}$.  It is exact at the aligned boundary.
Thus neither independent integration nor a uniformly
sharp complete H\"older splitting has been justified.
\end{proposition}
\begin{proof}
The three relevant ratios are exactly
\[
 \frac{\Phi(60)}{\Phi(10)^6},\qquad
 \frac1{\Phi(10)^6},\qquad \Phi(60).
\tag{V20.20}\label{r20:ratios}
\]
For the last one, rotation invariance makes every
factor of the equal-weight H\"older bound equal
to $[e^{60}\Phi(60)]^{1/6}$; their product is
$e^{60}\Phi(60)$.  The balanced integral is $e^{60}$.
All six functions coincide in the aligned case,
where H\"older is an equality.

Here is an exact reproducible enclosure of the
numbers, with no numerical Bessel routine.
Let $t_k=(z^2/4)^k/[k!(k+1)!]$,
$s_K=\sum_{k=0}^K t_k$, and
$r_K=z^2/[4(K+2)(K+3)]<1$.  The positive terms
and their decreasing ratios give
\[
       s_K\le\Phi(z)\le s_K+\frac{t_{K+1}}{1-r_K}.
\tag{V20.21}\label{r20:Phi-bound}
\]
Use $(z,K)=(10,80)$ and $(60,150)$ and perform
the rational divisions in \eqref{r20:ratios} with
outward endpoints.  Direct integer comparisons
give \eqref{r20:star-bounds} and
$10^{23}<\Phi(60)<2\cdot10^{23}$.
The accompanying exact checker executes these
comparisons, not decimal rounding.
\end{proof}

The same test is evaluated independently in the
degree-50 polynomial model actually proposed in V19.
Put $p(t)=p_{50}(10(1+t))$ and
\[
 c=\int p(u_1)\,du,\quad
 J_+=\int p(u_1)^6\,du,\quad
 J_\pm=\int p(u_1)^3p(-u_1)^3\,du.
\tag{V20.22}\label{r20:poly-star}
\]
Expanding the degree-300 polynomials and applying
\eqref{r19:Haar} computes all three as exact
rationals.  The resulting ratios, shown only
for orientation, are
\[
 J_+/c^6\simeq8.383645493\times10^6,\qquad
 J_\pm/c^6\simeq4.303166115\times10^{-17},\qquad
 J_+/J_\pm\simeq1.948250490\times10^{23}.
\tag{V20.23}\label{r20:poly-star-values}
\]
With V19's certified $\overline\delta_{50}$,
the ratio of a six-factor integral to its independent
surrogate changes by factors between
$(1-\overline\delta_{50})^6$ and its reciprocal
when the polynomial model is replaced by Wilson.
This follows by applying the local relative
bound to the joint integral and to each single
integral.  The rational polynomial calculation
therefore independently certifies the lower
aligned bound, the upper balanced bound, and
the $10^{23}$ loss.  It is not a simulation of
the entire $L=4$ periodic system.

\paragraph{Resulting direction.}
The positive finite factor graph and its exact
physical spectrum are now available.  We have
also acquired an evaluated native local test
showing that shared-link geometry matters greatly
outside the strong-coupling comparison range.
The next useful contraction must retain connected
staple groups, or prove that its boundary-dependent
compression error is small in the appropriate
joint integral.  A uniform worst-boundary splitting
cannot be advertised as nearly exact on the basis
of positivity alone.  The balanced boundary might
have small Gibbs probability; no estimate of that
probability, and no averaged-error bound, has been
proved here.  This is the next upstream issue,
not another improvement of the already tiny scalar
Taylor tail.

\paragraph{Verification and preservation.}
\path{scripts/verify_ym_restart_v20_positive_cubature.py}
checks all 330 monomials through degree seven for
the 48-node rule in exact $\mathbb Q(\sqrt2)$,
and checks its degree-eight failure.  It tests
the separated Fourier/Hopf moments and Gaussian
polynomial-remainder identities, a one-link
conjugation-projected Gram with a nontrivial
magnetic multiplier, two native lattice incidence
counts, the degree-50 conditional star, the
independent positive-series Wilson enclosures,
and message certificates on finite binary graphs.
The latter fixtures are not substitutes for
an interacting four-dimensional contraction.
The earlier native strong-coupling arithmetic is
also rerun.  General theorems rest on the proofs
above, not on the number of test cases.

The full V19 body is preserved except for its
displayed title.  Reversing that title and removing
this append reconstructs its SHA--256
\path{6402965FB8FB706FC1255F1954803EC7BF541DC4FA323A942C5DC9C950B1804B}.
Only active V19 is replaced; the other 46 sources
are unchanged.  Compilation and visual inspection
use artifacts outside \path{latex_notes}.

\paragraph{Status.}
No useful weak-coupling global purity enclosure
has yet been evaluated.  Neither a fixed-period
finite representation nor the local stress test
constructs the interacting continuum, calibrates
physical time, or proves a positive Hamiltonian
mass gap.  The full Yang--Mills goal remains open.

% END CONSOLIDATED V20 APPEND

% BEGIN CONSOLIDATED V21 APPEND -- 2026-09-19
\clearpage
\section{V21: native cutoff budgets and normalized purity stability}
\label{r21:main}

V20 identified attainable boundary staples at which independent
plaquette splitting and its H\"older upper bound have enormous relative
errors.  The next question is not merely whether such boundaries are
rare.  It is whether removing or approximating them leaves the
\emph{normalized endpoint affinity} sufficiently unchanged to control
the original periodic purity.  This appendix proves a one-sided
filter estimate with a total-discarded-probability budget, including
sequential positive elimination.  It also evaluates a nontrivial
native Wilson cutoff budget at an explicitly declared, very large
coupling and fixed volume.  It does not estimate the affinity of the
retained law itself.

\subsection{Archive audit and the scope of the literature transfer}
\label{r21:scope}

The f6 V28 and V41 archive already contains the small-link partition
lower bound, Wilson reflection-positivity arguments, identity-defect
chessboard suppression, and a reduction from misaligned stars to
plaquette defects.  Its f6 V36 bounds concern a pinned reduced
$\SU(3)$ fibre and cannot be silently applied to the full periodic
$\SU(2)$ measure.  The f5 V4 rare-error result controls bounded
cumulants under an $L^1$ modification on a common probability space.
V17--V19 already use square-root densities and Hellinger
data processing.  These general methods, and the idea that a bad
star costs plaquette action, are not new results of V21.

The additional result here is the precise transport from a
\emph{normalized one-sided history filter} to V18's full-collar
purity, together with its sequential-loss accounting and an evaluated
continuous-Wilson example.  Neither the archived bounded-cumulant
estimate nor a partition-function error before normalization is that
transport estimate.

A targeted check of Biskup,
\href{https://arxiv.org/abs/math-ph/0610025}
{\emph{Reflection Positivity and Phase Transitions in Lattice Spin Models}},
Section 5.3, Theorem 5.8, confirms that chessboard estimates require
reflection positivity for the specified block reflections and
observables.  We do not assume that a history cutoff, an approximate
elimination message, or an arbitrarily frozen exterior preserves
these hypotheses.  The evaluated bound below uses only positive
integration and the original Wilson action; it is deliberately
volume-dependent.  A sharper chessboard transfer must retain its own
reflection and normalization proof.  This targeted check does not
reset the scheduled broad literature sweep: the last was V20 and
the next is V30.  The next companion-programme review is V25.

\subsection{Why a rare boundary alone is not an error certificate}
\label{r21:upward}

\begin{proposition}[Upward reweighting can concentrate on a rare set]
\label{r21:rare-counterexample}
For every $0<\delta<1$, let $\mu$ put probabilities $1-\delta,\delta$
on two points, and let an unnormalized approximation multiply their
weights by $1,\delta^{-2}$.  Its normalized law $\nu$ has
\[
 \nu(\{2\})=\frac{1}{1+\delta-\delta^2}\longrightarrow1
 \quad\hbox{as }\delta\downarrow0,
 \qquad \mu(\{2\})=\delta\longrightarrow0.
\tag{V21.1}\label{r21:upward-eq}
\]
In particular no estimate tending to zero with $\mu(\{2\})$ alone
can bound the total-variation change for unrestricted upward errors.
\end{proposition}
\begin{proof}
The new total weight is $1-\delta+\delta^{-1}$.
Dividing its second weight by this total gives
\eqref{r21:upward-eq}.  For $\delta=1/100$, the new probability is
$10000/10099$ and its difference from $1/100$ exceeds $98/100$.
\end{proof}

This does not assert that the Wilson H\"older surrogate has these
particular probabilities.  It explains why multiplying V20's
worst-boundary loss by a claimed rarity without controlling the
reweighted normalizer is invalid.  A lower envelope has a different
and useful structure.

\subsection{One-sided filtering and the square root of the purity defect}
\label{r21:filter}

For probability laws $\rho,\sigma$, write
$H(\rho,\sigma)=\int\sqrt{d\rho\,d\sigma}$ for their Hellinger affinity.
On an output space with a measure-preserving involution $S$, put
\[
 A(\rho)=H(\rho,S_*\rho),\qquad d_S(\rho)=\sqrt{1-A(\rho)}.
\tag{V21.2}\label{r21:defect}
\]
All applications use products of Haar measures and the exchange of
two entire retained spatial slices.  Zeros in a filtered density
are allowed.

\begin{theorem}[Normalized one-sided filter bound]
\label{r21:filter-thm}
Let $\mu$ be a probability law and $0\le r\le1$ a measurable filter.
Suppose
\[
 z=\E_\mu r=1-\alpha>0,\qquad d\nu=\frac r z\,d\mu.
\]
Let $K$ be any common Markov kernel, including deterministic
marginalization, from the original space to the output space.
Then
\[
 H(K\mu,K\nu)\ge\sqrt{1-\alpha},
\qquad
 |d_S(K\mu)-d_S(K\nu)|
 \le 2\sqrt{1-\sqrt{1-\alpha}}\le2\sqrt{\alpha}.
\tag{V21.3}\label{r21:filter-bound}
\]
Consequently, if $d_S(K\nu)\le b$ and $c$ is the middle upper bound
in \eqref{r21:filter-bound}, then
\[
 A(K\mu)\ge\max\{0,1-(b+c)^2\}.
\tag{V21.4}\label{r21:filter-export}
\]
The same statement is valid with $c=2\sqrt{\alpha}$ or with any
certified upper bound for it.
\end{theorem}
\begin{proof}
Since $\sqrt r\ge r$ on $[0,1]$,
\[
 H(\mu,\nu)=\frac{\E_\mu\sqrt r}{\sqrt z}\ge\sqrt z.
\]
To see data processing directly, apply Cauchy--Schwarz on the
original-variable fibre of the joint laws
$\mu(dx)K(x,dy)$ and $\nu(dx)K(x,dy)$, and then integrate over $y$.
It gives $H(K\mu,K\nu)\ge H(\mu,\nu)$.

Use a common $S$-invariant dominating measure on the output space;
one can always add its $S$-pushforward if necessary.  The pullback
$Sf=f\circ S$ is a self-adjoint unitary.  If $f,g$ are the
nonnegative unit square-root densities of $K\mu,K\nu$, then
\[
 d_S(K\mu)=\frac{\|(I-S)f\|_2}{\sqrt2},\qquad
 \|f-g\|_2^2=2(1-H(K\mu,K\nu)).
\]
The reverse triangle inequality and $\|I-S\|\le2$ give
$|d_S(K\mu)-d_S(K\nu)|\le\sqrt2\|f-g\|_2$.
Insert the affinity bound.  Finally,
$1-\sqrt{1-\alpha}\le\alpha$ and
$d_S(K\mu)\le b+c$ prove the remaining assertions.
\end{proof}

\paragraph{The native purity is recovered only after the right marginal.}
Let $\mathbb P_q$ be two independent full Wilson histories of period
$2q$, with all spatial and temporal links, as in V17.
Let $K_q$ keep precisely their spatial slices at times $0,q$.
On these four slices let $S$ exchange the two time-$q$ slices.
Theorem~\ref{r18:collar-thm} gives
\[
 A(K_q\mathbb P_q)=R_{2q}=\frac{Z_{4q}}{Z_{2q}^2}.
\tag{V21.5}\label{r21:native-purity}
\]
Thus Theorem~\ref{r21:filter-thm} transfers a bound on the affinity
of the \emph{normalized, filtered endpoint law} to the original
purity.  It does not replace this endpoint affinity by the
microscopic-history swap, which has V17's counterterm and floor.
The filtered law need not be reflection positive or generated by
one transfer operator.  We therefore do not call its affinity a
new Hamiltonian's purity, and use no spectral interpretation for it.

\subsection{Sequential lower envelopes pay total probability loss}
\label{r21:sequential}

\begin{theorem}[Accumulated rejection, including elimination lifts]
\label{r21:sequential-thm}
Let $\mu_0=\mu$, and successively set
\[
 z_j=\E_{\mu_{j-1}}r_j=1-\alpha_j>0,\qquad
 d\mu_j=\frac{r_j}{z_j}\,d\mu_{j-1},\qquad 0\le r_j\le1.
\]
Then
\[
 d\mu_k=\frac{\prod_{j=1}^k r_j}
                   {\prod_{j=1}^k(1-\alpha_j)}\,d\mu,\qquad
 \alpha_{\rm tot}:=1-\E_\mu\prod_{j=1}^k r_j
 =1-\prod_{j=1}^k(1-\alpha_j)\le\sum_{j=1}^k\alpha_j.
\tag{V21.6}\label{r21:total}
\]
For any final common marginal $K$ and swap $S$, the defect change
is at most $2\sqrt{\alpha_{\rm tot}}$.  If certified estimates
$\alpha_j\le a_j$ have $\sum_j a_j<1$, a sufficient bound is
$2\sqrt{\sum_j a_j}$, not a sum of their individual square roots.
\end{theorem}
\begin{proof}
Multiply the Radon--Nikodym derivatives, and integrate the
resulting probability density to obtain the equality.
The inequality follows inductively from
$1-(1-a)(1-b)\le a+b$ for $a,b\in[0,1]$.
Apply Theorem~\ref{r21:filter-thm} once with filter $\prod_j r_j$.
\end{proof}

Here is its exact connection to positive elimination.  In the
\emph{current} graph let $y$ be the remaining variables,
$h(y)>0$ the exact outgoing bucket integral, and $F(y)\ge0$
the product of all other factors.  The current marginal has
density proportional to $Fh$.  If
$0\le\widehat h\le h$, replacing $h$ by $\widehat h$ is exactly
filtering this marginal by
\[
 r(y)=\frac{\widehat h(y)}{h(y)},\qquad
 \alpha=\E_{\rm current}\left(1-\frac{\widehat h}{h}\right).
\tag{V21.7}\label{r21:message}
\]
Lift $r(y)$ to a function of the original full history by its
remaining-coordinate map.  Inductively, exact elimination is
marginalization of the already filtered full law, and the next
replacement is multiplication of that full law by the lifted
$r$.  This proves the full-history representation used in
\eqref{r21:total}.  To compute the final endpoint affinity one
must retain the endpoint coordinates during elimination, or
retain their exact induced marginal; a scalar partition-function
contraction alone is insufficient.

If $1-\widehat h/h\le\epsilon$ off a bad set $B$, then
\[
 \alpha\le\epsilon+\mu_{\rm current}(B).
\tag{V21.8}\label{r21:good-bad}
\]
This follows by integrating the loss separately on $B$ and $B^c$.
The bad-set probability must be for the \emph{current normalized
law}, not automatically for the initial Gibbs measure.  There is
also an honest initial-law fallback.  With
$R_{j-1}=\prod_{i<j}r_i$ and
$z_{<j}=\E_\mu R_{j-1}>0$,
\[
 \mu_{j-1}(B)=\frac{\E_\mu[R_{j-1}{\bf1}_B]}{z_{<j}}
 \le\frac{\mu(B)}{z_{<j}},
\tag{V21.9}\label{r21:initial-fallback}
\]
because $R_{j-1}\le1$.  Its denominator cannot be omitted.

These statements hold for either a continuous graph or a finite
positive graph on its own probability space.  V20's atomic
cubature law is singular with respect to Haar and does not
automatically inherit probabilities of nonpolynomial cutoff events.
In particular the continuous event estimate proved next is not
asserted for the cubature nodes.  V19's separately normalized
continuous polynomial-to-Wilson comparison remains a different
available bridge.

\subsection{An elementary native probability budget for frustrated stars}
\label{r21:star}

Consider the original $\SU(2)$ Wilson law on a four-dimensional
periodic torus with all side lengths at least four.  Write
$V$ for its number of sites, $E=4V$ for its independent oriented
links, and $P=6V$ for its plaquettes.  With probability Haar
reference measure,
\[
 d\mu_\gamma=Z_\gamma^{-1}e^{-\gamma\mathcal S(U)}\,dU,\qquad
 \mathcal S(U)=\sum_p d_p(U),\qquad
 d_p=1-\tfrac12\Tr U_p\in[0,2].
\tag{V21.10}\label{r21:wilson}
\]
Here $E$ is a cardinality, not expectation, and the action
normalization is exactly that of V19--V20.

For a fixed link with quaternion $u\in S^3$, cyclically orient
each of its six plaquette traces so that
$\tfrac12\Tr U_p=u\cdot a_{e,p}$ for a unit quaternion
$a_{e,p}$ independent of this link.  Set
\[
 H_e=\sum_{p\ni e}a_{e,p},\qquad r_e=|H_e|\le6,\qquad
 B_\epsilon=\{\hbox{some link }e\hbox{ has }r_e\le6-\epsilon\},
 \quad 0<\epsilon\le6 .
\tag{V21.11}\label{r21:bad}
\]
The norm is gauge invariant: gauge transformations act on $u$
and all six dual staples by the same orthogonal transformation.

\begin{lemma}[Star frustration forces total plaquette action]
\label{r21:energy}
For every configuration and every link,
\[
 \sum_{p\ni e}d_p=6-u\cdot H_e\ge6-r_e,\qquad
 36-r_e^2=\sum_{p<p':\,p,p'\ni e}|a_{e,p}-a_{e,p'}|^2.
\tag{V21.12}\label{r21:energy-eq}
\]
Consequently $B_\epsilon\subset\{\mathcal S\ge\epsilon\}$.
\end{lemma}
\begin{proof}
The first identity follows by summing the six plaquette traces,
and its inequality is Cauchy--Schwarz with $|u|=1$.
Expand the square of the sum of the six unit staples for the
second identity.  If one star is bad, its incident action is
at least $\epsilon$; all other plaquette terms are nonnegative.
The six incident plaquettes are distinct on the stated torus.
\end{proof}

This inclusion controls the union over \emph{all} links without
an additional union factor.  Its cost will instead be the
extensive partition lower bound.  The balanced centre boundary
of V20 has $r_e=0$ and is included in every such $B_\epsilon$.

\begin{lemma}[Explicit $\SU(2)$ small-ball lower bound]
\label{r21:haar}
Let $h(\eta)$ be the Haar probability that a unit quaternion has
geodesic angle at most $\eta$ from $1$.  For $0<\eta\le1$,
\[
 h(\eta)=\frac2\pi\int_0^\eta\sin^2\theta\,d\theta
 \ge\frac{25}{54\pi}\eta^3>\frac{\eta^3}{8}.
\tag{V21.13}\label{r21:haar-eq}
\]
If all links lie in that ball, then $d_p\le8\eta^2$ for every
plaquette, and hence
\[
 Z_\gamma\ge(\eta^3/8)^{4V}\exp(-48\gamma V\eta^2).
\tag{V21.14}\label{r21:partition}
\]
\end{lemma}
\begin{proof}
The radial Haar density on $S^3$ is $2\sin^2\theta/\pi$.
For $0\le\theta\le1$, $\sin\theta\ge\theta-\theta^3/6
\ge5\theta/6$.  Integration gives the first inequality;
$\pi<22/7$ gives $25/(54\pi)>175/1188>1/8$.
The quaternion chord distance of a link or its inverse from $1$
is at most its geodesic angle.  Multiplication by a unit
quaternion is an isometry, so a four-link plaquette product
satisfies $|U_p-1|\le4\eta$ by telescoping.  Since
$d_p=|U_p-1|^2/2$, its defect is at most $8\eta^2$.
Integrate the Wilson weight over the product of these link balls.
\end{proof}

\begin{theorem}[Unconditional finite-volume star cutoff budget]
\label{r21:tail}
For $\gamma\ge1/8$ define
\[
 Q_{V,\epsilon}(\gamma)
 =2^{30V}\gamma^{6V}\exp(6V-\gamma\epsilon).
\tag{V21.15}\label{r21:Q}
\]
Then the original continuous periodic Wilson law obeys
\[
 \mu_\gamma(B_\epsilon)\le\min\{1,Q_{V,\epsilon}(\gamma)\}.
\tag{V21.16}\label{r21:tail-eq}
\]
For two independent histories, each with $V=mL^3$ sites, let
$G_\epsilon$ require $r_e>6-\epsilon$ at every link of both
histories.  Its discarded probability satisfies
\[
 \alpha:=\mathbb P(G_\epsilon^c)
 \le\min\{1,\,2Q_{V,\epsilon}(\gamma)\}.
\tag{V21.17}\label{r21:pair-tail}
\]
\end{theorem}
\begin{proof}
On $B_\epsilon$, Lemma~\ref{r21:energy} bounds the unnormalized
weight by $e^{-\gamma\epsilon}$, and the reference measure has
total mass one.  Divide by \eqref{r21:partition} and take
$\eta^2=(8\gamma)^{-1}\le1$.  The reciprocal partition bound is
$2^{30V}\gamma^{6V}e^{6V}$, proving \eqref{r21:tail-eq}.
Use the two-history union bound for \eqref{r21:pair-tail}.
No reflection estimate, conditional exterior bound, gauge
fixing, or continuum approximation enters this proof.
\end{proof}

For a prescribed $0<a<1$, a sufficient numerical criterion
for the two-history loss to be at most $a$ is
\[
 \gamma\epsilon
 \ge V(30\log2+6\log\gamma+6)+\log(2/a).
\tag{V21.18}\label{r21:volume-budget}
\]
This estimate is not volume-uniform.  At fixed $\gamma$ its
displayed right side grows with $V$; it gives no cofinal
continuum result.  At the earlier test coupling $\gamma=10$,
with $V=32\cdot4^3$ and $\epsilon=3$, its probability upper
bound is simply one.  A useful single-point evaluation at a
different coupling must not be reported as resolving that test.

There is also a temporal limitation worth making explicit.  If
$m\ge c\gamma$ for some fixed $c>0$ and $L$ is fixed, then, for
every $0<\epsilon\le6$,
\[
 \frac{\log Q_{mL^3,\epsilon}(\gamma)}{\gamma}
 \ge 6cL^3\log\gamma-6\longrightarrow+\infty .
\]
Thus this particular estimate eventually becomes vacuous on
such growing-period sequences.  We do not assume that this
sequence is the required physical clock; the calculation
simply prevents a fixed-period success from being extrapolated
to arbitrary ground-state projection times.

\subsection{An evaluated cutoff cost, but not an evaluated purity}
\label{r21:evaluation}

\begin{corollary}[Exact arithmetic at a declared native point]
\label{r21:point}
Take the original Wilson model with
\[
 \gamma=2^{16}=65536,\qquad L=4,\qquad m=2q=32,\qquad
 V=2048,\qquad \epsilon=3.
\tag{V21.19}\label{r21:point-eq}
\]
For the two-history good event $G_3$,
\[
 \alpha=\mathbb P_q(G_3^c)<2^{-2048}.
\tag{V21.20}\label{r21:point-alpha}
\]
Define $\nu_G=\mathbb P_q(\,\cdot\,\mid G_3)$ and its
\emph{endpoint affinity} $A_G=A(K_q\nu_G)$.
Then
\[
 \left|\sqrt{1-R_{32}}-\sqrt{1-A_G}\right|<2^{-1023}.
\tag{V21.21}\label{r21:point-affinity}
\]
In particular, a separately proved $A_G\ge1-b^2$ would imply
$R_{32}\ge\max\{0,1-(b+2^{-1023})^2\}$.  No such nontrivial
bound on $A_G$ is proved here.
\end{corollary}
\begin{proof}
Substitution in \eqref{r21:pair-tail} gives
\[
 2Q_{2048,3}(65536)=2^{258049}e^{-184320}.
\]
The exponential series gives $e>8/3$.  Also
$184320=17\cdot10842+6$,
$3^{17}=129140163<134217728=2^{27}$ and $3^6<2^{10}$.
Thus
\[
 2^{258049}e^{-184320}
 <\frac{3^{184320}}{2^{294911}}
 <2^{292744-294911}=2^{-2167}<2^{-2048}.
\]
The good event has positive probability by this bound.
Apply Theorem~\ref{r21:filter-thm} with $r={\bf1}_{G_3}$,
the marginal $K_q$, and \eqref{r21:native-purity}.
Its error is at most $2\sqrt\alpha<2^{-1023}$.
\end{proof}

\paragraph{What this calculation does and does not remove.}
It certifies the normalized cost of discarding a specified
family of frustrated-star histories at one finite weak-coupling
point, for the original continuous measure.  It is not a
simulation, a probability assigned to one measure-zero centre
configuration, or a proof that the remaining connected graph
has small purity defect.  The threshold $\gamma=65536$ is a
conservative evaluated choice, not a physical transition.

Moreover, the retained region still includes the aligned
$r_e=6$ boundary, where V20's independent splitting was already
grossly inaccurate.  It also includes the familiar distinct flat
holonomy sectors: take commuting constant links
$U_{x,\mu}=\exp(i\theta_\mu\sigma_3/\ell_\mu)$ on a torus of
side lengths $\ell_\mu$.  Every plaquette is $I$, every star
has $r_e=6$, but the winding-loop traces
$2\cos\theta_\mu$ vary and are gauge invariant.
Indeed $u\cdot a_{e,p}=1$ forces every staple to equal $u$.
This standard flat-connection example is a scope check,
not a new claimed obstruction theorem.
Discarding $B_3$ therefore neither produces a single gauge-fixed
convex vacuum chart nor justifies independent plaquette contraction.

\subsection{The next quantitative target and verification ledger}
\label{r21:next}

The new usable input is a probability-to-normalized-affinity error
bound.  The outstanding input is now explicitly the affinity
of the connected retained endpoint law, or lower-envelope
messages whose \emph{current-law} expected losses can be bounded
while those endpoints are retained.  A useful next calculation
must preserve the aligned shared-staple dependence and the flat
holonomy sectors; another proof that only severely frustrated
stars are rare will not supply it.  If one uses a sharper
chessboard tail, it must improve the quantitative volume/coupling
budget without assuming that the filtered law retains reflection
positivity.

\paragraph{Verification.}
\path{scripts/verify_ym_restart_v21_cutoff_purity.py}
checks filters, marginalization, swaps, sequential probabilities,
elimination lifts, quaternion identities, native incidence and
the displayed constants, using exact rational arithmetic and
certified square-root intervals.  V20's mathematical checks are
rerun.  These fixtures do not evaluate $A_G$; the native
probability bound rests on the proofs above.

Removing this append and restoring V20's displayed title
recovers its complete 484826-byte source and SHA--256
\path{2A3D3BC2C29E8AF7043C4C9CB254D339D9AE113589A62D5DA94CBAD82FA38678}.
Only active V20 is replaced; the other 46 sources are unchanged.
All non-source artifacts remain outside \path{latex_notes}.

\paragraph{Status.}
The finite-volume cutoff cost is now evaluated, but a useful
weak-coupling global purity lower bound remains unproved.
Cofinal volume/cutoff estimates, physical-time calibration,
the interacting four-dimensional continuum construction and
its positive physical Hamiltonian mass gap remain open.
No part of this appendix is a claim to the full Yang--Mills
mass-gap theorem.

% END CONSOLIDATED V21 APPEND

% BEGIN CONSOLIDATED V22 APPEND -- 2026-09-19
\clearpage
\section{V22: connected pair messages and a certified native decimation}
\label{r22:main}

V21 controls the normalized cost of a one-sided approximation, but
does not construct a useful connected approximation.  This appendix
does so for a specified first elimination pass of the original
Wilson lattice.  A shared plaquette is integrated jointly with the
ten private plaquettes at its two opposite links.  The resulting
message keeps three exterior invariants, including the angle between
the transported staple fields.  A positive beta-moment formula gives
an all-exterior relative lower approximation.  At $\gamma=10$ its
error is below $10^{-18}$ per pair, and an actual endpoint-preserving
packing has a certified root-affinity error below $3\cdot10^{-8}$.

This evaluates the error of a nontrivial connected decimation,
not the endpoint affinity of the remaining graph.  The generated
interactions have not been shown to remain in a closed, efficiently
contractible family under further elimination.

\subsection{What the archive already supplies}
\label{r22:scope}

The f16.2 V62 archive explicitly treats the joint conditional
integral on a ladder of parallel links.  Its one-rung case already
has five private staples at each free link.  It proves
all-exterior curvature bounds and the coherent Hessian, retaining
the shared interaction; we do not claim that integral or those
geometric facts as new.  The f15 V69 archive likewise integrates
six parallel neighbours through the induced factors
$\Phi(\gamma|c_i+v|)$.  The f6 V23 one-plaquette witness has a
different scope from the eleven-plaquette conditional cell here.
V15's log-rectangle estimates and V20--V21's normalization
certificates also remain in force.

What is added is a positive finite formula with a uniform relative
error, its evaluated native centre-sector contrast, and its
application to a concrete simultaneous pair elimination with the
purity endpoints retained.  It is not another coherent Hessian
calculation or an inference of normalized correlations from an
action lower bound.

For the targeted literature check, Koyama, Nakayama, Nishiyama and
Takayama,
\href{https://arxiv.org/abs/1201.3239}
{\emph{Holonomic Gradient Descent for the Fisher--Bingham
Distribution on a $d$-dimensional Sphere}},
Section 4, Theorem 2, supplies a published example of spherical
normalizer series with explicit truncation estimates.  Their
single-sphere problem is not our product of two $S^3$ spheres.
We derive the required formula directly below.  The numerical
error propagation discussed in their Section 5 is not imported as
a deterministic enclosure.  No holonomic solver or statistical
confidence interval is used in the certificate.
The next broad sweep remains V30, and the next companion review V25.

\subsection{The literal two-link message and its three invariants}
\label{r22:cell}

Work on the original four-dimensional periodic $\SU(2)$ lattice,
with every side length at least four.  Free two parallel,
positively oriented $\mu$ links $e=(x,\mu)$ and
$f=(x+\widehat\nu,\mu)$, $\mu\ne\nu$, and freeze all other links.
Exactly one plaquette contains both free links.  Each has five
other incident plaquettes, and all eleven are distinct.
In quaternion coordinates $u,v\in S^3$ their conditional factor is
\[
 e^{-11\gamma}
 \exp\{\gamma(\mathbf a\cdot u+\mathbf b\cdot v+u\cdot Rv)\},
 \qquad |\mathbf a|,|\mathbf b|\le5,\quad R\in\mathrm{SO}(4).
\tag{V22.1}\label{r22:raw}
\]
Here $\mathbf a,\mathbf b$ are the sums of the five oriented
private staples.  The shared plaquette supplies
$R(v)=\ell v r^{-1}$ for its two retained complementary links.
Quaternion multiplication is orthogonal; the scalar Wilson
defect of each of the eleven plaquettes contributes its
$e^{-\gamma}$ factor.  All other action factors are left outside
the message.

Changing variable $w=Rv$ and putting $\mathbf c=R\mathbf b$ gives
the exact outgoing message
\[
 \begin{split}
 h_\gamma(\mathbf a,\mathbf b,R)
     &=e^{-11\gamma}J_\gamma(\mathbf a,\mathbf c),\\
 J_\gamma(\mathbf a,\mathbf c)
     &=\int_{S^3\times S^3}
        e^{\gamma(\mathbf a\cdot u+\mathbf c\cdot w+u\cdot w)}
                           \,dH(u)\,dH(w)\\
     &=\int_{S^3}e^{\gamma\mathbf c\cdot w}
                         \Phi(\gamma|\mathbf a+w|)\,dH(w),\\
 \Phi(z)&=\int_{S^3}e^{zu_0}\,dH(u)
       =\sum_{k=0}^{\infty}\frac{(z^2/4)^k}{k!(k+1)!}.
 \end{split}
\tag{V22.2}\label{r22:J}
\]
The last series is V20's one-link normalizer.  The first equality
in the last line for $J$ integrates $u$ exactly, not independently
of $w$.

\begin{proposition}[Three-invariant reduction]
\label{r22:invariants}
The message depends only on
$|\mathbf a|^2,|\mathbf c|^2,\mathbf a\cdot\mathbf c$.
For $A=|\mathbf a|>0$ define
\[
 B=\frac{\mathbf a\cdot\mathbf c}{A},\qquad
 C=\sqrt{|\mathbf c|^2-B^2}.
\]
If $A=0$, choose $B=|\mathbf c|$, $C=0$.
Let $\sigma=1$ for $B\ge0$ and $\sigma=-1$ for $B<0$;
put $t_\sigma=t$ or $1-t$ respectively, and let
$\mathbb E_\beta$ denote expectation under
$\mathrm{Beta}(3/2,3/2)$ on $[0,1]$.
Then, with $\mathcal C(z)=\sinh(z)/z$ and $\mathcal C(0)=1$,
\[
 \begin{split}
 J_\gamma={}&e^{-\gamma|B|}\mathbb E_\beta\left[
 e^{2\gamma|B|t}\,
 \mathcal C\!\left(2\gamma C\sqrt{t(1-t)}\right)
 \Phi\!\left(\gamma\sqrt{(A-1)^2+4At_\sigma}\right)
 \right].
 \end{split}
\tag{V22.3}\label{r22:beta}
\]
\end{proposition}
\begin{proof}
Simultaneous orthogonal transformations of $\mathbf a,\mathbf c,u,w$
preserve the first integral in \eqref{r22:J}.  Pairs of vectors
with the same Gram matrix are related by such a transformation.
Choose an axis along $\mathbf a$ when it is nonzero, and write
$x$ for the component of $w$ along that axis.  Its density is
$2\sqrt{1-x^2}/\pi$ on $[-1,1]$; conditionally, the other
three components are uniform on the sphere of radius
$\sqrt{1-x^2}$.  A coordinate of the unit $S^2$ is uniform on
$[-1,1]$, so its exponential integral is $\mathcal C$.
The longitudinal field contributes $e^{\gamma Bx}$ and
$|\mathbf a+w|^2=A^2+1+2Ax$.
Substitute $x=\sigma(2t-1)$.  The normalized density becomes
$8\sqrt{t(1-t)}/\pi$, and the stated factors follow.
At $A=0$ choose the axis along $\mathbf c$ (arbitrarily if it
also vanishes); the same calculation applies.
\end{proof}

In particular the norms alone are insufficient; the third
invariant is a genuine connected interaction.  Under actual
gauge transformations the two fields and $R$ change frames
compatibly, leaving these three scalar invariants unchanged.
They parametrize one message, not all exterior links of the
remaining lattice.

\subsection{A finite positive beta-moment lower envelope}
\label{r22:positive}

For nonnegative arguments define
\[
 E_N(x)=\sum_{n=0}^N\frac{x^n}{n!},\quad
 C_M(z)=\sum_{\ell=0}^M\frac{z^{2\ell}}{(2\ell+1)!},\quad
 \Phi_K(z)=\sum_{k=0}^K\frac{(z^2/4)^k}{k!(k+1)!}.
\tag{V22.4}\label{r22:trunc}
\]
Let $\widehat J$ be the expression \eqref{r22:beta} with its
three functions replaced by $E_N,C_M,\Phi_K$, respectively.
Thus $0<\widehat J\le J_\gamma$.  The positivity assertion is
made after the exact conditional $S^2$ integration, not by
claiming that an arbitrary odd exponential Taylor polynomial
is a lower approximation at a negative argument.

Here is an explicit finite formula requiring no remaining
quadrature.  Set
\[
 \begin{split}
 e_n&=\frac{(2\gamma|B|)^n}{n!},\qquad
 s_\ell=\frac{(4\gamma^2 C^2)^\ell}{(2\ell+1)!},\\
 W_j&=\sum_{k=j}^K
   \frac{(\gamma^2/4)^k}{k!(k+1)!}
   \binom{k}{j}(A-1)^{2(k-j)}(4A)^j,\\
 M_{p,r}&=\frac{(3/2)_p(3/2)_r}{(3)_{p+r}},
 \end{split}
\tag{V22.5}\label{r22:coeff}
\]
where $(a)_n=a(a+1)\cdots(a+n-1)$ and $(a)_0=1$.
The value $M_{p,r}$ is the beta moment
$\mathbb E_\beta[t^p(1-t)^r]$.

\begin{theorem}[Positive finite pair formula]
\label{r22:finite}
For every exterior, including vanishing fields,
\[
 \widehat J=e^{-\gamma|B|}
 \sum_{n=0}^N\sum_{\ell=0}^M\sum_{j=0}^K e_ns_\ell W_j
 \begin{cases}
 M_{n+\ell+j,\ell},&\sigma=1,\\
 M_{n+\ell,\ell+j},&\sigma=-1.
 \end{cases}
\tag{V22.6}\label{r22:finite-eq}
\]
All displayed summands are nonnegative.  If
$\gamma,A,B,C^2$ are rational, the sum after removing its
explicit exponential prefactor is rational.
\end{theorem}
\begin{proof}
Expand each of the three finite factors in \eqref{r22:trunc}.
The binomial expansion of
$[(A-1)^2+4At_\sigma]^k$ has the nonnegative coefficients
in \eqref{r22:coeff}; this uses a squared constant rather
than a signed expansion about $x=0$.
The remaining monomials are precisely those whose moments
appear in \eqref{r22:finite-eq}.
The beta integral gives the rising-factorial expression for
$M_{p,r}$, which is rational at integer indices.
The zeroth summands also ensure strict positivity.
\end{proof}

This formula is not asserted to be a polynomial in all
original link coordinates.  Its invariant parameters include
square roots and a sign choice; it defines a gauge-invariant
measurable positive function with the bounds proved below.
In particular V20's polynomial cubature cannot simply be
applied to the resulting message as though its degree were
$N+2M+K$ in the original links.

\subsection{Uniform relative tails and a numerical certificate}
\label{r22:tails}

For any of the positive series
$f(z)=\sum_{j\ge0}a_jz^j$, its relative omitted tail beyond
index $d$ increases with $z\ge0$.  Indeed, writing $T$ for
the tail and $S$ for the partial sum,
$T'S-TS'$ is a sum of terms
$(j-i)a_ja_i z^{i+j-1}\ge0$, $j>d\ge i$.
The argument also applies to the even-power series by
using $z^2$ as parameter.
If successive omitted terms have ratio at most $\rho<1$,
and $u$ is the first omitted term divided by $1-\rho$,
then
\[
 \frac{f-S}{f}\le\frac{u}{S+u}.
\tag{V22.7}\label{r22:relative-tail}
\]
This is because $T\le u$ and $T/(S+T)$ increases with $T$.
These are the same elementary geometric-tail principles
used for V19's positive local truncation.

For the three functions, use their maximal arguments
\[
       X=10\gamma,\qquad Z_C=5\gamma,\qquad Z_\Phi=6\gamma.
\tag{V22.8}\label{r22:max-args}
\]
Let $\delta_E,\delta_C,\delta_\Phi$ be the bounds
\eqref{r22:relative-tail} at those arguments, with partial
sums $E_N(X)$, $C_M(Z_C)$ and $\Phi_K(Z_\Phi)$.  The corresponding
geometric tail data are
\[
 \begin{array}{c|c|c}
 &\hbox{first omitted term}&\rho\\ \hline
 E&X^{N+1}/(N+1)!&X/(N+2)\\
 C&Z_C^{2M+2}/(2M+3)!&Z_C^2/((2M+4)(2M+5))\\
 \Phi&(Z_\Phi^2/4)^{K+1}/((K+1)!(K+2)!)
       &Z_\Phi^2/(4(K+2)(K+3)).
 \end{array}
\tag{V22.9}\label{r22:tail-table}
\]
Choose the truncation indices so all three $\rho$ are less
than one.

\begin{theorem}[All-exterior relative message certificate]
\label{r22:relative}
Put
$\delta=1-(1-\delta_E)(1-\delta_C)(1-\delta_\Phi)$ and
$\widehat h=e^{-11\gamma}\widehat J$.  Then for every
actual frozen exterior,
\[
        (1-\delta)h_\gamma\le\widehat h\le h_\gamma,
        \qquad \delta\le\delta_E+\delta_C+\delta_\Phi.
\tag{V22.10}\label{r22:relative-eq}
\]
At $\gamma=10$, the choice
$N=200,\ M=80,\ K=80$ gives the exact rational certificates
\[
 \delta_E<5\cdot10^{-19},\qquad
 \delta_\Phi<5\cdot10^{-28},\qquad
 \delta_C<2\cdot10^{-36},\qquad
 \delta<10^{-18}.
\tag{V22.11}\label{r22:tail-values}
\]
There is no exceptional-boundary hypothesis.
\end{theorem}
\begin{proof}
We have $A\le5$, $|B|,C\le5$, $0\le t\le1$ and
$2\sqrt{t(1-t)}\le1$.  Also the radial argument of
$\Phi$ is at most $\gamma(A+1)\le6\gamma$.
Relative-tail monotonicity therefore bounds each of
the three positive factors pointwise in $t$ by the
corresponding retained fraction at its maximal argument.
Multiply these inequalities and integrate against the
positive beta density to obtain \eqref{r22:relative-eq}.
The final inequalities are finite rational comparisons
using \eqref{r22:tail-table}; the accompanying checker
evaluates the partial sums and their geometric tails
exactly.  No numerical exponential or logarithm is
required for those comparisons.
\end{proof}

The certificate controls the whole connected message.
It is not an assertion that the eleven plaquettes are
independent or that the remaining variables follow
their original Haar reference law.

\subsection{A native angular contrast and a separable-minorant obstruction}
\label{r22:contrast}

The two centre-sector messages at $\gamma=10$ will be denoted
\[
 J_+=J_{10}(5e_0,5e_0),\qquad
 J_-=J_{10}(5e_0,-5e_0).
\tag{V22.12}\label{r22:pm}
\]
They have the same field lengths.  These are attainable
conditional boundary values in the actual lattice.
To see this, set the two complementary shared-plaquette
links and all transverse connector links to $I$.
The five private plaquettes at $e$ have five distinct
opposite parallel links, and the five at $f$ have another
five, distinct from the first set.  Assign these ten
links the centres $\sigma I$ and $\tau I$ in their
respective sets.  Then
$\mathbf a=5\sigma e_0$, $\mathbf c=5\tau e_0$, $R=I$.
Side lengths at least four ensure the asserted distinctness.
The checker verifies these assignments on the literal
periodic plaquette list.

\begin{proposition}[Evaluated connected angular contrast]
\label{r22:contrast-prop}
For these native conditional messages,
\[
          130000000<\frac{J_+}{J_-}<150000000.
\tag{V22.13}\label{r22:ratio}
\]
The corresponding message matrix for the four centre
choices is $e^{-110}
\begin{pmatrix}J_+&J_-\\J_-&J_+\end{pmatrix}$.
Both updated-link staple norms are at least four,
for every value of the other integrated link, in
each of these four boundary sectors.
\end{proposition}
\begin{proof}
For $A=5,B=\pm5,C=0$, only $\ell=0$ contributes to
\eqref{r22:finite-eq}.  Let
\[
 S_\pm=\sum_{n=0}^{200}\sum_{j=0}^{80}
              \frac{100^n}{n!}W_j
 \begin{cases}M_{n+j,0},&+,\\M_{n,j},&-.\end{cases}
\]
Here $W_j$ is the rational value from
\eqref{r22:coeff} at $A=5,\gamma=10,K=80$.
The exponential $e^{-50}$ cancels in the ratio.
Equation \eqref{r22:relative-eq} implies
\[
 (1-\delta)\frac{S_+}{S_-}
       \le\frac{J_+}{J_-}
       \le\frac{S_+}{(1-\delta)S_-}.
\tag{V22.14}\label{r22:ratio-cert}
\]
Exact rational comparisons of these endpoints give
\eqref{r22:ratio}.  A decimal check gives
$S_+/S_-\simeq1.408889717\cdot10^8$; it is not the certificate.
Simultaneous reversal of $u,w$ exchanges the two
like-sign choices and the two unlike-sign choices,
giving the matrix.  Finally
$|5\sigma e_0+w|\ge4$ and $|5\tau e_0+u|\ge4$.
\end{proof}

\begin{corollary}[A product lower envelope must discard almost everything
at some boundary]
\label{r22:rank-one}
Suppose a positive separable approximation to this four-entry
message has entries $\widehat h_{\sigma\tau}=f_\sigma g_\tau$
and $\widehat h_{\sigma\tau}\le h_{\sigma\tau}$.
Then
\[
        \min_{\sigma,\tau}
          \frac{\widehat h_{\sigma\tau}}{h_{\sigma\tau}}
        \le\frac{J_-}{J_+}<\frac1{130000000}.
\tag{V22.15}\label{r22:rank-one-eq}
\]
The first inequality is sharp for this matrix.
\end{corollary}
\begin{proof}
Let $r_{\sigma\tau}=\widehat h_{\sigma\tau}/h_{\sigma\tau}$
and $\rho=J_+/J_->1$.  The rank-one identity gives
$r_{++}r_{--}=\rho^{-2}r_{+-}r_{-+}\le\rho^{-2}$.
If every ratio is at least $r_*$, then
$r_*^2\le\rho^{-2}$.  The constant minorant
$\widehat h_{\sigma\tau}=e^{-110}J_-$ attains the bound.
\end{proof}

This elementary cross-ratio test is not claimed as a new
general rank theorem.  Its significance is the evaluated
native cell, including all ten private plaquettes.
The lower bound four applies to the two updated stars only;
the calculation is not a probability estimate for these
exteriors, or a claim that every other star lies in V21's
global good event.  It does show why an unfrustrated
updated star is not enough to justify separating the
two connected exterior fields.

\subsection{An endpoint-preserving pair packing in the full Wilson law}
\label{r22:packing}

Let the spatial torus have side $L$, a multiple of four,
and let its temporal period be $2q\ge4$.  Use coordinates
$(x_1,x_2,x_3,t)$, and free only links in the first
spatial orientation.  For each anchor satisfying
\[
 x_2\equiv0\pmod4,\qquad x_3\equiv0\pmod2,\qquad
 t\equiv0\pmod2,\qquad t\notin\{0,q\},
\tag{V22.16}\label{r22:anchors}
\]
free the parallel pair with anchors $x$ and
$x+\widehat 2$.  Every other link is retained.
There are
\[
       J=\frac{L^3}{8}\left(q-1-{\bf1}_{\{q\ {\rm even}\}}\right)
\tag{V22.17}\label{r22:pair-count}
\]
pairs in each periodic history.  At $q=2$ the set is empty;
the nontrivial evaluation below uses $q=12$.

\begin{lemma}[Exact independent-block factorization conditional on the exterior]
\label{r22:packing-lemma}
No plaquette contains free links from two different pairs.
Each pair touches exactly eleven plaquettes, and these
eleven-plaquette sets are disjoint between pairs.
Consequently, after exact integration of all free pairs,
the retained density is proportional to
$F(\eta)\prod_{j=1}^J h_j(\eta)$, where $F$ is the product
of untouched Wilson factors.  Both spatial endpoint
slices at times $0,q$ are retained in full.
\end{lemma}
\begin{proof}
A plaquette contains at most two links in the first
spatial orientation.  When it contains two, their
anchors differ by one step in one of the other three
directions.  In direction 2, the selected anchor
positions form separated pairs $\{4k,4k+1\}$; a
nearest-neighbour connection joins selected positions
only inside one pair.  In direction 3 and in time,
selected positions have even parity, so they cannot
be nearest neighbours.  This remains true across
the periodic boundaries under the divisibility
assumptions.  Collinear first-direction links do
not share a plaquette.  The only two-free-link
plaquette is therefore the designated shared
plaquette of a pair.

Each link has six incident plaquettes and the two
links share one, giving eleven.  If plaquette sets
of two pairs intersected, that plaquette would
contain a free link from each, which was excluded.
Fubini's theorem now gives the conditional product.
This is independence of integrations \emph{given
the retained exterior}, not independence of the
messages after that exterior is integrated.
The last assertion follows from the excluded times
in \eqref{r22:anchors}.  Counting the $L$, $L/4$,
$L/2$ spatial choices and the permitted even times
proves \eqref{r22:pair-count}.
\end{proof}

\begin{theorem}[Normalized one-pass endpoint certificate]
\label{r22:one-pass}
Apply this packing to both independent histories in
$\mathbb P_q$, and replace each exact message by
the lower envelope \eqref{r22:finite-eq}.
Normalize the resulting positive retained density,
and call its four-endpoint marginal $\widehat m_q$.
For $\widehat A_q=H(\widehat m_q,S_*\widehat m_q)$,
where $S$ exchanges the two time-$q$ slices,
\[
 \left|\sqrt{1-R_{2q}}-\sqrt{1-\widehat A_q}\right|
 \le2\sqrt{1-(1-\delta)^{2J}}\le2\sqrt{2J\delta}.
\tag{V22.18}\label{r22:one-pass-eq}
\]
In the original Wilson model at
$\gamma=10,\ L=4,\ q=12$, there are $J=80$ pairs
per history, and
\[
 \left|\sqrt{1-R_{24}}-\sqrt{1-\widehat A_{12}}\right|
                  <3\cdot10^{-8},
 \qquad R_{24}=\frac{Z_{48}}{Z_{24}^2}.
\tag{V22.19}\label{r22:one-pass-point}
\]
\end{theorem}
\begin{proof}
Relative to the exact retained two-history law,
the unnormalized density is multiplied by
$r(\eta)=\prod_{j=1}^{2J}\widehat h_j(\eta)/h_j(\eta)$.
The all-exterior inequality \eqref{r22:relative-eq}
gives $(1-\delta)^{2J}\le r\le1$.
Its discarded probability is therefore at most
$1-(1-\delta)^{2J}\le2J\delta$ under this exact
normalized law, whatever the distribution of $\eta$.
Lift $r$ to the full histories and apply
Theorem~\ref{r21:filter-thm}, followed by the common
endpoint marginal and \eqref{r21:native-purity}.

At the stated point, \eqref{r22:pair-count} gives
$J=4^3(12-2)/8=80$.  The squared error upper
bound is less than $4\cdot160\cdot10^{-18}
=6.4\cdot10^{-16}<9\cdot10^{-16}$.
The incidence checker independently verifies
160 eliminated links and 880 touched plaquettes
per history out of its 9216 plaquettes.
\end{proof}

This is a continuous-Wilson statement.  It neither
assigns continuous cutoff probabilities to atomic
cubature nodes nor spends V19's plaquette-model
comparison error.  Exact pair integration itself
changes no endpoint law; the displayed error comes
only from approximating those connected messages.
No reflection positivity or single-transfer
representation is asserted for the approximate
retained density.

\subsection{The remaining normalized interaction is still the target}
\label{r22:next}

The one-pass certificate removes a specified set of
actual lattice variables with a paid small error.
It does not evaluate $\widehat A_{12}$, reduce all
remaining integrations to finitely many states, or
make the generated messages independent.
In particular the exterior links occur in several
different messages, even though the original
affected plaquette sets were disjoint.
The three-invariant representation is local to
one pair and is not a global three-dimensional
description of the retained measure.

The next elimination sees the generated
$\widehat h_j$ as well as the untouched plaquettes.
Its outgoing factor is not generally another
eleven-plaquette Wilson cell.  Reusing
\eqref{r22:relative-eq} for that next factor without
deriving its new dependence would be circular.
A useful continuation should control one such
overlapping generated-message integral, or directly
bound the four-endpoint affinity of the retained
law.  The native angular contrast shows why
discarding the interaction angle is not a viable
uniform lower-envelope shortcut.

\paragraph{Verification and preservation.}
\path{scripts/verify_ym_restart_v22_connected_pairs.py}
compares the positive beta formula with independent
$S^3$ monomial integration, checks all displayed
relative tails and the native ratio using exact
fractions, and enumerates the actual pair packing
for even and odd $q$.  It also reruns V21's filter
and normalization checks.  Finite fixtures test the
algebra; the general statements rely on the proofs
above.  No fixture evaluates the unknown endpoint
affinity.

Removing this append and restoring V21's title
recovers its entire 505063-byte source and SHA--256
\path{2A04A07D334AAE2228DB5070116A10B56B08D90CE8589703CB11BBCB2ABE9CB2}.
Only active V21 is replaced; the other 46 sources
are unchanged.  All non-source artifacts remain
outside \path{latex_notes}.

\paragraph{Status.}
A normalized error bound for an actual connected
weak-coupling decimation is now evaluated.
A useful weak-coupling global purity lower bound,
cofinal control and physical-time transport, the
interacting four-dimensional continuum construction,
and its positive Hamiltonian mass gap remain
unproved.

% END CONSOLIDATED V22 APPEND

% BEGIN CONSOLIDATED V23 APPEND -- 2026-09-19
\clearpage
\section{V23: ancestral reblocking and a certified shared-message overlap}
\label{r23:main}

V22's next obstruction is not the accuracy of an isolated pair
integral.  Two generated messages depend on the same retained
variables, and the next elimination must integrate their product.
Here we identify such a bucket on the actual periodic lattice,
recover its original 29 plaquettes without duplication, and give
a current-law lower replacement by recomputing those ancestors.
When every first-pass message is absorbed, the constant safety
factors cancel on normalization.  The final law can then be
compared directly with the original Wilson law, without retaining
an artificial first-pass error floor.

We also evaluate a genuinely shared-message integral at an
attainable conditional boundary.  Its overlap factor lies between
$1.0849$ and $1.0850$ at $\gamma=10$; replacing the shared average
by independent averages is not an identity.  This boundary
calculation is not a probability estimate for that boundary, and
the general remaining endpoint affinity is still unevaluated.

\subsection{Inherited tools and the targeted literature import}
\label{r23:inheritance}

The parallel-link ladder integrals in f16.2 V62 and induced
neighbour messages in f15 V69 already supply the geometric
framework.  V19 proves the positive local plaquette truncation,
V20 gives exact positive cubature and factor-graph elimination,
and V21 proves normalized filter stability.  None of these
general constructions is claimed anew.  The additional content
is the explicit second-pass ancestry partition, complete
absorption of the first pass, its normalization consequence, and
the evaluated overlap witness below.

For a targeted primary-source check we read Liu et al.,
\emph{Exact blocking formulas for spin and gauge models},
Phys.\ Rev.\ D 88 (2013), 056005,
\href{https://arxiv.org/html/1307.6543}{arXiv:1307.6543},
Sections II.1--II.2 and IV, especially equations (24)--(37).
Their character/Haar blocking is useful here.  Their discussion
also makes explicit that common exterior variables survive
partial integration, and that exact blocking does not remove
the growth of tensor indices.  We import that distinction, not
a certified truncation rate or a four-dimensional gap theorem.
The positive character contraction below is proved for its
specified central boundary; general recoupling tensors are not
assumed entrywise positive.

This is an intervening targeted check, not a replacement for
the user's broad ten-append survey.  The last broad survey was
V20 and the next is V30; the next companion-program checkpoint
remains V25.

\subsection{The actual second-pass bucket has six links and 29 ancestors}
\label{r23:geometry}

Work with normalized Haar measure on $\SU(2)$ and the original
four-dimensional periodic Wilson lattice
$(\mathbb Z/L\mathbb Z)^3\times\mathbb Z/(2q)\mathbb Z$,
where $4\mid L$, $L\ge4$, and $q\ge3$.
Coordinates are $(x_\mu,x_\nu,x_\rho,t)$, with the free links
oriented in the first spatial direction $\mu$.
The plaquette weight is
\[
 w(U_p)=e^{-\gamma(1-t_p)},\qquad
 t_p=\tfrac12\Tr U_p\in[-1,1].
\tag{V23.1}\label{r23:weight}
\]
Choose $a,b$ divisible by four and even $t\notin\{0,q\}$.
At fixed $x_\mu$ take the six $\mu$-links whose anchors have
$(x_\nu,x_\rho)=(a+c,b+r)$, $c=0,1$, $r=0,1,2$.
Write their group variables as
\[
 \begin{matrix}
  U_0&\!-\!&V_0\\[-2pt]
  |&&|\\[-2pt]
  u&\!-\!&v\\[-2pt]
  |&&|\\[-2pt]
  U_2&\!-\!&V_2 .
 \end{matrix}
\tag{V23.2}\label{r23:grid}
\]
The horizontal and vertical marks denote shared plaquettes,
not additional link variables.  The top and bottom pairs
are two of V22's eliminated pairs; the middle pair is retained
after V22 and is now to be eliminated.

\begin{lemma}[Exact ancestry and packed incidence]
\label{r23:incidence}
Precisely 29 distinct plaquettes touch these six links.
Seven touch two free links, and 22 touch one.  The two old
eleven-plaquette ancestor sets are disjoint; their complement
within the 29-plaquette set consists of seven bare plaquettes
touching the middle pair.

Taking all the displayed anchors gives
\[
 B=\frac{L^3}{16}\bigl(q-1-\mathbf1_{\{q\ {\rm even}\}}\bigr)
\tag{V23.3}\label{r23:block-count}
\]
blocks per history.  No plaquette touches free links in
different blocks, every old V22 pair belongs to exactly one
block, and every spatial link at times $0,q$ is retained.
\end{lemma}
\begin{proof}
A $\mu$-link meets six plaquettes, one in each signed transverse
direction.  The $2$ by $3$ rectangle has three horizontal
and four vertical adjacencies.  Thus the incidence total is
$6\cdot6=36$, with seven doubly counted plaquettes; the number
of distinct plaquettes is $36-7=29$.  Four corner vertices
have internal degree two and four private plaquettes each;
two middle vertices have degree three and three private
plaquettes each.  This gives $4\cdot4+2\cdot3=22$ private
plaquettes.  The old top and bottom rows are separated by
one row, so they share no plaquette.  Each has eleven
ancestors, leaving $29-22=7$: the middle rung and its six
private plaquettes.

Across blocks there are two unused columns between successive
two-column strips, one unused row between successive
three-row strips, and at least one unused time slice.
No pair of free parallel links in different blocks is
transversely adjacent.  Collinear $\mu$-links at successive
$x_\mu$ do not share a plaquette.  These observations also
cover periodic wraparound when $L=4$.  A plaquette has at
most two $\mu$-links, proving the disjointness assertion.

There are $L(L/4)(L/4)=L^3/16$ spatial anchors.  Among the
$q$ even time slices remove zero and also $q$ when $q$ is
even.  This gives \eqref{r23:block-count}.  The old even
$x_\rho$ rows are partitioned into rows $b,b+2$ with
$b\equiv0\pmod4$, so each old pair is absorbed exactly once.
The exclusion of times $0,q$ proves the endpoint claim.
\end{proof}

Let $\eta$ denote all links outside one six-link block.
Let $F_7(u,v;\eta)$ be the product of its seven new bare
plaquette weights, and $h_0,h_2$ its exact old pair messages.
Tonelli's theorem and the ancestry partition give
\[
 \begin{split}
 H(\eta)
 &=\int_{\SU(2)^2}F_7(u,v;\eta)h_0(u,v;\eta)
                         h_2(u,v;\eta)\,du\,dv\\
 &=\int_{\SU(2)^6}\prod_{p\in\mathcal P_{\rm block}}
                   w(U_p)\,dU_{\rm block}.
 \end{split}
\tag{V23.4}\label{r23:exact-bucket}
\]
All shared exterior variables have the same value in both
messages; there are no independently sampled exteriors.
The same identity is valid simultaneously for the packed
blocks conditional on the complete retained configuration.

\subsection{Recompute the ancestors and pay the current-law comparison}
\label{r23:recompute}

Suppose the first-pass messages satisfy
$\ell h_j\le\widehat h_j\le h_j$ for $j=0,2$, with one
constant $0<\ell\le1$ independent of $\eta,u,v$.
At $\gamma=10$, V22 permits $\ell=1-10^{-18}$.
The exact second-pass bucket of the \emph{current} approximate
law is
\[
 H_{\rm cur}(\eta)=\int F_7\widehat h_0\widehat h_2\,du\,dv,
 \qquad \ell^2H\le H_{\rm cur}\le H.
\tag{V23.5}\label{r23:current}
\]
This is not a bare two-link Wilson bucket.

For an integer $d\ge0$, recall V19's polynomial and its
pointwise lower plaquette weight:
\[
 p_d(z)=\sum_{k=0}^d\frac{z^k}{k!},\quad
 w_d(t)=e^{-2\gamma}p_d(\gamma(1+t)),\quad
 (1-\delta_d)w(t)\le w_d(t)\le w(t).
\tag{V23.6}\label{r23:poly}
\]
Here $\delta_d=1-e^{-2\gamma}p_d(2\gamma)$, or any certified
larger number smaller than one.  Monotonicity of the Poisson
upper tail in its mean gives the bound for every $t$.
Define a freshly computed ancestral integral, not a
polynomial fit to $\widehat h_0\widehat h_2$:
\[
 B_d(\eta)=\int_{\SU(2)^6}
                 \prod_{p\in\mathcal P_{\rm block}}w_d(t_p)
                     \,dU_{\rm block},
 \qquad r_d=(1-\delta_d)^{29}.
\tag{V23.7}\label{r23:B}
\]

\begin{theorem}[A normalized replacement with recoverable ancestry]
\label{r23:replacement}
For every exterior configuration,
\[
 r_dH\le B_d\le H,\qquad
 \widehat H_{\rm new}:=\ell^2B_d,\qquad
 \ell^2r_d H_{\rm cur}\le\widehat H_{\rm new}\le H_{\rm cur}.
\tag{V23.8}\label{r23:current-lower}
\]
Replace all the second-pass buckets in the packed graph by
$\widehat H_{\rm new}$.  Its normalized retained law is
exactly the law obtained directly from the original Wilson
model by replacing the 29 ancestral plaquettes of every
block by $w_d$ and integrating the six links.
In particular, the normalized final law does not depend on
the first-pass approximation error $\ell$.
\end{theorem}
\begin{proof}
Multiply the 29 inequalities in \eqref{r23:poly} and integrate
to obtain the bounds on $B_d$.  Then
$\ell^2B_d\le\ell^2H\le H_{\rm cur}$, while
$\ell^2B_d\ge\ell^2r_dH\ge\ell^2r_dH_{\rm cur}$.
This proves the one-sided current-law certificate; no
boundary-dependent division or unknown normalization is used.

Write $F_{\rm out}(\eta)$ for all untouched original
plaquettes.  By Lemma~\ref{r23:incidence}, all old messages
are absorbed and all ancestors occur once.  The final
unnormalized density for one history is
\[
 F_{\rm out}(\eta)\prod_{j=1}^B\ell^2B_{d,j}(\eta)
  =\ell^{2B}F_{\rm out}(\eta)\prod_{j=1}^B B_{d,j}(\eta).
\tag{V23.9}\label{r23:scalar}
\]
The scalar $\ell^{2B}$ cancels between the density and its
partition function.  The density on the right without that
scalar is precisely the directly recomputed ancestral model.
The same argument applies to the product of two histories.
\end{proof}

The conclusion concerns a \emph{different recomputation},
not automatic removal of an error from an existing
approximation.  It requires access to the original factors,
the exact ancestry partition, and absorption of every old
message.  If old approximate messages remain elsewhere,
their error is still present.  A boundary-dependent safety
factor also cannot be discarded as a scalar.  Nor does the
theorem make the $B_{d,j}$ independent after the common
exterior $\eta$ is integrated.

The integral $B_d$ is a finite polynomial calculation in
principle.  Each free link appears in six plaquettes and
has coordinate degree at most $6d$, so V20's positive
product cubature is exact.  At $d=80$ it uses
\[
 D=480,\qquad
 \bigl(\lfloor D/4\rfloor+1\bigr)(D+1)^2
       =121\cdot481^2=27\,994\,681
\tag{V23.10}\label{r23:cost}
\]
states per free link.  A dense six-fold product is enormous.
We have not performed the general exterior-dependent
$B_d$ contraction, compressed its retained boundary, or
proved an efficient algorithm for it.  Exact finite
representability was already V19--V20's result; the new
claim here is the correctly normalized ancestral
replacement for the actual overlapping bucket.

\subsection{An attainable zero-private-field boundary and its overlap}
\label{r23:theta}

There is one conditional boundary for which the exact
six-link integral simplifies considerably.  Set all
retained links not oriented in direction $\mu$ to $I$.
For each corner free link, set its external $\mu$-neighbour
in direction $\rho$ to $I$, its external $\mu$-neighbour
in direction $\nu$ to $I$, and its two temporal
$\mu$-neighbours to $-I$.  Its four private staples sum
to zero.  For each middle free link set its external
$\nu$-neighbour to $I$ and its temporal neighbours to
\[
 Q_\pm=-\tfrac12 I\ \mathord{\pm}\ \tfrac{\sqrt3}{2}\,\mathbf i,
 \qquad Q_\pm\in\SU(2),\qquad I+Q_++Q_-=0,
\tag{V23.11}\label{r23:boundary}
\]
where $\mathbf i$ is a unit imaginary quaternion.
All other retained $\mu$-links may be set to $I$.
These assignments are compatible: the only corner
external $\rho$-neighbours that coincide at $L=4$ receive
the same value $I$; the other specified neighbours are
distinct.  Since $2q\ge6$, the temporal neighbours are
distinct as well.  The checker enumerates this incidence.

The 22 private linear fields cancel exactly, and all seven
internal transports are the identity.  Put
\[
 k(g)=e^{\gamma\Tr(g)/2},\quad F=k*k*k,\quad
 z_1=\int k(g)\,dg,\quad Z_4=\int k(g)F(g)\,dg,
\tag{V23.12}\label{r23:k}
\]
where convolution uses normalized Haar measure,
$(f*g)(x)=\int f(xy^{-1})g(y)\,dy$.
The kernels are central and invariant under inversion.

\begin{proposition}[Exact native shared-message witness]
\label{r23:theta-id}
At boundary \eqref{r23:boundary} the exact original
29-plaquette integral is
\[
 H=e^{-29\gamma}Z_\theta,\qquad
 Z_\theta=\int k(g)F(g)^2\,dg.
\tag{V23.13}\label{r23:theta-Z}
\]
For the probability law $d\nu(g)=k(g)\,dg/z_1$, the ratio
between the shared-message average and the independent
replacement is
\[
 C_\gamma=\frac{\E_\nu F^2}{(\E_\nu F)^2}
          =\frac{z_1Z_\theta}{Z_4^2}
          =1+\frac{\Var_\nu F}{(\E_\nu F)^2}\ge1.
\tag{V23.14}\label{r23:overlap}
\]
\end{proposition}
\begin{proof}
Every original plaquette supplies its constant $e^{-\gamma}$.
The private exponents sum to zero, leaving exactly the seven
internal bonds of \eqref{r23:grid}.  Conditional on $u,v$,
the top row is a length-three path from $u$ to $v$; integrating
its two internal variables gives $F(uv^{-1})$.  The bottom
row gives the same function of the \emph{same} argument.
The middle rung contributes $k(uv^{-1})$.
Haar invariance reduces the last two integrals to the single
$g=uv^{-1}$ integral in \eqref{r23:theta-Z}.  Finally
$\E_\nu F=Z_4/z_1$ and $\E_\nu F^2=Z_\theta/z_1$;
the variance identity proves \eqref{r23:overlap}.
\end{proof}

This theta graph consists of three paths of lengths $1,3,3$
between the middle vertices.  Averaging its two length-three
messages independently would replace $Z_\theta$ by
$Z_4^2/z_1$.  The following evaluation measures the resulting
error at a native boundary.  It asserts neither positive
probability of that exact boundary nor prevalence of small
private fields in the full Wilson law.

\subsection{Finite positive character certificate for the overlap}
\label{r23:certificate}

Let $\chi_n$ be the $\SU(2)$ character of dimension $n+1$,
so $\chi_n(g)=U_n(\Tr(g)/2)$ with $U_n$ the Chebyshev
polynomial of the second kind.  For this calculation use
$\kappa_d(g)=p_d(\gamma(1+\Tr(g)/2))$; its common scalar
$e^{-\gamma}$ is omitted because it cancels from
\eqref{r23:overlap}.  Write
\[
 \kappa_d=\sum_{n=0}^d a_n\chi_n,\quad
 \lambda_n=\frac{a_n}{n+1},\quad
 f_d=\kappa_d^{*3}
       =\sum_{n=0}^d(n+1)\lambda_n^3\chi_n .
\tag{V23.15}\label{r23:characters}
\]
V19's coefficient recursion, recalled here to fix conventions,
starts at $b^{(0)}_n=\mathbf1_{\{n=0\}}$, sets
$b^{(k+1)}_n=b^{(k)}_n+
(b^{(k)}_{n-1}+b^{(k)}_{n+1})/2$ with $b^{(k)}_{-1}=0$,
and gives $a_n=\sum_{k=0}^d\gamma^kb^{(k)}_n/k!$.
Thus every coefficient is nonnegative.

\begin{theorem}[Evaluated overlap interval]
\label{r23:theta-certificate}
Define the finite positive sums
\[
 \begin{split}
 z_{1,d}&=\lambda_0,\qquad
 Z_{4,d}=\sum_{n=0}^d(n+1)^2\lambda_n^4,\\
 Z_{\theta,d}
 &=\sum_{i,j=0}^d(i+1)(j+1)\lambda_i^3\lambda_j^3
   \sum_{\substack{|i-j|\le n\le\min(i+j,d)\\
                    n\equiv i+j\ (2)}}(n+1)\lambda_n,\\
 C_d&=\frac{z_{1,d}Z_{\theta,d}}{Z_{4,d}^2}.
 \end{split}
\tag{V23.16}\label{r23:finite-theta}
\]
Then
\[
 (1-\delta_d)^8 C_d\le C_\gamma
                 \le (1-\delta_d)^{-8}C_d.
\tag{V23.17}\label{r23:theta-tail}
\]
Exact rational arithmetic at $\gamma=10,d=80$ gives
\[
 \delta_{80}<2\cdot10^{-24},\qquad
 \frac{10849}{10000}<C_{10}<\frac{10850}{10000}.
\tag{V23.18}\label{r23:theta-value}
\]
\end{theorem}
\begin{proof}
Haar orthogonality gives
$\chi_i*\chi_j=\mathbf1_{\{i=j\}}\chi_i/(i+1)$.
It follows either by integrating the representation matrix
coefficients or by Schur orthogonality and taking traces.
The identity
\[
 \chi_i\chi_j
   =\sum_{\substack{|i-j|\le n\le i+j\\ n\equiv i+j\ (2)}}\chi_n
\]
follows from the sine formula
$U_n(\cos\theta)=\sin((n+1)\theta)/\sin\theta$,
first away from the endpoints and then by polynomial
continuity.  Integrating against $\chi_n$ shows that
the triple character integral is one precisely on the
displayed parity/triangle range, and zero otherwise.
These identities prove all three sums in
\eqref{r23:finite-theta}.  For rational $\gamma$ they
involve only rational numbers and nonnegative summands.

The scaled polynomial kernel $e^{-\gamma}\kappa_d$ lies
between $(1-\delta_d)k$ and $k$.  Positivity of the graph
integrals gives relative retention $(1-\delta_d)^m$ for
any graph with $m$ bonds.  The numerator $z_1Z_\theta$
and denominator $Z_4^2$ each contain eight bonds.  Bounding
the numerator and denominator separately gives
\eqref{r23:theta-tail}.  The scale omitted in
\eqref{r23:characters} occurs eight times in each and cancels.

For a reproducible rational tail bound put
\[
 s=p_{80}(20),\qquad
 u=\frac{20^{81}}{81!}\frac1{1-20/82}.
\]
Successive omitted terms of $e^{20}$ have ratio at most
$20/82$, so $e^{20}-s\le u$ and
$\delta_{80}\le u/(s+u)<2/10^{24}$.
The checker evaluates \eqref{r23:finite-theta} with the
coefficient recursion, and verifies the two strict rational
comparisons
\[
 \frac{10849}{10000}
 <\left(1-\frac{u}{s+u}\right)^8 C_{80},
 \qquad
 \frac{C_{80}}{(1-u/(s+u))^8}<\frac{10850}{10000}.
\]
The full rational numerator and denominator are recorded in
the external check artifact.  For orientation only,
$C_{80}\simeq1.0849692753641804$ and
$u/(s+u)\simeq1.136971116945502\cdot10^{-24}$.
The decimal values are not the certificate.
\end{proof}

The seven-bond polynomial calculation here is performed
\emph{after exact cancellation of the private exponential
fields}.  It is not the same polynomial as the general
29-plaquette $B_d$, whose private polynomial factors need
not cancel.  The two constructions give valid, different
enclosures of the same exact conditional integral at this
boundary.  Keeping this distinction avoids assigning the
seven-bond truncation error to an arbitrary exterior.

\subsection{Endpoint-preserving second-pass error, with no inherited floor}
\label{r23:global-error}

Let $\widetilde{\mathbb P}_{q,d}$ be the normalized retained
two-history law obtained by the recomputation in
Theorem~\ref{r23:replacement}.  Let $K_q$ keep just the four
spatial endpoint slices at times $0,q$, and let $S$ swap
the two time-$q$ slices, as in V21.  Define
$\widetilde A_{q,d}=H(K_q\widetilde{\mathbb P}_{q,d},
S_*K_q\widetilde{\mathbb P}_{q,d})$.

\begin{corollary}[Actual packed second-pass comparison]
\label{r23:affinity}
For the original Wilson purity $R_{2q}=Z_{4q}/Z_{2q}^2$,
\[
 \left|\sqrt{1-R_{2q}}-\sqrt{1-\widetilde A_{q,d}}\right|
 \le 2\sqrt{1-(1-\delta_d)^{29B}}
 \le 2\sqrt{58B\delta_d}.
\tag{V23.19}\label{r23:affinity-bound}
\]
At $\gamma=10,L=4,q=12,d=80$,
\[
 B=40,\qquad
 \left|\sqrt{1-R_{24}}-\sqrt{1-\widetilde A_{12,80}}\right|
                  <2\cdot10^{-10}.
\tag{V23.20}\label{r23:affinity-point}
\]
There are 240 integrated links and 1160 replaced plaquettes
per history, with 9216 plaquettes in the original history.
\end{corollary}
\begin{proof}
By the scalar cancellation in \eqref{r23:scalar}, compare
the final law directly with the original Wilson law.
On two full histories the direct ancestral plaquette
replacement has density filter
\[
 r=\prod_{\text{two histories}}\ 
       \prod_{j=1}^B\prod_{p\in\mathcal P_j}\frac{w_d(t_p)}{w(t_p)},
 \qquad (1-\delta_d)^{58B}\le r\le1.
\]
Integration of the six-link blocks is exact and does not
change the endpoint law.  Thus
$\alpha=1-\E r\le1-(1-\delta_d)^{58B}$ under the original
normalized two-history law.  Apply
Theorem~\ref{r21:filter-thm} and the native identity
\eqref{r21:native-purity}.  This gives the first inequality
in \eqref{r23:affinity-bound}; its coarser last bound follows
from $1-(1-\delta_d)^{58B}\le58B\delta_d$.

At the stated point, $B=4^3(12-2)/16=40$.
The square of the coarse error is less than
$4\cdot2320\cdot2\cdot10^{-24}
=1.856\cdot10^{-20}< (2\cdot10^{-10})^2$.
The link and plaquette counts follow from
Lemma~\ref{r23:incidence}, and $6L^3(2q)=9216$.
\end{proof}

This is an error bound for a specified, well-defined
approximate endpoint law.  It is not an evaluation of
$\widetilde A_{12,80}$, an assertion of its closeness to
one, or an efficient computation of all $B_d$ messages.
No reflection positivity, new Hamiltonian purity, or
closed Wilson action is assigned to the approximating law.
The disappearance of the first-pass error from the direct
bound is exact bookkeeping after ancestral recomputation;
it does not establish contraction of the surviving
interaction.

\subsection{Verification and the next unresolved task}
\label{r23:next}

\path{scripts/verify_ym_restart_v23_ancestral_overlap.py}
checks the character fusion integral against independent
$S^3$ monomial moments in 343 cases, and the three theta
graph sums against direct polynomial integration in 21
cases.  The $\gamma=10$ overlap bounds are exact rational
comparisons.  Native incidence is enumerated at
$(L,q)=(4,12),(4,5),(8,4)$, including compatible boundary
assignments and absorption of every old pair.  Three
six-binary-spin positive factor-graph fixtures verify
the current-law inequalities and cancellation of a
constant, but not a boundary-dependent, multiplier.
V21's normalized filter tests are also rerun.
These fixtures supplement the general proofs; they
are not substitute four-dimensional Wilson calculations.

Deleting this append and restoring the displayed V22
title recovers its complete 527799-byte source and
SHA--256
\path{DE3C2190B08739650E955FE804B237D4B0E497398EAAE23495636C46215C1621}.
Only the immediate active V22 is replaced; the other 46
sources are unchanged.  Compilation products and check
artifacts remain outside \path{latex_notes}.

The useful direction is now more specific: obtain a
boundary-sensitive contraction or normalized endpoint
observable without enumerating the dense cubature product.
The theta witness shows an explicit dependence that must
survive such compression.  Its conditional ratio cannot
be multiplied over the lattice, because blocks retain
common exterior variables, and their test boundaries
have not been assigned any global Gibbs weight.
Repeating a general finite-dimensional existence argument
or making the scalar tail still smaller would not address
this computational and analytic bottleneck.

\paragraph{Status.}
We have a second-pass ancestry certificate, a normalized
recomputation theorem for the actual packing, and an
evaluated shared-message overlap at a native boundary.
The retained endpoint affinity, useful cofinal purity
control, physical-time transport, the interacting
four-dimensional continuum construction, and its
positive Hamiltonian mass gap remain unproved.

% END CONSOLIDATED V23 APPEND

% BEGIN CONSOLIDATED V24 APPEND -- 2026-09-19
\clearpage
\section{V24: normalized dependence and positive spectral compression}
\label{r24:main}

V23 identifies an actual overlapping bucket, but its scalar
overlap does not quantify the conditional operator or justify
compression in a changed exterior.  Here we compute that
operator at the zero-private-field boundary, distinguish a
small Hilbert--Schmidt error from a positive relative
certificate, and construct a finite-rank lower kernel that
survives arbitrary positive reweighting.  Its error also
controls the conditional operator with its \emph{own}
marginal normalizers.

This yields a useful extension beyond one boundary point:
the two middle private fields can vary throughout their
native radius-three balls while the corner fields remain
zero and internal transports remain the identity.  We
evaluate nonzero-field transverse correlations and reduce
the finite normalized operator to explicit Gram matrices.
The restriction on corner fields and transports is essential.
There is no claim that this boundary family carries positive
probability in the full Wilson exterior law, or that the
remaining four-endpoint affinity has been evaluated.

\subsection{Archive audit and the environment-sensitive literature input}
\label{r24:scope}

The identification of maximal correlation with a centered
conditional-expectation norm is inherited from f6 V18 and
f12 V36.  The normalized kernel
$J(u,v)/\sqrt{J_u(u)J_v(v)}$ and compact marginal-denominator
analysis occur in f16.1 V15.  The complete harmonic-band
axial decomposition is already proved in f15 V85.
We reuse those tools, not claim them as new results.
V19's character coefficients, V21's filter comparison,
and V23's actual theta-graph boundary are also retained.

We checked Zhao et al.,
\emph{Renormalization of tensor-network states}.
The source is
\href{https://arxiv.org/html/1002.1405}{arXiv:1002.1405};
the relevant passages are Section III.2,
equations (36)--(44), and Section III.3.
It distinguishes optimizing an isolated tensor from
controlling the contraction with its environment.
We use that distinction to select the required error
notion.  We do not import a numerical truncation heuristic
as a rigorous relative bound, or assume its environment
tensor is known for the present Wilson problem.
The broad review schedule remains V20, V30, and so on;
the next companion review is V25.

\subsection{The normalized native theta kernel and its full spectrum}
\label{r24:central}

Use V23's boundary with all six private fields zero,
at $\gamma=10$.  With $k(g)=e^{10\Tr(g)/2}$,
$F=k^{*3}$ and $Z_\theta=\int kF^2$, the marginal
density of the middle variables is
\[
 p(u,v)=p(uv^{-1}),\qquad
 p(g)=\frac{k(g)F(g)^2}{Z_\theta},\qquad
 (Pf)(u)=\int p(uv^{-1})f(v)\,dv .
\tag{V24.1}\label{r24:p}
\]
Its two marginals are Haar probability, so $P$ is the
actual conditional-expectation operator for this
specified native pair law, not a physical time transfer.
It is self-adjoint and compact, and $P1=1$.

Let $\kappa=p_{80}(10(1+t))$, $t=\Tr(g)/2$, and
$f=\kappa^{*3}$.  Use V23's rational character
coefficients $a_n$ and $\lambda_n=a_n/(n+1)$.
Define $b_n$ by finite character multiplication:
\[
 \kappa f^2=\sum_{n=0}^{240}b_n\chi_n,\qquad
 p_d=\frac{\kappa f^2}{b_0}
      =\sum_{n=0}^{240}(n+1)m_n\chi_n,\qquad
 m_n=\frac{b_n}{(n+1)b_0}.
\tag{V24.2}\label{r24:pd}
\]
Here $b_0=Z_{\theta,80}>0$, $m_0=1$, and every $b_n>0$.
Character fusion supplies a finite sum of positive
rationals for every coefficient.  The convolution
operator $P_d$ has eigenvalue $m_n$ on the complete
degree-$n$ harmonic space of dimension $(n+1)^2$,
and zero on the remaining degrees.

For clarity, fix a rational tail bound
\[
 s=p_{80}(20),\quad
 u=\frac{20^{81}}{81!(1-20/82)},\quad
 \bar\delta=\frac{u}{s+u}<2\cdot10^{-24},\quad
 r=(1-\bar\delta)^7.
\tag{V24.3}\label{r24:r}
\]
The seven-bond integral and its normalizer are both
approximated, giving
\[
 r p\le p_d\le r^{-1}p,\qquad
 \|P-P_d\|_{\rm op}\le2(1-r),\qquad
 \|P-P_d\|_{\rm HS}
       \le(r^{-1}-1)\|p_d\|_{L^2}.
\tag{V24.4}\label{r24:paid}
\]
Indeed, the unnormalized seven-bond replacement is a
filter bounded between $r$ and one.  Normalizing proves
the pointwise comparison; the normalized laws have
total variation at most $1-r$.  Young's convolution
inequality gives the operator bound from the $L^1$
bound.  The pointwise comparison gives
$|p-p_d|\le(r^{-1}-1)p_d$ and hence the last inequality.
For convolution on a probability compact group the
Hilbert--Schmidt norm is the $L^2$ norm of its kernel.

\begin{theorem}[Computer-assisted normalized dependence certificate]
\label{r24:spectrum}
The maximal correlation of the pair law \eqref{r24:p}
satisfies
\[
 0.90847<\rho_{\max}(u,v)<0.90848.
\tag{V24.5}\label{r24:rho}
\]
It is attained on the degree-one space, of dimension
four.  Every multiplier in degree $n\ge2$ has absolute
value below $0.775$.
\end{theorem}
\begin{proof}
Character orthogonality diagonalizes the central
convolution.  For a joint law with Haar marginals,
the supremum of $\E[g(u)h(v)]$ over centered unit
$L^2$ functions is $\|P\|_{L^2_0\to L^2_0}$, by
Cauchy--Schwarz and the definition of operator norm.
This is the inherited maximal-correlation identity.

For the explicit rational data \eqref{r24:pd}, the
checker verifies
\[
 \begin{split}
 0.90847&<m_1-2(1-r)<m_1+2(1-r)<0.90848,\\
 \max_{2\le n\le240}m_n+2(1-r)&<0.775,\qquad
 2(1-r)<0.775 .
 \end{split}
\tag{V24.6}\label{r24:rho-check}
\]
Outside the polynomial band $P_d=0$.
Both $P$ and $P_d$ are scalar on each complete
harmonic degree, so \eqref{r24:paid} bounds each
multiplier error by $2(1-r)$, including the infinite
complement.  The displayed separation proves the
theorem.  These are exact rational comparisons,
not a floating-point eigensolver calculation.
The diagnostic value is $m_1\simeq0.9084710694871205$.
\end{proof}

In particular, replacing the middle pair by independent
Haar variables discards a large normalized mode,
despite the modest scalar overlap in V23.  Conversely,
$\rho_{\max}<1$ for this local pair does not give a
volume-uniform physical mass gap.

\subsection{A small norm tail is not the required positive certificate}
\label{r24:tails}

For $0\le N\le240$, put
\[
 p_{d,N}(g)=\sum_{n=0}^N(n+1)m_n\chi_n(g),\quad
 T_N=\sum_{n>N}(n+1)^2m_n^2,\quad
 \varepsilon_N=\sum_{n>N}(n+1)^2m_n.
\tag{V24.7}\label{r24:tails-def}
\]
The rank is
$D_N=\sum_{n=0}^N(n+1)^2=(N+1)(N+2)(2N+3)/6$.
Orthogonality gives
$\|P_d-P_{d,N}\|_{\rm HS}^2=T_N$, whereas
$|\chi_n|\le n+1$ gives the pointwise bound
$|p_d-p_{d,N}|\le\varepsilon_N$.
These two errors must not be interchanged.

The exact checker proves
\[
 \|P-P_{d,32}\|_{\rm HS}<1.2\cdot10^{-12},
 \qquad D_{32}=12529,
\tag{V24.8}\label{r24:hs-small}
\]
using $\|p_d\|_2^2<30$,
$T_{32}<(1.2\cdot10^{-12}-(r^{-1}-1)11/2)^2$,
and \eqref{r24:paid}.  Nevertheless, it also verifies
\[
 p_{d,31}(-I)<0,\qquad
 p_{d,32}(-I)-\varepsilon_{32}<0 .
\tag{V24.9}\label{r24:negative}
\]
Thus one nearby small-rank truncation is not a density,
and the degree-32 tail-subtracted lower approximation
fails positivity.  The second assertion does not say
that every degree-32 construction fails; it tests
this explicit construction.

Here is a positive replacement with a certified
pointwise relative budget.  Let
\[
 m_*=\min_{g\in\SU(2)}p_d(g),\quad
 \tau_N=\frac{2\varepsilon_N}{m_*},\qquad
 q_N(g)=r\bigl(p_{d,N}(g)-\varepsilon_N\bigr).
\tag{V24.10}\label{r24:q}
\]

\begin{theorem}[Explicit positive finite-rank lower kernel]
\label{r24:positive}
At $N=52$,
\[
 0<\ell p(g)\le q_{52}(g)\le p(g),\qquad
 \ell=r^2(1-\tau_{52})>1-10^{-10}.
\tag{V24.11}\label{r24:lower}
\]
The integral operator with kernel $q_{52}(uv^{-1})$
has rank $D_{52}=51039$.  It is positive semidefinite
and pointwise strictly positive, but is not yet
normalized to preserve constants.
\end{theorem}
\begin{proof}
The tail bound gives
$p_{d,N}-\varepsilon_N\le p_d$ and
$p_{d,N}-\varepsilon_N\ge p_d-2\varepsilon_N
\ge(1-\tau_N)p_d$.
Multiplying by $r$ and using \eqref{r24:paid} proves
the two comparisons, provided $\tau_N<1$.

We specify how the finite minimum is certified.
Expand the known polynomial $f=\kappa^{*3}$ in
$s_0=1+t$:
$f(t)=\sum_{j=0}^{80}c_js_0^j$.
The checker calculates these coefficients exactly
from \eqref{r24:pd} and verifies all 81 inequalities
$c_j\ge0$.  This is a finite arithmetic certificate
at $\gamma=10,d=80$, not an unproved general
positivity assertion about shifted coefficients.
Since $\kappa$ also has nonnegative $s_0$ coefficients
and $\kappa(-1)=1$, the minimum of $\kappa f^2/b_0$
is attained at $t=-1$, with
\[
 m_*=\frac{c_0^2}{b_0},\qquad
 1.08\cdot10^{-16}<m_*<1.09\cdot10^{-16}.
\tag{V24.12}\label{r24:minimum}
\]
The remaining exact comparisons are
\[
 \tau_{52}<9.9\cdot10^{-11},\qquad
 r^2(1-\tau_{52})>1-10^{-10}.
\tag{V24.13}\label{r24:loss}
\]
They prove strict positivity and \eqref{r24:lower}.

The convolution multipliers of $q_{52}$ are
$r(1-\varepsilon_{52})$ in degree zero and
$r m_n$ in degrees $1,\ldots,52$.  All are positive,
and all higher ones vanish.  Their multiplicities
give the rank and positive semidefiniteness.
Its integral is $r(1-\varepsilon_{52})$, which
need not be one.
\end{proof}

The diagnostic values
$\varepsilon_{52}\simeq5.34885\cdot10^{-27}$ and
$1-\ell\simeq9.89510\cdot10^{-11}$ explain the
chosen degree.  The proof uses the rational
comparisons above.  Positivity of the kernel does
not assert nonnegative entries in every harmonic
tensor basis, nor an efficient global contraction.

\subsection{Positive reweighting and the correct conditional normalizers}
\label{r24:reweight}

Let $W(u,v)\ge0$ be measurable with
$0<\int Wp<\infty$.  Replacing $p$ by $q_{52}$ in
the unnormalized density $Wp$ gives a filter
between $\ell$ and one, for \emph{this same} $W$.
Consequently the two normalized joint laws satisfy
\[
 \|\mu_W-\widehat\mu_W\|_{\rm TV}\le1-\ell<10^{-10},
 \qquad H(\mu_W,\widehat\mu_W)\ge\sqrt{\ell}.
\tag{V24.14}\label{r24:weighted-law}
\]
There is no factor involving $\sup W/\inf W$.
Indeed, if $z=\E_{\mu_W}(q_{52}/p)\ge\ell$, then
$\mu_W=z\widehat\mu_W+(1-z)\mu_{\rm discarded}$.
This proves the total-variation bound; V21 proves
the affinity bound.  Bounded observables of norm
at most one differ in expectation by at most
$2(1-\ell)$.  A common marginal and swap also
have root-affinity-defect difference at most
$2\sqrt{1-\sqrt{\ell}}$.

There is a stronger operator consequence when
the two marginals are included rather than
silently held fixed.

\begin{lemma}[Relative lower kernels control normalized conditional operators]
\label{r24:operator-lemma}
Let $J,\widehat J$ be positive integrable joint
densities, with $\ell J\le\widehat J\le J$,
$0<\ell\le1$.  Write
$J_u(u)=\int J(u,v)\,dv$ and
$J_v(v)=\int J(u,v)\,du$, and define the
corresponding normalized kernels on the fixed
reference $L^2$ spaces by
\[
 C_J(u,v)=\frac{J(u,v)}{\sqrt{J_u(u)J_v(v)}},
 \qquad
 C_{\widehat J}(u,v)=
       \frac{\widehat J(u,v)}
            {\sqrt{\widehat J_u(u)\widehat J_v(v)}}.
\tag{V24.15}\label{r24:C}
\]
Then
\[
 \|C_{\widehat J}-C_J\|_{\rm op}\le\ell^{-1}-1.
\tag{V24.16}\label{r24:op-robust}
\]
For compact kernels every singular value, including
the second singular value representing maximal
correlation, changes by at most this amount.
\end{lemma}
\begin{proof}
Integrating the relative comparison gives the same
comparison for each marginal.  Therefore
$\ell C_J\le C_{\widehat J}\le\ell^{-1}C_J$
pointwise and
$|C_{\widehat J}-C_J|\le(\ell^{-1}-1)C_J$.
For any reference-space functions $a,b$,
Cauchy--Schwarz under the unnormalized measure
$J(u,v)\,du\,dv$ gives
\[
 \int C_J(u,v)|a(u)b(v)|\,du\,dv
 \le
 \left(\int J\frac{|a(u)|^2}{J_u(u)}\,du\,dv\right)^{1/2}
 \left(\int J\frac{|b(v)|^2}{J_v(v)}\,du\,dv\right)^{1/2}
 =\|a\|_2\|b\|_2 .
\]
Thus $C_J$ has norm one, attained on the square
roots of its normalized marginals.  The pointwise
domination proves \eqref{r24:op-robust}.  The
usual singular-value min--max characterization,
or the rank-approximation characterization,
gives the last assertion by the triangle inequality.
The total partition function cancels from
\eqref{r24:C}; neither marginal cancels.
\end{proof}

Applied to $J=Wp$, $\widehat J=Wq_{52}$, this gives
operator error below $2\cdot10^{-10}$ whenever
the conditional kernels are defined as above.
For general nonseparable $W$, the relative bound
survives but the rank need not.  For separate
weights $W=A(u)B(v)$, left and right multiplication
preserve rank, including after division by the
new marginal square roots.  This distinction is
important for the native application.

\subsection{A native nonzero-field family and explicit transverse correlations}
\label{r24:fields}

In V23's six-link block, keep the four corner
private fields zero and every internal transport
equal to the identity.  Vary the three private
staples at each middle link independently.  They
do not occur in any corner private plaquette.
Thus every pair of fields $a,b\in\mathbb R^4$
with $|a|,|b|\le3$ is attainable, and the normalized
middle law is
\[
 d\mu_{a,b}(u,v)
  =\frac{e^{10(a\cdot u+b\cdot v)}p(uv^{-1})}
         {\int e^{10(a\cdot u+b\cdot v)}p(uv^{-1})\,du\,dv}
                         \,du\,dv .
\tag{V24.17}\label{r24:family}
\]
To see attainability, the six exterior $\mu$-links
providing the middle private staples are distinct
and disjoint from the corner private neighbours.
Any vector $se_0$, $0\le s\le3$, is a sum of three
unit quaternions: choose the first as $e_0$ for
$s\ge1$, or $-e_0$ for $s<1$, and choose the other
two with common scalar component $(s-1)/2$ or
$(s+1)/2$ and opposite transverse components.
A common quaternion rotation supplies any direction.
V23's compatible corner assignment remains unchanged.

The degree-52 kernel therefore approximates the
\emph{own-marginal normalized conditional operator}
of every law \eqref{r24:family} within
$2\cdot10^{-10}$ in operator norm, at finite rank
at most 51039.  This is uniform over the whole
two-field family, not over all corner fields
and transports.

For a concrete dependence test take $a=3e_0$,
$b=3\sigma e_0$, $\sigma=1$ or $-1$, and
write $\mu_\sigma$ for the corresponding law.
Joint transverse rotations give
$\E u_1=\E v_1=0$ and
$\E u_1^2=\E v_1^2=:v_\sigma$.
Set $M_\sigma=\E u_0$ and $S_\sigma=\E u_0^2$.
Then $v_\sigma=(1-S_\sigma)/3$.

\begin{proposition}[Certified surviving dependence at nonzero fields]
\label{r24:field-values}
The transverse Pearson correlations satisfy
\[
 \begin{split}
 0.34097&<\operatorname{Corr}_{\mu_+}(u_1,v_1)<0.34098,\\
 0.84221&<\operatorname{Corr}_{\mu_-}(u_1,v_1)<0.84222.
 \end{split}
\tag{V24.18}\label{r24:field-corr}
\]
These are lower witnesses for the respective maximal
correlations, not evaluations of those full operator
norms.  If
$Z_\sigma=\int e^{30u_0+30\sigma v_0}p(uv^{-1})\,du\,dv$,
then
\[
 5.83\cdot10^{12}<Z_+/Z_-<5.84\cdot10^{12}.
\tag{V24.19}\label{r24:field-Z}
\]
\end{proposition}
\begin{proof}
The joint rotation generator in the $(0,1)$ plane
preserves Haar measure and $u\cdot v$.
Apply its integral identity to $u_1$ times the
density.  With the convention
$Lu_1=u_0$, $Lu_0=-u_1$, it gives
\[
 M_\sigma=30\,\E u_1^2+30\sigma\,\E u_1v_1 .
\]
Exchange symmetry, or $(u,v)\mapsto(-v,-u)$
when $\sigma=-1$, identifies the two transverse
variances.  Hence the exact Ward identities are
\[
 \operatorname{Corr}_{\mu_+}(u_1,v_1)
       =\frac{M_+}{30v_+}-1,\qquad
 \operatorname{Corr}_{\mu_-}(u_1,v_1)
       =1-\frac{M_-}{30v_-}.
\tag{V24.20}\label{r24:Ward}
\]
No asymptotic Gaussian approximation is used.

For reproducible finite arithmetic expand
$p_{160}(30(1+t))=\sum A_n\chi_n(t)$.
Let $A^{[1]}$ and $A^{[2]}$ be the character
coefficients after multiplying by $t$ and $t^2$,
respectively.  Multiplication by $t$ is the
recursion $A^{[1]}_n=(A_{n-1}+A_{n+1})/2$,
with $A_{-1}=0$ and coefficients outside the
finite range zero.  With the $m_n$ in
\eqref{r24:pd}, define
\[
 z_\sigma=\sum_n m_n A_n\sigma^n A_n,\quad
 \widehat M_\sigma=
    \frac{\sum_n m_n A^{[1]}_n\sigma^n A_n}{z_\sigma},
 \quad
 \widehat S_\sigma=
    \frac{\sum_n m_n A^{[2]}_n\sigma^n A_n}{z_\sigma}.
\tag{V24.21}\label{r24:axis-sums}
\]
Only indices common to the indicated finite arrays
are summed.  The identity
\[
 \int\chi_i(u)p_d(uv^{-1})\chi_j(v)\,du\,dv
       =m_i\mathbf1_{\{i=j\}}
\]
proves these formulas.
Although the opposed-field sums alternate, they
are evaluated as exact fractions and their
normalizers are positive integrals.

The exponential tilt approximation has relative
loss at most
\[
 \bar\delta_{\rm t}
 =\frac{u_{\rm t}}{p_{160}(60)+u_{\rm t}},\qquad
 u_{\rm t}=\frac{60^{161}}{161!(1-60/162)} .
\]
Together with the seven internal bonds, the
original unnormalized graph filter has loss
\[
 \alpha_*=1-(1-\bar\delta)^7
                     (1-\bar\delta_{\rm t})^2<2\cdot10^{-23}.
\tag{V24.22}\label{r24:axis-loss}
\]
Use this original-graph comparison, not an
uncorrected comparison of normalized factors.
The bounded moments $u_0,u_0^2$ differ from
\eqref{r24:axis-sums} by at most $2\alpha_*$.
Consequently
\[
 \frac{1-\widehat S_\sigma-2\alpha_*}{3}
 \le v_\sigma\le
 \frac{1-\widehat S_\sigma+2\alpha_*}{3}.
\]
Both lower variance endpoints are strictly
positive.  Rational interval division in
\eqref{r24:Ward} proves \eqref{r24:field-corr}.
The same original-graph bound puts $Z_+/Z_-$
between $(1-\alpha_*)z_+/z_-$ and
$(1-\alpha_*)^{-1}z_+/z_-$, proving
\eqref{r24:field-Z}.  Common exponential scalars
and the central normalizer cancel in this ratio.
\end{proof}

For orientation, the moment diagnostics are
$\widehat M_+\simeq0.96295324$,
$\widehat S_+\simeq0.92819004$,
$\widehat M_-\simeq0.72395228$, and
$\widehat S_-\simeq0.54116776$.
In particular, sizeable private fields do not
by themselves remove the surviving transverse
dependence.  None of these conditional moments
is a global Wilson expectation.

\subsection{The finite own-marginal problem and the archived axial reduction}
\label{r24:Gram}

For separate positive continuous weights $A,B$,
let $Q$ denote convolution by $q_{52}$ and choose
a real Haar-orthonormal harmonic basis $\phi_i$
of degrees at most 52.  Write its positive
multipliers as $\Lambda=\operatorname{diag}(\lambda_i)$,
so $q_{52}(uv^{-1})=\sum_i\lambda_i\phi_i(u)\phi_i(v)$.
The normalized conditional kernel of the
approximating joint law $A(u)q_{52}(uv^{-1})B(v)$ is
\[
 \widehat C=
 M_{\sqrt{A/(QB)}}\,Q\,M_{\sqrt{B/(QA)}} .
\tag{V24.23}\label{r24:weighted-C}
\]
Here $(QB)(u)=\int q_{52}(uv^{-1})B(v)\,dv$;
the symmetry of $Q$ gives the other marginal.
All denominators are strictly positive.

Define finite positive Gram matrices
\[
 \begin{split}
 (G_L)_{ij}
 &=\sqrt{\lambda_i\lambda_j}
      \int\frac{A(u)}{(QB)(u)}\phi_i(u)\phi_j(u)\,du,\\
 (G_R)_{ij}
 &=\sqrt{\lambda_i\lambda_j}
      \int\frac{B(v)}{(QA)(v)}\phi_i(v)\phi_j(v)\,dv .
 \end{split}
\tag{V24.24}\label{r24:Grams}
\]
The nonzero squared singular values of
$\widehat C$ are the eigenvalues of
$G_L^{1/2}G_R G_L^{1/2}$.  Indeed,
factor $\widehat C=XY^*$ using the two weighted
feature maps in \eqref{r24:weighted-C};
then $X^*X=G_L$, $Y^*Y=G_R$.
Polar decomposition of $X$ proves the assertion.
The leading singular value is exactly one,
with the correct marginal square-root vectors;
it must be removed when computing maximal correlation.

For the axial fields in \eqref{r24:field-corr},
$A(u)=e^{30u_0}$ and $B(v)=e^{30\sigma v_0}$.
The functions $A/(QB)$ and $B/(QA)$ depend only
on the scalar quaternion coordinate.  Apply
the already established f15 V85 basis
\[
 \psi_{\ell km}(t,\omega)
   =h_{\ell k}^{-1/2}(1-t^2)^{\ell/2}
          C_k^{(\ell+1)}(t)Y_{\ell m}(\omega),\qquad
 h_{\ell k}=\frac{(k+2\ell+1)!}
                  {4^\ell(k+\ell+1)(\ell!)^2k!}.
\]
The two Grams split into $\ell=0,\ldots,52$ blocks,
of sizes $53-\ell$, each repeated $2\ell+1$ times.
For example, before the multiplier square roots,
the left block entries are
\[
 \frac{2}{\pi\sqrt{h_{\ell k}h_{\ell j}}}
 \int_{-1}^1
       \frac{A(t)}{(QB)(t)}
       (1-t^2)^{\ell+1/2}
       C_k^{(\ell+1)}(t)C_j^{(\ell+1)}(t)\,dt.
\tag{V24.25}\label{r24:radial}
\]
Angular orthogonality proves the block structure,
and the archived dimension identity
$\sum_{\ell=0}^{52}(2\ell+1)(53-\ell)=51039$
shows that no harmonic sector is omitted.
Thus there are 53 matrix blocks, with largest
size 53, rather than one unrestricted dense
51039 by 51039 problem.

This is an application of the archived axial
basis to the \emph{new own-marginal weights}
$A/(QB)$ and $B/(QA)$, not a new derivation of
that basis or an identification with the
archived magnetic-weight Grams.
The entries in \eqref{r24:radial} contain
nonpolynomial denominators: V20's polynomial
cubature is not automatically exact for them.
Their rigorous interval evaluation and complete
noncentral maximal-correlation bounds are not
performed in this note.  The proved error
\eqref{r24:op-robust} would transfer such a
certified finite-block calculation to the true
native two-field law.

\subsection{Verification, limitations, and the next useful calculation}
\label{r24:next}

\path{scripts/verify_ym_restart_v24_normalized_compression.py}
contains the rational certificate.  It checks
character products against independent monomial
integration in 21 cases, computes the full
241-coefficient central polynomial, checks all
shifted-path coefficients and spectral tails,
and evaluates the nonzero-field moment intervals.
It also enumerates the six independently variable
middle private links, tests reweighting and
own-marginal kernel comparisons on finite
positive laws, and checks axial band dimensions.
As an independent check of the Ward computation,
it expands $u_1p_{160}(30(1+u_0))$ in transverse
harmonics using
$U_n=(C_n^{(2)}-C_{n-2}^{(2)})/(n+1)$.
The squared Haar norm of $u_1C_k^{(2)}(u_0)$ is
$(k+1)(k+3)/12$, and its degree is $k+1$.
Direct contraction with the multipliers $m_{k+1}$
gives correlation intervals agreeing with
\eqref{r24:field-corr}, without using the Ward identity.
The general statements follow from the proofs;
the numerical intervals are explicitly
computer-assisted finite rational results.

Deleting this append and restoring the V23 title
recovers its entire 550406-byte source with SHA--256
\path{CD0217DC9078BC0C42256F5FF919EF363637D2B43F26EBF27FF3BA9A5C29CDB1}.
Only the immediate active predecessor is replaced.
All other 46 sources remain unchanged, and
non-source artifacts stay outside \path{latex_notes}.

The next concrete calculation is to certify
the own-marginal axial Gram blocks, or to
extend the positive relative kernel construction
to nonzero corner fields and nontrivial internal
transports.  Either extension must keep its
actual exterior dependence.  Writing
$W=J_\eta/p$ for an arbitrary unknown bucket
does not solve that bucket: the quotient
may be nonseparable and as difficult as
$J_\eta$ itself.  Nor may the central
maximal correlation or the two-field
Pearson witnesses be multiplied over blocks
sharing retained exterior links.

\paragraph{Status.}
This note supplies a normalized local spectrum,
a certified positive finite-rank kernel,
an own-marginal perturbation bound, and native
nonzero-field dependence tests.  It does not
supply the four-dimensional endpoint affinity,
a contraction across general exteriors,
cofinal physical-time control, the interacting
continuum construction, or its positive
Hamiltonian mass gap.

% End the contents on its own page, including after future appends.
\AtEndDocument{\addtocontents{toc}{\protect\clearpage}}

% END CONSOLIDATED V24 APPEND

% BEGIN CONSOLIDATED V25 APPEND -- 2026-09-19
\clearpage
\section{V25: Certified axial dependence and open exterior neighbourhoods}
\label{r25:context}

V24 reduced the own-marginal conditional operator of a native six-link
bucket to finite positive-kernel Gram matrices, but evaluated only
coordinate-correlation lower witnesses at nonzero fields. This append
closes that finite spectral calculation. At $\gamma=10$, the full maximal
correlations of its aligned and opposed radius-three middle-field examples obey
\[
 \boxed{\quad
 0.34101369<\rho_+<0.34101371,\qquad
 0.87544428<\rho_-<0.87544430 .
 \quad}
 \tag{V25.1}\label{r25:main}
\]
The opposed value is appreciably larger than V24's transverse Pearson
witness. A selected coordinate is not an extremizing function.

We extend these certificates to open neighbourhoods of the actual
retained-link configurations, permitting small nonzero corner fields
and nontrivial internal transports. These neighbourhoods have an explicit
positive Wilson-probability lower bound. The bound is extremely small
and does not control their complement. Neither a local maximal correlation
nor the existence of favourable exteriors is the required four-endpoint
affinity or a physical Hamiltonian mass gap.

\subsection{Inherited objects, companion review, and nonduplication}
\label{r25:review}

Use precisely V23's native six-link geometry and V24's two-field family.
Four corner links are integrated out; the middle links are $u,v\in S^3$.
At the two reference exteriors their unnormalized density is
\[
 J_\sigma(u,v)=e^{30u_0}p(uv^{-1})e^{30\sigma v_0},
 \qquad \sigma\in\{1,-1\}.
 \tag{V25.2}\label{r25:joint}
\]
Here $p=k(k*k*k)^2/Z_\theta$ and $k(g)=e^{10\operatorname{Tr}(g)/2}$
are V24's central theta density and single-bond kernel.
For a positive density $J$ relative to product normalized Haar, set
\[
 C_J(u,v)=\frac{J(u,v)}{\sqrt{J_u(u)J_v(v)}},
 \qquad J_u(u)=\int J(u,v)\,dv,\quad
 J_v(v)=\int J(u,v)\,du .
 \tag{V25.3}\label{r25:own}
\]
Its first singular value is one and its second is the maximal
correlation. The total scalar normalizer cancels; the marginal functions do not.

The scheduled V25 companion review rehashed all four current companion
sources and reread their terminal arguments. All four are unchanged
since V20. NS V26's rank-one Oseen-kernel derivative norms concern a
different operator. SM V72's fixed-mode, fixed-time many-copy transport
is not a growing-cutoff physical continuum theorem. Shell-mixing V16
retains its noncommuting source-plane and nonperturbative remainder
obligations. Parallel YM V5's energy--entropy and tensorization warnings
do not provide a native endpoint-affinity estimate. No numerical constant
is imported from them. Both the next companion review and the next broad
literature checkpoint are V30.

The axial harmonic basis is inherited from f15 V85, as credited in V24.
The f16.1 V49 $LDL^{\mathsf T}$ calculations concern a different
constant-mode matrix Hamiltonian. General interval arithmetic,
finite-energy bounds, and spectral perturbation are standard tools,
not new general theorems here. The new evaluated objects are the
own-marginal spectra of \eqref{r25:joint}, their complete sector audit,
and their quantified retained-link neighbourhoods.

\subsection{Positive approximations and the common error ledger}
\label{r25:approx}

Retain V24's exact rational multipliers $m_n$, $0\leq n\leq240$,
from its degree-80 approximation of seven internal Wilson bonds. Put
\[
 \epsilon=\sum_{n=53}^{240}(n+1)^2m_n,\quad
 \lambda_0=1-\epsilon,\quad \lambda_n=m_n\ (1\leq n\leq52),\qquad
 q_0(g)=\sum_{n=0}^{52}(n+1)\lambda_n\chi_n(g).
 \tag{V25.4}\label{r25:q}
\]
Thus $q_0=q_{52}/r$ with $r=(1-\delta_{80})^7$ as in V24.
Removing a common positive scalar changes neither a normalized joint
law nor its conditional operator. All $\lambda_n$ and the pointwise
kernel $q_0$ are strictly positive by V24.

For the sources use
\[
 a(t)=p_{160}(30(1+t)),\qquad
 \widehat A(u)=e^{-30}a(u_0),\quad
 \widehat B_\sigma(v)=e^{-30}a(\sigma v_0).
 \tag{V25.5}\label{r25:source}
\]
Let $\delta_{\rm s}$ be V19's rational upper bound on the
$\operatorname{Poisson}(60)$ tail above 160. Then
$(1-\delta_{\rm s})e^{30t}\leq e^{-30}a(t)\leq e^{30t}$.
If $\ell_{24}$ is V24's degree-52 kernel retention, the approximating
joint density with $q_{52}$ lies between $LJ_\sigma$ and $J_\sigma$, where
\[
 L=\ell_{24}(1-\delta_{\rm s})^2>1-10^{-10},
 \qquad L^{-1}-1<2\cdot10^{-10}.
 \tag{V25.6}\label{r25:budget}
\]
These are exact rational checks. V24's own-marginal comparison gives
\[
 \|C_{J_\sigma}-\widehat C_\sigma\|_{\rm op}<2\cdot10^{-10}.
 \tag{V25.7}\label{r25:operator-error}
\]
One may equivalently compute $\widehat C_\sigma$ using
$\widehat A q_0\widehat B_\sigma$. Each side uses its own marginals.
No exponentially large source-condition-number factor enters this budget.

\subsection{A signed self-adjoint reduction for the opposed source}
\label{r25:reflection}

Let $Q_0$ be convolution by $q_0$ and let $R f(u)=f(-u)$.
This antipodal unitary is $(-1)^n$ on degree-$n$ harmonics. Define
\[
 s_\sigma(u)=
 \sqrt{\frac{\widehat A(u)}{(Q_0\widehat B_\sigma)(u)}} .
\]
The right multiplier in V24's factorization is $s_+$ in the aligned
case and $u\mapsto s_-(-u)$ in the opposed case. Therefore
\[
 \widehat C_+=M_{s_+}Q_0M_{s_+},\qquad
 \widehat C_-R=M_{s_-}Q_0R M_{s_-}.
 \tag{V25.8}\label{r25:twisted}
\]

\begin{proposition}[Complete finite signed spectrum]
\label{r25:signed-spectrum}
Let $X_\sigma$ have columns $\sqrt{\lambda_i}\,s_\sigma\phi_i$
in a real Haar-orthonormal basis of the complete degree-52 harmonic band,
and put $G_\sigma=X_\sigma^*X_\sigma$.
Let $D_{ii}=(-1)^{n_i}$. The nonzero eigenvalues of $\widehat C_+$
are those of $G_+$. The nonzero eigenvalues of the self-adjoint
$\widehat C_-R$ are those of $G_-D$. They are real, and their absolute
values are the singular values of $\widehat C_-$.
For each sign exactly one singular value is the Perron value one.
\end{proposition}
\begin{proof}
All multipliers and weights are strictly positive, so the independent
harmonic features give $X_\sigma$ full column rank.
Equation \eqref{r25:twisted} is respectively $X_+X_+^*$ and $X_-DX_-^*$.
Polar decomposition gives the finite self-adjoint matrices $G_+$
and $G_-^{1/2}DG_-^{1/2}$. The latter is similar to $G_-D$.
Right multiplication by $R$ preserves singular values.

The marginal square-root vectors give singular value one. For the
opposed law the right marginal is the antipodal reflection of the left,
so the left square-root vector is a $+1$ eigenvector of $\widehat C_-R$.
Both kernels in \eqref{r25:twisted} are strictly positive pointwise;
the second uses $q_0(-uv^{-1})>0$. Equality in conditional
Cauchy--Schwarz forces any unit singular vectors to be proportional to
the marginal square-root vectors, since the joint density is everywhere
positive. The unit singular value is consequently simple.
\end{proof}

An unvalidated double-precision eigenspectrum of $G_-D$ is not a
certificate. Although similar to a self-adjoint matrix, it need not be
well-conditioned in this basis. A diagnostic ordinary-precision
calculation returned magnitudes far above one. That failed check is
rejected, not interpreted as a physical effect.

\subsection{Radial enclosures and all 106 finite blocks}
\label{r25:computation}

Write $a(t)=\sum_{n=0}^{160}a_n\chi_n(t)$ with exact rational
coefficients and set
\[
 d(t)=\sum_{n=0}^{52}\lambda_n a_n\chi_n(t),\qquad
 w_\sigma(t)=\frac{a(t)}{d(\sigma t)}.
 \tag{V25.9}\label{r25:weight}
\]
The denominator is strictly positive on $[-1,1]$ because it is the
convolution of two pointwise positive functions. Its shifted-power
coefficients need not all be positive; no such assumption is used.
The common $e^{-30}$ cancels, so $w_\sigma$ is exactly the left
own-marginal Gram weight, not the magnetic weight $e^{30t}/Z$.

For each sign compute the 105 real moments
\[
 M_{\sigma,j}=\frac2\pi\int_0^\pi
   \sin^2\theta\,\cos^j\theta\,
   \frac{a(\cos\theta)}{d(\sigma\cos\theta)}\,d\theta,
 \qquad 0\leq j\leq104 .
 \tag{V25.10}\label{r25:moments}
\]
The substitution $t=\cos\theta$ identifies these with the required
weighted Haar moments. This trigonometric integrand is meromorphic,
with no square-root branch cut. Complex-rectangle evaluation is accepted
only when interval evaluation excludes zero from the denominator;
otherwise it returns a nonfinite enclosure and must be subdivided
or rejected.

For $\ell=0,\ldots,52$ and $0\leq i,j\leq52-\ell$ form
\[
 (G_{\sigma,\ell})_{ij}
   =\sqrt{\frac{\lambda_{\ell+i}\lambda_{\ell+j}}
                    {h_{\ell i}h_{\ell j}}}
          \sum_k c_{\ell ij,k}M_{\sigma,k},
 \qquad
 \sum_k c_{\ell ij,k}t^k
     =(1-t^2)^\ell C_i^{(\ell+1)}(t)C_j^{(\ell+1)}(t).
 \tag{V25.11}\label{r25:matrix}
\]
Here $h_{\ell i}$ is V24's archived Haar norm.
All $c_{\ell ij,k}$ are integers, and the maximum required degree is 104.
Each block has size $53-\ell$ and multiplicity $2\ell+1$.
Use $G_{+,\ell}$ for the aligned spectrum and $G_{-,\ell}D_\ell$
for the opposed signed spectrum, with $(D_\ell)_{jj}=(-1)^{\ell+j}$.

\begin{theorem}[Computer-assisted own-marginal spectral certificate]
\label{r25:certificate}
For the exact native laws \eqref{r25:joint}, \eqref{r25:main} holds.
The largest non-Perron singular cluster lies in the $\ell=1$ sector
for both signs, with angular multiplicity three.
\end{theorem}
\begin{proof}
The checker \path{scripts/verify_ym_restart_v25_axial_spectra.py}
starts with V24's exact rational coefficients and performs no coefficient
fitting. It uses python-flint 0.9.0, FLINT 3.6.0, at 640-bit precision.
Validated complex integration is requested with relative and absolute
tolerance $10^{-140}$. The actual returned balls, not the requested
tolerance, must pass finiteness and radius checks. Each of the 210
moment enclosures has radius less than $10^{-120}$.

Integer Gegenbauer recurrence and binomial expansion construct all
matrix entries by ball arithmetic. The certified eigensolver is called
with multiplicities allowed, never with its uncertified approximation
algorithm. Every block must return its full number of finite eigenvalue
enclosures, with imaginary part containing zero and real radius less
than $10^{-30}$. Reality follows independently from
Proposition~\ref{r25:signed-spectrum}, not from observing a small
imaginary part. There are 106 blocks and 2862 eigenvalues before angular
multiplicities; each sign has full rank 51039.

The finite-operator bounds checked on these enclosures are as follows.
Only the $\ell=0$ row excludes the exact unit mode. The first row
locates the leading values; interval comparisons reserve the full error
in \eqref{r25:operator-error}.
\begin{center}
\begin{tabular}{@{}lcc@{}}
\toprule
Finite spectral quantity & Aligned & Opposed\\
\midrule
Leading non-Perron magnitude, approximately
 & $0.341013700112$ & $0.875444289409$\\
$\ell=0$, excluding the unit mode & $<0.125$ & $<0.142$\\
$\ell=1$, excluding its leading eigenvalue & $<0.044$ & $<0.127$\\
Every eigenvalue in $\ell=2,\ldots,52$ & $<0.115$ & $<0.678$\\
\bottomrule
\end{tabular}
\end{center}
All comparisons are ball-enclosure inequalities against exact decimal
rationals. Sorting midpoints only nominates candidates; enclosure
tests certify the stated separations.

Singular-value perturbation by \eqref{r25:operator-error} proves
\eqref{r25:main} and controls the infinite-dimensional complement of
the finite-rank approximation. Simultaneous transverse rotations
commute with the exact and approximating operators. The strict
separation from the other angular sectors places the exact leading
non-Perron cluster in $\ell=1$ as well. The strict gap to the next
eigenvalue within that radial block and angular irreducibility give
multiplicity three.

Independent algebraic checks include 180 exact Haar-orthogonality
identities by monomial integration and the trace of each spectrum
against its input matrix. The addition theorem gives full-band checks:
\[
 \begin{split}
 \sum_{\ell=0}^{52}(2\ell+1)\Tr G_{\sigma,\ell}
     &=q_0(I)M_{\sigma,0},\\
 \sum_{\ell=0}^{52}(2\ell+1)\Tr(G_{\sigma,\ell}D_{\sigma,\ell})
     &=q_0(\sigma I)M_{\sigma,0},
 \end{split}
 \tag{V25.12}\label{r25:trace}
\]
where $D_{+,\ell}=I$, $D_{-,\ell}=D_\ell$.
The checker verifies both identities with enclosures, including the
strongly cancelling opposed trace. These supplementary checks do not
replace validated integration and spectral enclosure.
\end{proof}

Primary numerical-method sources checked here are the
\href{https://python-flint.readthedocs.io/en/latest/acb.html}
{python-flint integration contract},
\href{https://python-flint.readthedocs.io/en/latest/acb_mat.html}
{certified eigenvalue documentation}, and
\href{https://flintlib.org/doc/arb_mat.html}{FLINT's ball-matrix contract}.
The analytic-domain flag and the distinction between certified and
\texttt{approx} methods are retained. This targeted method review is
not the scheduled broad V30 sweep. No Yang--Mills theorem is imported
from these numerical libraries.

\subsection{Robustness under all retained-link perturbations}
\label{r25:open}

Fix either compatible reference exterior $\eta^\sigma$ from V24.
Let $E$ be the retained links occurring in the 29 plaquettes incident
to the six-link bucket. There are 22 plaquettes with one free link
and seven with two free links. Hence
\[
 |E|\leq22\cdot3+7\cdot2=80 .
 \tag{V25.13}\label{r25:incidence}
\]
This counts retained occurrences, not necessarily distinct links.
Enumeration at spatial sizes 4, 5, 6 and temporal period 8 gives
$|E|=78,80,80$, respectively, and 80 occurrences in each case.

Use the unit-quaternion chord metric and put
\[
 \mathcal U^\sigma_\varepsilon
  =\{\eta:|\eta_e-\eta^\sigma_e|<\varepsilon\ \text{for all }e\in E\}.
 \tag{V25.14}\label{r25:neighbourhood}
\]
For arbitrary $\eta$, let $J_\eta(u,v)$ be the true six-link bucket
with its four corner links integrated out, including all 29 original
incident plaquettes. There is no polynomial approximation in the
following comparison.

\begin{proposition}[An open boundary family]
\label{r25:robustness}
At $\gamma=10$, for every $\eta\in\mathcal U^\sigma_\varepsilon$,
\[
 \|C_{J_\eta}-C_{J_{\eta^\sigma}}\|_{\rm op}
       \leq e^{1600\varepsilon}-1 .
 \tag{V25.15}\label{r25:robust}
\]
In particular, at $\varepsilon=10^{-5}$,
\[
 \rho(J_\eta)<0.36\quad(\sigma=1),\qquad
 \rho(J_\eta)<0.90\quad(\sigma=-1).
 \tag{V25.16}\label{r25:open-bound}
\]
No equality constraint fixes the perturbed corner fields or internal transports.
\end{proposition}
\begin{proof}
Inversion and multiplication by unit quaternions preserve chord distance.
Telescope each oriented plaquette product with free links held fixed.
Its scalar part changes by at most the sum of the retained-factor
chord changes. The 80 retained occurrences therefore give
\[
 |\log W_\eta(U)-\log W_{\eta^\sigma}(U)|
     \leq80\gamma\varepsilon=800\varepsilon
\]
for every six-link assignment $U$. Integrate the four corners and put
$h=800\varepsilon$ to obtain
$e^{-h}J_{\eta^\sigma}\leq J_\eta\leq e^hJ_{\eta^\sigma}$.
Each marginal obeys the same comparison. The own-marginal kernel ratio
lies between $e^{-2h}$ and $e^{2h}$, whence
\[
 |C_{J_\eta}-C_{J_{\eta^\sigma}}|
      \leq(e^{2h}-1)C_{J_{\eta^\sigma}}.
\]
The positive kernel on the right has norm one. Apply the domination
to absolute values of test functions to prove \eqref{r25:robust}.
Singular-value perturbation bounds the change in maximal correlation.
At $\varepsilon=10^{-5}$, $e^x-1\leq x/(1-x)$ with $x=0.016$
gives a perturbation at most $2/123$.
Combining with \eqref{r25:main} proves \eqref{r25:open-bound};
these last inequalities are checked rationally.
\end{proof}

\subsection{Positive Wilson probability and its quantitative limitation}
\label{r25:probability}

\begin{proposition}[A finite-energy bound at the stated coupling]
\label{r25:positive-mass}
For the ordinary four-dimensional $\SU(2)$ Wilson Gibbs measure at
$\gamma=10$ on the admissible finite tori with the block geometry above,
\[
 \Pr\{\eta\in\mathcal U^\sigma_{10^{-5}}\}>10^{-6080},
 \qquad \sigma\in\{1,-1\}.
 \tag{V25.17}\label{r25:prob}
\]
The bound also holds conditional on all links outside $E$, uniformly
in their values, and after any further marginalization of those links.
\end{proposition}
\begin{proof}
The normalized Haar mass of an $S^3$ chord ball is independent of its
centre. For $0<\varepsilon\leq1$,
\[
 H_\varepsilon=\frac2\pi
   \int_{1-\varepsilon^2/2}^1\sqrt{1-t^2}\,dt
 \geq\frac{2}{3\pi\sqrt2}\varepsilon^3
 >\frac{\varepsilon^3}{10}.
 \tag{V25.18}\label{r25:cap}
\]
Use $1+t\geq1$ for the first inequality; $\pi<4$ and $\sqrt2<3/2$
suffice for the final one.

Each link belongs to six plaquettes, so at most $6|E|$ plaquettes
enter the density of the $E$ links given their complement.
Each normalized Wilson factor $e^{-10(1-t)}$ lies in $[e^{-20},1]$.
The conditional density relative to product Haar is at least
$e^{-120|E|}$: the numerator has this lower bound and its partition
function is at most one. Thus, using product Haar, not Gibbs independence,
\[
 \Pr(\mathcal U^\sigma_\varepsilon\mid E^c)
 \geq e^{-120|E|}H_\varepsilon^{|E|}
 >e^{-9600}10^{-1280}>10^{-6080}
\]
when $\varepsilon=10^{-5}$. The last comparison follows from $e^2<10$.
The estimates hold for every outside-link assignment, so integrating
them proves the remaining assertions.
\end{proof}

This conservative bound removes the limitation of an exact boundary
with unassigned probability. It does not show the neighbourhoods are
frequent enough for a useful multiscale contraction and is not uniform
as $\gamma\to\infty$ with the same constants. Different buckets share
retained links: no product estimate or independence of their events
is asserted. Most importantly, dependence between two middle links
within a conditioned bucket is not the response of its outgoing
message to all retained links. Neither estimate can silently replace
the actual four-endpoint sewing affinity.

\subsection{Verification and the upstream task now exposed}
\label{r25:next}

The numerical theorem is computer-assisted. The checker verifies the
relative approximation budget, every harmonic block and
multiplicity, unit-mode exclusion, trace identities, independent Haar
norm fixtures, and native incidence counts. It rejects nonfinite balls
or failed inequalities; ordinary numerical stabilization is not used
as proof. Its generated certificate and all build products remain
outside \path{latex_notes}.

Removing this append and restoring V24's title recovers its entire
574039-byte source with SHA--256
\path{50A451F89B1544879F1B20197116B804C6F5DEDB1A73A66173D4F814531129E8}.
Only this immediate active predecessor is replaced; the other
46 manuscript sources remain unchanged.

The next useful calculation is upstream of multiplying local numbers:
derive an outgoing-message response estimate for this same normalized
six-link integration under retained-link variations, or control a much
larger exterior set with an adequate bad-set budget. It must survive
shared retained variables and connect explicitly to the endpoint-affinity
ledger. The present calculation supplies a certified local spectral
primitive and a narrow robustness region, not that missing connection.

\paragraph{Status.}
The full axial own-marginal spectra are enclosed at the two native
exteriors. Strict local contraction persists on explicitly
positive-probability open neighbourhoods. General-exterior response,
four-endpoint affinity, cofinal volume and cutoff control, physical-time
transport, the interacting four-dimensional continuum construction,
and its positive Hamiltonian mass gap remain open.

% END CONSOLIDATED V25 APPEND

% BEGIN CONSOLIDATED V26 APPEND -- 2026-09-19
\clearpage
\section{V26: Native source response and a retained-metric obstruction}
\label{r26:context}

V25 paid the full conditional spectra at two native exteriors, but
explicitly did not identify them with outgoing-message response
constants. This append evaluates that missing distinction on the actual
six-link bucket. A perturbation of its three right private retained
links has Fisher-information transmission ratio approximately
$0.1162753$ in the aligned case and $0.74109005$ in the opposed case.
Both are smaller than one. Nevertheless, with unit speed in the
retained links' product round metric, the opposed perturbation has
outgoing Fisher response greater than 34 and Hellinger speed greater
than 2.9156. The integrated bucket also contributes a negative-curvature
direction to its retained effective action.

These are not merely equality-constrained boundary tests. On V25's
entire opposed chord neighbourhood of radius $10^{-5}$, a specified
unit retained-link direction still has Hellinger speed greater than
$1.5$, although the middle-pair maximal correlation is below $0.90$.
A specified unit direction has bucket effective-action curvature below
$-66$ there. This rejects a unit-Lipschitz or nonnegative-bucket-curvature
shortcut for the fixed native metric. It is not a counterexample to the
Yang--Mills mass gap: the total retained action has additional terms,
and a different, quantitatively matched source metric may be useful.

\subsection{Scope and already established machinery}
\label{r26:scope}

V5, V6, V11 and V12 already give normalized score, conditional-response
and score-Gram identities. The f12 V54 outer-source estimate already
differentiates both a conditional numerator and its normalizer.
The f16.2 V62 joint-ladder analysis gives general retained-curvature
bounds and a nonnegative coherent Hessian when all field slots align.
We reuse this calculus. We do not claim a new general information
contraction theorem or another proof of the fully coherent Hessian.
The zero corner fields in the present bucket are realized by cancelling
individual staples, and the opposed middle fields are not the fully
coherent point of that archive.

For primary-source context, Polyanskiy--Wu,
\href{https://arxiv.org/html/1508.06025}
{\emph{Strong data-processing inequalities for channels and Bayesian
networks}}, Section 2.1, equation (9), identifies fixed-law squared
maximal correlation with the chi-squared contraction coefficient.
The fixed-input qualifications in that section matter here.
We use only the elementary conditional-expectation proof below,
not a tensorization or a channel-composition theorem for the Wilson
graph. This is a targeted review; the next scheduled companion
review and broad literature sweep remain V30.

We keep $\gamma=10$ and V24--V25's two exact middle laws
\[
 \mu_\sigma(du\,dv)
  =Z_\sigma^{-1}e^{30u_0+30\sigma v_0}p(uv^{-1})\,du\,dv,
 \qquad \sigma\in\{1,-1\}.
 \tag{V26.1}\label{r26:law}
\]
Set
\[
 M_\sigma=\E u_0,\quad
 v_\sigma=\E u_1^2=\E v_1^2,\quad c_\sigma=\E u_1v_1,\quad
 m_\sigma(u)=\E[v_\perp\mid u],\quad
 \tau_\sigma=\E[m_{\sigma,1}(u)^2].
 \tag{V26.2}\label{r26:data}
\]
Here $u_\perp,v_\perp$ have three transverse components.
Their means, and those of $m_\sigma$, are zero by transverse rotations.
V24's exact Ward identity is
\[
 M_\sigma=30(v_\sigma+\sigma c_\sigma).
 \tag{V26.3}\label{r26:ward}
\]

\subsection{The actual six retained source links}
\label{r26:geometry}

At each middle link there are three private plaquettes, each containing
a single parallel retained link and two retained connector links.
V24 verified that these six parallel retained links are distinct and
do not occur in a corner private plaquette. Within the bucket each
occurs in exactly one incident plaquette. The verification here
re-enumerates this fact at spatial sizes 4, 5, 6 and temporal period 8.

With the other retained links fixed at the reference values, use
geodesic coordinates for their dual staple quaternions:
\[
 q_{Lj}(t)=e^{t x_j},\qquad
 q_{Rj}(t)=\sigma e^{t y_j},\qquad
 x_j,y_j\in\R^3,\quad j=1,2,3.
 \tag{V26.4}\label{r26:slots}
\]
Multiplication and inversion by fixed unit quaternions are isometries,
so these are genuine unit-scale coordinates of the corresponding
retained links, not independently invented field parameters.
Their product-metric speed squared is
$\sum_j(|x_j|^2+|y_j|^2)$.
Write $X=\sum_jx_j$ and $Y=\sum_jy_j$.

The first derivative of the log \emph{unnormalized} six-link density
at the reference point is
\[
 S=\gamma(X\cdot u_\perp+\sigma Y\cdot v_\perp).
 \tag{V26.5}\label{r26:score}
\]
Its mean vanishes there. The general normalized score is always
$S-\E S$. The second derivative of the log unnormalized density is
\[
 R=-\gamma\left(
     u_0\sum_j|x_j|^2+\sigma v_0\sum_j|y_j|^2\right).
 \tag{V26.6}\label{r26:contact}
\]
It cannot be discarded even when a zero-sum choice makes $S=0$.

There are two distinct outputs. First, integrate the four corners and
$v$ but retain $u$, obtaining a normalized left-middle record density
$p_L(u;\eta)$. Second, integrate all six free links, obtaining the
unnormalized outgoing bucket message $\mathcal H(\eta)$.
We compute the response of the first and curvature of the second.
Neither is the whole renormalized retained law or its four-endpoint
affinity.

\subsection{Fisher transmission versus the native parameter metric}
\label{r26:fisher}

For any bounded score $S(v)$ changing only the right weight, direct
differentiation, with both normalizers retained, gives
\[
 \partial_t\log p_L(u)
      =\E[S\mid u]-\E S .
 \tag{V26.7}\label{r26:conditional-score}
\]
Thus
\[
 I_{\rm out}=\Var(\E[S\mid u])
       \leq \rho_\sigma^2\Var(S)=\rho_\sigma^2 I_{\rm in}.
 \tag{V26.8}\label{r26:information}
\]
This follows from the definition of maximal correlation as the
centered conditional-expectation norm. It compares two statistical
norms under the same joint law. It does not bound $I_{\rm in}$ by the
squared speed of the retained parameter.

For the right-link variation in \eqref{r26:slots}, with $X=0$,
rotational symmetry gives the exact quadratic forms
\[
 I_{\rm in}(y)=\gamma^2v_\sigma|Y|^2,\qquad
 I_{\rm out}(y)=\gamma^2\tau_\sigma|Y|^2.
 \tag{V26.9}\label{r26:physical-Fisher}
\]
They each have rank three in the nine right-link tangent coordinates.
On the coherent subspace $y_1=y_2=y_3$ their nonzero eigenvalues in the
product round metric are $3\gamma^2v_\sigma$ and
$3\gamma^2\tau_\sigma$; on the six-dimensional zero-sum subspace both
vanish. The exact transmission factor for these source directions is
$\eta_\sigma=\tau_\sigma/v_\sigma$, not necessarily the full
$\rho_\sigma^2$.

Conditional Cauchy--Schwarz and maximal correlation give independent
checks:
\[
 \frac{c_\sigma^2}{v_\sigma}\leq\tau_\sigma
             \leq v_\sigma\rho_\sigma^2 .
 \tag{V26.10}\label{r26:sandwich}
\]
Indeed $\E[u_1m_{\sigma,1}(u)]=c_\sigma$ and $\E u_1^2=v_\sigma$.

If both triples move, the left record also has a direct source term:
\[
 I_L(x,y)=\gamma^2\bigl(
   v_\sigma|X|^2+2\sigma c_\sigma X\cdot Y+\tau_\sigma|Y|^2\bigr).
 \tag{V26.11}\label{r26:both-Fisher}
\]
This positive semidefinite form follows from the square of
$\gamma(X\cdot u_\perp+\sigma Y\cdot m_\sigma(u))$.
The right-only information inequality must not be applied to delete
the direct term.

\begin{theorem}[Certified native source response]
\label{r26:response-values}
For the exact laws \eqref{r26:law}, the following strict intervals hold.
The last two rows use a unit coherent right-link tangent and
$\gamma=10$.
\begin{center}
\begin{tabular}{@{}lcc@{}}
\toprule
Quantity & Aligned & Opposed\\
\midrule
$\tau_\sigma$
 & $(0.002783240,0.002783242)$ & $(0.113345333,0.113345336)$\\
$\eta_\sigma=\tau_\sigma/v_\sigma$
 & $(0.1162752,0.1162754)$ & $(0.7410900,0.7410901)$\\
$I_{\rm out}=300\tau_\sigma$
 & $(0.834972,0.834973)$ & $(34.003600,34.003601)$\\
Hellinger speed
 & $(0.456884,0.456885)$ & $(2.915630,2.915631)$\\
\bottomrule
\end{tabular}
\end{center}
We use Hellinger distance $H(p,q)=\|\sqrt p-\sqrt q\|_{L^2}$,
without an additional factor $1/\sqrt2$.
\end{theorem}
\begin{proof}
Here is the complete finite evaluation, including the nonlinear
conditional normalizer. Use V25's positive kernel $q_0$, its rational
multipliers $\lambda_n$, and
$a(t)=p_{160}(30(1+t))=\sum a_n\chi_n(t)$.
Let $d(t)=\sum_{n=0}^{52}\lambda_n a_n\chi_n(t)$ and set
\[
 h_k=\frac{a_k}{k+1}-\frac{a_{k+2}}{k+3},\qquad
 N_\sigma(t)=\sum_{k=0}^{51}\lambda_{k+1}\sigma^k h_k C_k^{(2)}(t),
 \qquad
 z_\sigma=\sum_{n=0}^{52}\lambda_n a_n^2\sigma^n .
 \tag{V26.12}\label{r26:polynomials}
\]
The identity
$U_n=(C_n^{(2)}-C_{n-2}^{(2)})/(n+1)$ expands
$v_1a(\sigma v_0)$ in degree-$(k+1)$ transverse harmonics.
Convolution by $q_0$ multiplies that degree by $\lambda_{k+1}$.
For the approximating joint law this proves
\[
 \widehat m_{\sigma,1}(u)
    =u_1\frac{N_\sigma(u_0)}{d(\sigma u_0)},\qquad
 \widehat\tau_\sigma
   =\frac{2}{3\pi z_\sigma}\int_0^\pi
       \frac{\sin^4\theta\,a(\cos\theta)
                   N_\sigma(\cos\theta)^2}
            {d(\sigma\cos\theta)}\,d\theta .
 \tag{V26.13}\label{r26:integral}
\]
Angular averaging gives $\E_{\omega\in S^2}\omega_1^2=1/3$.
Both denominators are positive; in particular, an alternating
expression for $z_-$ is evaluated exactly as a rational number.

V25 gives one-sided joint retention $L>1-10^{-10}$.
After conditioning on $u$ the same retention bound holds.
Put $\delta=1-L$. For $|v_1|\leq1$, the conditional means differ
by at most $2\delta$. Their squared means differ by at most
$4\delta$, and the $u$ marginals have total variation at most
$\delta$. Since a squared conditional mean lies in $[0,1]$,
\[
 |\tau_\sigma-\widehat\tau_\sigma|
          \leq5\delta<5\cdot10^{-10}.
 \tag{V26.14}\label{r26:tau-error}
\]
Both laws have zero transverse means. This comparison approximates a
fixed bounded regression functional. It does \emph{not} differentiate
a polynomial filter and pretend its derivative is the exponential score.

The checker \path{scripts/verify_ym_restart_v26_native_response.py},
run with \texttt{--certify}, evaluates \eqref{r26:integral} by validated
complex integration at 480-bit precision, using relative and absolute
tolerance $10^{-80}$. It requires each actual returned real radius
to be below $10^{-60}$. The callback is meromorphic, with division
returning a nonfinite ball if the denominator cannot exclude zero.
The numerator is built from exact rational polynomials. Applying
\eqref{r26:tau-error}, V24's much narrower moment enclosures for
$v_\sigma,c_\sigma,M_\sigma$, and interval arithmetic proves the table.
The two bounds in \eqref{r26:sandwich} are also checked.

An independent check removes the conditional denominator:
monomial Haar integration of
$(1-t^2)a(t)N_\sigma(t)/(3z_\sigma)$ agrees exactly with
\[
 \frac1{z_\sigma}\sum_{k=0}^{51}
    \lambda_{k+1}\sigma^k h_k^2\frac{(k+1)(k+3)}{12}.
 \tag{V26.15}\label{r26:cross-check}
\]
This is the transverse mixed moment of the polynomial law.
Its proximity to V24's true mixed moment is checked separately.

Finally a smooth positive density path with score $\dot\ell$ obeys
$\partial_t\sqrt p=\tfrac12\sqrt p\,\dot\ell$, hence
\[
 \lim_{t\to0}\frac{H(p_t,p_0)}{|t|}
        =\frac12\sqrt{I_{\rm out}} .
 \tag{V26.16}\label{r26:H-speed}
\]
Compactness and strictly positive smooth densities justify the $L^2$
differentiation here. This proves the final row.
\end{proof}

The input coherent Fisher eigenvalues are approximately
$7.18099594$ and $45.88322352$. Thus information contracts in both
examples, while the native geometric sensitivity differs substantially.
No rescaling of the physical retained parameters has been silently
included in \eqref{r26:information}.

\subsection{The complete 18-direction source Hessian of the bucket}
\label{r26:Hessian}

Keep all other retained links fixed and integrate all six free links.
The resulting $\mathcal H$ is, up to a constant independent of the
six selected links, the source normalizer of \eqref{r26:law}.
At either reference point $\nabla\log\mathcal H=0$ in these 18
source coordinates. By \eqref{r26:score}--\eqref{r26:contact},
\[
 \begin{split}
 Q_\sigma(x,y)
  :=D^2\log\mathcal H[(x,y),(x,y)]
  ={}&-\gamma M_\sigma\sum_j(|x_j|^2+|y_j|^2)\\
     &+\gamma^2\bigl(v_\sigma|X|^2+v_\sigma|Y|^2
                         +2\sigma c_\sigma X\cdot Y\bigr).
 \end{split}
 \tag{V26.17}\label{r26:Q}
\]
This is the Hessian of $\log\mathcal H$, not of $-\log\mathcal H$.

\begin{proposition}[Exact source-sector splitting]
\label{r26:Hessian-spectrum}
In the native product metric, the spectrum of \eqref{r26:Q} is
\[
 -\gamma M_\sigma\quad(12\text{ times}),\qquad
 0\quad(3\text{ times}),\qquad
 -6\sigma\gamma^2 c_\sigma\quad(3\text{ times}).
 \tag{V26.18}\label{r26:H-spectrum}
\]
At $\gamma=10$, the last eigenvalue lies in
$(-4.897074,-4.897072)$ for $\sigma=1$ and in
$(77.287400,77.287403)$ for $\sigma=-1$.
\end{proposition}
\begin{proof}
Zero-sum variations within either triple give the first eigenvalue.
On coherent triples write $x_j=x/\sqrt3$, $y_j=y/\sqrt3$.
The remaining matrix, for each transverse coordinate, is
\[
 \begin{pmatrix}
  -\gamma M_\sigma+3\gamma^2v_\sigma&3\sigma\gamma^2c_\sigma\\
  3\sigma\gamma^2c_\sigma&-\gamma M_\sigma+3\gamma^2v_\sigma
 \end{pmatrix}.
\]
Since $M_\sigma=3\gamma(v_\sigma+\sigma c_\sigma)$, its common
rotation eigenvalue is zero and its relative eigenvalue is
$-6\sigma\gamma^2c_\sigma$. The finite-dimensional decomposition has
dimensions $12+3+3=18$, so no selected tangent direction is omitted.
The numerical intervals follow directly from V24's moment enclosures.
The checker also tests all six algebraic eigenvectors in exact
rational fixtures for each transverse coordinate.
\end{proof}

In particular, the opposed bucket effective action
$-\log\mathcal H$ has three negative source-sector eigenvalues.
This is not in conflict with f16.2 V62's fully coherent result.
Nor does it establish negative curvature of the full retained action:
other original plaquettes and other generated messages remain present.
The claim concerns exactly the specified bucket contribution.

\subsection{A quantitative open-neighbourhood obstruction}
\label{r26:robust}

Let $\mathcal U^-_\varepsilon$ be V25's opposed retained-link
neighbourhood. Its set $E$ contains at most 80 retained links, with
80 retained plaquette occurrences. For every six-free-link assignment,
V25 proves an unnormalized density ratio between $e^{-h}$ and $e^h$,
where $h=800\varepsilon$. Multiplication by the irrelevant scalar
$e^{-h}$ makes this a one-sided filter with retention $e^{-2h}$.
Thus the full free-link laws, every relevant marginal, and their
conditional laws given $u$ have total-variation distance at most
\[
 \delta=1-e^{-1600\varepsilon}\leq1600\varepsilon .
 \tag{V26.19}\label{r26:TV}
\]

We need a response-functional estimate, since maximal-correlation
stability alone does not control the moving source.
For a law $\nu$ on the free links define
$F_\nu(S)=\Var_\nu(\E_\nu[S\mid u])$.

\begin{lemma}[Stability of the centered regression response]
\label{r26:F-stability}
Suppose the full, $u$-marginal, and $u$-conditional laws $\mu,\nu$
obey the preceding TV bound. If $|S|\leq K$ and $\E_\mu S=0$, then
\[
 |F_\nu(S)-F_\mu(S)|\leq K^2(5\delta+4\delta^2).
 \tag{V26.20}\label{r26:F-stable}
\]
If also $|S'|\leq K$ and $\|S'-S\|_\infty\leq d$, then
\[
 |F_\nu(S')-F_\nu(S)|\leq2Kd .
 \tag{V26.21}\label{r26:moving-score}
\]
\end{lemma}
\begin{proof}
Conditional means change by at most $2K\delta$, so their squares
change by at most $4K^2\delta$. Integrating the old squared conditional
mean under the new marginal costs at most $K^2\delta$.
The new mean has magnitude at most $2K\delta$, whereas the old mean
is zero. Subtracting their squared means adds at most $4K^2\delta^2$.
This proves the first bound.

The map $S\mapsto\E_\nu[S\mid u]-\E_\nu S$ is an orthogonal projection
followed by centering. Its $L^2$ norm on $S'-S$ is at most
$\sqrt{\Var_\nu(S'-S)}\leq d$. Each individual response norm is
at most $K$. Factor the difference of the squared norms to obtain
the second inequality.
\end{proof}

Choose at the reference point a unit coherent right-link direction
$y_j=e_1/\sqrt3$. At nearby exteriors use the same translated
Lie-algebra generators on these actual retained links. They remain
unit-speed geodesic directions. Each selected link occurs in one
private bucket plaquette with three retained factors. Telescoping
that differentiated plaquette product gives, for the unnormalized
score $S_\eta$,
\[
 |S_\eta|\leq K=10\sqrt3,\qquad
 |S_\eta-S_{\eta^-}|\leq3K\varepsilon .
 \tag{V26.22}\label{r26:score-variation}
\]
The inserted unit Lie-algebra generator has norm one; inversion and
the three retained unit factors preserve the same telescoping bound.
The free variables are held fixed in this comparison.

\begin{theorem}[Local dependence contraction can coexist with native expansion]
\label{r26:open-obstruction}
On the entire neighbourhood $\mathcal U^-_{10^{-5}}$:
\begin{enumerate}
\item The middle-pair maximal correlation is below $0.90$.
\item The selected unit right-link direction has outgoing Fisher
response greater than 9 and Hellinger speed greater than $1.5$.
\item A selected unit relative direction of the six source links
has $D^2\log\mathcal H>66$, hence bucket effective-action curvature
$D^2(-\log\mathcal H)<-66$.
\end{enumerate}
The neighbourhood has Wilson probability greater than $10^{-6080}$
at the fixed coupling, including conditionally on all links outside $E$.
\end{theorem}
\begin{proof}
The first and last assertions are V25's established bounds for this
same event; we do not multiply probabilities for overlapping events.
For the second, apply Lemma~\ref{r26:F-stability},
\eqref{r26:score-variation}, and Theorem~\ref{r26:response-values}.
With $\delta\leq0.016$, $\varepsilon=10^{-5}$ and $K^2=300$,
\[
 F_{\mu_\eta}(S_\eta)
 >34.003600-300(5\delta+4\delta^2+6\varepsilon)
 \geq34.003600-24.3252>9.
 \tag{V26.23}\label{r26:open-Fisher}
\]
Equation \eqref{r26:H-speed} gives Hellinger speed greater than $1.5$.

For the third choose $x_j=e_1/\sqrt6$, $y_j=-e_1/\sqrt6$ at the
reference point and transport the generators as above. Its total
product-metric speed is one. Now $|S_\eta|\leq K=10\sqrt6$ and
$|S_\eta-S_{\eta^-}|\leq3K\varepsilon$.
The second log-density derivative obeys
$|R_\eta|\leq10$ and
$|R_\eta-R_{\eta^-}|\leq30\varepsilon$.
Each selected plaquette is linear in its selected unit link, so its
second geodesic derivative is minus the squared speed times its trace.

At the reference point $\E S_{\eta^-}=0$. A squared bounded variable
takes values in $[0,K^2]$, so
\[
 |\Var_{\mu_\eta}S_{\eta^-}
             -\Var_{\mu_{\eta^-}}S_{\eta^-}|
       \leq K^2(\delta+4\delta^2).
\]
Moving the score changes its variance by at most $2K(3K\varepsilon)$,
by the same centered $L^2$ argument. Moving the expectation of $R$
costs at most $20\delta+30\varepsilon$.
Since $D^2\log\mathcal H=\E R+\Var S$, the full loss is at most
\[
 10(2\delta+3\varepsilon)
       +600(\delta+4\delta^2+6\varepsilon)=10.5707.
 \tag{V26.24}\label{r26:open-curvature}
\]
The reference lower bound $77.287400$ leaves a value greater than 66.
All decimal inequalities here are exact rational checks.
\end{proof}

This theorem does not refute a Wilson spectral gap or the existence of
a weighted contraction scheme. It refutes the specific inference that
V25's conditional $\rho<1$ makes the normalized record map unit-Lipschitz
in the unweighted retained-link metric, or makes this bucket contribution
convex. The rare-event lower probability is far too small to establish
the importance of this region in the actual cofinal Gibbs law.

\subsection{Verification and the changed direction}
\label{r26:next}

Checks cover source incidence, V25's approximation budget, both regression
variances, 18 polynomial identities, independent mixed moments, 24 Hessian
eigenvector fixtures, and 16 filtered-law response fixtures. Numerical
enclosures are computer-assisted; geometry and stability are proved above.
All generated data and build products stay outside \path{latex_notes}.

Deleting this append and restoring V25's title recovers exactly its
593793-byte source, SHA--256
\path{C2C4BA72D4CF65DFDB1118FF70FAA1D4EEF0EC6FBDC0E2D57EE021E35E5993D1}.
Only the immediate active predecessor is replaced; all 46 other
manuscript sources are unchanged.

Next, test compatibility of adjacent input/output source metrics under
their \emph{actual shared retained law}, or quantify compensation from
the other retained-action terms and the exceptional-set budget.
Local maximal correlation cannot replace the unweighted derivative
constant. Metric compatibility and four-endpoint affinity remain unpaid.

\paragraph{Status.}
A genuine native source-response coefficient is now evaluated, and a
retained-metric expansion and bucket-curvature obstruction are proved
on an open positive-probability exterior set. The full global response
estimate, cofinal physical-time and cutoff control, interacting
four-dimensional continuum construction, and positive physical
Hamiltonian mass gap remain unproved.

% END CONSOLIDATED V26 APPEND

% BEGIN CONSOLIDATED V27 APPEND -- 2026-09-19
\clearpage
\section{V27: Full retained curvature after exact packed decimation}
\label{r27:context}

V26 deliberately stopped at the curvature of one generated bucket
message. Its negative sign alone did not decide the curvature of the
retained Wilson action: original plaquettes and other generated
messages could have compensated it. This append resolves that
specific alternative for V23's actual simultaneous six-link packing.
We prove a three-dimensional negative-curvature bound for the
\emph{full exact retained marginal}, uniformly in every other retained
link, on the same opposed neighbourhood used in V25--V26.

This is a finite-lattice, fixed-coupling result for $\SU(2)$ in the
product round metric. It does not establish, or disprove, the
Yang--Mills mass gap. It rules out a particular pointwise positive
Bakry--Emery curvature route after this decimation. Neither a
Poincare inequality nor a multiscale method requires that route.

\subsection{Inherited inputs and the exact retained action}
\label{r27:retained}

We use the periodic lattice
\[
 \Lambda=(\mathbb Z/L\mathbb Z)^3\times
                     (\mathbb Z/(2q)\mathbb Z),\qquad
 L\in4\mathbb N,\quad q\geq3,\quad \gamma=10 .
 \tag{V27.1}\label{r27:lattice}
\]
Coordinates are $(x_\mu,x_\nu,x_\rho,t)$, and link variables have the
round unit-$S^3$ metric. The Wilson plaquette action is
$V_p=\gamma(1-t_p)$, where
$t_p=\tfrac12\operatorname{Re}\Tr U_p\in[-1,1]$.
Use exactly the free blocks of Proposition~\ref{r23:incidence}:
all six $\mu$-links in two $\nu$-columns and three $\rho$-rows,
anchored at $\nu,\rho$ multiples of four and at even
$t\notin\{0,q\}$, for every $x_\mu$.
Write $\mathcal B$ for this packing, $F_b$ for a free block,
$F=\bigcup_bF_b$, and $R$ for all remaining links.

Let $\mathcal P_b$ be the 29 plaquettes meeting $F_b$, and
$\mathcal P_0$ those meeting no free link. V23 proves that these
sets partition all plaquettes: no plaquette meets two distinct free
blocks. Therefore, for a retained assignment $\eta\in\SU(2)^R$,
\[
 \begin{split}
 V_b(z_b,\eta)&=\sum_{p\in\mathcal P_b}V_p(z_b,\eta),\\
 H_b(\eta)&=\int_{\SU(2)^{F_b}}e^{-V_b(z_b,\eta)}\,dz_b,\\
 V_R(\eta)&=\sum_{p\in\mathcal P_0}V_p(\eta)
                          -\sum_{b\in\mathcal B}\log H_b(\eta),\\
 d\mu_R(\eta)&=Z^{-1}e^{-V_R(\eta)}\,d\eta .
 \end{split}
 \tag{V27.2}\label{r27:exact-action}
\]
Haar measures are normalized. The additive constant in $V_R$ is
irrelevant. This formula uses the original exponential Wilson
weights, not the polynomial approximation to them.

Conditional on $\eta$, the free fibres factorize; the functions
$H_b$ generally share retained arguments. Thus
\eqref{r27:exact-action} does \emph{not} assert independence of the
messages under $\mu_R$. All derivatives below keep those shared
arguments and use the actual normalized conditional fibre laws.

The numerical input is the already certified V26 bound
\[
 D^2\log H_{b_0}(\xi,\xi)>66
 \tag{V27.3}\label{r27:inherited}
\]
for its selected unit relative-source direction on the opposed
chord-$10^{-5}$ exterior neighbourhood. We extend it to a
three-dimensional subspace below by inspecting its isotropic proof,
not by assuming operator-norm control of uncomputed directions.
The V23 packing, the V25 event probability, the V26 moment
certificate and the general score-covariance identity are inherited.
The new ingredient is the complete sparse incidence budget which
controls \emph{all other} terms in \eqref{r27:exact-action}.
The coherent-background Hessian in f16.2 V62 is not the opposed
background here and is not being reproved.

\subsection{All exterior links remain retained; the outside incidence budget}
\label{r27:incidence}

Fix $b_0\in\mathcal B$. Write $x$ for its left middle free-link base,
so the right middle base is $x+e_\nu$. The six retained source
$\mu$-links, specified by their base points, are
\[
 \begin{array}{lll}
 s_{L0}=x-e_\nu,&s_{L+}=x+e_t,&s_{L-}=x-e_t,\\
 s_{R0}=x+2e_\nu,&s_{R+}=x+e_\nu+e_t,&
 s_{R-}=x+e_\nu-e_t .
 \end{array}
 \tag{V27.4}\label{r27:sources}
\]
Denote their set by $S$. Let $E_{b_0}$ be all retained links in the
plaquettes $\mathcal P_{b_0}$, including the non-$\mu$ transport
links. It has at most 80 distinct links and 80 plaquette occurrences.

\begin{lemma}[Complete outside count]
\label{r27:count}
Every link in $E_{b_0}$ belongs to $R$ even when \emph{all} blocks
are eliminated. Each link of $S$ belongs to exactly one plaquette
of $\mathcal P_{b_0}$ and five other plaquettes. Among the other
plaquettes meeting $S$, precisely
\[
 m_L=\begin{cases}3,&L=4,\\2,&L\geq8\end{cases}
 \tag{V27.5}\label{r27:m}
\]
contain two selected links, and $30-2m_L$ contain one. None contains
more than two. Every double-selected plaquette belongs to
$\mathcal P_0$.
\end{lemma}
\begin{proof}
Only $\mu$-links are eliminated. Non-$\mu$ links in $E_{b_0}$ are
therefore retained. The private $\mu$-neighbours of a block lie
either in the buffer columns or buffer rows (residues $3$ or $2$
for the former, residue $3$ for the latter), or on an odd time
slice. They are outside every free block. This proves
$E_{b_0}\subset R$, including spatial periodic wraparound.

The six links in \eqref{r27:sources} are distinct. Their only
plaquettes meeting $F_{b_0}$ are the three private plaquettes of
each middle free link. Each selected link has the usual six
plaquette occurrences in four dimensions, leaving five outside
occurrences and hence 30 in total.

A plaquette can contain two selected links only when their
$\mu$-oriented bases differ by one transverse lattice step.
At time $t+1$ the pair is $(s_{L+},s_{R+})$, and at time $t-1$
it is $(s_{L-},s_{R-})$. At time $t$ the two side bases differ
by $3e_\nu$ and are adjacent precisely when $L=4$. No other pair
in \eqref{r27:sources} is transversely adjacent. The two time
offsets cannot become adjacent since $2q\geq6$.
All these plaquettes have their two $\mu$-links retained, and
their other links are not $\mu$-links. Thus they meet no free
block. Subtracting two occurrences per double-selected plaquette
from 30 gives the claimed single count.
\end{proof}

Let $\xi(a)$, $a\in\mathbb R^3$, vary the three left source slots
with velocities $a/\sqrt6$ and the three right source slots with
velocities $-a/\sqrt6$ in V26's reference dual-staple coordinates.
Realize these as velocities of the actual six retained links.
At the reference exterior, the coordinate identifications use
only multiplication by fixed unit quaternions and inversion;
they are isometries. At any nearby exterior translate the same
actual-link Lie-algebra generators, as in V26. This defines
a linear isometric embedding
\[
 a\longmapsto\xi_\eta(a)\in T_\eta\SU(2)^R,\qquad
 |\xi_\eta(a)|^2=|a|^2 .
 \tag{V27.6}\label{r27:subspace}
\]
All other link velocities vanish. For each fixed $a$ the resulting
product curve is geodesic. We write $W_\eta$ for the image of
this three-dimensional embedding.

\begin{lemma}[Plaquette curvature budget]
\label{r27:plaquette-bound}
For any product-geodesic retained direction $\xi$, holding all
free links fixed, set
$b_p(\xi)=\sum_{e\in p\cap R}|\xi_e|$. Then
\[
 |D^2t_p(\xi,\xi)|\leq b_p(\xi)^2,\qquad
 D^2V_p(\xi,\xi)\leq\gamma b_p(\xi)^2 .
 \tag{V27.7}\label{r27:local-bound}
\]
Consequently, along \eqref{r27:subspace},
\[
 \sum_{p\notin\mathcal P_{b_0}}\gamma b_p(\xi_\eta(a))^2
  =\gamma\left(5+\frac{m_L}{3}\right)|a|^2
  =C_L|a|^2,\qquad
 C_L=\begin{cases}60,&L=4,\\170/3,&L\geq8.\end{cases}
 \tag{V27.8}\label{r27:outside-budget}
\]
\end{lemma}
\begin{proof}
Along a round-$S^3$ geodesic of speed $v$, a unit quaternion
factor, or its inverse, has first derivative of norm $v$
and second derivative $-v^2$ times itself. Differentiate the
four-factor plaquette word twice. The diagonal terms have
norms at most $|\xi_e|^2$ and the two ordered mixed terms for
$e\ne f$ have combined norm at most $2|\xi_e||\xi_f|$.
Taking its scalar part cannot increase the norm. Their sum is
$b_p^2$. Multiplying by $-\gamma$ gives the upper bound for $V_p$.

Each selected speed in \eqref{r27:subspace} is $|a|/\sqrt6$.
Lemma~\ref{r27:count} therefore gives
\[
 \frac{\gamma|a|^2}{6}
       \bigl[(30-2m_L)+4m_L\bigr]
       =\gamma(5+m_L/3)|a|^2 .
\]
No sign assumption on the remaining link variables was used.
\end{proof}

\subsection{Integrating the other blocks cannot exceed that budget}
\label{r27:other-messages}

For a fixed retained product-geodesic, smoothness on the compact
free fibre permits differentiating $H_b$ twice. With
$\nu_{b,\eta}(dz_b)=H_b(\eta)^{-1}e^{-V_b(z_b,\eta)}dz_b$,
the standard normalized marginal identity gives
\[
 \begin{split}
 D(-\log H_b)&=\E_{\nu_{b,\eta}}[DV_b],\\
 D^2(-\log H_b)
   &=\E_{\nu_{b,\eta}}[D^2V_b]
                         -\Var_{\nu_{b,\eta}}(DV_b).
 \end{split}
 \tag{V27.9}\label{r27:marginal-Hessian}
\]
For completeness, $H_b'/H_b=-\E V_b'$ and
$H_b''/H_b=\E[(V_b')^2-V_b'']$; substitution into the
second logarithmic derivative proves the identity. The minus
sign on the variance is important: integrating a fibre cannot
produce a larger directional upper bound than the expected
bare directional Hessian.

\begin{proposition}[Full compensation bound]
\label{r27:compensation}
For every retained assignment $\eta$, along the directions
\eqref{r27:subspace},
\[
 \operatorname{Hess}V_R(\xi_\eta(a),\xi_\eta(a))
 \leq C_L|a|^2-
       \operatorname{Hess}\log H_{b_0}(\xi_\eta(a),\xi_\eta(a)).
 \tag{V27.10}\label{r27:compensation-bound}
\]
\end{proposition}
\begin{proof}
In \eqref{r27:exact-action}, retain the target term
$-\log H_{b_0}$ exactly. Bound the original $\mathcal P_0$
terms by \eqref{r27:local-bound}. For every $b\ne b_0$,
drop the nonnegative variance in
\eqref{r27:marginal-Hessian} and bound $D^2V_b$ by the
sum of $\gamma b_p^2$ over $\mathcal P_b$.
These bounds are uniform in the free variables and can be
integrated without change. The plaquette partition allocates
each outside plaquette exactly once, so their sum is
\eqref{r27:outside-budget}. The retained curve is geodesic;
its second derivative is the Riemannian Hessian in that direction.
\end{proof}

This proof includes other buckets sharing selected links. It
does not replace their arguments by frozen reference values,
sum independent-event probabilities, or assume a product retained
law. The same upper bound holds if only $b_0$ is eliminated;
\eqref{r27:compensation-bound} establishes the stronger statement
for the full V23 packing.

\subsection{A three-dimensional obstruction for the full marginal}
\label{r27:full-theorem}

Let $\eta^-|_{E_{b_0}}$ be V25--V26's attainable opposed exterior:
the four corner private fields cancel, the left middle private
staples are $I$, the right ones are $-I$, and the internal
transports are identity. Define
\[
 \mathcal U^-_{b_0}
 =\{\eta:|\eta_e-\eta^-_e|<10^{-5}
                       \text{ for every }e\in E_{b_0}\}.
 \tag{V27.11}\label{r27:event}
\]
No condition is imposed on $R\setminus E_{b_0}$.

\begin{lemma}[Uniformity over the three relative colours]
\label{r27:three-colours}
For every $\eta\in\mathcal U^-_{b_0}$ and $a\ne0$,
\[
 \operatorname{Hess}\log H_{b_0}
                 (\xi_\eta(a),\xi_\eta(a))>66|a|^2 .
 \tag{V27.12}\label{r27:three-bound}
\]
\end{lemma}
\begin{proof}
At $\eta^-$, Proposition~\ref{r26:Hessian-spectrum} gives the
same relative eigenvalue greater than $77.287400$ in all three
transverse colours. Equivalently the restriction to $W_{\eta^-}$
is that scalar times the identity; hence the bound holds for
every unit $a$, not just the coordinate vectors.

For each unit $a$, V26's comparison with the reference fibre
uses $\delta\leq0.016$, the score bounds
$|S|\leq10\sqrt6$, $|S_\eta-S_{\eta^-}|
\leq30\sqrt6\,10^{-5}$, and the contact bounds
$|R|\leq10$, $|R_\eta-R_{\eta^-}|\leq30\,10^{-5}$.
Quaternion norm bounds depend only on $|a|$, not its direction.
Thus the same proof and the same constants in
\eqref{r26:open-curvature} apply simultaneously to every unit
$a$: the Hessian loss is at most $10.5707$. The remaining
lower bound is $66.716700>66$. Homogeneity gives the result.
This argument is not an assertion that the perturbed Hessian
retains rotational symmetry.
\end{proof}

\begin{theorem}[Negative curvature survives the entire exact pass]
\label{r27:negative}
For every $L,q$ in \eqref{r27:lattice}, every target block $b_0$,
and every $\eta\in\mathcal U^-_{b_0}$, the full retained Wilson
action in \eqref{r27:exact-action} satisfies
\[
 \operatorname{Hess}V_R(\xi,\xi)<
 \begin{cases}
   -6|\xi|^2,&L=4,\\
   -\dfrac{28}{3}|\xi|^2,&L\geq8
 \end{cases}
 \quad\text{for all }0\ne\xi\in W_\eta .
 \tag{V27.13}\label{r27:negative-bound}
\]
In particular its full retained Hessian has at least three
eigenvalues below the respective displayed threshold.
The event has $\mu_R$-probability greater than $10^{-6080}$,
also conditionally on arbitrary retained assignments outside
$E_{b_0}$.
\end{theorem}
\begin{proof}
Subtract \eqref{r27:three-bound} from
\eqref{r27:compensation-bound}: $60-66=-6$ and
$170/3-66=-28/3$. The min--max characterization of eigenvalues
applied to the three-dimensional subspace $W_\eta$ proves
the eigenvalue assertion. We do not assume $W_\eta$ is an
invariant subspace of the full Hessian.

The finite-energy proof of Proposition~\ref{r25:positive-mass}
bounds the original Wilson conditional probability of
\eqref{r27:event}, uniformly in every original link outside
$E_{b_0}$. Explicitly, if $m=|E_{b_0}|\leq80$, its density
relative to product Haar is bounded below by $e^{-120m}$,
and each chord cap has Haar mass greater than
$(10^{-5})^3/10$. Thus a common lower bound is
\[
 e^{-9600}10^{-1280}>10^{-6080}.
 \tag{V27.14}\label{r27:probability}
\]
Since $E_{b_0}\subset R$, conditional expectation over the
free links, with all remaining retained links fixed, transfers
this same lower bound to $\mu_R$. This is the tower property,
not an assumption that the event is independent of the other
buckets. Smooth positive finite-volume densities give the
stated conditional version for every fixed exterior.
\end{proof}

\subsection{What this excludes, including constant linkwise reweighting}
\label{r27:BE}

Consider only the auxiliary reversible diffusion on the retained
link manifold with generator
\[
 \mathcal L_R f=\Delta f-\langle\nabla V_R,\nabla f\rangle .
 \tag{V27.15}\label{r27:diffusion}
\]
Its time is not physical Euclidean transfer time. On the unit
round $S^3$, $\operatorname{Ric}=2g$, hence on the product,
$\operatorname{Ric}(\xi,\xi)=2|\xi|^2$.
The weighted Bochner identity is
\[
 \Gamma_2(f)=\|\operatorname{Hess}f\|_{\mathrm{HS}}^2+
   (\operatorname{Ric}+\operatorname{Hess}V_R)
                                  (\nabla f,\nabla f).
 \tag{V27.16}\label{r27:Bochner}
\]
This classical identity and the drift sign are consistent with
Sturm's Proposition 8.4; see
\href{https://arxiv.org/html/1401.0687v3}
{Ricci Tensor for Diffusion Operators}, Sections 5 and 8.1.
No new curvature-dimension theorem is claimed here.

\begin{corollary}[Failure of pointwise positive curvature]
\label{r27:BE-failure}
On the same event, the weighted Ricci form has a three-dimensional
subspace on which
\[
 (\operatorname{Ric}+\operatorname{Hess}V_R)(\xi,\xi)<
 \begin{cases}
  -4|\xi|^2,&L=4,\\
  -\dfrac{22}{3}|\xi|^2,&L\geq8.
 \end{cases}
 \tag{V27.17}\label{r27:weighted-bound}
\]
Consequently the global pointwise inequality
$\Gamma_2(f)\geq\kappa\Gamma(f)$ for all smooth $f$ fails
for every $\kappa\geq-4$ in this metric, in particular for
every nonnegative $\kappa$.
\end{corollary}
\begin{proof}
Add $2|\xi|^2$ to \eqref{r27:negative-bound}. At any one such
point choose a smooth function whose gradient there equals
a chosen nonzero $\xi\in W_\eta$ and whose Hessian there is zero.
A linear function in normal coordinates, multiplied by a
cutoff equal to one near the point, has these prescribed jets.
Equation \eqref{r27:Bochner} then contradicts the proposed
pointwise inequality. The strict violation persists in an
open neighbourhood for this fixed smooth function.
\end{proof}

\begin{corollary}[Constant per-link metric scales do not remove the sign]
\label{r27:constant-metric}
Replace the product metric by
$g_{\mathbf a}=\bigoplus_{e\in R}a_e^2 g_{S^3}$ with arbitrary
positive constants $a_e$, and use the corresponding
$\Delta_{g_{\mathbf a}}-\nabla_{g_{\mathbf a}}V_R\cdot
\nabla_{g_{\mathbf a}}$ diffusion.
It still fails every nonnegative pointwise
Bakry--Emery curvature lower bound.
\end{corollary}
\begin{proof}
Constant rescaling of each product factor leaves its
Levi--Civita connection unchanged. Thus the covariant
$(0,2)$ Hessian of $V_R$ is unchanged. The covariant Ricci
tensor is also unchanged:
$\operatorname{Ric}_{g_{\mathbf a}}
=\bigoplus_e2g_{S^3}$, not $2g_{\mathbf a}$.
The weighted volume differs from product Haar by a constant,
so its potential differs only by an additive constant.
The negative quadratic form in \eqref{r27:weighted-bound}
therefore remains negative on exactly the same vectors.
Normalizing a vector in the new metric divides its value
by a positive number and cannot change the sign.
The preceding prescribed-jet argument applies in the new metric.
\end{proof}

This corollary concerns constant scalar weights on link factors.
It does not address configuration-dependent metrics, metrics
coupling different factors, gauge-quotient geometry, integrated
curvature estimates or a scale-dependent entropy argument.
It also does not identify the auxiliary diffusion gap with
the physical Hamiltonian gap. Negative pointwise curvature
by itself supplies no counterexample to either spectral gap.

\subsection{Literature check, verification and next obligation}
\label{r27:next}

The relevant literature alternative was already listed in f16.2
V12: Bauerschmidt--Bodineau's
\href{https://arxiv.org/html/1907.12308}
{Log-Sobolev inequality for the continuum sine-Gordon model}.
We rechecked the actual hypotheses of Theorem 1.2 and
equations (1.1)--(1.9). It allows negative scale-dependent
Hessian lower bounds, provided the resulting integrated
bound is finite. Its stated construction uses a positive
Gaussian precision matrix on a linear field space and
Gaussian convolutions of the potential. Those hypotheses
are not supplied by the present compact-group decimation.
The justified import is the need to quantify a
\emph{scale-dependent negative-curvature budget}; it is not
a ready-made Yang--Mills log-Sobolev inequality.
This targeted check does not replace the scheduled broad
literature sweep and companion review at V30.

The exact-incidence checker visits all 760 target blocks in
nine lattices: $L=4$ with $q\in\{3,4,6,12\}$;
$L=8$ with $q\in\{3,4,6\}$; and $L=12$ with $q\in\{3,4\}$.
It verifies the complete free ownership
partition, $E_{b_0}\subset R$, the 30 outside occurrences,
every double-selected pair, and the constants $60,170/3$.
The coordinate argument above covers all admissible lattice
sizes; the finite tests do not substitute for that proof.
Independent rational finite-law fixtures check the marginal
Hessian sign and shared-parameter additivity. V26's validated
moment/response computation is rerun, rather than replacing
its enclosure by floating-point diagnostics.
\[
 \text{\path{scripts/verify_ym_restart_v27_retained_curvature.py}}
\]
contains the checks; certificates and all build products
remain outside \path{latex_notes}.

Deleting this append and restoring V26's title reconstructs its
615073-byte source with SHA--256
\path{C508AFFF054200A24413996A8B6B86B09F176449C59E3B630A885A43481E3947}.
All 46 other manuscript sources are preserved. Only the immediate
active predecessor is replaced.

The next useful step is no longer to ask whether the untouched
plaquettes restore pointwise convexity for this exact pass:
\eqref{r27:negative-bound} settles that question negatively.
Instead, quantify the actual law's negative-curvature weight
or an adjacent-scale source-metric comparison which can
tolerate this region. A lower probability bound as small as
\eqref{r27:probability} gives neither a useful upper exceptional-set
bound nor a cofinal entropy estimate. Reblocking again may change
the situation, but its curvature and metric transport must be
computed from the actual new normalized marginal.

\paragraph{Status.}
The opposed native-source obstruction now survives the complete
exact packed decimation, with three negative full-action and
weighted-Ricci directions on an open positive-probability event.
Constant linkwise metric rescalings cannot repair this pointwise
sign. Global endpoint affinity, cofinal physical-time and
cutoff estimates, an interacting four-dimensional continuum
construction, and a positive physical Hamiltonian mass gap
remain unproved.

% END CONSOLIDATED V27 APPEND

% BEGIN CONSOLIDATED V28 APPEND -- 2026-09-19
\clearpage
\section{V28: Gauge-horizontal escape and flat critical families}
\label{r28:context}

V27 rules out pointwise positive curvature for one exact retained-link
diffusion in its product metric. Before choosing a replacement metric
or an escape-from-saddles argument, we must determine whether its bad
direction is merely gauge redundancy. This append performs that test
on the \emph{same exact packed marginal}, not on an independent
surrogate model.

The answer is quantitative. At a specified central retained
configuration, the direction survives gauge projection and changes
retained Wilson-loop traces. On the $4^3\times8$ lattice the full
action curvature in its normalized horizontal projection belongs to
\[
 (-125.944223,-125.944222).
 \tag{V28.1}\label{r28:main-value}
\]
This is an action-Hessian escape direction, not a transfer eigenvalue.
We also prove that the exact retained action has a continuous family
of gauge-inequivalent flat-holonomy critical points. Thus a method
requiring all critical points to be isolated cannot be applied without
modification. The useful next object is transverse coercivity near
that family, together with its actual induced probability law.

\subsection{Scope, source audit and the critical-point gauge identity}
\label{r28:gauge}

Use V27's lattice, packing and notation:
$\gamma=10$, $4\mid L$, $L\geq4$, $q\geq3$,
$M=\SU(2)^R$, and $V_R$ as in \eqref{r27:exact-action}.
Let $\mathcal G=\SU(2)^\Lambda$ act by
\[
 (g\eta)_e=g_{s(e)}\eta_e g_{t(e)}^{-1}.
 \tag{V28.2}\label{r28:gauge-action}
\]
Each bucket $H_b$ is invariant under this action: apply the same
transformation to its free integration variables and use Haar
invariance and the trace invariance of every plaquette. Therefore
$V_R$ is invariant. This does not require any factorization of
the retained law.

The graph with vertex set $\Lambda$ and retained edges $R$ is
connected. Indeed, every non-$\mu$ edge is retained and connects
the full transverse three-torus in each fixed $x_\mu$ layer;
the retained $\mu$-links at time zero connect consecutive layers.
Fix a root $o$. The based group
$\mathcal G_*=\{g:g_o=I\}$ acts freely: if it fixes a link
configuration, propagation along a retained path from the root
forces every $g_x=I$.
Its action is isometric in the product round metric.
Consequently the based quotient $B=M/\mathcal G_*$ is a smooth
compact Riemannian quotient.

For clarity, the last assertion can also be seen using a retained
spanning tree. Successively gauging its links to identity gives
unique based gauge coordinates and free chord coordinates.
Averaging a local section yields the usual quotient metric whose
tangent vectors are horizontal lifts. Tree gauge and cycle-space
projection are established machinery: compare f15 V99 and f16.2
V62, V71 and V86. We use them on V27's newly retained graph.
The winding example in f16.2 V71 is also inherited; its persistence
as a \emph{critical family after this exact marginalization} is
proved below.

\begin{lemma}[Central critical points and gauge-null Hessians]
\label{r28:critical}
Every central assignment $\eta_e\in\{I,-I\}$ is a critical point
of $V_R$. At a critical point, the Hessian of any smooth
gauge-invariant function annihilates every gauge-orbit tangent,
including all mixed pairings with that tangent.
\end{lemma}
\begin{proof}
Simultaneous conjugation by one $h\in\SU(2)$ fixes a central
assignment and rotates each Lie-algebra tangent by $\Ad_h$.
There is no nonzero rotation-invariant covector on $\mathbb R^3$.
Applying this to variations supported on each individual retained
link proves $dV_R=0$.

For the second assertion let $K$ be a fundamental gauge vector
field and $f$ invariant. The identity $Kf=0$, differentiated in
an arbitrary direction $X$, gives
\[
 \operatorname{Hess}f(X,K)+df(\nabla_XK)=0 .
\]
At a critical point the second term vanishes. Symmetry of the
Hessian proves the claim in either argument.
\end{proof}

Thus, if $P_\eta$ is orthogonal projection onto the horizontal
subspace at such a point,
\[
 \operatorname{Hess}V_R(\xi,\xi)
 =\operatorname{Hess}V_R(P_\eta\xi,P_\eta\xi)
 =\operatorname{Hess}_B\overline V_R(d\pi\xi,d\pi\xi).
 \tag{V28.3}\label{r28:H-horizontal}
\]
The last equality follows by lifting a quotient curve: at a
critical point its acceleration contributes no first-derivative
term to the second variation. These standard critical-point
identities agree with Propositions A.60--A.62 of Capel et al.,
\href{https://arxiv.org/abs/2606.13453v2}
{Rapid mixing for Gibbs measures in Riemannian manifolds},
revision 17 September 2026. The native calculation below, not
that general identity, is the new result.

The full quotient $M/\mathcal G$ is different from $B$.
At a central assignment it has a nontrivial stabilizer:
the residual global conjugation group. We will not treat that
singular quotient as a smooth manifold or count three colour
directions as three independent physical modes.

\subsection{The native horizontal metric and observable witnesses}
\label{r28:metric}

Choose a target block $b_0$. Complete its opposed reference
exterior from V26 to a \emph{fully central} assignment $\eta^c$:
all non-$\mu$ links are $I$; the corner private staples consist
of two $I$ and two $-I$; all three left middle staples are $I$,
and all three right middle staples are $-I$.
Every unspecified retained $\mu$-link is $I$.
The private-link construction of V23--V26 proves compatibility,
including the repeated retained neighbours when $L=4$.

In signed Lie coordinates $\delta\eta_e=\eta_e A_e$, the six
selected source links of \eqref{r27:sources} have the same
dual-staple coordinates as their actual parallel links: with
identity connectors the private trace is
$\tfrac12\Tr(U_{\rm free}\eta_e^{-1})$.
Write $J\in\mathbb R^R$ for the cochain equal to $+1$ on
$s_{L0},s_{L+},s_{L-}$, $-1$ on the right triple, and zero
elsewhere. Then
\[
 \xi(a)_e=J_ea/\sqrt6,\qquad |\xi(a)|^2=|a|^2 .
 \tag{V28.4}\label{r28:J}
\]

Let $d$ be the oriented incidence map from vertex functions
to edge cochains,
$(dz)_e=z_{s(e)}-z_{t(e)}$. At a central configuration
the adjoint transports are identity, so the vertical space
is $(\operatorname{im}d)\otimes\mathbb R^3$.
The horizontal projection is
\[
 P=\bigl[I-d(d^*d)^\dagger d^*\bigr]\otimes I_3,\qquad
 h(a)=P\xi(a).
 \tag{V28.5}\label{r28:projection}
\]
The dagger is the inverse on mean-zero vertex functions.
Fixing the root gives the same edge-space projection.
In particular
\[
 |h(a)|^2=\alpha_{b_0,L,q}|a|^2,\qquad
 \alpha_{b_0,L,q}
 =1-\frac16\langle d^*J,(d^*d)^\dagger d^*J\rangle .
 \tag{V28.6}\label{r28:alpha}
\]
This is the native metric comparison; the gauge projection
is not assumed to preserve the unweighted source norm.

\begin{proposition}[Uniform nonvanishing and Wilson-loop detection]
\label{r28:witness}
For every permitted lattice and target block,
\[
 \frac{m_L}{6}\leq\alpha_{b_0,L,q}\leq1,\qquad
 m_L=\begin{cases}3,&L=4,\\2,&L\geq8.\end{cases}
 \tag{V28.7}\label{r28:alpha-lower}
\]
There are $m_L$ edge-disjoint, entirely retained plaquette
cycles detecting $h(a)$ whenever $a\ne0$.
For each of them, along
$\eta_e(s)=\eta^c_e\exp(sh(a)_e)$,
\[
 \tfrac12\Tr U_C(\eta(s))
       =-\cos\bigl(2s|a|/\sqrt6\bigr),\qquad
 \left.\frac{d^2}{ds^2}\tfrac12\Tr U_C(\eta(s))\right|_{s=0}
       =\frac23|a|^2 .
 \tag{V28.8}\label{r28:Wilson}
\]
\end{proposition}
\begin{proof}
Use the two $\mu\nu$ plaquettes joining
$(s_{L+},s_{R+})$ and $(s_{L-},s_{R-})$.
When $L=4$, also use the wraparound plaquette joining the two
side links. V27's incidence proof shows they are entirely
retained. They are edge-disjoint, since their times are
$t+1,t-1,t$, respectively. Each oriented cycle cochain $c_C$
has norm squared four and $d^*c_C=0$.
Moreover $|\langle J,c_C\rangle|=2$.
Orthogonal projection and Bessel's inequality therefore give
\[
 |P(J/\sqrt6)|^2
 \geq\sum_C\frac{|\langle J/\sqrt6,c_C\rangle|^2}{4}
 =\frac{m_L}{6}.
\]
The upper bound is the contractivity of an orthogonal projection.

Every edge component of $h(a)$ is a real multiple of the same
$a$. All its exponentials commute. Subtracting a gradient
does not change cycle circulation, and each witness cycle has
central holonomy $-I$ at $\eta^c$.
Its holonomy at time $s$ is therefore
$-\exp(\pm2sa/\sqrt6)$, proving \eqref{r28:Wilson}.
Its trace is gauge invariant and changes, so this curve cannot
lie in one full gauge orbit.
\end{proof}

Combining \eqref{r28:H-horizontal} with V27's complete outside
budget and V26's reference log-message bound gives
\[
 \operatorname{Hess}_B\overline V_R(d\pi\xi(a),d\pi\xi(a))
 <-d_L|a|^2,\qquad
 d_L=\begin{cases}
  17.287400,&L=4,\\
  309311/15000,&L\geq8 .
 \end{cases}
 \tag{V28.9}\label{r28:escape-lower}
\]
Here $d_L=77.287400-C_L$; $309311/15000>20.6207$.
The three horizontal colour directions are linearly independent
before residual conjugation by
\eqref{r28:alpha-lower}. Their quotient-normalized Rayleigh
values are below $-d_L/\alpha_{b_0,L,q}$.
No full-Hessian invariant-subspace claim is needed.

All directions $h(a)$ with the same $|a|$ are related by global
conjugation; they give at least one physical radial escape,
not a count of three physical excitations.
At a stationary point the Hessian is an intrinsic quadratic
form independent of the connection. Hence changing to any
smooth positive Riemannian metric cannot make this action
Hessian positive semidefinite there. It can change its
normalization and the Ricci contribution to a diffusion
estimate; those are separate questions.

\subsection{Exact rational projection and full curvature on the test lattice}
\label{r28:certificate}

We now evaluate the metric comparison rather than only bound it.
Take $L=4$ and the target anchored at $(0,0,0,2)$.
The retained pattern is translation invariant in $x_\mu$.
On the transverse torus
$T=(\mathbb Z/4\mathbb Z)^2\times\mathbb Z/(2q)\mathbb Z$,
let $\mathsf L_T$ be the usual unweighted graph Laplacian
and $D$ the diagonal indicator that the $\mu$-link at that
transverse site is retained. Let $j\in\mathbb R^T$ have
entries $+1$ at
\[
 (3,1,2),\quad(0,1,3),\quad(0,1,1)
\]
and entries $-1$ at
\[
 (2,1,2),\quad(1,1,3),\quad(1,1,1).
\]
These are the six sources of \eqref{r28:J}.
The graph Laplacian decomposes along the longitudinal
four-cycle as
\[
 d^*d=I\otimes\mathsf L_T+\mathsf L_{C_4}\otimes D .
 \tag{V28.10}\label{r28:graph-L}
\]
The divergence of $J$ is
$(\delta_0-\delta_1)\otimes j$.

\begin{proposition}[Rational horizontal norm]
\label{r28:rational-alpha}
For this target,
\[
 \alpha_{4,q}
 =1-\frac16\left[
  j^{\mathsf T}(\mathsf L_T+2D)^{-1}j+
  j^{\mathsf T}(\mathsf L_T+4D)^{-1}j\right].
 \tag{V28.11}\label{r28:alpha-rational}
\]
The following intervals have rational endpoints and are certified
by exact linear solves:
\[
\begin{array}{c|c|c}
q&|T|&\alpha_{4,q}\\ \hline
3&96 &(0.77190759,\,0.77190760)\\
4&128&(0.77246418,\,0.77246419)\\
6&192&(0.77196414,\,0.77196415)
\end{array}
 \tag{V28.12}\label{r28:alpha-table}
\]
In particular,
\[
 \alpha_{4,4}
 =\frac{13060748345027910925208597681161}
        {16907901496085842905045272675304}.
 \tag{V28.13}\label{r28:alpha-exact}
\]
\end{proposition}
\begin{proof}
The longitudinal eigenvalues are $0,2,4,2$.
In the unitary Fourier basis the squared component of
$\delta_0-\delta_1$ at eigenvalue $\lambda_k$ is $\lambda_k/4$.
The zero mode contributes nothing, the two eigenvalue-two
modes have total weight one, and the eigenvalue-four mode
has weight one. Substitution in \eqref{r28:alpha} gives
\eqref{r28:alpha-rational}.

For $\lambda>0$,
$\mathsf L_T+\lambda D$ is positive definite:
its quadratic form can vanish only on a constant vector,
and $D$ is nonzero. Its entries are integers, so the solution
with integer right side $j$ is rational.
The checker solves both systems exactly, verifies each
rational residual, reconstructs the vertex potential on the
full four-dimensional retained graph, and verifies
$d^*(J-d\phi)=0$ and
\[
 |J-d\phi|^2=\langle J,J-d\phi\rangle=6\alpha .
\]
Thus the Fourier reduction is checked against the original
retained graph, not merely against itself. Rational comparisons
give \eqref{r28:alpha-table} and \eqref{r28:alpha-exact}.
\end{proof}

The $q=4$ lattice also permits an exact evaluation of the
\emph{whole} directional action Hessian before gauge projection.

\begin{proposition}[Evaluated gauge-horizontal escape]
\label{r28:exact-escape}
At the specified central assignment on $4^3\times8$,
\[
 \operatorname{Hess}V_R(\xi(a),\xi(a))
       =-(20+600c_-)|a|^2 ,
 \qquad c_-=\E_{\mu_-}(u_1v_1)
 \tag{V28.14}\label{r28:raw-escape}
\]
with exactly V26's opposed native middle-pair law.
On the normalized horizontal three-dimensional subspace
the Hessian quadratic form is the scalar
\[
 -\frac{20+600c_-}{\alpha_{4,4}},
 \tag{V28.15}\label{r28:normalized-escape}
\]
which lies in \eqref{r28:main-value}.
\end{proposition}
\begin{proof}
At $q=4$ the even free time slices are 2 and 6; slices 0 and 4
are retained. No selected source link meets a free block other
than the target, so all other $\log H_b$ have zero derivatives
along the original six-link direction $\xi(a)$.

For every remaining incident plaquette, all links are central.
If its central trace is $\epsilon_p\in\{1,-1\}$ and its oriented
$J$ circulation is $n_p$, its action Hessian is
$\gamma\epsilon_p n_p^2|a|^2/6$.
The outside single-source plaquettes can be counted as follows:
\[
 \begin{array}{c|rr}
 \text{source type}&\epsilon_p=+1&\epsilon_p=-1\\ \hline
 \text{left side link}&4&0\\
 \text{right side link}&0&4\\
 \text{four time-offset links together}&8&8
 \end{array}
\]
For a side link the four single neighbours are its two
$\rho$-neighbours and two time neighbours. For a time-offset
link they are its side neighbour, two corner neighbours and
its next even-time neighbour, giving two of each sign.
There are also three double-source plaquettes with
$\epsilon_p=-1$ and $|n_p|=2$.
Hence the exact outside curvature is
\[
 \frac{10}{6}\bigl(12-12-3\cdot4\bigr)|a|^2
       =-20|a|^2 .
\]
V26 gives target contribution $-600c_-|a|^2$.
This proves \eqref{r28:raw-escape}.

Now apply the full-action identity \eqref{r28:H-horizontal}
and divide by $\alpha_{4,4}$. The projected vector generally
has many nonzero link components. We do \emph{not} drop the
other bucket derivatives in that projected direction:
their combined effect is already covered by the exact
gauge-null Hessian identity.
The rational value \eqref{r28:alpha-exact} and the inherited
certified native moment give \eqref{r28:main-value}.
For verification the checker recomputes the moment from the
exact positive-polynomial comparison; its residual error
is smaller than the interval margin.
\end{proof}

This is a strict descent of the invariant action on nearby
inequivalent configurations. It need not be an eigenvector of
the entire quotient Hessian, and the displayed number is not
a diffusion mixing rate or a physical mass.

\subsection{The exact marginal retains a flat critical family}
\label{r28:flat}

The nontrivial winding configurations themselves are not new;
f16.2 V71 gives an explicit gauge-invariant winding witness.
What needs checking here is that integrating the entire V23
packing neither lifts their action nor introduces a force
transverse to their restricted family.

\begin{theorem}[Gauge-inequivalent nonisolated critical points]
\label{r28:flat-critical}
For arbitrary $g_0,\ldots,g_{L-1}\in\SU(2)$ define a retained
configuration by
\[
 \eta_e(\mathbf g)=
 \begin{cases}
    g_k,&e\text{ is a retained }\mu\text{-link with }x_\mu=k,\\
    I,&e\text{ has another orientation}.
 \end{cases}
 \tag{V28.16}\label{r28:flat-family}
\]
Then for the exact action \eqref{r27:exact-action},
\[
 V_R(\eta(\mathbf g))=V_R(I),\qquad
 dV_R|_{\eta(\mathbf g)}=0 .
 \tag{V28.17}\label{r28:flat-stationary}
\]
Its based quotient, and also its set of full gauge orbits,
therefore contain nonisolated critical points.
This statement holds for every coupling $\gamma\geq0$;
it does not assert that these points are global minima.
\end{theorem}
\begin{proof}
Every untouched plaquette equals identity: a plaquette using
$\mu$ contains $g_k$ and $g_k^{-1}$ at the same longitudinal
base layer, and every other factor is $I$.
Its action is zero and all its first derivatives vanish.

Fix any bucket $b$ at longitudinal layer $k$. Every retained
$\mu$-link in its 29 plaquettes has that same base layer.
Its remaining retained links are transverse, lying at layers
$k$ or $k+1$. Apply the gauge transformation equal to $I$ on
layer $k$, $g_k$ on layer $k+1$, and identity on all other
layers. These layers are distinct, including wraparound.
It sends every retained argument of $H_b$ in
\eqref{r28:flat-family} to $I$.
By invariance of the free integral, $H_b$ has the same value
as at its all-identity exterior.

At that exterior the differential of $H_b$ vanishes:
simultaneous conjugation fixes its arguments, so the argument
of Lemma~\ref{r28:critical} applies to this smooth invariant
function itself. The preceding gauge transformation is a
diffeomorphism of its retained argument space; therefore
$dH_b=0$ at \eqref{r28:flat-family}.
The transformations may differ from bucket to bucket.
They are used to prove the value and differential of each
summand separately, not to assert one global gauge fixing
of a nontrivial winding configuration.
Summing the untouched actions and all $-\log H_b$ proves
both parts of \eqref{r28:flat-stationary}.

At any odd time slice every $\mu$-link is retained.
Its winding loop has holonomy
$g_0g_1\cdots g_{L-1}$.
For $g_k=\exp(\theta\mathbf i/L)$ its normalized trace is
$\cos\theta$. Distinct $\theta\in(0,\pi)$ give distinct full
gauge orbits and a continuous family of critical points.
The based quotient is smooth, so differentiating its
stationary gradient along this family supplies a nonzero
Hessian-kernel direction there. This is not a gauge
direction: the derivative of $\cos\theta$ is nonzero on
that interval.
\end{proof}

Flatness of the original plaquettes alone would not have
proved $dV_R=0$ after marginalization. The separate
bucket gauge transformations establish that stronger fact.
Nor does a flat critical family prove a zero eigenvalue
above the physical vacuum. It is a property of a finite
retained action, not a spectrum of the transfer Hamiltonian.

\subsection{The probability measure on a quotient still has to be paid}
\label{r28:measure}

The preceding action statements must not be substituted for
a curvature calculation of the pushed-forward law.
For the based Riemannian quotient $\pi:M\to B$, write
$v(b)$ for its orbit volume in the induced metric.
The coarea formula gives, up to a constant,
\[
 \pi_*\mu_R(db)
   \propto e^{-\overline V_R(b)}v(b)\,d\operatorname{vol}_B(b),
 \qquad
 V_B^{\rm prob}=\overline V_R-\log v .
 \tag{V28.18}\label{r28:orbit-volume}
\]
Thus $\operatorname{Ric}_B+
\operatorname{Hess}_B V_B^{\rm prob}$ is not obtained merely
by deleting the vertical rows of V27's weighted Ricci form.
Tree-gauge Haar coordinates can instead absorb this volume
factor into their reference measure, but their product metric
is not automatically the induced quotient metric.
Neither coordinate convention permits dropping a Jacobian
from the corresponding Dirichlet form.

Equation \eqref{r28:orbit-volume} follows by integrating first
over each orbit in the ordinary Riemannian coarea formula.
The isometric action makes the integrand $e^{-V_R}$ constant
on that orbit, but does not imply that \emph{different}
orbits have equal volume.
No bound on $v$, its Hessian, or quotient weighted Ricci
is proved here. The action-Hessian result remains valid
regardless of this missing measure estimate.

\subsection{Literature import, nonduplication and the next positive estimate}
\label{r28:next}

The targeted search found Capel et al.'s revised
\href{https://arxiv.org/abs/2606.13453v2}
{rapid-mixing preprint}. We inspected its main hypotheses,
the coarea discussion and Appendix A.5, including the
rendered hypothesis and Hessian pages.
Its strategy uses escape from saddles, not convexity everywhere.
However, the stated main theorem assumes a free isometric
action, additional quotient geometry, a unique minimum,
isolated critical points, escape bounds and gradient control.
Theorem~\ref{r28:flat-critical} contradicts the isolated-point
hypothesis for the direct choice $F=V_R$ on the based gauge
quotient. A corrected probability potential as in
\eqref{r28:orbit-volume} is a different function and must be
computed before testing any of those hypotheses.

We therefore import a direction of attack, not the theorem's
mixing conclusion: treat the negative Hessian as a verified
escape direction, while controlling transverse fluctuations
near critical families and accounting for their measure.
A Morse--Bott analysis would additionally require that the
Hessian kernel equal the tangent space of a smooth critical
manifold; the existence of the family proved here does
\emph{not} establish that condition.

The next concrete calculation is the native Hessian transverse
to \eqref{r28:flat-family}. The coherent slot-and-transport
expansion in f16.2 V62 is relevant, but must be pulled back
to the actual retained Wilson links, with the full global
gauge and winding kernel kept. Its auxiliary slot-space
kernel is not automatically the native kernel. A useful
estimate must then retain the actual shared probability law,
control its volume and coupling dependence, and connect
to the unchanged physical endpoint-affinity target.
Neither saddle escape at one configuration nor a fixed-box
transverse estimate would complete those steps.

The checker verifies connectivity, the disjoint witness cycles,
their incidences and the flat-family local gauge supports for
all 160 target blocks in five lattices. It solves six rational
transverse systems and checks 1664 full-graph divergence
residuals, plus the exact projection energies.
The $4^3\times8$ outside-curvature count is independently
enumerated from the original plaquette words.
\begin{center}\small
\path{scripts/verify_ym_restart_v28_horizontal_escape.py}
\end{center}
contains these checks. Geometry and gauge arguments are proved
above for all stated sizes; finite fixtures are not replacements
for those proofs.

Removing this append and restoring V27's title reconstructs its
635865-byte source, SHA--256
\path{8E1171F239801510633D4EBBD96777FB0EC4229C52291027DF9A6D7EBA5A2826}.
All 46 other manuscript sources are unchanged.
The next scheduled companion review and broad literature sweep
remain V30. Build products and the inspected reference PDF
are kept outside \path{latex_notes}.

\paragraph{Status.}
The native nonconvex direction is gauge-observable, has a
certified horizontal norm, and has an evaluated quadratic
descent coefficient on one exact finite lattice.
The same exact marginal retains nonisolated flat critical
families. A quotient probability-curvature estimate,
native transverse coercivity with the required uniformity,
global physical contraction, continuum construction and the
Yang--Mills mass gap remain unproved.

% END CONSOLIDATED V28 APPEND

% BEGIN CONSOLIDATED V29 APPEND -- 2026-09-19
\clearpage
\section{V29: native Maxwell coercivity and the exact flat minimum}
\label{r29:start}

\subsection{The positive estimate selected after V28}

V28 exhibited a genuine gauge-horizontal descent direction at one
nonminimal critical configuration, but also a continuous flat critical
family. This append treats the complementary, positive question:
what does the \emph{full exact retained action} control at the
all-identity configuration? The answer is a discrete Maxwell energy
with an explicit, volume-independent comparison coefficient.
Its full native kernel can be determined. We also identify all global
minima of this marginal action. These are finite-regulator statements
for the same packed $\SU(2)$ Wilson model, not statements about its
physical transfer Hamiltonian.

The coherent positive diagram expansion and its local-frame kernel
were already proved in archive $C$, V62. They are credited inputs,
not new results here. We audited the compact Haar expansion, the
marked-component insertion and the native endpoint-plaquette
interpretation in that section. What is new in this append is their
assembly on the \emph{entire actual retained graph}: an explicit
extension to original link cochains, the Maxwell comparison, the
complete native kernel, and a global equality characterization.
V28 proved flat critical families but did not prove that these
exhaust the global minima.

The resulting kernel has more linear directions than can be tangent
to a smooth family of minima. Consequently the action at the identity
is not Morse--Bott, even on the smooth based gauge quotient.
This directs the next work toward the nonlinear winding sector and
its actual probability measure, not an assumed nondegenerate
Gaussian minimum. The constant-field quartic comparisons in $C$,
V44 and the normal chart in $C$, V41 are relevant predecessors;
a quartic model is not thereby the native physical Hamiltonian.

\subsection{Literal packed geometry and normalizations}
\label{r29:geometry}

Use the V23--V28 packing on the periodic lattice
\[
 \Lambda=(\mathbb Z/L\mathbb Z)^3\times
                 (\mathbb Z/(2q)\mathbb Z),\qquad
 L\in4\mathbb N,\quad L\ge4,\quad q\ge3 .
 \tag{V29.1}\label{r29:lattice}
\]
The coordinate directions are $\mu,\nu,\rho,t$.
For every $x_\mu$, place a $2$-column by $3$-row rectangle of
free $\mu$-links at each anchor with
$x_\nu,x_\rho\equiv0\pmod4$ and even $t\notin\{0,q\}$.
The column and row coordinates are $\nu$ and $\rho$.
Write $\mathcal F_b$ for its six links, $\mathcal F$ for their union,
and $\mathcal R$ for all other positively oriented links.
The internal graph of a block has seven edges and two square faces.
There are $22$ private plaquettes and seven two-free-link plaquettes;
every incident original plaquette belongs to just one block.
Every non-$\mu$ link is retained. The retained graph is connected,
and contains a complete winding cycle in each coordinate direction:
the $\mu$ cycle can be taken at $t=0$.

Let $\mathcal P_b$ be the $29$ plaquettes incident to block $b$ and
$\mathcal P_0$ the untouched plaquettes. With
$w(U)=\tfrac12\Tr U=\operatorname{Re}U$ and normalized Haar measures,
the exact marginal action, up to an irrelevant constant, is
\[
 V_{\mathcal R}(\eta)
   =\gamma\sum_{p\in\mathcal P_0}(1-w(U_p(\eta)))
           -\sum_b\log H_b(\eta),\qquad
 H_b=e^{-29\gamma}Z_b,\qquad \gamma>0 .
 \tag{V29.2}\label{r29:action}
\]
Identify $\SU(2)$ with the unit quaternions $S^3$.
Each private plaquette at free link $i$ gives a retained
three-link transport $p_{i\alpha}\in S^3$ between its endpoints;
write $n_i$ for the number of these slots.
The four corner vertices have $n_i=4$, the other two have $n_i=3$,
and $n_i+\deg(i)=6$.

An internal plaquette gives
\nopagebreak[4]
\[
 \begin{gathered}
 R_{ij}\in\operatorname{SO}(4),\qquad
 R_{ij}(x)=\ell_{ij}x r_{ij}^{-1},\\
 Z_b=\int_{(S^3)^6}
 \exp\!\left\{\gamma\sum_{\{i,j\}}x_i\cdot R_{ij}x_j
       +\gamma\sum_{i,\alpha}x_i\cdot p_{i\alpha}\right\}
                  \prod_i d\sigma(x_i).
 \end{gathered}
 \tag{V29.3}\label{r29:compact}
\]
The two retained links $\ell_{ij},r_{ij}$ connect the respective
left and right endpoints of $i,j$.
These slots and transports are dependent functions of the original
retained links. No independence of their law is asserted.

Use the unit-round $S^3$ metric on every retained link, with imaginary
quaternions having their Euclidean norm. A native tangent at
$\eta=1$ is a cochain
$a=(a_e)_{e\in\mathcal R}\in(\mathbb R^3)^{\mathcal R}$.
Let $b_{i\alpha}$ be the derivative of its private three-link path.
If $i$ is the $\mu$ link at $x$ and the private plaquette extends
in the positive $\nu$ direction, for example,
\[
 b_{i\alpha}=a_\nu(x)+a_\mu(x+\widehat\nu)
                               -a_\nu(x+\widehat\mu).
 \tag{V29.4}\label{r29:private}
\]
Reverse transverse plaquettes use the corresponding oriented path.
For an internal edge directed from $i$ at $x$ to $j$ at
$x+\widehat\nu$,
\[
 A_{ij}z=a_\nu(x)z-z a_\nu(x+\widehat\mu),\qquad
 t_{ij}:=A_{ij}1=a_\nu(x)-a_\nu(x+\widehat\mu).
 \tag{V29.5}\label{r29:transport}
\]
Changing this orientation negates both the mismatch and its
transport tangent. All subsequent norms are independent of that
choice. Define the retained action Hessian quadratic form
\[
 \mathcal H_{\mathcal R}(a)
       =\operatorname{Hess}V_{\mathcal R}(1)[a,a].
 \tag{V29.6}\label{r29:hessian}
\]

\subsection{The inherited coherent identity and its exact scope}
\label{r29:coherent}

Here is the part of $C$, V62 used below, with the normalization made
explicit. At $p_{i\alpha}=1$, $R_{ij}=I$, write $Q_b$ for the Hessian
of $-\log Z_b$ in arbitrary unit-slot and orthogonal-transport
directions $(b,A)$. Its first derivative is zero. With
$D_\gamma=4+6\gamma$, the audited bound is
\[
 \begin{split}
 Q_b(b,A)\ \ge{}&
  \frac{\gamma^2}{D_\gamma}
       \sum_i\sum_{\alpha<\beta}|b_{i\alpha}-b_{i\beta}|^2\\
 &+\frac{\gamma^3}{D_\gamma^2}
       \sum_{\{i,j\}}\sum_{\alpha,\beta}
               |b_{i\alpha}-b_{j\beta}-t_{ij}|^2\\
 &+\frac{\gamma^4}{D_\gamma^4}
       \sum_{C\text{ square}}\|A(C)\|_{\rm F}^2 .
 \end{split}
 \tag{V29.7}\label{r29:paid}
\]
It is one simultaneous lower bound, not a sum of three independently
spent estimates. In particular $Q_b\ge0$.

For clarity, the exact identity behind this import is
\[
 Q_b=\mathbb E_0\!\left[
  \sum_{P\text{ open}}
       |b_{i(P)\alpha(P)}-b_{j(P)\beta(P)}-A(P)1|^2
       +\frac14\sum_{C\text{ closed}}\|A(C)\|_{\rm F}^2\right].
 \tag{V29.8}\label{r29:diagram-hess}
\]
To recall why this has the required sign, expand every exponential
in \eqref{r29:compact}. At a vertex with $N_i$ half-edges the
even Haar moment pairs them with coefficient
$u(N_i)=(4\cdot6\cdots(N_i+2))^{-1}$, with $u(0)=1$;
odd moments vanish. A pairing diagram has positive coherent weight
\[
 a_D=\frac{\gamma^{\sum_e k_e+\sum_{i,\alpha}\ell_{i\alpha}}}
              {\prod_e k_e!\prod_{i,\alpha}\ell_{i\alpha}!}
                 \prod_i u(N_i)\,4^{\#\text{ closed strands}}.
 \tag{V29.9}\label{r29:weights}
\]
These weights sum to $Z_b(1)$, defining $\mathbb P_0$.
For general exterior data its relative factor is
\[
 f_D=
 \prod_{P\text{ open}}p_{i(P)\alpha(P)}^{\mathsf T}
                   R(P)p_{j(P)\beta(P)}
       \prod_{C\text{ closed}}\frac{\Tr R(C)}4 .
 \tag{V29.10}\label{r29:factor}
\]
At coherence every factor is one and its first derivative vanishes.
The negative second derivative of an open factor is its squared
tangent mismatch; that of a closed factor is
$\|A(C)\|_{\rm F}^2/4$.
The series is absolutely dominated by $\sum_Da_D=Z_b(1)$,
and derivatives by its finite occurrence moments, proving
\eqref{r29:diagram-hess} and stationarity.

The insertion argument in $C$, V62 gives the expected numbers of
zero-edge open strands, one-edge open strands and elementary
square closed strands as, respectively,
\[
 \gamma^2\mathbb E_0(4+N_i)^{-1},\quad
 \gamma^3\mathbb E_0[(4+N_i)(4+N_j)]^{-1},\quad
 4\gamma^4\mathbb E_0\prod_{i\in C}(4+N_i)^{-1}.
\]
Inserting a marked strand leaves every unmarked component in place;
occurrence-label choices cancel the exponential factorial ratios.
Moreover $\mathbb E_0N_i\le6\gamma$.
Convexity of the exponential and concavity of the logarithm bound
each displayed expectation below by the corresponding power of
$D_\gamma^{-1}$. Selecting these disjoint component types in
\eqref{r29:diagram-hess} yields \eqref{r29:paid}.
There is no empty-patch event or exponential isolation cost.

The other audited consequence is the exact local kernel:
on a rectangular internal graph,
\[
 Q_b=0\quad\Longleftrightarrow\quad
 b_{i\alpha}=\Omega_i1,\quad
 A_{ij}=\Omega_i-\Omega_j,\qquad
                 \Omega_i\in\mathfrak{so}(4).
 \tag{V29.11}\label{r29:aux-kernel}
\]
Indeed \eqref{r29:paid} forces equal slots, compatible edge
differences and zero square sums of $A$; integrating the latter
on the rectangle constructs $\Omega_i$.
Conversely the change of variables $x_i=O_i z_i$ for
$p_{i\alpha}=O_i1$, $R_{ij}=O_iO_j^{\mathsf T}$ leaves $Z_b$
unchanged. Differentiating gives the stated kernel.
Its unrestricted auxiliary dimension is not the native dimension.

Because the first derivative of $-\log Z_b$ vanishes at coherence,
pullback by the actual Wilson slot maps has no acceleration term.
Consequently the \emph{full} native quadratic form is exactly
\[
 \mathcal H_{\mathcal R}(a)
    =\gamma\sum_{p\in\mathcal P_0}|(d_1a)_p|^2
                         +\sum_b Q_b(b(a),A(a)),
 \tag{V29.12}\label{r29:full}
\]
where $d_1$ is the ordinary oriented plaquette curl at the identity.

\subsection{A canonical extension proves native Maxwell coercivity}
\label{r29:maxwell}

Define the linear extension $E$ to all original positive links by
\[
 (Ea)_e=a_e\quad(e\in\mathcal R),\qquad
 (Ea)_i=\overline b_i:=\frac1{n_i}\sum_\alpha b_{i\alpha}
                                      \quad(i\in\mathcal F).
 \tag{V29.13}\label{r29:extension}
\]
This is a declared averaging operator, not a minimizer over free
link configurations and not a probabilistic conditional mean.

\begin{theorem}[Maxwell comparison for the actual retained Hessian]
\label{r29:maxwell-thm}
For every native tangent, every packing in \eqref{r29:lattice}, and
every $\gamma>0$,
\[
 \boxed{\quad
 \mathcal H_{\mathcal R}(a)\ \ge\
       c_\gamma\sum_{p\text{ original}}|(d_1Ea)_p|^2,\qquad
 c_\gamma=\frac{9\gamma^3}{(4+6\gamma)^2}.
       \quad}
 \tag{V29.14}\label{r29:maxwell-bound}
\]
The coefficient is independent of $L,q$ and the number of blocks.
It satisfies $c_\gamma\sim\gamma/4$ as $\gamma\to\infty$.
\end{theorem}

\begin{proof}
Put $v_i=\sum_\alpha|b_{i\alpha}-\overline b_i|^2$.
For a private plaquette at $i$, the original curl of $Ea$ is,
up to sign, $\overline b_i-b_{i\alpha}$. For an internal
plaquette it is, up to sign,
$\overline b_i-\overline b_j-t_{ij}$.
The elementary variance identities give
\[
 \begin{split}
 \sum_{\alpha<\beta}|b_{i\alpha}-b_{i\beta}|^2&=n_i v_i,\\
 \sum_{\alpha,\beta}|b_{i\alpha}-b_{j\beta}-t_{ij}|^2
   &=n_i n_j|\overline b_i-\overline b_j-t_{ij}|^2
                                      +n_jv_i+n_iv_j .
 \end{split}
 \tag{V29.15}\label{r29:variance}
\]
They follow by expanding the squares and using the zero sum of
the deviations from each mean.

Write $A_\gamma=\gamma^2/D_\gamma$ and
$B_\gamma=\gamma^3/D_\gamma^2$.
The first two lines of \eqref{r29:paid}, after dropping the
nonnegative variance terms in its second line, imply
\[
 Q_b\ge3A_\gamma\sum_i v_i
        +9B_\gamma\sum_{\{i,j\}}
               |\overline b_i-\overline b_j-t_{ij}|^2
       \ge c_\gamma\sum_{p\in\mathcal P_b}|(d_1Ea)_p|^2.
\]
Here $n_i\ge3$, $n_i n_j\ge9$, and
$3A_\gamma\ge9B_\gamma=c_\gamma$ since $D_\gamma>3\gamma$.
Also $\gamma\ge c_\gamma$.
The plaquettes $\mathcal P_0,\mathcal P_b$ partition the original
plaquettes, so \eqref{r29:full} proves the assertion.
The cycle line of \eqref{r29:paid} is not needed for this lower
bound, but it remains part of the exact Hessian.
\end{proof}

In particular $c_{10}=1125/512=2.197265625$ in the metric used
here. This is an evaluated \emph{action-Hessian comparison
coefficient}; its units and meaning are not those of a physical
particle mass.

\subsection{The full native kernel and the finite-box constant}
\label{r29:kernel}

Let $\mathcal K=\ker d_1$ in the full real three-colour link
cochain space, and
$\mathcal K_{\mathcal R}=\{k|_{\mathcal R}:k\in\mathcal K\}$.
On the periodic lattice,
\[
 k_\alpha(x)=\phi(x)-\phi(x+\widehat\alpha)+h_\alpha,\qquad
       h_\alpha\in\mathbb R^3
 \tag{V29.16}\label{r29:closed}
\]
describes all closed cochains. The sign of the gradient is chosen
to agree with the infinitesimal gauge action.

\begin{theorem}[Exact native kernel]
\label{r29:kernel-thm}
At the identity,
\[
 \ker\operatorname{Hess}V_{\mathcal R}(1)=\mathcal K_{\mathcal R},
 \qquad
 \dim\mathcal K_{\mathcal R}=3(|\Lambda|-1)+12.
 \tag{V29.17}\label{r29:kernel-eq}
\]
The restriction $\mathcal K\to\mathcal K_{\mathcal R}$ is injective.
After passing to the smooth based gauge quotient, the action Hessian
therefore has a twelve-dimensional kernel at this point.
\end{theorem}

\begin{proof}
The form is positive semidefinite by \eqref{r29:full}.
If it vanishes on $a$, \eqref{r29:maxwell-bound} gives $d_1Ea=0$,
hence $a\in\mathcal K_{\mathcal R}$.

Conversely let $a$ be the restriction of a full closed cochain $k$,
and write $z_i=k_i$ on the free links.
The original private and internal curl equations say
\[
 b_{i\alpha}=z_i,\qquad z_i-z_j=A_{ij}1.
 \tag{V29.18}\label{r29:closed-local}
\]
For every elementary internal square $C$,
$A(C)z=\lambda_Cz-z\rho_C$.
The two coefficients are the curls of $k$ on the pure transverse
plaquettes at the two endpoints of the free $\mu$ links.
Those plaquettes are untouched; both curls are zero.
Thus $A(C)=0$. The internal rectangular graph has no additional
winding cycle, so its square faces generate its cycle space.
Choose $\Omega_*$ skew with $\Omega_*1=z_*$ and integrate $A$
from that root. This constructs
$A_{ij}=\Omega_i-\Omega_j$, and \eqref{r29:closed-local} implies
$\Omega_i1=z_i$ everywhere. Equation \eqref{r29:aux-kernel}
gives $Q_b=0$ for each block. The untouched curls are also zero,
so \eqref{r29:full} vanishes.

For completeness, the standard periodic closed-cochain
description \eqref{r29:closed} follows by Fourier transform:
at a nonzero frequency the curl equation says that the four
components are proportional to
$(e^{ik_\alpha}-1)_\alpha$, hence are a gradient; at zero
frequency they are the four constants $h_\alpha$.
This also proves the full-space dimension
$3(|\Lambda|-1)+12$.

If such a cochain restricts to zero on $\mathcal R$, its circulation
around each wholly retained coordinate winding cycle is zero.
The gradient telescopes and the circulation is the corresponding
period times $h_\alpha$, so every $h_\alpha=0$.
The retained graph is connected, so a gradient vanishing on it has
constant $\phi$. Thus the full cochain is zero, proving injectivity
and the dimension statement.
The based gauge action is free on the retained configuration
space, as proved in V28. Its vertical dimension is
$3(|\Lambda|-1)$, and its stationary Hessian vanishes on those
vertical vectors, giving the quotient assertion.
\end{proof}

The missing square rotations of the \emph{unrestricted} auxiliary
slot problem have not been silently dropped: their vanishing in
the converse used actual untouched endpoint plaquettes. This is
the key extra relation supplied by the native Wilson geometry.

\begin{corollary}[Explicit transverse coercivity at the identity]
\label{r29:finite-coercivity}
In the retained product metric, put
$\ell_{\max}=\max\{L,2q\}$. Then
\[
 \boxed{\quad
 \mathcal H_{\mathcal R}(a)\ \ge\
 \frac{9\gamma^3}{(4+6\gamma)^2}
       4\sin^2\!\left(\frac{\pi}{\ell_{\max}}\right)
          \operatorname{dist}(a,\mathcal K_{\mathcal R})^2 .
       \quad}
 \tag{V29.19}\label{r29:coercivity}
\]
\end{corollary}

\begin{proof}
For a Fourier frequency $k$, put
$\omega_\alpha=e^{ik_\alpha}-1$.
The curl symbol obeys
\[
 \sum_{\alpha<\beta}
       |\omega_\alpha u_\beta-\omega_\beta u_\alpha|^2
   =|\omega|^2|u|^2
                -\left|\sum_\alpha\overline\omega_\alpha u_\alpha\right|^2.
\]
Thus its nonzero squared singular values equal $|\omega|^2$.
Their minimum over the periodic frequencies is
$\lambda_\Lambda=4\sin^2(\pi/\ell_{\max})$.
Parseval therefore yields
$\|d_1u\|^2\ge\lambda_\Lambda\operatorname{dist}(u,\mathcal K)^2$.
For every $k\in\mathcal K$,
$\|Ea-k\|_{\rm full}\ge\|a-k|_{\mathcal R}\|_{\mathcal R}$;
taking infima shows
$\operatorname{dist}(Ea,\mathcal K)\ge
\operatorname{dist}(a,\mathcal K_{\mathcal R})$.
Apply \eqref{r29:maxwell-bound}.
\end{proof}

For $L=4,q=4,\gamma=10$ the displayed transverse coefficient is
$(1125/512)(2-\sqrt2)$. The comparison coefficient
$c_\gamma$ is volume-independent, but
$\lambda_\Lambda$ tends to zero with the longest lattice period.
This is not a volume-uniform positive Hessian bound, much less
a continuum mass gap. Conversely its decay is not a proof of
physical gaplessness: the action Hessian at one configuration is
not the interacting transfer Hamiltonian.

\subsection{Every global minimum is a restricted flat connection}
\label{r29:minima}

\begin{theorem}[Exact minimum and unique flat extension]
\label{r29:minima-thm}
For the packing above and $\gamma>0$,
\[
 V_{\mathcal R}(\eta)\ge V_{\mathcal R}(1).
 \tag{V29.20}\label{r29:min-bound}
\]
Equality holds if and only if $\eta$ is the restriction of a
full lattice connection with every original plaquette holonomy
equal to $1$. If equality holds, the free-link values in that
flat extension are unique.
The extension commutes with lattice gauge transformations.
\end{theorem}

\begin{proof}
Each factor in \eqref{r29:factor} lies in $[-1,1]$:
it is an inner product of unit vectors or the normalized trace
of a real orthogonal matrix. Hence $f_D\le1$, although it need
not be nonnegative. The positive, summable weights give
\[
 \frac{Z_b(\eta)}{Z_b(1)}=\mathbb E_0 f_D\le1.
 \tag{V29.21}\label{r29:Zmax}
\]
Every untouched plaquette action is nonnegative. Equation
\eqref{r29:action} proves \eqref{r29:min-bound}.
If equality holds, every untouched plaquette is flat and every
block satisfies equality in \eqref{r29:Zmax}.
Since $1-f_D\ge0$ and each allowed diagram has positive weight
for $\gamma>0$, equality forces $f_D=1$ for every such diagram.

Select a diagram consisting of just a zero-edge open strand
joining two distinct slots at a vertex, with no other activity.
Its factor is $p_{i\alpha}\cdot p_{i\beta}$, so these two unit
vectors must coincide. All slots at $i$ therefore have a common
value $p_i$. Next select a single-edge open strand joining any
slot at $i$ to any slot at $j$. It forces
$p_i\cdot R_{ij}p_j=1$, hence $p_i=R_{ij}p_j$.
Set the free link $U_i=p_i$. All private and all internal
plaquettes are then flat, by their literal three-link path and
Spin$(4)$ formulas. Along with the untouched plaquettes this
constructs a flat full connection. Every free link has private
slots, so its value is forced by any one of them. This proves
uniqueness.

For the converse, take a full flat connection. Private
plaquettes give $p_{i\alpha}=U_i$ and internal plaquettes give
$U_i=R_{ij}U_j$. Around an elementary square in the internal
rectangle, the composed transport is
$R(C)z=L_C z R_C^{-1}$, where $L_C,R_C$ are the two untouched
endpoint plaquette holonomies. They both equal $1$.
The internal rectangle is contractible, so one can choose
$O_i\in\operatorname{SO}(4)$ with
$R_{ij}=O_iO_j^{\mathsf T}$ and $O_i1=U_i$:
choose a root frame and propagate along paths, whose
independence follows from the square identities.
The change of variables $x_i=O_i z_i$ in \eqref{r29:compact}
then makes all slots equal to $1$ and all internal transports
equal to $I$. Thus $Z_b(\eta)=Z_b(1)$ for every block, and the
untouched actions vanish. This proves equality in
\eqref{r29:min-bound}.

Finally each private retained path transforms between the
same two endpoints as its free link. The unique extension is
therefore gauge equivariant. Its formula by a selected private
path is smooth on the minimum set (and is the restriction of
a smooth map on the ambient retained space).
\end{proof}

Selecting isolated diagrams here is used only for an exact
equality characterization. No useful quantitative bound on
their small probabilities is claimed. The quantitative
coercivity proof instead used the insertion intensities,
which retain all background activity.

This theorem strictly strengthens the flat-family result in
V28. Nontrivial commuting winding holonomies are allowed;
flat does not mean globally pure gauge on the torus.
The condition $\gamma>0$ matters: at $\gamma=0$ the entire
retained action is constant.

\subsection{Why the identity minimum is not Morse--Bott}
\label{r29:singular}

\begin{proposition}[A nonlinear obstruction inside the native kernel]
\label{r29:not-MB}
The action $V_{\mathcal R}$ is not Morse--Bott at the identity,
either before quotienting or on the smooth based gauge quotient.
In particular its twelve quotient Hessian zero modes cannot be
identified with the tangent space of a smooth minimum manifold.
\end{proposition}

\begin{proof}
Choose two different coordinate directions, say $\mu,\nu$,
and the full constant cochain
$h_\mu=\mathbf i$, $h_\nu=\mathbf j$, with the other two
components zero. Its restriction $a$ belongs to the kernel in
Theorem~\ref{r29:kernel-thm}.
Suppose a twice differentiable curve of minima through the
identity had this tangent, up to an infinitesimal gauge vector.
Theorem~\ref{r29:minima-thm} gives a smooth flat extension along
that curve. Let $W_\mu(s),W_\nu(s)$ be its two coordinate
winding holonomies based at one fixed vertex.
Flat plaquettes allow elementary homotopies of lattice loops;
the two torus generators commute, so
$W_\mu(s)W_\nu(s)=W_\nu(s)W_\mu(s)$ for every $s$.
If their lengths are $\ell_\mu,\ell_\nu$, their derivatives are
\[
 W_\mu'(0)=\ell_\mu\mathbf i,\qquad
 W_\nu'(0)=\ell_\nu\mathbf j.
\]
An infinitesimal gauge gradient telescopes around each loop
and does not change these derivatives at the identity.
But Taylor expansion of the commutation identity at order
$s^2$ requires
\[
 0=[W_\mu'(0),W_\nu'(0)]
     =\ell_\mu\ell_\nu[\mathbf i,\mathbf j]
     =2\ell_\mu\ell_\nu\mathbf k,
 \tag{V29.22}\label{r29:commutator}
\]
a contradiction. The unknown second derivatives of the
holonomies cancel from this identity.

For a Morse--Bott critical point, the critical set near it is
a smooth submanifold whose tangent space equals the Hessian
kernel. Every vector in that tangent space is the derivative
of a smooth curve in the submanifold. Along such a curve the
action is constant, because its differential vanishes.
The component through the identity would consequently consist
of global minima, by \eqref{r29:min-bound}. The preceding
kernel vector supplies a contradiction.

On the based quotient the same reasoning applies. The action
descends smoothly because the based action is free and proper.
A curve there has a local smooth lift, and its tangent differs
from a representative by an infinitesimal gauge vector.
The winding argument is unchanged. Thus passing to this
smooth quotient does not repair the non-Morse--Bott minimum.
\end{proof}

This is the familiar nonlinear commuting-holonomy obstruction,
now established for the \emph{exact packed marginal}.
For background on the space of commuting tuples and the distinction
between $\operatorname{Hom}(\mathbb Z^n,G)$ and its conjugation
quotient, see Adem--Cohen,
\href{https://arxiv.org/html/math/0603197v2}
{Commuting Elements and Spaces of Homomorphisms}, Introduction
and Section 2. Their topological computations are not imported
as analytic coercivity estimates; the obstruction just used
was proved directly.

The proposition does not say that each harmonic direction is
spurious. Any single direction, and any mutually commuting
tuple of directions, does integrate to flat connections.
It is their noncommuting combinations which obstruct a smooth
minimum manifold with the entire linear kernel as tangent.
Nor do the twelve linear modes count physical particles.

\subsection{Exact checks, preserved history, and the next import}
\label{r29:checks}

The checker
\begin{center}\small
\path{scripts/verify_ym_restart_v29_native_maxwell.py}
\end{center}
constructs the original oriented plaquettes, their block ownership,
the literal private path rows and internal transport rows.
It checks the variance and extension identities using rational
arithmetic on two independent deterministic fixtures per block,
at four rational couplings. More importantly, exact row boundaries
and directional sums certify that \emph{all} gradients and all
constant one-forms are annihilated, and that the extension preserves
every closed cochain. Connectivity and four wholly retained winding
cycles are checked directly.

\begin{center}
\begin{tabular}{@{}rrrrr@{}}
\toprule
$(L,q)$ & $|\Lambda|$ & $|\mathcal R|$ & blocks
                                      & scalar kernel dimension\\
\midrule
$(4,3)$ & 384  & 1488  & 8  & 387\\
$(4,4)$ & 512  & 2000  & 8  & 515\\
$(4,6)$ & 768  & 2976  & 16 & 771\\
$(8,3)$ & 3072 & 11904 & 64 & 3075\\
\bottomrule
\end{tabular}
\end{center}
There are $192$ rational block fixtures in these four geometries.
For $(L,q)=(4,3)$, the $3024$ by $1488$ integer matrix consisting
of untouched curls, same-site slot differences and one-edge
mismatches has rank $1101$ modulo the prime $1000003$.
This is a rigorous lower bound on its rational rank.
The independently verified $387$-dimensional closed-cochain
kernel gives the matching upper bound, so its rational rank
is exactly $1101$. This is an incidence certificate, not a
numerical diagonalization of the compact integral Hessian.
The proofs above apply to all stated lattice sizes.

Removing this append and restoring V28's title reconstructs its
entire 659044-byte source with SHA--256
\begin{center}\small
\path{9979B463EDCAFE8DEBD1E06B97129F8B2645ABED35200597F46CCB0940B3D394}.
\end{center}
All 46 other manuscript sources remain unchanged.
Only the latest active \texttt{.tex} is retained in its active
lineage; builds and checks are outside \path{latex_notes}.

\paragraph{Direction after this positive step.}
The local positive coefficient is now explicit, and the native
quadratic kernel is complete. The next question is a
\emph{nonlinear, measure-accounted} estimate near the exact flat
minimum set, including its singular winding configurations.
A regular Morse--Bott Gaussian theorem cannot simply be applied
there. The proof must retain the shared exact marginal and the
orbit-volume or equivalent reference-measure correction recorded
in \eqref{r28:orbit-volume}; this append has not computed that
correction.

V30 is the next scheduled companion review and broad
primary-literature sweep. The search should test methods for
degenerate Gibbs minima, commuting-holonomy singularities and
nonlinear coercivity against the already identified physical
endpoint-affinity bottleneck. A result for a fixed finite
quartic model, an auxiliary sampler or a regular stratum is
not automatically a bound on the native global law.
The V28 saddle escape and the present positive transverse
comparison are complementary finite-regulator inputs.
No all-saddle estimate, global gradient control, normalized
weak-coupling contraction, cofinal continuum construction or
positive physical Yang--Mills mass gap has been proved here.

% END CONSOLIDATED V29 APPEND

% BEGIN CONSOLIDATED V30 APPEND -- 2026-09-19
\clearpage
\section{V30: global nonlinear filling and the exact return measure}
\label{r30:start}

\subsection{Context, nonduplication, and the scheduled reviews}

V29 established the exact flat minimum and a native Maxwell lower
form at the identity. Its noncommuting harmonic directions also
showed why a regular Morse--Bott argument does not apply there.
The present append proves a \emph{global nonlinear} comparison,
valid for every retained configuration, rather than extending an
identity Hessian estimate by assertion. It constructs full original
links from retained paths and pays every original plaquette with
a volume-independent coefficient. It then restores the integrated
fluctuations exactly, so that an action comparison is not mistaken
for a comparison of normalized measures.

The all-exterior Ward bound of archive $C$, V62 is an inherited
input, rechecked below in the present normalization. The new parts
are the native filling map, its ordered cube identity, the global
face count and the resulting full-curvature comparison.
The radial comparison in $A$, V69 and the four-free-link cell in
$C$, V63 are different constructions. The latter has a cyclic
four-spin interaction and does not inherit the ferromagnetic
pair-graph argument used here. The generic constant-field quartic
phenomenon already appears in $C$, V44; only its explicit
realization in this exact packed marginal is added below.

\paragraph{Companion checkpoint at V30.}
Fresh hashes and terminal arguments were checked for NS V26,
SM--Einstein V72, shell-mixing V16 and parallel YM V5. All four
are unchanged since V25. NS's rank-one derivative reductions
suggest optimizing how shared terms are paid, but their constants
do not transfer to Yang--Mills. SM's copied-system limit is at
fixed momentum modes and supplies no uniform ultraviolet estimate.
The shell-mixing flat-holonomy calculation does not bound the
noncommuting curvature or a common nonperturbative remainder.
The parallel YM warning against global extensive-charge tightness
supports a local, normalized-law target. No companion result is
being treated as the missing physical contraction estimate.
The next companion checkpoint is V35.

\paragraph{Broad primary-literature checkpoint at V30.}
The user's sweep was reconsidered alongside searches in flat
connection geometry, degenerate Gibbs analysis, plaquette
identities, loop and spectral bootstraps, Gibbs factorization,
spectral renormalization and stochastic quantization.
The following is a hypothesis audit, not a claim to have
independently verified every proof in each paper. The first
three entries include the specific full-text passages indicated;
the other entries are abstract/status checks unless otherwise
stated. No new external spectral-gap theorem is invoked.

\begin{longtable}{@{}>{\raggedright\arraybackslash}p{0.30\textwidth}
                    >{\raggedright\arraybackslash}p{0.65\textwidth}@{}}
\toprule
Primary source and scope read & Import decision and outstanding input\\
\midrule
\endhead
\href{https://arxiv.org/html/1906.03954}{Feehan,
\emph{Morse theory for the Yang--Mills energy function near flat
connections}}, Introduction, Theorem 10 and surrounding remarks.
&
The corrected distance estimate can have a fractional curvature
exponent; a linear estimate requires additional structure such
as Morse--Bott regularity. The torus product-connection example
is particularly relevant to V29. Do not import a uniform
lattice exponent or constant. A classical stationary-energy
gap is not a quantum Hamiltonian gap.\\[4pt]
\href{https://arxiv.org/html/2410.21899}{Delande,
\emph{Sharp spectral gap for some degenerate Witten Laplacians}},
Assumptions 1--2 and Remark 1.3.
&
The local power normal form still has isolated critical points.
Our continuous, singular flat family does not meet that
hypothesis. Nonquadratic local analysis remains useful; the
stated theorem cannot simply be applied to this marginal.\\[4pt]
\href{https://arxiv.org/html/hep-lat/0508003}{Borisenko--Voloshin--Faber,
\emph{Plaquette representation for 3D lattice gauge models}},
Section 2.1, (14)--(17), and the Section 3 connector discussion.
&
Ordered non-Abelian cube constraints retain conjugating
connectors. We use this bookkeeping principle, proving a
short, literal four-dimensional cube identity below. No
three-dimensional perturbative or confinement conclusion
is imported.\\[4pt]
\href{https://arxiv.org/abs/2605.20509}{Abbott et al.,
\emph{The Causal Bootstrap}}; abstract, 2026.
&
Smeared spectral bounds from finite Euclidean data are a
potential certificate-discovery tool. They do not supply our
unproved uniform, all-depth native endpoint estimate; the
existing moment route remains conditional on that input.\\[4pt]
\href{https://arxiv.org/abs/2404.16925}{Kazakov--Zheng,
\emph{Bootstrap for Finite $N$ Lattice Yang--Mills Theory}};
abstract, revised 2024.
&
Finite-$N$ and specifically $\SU(2)$ loop constraints fit the
model better than a planar replacement. Finite constraint
bounds or string-tension estimates are not an all-scale
continuum mass-gap proof. Retain for exact local certificates.\\[4pt]
\href{https://arxiv.org/abs/2505.16585}{Cao--Nissim--Sheffield,
\emph{Expanded regimes of area law for lattice Yang--Mills
theories}}; abstract, revised 2025.
&
Truncated loop-equation control and restoration of omitted
terms are useful methodological leads. Their large-$N$
$U(N)$ regimes are not the present fixed-$\SU(2)$
weak-coupling return-measure estimate.\\[4pt]
\href{https://arxiv.org/abs/2208.11299}{Qin--Wang,
\emph{Spectral Telescope}}; abstract.
&
Conditional correlations and influence bounds can organize
a Gibbs-sampler gap. Such bounds for this actual joint law
are still required, as is a separate physical-time transfer
argument. An auxiliary sampler gap is not the target gap.\\[4pt]
\href{https://arxiv.org/abs/2508.07643}{Bach--Ballesteros--Geisler,
\emph{Spectral Renormalization Flow Based on Smooth
Feshbach--Schur}}; abstract, with the V20 hypothesis audit.
&
Retain spectral-parameter tracking if a physical operator
reduction becomes available. An iterated classical action
bound does not provide the domain, inverse and operator-norm
controls required for that reduction.\\[4pt]
\href{https://arxiv.org/abs/2503.03060}{Chevyrev--Shen,
\emph{Uniqueness of gauge covariant renormalisation of
stochastic 3D Yang--Mills--Higgs}}; abstract/status, revised 2026.
&
Gauge-covariant renormalization is relevant context, but
three-dimensional local stochastic dynamics and short-time
observables do not construct the interacting four-dimensional
continuum vacuum or its positive physical mass gap.\\
\bottomrule
\end{longtable}

The selected import is therefore \emph{exact noncommutative
curvature bookkeeping}, with singular-minimum and measure
hypotheses kept visible. The broad sweep is scheduled again
at V40, with the count continuing across any archive restart.

\subsection{The inherited global bound in the exact marginal}
\label{r30:ward}

Use precisely the lattice, packing and action of
\eqref{r29:lattice}--\eqref{r29:compact}. Let $B$ be the number
of six-link blocks and set
\[
 \begin{gathered}
 e(U)=1-w(U),\qquad E_W(U)=\sum_{p\in\mathcal P} e(U_p),\\
 W_\gamma(\eta)=V_{\mathcal R}(\eta)-V_{\mathcal R}(1),\qquad
 S_0(\eta)=\sum_{p\in\mathcal P_0}e(U_p(\eta)),\\
 S_{\rm slot}(\eta)=
 \sum_{i\in\mathcal F}\sum_{\alpha<\beta}
                  (1-p_{i\alpha}\cdot p_{i\beta}),\qquad
 A_\gamma=\frac{\gamma^2}{4+6\gamma}.
 \end{gathered}
 \tag{V30.1}\label{r30:definitions}
\]
All energies are dimensionless. V29 identified $V_{\mathcal R}(1)$
as the global minimum. All measures in this append are the
finite-regulator Wilson law or its exact marginal, with
normalized product Haar as reference.

\begin{lemma}[Inherited all-exterior Ward comparison]
\label{r30:inherited}
For every $\gamma>0$ and every retained configuration,
\[
 W_\gamma(\eta)\ge
       \gamma S_0(\eta)+A_\gamma S_{\rm slot}(\eta).
 \tag{V30.2}\label{r30:ward-bound}
\]
This includes zero aggregate staples and arbitrary transport
holonomies within each block.
\end{lemma}
\begin{proof}
This is $C$, V62's global graph bound, not a new theorem.
Here is its application with the constants recorded. Put
$r_i=|\sum_\alpha p_{i\alpha}|\le n_i$ in each block.
The positive Haar-pairing expansion bounds its partition
function by that with fields $\gamma r_i e_0$ and all pair
transports equal to the identity. Indeed each open strand
has endpoint contraction at most the product of endpoint
norms, and each closed strand has orthogonal-product trace
at most $4$; all scalar pairing weights are nonnegative.

For aligned fields $\gamma u_i e_0$, $0\le u_i\le n_i$,
let $m_i=\langle x_i^0\rangle$ and
$c_{ij}=\langle x_i^a x_j^a\rangle$ for a transverse coordinate.
Conditioning on absolute transverse coordinates gives a
zero-field ferromagnetic Ising sign law, so $c_{ij}\ge0$.
The one-spin raw second-moment bound from $C$, V62 gives
$c_{ii}\ge(4+6\gamma)^{-1}$, since every conditional field
has norm at most $\gamma(n_i+\deg i)=6\gamma$.
This uses a raw moment, not a longitudinal variance.
A simultaneous rotation in the $(0,a)$ plane gives the
exact Ward identity
\[
 m_i=\gamma\sum_j u_jc_{ij}
             \ge\frac{\gamma u_i}{4+6\gamma}.
\]
Interpolate $u_i$ from $r_i$ to $n_i$. The derivative of
$\log Z$ is $\gamma\sum_i(n_i-r_i)m_i$, and integration
bounds $\log Z(n,I)-\log Z(r,I)$ below by
\[
 \frac{\gamma^2}{2(4+6\gamma)}
        \sum_i(n_i^2-r_i^2)
   = A_\gamma\sum_i\sum_{\alpha<\beta}
                      (1-p_{i\alpha}\cdot p_{i\beta}).
\]
The original block factors are conditional on the same
retained configuration. Summing their logarithms and adding
the untouched plaquette action proves \eqref{r30:ward-bound}.
It does not factorize the retained probability law.
\end{proof}

\subsection{A global retained-path filling}
\label{r30:filling}

For each free $\mu$ link $i$ at $x$, choose its forward-time
private staple
\[
 s_i(\eta)=U_t(x)\,U_\mu(x+\widehat t)\,
                         U_t(x+\widehat\mu)^{-1}.
 \tag{V30.3}\label{r30:staple}
\]
All three factors are retained: free links occur only in the
$\mu$ direction at even time, and $x+\widehat t$ is at odd
time. Define a full-link configuration by
\[
 \Phi(\eta)_e=
 \begin{cases}
 \eta_e,&e\in\mathcal R,\\
 s_e(\eta),&e\in\mathcal F .
 \end{cases}
 \tag{V30.4}\label{r30:Phi}
\]

\begin{lemma}[Global section and private energy]
\label{r30:section}
The map $\Phi$ is smooth, gauge-equivariant and injective;
restriction to $\mathcal R$ is its left inverse. Every selected
forward $\mu t$ plaquette is flat. If $E_{\rm priv}$ denotes
the sum of all filled private-plaquette energies, then
\[
 E_{\rm priv}(\Phi\eta)
   =\sum_i\sum_\alpha(1-s_i\cdot p_{i\alpha})
   \le S_{\rm slot}(\eta).
 \tag{V30.5}\label{r30:private}
\]
\end{lemma}
\begin{proof}
Multiplication and inversion in $\SU(2)$ are smooth. Under
a vertex gauge transformation, the internal factors of the
path cancel, leaving
$s_i\mapsto g_xs_i g_{x+\widehat\mu}^{-1}$, the transformation
of the missing link. Restriction proves injectivity. Substituting
\eqref{r30:staple} into its selected plaquette makes that
holonomy the identity. A private plaquette has energy
$1-w(s_i p_{i\alpha}^{-1})=1-s_i\cdot p_{i\alpha}$.
Since $s_i$ is one of the slots, the nonzero terms in this
sum are a subset of the nonnegative unordered slot-pair
terms in \eqref{r30:definitions}. Sum over free links.
\end{proof}

There is no normalized staple sum, exceptional denominator,
local gauge choice or small-field assumption. This is a
deterministic section, not a conditional expectation or a
physical renormalization map. Its one-sided time choice
does not establish reflection compatibility.

\subsection{The ordered cube identity and its global face budget}
\label{r30:cube}

Consider an internal $\mu\alpha$ plaquette, where
$\alpha\in\{\nu,\rho\}$, based at $x$. Write its bottom vertices
as $A=x$, $B=x+\widehat\mu$, $D=x+\widehat\alpha$,
$C=x+\widehat\mu+\widehat\alpha$, and their vertical links
as $v_A,v_B,v_D,v_C$. At the upper, odd time layer put
\[
 a=U_\mu(x+\widehat t),\quad
 b=U_\mu(x+\widehat\alpha+\widehat t),\quad
 c=U_\alpha(x+\widehat t),\quad
 d=U_\alpha(x+\widehat\mu+\widehat t).
\]
At the lower layer put $c_0=U_\alpha(x)$ and
$d_0=U_\alpha(x+\widehat\mu)$. The filled bottom $\mu$ links
are $v_Aav_B^{-1}$ and $v_Db v_C^{-1}$. Define the top and
two endpoint side holonomies by
\[
 \begin{split}
 P&=a d b^{-1}c^{-1},\\
 S_L&=c_0v_Dc^{-1}v_A^{-1},\\
 S_R&=d_0v_Cd^{-1}v_B^{-1}.
 \end{split}
 \tag{V30.6}\label{r30:faces}
\]
All three faces belong to $\mathcal P_0$.

\begin{lemma}[Short non-Abelian cube identity]
\label{r30:cube-lemma}
If $W$ is the filled bottom plaquette holonomy, then
\[
 v_A^{-1}Wv_A
  =(a v_B^{-1}S_Rv_Ba^{-1})\,
             P\,(v_A^{-1}S_L^{-1}v_A).
 \tag{V30.7}\label{r30:cube-identity}
\]
Consequently,
\[
 e(W)\le 5e(P)+\frac52\{e(S_L)+e(S_R)\}.
 \tag{V30.8}\label{r30:weighted-cube}
\]
\end{lemma}
\begin{proof}
The left side of \eqref{r30:cube-identity} is
\[
 a v_B^{-1}d_0v_C b^{-1}v_D^{-1}c_0^{-1}v_A.
\]
Substitute the three expressions in \eqref{r30:faces} into
the right side and cancel adjacent inverse factors. The
same ordered word results; no factors have been commuted.

For unit quaternions,
$\delta(U)=\sqrt{e(U)}=|1-U|/\sqrt2$ is the chordal distance
to the identity. It is conjugation invariant, unchanged by
inversion, and satisfies $\delta(UV)\le\delta(U)+\delta(V)$
by the Euclidean triangle inequality and invariance of norm
under multiplication by a unit quaternion. Thus
\[
 \sqrt{e(W)}\le \sqrt{e(P)}+\sqrt{e(S_L)}+\sqrt{e(S_R)}.
\]
Weighted Cauchy--Schwarz with weights $1,2,2$ proves
\eqref{r30:weighted-cube}. The unequal weights will exactly
compensate for the shared side faces.
\end{proof}

\begin{lemma}[Global incidence count]
\label{r30:incidence}
Over all internal plaquettes in the packing, each used top
face occurs once and each used side face occurs twice.
The two sets are disjoint and each contains $7B$ distinct
untouched faces. Let $\Gamma(\eta)$ be the sum of the
energies over their union. Then
\[
 0\le\Gamma\le S_0,\qquad
 E_{\rm int}(\Phi\eta)\le5\Gamma,\qquad
 E_W(\Phi\eta)\le S_0+S_{\rm slot}+5\Gamma
                                      \le6S_0+S_{\rm slot}.
 \tag{V30.9}\label{r30:budget}
\]
\end{lemma}
\begin{proof}
The seven internal grid edges in a block consist of three
$\nu$ edges and four $\rho$ edges. Moving an internal
$\mu\alpha$ face up by one time step is injective on the
entire packing, so its top face is used once. It is untouched
because its $\mu$ links are at odd time.

Fix an internal transverse grid edge and its even time.
The associated $\alpha t$ side face at a specified
$\mu$ coordinate is used by the two adjacent $\mu$ intervals:
the one starting there and the one ending there. Both
intervals occur because the packing includes every
$x_\mu$. No other internal edge yields that same face.
The transverse anchors are separated by four and the
rectangles have widths two and three, so distinct transverse
rectangles cannot add a third use. Distinct even-time
layers also give distinct side faces. Periodic closure in
the $\mu$ direction preserves exactly these two uses.
Every side face is untouched since neither of its directions
is $\mu$. The top and side sets have different direction
pairs and are disjoint.

There are $7B$ top-face occurrences and $14B$ side-face
occurrences, hence $7B$ distinct faces in each set.
Summing \eqref{r30:weighted-cube} gives coefficient $5$
for every used face: $5$ for a top and $2(5/2)$ for a
side. This proves the internal bound. The original
plaquettes partition into untouched, private and internal
ones. Apply \eqref{r30:private} to their sum to obtain
\eqref{r30:budget}.
\end{proof}

\subsection{Global coercivity for the filled original curvature}
\label{r30:coercivity}

\begin{theorem}[All-configuration nonlinear comparison]
\label{r30:main}
For every lattice in \eqref{r29:lattice}, every $\gamma>0$
and every retained configuration, the exact marginal obeys
\[
 \boxed{\quad
 W_\gamma(\eta)\ge
   \frac{\gamma^2}{4+6\gamma}\,E_W(\Phi\eta)
       +\frac{4\gamma}{4+6\gamma}\,S_0(\eta).
 \quad}
 \tag{V30.10}\label{r30:main-bound}
\]
More precisely, before replacing $\Gamma$ by $S_0$,
\[
 W_\gamma(\eta)\ge
 A_\gamma E_W(\Phi\eta)
             +(\gamma-A_\gamma)S_0(\eta)-5A_\gamma\Gamma(\eta).
 \tag{V30.11}\label{r30:refined-bound}
\]
The coefficients are independent of the volume. This is
a finite-regulator action inequality, not a spectral-gap
or normalized-measure theorem.
\end{theorem}
\begin{proof}
Multiply the first filled-energy bound in
\eqref{r30:budget} by $A_\gamma$ and rearrange:
\[
 \gamma S_0+A_\gamma S_{\rm slot}
 \ge A_\gamma E_W(\Phi\eta)
                +(\gamma-A_\gamma)S_0-5A_\gamma\Gamma.
\]
Use \eqref{r30:ward-bound} to get
\eqref{r30:refined-bound}. Since $\Gamma\le S_0$ and
$\gamma-6A_\gamma=4\gamma/(4+6\gamma)>0$,
\eqref{r30:main-bound} follows.
\end{proof}

At $\gamma=10$ the filled-energy coefficient is $25/16$.
As $\gamma\to\infty$, $A_\gamma/\gamma\to1/6$.
Thus there is no exponential empty-block isolation loss
in this estimate. The optimized face payment improves
the provisional unweighted cost seven to cost six; no
optimality among all fillings or inequalities is asserted.
Unlike V29's Hessian form, this inequality concerns every
retained configuration and the full nonlinear original
plaquette energy of its filled extension.

\begin{corollary}[An exact quartic native benchmark]
\label{r30:quartic}
On retained links put $U_\mu=e^{s\mathbf i}$,
$U_\nu=e^{s\mathbf j}$ and $U_\rho=U_t=1$, and denote this
retained configuration by $\eta(s)$. At fixed lattice
and fixed $\gamma>0$,
\[
 \begin{gathered}
 E_W(\Phi\eta(s))=2|\Lambda|\sin^4s,\\
 W_\gamma(\eta(s))
   \ge2\gamma(|\Lambda|-9B)\sin^4s,\qquad
 |\Lambda|-9B\ge\frac{23}{32}|\Lambda| .
 \end{gathered}
 \tag{V30.12}\label{r30:quartic-bound}
\]
Moreover $W_\gamma(\eta(s))=\Theta(s^4)$ as $s\to0$,
with constants allowed to depend on the regulator and
$\gamma$.
\end{corollary}
\begin{proof}
The filling restores the same constant $\mu$ links.
Quaternion multiplication gives
\[
 w(e^{s\mathbf i}e^{s\mathbf j}
                   e^{-s\mathbf i}e^{-s\mathbf j})
       =1-2\sin^4s.
\]
Only the $\mu\nu$ planes have nonzero plaquette energy.
Each block contains three internal and six private
$\mu\nu$ plaquettes, and the incident plaquette sets of
different blocks are disjoint. Hence precisely
$|\Lambda|-9B$ such plaquettes are untouched. Equation
\eqref{r30:ward-bound}, with its slot term discarded, gives
the lower bound. The packing has
$B\le L(L/4)^2q=|\Lambda|/32$, proving the last estimate.

At fixed regulator the compact integral and its logarithm
are real analytic along this curve. The tangent is the
restriction of a constant full-lattice cochain, so its
native Hessian vanishes by V29. The gradient vanishes at
the minimum. Conjugation at every vertex by $\mathbf k$
sends $\eta(s)$ to $\eta(-s)$, so the action is even.
Its Taylor expansion therefore starts at order at least
four. This proves the $O(s^4)$ upper bound; the displayed
positive lower bound proves the matching order.
\end{proof}

This calculation includes an actual noncommuting direction,
rather than assuming the flat set is smooth. It gives
neither a cutoff-uniform quartic normal form nor a physical
quartic Hamiltonian. The general quartic obstruction was
already known in the earlier notes; the native lower
constant and exact untouched-face count are the additions.

\subsection{Restoring the measure: the exact return integral}
\label{r30:return}

Let $\mu_W$ be the normalized full Wilson law and
$\mu_{\mathcal R}$ its exact retained marginal.

\begin{proposition}[The deterministic filling is not the full law]
\label{r30:singular}
If $B>0$ and $\gamma<\infty$, then
$\Phi_*\mu_{\mathcal R}$ and $\mu_W$ are mutually singular.
\end{proposition}
\begin{proof}
The image of $\Phi$ lies in the event that every chosen
forward private plaquette is exactly the identity. One
such plaquette suffices: conditional on all links except
its free $\mu$ link, that event prescribes a single group
element. Haar measure has no atoms. The full Wilson law
has a positive continuous density with respect to finite
product Haar, so the event has $\mu_W$-measure zero.
It has $\Phi_*\mu_{\mathcal R}$-measure one.
\end{proof}

An energy certificate must therefore retain the discarded
fluctuations if it is to be used for probabilities. Introduce
$z=(z_i)_{i\in\mathcal F}\in\SU(2)^{\mathcal F}$ and put
\[
 F(\eta,z)_e=
 \begin{cases}
 \eta_e,&e\in\mathcal R,\\
 s_e(\eta)z_e,&e\in\mathcal F.
 \end{cases}
 \tag{V30.13}\label{r30:F}
\]
Write $E_\Phi(\eta)=E_W(\Phi\eta)$ and define
\[
 \begin{split}
 J_\gamma(\eta)
   &=\int\exp\{-\gamma[
             E_W(F(\eta,z))-E_\Phi(\eta)]\}\,dz,\\
 R_\gamma(\eta)&=\frac{J_\gamma(\eta)}{J_\gamma(1)},\qquad
 d\nu_\gamma(\eta)
   =\frac{e^{-\gamma E_\Phi(\eta)}\,d\eta}
          {\int e^{-\gamma E_\Phi(\xi)}\,d\xi}.
 \end{split}
 \tag{V30.14}\label{r30:J}
\]
Both integrals use normalized product Haar.

\begin{theorem}[Exact return measure, including its normalization]
\label{r30:return-theorem}
The map $F$ is a smooth bijection of the product configuration
spaces and preserves product Haar in the sense $dU=d\eta\,dz$.
The functions $J_\gamma$ and $R_\gamma$ are strictly positive,
and
\[
 \boxed{\quad
 W_\gamma=\gamma E_\Phi-\log R_\gamma,\qquad
 \frac{d\mu_{\mathcal R}}{d\nu_\gamma}
       =\frac{R_\gamma}{\nu_\gamma[R_\gamma]}.
 \quad}
 \tag{V30.15}\label{r30:exact-return}
\]
The integral $J_\gamma$ factorizes into one factor per
free block conditional on $\eta$. This is not a product
decomposition of $\mu_{\mathcal R}$.
\end{theorem}
\begin{proof}
The inverse is the retained restriction of $U$ together
with $z_i=s_i(\eta)^{-1}U_i$. For each fixed $\eta$, left
translation preserves Haar on every free link; Fubini
therefore proves the asserted reference-measure identity,
without a coordinate Jacobian. Gauge transformation
conjugates $z_i$ at the right endpoint of its free link,
which also preserves Haar.

The unnormalized retained density is exactly
\[
 Q_\gamma(\eta)
    :=\int e^{-\gamma E_W(F(\eta,z))}\,dz
     =e^{-\gamma E_\Phi(\eta)}J_\gamma(\eta).
\]
Compactness makes this finite and positive. Since
$E_\Phi(1)=0$,
$-\log Q_\gamma(\eta)+\log Q_\gamma(1)$ is precisely
the relative action $W_\gamma(\eta)$, proving the first
identity. Dividing $Q_\gamma$ by its retained integral
and comparing to the definition of $\nu_\gamma$ proves
the second, including its denominator.

Every original plaquette has free links in at most one
block. Its energy difference between $F(\eta,z)$ and
$\Phi\eta$ is consequently either zero or a function of
that block's $z_b$ and the fixed common exterior $\eta$.
Thus the exponent is a sum of block functions and
$J_\gamma=\prod_bJ_{\gamma,b}$. The block functions depend
on overlapping retained links. Integrating over $\eta$
does not make these factors independent.
\end{proof}

On every restricted full flat configuration V29 gives
$W_\gamma=0$, and the present filling is flat there, so
$E_\Phi=0$ and $R_\gamma=1$. This statement concerns this
particular exact marginal and these coordinates; it is
not a replacement for a different blocking scheme's
flat-holonomy effective potential.

Combining \eqref{r30:main-bound} and
\eqref{r30:exact-return} gives the rigorous one-sided estimate
\[
 \log R_\gamma(\eta)
   \le(\gamma-A_\gamma)E_\Phi(\eta)
               -\frac{4\gamma}{4+6\gamma}S_0(\eta).
 \tag{V30.16}\label{r30:return-upper}
\]
It gives no corresponding small oscillation or useful
uniform lower bound for $R_\gamma$. In particular,
$(\gamma-A_\gamma)/\gamma\to5/6$, not zero. One cannot
erase $R_\gamma$, its normalizer, or its correlations,
nor conclude that the two normalized laws are close.
This is the exact location of the remaining return-measure
problem after the positive global action estimate.

\subsection{Certificates, preservation, and the next upstream task}
\label{r30:checks}

The checker
\begin{center}\small
\path{scripts/verify_ym_restart_v30_nonlinear_filling.py}
\end{center}
constructs original oriented plaquettes and all retained
paths. It compares the two sides of each cube identity
as reduced words in noncommuting link symbols, before
assigning any numerical link values.

\begin{center}
\begin{tabular}{@{}rrrrrr@{}}
\toprule
$(L,q)$ & $|\Lambda|$ & $|\mathcal R|$ & $B$
                           & free links & internal cubes\\
\midrule
$(4,3)$ & 384  & 1488  & 8  & 48  & 56\\
$(4,4)$ & 512  & 2000  & 8  & 48  & 56\\
$(4,6)$ & 768  & 2976  & 16 & 96  & 112\\
$(8,3)$ & 3072 & 11904 & 64 & 384 & 448\\
$(8,4)$ & 4096 & 16000 & 64 & 384 & 448\\
\bottomrule
\end{tabular}
\end{center}

These are $1120$ exact cube identities and $960$ selected
private paths. Top multiplicity one, side multiplicity
two and disjointness are checked on all five geometries.
Four complete rational quaternion configurations on the
two smaller geometries verify the private-pair comparison,
every weighted cube bound and their full sums, including
noncentral gauge covariance and the inverse fluctuation
coordinates. Five rational couplings check the final
coefficient arithmetic. Exact constant noncommuting
backgrounds also verify the quartic plaquette count.
The checker does not numerically evaluate the compact
Ward integral, normalized-law mixing or a physical gap.
The proofs above, rather than these finite fixtures,
supply the all-size statements.

Removing this append and restoring V29's title reconstructs
its full 686681-byte source with SHA--256
\begin{center}\small
\path{3E19488D46FFF85A7599D42A65238028BAC70B65366F3A2A271F9CBA20A247D5}.
\end{center}
The 46 other manuscript sources remain unchanged.
Only the active predecessor is replaced; archived and
parallel sources are preserved. Builds and certificates
remain outside \path{latex_notes}, which contains only
\texttt{.tex} files.

\paragraph{Direction after the literature and nonlinear checks.}
The next target is not another unnormalized lower action
bound or a regular-minimum approximation. It is the local
interaction $\log J_{\gamma,b}$ in the exact common
exterior law: its conditional derivatives, cumulants or
overlap estimates, with the return normalizer paid.
Local loop constraints and finite-block analytic bounds
are candidates for this step; their hypotheses must be
proved for this same joint law. Any exceptional-set
argument must include its probability under that law.

This append does not evaluate the based-quotient metric
and orbit-volume correction, nor prove global normalized
weak-coupling contraction. Uniform physical-time transport,
cofinal volume/cutoff/coupling control, the interacting
four-dimensional continuum construction and a positive
physical Hamiltonian mass gap remain open.
The new theorem removes the purely nonlinear
\emph{action-comparison} obstruction for this packing;
it does not remove those measure and continuum obligations.

% END CONSOLIDATED V30 APPEND

% BEGIN CONSOLIDATED V31 APPEND -- 2026-09-19
\clearpage
\section[V31: native local minima and a normalized return approximation]
{V31: native local minima\\and a normalized return approximation}
\label{r31:start}

\subsection{Context and the nonduplication audit}

V30 pays the nonlinear curvature of a filled configuration but
leaves its exact return measure. This append connects a native
retained-frame defect, a $45$-plaquette probability witness,
and the integral around the six-link block's \emph{moving}
conditional minimum. Both its relaxed energy and fluctuation
determinant are retained. Replacing $M$ selected block factors
then changes the normalized retained law by $O(M/\gamma)$,
plus a paid bad-event term, uniformly in surrounding volume
for fixed $M$. This is not an all-block or physical-gap result.

V5 already proves parameter-dependent conditional Laplace
control in a declared fixed flat-boundary box. Archive $C$,
V63--V64 already treats pinned one- and two-layer limits
with their conditional normalizers. $C$, V65 and V76 already
prove the periodic Wilson chessboard input and a
volume-uniform marked-plaquette tail. V21 already proves
normalized one-sided filter stability. These are credited
methods, not new discoveries. The new work is their native
geometric connection, the explicit conditional-minimum
neighborhood, determinant variation, and normalized local
replacement under the unrestricted periodic law.

We rechecked Biskup,
\href{https://arxiv.org/html/math-ph/0610025v2}
{\emph{Reflection Positivity and Phase Transitions in Lattice
Spin Models}}, Section 5.3, event inequality (5.35).
The HTML labels that chessboard theorem 5.2, whereas earlier
notes cite 5.8; the formula and block hypotheses identify
the result without that numbering ambiguity.
No filtered or arbitrarily conditioned law is assumed
reflection positive. The next companion and broad-literature
checkpoints remain V35 and V40 respectively.

\subsection{The exact block frame and its curvature defect}
\label{r31:frame}

Keep V29--V30's geometry. Index a block's six free links
by $i=(c,r)$, $c=0,1$, $r=0,1,2$, and use the forward
staples $s_i$ of \eqref{r30:staple}. The left endpoints
form a retained $\nu,\rho$ grid. Choose its tree consisting
of three horizontal edges and the two left-column vertical
edges. Let $k_i$ be the ordered tree transport from the
root $(0,0)$ to $i$, with $k_{(0,0)}=1$.

For an internal edge $i\to j$, use V29's retained left-
and right-endpoint links $\ell_{ij},r_{ij}$. Set
\[
 \begin{aligned}
 v_{i\alpha}&=k_i p_{i\alpha}s_i^{-1}k_i^{-1},&
 \ell'_{ij}&=k_i\ell_{ij}k_j^{-1},\\
 r'_{ij}&=k_i s_i r_{ij}s_j^{-1}k_j^{-1},&
 T_{ij}z&=\ell'_{ij}z(r'_{ij})^{-1},\\
 z_i&=k_iU_is_i^{-1}k_i^{-1}.&&
 \end{aligned}
 \tag{V31.1}\label{r31:frame-formula}
\]
Here $v_{i\alpha}\in S^3$, $T_{ij}\in\operatorname{SO}(4)$,
and the change of integration coordinates preserves Haar.
It differs from V30's fluctuation coordinate by a conjugation
depending only on the fixed exterior.

For $\vartheta=(v,T)$ define
\[
 \begin{split}
 {\cal E}_\vartheta(z)
   &=\sum_{i,\alpha}(1-z_i\cdot v_{i\alpha})
       +\sum_{\{i,j\}}(1-z_i\cdot T_{ij}z_j),\\
 Q_\gamma(\vartheta)&=\int_{(S^3)^6}
                              e^{-\gamma{\cal E}_\vartheta(z)}\,dz,\\
 d(\vartheta)^2&=\sum_{i,\alpha}|v_{i\alpha}-1|^2
                   +\sum_{\{i,j\}}\|T_{ij}-I\|_{\rm op}^2 .
 \end{split}
 \tag{V31.2}\label{r31:energy}
\]
The symbol $0$ denotes coherent coefficients $v=1,T=I$.
For native $\vartheta_b(\eta)$, this $Q_\gamma$ is
exactly $H_b(\eta)$ in \eqref{r29:action}.

Let $L_1,L_2$ be the two original $\nu\rho$ faces in
the left-endpoint grid. Define the retained, gauge-invariant
quantity
\[
 D_b^2=2E_{{\rm priv},b}(\Phi\eta)
       +4E_{{\rm int},b}(\Phi\eta)
       +16\{e(L_1)+e(L_2)\}.
 \tag{V31.3}\label{r31:D}
\]

\begin{lemma}[Exact frame and coefficient control]
\label{r31:frame-lemma}
The $29$ original block-plaquette energies become
$\mathcal E_\vartheta(z)$ in \eqref{r31:frame-formula}.
The filled configuration becomes $z_i=1$, and
\[
 d(\vartheta_b(\eta))^2\le D_b(\eta)^2 .
 \tag{V31.4}\label{r31:d-D}
\]
\end{lemma}
\begin{proof}
Multiplication by unit quaternions preserves the real
inner product. Private traces become $z_i\cdot v_{i\alpha}$.
An internal holonomy becomes
$z_i r'_{ij}z_j^{-1}(\ell'_{ij})^{-1}$, with trace
$z_i\cdot T_{ij}z_j$. Substituting $U_i=s_i$ gives
$z_i=1$ and
$\sum_{i,\alpha}|v_{i\alpha}-1|^2
=2E_{{\rm priv},b}(\Phi\eta)$.

The five tree edges have $\ell'=1$. Each of the two
remaining right-column vertical edges gives a conjugate
of the corresponding single left-grid square holonomy.
Thus $\sum_{\{i,j\}}|\ell'_{ij}-1|^2
=2(e(L_1)+e(L_2))$.
Also $|r'_{ij}-\ell'_{ij}|^2$ is twice its filled
internal-plaquette energy. The operator norm of
quaternion multiplication is its Euclidean norm, so
\[
 \|T_{ij}-I\|_{\rm op}
 \le2|\ell'_{ij}-1|+|r'_{ij}-\ell'_{ij}|,\qquad
 \|T_{ij}-I\|_{\rm op}^2
 \le8|\ell'_{ij}-1|^2+2|r'_{ij}-\ell'_{ij}|^2.
\]
Summing proves the bound. Gauge invariance of $D_b$
follows from equivariance of $\Phi$ and plaquette traces.
\end{proof}

These frames belong to individual conditional integrals.
Frames for adjacent blocks need not assemble into one
global gauge transformation, and no such assembly is used.

\subsection{A native 45-plaquette probability witness}
\label{r31:witness}

Let $\mathcal W_b$ contain the $22$ original private
plaquettes, seven top and fourteen endpoint side faces
of the block's V30 cubes, and $L_1,L_2$. These four
collections are disjoint by their direction pairs
or time layers; each contains distinct faces.
Thus $|\mathcal W_b|=45$.

\begin{proposition}[Original action pays the retained defect]
\label{r31:witness-prop}
For every full configuration $U$, with retained restriction
$\eta$,
\[
 D_b(\eta)^2\le20\sum_{p\in\mathcal W_b}e(U_p).
 \tag{V31.5}\label{r31:witness-bound}
\]
\end{proposition}
\begin{proof}
For a selected slot $s$ and another private slot $p$,
the chordal triangle inequality through the actual free
link gives
$e(sp^{-1})\le2e(U_is^{-1})+2e(U_ip^{-1})$.
There are at most three unselected slots per site.
Hence
$E_{{\rm priv},b}(\Phi\eta)\le6\sum_{\rm private}e(U_p)$.
V30's weighted cube inequality, before sharing between
blocks, gives
\[
 E_{{\rm int},b}(\Phi\eta)
   \le5\sum_{\rm top}e(U_p)+\tfrac52\sum_{\rm side}e(U_p).
\]
In \eqref{r31:D} the four collections therefore receive
coefficients $12,20,10,16$, all at most $20$.
\end{proof}

For probability statements take the cofinal dyadic tori
$L=2^a$, $2q=2^h$, $a,h\ge3$, the original isotropic
periodic Wilson law $\mu_W$, and $\gamma\ge1$.
Its exact retained marginal is $\mu_{\mathcal R}$.
The analytic block estimates do not require dyadic sizes.
No arbitrary exterior conditioning is assigned the
following tail estimate.

Archive $C$, V76 proves, using V65's coordinate-plane
reflections,
\[
 \begin{gathered}
 \mu_W\{e(U_p)\ge t\}\le\min\{1,e^{b_\gamma-\gamma t}\},
                   \qquad 0\le t\le2,\\
 c_H=\frac8{3\pi^3},\qquad
 b_\gamma=384\left(8+\log\gamma-\frac23\log c_H\right).
 \end{gathered}
 \tag{V31.6}\label{r31:inherited-tail}
\]
To record its normalization, a Haar cap of angular
radius $\gamma^{-1/2}$ has probability at least
$c_H\gamma^{-3/2}$. Restricting all links to these caps
gives
$-\log Z_\gamma\le
6|\Lambda|(8+\log\gamma-\frac23\log c_H)$.
The marked plaquette lies strictly inside a closed
$(4,4,2,2)$ reflection block, of volume $64$.
Dissemination forces action at least $t$ per block;
the normalized chessboard root pays $6\cdot64=384$
times the cap constant. This uses the original law
and its normalizer, not a cutoff law.

\begin{corollary}[Probability of the controlled coefficient set]
\label{r31:probability}
For $0<u\le1$,
\[
 \mu_{\mathcal R}\{D_b^2>u\}
 \le p_\gamma(u):=\min\{1,45e^{b_\gamma-\gamma u/900}\}.
 \tag{V31.7}\label{r31:bad}
\]
For fixed $p>0$ set
\[
 u_\gamma=\frac{900}{\gamma}
                 (b_\gamma+p\log\gamma+\log45).
 \tag{V31.8}\label{r31:u}
\]
When $u_\gamma\le1$, the bad probability is at most
$\gamma^{-p}$, uniformly in surrounding volume and
block position.
\end{corollary}
\begin{proof}
Equation \eqref{r31:witness-bound} makes $D_b^2>u$
imply that one of its $45$ witness energies exceeds
$u/900$. Apply \eqref{r31:inherited-tail} and the union
bound under $\mu_W$. The left event depends only on
retained links, so it has the same probability under
$\mu_{\mathcal R}$. Substitute \eqref{r31:u}.
\end{proof}

No dependent probabilities were multiplied.
The sufficient constants are large and can give no
useful moderate-coupling bound, but
$u_\gamma=O(\log\gamma/\gamma)\to0$ independently of
surrounding volume. The original witness contains free
links; the event whose probability is bounded does not.

\subsection{An explicit unique conditional minimum}
\label{r31:minimum}

Allow arbitrary unit slots and orthogonal pair transports
on this same graph, and use the product unit-round
$S^3$ metric.

\begin{theorem}[Uniform pinned neighborhood]
\label{r31:minimum-theorem}
If $d(\vartheta)\le\delta_0:=1/32$, the block energy has
a unique global minimum $z_*(\vartheta)$, smooth in the
coefficients on a neighborhood of this set. Put
$m(\vartheta)=\mathcal E_\vartheta(z_*)$ and
$M(\vartheta)=\operatorname{Hess}\mathcal E_\vartheta(z_*)$.
Then
\[
 \begin{gathered}
 \left(\sum_i|z_{*,i}-1|^2\right)^{1/2}\le2d(\vartheta),
 \qquad
 0\le m(\vartheta)\le\mathcal E_\vartheta(1)
                                  \le\tfrac12d(\vartheta)^2,\\
 2I_{18}\preceq M(\vartheta)\preceq9I_{18}.
 \end{gathered}
 \tag{V31.9}\label{r31:minimum-bounds}
\]
On the product of angular caps of radius $1/4$ the
Hessian is at least $(75/32)I$. Outside those caps,
\[
 \mathcal E_\vartheta(z)-m(\vartheta)\ge33/800>1/32.
 \tag{V31.10}\label{r31:barrier}
\]
\end{theorem}
\begin{proof}
Put $a_i=|z_i-1|$, $a^2=\sum_i a_i^2$, and $d=d(\vartheta)$.
Since $n_i\ge3$, the coherent energy satisfies
$\mathcal E_0(z)\ge(3/2)a^2$. Subtracting the energy at
$z=1$ cancels constant coefficient perturbations.
Field differences cost at most
$\sum_{i,\alpha}|v_{i\alpha}-1|a_i$.
For an edge,
\[
 |z_i\cdot(T-I)z_j-1\cdot(T-I)1|
                   \le\|T-I\|_{\rm op}(a_i+a_j).
\]
Cauchy--Schwarz and $n_i+2\deg(i)\le9$ therefore give
\[
 \begin{split}
 |(\mathcal E_\vartheta(z)-\mathcal E_\vartheta(1))
                                      -\mathcal E_0(z)|
 &\le d\left\{\sum_i n_i a_i^2+
                     \sum_{\{i,j\}}(a_i+a_j)^2\right\}^{1/2}\\
 &\le3da .
 \end{split}
 \tag{V31.11}\label{r31:localization}
\]
A global minimizer exists by compactness and has
$\mathcal E_\vartheta(z_*)\le\mathcal E_\vartheta(1)$.
Thus $(3/2)a^2-3da\le0$ there and $a\le2d$.
All energy summands are nonnegative, while
$1-1\cdot v=|v-1|^2/2$ and
$1-1\cdot T1=|T1-1|^2/2$, proving the energy bounds.

For tangent vectors $\xi_i\perp z_i$, direct differentiation
along product geodesics gives
\[
 \begin{split}
 \operatorname{Hess}\mathcal E_\vartheta(z)[\xi,\xi]
 ={}&\sum_{i,\alpha}(z_i\cdot v_{i\alpha})|\xi_i|^2\\
 &+\sum_{\{i,j\}}\left\{
 (z_i\cdot T_{ij}z_j)(|\xi_i|^2+|\xi_j|^2)
                              -2\xi_i\cdot T_{ij}\xi_j\right\}.
 \end{split}
 \tag{V31.12}\label{r31:hessian}
\]
Within the caps, $\rho=1/4$ gives
$z_i\cdot v_{i\alpha}\ge1-\rho^2/2-d$ and
$z_i\cdot T_{ij}z_j\ge1-2\rho^2-d$.
Use $2|\xi_i||\xi_j|\le|\xi_i|^2+|\xi_j|^2$.
The resulting coefficient at each site is at least
\[
 n_i(1-\rho^2/2)-2\deg(i)\rho^2-6d
 \ge3(1-1/32)-3/8-3/16=75/32.
\]
The upper bound from \eqref{r31:hessian} is
$\max_i(n_i+2\deg i)=9$.

All global minimizers, already localized by $a\le1/16$,
lie strictly within the caps. This product of caps is
geodesically convex; the positive Hessian excludes two
distinct critical minima there. Hence the global
minimum is unique. Its positive Hessian and the
implicit-function theorem give smooth dependence.
The same strict arguments work on a slightly larger
neighborhood, for example $d<1/24$, with a slightly
weaker Hessian lower coefficient.

Outside the caps some $a_i\ge2\sin(1/8)>1/5$.
By \eqref{r31:localization} and monotonicity of
$(3/2)a^2-3a/32$ for $a\ge1/5$,
\[
 \mathcal E_\vartheta(z)-m(\vartheta)
 \ge\mathcal E_\vartheta(z)-\mathcal E_\vartheta(1)
 \ge\frac32(1/5)^2-\frac3{32}\frac15=\frac{33}{800}.
\]
\end{proof}

This does not contradict V29's singular full flat
minimum: the exterior here is fixed and the six
integrated links are pinned. The positive conditional
Hessian is not a positive Hessian on all retained
or physical directions.

\subsection{Uniform Haar integration and the determinant correction}
\label{r31:laplace-section}

\begin{theorem}[Uniform relative block integral]
\label{r31:laplace-theorem}
There is a finite $C_L\ge1$, depending only on this
graph and $\delta_0$, such that for $d(\vartheta)\le\delta_0$
and $\gamma\ge1$,
\[
 \begin{split}
 Q_\gamma(\vartheta)
 &=e^{-\gamma m(\vartheta)}\gamma^{-9}
       \frac8{\pi^3\sqrt{\det M(\vartheta)}}
                       (1+\varepsilon_\gamma(\vartheta)),\\
 |\varepsilon_\gamma(\vartheta)|&\le C_L/\gamma.
 \end{split}
 \tag{V31.13}\label{r31:laplace}
\]
The constant is uniform in the admissible exterior,
block position and surrounding volume. Its numerical
value is not evaluated here.
\end{theorem}
\begin{proof}
Use product Riemannian normal coordinates
$y\in\mathbb R^{18}$ centered at $z_*(\vartheta)$.
A common positive radius exists since the minimizers
lie in a fixed smaller cap. Taylor's formula gives
\[
 \mathcal E_\vartheta(\exp_{z_*}y)-m(\vartheta)
 =\tfrac12 y^{\mathsf T}M(\vartheta)y
                         +P_{3,\vartheta}(y)+R_{4,\vartheta}(y),
\]
where $P_3$ is homogeneous cubic,
$|P_3|\le C_3|y|^3$ and $|R_4|\le C_4|y|^4$,
uniformly. These finite constants follow from
compactness of the coefficient set, smooth dependence
of its minimum and bounded derivatives of the finite
dot-product energy.

Normalized product Haar is
\[
 (2\pi^2)^{-6}j(y)\,dy,\qquad
 j(y)=\prod_{i=1}^6
        \left(\frac{\sin|y_i|}{|y_i|}\right)^2
       =1+O(|y|^2).
\]
It is even. Choose the common radius small enough
that Taylor errors absorb at most a fixed fraction
of the positive quadratic form. Using
$|e^{-s}-1+s|\le s^2e^{|s|}/2$, the error after
subtracting the linear cubic term is bounded by
\[
 C e^{-c\gamma|y|^2}
       \{\gamma|y|^4+\gamma^2|y|^6+|y|^2\}.
\]
The subtracted cubic times the centered Gaussian
integrates to zero over the symmetric ball.
The displayed bound has integral $O(\gamma^{-10})$
by $y=\gamma^{-1/2}Y$.

Inside the caps strong convexity gives a uniform
positive energy cost outside the normal ball;
outside the caps use \eqref{r31:barrier}.
Both that omitted compact integral and the Gaussian
tail are exponentially small, uniformly after
multiplication by any fixed power of $\gamma$.
The full Gaussian integral is
$(2\pi)^9\gamma^{-9}/\sqrt{\det M}$.
Multiplying by $(2\pi^2)^{-6}$ gives $8/\pi^3$.
Its leading coefficient is uniformly bounded below
because $M\preceq9I$. Division gives the relative
$O(\gamma^{-1})$ bound. Increasing its constant on
a compact initial interval covers every $\gamma\ge1$.
\end{proof}

The proof establishes a finite uniform constant,
not a certified moderate-coupling threshold. Conditions
involving $C_L$ below must not be evaluated with a
guessed numerical value.

\begin{proposition}[Coherent determinant and zero linear variation]
\label{r31:determinant}
Let $\mathsf A=6I_6-\operatorname{Adj}(P_2\times P_3)$.
Then
\[
 M(0)=\mathsf A\otimes I_3,\qquad
 \det\mathsf A=37835,\qquad
 \frac8{\pi^3\sqrt{\det M(0)}}=
                       \frac8{\pi^3\,37835^{3/2}}.
 \tag{V31.14}\label{r31:det0}
\]
For some finite graph constant $C_D$,
\[
 |\log\det M(\vartheta)-\log\det M(0)|
                    \le C_D d(\vartheta)^2
           \quad(d(\vartheta)\le\delta_0).
 \tag{V31.15}\label{r31:det-bound}
\]
\end{proposition}
\begin{proof}
At coherence \eqref{r31:hessian} has diagonal $6I_3$
and off-diagonal $-I_3$ on graph edges.
The scalar eigenvalues are
$5,5\pm\sqrt2,7,7\pm\sqrt2$; their product is
$5(25-2)\,7(49-2)=37835$.

Take any smooth coefficient curve through coherence,
and parallel-transport orthonormal tangent frames
along its moving minimizer. Unit-slot derivatives
are perpendicular to $1$; transport derivatives
are $K_{ij}\in\mathfrak{so}(4)$.
The diagonal contact coefficients in
\eqref{r31:hessian} have zero first derivatives.
An off-diagonal derivative is minus the transverse
$3$ by $3$ block of $K_{ij}$: the moving-frame
contributions have zero first-order inner product
at coherence. That transverse block is skew.
Every block of $M(0)^{-1}=\mathsf A^{-1}\otimes I_3$
is scalar in color, so
\[
 \left.\frac{d}{ds}\log\det M(\vartheta(s))\right|_{s=0}
                   =\Tr(M(0)^{-1}\dot M(0))=0.
\]
The second derivatives of $\log\det M$ are bounded
on a slightly larger compact parameter neighborhood,
by smoothness and the Hessian lower bound.
Spherical and orthogonal-group logarithms have squared
coordinate norm bounded by a constant times
$d(\vartheta)^2$ there. Taylor's formula with the
zero first derivative proves \eqref{r31:det-bound}.
\end{proof}

Write $E_{\Phi,b}=\mathcal E_{\vartheta_b}(1)$.
V30's exact factor is
$J_{\gamma,b}=e^{\gamma E_{\Phi,b}}Q_\gamma(\vartheta_b)$.
For $d(\vartheta_b)\le\delta_0$, $\gamma\ge2C_L$,
\[
 \begin{split}
 \log\frac{J_{\gamma,b}(\eta)}{J_{\gamma,b}(1)}
 &=\gamma(E_{\Phi,b}-m_b)
       +\tfrac12\log\frac{\det M(0)}{\det M_b}
       +r_{\gamma,b},\\
 |r_{\gamma,b}|&\le4C_L/\gamma,\qquad
 0\le E_{\Phi,b}-m_b\le\tfrac12D_b^2 .
 \end{split}
 \tag{V31.16}\label{r31:return}
\]
The remainder estimate uses
$|\log(1+\varepsilon)|\le2|\varepsilon|$ for
$|\varepsilon|\le1/2$ at both exteriors.
The energy relaxation can be of order $\gamma D_b^2$;
it cannot be discarded simply because $D_b\to0$.

\subsection{A normalized replacement in the common exterior}
\label{r31:normalized}

Let $\mathcal S$ be any $M$ blocks of the packing,
possibly with overlapping retained supports.
For $0<u\le\delta_0^2$, put $G_b(u)=\{D_b^2\le u\}$
and, on this set, define
\[
 \mathcal A_{\gamma,b}=
 e^{-\gamma m_b}\gamma^{-9}
                  \frac8{\pi^3\sqrt{\det M_b}},\qquad
 \mathcal L_{\gamma,b}=
 (1-C_L/\gamma)\mathcal A_{\gamma,b}\mathbf1_{G_b(u)}.
 \tag{V31.17}\label{r31:lower}
\]
Outside $G_b$, define the lower factor to be zero
without evaluating the approximation there.
Assume $\gamma\ge2C_L$.
These functions are gauge invariant: changing the
exterior gauge acts by isometries in the conditional
integral and preserves its minimum and intrinsic
Hessian determinant.

\begin{theorem}[Native local normalized replacement]
\label{r31:replacement}
In the exact unnormalized retained density
\[
 e^{-\gamma S_0(\eta)}\prod_b Q_{\gamma,b}(\eta),
 \qquad Q_{\gamma,b}=H_b,
\]
replace $Q_{\gamma,b}$ by $\mathcal L_{\gamma,b}$
only for $b\in\mathcal S$, leaving every other
factor unchanged. Put
\[
 \overline\alpha
 =M\{p_\gamma(u)+2C_L/\gamma\}.
 \tag{V31.18}\label{r31:alpha}
\]
On the probability-compatible tori above, if
$\overline\alpha<1$ the new normalizer is positive
and the resulting probability law satisfies
\[
 \|\widehat\mu_{\mathcal R}-\mu_{\mathcal R}\|_{\rm TV}
 \le\overline\alpha,\qquad
 H(\widehat\mu_{\mathcal R},\mu_{\mathcal R})
                    \ge\sqrt{1-\overline\alpha}.
 \tag{V31.19}\label{r31:TV}
\]
The same bounds hold after a common Markov operation,
including marginalization. With
$u=u_\gamma\le\delta_0^2$, one has
\[
 \overline\alpha\le
               M\{\gamma^{-p}+2C_L/\gamma\}.
 \tag{V31.20}\label{r31:rate}
\]
For fixed $M$ this tends to zero independently of
surrounding volume; the other native factors and
their common correlations remain in the measure.
\end{theorem}
\begin{proof}
Put $c=C_L/\gamma\le1/2$.
The relative integral estimate gives
$0\le\mathcal L_b\le Q_b$ everywhere and, on $G_b$,
\[
 \mathcal L_b/Q_b\ge(1-c)/(1+c)\ge1-2c.
\]
Set $r=\prod_{b\in\mathcal S}\mathcal L_b/Q_b\in[0,1]$.
Using $1-\prod r_b\le\sum(1-r_b)$ under the original
retained law,
\[
 1-\mu_{\mathcal R}[r]
 \le\sum_{b\in\mathcal S}
          \{\mu_{\mathcal R}(G_b^c)+2c\}
 \le\overline\alpha .
\]
This pays the relative normalizer, not merely the
individual integrals.

Let $z=\mu_{\mathcal R}[r]>0$. The new law is
$r\,d\mu_{\mathcal R}/z$; the old law is
$z\widehat\mu_{\mathcal R}$ plus a positive measure
of mass $1-z$. This proves the total-variation bound.
Also $H=\mu_{\mathcal R}[\sqrt r]/\sqrt z\ge\sqrt z$.
Integration proves total-variation contraction and
conditional Cauchy--Schwarz proves affinity increase
under a common Markov operation. These are V21's
inherited filter facts, now with a native local
error budget. Equation \eqref{r31:bad} and then
\eqref{r31:u} give the remaining assertions.
\end{proof}

Using only the conditional minimum and the constant
coherent prefactor in \eqref{r31:det0} would add
relative logarithmic error at most $C_Du/2$ on $G_b$.
A corresponding lower envelope would give
$O(M(u+\gamma^{-1}))+Mp_\gamma(u)$ normalized loss,
hence $O(M\log\gamma/\gamma)+M\gamma^{-p}$ at
\eqref{r31:u}. Keeping the determinant gives the
stronger bound \eqref{r31:rate}, selected here.
Neither version substitutes the filled configuration
for the moving minimum.

\paragraph{Remaining scope boundary.}
Taking every block makes this error extensive.
The theorem proves no intensive accumulated error,
convergent cluster expansion, uniform conditional
mixing, or physical reflection positivity of the
approximate law. One cannot reuse the original
chessboard estimate after replacement without
proving its hypotheses for the changed law.
A local normalized approximation is not an
all-scale mass-gap proof.

\subsection{Checks, preservation, and the next obligation}
\label{r31:checks}

The checker
\begin{center}\small
\path{scripts/verify_ym_restart_v31_local_return.py}
\end{center}
constructs the $45$-face patch on
$(L,q)=(4,3),(4,4),(8,4),(8,8)$, with respectively
$8,8,64,192$ blocks. Four rational quaternion
configurations on the smaller two geometries give
$32$ block fixtures. They check all $29$ original
block energies against their framed dot-product
sum, both ordered tree-square identities per
fixture, and \eqref{r31:witness-bound}.
All principal minors of the rational $4$ by $4$
Gram-difference matrix additionally certify the
transport operator-norm domination on each fixture.

The determinant $37835$, Hessian lower coefficient
$75/32$, upper coefficient $9$ and barrier $33/800$
are exact arithmetic checks. All $42$ transport
tangent generators have zero logarithmic determinant
variation. The $66$ unit-slot tangent generators
have zero first Hessian variation by the formula
in the proof. The script does not numerically
evaluate $C_L$, $C_D$, a global mixing rate or a
physical Hamiltonian spectrum. The all-parameter
and uniform remainder arguments are proved in
the text, not inferred from finite fixtures.

Removing this append and restoring V30's title
recovers the entire 713901-byte predecessor with
SHA--256
\begin{center}\small
\path{E4883258E4102B97705D0BF282EB37FEC30F6A2385B8583AFD45E7776B504375}.
\end{center}
All 46 other manuscript sources remain unchanged.
Only the active predecessor filename is replaced;
its mathematical body is preserved. The folder
\path{latex_notes} contains only \texttt{.tex};
checks and typesetting outputs remain outside it.

\paragraph{Direction.}
The next target is nonextensive control of the
\emph{interacting retained minimum action}, its determinant
corrections and physical endpoint influence. Summing local
errors will not resolve it: a global argument must retain
connected dependencies and the changing normalizer.
The interacting four-dimensional continuum construction,
cofinal scale control and positive physical Hamiltonian
mass gap remain unproved.

% END CONSOLIDATED V31 APPEND

% BEGIN CONSOLIDATED V32 APPEND -- 2026-09-19
\clearpage
\section[V32: differentiable local remainders and finite-range comparison]
{V32: differentiable local remainders\\and finite-range comparison}
\label{r32:start}

\subsection{Context: which accumulation problem is addressed}

V31's all-block total-variation estimate would pay the number
of blocks. We now control a different object: the actual
logarithmic remainder as a finite-range interaction, including
two derivatives, uniformly in ambient volume. This does not
prove that an extensive perturbation has a small total
variation. It gives a local conditional comparison at every
exterior, without assuming independence or retained-field mixing.

The moving minimum and its determinant remain in the leading
factor. A coherent scalar normalization removes a field-independent
remainder. Its first coefficient derivative vanishes exactly,
so the field-dependent remainder is quadratic in the local
coefficient defect, with a prefactor \(O(\gamma^{-1})\).
A smooth cutoff can shrink with coupling without making its
second derivatives large. Rough factors are kept \emph{exact},
not set to zero or declared negligible.

Archive \(A\), V56 already transports covariance matrices for
growing separated families and explicitly distinguishes a
response norm from spatial decay. Archive \(C\), V28 already
removes an extensive loss from rough \emph{update forms}.
Archive \(C\), V31 already constructs an exponentially summable
interaction for a different, window-restricted bridge law,
under its fast-field clustering hypotheses. Those methods and
results are not rediscovered here. Our object is V31's exact
six-link elimination in the unrestricted compact retained
variables. The new derivative bound is proved before any
retained integration or conditional mixing argument.

A targeted primary-source check of Katsevich,
\href{https://arxiv.org/abs/2406.12706v3}
{\emph{The Laplace asymptotic expansion in high dimensions}},
confirms that growing integration dimension requires additional
derivative control. We use no theorem from that paper and do
not treat all free blocks as one uniformly nondegenerate saddle.
The proof below is a fixed \(18\)-dimensional, parameter-smooth
calculation followed by exact finite-range bookkeeping.
The scheduled companion and broad-literature checks remain
V35 and V40.

\subsection{A twice differentiable logarithmic remainder}

Keep the coefficient manifold
\(\mathcal P=(S^3)^{22}\times\mathrm{SO}(4)^7\), with coherent
point \(o=(1,\ldots,1;I,\ldots,I)\).
For slots use spherical logarithms at \(1\); for transports
use principal skew-matrix logarithms, with Frobenius norm.
Their product coordinates are denoted \(a\).
Define the globally smooth squared chord defect
\[
 \rho(\vartheta)^2
 =\sum_{i,\alpha}|v_{i\alpha}-1|^2
                  +\sum_{\{i,j\}}\|T_{ij}-I\|_{\mathrm F}^2 .
 \tag{V32.1}\label{r32:rho}
\]
In the neighborhood used below,
\[
 d(\vartheta)\le\rho(\vartheta),\qquad
 \rho(\vartheta)^2\le4d(\vartheta)^2,\qquad
 \rho(\vartheta)\le |a|\le2\rho(\vartheta).
 \tag{V32.2}\label{r32:comparison}
\]
The first two inequalities use
\(\|B\|_{\rm op}\le\|B\|_{\rm F}\le2\|B\|_{\rm op}\)
for \(4\)-by-\(4\) matrices. The last two follow componentwise
from \(2\sin(t/2)\le t\le2[2\sin(t/2)]\) for the small
principal angles. All derivative norms below use these fixed
finite-dimensional coordinates.

Set \(r_0=1/256\). On \(|a|<1/32\) V31 supplies the unique
smooth minimum \(m(a)\) and positive Hessian \(M(a)\).
Define, without asymptotic notation,
\[
 f_\gamma(a)=\log Q_\gamma(a)+\gamma m(a)+9\log\gamma
             -\log(8/\pi^3)+\tfrac12\log\det M(a).
 \tag{V32.3}\label{r32:f}
\]
Thus \(f_\gamma=\log(1+\varepsilon_\gamma)\) wherever V31's
notation is used. Positivity of \(Q_\gamma\), not a guessed
bound on \(\varepsilon_\gamma\), makes this definition valid.

\begin{theorem}[A differentiable, coherence-subtracted remainder]
\label{r32:derivatives}
There is a finite graph constant \(C\), independent of
\(\gamma\ge1\), such that on \(|a|\le1/64\)
\[
 \|f_\gamma\|_{C^2}\le C/\gamma,\qquad
 D_af_\gamma(0)=0.
 \tag{V32.4}\label{r32:c2}
\]
Consequently \(h_\gamma(a)=f_\gamma(a)-f_\gamma(0)\) satisfies
\[
 |h_\gamma(a)|\le C|a|^2/\gamma,\qquad
 |D_ah_\gamma(a)|\le C|a|/\gamma,\qquad
 \|D_a^2h_\gamma(a)\|\le C/\gamma.
 \tag{V32.5}\label{r32:h}
\]
Constants can be enlarged between displays. They are proved
finite but are not numerically evaluated.
\end{theorem}
\begin{proof}
The closed coordinate ball has a larger compact neighborhood
inside V31's strict uniqueness and Hessian region.
Use smoothly chosen orthonormal tangent frames at the moving
minimum; these exist over this small parameter ball.
In product normal coordinates write
\[
 F_a(y)=\mathcal E_a(\exp_{z_*(a)}y)-m(a).
\]
Its value and first \(y\)-derivative vanish at zero, and
\(D_y^2F_a(0)=M(a)\). On one common small ball,
\[
 B_a(y)=2\int_0^1(1-s)D_y^2F_a(sy)\,ds,\qquad
 F_a(y)=\tfrac12y^{\mathsf T}B_a(y)y .
\]
The matrices \(B_a(y)\) are uniformly positive. The map
\(u=B_a(y)^{1/2}y\) has derivative \(M(a)^{1/2}\) at zero.
Uniform smoothness and the inverse-function theorem give
a common smaller \(u\)-ball, a smooth inverse \(y_a(u)\),
and bounded derivatives of every fixed order. In these
coordinates the phase is exactly \(|u|^2/2\).

Choose an even smooth cutoff \(\psi(u)\), equal to one near
zero and supported inside that common ball. Pull it back
to the original integration variables, extending it by zero.
The transformed normalized Haar density is a smooth
\(A(a,u)\,du\), with
\[
 A(a,0)=(2\pi^2)^{-6}(\det M(a))^{-1/2}.
\]
Both parameter and \(u\)-derivatives are bounded on its
compact domain. Taylor's formula in \(u\), also after up to
two \(a\)-derivatives, gives
\[
 A(a,u)=A(a,0)+D_uA(a,0)u+R(a,u),\qquad
 \|R(\,\cdot\,,u)\|_{C^2_a}\le C|u|^2.
\]
The linear term integrates to zero against
\(\psi(u)e^{-\gamma|u|^2/2}\). The remainder has integral
\(O_{C^2_a}(\gamma^{-10})\), while the leading Gaussian
integral is \((2\pi)^9\gamma^{-9}A(a,0)\).

The complementary original integral has a uniform positive
energy barrier. This follows from uniqueness on a compact
parameter set, local positivity, and V31's exterior-cap
barrier. Differentiating it twice in \(a\) costs at most
\(C(1+\gamma^2)e^{-c\gamma}\): the moving cutoff and
\(F_a=\mathcal E_a-m(a)\) have bounded derivatives there.
The same estimate treats the Gaussian cutoff tail.
It follows that
\[
 Q_\gamma(a)e^{\gamma m(a)}\gamma^9
 =\frac8{\pi^3\sqrt{\det M(a)}}
                         \{1+O_{C^2_a}(\gamma^{-1})\}.
\]
The leading coefficient is bounded below. Taking its
logarithmic ratio proves the first assertion for all
sufficiently large \(\gamma\). On the remaining compact
interval \(1\le\gamma\le\gamma_0\), the positive integral
and all displayed smooth functions have bounded \(C^2\)
norms. Enlarging \(C\) proves the bound for every \(\gamma\ge1\).
This is a proof of differentiable remainders; differentiating
V31's scalar big-O alone would not be legitimate.

For the exact first derivative, at coherence the spin law is
invariant under simultaneous \(R\in\mathrm{SO}(3)\) acting
on the three transverse quaternion components. Hence
\(\E z_i\) is a multiple of \(1\), and
\(\E(z_jz_i^{\mathsf T})\) is block diagonal with a scalar
transverse block. It is symmetric. A slot variation is
orthogonal to \(1\), and a transport variation is a
skew matrix \(K\). Therefore
\[
 \E(z_i\cdot \dot v)=0,\qquad
 \E(z_i\cdot Kz_j)=0.
\]
Differentiation of the compact integral gives
\(D_a\log Q_\gamma(0)=0\). The envelope formula at \(z_i=1\)
gives \(D_am(0)=0\), and V31's determinant calculation gives
\(D_a\log\det M(0)=0\). Thus \(D_af_\gamma(0)=0\) for every
\(\gamma>0\), not just asymptotically. Integrating the
\(C^2\) bound along the straight segment from zero to \(a\)
proves \eqref{r32:h}.
\end{proof}

\subsection{A cutoff with no inverse-radius Hessian loss}

Fix a \(C^\infty\) function \(\chi:[0,\infty)\to[0,1]\)
equal to one on \([0,1]\) and zero on \([4,\infty)\).
For \(0<r\le r_0\), in the native frame of block \(b\), set
\[
 w_{\gamma,r,b}(\eta)=
 \chi\bigl(\rho(\vartheta_b(\eta))^2/r^2\bigr)
                         h_\gamma(\vartheta_b(\eta)).
 \tag{V32.6}\label{r32:w}
\]
Interpret this as zero outside \(\rho<2r\); no minimum or
logarithm is evaluated there. On its support \(|a|\le4r\le1/64\).
Vanishing of the cutoff with all derivatives makes the
extension globally smooth on the compact retained link space.

Let \(S_b\) be the union of retained links appearing in
the slots, selected forward staples, left-grid tree paths,
and the two endpoint links of each internal interaction.
For a retained link \(e\), \(\nabla_e\) denotes its gradient
in the unit-round \(S^3\) metric; second derivatives are
covariant product derivatives.

\begin{proposition}[Native finite support and derivative bounds]
\label{r32:native}
Each \(w_b=w_{\gamma,r,b}\) is gauge invariant, depends only
on \(\eta_{S_b}\), and obeys
\[
 \|w_b\|_\infty\le Cr^2/\gamma,\qquad
 \|\nabla_e w_b\|_\infty\le Cr/\gamma,\qquad
 \|\nabla_e\nabla_f w_b\|_\infty\le C/\gamma
 \quad(e,f\in S_b).
 \tag{V32.7}\label{r32:local-norms}
\]
The same \(C\) works for all blocks, tori, and
\(0<r\le r_0\).
With link distance defined as torus \(\ell^1\) distance
between initial vertices plus one when directions differ,
\[
 |S_b|\le480=:s_*,\qquad
 \sup_e\#\{b:e\in S_b\}\le120=:\kappa_*,
 \qquad\operatorname{diam}S_b\le11.
 \tag{V32.8}\label{r32:incidence}
\]
These are conservative combinatorial bounds.
\end{proposition}
\begin{proof}
Write \(g=\rho^2\). On the coefficient chart,
\(g=O(|a|^2)\), \(Dg=O(|a|)\), and \(D^2g=O(1)\).
Where derivatives of \(\chi(g/r^2)\) occur, \(|a|\le4r\).
The product rule and \eqref{r32:h} bound its second derivative
by sums of terms of sizes
\[
 \gamma^{-1},\quad
 r^{-2}r(r/\gamma),\quad
 r^{-2}(r^2/\gamma),\quad
 r^{-4}r^2(r^2/\gamma).
\]
Every term is \(O(\gamma^{-1})\); the zeroth and first
derivatives give \(Cr^2/\gamma\) and \(Cr/\gamma\).
The native coefficient maps are fixed finite products and
inverses of compact group matrices. Their first two
derivatives, including the coordinate logarithms on the
cutoff support, are bounded independently of volume.
The chain rule proves \eqref{r32:local-norms}.

A gauge transformation conjugates every \(v_{i\alpha}\)
by the root gauge and conjugates every \(T_{ij}\) by its
orthogonal quaternion action, which fixes \(1\).
Thus \(\rho\), \(Q\), \(m\), and \(\det M\) are unchanged.
Their combination \(w_b\) is gauge invariant; the local
tree frames need not assemble into one global gauge.

For a block anchored at \(x\), every initial vertex of
a link in \(S_b\) lies in the image of
\[
 x+\{0,1\}\times\{-1,0,1,2\}
                \times\{-1,0,1,2,3\}\times\{-1,0,1\}.
 \tag{V32.9}\label{r32:box}
\]
This is checked directly for the three-edge plaquette
staples, forward staple, five-edge grid tree and endpoint
links. There are \(120\) possible vertices and four
directions. For a fixed initial vertex, at most \(120\)
anchor displacements can contain it. The actual packing
uses a subset of those anchors. Periodic identifications
only lower these counts. Coordinate spans sum to \(10\);
the direction allowance gives \(11\).
\end{proof}

For \(a_*>0\), define a supremum interaction norm and a
weighted Hessian row norm by
\[
 \begin{aligned}
 \|w\|_{{\rm int},a_*}
 &=\sup_e\sum_{b:e\in S_b}|S_b|e^{a_*\operatorname{diam}S_b}
                                                     \|w_b\|_\infty,\\
 \|w\|_{{\rm Hess},a_*}
 &=\sup_e\sum_f e^{a_*\operatorname{dist}(e,f)}
                  \sum_{b:e,f\in S_b}
                              \|\nabla_e\nabla_f w_b\|_\infty .
 \end{aligned}
\]
Finite support, not cancellation or mixing, proves
\[
 \|w\|_{{\rm int},a_*}\le
      \kappa_*s_*e^{11a_*}Cr^2/\gamma,\qquad
 \|w\|_{{\rm Hess},a_*}\le
      \kappa_*s_*e^{11a_*}C/\gamma.
 \tag{V32.10}\label{r32:int}
\]
There is no ambient block count in these norms.


\subsection{All factors at once, but only a local comparison}

Define the strictly positive hybrid factor everywhere by
\[
 \overline Q_{\gamma,r,b}(\eta)
       =Q_{\gamma,b}(\eta)\exp\{-w_{\gamma,r,b}(\eta)\}.
 \tag{V32.11}\label{r32:hybrid}
\]
On \(\rho_b\le r\), this is exactly
\[
 \overline Q_{\gamma,r,b}
 =e^{f_\gamma(0)}e^{-\gamma m_b}\gamma^{-9}
                           \frac8{\pi^3\sqrt{\det M_b}},
\]
and on \(\rho_b\ge2r\) it is exactly \(Q_{\gamma,b}\).
In between it is the geometric interpolation prescribed by
\(\chi\). The scalar \(e^{f_\gamma(0)}\) is the same coherent
normalization for every block. It is retained in this
definition; dropping it only on good blocks would introduce
a field-dependent cutoff term of a different size.

Normalize the full hybrid weight on the same product Haar
space:
\[
 d\overline\mu_R
  =\overline Z^{-1}e^{-\gamma S_0}
                         \prod_b\overline Q_{\gamma,r,b}\,d\eta.
\]
All factors are smooth and positive on a compact space.
The actual law satisfies the exact identity
\[
 \frac{d\mu_R}{d\overline\mu_R}
  =\frac{\exp\{\sum_b w_b\}}
           {\overline\mu_R(\exp\{\sum_b w_b\})}.
 \tag{V32.12}\label{r32:density}
\]
This identity does not factor the retained law.

For a set \(D\) of retained links, condition both laws on
the same exterior \(\eta_{D^c}=\xi\). Positive densities
specify these kernels for every \(\xi\), not only almost
every exterior. Write them \(K_D^\xi,\overline K_D^\xi\).
Let \(n_D=\#\{b:S_b\cap D\ne\varnothing\}\).

\begin{theorem}[Uniform normalized local specifications]
\label{r32:kernels}
For every \(D,\xi\), with
\(\omega_D=2Cr^2 n_D/\gamma\) and
\(n_D\le\kappa_*|D|\),
\[
 \begin{gathered}
 e^{-\omega_D}
 \le\frac{dK_D^\xi}{d\overline K_D^\xi}
 \le e^{\omega_D},\qquad
 d_{\rm TV}(K_D^\xi,\overline K_D^\xi)
                         \le\min\{1,\omega_D/4\},\\
 e^{-\omega_D}\Var_{\overline K_D^\xi}F
 \le\Var_{K_D^\xi}F
 \le e^{\omega_D}\Var_{\overline K_D^\xi}F .
 \end{gathered}
 \tag{V32.13}\label{r32:kernel-bounds}
\]
The variance assertion holds for any square-integrable
real \(F\). All constants are independent of the number
of blocks outside \(D\).
\end{theorem}
\begin{proof}
In the conditional density all factors with
\(S_b\cap D=\varnothing\) cancel, including their
normalization contributions. Thus the tilt is
\(\Phi_D=\sum_{b:S_b\cap D\ne\varnothing}w_b\),
with \(\operatorname{osc}\Phi_D\le\omega_D\).
Its normalized exponential lies between
\(e^{-\omega_D}\) and \(e^{\omega_D}\).
The variance bounds follow by applying these bounds
to \((F-c)^2\) and taking the infimum over constants \(c\).

For completeness, let \(P_t\) be the normalized tilt of
\(\overline K_D^\xi\) by \(t\Phi_D\), \(0\le t\le1\).
For an event \(A\), compact boundedness permits
\[
 \frac{d}{dt}P_t(A)=\Cov_{P_t}(1_A,\Phi_D).
\]
Subtracting a range midpoint bounds the variance of a
variable of range length \(l\) by \(l^2/4\).
The Cauchy--Schwarz inequality bounds the last covariance
by \(\operatorname{osc}\Phi_D/4\). Integrate in \(t\)
and take the supremum over events.
Finally each touching block contains at least one link
of \(D\), so \(n_D\le\sum_{e\in D}\#\{b:e\in S_b\}\).
\end{proof}

In particular, every fixed-size conditional update has
error \(O(r^2/\gamma)\), even when every block factor
has been replaced. This is stronger in its local scope
than a union bound over the whole lattice. It is not
a bound on the unconditional marginal on \(D\):
the two kernels would then be integrated against
different exterior laws. Likewise, integrating the
variance comparison against different exterior laws
does not compare the two global Dirichlet forms or gaps.

\subsection{A shrinking thermal window and its exact limitations}

For the original periodic law on V31's admitted dyadic tori,
let \(p>0\), and use its already proved constants
\[
 \begin{aligned}
 c_H&=8/(3\pi^3),&
 b_\gamma&=384\{8+\log\gamma-(2/3)\log c_H\},\\
 u_\gamma&=\frac{900}{\gamma}
          \{b_\gamma+p\log\gamma+\log45\},&
 r_\gamma^2&=4u_\gamma .
 \end{aligned}
 \tag{V32.14}\label{r32:thermal}
\]
For sufficiently large \(\gamma\), \(r_\gamma\le r_0\).
V31's native witness and \eqref{r32:comparison} give
\[
 \mu_R(\rho_b>r_\gamma)
 \le\mu_R(D_b^2>u_\gamma)\le\gamma^{-p}.
 \tag{V32.15}\label{r32:probability}
\]
There is no use of reflection positivity for
\(\overline\mu_R\). Under the actual \(\mu_R\), each block
is therefore in the pure moving-minimum/determinant
region with probability at least \(1-\gamma^{-p}\).
The deterministic results simultaneously give
\[
 \begin{aligned}
 \|w\|_{{\rm int},a_*}
      &=O_{p,a_*}(\log\gamma/\gamma^2),\\
 \|w\|_{{\rm Hess},a_*}
      &=O_{a_*}(\gamma^{-1}),\\
 \sup_\xi d_{\rm TV}(K_D^\xi,\overline K_D^\xi)
      &=O_p(|D|\log\gamma/\gamma^2).
 \end{aligned}
 \tag{V32.16}\label{r32:thermal-rates}
\]
The conditions on \(r_\gamma\) are asymptotic and have
very large constants. No moderate-coupling numerical
certificate is asserted.

Small local changes do not by themselves give global
total-variation closeness. An elementary exact example
is a product of \(N\) Bernoulli variables with parameters
\(1/2\) and \(1/2+\epsilon\), \(0<\epsilon<1/2\).
Every one-site conditional difference is \(\epsilon\).
The empirical mean has variance at most \(1/(4N)\);
Chebyshev's inequality at threshold \(1/2+\epsilon/2\)
shows that the total variation of the two product laws
tends to one as \(N\to\infty\). The example is not a YM
counterexample, but rules out the unsupported inference.

The remaining obstacle is now sharper. The small term
controlled here is the \emph{centered local Laplace
remainder}, not the whole retained action. On good
blocks that action still contains \(\gamma m_b\) and
\(\tfrac12\log\det M_b\), and on rough blocks it contains
the original exact integral. Their long-distance,
normalized retained response is not bounded here.
Local specifications can stay close even when a reference
law has no uniform susceptibility or uniqueness bound.
Neither a physical-time contraction nor a continuum
Hamiltonian gap follows from \eqref{r32:int}.

\subsection{Verification and the next upstream input}

The accompanying checker verifies the native coefficient
carriers on four finite tori, including periodic seams,
and checks the stated support, diameter and overlap
bounds. It checks the algebraic symmetry responsible for
the exact first derivative and the scale cancellation
in the cutoff Hessian. These are finite bookkeeping
tests, not numerical evaluations of the compact-integral
constant \(C\), retained correlations, or physical spectra.

The complete predecessor V31 is recoverable byte-for-byte
by removing this append and restoring the displayed title.
Other manuscript sources remain unchanged; build and
verification files are outside \path{latex_notes}.
The next substantive task is a spatially resolved estimate
for the \emph{leading} hybrid retained law, or an exact
physical-endpoint comparison which avoids requiring it.
Reproving scalar Laplace asymptotics, replacing all-block
total variation by an unproved mixing assumption, or
suppressing rough factors would not close that obligation.
The interacting four-dimensional continuum construction
and its positive physical mass gap remain open.

% END CONSOLIDATED V32 APPEND

% BEGIN CONSOLIDATED V33 APPEND -- 2026-09-19
\clearpage
\section[V33: energy-relative elimination and the native Schur limit]
{V33: energy-relative elimination\\and the native Schur limit}
\label{r33:start}

\subsection{Context and a genuinely global comparison}

V32 controls a small finite-range remainder but retains exact
rough factors. We now go upstream to the leading action:
minimize the original Wilson energy over precisely the links
that were integrated out. A uniform cap estimate covers
\emph{all} conditional minima, including degenerate and
nonunique ones. Near coherence, V31--V32 remove the constant
normalization loss. The native curvature witness then pays
the remaining error with the minimized energy itself.
There is no additive block-count term and no bad-event deletion.

The result is a global pointwise comparison, uniform in
the torus size and in every retained configuration:
the exact action differs from coupling times the partial
minimum by a relative \(O(\log\gamma/\gamma)\).
Near identity the logarithmic loss disappears. This
identifies the exact Maxwell Schur form as the leading
native Hessian, with relative \(O(\gamma^{-1})\) error,
even on arbitrarily long tori.

V29 already proves Maxwell coercivity, the native Hessian
kernel and the full flat minimum. V30 already constructs
the nonlinear filling. V31 already proves the \(45\)-face
witness and the near-coherent minimum; V32 already proves
the differentiable remainder. These are inputs, not
rediscoveries. Archive \(C\), V33 distinguishes relative
curvature control from an absolute Hessian error, and
\(C\), V35 obtains a relative bound only on a
duration-dependent comparison ball for a different
endpoint-fixed law. Neither supplies the present
all-configuration partial-minimum comparison.

A targeted literature check also distinguishes this
pointwise question from free-energy asymptotics.
Chatterjee,
\href{https://arxiv.org/abs/1602.01222}
{\emph{The leading term of the Yang--Mills free energy}},
treats the leading weak-coupling free energy.
Brennecke's 2026 paper,
\href{https://doi.org/10.1007/s10955-026-03649-4}
{\emph{On the Leading Order Term of the Lattice
Yang--Mills Free Energy}}, explicitly evaluates its
lattice-Maxwell constant. Their primary abstracts and
the latter's introductory setup were checked. No
conditional-law, pointwise-action or mass-gap theorem
is imported from those statements. The next scheduled
companion and broad-literature reviews remain V35 and V40.

\subsection{The exact partial minimum and its paid defect}

Use V29's packing for every \(L\in4\mathbb N\), \(L\ge4\),
\(q\ge3\), not only dyadic sizes. Define
\[
 \mathcal I(\eta)
   =\min_{U:\,U|_{\mathcal R}=\eta}E_W(U)
   =S_0(\eta)+\sum_b m_b(\eta),\qquad
 m_b(\eta)=\min_{z\in(S^3)^6}\mathcal E_{\vartheta_b(\eta)}(z).
 \tag{V33.1}\label{r33:I}
\]
Every minimum exists by compactness. The second equality
holds because no original plaquette contains free links
from two blocks. It does not make the retained measure
a product. Globally \(m_b\) means the minimum \emph{value};
no smooth minimizing selection is asserted outside the
near-coherent region.

Write \(\mathcal W_b\) for V31's \(45\)-plaquette witness
and \(\rho_b\) for V32's Frobenius coefficient defect.
Every witness plaquette's initial vertex lies in
V32's \(120\)-vertex box \eqref{r32:box}. Therefore
\[
 \kappa_W:=\sup_p\#\{b:p\in\mathcal W_b\}\le120 .
 \tag{V33.2}\label{r33:coverage}
\]
Indeed a fixed initial vertex permits at most the \(120\)
anchor displacements in that box, and a fixed oriented
plaquette has only that one initial vertex. The packing
uses a subset of those anchors; periodic identifications
can only reduce the count.

\begin{proposition}[The minimized energy pays all local defects]
\label{r33:budget}
For every retained configuration,
\[
 \sum_b\rho_b^2\le80\kappa_W\mathcal I(\eta),\qquad
 \frac16 E_W(\Phi\eta)\le\mathcal I(\eta)\le E_W(\Phi\eta).
 \tag{V33.3}\label{r33:paid}
\]
The partial minimum is continuous and gauge invariant.
Its zero set consists exactly of retained restrictions
of full flat configurations.
\end{proposition}
\begin{proof}
For \emph{every} full extension \(U\), V31 and the
matrix norm comparison give
\[
 \rho_b^2\le4D_b^2
                    \le80\sum_{p\in\mathcal W_b}e(U_p).
\]
Summing and using \eqref{r33:coverage}, then minimizing
over \(U\), proves the first assertion. No individual
witness is minimized separately.

For the second, the filled configuration is an admissible
extension, which gives the upper bound. At a site with
\(n_i\le4\) private slots, their unit vectors satisfy
\[
 \sum_{\alpha<\beta}(1-p_\alpha\cdot p_\beta)
 =\frac{n_i}{2}\sum_\alpha|p_\alpha-\overline p|^2
 \le\frac{n_i}{2}\sum_\alpha|p_\alpha-U_i|^2
 \le4\sum_\alpha e(U_i p_\alpha^{-1}).
\]
V30 gives \(E_W(\Phi\eta)\le6S_0+S_{\rm slot}\).
Thus for every extension,
\[
 E_W(\Phi\eta)\le6S_0+4E_{\rm private}(U)
                                      \le6E_W(U).
\]
Take the minimum. Continuity of a minimum over a fixed
compact fibre follows directly from uniform continuity
of the energy. Gauge transformations biject the fibres
and preserve energy. Finally a sum of nonnegative
plaquette energies vanishes exactly when every plaquette
holonomy is \(1\); attainment of the minimum proves the
zero-set assertion.
\end{proof}

\subsection{All-exterior integration without a nondegenerate saddle}

For arbitrary unit slots and orthogonal transports, put
\[
 Q_\gamma(\vartheta)=\int e^{-\gamma\mathcal E_\vartheta(z)}
                                      \prod_{i=1}^6d\sigma(z_i),
 \qquad m(\vartheta)=\min_z\mathcal E_\vartheta(z).
\]
Let \(o\) denote coherent coefficients and set
\[
 c_H=\frac8{3\pi^3},\quad
 \kappa=c_H^6e^{-27},\quad
 C_0=\left[\frac1{2\sqrt\pi}
                     \left(\frac{\pi^2}{6}\right)^{3/2}\right]^6 .
 \tag{V33.4}\label{r33:constants}
\]
These explicit constants concern only the next cap bounds.

\begin{lemma}[A uniform cap at any minimizing configuration]
\label{r33:cap}
For every \(\vartheta\) and \(\gamma\ge1\),
\[
 \kappa\gamma^{-9}e^{-\gamma m(\vartheta)}
       \le Q_\gamma(\vartheta)\le e^{-\gamma m(\vartheta)},
 \qquad
 \kappa\gamma^{-9}\le Q_\gamma(o)\le C_0\gamma^{-9}.
 \tag{V33.5}\label{r33:caps}
\]
No uniqueness, positivity of the minimum Hessian, or
restriction on the exterior is required.
\end{lemma}
\begin{proof}
The exact product-geodesic Hessian formula
\eqref{r31:hessian} has the global upper bound
\[
 \operatorname{Hess}\mathcal E_\vartheta(z)[\xi,\xi]
 \le\sum_i(n_i+2\deg i)|\xi_i|^2\le9\sum_i|\xi_i|^2.
\]
Here every contact scalar is at most one and
\(2|\xi_i||\xi_j|\le|\xi_i|^2+|\xi_j|^2\).
At any global minimizer the first derivative is zero.
Taylor's formula along product geodesics therefore gives
\(\mathcal E_\vartheta(z)-m\le(9/2)\sum_i\theta_i^2\).
Restrict each angular displacement to
\(\theta_i\le\gamma^{-1/2}\le1\); the excess is at most
\(27/\gamma\). A normalized Haar cap has mass
\[
 \frac2\pi\int_0^r\sin^2\theta\,d\theta
 \ge\frac8{3\pi^3}r^3\quad(0\le r\le1),
\]
using \(\sin\theta\ge2\theta/\pi\).
The six caps prove the lower bound. The upper bound
uses only \(\mathcal E_\vartheta\ge m\) and normalized Haar.

At coherence \(m(o)=0\), and all seven internal energies
are nonnegative. With \(\theta_i\) the angle from \(1\),
\[
 \mathcal E_o(z)\ge3\sum_i(1-\cos\theta_i)
                     \ge\frac6{\pi^2}\sum_i\theta_i^2.
\]
Since \(\sin^2\theta\le\theta^2\), each radial integral
is at most
\[
 \frac2\pi\int_0^\infty
       \theta^2e^{-(6\gamma/\pi^2)\theta^2}\,d\theta
 =\frac1{2\sqrt\pi}
                     \left(\frac{\pi^2}{6\gamma}\right)^{3/2}.
\]
Taking the sixth power proves the final upper bound.
In particular \(0<\kappa\le C_0\).
\end{proof}

Define the centered, exact elimination error
\[
 g_\gamma(\vartheta)
   =-\log\frac{Q_\gamma(\vartheta)}{Q_\gamma(o)}
                                       -\gamma m(\vartheta).
 \tag{V33.6}\label{r33:g}
\]
The cap bounds immediately give
\[
 -|\log\kappa|-9\log\gamma
        \le g_\gamma(\vartheta)\le\log(C_0/\kappa).
 \tag{V33.7}\label{r33:rough}
\]
An additive use of these inequalities over all blocks
would still be extensive. The next step removes that loss.

\begin{lemma}[A curvature-relative error at every exterior]
\label{r33:relative-block}
There is a finite constant \(C\), depending only on the
six-spin graph, such that for every \(\gamma\ge1\)
and every coefficient configuration,
\[
 -C(1+\log\gamma)\rho(\vartheta)^2
       \le g_\gamma(\vartheta)\le C\rho(\vartheta)^2 .
 \tag{V33.8}\label{r33:relative-block-eq}
\]
On \(\rho\le1/128\), the stronger
\(\lvert g_\gamma\rvert\le C\rho^2\) holds.
This constant is proved finite, not numerically evaluated.
\end{lemma}
\begin{proof}
On that small neighborhood, \(|a|\le2\rho\le1/64\),
and V32's exact identity gives
\[
 g_\gamma(\vartheta)
 =\tfrac12\log\frac{\det M(\vartheta)}{\det M(o)}
                                      -h_\gamma(\vartheta).
\]
V31's stationary determinant and V32's centered
\(C^2\) remainder bound each term by \(C\rho^2\),
uniformly for \(\gamma\ge1\). Outside that neighborhood,
\(\rho^2>128^{-2}\), so \eqref{r33:rough} implies
\eqref{r33:relative-block-eq} after enlarging \(C\).
No differentiability of \(m\) outside the small
neighborhood is used.
\end{proof}


\subsection{The volume-uniform nonlinear action comparison}

As before let \(W_\gamma(\eta)=V_{\mathcal R}(\eta)-V_{\mathcal R}(1)\).
All blocks have the same coherent factor, so exactly
\[
 W_\gamma(\eta)-\gamma\mathcal I(\eta)
                         =\sum_b g_\gamma(\vartheta_b(\eta)).
 \tag{V33.9}\label{r33:exact-difference}
\]

\begin{theorem}[Energy-relative elimination without exceptional sets]
\label{r33:global}
There is a finite \(C_*\), independent of \(L,q,\gamma\),
such that for every retained configuration and
\(\gamma\ge1\),
\[
 \boxed{\quad
 [\gamma-C_*(1+\log\gamma)]\mathcal I(\eta)
 \le W_\gamma(\eta)
 \le(\gamma+C_*)\mathcal I(\eta).
 \quad}
 \tag{V33.10}\label{r33:global-bound}
\]
On configurations with \(\rho_b\le1/128\) for every block,
\[
 |W_\gamma(\eta)-\gamma\mathcal I(\eta)|
                                      \le C_*\mathcal I(\eta).
 \tag{V33.11}\label{r33:near-bound}
\]
In particular, uniformly over all admitted tori and
all configurations with \(\mathcal I(\eta)>0\),
\[
 \left|\frac{W_\gamma(\eta)}{\gamma\mathcal I(\eta)}-1\right|
                       \le\frac{C_*(1+\log\gamma)}{\gamma}.
 \tag{V33.12}\label{r33:uniform-ratio}
\]
\end{theorem}
\begin{proof}
Sum \eqref{r33:relative-block-eq} in
\eqref{r33:exact-difference} and apply
\(\sum_b\rho_b^2\le80\kappa_W\mathcal I\).
Absorb the fixed factor \(80\kappa_W\) in \(C_*\).
If all coefficients are in the smaller region, use the
stronger two-sided local bound without \(\log\gamma\).
Dividing by \(\gamma\mathcal I>0\) proves the ratio bound.
If \(\mathcal I=0\), the defect budget forces every
\(\rho_b=0\), and both actions are exactly zero.
\end{proof}

The first lower coefficient need not be positive at
moderate coupling; V29's nonnegative action remains
available there. The theorem is not obtained by asserting
that most blocks are good. It covers configurations with
arbitrarily many rough blocks and all winding sectors.
The conservative constants are large and unevaluated
because their near-coherent part uses V31--V32.
This is uniformity of an \emph{energy ratio}, not
uniformity of normalized correlations.

\subsection{The true Maxwell Schur form, with relative error}

Let \(d_1\) be the original plaquette curl on real
three-color link cochains. Split its columns as
\(d_1=(D_R,D_F)\) according to retained and free links.
Define the partial-minimum quadratic form
\[
 \begin{aligned}
 \mathcal S(a)&=\min_z\|D_Ra+D_Fz\|_2^2
                           =\langle a,\mathsf S_Ra\rangle,\\
 \mathsf S_R
 &=D_R^*D_R-D_R^*D_F(D_F^*D_F)^{-1}D_F^*D_R .
 \end{aligned}
 \tag{V33.13}\label{r33:Schur}
\]
The inverse is an ordinary inverse, not a
pseudoinverse. Up to the color factor,
\[
 D_F^*D_F=\bigoplus_b\mathsf A,\qquad
 \mathsf A=6I_6-\operatorname{Adj}(P_2\times P_3),
 \qquad\det\mathsf A=37835>0 .
 \tag{V33.14}\label{r33:A}
\]
Every free link meets six plaquettes and internal
two-free-link plaquettes give off-diagonal \(-1\).
Different blocks have no shared plaquette.

\begin{theorem}[A sharp relative coherent Hessian limit]
\label{r33:Hessian}
The partial minimum is smooth near the all-identity
retained configuration and
\[
 \operatorname{Hess}\mathcal I(1)[a,a]=\mathcal S(a).
 \tag{V33.15}\label{r33:I-Hess}
\]
For a finite \(C_*\) independent of the torus size,
\[
 (\gamma-C_*)\mathcal S(a)
 \le\operatorname{Hess}V_{\mathcal R}(1)[a,a]
 \le(\gamma+C_*)\mathcal S(a)
 \qquad(\gamma\ge1).
 \tag{V33.16}\label{r33:relative-Hess}
\]
No condition such as \(\gamma\gg\ell_{\max}^2\)
is required for this relative comparison.
\end{theorem}
\begin{proof}
V31 gives a unique smooth minimizer for each free block
in a neighborhood of identity. Their product is the
global minimizer in the free fibre, so \(\mathcal I\)
is smooth there. At identity its free Hessian is
\(D_F^*D_F\), positive by \eqref{r33:A}.
Differentiating the stationary equation for the free
minimum gives the first-order minimizing extension
\[
 z_*(a)=-(D_F^*D_F)^{-1}D_F^*D_Ra.
\]
Since the full Wilson Hessian is \(d_1^*d_1\),
substitution gives \eqref{r33:I-Hess}; equivalently
one completes the positive free-variable square.
The minimum and first derivatives at identity vanish.

For any fixed tangent \(a\), the curve
\(\eta_e(s)=\exp(sa_e)\) lies in
\(\rho_b\le1/128\) for every block when \(s\) is
sufficiently small. That permissible interval may depend
on \(a\) and the finite lattice; the coefficient in
\eqref{r33:near-bound} does not.
Both sides vanish to second order at identity.
Divide that inequality by \(s^2/2\) and let \(s\to0\).
This gives the relative bound with exactly the same
volume-independent constant. An absolute operator
error followed by division by the smallest eigenvalue
has not been used.
\end{proof}

This identifies the minimizing extension, unlike V29's
declared slot-average extension. It sharpens the
large-\(\gamma\) comparison coefficient to the correct
leading coefficient one against the actual Schur form.
It is still a Hessian assertion at one configuration,
not a curvature lower bound throughout a retained
probability space.

\begin{proposition}[The exact soft scale of the Schur form]
\label{r33:soft}
Let \(\ell_{\max}=\max\{L,2q\}\), and let
\(\mathcal K_R\) be V29's restrictions of closed full
cochains. Then
\[
 \ker\mathsf S_R=\mathcal K_R,\qquad
 \lambda_{\min}^{+}(\mathsf S_R)
                   =4\sin^2(\pi/\ell_{\max})=:\lambda_\Lambda .
 \tag{V33.17}\label{r33:soft-eq}
\]
Consequently, for \(\gamma>C_*\), the smallest positive
eigenvalue of the actual identity Hessian obeys
\[
 (\gamma-C_*)\lambda_\Lambda
 \le\lambda_{\min}^{+}(\operatorname{Hess}V_{\mathcal R}(1))
 \le(\gamma+C_*)\lambda_\Lambda .
 \tag{V33.18}\label{r33:Hsoft}
\]
These eigenvalues use the retained product metric and
remove its full closed-cochain kernel.
\end{proposition}
\begin{proof}
The minimum squared curl vanishes exactly when \(a\)
has a closed full extension, proving the kernel.
For any extension \(u\), periodic Fourier decomposition
gives V29's inequality
\[
 \|d_1u\|^2\ge
 \lambda_\Lambda\operatorname{dist}(u,\ker d_1)^2
 \ge\lambda_\Lambda\operatorname{dist}(a,\mathcal K_R)^2 .
\]
Minimizing over the free components proves the lower
spectral bound.

For the matching upper bound choose a longest direction
\(k\) among \(\nu,\rho,t\), of period \(\ell_{\max}\),
and choose a different \(j\) among those three directions.
All links in direction \(j\) are retained. For a fixed
unit color vector \(v\), set
\[
 a_j(x)=v\cos(2\pi x_k/\ell_{\max}),\qquad
 a_\alpha(x)=0\quad(\alpha\ne j).
\]
Its full extension with free \(\mu\)-components zero
has curl only in the \(j,k\) plane. Free columns involve
only \(\mu\)-plaquettes, so \(D_F^*D_Ra=0\):
zero is exactly the optimal free extension.
The one-dimensional difference identity gives
\(\mathcal S(a)=\lambda_\Lambda\|a\|^2\).
Discrete summation by parts in direction \(j\) shows
orthogonality to every restricted gradient, and the
zero mean in direction \(k\) shows orthogonality to
the four harmonic constants. Thus \(a\perp\mathcal K_R\),
and the lower bound is attained. Finally use
\eqref{r33:relative-Hess} on \(\mathcal K_R^\perp\);
V29 identifies the actual Hessian's kernel.
\end{proof}

At any fixed sufficiently large \(\gamma\), this
coherent action-Hessian scale goes to zero as the
longest period grows. That is not a proof of a
gapless physical theory. It does show why improving
the local Hessian approximation cannot by itself
supply a volume-uniform physical mass.

\subsection{What the global comparison still does not transport}

Write \(\Delta_\gamma=W_\gamma-\gamma\mathcal I\)
and normalize
\(d\nu_\gamma=Z_{\mathcal I}(\gamma)^{-1}
e^{-\gamma\mathcal I}\,d\eta\).
The exact probability identity is
\[
 \frac{d\mu_R}{d\nu_\gamma}
   =\frac{e^{-\Delta_\gamma}}
                    {\nu_\gamma(e^{-\Delta_\gamma})}.
 \tag{V33.19}\label{r33:law}
\]
Although \(\Delta_\gamma\) is small \emph{relative to}
\(\gamma\mathcal I\), it can be extensive.
The normalizer in \eqref{r33:law} cannot be omitted.
Pointwise ordered actions do not generally order
connected correlations or transfer-operator gaps.

Thus rough conditional minima no longer obstruct the
energy comparison. The remaining task is a normalized,
spatially or temporally resolved estimate for the
nonlinear partial-minimum law and its paid correction,
or a direct physical-transfer argument avoiding such
a comparison. Noncommuting winding modes, their
nonquadratic energy, and the compact measure remain
present. Neither the uniform energy ratio nor the
sharp Schur limit constructs the interacting
four-dimensional continuum theory or its mass gap.

The diagnostic checks the witness coverage and exact
linear minimizing equations on four finite tori,
with rational retained cochains. It also checks the
filled-energy comparison and the soft-mode difference
identity in exact cyclotomic arithmetic. These are
geometric and algebraic checks, not numerical
evaluation of \(C_*\) or certification of a physical
spectrum. V32 is recoverable byte-for-byte after
removing this append and restoring the title.
All other manuscript sources are preserved, and
all build outputs remain outside \path{latex_notes}.

% END CONSOLIDATED V33 APPEND

% BEGIN CONSOLIDATED V34 APPEND -- 2026-09-19
\clearpage
\section{V34: Strict saddles, normalized conditional response, and native innovations}
\label{r34:section}

\subsection{Purpose, provenance, and the remaining distinction}

V33 compared the actual retained action with the partial minimum of
the original Wilson energy. Its relative error is small, but can
still be extensive. Neither that comparison nor a positive Hessian
at a conditional minimum bounds variance under the whole compact
conditional law: a competing local minimum could produce a slow
mode. This append supplies the missing global conditional estimate
for the actual six-link block, without deleting any spin configurations.
It then controls arbitrary sums of \emph{conditionally centered}
block fluctuations under the original normalized Wilson law.
The retained conditional predictors are not centered away from
the physical problem; their unresolved covariance is displayed explicitly.

The graph, frames, coefficient witness, and local convexity are
inherited from V29--V33. The new graph-specific input is that every
critical point other than the pinned minimum has a uniformly
negative one-site Hessian direction. A ground-state transformation
and a partition-of-unity identity turn this into a Poincare inequality
on the full compact block. These analytic tools are standard, not
new claims about Witten theory. In particular, the targeted sources
\[
 \text{\url{https://arxiv.org/abs/1607.08714}}
 \qquad\text{and}\qquad
 \text{\url{https://arxiv.org/html/2410.21899}}
\]
are respectively Le Peutrec's boundary Witten/Brascamp--Lieb
framework and Delande's degenerate Witten analysis.
The former supplies methodological context; the latter's isolated
critical-point hypotheses are not assumed here. The needed estimates
are proved below, including the boundary sign and localization error.
No external spectral asymptotics are imported.

Archive $A$, V56 already recorded conditional-predictor covariance
decomposition; archive $C$ contains conditional and Witten routes.
We use that decomposition rather than claim it anew.
The new result is the verified native hypothesis and its uniform
application to normalized six-spin laws and to the original joint
law. The next scheduled companion review is V35; the next broad
primary-literature sweep is V40.

\subsection{The compact coefficient family and a strict-saddle theorem}

Let $\mathcal M=(S^3)^6$ carry the product unit-round metric and
normalized product Haar measure $\sigma$. Write $e_0=1\in S^3$.
The graph is $P_2\times P_3$, with vertices $(c,r)$,
$c\in\{0,1\}$, $r\in\{0,1,2\}$. There are four corner vertices,
each with degree two and four unit slots, and two middle vertices,
each with degree three and three unit slots. Each middle vertex
is adjacent to two corners and the other middle vertex.
There are $22$ slots and seven edges. Thus
\[
 n_i+\deg(i)=6,\qquad n_i\ge3.
\]
For unit slots $v_{i\alpha}$ and orthogonal edge transports $T_{ij}$
use V31's energy and defect
\[
 \begin{aligned}
 E_\vartheta(z)
   &=\sum_{i,\alpha}(1-z_i\cdot v_{i\alpha})
       +\sum_{\{i,j\}}(1-z_i\cdot T_{ij}z_j),\\
 d(\vartheta)^2
   &=\sum_{i,\alpha}|v_{i\alpha}-e_0|^2
       +\sum_{\{i,j\}}\|T_{ij}-I\|_{\rm op}^2,
 \qquad \delta_0=\frac1{32}.
 \end{aligned}
 \tag{V34.1}\label{r34:energy}
\]
The native quaternion transports form a subfamily of this compact
orthogonal family. An edge read at its other endpoint uses $T_{ij}^*$.
All statements in this subsection hold for every coefficient family
with $d(\vartheta)\le\delta_0$, not just for almost every exterior.

\begin{theorem}[No competing conditional local minima]
\label{r34:saddle}
For $d(\vartheta)\le1/32$, V31's unique global minimum $z_*(\vartheta)$
is the only local minimum and the only critical point in the
product of angular caps of radius $1/4$ about $e_0$.
Every other critical point has a unit tangent vector supported at
one vertex such that
\[
 \operatorname{Hess} E_\vartheta(z)[\xi,\xi]\le-\frac{29}{16}.
 \tag{V34.2}\label{r34:negative}
\]
The other critical points need not be isolated or nondegenerate.
\end{theorem}
\begin{proof}
Define the local field, with the endpoint convention above, by
\[
 h_i=\sum_\alpha v_{i\alpha}+\sum_{j\sim i}T_{ij}^{(i)}z_j,
 \qquad h_i^0=n_i e_0+\sum_{j\sim i}z_j .
\]
Each individual coefficient defect is at most $d$, and hence
\[
 |h_i-h_i^0|\le6d\le\frac3{16},\qquad |h_i|\le6.
 \tag{V34.3}\label{r34:field}
\]
Holding the other spins fixed, the nonconstant energy is $-z_i\cdot h_i$.
At a critical point, $h_i$ is parallel to $z_i$ or zero; the
one-site Hessian is $(z_i\cdot h_i)|\xi_i|^2$.

For a corner,
\[
 h_i\cdot e_0\ge4-2-\frac3{16}=\frac{29}{16}.
\]
Its field cannot vanish. If any corner is anti-aligned with its
field, the one-site Hessian is $-|h_i|\le-29/16$, as required.
Otherwise all four corners are aligned and
\[
 z_i\cdot e_0=\frac{h_i\cdot e_0}{|h_i|}
                       \ge\frac{29}{96}
       \quad\hbox{at every corner}.
\]
Each middle vertex now satisfies
\[
 h_i\cdot e_0
   \ge3+2\frac{29}{96}-1-\frac3{16}=\frac{29}{12}.
\]
Its field cannot vanish either. If it is anti-aligned, its
one-site Hessian is at most $-29/12$. We are therefore done
unless all six spins are aligned with their nonzero fields.
In this remaining case all $z_i\cdot e_0$ are positive.

Choose a vertex whose height $c=z_i\cdot e_0$ is smallest,
and set $u=e_0-cz_i\in T_{z_i}S^3$. If $c=1$ there is
nothing to localize. Otherwise stationarity gives $h_i\cdot u=0$,
whereas
\[
 \begin{aligned}
 h_i^0\cdot u
  &=n_i(1-c^2)+\sum_{j\sim i}(z_j\cdot e_0-c\,z_i\cdot z_j)\\
  &\ge n_i(1-c^2).
 \end{aligned}
\]
Here $z_j\cdot e_0\ge c>0$ and $z_i\cdot z_j\le1$.
Using \eqref{r34:field} and $|u|=\sqrt{1-c^2}$ gives
\[
 \sqrt{1-c^2}\le\frac{6d}{n_i}\le2d\le\frac1{16}.
\]
The maximum polar angle is thus at most $\arcsin(1/16)<1/4$.
V31's Hessian lower bound $75/32$ throughout this geodesically
convex product of caps makes the gradient have at most one zero
there. Indeed, two distinct zeros would contradict strict positivity
of the second derivative along their connecting geodesic.
V31 locates $z_*$ inside these caps. Thus the all-aligned critical
point is exactly $z_*$. All other cases have already supplied
\eqref{r34:negative}.
\end{proof}

This theorem concerns the \emph{pinned conditional block}, not the
full retained energy. It does not remove V29's flat-family
degeneracies or V33's long-wavelength Schur modes.

\subsection{A full-law Poincare estimate without a spin-domain cutoff}

Let $m_\vartheta=\min E_\vartheta$ and define the probability measure
\[
 d\lambda_{\gamma,\vartheta}
   =Q_\gamma(\vartheta)^{-1}e^{-\gamma E_\vartheta}\,d\sigma,
 \qquad Q_\gamma(\vartheta)=\int_{\mathcal M}e^{-\gamma E_\vartheta}\,d\sigma .
 \tag{V34.4}\label{r34:lambda}
\]
The domain of this law is always all of $\mathcal M$.

\begin{theorem}[Uniform normalized conditional variance]
\label{r34:poincare}
There are finite constants $\Gamma_*\ge1$ and $C_P$, depending
only on the six-vertex graph and $\delta_0$, such that for
$d(\vartheta)\le\delta_0$, $\gamma\ge1$, and real
$F\in H^1(\mathcal M)$,
\[
 \operatorname{Var}_{\lambda_{\gamma,\vartheta}}F
       \le \frac{C_P}{\gamma}
                         \lambda_{\gamma,\vartheta}(|\nabla F|^2).
 \tag{V34.5}\label{r34:poincare-eq}
\]
For $\gamma\ge\Gamma_*$ one can replace $C_P$ by $4$.
These constants are independent of exterior volume and position.
Their finite existence is proved; no numerical value of $\Gamma_*$
or moderate-coupling guarantee is asserted.
\end{theorem}
\begin{proof}
We give the normalization and localization details. Put
\[
 R(z)=\left(\sum_i\theta_i(z_i)^2\right)^{1/2},
 \qquad \theta_i(z_i)=\operatorname{dist}_{S^3}(z_i,e_0),
 \qquad \Omega=\{R<1/4\}.
\]
Only $R^2$ in the injectivity neighborhood and smooth functions
constant near its center are used as smooth functions.
V31 gives, with $a=(\sum_i|z_i-e_0|^2)^{1/2}$,
\[
 R(z_*)\le\frac{\pi}{2}a(z_*)\le\frac{\pi}{32}<\frac18,
 \qquad
 E_\vartheta(z)-m_\vartheta\ge\frac32a^2-\frac3{32}a.
\]
For $R\ge1/8$ we have $a\ge(2/\pi)R\ge1/(4\pi)$.
The displayed quadratic is increasing on that interval, so
\[
 E_\vartheta(z)-m_\vartheta\ge b_*:=\frac{3(4-\pi)}{128\pi^2}>0
                      \quad(R\ge1/8).
\]
V33's global cap lower bound, with
$\kappa=(8/(3\pi^3))^6e^{-27}$, proves
\[
 \lambda_{\gamma,\vartheta}\{R\ge1/8\}
       \le t_\gamma:=\min\{1,\kappa^{-1}\gamma^9e^{-b_*\gamma}\}.
 \tag{V34.6}\label{r34:tail}
\]
The denominator $Q_\gamma$ has been included. This is a
conditional probability estimate obtained from the block integral,
not an application of the original-law chessboard estimate to a
changed measure.

On $\Omega$, $\operatorname{Hess}E_\vartheta\ge2g$ and
$\operatorname{Ric}=2g$. Its smooth boundary is convex:
on $S^3$, the eigenvalues of
$\operatorname{Hess}(\theta^2/2)$ are $1$ radially and
$\theta\cot\theta$ transversely. Their product-space sum
is positive for $R\le1/4$, making the outward second
fundamental form of the level set nonnegative.
For the probability law $\lambda^\Omega$ obtained by restricting
and normalizing $\lambda_{\gamma,\vartheta}$, the Neumann
operator $L=\Delta-\gamma\nabla E_\vartheta\cdot\nabla$ obeys
the integrated Bochner identity
\[
 \begin{aligned}
 \lambda^\Omega((LF)^2)
  ={}&\lambda^\Omega\!\left(
     |\operatorname{Hess}F|^2+
     (\operatorname{Ric}+\gamma\operatorname{Hess}E_\vartheta)
                                      [\nabla F,\nabla F]\right)\\
     &+\int_{\partial\Omega}
          \mathrm{II}(\nabla_{\partial}F,\nabla_{\partial}F)\,
                                      d\lambda^\Omega_{\partial}\\
  \ge{}&(2\gamma+2)\lambda^\Omega(|\nabla F|^2).
 \end{aligned}
 \tag{V34.7}\label{r34:bochner}
\]
Here $d\lambda^\Omega_{\partial}$ is the boundary density
with the same normalization, and $F$ initially satisfies
$\partial_nF=0$. The identity follows by integrating
$\tfrac12 L|\nabla F|^2
 =|\operatorname{Hess}F|^2+
 \langle\nabla F,\nabla LF\rangle+
 (\operatorname{Ric}+\gamma\operatorname{Hess}E_\vartheta)
 [\nabla F,\nabla F]$ and using
$\partial_n|\nabla F|^2=-2\mathrm{II}(\nabla_\partial F,\nabla_\partial F)$.
For a nonconstant Neumann eigenfunction $-LF=\ell F$,
integration by parts gives
$\lambda^\Omega(|\nabla F|^2)=\ell\lambda^\Omega(F^2)$.
Consequently \eqref{r34:bochner} implies $\ell\ge2\gamma+2$.
The Neumann eigenfunction expansion on this smooth bounded
domain gives the same Poincare inequality for every $H^1$
function on $\Omega$. This is a local ingredient, not a
replacement of the full law by the restricted one.

Let
\[
 \psi=Q_\gamma(\vartheta)^{-1/2}e^{-\gamma E_\vartheta/2}.
\]
On $L^2(\sigma)$, with no boundary, introduce the nonnegative form
\[
 \begin{aligned}
 q_\gamma(u)
  &=\int\left|\nabla u+\frac{\gamma}{2}u\nabla E_\vartheta\right|^2d\sigma\\
  &=\int\left\{|\nabla u|^2+
       \left(\frac{\gamma^2}{4}|\nabla E_\vartheta|^2
                         -\frac{\gamma}{2}\Delta E_\vartheta\right)
                                      |u|^2\right\}d\sigma .
 \end{aligned}
 \tag{V34.8}\label{r34:form}
\]
Then $\|\psi\|_2=1$, $q_\gamma(\psi F)=\lambda(|\nabla F|^2)$,
and $\psi F\perp\psi$ exactly when $\lambda(F)=0$.

We now justify uniform localization outside $R<1/8$.
The coefficient family $K=\{\vartheta:d(\vartheta)\le1/32\}$
and $K\times\{R\ge1/8\}$ are compact. At every critical point
in the latter set, Theorem~\ref{r34:saddle} supplies a unit
direction of Hessian at most $-29/16$. Extend that direction
to a smooth unit vector field $V$ in a coordinate neighborhood.
At the critical point,
\[
 \operatorname{div}((VE_\vartheta)V)
                   =\operatorname{Hess}E_\vartheta[V,V]\le-29/16.
\]
By continuity in position and coefficients, shrink the neighborhood
so that the divergence is at most $-1$. At every other point
a neighborhood has $|\nabla E_\vartheta|$ bounded below by a
positive number. A finite subcover therefore consists of
saddle patches of the first type and gradient patches on
which $|\nabla E_\vartheta|\ge g_0>0$. Neither isolated
critical points nor a nonsingular full Hessian is needed.

For $u$ supported in a saddle patch, $|V|\le1$ and
integration by parts give
\[
 \begin{aligned}
 q_\gamma(u)
 &\ge\int\left|Vu+\frac{\gamma}{2}(VE_\vartheta)u\right|^2d\sigma\\
 &=\int\left\{|Vu|^2+
       \frac{\gamma^2}{4}(VE_\vartheta)^2u^2
       -\frac{\gamma}{2}\operatorname{div}((VE_\vartheta)V)u^2\right\}d\sigma\\
 &\ge\frac{\gamma}{2}\|u\|_2^2 .
 \end{aligned}
 \tag{V34.9}\label{r34:saddle-form}
\]
The unit sphere coordinate identity $\Delta z_a=-3z_a$
shows $|\Delta E_\vartheta|\le3\cdot22+6\cdot7=108$ globally.
Thus on a gradient patch,
\[
 q_\gamma(u)\ge
       \left(\frac{\gamma^2g_0^2}{4}-54\gamma\right)\|u\|_2^2
       \ge\frac{\gamma}{2}\|u\|_2^2
 \tag{V34.10}\label{r34:gradient-form}
\]
once $\gamma$ exceeds a fixed finite threshold.

Choose a smooth $0\le\chi_0\le1$, equal to one for $R\le1/8$
and supported in $R<3/16$, and complete it to a real smooth
partition $\sum_j\chi_j^2=1$ with all other terms supported
in the finite patch cover. This can be done with uniformly
bounded position derivatives for all $\vartheta\in K$.
For example choose $\chi_0=\cos\alpha(R)$, with $\alpha=0$
on the inner ball and $\alpha=\pi/2$ beyond $3/16$,
and multiply a normalized squared partition on the exterior
cover by $\sin\alpha(R)$. Smooth parameter charts on a
neighborhood of the compact coefficient set and a finite
subcover give one common bound
\[
 \sum_j|\nabla\chi_j|^2\le C_{\rm IMS}<\infty .
\]
If $u\perp\psi$, the local Poincare inequality on $\Omega$,
applied to $(\chi_0u)/\psi$, gives
\[
 \begin{aligned}
 q_\gamma(\chi_0u)
 &\ge2\gamma\left(\|\chi_0u\|_2^2
             -\frac{|\langle\chi_0u,\psi\rangle|^2}{\lambda(\Omega)}\right)\\
 &\ge2\gamma\|\chi_0u\|_2^2-4\gamma t_\gamma\|u\|_2^2
                                      \quad(t_\gamma\le1/2).
 \end{aligned}
 \tag{V34.11}\label{r34:projection}
\]
Indeed $\langle\chi_0u,\psi\rangle
=\langle u,(\chi_0-1)\psi\rangle$, whose squared magnitude
is at most $t_\gamma\|u\|_2^2$, while
$\lambda(\Omega)\ge1-t_\gamma$.

Expanding the gradients and using
$\sum_j\chi_j\nabla\chi_j=0$ proves the exact localization identity
\[
 q_\gamma(u)=\sum_j q_\gamma(\chi_ju)
                  -\int\sum_j|\nabla\chi_j|^2|u|^2d\sigma .
\]
Combining it with \eqref{r34:saddle-form}--\eqref{r34:projection}
yields
\[
 q_\gamma(u)\ge
       \left(\frac{\gamma}{2}-C_{\rm IMS}-4\gamma t_\gamma\right)\|u\|_2^2
       \ge\frac{\gamma}{4}\|u\|_2^2,\qquad u\perp\psi,
 \tag{V34.12}\label{r34:gap}
\]
for all $\gamma\ge\Gamma_*$, with one finite $\Gamma_*$.
Take $u=\psi(F-\lambda F)$ to obtain the claimed constant $4$.

For completeness, the initial interval costs only a finite
uniform constant. Every energy summand lies in $[0,2]$,
so $\operatorname{osc}E_\vartheta\le58$.
Product Haar has Poincare gap $3$, from the first spherical
eigenvalue and tensorization. The ratio between the largest
and smallest density of $\lambda$ relative to Haar is at
most $e^{58\gamma}$. Comparing the variance with its Haar
centering and then comparing Dirichlet integrals gives
$\operatorname{Var}_{\lambda}F
 \le e^{58\gamma}\lambda(|\nabla F|^2)/3$.
Thus a possible finite choice is
\[
 C_P=\max\left\{4,\frac{\Gamma_*}{3}e^{58\Gamma_*}\right\}.
 \tag{V34.13}\label{r34:CP}
\]
Smooth approximation proves the asserted $H^1$ statement.
\end{proof}

The gap in \eqref{r34:gap} is the gap of an
\emph{auxiliary block diffusion form}. It is not a physical
transfer eigenvalue or a Hamiltonian mass, and its time
parameter is not the physical Euclidean time.

\subsection{Normalized exterior response with no coupling-factor loss}

For a differentiable coefficient curve $\vartheta(s)$ in $K$ put
\[
 B(s)^2=\sum_{i,\alpha}|\dot v_{i\alpha}|^2
                   +\sum_{\{i,j\}}\|\dot T_{ij}\|_{\rm op}^2.
 \tag{V34.14}\label{r34:speed}
\]
All spin gradients here are in the fixed framed coordinates.

\begin{corollary}[Uniform normalized response]
\label{r34:response}
For smooth $F_s$ on $\mathcal M$, at every point of such a curve,
\[
 \left|\frac{d}{ds}\lambda_{\gamma,\vartheta(s)}(F_s)
                     -\lambda_{\gamma,\vartheta(s)}(\dot F_s)\right|
 \le\sqrt{12}\,C_P B(s)
                      \left[\lambda_{\gamma,\vartheta(s)}
                                         (|\nabla F_s|^2)\right]^{1/2}.
 \tag{V34.15}\label{r34:response-eq}
\]
For $\gamma\ge\Gamma_*$, $C_P$ can be replaced by $4$.
In particular, with the product geodesic distance,
\[
 W_1(\lambda_{\gamma,\vartheta(0)},\lambda_{\gamma,\vartheta(1)})
           \le\sqrt{12}\,C_P\int_0^1 B(s)\,ds .
 \tag{V34.16}\label{r34:wasserstein}
\]
The curve must remain in $K$. This is a stability bound,
not a coefficient less than one or a mixing theorem.
\end{corollary}
\begin{proof}
Differentiating the numerator \emph{and} the normalizer gives
\[
 \frac{d}{ds}\lambda(F_s)
      =\lambda(\dot F_s)-\gamma\operatorname{Cov}_\lambda(F_s,S_s),
          \qquad S_s=\partial_sE_{\vartheta(s)} .
\]
At vertex $i$ the spin gradient of $S_s$ is the tangent
projection of minus a sum of $n_i+\deg(i)=6$ terms:
the slot derivatives and the incident transport derivatives
applied to neighboring unit spins. Cauchy--Schwarz at each
vertex, followed by counting each edge twice, gives
\[
 |\nabla S_s|^2
       \le6\sum_{i,\alpha}|\dot v_{i\alpha}|^2
                +12\sum_{\{i,j\}}\|\dot T_{ij}\|_{\rm op}^2
       \le12B(s)^2 .
 \tag{V34.17}\label{r34:score-gradient}
\]
Cauchy--Schwarz for covariance and two applications of
\eqref{r34:poincare-eq} now cancel the explicit $\gamma$:
\[
 \gamma|\operatorname{Cov}_\lambda(F_s,S_s)|
    \le C_P\sqrt{\lambda(|\nabla F_s|^2)\lambda(|\nabla S_s|^2)}.
\]
This proves \eqref{r34:response-eq}. Apply it to any fixed
one-Lipschitz test function (by smooth approximation), integrate
along the curve, and use the dual formula for $W_1$ on the compact
metric space to obtain \eqref{r34:wasserstein}.
\end{proof}

In particular the centered normalized score has Fisher bound
$\gamma^2\operatorname{Var}_\lambda S_s\le12C_P\gamma B(s)^2$.
Moving native frames contribute to $\dot F_s$; they are not
silently omitted. Nor does \eqref{r34:response-eq} control
the covariance of retained predictors under their own interacting
marginal. The constant is uniform but not contractive.

\subsection{All-block conditional innovations under the original law}

Return to the exact V29 packing and its original normalized
Wilson measure $\mu_W$. Denote all retained links by $\eta$.
Conditioned on $\eta$, the free six-link blocks are independent,
and the framed law of block $b$ is exactly
$\lambda_{\gamma,\vartheta_b(\eta)}$. This factorization is
under the original measure, without replacing any factor
by its Laplace approximation.

Let $\mathcal B$ be any finite set of these blocks.
Allow $F_b=F_b(\eta,z_b)$ to depend on retained links as well
as its own free block, but not on other free blocks.
Assume real $F_b$ are square integrable, have $H^1$
spin dependence, and for every retained configuration satisfy
\[
 \lambda_{\gamma,\vartheta_b(\eta)}(|\nabla_{z_b}F_b|^2)\le L_b^2,
 \qquad
 \operatorname{osc}_{z_b}F_b(\eta,\cdot)\le O_b .
 \tag{V34.18}\label{r34:observable-bounds}
\]
For example uniform spin Lipschitz and oscillation bounds suffice.
Define
\[
 M_b(\eta)=\mathbb E_{\mu_W}[F_b\mid\eta],
        \qquad X_b=F_b-M_b(\eta).
 \tag{V34.19}\label{r34:innovation}
\]

\begin{theorem}[Native innovation covariance without a volume union bound]
\label{r34:innovation-theorem}
On the dyadic periodic regulator class for which V31's original-law
witness tail holds, every real coefficient family $(a_b)$ obeys
\[
 \mathbb E_{\mu_W}\left|\sum_{b\in\mathcal B}a_bX_b\right|^2
 \le\sum_{b\in\mathcal B}a_b^2
       \left\{\frac{C_P L_b^2}{\gamma}
                     +\frac{O_b^2}{4}p_\gamma(1/1024)\right\}.
 \tag{V34.20}\label{r34:innovation-bound}
\]
Here, exactly as in V31,
\[
 \begin{aligned}
 p_\gamma(1/1024)
    &=\min\{1,45\exp(b_\gamma-\gamma/(900\cdot1024))\},\\
 b_\gamma
    &=384\left(8+\log\gamma-\frac23\log\frac8{3\pi^3}\right).
 \end{aligned}
 \tag{V34.21}\label{r34:original-tail}
\]
For $\gamma\ge\Gamma_*$ replace $C_P$ by $4$.
No event requiring all blocks to be good is imposed.
In particular, for uniformly bounded $L_b,O_b$ and
$\sum_ba_b^2=1$, the bound is independent of $|\mathcal B|$
and of surrounding volume.
\end{theorem}
\begin{proof}
Conditional centering and conditional independence give
\[
 \mathbb E[X_bX_c\mid\eta]=0\quad(b\ne c),\qquad
 \mathbb E\left|\sum_ba_bX_b\right|^2
        =\sum_ba_b^2\,\mathbb E\operatorname{Var}(F_b\mid\eta).
 \tag{V34.22}\label{r34:orthogonality}
\]
Overlap of retained carriers does not spoil this identity,
since they are fixed by the conditioning.

On $G_b=\{D_b(\eta)^2\le1/1024\}$,
\eqref{r31:d-D} gives $d(\vartheta_b)\le1/32$.
Theorem~\ref{r34:poincare} bounds the conditional variance
there by $C_PL_b^2/\gamma$. On $G_b^c$, the elementary bound
$\operatorname{Var}Y\le(\operatorname{osc}Y)^2/4$
bounds it by $O_b^2/4$. Integrate these two bounds and use
$\mu_{\mathcal R}(G_b^c)\le p_\gamma(1/1024)$ from
\eqref{r31:bad}. Sum the exact orthogonality identity.
No dependent event probabilities have been multiplied or
conditioned on an arbitrary exterior, and no union over
$\mathcal B$ has been taken.
\end{proof}

This bound controls a whole family in its natural $\ell^2$
normalization. It does not say that an unnormalized extensive
sum has volume-independent variance. The constants in
\eqref{r34:original-tail} are deliberately conservative and
do not certify useful moderate coupling.

For clarity, the exact covariance bookkeeping is
\[
 \operatorname{Cov}_{\mu_W}(F_b,F_c)
  =\operatorname{Cov}_{\mu_{\mathcal R}}(M_b,M_c)
       +\mathbf1_{\{b=c\}}\,
                       \mathbb E_{\mu_{\mathcal R}}
                              \operatorname{Var}(F_b\mid\eta).
 \tag{V34.23}\label{r34:covariance}
\]
It follows simply by writing $F_b=M_b+X_b$ and conditioning
each mixed term. This is the inherited predictor decomposition,
not a newly discovered identity. V34 bounds its second term.
The first term still contains all correlations propagated
through the retained field.

\subsection{ESM V23 transfer: an exact plaquette derivative floor}
\label{r34:ESM}

During this iteration the user supplied
\path{Renewal_Perturbative_ESM_Continuum_Research_Notes_V23.tex}
and a transfer memorandum. We read the whole V23 append,
its domain and open ledger, and the earlier source-composition,
rooted-contraction, sector-rarity, and conditional-guard arguments.
The manuscript's detailed BV coefficients are specific to its
completed gravitational source, paired odd grids, finite loop/source
order, and fixed positive lower scale. None is substituted for a
nonperturbative Yang--Mills estimate.

The transferable discipline is to test whether the lowest
coefficient is actually present before paying a generic bound.
Field-amplitude order, external-momentum order, and coupling order
must remain distinct. V32 already proved a vanishing linear
coefficient of its centered remainder; V33 already made the
identity-Hessian error relative to the Maxwell Schur form.
Those results are inherited, not reproved as new momentum
cancellations. In particular the Schur form still has the exact
soft scale $4\sin^2(\pi/\ell_{\max})$. An extra factor of momentum
in an error does not itself create a physical mass.

A different, useful cancellation occurs in the native observable.
For a plaquette $p$ meeting block $b$, put
$e_p=e(U_p)=1-\tfrac12\operatorname{Re}\operatorname{Tr}U_p$.
Exactly one or two of its links are free in that block;
write this number as $s_p\in\{1,2\}$.
At fixed exterior, the framed function is respectively
$1-z_i\cdot v$ or $1-z_i\cdot Tz_j$.
Consequently
\[
 |\nabla_{z_b}e_p|^2=s_p(2e_p-e_p^2)\le2s_p e_p .
 \tag{V34.24}\label{r34:plaquette-floor}
\]
Indeed the tangent projection of a unit vector $w$ at $z$
has squared norm $1-(z\cdot w)^2$; for an internal edge
both endpoint gradients have this norm because $T$ is orthogonal.
Thus the gradient vanishes exactly at a flat plaquette, as well
as at its antipodal maximum. Equation~\eqref{r34:plaquette-floor}
is an elementary spherical identity. Its new use here is in
the full native conditional variance theorem, not a novelty
claim for the identity itself.

The original single-plaquette tail \eqref{r31:inherited-tail}
also gives the volume-independent mean bound
\[
 \mathbb E_{\mu_W}e_p\le
       \epsilon_\gamma:=\min\left\{2,\frac{b_\gamma+1}{\gamma}\right\}.
 \tag{V34.25}\label{r34:mean-energy}
\]
To see this, integrate the tail over $0\le t\le2$ and bound
it by the integral over $[0,\infty)$ of
$\min\{1,e^{b_\gamma-\gamma t}\}$. Since $b_\gamma\ge0$,
that integral is $(b_\gamma+1)/\gamma$.
This uses the original normalized law, not a prescribed
conditional exterior.

\begin{corollary}[Improved native plaquette-innovation scale]
\label{r34:plaquette-innovation}
For one chosen plaquette $p_b$ meeting each block $b$, set
$X_b=e_{p_b}-\mathbb E[e_{p_b}\mid\eta]$.
On the same regulator class as Theorem~\ref{r34:innovation-theorem},
\[
 \mathbb E_{\mu_W}\left|\sum_b a_bX_b\right|^2
 \le\sum_ba_b^2\left\{
       \frac{2s_{p_b}C_P\epsilon_\gamma}{\gamma}
                           +p_\gamma(1/1024)\right\}.
 \tag{V34.26}\label{r34:improved-single}
\]
More generally let
$F_b=\sum_{p\in\mathcal P_b}c_{b,p}e_p$ with real fixed coefficients,
where $\mathcal P_b$ consists of plaquettes meeting that block,
and set $A_b=\sum_p|c_{b,p}|$.
For $X_b=F_b-\mathbb E[F_b\mid\eta]$,
\[
 \mathbb E_{\mu_W}\left|\sum_b a_bX_b\right|^2
 \le\sum_ba_b^2A_b^2
       \left\{\frac{4C_P\epsilon_\gamma}{\gamma}
                         +p_\gamma(1/1024)\right\}.
 \tag{V34.27}\label{r34:improved-panel}
\]
Replace $C_P$ by $4$ when $\gamma\ge\Gamma_*$.
Thus a uniformly bounded panel with $\sum_ba_b^2=1$ has
variance $O(\log\gamma/\gamma^2)$ as $\gamma\to\infty$,
uniformly in ambient volume. These are conditional innovations,
not the uncentered plaquette correlations.
\end{corollary}
\begin{proof}
On $G_b$, \eqref{r34:poincare-eq} and
\eqref{r34:plaquette-floor} give
\[
 \operatorname{Var}(e_p\mid\eta)
          \le\frac{2s_p C_P}{\gamma}
                              \mathbb E[e_p\mid\eta].
\]
Multiply by $\mathbf1_{G_b}$ and average. Since $e_p\ge0$,
discarding this indicator increases the right side, so
\eqref{r34:mean-energy} applies. On $G_b^c$, $0\le e_p\le2$
gives variance at most one, paid by the original-law probability
$p_\gamma(1/1024)$. The orthogonality
\eqref{r34:orthogonality} proves \eqref{r34:improved-single}.

For a panel, weighted Cauchy--Schwarz and $s_p\le2$ give
\[
 |\nabla F_b|^2
       \le A_b\sum_p|c_{b,p}|\,|\nabla e_p|^2
       \le4A_b\sum_p|c_{b,p}|\,e_p.
\]
The same good-event argument costs at most
$4C_PA_b^2\epsilon_\gamma/\gamma$.
Its spin oscillation is at most $2A_b$, giving bad-event
variance at most $A_b^2$. Average and use orthogonality.
Finally $b_\gamma=O(\log\gamma)$, while the fixed-threshold
$p_\gamma(1/1024)$ decays faster than any inverse power of
$\gamma$. This proves the stated asymptotic scale without
claiming a useful explicit threshold.
\end{proof}

This transfer improves the \emph{actual normalized source variance}
by exploiting a proved missing linear spin response, rather than
assuming an extra spatial derivative of an unknown retained kernel.
It leaves the predictor term in \eqref{r34:covariance} untouched.

Two further audit points from the memorandum matter. First, the
native factor is $J_{\gamma,b}=e^{\gamma E_{\Phi,b}}Q_\gamma(\vartheta_b)$.
At a controlled exterior, its coherent logarithmic ratio is
\[
 \log\frac{J_{\gamma,b}(\eta)}{J_{\gamma,b}(1)}
   =\gamma(E_{\Phi,b}-m_b)
       -\frac12\log\frac{\det M_b}{\det M(0)}
       +h_{\gamma,b},
 \tag{V34.28}\label{r34:correct-factor}
\]
with V32's coherence-subtracted remainder $h_{\gamma,b}$.
A shorthand keeping only $-\gamma m_b$ would omit the
filled-energy term. We retain the exact convention.

Second, ESM V1's source-cumulant tower says that a second-order
source coefficient contains both the averaged conditional
variance and the covariance of conditional means. This agrees
with \eqref{r34:covariance}; keeping the mean alone does not
close source transport. The ESM warning that rare sectors need
not have decaying connected activities also survives unchanged.
We therefore select this derivative-floor improvement and
explicit return bookkeeping, but do not import a claimed
contractive RG trajectory, BV restoration, or full continuum
measure. Paired source jets remain candidates only after their
contact annihilation is checked against the actual YM Gram
form. Fixed-order coefficientwise limits remain distinct from
the desired interacting positive-time spectral theorem.

\subsection{Verification and the next genuinely physical obligation}

The exact-arithmetic diagnostic
\begin{center}\small
\path{scripts/verify_ym_restart_v34_conditional_variance.py}
\end{center}
checks the graph counts, the corner/middle constants, all $64$
coherent collinear critical configurations, the incident derivative
count, the exact plaquette-gradient identity, and a rational
correlated-retained conditional-product
fixture. In that fixture the innovations are orthogonal but the
unconditional observables retain nonzero off-diagonal covariance.
It therefore checks both the useful identity and the reason it
must not be mistaken for full mixing. These finite checks are
not substitutes for the compactness, Bochner, and localization
proof above, and do not numerically certify $\Gamma_*$.

The achieved input is a normalized full-domain block variance
bound and a coefficient-response bound, with a non-extensive
$\ell^2$ assembly for native conditional innovations.
The next obligation is to control the first term in
\eqref{r34:covariance}, or an equivalent retained-field/physical-time
influence. Iterating the block estimate without paying the retained
predictors would be circular: those predictors are precisely what
survive the conditional elimination. Any further block iteration
must track the actual changed interaction and its normalized law.

V33's global action comparison, V34's conditional inequality,
and the original-law rare-defect estimate are compatible inputs,
but they do not yet imply contractivity, exponential clustering,
a cofinal interacting four-dimensional continuum measure, or a
positive physical Hamiltonian mass gap. Those obligations remain
open. V33's entire body is preserved; only the displayed version
title changes. Build outputs and diagnostic artifacts remain
outside \path{latex_notes}.

% END CONSOLIDATED V34 APPEND

% BEGIN CONSOLIDATED V35 APPEND -- 2026-09-19
\clearpage
\section{V35: Joint native coefficient defects and shrinking rough components}
\label{r35:section}

\subsection{Companion checkpoint and the upstream choice}

V34 controls conditional fluctuations, but not the covariance of their
retained predictors. Its exceptional-exterior estimate concerns one
block at a time. We now control the joint geometry of those very
coefficient defects under the original Wilson law, including overlapping
translated conditional blocks. The result is an arbitrary-support
one-sided source bound with explicit constants, not merely a union
bound for a fixed list of observables.

This is not a new general bad-cluster method. Archive $A$, V20--V21
already developed full-law bad-patch and conditional-confidence
estimates, and V32 there improved their strip-energy witness.
Those statements explicitly depend on their selected-coupling
cell-source input. Archive $A$, V71 treats a different integrated
six-neighbour conditional. Here the input is instead the original
periodic Wilson reflection inequality of archive $C$, V65, the
six-link coefficient witness of current V31, and the actual cap
normalizer. No susceptibility-selected coupling or independently
postulated source theorem is assumed. The unsplit witness and its
32-colour native packing give the quantitative new source budget.

The scheduled companion checkpoint reread the terminal proofs and
scope statements of NS V26, SM V72, shell-mixing V16 and parallel
YM V5. Their complete hashes are unchanged from V30.
NS supplies no new YM spatial estimate; its reference to a future
V27 certificate is not taken as evidence. SM's many-copy theorem
keeps its physical mode bank fixed. Shell mixing's measured
flat-holonomy coefficient does not control a common nonperturbative
integral remainder or a physical mass. Parallel YM's extensive-charge
warning still favors local geometry over global defect absence.
The supplied perturbative ESM V23 also remains unchanged.
Its derivative-floor and explicit-return discipline is already
used in V34; its rarity-versus-connectedness warning remains relevant.

For the reflection step we checked the primary closed-block
statement in M.~Biskup,
\emph{Reflection Positivity and Phase Transitions in Lattice Spin Models},
\url{https://arxiv.org/html/math-ph/0610025v2},
Section 5.3, equations (5.34)--(5.35).
The HTML numbers the theorem 5.2; earlier archive references use
5.8. Blocks share boundary variables. We use the reflection
inequality, not a spin-model clustering conclusion.
This targeted check does not reset the broad-sweep cadence:
the next broad literature review and companion review are both V40.

\subsection{One local conditional at every anchor}

Fix the coordinate order $(\mu,\nu,\rho,\tau)$ used in V29--V34.
For any fine vertex $x$ let $\mathcal F_x$ be the six positive
$\mu$ links starting at
\[
       x+c\widehat\nu+r\widehat\rho,\qquad
                         c\in\{0,1\},\quad r\in\{0,1,2\}.
 \tag{V35.1}\label{r35:local-block}
\]
Fix all links outside $\mathcal F_x$ and integrate just these six
links. There are exactly $22$ private plaquettes and seven internal
plaquettes. Thus this \emph{single-block} conditional is V31's
$22$-slot, seven-transport law on $(S^3)^6$, at every anchor.
It does not require the anchor to belong to the sparse packing,
or its time coordinate to avoid the two retained boundary slices.
Different such conditionals overlap; they are not asserted independent.

Use the same upper-$\tau$ three-link filling for each free link,
the same left-column frame tree, and the same two left-grid
$\nu\rho$ squares as in V31. Define the nonnegative, exterior-measurable
coefficient defect
\[
 \mathcal D_x
   =2E_{{\rm priv},x}(\Phi_x\eta_x)
       +4E_{{\rm int},x}(\Phi_x\eta_x)
       +16\{e(L_{x,1})+e(L_{x,2})\},
       \qquad \eta_x=U|_{\mathcal F_x^c}.
 \tag{V35.2}\label{r35:D}
\]
Thus $\mathcal D_x$ is the quantity previously denoted $D_b^2$,
not its square root. It is gauge invariant. The local coefficient
defect obeys $d(\vartheta_x)^2\le\mathcal D_x$.
All these facts are the translated literal-word proof of
\eqref{r31:d-D}; every word in \eqref{r35:D} uses links outside
$\mathcal F_x$. Translation introduces no new analytic assumption.
In particular V34's full-domain conditional Poincare theorem
applies wherever $\mathcal D_x\le1/1024$, uniformly in $x$.

Let $\mathcal W_x$ consist of the same four face collections as V31:
$22$ private faces, seven top faces, fourteen endpoint side faces
and the two left-grid squares. They are distinct, and
\[
 \mathcal D_x\le20\sum_{p\in\mathcal W_x}e(U_p)
       \quad\hbox{for every full configuration }U.
 \tag{V35.3}\label{r35:witness}
\]
This is the inherited nonlinear witness inequality, now used
before splitting its $45$ faces into separate threshold events.

\begin{lemma}[A strictly interior reflection carrier]
\label{r35:carrier}
Every vertex of every face of $\mathcal W_x$ belongs to
\[
 x+\{0,1\}\times\{-1,0,1,2\}\times
              \{-1,0,1,2,3\}\times\{-1,0,1\}.
 \tag{V35.4}\label{r35:vertex-box}
\]
The $45$ faces therefore lie strictly inside the closed
rectangle with lower vertex
$x-(1,2,2,2)$ and side lengths
\[
             h=(4,8,8,4),\qquad v_h=1024.
 \tag{V35.5}\label{r35:h}
\]
The assertion includes all endpoints of supporting links,
not just plaquette base vertices.
\end{lemma}
\begin{proof}
The private $\mu\nu$ faces are the six faces just outside the
two-column grid; their $\nu$ coordinates lie between $-1$ and $2$.
The private $\mu\rho$ faces are the four faces outside the
three-row grid; their $\rho$ coordinates lie between $-1$ and $3$.
The twelve private $\mu\tau$ faces are the two temporal neighbours
of each free link, with $\tau$ coordinates in $[-1,1]$.
Every one has $\mu$ coordinate zero or one relative to $x$.
The seven top faces are the three internal $\mu\nu$ faces and
four internal $\mu\rho$ faces translated by $+\widehat\tau$.
The fourteen endpoint side faces have directions $\nu\tau$ or
$\rho\tau$, at $\mu=0$ or $1$, with $\tau$ between zero and one.
The two left-grid squares have directions $\nu\rho$, at
$\mu=\tau=0$. These listings prove \eqref{r35:vertex-box}.
Relative to the stated lower rectangle vertex, its coordinate
ranges are respectively $[1,2]$, $[1,4]$, $[1,5]$, $[1,3]$,
all strictly between zero and the corresponding side length.
\end{proof}

Use the cofinal subfamily of V31's dyadic tori with
$L=2^j$, $j\ge4$, and $2q=2^\ell$, $\ell\ge3$.
Each side divided by its corresponding length in
\eqref{r35:h} is a power of two at least two.
This mildly larger minimum spatial size is needed by the
closed-block reflection theorem. Set
\[
 B_\gamma=8+\log\gamma-\frac23\log c_H,\qquad
 c_H=\frac8{3\pi^3},\qquad \gamma\ge1.
\]
The original full measure is $\mu_W=Z^{-1}e^{-\gamma E_W}dH$,
and the inherited cap estimate is
\[
                 -\log Z\le6|\Lambda|B_\gamma.
 \tag{V35.6}\label{r35:Z}
\]

\subsection{Joint tails with the original normalizer}

Two anchor families will be useful. For all fine anchors,
colour $x$ by its four residues modulo $(4,8,8,4)$;
there are $K=1024$ colours. For the native sparse packing
of V29, $x_\nu,x_\rho$ are multiples of four and $x_\tau$
is even, with the two excluded temporal slices still omitted.
At most $4\cdot2\cdot2\cdot2=32$ of these residue colours occur;
use $K=32$ for that family. All members of one colour occupy
distinct rectangles of one translated reflection tiling.
Missing anchors do not affect this property.

For either family define
\[
 \begin{aligned}
 \lambda_{\gamma,K}&=\frac{\gamma}{20K},&
 a_\gamma&=\frac{122880\,B_\gamma}{\gamma},\\
 q_{\gamma,K}(u)
   &=\min\left\{1,
         \exp\left(\frac{6144B_\gamma-\gamma u/20}{K}\right)\right\}
     =\exp\{-\lambda_{\gamma,K}(u-a_\gamma)_+\}.
 \end{aligned}
 \tag{V35.7}\label{r35:q}
\]

\begin{theorem}[Arbitrary-support joint coefficient defects]
\label{r35:joint}
For every finite set of distinct anchors $A$ in the selected family
and every collection $u_x\ge0$,
\[
 \mu_W\{\mathcal D_x>u_x\ \hbox{for every }x\in A\}
              \le\prod_{x\in A}q_{\gamma,K}(u_x).
 \tag{V35.8}\label{r35:joint-eq}
\]
For the native packing the same assertion holds under the exact
retained marginal $\mu_{\mathcal R}$, since all these defects
are retained-measurable. No independence, selected coupling,
changed reference law, or all-good event is assumed.
\end{theorem}
\begin{proof}
Fix a residue colour. By \eqref{r35:witness}, the event at $x$
is contained in
\[
        A_x(u_x)=\left\{\sum_{p\in\mathcal W_x}e(U_p)>u_x/20\right\}.
\]
This event belongs to the closed-block algebra of its rectangle.
Pull it back by the alternating reflections to the base rectangle
of that colour's tiling. Fully disseminate this pulled-back event.
All reflected witness faces remain strictly interior, by
Lemma~\ref{r35:carrier}. Witnesses in different rectangles
therefore contain distinct original plaquettes, including at
periodic seams. If $N_h=|\Lambda|/1024$, the disseminated
event has total action at least $N_hu_x/20$.
The product Haar mass of any event is at most one, so its
\emph{normalized} probability is at most
$Z^{-1}e^{-\gamma N_hu_x/20}$. Taking its chessboard root
and using \eqref{r35:Z} gives
\[
 \mu_W(\hbox{disseminated }A_x)^{1/N_h}
       \le\min\{1,\exp(6144B_\gamma-\gamma u_x/20)\}.
 \tag{V35.9}\label{r35:root}
\]
The closed-block chessboard inequality of archive $C$, V65
then bounds any same-colour intersection by the product
of \eqref{r35:root}. The labels $u_x$ and the pulled-back
local events may differ from rectangle to rectangle; the
inequality explicitly allows distinct block functions.
Reflection of an energy witness need not preserve the original
sparse elimination pattern, and no such property is used.

For each of the $K$ colours let $I_c$ be the indicator of
the corresponding intersection. Empty colours have $I_c=1$.
Generalized Holder inequality, using $I_c^K=I_c$, gives
\[
 \mu_W\left(\prod_c I_c\right)
                  \le\prod_c\mu_W(I_c)^{1/K}.
\]
Insert the same-colour bounds. Each anchor is counted exactly
once, yielding \eqref{r35:joint-eq}. The retained-marginal
version follows by the pushforward identity.
The proof reflects only the original Wilson law; it never
assigns reflection positivity to the retained marginal.
\end{proof}

For clarity, the two useful probabilities are
\[
 \begin{aligned}
 q_{\gamma,32}(u)&=\min\{1,e^{192B_\gamma-\gamma u/640}\},\\
 q_{\gamma,1024}(u)&=\min\{1,e^{6B_\gamma-\gamma u/20480}\}.
 \end{aligned}
 \tag{V35.10}\label{r35:two-q}
\]
These constants are explicit and conservative. Both tend to zero
for fixed $u>0$ as $\gamma\to\infty$, without specifying a
susceptibility-selected trajectory. The larger all-anchor colour
loss pays overlapping local conditionals, not a hidden independence
assumption.

\subsection{A native one-sided multisource estimate}

The joint tail has an integrated form which can be used before
choosing a particular defect threshold.

\begin{theorem}[Positive source budget at arbitrary support]
\label{r35:source}
For a finite anchor family $A$ and $0\le t_x<\lambda_{\gamma,K}$,
\[
 \mathbb E_{\mu_W}
       \exp\left(\sum_{x\in A}t_x\mathcal D_x\right)
 \le \prod_{x\in A}
                \frac{e^{a_\gamma t_x}}{1-t_x/\lambda_{\gamma,K}} .
 \tag{V35.11}\label{r35:mgf}
\]
The same estimate holds under $\mu_{\mathcal R}$ for the native
packing. In particular, with the explicit upper centering
$m_{\gamma,K}^+=a_\gamma+\lambda_{\gamma,K}^{-1}$,
\[
 \log\mathbb E\exp\left(\sum_x t_x
                           (\mathcal D_x-m_{\gamma,K}^+)\right)
 \le\sum_x\left\{-\log(1-t_x/\lambda_{\gamma,K})
                                     -t_x/\lambda_{\gamma,K}\right\}.
 \tag{V35.12}\label{r35:center}
\]
This centering is not identified with the actual mean.
The estimate is one-sided in every source.
\end{theorem}
\begin{proof}
For a nonnegative random variable $X$ and $t\ge0$,
\[
 e^{tX}=1+\int_0^\infty t e^{tu}\mathbf1_{\{X>u\}}\,du.
\]
Multiply these identities over $A$, expand over subsets, and
apply Tonelli. Every term has nonnegative coefficient, so
\eqref{r35:joint-eq} bounds it by the corresponding product
of one-dimensional integrals. The function
$q_{\gamma,K}(u)$ is exactly the survival function of
$Y=a_\gamma+Z$, where $Z$ is exponential of rate
$\lambda_{\gamma,K}$. Thus the resulting product is
$\prod_x\mathbb E e^{t_xY}
=\prod_x e^{a_\gamma t_x}/(1-t_x/\lambda_{\gamma,K})$.
Subtract the displayed upper centering and take logarithms
to get \eqref{r35:center}.
\end{proof}

The same nonnegative-integral argument gives, for positive
integers $r_x$,
\[
 \mathbb E\prod_{x\in A}\mathcal D_x^{r_x}
       \le\prod_{x\in A}
          \left\{\sum_{j=0}^{r_x}
               \binom{r_x}{j}a_\gamma^{\,r_x-j}
                          \frac{j!}{\lambda_{\gamma,K}^{\,j}}\right\}.
 \tag{V35.13}\label{r35:moments}
\]
Hence every fixed collection of powers costs
$O((\log\gamma/\gamma)^{\sum r_x})$, with the dependence
on the number and orders of marks explicit in the product.
The estimate also applies to the whole finite lattice at
a regulator; its logarithm is then properly extensive in
the number of nonzero source marks.

One must not differentiate \eqref{r35:center} as if it
were a two-sided centered Gaussian bound. Its upper centering
may strictly exceed $\mathbb E\mathcal D_x$, and negative
sources have not been controlled. It therefore does not
bound separated connected covariances by a decaying kernel.
Nor is it an interacting analytic RG self-map. These are
the same source-versus-connectedness distinctions retained
in the ESM audit.

For a fixed $\sigma>0$, the explicit moving threshold
\[
 u_{\gamma,K}(\sigma)
       =\frac{122880B_\gamma+20K\sigma\log\gamma}{\gamma}
 \tag{V35.14}\label{r35:moving-threshold}
\]
obeys $q_{\gamma,K}(u_{\gamma,K})=\gamma^{-\sigma}$ and
tends to zero. Eventually it is smaller than $1/1024$.
This provides simultaneous polynomial joint failure bounds
for the actual local conditional entry conditions. It does
not claim absence of defects in a much larger physical
volume or after a changed interaction has been integrated.

\subsection{Rooted rough components and a physical-diameter limit}

Here is the geometric consequence relevant to the next
retained-field decomposition. Fix $u>0$, for example
$u=1/1024$, and call $x$ bad when $\mathcal D_x>u$.
Use the graph whose vertices are the chosen anchors and
whose edges join distinct anchors when their vertex boxes
\eqref{r35:vertex-box} intersect on the periodic lattice.
For all anchors, the maximum degree is at most
\[
             \Delta=3\cdot7\cdot9\cdot5-1=944.
\]
Indeed the four relative anchor coordinates must have
periodic magnitudes at most $1,3,4,2$, respectively.
For the native packing these restrictions give at most
$3\cdot1\cdot3\cdot3-1=26$ neighbours, so use $\Delta=26$.
The omitted temporal anchors only decrease degrees.
Every graph edge moves each fine coordinate by at most four.

Let $\mathcal C_x$ be the connected component of bad anchors
containing $x$, or the empty set when $x$ is good.
Write $q_*=q_{\gamma,K}(u)$ and $r_*=\Delta^2q_*$.

\begin{corollary}[Rooted cluster probability and exponential cost]
\label{r35:clusters}
For every $n\ge1$,
\[
       \mu_W\{|\mathcal C_x|\ge n\}
                   \le q_* r_*^{\,n-1}.
 \tag{V35.15}\label{r35:cluster-tail}
\]
For $\alpha\ge0$ with $e^\alpha r_*<1$, one has
\[
 \begin{aligned}
 \mathbb E e^{\alpha|\mathcal C_x|}
       &\le1+\frac{q_*(e^\alpha-1)}{1-e^\alpha r_*},\\
 \sum_{\substack{A\ni x\\ A\ {\rm connected}}}
     e^{\alpha|A|}
          \mu_W\{\mathcal D_y>u\text{ for all }y\in A\}
       &\le\frac{e^\alpha q_*}{1-e^\alpha r_*}.
 \end{aligned}
 \tag{V35.16}\label{r35:rooted}
\]
The bounds are volume independent, and hold for the retained
marginal in the packing case.
\end{corollary}
\begin{proof}
The standard rooted-animal count, already used in archive $A$,
is at most $\Delta^{2(n-1)}$ for connected $n$-vertex sets
containing a fixed root. To recall its elementary justification,
choose a deterministic rooted spanning tree of each such set
and its deterministic depth-first traversal. The resulting
walk has length $2(n-1)$, its visited set recovers the original
set, and there are at most $\Delta^{2(n-1)}$ such walks.
If the bad root component has at least $n$ vertices, it
contains a connected $n$-vertex set containing the root.
Theorem~\ref{r35:joint} and a union bound give
\eqref{r35:cluster-tail}. Summing the tail identity for
$e^{\alpha|\mathcal C_x|}$ proves the first inequality
in \eqref{r35:rooted}; summing the rooted-animal count
with its event bound proves the second.
These are sums of probabilities, not constructed polymer
activities.
\end{proof}

\begin{corollary}[Rough coefficient components shrink in physical units]
\label{r35:physical-diameter}
Give each fine edge physical length $a\downarrow0$ and let
$\gamma(a)\to\infty$ along any of the admitted dyadic tori.
Keep $u>0$ fixed. For any fixed bounded physical window $W$
and any $\rho>0$, the probability that a bad component
with an anchor in $W$ has physical anchor diameter at least
$\rho$ tends to zero. No particular relation between
$\gamma(a)$ and $a$ is needed for this statement.
\end{corollary}
\begin{proof}
Use the periodic physical $\ell^\infty$ distance, or a
nonwinding lift of the window. The number of anchors in $W$
is at most $C_W(1+a^{-1})^4$. A connected component of $n$
anchors has diameter at most $4a(n-1)$: connect any two
anchors by a simple path in a spanning tree. Therefore
diameter at least $\rho$ requires at least
\[
                 n_a=\left\lceil\frac{\rho}{4a}\right\rceil+1
\]
anchors. A root union bound and \eqref{r35:cluster-tail}
give an upper bound
\[
      C_W(1+a^{-1})^4q_* r_*^{\,n_a-1}.
 \tag{V35.17}\label{r35:physical-probability}
\]
Since $\gamma(a)\to\infty$, $q_*\to0$, and eventually
$r_*\le1/2$. The displayed bound then tends to zero.
The union of the corresponding witness boxes changes the
diameter by at most a constant times $a$, so it has the
same vanishing macroscopic-diameter conclusion.
\end{proof}

This is a cofinal statement about \emph{coefficient-defect geometry}.
It neither constructs a continuum measure nor says that all
plaquette curvature is uniformly small. The witness inequality
is one-sided. It also does not require or imply absence of all
bad anchors: their total number may grow with volume and
refinement. The physical diameter here is a geometric length,
not the reciprocal of a proved particle mass.

\subsection{What this permits next, and what still cannot be inferred}

The original-law coefficient defects now have explicit joint
tails, mixed positive moments, and rooted rough-component
costs without a selected-coupling source hypothesis. V34's
normalized conditional theorem is available at every good
translated six-link block, not only a member of one sparse
packing. These are actual-law inputs for treating rough
components locally while analyzing propagation through their
complement.

The next missing estimate is still the propagation itself:
one must bound the retained conditional-predictor covariance,
or a normalized physical-time operator, across the good
region and across the paid rough components. V34's response
constant is not a contraction constant. Overlapping good
conditionals cannot be multiplied as independent kernels.
The second inequality of \eqref{r35:rooted} is not a
connected-activity norm for the retained action. The archive
and ESM counterexamples already explain why rare sectors
alone do not supply that norm; no further toy obstruction
is claimed as a new result here.

In particular a proof that uses \eqref{r35:physical-probability}
to discard the retained predictor term in
\eqref{r34:covariance} would be invalid. Collective smooth
modes survive entirely in the good region, and the exact
native Schur scale in V33 remains. Conversely that soft
action-Hessian scale is not a proof of physical gaplessness.
Neither issue is settled by the present probability geometry.

The checker
\begin{center}\small
\path{scripts/verify_ym_restart_v35_joint_defects.py}
\end{center}
verifies the literal $45$-face carrier, all supporting vertices,
reflection copies including periodic seams, residue colour
counts, native packing occupancy and overlap degrees. Exact
rational tests cover the source integral coefficients and
the root-tail summation identities. Those checks validate
geometry and algebra; the reflection theorem and the analytic
probability proofs above carry the infinite-family statements.
The full V34 body is preserved, with only its displayed
version title changed. All build files remain outside
\path{latex_notes}. The interacting four-dimensional
continuum construction and physical Yang--Mills mass gap
remain open.

% END CONSOLIDATED V35 APPEND

% BEGIN CONSOLIDATED V36 APPEND -- 2026-09-19
\clearpage
\section[V36: co-moving transport and retained source response]
{V36: co-moving conditional transport\\and retained source response}
\label{r36:start}

\subsection{The upstream choice and its nonduplication boundary}

V35 pays joint rough-coefficient geometry under the original Wilson
law. The remaining correlations can live entirely in the smooth
retained variables. We therefore return to the derivative of an
\emph{actual normalized conditional predictor}, and separate motion
of its conditional minimum from deformation of its thermal cloud.
This gives an exact differential transport, not a replacement
probability measure. Its deformation error is small; the motion
of the minimum is deliberately retained.

The new input is a quadratic cancellation in a co-moving score.
It upgrades V34's order-one coefficient stability to an
\(O(\gamma^{-1/2})\) error after the classical motion is removed.
For a native plaquette predictor, its vanishing spin derivative
at low curvature gives a further factor. A finite-incidence
estimate then pays the response error of an arbitrary
\(\ell^2\)-normalized family under the true retained marginal,
without assuming all blocks are good.

The general normalized score identity is not new. Current V5
already uses conditional minimizers and native marginal scores
on a different fixed history box. Archive \(C\), V32 proves
a Poisson/one-form response theorem for a restricted, rescaled
bridge law; \(C\), V35 compares normalized laws on convex balls.
Archive \(C\), V61 already has a coherent ribbon's exact
barycentric response, including its common-rotation eigenvalue
one. Its V71--V72 retain the true predictor, gauge-aligned
replicas and surviving singlets rather than deleting induced
correlations. These arguments do not provide the present
full compact six-link co-moving transport, but they rule out
claiming that its classical part is newly proved contractive.

A targeted primary-source check used Courtade, Fathi and
Pananjady,
\href{https://arxiv.org/html/1703.07707}
{\emph{Existence of Stein Kernels under a Spectral Gap, and
Discrepancy Bounds}}, Section 2.1 and the proof of Theorem 2.4.
Its variational construction motivates solving a centered
Poisson equation using a Poincare bound. We prove the required
compact-manifold equation directly below, not by identifying
our conditional law with a Euclidean Stein-kernel example.
Menz,
\href{https://arxiv.org/html/1402.5160}
{\emph{A Brascamp--Lieb type covariance estimate}},
Theorem 2.3, separately requires a positive intersite comparison
matrix to infer spatial covariance decay. We have not verified
such a matrix for the retained Wilson law. No decay theorem is
imported. The scheduled companion and broad-literature
checkpoints remain V40.

\subsection{Uniform moments around the actual conditional minimum}

Use the same six-spin compact manifold
\(\mathcal M=(S^3)^6\), product round metric, coefficient family
\(K=\{\vartheta:d(\vartheta)\le1/32\}\), and energy
\(E_\vartheta\) as V31 and V34. Write
\[
 \lambda_{\gamma,\vartheta}(dz)
   =Q_\gamma(\vartheta)^{-1}e^{-\gamma E_\vartheta(z)}\,dz,
 \qquad z_*=z_*(\vartheta),\qquad m=E_\vartheta(z_*).
 \tag{V36.1}\label{r36:law}
\]
Here \(dz\) is normalized product Haar measure and \(\gamma\ge1\).
The minimum is unique and smooth on a neighborhood of \(K\),
and \(M=\operatorname{Hess}E_\vartheta(z_*)\succeq2I\).
All constants \(C\) below are finite constants for this fixed
graph and coefficient set; they may increase between displays.
They are independent of coupling, exterior, torus and marked
support, but are not numerically evaluated.

Put \(r_\vartheta(z)=\operatorname{dist}_{\mathcal M}(z,z_*)\).
\begin{lemma}[Full compact displacement moments]
\label{r36:moments}
For \(k=1,2\),
\[
 \lambda_{\gamma,\vartheta}(r_\vartheta^{2k})
       \le C_k\gamma^{-k},\qquad
 0\le\lambda_{\gamma,\vartheta}(E_\vartheta-m)\le C/\gamma .
 \tag{V36.2}\label{r36:moment}
\]
If \(F\) is any one of the 29 plaquette-energy summands,
then \(0\le F\le2\), and
\[
 \lambda(F)\le \tfrac12d(\vartheta)^2+C/\gamma,\qquad
 \lambda(|\nabla_zF|^2)
       \le C\{d(\vartheta)^2+\gamma^{-1}\}.
 \tag{V36.3}\label{r36:plaquette-moment}
\]
\end{lemma}
\begin{proof}
V31's positive Hessian and compactness of \(K\) give one normal
coordinate neighborhood of the minimum in which
\(E_\vartheta-m\ge c r_\vartheta^2\), with uniform positive
\(c\). Outside a smaller such neighborhood the energy
difference has a strictly positive uniform lower bound:
otherwise a convergent sequence in the compact parameter-spin
product would give a second minimum. Haar density and its
normal-coordinate Jacobian have uniform upper and lower bounds
on the chosen small neighborhood.

A ball of radius proportional to \(\gamma^{-1/2}\) at \(z_*\),
and a uniform upper quadratic Taylor bound on its energy,
give \(Q_\gamma e^{\gamma m}\ge c_0\gamma^{-9}\).
The local numerator for \(r^{2k}\) is bounded by the
18-dimensional Gaussian integral \(C_k\gamma^{-9-k}\).
The complementary numerator is at most \(C_ke^{-c_1\gamma}\).
Division by the lower normalizer proves the first assertion,
enlarging constants for a bounded initial coupling interval.
No spin-domain restriction is made in the resulting law.

The global Hessian norm of \(E_\vartheta\) is uniformly bounded.
Taylor's formula along a minimizing product geodesic from
\(z_*\), whose initial energy derivative is zero, gives
\(E_\vartheta(z)-m\le C r_\vartheta(z)^2\).
The first moment proves the energy estimate.
Every summand is nonnegative and \(m\le E_\vartheta(e_0)
\le d(\vartheta)^2/2\), proving the first inequality in
\eqref{r36:plaquette-moment}. For a summand with
\(s=1\) or \(2\) free spins, V34's exact identity is
\[
              |\nabla_zF|^2=s(2F-F^2)\le4F.
\]
The second inequality follows.
\end{proof}

\subsection{A co-moving score with its constant and linear terms removed}

Let \(\vartheta(s)\) be a smooth coefficient curve in \(K\),
at fixed \(\gamma\). A dot denotes differentiation at the
parameter under consideration. Use V34's speed
\[
 B^2=\sum_{i,\alpha}|\dot v_{i\alpha}|^2+
                  \sum_{\{i,j\}}\|\dot T_{ij}\|_{\rm op}^2,
 \qquad S(z)=\partial_s E_{\vartheta(s)}(z).
 \tag{V36.4}\label{r36:speed}
\]
The metric and the spin coordinates are fixed during this
partial derivative. Define \(k=\dot z_*\) covariantly.
Differentiating the critical-point equation gives
\[
 M k=-\nabla S(z_*),\qquad |k|\le\sqrt3 B ,
 \tag{V36.5}\label{r36:velocity}
\]
where V34's \(|\nabla S|^2\le12B^2\) and \(M\succeq2I\)
were used. Thus \(k\) is linear in the coefficient tangent.

For each spin, regard \(z_{*,i},k_i\) as vectors in
\(\mathbb R^4\), and put
\[
 A_i=k_i z_{*,i}^{\mathsf T}-z_{*,i}k_i^{\mathsf T},
 \qquad W_i(z_i)=A_i z_i,\qquad
 R(z)=S(z)+W E_\vartheta(z).
 \tag{V36.6}\label{r36:residual}
\]
The field \(W\) is globally smooth and tangent, preserves Haar
measure infinitesimally, and satisfies
\[
 W(z_*)=k,\qquad \sup_z|W(z)|\le|k|\le\sqrt3 B .
\]
These are product orthogonal rotations for an auxiliary
change of integration coordinates, not physical time
evolution or an asserted lattice gauge transformation.

\begin{lemma}[Quadratic residual cancellation]
\label{r36:cancellation}
Exactly,
\[
 R(z_*)=\dot m,\qquad \nabla R(z_*)=0,\qquad
 |R(z)-\dot m|\le C B\,r_\vartheta(z)^2 .
 \tag{V36.7}\label{r36:quadratic}
\]
Consequently the score's actual centered variance satisfies
\[
             \operatorname{Var}_\lambda R
                         \le C B^2/\gamma^2 .
 \tag{V36.8}\label{r36:score-var}
\]
\end{lemma}
\begin{proof}
At the minimum, \(WE=0\) and the envelope formula gives
\(S(z_*)=\dot m\). Covariant differentiation of \(WE\) gives
\(\nabla(WE)(z_*)=M k\): the derivative of \(W\) multiplies
\(\nabla E(z_*)=0\). Equation \eqref{r36:velocity} cancels
the remaining \(\nabla S(z_*)\).

All spin derivatives of the finite polynomial energy are
uniformly bounded on the compact coefficient-spin product.
The derivatives of \(S\) are bounded by constants times \(B\);
the derivatives of \(W\) are bounded by constants times \(|k|\).
In particular \(\|\operatorname{Hess}R\|_\infty\le CB\).
Taylor's integral remainder along a minimizing geodesic proves
\eqref{r36:quadratic} globally. Finally
\(\operatorname{Var}_\lambda R\le\lambda[(R-\dot m)^2]\);
apply the fourth displacement moment in \eqref{r36:moment}.
The centering in the variance is \(\lambda R\), not \(\dot m\).
\end{proof}

\subsection{Exact normalized transport and a small deformation velocity}

Let the positive weighted Laplacian be
\[
 \mathcal L_{\gamma,\vartheta}
       =-\Delta+\gamma\nabla E_\vartheta\cdot\nabla,\qquad
 \lambda(f\mathcal L g)=\lambda(\nabla f\cdot\nabla g).
 \tag{V36.9}\label{r36:generator}
\]
This is the auxiliary block diffusion operator, never a
physical Hamiltonian.

\begin{theorem}[Co-moving conditional transport]
\label{r36:transport}
There is a unique mean-zero weak solution \(u\) of
\[
 \mathcal L_{\gamma,\vartheta}u
                 =\gamma(R-\lambda R),\qquad \lambda u=0,
 \tag{V36.10}\label{r36:poisson}
\]
and it is smooth for smooth coefficients and tangents.
The correction \(V_{\rm th}=-\nabla u\) obeys
\[
               \lambda(|V_{\rm th}|^2)\le C B^2/\gamma.
 \tag{V36.11}\label{r36:correction}
\]
For any smooth \(s\)-dependent observable \(F_s\),
\[
 \frac{d}{ds}\lambda_{\gamma,\vartheta(s)}(F_s)
       =\lambda\bigl(\dot F_s+W F_s+
                              V_{\rm th}\cdot\nabla F_s\bigr).
 \tag{V36.12}\label{r36:exact}
\]
In particular,
\[
 \left|\frac{d}{ds}\lambda(F_s)
                       -\lambda(\dot F_s+WF_s)\right|
 \le \frac{CB}{\sqrt\gamma}
                     \{\lambda(|\nabla F_s|^2)\}^{1/2}.
 \tag{V36.13}\label{r36:response}
\]
The formulas retain the exact conditional normalizer.
\end{theorem}
\begin{proof}
On the mean-zero weighted Sobolev space, V34's full compact
Poincare inequality implies
\(\|f\|_2^2\le(C_P/\gamma)\lambda(|\nabla f|^2)\).
Thus the Dirichlet form is coercive and the centered right
side in \eqref{r36:poisson} is a continuous linear functional.
The Hilbert-space representation theorem gives the unique
solution. At fixed finite \(\gamma\), the smooth positive
density on the closed compact manifold gives the usual
elliptic regularity; alternatively the weak solution suffices
for all integrated identities. Testing with \(u\) yields
\[
 \begin{aligned}
 \lambda(|\nabla u|^2)
  &=\gamma\lambda[u(R-\lambda R)]\\
  &\le\gamma(C_P/\gamma)^{1/2}
       \lambda(|\nabla u|^2)^{1/2}
                  \{\operatorname{Var}_\lambda R\}^{1/2}.
 \end{aligned}
\]
Consequently
\(\lambda(|\nabla u|^2)\le C_P\gamma
\operatorname{Var}_\lambda R\), proving
\eqref{r36:correction}. Linearity of this mean-zero solution
in the tangent follows from the same uniqueness.

Differentiate the compact normalized integral:
\[
 \frac{d}{ds}\lambda(F_s)
               =\lambda(\dot F_s)-\gamma\operatorname{Cov}_\lambda(F_s,S).
\]
Since \(W\) has zero Haar divergence, integration of
\(W(F_s e^{-\gamma E})\) gives
\(\lambda(WF_s)=\gamma\lambda(F_s WE)\).
Taking \(F_s=1\) gives \(\lambda(WE)=0\), and hence
\[
 \frac{d}{ds}\lambda(F_s)
   =\lambda(\dot F_s+WF_s)
                         -\gamma\operatorname{Cov}_\lambda(F_s,R).
\]
Equation \eqref{r36:poisson} changes the last term into
\(-\lambda(\nabla F_s\cdot\nabla u)\), including its sign.
Cauchy--Schwarz proves \eqref{r36:response}.
No differentiation of a scalar asymptotic remainder is used.
\end{proof}

There is also a full-law finite-path consequence. Solve
\(O_i'=A_iO_i\), \(O_i(0)=I\), and let
\(\mathcal O_s z=(O_i(s)z_i)_i\). Then
\(\mathcal O_s z_*(0)=z_*(s)\).
Define the co-moving normalized law
\(\widetilde\lambda_s=(\mathcal O_s^{-1})_\#\lambda_s\).
With product geodesic distance and the quadratic Wasserstein
distance,
\[
 W_2(\widetilde\lambda_0,\widetilde\lambda_1)
              \le\frac{C}{\sqrt\gamma}\int_0^1 B(s)\,ds .
 \tag{V36.14}\label{r36:W2}
\]
To justify this without postulating a coupling, apply
\eqref{r36:exact} to \(f\circ\mathcal O_s^{-1}\).
Its explicit derivative cancels \(W\). The remaining
continuity velocity is
\(O(s)^{-1}V_{\rm th}(\mathcal O_s z)\), whose squared
integral is bounded by \(CB(s)^2/\gamma\).
For a smooth coefficient path, the parameter-dependent elliptic
solution is smooth, so this vector field has a global smooth
flow on the compact manifold. Uniqueness of the continuity
equation identifies its transported initial density with
\(\widetilde\lambda_s\). Along this flow, distance is at most
integrated speed. Minkowski's inequality in the resulting
path probability space proves \eqref{r36:W2}.
Piecewise smooth paths follow by concatenation.
The curve must remain in \(K\); nothing here controls a path
through arbitrary rough coefficients.

This is a statement about deformation \emph{after} the
minimum is moved. The laboratory-frame motion \(W\) is
order \(B\), not \(B/\sqrt\gamma\). Even at coherence a common
rotation of all slots is reproduced by the minimum.
For that curve with transports fixed at identity, taking
the common plane rotation as \(W\) gives \(R\equiv0\)
and \(V_{\rm th}=0\); the entire response is \(W\).
A small correction therefore cannot be called a contraction
of the complete predictor.

\subsection{Pullback to every native retained link}

For the V29 packing, denote the actual retained links by
\(\eta\), their marginal by \(\mu_{\mathcal R}\), and the exact
coefficient map by \(\vartheta_b(\eta)\).
Let \(h=(h_e)_{e\in\mathcal R}\) be a tangent to those links,
with the product unit-round norm
\(\|h\|^2=\sum_e|h_e|^2\).
Write \(B_b(h)\) for \eqref{r36:speed} with coefficient
derivatives induced by \(h\) using precisely V31's frames.

\begin{lemma}[Volume-independent coefficient incidence]
\label{r36:incidence}
For every exterior, tangent, and subset of packed blocks,
\[
                \sum_b B_b(h)^2\le716160\,\|h\|^2.
 \tag{V36.15}\label{r36:incidence-bound}
\]
No derivative of a frame is omitted.
\end{lemma}
\begin{proof}
The tree word \(k_i\) in \eqref{r31:frame-formula} has
length at most three. Both \(p_{i\alpha}\) and \(s_i\)
have length three. Thus the word
\(k_ip_{i\alpha}s_i^{-1}k_i^{-1}\) has at most 12
retained-link occurrences. The two words defining
\(\ell'_{ij}\) and \(r'_{ij}\) have lengths at most seven
and thirteen; differentiation of
\(z\mapsto\ell'_{ij}z(r'_{ij})^{-1}\) in operator norm
therefore costs at most 20 occurrences in total.

Quaternion multiplication is isometric and inversion
preserves tangent norm. The squared derivative of a word
with \(l\) occurrences is at most
\(l\sum_{\rm occurrences}|h_e|^2
 \le l^2\sum_{e\in S_b}|h_e|^2\).
The same argument applies to the sum of the two transport
word derivatives. With 22 slots and seven transports this gives
\[
 B_b(h)^2\le(22\cdot12^2+7\cdot20^2)
                         \sum_{e\in S_b}|h_e|^2
                 =5968\sum_{e\in S_b}|h_e|^2.
\]
All these words use retained edges whose two endpoints lie
in V35's 120-vertex bounding box at the block anchor.
For a fixed positive edge, its base point allows at most
120 such anchors. Summing gives \(5968\cdot120=716160\).
Periodic identifications can only reduce this upper bound.
The estimate is intentionally generous, not an optimized
native Jacobian norm.
\end{proof}

Fix, for each marked block \(b\), one of its 29 original
plaquette energies \(F_b(\eta,z_b)\), expressed in the exact
frame, and define its predictor
\[
          M_b(\eta)=\lambda_{\gamma,\vartheta_b(\eta)}
                                      (F_b(\eta,\cdot)).
 \tag{V36.16}\label{r36:predictor}
\]
The derivative \(\partial_hF_b\) below is at fixed framed spin
coordinates and includes the explicit variation of that frame.

Write \(\mathcal D_b=D_b^2\) as in V35 and
\(G_b=\{\mathcal D_b\le1/1024\}\).
Define a pointwise differential approximation by
\[
 \mathcal A_b(h)=
 \begin{cases}
 \lambda_b(\partial_hF_b+W_{b,h}F_b),&\eta\in G_b,\\
 \lambda_b(\partial_hF_b),&\eta\notin G_b,
 \end{cases}
 \qquad
 \varepsilon_b(h)=\partial_h M_b-\mathcal A_b(h).
 \tag{V36.17}\label{r36:native-error}
\]
On \(G_b\), \(W_{b,h}\) is the moving-minimum field already
constructed; no minimum is selected outside \(G_b\).
The approximation is a tangent-linear expression, not
asserted to be the derivative of a scalar effective source.

\begin{proposition}[Actual-law source-response error]
\label{r36:native}
Use V35's dyadic tori \(L=2^j\), \(j\ge4\), and
\(T=2^\ell\), \(\ell\ge3\), and their original Wilson law.
Set
\[
 \begin{aligned}
 \mathcal B_\gamma&=8+\log\gamma-\tfrac23\log c_H,
                         &c_H&=\frac8{3\pi^3},\\
 m_\gamma^+&=\frac{122880\mathcal B_\gamma+640}{\gamma},
 &q_\gamma^*&=\min\{1,\exp(192\mathcal B_\gamma
                                  -\gamma/(640\cdot1024))\}.
 \end{aligned}
 \tag{V36.18}\label{r36:probability}
\]
For deterministic real source coefficients \(c_b\) and any
smooth retained tangent field \(h(\eta)\) with
\(\sup_\eta\|h(\eta)\|\le H\),
\[
 \int\left|\sum_b c_b\varepsilon_b(h(\eta))\right|^2
                                    d\mu_{\mathcal R}(\eta)
 \le C H^2\sum_b c_b^2
       \left\{\frac{m_\gamma^+}{\gamma}
                    +\frac1{\gamma^2}+\gamma^2q_\gamma^*\right\}.
 \tag{V36.19}\label{r36:native-bound}
\]
In particular its root-mean-square bound is
\[
 C H\|c\|_{\ell^2}
       \left\{\frac{\sqrt{1+\log\gamma}}{\gamma}
                                  +\gamma\sqrt{q_\gamma^*}\right\}.
 \tag{V36.20}\label{r36:rate}
\]
The constants are uniform in the number of marked blocks.
The averaging law is the full unmodified retained marginal,
not its conditioning to the simultaneous good event.
\end{proposition}
\begin{proof}
At a good exterior, \eqref{r36:response},
\eqref{r36:plaquette-moment}, and \(d^2\le\mathcal D_b\)
give the pointwise inequality
\[
 |\varepsilon_b(h)|^2
       \le C B_b(h)^2
               \left(\frac{\mathcal D_b}{\gamma}+
                                      \frac1{\gamma^2}\right).
\]
At any exterior the normalized score formula is still exact,
whether or not a unique minimum exists. There are 29 score
terms, so \(\|S_b\|_\infty\le\sqrt{29}B_b(h)\) and
\(\operatorname{osc}S_b\le2\sqrt{29}B_b(h)\).
Since \(0\le F_b\le2\), the elementary range bound for
covariance gives
\[
 |\gamma\operatorname{Cov}_{\lambda_b}(F_b,S_b)|^2
                          \le29\gamma^2B_b(h)^2.
\]
Thus, enlarging \(C\) once,
\[
 |\varepsilon_b(h)|\le B_b(h) K_b,\qquad
 K_b^2=C{\bf1}_{G_b}
       \left(\frac{\mathcal D_b}{\gamma}+\frac1{\gamma^2}\right)
                     +29\gamma^2{\bf1}_{G_b^c}.
\]
Cauchy--Schwarz and \eqref{r36:incidence-bound} give, for
each exterior separately,
\[
 \left|\sum_b c_b\varepsilon_b(h)\right|^2
           \le716160\,\|h\|^2\sum_b c_b^2 K_b^2.
\]
V35's native \(K=32\) tail implies
\(\mu_{\mathcal R}(G_b^c)\le q_\gamma^*\) and
\(\int\mathcal D_b\,d\mu_{\mathcal R}\le m_\gamma^+\).
Taking expectation proves \eqref{r36:native-bound}.
Equation \eqref{r36:rate} follows from
\(m_\gamma^+=O((1+\log\gamma)/\gamma)\) and
\(\sqrt{x+y}\le\sqrt x+\sqrt y\).

Only one-point consequences of V35 are needed for this
first-derivative estimate; its stronger joint tails remain
available but are not falsely presented as necessary here.
The bad-event indicator is \emph{not differentiated}:
\eqref{r36:native-error} splits the already defined
pointwise derivative, and does not differentiate a cutoff
predictor. The coefficient of the rare event is paid under
the same measure as the left-hand side.
\end{proof}

For fixed \(\gamma\), the marked predictor \(M_b\) is a
smooth function of the retained links by compact integration,
even across the good-set boundary. The approximation
\(\mathcal A_b\) need not be continuous there, and no
integrability or curl-free property of that differential
approximation has been proved. Exact co-moving transport
along a wholly good coefficient path uses
\eqref{r36:exact}, not differentiation of this piecewise
approximation.

\subsection{What has moved, and the remaining propagation problem}

The new statement controls the derivative of a genuine
retained predictor, including its normalizer. It separates
\[
 \begin{aligned}
 &\text{explicit source/frame derivative}
       +\text{motion of the conditional minimum}\\
 &\hspace{4em}+\text{paid thermal deformation}.
 \end{aligned}
\]
The last term is now estimated, pointwise on a good block
and in \(L^2(\mu_{\mathcal R})\) for an arbitrary marked
plaquette bank. The first two terms are not set to zero.
All induced interactions defining \(\mu_{\mathcal R}\)
are unchanged.

At coherence, \eqref{r36:velocity} is the harmonic
response equation on the actual six-vertex graph:
\(M=(6I-\operatorname{Adj}(P_2\times P_3))\otimes I_3\).
For a common slot rotation with fixed transports,
\((6I-\operatorname{Adj}){\bf1}=(n_i)_i\), so its
minimizer velocity reproduces that rotation exactly.
This checks consistency with the inherited coherent
Ward/barycentric identities, not a new spectral obstruction.
Away from coherence the same equation uses the actual
conditional Hessian and signed score gradient.
Its low-frequency retained propagation, under the full
marginal, remains to be controlled.

In particular, \eqref{r36:native-bound} is not an estimate of
\(\operatorname{Cov}_{\mu_{\mathcal R}}(M_b,M_c)\).
Converting a small differential error into a centered
\(L^2\) error of the predictors would already require an
appropriate retained-law functional inequality or a
different justified transport argument. Inserting the
unknown mass-gap estimate at that point would be circular.
Nor does multiplying physical source normalizations into
\eqref{r36:rate} preserve its decay automatically:
any powers of the cutoff and any growing physical support
have to be kept explicitly. The theorem currently controls
one specified native elimination, not an arbitrary number
of RG iterations with changed actions.

The next calculation should estimate the retained
classical differential \(\mathcal A_b\), or construct its
compatible finite-path/source interaction with all generated
terms paid. Repeating single-block concentration or merely
writing another formal score identity would not accomplish
this. The needed new input concerns collective smooth modes
and physical-time propagation; V35's rough geometry and
the present thermal correction are useful inputs, not that
conclusion.

The exact diagnostic
\begin{center}\small
\path{verify_ym_restart_v36_comoving_transport.py}
\end{center}
checks four rational coefficient families with genuinely
nonidentity stabilizer transports, the 18-dimensional
implicit velocity, and a direct independent differentiation
of the vanishing residual gradient. It also checks global
common-rotation cancellation on rational spin configurations;
every native framed word and its retained support on three
periodic lattices; and a normalized Gaussian precision-change
fixture, where both the Poisson centering and the sign of
the response can be evaluated exactly. In the last fixture
the thermal velocity is nonzero and has precisely
\(\gamma^{-1/2}\) scale; no better generic rate is silently
assumed. These diagnostics do not numerically evaluate the
analytic constants, formalize the proof in a proof assistant,
or certify a physical spectrum.

\paragraph{V36 ledger.}
New: full compact displacement moments used with a co-moving
score whose first spin derivative vanishes exactly; an exact
centered Poisson transport with \(O(\gamma^{-1/2})\) thermal
velocity; a normalized co-moving \(W_2\) path estimate;
an explicit retained-link incidence bound; and a
volume-independent actual-law plaquette-predictor response
error of order \(\sqrt{\log\gamma}/\gamma\), with its rare
exterior cost displayed.
Inherited: V31's minimum, V34's full conditional Poincare
bound and plaquette derivative identity, V35's original-law
tail, and the general variational Poisson method.
Still open: collective retained covariance, a uniform
physical-time contraction, a noncollapsing interacting
four-dimensional continuum construction, and the positive
physical Yang--Mills mass gap.
The complete V35 body is preserved below the changed title;
removing this appendix and restoring that title recovers
its exact original bytes.

% END CONSOLIDATED V36 APPEND

% BEGIN CONSOLIDATED V37 APPEND -- 2026-09-19
\clearpage
\section[V37: coherent reconstruction and reflected source comparison]
{V37: coherent reconstruction\\and native reflected source comparison}
\label{r37:start}

\subsection{A finite-path target, not another differential approximation}

V36 isolates the small thermal part of a normalized conditional
response, but its piecewise differential approximation need not
integrate to a scalar source. We now choose one definite coefficient
path and integrate the \emph{exact} transport. This constructs an
actual normalized stochastic reconstruction, retaining the entire
original marginal on the unintegrated links. Good blocks use a
rotated copy of the same coherent six-spin law at coupling
\(\gamma\); rough blocks keep their exact conditional law.

The second step checks time reflection literally. Matching the
constructions on the two half-tori gives a reflection-positive
joint coupling. More importantly, it gives a quantitative
comparison of positive-time sources \emph{under the original
Wilson law}: their conditional noise is null in the reflected
form, and their remaining conditional-mean bias is paid. This
does not assume a retained Poincare inequality.

The distinctions from earlier results are essential. V30's
deterministic filling is singular relative to the original
measure. Our reconstruction keeps compact thermal randomness.
V31 changes selected conditional partition functions and hence
changes the retained marginal; here that marginal is unchanged.
V34's innovation decomposition and archive \(C\), V66's
boundary-predictor identities already distinguish ordinary
variance from reflected norm. Conditional factorization and
the resulting projection identity are standard tools, not
new abstract discoveries. The new estimates concern the
specified six-link reconstruction and its finite-path bias.
Archive \(C\), V57's exact source Feshbach return was also
reread: identifying a coarse source never licenses dropping
its nonreducing physical-time return.

The primary-source check distinguishes two classical tools.
Moser's \href{https://web.ma.utexas.edu/users/dafr/M392C-2018-MorseTheory/Readings/Moser.pdf}
{\emph{On the Volume Elements on a Manifold}}, Section 4,
constructs equivalence of positive volume forms through a
time-dependent vector field. We use V36's already proved
weighted-Poisson velocity, with its quantitative energy bound,
rather than claim a new general transport principle.
Biskup's \href{https://arxiv.org/html/math-ph/0610025v2}
{\emph{Reflection Positivity and Phase Transitions in Lattice
Spin Models}}, Section 5.1, makes the positive-half algebra
and conditional-factorization argument explicit.
The original Wilson reflection positivity remains the
audited input of archive \(C\), V65. No spin-model decay
theorem is transferred. Both scheduled reviews remain V40.

\subsection{A canonical radial coefficient path and its actual coupling}

Retain V36's \(\mathcal M=(S^3)^6\) and
\(K=\{\vartheta:d(\vartheta)\le1/32\}\).
Write \(o\) for coherent coefficients and
\(\lambda_0=\lambda_{\gamma,o}\), at the \emph{same}
coupling \(\gamma\ge1\) as the target law.
For \(\vartheta=(v,T)\in K\), choose the radial path
\[
 v_{i\alpha}(s)=\exp_{e_0}
                   \bigl(s\log_{e_0}v_{i\alpha}\bigr),
 \qquad T_{ij}(s)=\exp(s\log T_{ij}),\qquad 0\le s\le1.
 \tag{V37.1}\label{r37:path}
\]
The logs are the unique small spherical and principal
orthogonal logs. No native retained-link curve realizing
these auxiliary coefficients is assumed or needed.

\begin{lemma}[The radial path remains in the proved coefficient set]
\label{r37:path-lemma}
Put \(d=d(\vartheta)\). For V36's speed \(B(s)\),
\[
       d(\vartheta(s))\le d,\qquad B(s)\le2d,\qquad
       \int_0^1B(s)\,ds\le2d .
 \tag{V37.2}\label{r37:path-bounds}
\]
\end{lemma}
\begin{proof}
A slot with angle \(\alpha\) has chord
\(2\sin(\alpha/2)\), and its intermediate chord is
\(2\sin(s\alpha/2)\), no larger. Its speed is
\(\alpha\le2[2\sin(\alpha/2)]\) in this neighborhood.
For an orthogonal transport, diagonalize into real
rotation planes. The operator-norm chord is
\(2\max_j|\sin(\alpha_j/2)|\), and the logarithm norm
is \(\max_j|\alpha_j|\). The same two inequalities hold
plane by plane. Squaring and summing proves the claims.
\end{proof}

Along this path let \(W_s\), \(V_{{\rm th},s}\) and
\(\mathcal O_s\) be exactly V36's fields and product
rotation, with \(\mathcal O_0=I\).
Let \(\Psi_s\) be the flow of the co-moving thermal
velocity
\[
 \widetilde V_s(y)=O(s)^{-1}
                      V_{{\rm th},s}(\mathcal O_s y),
 \qquad \Psi_0=I,
 \qquad
 \mathcal T_\vartheta=\mathcal O_1\Psi_1,\quad
 \mathcal O_\vartheta=\mathcal O_1.
 \tag{V37.3}\label{r37:maps}
\]
Here \(O(s)^{-1}\) acts on the tangent vector spinwise.
The flows are smooth diffeomorphisms at each finite
\(\gamma\). Smooth dependence on the small coefficients
follows from the unique centered elliptic solution and
the ordinary differential equation.

\begin{theorem}[Normalized coherent reconstruction on a good block]
\label{r37:coupling}
For \(Y\sim\lambda_0\), put
\(Z=\mathcal T_\vartheta Y\) and
\(\overline Z=\mathcal O_\vartheta Y\). Then
\[
 \begin{gathered}
 Z\sim\lambda_{\gamma,\vartheta},\qquad
 \overline Z\sim(\mathcal O_\vartheta)_\#\lambda_0,\qquad
 \mathcal O_\vartheta(e_0,\ldots,e_0)=z_*(\vartheta),\\
 \mathbb E\,\operatorname{dist}_{\mathcal M}
                         (Z,\overline Z)^2
                   \le C d(\vartheta)^2/\gamma .
 \end{gathered}
 \tag{V37.4}\label{r37:coupling-bound}
\]
At \(\vartheta=o\), both maps are the identity.
The construction commutes with simultaneous orthogonal
conjugations of the coefficients fixing \(e_0\).
\end{theorem}
\begin{proof}
V36's continuity equation says
\((\Psi_s)_\#\lambda_0
       =(\mathcal O_s^{-1})_\#\lambda_{\gamma,\vartheta(s)}\).
Applying \(\mathcal O_s\) gives the first assertion,
including the true normalizer, and its minimum-motion
equation gives the third. Since \(\mathcal O_1\) is an
isometry, along the flow of \(\Psi_s\),
\[
 \begin{aligned}
 \{\mathbb E\operatorname{dist}(Z,\overline Z)^2\}^{1/2}
 &=\{\mathbb E\operatorname{dist}(\Psi_1Y,Y)^2\}^{1/2}\\
 &\le\int_0^1
       \{\mathbb E|\widetilde V_s(\Psi_sY)|^2\}^{1/2}ds\\
 &\le C\gamma^{-1/2}\int_0^1 B(s)\,ds .
 \end{aligned}
\]
This is a specified coupling, not just existence of a
Wasserstein minimizer. Equation \eqref{r37:path-bounds}
proves the bound.

At coherence the path is constant, so both vector fields
vanish. An orthogonal transformation fixing \(e_0\)
commutes with the radial logs, the coefficient energy,
the intrinsic metric and the coherent law. Uniqueness
of the minimum, of the centered Poisson solution and
of the flow proves equivariance of both maps.
\end{proof}

The full rotation \(\mathcal O_\vartheta\), not just the
location \(z_*\), is part of this definition. Separate
rotations fixing individual spin minima need not preserve
the interacting coherent law. Thus no unrecorded choice
of six independent residual frames is allowed.
The radial coefficient path fixes that choice.

\subsection{An improved finite-path plaquette bias}

Let \(F_\vartheta(z)\) be any one of the 29 energy summands
of the \emph{terminal} block energy \(E_\vartheta\).
Define two actual expectations
\[
 M_\vartheta=\lambda_{\gamma,\vartheta}(F_\vartheta),
 \qquad
 N_\vartheta=\lambda_0(F_\vartheta\circ\mathcal O_\vartheta),
 \qquad b_\vartheta=M_\vartheta-N_\vartheta .
 \tag{V37.5}\label{r37:means}
\]

\begin{proposition}[Finite-path bias with the source derivative floor]
\label{r37:bias}
Uniformly for \(\vartheta\in K\) and \(\gamma\ge1\),
\[
 \begin{aligned}
 |b_\vartheta|
    &\le\frac{C d}{\sqrt\gamma}
                              \sqrt{d^2+\gamma^{-1}},\\
 |b_\vartheta|^2
    &\le C(d^4/\gamma+d^2/\gamma^2),\\
 \mathbb E|F_\vartheta(Z)-F_\vartheta(\overline Z)|^2
    &\le C d^2/\gamma .
 \end{aligned}
 \tag{V37.6}\label{r37:bias-bounds}
\]
These compare values, not merely differentials.
\end{proposition}
\begin{proof}
Freeze the terminal observable and set
\[
        f_s(z)=F_\vartheta(\mathcal O_1\mathcal O_s^{-1}z).
\]
Its explicit derivative satisfies
\(\partial_s f_s+W_s f_s=0\).
Moreover \(\mathcal O_1\mathcal O_s^{-1}z_*(\vartheta(s))
=z_*(\vartheta)\). At that final minimum,
nonnegativity of all summands and V31's energy bound imply
\(F_\vartheta(z_*)\le m(\vartheta)\le d^2/2\).
The exact plaquette gradient identity therefore gives
\(|\nabla F_\vartheta(z_*)|^2\le2d^2\).
All second spin derivatives of this summand are uniformly
bounded. Transport its gradient along a minimizing
geodesic from the current minimum. Since
\(\mathcal O_1\mathcal O_s^{-1}\) is an isometry,
\[
 |\nabla f_s(z)|\le C\{d+
          \operatorname{dist}(z,z_*(\vartheta(s)))\}.
\]
V36's displacement moment yields
\(\lambda_s(|\nabla f_s|^2)\le C(d^2+\gamma^{-1})\).

Use the exact normalized transport of V36:
\[
       \frac{d}{ds}\lambda_s(f_s)
                 =\lambda_s(V_{{\rm th},s}\cdot\nabla f_s).
\]
The endpoint difference is precisely \(b_\vartheta\).
Cauchy--Schwarz, V36's thermal velocity bound and
\eqref{r37:path-bounds} give the first two inequalities.
Finally the global spin Lipschitz norm of a summand is
at most \(\sqrt2\), since
\(|\nabla F|^2=s(2F-F^2)\le s\le2\).
Apply \eqref{r37:coupling-bound} to obtain the last bound.
\end{proof}

This proof does not integrate the piecewise differential
approximation in V36. It integrates a smooth full
conditional law along one path entirely inside \(K\).
No retained functional inequality or pathwise good-event
probability has been inserted.

\subsection{A genuine full-law coupling with the marginal unchanged}

Use V35's dyadic regulators \(L\ge16\), \(T=2q\ge8\).
Let \(\eta\) be all retained links of the native packing.
For now use V31's frames and defects
\(\mathcal D_b=D_b^2\); a reflection-matched modification
will be specified in the next subsection.
At a good exterior
\(G_b=\{\mathcal D_b\le1/1024\}\), use the coupling
\eqref{r37:coupling-bound} with
\(\vartheta=\vartheta_b(\eta)\).
At a rough exterior sample
\(Z_b\sim\lambda_{\gamma,\vartheta_b(\eta)}\) and set
\(\overline Z_b=Z_b\) exactly.
Conditioned on \(\eta\), make these paired draws independent
between distinct free blocks, and draw \(\eta\) from
its \emph{original} marginal \(\mu_{\mathcal R}\).
Call the resulting joint law \(\Gamma\).

Undo the exact Haar-preserving block frames to obtain full
configurations \(U,\overline U\), with their retained links
both equal to \(\eta\). Its first marginal is exactly
\(\mu_W\); denote the second by \(\overline\mu_W\).
The second marginal is normalized and has exactly
\(\mu_{\mathcal R}\) as retained marginal. Its conditional
block laws are
\[
 \overline\lambda_b=
 \begin{cases}
  (\mathcal O_{\vartheta_b})_\#\lambda_0,&\eta\in G_b,\\
  \lambda_{\gamma,\vartheta_b},&\eta\notin G_b.
 \end{cases}
 \tag{V37.7}\label{r37:hybrid}
\]
Both laws have positive densities relative to full product
Haar at every finite coupling; unlike deterministic
filling, this construction is not singular.

All normalization factors remain present. In framed
coordinates the exact density ratio is
\[
 \frac{d\overline\mu_W}{d\mu_W}(\eta,z)
 =\prod_{b:\eta\in G_b}
   \frac{Q_\gamma(\vartheta_b)}{Q_\gamma(o)}
       \exp\{\gamma E_{\vartheta_b}(z_b)
              -\gamma E_o(\mathcal O_{\vartheta_b}^{-1}z_b)\}.
 \tag{V37.8}\label{r37:density}
\]
Every factor has conditional expectation one under its
original block law. This proves normalization and unchanged
retained marginal without assuming a small global
likelihood ratio. The products are over a finite regulator.
No volume-independent total-variation bound is asserted.

These kernels are measurable at the threshold and smooth
on the interiors of the two regions. Native gauge
transformations conjugate all framed spins and coefficients
by the common root conjugation fixing \(e_0\).
The equivariance in Theorem~\ref{r37:coupling}, and gauge
invariance of \(\mathcal D_b\), show that the reconstructed
law and coupling are gauge covariant. The original
configuration-space measure is not gauge-fixed or given
a missing Jacobian.

For a marked plaquette \(F_b\) from each of distinct blocks,
write
\[
 \begin{gathered}
 M_b=\mathbb E_\Gamma[F_b(U)\mid\eta],\qquad
 N_b=\mathbb E_\Gamma[F_b(\overline U)\mid\eta],\\
 b_b=M_b-N_b,\qquad
 \Delta_b=F_b(U)-F_b(\overline U),\qquad
 \xi_b=\Delta_b-b_b .
 \end{gathered}
 \tag{V37.9}\label{r37:decomposition}
\]
On rough exteriors \(\Delta_b=b_b=\xi_b=0\).
Define the explicit moment budgets
\[
 \begin{aligned}
 A_\gamma&=\frac{122880\mathcal B_\gamma}{\gamma},
 &\ell_\gamma&=\frac{\gamma}{640},
 &\mathcal B_\gamma&=8+\log\gamma-\tfrac23\log\frac8{3\pi^3},\\
 m_{1,\gamma}&=A_\gamma+\ell_\gamma^{-1},
 &m_{2,\gamma}&=A_\gamma^2+
                  2A_\gamma\ell_\gamma^{-1}+2\ell_\gamma^{-2},\\
 \alpha_\gamma&=\frac{m_{1,\gamma}}{\gamma},
 &\beta_\gamma^2&=\frac{m_{2,\gamma}}{\gamma}
                                     +\frac{m_{1,\gamma}}{\gamma^2}.
 \end{aligned}
 \tag{V37.10}\label{r37:budgets}
\]
Thus \(\alpha_\gamma=O((1+\log\gamma)/\gamma^2)\) and
\(\beta_\gamma=O((1+\log\gamma)/\gamma^{3/2})\).
V35 gives the unconditioned bounds
\(\mathbb E_{\mu_{\mathcal R}}\mathcal D_b\le m_{1,\gamma}\)
and
\(\mathbb E_{\mu_{\mathcal R}}\mathcal D_b^2\le m_{2,\gamma}\).

\begin{theorem}[Innovations and retained bias are paid separately]
\label{r37:ordinary}
For arbitrary deterministic complex coefficients \(c_b\),
\[
 \begin{aligned}
 \mathbb E_\Gamma\left|\sum_b c_b\xi_b\right|^2
                         &\le C\alpha_\gamma\|c\|_2^2,\\
 \mathbb E_{\mu_{\mathcal R}}\left|\sum_b c_b b_b\right|^2
                         &\le C\beta_\gamma^2\|c\|_1^2,\\
 \mathbb E_\Gamma\left|\sum_b c_b\Delta_b\right|^2
   &=\mathbb E_\Gamma\left|\sum_b c_b\xi_b\right|^2
       +\mathbb E_{\mu_{\mathcal R}}\left|\sum_b c_b b_b\right|^2\\
   &\le C\{\alpha_\gamma\|c\|_2^2+
                              \beta_\gamma^2\|c\|_1^2\}.
 \end{aligned}
 \tag{V37.11}\label{r37:ordinary-bounds}
\]
There is no all-good event, discarded rough block, or
multiplication of dependent defect probabilities.
\end{theorem}
\begin{proof}
Given \(\eta\), the paired block draws are independent and
each \(\xi_b\) has zero mean. Thus their conditional
cross covariances vanish. Their variances are bounded
by \(\mathbb E[|\Delta_b|^2\mid\eta]\).
On a good block \eqref{r37:bias-bounds} and
\(d^2\le\mathcal D_b\) bound this by
\(C\mathcal D_b/\gamma\); on a rough block it is zero.
This proves the first assertion. The same proposition gives
\[
 |b_b|^2\le C{\bf1}_{G_b}
        (\mathcal D_b^2/\gamma+\mathcal D_b/\gamma^2),
\]
so \(\|b_b\|_{L^2(\mu_{\mathcal R})}\le C\beta_\gamma\).
The triangle inequality in that actual \(L^2\) space
proves the second assertion, without a covariance assumption.
Conditional centering makes the innovation sum orthogonal
to every retained function, giving the exact third identity.
\end{proof}

The \(\ell^1\) bias cost has not been silently changed
to \(\ell^2\). Conditional independence of the noises does
not make the retained biases independent. Keeping rough
blocks exact removes a \emph{comparison error} there;
it does not remove their contribution to \(N_b\) or
their remaining correlations.

\subsection{Reflection-matched kernels and the native quotient}

Let \(\theta\) be site time reflection \(t\mapsto-t\)
on the periodic time circle, with fixed planes \(0,q\).
It acts on oriented temporal links with inversion and
on spatial links by reflection of their base point.
Write \(\Theta F(U)=\overline{F(\theta U)}\).
The positive open half is \(0<t<q\).

\begin{lemma}[The actual block conditionals do not cross a reflection plane]
\label{r37:half-geometry}
The native free-link packing is invariant under \(\theta\)
and has no free link on either fixed plane.
Every block belongs to one open half. Its 29 incident
plaquettes and the exterior links in its conditional
law stay in that same open half.
For a positive-half block, V31's frame and V35's
45-face witness also stay in that half.
\end{lemma}
\begin{proof}
Free block times are even and are not \(0,q\), with \(q\)
even. Positive block times consequently lie between
2 and \(q-2\); reflection pairs them with negative blocks.
All incident plaquettes have vertices only at block time
or its two neighbors. The frame tree is spatial, and
the forward staple reaches time \(t+1\). V35's full
witness carrier has temporal offsets in \(\{-1,0,1\}\).
Its supporting vertices therefore have times in
\(\{1,\ldots,q-1\}\). The same conclusion for the
actual negative conditional follows by reflection.
No plaquette contains free links from different blocks.
\end{proof}

For reflection positivity use the following \emph{matched}
convention, rather than the forward-time frame in both halves.
On every positive block use the construction above.
Define the frame, good event and paired kernel on its
negative partner by reflecting that positive construction.
The negative frame then uses the backward-time staple.
This changes no conditional law of \(U\); it only fixes
the reconstruction rule for \(\overline U\).
One-point defect moments on negative blocks equal those
on their reflected positive partners by invariance of
\(\mu_W\). Thus \eqref{r37:ordinary-bounds} still holds.
A joint V35 tail for mixed forward/backward witnesses
is neither used nor asserted.

The retained marginal is reflection positive on its
retained positive-half algebra: extend such an observable
to the full configuration without using free links
and apply the original Wilson reflection positivity.

\begin{proposition}[Positive joint reconstruction]
\label{r37:RP}
With the matched convention, the joint law \(\Gamma\)
and its second marginal \(\overline\mu_W\) are reflection
positive for this fixed reflection.
For any bounded positive-half joint observable \(A\), put
\(\mathcal K A(\eta)=\mathbb E_\Gamma[A\mid\eta]\).
It is a retained positive-half observable, and
\[
       \mathbb E_\Gamma[(\Theta A)A]
          =\mathbb E_{\mu_{\mathcal R}}
                    [(\Theta\mathcal K A)\mathcal K A]\ge0.
 \tag{V37.12}\label{r37:RP-identity}
\]
\end{proposition}
\begin{proof}
Given \(\eta\), paired noises in the two halves are
independent. All their kernels depend only on retained
links in the corresponding half, and reflection takes
one half's kernels exactly to the other's. Conditional
expectation of \((\Theta A)A\) therefore factors into
\((\Theta\mathcal K A)\mathcal K A\).
Reflection positivity of \(\mu_{\mathcal R}\) proves
the result. Reflection invariance follows from the same
kernel covariance. A function of the second full
configuration alone is a joint observable, proving
the marginal assertion.
\end{proof}

For the first, original marginal the same proof gives
\[
  \mathbb E_{\mu_W}[(\Theta F)F]
       =\mathbb E_{\mu_{\mathcal R}}
                   [(\Theta P F)PF],\qquad
             PF=\mathbb E_{\mu_W}[F\mid\eta],
 \tag{V37.13}\label{r37:projection}
\]
for every bounded positive-half \(F\). This is not
a claim that an arbitrary spatial sigma-field is a
Markov separator: it uses the verified free-block
factorization on the entire retained history.
In particular \(F-PF\) is null in this fixed reflected
form, although it usually has positive ordinary variance.
Conditional projection identifies the corresponding
fixed-reflection quotients with that of \(\mu_{\mathcal R}\);
every retained positive observable is already in its range.

Neither this identification nor \eqref{r37:RP-identity}
intertwines one-step physical time translation.
The packing and matched frames single out the planes
and are not invariant under arbitrary time shifts.
No new Hamiltonian or equality of physical transfer
spectra is inferred from equality of these quotients.

\subsection{A quantitative source comparison in the original reflected form}

Choose distinct marked blocks in the positive open half,
and for each choose one incident plaquette energy \(F_b(U)\).
Let \(N_b(\eta)\) be its hybrid conditional mean from
\eqref{r37:decomposition}, using the matched convention.
It is a bounded, gauge-invariant, retained positive-half
observable, including its original conditional predictor
on rough exteriors. For \(c=(c_b)\), put
\[
 F_c=\sum_b c_bF_b,\qquad N_c=\sum_b c_bN_b,\qquad
 B_c=\sum_b c_b b_b .
\]
Let \(H^\circ=H-\mathbb E_{\mu_W}H\) for native
observables, and use the same original marginal for
retained observables. Define
\(\|H\|_{\rm OS}^2=\mathbb E_{\mu_W}[(\Theta H)H]\).

\begin{theorem}[Native reflected source error is the paid bias]
\label{r37:OS}
Exactly,
\[
 \begin{aligned}
 \|F_c-N_c\|_{\rm OS}^2
       &=\mathbb E_{\mu_{\mathcal R}}
                                     [(\Theta B_c)B_c],\\
 \|F_c^\circ-N_c^\circ\|_{\rm OS}^2
       &=\mathbb E_{\mu_{\mathcal R}}
                     [(\Theta B_c)B_c]
                         -|\mathbb E_{\mu_{\mathcal R}}B_c|^2 .
 \end{aligned}
 \tag{V37.14}\label{r37:OS-exact}
\]
Both are nonnegative, and
\[
 \boxed{\quad
 \|F_c^\circ-N_c^\circ\|_{\rm OS}
       \le C\beta_\gamma\|c\|_1
       \le C\frac{1+\log\gamma}{\gamma^{3/2}}\|c\|_1 .
 \quad}
 \tag{V37.15}\label{r37:OS-bound}
\]
The same bound without centering also holds.
These are comparisons under the original Wilson law,
not a replacement law's spectral assertion.
\end{theorem}
\begin{proof}
The conditional projection of \(F_c-N_c\) is \(B_c\).
Equation \eqref{r37:projection} proves the first identity.
Subtracting its native mean is orthogonal projection
off the constant vector of reflected norm one, proving
the second. Reflection positivity gives both signs.
By ordinary Cauchy--Schwarz and reflection invariance,
\[
  \left|\mathbb E_{\mu_{\mathcal R}}
                      [(\Theta B_c)B_c]\right|
          \le\mathbb E_{\mu_{\mathcal R}}|B_c|^2
          \le C\beta_\gamma^2\|c\|_1^2 .
\]
Theorem~\ref{r37:ordinary} gives the last step.
Centering can only decrease the squared reflected norm.
\end{proof}

The coupled innovation does not appear in
\eqref{r37:OS-bound}. In fact its reflected pairing is
zero by conditional centering and the half-factorization.
This cancellation is specific to the verified reflection
placement. It is not an assertion that all ordinary
conditional covariances vanish.

For a bank of \(M\) marked blocks with \(\|c\|_2=1\),
the sufficient absolute error is \(C\sqrt M\,\beta_\gamma\).
Thus \(M(1+\log\gamma)^2/\gamma^3\to0\) suffices for
this bare reflected comparison to vanish. This is not
a physical-window criterion until the relation between
\(M\), lattice spacing and the source normalization is
specified. Multiplying a source by \(Z(a)\) multiplies
the error by \(|Z(a)|\); a small bare error need not
survive the intended continuum normalization.

On the fixed-reflection quotient let \(V,\widehat V\)
be the centered source syntheses for \(F_b,N_b\).
The estimate gives \(\|V-\widehat V\|\le C\sqrt M\,\beta_\gamma\).
Consequently, if \(\epsilon=C\sqrt M\,\beta_\gamma\),
\[
 \|V^*V-\widehat V^*\widehat V\|
                    \le2\|V\|\epsilon+\epsilon^2 .
 \tag{V37.16}\label{r37:Gram}
\]
This follows just by expanding the Gram difference.
The same generic Gram estimate is already used in the
archives; the new input is the evaluated native error
feeding it. There is still no lower bound on the
positive-time Gram or contraction of the resulting
retained source family.

\subsection{Explicit coherent moments and a native minimum-energy source}

The retained mean need not remain an unevaluated
conditional integral. Write the coherent sample as
\(Y_i=(Y_i^0,Y_i^1,Y_i^2,Y_i^3)\), and define, for
an internal edge \((i,j)\),
\[
 m_{i,\gamma}=\mathbb E_{\lambda_0}Y_i^0,\qquad
 a_{ij,\gamma}=\mathbb E_{\lambda_0}Y_i^0Y_j^0,\qquad
 c_{ij,\gamma}=\mathbb E_{\lambda_0}Y_i^1Y_j^1 .
 \tag{V37.17}\label{r37:coherent-moments}
\]
These numbers depend on the fixed six-vertex graph
and on \(\gamma\), not on the exterior.
The diagonal transverse \(SO(3)\) symmetry gives
\[
 \mathbb E_{\lambda_0}Y_i=m_{i,\gamma}e_0,\qquad
 \mathbb E_{\lambda_0}[Y_jY_i^{\mathsf T}]
   =c_{ij,\gamma}I_4+
                 (a_{ij,\gamma}-c_{ij,\gamma})e_0e_0^{\mathsf T}.
 \tag{V37.18}\label{r37:coherent-tensor}
\]
Mixed longitudinal--transverse moments vanish and
the three transverse diagonal moments agree.

Write \(O_i\) for the \(i\)-th component of
\(\mathcal O_\vartheta\), so \(O_i e_0=z_i^*\).
For a private slot and an internal edge respectively,
put
\[
 \begin{aligned}
 e_{i\alpha}^*&=1-z_i^*\cdot v_{i\alpha},&
 h_{i,\gamma}&=1-m_{i,\gamma},\\
 e_{ij}^*&=1-z_i^*\cdot T_{ij}z_j^*,&
 h_{ij,\gamma}&=1-a_{ij,\gamma}-3c_{ij,\gamma},\\
 S_{ij}&=O_i^{\mathsf T}T_{ij}O_j,&
 \kappa(S)&=4-\operatorname{Tr}S
                       =\tfrac12\|S-I_4\|_{\mathrm F}^2 .
 \end{aligned}
 \tag{V37.19}\label{r37:offsets}
\]
The last identity holds for every \(S\in SO(4)\).
The constants \(h_{i,\gamma}\) and \(h_{ij,\gamma}\)
are the corresponding coherent plaquette means.

\begin{proposition}[Exact coherent writers and their classical part]
\label{r37:writers}
On a good exterior the hybrid conditional means satisfy
\[
 \begin{aligned}
 N_{i\alpha}-h_{i,\gamma}&=m_{i,\gamma}e_{i\alpha}^*,\\
 N_{ij}-h_{ij,\gamma}
       &=(a_{ij,\gamma}-c_{ij,\gamma})e_{ij}^*
                                  +c_{ij,\gamma}\kappa(S_{ij}).
 \end{aligned}
 \tag{V37.20}\label{r37:writers-exact}
\]
All three coefficients \(m_{i,\gamma}\),
\(c_{ij,\gamma}\), and \(a_{ij,\gamma}-c_{ij,\gamma}\)
are strictly positive for \(\gamma>0\).
Uniformly for \(\gamma\ge1\),
\[
 \begin{gathered}
 0\le1-m_{i,\gamma}\le C/\gamma,\qquad
 0\le1-(a_{ij,\gamma}-c_{ij,\gamma})\le C/\gamma,\\
 0<c_{ij,\gamma}\le C/\gamma,\qquad
 0\le h_{i,\gamma},h_{ij,\gamma}\le C/\gamma .
 \end{gathered}
\]
For any one of the 29 incident plaquettes, with its
corresponding constant \(h_{b,\gamma}\),
\[
 0\le N_b-h_{b,\gamma},\qquad
 |N_b-h_{b,\gamma}-e_b^*|
                              \le C d(\vartheta)^2/\gamma .
 \tag{V37.21}\label{r37:writer-error}
\]
\end{proposition}
\begin{proof}
For a private slot use \(\mathbb E O_iY_i=m_{i,\gamma}z_i^*\).
For an edge, the expected dot product is
\[
 \mathbb E[Y_i^{\mathsf T}S_{ij}Y_j]
   =c_{ij,\gamma}\operatorname{Tr}S_{ij}
      +(a_{ij,\gamma}-c_{ij,\gamma})
                              e_0^{\mathsf T}S_{ij}e_0 .
\]
Subtracting its coherent value proves
\eqref{r37:writers-exact} exactly.

For clarity, positivity is not deduced merely from
rotational symmetry. Use the classical sphere-pairing
expansion already employed in the archive and in V29.
Apart from a constant, the coherent density is
\[
 \exp\!\left(\gamma\sum_i n_iY_i^0+
                         \gamma\sum_{(i,j)}Y_i\cdot Y_j\right),
\]
where every private multiplicity \(n_i\) is positive.
Expand each field and edge exponential in its power
series, and integrate the sphere monomials by the
even-moment pairing formula. Every pairing has a
positive weight and joins strand endpoints in pairs;
odd moments vanish. The resulting components are
closed loops or paths whose endpoints are field
insertions or the displayed observable marks.
Absolute convergence follows from compactness,
the factorial denominators and the finiteness of the
graph, so these contractions can be grouped by the
components containing the marks.

For one mark the open path ends at a field insertion,
giving a nonnegative multiple of \(e_0\). The term
with one field at that vertex is strictly positive.
For two marks, either their path joins them to each
other, giving a nonnegative multiple of \(I_4\),
or their separate paths end at field insertions,
giving a nonnegative multiple of \(e_0e_0^{\mathsf T}\).
For adjacent vertices the single-edge term makes the
first coefficient strictly positive, and one field
at each marked vertex makes the second strictly
positive. All remaining loops and field-to-field
paths have nonnegative contractions. Dividing by
the positive partition function and comparing with
\eqref{r37:coherent-tensor} proves the three asserted
signs. This is a finite-graph ferromagnetic pairing
argument, not a new mass-gap inequality.

The coherent V36 displacement estimate yields
\(1-m_{i,\gamma}\le C/\gamma\) and
\(\mathbb E(Y_i^r)^2\le C/\gamma\) for \(r=1,2,3\).
Thus \(c_{ij,\gamma}\le C/\gamma\) by
Cauchy--Schwarz. For \(x,y\in[-1,1]\),
\(1-xy\le(1-x)+(1-y)\), so
\(1-a_{ij,\gamma}\le
 (1-m_{i,\gamma})+(1-m_{j,\gamma})\).
Since \(a_{ij,\gamma}\le1\) and \(c_{ij,\gamma}>0\),
the stated estimate for \(1-(a-c)\) follows.
The \(h\)'s are nonnegative plaquette energies and
are bounded by the coherent expected total energy,
which V36 bounds by \(C/\gamma\).

Along the prescribed radial path,
\(\|O_i-I_4\|_{\rm op}\le\int|k_i(s)|\,ds\le C d(\vartheta)\).
The coefficient defect also bounds
\(\|T_{ij}-I_4\|_{\rm op}\), hence
\(\kappa(S_{ij})\le C d(\vartheta)^2\).
Each individual minimum energy is at most the
nonnegative total minimum, which is at most
\(d(\vartheta)^2/2\). Applying the coefficient
estimates to \eqref{r37:writers-exact} proves the
absolute error in \eqref{r37:writer-error}.
Its sign follows from the positive coefficients,
\(e_b^*\ge0\), and \(\kappa\ge0\).
\end{proof}

The trace term records the full relative rotations
of the canonical transport, not only the locations
of the minima. It may be bounded as above, but must
not be silently identified with \(e_{ij}^*\).

For each chosen native positive-half plaquette define
the retained observable
\[
 J_b(\eta)=
 \begin{cases}
   e_b^*(\vartheta_b(\eta)),&\eta\in G_b,\\
   M_b(\eta)-h_{b,\gamma},&\eta\notin G_b ,
 \end{cases}
 \qquad J_c=\sum_b c_bJ_b .
 \tag{V37.22}\label{r37:J}
\]
This definition chooses no minimum on rough exteriors.
It is bounded, measurable, gauge invariant and local
to the same retained positive half; smoothness at
the threshold is neither asserted nor needed.

\begin{corollary}[Native minimum-energy source comparison]
\label{r37:minimum-source}
With the unchanged original marginal and centering,
\[
 \begin{aligned}
 \|N_c^\circ-J_c^\circ\|_{\rm OS}
       &\le C\frac{\sqrt{m_{2,\gamma}}}{\gamma}\|c\|_1,\\
 \|F_c^\circ-J_c^\circ\|_{\rm OS}
       &\le C\left(\beta_\gamma+
                         \frac{\sqrt{m_{2,\gamma}}}{\gamma}\right)\|c\|_1
        \le C\frac{1+\log\gamma}{\gamma^{3/2}}\|c\|_1 .
 \end{aligned}
 \tag{V37.23}\label{r37:minimum-source-bound}
\]
\end{corollary}
\begin{proof}
Let \(R_b=N_b-h_{b,\gamma}-J_b\).
It vanishes exactly on rough exteriors and is bounded
by \(C\mathcal D_b/\gamma\) on good exteriors.
The V35 second moment gives
\(\|R_b\|_{L^2(\mu_{\mathcal R})}\le C\sqrt{m_{2,\gamma}}/\gamma\).
A sum is bounded by the triangle inequality with
\(\|c\|_1\). The centered reflected norm is bounded
by the ordinary \(L^2\) norm, by reflection invariance,
Cauchy--Schwarz and projection off the constant.
The constants \(h_{b,\gamma}\) disappear under global
centering, proving the first line. Theorem~\ref{r37:OS}
and the triangle inequality prove the second.
Finally \(m_{2,\gamma}=O((1+\log\gamma)^2/\gamma^2)\) and
\(\gamma\ge1\) give the displayed rate.
\end{proof}

The coherent offset is subtracted globally, not only
on the good event: dropping
\(h_{b,\gamma}\mathbf1_{G_b}\) would leave a random
indicator term. The rough definition in \eqref{r37:J}
is what makes the global subtraction legitimate.
Nor does this corollary contradict V30's singularity
obstruction to deterministic minimum filling.
No probability measure has been replaced by a
point mass. Instead an observable is compared in
the original reflected form, its conditional noise
projects to a null vector, a constant is centered
away, and its remaining mean error is explicitly paid.

\subsection{Checks and the remaining nonperturbative task}

The accompanying exact diagnostic is
\begin{center}\small
\path{verify_ym_restart_v37_reflected_reconstruction.py}.
\end{center}
It verifies the temporal orientation reversal, the
29-plaquette block conditionals, the full 45-face
witness support and reflection pairing on four tori.
A 64-atom rational probability fixture keeps the
retained law unchanged while changing one good
conditional and retaining a rough one. All 255
principal minors of its complete positive-half joint
Gram are checked, together with its positive
factorization. Its ordinary innovation variance is
\(15/32\), while its reflected innovation pairing is
exactly zero. The centered reflected source error is
\(7/256\), entirely from the retained bias. Repeated
conditionally independent blocks share a correlated
bias, exhibiting its quadratic block-count cost.
Shifted-exponential moments and a Gaussian scale
coupling are separate exact checks of the budgets.
A rational isotropic moment fixture independently
checks both coherent writer identities and the
rotation-trace identity. It is an algebraic fixture,
not an evaluation of the coherent Gibbs moments.
None numerically solves the compact conditional flow,
evaluates the analytic constants, or certifies a
continuum spectrum.

The remaining native quantity is now concrete:
the reflected correlations of \(J_b(\eta)\), or its
explicit coherent refinement \(N_b(\eta)\), under
the unchanged interacting \(\mu_{\mathcal R}\),
including the exact rough-sector predictors.
The reconstruction does not turn this marginal into
a product, and the common coherent noise is not a
vacuum state for the full theory. The previous
classical collective motion survives inside
\(\mathcal O_{\vartheta_b(\eta)}\) and the induced
marginal. To infer the full goal one still needs
noncollapsing physically normalized sources,
control of their collective long-distance return,
a valid physical-time contraction, and the
interacting four-dimensional continuum construction.

A next useful calculation should attack these
retained reflected means and their actual induced
interaction, or prove a compatible physical-time
return estimate. It should not reprove that the
conditional noises are small or discard the
\(\ell^1\) bias cost. If a proposed argument invokes
a retained spectral gap to upgrade the present
comparison, it has imported the missing conclusion.

\paragraph{V37 ledger.}
New: a fixed radial coefficient path and actual
normalized coherent coupling; a finite-path
plaquette-mean bias
\(O(d^2/\sqrt\gamma+d/\gamma)\); a full stochastic
reconstruction retaining the original marginal
and every rough conditional; reflection-matched
kernels; and a native positive-time source
comparison with absolute error
\(O((1+\log\gamma)\gamma^{-3/2}\|c\|_1)\).
Also new: explicit coherent moment writers and a
comparison, at the same reflected rate, to the
native classical minimum-energy source with its
rough-sector predictor kept exact.
Inherited and credited: transport by continuity
flows, conditional-factorization reflection
positivity, V34's source derivative floor,
V35's one-point moments and V36's exact velocity.
Not proved: a dimension-free bound for correlated
bias sums, translation-compatible spectral
equivalence, a positive physical source floor,
or the continuum Yang--Mills mass gap.
Removing this appendix and restoring the displayed
V36 title recovers the entire predecessor exactly.

% END CONSOLIDATED V37 APPEND

% BEGIN CONSOLIDATED V38 APPEND -- 2026-09-19
\clearpage
\section[V38: thermal cancellation and native block-energy sources]
{V38: thermal cancellation\\and native block-energy sources}
\label{r38:start}

\subsection{The upstream coefficient test and its precise scope}

V37 provides a source comparison under the original Wilson law,
but its general plaquette estimate does not use all cancellations
of a \emph{whole conditional-energy insertion}. V33 already
identifies the Maxwell Schur scale of the leading retained action.
Repeating that Hessian computation would not control its
long-distance correlations. Instead we now test a potentially
dangerous source coefficient before bounding it.

This is the transferable part of the ESM V23 suggestion:
identify a forbidden or field-independent low-order term in the
actual normalized quantity. No ESM coupling power, BV identity,
or perturbative continuum theorem is imported. The scalar
source used here is the sum of the 29 plaquette energies incident
to one specified six-link block. It is not an arbitrary single
plaquette, not the entire lattice action, and not asserted to
span all physical spectral sectors.

Equipartition itself is classical. Archive f1, V122 already uses
Laplace equipartition for a single \(SU(3)\) Wilson well; archive
f16.1, V83 uses a global quadratic plaquette check. Here the
number \(9\) is half the \emph{18 real conditional spin dimensions};
it is not a count of global physical polarizations. The new
work is a uniform, twice coefficient-differentiable and twice
coupling-differentiable remainder for the native compact block,
its exact coherence subtraction, an explicit first correction,
and the resulting original-law reflected source bounds.
V32's first-derivative cancellation and V37's reflected projection
are credited inputs, not new abstract principles.

As a targeted primary-source check, Katsevich's
\href{https://arxiv.org/html/2406.12706v3}
{\emph{The Laplace asymptotic expansion in high dimensions}},
Theorem 2.12 and its Gaussian coefficient formula, confirm the
standard even-power structure and the need to control remainders.
We prove the mixed derivative bounds below directly in fixed
dimension; that theorem is not used to justify an unproved
coupling derivative or a growing-volume saddle expansion.
The scheduled broad literature and companion reviews remain V40.

\subsection{Mixed coupling derivatives of the normalized remainder}

Use V31's \(K=\{\vartheta:d(\vartheta)\le1/32\}\),
minimum \(m(\vartheta)\), Hessian \(M(\vartheta)\), and
\(Q_\gamma(\vartheta)=\int e^{-\gamma E_\vartheta}\,dH\).
Let \(o\) be coherence and put
\[
 f_\gamma(\vartheta)=
  \log Q_\gamma(\vartheta)+\gamma m(\vartheta)+9\log\gamma
   -\log(8/\pi^3)+\tfrac12\log\det M(\vartheta).
 \tag{V38.1}\label{r38:f}
\]
This is the same exact logarithmic remainder as V32.
Norms in coefficient space use its small spherical and
principal-orthogonal logarithm coordinates \(a\).
On \(K\), \(|a|\le4d(\vartheta)\); the radial segment
\(sa\) remains in \(K\) by Lemma~\ref{r37:path-lemma}.
A \(C^2_a\) bound on \(K\) means the restriction of
coordinate derivatives from an open neighborhood.

\begin{theorem}[A mixed-derivative Laplace remainder]
\label{r38:mixed}
There is a smooth scalar \(a_1(\vartheta)\) on a neighborhood
of \(K\) and finite graph constants \(C_j\), independent of
\(\gamma\ge1\), such that, for \(j=0,1,2\),
\[
 \left\|\partial_\gamma^j
       \{f_\gamma-a_1/\gamma\}\right\|_{C^2_a(K)}
       \le C_j\gamma^{-2-j},\qquad
 \|\partial_\gamma^j f_\gamma\|_{C^2_a(K)}
       \le C_j\gamma^{-1-j}.
 \tag{V38.2}\label{r38:mixed-bounds}
\]
Moreover \(D_af_\gamma(o)=D_aa_1(o)=0\).
Writing \(\Delta a_1=a_1(\vartheta)-a_1(o)\),
\[
 \begin{aligned}
 |\Delta a_1|&\le C d(\vartheta)^2,\\
 \left|\partial_\gamma^j
  \left\{f_\gamma(\vartheta)-f_\gamma(o)
                           -\frac{\Delta a_1}{\gamma}\right\}\right|
     &\le C_j d(\vartheta)^2\gamma^{-2-j}.
 \end{aligned}
 \tag{V38.3}\label{r38:centered-remainder}
\]
\end{theorem}
\begin{proof}
V31 gives smooth unique minima on a neighborhood of the
compact set \(K\), a uniformly positive conditional Hessian
and a uniform barrier outside a fixed minimum neighborhood.
Thus the exact-phase construction in V32 applies on all of
\(K\), not only its smaller coordinate ball. To be explicit,
choose smooth tangent frames at the minima, which lie in
one small product cap, and product normal coordinates \(y\).
For \(F_a(y)=E_a(\exp_{z_*(a)}y)-m(a)\), set
\[
 B_a(y)=2\int_0^1(1-s)D_y^2F_a(sy)\,ds,\qquad
 u=B_a(y)^{1/2}y .
\]
On one sufficiently small common ball this is smoothly
invertible, \(F_a=|u|^2/2\), and all derivatives of each
fixed order are bounded uniformly in \(a\).
Its density \(A(a,u)\) satisfies
\(A(a,0)=(2\pi^2)^{-6}(\det M(a))^{-1/2}\).
Use a fixed even cutoff \(\psi(u)\), equal to one near zero.

Taylor-expand \(A\) to degree three. The odd terms have
zero Gaussian integral. After multiplication by \(\psi\),
the degree-four remainder and its first two coefficient
derivatives are bounded by \(C|u|^4\). Replacing the
cutoff polynomial terms by their full Gaussian integrals
has exponentially small error. For \(j=0,1,2\),
\[
 \partial_\gamma^j\{\gamma^9e^{-\gamma|u|^2/2}\}
 =\gamma^{9-j}p_j(\gamma|u|^2/2)e^{-\gamma|u|^2/2},
 \tag{V38.4}\label{r38:kernel-derivative}
\]
where \(p_j\) is a fixed polynomial. Integrating
\(|u|^4\) against these absolute kernels gives
\(C_j\gamma^{-2-j}\). The same proof works after two
coefficient derivatives. This differentiates the actual
integral with an explicit dominating bound, not a scalar
big-O assertion.

The complementary original integral, after extracting
\(e^{-\gamma m}\), has a positive uniform phase barrier.
Up to two coefficient derivatives cost a polynomial in
\(\gamma\); up to two coupling derivatives insert bounded
powers of \(E-m\), and derivatives of \(\gamma^9\).
All these errors are polynomial times \(e^{-c\gamma}\).
We have therefore proved
\[
 \frac{Q_\gamma(a)e^{\gamma m(a)}\gamma^9}
      {8/(\pi^3\sqrt{\det M(a)})}
   =1+\frac{a_1(a)}{\gamma}+r_\gamma(a),\qquad
 \|\partial_\gamma^j r_\gamma\|_{C^2_a}\le C_j\gamma^{-2-j}.
 \tag{V38.5}\label{r38:relative-integral}
\]
The leading factor is uniformly positive. The chain and
product rules for the logarithm give
\eqref{r38:mixed-bounds} at sufficiently large \(\gamma\).
On the remaining compact coupling interval, positivity
and smoothness of the original integral give the bounds
after increasing the constants.

V32 proves \(D_af_\gamma(o)=0\) for every positive coupling.
The \(C^2\) convergence
\(\gamma f_\gamma\longrightarrow a_1\) implies
\(D_aa_1(o)=0\). Smooth mixed differentiation is legitimate
for the compact integral and all terms in \eqref{r38:f}.
Consequently the expression in braces in
\eqref{r38:centered-remainder}, and its indicated coupling
derivatives, vanish together with their first coefficient
derivative at \(o\). Taylor's formula on \(sa\), with the
proved \(C^2_a\) bounds and \(|a|\le4d\), proves
\eqref{r38:centered-remainder}.
\end{proof}

\subsection{The first coefficient, including the Haar term}

Let \(X\) be the centered Gaussian in the tangent space at
\(z_*(\vartheta)\) with covariance \(G=M(\vartheta)^{-1}\).
In the product normal coordinates above, write
\[
 E_\vartheta(\exp_{z_*}y)-m
 =\tfrac12\langle y,My\rangle+
       \tfrac16 T_3[y,y,y]+\tfrac1{24}T_4[y,y,y,y]+O(|y|^5).
\]
The tensors here are derivatives of the normal-coordinate
phase; this convention removes ambiguities between arbitrary
coordinate and covariant higher derivatives.

\begin{proposition}[An intrinsic finite Gaussian contraction]
\label{r38:a1}
With \(X_i\in\mathbb R^3\) the tangent variable of spin \(i\),
\[
 a_1(\vartheta)=
 \mathbb E_G\left[
    -\frac13\sum_{i=1}^6|X_i|^2
    -\frac1{24}T_4[X,X,X,X]
    +\frac1{72}\{T_3[X,X,X]\}^2
 \right].
 \tag{V38.6}\label{r38:a1formula}
\]
At coherence put \(g=\mathsf A^{-1}\), where
\(\mathsf A=6I-\operatorname{Adj}(P_2\times P_3)\).
Then
\[
 \begin{aligned}
 a_1(o)={}&-\sum_i g_{ii}
       +\frac{15}{4}\sum_i g_{ii}^2\\
 &+\sum_{\{i,j\}}
 \left\{\frac94 g_{ii}g_{jj}+\frac32g_{ij}^2
                   -\frac52(g_{ii}+g_{jj})g_{ij}\right\}\\
 ={}&-\frac{20555763}{572594890}.
 \end{aligned}
 \tag{V38.7}\label{r38:coherent-a1}
\]
\end{proposition}
\begin{proof}
For each unit-round \(S^3\), normal-coordinate Haar
density has the factor \((\sin|y_i|/|y_i|)^2\).
Thus its quadratic relative term is
\(-\frac13\sum_i|y_i|^2\); it has no linear term.
Rescale \(y=\gamma^{-1/2}X\) in the Gaussian integral.
The terms of order \(\gamma^{-1}\) in the exponential
are \(-T_4[X^4]/24\) and \(T_3[X^3]^2/72\).
The cubic term at order \(\gamma^{-1/2}\) is odd.
There is no linear-density/cubic cross term.
This gives \eqref{r38:a1formula}. A fixed small cutoff,
Taylor remainders, and Gaussian domination identify this
coefficient with the uniquely determined coefficient
in \eqref{r38:relative-integral}. Orthogonal changes
of tangent frames do not change any contraction.

At coherence all cubic terms vanish. If \(r_i=|y_i|\),
the homogeneous quartic part of the phase is
\[
 E_4(y)=-\frac14\sum_i r_i^4
       -\frac14\sum_{\{i,j\}}r_i^2r_j^2
       +\frac16\sum_{\{i,j\}}(r_i^2+r_j^2)y_i\cdot y_j .
 \tag{V38.8}\label{r38:quartic}
\]
Indeed each private term is \(1-\cos r_i\), and each
edge dot product is
\(\cos r_i\cos r_j+
 (\sin r_i/r_i)(\sin r_j/r_j)y_i\cdot y_j\);
the private multiplicity plus degree at every vertex is six.
For the covariance \(g\otimes I_3\), Wick contraction gives
\[
 \begin{gathered}
 \mathbb E|X_i|^2=3g_{ii},\qquad
 \mathbb E|X_i|^4=15g_{ii}^2,\\
 \mathbb E(|X_i|^2|X_j|^2)=9g_{ii}g_{jj}+6g_{ij}^2,\qquad
 \mathbb E(|X_i|^2X_i\cdot X_j)=15g_{ii}g_{ij}.
 \end{gathered}
\]
Substitution in the Haar term minus \(E_4\) gives the
displayed sum. Exact inversion of the displayed six-by-six
integer matrix, with determinant \(37835\), gives its
stated rational value. The accompanying checker evaluates
both the Wick contractions and the reduced sum, and checks
the same scalar in the independent chart
\(z_i=(\sqrt{1-|x_i|^2},x_i)\).
\end{proof}

The coefficient is not generically constant in the exterior
coefficient family. For example replace two private slots
at one vertex by
\((\cos\alpha,\pm\sin\alpha,0,0)\), leaving all other slots
and transports coherent. With
\(\delta=2(1-\cos\alpha)\), the unique minimum is still
\(z_i=e_0\), its energy is \(m=\delta\), and
\(M=(\mathsf A-\delta e_ie_i^{\mathsf T})\otimes I_3\).
Here \(d^2=2\delta\). Formula \eqref{r38:coherent-a1}
continues to hold with diagonal entries \(D_i=6-\delta\)
at the selected vertex, \(D_j=6\) elsewhere, and with
\(\frac{15}{4}\sum g_{jj}^2\) replaced by
\(\frac58\sum D_jg_{jj}^2\).
Four exact rational fixtures in the checker have
\(a_1(\vartheta)\ne a_1(o)\).
These are admissible coefficient examples, not a claimed
realization by particular native retained link histories.
They preclude deleting the correction from a theorem
uniform on the entire coefficient set.

\subsection{Whole-block thermal means and contact cumulants}

Let
\[
 H_\gamma(\vartheta)=\lambda_{\gamma,\vartheta}(E_\vartheta),
 \qquad h_\gamma=H_\gamma(o),\qquad
 V_\gamma(\vartheta)=\operatorname{Var}_{\lambda_{\gamma,\vartheta}}
                                                   (E_\vartheta).
\]
These symbols refer to conditional quantities, not to the
full-lattice mean or variance.

\begin{theorem}[A constant leading thermal term and a quadratic remainder]
\label{r38:thermal}
For \(\vartheta\in K\) and \(\gamma\ge1\),
\[
 \begin{aligned}
 H_\gamma(\vartheta)
    &=m(\vartheta)+\frac9\gamma
                         +\frac{a_1(\vartheta)}{\gamma^2}
                                      +O_{C^2_a}(\gamma^{-3}),\\
 V_\gamma(\vartheta)
    &=\frac9{\gamma^2}+\frac{2a_1(\vartheta)}{\gamma^3}
                                      +O_{C^2_a}(\gamma^{-4}),\\
 \left|H_\gamma(\vartheta)-h_\gamma-m(\vartheta)
                         -\frac{\Delta a_1}{\gamma^2}\right|
    &\le C d(\vartheta)^2\gamma^{-3},\\
 \left|V_\gamma(\vartheta)-V_\gamma(o)
                         -\frac{2\Delta a_1}{\gamma^3}\right|
    &\le C d(\vartheta)^2\gamma^{-4}.
 \end{aligned}
 \tag{V38.9}\label{r38:thermal-bounds}
\]
In particular
\(\left|H_\gamma(\vartheta)-h_\gamma-m(\vartheta)\right|
 \le Cd(\vartheta)^2/\gamma^2\).
\end{theorem}
\begin{proof}
At fixed coefficients, differentiation of the full compact
integral gives
\(H_\gamma=-\partial_\gamma\log Q_\gamma\) and
\(V_\gamma=\partial_\gamma^2\log Q_\gamma\), including the
normalizing denominator. In \eqref{r38:f}, the minimum and
Hessian determinant do not depend on \(\gamma\).
The two derivatives of \(-9\log\gamma\) give \(9/\gamma\)
and \(9/\gamma^2\), respectively. The two derivatives of
\(a_1/\gamma\), with their indicated signs, give
\(a_1/\gamma^2\) and \(2a_1/\gamma^3\).
Theorem~\ref{r38:mixed}, including its coherence-subtracted
bounds for \(j=1,2\), proves every assertion.
\end{proof}

The determinant correction disappears from this particular
coupling derivative because the Hessian is that of
\(E_\vartheta\), not of \(\gamma E_\vartheta\).
The latter scaling has already been accounted for by
\(\gamma^{-9}\). There is no missing differentiation of a
field-dependent normalizer. Individual plaquette insertions
vary individual action weights, so their leading thermal
coefficient need not be constant. The theorem cannot be
assigned to them term by term.

\subsection{A finite-source identity before retained integration}

For \(\gamma\ge2\), \(|s|\le\gamma/2\), define
\[
 \begin{aligned}
 Z_{\gamma,\vartheta}(s)
  &=\frac{Q_{\gamma+s}(\vartheta)}{Q_\gamma(\vartheta)}
                       \frac{Q_\gamma(o)}{Q_{\gamma+s}(o)},\\
 \varepsilon_{\gamma,\vartheta}(s)
  &=f_{\gamma+s}(\vartheta)-f_{\gamma+s}(o)
                           -f_\gamma(\vartheta)+f_\gamma(o).
 \end{aligned}
 \tag{V38.10}\label{r38:source}
\]
The first definition makes sense at every exterior, including
rough ones. The second is used only on \(K\).

\begin{proposition}[Normalized whole-block source transport]
\label{r38:finite-source}
On \(K\),
\[
 \begin{aligned}
 \log Z_{\gamma,\vartheta}(s)&=-s\,m(\vartheta)
                                      +\varepsilon_{\gamma,\vartheta}(s),\\
 |\varepsilon_{\gamma,\vartheta}(s)|
       &\le C|s|d(\vartheta)^2/\gamma^2,\\
 \left|\varepsilon_{\gamma,\vartheta}(s)
       -\Delta a_1\left(\frac1{\gamma+s}-\frac1\gamma\right)\right|
       &\le C|s|d(\vartheta)^2/\gamma^3 .
 \end{aligned}
 \tag{V38.11}\label{r38:source-bound}
\]
For distinct blocks of the actual native packing, let
\(\mathcal E_b(U)\) be the sum of their incident plaquette
energies and let \(s_b\) be deterministic real sources.
Exactly,
\[
 \frac{\mathbb E_{\mu_W}
                     \exp\{-\sum_b s_b\mathcal E_b(U)\}}
      {\prod_b[Q_{\gamma+s_b}(o)/Q_\gamma(o)]}
   =\mathbb E_{\mu_{\mathcal R}}
                         \prod_b Z_{\gamma,\vartheta_b(\eta)}(s_b).
 \tag{V38.12}\label{r38:source-functional}
\]
No good-event conditioning occurs in this identity.
\end{proposition}
\begin{proof}
Insert \eqref{r38:f} in the four logarithms defining \(Z\).
The determinant terms cancel between the two couplings,
and the universal \(9\log\gamma\) terms cancel between
the exterior and coherence. Since \(m(o)=0\), the first
line of \eqref{r38:source-bound} follows exactly.
Integrating \(\partial_\gamma(f_\gamma(\vartheta)-f_\gamma(o))\)
over the coupling interval gives the second line;
using \eqref{r38:centered-remainder} gives the third.
Every coupling in that interval is at least \(\gamma/2\).

Given the original retained links, the free blocks are
independent and \(\mathcal E_b=E_{\vartheta_b}\) exactly
in their Haar-preserving frames. Its conditional Laplace
transform is \(Q_{\gamma+s_b}(\vartheta_b)/Q_\gamma(\vartheta_b)\).
Multiply these transforms, integrate the unchanged original
marginal and divide by the displayed coherent constants.
\end{proof}

The right-hand side remains a genuinely interacting retained
expectation. On rough exteriors its \(Z\) is kept exact.
The bounds in \eqref{r38:source-bound} are pointwise local
bounds, not permission to apply V35's original-law moments
under a source-tilted law. Taking a logarithm of
\eqref{r38:source-functional} does not remove the retained
normalizer or manufacture connected spatial decay.

The first and second derivatives of \(\log Z\) at zero
are \(-H_\gamma(\vartheta)+h_\gamma\) and
\(V_\gamma(\vartheta)-V_\gamma(o)\).
Thus the previous theorem also evaluates the residual
same-block contact term. Different-block connected
derivatives still involve covariances of the retained
predictors; the contact subtraction does not delete them.

\subsection{Stronger original-law reflected block-energy sources}

Use V37's reflection-matched native packing, and select
distinct blocks in its positive open half. Let
\[
 \mathcal E_c=\sum_b c_b\mathcal E_b,\qquad
 G_b=\{\mathcal D_b\le1/1024\},\qquad
 H_b=\mathbb E_{\mu_W}[\mathcal E_b\mid\eta].
\]
The notation \(H_b\) here is a conditional mean, not the
conditional partition factor of earlier notes.
Define retained sources
\[
 \begin{aligned}
 K_b^{(0)}&=
 \begin{cases}
   m(\vartheta_b),&G_b,\\
   H_b-h_\gamma,&G_b^c ,
 \end{cases}\\
 K_b^{(1)}&=
 \begin{cases}
   m(\vartheta_b)+
        [a_1(\vartheta_b)-a_1(o)]/\gamma^2,&G_b,\\
   H_b-h_\gamma,&G_b^c .
 \end{cases}
 \end{aligned}
 \tag{V38.13}\label{r38:K}
\]
These are bounded, measurable, gauge-invariant local
retained observables. No minimum is chosen on \(G_b^c\);
there the original conditional predictor remains exact.
Let \(K_c^{(j)}=\sum_b c_bK_b^{(j)}\), and center all
observables with the unchanged original Wilson marginal.
In fact \(K_b^{(0)}\) is the sum of the 29 instances of
V37's source \eqref{r37:J} in this block, using their
common good event: the minimum plaquette energies sum
to \(m\), and their coherent offsets sum to \(h_\gamma\).
Thus the improved estimate exploits a cancellation in
the same whole-energy comparator, not a changed measure.

\begin{theorem}[Thermally centered native energy comparison]
\label{r38:OS}
For arbitrary deterministic complex coefficients \(c_b\),
\[
 \begin{aligned}
 \|\mathcal E_c^\circ-(K_c^{(0)})^\circ\|_{\rm OS}
    &\le C\frac{\sqrt{m_{2,\gamma}}}{\gamma^2}\|c\|_1
     \le C\frac{1+\log\gamma}{\gamma^3}\|c\|_1,\\
 \|\mathcal E_c^\circ-(K_c^{(1)})^\circ\|_{\rm OS}
    &\le C\frac{\sqrt{m_{2,\gamma}}}{\gamma^3}\|c\|_1
     \le C\frac{1+\log\gamma}{\gamma^4}\|c\|_1 .
 \end{aligned}
 \tag{V38.14}\label{r38:OS-bound}
\]
Here \(m_{2,\gamma}\) is the original-law defect second-moment
budget in \eqref{r37:budgets}. The comparison is under
\(\mu_W\), not a modified minimum-action law.
\end{theorem}
\begin{proof}
Set \(r_b^{(j)}=H_b-h_\gamma-K_b^{(j)}\).
On rough exteriors it vanishes exactly. On good exteriors
\(d(\vartheta_b)^2\le\mathcal D_b\), so
Theorem~\ref{r38:thermal} gives
\[
 |r_b^{(0)}|\le C\mathcal D_b/\gamma^2,\qquad
 |r_b^{(1)}|\le C\mathcal D_b/\gamma^3 .
\]
Consequently their ordinary \(L^2(\mu_{\mathcal R})\)
norms are bounded by \(C\sqrt{m_{2,\gamma}}/\gamma^2\)
and \(C\sqrt{m_{2,\gamma}}/\gamma^3\), respectively.

All 29 plaquettes and the retained predictor of a selected
block stay in the positive half. V37's exact conditional
projection, equation \eqref{r37:projection}, therefore
applies to \(\mathcal E_c-K_c^{(j)}\); it does not require
independence of plaquettes inside one block.
Its projected observable is
\(h_\gamma\sum_b c_b+\sum_b c_br_b^{(j)}\).
Global centering removes the first, deterministic term.
Ordinary Cauchy--Schwarz and reflection invariance bound
the remaining centered reflected norm by
\(\|\sum_b c_br_b^{(j)}\|_{L^2(\mu_{\mathcal R})}\).
The triangle inequality and the V35 moment bound prove
\eqref{r38:OS-bound}.
\end{proof}

Keeping \(H_b-h_\gamma\) on rough exteriors is essential:
subtracting \(h_\gamma\) only on \(G_b\) would leave a
random indicator. Neither theorem changes the measure,
drops rough-sector correlations, or replaces the
correlated \(\ell^1\) cost by an \(\ell^2\) cost.
For an \(\ell^2\)-normalized bank of \(M\) marked blocks,
the two sufficient absolute errors are respectively
\(C\sqrt M(1+\log\gamma)/\gamma^3\) and
\(C\sqrt M(1+\log\gamma)/\gamma^4\).
Any additional physical normalization \(Z(a)\) multiplies
these errors by \(|Z(a)|\).

\subsection{What this improves, and what remains genuinely missing}

The result replaces V37's generic plaquette bias estimate
by a stronger, proven comparison for a specified
whole-block source. The new cancellation is not that all
thermal fluctuations vanish: their leading conditional
variance is \(9/\gamma^2\). They are null only after the
verified reflected conditional projection; their conditional
mean has a constant leading term that can be centered away.

The explicit \(a_1\) correction is finite and local.
It is not a beta-function coefficient of the full
four-dimensional theory and not an interacting continuum
counterterm theorem. The native retained action, its
nonlinear collective modes, and the physical-time return
have not become contractive. In particular V33's soft
coherent Schur scale is unchanged.

There is also a simple scale warning. For any fixed
\(p,r>0\), \(a\downarrow0\), and \(\gamma(a)\to\infty\) with
\(\gamma(a)\le C\log(1/a)\),
\[
 a^{-p}\frac{1+\log\gamma(a)}{\gamma(a)^r}\longrightarrow\infty.
 \tag{V38.15}\label{r38:scale-warning}
\]
Indeed this is bounded below by
\(a^{-p}/[C\log(1/a)]^r\), which diverges.
This is a conditional arithmetic statement, not a proof
that the actual source error diverges or that such a
trajectory has been constructed. It explains why even a
higher inverse-coupling local bound does not automatically
survive a normalization containing a negative power of
lattice spacing. Further finite-order local improvements
alone are not a substitute for positive-time control.

The next calculation should therefore use these concrete
sources to estimate the \emph{retained} reflected
correlations, their low-energy component after physical
time damping, or an actual scale-generating return.
It must not iterate a local asymptotic series merely
because each additional coefficient is computable.
A noncollapsing physical source bank, compatible physical-time
contraction, the interacting four-dimensional continuum
construction and a positive Hamiltonian mass gap remain open.

\paragraph{Verification and ledger.}
The exact checker
\begin{center}\small
\path{verify_ym_restart_v38_energy_source.py}
\end{center}
checks the \(18\)-dimensional radial Gaussian moments,
normalizer-sensitive energy-ratio signs in a nonlinear
one-dimensional fixture, the finite-source shift algebra,
and the native coherent contraction \eqref{r38:coherent-a1}.
Four noncoherent paired-slot fixtures verify the same
coefficient in two coordinate charts and via direct Wick
contractions. These checks do not evaluate the uniform
analytic constants, prove the remainders by computation,
or certify a physical spectrum.
The uniform analytic statements are proved above.

New: mixed coupling derivatives of V32's remainder;
the field-independent \(9/\gamma\) energy coefficient;
coherence-subtracted mean and contact-variance bounds;
an evaluated native first correction; an exact sourced
retained functional with rough factors unchanged; and
the two original-law reflected block-energy comparisons.
Inherited and credited: finite-dimensional equipartition,
V31's pinned minimum, V32's symmetry floor, V35's native
moments and V37's reflected projection.
Removing this append and restoring the V37 title recovers
the entire predecessor exactly. All other source files
are preserved; build and verification outputs stay outside
\path{latex_notes}.

% END CONSOLIDATED V38 APPEND

% BEGIN CONSOLIDATED V39 APPEND
\clearpage
\section{V39: Native slab insertions on the unchanged physical transfer}
\label{r39:context}

V38 produced a sharper source comparison for the whole native
29-plaquette energy, but its displayed norm used a fixed reflected
form. The next issue is compatibility with physical time: can that
specific comparison be realized on the original Wilson Hilbert space,
with the same transfer operator at every later time? This append
answers that question at a finite regulator. It does not estimate the
remaining retained correlations or prove a continuum mass gap.

The construction keeps both complete boundary slices of a two-step
slab. Only the six spatial links of each selected native block in its
middle slice are integrated out. Consequently neither endpoint Hilbert
space nor the physical transfer is replaced by a process on retained
variables. A time-reflection-symmetric retained source is needed.
A one-sided cutoff is not enough; below a union of forward and backward
controlled frames supplies such a source, without discarding rough
configurations.

\paragraph{Inherited results and the literature check.}
We use V14's original gauge-projected Wilson transfer
\eqref{r14:T}, V31's native conditional energy and controlled minima,
V35's original-law defect moments \eqref{r35:moments}, and V38's
coherence-subtracted energy expansion. Transfer insertions themselves
are not new: current V14 and archive \(C\), V8 already use them.
Archive \(C\), V60--V63 already distinguishes exact retained
conditionals from a physical Markov quotient; its V57 retains
nonreducing returns. The elementary source-Gram transport inequality
used below is also inherited. The new input is a reflection-even,
native, two-step realization of the V38 sources with quantitative
errors on this \emph{same} physical Hilbert space.

The targeted primary-source check was M.~L\"uscher,
\href{https://bib-pubdb1.desy.de/record/396349/files/7611148.pdf}
{\emph{Construction of a Selfadjoint, Strictly Positive Transfer Matrix
for Euclidean Lattice Gauge Theories}}, DESY 76/54,
Commun.\ Math.\ Phys.\ \textbf{54} (1977), 283--292,
Sections II.B--II.C. Its pure-gauge transfer positivity and
reconstruction provide the classical setting, not a continuum gap
estimate. The relevant text was read; this scan was not
visually verified locally. We do not import a fermionic vacuum
uniqueness assertion. The scheduled broad literature and companion
reviews both remain due at V40. The ESM comparison continues to guide
the separation of exact identities, uniform error estimates and
unproved nonperturbative limits; no ESM scaling exponent is transferred.

\subsection{Full endpoints and exact two-step insertion kernels}

Fix a spatial torus of dyadic side \(L\ge16\) and a coupling
\(\gamma\ge\gamma_*\), where \(\gamma_*\ge2\) is large enough for V38.
All Haar measures are probability measures. Set
\[
 \mathcal X=\SU(2)^{3L^3},\qquad
 \mathcal H=L^2(\mathcal X,dx)_{\mathrm{gauge}} .
\]
For full spatial slices \(x,y\) and all temporal links \(v\)
between them, use the original weight
\[
 \begin{split}
 w_\gamma(x,y;v)
  &=\exp\{-\gamma[S_s(x)/2+S_t(x,y;v)+S_s(y)/2]\},\\
 T(x,y)&=\int w_\gamma(x,y;v)\,dv,\qquad
 T\Omega=\Lambda\Omega,\qquad \|\Omega\|_2=1,\\
 P&=T/\Lambda,\qquad \Pi_\Omega=|\Omega\rangle\langle\Omega|.
 \end{split}
 \tag{V39.1}\label{r39:transfer}
\]
Here \(P\) denotes the normalized \emph{physical} transfer, called
\(K\) in V14; it is not V14's ground-state Markov operator.
The inherited finite-regulator facts are
\[
 0<P\le I,\qquad P\Omega=\Omega,\qquad \Omega>0 .
\]
The operator is compact, trace class and smoothing, and its top
eigenvalue is simple. It is injective on \(\mathcal H\).
On unreduced \(L^2(dx)\) the gauge projection has a nullspace;
we never define a physical logarithm on that nullspace.
For \(a>0\), \(H_a=-a^{-1}\log P\) is defined by spectral calculus
on \(\mathcal H\). No lower bound on its first excited eigenvalue,
uniform in \(L,\gamma,a\), is assumed.

A slab configuration consists of \(x,y,z,v_0,v_1\), at times
\(0,1,2\). For a bounded gauge-invariant slab observable \(F\), define
\[
 B_F(x,z)=\Lambda^{-2}\int
 w_\gamma(x,y;v_0)w_\gamma(y,z;v_1)
 F(x,y,z,v_0,v_1)\,dy\,dv_0\,dv_1 .
 \tag{V39.2}\label{r39:insertion}
\]
Both endpoints \(x,z\) are full slices. In particular
\[
 B_1=P^2,\qquad
 |B_F(x,z)|\le\|F\|_\infty P^2(x,z),\qquad
 \|B_F\|\le\|F\|_\infty .
 \tag{V39.3}\label{r39:basic-domination}
\]
The last assertion follows by taking absolute values in the integral
kernel, then using \(\|P^2\|\le1\). Absolute values of physical
wavefunctions remain gauge invariant.

Let \(\iota\) reverse the slab, exchanging \(x,z\), leaving the middle
spatial links in place, and replacing
\((v_0,v_1)\) by \((v_1^{-1},v_0^{-1})\).
The action and Haar measure are invariant under this change.
Taking the adjoint kernel gives
\[
 B_F^*=B_{\overline{F\circ\iota}} .
 \tag{V39.4}\label{r39:adjoint}
\]
Thus real, reversal-even \(F\) yields a selfadjoint insertion.
A nonnegative observable gives a nonnegative \emph{kernel};
it does not in general give a positive operator in the Loewner order.

\subsection{The native free block survives the physical slab cut}

Choose spatial directions \(\mu,\nu,\rho\). At time \(1\), put
six free \(\mu\)-links at a \(2\)-by-\(3\) grid in the
\((\nu,\rho)\) plane. Use V35's spatial anchors:
both transverse coordinates are multiples of \(4\), while the
\(\mu\) coordinate is arbitrary. Select any finite subset of
these blocks. All other links, including all temporal links and
both complete endpoint slices, belong to the retained variable
\(\eta\). Denote the sum of the 29 plaquette energies incident
to the six free links of block \(b\) by \(\mathcal E_b\).

\begin{lemma}[No half-weight incident plaquette]
\label{r39:native-law}
Every plaquette incident to a selected free block has coefficient
one in the two-step slab action. There are 22 private plaquettes
and seven with two free links, and no plaquette contains free
links from different selected blocks. Given \(\eta\), the free
block law is precisely V31's native law
\[
 \lambda_{\gamma,\vartheta_b}(dU_b)
 =Q_\gamma(\vartheta_b)^{-1}
             e^{-\gamma E_{\vartheta_b}(U_b)}\,dH(U_b).
 \tag{V39.5}\label{r39:conditional}
\]
In particular \(E_{\vartheta_b}\) is the original
\(\mathcal E_b\), with no omitted additive energy.
\end{lemma}
\begin{proof}
The product of the two weights in \eqref{r39:insertion} has spatial
action \(S_s(x)/2+S_s(y)+S_s(z)/2\) and the full temporal
action in each adjacent slab. A spatial plaquette touching a free
link lies at time \(1\), hence has weight one. A temporal plaquette
touching one lies between times \(0,1\) or \(1,2\), also with weight
one. None touches a half-weight endpoint spatial plaquette.

Each of the six free links has six incident plaquettes, giving
36 incidences. The rectangular grid has
\(3+4=7\) internal adjacencies, each counted twice; hence there
are 29 distinct plaquettes, of which \(36-2\cdot7=22\) are private.
The transverse spacing four leaves at least one intervening
retained link between distinct grids. Since an incident plaquette
can connect only neighboring parallel \(\mu\)-links, it cannot
contain free links from two blocks. The actual Haar-preserving
native change of variables is then the one in V31.
All remaining slab factors are independent of these free links
and cancel in conditional normalization. Neither an endpoint
wavefunction nor a Perron factor depends on a free middle link.
\end{proof}

Define the bounded retained predictor
\[
 H_b(\eta)=\E[\mathcal E_b\mid\eta].
\]
It is real, gauge invariant and reversal even. The letter \(H_b\)
denotes this conditional mean, not the physical Hamiltonian \(H_a\).
Since \(0\le\mathcal E_b\le58\), also \(0\le H_b\le58\).

\begin{proposition}[An identity of full-endpoint insertion kernels]
\label{r39:predictor}
For each selected block,
\[
 B_{\mathcal E_b}=B_{H_b},\qquad
 B_{\mathcal E_b^2}-B_{H_b^2}
       =B_{\Var(\mathcal E_b\mid\eta)} .
 \tag{V39.6}\label{r39:projection}
\]
The first identity holds pointwise in both endpoints and therefore
between arbitrary endpoint wavefunctions, and after sewing to
any exterior that does not use the eliminated middle links.
\end{proposition}
\begin{proof}
At fixed \(x,z\), first integrate the six free links in
\eqref{r39:insertion}. Lemma~\ref{r39:native-law} identifies their
normalized conditional distribution. Replacing
\(\mathcal E_b\) by its conditional mean leaves that integral
unchanged. Replacing its square uses
\(\E[\mathcal E_b^2\mid\eta]=H_b^2+\Var(\mathcal E_b\mid\eta)\).
Fubini applies because all spaces are compact and the integrands
are bounded. The same argument allows additional factors
measurable with respect to the retained exterior.
\end{proof}

In particular \(B_{\mathcal E_b-H_b}=0\), whereas its ordinary
same-slab square need not vanish. Equation~\eqref{r39:projection}
is not a multiplicative map of observables. The contact-variance
kernel is nonnegative, but no operator-order assertion about it
is needed or made. Nor does the identity construct an intertwiner
between the physical transfer and an invented retained transfer.

\subsection{A reflection-even source from intrinsic frame overlap}

The forward frame of V35 has a retained defect
\(\mathcal D_b^+\); reflect its entire construction to obtain
\(\mathcal D_b^-\). Both are measurable functions of the retained
links in this same slab. The 45-face witness for the forward frame
lies at times \(0,1,2\), as does its reflected witness.
The reflected free set is the same middle grid. Write
\[
 G_b^\pm=\{\mathcal D_b^\pm\le1/1024\},\qquad
 \widehat G_b=G_b^+\cup G_b^- .
 \tag{V39.7}\label{r39:union}
\]
The two events are exchanged by \(\iota\).
The conditional energy \(E\), its law, and its predictor \(H_b\)
do not depend on which frame is used to describe them.

\begin{lemma}[Intrinsic agreement of the good-frame coefficients]
\label{r39:intrinsic}
If two native frames of the same block are both in V31's controlled
coefficient set, they have the same minimum value \(m_b\), the same
full Hessian determinant at the minimum, and the same first
Laplace coefficient \(a_{1,b}\) of V38.
\end{lemma}
\begin{proof}
For a fixed retained configuration, each frame is a product of
left and right multiplications on the six copies of \(\SU(2)\),
with an inversion if an orientation is reversed. These are
isometries of the round three-sphere and preserve Haar measure.
The resulting energy is exactly the original 29-plaquette sum
in either coordinate system. Hence the change between frames
preserves the minimum value and carries the full Hessian at
the unique controlled minimum by an orthogonal conjugation.

The partition integrals \(Q_\gamma\) coincide for every \(\gamma\).
The minimum and the Hessian determinant just shown to coincide
make their leading Laplace factors identical. The expansions
of V38 in the two controlled frames therefore read
\(f_\gamma=a_1/\gamma+O(\gamma^{-2})\) for the very same
function \(f_\gamma\). Multiply the difference by \(\gamma\)
and let \(\gamma\to\infty\); their \(a_1\)'s coincide.
This argument also accounts for the Haar factor, rather than
identifying only the phase contribution.
\end{proof}

Let \(h_\gamma=H_\gamma(o)\) be V38's exact coherent conditional
energy mean. On \(\widehat G_b\), use either available good frame
to define \(m_b,a_{1,b}\); they agree on the overlap.
For \(j=0,1\), define
\[
 \widehat K_b^{(j)}=
 \begin{cases}
 m_b+j[a_{1,b}-a_1(o)]/\gamma^2,&\widehat G_b,\\
 H_b-h_\gamma,&\widehat G_b^c,
 \end{cases}
 \qquad
 r_b^{(j)}=H_b-h_\gamma-\widehat K_b^{(j)} .
 \tag{V39.8}\label{r39:writer}
\]
One may choose the forward frame first on the overlap to make
measurability explicit; the resulting scalar is independent of
that choice. No minimum or asymptotic coefficient is evaluated
when both frames are rough.

\begin{proposition}[Reflection-even native writer with exact rough part]
\label{r39:writer-bound}
The observables \(\widehat K_b^{(j)}\) and \(r_b^{(j)}\) are bounded,
measurable, gauge invariant and reversal even. Uniformly over
the native exteriors and spatial volumes,
\[
 \begin{split}
 |r_b^{(j)}|
 &\le C\gamma^{-2-j}
       \min(\mathcal D_b^+,\mathcal D_b^-)
                      {\bf1}_{\widehat G_b}\\
 &\le \frac{C}{1024\,\gamma^{2+j}} .
 \end{split}
 \tag{V39.9}\label{r39:pointwise}
\]
\end{proposition}
\begin{proof}
In either good frame V38 gives, with its coefficient defect \(d\),
\[
 \begin{split}
 |H_b-h_\gamma-m_b|&\le C d^2\gamma^{-2},\\
 |H_b-h_\gamma-m_b-[a_{1,b}-a_1(o)]/\gamma^2|
                       &\le C d^2\gamma^{-3}.
 \end{split}
\]
The corresponding V35 defect bounds \(d^2\le\mathcal D_b^\pm\).
On the union, the smaller of the two defects is a good one,
so apply that frame's estimate. Lemma~\ref{r39:intrinsic} makes
the estimated scalar the one in \eqref{r39:writer}, including
when both frames are good. Off the union the residual is exactly
zero. This proves \eqref{r39:pointwise}.

Gauge changes are native Haar-preserving frame changes and preserve
the defect events. Reversal exchanges the two events and the two
descriptions of the identical conditional energy. Their intrinsic
agreement proves reversal symmetry on the union; off it the exact
predictor is already reversal even. Boundedness follows from the
compact controlled coefficient set, \(0\le H_b\le58\), and bounded
\(h_\gamma\).
\end{proof}

This construction uses a \emph{union}, not an intersection or the
arithmetic average of two one-sided writers. It exploits agreement
where both descriptions are controlled, and keeps the exact rough
predictor where neither is. It requires no joint-tail estimate for
the forward and backward defects and no independence between them.

\subsection{Uniform operator bounds without a boundary-overlap penalty}

For deterministic complex coefficients \(c_b\), put
\(\mathcal E_c=\sum_b c_b\mathcal E_b\),
\(\widehat K_c^{(j)}=\sum_b c_b\widehat K_b^{(j)}\), and
\(r_c^{(j)}=\sum_b c_br_b^{(j)}\). Define
\[
 \epsilon_{j,\gamma}(c)
       =\frac{C\|c\|_1}{1024\,\gamma^{2+j}} .
\]
All operator norms in this append act on the original \(\mathcal H\).

\begin{theorem}[Native full-endpoint insertion comparison]
\label{r39:operator}
For every such finite bank,
\[
 \begin{split}
 D_c^{(j)}
 &:=B_{\mathcal E_c}
     -h_\gamma\Big(\sum_b c_b\Big)P^2
     -B_{\widehat K_c^{(j)}}=B_{r_c^{(j)}},\\
 |D_c^{(j)}(x,z)|
 &\le\epsilon_{j,\gamma}(c)P^2(x,z),\qquad
 \|D_c^{(j)}\|\le\epsilon_{j,\gamma}(c).
 \end{split}
 \tag{V39.10}\label{r39:operator-bound}
\]
For a slab observable define its stationary mean and centered
insertion by
\[
 \overline F=\langle\Omega,B_F\Omega\rangle,\qquad
 B_F^\circ=B_F-\overline F\,P^2 .
 \tag{V39.11}\label{r39:center}
\]
Then
\[
 \|B_{\mathcal E_c}^\circ
            -B_{\widehat K_c^{(j)}}^\circ\|
            \le2\epsilon_{j,\gamma}(c).
 \tag{V39.12}\label{r39:centered-operator}
\]
\end{theorem}
\begin{proof}
Proposition~\ref{r39:predictor} and \eqref{r39:writer} give the
exact first identity. The triangle inequality in
\eqref{r39:pointwise} bounds \(|r_c^{(j)}|\) by
\(\epsilon_{j,\gamma}(c)\), so the positive two-step kernel
gives the pointwise estimate. For arbitrary physical \(f,g\),
\[
 |\langle f,D_c^{(j)}g\rangle|
 \le\epsilon_{j,\gamma}(c)\langle|f|,P^2|g|\rangle
 \le\epsilon_{j,\gamma}(c)\|f\|_2\|g\|_2 .
\]
This proves the norm estimate, including complex \(c\).
Taking its vacuum expectation bounds
\(|\overline{\mathcal E_c}
-h_\gamma\sum c_b-\overline{\widehat K_c^{(j)}}|\)
by the same \(\epsilon_{j,\gamma}(c)\).
Subtracting that scalar times \(P^2\), whose norm is one,
proves \eqref{r39:centered-operator}.
\end{proof}

The centering in \eqref{r39:center} subtracts
\(\overline F P^2\), not \(\overline F I\): the observable occupies
a two-step slab. The Perron normalization cancels through domination
by that same \(P^2\). No separately estimated lower bound on
\(\Lambda\), an endpoint density, or a vacuum overlap enters
\eqref{r39:operator-bound}.

\subsection{Original-law moments in the stationary vacuum slab}

The stronger V38 rate uses defect moments, not the worst-case
pointwise cutoff. We must therefore identify the correct law.
Define the augmented stationary slab measure by
\[
 d\nu_\gamma=
 \Lambda^{-2}\Omega(x)\Omega(z)
 w_\gamma(x,y;v_0)w_\gamma(y,z;v_1)
 \,dx\,dy\,dz\,dv_0\,dv_1 .
 \tag{V39.13}\label{r39:vacuum-law}
\]
This is a probability measure with endpoint marginals
\(\Omega^2dx\). It includes all temporal links, and its conditional
free middle-link law is still \eqref{r39:conditional}.
It is not a new retained Gibbs measure.

\begin{lemma}[Passing the inherited moments to the vacuum law]
\label{r39:moment-limit}
At fixed \(L,\gamma\), the local two-step slab marginal of the
original periodic Wilson measure with time extent \(N\)
converges in total variation to \(\nu_\gamma\), along dyadic
\(N\to\infty\). Consequently
\[
 \E_{\nu_\gamma}(\mathcal D_b^\pm)^2\le m_{2,\gamma},
 \quad
 m_{2,\gamma}=A_\gamma^2+2A_\gamma/\ell_\gamma
                                      +2/\ell_\gamma^2,
 \tag{V39.14}\label{r39:vacuum-moments}
\]
where \(A_\gamma,\ell_\gamma\) are exactly
\eqref{r37:budgets}. In particular
\(m_{2,\gamma}=O((1+\log\gamma)^2/\gamma^2)\),
with constants independent of \(L\).
\end{lemma}
\begin{proof}
After integrating the complementary cylinder, the periodic
augmented slab density relative to product Haar is
\[
 \Lambda^{-2}w_\gamma(x,y;v_0)w_\gamma(y,z;v_1)
       \frac{P^{N-2}(z,x)}{\Tr_{\mathcal H}P^N}.
 \tag{V39.15}\label{r39:periodic-density}
\]
At this \emph{fixed} regulator, compactness and the simple positive
vacuum give
\(\rho_{L,\gamma}:=\|P-\Pi_\Omega\|<1\).
This assertion has no uniformity in \(L,\gamma\).
Since \(P\) is positive trace class,
\(\Tr P^N\to1\), by domination of its eigenvalue powers
by the summable eigenvalues of \(P\).

For completeness, kernel convergence is uniform. Let
\(C_{L,\gamma}=\sup_x\|P(x,\cdot)\|_{L^2(dx)}<\infty\),
using the continuous smoothing kernel on compact \(\mathcal X\).
For \(N\ge5\),
\[
 \begin{split}
 P^{N-2}-\Pi_\Omega
     &=P(P^{N-4}-\Pi_\Omega)P,\\
 \sup_{x,z}|P^{N-2}(z,x)-\Omega(z)\Omega(x)|
     &\le C_{L,\gamma}^2\rho_{L,\gamma}^{\,N-4}.
 \end{split}
\]
One can use the gauge-averaged kernels on unreduced \(L^2\)
in this calculation; their extra nullspace contributes zero.
Together with \eqref{r39:periodic-density}, uniform convergence
proves convergence of the augmented densities in \(L^1\),
and hence total variation. All local defect functions are bounded
on the finite product group, including their measurable cutoff
indicators.

V35's one-point moment estimate is for the original periodic law,
uniformly in the allowed spatial and temporal extents. Translate
its native packing so the selected middle block is at time \(1\).
The forward witness fits in this slab, and reflection gives the
same estimate for the backward witness. The unchanged periodic
moments pass to the total-variation limit, proving
\eqref{r39:vacuum-moments}. Only after this fixed-regulator limit
are the constants observed to be independent of \(L\).
No uniform rate in \(N\) has been inferred.
\end{proof}

\begin{theorem}[Native comparison of physical vacuum source vectors]
\label{r39:vacuum-vector}
Let \(\psi_F=B_F^\circ\Omega\). For every finite coefficient bank,
\[
 \begin{split}
 \|\psi_{\mathcal E_c}
             -\psi_{\widehat K_c^{(j)}}\|
 &\le C\sqrt{m_{2,\gamma}}\,
                       \gamma^{-2-j}\|c\|_1\\
 &=O\bigl((1+\log\gamma)
                       \gamma^{-3-j}\|c\|_1\bigr),
 \qquad j=0,1 .
 \end{split}
 \tag{V39.16}\label{r39:vector-bound}
\]
The vectors lie in the same physical space and are orthogonal
to its vacuum.
\end{theorem}
\begin{proof}
For every bounded slab observable \(R\), disintegration of
\eqref{r39:vacuum-law} at its left endpoint gives
\[
 \frac{(B_R\Omega)(x)}{\Omega(x)}
                   =\E_{\nu_\gamma}[R\mid x].
 \tag{V39.17}\label{r39:Jensen}
\]
Indeed the conditional normalizing factor is
\(\Lambda^2\Omega(x)\), since
\(T^2\Omega=\Lambda^2\Omega\).
Furthermore \(\overline R=\E_{\nu_\gamma}R\).
Thus conditional Jensen, also for complex \(R\), implies
\[
 \|B_R\Omega-\overline R\Omega\|_2^2
       \le \E_{\nu_\gamma}|R-\overline R|^2
       \le \E_{\nu_\gamma}|R|^2 .
 \tag{V39.18}\label{r39:vector-Jensen}
\]
Apply this to \(R=r_c^{(j)}\), using the exact centered
identity of Theorem~\ref{r39:operator}. From
\eqref{r39:pointwise} and \eqref{r39:vacuum-moments},
\[
 \|r_b^{(j)}\|_{L^2(\nu_\gamma)}
       \le C\gamma^{-2-j}\sqrt{m_{2,\gamma}} .
\]
The \(L^2\) triangle inequality supplies \(\|c\|_1\),
without decorrelating retained blocks. Finally
\(\langle\Omega,\psi_F\rangle=0\) follows from
\eqref{r39:center} and \(P^2\Omega=\Omega\).
\end{proof}

The rough predictor remains in each retained source.
Neither the operator bound nor the vector bound assumes independent
defects, a small probability of an all-good configuration, or
a lower bound for a source Gram matrix.

\subsection{Every later physical time and spectral window}

The construction of \(\nu_\gamma\) extends consistently to any
finite time interval by the same endpoint vacuum factors.
Equivalently, this is the stationary ground-state path law of
the \emph{original} Wilson transfer, with its temporal-link
variables retained in each slab.

\begin{proposition}[Two nonoverlapping slab windows]
\label{r39:time-correlation}
Let \(F,G\) be bounded real gauge-invariant reversal-even
two-step observables. Put the first in times \([0,2]\) and
the second in \([n+2,n+4]\), \(n\ge0\). Then
\[
 \E\bigl[(F_{[0,2]}-\overline F)
              (G_{[n+2,n+4]}-\overline G)\bigr]
       =\langle\psi_F,P^n\psi_G\rangle .
 \tag{V39.19}\label{r39:time}
\]
There is no even-\(n\) restriction.
\end{proposition}
\begin{proof}
Integrate the first two-step window, then the \(n\) unmarked
steps, then the second window. Sewing the full boundary slices
gives
\(\langle\Omega,B_F^\circ P^nB_G^\circ\Omega\rangle\).
Equation~\eqref{r39:adjoint} makes the left insertion selfadjoint,
so this equals the displayed inner product. For \(n=0\)
the windows share a boundary slice, not an eliminated middle
link or a temporal plaquette, and the same composition holds.
\end{proof}

The slab centers are separated by \(n+2\) lattice steps,
whereas \(P^n\) counts the unmarked steps between their boundaries.
That two-step support must not be suppressed in physical-time
bookkeeping. For complex linear combinations use the sesquilinear
extension of the real columns, with conjugation on the first bank.

Let \(M\) marked blocks define synthesis maps
\(V,W_j:\mathbb C^M\to\mathcal H\) whose columns are
\(\psi_{\mathcal E_b}\) and \(\psi_{\widehat K_b^{(j)}}\).
Then Theorem~\ref{r39:vacuum-vector} gives
\[
 \|V-W_j\|\le\delta_{j,\gamma,M},
 \qquad
 \delta_{j,\gamma,M}
       =C\sqrt M\,\sqrt{m_{2,\gamma}}\gamma^{-2-j}.
 \tag{V39.20}\label{r39:synthesis}
\]
For every positive contraction \(Q\) on this same space,
\[
 \|V^*QV-W_j^*QW_j\|
       \le 2\|V\|\delta_{j,\gamma,M}
                            +\delta_{j,\gamma,M}^2 .
 \tag{V39.21}\label{r39:Gram}
\]
Indeed set \(D=W_j-V\), expand
\(W_j^*QW_j-V^*QV=V^*QD+D^*QV+D^*QD\), and take norms.
This familiar perturbation inequality is stated to specify
the native input \(\delta_{j,\gamma,M}\), not as a new
abstract theorem.

In particular it applies simultaneously to
\[
 Q=P^n,\quad Q=e^{-tH_a},\quad
 Q={\bf1}_I(H_a),\quad
 Q={\bf1}_I(H_a)e^{-tH_a},
 \qquad n\ge0,\quad t\ge0,
 \tag{V39.22}\label{r39:spectral}
\]
for every Borel energy set \(I\subset[0,\infty)\).
The error is uniform in these choices. Also
\(\|Q(V-W_j)\|\le\delta_{j,\gamma,M}\).
For \(Q=P^n\), the matrices are exactly the connected
two-window correlations of \eqref{r39:time}.
All these comparisons use \(H_a=-a^{-1}\log P\), not the
spectrum of a retained conditional sampler.

\subsection{The remaining low-energy problem is explicit}

The new step is a time-compatible physical realization of the native
V38 energy comparison, not a decay estimate for
\(\langle\psi_{\widehat K_b},P^n\psi_{\widehat K_{b'}}\rangle\).
In particular it does not turn an absolute source error into a
relative spectral estimate. Archive \(C\), V84 already explains
how arbitrarily small low-energy spectral weight can defeat such
an inference; that obstruction is not rediscovered here.

For example, with \(D=V-W_j\) and a physical threshold \(E_*>0\),
spectral splitting only yields
\[
 \|e^{-tH_a}D\|
 \le \|{\bf1}_{[0,E_*]}(H_a)D\|
                   +e^{-tE_*}\delta_{j,\gamma,M}.
 \tag{V39.23}\label{r39:low-energy}
\]
The low-energy term on the right is not shown to be smaller
than \(\delta_{j,\gamma,M}\). Although every column is vacuum
orthogonal, its nonvacuum spectrum can still approach zero
along a regulator sequence. Increasing \(t\) suppresses the
high-energy term only. A physical normalization \(Z(a)\)
multiplies all source errors by \(|Z(a)|\);
V38's scale warning \eqref{r38:scale-warning} therefore remains.

A productive next target is a native, scale-sensitive estimate for
the low-energy residual or the retained connected correlations,
including their rough part and nonreducing returns.
It should use the physical insertion representation just proved,
rather than replace it by a Markov assertion about a retained
history. A noncollapsing source bank, coverage of the needed
physical sectors, nonperturbative scale generation, construction
of the interacting four-dimensional continuum, and a positive
physical Hamiltonian gap remain unproved.

\paragraph{Verification and scope ledger.}
The exact checker
\begin{center}\small
\path{verify_ym_restart_v39_slab_insertions.py}
\end{center}
enumerates the open-slab geometry at \(L=8,16,32\);
\(L=8\) is a geometric check only, not an enlargement of the
moment theorem's hypotheses. It checks 2336 middle blocks and
67744 full-weight incident faces, confirms both complete boundary
slices are retained, and checks reversal and the forward/backward
witness support. The native conditional integral is proved
analytically above, not evaluated by this enumeration.

An independent four-state rational positive-transfer fixture
checks the predictor kernel identity, a failure of selfadjointness
for a one-sided cutoff, and its repair by an intrinsic union
writer. It retains a nonzero conditional innovation contact
while its insertion vanishes. All principal minors needed for
the uncentered and centered norm bounds are checked exactly.
Direct path sums at \(n=0,1\), transfer powers at
\(n=0,1,2,7,20\), and all eight subsets of the three excited
spectral modes verify the same-operator identities.
This fixture is neither Wilson quadrature nor a spectral
certificate for the gauge theory. Uniform analytic constants
are not numerically evaluated and the proofs are not
machine-formalized.

New: the full-weight native two-step embedding; intrinsic
forward/backward frame agreement and the reflection-even union
writer; exact full-endpoint predictor insertions with the contact
term retained; uniform operator and physical-vacuum-vector
comparisons; and their use with the unchanged physical time and
spectral projections. Inherited: Wilson transfer positivity,
conditional expectation identities, the abstract Gram inequality,
V35 moments, and V38's local expansion and rates.
Removing this append and restoring the V38 title recovers the
entire predecessor exactly. Other sources remain unchanged.
All checks and build outputs stay outside \path{latex_notes}.

% END CONSOLIDATED V39 APPEND

% BEGIN CONSOLIDATED V40 APPEND
\clearpage
\section{V40: Native electric curvature, a source vertex, and the normalization checkpoint}
\label{r40:context}

V39 embeds the native source comparison in the unchanged physical
Wilson Hilbert space. Its error is absolute, however, and cannot by
itself control a physically normalized low-energy source. This checkpoint
asks a concrete upstream question: does the local residual have an extra
power of curvature, or can it mix with the electric energy already at
quadratic order in an actual link background?

The answer is not obtained by computing another inverse-coupling
coefficient. We realize V38's formerly abstract paired-pin test by
literal endpoint links, compute the \emph{curvature derivative of the
existing first coefficient}, and identify the corresponding exact
finite-coupling response. It is nonzero, with two distinct site weights.
A single scalar renormalization of the minimum energy cannot cancel all
these electric profiles. A profile-resolved response subtraction does
remove their quadratic term. Its proved scope is the specified native
test family, not a replacement of the full retained measure.

The interacting four-dimensional continuum and its mass gap remain
unproved. In particular a smooth classical background is not a sample
of the continuum quantum law and is not a low-energy spectral vector.

\subsection{Scheduled companion and broad primary-source reviews}

The V39 predecessor was verified against its recorded hash. Fresh
hashes show that the repository companions NS V26, SM--Einstein V72,
Shell Mixing V16, and Parallel YM V5 are unchanged since V35.
Their terminal proofs and scope statements were reconsidered:
NS's rank-one tensor improvement supplies no YM covariance estimate;
SM's many-copy, fixed-mode transport is not a cutoff-uniform physical
theory; Shell Mixing's measured flat-holonomy calculation still leaves
the noncommuting mixed-curvature and nonperturbative remainders; and
Parallel YM's global extensive-charge tightness target remains excluded.

The Downloads directory contains newer related notes than those in the
last explicit attachment list. They were read as research data, not
instructions. The relevant new statements inspected are:
\begin{itemize}
\item Source Selection V35: the stationary reverse is normalized by the
actual faithful state. Its contextual identity keeps both boundary
weights; a GNS adjoint, a trace adjoint and a CP reverse are different
objects. This reinforces V39's use of the actual physical adjoint, not
an arbitrarily chosen retained reversal.
\item Native Regenerative Stationarity V38: a nonzero complex-square
normal coefficient excludes the proposed nearby three-profile branch,
and an approximate solve pays an explicit conditioning loss. This
suggests testing a supposedly absent source coefficient directly,
rather than repeatedly shrinking a residual.
\item Exact-Action Einstein Bridge V28: enlarging the first-field family
removes an old Poisson cokernel, but three cubic means remain and a
differential coefficient lift is not the original physical-time
multiplier rate. No rank or lifespan constant transfers to YM.
\end{itemize}
Perturbative ESM V23 is unchanged. Its relevant coefficient and
normalization discipline is retained, not its perturbative exponents.
None of these companions proves the missing YM retained correlation
or physical low-energy estimate. These are targeted readings of the
relevant new statements and their proofs, not a recertification of
every companion manuscript.

The broad sweep reconsidered the user's literature suggestions and
searched constructive gauge theory, composite-source normalization,
critical SPDEs, multiscale functional inequalities, moment/loop
bootstraps, spectral RG, and local-to-global gap certificates.
The principal decisions are as follows.
\begin{enumerate}
\item \textbf{Composite operators.}
Suzuki,
\href{https://arxiv.org/html/1304.0533}
{\emph{Energy--momentum tensor from the Yang--Mills gradient flow}},
Sections 2--3, separates vacuum subtraction from operator mixing,
notably in (2.10)--(2.11) and (3.7)--(3.8).
This motivates the present native source-vertex test.
We import neither its renormalized continuum assumptions nor its
matching coefficients. Flow time is not physical transfer time.
Rothe's
\href{https://arxiv.org/abs/hep-lat/9508005}
{energy sum-rule and trace-anomaly abstract} was also screened;
its warning against identifying an action insertion with a physical
energy without the full normalization is relevant.
\item \textbf{Critical weak-coupling expansions.}
Deng--Shen,
\href{https://arxiv.org/html/2607.10105v2}
{\emph{The four-dimensional Anderson model: a case study for critical SPDEs}},
Theorem 1.1 and Section 1.2, treats a weakly coupled four-dimensional
Anderson model, not Yang--Mills. Its high-order, cutoff-dependent
remainder control is a useful future method; its limit is Gaussian.
It gives no nonperturbative YM spectrum or permission to replace
a uniform remainder by finitely many computed coefficients.
\item \textbf{Multiscale functional inequalities.}
Bauerschmidt--Bodineau--Dagallier,
\href{https://arxiv.org/html/2307.07619}
{\emph{Stochastic dynamics and the Polchinski equation}},
Theorem 3.6, requires a lower bound on the Hessian of the
\emph{renormalized} potential along a covariance decomposition.
The nonflat discussion retains a separate terminal convergence
assumption. Our native conditional minimum does not supply these
global hypotheses. The resulting dynamics is not automatically \(P\).
\item \textbf{Physical spectral moments.}
The
\href{https://arxiv.org/html/2605.20509}
{causal bootstrap of Abbott and collaborators}, Section II,
optimizes smeared spectral quantities against actual Euclidean
correlators and positivity. The
\href{https://arxiv.org/abs/2508.01377}
{matrix moment paper of Abbott--Jay--Oare} was checked at abstract
level. V39 now supplies compatible physical source vectors, but
not a certified growing-depth correlator bank. Finite absolute
uncertainty still does not exclude a tiny slow component.
\item \textbf{Loops and strong coupling.}
The current abstracts of
\href{https://arxiv.org/abs/2404.16925}
{Kazakov--Zheng's finite-\(N\) bootstrap},
\href{https://arxiv.org/abs/2505.16585}
{Cao--Nissim--Sheffield's expanded area-law regimes}, and
\href{https://arxiv.org/abs/2204.12737}
{Shen--Zhu--Zhu's strong-coupling stochastic analysis}
were screened. Exact loop constraints and controlled restoration
of a removed interaction remain useful architectures. Their
available conclusions do not provide our fixed-\(\SU(2)\)
weak-coupling continuum estimate.
\item \textbf{Spectral RG and certificate hierarchies.}
The current
\href{https://arxiv.org/abs/2508.07643}
{Bach--Ballesteros--Geisler abstract} confirms the smooth
Feshbach--Schur semigroup construction; it does not pay our
uncontrolled complementary resolvent. Frustration-free gap
hierarchies were searched, but no theorem was imported: the
proposed journal page was unavailable, and the required
frustration-free decomposition of the physical Wilson generator
has not been established.
\end{enumerate}
The full search and access ledger is stored outside
\path{latex_notes}. The next companion review is V45; the next
regular broad sweep is V50, retaining the count across archives.

\paragraph{Nonduplication.}
Archive \(A\), V58 already identifies the uniform action insertion
with a weighted derivative of the vacuum. It also explains why
this does not bound the unweighted response. V60--V62 continue that susceptibility
route. We do not rename those identities a new low-energy theorem.
V38 already computes \(a_1\) and its abstract paired-slot examples;
it explicitly does not realize those examples by native link
histories. The present advance is that realization, its electric
profile derivative, the exact finite-coupling response, and the
resulting distinction between smooth-field normalization and a
quantum-law estimate.

\subsection{Literal two-sided temporal endpoints}

Keep V39's two-step slab and one middle six-link block. Number its
vertices by
\[
 (0,0),(0,1),(0,2),(1,0),(1,1),(1,2).
\]
Let \(\mathcal G=P_2\times P_3\) be its seven-edge grid, and let
\(n_i=6-\deg_{\mathcal G}(i)\) be the number of private slots:
\(n_i=4\) at the four corners and \(n_i=3\) at the two middle
vertices. Exactly two private plaquettes at each vertex are temporal.

Choose the following \emph{actual} retained links in the original slab.
Every temporal link is the identity, as is every retained middle
spatial link. At the spatial location of free link \(i\), set the
corresponding \(\mu\)-link in the left endpoint slice to
\[
 q_i=\exp(\boldsymbol\theta_i\cdot\boldsymbol e)
       \in\SU(2),
\]
and in the right endpoint slice to \(q_i^{-1}\).
Here \(\boldsymbol e\) is the unit imaginary quaternion basis, so
the scalar part of \(q_i\) is \(\cos|\boldsymbol\theta_i|\).
Other endpoint links may be set to the identity.
The axes \(\boldsymbol\theta_i/|\boldsymbol\theta_i|\) need not
coincide or commute. Let \(z_i=(z_{i0},\boldsymbol z_i)\in S^3\)
be the six free middle links, and put
\[
 \delta_i=2(1-\cos|\boldsymbol\theta_i|),\qquad
 \Delta=\sum_i\delta_i .
 \tag{V40.1}\label{r40:delta}
\]

\begin{proposition}[An actual native electric deformation]
\label{r40:native}
The entire 29-plaquette conditional energy on these retained
configurations is exactly
\[
 E_\delta(z)
   =29-\sum_i(n_i-\delta_i)z_{i0}
                  -\sum_{\{i,j\}\in\mathcal G}z_i\cdot z_j
   =E_0(z)+\sum_i\delta_i z_{i0}.
 \tag{V40.2}\label{r40:energy}
\]
If \(0\le\delta_i< n_i\), its unique global minimum is
\(z_i=e_0\) for every \(i\), with value \(m(\delta)=\Delta\).
No endpoint spatial action is included in this conditional energy;
such factors remain in the full slab kernel and are independent
of the free middle links.
\end{proposition}
\begin{proof}
For a free link \(z_i\), its two temporal plaquettes contribute
\[
 [1-\operatorname{Re}(z_iq_i^{-1})]
       +[1-\operatorname{Re}(z_iq_i)]
 =2-2\cos|\boldsymbol\theta_i|\,z_{i0}.
\]
The imaginary terms cancel for each \(i\) separately; different
\(q_i\)'s do not need to commute. The other \(n_i-2\) private
spatial plaquettes contribute \(1-z_{i0}\).
An internal spatial plaquette has all transverse links equal
to the identity and contributes \(1-z_i\cdot z_j\).
The 12 temporal and 17 spatial plaquettes are exactly those
counted in V39, giving \eqref{r40:energy}.
Subtracting \(\Delta\) leaves
\[
 E_\delta(z)-\Delta
 =\sum_i(n_i-\delta_i)(1-z_{i0})
                    +\sum_{\{i,j\}\in\mathcal G}(1-z_i\cdot z_j).
 \tag{V40.3}\label{r40:positive}
\]
Every term is nonnegative and each private coefficient is
positive. Equality therefore forces all \(z_i=e_0\).
\end{proof}

This is a specification of an original link configuration, not
a change of the Wilson action or measure. Endpoint spatial
plaquettes can be nontrivial for inhomogeneous \(q_i\)'s and are
not silently deleted from that measure. Only their independence
of the free middle links is used.

The parameters have a gauge-invariant description on this family.
Let \(W_{2,i}\) be the half-trace of the temporal rectangle with
spatial width one and temporal length two over link \(i\).
On the displayed representative \(W_{2,i}=\cos(2|\boldsymbol\theta_i|)\).
Consequently, if \(|\boldsymbol\theta_i|<\pi/2\),
\[
 \delta_i=2-\sqrt{\,2+2W_{2,i}\,}.
 \tag{V40.4}\label{r40:rectangle}
\]
The sign of the square root is fixed by proximity to the identity.
Thus the electric test and its profile weights are not artifacts
of a gauge-coordinate choice. The right side is defined more
generally, but \eqref{r40:energy} is asserted only for this
paired-endpoint family and its gauge orbit.

As \(\boldsymbol\theta\to0\), all finite words in the forward and
backward V35 witnesses approach the identity. Their defects are
\(O(\sum_i|\boldsymbol\theta_i|^2)\), by telescoping the finitely
many link products. Hence both V39 good events hold in a fixed
small neighborhood of this family, independently of \(\gamma\).
The V39 writer there equals \(m\) or
\(m+[a_1(\delta)-a_1(0)]/\gamma^2\), using its intrinsic
frame agreement.

\subsection{A uniform radial-parameter expansion, not a differentiated big-O}

For purposes of differentiation extend \(\delta\) to the real box
\(\mathcal D=[-1/4,1/4]^6\), with energy \eqref{r40:energy}.
Negative \(\delta_i\)'s in this auxiliary extension need not be
realizable by the paired links. All statements about actual
exteriors below use \(\delta_i\ge0\).
The extension leaves a uniformly pinned minimum, since
\(n_i-\delta_i\ge11/4\).
Let
\[
 \begin{gathered}
 D_i=6-\delta_i,\qquad
 M_\delta=\operatorname{diag}(D_i)-\operatorname{Adj}(\mathcal G),
 \qquad G_\delta=M_\delta^{-1},\\
 Q_\gamma(\delta)=\int e^{-\gamma E_\delta(z)}\,dH(z),
 \qquad H_\gamma(\delta)=-\partial_\gamma\log Q_\gamma(\delta),
 \qquad h_\gamma=H_\gamma(0).
 \end{gathered}
 \tag{V40.5}\label{r40:partition}
\]
The full Hessian is \(M_\delta\otimes I_3\), and
\(11I/4\le M_\delta\le37I/4\).
These are six-spin conditional matrices, not a Hessian bound
for the collective retained field.

\begin{lemma}[Uniform differentiated radial remainder]
\label{r40:radial}
For \(\gamma\ge\gamma_*\), with a constant \(C_0>0\)
independent of \(\delta,\gamma\),
\[
 \log Q_\gamma(\delta)
 =-\gamma\Delta-9\log\gamma+\log C_0
       -\tfrac32\log\det M_\delta
       +\frac{A(\delta)}{\gamma}+R_\gamma(\delta),
 \tag{V40.6}\label{r40:Laplace}
\]
where
\[
 \|R_\gamma\|_{C^2(\mathcal D)}
          \le C\gamma^{-2},\qquad
 \|\partial_\gamma R_\gamma\|_{C^2(\mathcal D)}
          \le C\gamma^{-3}.
 \tag{V40.7}\label{r40:remainder}
\]
The first coefficient \(A(\delta)\) is V38's intrinsic
\(a_1\) for these energies, and is
\[
 A(\delta)
 =\frac32\sum_i (G_\delta)_{ii}
 -\frac{15}{8}\sum_iD_i(G_\delta)_{ii}^2
 +\sum_{\{i,j\}\in\mathcal G}
    \left\{\frac94(G_\delta)_{ii}(G_\delta)_{jj}
                            +\frac32(G_\delta)_{ij}^2\right\}.
 \tag{V40.8}\label{r40:A}
\]
\end{lemma}
\begin{proof}
The global positive representation \eqref{r40:positive}
holds throughout a neighborhood of \(\mathcal D\) with a
strictly positive pinning constant. Outside a fixed neighborhood
of \(e_0^6\) it gives a uniform positive phase barrier.
In a fixed neighborhood the phase has a uniformly positive
Hessian and smooth parameter dependence of every finite order.
The parameter-dependent Morse change of variables used in
V38 therefore reduces the phase \emph{exactly} to a positive
quadratic form, with a smooth compactly supported amplitude
and uniform bounds on as many derivatives as needed.

Expand that amplitude to even degree two with the degree-four
Taylor remainder retained, and rescale by \(\gamma^{-1/2}\).
This gives a normalized integral
\(1+A(\delta)/\gamma+O_{C^2(\delta)}(\gamma^{-2})\).
For its actual \(\gamma\) derivative, differentiate the normalized
Gaussian kernel as in \eqref{r38:kernel-derivative}; the resulting
polynomial kernel integrates the constant exactly to zero and
differentiates the quadratic term exactly. Applying it to the
degree-four amplitude remainder costs \(O_{C^2}(\gamma^{-3})\).
The same bounds apply after two \(\delta\) derivatives because
the Morse amplitude and its parameter derivatives have the
uniform Taylor bounds just described. Differentiated tails
remain exponentially small. Taking the logarithm, whose
argument is uniformly bounded away from zero, proves
\eqref{r40:Laplace}--\eqref{r40:remainder}.
This proves the derivative bound on the integral, rather than
formally differentiating scalar asymptotic notation.

To evaluate the coefficient independently, use the local
Cartesian chart \(z_i=(\sqrt{1-|x_i|^2},x_i)\).
The relative Haar density is
\(\prod_i(1-|x_i|^2)^{-1/2}\); its quadratic term is
\(\frac12\sum_i|x_i|^2\). The phase is even, with quadratic
form \(M_\delta\otimes I_3\) and quartic part
\[
 \frac18\sum_iD_i|x_i|^4
                  -\frac14\sum_{\{i,j\}\in\mathcal G}|x_i|^2|x_j|^2 .
\]
There is no cubic term. The Gaussian covariance is
\(G_\delta\otimes I_3\).
The Wick identities
\(\E|X_i|^2=3G_{ii}\),
\(\E|X_i|^4=15G_{ii}^2\), and
\(\E(|X_i|^2|X_j|^2)=9G_{ii}G_{jj}+6G_{ij}^2\)
give \eqref{r40:A}. Uniqueness of the normalized Laplace
coefficient identifies it with V38's intrinsic \(a_1\).
\end{proof}

Set \(r_\gamma(\delta)=H_\gamma(\delta)-h_\gamma-\Delta\),
the uncorrected V39 residual on this family.
Differentiating the \emph{proved} logarithmic formula gives
\[
 \begin{split}
 r_\gamma(\delta)
  &=\frac{A(\delta)-A(0)}{\gamma^2}
          -\partial_\gamma[R_\gamma(\delta)-R_\gamma(0)],\\
 \left|r_\gamma(\delta)
       -\frac{A(\delta)-A(0)}{\gamma^2}\right|
  &\le C\gamma^{-3}\|\delta\|_1,\qquad
 \|r_\gamma\|_{C^2(\mathcal D)}\le C\gamma^{-2}.
 \end{split}
 \tag{V40.9}\label{r40:residual}
\]
The determinant has no \(\gamma\) derivative.
As before, the common \(9/\gamma\) cancels through the exact
coherent mean \(h_\gamma\).

\subsection{The nonzero, profile-dependent electric vertex}

Put \(G=M_0^{-1}\). For the parameter \(\delta_p\), inverse
differentiation gives
\[
 J^{(p)}_{ij}
 =\left.\partial_{\delta_p}(G_\delta)_{ij}\right|_0
 =G_{ip}G_{pj}.
 \tag{V40.10}\label{r40:J}
\]
Define \(\kappa_p=\partial_{\delta_p}A(0)\).

\begin{proposition}[Exact native electric response coefficient]
\label{r40:kappa}
With the grid ordering above,
\[
 \begin{split}
 \kappa_p={}&
 \frac32\sum_iJ^{(p)}_{ii}
 +\frac{15}{8}G_{pp}^2
 -\frac{45}{2}\sum_iG_{ii}J^{(p)}_{ii}\\
 &+\sum_{\{i,j\}\in\mathcal G}
    \left\{\frac94(J^{(p)}_{ii}G_{jj}+G_{ii}J^{(p)}_{jj})
                                    +3G_{ij}J^{(p)}_{ij}\right\}.
 \end{split}
 \tag{V40.11}\label{r40:kappa-formula}
\]
There are exactly two values:
\[
 \begin{gathered}
 \kappa_{\mathrm{corner}}
      =\frac{43168352898}{10832063831575},\qquad
 \kappa_{\mathrm{middle}}
      =\frac{150461307}{12632144410}
         >\kappa_{\mathrm{corner}}>0,\\
 \kappa_{\mathrm{tot}}
       =4\kappa_{\mathrm{corner}}+2\kappa_{\mathrm{middle}}
       =\frac{430714553097}{10832063831575}.
 \end{gathered}
 \tag{V40.12}\label{r40:kappa-values}
\]
Thus
\[
 r_\gamma(\delta)
 =\gamma^{-2}\sum_i\kappa_i\delta_i
      +O\!\left(\gamma^{-2}\|\delta\|_1^2
                      +\gamma^{-3}\|\delta\|_1\right).
 \tag{V40.13}\label{r40:vertex}
\]
The constants are uniform near this family, including in the
endpoint color directions.
\end{proposition}
\begin{proof}
Differentiate \(M_\delta G_\delta=I\), using
\(\partial_{\delta_p}M_\delta=-e_pe_p^{\mathsf T}\), to get
\eqref{r40:J}. Differentiate each term of
\eqref{r40:A}; its diagonal coefficient derivative contributes
\(15G_{pp}^2/8\), and \(D_i=6\) at the origin supplies the
factor \(-45/2\). This gives \eqref{r40:kappa-formula}.
Inverting the displayed integer matrix \(M_0\), of determinant
\(37835\), and substituting its entries gives the rational
values \eqref{r40:kappa-values}. The grid symmetries identify
the four corner values and the two middle values. Positivity
and their strict inequality are rational comparisons.

An independent check differentiates the formula obtained in V38
from normal-coordinate Haar factors and Wick contractions,
whose separate Haar and quartic terms differ from the
Cartesian ones. Finally Taylor's theorem for the smooth
rational function \(A\), followed by \eqref{r40:residual},
proves \eqref{r40:vertex}. Exact formula
\eqref{r40:energy} removes the color directions from the
conditional law, so the constants do not depend on them.
\end{proof}

Since
\(\delta_i=|\boldsymbol\theta_i|^2
+O(|\boldsymbol\theta_i|^4)\),
the displayed leading term is a nonzero quadratic electric
vertex. Coherence's vanishing first derivative in the link
angles does \emph{not} force this quadratic term to vanish.
This is an actual native exterior test, unlike the explicitly
abstract examples in V38.

\subsection{An exact finite-coupling response and its subtraction}

All expectations below use the same coherent six-spin law
\[
 \lambda_{\gamma,0}(dz)=Q_\gamma(0)^{-1}e^{-\gamma E_0(z)}dH(z).
\]
Write \(M_i(\gamma)=\E_{\gamma,0}z_{i0}\).

\begin{theorem}[Normalized response at finite coupling]
\label{r40:finite-response}
For each \(\gamma>0\), define
\[
 \begin{split}
 \rho_i(\gamma)
   &:=\left.\partial_{\delta_i}r_\gamma(\delta)\right|_0\\
   &=M_i(\gamma)-\gamma\Cov_{\gamma,0}(E_0,z_{i0})-1\\
   &=M_i(\gamma)+\gamma M_i'(\gamma)-1
     =-\frac{d}{d\gamma}
                     \bigl[\gamma\,\E_{\gamma,0}(1-z_{i0})\bigr].
 \end{split}
 \tag{V40.14}\label{r40:rho}
\]
It is an exact conditional response, including the derivative
of the normalizer. Uniformly for \(\gamma\ge\gamma_*\),
\[
 \rho_i(\gamma)=\frac{\kappa_i}{\gamma^2}
                                      +O(\gamma^{-3}).
 \tag{V40.15}\label{r40:rho-asymptotic}
\]
On the paired-endpoint family the test-source correction
\[
 K_{\gamma}^{\mathrm{pair}}(\delta)
             =\Delta+\sum_i\rho_i(\gamma)\delta_i
 \tag{V40.16}\label{r40:pair-source}
\]
satisfies
\[
 \left|H_\gamma(\delta)-h_\gamma
                           -K_{\gamma}^{\mathrm{pair}}(\delta)\right|
             \le C\gamma^{-2}\|\delta\|_1^2 .
 \tag{V40.17}\label{r40:pair-error}
\]
\end{theorem}
\begin{proof}
All integrations are on a fixed compact space with bounded
energy and insertions. Differentiating the normalized integral
at fixed \(\gamma\), using
\(\partial_{\delta_i}E_\delta=z_{i0}\), gives
\[
 \left.\partial_{\delta_i}H_\gamma(\delta)\right|_0
       =\E_{\gamma,0}z_{i0}
                           -\gamma\Cov_{\gamma,0}(E_0,z_{i0}).
\]
The derivative of \(m=\Delta\) is one, proving the first
two lines of \eqref{r40:rho}. Independently,
\(M_i'=-\Cov_{\gamma,0}(z_{i0},E_0)\), proving the remaining
identities. Equation~\eqref{r40:residual}, with its
\(C^2(\delta)\) remainder, gives \eqref{r40:rho-asymptotic}.
Taylor's theorem at \(\delta=0\), using \(r_\gamma(0)=0\)
and its \(C^2\) bound, gives \eqref{r40:pair-error}.
\end{proof}

The constants \(\rho_i(\gamma)\) are given by fixed
six-spin integrals; they have not been replaced by an
unverified numerical fit. The formula is valid at finite
coupling, not merely to a chosen inverse-coupling order.
Its exact normalizer subtraction is essential.
There is no claim that \eqref{r40:pair-source} already
renormalizes a general YM composite operator.

\begin{corollary}[A scalar cell-energy rescaling fails this native test]
\label{r40:scalar-no}
For all sufficiently large \(\gamma\), no scalar \(c_\gamma\)
can satisfy, for every small paired electric profile,
\[
 H_\gamma(\delta)-h_\gamma
       -(1+c_\gamma)m(\delta)
           =O\!\left(\left(\sum_i|\boldsymbol\theta_i|^2\right)^2\right) .
 \tag{V40.18}\label{r40:no-scalar}
\]
A scalar can match the uniform profile but not all six
independent profiles.
\end{corollary}
\begin{proof}
Allow only \(\boldsymbol\theta_p\) to be nonzero and let it
tend to zero. The coefficient of its squared norm in
\eqref{r40:no-scalar} forces \(c_\gamma=\rho_p(\gamma)\).
Requiring this for a corner and a middle vertex contradicts
\eqref{r40:rho-asymptotic} and
\(\kappa_{\mathrm{middle}}>\kappa_{\mathrm{corner}}\), once
\(\gamma\) is sufficiently large. In the limiting coefficient
problem the unavoidable worst error is at least
\[
 \inf_c\max_i|\kappa_i-c|
       =\frac{\kappa_{\mathrm{middle}}
                         -\kappa_{\mathrm{corner}}}{2}
       =\frac{171704435709}{43328255326300}>0 .
\]
For the uniform profile only the sum of the six coefficients
is tested, so \(c_\gamma=\frac16\sum_i\rho_i(\gamma)\)
can match that one direction.
\end{proof}

This obstruction concerns one scalar multiple of the
\emph{specified block minimum}. It is not a no-go theorem
for a properly chosen multi-operator counterterm, an
averaged translation-covariant source, or an interacting
continuum construction. It tells the next calculation
which profile information a scalar calibration loses.

\subsection{What the calculation says about lattice-spacing normalization}

There is a particularly transparent smooth background.
Use the same \(q=\exp(a^2 E\,e_1)\) on every \(\mu\)-link
of the full left slice and \(q^{-1}\) on the full right
slice, with all other links as before and the minimizing
middle slice equal to the identity. Every spatial
plaquette is then the identity, including across periodic
spatial seams. Locally in physical time \(s\in[-a,a]\),
this is sampled from
\[
 A_\mu(s)=-sE\,e_1,\qquad A_\nu=A_\rho=A_0=0 ,
\]
using link convention \(\exp(aA_\mu)\).
Its classical electric curvature is constant \(-E e_1\).
This is a local background test; no globally periodic
linear-in-time continuum field is required or asserted.

All six \(\delta_i\)'s now equal
\(\delta(a)=2(1-\cos(a^2E))=a^4E^2+O(a^8E^4)\).
For bounded fixed \(E\), \eqref{r40:vertex} yields, uniformly
for small \(a\) and large \(\gamma\),
\[
 \begin{split}
 a^{-4}r_\gamma(\delta(a)\mathbf1)
  &=\frac{\kappa_{\mathrm{tot}}E^2}{\gamma^2}
           +O\!\left(\frac{a^4E^4}{\gamma^2}
                               +\frac{E^2}{\gamma^3}\right),\\
 \gamma^2a^{-4}r_\gamma(\delta(a)\mathbf1)
  &=\kappa_{\mathrm{tot}}E^2+O(a^4E^4+E^2/\gamma).
 \end{split}
 \tag{V40.19}\label{r40:physical-scaling}
\]
In particular for any \(\gamma(a)\to\infty\),
the second line tends to the explicitly positive
\(\kappa_{\mathrm{tot}}E^2\) when \(E\ne0\).
There is no extra positive power of \(a\) after factoring
the electric-energy scaling and the proved
\(\gamma^{-2}\) coefficient. In contrast,
\eqref{r40:pair-error} gives
\[
 \gamma^2a^{-4}
 \left|H_\gamma(\delta(a)\mathbf1)-h_\gamma
                  -K_\gamma^{\mathrm{pair}}(\delta(a)\mathbf1)\right|
                    \le C a^4E^4 .
 \tag{V40.20}\label{r40:improved-scaling}
\]
This is a real improvement in \emph{curvature order},
not another finite term in an inverse-coupling series.

These formulas also refine the interpretation of V38's
scale warning. Its coarse unconditional defect bound alone
cannot certify a negative-power-of-\(a\) normalization.
On a specified smooth background, however, the actual defect
has its own \(a^4\) factor and the uncorrected \(a^{-4}\)
residual in \eqref{r40:physical-scaling} tends to zero as
\(\gamma\to\infty\). There is no contradiction: the hypotheses
are different. Neither formula proves that quantum exteriors
have such smooth curvature scaling, or that their normalized
errors vanish in a physical vacuum norm.

For example a generalization of \eqref{r40:pair-error} would
need a proof for the full noncommuting native coefficient
family, including magnetic and mixed responses, and a bound
on the resulting fourth-order curvature remainder under the
\emph{original} retained law. V35 controls bare defects; it
does not give a cutoff-uniform moment for plaquette angles divided
by \(a^2\), or
the required renormalized composite field. Evaluating a
classical constant-curvature background cannot replace that
missing quantum estimate. No spectral statement follows
merely because the test profile is smooth.

\subsection{Result, verification, and the next nonduplicating target}

The new exact checker
\begin{center}\small
\path{verify_ym_restart_v40_electric_vertex.py}
\end{center}
evaluates all 29 original Wilson words in five rational
noncommuting free-spin and endpoint-color fixtures and
compares them with \eqref{r40:energy}. It checks the
two-step rectangle identity and positivity of the displayed
pin coefficients. A separate full periodic endpoint check
tests 384 spatial faces for the uniform electric profile.
The matrix inverse derivative is checked entrywise for
each of the six directions, and two independent Haar
coordinate formulas give the same rational
\(\kappa_i\)'s and their strict inequalities.
A finite normalized-law fixture checks the covariance-response
sign; it is not Wilson quadrature. Uniform analytic remainder
constants and a numerical large-\(\gamma\) threshold are
not evaluated. The proofs are written analytic arguments,
not proof-assistant formalizations.

The native electric vertex rules out a specific shortcut:
coherent centering and a single block-energy normalization
do not remove all quadratic local electric response.
The exact response subtraction supplies a benchmark that
a general source-renormalization argument should reproduce.
V38's already corrected \(K^{(1)}\) subtracts the leading
\(A(\delta)-A(0)\) term; it does not remove this exact
finite-coupling response to all orders. The calculation
does not alter V39's physical transfer or its rough-sector
prescription, and does not claim a new global writer.

The next target is to identify the complete native quadratic
response on actual retained-link directions modulo gauge,
including electric, magnetic and mixed components, with a
remainder suitable for the original-law physical source
norm. The first test should determine which components are
forced by symmetry and which must be retained, not assume
that the two electric coefficients above exhaust the map.
That target must remain coupled to positive-time correlations,
source noncollapse and scale generation; a sequence of
ever higher local coupling coefficients is not a substitute.

The full interacting four-dimensional continuum construction
and its physical mass gap remain open in this programme.
Removing this append and restoring the V39 title recovers
the predecessor exactly. All other sources are preserved,
only the latest active main \(\textrm{.tex}\) is retained,
and build/check files remain outside \path{latex_notes}.

% END CONSOLIDATED V40 APPEND

% BEGIN CONSOLIDATED V41 APPEND -- 2026-09-19
\clearpage
\section[V41: native cycle response and magnetic source curvature]
{V41: Native cycle response\\and magnetic source curvature}
\label{r41:start}

\subsection{The upstream test and the nonduplication boundary}

V40's paired electric family does not exhaust the retained variables.
We now identify the complete quadratic response of the same
six-free-link conditional energy, using actual retained links rather
than freely prescribed slot coefficients. There are 25 retained cycles
per color after based gauge removal. The classical partial minimum
sees only 23 of them at quadratic order. The other two are transverse
magnetic cycles, and we compute a strictly positive thermal source
response on both. A source correction based only on that block's
minimum therefore misses actual native directions.

The positive result is an exact finite-coupling quadratic source
correction in all 25 cycle variables. Its cubic remainder is bounded
under the original interacting law, and hence in V39's physical vacuum
source norm. This improves the previously proved error; it does not
give decay of retained correlations, exclude low-energy spectral
weight, or construct the interacting continuum theory.

Tree gauge, Schur complements, normalized cumulant differentiation,
and Taylor subtraction are standard methods. V28 already separates
based gauge geometry from the quotient probability measure; V33
already proves the full-action Maxwell Schur formula; V38 supplies
the first Laplace coefficient; V39 supplies the physical source map.
Archive $A$, V99 already studies a moving-frame magnetic response and
a fixed-volume Gaussian comparison on a different five-chord cube
packing. Its magnetic carrier and soft-mode warning are not new
results here. The new items are this block's 25-cycle geometry,
its two-dimensional magnetic response, the full finite-coupling
response formula, and its native-defect-controlled remainder.

The targeted primary-source check is I.~M.~Burbano and C.~W.~Bauer,
\emph{Gauge Loop-String-Hadron Formulation on General Graphs and
Applications to Fully Gauge Fixed Hamiltonian Lattice Gauge Theory},
\url{https://arxiv.org/html/2409.13812v2}, Sections III.1 and C.3.2.
Their tree construction reduces gauge variables to based loop
holonomies; dynamics still requires the transformed operators.
We use only this standard coordinate viewpoint, proving our finite
graph and estimates below. We import neither a Hamiltonian truncation
nor a spectral conclusion. The regular companion and broad literature
reviews were completed at V40; the next are V45 and V50.

\subsection{The actual 80 retained links and their 25 cycles}

Place the six free $\mu$ links at middle time $1$, at positions
$(0,c,r,1)$ with $c=0,1$ and $r=0,1,2$, as in V39.
Coordinates have order $(\mu,\nu,\rho,\tau)$.
Let $\mathcal P$ be their 29 incident plaquettes and let $\mathcal R_b$
contain every other link in these plaquettes. Include no additional
exterior action in the conditional energy
\[
 E(z;\eta)=\sum_{p\in\mathcal P}\{1-\operatorname{Re}U_p(z,\eta)\},
 \qquad
 Q_\gamma(\eta)=\int_{(S^3)^6}e^{-\gamma E(z;\eta)}\,dH(z).
 \tag{V41.1}\label{r41:native}
\]
Untouched plaquettes remain in the full Wilson law, as before.
Write $H_\gamma=-\partial_\gamma\log Q_\gamma$,
$h_\gamma=H_\gamma(1)$ and $m(\eta)=\min_zE(z;\eta)$ wherever the
controlled smooth minimum is used.

\begin{proposition}[Exact local ranks]
\label{r41:ranks}
The supporting graph has 56 vertices and 86 links, six free and
80 retained. The retained graph is connected. For one real color,
let $C_F$ and $C_R$ be the signed plaquette incidence matrices with
six and 80 columns, and let $B$ be the retained vertex-gradient
matrix. Then
\[
 \begin{gathered}
 C_F^{\mathsf T}C_F=\mathsf A
       =6I-\operatorname{Adj}(P_2\times P_3),\qquad
 \det\mathsf A=37835,\\
 \operatorname{rank}B=55,\qquad \operatorname{rank}C_R=29,\\
 \Pi_F^\perp=I-C_F\mathsf A^{-1}C_F^{\mathsf T},\qquad
 S=C_R^{\mathsf T}\Pi_F^\perp C_R,\\
 \operatorname{rank}S=23,\qquad SB=0 .
 \tag{V41.2}\label{r41:ranks-eq}
 \end{gathered}
\]
Thus the based retained quotient has 25 scalar tangent coordinates;
$\ker S/\operatorname{im}B$ has dimension two. The color factors
multiply these dimensions by three.
\end{proposition}
\begin{proof}
There are two six-vertex transverse grids at the ends of the free
links. Each of the 22 private plaquettes adds a three-link retained
path between a pair of these endpoints, with two distinct new
vertices. The seven internal plaquettes add one grid edge on each
side. Hence the counts are $12+44=56$ and
$6+66+14=86$. The grid edges and private paths make the retained
graph connected, giving $\operatorname{rank}B=56-1$.

Every private row has its own outer retained $\mu$ link.
Each internal row has its own left transverse grid link.
These 29 columns give a signed diagonal $29$-column submatrix.
Therefore $C_R$ is onto $\mathbb R^{29}$.
The free incidence gives the displayed $\mathsf A$, which is positive:
its diagonal minus off-diagonal absolute row sum is at least three.
Thus $\Pi_F^\perp$ is an orthogonal projection of rank $29-6=23$.
Surjectivity of $C_R$ gives $\operatorname{rank}S=23$.
For every full vertex gradient the plaquette curl vanishes, so
$C_RB$ lies in the range of $C_F$, proving $SB=0$.
Finally $80-55=25$ and $80-23-55=2$.
\end{proof}

This $S$ is the local incident-energy Schur form.
It is not V33's entire retained-action Schur form: the latter
also includes untouched plaquettes. No global flatness or
physical zero mode follows from its two extra local null directions.

Here is a particular retained tree $\mathcal T_+$.
On the left grid choose its two left-column $\rho$ edges and
all three $\nu$ edges. Include all three edges of each of the six
upper-temporal private paths. For the other 16 private paths,
include their two transverse legs, leaving their outer $\mu$ link
out. These $5+18+32=55$ links form a tree on all 56 vertices.
Its chords are the 16 outer $\mu$ links, seven right-grid edges,
and two left-grid $\rho$ edges. Fix a root on the left grid,
gauge the tree links to identity, and denote the remaining
holonomies by $q_1,\ldots,q_{25}$. Root gauge acts by simultaneous
conjugation. Locally write
\[
 q_k=\exp(\xi_k\cdot{\bf e}),\qquad
 \xi=(\xi_1,\ldots,\xi_{25})\in(\mathbb R^3)^{25}.
 \tag{V41.3}\label{r41:coordinates}
\]
No gauge fixing is imposed on the probability law. These are
coordinates for evaluating gauge-invariant scalar functions of
the original exterior. Haar translations of the six free links
preserve \eqref{r41:native}. A quotient metric, Jacobian, or Laplacian
is not being substituted for a physical one.

\begin{lemma}[The original defect pays all cycle coordinates]
\label{r41:cycle-defect}
Let $\mathcal D^+$ be exactly V35's defect, using the upper filling.
For every exterior,
\[
 \sum_{k=1}^{25}|q_k-1|^2
 \le 2E_{\rm priv}(\Phi\eta)+4E_{\rm int}(\Phi\eta)
               +6\{e(L_1)+e(L_2)\}
 \le\mathcal D^+ .
 \tag{V41.4}\label{r41:chord-bound}
\]
If the principal logarithms satisfy $|\xi|\le\rho:=1/256$, then
$|\xi|^2\le4\mathcal D^+$. Conversely
$\mathcal D^+\le\rho^2/4$ ensures that these logarithms exist and
$|\xi|\le\rho$.
\end{lemma}
\begin{proof}
In this tree gauge all six upper filling paths are identity.
The six reference private faces therefore have zero filled energy.
Each other private face is, up to inversion, its outer $\mu$ chord,
so these 16 squared chords sum to $2E_{\rm priv}(\Phi\eta)$.
The two left chords have squared norms summing to
$2\{e(L_1)+e(L_2)\}$; their fundamental loops are the two left
grid squares, up to conjugation and orientation.

For an internal edge write $l$ and $r$ for the left and right
transverse links in this gauge. Its filled plaquette has energy
$e(lr^{-1})$. The triangle inequality gives
\[
 |r-1|^2\le2|l-1|^2+2|r-l|^2
               =2|l-1|^2+4e(lr^{-1}).
\]
Only the two left chords are nonidentity among the seven $l$'s.
Summing these inequalities and the previous contributions gives
\eqref{r41:chord-bound}. V35 has coefficient 16, rather than six,
on the two left squares. For small unit quaternions,
$|\log q|\le2|q-1|$ in the stated neighborhood. If the defect is
at most $\rho^2/4$, every chord is already in that neighborhood,
and the same inequality proves the converse.
\end{proof}

In this gauge the private slots are the outer chords or identity.
An internal transport uses one left and one right transverse link;
its distance from identity is at most $|l-1|+|r-1|$.
Consequently its full coefficient distance satisfies
\[
 d(\vartheta(\xi))^2\le2\sum_k|q_k-1|^2\le2|\xi|^2 .
 \tag{V41.5}\label{r41:coefficient}
\]
In particular $|\xi|\le2\rho$ is strictly within V31's controlled
coefficient set. This verifies the minimum and Laplace hypotheses
for all 75 coordinates, not just an electric subfamily.

\subsection{The complete normalized quadratic response}

All expectations here use the same coherent law
$\lambda_{\gamma,0}=Q_\gamma(0)^{-1}e^{-\gamma E_0}\,dH$.
Let $I=(k,a)$ denote a chord and color coordinate, and set
\[
 F_I=\partial_I E(z;0),\qquad
 K_{IJ}=\partial_I\partial_J E(z;0),\qquad
 S^{\mathcal T}_{IJ}=\partial_I\partial_J m(0).
 \tag{V41.6}\label{r41:derivatives}
\]
These are literal derivatives of the 29 original plaquette words
with tree links fixed to identity. $S^{\mathcal T}$ is the restriction
of $S\otimes I_3$ to chord columns. Define
\[
 \begin{split}
 \mathsf B_{\gamma,IJ}^{(0)}
 ={}&\E K_{IJ}-\gamma\Cov(E_0,K_{IJ})
       -2\gamma\Cov(F_I,F_J)\\
    &+\gamma^2\operatorname{cum}(E_0,F_I,F_J)
       -S^{\mathcal T}_{IJ},
 \end{split}
 \tag{V41.7}\label{r41:response}
\]
where the third cumulant is the expectation of the product of
the three centered variables. Put, for $j=0,1$,
\[
 \begin{split}
 r_\gamma^{(j)}(\xi)
   &=H_\gamma(\xi)-h_\gamma-m(\xi)
           -j\gamma^{-2}\{a_1(\xi)-a_1(0)\},\\
 \mathsf B_\gamma^{(j)}
   &=\mathsf B_\gamma^{(0)}-j\gamma^{-2}D^2a_1(0).
 \end{split}
 \tag{V41.8}\label{r41:residual}
\]
Here $a_1$ is V38's intrinsic Gaussian contraction, pulled back
along these actual retained-link coordinates.

\begin{theorem}[All directions, with the normalizer retained]
\label{r41:full-response}
For every positive $\gamma$,
$\mathsf B_\gamma^{(j)}=D^2r_\gamma^{(j)}(0)$.
There is a real symmetric $25$ by $25$ matrix $b_\gamma^{(j)}$ with
\[
 \mathsf B_{\gamma,(k,a)(l,b)}^{(j)}
            =b_{\gamma,kl}^{(j)}\delta_{ab}.
 \tag{V41.9}\label{r41:color}
\]
In unreduced retained-link coordinates the same Hessian annihilates
every vertex-gradient direction. For sufficiently large $\gamma$,
uniformly on $|\xi|\le2\rho$,
\[
 \begin{split}
 \|r_\gamma^{(j)}\|_{C^3}&\le C\gamma^{-2-j},\\
 \left|r_\gamma^{(j)}(\xi)
      -\tfrac12\sum_{k,l}b_{\gamma,kl}^{(j)}
                              \xi_k\cdot\xi_l\right|
       &\le C\gamma^{-2-j}|\xi|^3 .
 \end{split}
 \tag{V41.10}\label{r41:cubic}
\]
The constants are independent of surrounding volume and location.
\end{theorem}
\begin{proof}
Differentiate the actual compact integral at fixed $\gamma$:
\[
 \partial_I\partial_J\log Q_\gamma
       =-\gamma\E K_{IJ}+\gamma^2\Cov(F_I,F_J).
\]
The derivative of any fixed observable's expectation with respect
to $\gamma$ is minus its covariance with $E_0$. The derivative of
$\Cov(F_I,F_J)$ is $-\operatorname{cum}(E_0,F_I,F_J)$.
Applying $-\partial_\gamma$ gives the first four terms of
\eqref{r41:response}. The free-variable stationary equation and
Schur completion give the last one after subtracting $m$.
This proves the exact formula, including the normalizer derivatives.

The scalars $H_\gamma,m,a_1$ are gauge invariant; for $a_1$ this
also follows from uniqueness of the normalized Laplace coefficient.
Simultaneous conjugation fixes $\xi=0$ and rotates each $\xi_k$.
There is no invariant linear functional on this representation,
and its invariant bilinear forms are multiples of the Euclidean
inner product. Hence $Dr_\gamma^{(j)}(0)=0$ and
\eqref{r41:color} holds. Also $r_\gamma^{(j)}(0)=0$.
In unreduced coordinates differentiate
$Dr_\gamma^{(j)}(\eta)[X_\chi(\eta)]=0$ at identity.
The term involving the first derivative vanishes; the remaining
term says that the Hessian kills the gauge vector $B\chi$.

For the analytic estimate repeat V38's proof with three, rather
than two, retained-coordinate derivatives. By
\eqref{r41:coefficient}, the minimum, Hessian bounds and phase
barrier hold on a neighborhood of the compact ball. The
parameter-dependent Morse map has bounded derivatives of every
fixed order there. Expand its density through cubic degree in
the Gaussian variable. The remainder and its first three parameter
derivatives are bounded by $C|u|^4$. The normalized Gaussian
kernel derivative in \eqref{r38:kernel-derivative} bounds the
integrated remainder by $C\gamma^{-2}$ and its coupling derivative
by $C\gamma^{-3}$, also after three parameter derivatives.
The complementary integral costs a polynomial in $\gamma$ times
the uniform exponential barrier. Taking the logarithm is harmless
because the leading normalized integral is uniformly positive.
Thus, in $C^3$,
\[
 H_\gamma(\xi)
   =m(\xi)+9/\gamma+a_1(\xi)/\gamma^2+O(\gamma^{-3}).
\]
Subtract its value at zero to get both $C^3$ bounds.
Taylor's formula with the integral third derivative and the
vanishing value and first derivative proves \eqref{r41:cubic}.
\end{proof}

Formula \eqref{r41:response} specifies every electric, magnetic,
and mixed entry by fixed coherent six-spin integrals. It does not
claim that all 325 scalar entries have been evaluated numerically.
Estimating its large covariance terms separately would lose the
cancellation proved by the normalized integral argument.
A cubic invariant $\xi_k\cdot(\xi_l\times\xi_m)$ is allowed by
simultaneous conjugation; we have not proved its absence for this
whole native family. V40's quartic paired-family remainder is
therefore not extended to all retained directions.

\subsection{The two magnetic directions have a nonzero response}

Number grid vertices as in V40: $0,1,2$ on the left column and
$3,4,5$ on the right. For every oriented grid edge $i<j$, assign
the same $U_{ij}$ to its transverse link at $\mu=0$ and $\mu=1$.
Set every other retained link to identity. These are actual
middle-slice spatial links in the original slab. Direct multiplication
of the original plaquettes gives
\[
 E_U(z)=\sum_i n_i(1-z_{i0})
       +\sum_{i<j:\,i\sim j}
            \{1-z_i\cdot\Ad(U_{ij})z_j\},
       \qquad n_i=6-\deg(i).
 \tag{V41.11}\label{r41:magnetic-energy}
\]
Conjugation on quaternions fixes their scalar part and rotates
their three vector components. Each summand is nonnegative,
and $n_i\ge3$, so the unique global minimum is $z_i=1$, with
\[
                    m(U)=0
 \tag{V41.12}\label{r41:magnetic-min}
\]
for this whole family. This does not set untouched $\nu\rho$
plaquettes to identity or remove their action.

At first order the seven simultaneous left/right link changes
have zero incident plaquette curl. Modulo gauge they have dimension
two: the connected grid has seven edges and six vertices.
More explicitly, gradient changes extend to gauge changes by taking
the same value at paired left/right vertices and along their
private paths; the two grid curls distinguish the remaining
directions. This realizes exactly $\ker S/\operatorname{im}B$
in Proposition~\ref{r41:ranks}.

\begin{proposition}[An evaluated magnetic source quadratic form]
\label{r41:magnetic}
Take $U_{ij}(t)=\exp(tc_{ij}{\bf e}_3)$, with real $c_{ij}$, and put
\[
 \begin{split}
 \beta_1&=c_{03}+c_{34}-c_{14}-c_{01},\\
 \beta_2&=c_{14}+c_{45}-c_{25}-c_{12}.
 \end{split}
 \tag{V41.13}\label{r41:beta}
\]
Then V38's same first coefficient satisfies
\[
 \begin{gathered}
 a_1(U(t))-a_1(0)=t^2\beta^{\mathsf T}\mathsf L\beta+O(t^4),\\
 \mathsf L=\frac1{10832063831575}
 \begin{pmatrix}
 25719719164&1035710444\\
 1035710444&25719719164
 \end{pmatrix}\succ0 .
 \end{gathered}
 \tag{V41.14}\label{r41:L}
\]
For fixed $c$, uniformly at small $t$ and large $\gamma$,
\[
 H_\gamma(U(t))-h_\gamma
   =\frac{t^2}{\gamma^2}\beta^{\mathsf T}\mathsf L\beta
       +O(\gamma^{-2}t^4+\gamma^{-3}t^2).
 \tag{V41.15}\label{r41:magnetic-response}
\]
Thus no estimate $|H_\gamma-h_\gamma-m|\le C_\gamma m$
can hold on a whole native neighborhood of coherence when
$\gamma$ is sufficiently large.
\end{proposition}
\begin{proof}
We give the finite matrix evaluation, including its normalization.
Put $R_{ij}=\Ad(U_{ij})|_{\mathbb R^3}$ and form the positive
$18$ by $18$ matrix
\[
 M_{ii}=6I_3,\qquad M_{ij}=-R_{ij},\qquad
 M_{ji}=-R_{ij}^{\mathsf T},\qquad G=M^{-1}.
\]
The block diagonal-dominance bound gives $3I\preceq M\preceq9I$.
In the chart $z_i=(\sqrt{1-|x_i|^2},x_i)$ the Haar density has
quadratic part $\tfrac12\sum_i|x_i|^2$.
The phase has no cubic term and quartic part
\[
             \frac68\sum_i|x_i|^4
                  -\frac14\sum_{i<j:\,i\sim j}|x_i|^2|x_j|^2 .
\]
Set $T_i=\operatorname{tr}G_{ii}$. Gaussian pairing gives
\[
 \begin{split}
 a_1(U)={}&\tfrac12\sum_iT_i
     -\tfrac68\sum_i\{T_i^2+2\operatorname{tr}(G_{ii}^2)\}\\
     &+\tfrac14\sum_{i<j:\,i\sim j}
          \{T_iT_j+2\operatorname{tr}(G_{ij}G_{ji})\}.
 \end{split}
 \tag{V41.16}\label{r41:magnetic-a1}
\]
This includes the full Haar term, not a determinant alone.

For a rational evaluation write
$M(t)=M_0+tM_1+t^2M_2+O(t^3)$, where
$M_0=\mathsf A\otimes I_3$. With
\[
 J=\begin{pmatrix}0&-1&0\\1&0&0\\0&0&0\end{pmatrix},
 \qquad \Pi=\operatorname{diag}(1,1,0),
\]
the upper off-diagonal edge blocks are
$(M_1)_{ij}=-2c_{ij}J$ and $(M_2)_{ij}=2c_{ij}^2\Pi$;
the lower blocks are their transposes. For $G_0=M_0^{-1}$,
\[
 G_1=-G_0M_1G_0,\qquad
 G_2=G_0M_1G_0M_1G_0-G_0M_2G_0 .
 \tag{V41.17}\label{r41:inverse-jets}
\]
The entries of $G_0$ have denominator dividing 37835.
Substitute $G_0+tG_1+t^2G_2$ into \eqref{r41:magnetic-a1}.
In the edge order $(01,03,12,14,25,34,45)$ the coefficient of
$t^2$ is $c^{\mathsf T}D^{\mathsf T}\mathsf LDc$, with
\[
 D=\begin{pmatrix}
 -1&1&0&-1&0&1&0\\
 0&0&-1&1&-1&0&1
 \end{pmatrix}.
 \tag{V41.18}\label{r41:magnetic-curl}
\]
These are finite rational identities: matrix inversion and
multiplication give the two numerators in \eqref{r41:L}.
The exact checker verifies every entry, the five-dimensional
gradient kernel, and rank two. It also differentiates the independent
product-normal-coordinate Haar and quartic-phase formula, obtaining
the same coefficient. The eigenvalues of $\mathsf L$ are its
diagonal entry plus and minus its off-diagonal entry, both positive.

Conjugation by a fixed quaternion sending $\mathbf e_3$ to
$-\mathbf e_3$ sends every $U_{ij}(t)$ to $U_{ij}(-t)$.
Thus these scalar functions are even in $t$.
Smoothness gives the $O(t^4)$ remainder in \eqref{r41:L}.
The differentiated Laplace remainder centered at zero is
$O(\gamma^{-3}t^2)$, by its $C^2$ bound and zero first derivative.
This proves \eqref{r41:magnetic-response}.
Choose $\beta\ne0$, take $\gamma$ sufficiently large, and then
$t\ne0$ sufficiently small. The response is strictly positive
although \eqref{r41:magnetic-min} is exactly zero, excluding
the proposed minimum-relative bound.
\end{proof}

This is a source-basis obstruction, not a massless-spectrum
argument. The full action already contains the transverse
plaquettes whose curvature is varied. Conversely the positive
$2$ by $2$ matrix is not a lower bound for the full $25$ by $25$
response or for a physical Hamiltonian. V40's electric constants
and \eqref{r41:L} test different native curvature components.

\subsection{A reflected full-cycle writer and an original-law bound}

For $\mathcal T_+$ let $G_+=\{|\xi^+|\le\rho\}$, with principal
logarithms defined there. Outside this set no logarithm is used.
Define the bounded exterior-measurable source
\[
 J_+^{(j)}(\eta)=
 \begin{cases}
 m(\eta)+j\gamma^{-2}\{a_1(\eta)-a_1(0)\}
   +\tfrac12\sum_{k,l}b_{\gamma,kl}^{(j)}
                                  \xi_k^+\cdot\xi_l^+,&G_+,\\
 H_\gamma(\eta)-h_\gamma,&G_+^c .
 \end{cases}
 \tag{V41.19}\label{r41:writer}
\]
The finite-coupling coefficients of \eqref{r41:response} are
evaluated only in the fixed coherent six-spin law, not in the
varying exterior. No unverified numerical fit is used.

Let $\mathcal T_-$ be the time-reflected lower-path tree, set
$J_-^{(j)}(\eta)=J_+^{(j)}(\theta\eta)$ with reflected coordinate
labels, and put
\[
                  \overline J^{(j)}
                    =\tfrac12(J_+^{(j)}+J_-^{(j)}).
 \tag{V41.20}\label{r41:reflected-writer}
\]
Each scalar is gauge invariant; the last is real and reflection
even. Unlike V39's intrinsic coefficients, quadratic polynomials
in different tree coordinates need not agree exactly on overlaps.
We have deliberately defined their reflection average. When one
chart is rough its term keeps the exact predictor; when both are
rough the average is $H_\gamma-h_\gamma$. This is a new explicit
source convention, not an assertion that the previous writers agree.

\begin{theorem}[Cubic native remainder in the physical source norm]
\label{r41:vacuum}
For $j=0,1$ and all sufficiently large $\gamma$, the unreflected
error is zero off $G_+$ and satisfies
\[
 |H_\gamma-h_\gamma-J_+^{(j)}|
       \le C\gamma^{-2-j}(\mathcal D^+)^{3/2}\mathbf1_{G_+}.
 \tag{V41.21}\label{r41:pointwise}
\]
For $A_\gamma,\ell_\gamma$ of \eqref{r37:budgets}, set
\[
 m_{3,\gamma}=A_\gamma^3+
       3A_\gamma^2/\ell_\gamma+
       6A_\gamma/\ell_\gamma^2+6/\ell_\gamma^3.
 \tag{V41.22}\label{r41:m3}
\]
Under the unchanged stationary slab law $\nu_\gamma$ of V39,
\[
 \|H_\gamma-h_\gamma-\overline J^{(j)}\|_{L^2(\nu_\gamma)}
      \le C\gamma^{-2-j}\sqrt{m_{3,\gamma}} .
 \tag{V41.23}\label{r41:L2}
\]
For a finite bank of the same middle blocks and deterministic
complex coefficients $c_b$, the physical source vectors satisfy
\[
 \begin{split}
 \|\psi_{\mathcal E_c}-\psi_{\overline J_c^{(j)}}\|
 &\le C\gamma^{-2-j}\sqrt{m_{3,\gamma}}\|c\|_1\\
 &=O\bigl((1+\log\gamma)^{3/2}
                           \gamma^{-7/2-j}\|c\|_1\bigr).
 \end{split}
 \tag{V41.24}\label{r41:physical}
\]
These vectors use the original physical Hilbert space and
transfer $P$, without a retained Markov quotient.
\end{theorem}
\begin{proof}
Theorem~\ref{r41:full-response} and
Lemma~\ref{r41:cycle-defect} prove \eqref{r41:pointwise};
the complementary source is exact by definition.
Use the same estimate after reflection and the $L^2$ triangle
inequality for \eqref{r41:reflected-writer}.

V35's moment formula with $r=3$ bounds the third moment of
each original defect by \eqref{r41:m3}. V39's total-variation
limit from periodic cylinders to the augmented vacuum slab
applies equally to the third moment: each defect is a bounded
function on the same finite product group. Thus both
$\mathcal D^+$ and $\mathcal D^-$ have this third-moment bound
under $\nu_\gamma$. Squaring \eqref{r41:pointwise} proves
\eqref{r41:L2}.

V39's conditional insertion identity sends the original block
energy to its exact predictor $H_\gamma$. The constant $h_\gamma$
has zero centered vacuum source since $B_1=P^2$ and
$P^2\Omega=\Omega$. Apply the unchanged Jensen inequality
\eqref{r39:vector-Jensen} to the residual sum and use the
$L^2$ triangle inequality over $b$. This gives
\eqref{r41:physical}, with no independence or separated
correlation decay assumed.
\end{proof}

The $j=0$ bound improves V39's $\gamma^{-3}(1+\log\gamma)$
rate by one defect half-power; $j=1$ similarly improves its
$\gamma^{-4}(1+\log\gamma)$ rate. More importantly for V40's
normalization question, every quadratic native cycle response
has been removed at finite coupling. On a specified smooth
background with $\xi=O(a^2)$, the local remainder is
$O(\gamma^{-2-j}a^6)$; multiplication by $a^{-4}$ leaves
$O(\gamma^{-2-j}a^2)$. This is only a smooth-background statement.
The proved quantum estimate is \eqref{r41:physical}, at its
displayed inverse-coupling scale. It still does not pay arbitrary
negative powers of $a$ along a continuum trajectory.

\subsection{Verification and the next physical obstruction}

The exact diagnostic is
\begin{center}\small
\path{scripts/verify_ym_restart_v41_full_response.py}.
\end{center}
It checks the native incidence ranks and Ward null directions,
three noncommuting 29-face magnetic identities, the stated
55-edge tree with its 25 chords, and three noncommuting
full-retained fixtures for the native filling and chord-defect
inequality. It evaluates the seven-edge magnetic quadratic form
by rational inverse jets, verifies its rank and curl factorization,
and independently checks its Haar-chart formula. A normalized
four-state fixture verifies the second-response cumulant formula
by a direct Taylor quotient; it is not Wilson quadrature.
Uniform analytic constants, finite-coupling Wilson integrals,
and a physical spectrum are not numerically certified.
The remainder and original-law statements are analytic arguments,
not proof-assistant formalizations.

The new writer gives a stronger absolute source comparison.
V39's all-time and spectral-window comparisons can use the
smaller vector error, but arbitrarily small nonzero spectral
weight can still lie near zero energy. This substitution proves
no positive lower bound on a physical source Gram, no exhaustive
sector coverage, and no connected correlation decay.

The next step must examine actual connected correlations of
these retained cycle sources and their rough-exterior predictors,
or a source-relative estimate at low physical energy.
Further local thermal coefficients do not supply that estimate.
The two magnetic cycles must remain in the source bank;
bounding every residual by the classical block minimum is now
explicitly excluded. The full interacting four-dimensional
continuum construction and physical mass gap remain open.

The complete V40 body is preserved. Removing this appendix and
restoring its title recovers that version exactly. Only the latest
active main source is retained, all other sources are preserved,
and build and check artifacts remain outside \path{latex_notes}.
% END CONSOLIDATED V41 APPEND

% BEGIN CONSOLIDATED V42 APPEND -- 2026-09-20
\clearpage
\section[V42: finite magnetic holonomies in the compact strand law]
{V42: Finite magnetic holonomies\\in the compact strand law}
\label{r42:start}

\subsection{What is new, and the transfer from the ESM notes}

We go upstream from V41's absolute source error to the actual
conditional normalizer. The two magnetic cycles have zero classical
partial minimum, but this does not make their integrated interaction
zero. We prove a global, finite-holonomy lower bound for that interaction
on the native shared-transport family. Its exact compact representation
has nonnegative loop factors. A lower bound valid at every positive
coupling has only a polynomial loss at large coupling; a separate
uniform Laplace argument removes that loss for this fixed six-spin
block. Neither statement is a bound on correlations between blocks.

The supplied \emph{Perturbative ESM Continuum Research Notes V23},
especially its derivative-floor theorem and first-background return
discussion, recommends testing the actual forbidden coefficients
before estimating them. We use that discipline, not its model-specific
coefficient or power of the mesh. In the present magnetic sector the
response is a curvature response, but its curvature-quadratic term
is \emph{not} forbidden. The bounds below prove this without computing
another inverse-coupling coefficient. They also identify a return that
cannot be suppressed: general left/right retained transports destroy
the nonnegative termwise representation, even arbitrarily close to
coherence at sufficiently high strand multiplicity.

Nonduplication is important here. Archive $C$
(\path{Renewal_Yang_Mills_Mass_Gap_Research_Notes_V89_f16.2.tex}),
V61--V62, already proves alignment domination, the exact compact
Haar-pairing identity, marked simple-cycle insertion, and a positive
coherent Hessian with a holonomy term. Its equations
\texttt{f16v62:diagram-weight}, \texttt{f16v62:insert-cycle} and
\texttt{f16v62:count-Jensen} are inputs below, not new results.
Its literal $\mathrm{Spin}(4)$ holonomy formula already distinguishes
left and right Wilson loops and their simultaneous center.
Archive $A$, V99 already supplies a Gaussian closed-walk normalizer
on a different packing. Archive f14 V63 already disproves a global
convexity route. We do not repeat any of these as progress.

The new passage is from the coherent Hessian to an actual
\emph{finite-holonomy} compact ratio, with a quantitative probability
of a marked cycle, a global equality set, a relative finite-coupling
comparison on the full compact magnetic family, and a literal
small-exterior sign obstruction to extending the positive expansion.
The strand method is standard; a primary-source reference is
B.~Lees and L.~Taggi, \emph{Exponential decay of transverse
correlations for $O(N)$ spin systems and related models},
\url{https://doi.org/10.1007/s00440-021-01053-5}, Sections 1.1 and 2.
Their transverse-correlation theorem concerns their specified
external-field spin model. It is not a Wilson retained-law
correlation theorem and is not imported here. This targeted check
does not replace the scheduled V50 broad review; the next regular
companion review remains V45.

\subsection{The same native compact integral and its strand law}

Use exactly the graph, pins and actual shared left/right links
of \eqref{r41:magnetic-energy}. Thus
\[
 V=\{0,1,2,3,4,5\},\quad
 E=\{01,03,12,14,25,34,45\},\quad
 n_i=6-\deg(i),\quad \sum_i n_i=22 .
 \tag{V42.1}\label{r42:graph}
\]
For an oriented edge set $U_{ji}=U_{ij}^{-1}$ and
\[
 T_{ij}=\operatorname{diag}(1,R_{ij}),\qquad
 R_{ij}=\operatorname{Ad}(U_{ij})|_{\mathbb R^3}.
\]
The spin is a real unit quaternion $z_i\in S^3$, and
\[
 \begin{split}
 Z_\gamma(U)
 &=\int\exp\left\{\gamma\sum_i n_i z_{i0}
             +\gamma\sum_{ij\in E}z_i\cdot T_{ij}z_j\right\}dH(z),\\
 Q_\gamma(U)&=e^{-29\gamma}Z_\gamma(U),\qquad
 F_\gamma(U)=-\log\frac{Q_\gamma(U)}{Q_\gamma(1)} .
 \end{split}
 \tag{V42.2}\label{r42:Z}
\]
All Haar measures are normalized. This $F_\gamma$ is the
coherence-subtracted \emph{action contribution} of the integrated
block, not V41's energy-source correction. The untouched Wilson
plaquettes still belong to the surrounding action.

For precision recall the inherited strand law in these conventions.
Let $m_e$ count labelled bond occurrences and $k_i$ count labelled
occurrences of the combined pin $n_i e_0$. Put
$N_i=\sum_{e\ni i}m_e+k_i$. Only even $N_i=2r_i$ contribute.
Pair all local occurrences using
\[
 \int_{S^3}z_{a_1}\cdots z_{a_{2r}}\,dH(z)
 =u(2r)\sum_{\text{pairings }\pi}\prod_{\{s,t\}\in\pi}\delta_{a_s a_t},
 \qquad
 u(2r)=\frac1{2^r(r+1)!}.
 \tag{V42.3}\label{r42:moment}
\]
The convention is $u(0)=1$; odd moments vanish.
A diagram $D$ consists of the multiplicities and these local pairings.
Its components are closed bond strands and open strands ending
at two pins, including a zero-bond pin pair. Write $L(D)$ for
its number of closed components. The coherent diagram probability is
\[
 \mathbb P_{\gamma,0}(D)=
 \frac1{Z_\gamma(1)}
 \left(\prod_e\frac{\gamma^{m_e}}{m_e!}\right)
 \left(\prod_i\frac{(\gamma n_i)^{k_i}}{k_i!}\right)
 \left(\prod_i u(N_i)\right)4^{L(D)} .
 \tag{V42.4}\label{r42:law}
\]
Combining the $n_i$ identical slots simply sums their multinomial
labels, so this is the same law used in archive $C$, not a new
reference measure on retained links.

For a closed strand $C$, let $U(C)$ be its ordered quaternion
transport, with inverse factors on reversed traversals. Choosing
another start conjugates this product; reversing orientation
inverts it. Both leave its real part squared unchanged.

\begin{theorem}[Exact nonnegative finite-holonomy factors]
\label{r42:positive}
For all actual shared-transport exteriors and every $\gamma>0$,
\[
 \frac{Q_\gamma(U)}{Q_\gamma(1)}
 =\mathbb E_{\gamma,0}
        \prod_{C\text{ closed in }D}\bigl(\operatorname{Re}U(C)\bigr)^2,
 \qquad 0<\frac{Q_\gamma(U)}{Q_\gamma(1)}\le1 .
 \tag{V42.5}\label{r42:ratio}
\]
The series is absolutely convergent. No small-link, small-coupling
or Gaussian assumption is used.
\end{theorem}
\begin{proof}
The inherited compact pairing expansion gives an open component
factor $e_0^{\mathsf T}T(P)e_0$ and a closed relative factor
$\operatorname{Tr}_{\mathbb R^4}T(C)/4$. Every shared transport
fixes $e_0$, so all open factors are exactly one. The exact
quaternion conjugation trace is
\[
 \operatorname{Tr}_{\mathbb R^4}\operatorname{Ad}(q)
       =1+\operatorname{tr}_{\mathbb R^3}\operatorname{Ad}(q)
       =4(\operatorname{Re}q)^2 .
\]
This proves the equality and puts every factor in $[0,1]$.
The absolute diagram sum is bounded by the coherent sum
$Z_\gamma(1)<\infty$. Equivalently one may first expand the
finite compact exponential in components; absolute component
series are dominated by an exponential with a finite multiple
of $\gamma$, which justifies integration and regrouping.
Finally $Q_\gamma(U)>0$ follows from its strictly positive
compact integrand. The upper bound follows from the expectation.
\end{proof}

This is an identity for a conditional normalizer at prescribed
exterior data. The shared-transport family has positive codimension
inside the general retained space. It has not been substituted
for that space in a probability integral.

\subsection{A cycle-occurrence estimate gives a nonlinear action floor}

Let $C_1=(0,3,4,1)$ and $C_2=(1,4,5,2)$ be the two elementary
unoriented squares, with cyclic closure understood. Set
\[
 W_p=U(C_p),\qquad
 \chi_p=1-(\operatorname{Re}W_p)^2,\qquad
 X(U)=\chi_1+\chi_2,\qquad D_\gamma=4+6\gamma .
 \tag{V42.6}\label{r42:flux}
\]
Thus $0\le\chi_p\le1$. These are adjoint flux defects: they
vanish at both $W_p=1$ and $W_p=-1$.

\begin{lemma}[A polynomial price for seeing a specified square]
\label{r42:probability}
Let $K_p$ count closed components traversing $C_p$ exactly once
under \eqref{r42:law}. Then
\[
 \begin{split}
 \mathbb E_{\gamma,0}K_p
 &=4\gamma^4\,\mathbb E_{\gamma,0}
                  \prod_{i\in C_p}(4+N_i)^{-1}
       \ge \frac{4\gamma^4}{D_\gamma^4},\\
 \mathbb P_{\gamma,0}\{K_p\ge1\}
 &\ge \frac{16\gamma^8}{D_\gamma^8(\gamma^2+\gamma)} .
 \end{split}
 \tag{V42.7}\label{r42:cycle-probability}
\]
\end{lemma}
\begin{proof}
The equality is archive $C$ V62's marked simple-cycle insertion:
adding one labelled occurrence on each of its four bonds cancels
the exponential factorials, adds a factor $(4+N_i)^{-1}$ at
each visited vertex, and adds one trace factor four. The bound
$\mathbb E N_i\le6\gamma$ and two applications of Jensen give
the first inequality, at the inherited scope.

Here is the additional probability step. For any edge $e$ on
$C_p$, distinct counted components use distinct occurrences, so
$K_p\le m_e$. Give that one bond an independent coupling $s$.
Differentiating the coherent positive series and its compact
integral, then setting $s=\gamma$, gives
\[
 \mathbb E m_e=\gamma\langle z_i\cdot z_j\rangle_0\le\gamma,\qquad
 \mathbb E[m_e(m_e-1)]
       =\gamma^2\langle(z_i\cdot z_j)^2\rangle_0\le\gamma^2 .
\]
Hence $\mathbb E K_p^2\le\gamma^2+\gamma$.
Cauchy--Schwarz applied to
$K_p\mathbf1_{\{K_p\ge1\}}$ gives
$\mathbb P(K_p\ge1)\ge(\mathbb E K_p)^2/\mathbb E K_p^2$.
This proves the second line. No independence of strand counts
or empty-background event is used.
\end{proof}

\begin{theorem}[Global finite-coupling magnetic coercivity]
\label{r42:coercivity}
For every $\gamma>0$ and every shared-transport exterior,
\[
 F_\gamma(U)\ge\kappa_\gamma X(U),\qquad
 \kappa_\gamma=
 \frac{8\gamma^8}{(4+6\gamma)^8(\gamma^2+\gamma)}>0 .
 \tag{V42.8}\label{r42:floor}
\]
Moreover
\[
 F_\gamma(U)=0
 \quad\Longleftrightarrow\quad
 W_1,W_2\in\{1,-1\}
 \quad\Longleftrightarrow\quad
 R_{ij}\text{ is a flat graph connection}.
 \tag{V42.9}\label{r42:equality}
\]
In particular the large-coupling loss in this explicit bound is
$\kappa_\gamma\sim8/(6^8\gamma^2)$, not exponential.
\end{theorem}
\begin{proof}
Write $q(D)$ for the product in \eqref{r42:ratio}. Its factors
are nonnegative. If $K_p\ge1$, at least one factor is
$1-\chi_p$, so $1-q(D)\ge\chi_p$ on this event. Consequently
\[
 1-q(D)\ge\frac12\sum_{p=1}^2
                  \chi_p\mathbf1_{\{K_p\ge1\}} .
\]
Take expectations, use \eqref{r42:cycle-probability}, and use
$-\log r\ge1-r$ for $0<r\le1$. This proves
\eqref{r42:floor}.

Its strict coefficient proves that equality forces both
$\chi_p=0$. Conversely tree-gauge the adjoint connection using
the five edges $01,12,03,14,25$. The two remaining edges
$34,45$ are the two square holonomies in this gauge. If both
are identity in $\mathrm{SO}(3)$, every transport can be
removed by vertex rotations fixing $e_0$. Haar invariance then
gives $Q_\gamma(U)=Q_\gamma(1)$. Equivalently every closed
quaternion holonomy is central, so every factor in
\eqref{r42:ratio} is one. This also proves both flatness
equivalences. The asymptotic coefficient is immediate.
\end{proof}

The center ambiguity is real, not a defect in this proof.
Simultaneous left/right central plaquettes can be invisible
to the incident conditional integral. Their untouched
fundamental Wilson action is not invisible and must be kept.
Nor does a lower action bound order normalized measures or
bound their connected correlations.

\subsection{The fixed-block weak-coupling floor does not collapse}

For this fixed graph a global compact comparison improves the
polynomial coefficient at large $\gamma$. Define the same
$18$ by $18$ matrix as in V41,
\[
 M(U)_{ii}=6I_3,\qquad
 M(U)_{ij}=-R_{ij},\qquad
 M(U)_{ji}=-R_{ij}^{\mathsf T},
 \qquad
 F_\infty(U)=\frac12\log\frac{\det M(U)}{\det M(1)} .
 \tag{V42.10}\label{r42:determinant}
\]
For all connections $3I\preceq M(U)\preceq9I$.

\begin{proposition}[Closed-walk bounds at finite holonomy]
\label{r42:walks}
Let $\mathcal W_k$ be the rooted oriented closed walks of length
$k$ on this six-vertex graph. Then
\[
 F_\infty(U)=\frac12\sum_{k\ge1}\frac1{k6^k}
       \sum_{w\in\mathcal W_k}\{3-\operatorname{tr}R(w)\},
 \qquad
 \frac1{324}X(U)\le F_\infty(U)\le24X(U).
 \tag{V42.11}\label{r42:walk-bound}
\]
Every summand is nonnegative and the series converges absolutely.
\end{proposition}
\begin{proof}
Write $M=6I-A_R$. The connection adjacency has norm at most
three, by the maximum scalar degree. Thus the logarithm series
converges with norm ratio at most $1/2$, and its trace is the
rooted closed-walk sum. Subtract the identity connection.
A rotation in $\mathrm{SO}(3)$ has trace at most three, so
all displayed deficits are nonnegative. At length four each
elementary square occurs eight times, from four starts and
two orientations; backtracking walks have identity holonomy.
Since $3-\operatorname{tr}R(C_p)=4\chi_p$, this term is
$X/324$. Keeping it proves the lower bound.

For the upper bound use the five-edge tree of the preceding
proof. Only two edge transports remain, say $A,B$, and
\[
 \|A-I\|_{\mathrm F}^2=8\chi_1,\qquad
 \|B-I\|_{\mathrm F}^2=8\chi_2 .
\]
Telescoping a product of at most $k$ such orthogonal factors,
including inverses, gives
$\|R(w)-I\|_{\mathrm F}\le k\sqrt{8X}$.
Therefore $3-\operatorname{tr}R(w)\le4k^2X$.
There are at most $6\cdot3^k$ rooted walks. Inserting these
bounds in the convergent series gives
$12X\sum_{k\ge1}k2^{-k}=24X$.
\end{proof}

\begin{theorem}[Uniform relative comparison on the entire magnetic family]
\label{r42:relative}
There are finite constants $C,\gamma_0$, depending only on this
fixed six-spin block, such that for all connections $U$ and
all $\gamma\ge\gamma_0$,
\[
 |F_\gamma(U)-F_\infty(U)|\le\frac C\gamma X(U).
 \tag{V42.12}\label{r42:relative-bound}
\]
After increasing $\gamma_0$ if necessary,
\[
                  \frac1{648}X(U)\le F_\gamma(U)\le25X(U).
 \tag{V42.13}\label{r42:uniform-floor}
\]
The threshold is not evaluated numerically. No volume-uniform
claim for growing free blocks is made.
\end{theorem}
\begin{proof}
For all these exteriors the unique minimum is $z_i=1$,
and its energy is zero. More strongly,
\[
 E_U(z)\ge\sum_i n_i(1-z_{i0})\ge3\sum_i(1-z_{i0}),
\]
which supplies a uniform barrier outside any fixed sufficiently
small neighborhood of that point. The Hessian bounds above
are uniform on the full compact connection space.

In the positive scalar-part chart
$z_i=(\sqrt{1-|x_i|^2},x_i)$, the quadratic phase is
$x^{\mathsf T}M(U)x/2$. The remaining phase starts at degree
four and the Haar density has constant term
$(2\pi^2)^{-6}$. The parameter-dependent Laplace argument of
V38, now on a finite atlas of the compact connection space,
therefore gives
\[
 Q_\gamma(U)=\frac8{\pi^3}\gamma^{-9}\det M(U)^{-1/2}
                        \{1+\varepsilon_\gamma(U)\},
 \qquad
 \|\varepsilon_\gamma\|_{C^2}\le C\gamma^{-1}.
 \tag{V42.14}\label{r42:Laplace}
\]
Here the $C^2$ norm is on a fixed finite atlas, not on a growing
graph. To justify the uniformity explicitly, the minimum is
fixed, all phase and density derivatives have compact uniform
bounds, and $M\ge3I$. Rescaling $x=\gamma^{-1/2}y$ makes the
quartic phase and quadratic density errors integrable
$O(\gamma^{-1})$ Gaussian moments. Two parameter derivatives
introduce only fixed polynomial moments of $y$. A smaller
chart gives a common Gaussian majorant; the complementary
integral and its parameter derivatives are a polynomial in
$\gamma$ times the uniform exponential barrier. This proves
the stated differentiated remainder. The constant is the
$18$-dimensional Gaussian factor divided by the six normalized
sphere areas.

Tree gauge reduces the adjoint connection to
$(A,B)\in\mathrm{SO}(3)^2$. Both $F_\gamma$ and $F_\infty$
are invariant under this reduction and simultaneous
conjugation. Taking the ratio with $U=1$ and its logarithm in
\eqref{r42:Laplace} gives
$R_\gamma:=F_\gamma-F_\infty$ with
$\|R_\gamma\|_{C^2}\le C/\gamma$ and $R_\gamma(I,I)=0$.
In logarithmic coordinates at $(I,I)$, simultaneous
conjugation rotates the two three-vectors. It has no invariant
linear functional, so $DR_\gamma(I,I)=0$. Taylor's formula
then bounds $|R_\gamma|$ by $C/\gamma$ times the squared
coordinate norm, which is comparable to $X$ near $(I,I)$.
Outside that neighborhood $X$ has a positive lower bound on
the compact space; the $C^0$ estimate gives the same inequality.
This proves \eqref{r42:relative-bound} globally, including
the central lifts that are all the same adjoint flat point.

Choose $\gamma_0$ so that $C/\gamma_0\le1/648$ and
$C/\gamma_0\le1$. Combine with \eqref{r42:walk-bound} to obtain
\eqref{r42:uniform-floor}.
\end{proof}

This comparison is relative to the actual magnetic defect, not
an additive $O(\gamma^{-1})$ error which could dominate an
arbitrarily small flux. In particular, if
$U_e(t)=\exp(tc_e\mathbf e_3)$ with V41's two curls $\beta_p$,
then
\[
 X(U(t))=t^2(\beta_1^2+\beta_2^2)+O(t^4).
 \tag{V42.15}\label{r42:quadratic}
\]
Thus the magnetic action's quadratic term cannot be removed by
a proposed extra symmetry cancellation. Gauge-gradient directions
still have zero curl, as required. This confirms the existing
coherent-holonomy floor rather than inventing a new Ward identity.

Also, directly from the actual normalized integral,
\[
            \partial_\gamma F_\gamma(U)=H_\gamma(U)-h_\gamma .
 \tag{V42.16}\label{r42:source}
\]
The leading determinant cost is independent of $\gamma$ and
therefore does not itself produce an order-one energy-source
response. V41 already evaluates the first nonzero source term.
We do not infer monotonicity in $\gamma$ for all $U$ from
positivity of $F_\gamma$.

\subsection{The positive strand mechanism does not survive general returns}

\begin{proposition}[Negative native strand terms arbitrarily near coherence]
\label{r42:sign}
For general actual retained exteriors of this same six-link
block, the compact pairing expansion need not have nonnegative
relative factors. There are such exteriors arbitrarily close
to the all-identity exterior with a strictly negative closed
strand factor.
\end{proposition}
\begin{proof}
An internal edge with independent retained left and right
links $\ell,r$ acts on a quaternion by
$T(q)=\ell q r^{-1}$. Its four-dimensional trace is
$4\operatorname{Re}\ell\,\operatorname{Re}r$, as proved for
literal Wilson words in archive $C$ V62.

Choose the actual left transverse retained link belonging to
edge $34$, set it to
$u_w=\exp(\pi\mathbf e_1/w)$ for an integer $w\ge2$, and leave
its right partner and all other retained links identity.
The 22 private slots stay identity; this link belongs to an
internal, not a private, face in the incident energy.
A closed strand traversing $C_1$ exactly $w$ times has transport
$q\mapsto u_w^wq=-q$. Its relative trace is $-1$.

Such a diagram really occurs: use $w$ labelled copies of every
edge of $C_1$, zero other multiplicities, and pair consecutive
edges at its vertices. At one chosen vertex connect the end
of copy $j$ to the start of copy $j+1$ cyclically, so there is
one component winding $w$ times, not $w$ separate components.
All local degrees are $2w$, all Haar pairing coefficients and
bond factorial weights are strictly positive, and its only
closed relative factor is $-1$. Finally $u_w\to1$ as
$w\to\infty$, proving the small-exterior assertion.
\end{proof}

The negative term is not a negative Wilson probability:
the full compact integral is strictly positive, and physical
reflection positivity is unchanged. Nor does this rule out
a different regrouping with positive weights. It rules out
using \eqref{r42:ratio}'s termwise nonnegative domination on
the whole retained neighborhood without further work.
A fixed low-degree expansion can miss this obstruction, since
the required strand multiplicity grows as the link angle shrinks.
The ESM warning about omitted returns therefore has a concrete
native counterpart here.

\subsection{Verification, interpretation and the next estimate}

The reproducible diagnostic is
\begin{center}\small
\path{scripts/verify_ym_restart_v42_magnetic_strands.py}.
\end{center}
Exact diagnostics verify sphere moments, quaternion traces, literal
native Wilson words, matrix/walk coefficients, central blindness and
a negative strand fixture. Probability and Laplace estimates are
analytic proofs, not numerical integrations or spectral certificates.

The new nonlinear magnetic action floor is local and center-blind;
it does not define autonomous retained dynamics. The next bound must
control signed returns or a justified positive regrouping for independent
left/right transports and private slots, including rough-exterior
predictors, under the original retained Wilson law. Local bounds
cannot be multiplied as though blocks were independent.

The interacting four-dimensional continuum theory, cutoff-stable
nonzero physical source norms, connected decay at physical separation
and a positive physical Hamiltonian mass gap remain unproved.
The goal remains active.

V41 is recovered exactly by removing this appendix and restoring its
title. Only the newest active main source is kept; archives and companion
sources are unchanged. All build/check artifacts remain outside
\path{latex_notes}.
% END CONSOLIDATED V42 APPEND

% BEGIN CONSOLIDATED V43 APPEND -- 2026-09-20
\clearpage
\section[V43: long-strand tails and normalized signed returns]
{V43: Long-strand tails\\and normalized signed returns}
\label{r43:start}

\subsection{From the sign obstruction to a quantitative return estimate}

V42 found negative strand terms at exteriors arbitrarily close to
coherence. Their multiplicity, however, grows as the exterior angle
shrinks. This section bounds that effect rather than imposing
nonnegativity on each term. For the actual six-spin conditional,
an anisotropic compact integral controls long closed strands with
a coupling-independent constant and long open strands with an
$O(\gamma)$ constant. The resulting estimate pays the actual
normalizing denominator and its coupling derivative. It is not
another local Taylor coefficient.

On a native defect threshold of order $\gamma^{-2/3}$, replacing
signed diagram factors by their absolute values changes both the
block log normalizer and its conditional energy source by at most
$C(1+\gamma)e^{-\gamma^{1/3}}$. Keeping the exact predictor outside
that threshold gives the same uniform error in the original
physical vacuum source norm. General private slots and independent
left/right transports are included. No shared-transport hypothesis
is made in this comparison.

The Haar-pairing representation is inherited from archive $C$
V61--V62 and V42. Colour switching is an established method:
B.~Lees and L.~Taggi, \emph{Exponential decay of transverse
correlations for $O(N)$ spin systems and related models},
\url{https://doi.org/10.1007/s00440-021-01053-5},
Section 2.2, Lemma 2.5. We prove below the weighted, finite-block
identity needed here, including coincident endpoints. We do not
import their infinite-graph decay conclusion as a theorem for
the retained Wilson law. Archive $C$ V72's long physical
Wilson-path singlet calculation concerns a different object,
not the internal pairing strands estimated here.

The new estimate removes a particular local signed-return
obstruction. It does not estimate connected correlations of the
resulting positive diagram interaction, prove a positive source
Gram, or reach the nonperturbative physical mass scale.
The next companion and broad literature checkpoints remain
V45 and V50.

\subsection{An anisotropic compact integral controls component lengths}

Keep the graph and pins \eqref{r42:graph}, and write
$z_i=(z_{i0},x_i)\in S^3$, $x_i\in\mathbb R^3$.
For $1\le s\le5/4$ define the auxiliary integral
\[
 \begin{split}
 Z_{\gamma,s}
 &=\int\exp\left\{\gamma\sum_i n_i z_{i0}
       +\gamma\sum_{ij\in E}
                    (z_{i0}z_{j0}+s\,x_i\cdot x_j)\right\}dH,\\
 Q_{\gamma,s}&=e^{-29\gamma}Z_{\gamma,s},\qquad
 E_s=-\gamma^{-1}\log(\text{integrand of }Q_{\gamma,s}).
 \end{split}
 \tag{V43.1}\label{r43:anisotropic}
\]
Thus $Z_{\gamma,1}=Z_\gamma(1)$ is the old coherent integral.
The auxiliary anisotropy is only a generating device; it
does not replace a Wilson action or a physical evolution.

In V42's coherent uncoloured diagram law
$\mathbb P_{\gamma,0}$, let $\mathcal C(D)$ and $\mathcal O(D)$
be the closed and open components. Their length is the number
of graph-bond occurrences, excluding the two pin occurrences
of an open component. A zero-bond pin pair has length zero.
Set
\[
 G_s(D)=\prod_{C\in\mathcal C(D)}
                      \frac{1+3s^{|C|}}4,\qquad
 Y(z)=\sum_i n_i z_i^1 .
 \tag{V43.2}\label{r43:Gs}
\]

\begin{lemma}[The two generating identities, with the endpoint factor]
\label{r43:generators}
For every $\gamma>0$ and $1\le s\le5/4$,
\[
 \begin{split}
 \mathbb E_{\gamma,0}G_s
     &=\frac{Z_{\gamma,s}}{Z_{\gamma,1}},\\
 \mathbb E_{\gamma,0}
       \left[G_s\sum_{P\in\mathcal O(D)}s^{|P|}\right]
     &=\frac{\gamma^2}{2}
             \frac{Z_{\gamma,s}}{Z_{\gamma,1}}
                              \langle Y^2\rangle_{\gamma,s}.
 \end{split}
 \tag{V43.3}\label{r43:generator-identities}
\]
Each open component is counted once, not once per orientation.
\end{lemma}
\begin{proof}
Expand the compact integral in its four coordinate colours.
Pins have colour zero. An open component must therefore have
colour zero along its whole path. A closed component can have
colour zero, with weight one, or one of the three transverse
colours, with weight $s^{|C|}$. Summing its colours changes the
old trace factor four to $1+3s^{|C|}$. This proves the first
identity using the same labelled-occurrence factorials as
\eqref{r42:law}.

For the second identity differentiate twice at $t=0$ the
integral with additional source $tY$. A surviving diagram
contains two labelled transverse source endpoints joined by
one open strand of colour one. It cannot end at an ordinary
colour-zero pin. Recolour this entire marked path zero and
replace its two source endpoints by ordinary pins. The bond
weights lose exactly $s^{|P|}$, and the endpoint coefficients
gain $\gamma^2$. All other open and closed components are
unchanged. Insertion positions among repeated pin labels
cancel their exponential factorial ratios, just as in the
inherited component-insertion proof.

Conversely a chosen open component has two ways to receive
the two differentiated source labels. This remains true for
coincident endpoint vertices: its two pin occurrences are
distinct labelled objects. Thus differentiation gives
$2\gamma^{-2}$ times the marked-open-component sum.
The source derivative is $Z_{\gamma,s}\langle Y^2\rangle_{\gamma,s}$,
which proves the formula and its factor $1/2$.
All regroupings are justified by the absolutely convergent
coordinate expansion of a compact exponential integral,
also after its two source derivatives.
\end{proof}

\begin{lemma}[Uniform compact integral stocks for this fixed block]
\label{r43:stocks}
There are finite constants $K_0,K_1$, independent of $\gamma\ge1$
and $1\le s\le5/4$, such that
\[
 \frac{Z_{\gamma,s}}{Z_{\gamma,1}}\le K_0,\qquad
 \frac{\gamma^2}{2}\frac{Z_{\gamma,s}}{Z_{\gamma,1}}
                  \langle Y^2\rangle_{\gamma,s}\le K_1\gamma .
 \tag{V43.4}\label{r43:stock-bounds}
\]
These constants concern six spins, not an increasing free graph.
\end{lemma}
\begin{proof}
Put $S(z)=\sum_i(1-z_{i0})$. Since $|x_i|^2\le2(1-z_{i0})$,
\[
 \sum_{ij\in E}x_i\cdot x_j
       \le\frac12\sum_i\deg(i)|x_i|^2
       \le\sum_i\deg(i)(1-z_{i0}).
\]
The coherent bond energies are nonnegative. Consequently
\[
 E_s\ge\sum_i\{n_i-(s-1)\deg(i)\}(1-z_{i0})
                                  \ge\frac94 S(z).
 \tag{V43.5}\label{r43:barrier}
\]
This is a global compact inequality, including negative scalar
spin components.

Here are elementary constants if desired. Set $a_*=9/4$,
\[
 A_0=\frac{\pi^{5/2}}{2(2a_*)^{3/2}},\quad
 A_1=\frac{3\pi^{9/2}}{4(2a_*)^{5/2}},\quad
 c_0=e^{-27}\left(\frac8{3\pi^3}\right)^6 .
\]
The normalized radial Haar density is $(2/\pi)\sin^2\theta\,d\theta$.
Use $\sin\theta\le\theta$ and
$1-\cos\theta\ge2\theta^2/\pi^2$ for $0\le\theta\le\pi$.
Extending the Gaussian radial integral to infinity gives
\[
 \int e^{-a_*\gamma(1-z_0)}dH\le A_0\gamma^{-3/2},\qquad
 \int |x|^2e^{-a_*\gamma(1-z_0)}dH\le A_1\gamma^{-5/2}.
\]
Thus $Q_{\gamma,s}\le A_0^6\gamma^{-9}$ and
\[
 \int\sum_i|x_i|^2e^{-\gamma E_s}dH
                              \le6A_1A_0^5\gamma^{-10}.
\]

For the denominator restrict each spin to
$\theta_i\le\gamma^{-1/2}\le1$. The triangle inequality on
$S^3$ and $1-\cos t\le t^2/2$ give
\[
 E_1\le\sum_i(n_i/2+\deg(i))\theta_i^2\le27/\gamma .
\]
Each cap has Haar mass at least
$[8/(3\pi^3)]\gamma^{-3/2}$, using
$\sin\theta\ge2\theta/\pi$ on this cap. Therefore
$Q_{\gamma,1}\ge c_0\gamma^{-9}$.
Finally $\sum_i n_i^2=82$, so $Y^2\le82\sum_i|x_i|^2$.
We may take $K_0=A_0^6/c_0$ and
$K_1=246A_1A_0^5/c_0$. Their large numerical size is not
optimized or hidden as a smallness assertion.
\end{proof}

\begin{theorem}[Exponential tails of internal strand lengths]
\label{r43:tails}
Fix $s=5/4$. Uniformly for $\gamma\ge1$ and integers $L\ge1$,
\[
 \begin{split}
 \mathbb E_{\gamma,0}\#\{C\in\mathcal C:|C|\ge L\}
                                      &\le C_0s^{-L},\\
 \mathbb E_{\gamma,0}\#\{P\in\mathcal O:|P|\ge L\}
                                      &\le C_1\gamma s^{-L}.
 \end{split}
 \tag{V43.6}\label{r43:tail-bound}
\]
In particular the probability of any component of length at
least $L$ is at most $C(1+\gamma)s^{-L}$.
\end{theorem}
\begin{proof}
Since every factor of $G_s$ is at least one,
\[
 G_s-1\ge\frac34\sum_{C\in\mathcal C}(s^{|C|}-1).
\]
The first generator and \eqref{r43:stock-bounds} bound the
expectation of this sum by $4(K_0-1)/3$.
For $|C|\ge L\ge1$,
$s^{|C|}-1\ge(1-s^{-1})s^L$.
This proves the first line. For open components use $G_s\ge1$
in the second generator to get
$\mathbb E\sum_Ps^{|P|}\le K_1\gamma$.
Keeping only lengths at least $L$ proves the second line.
The probability assertion follows by bounding an indicator
by the total number of long components.
\end{proof}

These are tails under the coherent internal diagram law.
They are not spatial clustering estimates under the interacting
retained measure. A long internal strand may repeatedly traverse
the same short square, exactly as in V42's sign example.

\subsection{General exterior data and the normalizing denominator}

Now allow every actual private slot $p_{i\alpha}\in S^3$,
$\alpha=1,\ldots,n_i$, and every internal transport
$T_{ij}\in\mathrm{SO}(4)$, including independent left/right
Wilson links. Define
\[
 Z_\gamma(p,T)=\int
 e^{\gamma I_{p,T}(z)}dH,\qquad
 I_{p,T}=\sum_{i,\alpha}p_{i\alpha}\cdot z_i
                    +\sum_{ij\in E}z_i\cdot T_{ij}z_j .
 \tag{V43.7}\label{r43:general}
\]
The original conditional energy is $29-I_{p,T}$.
The inherited uncoloured expansion is
\[
 R_\gamma:=\frac{Z_\gamma(p,T)}{Z_{\gamma,1}}
     =\mathbb E_{\gamma,0}\Phi_{p,T}(D),
 \quad
 \Phi_{p,T}=
 \prod_{P\in\mathcal O}
    p_{i(P)\alpha(P)}^{\mathsf T}T(P)p_{j(P)\beta(P)}
 \prod_{C\in\mathcal C}\frac{\operatorname{Tr}T(C)}4 .
 \tag{V43.8}\label{r43:Phi}
\]
For distinct slot vectors keep their individual labels in
the coherent law; summing them recovers \eqref{r42:law}.
Their introduction changes no strand-length estimate above.
All component factors lie in $[-1,1]$, so $|\Phi|\le1$.
Define the positive absolute-diagram normalizer
\[
 A_\gamma=\mathbb E_{\gamma,0}|\Phi_{p,T}|,\qquad
 Z_\gamma^{\rm abs}=Z_{\gamma,1}A_\gamma .
 \tag{V43.9}\label{r43:absolute}
\]
It is a full convergent series, not a fixed-degree truncation.
Since $Z_\gamma(p,T)>0$,
$0<R_\gamma\le A_\gamma\le1$.
Open inner products and closed traces are unchanged by
consistent vertex frame rotations; so is their absolute
product. Thus $Z_\gamma^{\rm abs}$ is an invariant scalar
under the same free-spin frame changes. It is not assigned
reflection positivity or an autonomous transfer operator.

Measure closeness in a specified coherent frame by
\[
 \begin{split}
 \epsilon&=\max\left\{\max_{i,\alpha}|p_{i\alpha}-e_0|,
                         \max_{ij}\|T_{ij}-I\|_{\rm op}\right\},\\
 \delta^2&=\sum_{i,\alpha}|p_{i\alpha}-e_0|^2
                         +\sum_{ij}\|T_{ij}-I\|_{\rm op}^2 .
 \end{split}
 \tag{V43.10}\label{r43:distances}
\]

\begin{lemma}[A quadratic, normalized lower bound]
\label{r43:Jensen}
For arbitrary unit slots and orthogonal transports,
\[
                       R_\gamma\ge e^{-\gamma\delta^2/2}.
 \tag{V43.11}\label{r43:denominator}
\]
\end{lemma}
\begin{proof}
Use the coherent spin law, not the diagram law, to write
$R_\gamma=\langle e^{\gamma(I_{p,T}-I_{e_0,I})}\rangle_0$.
By transverse rotation invariance
$\langle z_i\rangle_0=m_i e_0$ and
$\langle z_jz_i^{\mathsf T}\rangle_0
=\operatorname{diag}(a_{ij},b_{ij},b_{ij},b_{ij})=:C_{ij}$.
All $m_i,a_{ij},b_{ij}$ are nonnegative. One direct verification
expands the coherent positive-coordinate monomials: every
nonzero product-sphere moment in each of these numerators
is nonnegative. Also $m_i\le1$ and
$\operatorname{tr}C_{ij}=\langle z_i\cdot z_j\rangle_0\le1$.
Hence $C_{ij}\succeq0$.

A pin term has mean change
$-m_i|p_{i\alpha}-e_0|^2/2$.
Since $C_{ij}$ is symmetric, orthogonality gives
\[
 \operatorname{tr}((T-I)C_{ij})
   =-\tfrac12\operatorname{tr}((T-I)^{\mathsf T}(T-I)C_{ij})
   \ge-\tfrac12\|T-I\|_{\rm op}^2 .
\]
Sum these bounds and apply Jensen's inequality to the
exponential. This proves \eqref{r43:denominator}.
\end{proof}

\begin{lemma}[Negative factors require a long component]
\label{r43:negative}
For $0<\epsilon\le1/8$, put
$L_\epsilon=\lfloor(2\epsilon)^{-1}\rfloor$ and
$b=\tfrac12\log(5/4)$. Then
\[
 p_-:=\mathbb P_{\gamma,0}\{\Phi_{p,T}<0\}
        \le C(1+\gamma)e^{-b/\epsilon},\qquad \gamma\ge1 .
 \tag{V43.12}\label{r43:negative-tail}
\]
\end{lemma}
\begin{proof}
For a path of length $\ell$,
$\|T(P)-I\|_{\rm op}\le\ell\epsilon$ by telescoping orthogonal
products, also for reversed traversals. Thus
\[
 |p_i-T(P)p_j|\le(2+\ell)\epsilon .
\]
If $\ell<L_\epsilon$, this is at most $2\epsilon+1/2\le3/4$.
Its open inner product, equal to one minus half this squared
distance, is strictly positive.
For a closed component of that length,
\[
 \frac{\operatorname{Tr}T(C)}4
 =1-\frac18\|T(C)-I\|_{\mathrm F}^2
 \ge1-\frac12\ell^2\epsilon^2>0 .
\]
Therefore a negative product has a component of length at
least $L_\epsilon$. Apply \eqref{r43:tail-bound} and
$s^{-L_\epsilon}\le s\,e^{-b/\epsilon}$.
At $\epsilon=0$ every factor is one and no negative term occurs.
\end{proof}

\subsection{The coupling derivative is controlled as well}

Let
\[
 \mathcal L_\gamma=\log\frac{Z_\gamma^{\rm abs}}{Z_\gamma(p,T)}
                         =\log(A_\gamma/R_\gamma)\ge0,\quad
 H_\gamma=29-\partial_\gamma\log Z_\gamma,\quad
 H_\gamma^{\rm abs}=29-\partial_\gamma\log Z_\gamma^{\rm abs}.
 \tag{V43.13}\label{r43:L}
\]
All derivatives here keep the exterior and its frame fixed.
The absolute values are independent of $\gamma$.

\begin{theorem}[Normalized signed-return and source bounds]
\label{r43:normalized}
For $\gamma\ge1$ and $0<\epsilon\le1/8$,
\[
 \begin{split}
 0\le\mathcal L_\gamma
   &\le C(1+\gamma)
             \exp\{\gamma\delta^2/2-b/\epsilon\},\\
 |H_\gamma-H_\gamma^{\rm abs}|
     =|\partial_\gamma\mathcal L_\gamma|
   &\le C(1+\gamma)
             \exp\{\gamma\delta^2-b/(2\epsilon)\}.
 \end{split}
 \tag{V43.14}\label{r43:normalized-bounds}
\]
Both differences vanish at $\epsilon=0$.
\end{theorem}
\begin{proof}
Let $M(D)$ be the total number of bond and pin occurrences,
and put $\Phi_-=\max\{-\Phi,0\}$,
$N=\mathbb E_0\Phi_-$ and $N_1=\mathbb E_0(M\Phi_-)$.
Then $A-R=2N$ and $N\le p_-$, so
\[
 0\le\log(A/R)=\log(1+2N/R)\le2N/R .
\]
The denominator lemma and \eqref{r43:negative-tail} give the
first bound.

Here the derivative of the normalizer cannot be discarded.
The coherent weights before normalization are proportional
to $\gamma^M$. Compact differentiation gives
\[
 \mathbb E_0M=\gamma\langle I_{e_0,I}\rangle_0\le29\gamma,\qquad
 \mathbb E_0M^2
 =\gamma^2\langle I_{e_0,I}^2\rangle_0
           +\gamma\langle I_{e_0,I}\rangle_0
 \le841\gamma^2+29\gamma .
\]
Define $A_1=\mathbb E_0(M|\Phi|)$ and
$R_1=\mathbb E_0(M\Phi)$ in this paragraph only.
Termwise differentiation of the two entire positive-majorized
series yields the exact identity
\[
 \partial_\gamma\mathcal L_\gamma
       =\frac{A_1}{\gamma A}-\frac{R_1}{\gamma R}
       =\frac{2(N_1R-NR_1)}{\gamma AR}.
 \tag{V43.15}\label{r43:derivative}
\]
The coherent normalizer derivative cancels between the two
logarithms; it was not omitted from either one.
Cauchy--Schwarz gives
$N_1\le(841\gamma^2+29\gamma)^{1/2}p_-^{1/2}$.
Also $|R_1|\le A_1\le29\gamma$ and $A\ge R$.
Consequently
\[
 |\partial_\gamma\mathcal L_\gamma|
 \le 2e^{\gamma\delta^2/2}N_1/\gamma
                         +58e^{\gamma\delta^2}N .
\]
Insert \eqref{r43:negative-tail} and enlarge the fixed constant.
The displayed second estimate bounds both terms. The equality
with the energy-source difference follows from \eqref{r43:L}.
Absolute convergence with one or two occurrence marks follows
from the same compact exponential majorant.
\end{proof}

The estimate is genuinely normalized. A small negative
unnormalized weight alone would not give it, since $R_\gamma$
can shrink with the exterior. The explicit quadratic
denominator cost is what permits the next native threshold.

\subsection{The original defect pays the return estimate}

Use the actual upper-path retained tree of V41, without
restricting to the shared magnetic family. Gauge its links
to identity and form the 22 private slot vectors and seven
$\mathrm{SO}(4)$ transports in that frame.

\begin{lemma}[Global coefficient distances paid by the native defect]
\label{r43:defect}
With the distances \eqref{r43:distances} in this tree frame,
\[
                      \epsilon^2\le\delta^2\le2\mathcal D^+ .
 \tag{V43.16}\label{r43:defect-bound}
\]
No principal logarithm or smallness assumption is needed.
\end{lemma}
\begin{proof}
The six upper reference slots are identity; the other 16 are
outer quaternion chords, possibly inverted. Their squared
distances sum to the squared distances of those chords.
An internal edge acts by $T(q)=lqr^{-1}$, with $l,r$ its
left and right tree-gauged retained links. Orthogonality of
quaternion multiplication gives
\[
 \|T-I\|_{\rm op}\le|l-1|+|r-1|,\qquad
 \|T-I\|_{\rm op}^2\le2(|l-1|^2+|r-1|^2).
\]
Among left grid links only two are chords; all seven right
links are chords. Each occurs in exactly one of the seven
internal transports. Thus $\delta^2$ is at most twice the
sum of all 25 squared chord distances. The global inequality
\eqref{r41:chord-bound} bounds that sum by $\mathcal D^+$.
The maximum squared distance is no larger than their sum.
\end{proof}

The chosen retained tree is globally defined by path products.
Its residual root gauge acts by simultaneous conjugation,
fixing $e_0$ and preserving both distance norms. In particular
the good event and scalar construction below are native
gauge-invariant functions, not a change of quotient measure.

Set
\[
 d_\gamma=\frac{\gamma^{-2/3}}{2048},\qquad
 G_\gamma^+=\{\mathcal D^+\le d_\gamma\},\qquad
 e_\gamma=C(1+\gamma)e^{-\gamma^{1/3}} .
 \tag{V43.17}\label{r43:threshold}
\]
On $G_\gamma^+$ both $\epsilon^2$ and $\delta^2$ are at most
$\gamma^{-2/3}/1024$. Since
$b=\log(5/4)/2\ge1/10$,
\[
 \gamma\delta^2-\frac b{2\epsilon}
 \le\left(\frac1{1024}-\frac85\right)\gamma^{1/3}
                                                  <-\gamma^{1/3}.
\]
At zero distance use the zero-error convention.
Both bounds in \eqref{r43:normalized-bounds} are therefore
at most $e_\gamma$ on this event.

Define the retained source at a fixed coupling by
\[
 W_\gamma^+(\eta)=
 \begin{cases}
 H_\gamma^{\rm abs}(\eta)-h_\gamma,&\eta\in G_\gamma^+,\\
 H_\gamma(\eta)-h_\gamma,&\eta\notin G_\gamma^+,
 \end{cases}
 \quad
 W_\gamma^-=W_\gamma^+\circ\theta,\quad
 \overline W_\gamma=\tfrac12(W_\gamma^++W_\gamma^-).
 \tag{V43.18}\label{r43:writer}
\]
The lower construction uses reflected tree and slot labels.
The scalar $Z^{\rm abs}$ itself is frame invariant.
The cutoffs need not agree; their reflection average is
deliberate. This is a source definition at fixed $\gamma$:
no derivative of the moving indicator in
\eqref{r43:threshold} is asserted to vanish.

\begin{theorem}[Uniform source error in the original physical space]
\label{r43:physical}
For every exterior,
\[
 |H_\gamma-h_\gamma-W_\gamma^+|
                \le e_\gamma\mathbf1_{G_\gamma^+},\qquad
 |H_\gamma-h_\gamma-\overline W_\gamma|\le e_\gamma .
 \tag{V43.19}\label{r43:pointwise}
\]
For any finite bank of the original middle blocks and
deterministic complex coefficients $c_b$,
\[
       \|\psi_{\mathcal E_c}-\psi_{\overline W_{\gamma,c}}\|
                          \le e_\gamma\|c\|_1 .
 \tag{V43.20}\label{r43:physical-bound}
\]
The physical Hilbert space, vacuum, transfer and source
centering are exactly those of V39.
\end{theorem}
\begin{proof}
The good-set estimate was just proved; the rough-set
predictor is exact by definition. The same estimate holds
after reflection, proving the average bound.
The source is bounded at fixed coupling: on the good set
$H^{\rm abs}$ differs by at most $e_\gamma$ from the original
$0\le H_\gamma\le58$, and on the complement it equals that
original predictor. V39's exact conditional insertion
therefore applies. Its residual for a coefficient bank is
bounded by $e_\gamma\|c\|_1$ on the entire augmented slab.
Apply \eqref{r39:vector-Jensen}; subtracting the original
vacuum expectation cannot increase the $L^2$ variance.
The constant $h_\gamma$ has zero centered vacuum source
because $B_1=P^2$ and $P^2\Omega=\Omega$.
No independence or correlation decay was used.
\end{proof}

There is also a precise original-law size for the rough event.
On V35's cofinal dyadic regulators, its native-packing bound
and V39's fixed-regulator vacuum-slab limit give
\[
 \nu_\gamma((G_\gamma^+)^c)
 \le \exp\left\{-\frac{\gamma}{640}(d_\gamma-A_\gamma)_+\right\},
 \qquad A_\gamma=\frac{122880B_\gamma}{\gamma}.
 \tag{V43.21}\label{r43:rough}
\]
For periodic laws, the corresponding joint bound holds on
every selected set of distinct native anchors by V35.
For the slab statement, take the established total-variation
limit at fixed coupling; the event is a bounded measurable
function of the same finite slab variables.
Since $A_\gamma/d_\gamma\to0$, eventually
\[
 \nu_\gamma((G_\gamma^+)^c)
             \le\exp\{-\gamma^{1/3}/2621440\}.
 \tag{V43.22}\label{r43:rough-asymptotic}
\]
The sufficient threshold $d_\gamma\ge2A_\gamma$ is explicit
and extremely conservative. The uniform source estimate
\eqref{r43:physical-bound} does not need to discard this
rough event or assume that an entire lattice is good.

\subsection{What this does not close, and why threshold tuning is insufficient}

The absolute-diagram reference retains all strand orders.
It has not been reduced to a finite source bank or a practical
quadrature algorithm. Its absolute values can also be
nonsmooth in exterior coordinates; no uniform retained-field
derivative theorem is implicit in \eqref{r43:normalized-bounds}.
The derivative proved here is only the coupling derivative.

Most importantly, positivity of the local diagram weights
does not give connected decay of the original retained law.
Changing many block factors can accumulate an extensive
normalizer error, and no reflection-positive transfer for
the changed factors has been constructed. Equation
\eqref{r43:physical-bound} keeps the original transfer
precisely to avoid that substitution.

The error is smaller than every inverse power of $\gamma$,
but is not a cutoff-power estimate on a proposed trajectory
$\log(1/a)\asymp\gamma$. Multiplication of this upper bound
by $a^{-p}$, $p>0$, does not tend to zero there. This is a
limitation of the proved bound, not a lower bound on the
actual error. The basic competition between a long-strand
cost and a defect tail is already visible in
\[
          \sup_{\epsilon>0}
                  \min\{\epsilon^{-1},\gamma\epsilon^2\}
                                      =\gamma^{1/3}.
 \tag{V43.23}\label{r43:balance}
\]
Indeed the second entry is at most $\gamma^{1/3}$ below
$\epsilon=\gamma^{-1/3}$ and the first is at most that value
above it; equality holds at the crossing.
Thus retuning this elementary two-budget argument alone
does not reach an $e^{-c\gamma}$ error.

The next task is an actual connected estimate for the
positive absolute-diagram interaction with the exact
rough predictors and original retained normalizer, or a
source-relative physical low-energy estimate. Further
pointwise sign examples or higher local Taylor coefficients
do not supply that estimate. The continuum construction
and full Yang--Mills mass gap remain unproved and active.

\subsection{Checks and preservation}

The diagnostic
\begin{center}\small
\path{scripts/verify_ym_restart_v43_signed_returns.py}
\end{center}
compares a labelled-pairing enumeration with an independent
coordinate-moment expansion for the anisotropic generator
and its marked-open source, including coincident endpoints.
It checks exact compact barrier fixtures, the denominator
matrix identity, native tree-distance bounds, the normalized
signed derivative identity and the displayed exponent budgets.
These are finite algebraic diagnostics, not evaluations of
an infinite diagram sum or certificates of physical mixing.

The complete V42 body is preserved. Removing this appendix
and restoring its title recovers V42 exactly. Only the
newest active main source is retained; all archived and
companion sources are unchanged. Build and check files
remain outside \path{latex_notes}.
% END CONSOLIDATED V43 APPEND

% BEGIN CONSOLIDATED V44 APPEND -- 2026-09-20
\clearpage
\section{V44: local specification returns and an untouched-link obstruction}
\label{r44:start}

V43 controlled a conditional normalizer and its energy source.
It did not compare the resulting retained Gibbs specifications.
This appendix makes that comparison with boundary normalization
included. On the actual sparse packing, each retained link enters
at most two changed factors. Every finite-window conditional law
is consequently close to the original law, uniformly in the ambient
volume and in all boundary values.

We then test single-link Dobrushin contraction of the positive
reference. It fails at large coupling for this native construction:
some retained links lie outside every integrated factor.
Their six-staple influences tend to one even along boundary changes
tending to the identity, with every selected coefficient defect zero.
This embeds an inherited separated-minimizer mechanism in the
present packing; it is not a general no-go theorem for Yang--Mills.
It excludes this microscopic contraction route, not physical
clustering, larger-block criteria or the mass gap.

\paragraph{Nonduplication and targeted reading.}
Archive f6 V23 already proves compact separated-minimizer
concentration and an isolated one-plaquette example. Its warning
requires the other staples and generated interactions to be
included before applying that example in four dimensions.
The present test includes all six staples and proves that the
generated factors are absent from the selected row.
Archive f12 V71 also uses a cleared corridor, but for a
product-reference oscillation obstruction in a different packing.
The new claims here are the exact overlap-two return and the
six surviving influences of the present hybrid specification.
V16 and archive f6 V19--V22 already contain strong-coupling
Dobrushin comparison. Those are not new weak-coupling results.
A targeted recheck of Rebeschini--van Handel,
\href{https://arxiv.org/html/1308.4117}
{\emph{Comparison Theorems for Gibbs Measures}},
Corollary 2.6 and the subsequent block discussion, confirms that
closeness of specifications is not itself contraction.
All estimates used below are proved directly. The next companion
review is V45; the next broad review is V50.

\subsection{A positive reference with rough factors unchanged}

Fix a finite periodic regulator and V29's sparse packing
\(\mathcal B\), using coordinate order \((\mu,\nu,\rho,\tau)\).
Anchor residues are \(x_\nu,x_\rho\in4\mathbb Z\),
\(x_\tau\in2\mathbb Z\); the excluded temporal layers remain
excluded. No plaquette meets two distinct free six-link blocks.
Let \(\mathcal R\) be the retained links and \(E_{\rm u}\) the
sum of untouched plaquette energies. The exact marginal is
\[
 d\mu_\gamma(\eta)=
 \frac{e^{-\gamma E_{\rm u}(\eta)}
             \prod_{b\in\mathcal B}Q_{\gamma,b}(\eta)}
      {\mathcal Z_\gamma}\,dH_{\mathcal R}(\eta).
 \tag{V44.1}\label{r44:original}
\]
Here \(Q_{\gamma,b}=e^{-29\gamma}Z_{\gamma,b}\) is V41's original
compact integral. No all-good conditioning is imposed.

Fix a base coupling \(\gamma_0\ge1\). During coupling
differentiations freeze the indicators
\(\chi_b=\mathbf1_{\{\mathcal D_b^+\le d_{\gamma_0}\}}\).
For \(\gamma\) near \(\gamma_0\), put
\[
 \ell_{\gamma,b}=\chi_b
       \log\frac{Q_{\gamma,b}^{\rm abs}}{Q_{\gamma,b}},
 \qquad
 \widehat Q_{\gamma,b}=Q_{\gamma,b}e^{\ell_{\gamma,b}},
 \qquad V_\gamma=\sum_b\ell_{\gamma,b}.
 \tag{V44.2}\label{r44:hybrid}
\]
The hybrid factor equals \(Q^{\rm abs}\) on this fixed good
set and \(Q\) elsewhere. V43 proves, at \(\gamma_0\),
\[
 0\le\ell_{\gamma_0,b}\le e_{\gamma_0},\qquad
 |\partial_\gamma\ell_{\gamma,b}|_{\gamma=\gamma_0}
                 \le e_{\gamma_0},\qquad
 e_\gamma=C(1+\gamma)e^{-\gamma^{1/3}}.
 \tag{V44.3}\label{r44:small}
\]
Thus \(-\partial_\gamma\log\widehat Q\) equals \(H^{\rm abs}\)
on the fixed good set and \(H\) off it. No derivative of a
moving indicator is suppressed. Exterior smoothness is not assumed.

For \(0\le t\le1\), the normalized interpolation is
\[
 d\mu_{\gamma,t}
    =\frac{e^{tV_\gamma}}{\mu_\gamma(e^{tV_\gamma})}\,d\mu_\gamma,
 \qquad \widehat\mu_\gamma=\mu_{\gamma,1}.
 \tag{V44.4}\label{r44:interpolation}
\]
These positive finite-volume laws are auxiliary.
No reflection-positive transfer or original physical vacuum is
assigned to them. V43's physical source theorem still uses the
unchanged original transfer.

\subsection{The changed support has overlap at most two}

Let \(S_b=\mathcal R_b\) be V41's 80 retained supporting links.
The upper filling, frame tree and two left squares use only \(S_b\).
Hence \(Q_b,Q_b^{\rm abs},\mathcal D_b^+\) and \(\ell_b\)
are \(S_b\)-local. The additional faces of V35's energy witness
are not needed in this functional support.

\begin{lemma}[Native support multiplicity and corridor]
\label{r44:overlap}
For this packing,
\[
 k_i:=\#\{b:i\in S_b\}\le2,\qquad
 n_U:=\#\{b:S_b\cap U\ne\varnothing\}\le2|U|
 \tag{V44.5}\label{r44:counts}
\]
for each retained link \(i\) and finite \(U\subset\mathcal R\).
The corridor
\[
 \mathcal C=\{(x,\nu):x_\nu=2\pmod4\}
 \tag{V44.6}\label{r44:corridor}
\]
misses every \(S_b\) and every free block.
\end{lemma}
\begin{proof}
Translate an anchor to zero. The 22 positive retained \(\mu\)
links are the private outer links with bases
\[
 \begin{array}{c|c|c|c}
 \mu&\nu&\rho&\tau\\ \hline
 0&-1,2&0,1,2&0\\
 0&0,1&-1,3&0\\
 0&0,1&0,1,2&-1,1
 \end{array}
\]
and the other retained links have bases
\[
 \begin{array}{c|c|c|c|c|c}
 \text{direction}&\mu&\nu&\rho&\tau&\text{count}\\ \hline
 \nu&0,1&-1,0,1&0,1,2&0&18\\
 \rho&0,1&0,1&-1,0,1,2&0&16\\
 \tau&0,1&0,1&0,1,2&-1,0&24 .
 \end{array}
\]
These are the retained edges of the 29 incident plaquettes;
they contain the upper paths and left squares.

For a transverse link, the \(\nu,\rho,\tau\) offsets determine
at most one residue-compatible anchor. Its two \(\mu\) offsets
allow at most two anchors. For a \(\mu\) link, the three rows
of the first table have disjoint residue patterns.
The first allows at most one anchor. The \(\rho\) offsets
\(-1,3\) in the second can give two anchors, as can the
\(\tau\) offsets \(-1,1\) in the third; other coordinates
are fixed. Deleting anchors cannot increase these counts.
The cofinal periods have no local self-identifications, so the
same argument works at seams.
Thus \(k_i\le2\), and
\(n_U\le\sum_{i\in U}k_i\).
Supported \(\nu\) links have base residue \(3,0\), or \(1\),
never \(2\); all free links have direction \(\mu\).
\end{proof}

The multiplicity two is specific to this packing and coordinate
order. Arbitrary translates of this template only give the bound
24 from the direction counts. Superposing orientations or new
eliminations requires a fresh count.

\subsection{An all-boundary local return theorem}

Let \(\pi_U^\zeta,\widehat\pi_U^\zeta\) be the original and
hybrid conditional laws on \(U\) at \(\gamma_0\), with all
other retained links frozen to the same \(\zeta\).
The positive factor formulas define these kernels for every
boundary, including zero-probability configurations.

\begin{lemma}[Bounded-tilt estimate]
\label{r44:tilt}
If \(d\nu=e^v\,d\mu/\mu(e^v)\) and
\(\operatorname{osc}v\le d\), then
\[
 \|\nu-\mu\|_{\rm TV}\le\tanh(d/4),\qquad
 \|\nu-\mu\|_{\rm TV}:=\sup_A|\nu(A)-\mu(A)|.
 \tag{V44.7}\label{r44:tilt-bound}
\]
\end{lemma}
\begin{proof}
Set \(r=e^v\), \(m=\inf r\), \(M=\sup r\), \(z=\mu r\).
For \(M>m\), the convex chord bound for \(|r-z|\) gives
\[
 \|\nu-\mu\|_{\rm TV}
 \le\frac{(M-z)(z-m)}{z(M-m)}
 \le\frac{\sqrt M-\sqrt m}{\sqrt M+\sqrt m}.
\]
The maximum occurs at \(z=\sqrt{Mm}\).
The final expression is \(\tanh(\log(M/m)/4)\).
Constant tilts are immediate; essential bounds suffice.
\end{proof}

\begin{theorem}[Native finite-window comparison]
\label{r44:return}
Uniformly in the ambient regulator and every common boundary,
\[
 \frac{d\widehat\pi_U^\zeta}{d\pi_U^\zeta}
 =\frac{e^{V_U^\zeta}}{\pi_U^\zeta(e^{V_U^\zeta})},
 \qquad
 V_U^\zeta=\sum_{b:S_b\cap U\ne\varnothing}\ell_{\gamma_0,b},
 \tag{V44.8}\label{r44:local-ratio}
\]
and
\[
 \|\widehat\pi_U^\zeta-\pi_U^\zeta\|_{\rm TV}
 \le t_U:=\tanh(n_Ue_{\gamma_0}/4)
 \le\tanh(|U|e_{\gamma_0}/2).
 \tag{V44.9}\label{r44:TV}
\]
The bound with \(|t-s|e_{\gamma_0}\) compares interpolation
parameters \(s,t\). For real bounded \(F,G\) on the window,
\[
 |\Cov_{\widehat\pi_U^\zeta}(F,G)-\Cov_{\pi_U^\zeta}(F,G)|
 \le2\,\operatorname{osc}F\,\operatorname{osc}G\,t_U .
 \tag{V44.10}\label{r44:covariance-return}
\]
\end{theorem}
\begin{proof}
Factors missing \(U\) are constants in its conditional integral
and cancel from the normalized ratio. The remaining tilt lies
between zero and \(n_Ue_{\gamma_0}\). The preceding lemma and
support count prove the first assertions.

For covariance use
\[
 \Cov_\lambda(F,G)=\frac12\int
 (F(x)-F(y))(G(x)-G(y))\,d\lambda(x)d\lambda(y).
\]
The full integrand, including \(1/2\), has oscillation at most
\(\operatorname{osc}F\,\operatorname{osc}G\).
Telescoping products gives
\(\|\lambda^{\otimes2}-\lambda'^{\otimes2}\|_{\rm TV}
\le2\|\lambda-\lambda'\|_{\rm TV}\).
The bounded-function total-variation inequality finishes the proof.
\end{proof}

This compares conditionals, not unconditional local marginals of
different Gibbs states, whose boundary distributions also change.
For example set
\(B_U(F)=\sup_{\zeta,\zeta'}|\pi_U^\zeta F-\pi_U^{\zeta'}F|\),
with \(\widehat B_U\) defined similarly. The triangle inequality gives
\[
 |B_U(F)-\widehat B_U(F)|\le2\operatorname{osc}F\,t_U ,
 \tag{V44.11}\label{r44:boundary}
\]
not smallness of either influence.

For windows with \(O(r^4)\) links the return tends to zero if
\(r^4e_\gamma\to0\). Thus \(r=\exp(\alpha\gamma^{1/3})\),
\(0<\alpha<1/4\), is allowed. This does not certify physical-sized
windows on a trajectory \(\log(1/a)\asymp\gamma\).
It is a mesoscopic comparison, not physical correlation decay.

\subsection{The connected and source normalizers remain explicit}

At fixed \(\gamma_0\), for fixed bounded observables,
\[
 \begin{split}
 \widehat\mu F-\mu F
   &=\int_0^1\sum_b\Cov_{\mu_t}(F,\ell_b)\,dt,\\
 \Cov_{\widehat\mu}(F,G)-\Cov_\mu(F,G)
   &=\int_0^1\sum_b\kappa_{\mu_t}(F,G,\ell_b)\,dt ,
 \end{split}
 \tag{V44.12}\label{r44:connected}
\]
where
\(\kappa_\lambda(F,G,L)=
\lambda[(F-\lambda F)(G-\lambda G)(L-\lambda L)]\).
Indeed, differentiating the normalized expectation gives
\(\partial_t\mu_t F=\mu_t(FV)-\mu_t F\,\mu_t V\).
Apply this to \(\mu_t(FG)-\mu_t F\,\mu_t G\);
its five terms are the centered triple product.
Boundedness and finite volume justify differentiation,
and linearity in \(V=\sum_b\ell_b\) gives the formulas.

Locality does not make a distant covariance vanish.
Using only
\(|\kappa(F,G,\ell_b)|\le
\operatorname{osc}F\,\operatorname{osc}G\,e_{\gamma_0}\)
still pays the total number of blocks. Likewise, at the base coupling,
\[
 0\le\log(\widehat{\mathcal Z}_\gamma/\mathcal Z_\gamma)
                                      \le|\mathcal B|e_\gamma .
 \tag{V44.13}\label{r44:pressure}
\]
For a coupling-source check put
\(K_\gamma=E_{\rm u}+\sum_bH_{\gamma,b}\).
With the cutoff frozen, exact differentiation at \(\gamma_0\) yields
\[
 \partial_\gamma\log
       \frac{\widehat{\mathcal Z}_\gamma}{\mathcal Z_\gamma}
 =\widehat\mu_\gamma\!\left(\sum_b\partial_\gamma\ell_{\gamma,b}\right)
       -\{\widehat\mu_\gamma K_\gamma-\mu_\gamma K_\gamma\}.
 \tag{V44.14}\label{r44:pressure-derivative}
\]
The original log-density derivative is \(-K_\gamma\);
the changed one is \(-K_\gamma+\sum_b\partial_\gamma\ell_{\gamma,b}\).
The second term is a genuine change of expectation. The bounded
local source error alone is not the derivative of the full changed
normalizer.

\subsection{A native all-good obstruction to single-link contraction}

For either law define
\[
 C_{ij}=\sup_{\zeta,\zeta':\,\zeta_k=\zeta'_k\ (k\ne j)}
 \|\pi_{\{i\}}^\zeta-\pi_{\{i\}}^{\zeta'}\|_{\rm TV},\qquad i\ne j .
 \tag{V44.15}\label{r44:influence}
\]
The usual row criterion requires
\(\sup_i\sum_{j\ne i}C_{ij}<1\).
The single-link return theorem implies
\[
 |\widehat C_{ij}-C_{ij}|\le2\tanh(e_\gamma/2),
 \tag{V44.16}\label{r44:coefficient-return}
\]
but supplies no contraction margin. For the following row,
the two specifications are exactly the same.

\begin{theorem}[Six surviving native influences]
\label{r44:obstruction}
Choose \(i=(x,\nu)\in\mathcal C\) and its six opposite
parallel neighbours
\[
 j_{d,\sigma}=(x+\sigma\widehat d,\nu),\qquad
 d\in\{\mu,\rho,\tau\},\quad \sigma=\pm1 .
\]
The kernels at \(i\) are identical for the original, hybrid,
and every interpolated law, at every boundary. For these six \(j\),
\[
 C_{ij}=\widehat C_{ij}
 \ge1-2K\gamma^{3/2}e^{-5\gamma^{1/3}/2592},
 \qquad K=e^3/c_H,\quad c_H=8/(3\pi^3).
 \tag{V44.17}\label{r44:lower}
\]
In particular
\[
 \liminf_{\gamma\to\infty}\sum_{j\ne i}C_{ij}\ge6 .
 \tag{V44.18}\label{r44:row}
\]
Thus the unweighted single-link row criterion fails for all
sufficiently large \(\gamma\). The witnesses have every selected
\(\mathcal D_b^+=0\), and the changed links tend to the identity.
\end{theorem}
\begin{proof}
No generated factor or cutoff depends on \(i\), by the support
lemma. No plaquette at \(i\) contains a free link: its
\(\mu\nu\) plaquettes have parallel \(\mu\) links at \(\nu\)
residues 2 and 3; the others contain no \(\mu\) link.
Its conditional action is exactly its six untouched Wilson
plaquettes. All generated factors cancel at every boundary.

For a specified \(j\), compare the all-identity boundary to the
boundary with only \(U_j=q_\theta=\cos\theta+\mathbf e_1\sin\theta\).
This opposite parallel link occurs in exactly one incident plaquette.
Writing \(u=U_i\), their full energy is
\[
 6-5\operatorname{Re}u-\operatorname{Re}(u q_\theta^{-1})
 =6-h_\theta\cdot u,\qquad
 h_\theta=(5+\cos\theta)e_0+\sin\theta\,e_1 .
 \tag{V44.19}\label{r44:staple}
\]
Both plaquette orientations give this same real part.
The two laws have densities proportional to
\(e^{6\gamma e_0\cdot u}\) and \(e^{\gamma h_\theta\cdot u}\).
For \(0<\theta\le1\), \(5\le|h_\theta|\le6\), and their centers
satisfy
\[
 d_\theta:=|h_\theta/|h_\theta|-e_0|
                  \ge\sin\theta/6\ge\theta/12 .
 \tag{V44.20}\label{r44:separation}
\]

For \(5\le h\le6\), a unit vector \(v\), and \(0<r<2\),
the law \(d\lambda_{h,v}\propto e^{\gamma h v\cdot u}dH(u)\)
has the normalized tail bound
\[
 \lambda_{h,v}\{|u-v|\ge r\}
              \le K\gamma^{3/2}e^{-\gamma h r^2/2}.
 \tag{V44.21}\label{r44:cap}
\]
Factor out \(e^{\gamma h}\). The geodesic cap of radius
\(\gamma^{-1/2}\) has Haar mass at least \(c_H\gamma^{-3/2}\),
and on it \(\gamma h(1-v\cdot u)\le h/2\le3\).
This gives denominator at least \(e^{-3}c_H\gamma^{-3/2}\).
Outside the Euclidean radius \(r\), the exponent costs
\(\gamma h|u-v|^2/2\ge\gamma h r^2/2\).
The remaining Haar mass is at most one, proving the bound.

The event \(|u-e_0|<d_\theta/3\) has probability at least
\(1-K\gamma^{3/2}e^{-5\gamma d_\theta^2/18}\) under the first
law, and at most the same error under the second: its points
are at distance greater than \(2d_\theta/3\) from the second center.
Subtract the probabilities, use \(d_\theta\ge\theta/12\),
and choose \(\theta=\gamma^{-1/3}\).
This proves \eqref{r44:lower}. The six neighbours are distinct
on the cofinal regulators, so their sum gives \eqref{r44:row}.

Both \(i\) and \(j\) miss every \(S_b\). Every native block
therefore sees its coherent exterior, for every value of \(u\),
and its defect is zero. The boundary change is at most
\(\theta\to0\). This is an exact normalized four-dimensional
native conditional, not an isolated-plaquette model.
\end{proof}

\begin{corollary}[Positive site weights do not repair this row test]
\label{r44:weighted}
For every assignment \(w_i>0\) to retained links,
\[
 \max_i\frac1{w_i}\sum_{j\ne i}C_{ij}w_j
 \ge6\bigl(1-2K\gamma^{3/2}e^{-5\gamma^{1/3}/2592}\bigr).
 \tag{V44.22}\label{r44:weighted-row}
\]
The same holds for every interpolated reference.
\end{corollary}
\begin{proof}
Fix one \(\nu\) coordinate congruent to two modulo four.
Its corridor links form a three-dimensional periodic graph
in the \(\mu,\rho,\tau\) directions, with six neighbours per vertex.
Choose a link minimizing \(w_i\) on this finite graph.
Each of its six neighbours has weight at least \(w_i\);
apply \eqref{r44:lower}. If the displayed lower bound is negative,
nonnegativity of the row already proves the assertion.
\end{proof}

The changed boundary is not globally flat: untouched plaquettes
can have defect \(1-\cos\theta\).
Zero selected defects do not control these plaquettes.
There is no contradiction with V35's tail theorem. The conclusion
is that making the integrated blocks all good does not pay
microscopic retained contraction.

This does not prove a massless physical excitation.
A supremum over pinned link coordinates, including gauge-dependent
tests, is stronger than decay of gauge-invariant correlations
under the original law. Larger blocks, adapted metrics, averaged
couplings and physical spectral arguments remain possible,
with their hypotheses still to be proved.

\subsection{Next direction and checks}

The positive hybrid reference now has an all-boundary local return
with overlap two, not an ambient-volume error.
Its single-link contraction route is nevertheless excluded by a
literal native six-staple row, even on the selected all-good set.
The next attempt must include untouched plaquettes and their links
in a block or gauge-invariant correlation estimate.
Equation \eqref{r44:TV} supplies the return comparison for such
blocks, not their contraction or cumulant summability.

The diagnostic
\begin{center}\small
\path{scripts/verify_ym_restart_v44_specification_returns.py}
\end{center}
independently enumerates the template, periodic overlap and corridor,
and checks all six staple positions with rational quaternions.
Finite-state fixtures verify local normalizer cancellation,
the sharp bounded-tilt constant and the connected derivative
through a separate three-source logarithmic expansion.
These do not certify Wilson mixing or a physical spectrum.

All V43 content is preserved; removing this appendix and restoring
its title recovers the predecessor exactly. All other sources are
unchanged. Builds stay outside \path{latex_notes}.
The interacting continuum construction and full Yang--Mills
mass gap remain unproved; the goal stays active.
% END CONSOLIDATED V44 APPEND


% BEGIN CONSOLIDATED V45 APPEND -- 2026-09-20
\clearpage
\section{V45: joint corridor integration and collective strand moments}
\label{r45:start}

V44 identified links untouched by the sparse six-link elimination.
The next step is to release an entire connected layer of these links,
including all plaquettes joining them, instead of treating their
one-link conditionals separately. The resulting normalized conditional
is an unpinned three-dimensional \(O(4)\) connection model.
At coherent exterior data it is the ordinary ferromagnetic model.
There are no private pins: a neighbour released into the same block
cannot still be counted as a fixed aligning field.

We prove an exact squared-strand-length identity for this joint model.
For any fixed connected finite graph with \(m>0\) edges it implies
\[
 \lim_{\gamma\to\infty}
  \frac{\mathbb E_\gamma\sum_C|C|^2}{(\gamma m)^2}=\frac13,
 \qquad
 \liminf_{\gamma\to\infty}
  \mathbb P_\gamma\!\left\{\max_C|C|\ge\frac{\gamma m}{6}\right\}
                                                        \ge\frac1{36}.
 \tag{V45.1}\label{r45:headline}
\]
Here strand length counts bond occurrences, not spatial diameter.
The limit fixes the regulator first. It is neither a thermodynamic
loop theorem nor a physical low-energy statement.
It does prove that V43's coupling-uniform exponential strand tail
cannot simply be extended to this unpinned native block.

\subsection{Scheduled companion checkpoint and inherited scope}

Fresh hashes show that repository NS V26, SM--Einstein V72,
Shell Mixing V16 and Parallel YM V5 are unchanged since V40.
Their terminal proofs and limitations were reconsidered.
NS's rank-one refinement gives no retained YM covariance estimate;
SM's many-copy fixed-mode transport is not a cutoff-uniform
finite-species construction; Shell Mixing still leaves its
noncommuting mixed-curvature calculation and nonperturbative
remainder open; Parallel YM's extensive-charge tightness target
remains excluded.

The newer Downloads Source Selection V35, Regenerative Stationarity
V38 and Exact-Action Bridge V28 also have unchanged hashes.
Their respective stationary-adjoint convention, boundary-root
obstruction and enlarged-profile rank result do not supply
a physical YM contraction. Their terminal scope statements were
rechecked, not every old proof recertified.

Perturbative ESM has advanced from the supplied V23 to V27.
The V24--V26 closing ledgers describe a degree-closed one-loop
pure-gravity source module at a retained positive lower scale.
The relevant new V27 construction and proof were read:
its metric--Higgs hard Hessian is a joint block matrix, with
a scalar-degree-two coefficient containing both
\[
 \frac12\Tr(G_gU)
 \quad\hbox{and}\quad
 -\frac12\Tr(G_gLG_HL^{\mathsf T}).
\]
Adding isolated sector determinants would omit the mixed term.
This is a methodological warning to assemble the simultaneous
operands and their normalization. No ESM coefficient, contour
theorem, infrared removal or nonperturbative exponent is imported
into Yang--Mills. Here the simultaneous compact corridor integral
is obtained directly from Wilson plaquettes, without a Gaussian
approximation or a new ultraviolet regulator.

\paragraph{Targeted literature and nonduplication.}
The random-wire representation is established prior work; see
Benassi--Ueltschi,
\href{https://arxiv.org/html/1807.06564v4}
{\emph{Loop correlations in random wire models}},
Section 6, Proposition 6.1.
The revised primary text and the representation's scope were checked.
We do not import its XY-specific loop-distribution results as an
\(O(4)\) theorem. Nor do we claim a Poisson--Dirichlet law from
the single moment in \eqref{r45:headline}.
Archive \(C\), V61--V62 and V42--V43 already provide the Haar-pairing
representation. The new work is the actual joint corridor
identification, its normalized tensor moment and the resulting
failure of the proposed pinned-tail extension.
The next companion and broad literature checkpoints are both V50.

\subsection{The exact joint native conditional}

Keep V44's original sparse packing and coordinate order.
For a fixed \(k=2\pmod4\), let
\[
 \mathcal C_k=\{(x,\nu):x_\nu=k\}.
\]
Release all links in this one layer and fix every other retained
link. Write \(z_x=U_\nu(x)\in S^3\) for its variables.
V44 proves that no generated factor \(Q_b\), hybrid factor
\(\widehat Q_b\), or native good-set cutoff depends on them.
The remaining graph \(G_k\) has coordinates \((x_\mu,x_\rho,x_\tau)\).
On periods \(L,L,2q\), it has
\[
 n=L^2(2q),\qquad m=3n,
\]
with periodic nearest-neighbour edges in the \(\mu,\rho,\tau\)
directions.

\begin{theorem}[Joint corridor law, without retained pins]
\label{r45:joint}
For arbitrary fixed exterior data \(\eta\), the normalized law
on \(\mathcal C_k\), under the original retained law and every
V44 interpolation, is
\[
 d\lambda_{\gamma,\eta}(z)=
 \frac{\exp\{\gamma\sum_{x,d}z_x\cdot T_{x,d}z_{x+\widehat d}\}}
      {Z_{\gamma,G_k}(T)}
                       \prod_{x\in G_k}dH(z_x),
 \qquad d\in\{\mu,\rho,\tau\},
 \tag{V45.2}\label{r45:law}
\]
where
\[
 T_{x,d}z=l_{x,d}zr_{x,d}^{-1},\qquad
 l_{x,d}=U_d(x),\quad r_{x,d}=U_d(x+\widehat\nu).
 \tag{V45.3}\label{r45:transport}
\]
Every \(T_{x,d}\) is the actual orthogonal left/right transport.
There is no one-spin field term.
\end{theorem}
\begin{proof}
Each plaquette containing a released \(\nu\) link has its
opposite parallel \(\nu\) link in this same layer. Counting it
once gives exactly the displayed three positive graph directions.
For either plaquette orientation its real trace is
\[
 \operatorname{Re}
 (z_x r_{x,d}z_{x+\widehat d}^{-1}l_{x,d}^{-1})
       =z_x\cdot(l_{x,d}z_{x+\widehat d}r_{x,d}^{-1}).
\]
No other Wilson plaquette depends on the released links.
The sum of incident energies is consequently
\(m-\sum_{x,d}z_x\cdot T_{x,d}z_{x+\widehat d}\).
Its scalar \(e^{-\gamma m}\), all exterior plaquettes and every
generated factor cancel from the normalized conditional.
This proves the formula, including its denominator.
\end{proof}

At the coherent exterior \(l=r=1\), \eqref{r45:law} is the
zero-field \(O(4)\) model on the whole periodic graph.
Every selected six-link defect remains zero for every
corridor configuration: those defects see no corridor links.
The corridor action itself is not zero or discarded.
The coherent exterior is a special boundary configuration,
not a positive-probability phase of the full Wilson law.

\subsection{Closed strands and their anisotropic generator}

For this subsection take any fixed connected finite simple graph
\(G=(V,E)\), with \(n=|V|\), \(m=|E|>0\), and set
\[
 B(z)=\sum_{\{i,j\}\in E}z_i\cdot z_j,\qquad
 Z_\gamma=\int e^{\gamma B(z)}\,dH(z).
\]
Use \(\langle\cdot\rangle_\gamma\) for this spin law.
In the inherited pairing expansion there are no pin occurrences.
A diagram has bond multiplicities \(m_e\), even endpoint degrees
\(N_i=\sum_{e\ni i}m_e\), and a pairing at each vertex.
Its components are all closed. Its normalized weight is
\[
 \frac1{Z_\gamma}
 \left(\prod_e\frac{\gamma^{m_e}}{m_e!}\right)
 \left(\prod_i u(N_i)\right)4^{\#\mathcal C},
 \qquad u(2r)=\frac1{2^r(r+1)!}.
 \tag{V45.4}\label{r45:diagram}
\]
Let
\[
 M=\sum_em_e=\sum_{C\in\mathcal C}|C|,\qquad
 S_2=\sum_{C\in\mathcal C}|C|^2,\qquad L_*=\max_C|C|,
\]
with \(L_*=0\) for the empty diagram.
Expectation under \eqref{r45:diagram} is \(\mathbb E_\gamma\),
distinguished from spin expectation.

Write \(z_i=(z_i^0,\mathbf z_i)\) and define the auxiliary
anisotropic partition function
\[
 Z_{\gamma,s}=\int
 \exp\{\gamma\sum_{\{i,j\}}
        (z_i^0z_j^0+s\,\mathbf z_i\cdot\mathbf z_j)\}\,dH .
\]
Color summation in each closed component gives exactly
\[
 \frac{Z_{\gamma,s}}{Z_\gamma}
  =\mathbb E_\gamma G_s,\qquad
 G_s=\prod_C\frac{1+3s^{|C|}}4,\quad s\ge0 .
 \tag{V45.5}\label{r45:generator}
\]
One component is either scalar-colored, with weight one,
or has one of three transverse colors, with weight \(s^{|C|}\).
This is the inherited pairing argument with the pins removed.
Compact exponential majorants justify the full series and all
finitely many marks used below.

\begin{lemma}[Occurrence moments and the second anisotropic mark]
\label{r45:marks}
For every integer \(j\ge1\),
\[
 \mathbb E_\gamma(M)_j=\gamma^j\langle B^j\rangle_\gamma ,
 \qquad
 \mathbb E_\gamma M=\gamma\langle B\rangle_\gamma,\qquad
 \Var_\gamma M=\gamma^2\Var_{\rm spin}(B)
                         +\gamma\langle B\rangle_\gamma .
 \tag{V45.6}\label{r45:occurrences}
\]
Here \((M)_j=M(M-1)\cdots(M-j+1)\).
At \(s=1\),
\[
 \partial_s\log G_s=\frac34 M,\qquad
 \partial_s^2\log G_s=\frac3{16}S_2-\frac34M .
 \tag{V45.7}\label{r45:marks-eq}
\]
\end{lemma}
\begin{proof}
Differentiate the unnormalized positive series \(Z_\gamma\)
\(j\) times in its common coupling, multiply by \(\gamma^j\),
and divide by the same \(Z_\gamma\).
The compact integral derivative is \(\int B^je^{\gamma B}dH\).
This gives the first identity; using
\(M^2=(M)_2+M\) gives the variance identity.
For a component of length \(a\), differentiate
\(\log((1+3s^a)/4)\) twice at one. The answers are
\(3a/4\) and \(3a^2/16-3a/4\). Sum over components.
\end{proof}

\subsection{The normalized tensor identity}

Define the real symmetric \(4\)-by-\(4\) matrix
\[
 J(z)=\frac12\sum_{\{i,j\}\in E}
       (z_i z_j^{\mathsf T}+z_j z_i^{\mathsf T}),
 \qquad \Tr J=B,\qquad
 A(z)=4\Tr(J^2)-B^2\ge0 .
 \tag{V45.8}\label{r45:tensor}
\]
The nonnegativity follows from the Cauchy--Schwarz inequality
for the four real eigenvalues; \(J\) itself need not be positive.

\begin{theorem}[Exact closed-strand second moment]
\label{r45:second}
For every \(\gamma>0\),
\[
 \boxed{\displaystyle
 \mathbb E_\gamma S_2
 =\gamma\langle B\rangle_\gamma
       +\frac{\gamma^2}{9}
                      \langle4\Tr(J^2)-B^2\rangle_\gamma .}
 \tag{V45.9}\label{r45:identity}
\]
\end{theorem}
\begin{proof}
Put \(T(z)=\sum_{\{i,j\}}\mathbf z_i\cdot\mathbf z_j=B-J_{00}\).
Taking two logarithmic derivatives of \eqref{r45:generator}
at \(s=1\), with its full normalization, yields
\[
 \gamma^2\Var_{\rm spin}(T)
 =\frac3{16}\mathbb E_\gamma S_2-\frac34\mathbb E_\gamma M
                                      +\frac9{16}\Var_\gamma M .
 \tag{V45.10}\label{r45:variance}
\]
Indeed the logarithm of the expectation contributes the variance
of the first mark, not just the second mark's expectation.

To evaluate the spin variance, use simultaneous \(O(4)\) invariance.
For an independent uniform unit vector \(v\in S^3\), and any fixed
real symmetric \(J\), the sphere second and fourth moments give
\[
 \mathbb E_v(v^{\mathsf T}Jv)=\frac{\Tr J}{4},\qquad
 \mathbb E_v(v^{\mathsf T}Jv)^2
                      =\frac{(\Tr J)^2+2\Tr(J^2)}{24}.
 \tag{V45.11}\label{r45:rotation}
\]
For example the fourth moment tensor is
\((\delta_{ab}\delta_{cd}+\delta_{ac}\delta_{bd}
+\delta_{ad}\delta_{bc})/24\); contracting twice with \(J\)
proves the formula.
Randomly rotate the whole spin configuration, preserving its law.
Conditional on the unrotated configuration, \(T\) has mean \(3B/4\)
and variance \(A/48\). Total variance therefore gives
\[
 \Var_{\rm spin}(T)
       =\frac9{16}\Var_{\rm spin}(B)+\frac1{48}\langle A\rangle_\gamma .
 \tag{V45.12}\label{r45:total-variance}
\]
Insert this and \eqref{r45:occurrences} into
\eqref{r45:variance}. The two
\(9\gamma^2\Var_{\rm spin}(B)/16\) terms cancel exactly.
The remaining equation is
\(\gamma^2\langle A\rangle/48
 =3(\mathbb E S_2-\gamma\langle B\rangle)/16\),
which proves \eqref{r45:identity}.
\end{proof}

This is a normalized, all-orders identity at fixed graph size.
It does not identify strand moments with a physical transfer
spectrum. The anisotropy is a generating parameter for the proof,
not a permitted change of the Wilson action.

\subsection{A positive probability of coupling-scale strands}

\begin{lemma}[Fixed-graph compact concentration]
\label{r45:concentration}
For each fixed connected \(G\), as \(\gamma\to\infty\),
\[
 \langle B^j\rangle_\gamma\longrightarrow m^j\quad(j\ge1),
 \qquad \langle A\rangle_\gamma\longrightarrow3m^2 .
 \tag{V45.13}\label{r45:limits}
\]
\end{lemma}
\begin{proof}
The deficit
\(R=m-B=\frac12\sum_{\{i,j\}}|z_i-z_j|^2\) is nonnegative.
Its zero set consists exactly of \(z_i=v\) for all \(i\),
with \(v\in S^3\), by graph connectivity.

For \(\gamma\ge1\), restrict all spins to the geodesic cap of
radius \(\gamma^{-1/2}\) around \(e_0\).
Each normalized Haar cap has mass at least
\(c_H\gamma^{-3/2}\), \(c_H=8/(3\pi^3)\), and \(R\le2m/\gamma\)
there. Thus for every fixed \(\delta>0\),
\[
 \mathbb P_{\rm spin}\{R\ge\delta\}
 \le c_H^{-n}\gamma^{3n/2}e^{2m-\gamma\delta}
                                      \longrightarrow0 .
 \tag{V45.14}\label{r45:concentration-bound}
\]
On the zero set, \(B=m\), \(J=mvv^{\mathsf T}\) and \(A=3m^2\).
Compactness, continuity and \eqref{r45:concentration-bound}
give convergence of these bounded functions to their constant
values on the zero set. No uniform thermodynamic limit was used.
\end{proof}

\begin{theorem}[Long strands at fixed regulator]
\label{r45:long}
For every fixed connected \(G\) with \(m>0\),
\[
 \frac{\mathbb E_\gamma S_2}{(\gamma m)^2}\longrightarrow\frac13,
 \quad
 \frac{\mathbb E_\gamma M}{\gamma m}\longrightarrow1,\quad
 \frac{\mathbb E_\gamma M^4}{(\gamma m)^4}\longrightarrow1 .
 \tag{V45.15}\label{r45:moment-limits}
\]
For each \(0<a<1/3\),
\[
 \liminf_{\gamma\to\infty}
       \mathbb P_\gamma\{L_*\ge a\gamma m\}\ge(1/3-a)^2 .
 \tag{V45.16}\label{r45:probability}
\]
\end{theorem}
\begin{proof}
The first two limits follow from
\eqref{r45:identity}, \eqref{r45:occurrences}
and \eqref{r45:limits}.
For the third use
\(M^4=(M)_4+6(M)_3+7(M)_2+M\) and the same identities.
All graph parameters are fixed.

Pointwise, \(S_2\le L_*M\le M^2\). Let
\(E=\{L_*\ge a\gamma m\}\) and \(p_\gamma=\mathbb P_\gamma(E)\).
On its complement \(S_2\le a\gamma mM\); consequently
\[
 \mathbb E_\gamma S_2
 \le a\gamma m\,\mathbb E_\gamma M
                   +(\mathbb E_\gamma M^4)^{1/2}p_\gamma^{1/2}.
\]
Divide by \((\gamma m)^2\) and use \eqref{r45:moment-limits}.
Since \(1/3-a>0\), rearrangement and a lower limit give
\eqref{r45:probability}. Taking \(a=1/6\) proves
\eqref{r45:headline}.
\end{proof}

In particular there cannot be constants \(C_G,c_G>0\) and
finite \(p\ge0\), independent of \(\gamma\), such that
\[
 \mathbb P_\gamma\{L_*\ge L\}
                 \le C_G(1+\gamma)^p e^{-c_GL}
 \quad\hbox{for all }\gamma\ge1,\ L\ge1 .
 \tag{V45.17}\label{r45:no-tail}
\]
Taking \(L=\lceil\gamma m/6\rceil\) contradicts
\eqref{r45:probability}. This excludes the proposed form of
pinned-tail extension even when constants may depend on the
fixed corridor graph. It does not exclude a bound with decay
rate of order \(1/\gamma\), nor a differently reorganized expansion.

\subsection{The generator itself has an extensive low-temperature cost}

The failure can also be seen before taking its derivatives.
For every fixed \(s>1\),
\[
 \lim_{\gamma\to\infty}\frac1\gamma
     \log\frac{Z_{\gamma,s}}{Z_\gamma}=m(s-1).
 \tag{V45.18}\label{r45:generator-growth}
\]
Here is a direct compact proof. With
\(D_s=\operatorname{diag}(1,s,s,s)\), each bond
\(z_i^{\mathsf T}D_sz_j\) is at most \(s\).
That maximum is attained simultaneously when all spins are
the same transverse unit vector. In caps of geodesic radius
\(\gamma^{-1/2}\) around that vector, each bond deficit is at
most \(3s/\gamma\): use
\[
 s-z_i^{\mathsf T}D_sz_j
   =s(1-z_i\cdot z_j)+(s-1)z_i^0z_j^0
   \le2s/\gamma+(s-1)/\gamma .
\]
The cap mass is again at least
\(c_H^n\gamma^{-3n/2}\).
Together with the coherent cap bound this yields
\[
 \begin{split}
 m(s-1)\gamma-3sm+n\log c_H-\tfrac32n\log\gamma
 &\le\log(Z_{\gamma,s}/Z_\gamma),\\
 \log(Z_{\gamma,s}/Z_\gamma)
 &\le m(s-1)\gamma+2m-n\log c_H+\tfrac32n\log\gamma .
 \end{split}
 \tag{V45.19}\label{r45:generator-bounds}
\]
Division by \(\gamma\) proves \eqref{r45:generator-growth}.
Thus the bounded anisotropic normalizer used for V43's pinned
block is not available for this joint unpinned block.

\subsection{Physical meaning, next task, and checks}

The common orientation of the coherent minimum manifold is not
automatically a physical excitation. For a path in the corridor
graph, close a Wilson rectangle with the corresponding left
and right exterior paths. At the identity exterior its normalized
trace is \(z_x\cdot z_y\), equal to one for every common
orientation \(z_i=v\).
A gauge-invariant source therefore requires its own test;
the orbit variance used in the generating calculation cannot
be declared physical spectral weight.

Likewise, a strand can traverse the same finite graph many times.
The probability in \eqref{r45:probability} does not prove a strand
of large spatial diameter, long-range connected correlations of
the unconditioned Wilson law, or a zero physical mass.
The coherent exterior is not integrated with its original
probability here. The limits do not commute with a growing
regulator without additional estimates.

The concrete change of direction is to retain these collective
corridor fields, their left/right connections and gauge-invariant
sources in the next joint effective description.
One cannot replace the released neighbours by the old private
pins, or use V43's strand tail as a uniform input for the larger
block. The remaining problem is a connection-dependent,
normalized connected estimate after averaging the other retained
links under their original law. The mixed-term discipline of
ESM V27 supports this joint treatment, but does not provide that
missing estimate.

The diagnostic
\begin{center}\small
\path{scripts/verify_ym_restart_v45_corridor_strands.py}
\end{center}
enumerates the full native corridor incidence and checks its
literal plaquette identity on noncommuting quaternion fixtures.
Independent labelled-pairing and coordinate-moment expansions
agree on the anisotropic generator and squared-length mark:
a single edge through degree eight and a four-cycle through
degree six. Further exact checks verify the rotational fourth
moment, its centering, and the variance cancellation.
These finite checks are not proofs by sampling of an infinite
series, a thermodynamic loop law or a mass gap.

The complete V44 body is preserved, all 46 other source files
are unchanged, and only the latest active main source is retained.
Build and diagnostic files remain outside \path{latex_notes}.
The full interacting four-dimensional continuum and physical
Yang--Mills mass gap remain unproved and active.
% END CONSOLIDATED V45 APPEND


% BEGIN CONSOLIDATED V46 APPEND -- 2026-09-20
\clearpage
\section{V46: thermal exterior connections and signed Wilson-source returns}
\label{r46:start}

V45 forces us to retain the collective corridor orientation rather
than replace it by a private pin. The next question concerns actual
gauge-invariant sources, not strand length. We calculate their
normalized connected return when the exterior links differ from
identity on the thermal scale \(\gamma^{-1/2}\).
The result includes both a shifted Gaussian fluctuation and a
nontrivial integral over the common orientation.

The main new conclusion is a native counterexample to a proposed
coherent-source shortcut: two consecutive corridor plaquette energies
have a strictly negative conditional covariance at suitable
\(O(\gamma^{-1/2})\) exterior data, whereas their coherent-exterior
leading covariance is strictly positive. The exterior tends to
identity and eventually lies in every V43 native good set.
This is not a failure of physical reflection positivity, a negative
variance, or a mass-gap conclusion.

\subsection{Nonduplication and the inherited compact integral}

V1 already proves the full fixed-graph electric integral,
its Gaussian mean and covariance, its normalized root completion,
and two orders of its scalar correction:
Theorem~\ref{restart:main} and
Corollary~\ref{restart:root-expansion}.
After permuting coordinate names, its slab integral is precisely
V45's corridor conditional. A new leading normalizer or another
generic Laplace theorem would duplicate that work.

V2 already derives coherent connected plaquette insertions on
contractible four-dimensional boxes. Archive f16.2 V63 already
records negative transverse link covariances in a pinned four-link
cell; its V64 treats induced two-layer predictors.
Those are not the unpinned, thermal-exterior, gauge-invariant
plaquette covariance constructed here.
Archive f13's orbit disintegration and f16.1 V44's gauge tube are
also inherited techniques, not new claims.

For context, Bandeira, Singer and Spielman,
\href{https://arxiv.org/html/1204.3873}
{\emph{A Cheeger Inequality for the Graph Connection Laplacian}},
Section 1.3, describe the orthogonal-connection quadratic form.
Its spectral relaxation is not a physical Wilson transfer gap.
Only that mathematical context is used; all source formulas below
are proved from the original compact conditional and V1's already
controlled localization. The next broad and companion reviews
remain due at V50.

\subsection{The same conditional law and its collective orientation}

Fix a connected finite graph \(G\), with one orientation per edge,
\(n\ge2\) vertices and \(m\) edges. It will subsequently be the actual
periodic corridor graph. Set \(\varepsilon=\gamma^{-1/2}\).
For fixed imaginary quaternions \(\alpha_e,\beta_e\), prescribe
the actual exterior links and their transport by
\[
 l_e=e^{\varepsilon\alpha_e},\qquad
 r_e=e^{\varepsilon\beta_e},\qquad
 T_e(\varepsilon)z=l_ezr_e^{-1}=e^{\varepsilon A_e}z,\qquad
 A_ez=\alpha_ez-z\beta_e .
 \tag{V46.1}\label{r46:transport}
\]
Each \(A_e\) is a real skew-symmetric matrix. Left and right
multiplications commute, giving the last exponential exactly.
All constants below may depend on \(G\), the chosen finite paths,
and a fixed bound for \(\alpha,\beta\).
No uniformity in growing graph size is asserted.

Let \(D:\mathbb R^n\to\mathbb R^m\) have rows
\(b_e^{\mathsf T}\), with \(b_e=\mathbf1_{s(e)}-\mathbf1_{t(e)}\).
Define
\[
 L=D^{\mathsf T}D,\qquad
 R=DL^\dagger D^{\mathsf T},\qquad P=I-R .
 \tag{V46.2}\label{r46:projection}
\]
The dagger is the inverse on the mean-zero vertex space, zero
on constants. Thus \(R\) and \(P\) are orthogonal projections
onto gradients and cycles. They are finite graph matrices,
not physical spectral projections.
The conditional energy and its unnormalized compact integral are
\[
 E_\varepsilon(z)=\frac12\sum_e|z_{s(e)}-T_e(\varepsilon)z_{t(e)}|^2,
 \qquad
 Q_\varepsilon(A)=\int e^{-E_\varepsilon(z)/\varepsilon^2}\prod_i dH(z_i).
 \tag{V46.3}\label{r46:Q}
\]
This is exactly V45's original, hybrid and interpolated conditional,
with the scalar plaquette constants retained in \(E_\varepsilon\).

For \(v\in S^3\), let \(h_e(v)=A_ev\in v^\perp\), and put
\[
 \begin{split}
 K_A&=\sum_{e,f}P_{ef}A_e^{\mathsf T}A_f,\qquad
 V_A(v)=\tfrac12\|Ph(v)\|^2=\tfrac12v^{\mathsf T}K_Av,\\
 Z_A&=\int_{S^3}e^{-V_A(v)}dH(v),\qquad
 d\rho_A(v)=Z_A^{-1}e^{-V_A(v)}dH(v).
 \end{split}
 \tag{V46.4}\label{r46:rotor}
\]
Symmetry and idempotence of \(P\) show that \(K_A\) is symmetric
positive semidefinite.

\begin{proposition}[V1's leading root law in corridor variables]
\label{r46:inherited}
Uniformly on the declared bounded exterior sets,
\[
 \frac{Q_\varepsilon(A)}{Q_\varepsilon(0)}
                  =Z_A+O_{G,M}(\varepsilon).
 \tag{V46.5}\label{r46:ratio}
\]
In leading conditional moments, the common orientation has law
\(\rho_A\). Conditional on \(v\), the rescaled differences of the
spins have the distribution of the differences of
\[
 \xi=\mu(v)+\eta,\qquad
 \mu(v)=L^\dagger D^{\mathsf T}h(v),\qquad
 \mathbb E(\eta_i\eta_j^{\mathsf T}\mid v)
                 =L^\dagger_{ij}(I-vv^{\mathsf T}).
 \tag{V46.6}\label{r46:Gaussian}
\]
Here \(\sum_i\xi_i=0\). This is a description of the leading
conditional law, not a replacement of the compact law at finite
\(\varepsilon\).
\end{proposition}
\begin{proof}
Use the exact V1 root completion, writing \(z_i=b_iv\) and
\(b_o=1\). In its notation \(X=\alpha\), \(Y=\beta\), and the
root \(v\) changes the upper endpoint to \(\Ad_v\beta\).
Right multiplication by \(v\) is an isometry, so
\[
 \|P(\alpha-\Ad_v\beta)\|^2
                         =\|P(\alpha v-v\beta)\|^2=\|Ph(v)\|^2.
\]
This proves \eqref{r46:ratio} directly from the inherited expansion,
with its same compact complement and Haar factor.

For the Gaussian law, write \(z_i=v+\varepsilon\xi_i+O(\varepsilon^2)\)
in root coordinates. The leading energy is
\(\frac12\|D\xi-h(v)\|^2\).
Completing the square gives the displayed mean modulo constants.
The rooted inverse and \(L^\dagger\) give the same law of all
differences: subtracting the mean of the rooted field gives the
mean-zero Gaussian with this covariance. Equivalently both produce
the same edge covariance \(R\), since the incidence matrix with
one root column removed has the same range as \(D\).
The Gaussian determinant is independent of \(v\); its entire
normalizing factor cancels in \eqref{r46:ratio}.
The remaining root density is exactly \(\rho_A\).
\end{proof}

In particular fixing \(v=e_0\), discarding \(Z_A\), or averaging a
conditional source with unweighted Haar instead of \(\rho_A\)
would generally change the answer. If
\(A'_e=A_e+S_{s(e)}-S_{t(e)}\), with each \(S_i\) skew, then
\(Ph'=Ph\). This is the infinitesimal change induced by
\(T'_e=e^{\varepsilon S_s}T_ee^{-\varepsilon S_t}\).
It leaves \(K_A\) unchanged. The construction removes infinitesimal
pure-gauge gradients without deleting cycle or orientation variables.

\subsection{A normalized Wilson-path insertion theorem}

For a fixed oriented graph path \(p\) from \(i\) to \(j\), let
\(c_p\in\mathbb R^m\) count signed traversals. Then
\(D^{\mathsf T}c_p=b_p=\mathbf1_i-\mathbf1_j\).
Use the ordered transports, with inverses for reversed steps:
\[
 T_p(\varepsilon)=\prod_{e\in p}^{\rm ordered}T_e(\varepsilon),\qquad
 W_p=z_i\cdot T_p(\varepsilon)z_j,\qquad
 X_p^\varepsilon=\varepsilon^{-2}(1-W_p).
 \tag{V46.7}\label{r46:writer}
\]
For the corridor, \(W_p\) is the normalized Wilson trace of the
rectangle closed by its actual lower and upper exterior paths.
It is gauge invariant when those paths and the endpoint links
are transformed together. A one-edge path is an elementary
\(\nu d\) plaquette.

Define, allowing distinct paths with the same endpoints,
\[
 r_{pq}=b_p^{\mathsf T}L^\dagger b_q,\qquad
 C_p=\sum_e(Pc_p)_e A_e,\qquad f_p(v)=|C_pv|^2 .
 \tag{V46.8}\label{r46:source-data}
\]
Each \(C_p\) is skew. In particular its action is tangent to \(S^3\).
The path data depend on the cycle traversals through \(C_p\),
not just on their endpoints.

\begin{theorem}[Full leading connected source return]
\label{r46:source}
For the normalized compact conditional \eqref{r46:Q},
\[
 \begin{split}
 \mathbb E X_p^\varepsilon
 &=\frac32r_{pp}+\frac12\mathbb E_{\rho_A}f_p+O(\varepsilon),\\
 \Cov(X_p^\varepsilon,X_q^\varepsilon)
 &=\frac32r_{pq}^2
   +r_{pq}\,\mathbb E_{\rho_A}\bigl[(C_pv)\cdot(C_qv)\bigr]\\
 &\hspace{1.5em}+\frac14\Cov_{\rho_A}(f_p,f_q)+O(\varepsilon).
 \end{split}
 \tag{V46.9}\label{r46:covariance}
\]
The remainders are uniform on fixed bounded exterior sets.
Thus the expression on the right is the limit of
\(\gamma^2\Cov(W_p,W_q)\) as well.
The last term is a collective-orientation predictor covariance;
the middle term need not be positive.
\end{theorem}
\begin{proof}
The ordered product has expansion
\(T_p=I+\varepsilon\sum_e(c_p)_eA_e+O(\varepsilon^2)\).
Orthogonality gives the exact identity
\(1-W_p=\frac12|z_i-T_pz_j|^2\).
The leading rescaled source is therefore
\[
 X_p^0=\tfrac12|Y_p|^2,\qquad
 Y_p=\xi_i-\xi_j-\sum_e(c_p)_eA_ev .
\]
Under \eqref{r46:Gaussian},
\[
 \mathbb E(Y_p\mid v)=-C_pv,\qquad
 \Cov(Y_p,Y_q\mid v)=r_{pq}(I-vv^{\mathsf T}).
 \tag{V46.10}\label{r46:Y}
\]
Indeed \(b_p^{\mathsf T}\mu=c_p^{\mathsf T}Rh\), and subtracting
\(c_p^{\mathsf T}h\) gives \(-c_p^{\mathsf T}Ph=-C_pv\).

For two jointly Gaussian tangent vectors of dimension three,
with means \(d_p,d_q\) and cross covariance \(rI_3\), direct
Gaussian pairing gives
\[
 \Cov(\tfrac12|Y_p|^2,\tfrac12|Y_q|^2\mid v)
                          =\tfrac32r^2+r\,d_p\cdot d_q.
 \tag{V46.11}\label{r46:Wick}
\]
For completeness, the centered quadratic terms have covariance
\(2\Tr((rI_3)(rI_3)^{\mathsf T})/4=3r^2/2\);
the two linear terms give \(r\,d_p\cdot d_q\);
the mixed odd Gaussian terms vanish.
Their conditional means are
\((3r_{pp}+f_p(v))/2\).
Total expectation and total covariance now yield
\eqref{r46:covariance}, including the last normalization-sensitive
term.

We supply the insertion remainder rather than infer it by
differentiating an undifferentiated error.
V1's root-coordinate localization, uniformly also in \(v\),
bounds the rescaled compact density by
\(C e^{-c|\zeta|^2}\) after its common
\(\varepsilon^{3(n-1)}\) factor is removed.
On a sufficiently small fixed coordinate tube, Taylor's formula
bounds the error in the rescaled energy by
\(C\varepsilon(1+|\zeta|)^3\), and the error in each
\(X_p^\varepsilon-X_p^0\) by the same type of polynomial.
The source itself is bounded there by
\(C(1+|\zeta|)^2\). These bounds follow either from its exact
squared-distance identity or from the finite ordered product.
The two rescaled energies, exact and limiting, are bounded below
by \(c|\zeta|^2-C\). The inequality
\[
 |e^{-u}-e^{-w}|\le |u-w|(e^{-u}+e^{-w})
\]
therefore supplies an integrable polynomial times a Gaussian for
the errors in the partition, first source moment and product of
two source moments. The smooth Haar density obeys the same
dominated first-order estimate. Outside the fixed tube, the
source product is at most \(4\varepsilon^{-4}\), while V1's
compact gap makes its integral exponentially small.
The limiting divided denominator is bounded away from zero on
each bounded exterior set, since \(V_A\) is uniformly bounded.
Division and subtraction of the two exact means give
\(O(\varepsilon)\) for the normalized covariance. This proves
the asserted remainder with all normalizations present.
\end{proof}

When \(A=0\), this reduces to the familiar centered Gaussian
coefficient \(3r_{pq}^2/2\). That specialization is not claimed
new. What cannot be discarded at thermal exteriors are the
shifted-Gaussian and orientation-predictor terms.

\subsection{An actual all-good exterior with negative plaquette covariance}

Use the V45 corridor with fixed \(\nu\) coordinate \(k=2\pmod4\),
and the inherited periodic sizes \(L,L,2q\).
Choose one straight \(\mu\)-direction ring at fixed \(\rho,\tau\).
Let \(c\) be one on its positively oriented edges and zero
elsewhere. Thus \(D^{\mathsf T}c=0\) and \(Pc=c\).
Let \(\mathrm i\) be a fixed unit imaginary quaternion and choose
the literal retained links
\[
 U_\mu(x)=e^{2\varepsilon\mathrm i}
 \quad\hbox{on this lower ring in the plane }x_\nu=k;
 \qquad\hbox{all other exterior links are }1.
 \tag{V46.12}\label{r46:native}
\]
In particular all corresponding upper-plane links are identity.
These lower links are retained, since the eliminated sparse
blocks consist of \(\nu\) links. The generated factors still
cancel from the conditional on corridor links by V45.

If \(Jz=\mathrm i z\), then \(J^{\mathsf T}J=I\), and
\(A_e=2c_eJ\). Hence
\[
 K_A=4L I,\quad V_A=2L,\quad \rho_A=H,\quad
 C_e=2J\quad(e\hbox{ on the ring}).
 \tag{V46.13}\label{r46:ring}
\]
This is a genuine left/right mismatch, not a common conjugation
of the two exterior planes.

\begin{theorem}[Sign change inside a shrinking native good neighborhood]
\label{r46:negative}
Take two consecutive positively oriented edges
\(e=(x,x+\widehat\mu)\) and
\(f=(x+\widehat\mu,x+2\widehat\mu)\) of this ring.
For the elementary plaquette energies \(F_e=1-W_e\),
\(F_f=1-W_f\) in the normalized native conditional,
\[
 \begin{split}
 \gamma^2\Cov(F_e,F_f)\big|_{\eqref{r46:native}}
       &\longrightarrow \tfrac32r_{ef}^2+4r_{ef}<0,\\
 \gamma^2\Cov(F_e,F_f)\big|_{\text{identity exterior}}
       &\longrightarrow \tfrac32r_{ef}^2>0 .
 \end{split}
 \tag{V46.14}\label{r46:negative-eq}
\]
Both exteriors eventually belong to all native V43 good sets.
The inequalities hold at every fixed allowed corridor regulator.
\end{theorem}
\begin{proof}
First \(r_{ef}<0\), without any numerical inverse.
Set \(u=L^\dagger b_e\), so \(Lu=b_e\).
The discrete maximum principle puts the minimum at the sink
\(x+\widehat\mu\). To spell out the strictness needed here,
a minimum at any other harmonic vertex propagates to all its
neighbors. The remaining arc of the ring joins \(x+2\widehat\mu\)
to the source \(x\) without passing through the sink.
If \(u(x+2\widehat\mu)\) were equal to the sink minimum,
propagation along this arc would make the source a minimum,
contradicting \((Lu)(x)=1\).
Therefore
\[
 r_{ef}=b_f^{\mathsf T}u
           =u(x+\widehat\mu)-u(x+2\widehat\mu)<0 .
\]
The minimum assertion itself follows by the same propagation:
at a minimum the Laplacian is nonpositive, excluding the source;
a harmonic minimum propagates until reaching the sink.
Moreover \(|r_{ef}|\le1\), since \(R\) is an orthogonal projection.
It follows that \(r_{ef}(3r_{ef}/2+4)<0\).

For \eqref{r46:native}, both \(f_e(v)\) and \(f_f(v)\) are
the constant four, and
\((C_ev)\cdot(C_fv)=4\).
Insert these facts in \eqref{r46:covariance}; its orientation
covariance term is zero, not omitted. The identity exterior gives
the second line. The \(O(\gamma^{-1/2})\) scaled errors then give
the same strict signs at all sufficiently large \(\gamma\).

Finally each link of the exterior differs from identity by
\(O(\varepsilon)\). V35--V41's native defects are fixed finite
sums of squared chordal distances of bounded-length retained
words and frames, with value zero at identity. Telescoping each
word therefore gives
\(\mathcal D_b^+\le C_0\varepsilon^2=C_0/\gamma\), uniformly
over anchors of the fixed-orientation packing.
V43's threshold is \(\gamma^{-2/3}/2048\); the preceding bound
lies below it once \(\gamma^{1/3}\ge2048C_0\).
The constant depends on the fixed word templates, not on the
number of anchors. Every generated factor remains in its
declared good regime. Nothing in this step drops the untouched
corridor plaquettes.
\end{proof}

For an independent small-graph check, on the four-cycle
\(r_{ee}=3/4\), \(r_{ef}=-1/4\) for distinct positive cycle edges.
With the same amplitude two the limiting off-diagonal covariance
is \(-29/32\), compared with \(3/32\) at identity.
The diagonal is \(123/32\); hence this sign is compatible with
a positive semidefinite covariance matrix. This four-cycle is a
check of the formula, not the proof of its native torus realization.

\subsection{What the conditional result permits, and what it does not}

There is no contradiction with V44's comparison of the original
and hybrid conditional laws: they are exactly equal on this
corridor. The comparison that fails is replacement of the actual
thermal exterior by the coherent exterior in its leading connected
source coefficient. In particular no error
\(o(\gamma^{-2})\), uniform throughout the native good set, can
justify that replacement for these two sources.
The limit of the scaled difference is \(4r_{ef}\ne0\).

In this particular witness \(D^{\mathsf T}h=0\), so the leading
Gaussian drift is zero and the orientation remains uniform.
The sign change comes from keeping the actual twisted Wilson
insertions, whose linear Gaussian pieces have a signed cross
covariance. Thus we have not disproved using a coherent leading
fluctuation measure while retaining these perturbed insertions;
we have disproved replacing the complete source coefficient by
its zero-exterior value. For general left/right data the other
terms of \eqref{r46:covariance} must also remain.

A negative off-diagonal covariance is not a negative variance.
Nor is this pinned-exterior law a reflection-positive physical
vacuum with a freely selectable physical time. Nothing above
proves or disproves a positive physical mass.

For the original exterior marginal \(\nu_\gamma\), exact total
covariance still reads
\[
 \Cov_{\rm original}(F_e,F_f)
 =\mathbb E_{\nu_\gamma}\Cov(F_e,F_f\mid\eta)
   +\Cov_{\nu_\gamma}
           \bigl(\mathbb E[F_e\mid\eta],\mathbb E[F_f\mid\eta]\bigr).
 \tag{V46.15}\label{r46:outer}
\]
The conditional formulas supply actual operands, not the missing
law of \(\eta\). A prescribed exterior has Haar probability zero;
continuity makes the strict finite-\(\gamma\) sign persist on
an open neighborhood, but gives no useful lower probability
uniformly in coupling or regulator.
To integrate the thermal formulas, one still needs control of
the exterior's rescaled connections, the complementary region
and the source-weighted remainders under that original marginal.
The constants in the fixed-graph proof include its inverse
Laplacian and cannot be silently used on a growing corridor.

The next upstream task is thus to assemble the collective
connection and source predictors with their genuine exterior
weight. Repeating the coherent Gaussian coefficient, forcing
conditional positivity, or suppressing the orientation integral
would miss terms now calculated explicitly.
ESM V27's joint-operand discipline is respected, but supplies none
of these nonperturbative estimates.

The exact-arithmetic diagnostic
\begin{center}\small
\path{scripts/verify_ym_restart_v46_source_returns.py}
\end{center}
checks graph projections, root independence, noncommuting
quaternion tangent data, independent Gaussian polynomial
contractions, the cycle sign calculation, and native ring
incidence. These checks do not establish growing-volume
remainder bounds or the full physical gap.
The V45 body is retained exactly. Only the current main
\texttt{.tex} remains active; build and diagnostic files stay
outside \path{latex_notes}. The full interacting four-dimensional
continuum Yang--Mills construction and mass gap remain open.
% END CONSOLIDATED V46 APPEND


% BEGIN CONSOLIDATED V47 APPEND
\clearpage
\section{V47: Thermal joint defects from the full Wilson pressure ratio}
\label{r47:context}

V46 exposed the failure of coherent-source substitution on thermally
small exteriors. The next question is what can actually be averaged
under the original exterior law. We go upstream to the normalizer:
the absolute partition-function lower bound used in V35 loses a
factor of \(\log\gamma\) in every local defect budget. Keeping half
the Wilson action in the numerator replaces this loss by a
boundary-to-volume pressure mismatch. This gives stronger joint
defect, source-weighted remainder, and physical-source comparison
bounds on specified rectangular regulators.

This is not a new proof of the archived pressure mechanism.
Archive \path{Renewal_Yang_Mills_Mass_Gap_Research_Notes_V100_f14.tex},
V95, already proves the cubic comb-tree pressure comparison and a
thermal energy bound; its V98 proves cubic single-plaquette
exponential moments and an \(\ell^1\) source ball. Here the new
steps are the explicit rectangular boundary count, its use in
V35's \emph{joint native coefficient-defect} chessboard, and its
export to the original retained and physical vacuum source errors.
The cubic pressure argument is reproved only to make the rectangular
constants and hypotheses checkable. Neither an independent-plaquette
Wilson law nor an independent retained law is introduced.

A targeted literature check revisited P.\ A.\ Faria da Veiga and
M.\ O'Carroll, \emph{On Yang--Mills Stability Bounds and Plaquette
Field Generating Function}, Sections IV--V,
\url{https://arxiv.org/html/2205.07376}.
The relevant structural input is cancellation of the extensive
coupling power between normalized numerator and denominator.
Their displayed imaginary-trace source vanishes identically for
\(\SU(2)\); their Remark 16 also discusses action observables.
No generating-function theorem is imported with that source
silently changed. Our action and native defects are treated
directly below. Bounds on local moments, in that paper or here,
are not a construction of an interacting continuum theory or a
positive physical mass gap. The next regular broad literature
and companion-programme checkpoints remain V50.

\subsection{Rectangular pressure bounds without discarding torons}

Let
\[
 \Lambda_{\mathbf N}=\prod_{\mu=1}^4\mathbb Z/N_\mu\mathbb Z,
 \quad N_\mu\ge3,\quad V=\prod_\mu N_\mu,\quad
 \sigma_{\mathbf N}=\sum_\mu N_\mu^{-1}.
\]
There are \(4V\) positive links and \(6V\) plaquettes. On the full
periodic compact configuration space set
\[
 \begin{gathered}
 e(q)=1-\operatorname{Re}q,\qquad S(U)=\sum_p e(U_p),\\
 Z_{\mathbf N,\gamma}=\int e^{-\gamma S(U)}\,dH(U),\\
 d\mu_{\mathbf N,\gamma}=Z_{\mathbf N,\gamma}^{-1}
                         e^{-\gamma S(U)}dH(U),
 \end{gathered}
 \tag{V47.1}\label{r47:law}
\]
where \(q\) is a unit quaternion and every Haar factor is normalized.
The symbol \(\gamma\) has the same Wilson-action normalization as V35.

\begin{lemma}[Rectangular comb and triangular plaquettes]
\label{r47:comb}
A spanning-tree gauge leaves
\[
 D=3V+1
\]
independent Haar link coordinates. There is a collection of
\[
 n_\triangle=3V-V\sigma_{\mathbf N}+1,\qquad
 D-n_\triangle=V\sigma_{\mathbf N},
 \tag{V47.2}\label{r47:counts}
\]
distinct plaquettes whose holonomies, together with the remaining
link coordinates, are independent normalized Haar coordinates
under the \emph{reference} measure.
\end{lemma}
\begin{proof}
Choose coordinate representatives \(0\le x_\mu<N_\mu\).
Take as tree links \((x,\mu)\) those with
\(x_\mu<N_\mu-1\) and \(x_\nu=0\) for every \(\nu>\mu\).
There are
\(\sum_\mu(N_\mu-1)\prod_{\nu<\mu}N_\nu=V-1\) such links.
Every vertex has a unique route to the origin obtained by
decreasing its highest positive coordinate until it vanishes.
Thus this is a spanning tree, not a periodic gauge restriction.
Successive vertex gauge transformations set its links to \(1\).
Haar invariance gives product normalized Haar on the other
\(4V-(V-1)=D\) coordinates in every gauge-invariant integral.

Call a positive link open if it does not wrap its own coordinate.
There are \(4V-V\sigma_{\mathbf N}\) open links. For each open
off-tree link \((x,\mu)\), let
\(\nu=\max\{\nu'>\mu:x_{\nu'}>0\}\) and \(y=x-\widehat\nu\).
Select the \(\mu\nu\) plaquette based at \(y\). Its two
\(\nu\)-links belong to the tree: both have no positive coordinate
above \(\nu\), and their \(\nu\) coordinate is less than
\(N_\nu-1\). Consequently, in tree gauge its holonomy is
\[
 Q_{x,\mu}=k_\mu(y)k_\mu(x)^{-1}.
\]
Different selected links give different plaquettes. Order
off-tree open links of each direction by increasing
\(\sum_{\nu>\mu}x_\nu\). Its predecessor \(k_\mu(y)\) is
either a tree link or an earlier coordinate, so
\(k_\mu(x)=Q_{x,\mu}^{-1}k_\mu(y)\) is a triangular,
invertible Haar-preserving change of variables.
The number of selected links is the number of open links minus
\(V-1\), giving \eqref{r47:counts}.
All wrapping links remain among the unintegrated coordinates.
No winding holonomy, flat sector, or collective orientation
has been fixed. Independence is asserted only for this Haar
coordinate system, not for \(\mu_{\mathbf N,\gamma}\).
\end{proof}

\begin{proposition}[Full compact pressure ratio]
\label{r47:pressure}
Put \(c_H=8/(3\pi^3)\) and \(C_1=(\pi/2)^{5/2}\).
For \(\gamma\ge1\),
\[
 e^{-48V}c_H^D\gamma^{-3D/2}
 \le Z_{\mathbf N,\gamma}
 \le C_1^{n_\triangle}\gamma^{-3n_\triangle/2}.
 \tag{V47.3}\label{r47:two-sided}
\]
In particular, for \(\gamma\ge2\), define
\[
 B^*_{\mathbf N,\gamma}
       =66+\tfrac32\sigma_{\mathbf N}\log\gamma .
\]
Then
\[
 \frac1V\log\frac{Z_{\mathbf N,\gamma/2}}{Z_{\mathbf N,\gamma}}
            \le B^*_{\mathbf N,\gamma}.
 \tag{V47.4}\label{r47:ratio}
\]
\end{proposition}
\begin{proof}
For the class integral
\(z_\gamma=\int_{\SU(2)}e^{-\gamma e(q)}\,dq\), the class-angle
formula and \(1-\cos\theta\ge2\theta^2/\pi^2\) give
\[
 \begin{split}
 z_\gamma
  &=\frac2\pi\int_0^\pi
       e^{-\gamma(1-\cos\theta)}\sin^2\theta\,d\theta\\
  &\le\frac2\pi\int_0^\infty
          e^{-2\gamma\theta^2/\pi^2}\theta^2\,d\theta
    =C_1\gamma^{-3/2}.
 \end{split}
\]
Discard all but the selected nonnegative plaquette energies and
use Lemma~\ref{r47:comb}; this proves the upper bound.
For the lower bound restrict all \(D\) off-tree links to angular
distance at most \(r=\gamma^{-1/2}\le1\) from \(1\).
Each cap has Haar mass at least \(c_Hr^3\), since
\(\sin\theta\ge2\theta/\pi\) on \([0,1]\).
A product of at most four such links or inverses has chord
distance at most \(4r\) from \(1\); hence its energy is at most
\(8r^2\). This includes wrapping plaquettes. Therefore
\(\gamma S\le48V\) on this cap event, proving the lower bound.

Divide the upper bound at \(\gamma/2\) by the lower bound at
\(\gamma\). Its logarithm divided by \(V\) is at most
\[
 48+\frac DV\log(c_H^{-1})
      +\frac{n_\triangle}{V}\log C_1
      +\frac{3n_\triangle}{2V}\log2
      +\frac32\sigma_{\mathbf N}\log\gamma.
\]
Here \(D/V\le4\), \(n_\triangle/V\le3\).
The constant part is less than \(66\); for example
\(\log(c_H^{-1})<5/2\), \(\log C_1<23/20\), and
\(\log2<7/10\) bound it by \(64.6\).
This proves \eqref{r47:ratio}. The mismatch in the powers of
\(\gamma\) is exactly the last term, not an extensive
\(\log\gamma\) cost independent of shape.
\end{proof}

For the programme's regulators \(\mathbf N=(L,L,L,T)\), this gives
\[
 B^*_{L,T,\gamma}
       =66+\tfrac32(3/L+1/T)\log\gamma.
 \tag{V47.5}\label{r47:shape}
\]
It is uniformly bounded on every specified window
\[
       (3/L+1/T)\log\gamma\le A<\infty.
 \tag{V47.6}\label{r47:window}
\]
This includes sequences with all side lengths of order \(a^{-1}\)
and \(\gamma=O(\log(1/a))\), \emph{if} such a sequence is chosen;
we have not constructed or identified the physical continuum
trajectory. At fixed finite sides the displayed bound still
has a logarithmic loss. No all-regulator uniformity is claimed.

\subsection{The native joint tail with the numerator action retained}

Throughout the remainder of V47 assume \(\gamma\ge2\).
Use exactly V35's local defect \(\mathcal D_x\), forty-five-face
witness \(\mathcal W_x\), and carrier of side lengths
\(h=(4,8,8,4)\); see \eqref{r35:D}--\eqref{r35:h}.
Thus
\(\mathcal D_x\le20\sum_{p\in\mathcal W_x}e(U_p)\) pointwise
in the full configuration. Retain the dyadic restrictions
\(L=2^j\), \(j\ge4\), and \(T=2^\ell\), \(\ell\ge3\).
Use \(K=1024\) for all anchors, or \(K=32\) for the native sparse
packing with the same omitted temporal anchors as V35.
Set, with \(B^*=B^*_{L,T,\gamma}\),
\[
 \begin{split}
 \widetilde q_{\gamma,K}(u)
   &=\min\left\{1,\exp\left(
                \frac{1024B^*-\gamma u/40}{K}\right)\right\},\\
 a^*&=\frac{40960B^*}{\gamma},
 \qquad\lambda^*=\frac{\gamma}{40K},
 \qquad
 \widetilde q_{\gamma,K}(u)
       =\exp\{-\lambda^*(u-a^*)_+\}.
 \end{split}
 \tag{V47.7}\label{r47:q}
\]

\begin{theorem}[Thermal joint native coefficient defects]
\label{r47:joint}
For every finite set of distinct anchors \(\mathcal A\) in the
selected family and arbitrary thresholds \(u_x\ge0\),
\[
 \mu_{L,T,\gamma}
   \{\mathcal D_x>u_x\text{ for all }x\in\mathcal A\}
       \le\prod_{x\in\mathcal A}\widetilde q_{\gamma,K}(u_x).
 \tag{V47.8}\label{r47:joint-eq}
\]
For native anchors this is also the tail under the unchanged
retained marginal \(\mu_{\mathcal R}\).
\end{theorem}
\begin{proof}
Fix one residue colour. Pull the witness event of each anchor
back to the base rectangle of its translated reflection tiling.
The inherited closed-block chessboard inequality allows distinct
events and thresholds. In the fully disseminated event for one
threshold \(u\), all copies of the forty-five-face witness lie
strictly inside their rectangles; they are disjoint sets of
original plaquettes, also across periodic seams.
There are \(N_h=V/1024\) rectangles. Hence that disseminated
event, denoted \(A_{\rm dis}(u)\), implies \(S>N_hu/20\).

Unlike V35's absolute numerator estimate, keep half the action:
\[
 \begin{split}
 \mu_{L,T,\gamma}(A_{\rm dis}(u))
 &=Z_{\gamma}^{-1}\int
       \mathbf1_{A_{\rm dis}(u)}e^{-\gamma S/2}e^{-\gamma S/2}\,dH\\
 &\le e^{-\gamma N_hu/40}\frac{Z_{\gamma/2}}{Z_\gamma}.
 \end{split}
\]
Here both partition functions are for this same full periodic
regulator. Equation~\eqref{r47:ratio} and the probability bound
by \(1\) give the chessboard root
\[
 \mu(A_{\rm dis}(u))^{1/N_h}
       \le\min\{1,\exp(1024B^*-\gamma u/40)\}.
\]
Multiply these roots over anchors of the chosen colour.
Generalized Holder over all \(K\) colours, with empty colours
represented by the constant \(1\), takes a \(1/K\) power of
each such factor and gives \eqref{r47:joint-eq}.
Every reflected measure in this argument is the original full
Wilson law. The last assertion is its pushforward identity,
since native \(\mathcal D_x\) is retained-measurable.
\end{proof}

For clarity the two exponents are
\[
 \widetilde q_{\gamma,32}(u)=
       \min\{1,e^{32B^*-\gamma u/1280}\},\qquad
 \widetilde q_{\gamma,1024}(u)=
       \min\{1,e^{B^*-\gamma u/40960}\}.
 \tag{V47.9}\label{r47:two-q}
\]
The constants are deliberately conservative. For fixed \(K,A\)
and any desired \(s>0\), the threshold
\[
             u=\frac{40960B^*+40Ks}{\gamma}
\]
has tail at most \(e^{-s}\). Thus in the window
\eqref{r47:window} the defect scale is \(1/\gamma\), with
explicit constants and no selected reference coupling.

\subsection{Arbitrary-support thermal sources and weighted rough returns}

Write
\[
 Y_x=\gamma\mathcal D_x,\qquad
 A^*=40960B^*,\qquad
 \rho_K=(40K)^{-1}.
\]
For a finite nonnegative source bank \(t=(t_x)\) with
\(t_x<\rho_K\), define
\[
 \mathcal H(t)=A^*\sum_x t_x-\sum_x\log(1-40Kt_x),
 \quad
 M_r=\sum_{j=0}^r\binom rj(A^*)^{r-j}j!(40K)^j,
 \quad M_0=1.
 \tag{V47.10}\label{r47:budgets}
\]
The quantity \(M_r\) is the \(r\)-th moment of
\(A^*+\operatorname{Exp}(\rho_K)\). This is a comparison
variable, not a representation of the Wilson defects.

\begin{corollary}[Positive sources and mixed moments]
\label{r47:moments}
For every such source bank,
\[
 \log\E\exp\Big(\sum_x t_xY_x\Big)\le\mathcal H(t),
 \qquad
 \E\prod_{x\in\mathcal A}Y_x^{r_x}
                  \le\prod_{x\in\mathcal A}M_{r_x}
 \tag{V47.11}\label{r47:mgf}
\]
for positive integers \(r_x\). If \(40Kt_x\le1/2\), then
\(\mathcal H(t)\le(A^*+80K)\|t\|_1\).
Expectations are under the original full law, or its exact
retained marginal for native marks.
\end{corollary}
\begin{proof}
After multiplying \(\gamma\) into the thresholds,
\eqref{r47:joint-eq} is the product upper-tail bound with
one-coordinate survival
\(\exp\{-\rho_K(v-A^*)_+\}\).
Expand each exponential as
\(1+\int_0^\infty t e^{tv}\mathbf1_{\{Y>v\}}\,dv\),
multiply, and apply Tonelli to the nonnegative terms.
Each subset term is bounded by the corresponding product
integral for independent comparison variables. This proves
the first bound without assuming independence of the actual
variables. Applying the layer-cake formula
\(Y^r=\int_0^\infty r v^{r-1}\mathbf1_{\{Y>v\}}\,dv\)
proves the second. Integrating the shifted exponential gives
exactly \(M_r\). Finally \(-\log(1-z)\le2z\) on \([0,1/2]\).
\end{proof}

This is a support-independent coordinate radius with an
explicit extensive source cost. It is stronger at this interface
than merely controlling one \(\ell^1\) ball, but is still
one-sided and not centered at the actual means.
It yields neither a source log-normalizer zero-free region
uniform in arbitrary support nor a decaying connected kernel.

\begin{proposition}[Source-weighted original-law estimates]
\label{r47:weighted}
Let \(0\le2t_x<\rho_K\) and
\(\Psi=\sum_x t_xY_x\). For any finite marked event bank
\(\mathcal B\), with arbitrary overlap with the source bank,
\[
 \E\left[e^\Psi
       \mathbf1_{\{\mathcal D_b>u_b,\ b\in\mathcal B\}}\right]
 \le e^{\mathcal H(2t)/2}
          \prod_{b\in\mathcal B}
                    \widetilde q_{\gamma,K}(u_b)^{1/2}.
 \tag{V47.12}\label{r47:weighted-event}
\]
For every positive integer \(r\),
\[
 \E[\mathcal D_b^{2r}e^\Psi]
      \le\gamma^{-2r}M_{4r}^{1/2}e^{\mathcal H(2t)/2}.
 \tag{V47.13}\label{r47:weighted-moment}
\]
Consequently, if retained errors \(R_b\) vanish on the rough
exterior and satisfy
\(|R_b|\le C\mathcal D_b/\gamma^p\) on its complement, then
\[
 \left\|e^{\Psi/2}\sum_b c_bR_b\right\|_{L^2(\mu_{\mathcal R})}
   \le C\gamma^{-p-1}M_4^{1/4}
               e^{\mathcal H(2t)/4}\|c\|_1 .
 \tag{V47.14}\label{r47:weighted-error}
\]
Here \(c_b\) are arbitrary deterministic complex coefficients
and all marks belong to the same chosen anchor family.
\end{proposition}
\begin{proof}
Cauchy--Schwarz bounds the first left side by
\((\E e^{2\Psi})^{1/2}
  \mu(\mathcal D_b>u_b,\ b\in\mathcal B)^{1/2}\).
Use \eqref{r47:mgf} and \eqref{r47:joint-eq}.
The same inequality bounds the second left side by
\((\E\mathcal D_b^{4r})^{1/2}(\E e^{2\Psi})^{1/2}\).
Equation~\eqref{r47:mgf} gives \eqref{r47:weighted-moment}.
Apply that equation with \(r=1\) to each weighted error,
take a square root, and use the \(L^2\) triangle inequality.
No conditioning on an all-good event or change of marginal
has occurred.
\end{proof}

For V43's moving cutoff
\(d_\gamma=\gamma^{-2/3}/2048\), the rough-return factor in
\eqref{r47:weighted-event} is explicitly
\[
 \widetilde q_{\gamma,K}(d_\gamma)=
  \min\left\{1,\exp\left(
       \frac{1024B^*-\gamma^{1/3}/81920}{K}\right)\right\}.
 \tag{V47.15}\label{r47:moving}
\]
It is stretched exponentially small in the shape window,
per marked rough anchor, before the displayed source cost
is paid. If the source support grows, that cost must still
be paid; this formula does not say the whole lattice is good.
The weighted norm in \eqref{r47:weighted-error} is ordinary
\(L^2\). An arbitrary positive source tilt has not been shown
reflection positive, so it is not assigned an OS norm.

\subsection{Improved native and physical vacuum source errors}

Use the same energy observables, rough-sector definitions,
reflection-matched packing, and exact projected sources as
V38. The already-proved pointwise errors are
\[
 r_b^{(j)}=H_b-h_\gamma-K_b^{(j)},\qquad
 r_b^{(j)}=0\text{ on }G_b^c,\qquad
 |r_b^{(j)}|\le C\mathcal D_b\gamma^{-2-j}
             \text{ on }G_b,\qquad j=0,1.
 \tag{V47.16}\label{r47:inherited-error}
\]
No coefficient of that local Laplace expansion is recalculated
or claimed anew.

\begin{corollary}[Thermal-scale physical source comparison]
\label{r47:OS}
For these sources, with \(K=32\) in \eqref{r47:budgets},
\[
 \|\mathcal E_c^\circ-(K_c^{(j)})^\circ\|_{\rm OS}
      \le C\sqrt{M_2}\,\gamma^{-3-j}\|c\|_1,\qquad j=0,1.
 \tag{V47.17}\label{r47:OS-eq}
\]
In the original fixed-\(L,\gamma\) physical vacuum path law of
V39, put
\[
 B^*_{\infty}=66+\frac{9}{2L}\log\gamma
\]
and let \(M_{2,\infty}\) be \eqref{r47:budgets} with this value.
Then the unchanged physical vectors satisfy
\[
 \|\psi_{\mathcal E_c}
           -\psi_{\widehat K_c^{(j)}}\|
       \le C\sqrt{M_{2,\infty}}\,
                    \gamma^{-3-j}\|c\|_1 .
 \tag{V47.18}\label{r47:vacuum}
\]
Both comparisons have constants uniform on the corresponding
shape window; their powers are \(\gamma^{-3}\) and
\(\gamma^{-4}\), without the earlier factor \(1+\log\gamma\).
\end{corollary}
\begin{proof}
The second-moment bound is now
\(\E_{\mu_{\mathcal R}}\mathcal D_b^2\le M_2/\gamma^2\).
Insert it into \eqref{r47:inherited-error}.
V37's exact conditional projection \eqref{r37:projection}
and the proof of \eqref{r38:OS-bound} then bound the centered
OS norm by the ordinary \(L^2\) norm of \(\sum_b c_br_b^{(j)}\).
The deterministic offset is subtracted globally; rough
predictors are retained exactly. This proves
\eqref{r47:OS-eq} by the triangle inequality.

For the physical vacuum, take the dyadic temporal limit at
fixed \(L,\gamma\), precisely as in Lemma~\ref{r39:moment-limit}.
The augmented slab law converges in total variation.
The local defect squared is bounded at this fixed regulator,
so its expectation converges. The constants satisfy
\(B^*_{L,T,\gamma}\to B^*_\infty\); therefore
\(\E_{\nu_\gamma}\mathcal D_b^2\le M_{2,\infty}/\gamma^2\).
The same holds for the reflected witness.
The exact source identity and conditional Jensen bound in
\eqref{r39:vector-Jensen}, followed by
\eqref{r47:inherited-error}, prove \eqref{r47:vacuum}.
This uses no uniform temporal convergence rate.
\end{proof}

For an \(\ell^2\)-normalized bank of \(m\) physical columns,
V39's synthesis error may consequently be replaced by
\[
 \delta^*_{j,\gamma,m}
       =C\sqrt{mM_{2,\infty}}\,\gamma^{-3-j}.
 \tag{V47.19}\label{r47:synthesis}
\]
Its already-proved comparison \eqref{r39:Gram} then holds
for every positive contraction on the same physical Hilbert
space, including the original \(P^n\) and spectral projectors,
with this improved \emph{absolute} error.
It still gives no lower source Gram bound, relative
low-energy exclusion, or decay in \(n\).
A physical source normalization multiplies these errors
by its absolute magnitude. Inverse powers of a logarithmically
growing coupling do not automatically pay powers of \(a^{-1}\).

\subsection{What has moved, and what the next step cannot assume}

The new result is an original-law, arbitrary-support thermal
defect budget and its explicitly source-weighted remainder
estimate. It removes an avoidable normalizer loss in the
specified growing-regulator window and improves existing
physical source comparisons. It does not integrate the
whole V46 corridor formula for free: that expansion still has
graph-dependent remainders and assumes bounded rescaled
endpoint connections. Small local action does not control
all winding holonomies. In particular the flat toron sectors
retained in Lemma~\ref{r47:comb} cannot be replaced by a
near-identity global endpoint gauge.

The next useful target is a genuinely connected, centered
original-law estimate retaining the predictor covariance,
or an explicit treatment of compact holonomy sectors in
the growing-corridor expansion. Repeating the positive
moment argument cannot supply either one. V46's signed
conditional covariance is entirely consistent with
\eqref{r47:mgf}: a one-sided upper budget does not determine
the sign or spatial decay of centered covariances.

The exact diagnostic
\path{verify_ym_restart_v47_thermal_joint_budgets.py}
checks rectangular tree/face counts, triangular noncommuting
coordinate identities, the boundary mismatch, source-tail
constants, and shifted-exponential moment algebra.
These finite checks supplement the proofs; they do not
verify a continuum limit, reflection positivity of a tilted
law, or the infinite-volume mass gap.
The interacting four-dimensional continuum construction
and its strictly positive physical Hamiltonian gap remain open.

% BEGIN CONSOLIDATED V48 APPEND
\clearpage
\section{V48: Holonomy-resolved corridor normalization and source returns}
\label{r48:start}

V47 pays local defect moments under the original Wilson marginal,
but those moments do not localize winding holonomies. We now retain
such holonomies in V45's actual unpinned corridor integral.
The new result is a complete fixed-graph leading normalizer and
source return on matched left/right exteriors, including the change
in the dimension of the collective minimum. In particular, a flat
noncentral matched holonomy has an \(S^1\), not \(S^3\), family of
minima. Its normalized corridor integral is of order
\(\gamma^{-1}\) relative to coherence, even though all native
six-link coefficient defects can be exactly zero.
We then prove a uniform near-central crossover for the normalizer
and total-energy moments on the matched flat family, without
bounding the thermal parameter \(\gamma|\theta|^2\).

The geometry of parallel sections and the method of Laplace
expansion about a manifold of minima are established methods,
not new general theories here. The archive audit specifically
distinguishes f14 V54's regular and central strata of the
\emph{full flat Wilson connection}, f12 V9's charged flat Bloch
identity on a different restricted fibre, f12 V86's conditional
local-minimum certificate, and archive \(C\) V71's anchored
tree alignment and winding exception. V42's finite-holonomy
normalizer is for a \emph{pinned six-spin block}; it does not
give the unpinned corridor exponent below. V1's coherent-root
Laplace method and V46's bounded thermal-connection expansion
are also inherited, not repeated as new results.

For primary-source context, the quadratic connection form and
parallel-section kernel are discussed in Bandeira--Singer--Spielman,
\emph{A Cheeger Inequality for the Graph Connection Laplacian},
Section 1, \url{https://arxiv.org/html/1204.3873},
and Gao--Brodzki--Mukherjee,
\emph{The Geometry of Synchronization Problems and Learning Group Actions},
\url{https://arxiv.org/html/1610.09051v3}.
We use the finite-dimensional geometric formulation, with a direct
proof in the normalization needed here. No synchronization
algorithm, Cheeger estimate, or graph spectrum is identified
with the physical Yang--Mills Hamiltonian.
The regular companion and broad literature checkpoints remain V50.

\subsection{The complete minimum and its normal Hessian}

Let \(G=(V_G,E_G)\) be a fixed finite connected simple graph,
\(n=|V_G|\ge2\). Give each edge one orientation and an actual
orthogonal transport \(T_{ij}\), with
\(T_{ji}=T_{ij}^{-1}\). For unit real quaternions define
\[
 \begin{gathered}
 W_{ij}(z)=1-z_i\cdot T_{ij}z_j
                =\tfrac12|z_i-T_{ij}z_j|^2,\qquad
 S_T=\sum_{ij\in E_G}W_{ij},\\
 Q_{\gamma,G}(T)=\int_{(S^3)^n}e^{-\gamma S_T(z)}\,dH(z),
 \qquad
 d\lambda_{\gamma,T}=Q_{\gamma,G}(T)^{-1}e^{-\gamma S_T}dH .
 \end{gathered}
 \tag{V48.1}\label{r48:law}
\]
Every Haar factor has mass one; the round volume of \(S^3\)
is \(v_3=2\pi^2\). On \(G_k\), these are exactly V45's
conditional integral and incident original Wilson plaquettes,
with the scalar \(e^{-\gamma |E_G|}\) retained in \(Q\).

Choose a rooted spanning tree and remove its transports by
vertex orthogonal changes of variables. Write \(\widetilde T_e\)
for the remaining transports, so every tree edge has identity
transport. Put
\[
 F=\bigcap_{e\notin\mathcal T}
                  \ker(I-\widetilde T_e),\qquad
 r=\dim F,\qquad s=r-1,
 \tag{V48.2}\label{r48:F}
\]
and assume \(r\ge1\).
Equivalently \(F\) is the common fixed space of all based loop
transports. Changing the root gauge conjugates this space and
does not change its dimension. The assumption is automatic for
matched left/right transports \(T_e z=u_ezu_e^{-1}\), since
the real scalar quaternion is a parallel section.

Let \(L_0\) be the ordinary unweighted graph Laplacian and
\(\det' L_0\) its product of nonzero eigenvalues. In tree gauge
every \(\widetilde T_e\) is the identity on \(F\) and preserves
\(F^\perp\). Let \(L_\perp\) be its connection Laplacian on
\((F^\perp)^n\):
\[
 \langle x,L_\perp x\rangle
     =\sum_{ij\in E_G}|x_i-\widetilde T_{ij}x_j|^2.
 \tag{V48.3}\label{r48:Lperp}
\]
For \(r=4\), it is the empty operator with determinant one.

\begin{lemma}[All minima, not only a local critical orbit]
\label{r48:minima}
The complete zero set of \(S_T\), in tree gauge, is
\[
          \mathcal M=\{z_i=u\text{ for every }i:
                                      u\in F,\ |u|=1\}.
 \tag{V48.4}\label{r48:M}
\]
It is an embedded \(S^{r-1}\) with induced volume
\(n^{s/2}|S^s|\), where \(|S^0|=2\).
The Hessian has kernel exactly \(T\mathcal M\). Its
positive normal determinant is constant on \(\mathcal M\) and equals
\[
       \Delta_T=(\det' L_0)^s\det L_\perp>0.
 \tag{V48.5}\label{r48:det}
\]
Its normal dimension is \(d_T=3n-r+1=3n-s\).
\end{lemma}
\begin{proof}
Every summand of \(S_T\) is nonnegative. Vanishing forces
\(z_i=\widetilde T_{ij}z_j\) on every edge. Tree edges first
force all spins to equal one vector \(u\); the chords then
force \(u\in F\). Conversely every such unit vector makes
every summand zero. Thus there are no other global minima.

The diagonal embedding of the unit sphere in \(F\) multiplies
its tangent metric by \(n\), proving the stated volume.
For \(u\in F\cap S^3\), a tangent field \(x_i\perp u\) has
second variation
\[
 \operatorname{Hess}_u S_T[x,x]
             =\sum_{ij}|x_i-\widetilde T_{ij}x_j|^2.
\]
There is no additional first-derivative term: each squared
difference vanishes at the minimum. The orthogonal splitting
\[
 u^\perp=(F\cap u^\perp)\oplus F^\perp
\]
gives \(s\) copies of \(L_0\) and one copy of \(L_\perp\).
The latter has no kernel: a zero-energy field is constant on
tree edges and fixed by all chords, and hence lies in both
\(F\) and \(F^\perp\). The \(s\) constant neutral fields are
exactly the tangent space of \(\mathcal M\).
The other eigenvalues therefore give \eqref{r48:det}.
Rotating \(F\) while acting identically on \(F^\perp\)
commutes with all tree-gauged transports, so this determinant
does not depend on \(u\).
\end{proof}

For a matched adjoint connection,
\(F=\mathbb R1\oplus F_{\rm adj}\).
The common fixed space \(F_{\rm adj}\subset\mathbb R^3\)
has dimension \(0,1,\) or \(3\): two independent fixed
vectors would force every rotation to be the identity.
Accordingly the possible \(r\) are \(1,2,4\), with minimum
sets \(S^0,S^1,S^3\). In the \(S^0\) case both antipodal
minima must be included, not one selected branch.

\subsection{The measured normalizer and its energy insertions}

\begin{theorem}[Fixed-connection holonomy-stratum asymptotics]
\label{r48:laplace}
At fixed \(G,T\) satisfying the preceding assumptions,
as \(\gamma\to\infty\),
\[
 Q_{\gamma,G}(T)
 =\frac{n^{s/2}|S^s|}{v_3^n\sqrt{\Delta_T}}
           \left(\frac{2\pi}{\gamma}\right)^{d_T/2}
               \{1+O_{G,T}(\gamma^{-1/2})\}.
 \tag{V48.6}\label{r48:Q}
\]
Under the exact normalized law \eqref{r48:law},
\[
 \E_{\lambda_{\gamma,T}}S_T
     =\frac{d_T}{2\gamma}+O_{G,T}(\gamma^{-3/2}),
 \qquad
 \Var_{\lambda_{\gamma,T}}(S_T)
     =\frac{d_T}{2\gamma^2}+O_{G,T}(\gamma^{-5/2}).
 \tag{V48.7}\label{r48:energy}
\]
These are fixed-graph, fixed-connection statements. The
constants are not asserted uniform when the minimum dimension
changes, a positive normal eigenvalue tends to zero, or the
graph grows.
\end{theorem}
\begin{proof}
Lemma~\ref{r48:minima} identifies the entire compact minimum
set and a strictly positive Hessian on its normal bundle.
In a sufficiently small normal tube, with normal coordinate
\(x\) and induced volume \(d\operatorname{vol}_{\mathcal M}(u)\),
Taylor's theorem gives uniformly in \(u\)
\[
 S_T(u,x)=\tfrac12 x^{\mathsf T}H_u x+O_{G,T}(|x|^3),
 \qquad
 dH=v_3^{-n}\{1+O_{G,T}(|x|)\}\,
                         dx\,d\operatorname{vol}_{\mathcal M}(u).
\]
There are constants \(c,C>0\) with
\(c|x|^2\le S_T(u,x)\le C|x|^2\) on this tube.
Outside it, compactness and the complete zero-set
identification give a strictly positive action minimum.

Scale \(x=\gamma^{-1/2}y\).
The difference between the scaled integrand and its
Gaussian limit is bounded in absolute value by
\[
 C\gamma^{-1/2}(1+|y|^3)e^{-c'|y|^2}
\]
after decreasing the fixed tube radius if necessary.
For example apply the mean-value inequality to the two
exponents; both intermediate quadratic lower bounds are
positive, and the exponent difference is bounded by
\(C\gamma^{-1/2}|y|^3\). The volume-factor error is handled
the same way. Gaussian tails permit extension to the
whole normal fibre. The off-tube contribution is
exponentially small, including after multiplication by
any fixed power of \(\gamma\). Integrating the Gaussian
gives \((2\pi)^{d_T/2}/\sqrt{\Delta_T}\);
Lemma~\ref{r48:minima} supplies the volume factor.

For the two energy assertions apply this same argument
with insertions \(\gamma S_T\) and \((\gamma S_T)^2\),
before dividing by the same normalizer.
Their limiting variable is
\(\tfrac12 y^{\mathsf T}H_u y\), with mean \(d_T/2\)
and variance \(d_T/2\). Polynomial-weighted Gaussian
domination gives \(O_{G,T}(\gamma^{-1/2})\) errors in both
scaled moments. This proves \eqref{r48:energy}, without
differentiating an uncontrolled asymptotic remainder.
\end{proof}

For the identity connection, \(r=4\) and \(d_0=3(n-1)\).
Thus the known coherent-root coefficient is recovered.
At a matched \(r=2\) connection, \(d_T=3n-1=d_0+2\);
at \(r=1\), \(d_T=3n=d_0+3\).
The normalized mean incident energy consequently differs
from coherence by \(1/\gamma+O(\gamma^{-3/2})\) or
\(3/(2\gamma)+O(\gamma^{-3/2})\), respectively.
These additional energy terms come from formerly collective
directions that the loop transports no longer fix.

\subsection{Actual flat matched exteriors on the periodic corridor}

Take \(G_k\) with periods \(\mathbf N=(L,L,T)\).
For \(d\in\{\mu,\rho,\tau\}\), choose
\[
          l_{x,d}=r_{x,d}=u_d=e^{\theta_d\mathbf i}
          \quad\hbox{at every corridor edge}.
 \tag{V48.8}\label{r48:flat}
\]
These are actual retained Wilson links on its lower and upper
planes, not independently assigned abstract matrices.
They can be extended to the full four-dimensional regulator
by the same constant \(u_d\) on every \(d\)-link and identity
on every \(\nu\)-link. This full configuration is flat,
because all these quaternions commute.

The adjoint transport fixes the real two-plane
\(\operatorname{span}\{1,\mathbf i\}\); on the charged
plane \(\operatorname{span}\{\mathbf j,\mathbf k\}\)
it is multiplication by \(e^{2i\theta_d}\).
Define the complex magnetic incidence and Laplacian by
\[
 (B_\theta f)_{x,d}=f_x-e^{2i\theta_d}f_{x+\widehat d},
       \qquad L_\theta=B_\theta^*B_\theta .
 \tag{V48.9}\label{r48:magnetic}
\]
For the allowed momenta \(k_d=2\pi j_d/N_d\),
\[
 \lambda_\theta(k)
        =2\sum_{d=1}^3\{1-\cos(k_d+2\theta_d)\}.
 \tag{V48.10}\label{r48:symbol}
\]
This finite Fourier identity retains the periodic winding
phases; it does not gauge them away by a nonperiodic map.

\begin{corollary}[An explicit finite-holonomy determinant ratio]
\label{r48:ratio}
Assume \(N_d\theta_d\notin\pi\mathbb Z\) for at least one
direction. Then \(r=2\), \(L_\theta>0\), and
\[
 \frac{Q_{\gamma,G_k}(\theta)}{Q_{\gamma,G_k}(0)}
   =\frac{2\tau(G_k)}{\gamma\det_{\mathbb C}L_\theta}
               \{1+O_{G_k,\theta}(\gamma^{-1/2})\},
 \quad
 \tau(G_k)=\frac{\det' L_0}{n}.
 \tag{V48.11}\label{r48:ratio-eq}
\]
The number \(\tau(G_k)\) is equivalently any grounded
Laplacian cofactor. In particular
\[
 -\log\frac{Q_\gamma(\theta)}{Q_\gamma(0)}
   =\log\gamma+
          \log\frac{\det_{\mathbb C}L_\theta}{2\tau(G_k)}
             +O_{G_k,\theta}(\gamma^{-1/2}).
 \tag{V48.12}\label{r48:free-energy}
\]
\end{corollary}
\begin{proof}
Every summand in \eqref{r48:symbol} is nonnegative.
An eigenvalue is zero precisely when
\(k_d+2\theta_d\in2\pi\mathbb Z\) in every direction,
which is possible exactly when every
\(N_d\theta_d\in\pi\mathbb Z\).
Under the stated hypothesis the charged operator is positive.
There are exactly two real parallel sections, including after
periodic tree gauge, so \(r=2\).
The real determinant of the charged operator is
\((\det_{\mathbb C}L_\theta)^2\).
Insert this and \(s=1\) into \eqref{r48:Q}, divide by its
\(s=3\) identity-connection version, and use
\(|S^1|=2\pi\), \(|S^3|=2\pi^2\).
The coefficient simplifies to
\(2\det' L_0/(n\det_{\mathbb C}L_\theta)\).
This gives \eqref{r48:ratio-eq} and its logarithm.
\end{proof}

If all \(N_d\theta_d\in\pi\mathbb Z\), the adjoint winding
transports are trivial. A \emph{periodic orthogonal}
vertex change then removes the adjoint connection, and
\(Q_\gamma(\theta)=Q_\gamma(0)\) exactly.
Nontrivial fundamental central winding holonomies can
therefore be invisible to this matched conditional.
They have not been removed from the full Wilson theory.

\subsection{Conditional plaquette covariances with the holonomy retained}

At noncentral \(\theta\) as above, let
\[
 R_0=B_0L_0^\dagger B_0^{\mathsf T},
       \qquad R_\theta=B_\theta L_\theta^{-1}B_\theta^* .
 \tag{V48.13}\label{r48:projections}
\]
The first is a real edge projection of rank \(n-1\);
the second is a complex Hermitian edge projection of rank \(n\).
Write \((R_0)_{ef}=r_{ef}\), using the same chosen
orientations for the actual plaquette energies \(W_e\).

\begin{theorem}[Native conditional source coefficients]
\label{r48:sources}
For fixed edges \(e,f\), at fixed \(G_k,\theta\),
\[
 \begin{split}
 \gamma\,\E_\theta W_e
       &=\tfrac12 r_{ee}+(R_\theta)_{ee}
                       +O_{G_k,\theta}(\gamma^{-1/2}),\\
 \gamma^2\Cov_\theta(W_e,W_f)
       &=\tfrac12 r_{ef}^{\,2}+|(R_\theta)_{ef}|^2
                       +O_{G_k,\theta}(\gamma^{-1/2}).
 \end{split}
 \tag{V48.14}\label{r48:source-eq}
\]
The expectation is under the complete normalized compact
corridor law. These are not covariances of a replaced
exterior marginal.
\end{theorem}
\begin{proof}
At the parallel minimum \(z_x=1\), the tangent field consists
of a neutral real coordinate \(\xi_x\mathbf i\) and a charged
complex coordinate \(w_x\) on
\(\operatorname{span}\{\mathbf j,\mathbf k\}\).
The normal Hessian is \(L_0\) on mean-zero \(\xi\)
and the realification of \(L_\theta\) on \(w\).
Thus the Gaussian normal limit has independent neutral and
charged parts, with
\[
 \E[(B_0\xi)_e(B_0\xi)_f]=r_{ef},\quad
 \E[(B_\theta w)_e\overline{(B_\theta w)_f}]
                                   =2(R_\theta)_{ef},\quad
 \E[(B_\theta w)_e(B_\theta w)_f]=0 .
\]
The factor two in the complex covariance represents its two
real Gaussian coordinates. The leading insertion is
\[
       \gamma W_e
        \longrightarrow
       \tfrac12\{(B_0\xi)_e^2+|(B_\theta w)_e|^2\}.
\]
Every \(W_e\) is invariant under rotations of the common
fixed two-plane, so this calculation is identical at all
points of the minimum circle. Integrate the circle with
its measured volume instead of selecting one root.

The real Wick identity gives covariance
\(r_{ef}^2/2\) for the neutral squares.
For the proper complex Gaussian, the covariance of
the two half-squared moduli is
\(|(R_\theta)_{ef}|^2\). Their mixed covariance vanishes.
The normal-tube argument of Theorem~\ref{r48:laplace},
with these first and second source insertions, gives the
stated errors under the same normalized compact law.
\end{proof}

Summing over all edges, the two projection ranks recover
\(\gamma\E_\theta S_\theta\to(3n-1)/2\).
Likewise \(\sum_{e,f}|(R_\theta)_{ef}|^2=n\) and
\(\sum_{e,f}r_{ef}^2=n-1\) recover the total energy variance
in \eqref{r48:energy}. These are internal consistency
identities, not an inference of a spatial decay rate.

Here is an exact finite-graph check, not a replacement for
the three-dimensional corridor. On the four-cycle put
identity adjoint transport on three edges and
\(\operatorname{Ad}(\mathbf i)\) on the last.
Then \(\tau=4\), \(\det_{\mathbb C}L_\theta=4\),
\(R_\theta=I_4\), and \(R_0=I_4-\mathbf1\mathbf1^{\mathsf T}/4\).
Thus \(Q_\theta/Q_0\sim2/\gamma\), while the covariance
coefficient matrix has diagonal \(41/32\), off-diagonal
\(1/32\), and eigenvalues \(11/8,5/4,5/4,5/4\).
The coherent coefficients are instead \(27/32\) and
\(3/32\). All matrices remain positive semidefinite.
This comparison is distinct from V46's non-shared
thermal twist with negative conditional off-diagonal entries.

\subsection{A surviving harmonic source term near a central holonomy}

The failure of uniformity at \(\theta=0\) can be evaluated
without estimating a divergent determinant.
Fix a nonzero \(a=(a_1,a_2,a_3)\in\mathbb R^3\) and put
\(\theta=ta\). At a fixed graph, all sufficiently small
nonzero \(t\) have at least one noncentral adjoint winding
holonomy. On its oriented edges define
\[
             h_{x,d}=\frac{a_d}{\sqrt{n\sum_j a_j^2}},
                 \qquad \sum_e h_e^2=1 .
 \tag{V48.15}\label{r48:h}
\]

\begin{proposition}[The rank-one harmonic return]
\label{r48:harmonic}
At this fixed graph,
\[
                 R_{ta}\longrightarrow R_0+hh^{\mathsf T}
                    \quad(t\to0,\ t\ne0).
 \tag{V48.16}\label{r48:jump}
\]
Consequently the iterated source limits, with
\(\gamma\to\infty\) taken first, are
\[
 \begin{split}
 \lim_{t\to0,\ t\ne0}\lim_{\gamma\to\infty}
                     \gamma\E_{ta}W_e
        &=\tfrac32r_{ee}+h_e^2,\\
 \lim_{t\to0,\ t\ne0}\lim_{\gamma\to\infty}
                     \gamma^2\Cov_{ta}(W_e,W_f)
        &=\tfrac32r_{ef}^{\,2}
                  +2r_{ef}h_eh_f+h_e^2h_f^2 .
 \end{split}
 \tag{V48.17}\label{r48:iterated}
\]
At \(t=0\) the last two terms are absent.
\end{proposition}
\begin{proof}
Use the finite normalized Fourier basis on the vertex torus.
For each \(k\ne0\), the rank-one contribution
\((B_{ta}f_k)(B_{ta}f_k)^*/\lambda_{ta}(k)\)
converges to its ordinary zero-twist contribution.
There are finitely many momenta, and their sum tends to \(R_0\).
For \(k=0\), \(f_0=n^{-1/2}\), the edge vector is
\((1-e^{2ita_d})/\sqrt n\); its squared norm is
\(\lambda_{ta}(0)\).
Its normalized vector tends to \(-ih\), up to the immaterial
sign of \(t\). The corresponding projection tends to
\(hh^{\mathsf T}\). This proves \eqref{r48:jump}.
Substitute into \eqref{r48:source-eq}, retaining the
declared order of limits, to obtain \eqref{r48:iterated}.
\end{proof}

The vector \(h\) is divergence free and is orthogonal to
ordinary graph gradients. It is a winding harmonic mode,
not a local plaquette defect. The mean return sums to one.
The covariance return also sums to one: the term involving
\(r_{ef}\) sums to \(2h^{\mathsf T}R_0h=0\), and
\(\sum_{e,f}h_e^2h_f^2=1\).
Thus the two lifted real collective directions have been
tracked in both the normalizer and the connected source.

For fixed \(\gamma<\infty\), compact dominated convergence
instead gives continuity of \(Q_{\gamma,G}(ta)\), and of
the normalized first and second \(W_e\) moments, at \(t=0\).
In particular
\[
 \begin{split}
 \lim_{t\to0,\ t\ne0}\lim_{\gamma\to\infty}
    \gamma(\E_{ta}S_{ta}-\E_0S_0)&=1,\\
 \lim_{\gamma\to\infty}\lim_{t\to0}
    \gamma(\E_{ta}S_{ta}-\E_0S_0)&=0 .
 \end{split}
 \tag{V48.18}\label{r48:noncommuting}
\]
No single fixed-stratum asymptotic may be interchanged
with a holonomy integral merely by continuity at finite coupling.

\subsection{A uniform matched-flat crossover, beyond bounded thermal parameters}
\label{r48:crossover}

On this particular translation-invariant family the preceding
two orders of limits can be joined. This is stronger than
merely evaluating V1's root integral at bounded rescaled
endpoints: the parameter \(\gamma|\theta|^2\) below is
unrestricted. The graph is still fixed.

Set
\[
 \kappa=\gamma n\sum_d(1-\cos2\theta_d),\qquad
 I(\kappa)=\int_0^1e^{-\kappa x}\,dx
             =\frac{1-e^{-\kappa}}{\kappa},\quad I(0)=1.
 \tag{V48.19}\label{r48:kappa}
\]
\begin{theorem}[Uniform near-central normalizer and energy crossover]
\label{r48:uniform}
For every fixed corridor graph there are \(\delta_G>0\),
\(\gamma_G<\infty\), and \(C_G<\infty\) such that, uniformly
for \(\|\theta\|_1\le\delta_G\) and \(\gamma\ge\gamma_G\),
\[
 \frac{Q_{\gamma,G}(\theta)}{Q_{\gamma,G}(0)}
        =I(\kappa)\{1+\varepsilon_{G,\gamma,\theta}\},
 \qquad
 |\varepsilon_{G,\gamma,\theta}|
       \le C_G(\|\theta\|_1+\gamma^{-1/2}).
 \tag{V48.20}\label{r48:uniform-Q}
\]
With \(d_0=3(n-1)\), the same normalized compact law satisfies
\[
 \begin{split}
 \gamma\E_\theta S_\theta
     &=\frac{d_0}{2}+m(\kappa)
                     +O_G(\|\theta\|_1+\gamma^{-1/2}),\\
 \gamma^2\Var_\theta(S_\theta)
     &=\frac{d_0}{2}+v(\kappa)
                     +O_G(\|\theta\|_1+\gamma^{-1/2}),\\
 m(\kappa)&=1-\frac{\kappa}{e^\kappa-1},
 \qquad
 v(\kappa)=1-\frac{\kappa^2e^\kappa}{(e^\kappa-1)^2},
 \qquad m(0)=v(0)=0 .
 \end{split}
 \tag{V48.21}\label{r48:uniform-energy}
\]
There is no bound imposed on \(\kappa\).
\end{theorem}
\begin{proof}
Use a fixed normal tube about the \emph{coherent diagonal}
\(S^3\), not about the changing set of exact minima.
Its coordinates are \(u\in S^3\), \(y_i\perp u\),
\(\sum_i y_i=0\), and
\(z_i=\operatorname{Exp}_u(y_i)\).
Write \(X(u)=u_{\mathbf j}^2+u_{\mathbf k}^2\).
At the diagonal,
\[
 S_\theta(u,0)=n\sum_d(1-\cos2\theta_d)X(u).
\]
The first derivative in every normal direction is zero:
translation invariance makes the tangent gradient identical
at each vertex, and the normal field has zero sum.
The normal Hessian \(H(u,\theta)\) is therefore the actual
quadratic operator for the relative variables, without a
displaced saddle. At \(\theta=0\) it is three copies of
\(L_0\) restricted to mean-zero fields.

This Hessian stays positive uniformly in \(u\) when
\(\|\theta\|_1\) is small. To check the scale explicitly, put
\(c_d(u)=u\cdot T_du\). The quadratic form at the diagonal is
the restriction of
\[
 \sum_{x,d}
 \{c_d(u)(|y_x|^2+|y_{x+\widehat d}|^2)
                         -2y_x\cdot T_dy_{x+\widehat d}\}.
\]
Since \(\|T_d-I\|\le2|\theta_d|\) and
\(|c_d(u)-1|\le2|\theta_d|\), its difference from the
zero-twist form is at most
\(8\|\theta\|_1\sum_x|y_x|^2\).
Thus \(\|\theta\|_1\le\lambda_2(L_0)/16\) suffices
for a lower bound \(\lambda_2(L_0)/2\) on normal fields.

Shrink the fixed tube and \(\delta_G\) further if necessary.
Uniform Taylor bounds now give
\[
 S_\theta(u,y)=S_\theta(u,0)
             +\tfrac12 y^{\mathsf T}H(u,\theta)y+O_G(|y|^3)
\]
and a positive quadratic lower bound for the action
difference inside the tube. Outside the tube, the
zero-twist action has a positive minimum; uniform
continuity keeps the twisted action bounded below by a
positive constant for \(\|\theta\|_1\le\delta_G\).
The normal volume density at zero is unchanged.

Apply the Gaussian estimate from
Theorem~\ref{r48:laplace} \emph{pointwise in \(u\)}, factoring
out \(e^{-\kappa X(u)}\) first. Its relative error is
\(O_G(\gamma^{-1/2})\), uniformly in \(u,\theta\).
The determinant amplitude relative to zero twist is
\[
 a(u,\theta)=
 \sqrt{\frac{\det H(u,0)}{\det H(u,\theta)}}
                 =1+O_G(\|\theta\|_1),
\]
uniformly, with a positive lower bound.
Under normalized round measure on \(S^3\), the squared
length \(X\) in a specified real two-plane is uniform
on \([0,1]\). This follows directly from the polar
coordinates in \(\mathbb C^2\), whose normalized radial
factor is \(dX\). Integrating the normalizer over \(u\)
and dividing by its zero-twist version gives
\eqref{r48:uniform-Q}.
The off-tube error remains relatively negligible:
\(I(\kappa)\ge(1+\kappa)^{-1}\), and
\(\kappa\le C_G\gamma\), so its inverse costs only a
polynomial factor against the fixed exponential action gap.

For completeness the source estimates are proved by
insertions, not by differentiating
\eqref{r48:uniform-Q}. Conditional on \(u\), the leading
normal energy is a half chi-square variable with \(d_0\)
degrees of freedom, independently of \(u\), because its
covariance is \(H(u,\theta)^{-1}\).
The root energy is \(\kappa X\).
Under the root weight \(e^{-\kappa X}dX/I(\kappa)\),
this is an exponential variable truncated to \([0,\kappa]\).
All its fixed moments are uniformly bounded by the
corresponding full exponential moments. The factor
\(a(u,\theta)=1+O_G(\|\theta\|_1)\) consequently changes
its first two moments by \(O_G(\|\theta\|_1)\), uniformly
even when \(\kappa\) diverges. The polynomial-weighted
normal Gaussian remainders are \(O_G(\gamma^{-1/2})\)
by the same domination, including the mixed root-normal
terms. Direct integration of the truncated exponential
gives mean \(m(\kappa)\) and variance \(v(\kappa)\).
Adding the independent limiting normal energy proves
\eqref{r48:uniform-energy}.
\end{proof}

Thus \(m(\kappa),v(\kappa)\to0\) as \(\kappa\to0\) and
both tend to one as \(\kappa\to\infty\).
For \(\theta=ta\to0\),
\(\kappa\sim2\gamma n t^2|a|^2\).
Also \(\det_{\mathbb C}L_{ta}
\sim4t^2|a|^2\det' L_0\), by the finite Fourier product.
The large-\(\kappa\) factor \(I(\kappa)\sim\kappa^{-1}\)
therefore matches \eqref{r48:ratio-eq}.
This establishes a common normalization and total-energy
description through the rank change on the matched flat
family. The explicit appearance of \(\lambda_2(L_0)\)
in the proof prevents a claim of regulator-independent
chart size.

\subsection{The native zero-defect test and the next unpaid integration}

The full flat extension of \eqref{r48:flat} has every original
plaquette equal to \(1\), and hence every native
\(\mathcal D_b=0\), by the witness inequality and nonnegativity.
V44--V45 show that these local coefficient defects do not
depend on any released corridor link. Therefore, after fixing
this retained exterior, \(\mathcal D_b=0\) for \emph{every}
corridor configuration in the conditional integral, not only
at its minimum.

Nevertheless, at a fixed noncentral matched holonomy,
\[
 \gamma(\E_\theta S_\theta-\E_0S_0)\longrightarrow1,\qquad
 Q_{\gamma,G_k}(\theta)/Q_{\gamma,G_k}(0)\longrightarrow0 .
 \tag{V48.22}\label{r48:zero-defect}
\]
Thus a uniform coherent replacement whose error vanishes
solely with the native local defects is false for this whole
corridor. V47's stronger defect tails cannot pay the omitted
winding return. This does not invalidate the pinned local
comparisons of V35--V39: their operands, boundary conditions,
and minimum dimensions are different.

The matched flat exteriors form a special subset of the
retained space. No positive probability, dominating phase,
or weight under the full Wilson marginal is assigned to them
by this calculation. General unmatched left/right exteriors
need not have a parallel section at all, and are not covered
by Theorem~\ref{r48:laplace}. The normalizer and sources have
to be extended jointly to their neighborhoods before any
original-law averaging can be justified.

The crossover in Theorem~\ref{r48:uniform} now retains the compact
collective variable through small charged eigenvalues on this
matched flat family, without a bounded thermal-parameter restriction.
The next target is an extension to general left/right exterior
neighborhoods and growing graph scales, together with a bound on
their actual exterior weight and predictor covariance. V1 and V46
already supply the bounded-parameter root integral; recomputing it
alone is not progress. Normalized sources and off-chart contributions
must stay together in this extension.

The exact diagnostic
\path{verify_ym_restart_v48_holonomy_strata.py}
checks the tree-gauged fixed-space ranks, positive normal
operators, rational gauge covariance, four-cycle determinants,
projection source contractions, and the native Fourier zero-mode
counts, together with exact normal stationarity and the truncated
exponential moment algebra used in the crossover. These are finite
algebraic checks of the displayed
proofs, not numerical evidence for a continuum mass gap.
The original interacting four-dimensional construction,
noncollapsing physical source family, and strictly positive
physical Hamiltonian gap remain unproved.
% END CONSOLIDATED V48 APPEND

% BEGIN CONSOLIDATED V49 APPEND
\clearpage
\section{V49: Growing-corridor holonomy bounds and exact determinant cancellation}
\label{r49:start}

V48 leaves graph-dependent constants in its compact crossover.
We now pay two graph-size costs explicitly on the cubic cofinal
corridors \(G_N=(\mathbb Z/N\mathbb Z)^3\), \(N\ge3\); the native
dyadic subfamily has \(N\ge16\).
First, an energy-relative comparison improves the sufficient
normal-Hessian window from link angle \(O(N^{-2})\) to
\(O(N^{-1})\), that is, to a fixed winding-holonomy window.
Second, an exact transverse product cancels the extensive
determinant factors before any limit. Its finite-size factor
has a uniformly convergent infinite product with a proved
\(O(N^{-2})\) error, including two angular derivatives.

The determinant literature suggested this organization:
F.~Friedli, \emph{The bundle Laplacian on discrete tori},
\url{https://doi.org/10.4171/aihpd/66},
and Y.~Lin, S.~Wan, H.~Zhang,
\emph{Connection Laplacian on discrete tori with converging property},
Theorem 1.3, \url{https://arxiv.org/html/2403.06105}.
Their subject is a connection Laplacian, not an interacting
Wilson continuum measure. We prove the needed single-winding
case directly, including its zero-mode normalization and rate,
rather than importing a general determinant limit as a
nonlinear remainder estimate.

The nonduplication check also retains f12 V15's uniform
small-curvature theorem on its different, gapped restricted
fibre, f9 V68's relative inverse-jet identity, and the
unshifted hyperbolic product used in f16.1, V72.
None gives the growing \emph{unpinned} corridor estimate
below. The scalar cycle product itself is classical.
The new programme-specific results are its normalized
three-dimensional winding factor, angular-source return,
and their combination with an explicit relative Hessian
and Gaussian chart budget. No unchanged determinant
identity is counted as a mass-gap theorem.
The next companion review and regular broad primary-source
sweep are both due at V50.

\subsection{Energy-relative normal stability in a fixed holonomy window}

Use the matched flat transports of V48 with
\[
 \theta_d=\frac{\varphi_d}{2N},\qquad
 T_d=\operatorname{Ad}(e^{\theta_d\mathbf i}),\qquad
 \varphi=(\varphi_1,\varphi_2,\varphi_3).
 \tag{V49.1}\label{r49:angles}
\]
Here \(e^{i\varphi_d}\) is the \emph{charged adjoint}
winding phase; the fundamental quaternion winding angle
is \(\varphi_d/2\).
For each collective orientation \(u\in S^3\), consider tangent
fields \(y_x\perp u\) with \(\sum_x y_x=0\).
Write
\[
 E_0(y)=\sum_{x,d}|y_x-y_{x+\widehat d}|^2
                =\langle y,H_0y\rangle ,
\]
where \(H_0\) is three copies of the scalar graph Laplacian
on mean-zero fields. Let \(H(u,\theta)\) be the actual
normal Hessian at the coherent diagonal used in the
proof of Theorem~\ref{r48:uniform}. The diagonal need not
be a minimum in the root variable \(u\); it is stationary
in these relative variables.

\begin{theorem}[Regulator-independent relative Hessian comparison]
\label{r49:H}
Set
\[
 \eta(\varphi)=\tfrac12\sum_d|\varphi_d|
                          +\tfrac18\sum_d\varphi_d^2.
\]
For all \(N\ge3\), all \(u\in S^3\), and every such field,
\[
 |\langle y,(H(u,\theta)-H_0)y\rangle|
                   \le\eta(\varphi)E_0(y).
 \tag{V49.2}\label{r49:relative}
\]
In particular, if \(\|\varphi\|_1\le1\), then
\[
              \tfrac38H_0\preceq H(u,\theta)
                                     \preceq\tfrac{13}{8}H_0 .
 \tag{V49.3}\label{r49:window}
\]
This window is independent of \(N\) when expressed in
winding phases, not in link angles.
\end{theorem}
\begin{proof}
Let \(P\) be projection onto the charged real two-plane and
let \(J\) be its quarter-turn, extended by zero on its
orthogonal complement. Put
\[
 a_d=1-\cos2\theta_d,\qquad b_d=\sin2\theta_d,\qquad
 X=|Pu|^2.
\]
Then \(T_d=I-a_dP+b_dJ\) and \(u\cdot T_du=1-a_dX\).
Subtract the zero-twist quadratic form from the actual
form displayed in V48. The direction-\(d\) contribution is
\[
 a_d\left\{2\sum_x y_x\cdot Py_{x+\widehat d}
                            -2X\sum_x|y_x|^2\right\}
             -2b_d\sum_x y_x\cdot Jy_{x+\widehat d}.
\]
The absolute value of its symmetric bracket is at most
\(4\|y\|_2^2\).
Since \(y_x\cdot Jy_x=0\), its skew sum equals
\(\sum_x y_x\cdot J(y_{x+\widehat d}-y_x)\), whose absolute
value is at most \(\|y\|_2\sqrt{E_0(y)}\).
Mean-zero Poincare and the exact cubic gap give
\[
 \|y\|_2^2\le\lambda_2^{-1}E_0(y),\qquad
 \lambda_2=4\sin^2(\pi/N)\ge16/N^2.
\]
Thus the relative error is at most
\(\sum_d(4a_d/\lambda_2+2|b_d|/\sqrt{\lambda_2})\).
Use \(a_d\le\varphi_d^2/(2N^2)\) and
\(|b_d|\le|\varphi_d|/N\) to obtain \eqref{r49:relative}.
If \(\|\varphi\|_1\le1\), then
\(\sum_d\varphi_d^2\le1\) and \(\eta\le5/8\).
\end{proof}

The skew term is paid against a discrete gradient before
using Poincare. Bounding it first by a constant times
\(\|\theta\|_1\|y\|_2^2\) would recover V48's unnecessarily
smaller sufficient window. The theorem does not yet
control the normal coordinate volume or the compact
off-tube integral uniformly in \(N\).

\begin{corollary}[A genuine uniform bound for the normal Gaussian reference]
\label{r49:Gaussian}
Let \(Y\) be the centered real Gaussian normal field with
covariance \(\gamma^{-1}H(u,\theta)^{-1}\), under
\(\|\varphi\|_1\le1\). Then
\[
 \E|Y_x|^2\le\frac{13}{2\gamma},\qquad
 \Pr\{\max_x|Y_x|>r\}
                  \le6N^3e^{-\gamma r^2/13},\quad r>0.
 \tag{V49.4}\label{r49:Gaussian-eq}
\]
These estimates concern this explicitly declared Gaussian
reference, not the original compact corridor law.
\end{corollary}
\begin{proof}
Represent momenta by integers
\(\mathcal I_N=\{-\lfloor N/2\rfloor,\ldots,\lceil N/2\rceil-1\}\).
For \(p\ne0\) in \(\mathcal I_N^3\),
\[
 \lambda_0(p)=4\sum_d\sin^2(\pi p_d/N)\ge16|p|^2/N^2.
\]
The shell \(|p|_\infty=m\) has at most \(24m^2+2\le26m^2\)
points. Therefore the scalar mean-zero diagonal Green function
obeys
\[
 G_N(0)=N^{-3}\sum_{p\ne0}\lambda_0(p)^{-1}
       \le\frac1{16N}\sum_{m=1}^{\lfloor N/2\rfloor}26
       \le\frac{13}{16}.
\]
Inverting \eqref{r49:window}, each of the three orthonormal
tangent components at a site has variance at most
\((8/3)G_N(0)/\gamma\le13/(6\gamma)\).
Sum the component variances. If \(|Y_x|>r\), one component
exceeds \(r/\sqrt3\) in absolute value.
The scalar Gaussian tail and a union bound over three
components and \(N^3\) sites prove \eqref{r49:Gaussian-eq}.
No independence of sites is used.
\end{proof}

\subsection{An exact winding determinant with its zero mode removed correctly}

For the remainder of the determinant calculation put
\(\varphi=(\phi,0,0)\), without a small-\(\phi\) restriction.
Let \(L_{N,\phi}\) be the charged complex Laplacian with
link phase \(e^{i\phi/N}\) in direction one.
It is positive definite unless \(\phi\in2\pi\mathbb Z\).
For \(p=(p_2,p_3)\in\mathcal I_N^2\) define
\[
 \omega_N(p)=
   2\operatorname{arsinh}
       \sqrt{\sin^2(\pi p_2/N)+\sin^2(\pi p_3/N)} .
 \tag{V49.5}\label{r49:omega}
\]
Set
\[
 \mathcal P_N(\phi)=
 \prod_{\substack{p\in\mathcal I_N^2\\p\ne0}}
  \left\{1+\frac{1-\cos\phi}{\cosh(N\omega_N(p))-1}\right\}.
 \tag{V49.6}\label{r49:PN}
\]

\begin{theorem}[Exact normalized determinant factorization]
\label{r49:determinant}
For every \(N\ge3\) and real \(\phi\notin2\pi\mathbb Z\),
\[
 \frac{\det_{\mathbb C}L_{N,\phi}}{\det' L_{N,0}}
        =\frac{4\sin^2(\phi/2)}{N^2}\mathcal P_N(\phi).
 \tag{V49.7}\label{r49:det-eq}
\]
The right side extends continuously to zero at central phases.
\end{theorem}
\begin{proof}
At fixed transverse momentum \(p\), the \(N\) eigenvalues
in the winding direction are
\[
       2\cosh\omega_N(p)
                   -2\cos((2\pi j+\phi)/N),\qquad 0\le j<N .
\]
The roots and leading coefficient of the Chebyshev
polynomial give the scalar identity
\[
 \prod_{j=0}^{N-1}
       \{2\cosh\omega-2\cos((2\pi j+\phi)/N)\}
                  =2\{\cosh(N\omega)-\cos\phi\}.
\]
For \(p\ne0\), divide it by the \(\phi=0\) identity.
For \(p=0\), the twisted product is \(4\sin^2(\phi/2)\),
whereas the untwisted product with its one zero eigenvalue
omitted is \(N^2\). The last assertion follows either
from the cycle sine product or by dividing the
\(\phi=0\) identity by \(2\cosh\omega-2\) and taking
\(\omega\to0\). Multiply over all transverse momenta.
The omitted scalar zero mode is not counted as a finite
eigenvalue or as an extra determinant factor.
\end{proof}

This cancels the bulk factors \emph{at finite \(N\)}.
There is no subtraction of two separately divergent
numerical log-determinants.

\subsection{Uniform infinite product and a quantitative finite-size error}

Let
\[
 \mathcal P_\infty(\phi)=
 \prod_{p\in\mathbb Z^2\setminus\{0\}}
  \left\{1+\frac{1-\cos\phi}{\cosh(2\pi|p|)-1}\right\},
 \qquad
 C_*=\frac{16e^{-2}}{(1-e^{-2})^4}.
 \tag{V49.8}\label{r49:Pinf}
\]
\begin{theorem}[Uniform product limit, with angular marks]
\label{r49:limit}
The product is positive and converges uniformly for real
\(\phi\). For all \(N\ge3\),
\[
 1\le\mathcal P_N(\phi),\mathcal P_\infty(\phi)
                         \le e^{C_*(1-\cos\phi)}.
 \tag{V49.9}\label{r49:Pbound}
\]
There are universal finite constants \(C,C_2\) such that
\[
 \begin{split}
 |\log\mathcal P_N(\phi)-\log\mathcal P_\infty(\phi)|
                    &\le\frac{C(1-\cos\phi)}{N^2},\\
 \|\log\mathcal P_N-\log\mathcal P_\infty\|_{C^2(\mathbb R/2\pi\mathbb Z)}
                    &\le\frac{C_2}{N^2}.
 \end{split}
 \tag{V49.10}\label{r49:rate}
\]
The \(C^2\) norm is the maximum of the sup norms through
the second angular derivative.
\end{theorem}
\begin{proof}
Write \(a=1-\cos\phi\in[0,2]\) and
\(f(z,\phi)=\log(1+a/(\cosh z-1))\).
For centered \(p\), the sine lower bound and
\(\operatorname{arsinh}v\ge v/2\) on \(0\le v\le\sqrt2\)
give \(N\omega_N(p)\ge2|p|\).
Consequently
\[
 0\le f(N\omega_N(p),\phi)
        \le\frac{2a e^{-2|p|}}{(1-e^{-2})^2}.
\]
There are at most \(8m\) points with \(|p|_\infty=m\).
Summing the majorant gives \(C_*a\), and proves
\eqref{r49:Pbound} and uniform convergence of the infinite
log-product. Its exponential is strictly positive.

For the rate, put
\(v=\pi|p|/N\) and
\(w=\sqrt{\sum_{j=2}^3\sin^2(\pi p_j/N)}\).
The elementary estimates
\(|\sin x-x|\le|x|^3/6\) and
\(0\le w-\operatorname{arsinh}w\le w^3/6\), with \(w\le v\),
give
\[
 |N\omega_N(p)-2\pi|p||
                   \le\frac{2\pi^3|p|^3}{3N^2}.
\]
Both arguments are at least \(2|p|\). For \(z\ge2\),
\[
 |\partial_z f(z,\phi)|
    =\frac{a\sinh z}{(\cosh z-1)(\cosh z-\cos\phi)}
       \le\frac{2a e^{-z}}{(1-e^{-2})^4}.
\]
Hence the sum of the differences over \(\mathcal I_N^2\setminus0\)
is bounded by
\[
 \frac{4\pi^3a}{3(1-e^{-2})^4N^2}
            \sum_{p\ne0}|p|^3e^{-2|p|}.
\]
For a missing point outside \(\mathcal I_N^2\),
\(|p|\ge N/2\). Its tail contribution is bounded by
\[
 \frac{8a}{(1-e^{-2\pi})^2N^2}
                     \sum_{p\ne0}|p|^2e^{-2\pi|p|}.
\]
Both displayed series converge by the same shell count.
Their sum gives an explicit finite choice of \(C\).

For the angular marks, differentiation gives
\[
 \begin{split}
 \partial_\phi f(z,\phi)&=\frac{\sin\phi}{\cosh z-\cos\phi},\\
 \partial_\phi^2 f(z,\phi)&=
  \frac{\cos\phi}{\cosh z-\cos\phi}
               -\frac{\sin^2\phi}{(\cosh z-\cos\phi)^2}.
 \end{split}
\]
For each of these functions, its magnitude and its first
\(z\)-derivative are bounded by a universal constant times
\(e^{-z}\) on \(z\ge2\), uniformly in real \(\phi\).
For example each differentiated denominator is bounded below
by \(e^z(1-e^{-2})^2/2\); the numerator contains at most
one \(\sinh z\). The same sum and tail argument therefore
proves the second line of \eqref{r49:rate}, and justifies
termwise angular differentiation of the infinite product.
\end{proof}

\subsection{The growing coefficient in the actual corridor normalizer}

Write \(Q_{N,\gamma}(\phi)\) for the original compact
matched-flat corridor integral with
\(\theta=(\phi/(2N),0,0)\). Define the fixed-regulator
weak-coupling coefficient
\[
 A_N(\phi)=
       \lim_{\gamma\to\infty}
          \gamma\,\frac{Q_{N,\gamma}(\phi)}{Q_{N,\gamma}(0)}
       \quad(\phi\notin2\pi\mathbb Z).
\]
V48 proves existence of this limit. Since the corridor
has \(n=N^3\) vertices, \eqref{r48:ratio-eq} and
\eqref{r49:det-eq} give the exact coefficient
\[
 A_N(\phi)=
   \frac1{2N\sin^2(\phi/2)\mathcal P_N(\phi)}.
 \tag{V49.11}\label{r49:A}
\]
In particular
\[
 \frac{e^{-C_*(1-\cos\phi)}}{2\sin^2(\phi/2)}
       \le N A_N(\phi)\le\frac1{2\sin^2(\phi/2)},
 \qquad
 N A_N(\phi)\longrightarrow
 \frac1{2\sin^2(\phi/2)\mathcal P_\infty(\phi)}.
 \tag{V49.12}\label{r49:A-limit}
\]
The convergence has rate \(O(N^{-2})\) uniformly on every
closed angular set separated from \(2\pi\mathbb Z\).

Thus the fixed-regulator coefficient itself is of order
\(N^{-1}\), a dependence hidden inside V48's determinant.
The limit asserted here is
\[
 \lim_{N\to\infty}\lim_{\gamma\to\infty}
       \gamma N\,\frac{Q_{N,\gamma}(\phi)}{Q_{N,\gamma}(0)},
 \tag{V49.13}\label{r49:order}
\]
with the value in \eqref{r49:A-limit}.
This is \emph{not} a theorem along
\(\gamma=\gamma(N)\), and in particular not along a
logarithmic bare-coupling trajectory.

For small nonzero \(\phi\), the right side behaves as
\(2/\phi^2\), since \(\mathcal P_\infty(0)=1\).
This agrees with V48's root factor
\(I(\kappa)\sim\kappa^{-1}\), where
\(\kappa=\gamma N^3(1-\cos(\phi/N))
              \sim\gamma N\phi^2/2\).
Agreement of these coefficients is not a uniform
nonlinear error estimate.

\subsection{An angular source response with its variance term retained}

The finite-size factor also controls an actual normalized
conditional response in the same sequential regime.
Let \(S_\phi(z)\) be the corridor action and put
\(F_{N,\gamma}(\phi)=-\log(Q_{N,\gamma}(\phi)/Q_{N,\gamma}(0))\).
Direct differentiation of the finite compact integral gives
\[
 F_{N,\gamma}'=\gamma\E_\phi S_\phi',\qquad
 F_{N,\gamma}''=\gamma\E_\phi S_\phi''
                              -\gamma^2\Var_\phi(S_\phi').
 \tag{V49.14}\label{r49:response}
\]
The second derivative includes both the explicit second
action derivative and the covariance of the first derivative.
Neither can be discarded before taking the limit.

\begin{proposition}[A normalized angular return after bulk cancellation]
\label{r49:source}
At fixed \(N\), on every closed angular interval separated
from \(2\pi\mathbb Z\), as \(\gamma\to\infty\),
\[
 \begin{split}
 F_{N,\gamma}'(\phi)
     &=\cot(\phi/2)+\partial_\phi\log\mathcal P_N(\phi)
                                  +O_N(\gamma^{-1/2}),\\
 F_{N,\gamma}''(\phi)
     &=-\tfrac12\csc^2(\phi/2)
             +\partial_\phi^2\log\mathcal P_N(\phi)
                                  +O_N(\gamma^{-1/2}).
 \end{split}
 \tag{V49.15}\label{r49:angular}
\]
The interval dependence is included in the constants.
After this fixed-\(N\) limit, the replacement
\(\mathcal P_N\to\mathcal P_\infty\) costs \(O(N^{-2})\).
\end{proposition}
\begin{proof}
The minimum circle is the same fixed
\(\operatorname{span}\{1,\mathbf i\}\cap S^3\)
throughout such an interval. Its positive normal Hessian
has a uniform lower bound at this fixed \(N\).
Both \(S_\phi'\) and \(S_\phi''\) vanish to second order
in the normal displacement from that circle: the action
and its first normal derivative vanish there for every
\(\phi\). One and two angular derivatives of the compact
integral therefore introduce, after normal scaling,
only uniformly dominated quadratic and quartic Gaussian
polynomials. The normal-tube proof of
Theorem~\ref{r48:laplace} consequently gives its leading
coefficient with two angular derivatives and relative
\(O_N(\gamma^{-1/2})\) error. This is a differentiated
integral argument, not differentiation of an unspecified
pointwise remainder.

The leading negative log-ratio is
\[
       \log\gamma+\log(2N)
             +2\log|\sin(\phi/2)|+\log\mathcal P_N(\phi).
\]
Differentiate it twice to get \eqref{r49:angular}.
The final assertion is \eqref{r49:rate}.
\end{proof}

These sources are winding-angle responses of the
\emph{conditional} normalizer. They are not a local
physical spectral gap or an estimate of separated
correlations under the original exterior marginal.

\subsection{Which graph-size losses are now paid}

The relative Hessian window in \eqref{r49:window} is
uniform for fixed winding phases of total magnitude at
most one. The normal Gaussian diagonal and chart-tail
bounds are also uniform, with their explicit \(N^3\)
union cost. The entire determinant's bulk cancellation,
finite-size coefficient and two angular marks are
controlled before taking a regulator limit.

Three different unpaid steps must not be hidden in these
statements. First, \eqref{r49:Gaussian-eq} is not a bound
for the compact Wilson law. Second, uniform positivity
of the quadratic normal operator does not by itself pay
the nonlinear action, coordinate-volume corrections,
or the off-tube part of that law as \(N\) grows.
Third, the matched flat exteriors still have no assigned
probability under the original Wilson marginal.

For example, the Gaussian union bound becomes small only
under a declared inequality comparing \(\gamma r^2\) with
\(3\log N\); it does not establish that inequality on the
physical trajectory. Even after it holds, expanding an
extensive density term and bounding its absolute size
can incur a volume cost. A useful next step must estimate
the \emph{normalized difference or connected remainder},
with its winding background retained, rather than merely
repeat a local Gaussian moment or invoke the already
proved determinant formula.

The exact diagnostic
\path{verify_ym_restart_v49_winding_products.py}
checks rational cycle determinants, full signed-torus
determinants, the normal-Hessian decomposition and its
relative bound, and the diagonal Green-function constants.
The infinite-product convergence and its rate follow
from the analytic majorants above, not from numerical
extrapolation. The nonlinear growing-corridor remainder,
actual exterior predictor covariance, noncollapsing
physical source family, and full interacting continuum
Yang--Mills mass gap remain unproved.
% END CONSOLIDATED V49 APPEND

% BEGIN CONSOLIDATED V50 APPEND
\clearpage
\section[V50: Compact-law infrared bounds and macroscopic strands]{V50: Compact-law infrared control, marked chart tails, and macroscopic strands}
\label{r50:start}

The V50 checkpoint changes the next acquisition order. V49 controlled a
normal Gaussian reference, but its tail bound did not apply to the compact
corridor measure. On the coherent corridor, a classical finite-volume
infrared bound supplies precisely that missing probability input.
We transfer it with the actual normalization, retain the random collective
orientation, and obtain a joint plaquette-source exponential bound and
a source-weighted off-chart estimate. As a further consequence, V45's
long-strand conclusion now holds uniformly in the regulator at a fixed
explicit low-temperature threshold.

The corridor remains the \emph{actual conditional} of V45. This section
does not replace the original exterior distribution by coherent boundary
data. Noncentral winding, general left/right exteriors, and their Wilson
probabilities remain distinct obligations.

\subsection{Scheduled companion and broad literature checkpoint}

Fresh hashes of NS V26, SM--Einstein V72, Shell Mixing V16 and Parallel
YM V5 agree with V45. Their terminal statements and scope passages were
reread. The latest related downloads are also unchanged: Source Selection
V35, Stationarity Closure V38, Exact-Action Bridge V28 and Perturbative
ESM V27. NS's rank-one estimates and SM's fixed-mode many-copy limit
supply no corridor covariance bound. Shell Mixing still separates its
flat coefficient from the noncommuting mixed-curvature problem.
The ESM lesson retained here is to control marks and the normalizer in
one law; no ESM coefficient is imported.

The broad review covered low-temperature spin expansions, infrared
bounds and collective modes, stochastic gauge renormalization, lattice
area laws, multiscale functional inequalities, and physical spectral
bounds. The principal decisions are:
\begin{itemize}
\item A.~Giuliani and S.~Ott,
\emph{Low temperature asymptotic expansion for classical \(O(N)\)
vector models}, \url{https://arxiv.org/abs/2302.07299}:
use Proposition 1.1 and equation (3), which are finite-volume statements.
Their later asymptotic expansion uses a thermodynamic limit followed by
field removal. It is not a periodic winding-normalizer theorem.
\item M.~Biskup, \emph{Reflection Positivity and Phase Transitions
in Lattice Spin Models}, \url{https://arxiv.org/html/math-ph/0610025},
Section 5.2: the Gaussian-domination input requires the specified
ferromagnetic law. It must not be assigned to arbitrary transported
exteriors.
\item J.~Aru and A.~Korzhenkova,
\url{https://arxiv.org/abs/2405.04501}: their large-spin-dimension limit
is not a fixed four-component result. The collective mode cannot be
discarded by identifying the field with a centered GFF.
\item I.~Chevyrev and H.~Shen,
\url{https://arxiv.org/abs/2503.03060}: uniqueness of a gauge-covariant
renormalization for the stochastic three-dimensional theory is not
four-dimensional continuum construction or physical-time contraction.
\item S.~Cao, R.~Nissim and S.~Sheffield,
\url{https://arxiv.org/abs/2505.16585} and
\url{https://arxiv.org/abs/2509.04688}: retain the declared lattice
parameter regimes and the separate mass-gap-to-area-law input.
No weak-coupling continuum conclusion is imported.
\item R.~Bauerschmidt, T.~Bodineau and B.~Dagallier,
\url{https://arxiv.org/html/2307.07619}, and
R.~Abbott, W.~I.~Jay and P.~R.~Oare,
\url{https://arxiv.org/abs/2508.01377}: renormalized-Hessian hypotheses
and actual spectral-moment data remain required. A source bound at one
conditional exterior provides neither.
\end{itemize}
Only the first two references supply the classical input used below.
The remaining items are scope checks, not imported theorems.
The full-text spin-model pages containing the normalization and
Proposition 1.1 were checked; the later expansion statement was screened,
not reproved.

For nonduplication, f14 V62 already contains a sub-Gaussian and quadratic
exponential bound for its \emph{geometric reference law}. That Gaussian
randomization is reused, not claimed new. The reference law is not the
compact Wilson conditional used here. V45 supplies the normalized
strand identity; V49 supplies the explicit cubic Green-function bound.
The new point is their application with a regulator-uniform,
\emph{actual compact-law} input.

\subsection{The finite-volume input in the original corridor normalization}

Let \(N\ge4\) be even, \(n=N^3\), and \(m=3n\). The native dyadic
corridors are included. At coherent exterior the released variables
\(z_x\in S^3\subset\mathbb R^4\) have law
\[
 d\mu_{N,\gamma}(z)=Q_{N,\gamma}(0)^{-1}
       e^{-\gamma S_0(z)}\prod_xdH(z_x),\qquad
 S_0(z)=\frac12\sum_{\{x,y\}\in E}|z_x-z_y|^2 .
 \tag{V50.1}\label{r50:law}
\]
Edges are counted once. Write \(L=B^{\mathsf T}B\) for the scalar
nearest-neighbor Laplacian, \(C=L^\dagger\), and
\(c_N=C_{xx}\). Thus \(C\mathbf1=0\), and V49 proves
\[
                       c_N\le13/16.
 \tag{V50.2}\label{r50:green}
\]
The inverse is on the mean-zero space, not on a pinned graph.

\begin{lemma}[Classical finite-volume domination, with native units]
\label{r50:input}
For every \(f_x\in\mathbb R^4\) with \(\sum_x f_x=0\),
\[
 \E_{\mu_{N,\gamma}}
       e^{\sum_x f_x\cdot z_x}
 \le \exp\left\{\frac1{2\gamma}
                 \sum_{x,y}f_x\cdot C_{xy}f_y\right\}.
 \tag{V50.3}\label{r50:linear}
\]
\end{lemma}
\begin{proof}
This is the finite-volume infrared inequality in Proposition 1.1,
equation (3), of Giuliani--Ott, with spin dimension four, zero external
field, and \(\beta=\gamma\). Indeed,
\(e^{-\gamma S_0}=e^{-\gamma m}
 e^{\gamma\sum_{\{x,y\}}z_x\cdot z_y}\);
the constant and the area-versus-normalized-Haar factors cancel in
the same normalized expectation. The positive graph Laplacian there
is the present \(L\).

Equivalently, Biskup's Gaussian domination gives
\(Z(h)\le Z(0)\) for the shifted-gradient energy
\(\frac{\gamma}{2}\sum_e|(Bz)_e+(Bh)_e|^2\).
Expanding this square and taking \(h=-C f/\gamma\) gives
\eqref{r50:linear}. This uses ferromagnetic reflection positivity on
the even periodic torus, not a convex approximation to the sphere
or V49's normal Gaussian field.
\end{proof}

This lemma is an explicitly credited classical input. In particular,
there is no inference of it from a positive Hessian at a minimum.
The constraint on the source is essential: a constant source couples
to the compact collective variable and is not controlled by \(C\).

\subsection{Uniform compact-law control around a random collective orientation}

Put
\[
 M(z)=n^{-1}\sum_xz_x,\qquad v_x=z_x-M,\qquad
 D(z)=n^{-1}\sum_x|v_x|^2=1-|M|^2.
 \tag{V50.4}\label{r50:collective}
\]
Define \(U=M/|M|\) when \(M\ne0\). The exceptional set \(M=0\)
has probability zero: condition on all but one spin; the remaining
spin would have to equal one specified point of \(S^3\), which has
Haar measure zero. Give \(U\) any value there when making pointwise
statements.

\begin{theorem}[An actual-law collective chart budget]
\label{r50:chart}
For every \(\gamma>0\),
\[
 \begin{split}
 \E D&\le\frac{4c_N}{\gamma}\le\frac{13}{4\gamma},\\
 \E\left[n^{-1}\sum_x|z_x-U|^2\right]
                &\le\frac{8c_N}{\gamma}\le\frac{13}{2\gamma}.
 \end{split}
 \tag{V50.5}\label{r50:mean}
\]
For \(r>0\),
\[
 \mu_{N,\gamma}\{\max_x|v_x|>r\}
       \le8n\exp\{-\gamma r^2/(8c_N)\}.
 \tag{V50.6}\label{r50:v-tail}
\]
Consequently, for \(0<r<1\), the collective chordal chart
\(\mathcal A_r=\{\max_x|z_x-U|\le r\}\) satisfies
\[
 \mu_{N,\gamma}(\mathcal A_r^{\,c})
     \le p_{N,\gamma,r}
     :=\min\{1,\,8N^3e^{-\gamma r^2/26}\}.
 \tag{V50.7}\label{r50:chart-tail}
\]
The direction \(U\) is uniform on \(S^3\); it is not fixed to an
externally chosen orientation.
\end{theorem}
\begin{proof}
In \eqref{r50:linear} take
\(f_y=t(\mathbf1_{y=x}-1/n)e_a\), where \(e_a\) is any of the
four coordinate vectors. This gives
\(\E e^{t v_x^a}\le e^{t^2c_N/(2\gamma)}\).
The mean vanishes, and expansion at \(t=0\) bounds its second
moment by \(c_N/\gamma\). Sum the four components and average
over \(x\) for the first line of \eqref{r50:mean}.
The exact identity
\[
 n^{-1}\sum_x|z_x-U|^2=2(1-|M|)
                       \le2(1-|M|^2)
\]
proves the second line.

Chernoff optimization gives
\(\Pr(|v_x^a|>s)\le2e^{-\gamma s^2/(2c_N)}\).
If \(|v_x|>r\), some component exceeds \(r/2\) in absolute
value. Sum over the four components and all \(n\) sites.
For every configuration, the reverse triangle inequality gives
\(|1-|M||\le\max_x|z_x-M|\), hence
\(\max_x|z_x-U|\le2\max_x|v_x|\).
Apply \eqref{r50:v-tail} at \(r/2\), then use \eqref{r50:green}.
Finally the law and \(M\) are equivariant under simultaneous
\(O(4)\) rotations. The distribution of \(U\) is therefore the
unique rotation-invariant probability on \(S^3\).
\end{proof}

For example, \(\mu(\mathcal A_r^c)\le\epsilon\) is proved whenever
\[
 \gamma r^2\ge26\{3\log N+\log(8/\epsilon)\}.
 \tag{V50.8}\label{r50:scale}
\]
Unlike V49's Gaussian estimate, this includes the complete compact
law, its normalizer, and every configuration outside the chart.
It does not certify \eqref{r50:scale} for a physical bare-coupling
trajectory: its constants and the chosen radius must still be compared
with that trajectory.

The same input bounds two-site differences at every separation:
\[
 \E|z_x-z_y|^2\le\frac{4R_{xy}}{\gamma},\qquad
 R_{xy}=(e_x-e_y)^{\mathsf T}C(e_x-e_y).
 \tag{V50.9}\label{r50:resistance}
\]
This follows by taking four scalar sources \(t(e_x-e_y)\), as above.
In particular \(R_{xy}\le4c_N\) and
\(\E(z_x\cdot z_y)\ge1-13/(2\gamma)\).
This is orientational order of a conditional spin system, not a
physical massless mode or a statement about the full gauge-invariant
Wilson marginal.

\subsection{Joint native plaquette marks and a weighted off-chart estimate}

For an oriented edge \(e=\{x,y\}\), its actual corridor plaquette
deficit is
\[
 W_e=\tfrac12|z_x-z_y|^2,\qquad
 R=BCB^{\mathsf T}.
 \tag{V50.10}\label{r50:edge}
\]
Here \(R\) is an orthogonal edge projection of rank \(n-1\).
For nonnegative marks \(t_e\), let
\(T=\operatorname{diag}(t_e)\) and \(K=T^{1/2}RT^{1/2}\).

\begin{theorem}[Nonlinear compact-law source bound]
\label{r50:marks}
If \(\|K\|<1\), then
\[
 \E_\mu\exp\{\gamma\sum_e t_eW_e\}
                  \le\det(I-K)^{-2}.
 \tag{V50.11}\label{r50:quad}
\]
In particular, if \(\|K\|\le s<1\),
\[
 \log\E_\mu e^{\gamma\sum_e t_eW_e}
     \le \frac{2}{1-s}\sum_e t_eR_{ee}.
 \tag{V50.12}\label{r50:trace}
\]
If \(2\|K\|<1\), the same normalized law satisfies
\[
 \E_\mu\!\left[e^{\gamma\sum_e t_eW_e}
                         \mathbf1_{\mathcal A_r^c}\right]
       \le\det(I-2K)^{-1}\,p_{N,\gamma,r}^{1/2}.
 \tag{V50.13}\label{r50:marked-tail}
\]
\end{theorem}
\begin{proof}
Let \(g_e\) be independent standard Gaussian vectors in \(\mathbb R^4\).
Gaussian integration gives, for every compact spin configuration,
\[
 e^{\frac{\gamma}{2}\sum_e t_e|(Bz)_e|^2}
   =\E_g e^{\sqrt\gamma\sum_e\sqrt{t_e}\,g_e\cdot(Bz)_e}.
\]
For each \(g\), the site source
\(\sqrt\gamma B^{\mathsf T}T^{1/2}g\) has zero sum. Applying
\eqref{r50:linear}, followed by Tonelli and a finite Gaussian integral,
bounds the expectation by
\[
 \E_g e^{\frac12\sum_{a=1}^4(g^a)^{\mathsf T}K g^a}
                 =\det(I-K)^{-2}.
\]
This reuses the Gaussian randomization from f14 V62 with the new
compact-law input. It is not a Gaussian replacement of \(\mu\).
For eigenvalues \(0\le\lambda\le s\),
\(-\log(1-\lambda)\le\lambda/(1-s)\); sum and use
\(\Tr K=\sum_e t_eR_{ee}\).
Cauchy--Schwarz, \eqref{r50:quad} at marks \(2t_e\), and
\eqref{r50:chart-tail} prove \eqref{r50:marked-tail}.
\end{proof}

For marks on \(k\) edges with \(0\le t_e\le s<1/2\),
the determinant in \eqref{r50:marked-tail} is at most
\(\exp\{2ks/(1-2s)\}\), since \(R\preceq I\).
Thus its fixed-support prefactor is explicitly independent of \(N\).

For a single edge, \(K\) has its only possible nonzero eigenvalue
\(t_eR_{ee}\), so
\[
 \E e^{t\gamma W_e}\le(1-tR_{ee})^{-2},\qquad
                 \E(\gamma W_e)\le2R_{ee}.
 \tag{V50.14}\label{r50:one-edge}
\]
The last conclusion follows by the right derivative at zero, where
both normalizers equal one. Derivatives at nonzero sources or a
connected covariance estimate do \emph{not} follow merely by
differentiating an inequality.

A polynomial in finitely many scaled deficits also has a paid
off-chart bound: for \(a\ge0\),
\(a^k\le k!\delta^{-k}e^{\delta a}\).
Apply this to each factor and then \eqref{r50:marked-tail}, with
marks increased by the chosen positive \(\delta\)'s and with
\(2\|K\|<1\) verified for those increased marks. This provides actual
source-weighted tail control at fixed marked support. For marks on
the entire corridor the determinant can still cost volume;
the theorem has not made an extensive density remainder small.

If actual retained transports are a periodic pure gauge,
\(T_{xy}=g_xg_y^{-1}\) with \(g_x\in O(4)\), the substitution
\(z_x=g_xs_x\) preserves product sphere measure and sends every
\(|z_x-T_{xy}z_y|^2\) to \(|s_x-s_y|^2\).
All these conclusions hold in that transported frame.
This requires trivial transport holonomy on every loop.
It does not remove the noncentral winding treated in V48--V49.

\subsection{A regulator-uniform macroscopic-strand consequence}

Retain exactly V45's positive coherent diagram law and its variables
\(M_{\rm occ}=\sum_em_e\), \(S_2=\sum_C|C|^2\), \(L_*=\max_C|C|\).
The notation \(M_{\rm occ}\) distinguishes occurrences from the
spin average \(M\). Put
\[
 B_{\rm spin}=\sum_{\{x,y\}}z_x\cdot z_y,\qquad
 J=\frac12\sum_{\{x,y\}}(z_xz_y^{\mathsf T}+z_yz_x^{\mathsf T}),\qquad
 A=4\Tr(J^2)-B_{\rm spin}^2 .
\]
No diagram identity is changed.

\begin{proposition}[Uniform coherent tensor concentration]
\label{r50:tensor}
At every \(N,\gamma\) as above,
\[
 \E B_{\rm spin}\ge m\left(1-\frac{13}{2\gamma}\right),\qquad
 \E A\ge m^2\left(3-\frac{65}{\gamma}\right).
 \tag{V50.15}\label{r50:tensor-eq}
\]
\end{proposition}
\begin{proof}
Regularity and \(\sum_xv_x=0\) cancel the mixed average/fluctuation terms:
\[
 J=mMM^{\mathsf T}+J_v,\qquad
 J_v=\frac12\sum_{\{x,y\}}(v_xv_y^{\mathsf T}+v_yv_x^{\mathsf T}),\qquad
 \|J_v\|_F\le3\sum_x|v_x|^2=mD.
\]
With \(J_*=mUU^{\mathsf T}\), it follows that
\(\|J-J_*\|_F\le2mD\).
Also
\(m-B_{\rm spin}=S_0\le6\sum_x|v_x|^2=2mD\).
The elementary bounds \(\|J\|_F\le m\) and
\(|B_{\rm spin}|\le m\) give
\[
 |4\Tr J^2-4m^2|\le16m^2D,\qquad
 |B_{\rm spin}^2-m^2|\le4m^2D.
\]
Hence \(A\ge3m^2-20m^2D\). Take expectations and use
\eqref{r50:mean}. These inequalities also hold on \(M=0\) with any
chosen \(U\).
\end{proof}

\begin{theorem}[Macroscopic strands at a fixed coupling threshold]
\label{r50:strands}
For every even \(N\ge4\) and every \(\gamma\ge130\), under the
unchanged coherent diagram law of V45,
\[
       \mathbb P_\gamma\{L_*\ge\gamma m/6\}\ge1/162.
 \tag{V50.16}\label{r50:strands-eq}
\]
Both constants are independent of \(N\).
\end{theorem}
\begin{proof}
Write \(q=\gamma m\). V45's exact identities and
\eqref{r50:tensor-eq} yield
\[
 \frac{\mathbb E S_2}{q^2}
 =\frac{\gamma\E B_{\rm spin}}{q^2}
       +\frac{\E A}{9m^2}
 \ge\frac13-\frac{65}{9\gamma}\ge\frac5{18}.
\]
The first summand was discarded only after its nonnegativity was
verified by \(\gamma\ge130\) and \eqref{r50:tensor-eq}.
Furthermore
\[
 \mathbb E M_{\rm occ}\le q,\qquad
 \mathbb E M_{\rm occ}^4\le q^4+6q^3+7q^2+q\le2q^4.
\]
The first two inequalities are V45's factorial-moment identities
and \(|B_{\rm spin}|\le m\); the last uses \(q\ge8\).
On the complement of \(E=\{L_*\ge q/6\}\),
\(S_2\le(q/6)M_{\rm occ}\), while always \(S_2\le M_{\rm occ}^2\).
Cauchy--Schwarz therefore gives
\[
 \frac5{18}\le
 \frac{\mathbb E S_2}{q^2}
 \le\frac16+\sqrt{2\mathbb P_\gamma(E)}.
\]
Subtract \(1/6\) and square.
\end{proof}

This closes the graph-size loss in this particular V45 consequence.
It rules out a regulator-uniform short-strand description of the
coherent corridor even at a fixed sufficiently large \(\gamma\).
The strand measure is not the physical transfer spectrum, and the
coherent exterior has not acquired a probability under the full
Wilson marginal by this argument.

\subsection{Paid and unpaid nonlinear work after the checkpoint}

For the coherent, or periodically pure-gauge, \emph{actual} conditional,
the collective chart probability and fixed-support source-weighted
off-chart contribution are now bounded without a Gaussian
approximation. The random orientation is retained. The result holds
at finite \(N,\gamma\), with the explicit joint window
\eqref{r50:scale}; it is not an iterated-limit assertion.

What remains is a normalized \emph{on-chart} nonlinear estimate and
its connected source return, extended beyond the pure-gauge family.
The event \(\mathcal A_r\) uses the extrinsic spin average; a geodesic
normal-coordinate expansion must still supply its own coordinate
change and coarea density.
Noncentral winding invalidates the global change to coherent spins;
general unmatched left/right fields may also remove the collective
minimum. The original exterior weight and predictor covariance have
not been estimated here. In particular a linear-source moment bound
at zero source must not be renamed a uniform tilted-Hessian bound,
and orientational order must not be renamed a physical gap.

The next useful import from the low-temperature literature is its
integration-by-parts treatment of normalized local observables, with
the collective variable and winding kept explicitly. Before using
an infinite-volume selected-state expansion, its relation to these
finite periodic conditional sources must be proved. Recomputing a
fixed-volume coefficient or assuming that distant connected errors
are summable would evade the remaining problem.

The exact diagnostic
\path{verify_ym_restart_v50_compact_infrared.py}
checks the Laplacian/projection normalization, quadratic determinant
assembly, collective tensor inequalities and uniform strand constants.
It does not numerically certify the external infrared theorem or
a continuum limit. The next companion checkpoint is V55; the next
regular broad literature sweep is V60. The full interacting
four-dimensional Yang--Mills construction and positive physical
Hamiltonian mass gap remain unproved.
% END CONSOLIDATED V50 APPEND


% BEGIN CONSOLIDATED V51 APPEND
\clearpage
\section[V51: Collective identities and normalized energy returns]{V51: Chart-free collective identities, explicit normalizers, and an energy return}
\label{r51:start}

V50 paid compact-law infrared moments but did not determine their
three-tangent-component energy coefficient or a connected response.
We now use a globally smooth alignment vector field on the original
product of spheres. Its divergence includes the derivative of the
spin average. This term removes exactly three collective dimensions;
it would be wrong to treat that average as a fixed external field.

The result is a volume-uniform estimate for the mean energy density,
an explicit extensive logarithmic-normalizer remainder, and a bound
on the variance of the total energy. We also transfer the normalized
identity to finitely many positive plaquette marks. None requires a
chart cutoff, a chosen orientation, or an unproved coarea formula.
The connected estimate concerns the uniform energy source, not an
absolutely summable spatial covariance kernel.

The inputs are V50's actual coherent conditional, its credited
infrared inequality, and V1's fixed-graph leading compact integral.
The elementary spherical integration-by-parts rule is not claimed new.
V38 already proved mixed derivatives for a fixed 18-dimensional block,
and V48 already proved a fixed-graph winding crossover. Neither
provides the explicit growing-corridor estimates below. Searches of
the earlier sphere/Stein and equipartition sections found single-link,
fixed-block, or chart-reference results, which are not repeated here.
Biskup's Section 5.2,
\url{https://arxiv.org/html/math-ph/0610025}, was rechecked for the
inherited Gaussian-domination scope. No later selected-state expansion
from Giuliani--Ott is imported.

\subsection{A smooth vector field which retains the collective mode}

Keep \(N\ge4\) even, \(n=N^3\), \(m=3n\), and the law
\(\mu=\mu_{N,\gamma}\) of \eqref{r50:law}. Write
\[
 X_e=\gamma W_e,\qquad Y=\sum_e X_e=\gamma S_0,\qquad
 d=3(n-1),\qquad \alpha=d/2,\qquad c=c_N\le13/16.
 \tag{V51.1}\label{r51:data}
\]
The edge resistance \(R_{ee}\) is as in V50; in particular
\(\sum_eR_{ee}=n-1\). In this section \(L=B^{\mathsf T}B\) is
unrooted and \(C=L^\dagger\). Rooted matrices will be explicitly
identified when the leading normalizer is recalled.

Retain \(M=n^{-1}\sum_x z_x\), \(v_x=z_x-M\), and \(D=1-|M|^2\).
Define, at every configuration including \(M=0\),
\[
 a_x=z_x\cdot M,\qquad P_x=I-z_xz_x^{\mathsf T},\qquad
 \mathcal V_x=P_xM,\qquad
 \mathcal V F=\sum_x\mathcal V_x\cdot\nabla_x F.
 \tag{V51.2}\label{r51:field}
\]
Here \(\nabla_x\) is the intrinsic spherical gradient. The vector
field is polynomial and globally smooth. It aligns spins with
their actual mean; it does not fix the direction \(U\).

\begin{lemma}[Exact collective divergence and energy drift]
\label{r51:geometry}
Set
\[
 b_e=D+\frac{|v_x|^2+|v_y|^2}{2}\quad(e=\{x,y\}),\qquad
 T=\sum_eX_eb_e,\qquad \rho=T-3nD.
 \tag{V51.3}\label{r51:defect}
\]
Then \(b_e\ge0\), and the following identities hold pointwise:
\[
 \begin{split}
 a_x&=1-\tfrac12(D+|v_x|^2),\\
 \operatorname{div}\mathcal V&=-d+3nD,\\
 \mathcal V X_e&=-(2-b_e)X_e,\qquad
 \mathcal VY=-2Y+T.
 \end{split}
 \tag{V51.4}\label{r51:geometry-eq}
\]
Consequently, for every smooth real \(F\) on \((S^3)^n\),
\[
 \E\{\mathcal VF+(2Y-d-\rho)F\}=0.
 \tag{V51.5}\label{r51:ibp}
\]
\end{lemma}
\begin{proof}
The first formula follows by expanding \(|z_x-M|^2\).
For fixed \(M\), the divergence of \(P_xM\) is \(-3z_x\cdot M\).
When \(z_x\) varies, the actual mean has derivative \(I/n\);
its additional tangential trace is \(3/n\).
Thus the full divergence is
\(3-3\sum_xa_x=3-3n|M|^2=-d+3nD\).
This can also be checked by taking the trace, against \(P_x\),
of the ambient derivative of \(M-(z_x\cdot M)z_x\).

Since \(W_e=1-z_x\cdot z_y\), direct differentiation gives
\(\mathcal VW_e=-(a_x+a_y)W_e=-(2-b_e)W_e\).
Sum the edge identities. Finally integrate the divergence of
\(F\mathcal V e^{-Y}\) on the compact boundaryless product of
spheres and divide by its actual partition function.
No integration through a chart boundary occurs.
\end{proof}

Taking \(F=1\) gives the exact mean-energy identity
\[
             2\E Y=d-3n\E D+\E T.
 \tag{V51.6}\label{r51:energy-exact}
\]
For any smooth \(F\), subtracting its mean in \eqref{r51:ibp} gives
\[
       2\Cov(Y,F)=-\E\mathcal VF+\Cov(\rho,F).
 \tag{V51.7}\label{r51:connected}
\]
The final covariance is the explicit nonlinear return. It has not
been set to zero by a Ward identity.

\subsection{Moment bounds for the actual nonlinear defect}

For an integer \(p\ge1\) put
\[
                         A_p=(4\,p!\,2^p)^{1/p}.
 \tag{V51.8}\label{r51:Ap}
\]
All norms in the next lemma use the full normalized compact law.

\begin{lemma}[An extensive defect with explicit \(L^1\) and \(L^2\) costs]
\label{r51:moments}
For \(Z_x=\gamma|v_x|^2/2\),
\[
 \|X_e\|_p\le A_pR_{ee},\quad
 \|Z_x\|_p\le A_pc,\quad
 \|D\|_p\le\frac{2A_pc}{\gamma},\quad
 \|b_e\|_p\le\frac{4A_pc}{\gamma}.
 \tag{V51.9}\label{r51:moment-eq}
\]
In particular,
\[
 \begin{split}
 \E T&\le\frac{128c(n-1)}{\gamma},\\
 \|T\|_2&\le\frac{64\sqrt6\,c(n-1)}{\gamma},\\
 \|\rho\|_2&\le
   \frac{c}{\gamma}\{64\sqrt6(n-1)+24\sqrt2\,n\}
                  \le\frac{160n}{\gamma}.
 \end{split}
 \tag{V51.10}\label{r51:defect-bound}
\]
Also \(\E D\le4c/\gamma\), as already proved in V50.
\end{lemma}
\begin{proof}
V50 gives
\(\E e^{tX_e}\le(1-tR_{ee})^{-2}\).
The same four-component Gaussian randomization, now applied to the
zero-mean source defining \(v_x\), gives
\(\E e^{tZ_x}\le(1-tc)^{-2}\).
At \(t=1/(2R_{ee})\) and \(t=1/(2c)\), respectively, the right
sides equal four. The pointwise inequality
\(u^p\le p!e^u\) gives the first two bounds.
This uses a positive exponential estimate, not differentiation
of that inequality at a nonzero source.

Since \(D=2(\gamma n)^{-1}\sum_xZ_x\) and
\(b_e=D+(Z_x+Z_y)/\gamma\), Minkowski gives the remaining bounds.
Hölder and the resistance trace then give
\[
       \|T\|_p\le
       \frac{4A_{2p}^{\,2}c}{\gamma}\sum_eR_{ee}
       =\frac{4A_{2p}^{\,2}c(n-1)}{\gamma}.
\]
Here \(A_2^2=32\), \(A_4^2=16\sqrt6\), and
\(3n\|D\|_2\le24\sqrt2\,cn/\gamma\).
For the last numerical bound use
\(\sqrt6<5/2\), \(\sqrt2<3/2\), and
\(196(13/16)=637/4<160\).
\end{proof}

\begin{theorem}[Uniform mean energy in the original compact law]
\label{r51:mean}
For every allowed \(N\) and every \(\gamma>0\),
\[
 -\frac{6cn}{\gamma}
 \le \E Y-\alpha
 \le \frac{64c(n-1)}{\gamma},\qquad
                 |\E Y-\alpha|\le\frac{52n}{\gamma}.
 \tag{V51.11}\label{r51:mean-eq}
\]
By edge transitivity, each actual corridor plaquette satisfies
\[
 -\frac{2c}{\gamma}\le
 \E X_e-\frac{n-1}{2n}
 \le\frac{64c(n-1)}{3n\gamma}.
 \tag{V51.12}\label{r51:local-mean}
\]
\end{theorem}
\begin{proof}
In \eqref{r51:energy-exact}, use \(T,D\ge0\),
\(\E D\le4c/\gamma\), and \eqref{r51:defect-bound}.
Then \(64c\le52\) and \(6c<52\).
Translations and coordinate permutations act transitively on
unoriented edges, so \(\E X_e=\E Y/(3n)\).
\end{proof}

The coefficient \((n-1)/(2n)=\tfrac32R_{ee}\) is the
three-tangent-component coefficient. V50's four-ambient-component
upper bound \(2R_{ee}\) was valid but did not identify this coefficient.
The present error is uniform in \(N\); it uses neither
\(\gamma\gg\log N\) nor a simultaneous whole-corridor chart event.

\subsection{An explicit extensive logarithmic-normalizer remainder}

Write \(Q_{N,\gamma}=\int e^{-\gamma S_0}\prod_xdH(z_x)\).
Let \(L^{(o)}\) be the scalar rooted Laplacian and set
\[
 \kappa_N=
 \frac{(2\pi)^{3(n-1)/2}}
 {(2\pi^2)^{n-1}\{\det L^{(o)}\}^{3/2}},\qquad
 h_N(\gamma)=\log Q_{N,\gamma}-\log\kappa_N+\alpha\log\gamma.
 \tag{V51.13}\label{r51:normalizer}
\]
This is V1's exact leading factor, not a new choice of measure.
The global right-Haar integral is one. V1 proves, at each fixed
graph, \(h_N(\gamma)\to0\) as \(\gamma\to\infty\).

\begin{theorem}[Finite-regulator normalizer and uniform energy-source bound]
\label{r51:free-energy}
For every allowed \(N\) and \(\gamma>0\),
\[
 -\frac{6cn}{\gamma}\le h_N(\gamma)
 \le\frac{64c(n-1)}{\gamma},\qquad
                       |h_N(\gamma)|\le\frac{52n}{\gamma}.
 \tag{V51.14}\label{r51:free-energy-eq}
\]
For the uniform real energy mark \(t<1\),
\[
 \left|\log\E e^{tY}+\alpha\log(1-t)\right|
       \le\frac{52n}{\gamma}\frac{|t|}{1-t}.
 \tag{V51.15}\label{r51:uniform-mgf}
\]
\end{theorem}
\begin{proof}
Differentiating the compact integral is legitimate at finite \(N\);
it gives \(h_N'(\gamma)=(\alpha-\E_\gamma Y)/\gamma\).
Theorem~\ref{r51:mean} bounds this derivative by constants times
\(n/\gamma^2\). Using the fixed-\(N\) value at infinity,
\[
 h_N(\gamma)=\int_\gamma^\infty
             \frac{\E_u(uS_0)-\alpha}{u}\,du.
\]
Integrate the two bounds in \eqref{r51:mean-eq}.
This is a fixed-\(N\) identity with a uniform estimate, not an
interchange of an unproved thermodynamic limit with a saddle expansion.

The exact source normalizer is
\(\E_\gamma e^{tY}=Q_{N,\gamma(1-t)}/Q_{N,\gamma}\).
Substitute \eqref{r51:normalizer} and integrate
\(|h_N'(u)|\le52n/u^2\) between the two positive couplings.
The integral is \(52n|t|/\{\gamma(1-t)\}\).
\end{proof}

Thus the free-energy correction \emph{per vertex} is \(O(\gamma^{-1})\)
uniformly in the regulator. Relative approximation of the entire
normalizer by \(\kappa_N\gamma^{-\alpha}\) is proved in the stronger
window \(n/\gamma\to0\). These are different conclusions.
Equation \eqref{r51:uniform-mgf} has the harmonic exponent \(\alpha\);
it is not a license to differentiate a remainder inequality twice.

\subsection{A connected energy return obtained without differentiating a bound}

\begin{theorem}[Total-energy variance and its explicit window]
\label{r51:variance}
Let \(V=\Var_\mu(Y)\) and \(\delta=80n/\gamma\). Then
\[
 V=\alpha-\frac{3n}{2}\E D+\frac12\Cov(\rho,Y),
 \qquad
 \sqrt V\le\sqrt\alpha+\delta,
 \tag{V51.16}\label{r51:variance-exact}
\]
and
\[
 |V-\alpha|\le
       \frac{6cn}{\gamma}
       +\frac{80n\sqrt\alpha}{\gamma}
       +\frac{6400n^2}{\gamma^2}.
 \tag{V51.17}\label{r51:variance-error}
\]
In particular,
\[
 \Var(Y/n)\le
       \left(\frac{\sqrt\alpha}{n}+\frac{80}{\gamma}\right)^2.
 \tag{V51.18}\label{r51:density-variance}
\]
For any growing sequence with \(n\to\infty\) and
\(\gamma/\sqrt n\to\infty\), \(V/n\to3/2\).
\end{theorem}
\begin{proof}
Take \(F=Y\) in \eqref{r51:connected}. Since
\(-\mathcal VY=2Y-T\), and
\(\E Y-\E T/2=\alpha-3n\E D/2\), the first identity follows.
Cauchy--Schwarz and \eqref{r51:defect-bound} give
\[
 \tfrac12|\Cov(\rho,Y)|
 \le\tfrac12\|\rho\|_2\sqrt V\le\delta\sqrt V.
\]
Because \(D\ge0\), \(V\le\alpha+\delta\sqrt V\).
The positive root of this quadratic inequality is at most
\(\sqrt\alpha+\delta\). The absolute error now follows from
\(3n\E D/2\le6cn/\gamma\). Divide as stated for the remaining claims.
\end{proof}

For any edge \(f\), transitivity and the exact finite-volume sum give
\[
 \sum_e\Cov(X_e,X_f)=\Cov(Y,X_f)=\frac{V}{m}.
 \tag{V51.19}\label{r51:summed-return}
\]
This is a \emph{signed}, spatially summed connected response of the
actual law. Theorem~\ref{r51:variance} bounds it, and identifies its
limit \(1/2\) in the stated stronger window. It does not bound
\(\sum_e|\Cov(X_e,X_f)|\), a separated-support covariance, or a
physical transfer operator. The density variance tends to zero
whenever \(n,\gamma\to\infty\), but that fact alone gives none of
these missing spatial conclusions.

\subsection{Finite-support positive marks: the same normalizer and identity}

Let \(0\le t_e\le s<1/2\) be supported on at most \(k\) edges and put
\[
 J_t=\sum_et_eX_e,\quad
 d\mu_t=\frac{e^{J_t}}{\E_\mu e^{J_t}}\,d\mu,\quad
 Y_t=Y-J_t,\quad
 T_t=\sum_e(1-t_e)X_eb_e,\quad \rho_t=T_t-3nD.
 \tag{V51.20}\label{r51:tilt}
\]
The positive mark weakens the corresponding ferromagnetic bond.
All other bonds and the product-Haar measure are unchanged.

\begin{proposition}[A normalized marked identity with a paid defect]
\label{r51:marked-identity}
For every smooth \(F\),
\[
 \E_t\{\mathcal VF+(2Y_t-d-\rho_t)F\}=0,\qquad
 2\Cov_t(Y_t,F)=-\E_t\mathcal VF+\Cov_t(\rho_t,F).
 \tag{V51.21}\label{r51:marked-ibp}
\]
With \(H_{k,s}=\exp\{2ks/(1-2s)\}\),
\[
 \E_t|\rho_t|\le\frac{160H_{k,s}n}{\gamma},\qquad
            |\E_tY_t-\alpha|\le\frac{80H_{k,s}n}{\gamma}.
 \tag{V51.22}\label{r51:marked-defect}
\]
Consequently the absolute value of the last covariance in
\eqref{r51:marked-ibp} is at most
\(320H_{k,s}n\|F\|_\infty/\gamma\).
\end{proposition}
\begin{proof}
Integrate the divergence against the actual density \(e^{-Y_t}\).
The geometric identities give \(\mathcal VY_t=-2Y_t+T_t\).
This proves both equalities, including subtraction of the
\(\mu_t\) mean.

The denominator \(\E e^{J_t}\) is at least one. V50, applied to
the doubled marks, gives
\[
 (\E e^{2J_t})^{1/2}
 \le\det(I-2K_t)^{-1}\le H_{k,s},\qquad
 K_t=\operatorname{diag}(\sqrt t)\,R\,
                         \operatorname{diag}(\sqrt t).
\]
Here \(R\preceq I\) and its marked trace is at most \(ks\).
Since \(0\le T_t\le T\), Hölder and the proof of
\eqref{r51:defect-bound} give
\[
 \E_t|\rho_t|\le
 H_{k,s}(\|T\|_2+3n\|D\|_2)\le160H_{k,s}n/\gamma.
\]
Take \(F=1\) for the mean identity and use
\(|\Cov(A,F)|\le2\|F\|_\infty\E|A|\) for the last claim.
Reflection positivity of the defect-bond law was not assumed.
\end{proof}

There is also an explicit leading joint normalizer for these marks.
Let \(L_t=B^{\mathsf T}(I-\operatorname{diag}t)B\), and define
\(\kappa_N(t)\) by \eqref{r51:normalizer} with the rooted matrix
\(L_t^{(o)}\). All its conductances are at least \(1-s>0\).
At fixed \(N,t\), the compact action has exactly the coherent
diagonal as minima; after removing one right-Haar variable its
Hessian is \(L_t^{(o)}\otimes I_3>0\).
The same leading Laplace argument as V1 therefore gives
\(Q_{N,\gamma}(t)\gamma^\alpha\to\kappa_N(t)\).
This invocation needs only the fixed-graph leading limit; the
uniform remainder is supplied by the preceding proposition.

\begin{corollary}[Explicit joint marked normalizer remainder]
\label{r51:marked-normalizer}
With the same finite-support marks,
\[
 \left|\log\frac{Q_{N,\gamma}(t)}
                    {\kappa_N(t)\gamma^{-\alpha}}\right|
          \le\frac{80H_{k,s}n}{\gamma},
 \tag{V51.23}\label{r51:marked-normalizer-eq}
\]
and
\[
 \left|\log\E_\mu e^{J_t}
             +\frac32\log\det(I-K_t)\right|
          \le\frac{(80H_{k,s}+52)n}{\gamma}.
 \tag{V51.24}\label{r51:joint-normalizer}
\]
\end{corollary}
\begin{proof}
At fixed marks,
\(\partial_\gamma\log Q_{N,\gamma}(t)=-\E_tY_t/\gamma\).
Integrate \eqref{r51:marked-defect} from \(\gamma\) to infinity
using the fixed-\(N,t\) leading constant, exactly as above.
The matrix determinant lemma on the mean-zero site space gives
\[
 \frac{\det L_t^{(o)}}{\det L^{(o)}}=\det(I-K_t).
\]
Indeed each rooted determinant is the corresponding nonzero
eigenvalue product divided by \(n\); on the mean-zero space
factor \(L\) and then use Sylvester's determinant identity.
Hence \(\kappa_N(t)/\kappa_N(0)=\det(I-K_t)^{-3/2}\).
Subtract the two logarithmic remainder bounds.
\end{proof}

This is a statement about one normalized joint source function,
not independently chosen plaquette marginals. Its determinant
has three tangent components; V50's domination has four ambient
components. The former is an approximation with the displayed
error, whereas the latter remains a nonasymptotic upper bound.
Even at fixed marked support, our relative approximation still
requires \(n/\gamma\to0\); no smaller error is inferred by
cancelling two separately extensive bounds. Mixed derivatives
at nonzero marks require their own connected estimates.

\subsection{Where the next estimate must improve}

The new vector field yields exact identities and explicit
growing-volume bounds on the \emph{original coherent compact law}.
It replaces an unquantified fixed-graph normalizer error by
\(O(n/\gamma)\), and supplies one actual connected energy return.
It does not resolve general exteriors or noncentral winding.
As in V50, a periodic pure-gauge transport is removable exactly;
a nontrivial loop holonomy is not.

There are now two concrete upstream costs. First, the estimate
\(\|\rho\|_2\le160n/\gamma\) uses an absolute sum of local errors.
Improving its covariance with a specified local source, or its
centered norm, requires spatial cancellation that is not supplied
by one-point infrared moments. Second, the finitely marked
normalizer remainder is extensive. A connected comparison must
control its \emph{difference} before that extensive subtraction
can be replaced by a support-dependent error. Repeating scalar
Gaussian determinants or differentiating a zero-source inequality
would not pay either cost.

In particular, neither \(\gamma\gg\sqrt n\) nor
\(\gamma\gg n\) has been shown to describe the required physical
continuum trajectory. The full Wilson exterior distribution,
source-relative sector coverage, a noncollapsing physical source
Gram, physical-time contraction, and continuum construction remain
open. The next companion review is V55, the next broad literature
sweep V60, and the archive/restart checkpoint V100.

The exact diagnostic
\path{verify_ym_restart_v51_collective_identities.py}
checks the pointwise divergence and energy drift, normalized
Haar integration-by-parts coefficients, resistance determinant
normalization, and the rational constant bounds. It does not
certify spatial decay or the full physical mass gap.
% END CONSOLIDATED V51 APPEND


% BEGIN CONSOLIDATED V52 APPEND
\clearpage
\section[V52: Root-averaged local entropy and pair returns]{V52: Root-averaged compact entropy, a gradient frame, and uniform pair returns}
\label{r52:start}

V51's nonlinear defect bound was extensive. Subtracting two such
normalizer errors did not give a support-dependent comparison.
This section instead localizes the entropy of one actual compact
law. The key point is that Gaussian \emph{edge gradients} form a
bounded frame, even though the rooted site covariance is poorly
conditioned as the graph grows. Averaging the auxiliary root restores
translation symmetry before that localization.

We obtain a fixed-support comparison and a two-plaquette covariance
error uniform in the positions and separation of the plaquettes.
The sufficient window is now logarithmic in the number of vertices,
rather than \(\gamma\gg n\). This is an absolute pairwise error,
not an absolutely summable error kernel or physical exponential decay.

The direction was adjusted after testing a direct winding
reweighting: its exponential current is not bounded by V50's
one-point moments. That obstruction remains. The new argument
controls the coherent conditional and its periodic pure-gauge
transforms; it does not silently extend reflection positivity to
a frustrated transport.

For nonduplication, f16.1/f16.2 V6 already proves Gaussian
marginal-entropy localization by Fisher information and the
Ornstein--Uhlenbeck semigroup, as well as its use for unbounded
scores. We reuse that principle, not claim it anew. The old
condition-number estimate is replaced here by an explicit gradient
frame bound. The same archives treat anchored native laws; f14 V46
treats a declared hard energy-ball chart. Neither supplies the
root-averaged compact corridor law below. V1's root-Haar completion,
V50's original-law infrared moments, and V51's normalizer bounds
are retained exactly. The classical entropy-projection context is
E.~Carlen and D.~Cordero--Erausquin,
\url{https://arxiv.org/abs/0710.0870}; its primary record was
rechecked. No new external theorem about four-dimensional gauge
fields is imported.

\subsection{An exact root chart and its original-law probability}

Use \(N\ge4\) even, \(n=N^3\), \(m=3n\), and the actual coherent
law \(\mu=\mu_{N,\gamma}\). Let \(B\) have edge rows
\(e_x-e_y\), \(L=B^{\mathsf T}B\), \(C=L^\dagger\),
\(c=C_{xx}\le13/16\), and
\[
 R=BCB^{\mathsf T},\qquad r=R_{ee}=\frac{n-1}{3n}.
 \tag{V52.1}\label{r52:graph}
\]
All scalar edge diagonals coincide. For any vertices put
\(R_{xo}=(e_x-e_o)^{\mathsf T}C(e_x-e_o)\).
Then \(R_{xo}\le4c\) and
\[
                         \sum_x R_{xo}=2nc.
 \tag{V52.2}\label{r52:root-trace}
\]
The latter uses translation invariance and \(C\mathbf1=0\).

For a reference vertex \(o\), identify spins with unit quaternions
and write
\[
 R_{\rm root}=z_o,\quad b_x=z_xz_o^{-1},\quad b_o=1,\qquad
 \mathcal A_o=\{\operatorname{Re}b_x\ge1/2\text{ for every }x\}.
 \tag{V52.3}\label{r52:chart}
\]
The change of variables preserves product Haar:
\(R_{\rm root}\) is Haar and independent of the rooted variables
under \(\mu\). This removes an integration variable, not a
physical collective mode; reconstruction retains \(z_x=b_xR_{\rm root}\).

Let \(a=\mu(\mathcal A_o)\), which is independent of \(o\), and set
\[
                  p=8n e^{-\gamma/26}.
 \tag{V52.4}\label{r52:p}
\]
Statements containing \(1-p\) below assume \(p<1\).

\begin{lemma}[A paid root-chart event]
\label{r52:tail}
One has \(1-a\le\min(1,p)\).
On \(\mathcal A_o\) the coordinates
\[
 b_x=(s_x,x_x),\qquad s_x=\sqrt{1-|x_x|^2}\ge1/2,\qquad
 x_o=0,\qquad u_x=\sqrt\gamma\,x_x
 \tag{V52.5}\label{r52:coordinates}
\]
are single valued, and their exact Haar factor is
\[
              dH(b_x)=\frac{d^3x_x}{2\pi^2s_x}.
 \tag{V52.6}\label{r52:haar}
\]
\end{lemma}
\begin{proof}
The scalar part of \(b_x\) equals \(z_x\cdot z_o\).
If it is less than \(1/2\), then \(|z_x-z_o|>1\).
V50's infrared exponential inequality applied to \(e_x-e_o\)
and each of four ambient components gives
\[
 \mu\{|z_x-z_o|>1\}\le
           8e^{-\gamma/(8R_{xo})}
           \le8e^{-\gamma/(32c)}\le8e^{-\gamma/26}.
\]
The root term itself is zero; a union bound gives the claim.
The surface graph \((\sqrt{1-|x|^2},x)\) has area factor \(1/s\).
Divide by \(|S^3|=2\pi^2\).
\end{proof}

Let \(G_o\) be the three-component centered Gaussian on
\(u_o=0\) with covariance \((L^{(o)})^{-1}\otimes I_3\).
For every configuration in the chart,
\[
 \begin{split}
 \gamma S_0(b)
 &=\frac12\sum_e|u_x-u_y|^2+\mathcal Q(u),\\
 \mathcal Q(u)&=\frac{\gamma}{2}\sum_e(s_x-s_y)^2\ge0,\qquad
 \mathcal J(u)=-\sum_{x\ne o}\log s_x
                         \le2\sum_x|x_x|^2 .
 \end{split}
 \tag{V52.7}\label{r52:sign}
\]
For the last inequality use
\(-\log(1-t)\le t/(1-t)\) and \(t=|x_x|^2\le3/4\).
The favorable energy sign is specific to these positive
Cartesian coordinates, not to an arbitrary exponential chart.

\subsection{A normalized compact-to-Gaussian entropy budget}

Let \(P_o\) denote the law of \(u\) under
\(\mu(\,\cdot\,\mid\mathcal A_o)\), with the independent root
Haar variable integrated out. Define
\[
 H=\frac{16cn}{\gamma(1-p)}+\frac{6cn}{\gamma}
                       -\log(1-p),\qquad \epsilon=H/n.
 \tag{V52.8}\label{r52:H}
\]

\begin{theorem}[The actual chart likelihood, with its complement paid]
\label{r52:entropy}
For \(p<1\),
\[
                         D(P_o\Vert G_o)\le H.
 \tag{V52.9}\label{r52:entropy-eq}
\]
If \(p\le1/2\), then
\[
                     \epsilon\le\frac{247}{8\gamma}+\frac{2p}{n}.
 \tag{V52.10}\label{r52:epsilon}
\]
Here \(D\) is relative entropy of normalized probabilities.
\end{theorem}
\begin{proof}
Let \(h_N(\gamma)\) be exactly V51's logarithmic normalizer
remainder, including its rooted determinant and all Haar constants.
On the coordinate domain the likelihood is
\[
 \frac{dP_o}{dG_o}(u)
 =\frac{\mathbf1_{\mathcal A_o}(u)}
              {a\,e^{h_N(\gamma)}}\,
                  e^{\mathcal J(u)-\mathcal Q(u)}.
 \tag{V52.11}\label{r52:likelihood}
\]
Indeed the quadratic Gaussian integral and the factor
\((2\pi^2)^{-(n-1)}\gamma^{-3(n-1)/2}\)
give V1's \(\kappa_N\gamma^{-3(n-1)/2}\).
Thus no normalizer has been independently chosen.

Since \(\mathcal Q\ge0\),
\[
 D(P_o\Vert G_o)\le
             \E_{\mu(\cdot\mid\mathcal A_o)}\mathcal J
                         -h_N(\gamma)-\log a.
\]
Inside the chart,
\(|x_x|^2=1-(z_x\cdot z_o)^2\le|z_x-z_o|^2\).
V50 and \eqref{r52:root-trace} give
\(\E\sum_x|z_x-z_o|^2\le8nc/\gamma\).
Consequently the first term is at most \(16nc/(\gamma a)\).
V51 gives \(h_N(\gamma)\ge-6cn/\gamma\), and \(a\ge1-p\).
This proves \eqref{r52:entropy-eq}.
For \(p\le1/2\), use \(-\log(1-p)\le2p\) and
\(38c\le38(13/16)=247/8\).
\end{proof}

This is a chart-conditioned entropy statement with an explicit
original-law probability \(1-a\). The complete compact law has
not been identified with a Gaussian having support only in a chart.
The errors from removing this conditioning will be included in
the actual plaquette statements below.

\subsection{Root averaging and a volume-independent gradient frame}

For every root and every compact configuration, including those
outside its chart, define
\[
 U_e^{\,o}=\sqrt\gamma\,
       \operatorname{Im}\{(z_x-z_y)z_o^{-1}\}.
 \tag{V52.12}\label{r52:edge-field}
\]
Inside the chart this is \(u_x-u_y\). It is always a discrete
gradient. Let \(\mathcal H=\operatorname{ran}B\otimes\mathbb R^3\).
The image of \(G_o\) on \(\mathcal H\) is the same Gaussian \(G\)
for every root, with covariance \(R\otimes I_3\); on
\(\mathcal H\) this is standard Gaussian covariance.

Choose \(o\) uniformly from the \(n\) vertices, independently before
sampling \(\mu\), and condition the pair \((o,z)\) on \(\mathcal A_o\).
Let \(\overline P\) be the resulting law of \(U^{\,o}\) on
\(\mathcal H\). This is a mixture of chart laws, not a new vacuum law.
Translations act simultaneously on \(o,z\), and coordinate
permutations preserve the construction. Hence \(\overline P\)
is translation invariant and edge transitive. The probabilities
of the chart events are all \(a\), so the conditional root is
still uniform. Entropy convexity and the root-independent
Gaussian gradient image give
\[
                       D(\overline P\Vert G)\le H.
 \tag{V52.13}\label{r52:mixture}
\]
Both facts are needed: a fixed-root chart law alone is not
translation invariant.

\begin{theorem}[Support-dependent entropy, without a site condition number]
\label{r52:local}
For a nonempty set \(S\) of at most \(k\) distinct edges,
with arbitrary diameter,
\[
 D(\overline P_S\Vert G_S)
                  \le12k^2(k+1)\epsilon.
 \tag{V52.14}\label{r52:local-eq}
\]
Entropies are taken on the common linear image when cycle
constraints make the covariance singular. For a single edge the
sharper bound is
\[
                  D(\overline P_e\Vert G_e)\le H/(n-1).
 \tag{V52.15}\label{r52:single}
\]
\end{theorem}
\begin{proof}
We recall the projection form of the Gaussian localization principle
proved in f16 V6. For standardized Gaussian marginal maps with
orthogonal projections \(\Pi_j\), a bound
\(\sum_j\Pi_j\preceq A I\) implies
\[
                \sum_jD(P_j\Vert G_j)\le A D(P\Vert G).
 \tag{V52.16}\label{r52:projection-principle}
\]
Its proof is the same conditional Cauchy--Schwarz bound for the
marginal Fisher informations, integrated along the
Ornstein--Uhlenbeck semigroup. It works on any finite-dimensional
Euclidean subspace and its linear marginal images; zero redundant
coordinates are first removed. The truncation argument in f16 V6
requires only finite entropy, which \eqref{r52:mixture} supplies.

Here the scalar covariance on \(S\) is
\(R_S=B_SC B_S^{\mathsf T}\). Since \(\|L\|\le12\),
\[
       R_S\succeq\tfrac1{12}B_SB_S^{\mathsf T}.
\]
The two matrices have the same kernel. Each nontrivial connected
component of the graph consisting of \(S\) has at most \(k+1\)
vertices. Its positive Laplacian eigenvalues are at least
\(1/\{k(k+1)\}\): for a mean-zero vector, express each value as
the average of its differences from the other vertices, and
bound each difference along a path of at most \(k\) edges.
This yields \(\|v\|_2^2\le k(k+1)\|B_Sv\|_2^2\)
component by component. Therefore
\[
       \lambda_{\min}^{+}(R_S)\ge\frac1{12k(k+1)}.
 \tag{V52.17}\label{r52:gram-floor}
\]

For the \(n\) translates of \(S\), each edge coordinate occurs at
most \(k\) times. On the gradient space, the standardized marginal
projection is
\[
 \Pi_S=R\,P_S^{\mathsf T}R_S^\dagger P_S R
\]
in one color, with the identity in the three colors understood.
Here \(P_S\) extracts edge coordinates. Equation
\eqref{r52:gram-floor} gives
\(\sum_h\Pi_{S+h}\preceq12k^2(k+1)I_{\mathcal H}\).
All translated marginal entropies are equal by the root-averaged
translation symmetry. Divide \eqref{r52:projection-principle}
by \(n\) and use \eqref{r52:mixture}.

For the sharper one-edge statement,
\(\Pi_e=r^{-1}R P_e^{\mathsf T}P_eR\), so
\(\sum_e\Pi_e=r^{-1}I_{\mathcal H}\).
All \(m\) entropies are equal by edge transitivity, and
\(mr=n-1\). This proves \eqref{r52:single}.
\end{proof}

The proof does not bound the growing condition number of the
rooted \emph{site} covariance by a constant. It avoids that
quantity by computing the standardized gradient projections.

\subsection{Return from the auxiliary gradient to native plaquette energies}

Put \(X_e=\gamma W_e\) as in V51 and
\(q_e^{\,o}=\tfrac12|U_e^{\,o}|^2\). Under \(G\), put
\(q_e^G=\tfrac12|U_e|^2\). A useful exact relation, with no chart
restriction, is
\[
 X_e=q_e^{\,o}+\Delta_e^{\,o},\qquad
 \Delta_e^{\,o}=\frac{\gamma}{2}
       \{(z_x-z_y)\cdot z_o\}^2\ge0.
 \tag{V52.18}\label{r52:energy-return}
\]
In particular \(q_e^{\,o}\le X_e\) on the full compact law.

\begin{lemma}[A uniform nonlinear source error]
\label{r52:source-error}
With \(A_p\) from \eqref{r51:Ap}, at every root,
\[
 \E\Delta_e^{\,o}\le\frac{256cr}{\gamma},\qquad
 \|\Delta_e^{\,o}\|_2\le\frac{8cA_4^2r}{\gamma}.
 \tag{V52.19}\label{r52:source-error-eq}
\]
\end{lemma}
\begin{proof}
The difference \(z_x-z_y\) is orthogonal to \(z_x+z_y\).
Thus
\[
 \{(z_x-z_y)\cdot z_o\}^2
 \le\tfrac12|z_x-z_y|^2
                    (|z_x-z_o|^2+|z_y-z_o|^2).
\]
Let \(Z_{xo}=\gamma|z_x-z_o|^2/2\).
V51's positive-exponential moment argument, now applied to the
zero-mean source \(e_x-e_o\), gives
\(\|Z_{xo}\|_p\le A_pR_{xo}\).
Consequently
\[
 \Delta_e^{\,o}\le\gamma^{-1}X_e(Z_{xo}+Z_{yo}).
\]
Use Hölder with moments two for its expectation and moments four
for its \(L^2\) norm. Since \(A_2^2=32\),
\(R_{xo}+R_{yo}\le8c\), and \(\|X_e\|_p\le A_pr\),
the claims follow. A zero resistance means the variable is
identically zero, requiring no division by that resistance.
\end{proof}

\begin{theorem}[A joint bounded-source comparison in a logarithmic window]
\label{r52:bounded}
Let \(F:\mathbb R_+^S\to\mathbb R\) have \(|F|\le M_F\) and
Lipschitz constant \(L_F\) for the \(\ell^1\) metric.
For \(|S|\le k\) and \(p<1\),
\[
 \begin{split}
 |\E_\mu F((X_e)_{e\in S})-\E_G F((q_e^G)_{e\in S})|
 \le{}&\frac{256cL_Fkr}{\gamma}\\
 &+2M_F\{p+\sqrt{6k^2(k+1)\epsilon}\}.
 \end{split}
 \tag{V52.20}\label{r52:bounded-eq}
\]
The constants do not depend on the position or diameter of \(S\).
\end{theorem}
\begin{proof}
Under the original joint law of a uniform root and the compact
spins, \eqref{r52:source-error-eq} bounds the change from \(X\)
to \(q^{\,o}\). Conditioning that joint law on \(\mathcal A_o\)
changes an expectation of a function bounded by \(M_F\) by at
most \(2M_F(1-a)\le2M_Fp\).
Its conditional edge law is precisely \(\overline P\), not an
independently normalized collection of edge marginals.
Pinsker's inequality and \eqref{r52:local-eq} bound the remaining
change by \(2M_F\sqrt{6k^2(k+1)\epsilon}\).
\end{proof}

\subsection{Unbounded two-plaquette marks and an explicit covariance error}

\begin{theorem}[Actual coherent pair return, uniformly in separation]
\label{r52:covariance}
Assume \(\gamma\ge1\) and \(p\le1/2\). For any two edges,
including coincident edges,
\[
 \left|\Cov_\mu(X_e,X_f)-\frac32R_{ef}^{\,2}\right|
 \le512\{\gamma^{-1}+\sqrt p+\epsilon^{1/4}\}.
 \tag{V52.21}\label{r52:covariance-eq}
\]
In particular the error tends to zero along every sequence
\(n,\gamma\to\infty\) with
\(\gamma-26\log n\to+\infty\), uniformly over the two edge positions.
\end{theorem}
\begin{proof}
The Gaussian Wick value is
\[
 \E_Gq_e^G=\tfrac32r,\qquad
 \E_G(q_e^Gq_f^G)=\tfrac94r^2+\tfrac32R_{ef}^{\,2}.
\]
We give quantitative moment bounds to pay the unbounded marks.

First \eqref{r52:source-error-eq}, \(q_e^{\,o}\le X_e\), and
\(\|X_e\|_2\le A_2r\) imply
\[
 |\E X_eX_f-\E q_e^{\,o}q_f^{\,o}|
           \le\frac{16cA_4^2A_2r^2}{\gamma}.
 \tag{V52.22}\label{r52:product-return}
\]
This holds uniformly in the root and hence after root averaging.
Next \(\E(X_eX_f)^2\le A_4^4r^4=1536r^4\).
For \(F=q_e^{\,o}q_f^{\,o}\), Cauchy--Schwarz on the excluded event
and \(a\ge1/2\) show that conditioning costs at most
\[
                 100r^2\sqrt p.
 \tag{V52.23}\label{r52:conditioning-product}
\]
Indeed the two costs are bounded by
\(\sqrt{1536}\,r^2\sqrt p\) and
\(p\sqrt{1536/a}\,r^2\), whose sum is below the displayed bound.

For probabilities \(P,Q\) on the same space and a square-integrable
function \(F\), Cauchy--Schwarz against the measure \(P+Q\) gives
\[
 |\E_PF-\E_QF|
 \le\{2\|P-Q\|_{\rm TV}\}^{1/2}
                  \{\E_PF^2+\E_QF^2\}^{1/2}.
 \tag{V52.24}\label{r52:weighted-TV}
\]
Use their densities against \(P+Q\), and
\((p_0-q_0)^2/(p_0+q_0)\le|p_0-q_0|\), to verify the formula.
For \(P=\overline P_{\{e,f\}}\) and \(Q=G_{\{e,f\}}\),
the first second moment for \(F=q_eq_f\) is at most \(3072r^4\).
The second is at most \(945r^4/16\), by the fourth moment of
\(\tfrac12r\chi_3^2\) and Cauchy--Schwarz.
Their sum is less than \(56^2r^4\).
Pinsker and the \(k=2\) case of \eqref{r52:local-eq} therefore give
\[
 |\E_{\overline P}q_eq_f-\E_Gq_eq_f|
                         \le56r^2(288\epsilon)^{1/4}.
 \tag{V52.25}\label{r52:entropy-product}
\]

Finally V51 gives
\(|\E_\mu X_e-\tfrac32r|\le52/(3\gamma)\), while
\(\E_\mu X_e\le2r\). The difference of the two products of
means is at most \(182r/(3\gamma)\).
Combine this with \eqref{r52:product-return}--%
\eqref{r52:entropy-product}.
Since \(r\le1/3\), \(c\le13/16\),
\(A_4^2=16\sqrt6\), and \(A_2=4\sqrt2\),
the coefficient of \(\gamma^{-1}\) is at most \(390\)
(using \(\sqrt3\le2\)); the other two coefficients are each
less than \(512\). This proves \eqref{r52:covariance-eq}.
Under the stated sequence \(p\to0\), and
\eqref{r52:epsilon} gives \(\epsilon\to0\).
\end{proof}

This pays an actual compact-law, separated-support return, not
just the uniform-energy response of V51 or a fixed-graph
Gaussian coefficient. It is not permissible to sum the error
in \eqref{r52:covariance-eq} over all \(f\): that could cost the
number of edges. A spatially summable remainder, a relative
source-Gram bound, and an original-exterior average remain
different requirements.

\subsection{The remaining winding and physical bottlenecks}

The chart and its root averaging are analysis devices. The final
plaquette expectations use the full original coherent conditional;
its off-chart probability and unbounded source moments were paid.
There is no assumption that the auxiliary root is a physical
excitation, or that a Gaussian approximation is exact.

This also explains the limit of the attempted winding transfer.
The density for a transported action contains
\(e^{-\gamma(S_\theta-S_0)}\). It is not a uniformly bounded local
observable, and V50's positive deficit marks do not automatically
control its signed current. Nor does root averaging commute with
a nontrivial transported normalizer. The new local comparison
cannot be substituted for that missing estimate.

The sufficient logarithmic condition above has not been checked
against the physical bare-coupling trajectory; its explicit constant
must be respected. Even where it holds, a uniform absolute
pairwise error supplies neither exponential clustering of the
full Wilson law nor a positive physical Hamiltonian gap.
The next productive targets are a spatially weighted version of
the localized defect or a genuinely normalized winding-current
estimate, not another coherent determinant.

The diagnostic
\path{verify_ym_restart_v52_local_entropy.py}
checks root covariance identities, cyclic and separated carrier
Gram matrices, gradient-frame normalization, exact quaternion
source returns, Haar chart normalization, and the covariance
constants. Finite checks do not prove the imported infrared or
entropy principles. Companion review remains due V55; the next
regular broad literature sweep is V60. The interacting
four-dimensional continuum construction and its mass gap remain
unproved.
% END CONSOLIDATED V52 APPEND


% BEGIN CONSOLIDATED V53 APPEND
\clearpage
\section[V53: Full signed-chart entropy and local source normalizers]{V53: Full signed-chart entropy, without a logarithmic volume window}
\label{r53:start}

V52 localized the entropy of a root-chart conditional, then paid the
probability of leaving that chart. Its union bound imposed
\(\gamma-26\log n\to+\infty\). This section removes that restriction
for the coherent conditional: retain every scalar-sign sheet, and pay
the signs by their actual conditional Shannon entropy. A one-site
conditional density also controls the logarithmic Haar singularity
at the equator. These are per-site costs, not a claim that the
probability of any bad site is small.

The resulting entropy density is \(O(\gamma^{-1})\), uniformly over
all even \(N\ge4\). V52's gradient frame then gives the same
fixed-support comparison without conditioning. We also prove a
normalized comparison for signed, fixed-support exponential
plaquette sources, using their actual second exponential moments.
The support may have arbitrary diameter. Neither the covariance
error nor the source-normalizer error is asserted to be spatially
summable.

For nonduplication, f14 V54 already derived both scalar-sign Haar
sheets, and f14 V96 compared a one-site spherical conditional with
a Gaussian. Those results are reused as geometric context, not
claimed new. The information identities and warning about fixed-law
entropy versus functional inequalities in f15 V23 also remain in
force. What was missing here is a full correlated corridor
normalizer with all sheets retained and a volume-independent local
return. The input remains V50's classical finite-volume infrared
inequality, V51's scalar normalizer, and V52's gradient frame.
The primary records
\href{https://arxiv.org/abs/2302.07299}{Giuliani--Ott},
\href{https://arxiv.org/html/math-ph/0610025}{Biskup}, and
\href{https://arxiv.org/abs/0710.0870}{Carlen--Cordero--Erausquin}
were reopened for attribution; no new external theorem is imported.

\subsection{Both sign sheets in the same compact normalizer}

Use exactly the coherent law \eqref{r50:law}, with \(N\ge4\) even,
\(n=N^3\), \(m=3n\), \(\gamma>0\), and \(B,L,C,c,R,r\) as in
\eqref{r52:graph}. In particular \(c\le13/16\) and \(r<1/3\).
Choose a reference vertex \(o\). V1 and V52's root-Haar
factorization still writes \(z_x=b_xR_{\rm root}\), with
\(b_o=1\) and an independent Haar \(R_{\rm root}=z_o\).
No root orientation is physically pinned.

Off the equators, which have zero probability at finite \(N,\gamma\),
write
\[
 b_x=(\sigma_xs_x,x_x),\qquad
 s_x=\sqrt{1-|x_x|^2}>0,\qquad
 \sigma_x\in\{-1,1\},\qquad u_x=\sqrt\gamma\,x_x .
 \tag{V53.1}\label{r53:coordinates}
\]
At the root \(s_o=\sigma_o=1\), \(u_o=0\).
Every other sign is summed over. On each sheet the Haar factor is
\((2\pi^2s_x)^{-1}d^3x_x\), not twice that factor. The full
action identity is
\[
 \begin{split}
 \gamma S_0
 &=\frac12\sum_e|u_x-u_y|^2+\mathcal Q_\sigma(u),\\
 \mathcal Q_\sigma(u)
 &=\frac{\gamma}{2}\sum_e(\sigma_xs_x-\sigma_ys_y)^2\ge0,
 \qquad
 \mathcal J(u)=-\sum_{x\ne o}\log s_x .
 \end{split}
 \tag{V53.2}\label{r53:action}
\]
This identity holds for mixed signs too. It is not the
positive-sheet correction with the other sheets discarded.

Let \(\mathcal D=\{u_o=0,\ |u_x|<\sqrt\gamma\ (x\ne o)\}\),
let \(G_o\) be V52's normalized rooted Gaussian, and let \(P_o^*\)
be the image of the \emph{unconditioned} compact law under \(z\mapsto u\).
Put
\[
 Z_{\rm sign}(u)=\sum_{\sigma:\,\sigma_o=1}e^{-\mathcal Q_\sigma(u)}.
\]
Using the same \(h_N(\gamma)\) as in \eqref{r51:normalizer}, one has
\[
 \frac{dP_o^*}{dG_o}(u)
 =\mathbf1_{\mathcal D}(u)e^{-h_N(\gamma)}
                       e^{\mathcal J(u)}Z_{\rm sign}(u),\qquad
 \pi(\sigma\mid u)=\frac{e^{-\mathcal Q_\sigma(u)}}{Z_{\rm sign}(u)} .
 \tag{V53.3}\label{r53:likelihood}
\]
Indeed the rescaled one-sheet Haar constants and Gaussian integral
give precisely \(\kappa_N\gamma^{-3(n-1)/2}\); the remaining
integral is the original compact partition function. In particular
\(\int_{\mathcal D}e^{\mathcal J}Z_{\rm sign}\,dG_o=e^{h_N}\).
There is no added factor \(2^{-(n-1)}\) in the sign sum.

For a finite-valued random variable, write \(\mathsf H\) for
Shannon entropy with natural logarithms. The conditional distribution
in \eqref{r53:likelihood} gives the exact identity
\[
 \log Z_{\rm sign}(u)
   =\mathsf H(\sigma\mid u=u)-\E[\mathcal Q_\sigma\mid u].
\]
Consequently, once the integrability below is established,
\[
 D(P_o^*\Vert G_o)
   =\E_\mu\mathcal J-\E_\mu\mathcal Q_\sigma
                   +\mathsf H_\mu(\sigma\mid u)-h_N(\gamma).
 \tag{V53.4}\label{r53:entropy-identity}
\]
This is a finite conditional Shannon entropy, not a subtraction
of infinite differential entropies. Bounding the sign sum only by
\(2^{n-1}\) would lose \((n-1)\log2\). Instead,
\(\mathsf H(\sigma\mid u)\le\sum_{x\ne o}\mathsf H(\sigma_x)\);
no independence of the signs is assumed.

\subsection{A rare-event logarithmic estimate at the equator}

Define
\[
 \begin{split}
 \delta_\gamma&=\min\{1,8e^{-\gamma/26}\},\\
 b(t)&=-t\log t-(1-t)\log(1-t),\qquad b(0)=b(1)=0,\\
 a_\gamma&=\delta_\gamma
              \{1+6\gamma+\log(2/\delta_\gamma)\}
                 +b(\min\{1/2,\delta_\gamma\}),\\
 \eta_\gamma&=\frac{22c}{\gamma}+a_\gamma .
 \end{split}
 \tag{V53.5}\label{r53:budget}
\]
Unlike V52's \(p\), none of these quantities contains \(n\).

\begin{lemma}[The full-law Jacobian and sign costs]
\label{r53:costs}
At every root,
\[
 \begin{split}
 \E_\mu\mathcal J
 &\le\frac{16cn}{\gamma}
       +(n-1)\delta_\gamma\{1+6\gamma+\log(2/\delta_\gamma)\},\\
 \mathsf H_\mu(\sigma\mid u)
 &\le(n-1)b(\min\{1/2,\delta_\gamma\}).
 \end{split}
 \tag{V53.6}\label{r53:costs-eq}
\]
In particular the Jacobian logarithm is integrable.
\end{lemma}
\begin{proof}
Put \(w_x=z_x\cdot z_o=\sigma_xs_x\). For \(x\ne o\),
the event \(w_x<1/2\) implies \(|z_x-z_o|>1\).
V50's four-component sub-Gaussian bound gives
\[
 \mu(w_x<1/2)\le
 \min\{1,8e^{-\gamma/(8R_{xo})}\}\le\delta_\gamma .
 \tag{V53.7}\label{r53:site-tail}
\]
Thus both \(\{\sigma_x=-1\}\) and \(\{s_x<1/2\}\) have
probability at most \(\delta_\gamma\). For
\(\delta_\gamma\le1/2\), use monotonicity of binary entropy
on \([0,1/2]\); for larger \(\delta_\gamma\), use its maximum
\(\log2\). Conditional entropy is no larger than joint entropy,
and joint entropy is no larger than the sum of marginal entropies.
This proves the second assertion.

A small probability alone does not control \(-\log s_x\).
To pay that unbounded quantity, condition on every spin except
\(z_x\); since \(x\ne o\), the reference \(z_o\) is then fixed.
The exact one-site conditional density with respect to normalized
Haar on \(S^3\) is
\[
 \frac{\exp(\gamma z_x\cdot h_x)}
           {\int_{S^3}\exp(\gamma z\cdot h_x)\,dH(z)},
 \qquad h_x=\sum_{y:\,y\sim x}z_y,\qquad |h_x|\le6 .
\]
Jensen's inequality and \(\int z\,dH(z)=0\) put the denominator
above one. This density is therefore at most \(e^{6\gamma}\).
For a fixed unit vector \(z_o\), the Haar coordinate \(z\cdot z_o\)
has density \((2/\pi)\sqrt{1-t^2}\) on \((-1,1)\). Hence
\[
 \mu(s_x<\varepsilon)\le
               2e^{6\gamma}\varepsilon,\qquad 0<\varepsilon\le1/2.
 \tag{V53.8}\label{r53:equator}
\]
This is uniform in the conditioned exterior. It does not assert
that its density upper bound is small.

Here is the elementary integration step with all costs retained.
If \(0<S\le1\), \(\Pr(S<1/2)\le\delta\le1\), and
\(\Pr(S<e^{-t})\le Ke^{-t}\) for \(t\ge\log2\), where
\(K\ge2\delta>0\), put \(T=\log(K/\delta)\ge\log2\). Then
\[
 \begin{split}
 \E[-\log S;\ S<1/2]
 &=\log2\,\Pr(S<1/2)+\int_{\log2}^{\infty}\Pr(S<e^{-t})\,dt\\
 &\le\delta\log2+\int_{\log2}^{\infty}\min\{\delta,Ke^{-t}\}\,dt\\
 &=\delta(1+\log(K/\delta)).
 \end{split}
 \tag{V53.9}\label{r53:layer}
\]
The two integral pieces end and begin at \(T\), respectively.
Taking \(K=2e^{6\gamma}\), \(\delta=\delta_\gamma\) gives
the displayed rare-equator cost.

On \(s_x\ge1/2\), the same elementary bound used in V52 gives
\[
 -\log s_x\le2|x_x|^2,\qquad
 |x_x|^2=1-w_x^2\le|z_x-z_o|^2 .
\]
The second inequality holds for all \(w_x\in[-1,1]\).
Sum V50's bound
\(\E|z_x-z_o|^2\le4R_{xo}/\gamma\), and use
\(\sum_xR_{xo}=2nc\). This contributes at most \(16cn/\gamma\).
Adding the rare-equator contribution proves the first assertion.
\end{proof}

\begin{theorem}[Unconditioned compact entropy at every volume]
\label{r53:global}
For every \(\gamma>0\) and every even \(N\ge4\),
\[
 D(P_o^*\Vert G_o)
 \le\frac{22cn}{\gamma}+(n-1)a_\gamma
 \le n\eta_\gamma .
 \tag{V53.10}\label{r53:global-eq}
\]
Moreover, when \(\gamma\ge26\log16\),
\[
 \eta_\gamma\le
 \frac{143}{8\gamma}
       +8\left(2+\frac{79}{13}\gamma\right)e^{-\gamma/26}
                    =O(\gamma^{-1}),
 \tag{V53.11}\label{r53:rate}
\]
with constants independent of \(N\).
\end{theorem}
\begin{proof}
The finite sign space, the integrable \(\mathcal J\), and
\(0\le\mathcal Q_\sigma\le2\gamma m\) justify
\eqref{r53:entropy-identity}. Drop its nonpositive term
\(-\E\mathcal Q_\sigma\) and use V51's
\(h_N(\gamma)\ge-6cn/\gamma\). Lemma \ref{r53:costs} proves
\eqref{r53:global-eq}.

For the last assertion,
\(\delta_\gamma=8e^{-\gamma/26}\le1/2\), and
\(b(\delta)\le\delta(1-\log\delta)\). Thus
\[
 a_\gamma\le
 \delta_\gamma\left\{2+\frac{79}{13}\gamma-\log32\right\}
 \le8\left(2+\frac{79}{13}\gamma\right)e^{-\gamma/26}.
\]
Finally \(22c\le143/8\).
\end{proof}

The exponential \(e^{6\gamma}\) in the conditional density has
become a factor \(6\gamma\) multiplying a rare probability. This
is a logarithmic integration argument, not an invalid uniform
bound on the Haar Jacobian. Nor is the complement of a simultaneous
chart being discarded: that complement may have probability near
one at large volume.

\subsection{Localized gradients and native pair returns, without conditioning}

Choose a root uniformly and independently of \(z\sim\mu\), with
\emph{no subsequent conditioning}. Let \(\overline P^*\) be the
law of the full gradient
\[
 U_e^{\,o}=\sqrt\gamma\,\operatorname{Im}\{(z_x-z_y)z_o^{-1}\}.
\]
Each rooted Gaussian has the same gradient image \(G\) on
\(\mathcal H=\operatorname{ran}B\otimes\mathbb R^3\).
Entropy contraction and convexity, followed by
\eqref{r53:global-eq}, give
\(D(\overline P^*\Vert G)\le n\eta_\gamma\).
Simultaneous translations of root and spins make this mixture
translation invariant; coordinate permutations give edge
transitivity. Therefore the already-proved V52 frame calculation
applies without modification:
\[
 \begin{split}
 D(\overline P_S^*\Vert G_S)
       &\le C_k\eta_\gamma,\qquad C_k=12k^2(k+1),\quad |S|\le k,\\
 D(\overline P_e^*\Vert G_e)
       &\le\frac{n}{n-1}\eta_\gamma .
 \end{split}
 \tag{V53.12}\label{r53:local}
\]
For cyclic supports entropy is on the common linear image.
The V52 projection theorem and carrier Gram estimate, not
tensorization for independent edges, justify this use.
In particular separated supports require no new site
condition-number estimate.

Keep \(X_e=\gamma W_e\), \(q_e^{\,o}=|U_e^{\,o}|^2/2\), and
\(q_e^G=|U_e|^2/2\). V52's full-law identity and bounds remain
\[
 X_e=q_e^{\,o}+\Delta_e^{\,o},\quad
 \Delta_e^{\,o}\ge0,\quad
 \E\Delta_e^{\,o}\le\frac{256cr}{\gamma},\quad
 \|\Delta_e^{\,o}\|_2\le\frac{8cA_4^2r}{\gamma}.
 \tag{V53.13}\label{r53:return}
\]
These bounds already held outside every chart, so they are reused,
not reproved.

\begin{corollary}[Full coherent local observables and covariance]
\label{r53:observables}
If \(|F|\le M_F\) and \(F\) has \(\ell^1\)-Lipschitz constant
\(L_F\) on \(k\) nonnegative energies, then
\[
 |\E_\mu F(X_S)-\E_G F(q_S^G)|
 \le\frac{256cL_Fkr}{\gamma}
             +2M_F\sqrt{6k^2(k+1)\eta_\gamma}.
 \tag{V53.14}\label{r53:bounded}
\]
For every \(\gamma\ge1\), all even \(N\ge4\), and any two edges,
including a coincident pair,
\[
 \left|\Cov_\mu(X_e,X_f)-\frac32R_{ef}^2\right|
              \le512\{\gamma^{-1}+\eta_\gamma^{1/4}\}.
 \tag{V53.15}\label{r53:covariance}
\]
The covariance error is \(O(\gamma^{-1/4})\) uniformly in \(N,e,f\).
No comparison between \(\gamma\) and \(\log n\) is needed.
\end{corollary}
\begin{proof}
For the first assertion, use \eqref{r53:return}, then Pinsker with
\eqref{r53:local}. There is no chart-conditioning term.

For the second, the exact native-to-gradient product error is
still at most \(16cA_4^2A_2r^2/\gamma\), by
\eqref{r52:product-return}. Under \(\overline P^*\),
\(\E(q_eq_f)^2\le1536r^4\), since \(q^{\,o}\le X\).
Under \(G\) the same second moment is at most \(945r^4/16\).
The sum is less than \(56^2r^4\), even without V52's extra
conditioning factor. Its weighted total-variation inequality,
\eqref{r52:weighted-TV}, and
\(D(\overline P_{\{e,f\}}^*\Vert G_{\{e,f\}})\le144\eta_\gamma\)
give the product comparison
\[
 |\E_{\overline P^*}q_eq_f-\E_Gq_eq_f|
                 \le56r^2(288\eta_\gamma)^{1/4}.
\]
Use V51's mean error \(52/(3\gamma)\) and \(\E X_e\le2r\);
the products of means differ by at most \(182r/(3\gamma)\).
V52's constant calculation bounds the combined inverse-coupling
coefficient by \(390<512\), and the remaining coefficient is
also below \(512\). Wick's formula supplies
\(\Cov_G(q_e,q_f)=3R_{ef}^2/2\). This proves the assertion
without differentiating a normalizer remainder.
\end{proof}

\subsection{Signed exponential plaquette sources in their actual normalizer}

The new entropy also improves V51's extensive mixed-source error.
For a nonempty set \(S\) of at most \(k\) edges and real
\(|t_e|\le s<1/2\), put \(T=\operatorname{diag}(t_e:e\in S)\) and
\[
 \begin{split}
 Z_\mu(t)&=\E_\mu e^{\sum_{e\in S}t_eX_e},\\
 Z_G(t)&=\E_G e^{\sum_{e\in S}t_eq_e^G}
       =\det(I-R_S^{1/2}TR_S^{1/2})^{-3/2},\\
 \mathcal B_{k,s}(\gamma)
  &=(1-2s)^{-k}
       \left\{\frac{8cksA_4^2r}{\gamma}
                    +\sqrt2\,(2C_k\eta_\gamma)^{1/4}\right\}.
 \end{split}
 \tag{V53.16}\label{r53:marked-definitions}
\]
Zero covariance directions contribute determinant one. Since
\(0\preceq R_S\preceq I\) and \(\|T\|\le s\), the determinant
matrix is positive definite. The signed \(T\) is not replaced by
an undefined signed square root.

\begin{theorem}[A fixed-support joint log normalizer, uniform in volume]
\label{r53:marked}
In this setting,
\[
 |Z_\mu(t)-Z_G(t)|\le\mathcal B_{k,s}(\gamma),\qquad
 |\log Z_\mu(t)-\log Z_G(t)|
                  \le e^{2skr}\mathcal B_{k,s}(\gamma).
 \tag{V53.17}\label{r53:marked-eq}
\]
For fixed \(k,s<1/2\), these errors are \(O_{k,s}(\gamma^{-1/4})\)
uniformly in volume, support diameter, and the source signs.
They refer to the same baseline compact law and its normalized
joint source insertion.
\end{theorem}
\begin{proof}
For positive marks of size \(2s\) on \(S\), V50's determinant
bound, using rank at most \(k\), gives
\[
 \E_\mu e^{2s\sum_{e\in S}X_e}\le(1-2s)^{-2k}.
 \tag{V53.18}\label{r53:marked-moment}
\]
The corresponding three-color Gaussian moment is at most
\((1-2s)^{-3k/2}\). These are true second exponential moments;
bounded-observable total variation alone would not suffice.

At every root and configuration, the mean-value formula and
\(0\le q_e^{\,o}\le X_e\) imply
\[
 \left|e^{\sum t_eX_e}-e^{\sum t_eq_e^{\,o}}\right|
       \le s\sum_{e\in S}\Delta_e^{\,o}\,
                                 e^{s\sum_{e\in S}X_e}.
\]
Cauchy--Schwarz, \eqref{r53:return}, and
\eqref{r53:marked-moment} bound its expectation by
\(8cksA_4^2r\,(1-2s)^{-k}/\gamma\).
This is uniform in the root.

For the remaining comparison take
\(F(U)=e^{\sum t_e|U_e|^2/2}\) on the \(S\)-marginal.
Its second moment under \(\overline P^*\) is at most
\((1-2s)^{-2k}\), since \(q^{\,o}\le X\); under \(G\) it is
at most \((1-2s)^{-3k/2}\). Their sum is no larger than
\(2(1-2s)^{-2k}\). Weighted total variation and Pinsker with
\eqref{r53:local} give
\[
 |\E_{\overline P^*}F-\E_GF|
       \le\sqrt2\,(1-2s)^{-k}(2C_k\eta_\gamma)^{1/4}.
\]
This proves the first inequality.

Finally Jensen and the original mean bound give
\(Z_\mu(t)\ge e^{-2skr}\).
For the Gaussian mean \(3r/2\), Jensen gives
\(Z_G(t)\ge e^{-3skr/2}\ge e^{-2skr}\).
The logarithm is \(e^{2skr}\)-Lipschitz above this common
positive lower bound. This proves the second inequality.
\end{proof}

This pays a fixed-support \emph{relative source normalizer}
without the \(n/\gamma\) error of V51, and for either sign of
each source. It does not control growing supports at no cost:
both \(C_k\) and the exponential-moment constants are displayed.
Nor may derivatives of \eqref{r53:marked-eq} be taken as if its
remainder had already been differentiated. The covariance result
was proved separately using actual fourth moments.

\subsection{What has been removed, and what remains upstream}

The global root-chart logarithmic window is no longer an
obligation for the coherent comparisons \eqref{r53:bounded},
\eqref{r53:covariance}, and \eqref{r53:marked-eq}.
All signs, the equator, the collective Haar orientation, and the
joint source normalizer are included. The ESM transfer principle
retained in earlier reviews is implemented here as a common-law
normalizer and source estimate; no perturbative ESM coefficient
has been transplanted.

The spatial bottleneck is unchanged in kind: summing a uniform
absolute covariance error over all distant \(f\) may cost \(m\).
A summable spatial remainder or a relative quadratic source-form
estimate still requires additional structure. Likewise the
noncentral winding density contains a signed transported current,
not just the positive energy marks used in
\eqref{r53:marked-moment}. These estimates must not be assigned
to a general frustrated exterior without a new input.

The next useful direction is a spatially weighted connected
estimate, or an actual normalized noncentral current bound.
Small fixed-law entropy is not a log-Sobolev or Poincare
inequality for all test functions; f15 V23's warning remains
applicable. There is still no full Wilson exterior average,
verified physical source Gram and sector coverage, physical-time
contraction, or interacting four-dimensional continuum
construction. In particular no physical Yang--Mills mass gap
has been proved.

The diagnostic \path{verify_ym_restart_v53_signed_entropy.py}
checks exact finite signed-sheet probabilities and energy
identities, high-precision logarithmic identities, the
rare-event integration constants, and source determinant
fixtures. It is not an independent proof of the inherited
infrared theorem or the analytic entropy bounds.
V52's mathematical checks are also rerun. Companion review
is next due V55; the regular broad literature checkpoint is
V60, and the V100 archive/restart rule remains unchanged.
% END CONSOLIDATED V53 APPEND


% BEGIN CONSOLIDATED V54 APPEND
\clearpage
\section[V54: Transport likelihoods and native noncentral returns]{V54: Exact transport likelihood moments and a paid noncentral connection window}
\label{r54:start}

The first attempted route beyond V53 was to turn the signed
total-energy susceptibility into an absolutely summable covariance
row. That would require an additional correlation-sign inequality.
The targeted literature check did not supply one for the four-component
sphere law. We therefore go upstream of the transported-current
obstruction instead: bound the full likelihood of a connection
perturbation, including its normalizer and its second moment.
This allows unbounded native plaquette products to be transferred,
not just bounded tests.

The new input is an affine-matrix extension of an \emph{already
inherited} strand domination argument. The orthogonal partition
comparison itself occurs in f16.2 V62 and in V30's proof of
Lemma \ref{r30:inherited}; V45 already supplies the unpinned positive
pairing law. Neither those arguments nor V46's fixed-graph thermal
asymptotics gave the real-parameter likelihood bound below.
We do not claim either the pairing expansion or entropy-to-likelihood
bookkeeping anew: f15 V31 already used an order-two likelihood budget
for a different source problem. Here that budget is proved for the
actual noncommutative corridor transport, uniformly in graph size.

\subsection{Why covariance positivity is not an available shortcut}

Theorem 2.3 in
\href{https://arxiv.org/html/1708.00058}{Peled--Spinka,
\emph{Lectures on the Spin and Loop \(O(n)\) Models}},
Section 2.3, states the relevant nonnegative covariance of spin
inner products for spin dimensions one and two. Its stated
hypothesis does not include four. The special single-site
Laplace-transform condition in
\href{https://arxiv.org/html/hep-lat/0205028}{Orland,
\emph{Griffiths Inequalities for some \(O(n)\) Classical Spin Models
with \(n\ge3\)}}, equations (1.3)--(1.4), does not apply to
normalized Haar on \(S^3\) either. Indeed the inherited sphere
moments give, with \(v=|q|^2\),
\[
 \int_{S^3}e^{q\cdot z}\,dH(z)=1+\frac v8+\frac{v^2}{192}+O(v^3),
 \qquad
 \log\int_{S^3}e^{q\cdot z}\,dH(z)
                =\frac v8-\frac{v^2}{384}+O(v^3).
 \tag{V54.1}\label{r54:scope}
\]
The negative second coefficient excludes that sufficient
nonnegative-coefficient hypothesis. It does not disprove the
desired covariance inequality for this coherent law.
No such inequality is assumed below. In particular V51's signed
row sum has not been promoted to an absolute row sum.
The cited primary theorem and hypothesis passages were read;
the regular broad review remains due V60.

\subsection{The original connection law and a frame-dependent budget}

Let \(G=(V,E)\) be a finite connected graph, with edges counted
once and one orientation per edge. The later local comparison uses
the even cubic torus \(N\ge4\). For real orthogonal transports
\(T_e\in SO(4)\), set
\[
 \begin{split}
 S_T(z)&=\sum_{e=(x,y)}(1-z_x\cdot T_ez_y),\\
 Q_T&=\int e^{-\gamma S_T(z)}\prod_xdH(z_x),\qquad
 d\mu_T=Q_T^{-1}e^{-\gamma S_T}\prod_xdH(z_x).
 \end{split}
 \tag{V54.2}\label{r54:law}
\]
Here \(\gamma>0\), \(z_x\in S^3\). These are the actual left/right
corridor transports \(T_ez=l_ezr_e^{-1}\) of V45; they need not
commute or have trivial holonomy. Write \(Q_0,\mu_0\) for \(T_e=I\).
In the chosen frame define
\[
 \begin{split}
 d_e&=1-\tfrac14\Tr T_e=\tfrac18\|I-T_e\|_{\rm F}^2,\qquad
 a_e=\|I-T_e\|_{\rm op},\\
 \Gamma&=\gamma\sum_ed_e,\qquad
 \Phi_T(z)=\gamma\sum_{e=(x,y)}z_x\cdot(T_e-I)z_y
                         =-\gamma(S_T-S_0).
 \end{split}
 \tag{V54.3}\label{r54:budget}
\]
The Frobenius and operator norms are real four-dimensional norms.
Two orthogonal rotation planes show that
\[
                         a_e^2\le4d_e .
 \tag{V54.4}\label{r54:SO4}
\]
For if the rotation angles are \(\alpha,\beta\), then
\(a_e^2=2\max(1-\cos\alpha,1-\cos\beta)\), whereas
\(4d_e=2(1-\cos\alpha)+2(1-\cos\beta)\).

Vertex rotations \(g_x\in SO(4)\) change variables \(z_x=g_xw_x\)
and replace \(T_e\) by \(T_e^g=g_x^{\mathsf T}T_eg_y\).
Haar, \(Q_T\), and every native plaquette energy are preserved
under this simultaneous transformation. Every estimate below
can be used in any such specified frame. The minimum of
\(\Gamma(T^g)\) over the compact set of vertex frames is attained;
it may therefore be used for the invariant conclusions.
The comparison is then between the transformed \(\mu_T\) and
the coherent law in that frame, not between two unrelated
choices of normalizer. No averaging of exterior frames is
implicitly performed.

\subsection{An affine transport bound for every real likelihood parameter}

For real edge matrices \(M_e\), define
\(\mathcal Z(M)=\int\exp\{\gamma\sum z_x\cdot M_ez_y\}\,dH\).

\begin{lemma}[The inherited strand bound with nonorthogonal matrices]
\label{r54:matrix-domination}
If \(b_e\ge\|M_e\|_{\rm op}\), \(b_e\ge0\), then
\[
                 0<\mathcal Z(M)\le\mathcal Z((b_eI)_e).
 \tag{V54.5}\label{r54:matrix-eq}
\]
\end{lemma}
\begin{proof}
Use the existing Haar pairing expansion. At a vertex of degree
\(2j\) in the occurrence multigraph, each pairing has positive
coefficient \(1/\{2^j(j+1)!\}\). A closed strand contracts to
the trace of the ordered product of its edge matrices, with
transposes when traversed backwards. Its absolute value is at
most four times the product of their operator norms.
Replace those norms by \(b_e\). The resulting scalar diagram
is exactly bounded by the corresponding positive diagram for
\(b_eI\). This argument uses absolute values of diagram weights;
it does not assume transported traces are nonnegative.
All summed majorants are finite, bounded by the compact
scalar exponential majorant \(e^{\gamma\sum b_e}\).
Thus summation is legitimate. Positivity of \(\mathcal Z(M)\)
comes from its original real exponential integral.
\end{proof}

\begin{theorem}[Full transported-current exponential bound]
\label{r54:mgf}
Put \(q_\lambda=\max\{\lambda(\lambda-1),0\}\).
For every real \(\lambda\),
\[
 \begin{split}
 \E_{\mu_0}e^{\lambda\Phi_T}
 &\le
 \exp\!\left\{\gamma\sum_e
          \bigl(\sqrt{1+q_\lambda a_e^2}-1\bigr)\right\}\\
 &\le \exp(2q_\lambda\Gamma).
 \end{split}
 \tag{V54.6}\label{r54:mgf-eq}
\]
These are normalized expectations of the full compact law,
with arbitrary noncentral transports.
\end{theorem}
\begin{proof}
In the numerator of the expectation the edge matrix is
\(M_e(\lambda)=(1-\lambda)I+\lambda T_e\). Orthogonality gives
the exact matrix identity
\[
 M_e(\lambda)^{\mathsf T}M_e(\lambda)
   =I+\lambda(\lambda-1)(2I-T_e-T_e^{\mathsf T}).
 \tag{V54.7}\label{r54:affine}
\]
The last parenthesis is
\((I-T_e)^{\mathsf T}(I-T_e)\succeq0\), of norm \(a_e^2\).
Thus for \(0\le\lambda\le1\), \(\|M_e(\lambda)\|\le1\);
outside this interval its squared norm is \(1+q_\lambda a_e^2\).
In all cases choose \(b_e=\sqrt{1+q_\lambda a_e^2}\ge1\)
in Lemma \ref{r54:matrix-domination}.

Dividing its scalar majorant by the \emph{same} coherent
partition function gives
\[
 \frac{\mathcal Z((b_eI)_e)}{\mathcal Z(I)}
       =\E_{\mu_0}\exp\{\gamma\sum_e(b_e-1)z_x\cdot z_y\}
       \le e^{\gamma\sum_e(b_e-1)}.
\]
The constants \(e^{-\gamma|E|}\) cancel in the likelihood ratio.
Finally use \(\sqrt{1+t}-1\le t/2\) and
\(\gamma\sum a_e^2\le4\Gamma\).
\end{proof}

The estimate for \(\lambda=1\), \(Q_T\le Q_0\), is inherited
orthogonal domination. The new utility is that the \(\lambda=2\)
numerator is also controlled: \(2T-I\) is not orthogonal, but
its norm differs from one only quadratically in \(I-T\).
A crude linear bound \(\|2T-I\|\le1+2\|T-I\|\) would lose this
quadratic connection budget.

\subsection{The actual normalizer and both entropy directions}

Let \(\rho_e=\E_{\mu_0}(z_x\cdot z_y)\).
Global \(O(4)\) invariance gives
\(\E_{\mu_0}(z_yz_x^{\mathsf T})=(\rho_e/4)I\).
Also \(0\le\rho_e\le1\): the upper bound is pointwise, and
the lower bound follows by differentiating the positive
scalar-coupling pairing series in that edge coupling.
Set \(\Gamma_\rho=\gamma\sum_e\rho_ed_e\le\Gamma\).

\begin{corollary}[A normalized connection comparison with paid tails]
\label{r54:likelihood}
With \(L_T=d\mu_T/d\mu_0\) and
\(J_\Gamma=e^{6\Gamma}-1\),
\[
 \begin{gathered}
 e^{-\Gamma_\rho}\le Q_T/Q_0\le1,\qquad
 L_T=e^{\Phi_T}/\E_{\mu_0}e^{\Phi_T},\\
 D(\mu_0\Vert\mu_T)\le\Gamma_\rho,\qquad
 \E_{\mu_0}L_T^2\le e^{4\Gamma+2\Gamma_\rho}\le e^{6\Gamma},\qquad
 D(\mu_T\Vert\mu_0)\le6\Gamma .
 \end{gathered}
 \tag{V54.8}\label{r54:likelihood-eq}
\]
Consequently, for any \(F\in L^2(\mu_0)\),
\[
 |\E_{\mu_T}F-\E_{\mu_0}F|
                      \le\sqrt{J_\Gamma}\,\|F\|_{L^2(\mu_0)},
 \qquad
 \|\mu_T-\mu_0\|_{\rm TV}\le\min\{1,\sqrt{\Gamma/2}\}.
 \tag{V54.9}\label{r54:weighted}
\]
\end{corollary}
\begin{proof}
The preceding tensor identity gives
\(\E_{\mu_0}\Phi_T=-\Gamma_\rho\). Jensen supplies the lower
normalizer bound; \eqref{r54:mgf-eq} at one gives the upper.
Exactly,
\(D(\mu_0\Vert\mu_T)=\Gamma_\rho+\log(Q_T/Q_0)\).
For the second likelihood moment, use \eqref{r54:mgf-eq}
at two and the lower normalizer bound.
Jensen under \(\mu_T\) then gives
\[
 D(\mu_T\Vert\mu_0)=\E_{\mu_T}\log L_T
       \le\log\E_{\mu_T}L_T=\log\E_{\mu_0}L_T^2.
\]
Finally \(\E_{\mu_0}(L_T-1)^2=\E_{\mu_0}L_T^2-1\).
Cauchy--Schwarz proves the unbounded-test estimate, and
Pinsker applied to the opposite entropy direction proves
the stated total-variation bound.
\end{proof}

This does not use a Gaussian replacement for the transported law,
or differentiate V53's fixed-support source remainder.
The full current, all signs, and the transported normalizer
are present. The bounds are useful as small errors only when
the displayed connection budget is small.

\subsection{Return of the transported native plaquette energies}

Now let the graph be V53's even cubic torus. In the chosen frame put
\[
 X_e=\gamma(1-z_x\cdot z_y),\qquad
 Y_e=\gamma(1-z_x\cdot T_ez_y),\qquad
 \kappa_e=\gamma a_e^2 .
\]
The \(Y_e\) are the actual transported plaquette deficits, not
the unchanged reference \(X_e\) inserted into a new law.
Recall \(r=(n-1)/(3n)\) and \(\|X_e\|_p\le A_pr\) from V51.

\begin{lemma}[The source itself changes, with an explicit cost]
\label{r54:source}
Under the coherent law,
\[
 \begin{gathered}
 0\le Y_e\le2X_e+\kappa_e,\qquad
 \sum_e\kappa_e\le4\Gamma,\\
 \|Y_e-X_e\|_2
       \le2\sqrt{r\kappa_e}+\kappa_e/2
       \le d_\Gamma:=4\sqrt{r\Gamma}+2\Gamma .
 \end{gathered}
 \tag{V54.10}\label{r54:source-eq}
\]
For a set \(S\) of at most \(k\) edges,
\[
       \sum_{e\in S}\E_{\mu_0}|Y_e-X_e|
                     \le4\sqrt{kr\Gamma}+2\Gamma .
 \tag{V54.11}\label{r54:sum-source}
\]
\end{lemma}
\begin{proof}
Write \(v=z_x-z_y\), \(w=(I-T_e)z_y\).
Then \(Y_e=\gamma|v+w|^2/2\),
\(X_e=\gamma|v|^2/2\), and \(|w|\le a_e\).
The elementary square inequality gives \(Y_e\le2X_e+\kappa_e\).
Expanding the square also gives
\[
 |Y_e-X_e|\le\sqrt{2\kappa_e X_e}+\kappa_e/2.
\]
Use \(\E X_e\le2r\) to bound the \(L^2\) norm of the square-root
term. The budget bound follows from \eqref{r54:SO4};
summing and applying Cauchy--Schwarz over \(S\) proves
\eqref{r54:sum-source}.
\end{proof}

\begin{theorem}[Native noncentral source and covariance comparison]
\label{r54:return}
Let \(|F|\le M_F\) with \(\ell^1\)-Lipschitz constant \(L_F\),
on at most \(k\) plaquette energies. Then
\[
 \begin{split}
 |\E_{\mu_T}F(Y_S)-\E_G F(q_S^G)|
 \le{}&
  2M_F\min\{1,\sqrt{\Gamma/2}\}
       +L_F\{4\sqrt{kr\Gamma}+2\Gamma\}\\
 &+\frac{256cL_Fkr}{\gamma}
       +2M_F\sqrt{6k^2(k+1)\eta_\gamma}.
 \end{split}
 \tag{V54.12}\label{r54:bounded}
\]
Here \(G,\eta_\gamma,c\) are exactly V53's coherent reference data.
For an unbounded pair, define
\[
 \begin{split}
 u_\Gamma&=\sqrt{J_\Gamma}(2A_2r+4\Gamma)+d_\Gamma,\\
 \mathcal C(\Gamma)&=
       \sqrt{J_\Gamma}(2A_4r+4\Gamma)^2
       +d_\Gamma(3A_2r+4\Gamma)
       +u_\Gamma(4r+u_\Gamma).
 \end{split}
 \tag{V54.13}\label{r54:C}
\]
For \(\gamma\ge1\), including coincident edges,
\[
 \left|\Cov_{\mu_T}(Y_e,Y_f)-\frac32R_{ef}^2\right|
       \le \mathcal C(\Gamma)
                       +512\{\gamma^{-1}+\eta_\gamma^{1/4}\}.
 \tag{V54.14}\label{r54:covariance}
\]
Since \(r\le1/3\), \(\mathcal C(\Gamma)=O(\sqrt\Gamma)\)
as \(\Gamma\downarrow0\), uniformly in volume and edge positions.
\end{theorem}
\begin{proof}
For the bounded assertion, first change from \(\mu_T\) to
\(\mu_0\) with the observable \(F(Y_S)\) unchanged. Its error is
bounded by twice its supremum norm times total variation.
Then change the source from \(Y\) to \(X\) using
\eqref{r54:sum-source}, and finally apply V53's
\eqref{r53:bounded}.

For the unbounded assertion use, under \(\mu_0\),
\[
 \|Y_e\|_p\le2A_pr+\kappa_e\le2A_pr+4\Gamma .
\]
The likelihood estimate \eqref{r54:weighted}, with
\(F=Y_eY_f\), and Hölder with fourth moments give
\[
 |\E_{\mu_T}Y_eY_f-\E_{\mu_0}Y_eY_f|
          \le\sqrt{J_\Gamma}(2A_4r+4\Gamma)^2 .
\]
The pointwise decomposition
\(Y_eY_f-X_eX_f=(Y_e-X_e)Y_f+X_e(Y_f-X_f)\)
and Cauchy--Schwarz cost at most
\(d_\Gamma(3A_2r+4\Gamma)\).
For each mean, the same likelihood and source bounds give
\[
 |\E_{\mu_T}Y_e-\E_{\mu_0}X_e|\le u_\Gamma .
\]
Both coherent means are at most \(2r\); therefore the two
products of means differ by at most \(u_\Gamma(4r+u_\Gamma)\).
Combining these three errors compares the two actual covariances.
Apply \eqref{r53:covariance} for the coherent one. All moments used
above are genuine compact-law moments, and all constants are
independent of separation. The formula also covers \(e=f\).
\end{proof}

\subsection{A genuinely noncentral winding example and its quantitative limit}

For the matched flat connection with one winding direction, set
\[
 T_{x,x+\hat1}z=e^{(\theta/N)v}z e^{-(\theta/N)v},
 \qquad |v|=1,\quad v\in\operatorname{Im}\mathbb H,
 \qquad T_e=I\ \hbox{in the other directions}.
\]
The real transport has two fixed directions and a rotation by
\(2\theta/N\) on the other plane. There are \(N^3\) edges in
the winding direction, so this explicit frame gives
\[
 \Gamma=\frac{\gamma N^3}{2}\{1-\cos(2\theta/N)\}
                       \le\gamma N\theta^2 .
 \tag{V54.15}\label{r54:winding}
\]
Its loop transport is nontrivial if \(\theta\notin\pi\mathbb Z\).
No periodic gauge transformation makes it coherent in that case.
Nevertheless, along any sequence with \(\gamma\to\infty\) and
\(\gamma N\theta^2\to0\), \eqref{r54:covariance} proves convergence
of the actual transported native pair to \(3R_{ef}^2/2\), uniformly
in edge positions. The conclusion includes noncentral holonomies
at every finite stage.

This is a proved sufficient shrinking window, not a bound on
all fixed noncentral winding angles. At fixed nonzero \(\theta\),
the displayed budget grows like \(\gamma N\); the likelihood
bounds remain valid but are not a small comparison. Nor is small
individual link angle alone enough to make the sum small.
A general left/right connection may be tested in any declared
vertex frame by the same \(\Gamma\). Its probability under the
original Wilson exterior distribution has not been estimated.

For example, order \(\gamma^{-1/2}\) rotations on an extensive
family of edges can give an order-\(|E|\) budget, rather than
a vanishing one. Thus this shrinking window is not proposed
as a substitute for the typical thermal exterior problem.
That problem may require localized likelihood control or
a connection-adapted reference retaining V46's shifted fields
and collective orientation, rather than another comparison
to the coherent reference with a smaller declared window.

\subsection{Remaining obligations and the next checkpoint}

The new result pays an actual normalized transported-current
estimate, including unbounded-source return, on a quantitative
connection budget. It is not merely a coherent Gaussian determinant.
It also does not resolve V53's spatial summability problem:
summing the uniform error in \eqref{r54:covariance} can still
cost the number of plaquettes. Correlation positivity was not
silently imported to avoid that cost.

The next upstream issue is to obtain useful connection budgets
or a stronger fixed-holonomy comparison under the actual exterior
law. This must retain its predictor covariance and the physical
source family. A finite-dimensional likelihood comparison gives
no physical-time contraction or interacting four-dimensional
continuum construction. The full Yang--Mills mass gap remains
unproved.

The diagnostic \path{verify_ym_restart_v54_transport_likelihood.py}
checks noncommuting rational orthogonal transports, affine norm
identities, finite signed diagram contractions, exact discrete
likelihood normalization, and native source-return constants.
It is not a numerical proof of the infinite pairing expansion
or continuum claims. V53's algebraic regression and exact source
preservation are checked separately. Companion review is due
on the next append, V55; the regular broad literature sweep is
due V60, with the cadence retained across archive restarts.
% END CONSOLIDATED V54 APPEND


% BEGIN CONSOLIDATED V55 APPEND
\clearpage
\section{V55: Spatially summable winding-source comparison and the remaining native defect}
\label{r55:start}

V54 pays the actual transported likelihood on a global connection
budget. That budget is extensive on many thermal exteriors and grows
even for a fixed winding angle. We do not tighten its shrinking
window again. Instead, this installment proves that the
\emph{connection-adapted Gaussian source covariance} of a fixed
commuting winding background has an absolute row-sum difference
\(O(N^{-1})\) from the coherent coefficient, on the three-dimensional
cubic corridor. The bound allows arbitrary signed source profiles
and arbitrary support size. A global linear-field likelihood is
not needed; in fact the two Gaussian field laws have singular
supports. The nonlinear compact-law remainder is kept as a
separate, explicitly defined source operator.

\subsection{Scheduled companion review and the precise transferable input}

At the V55 checkpoint the repository NS V26, SM--Einstein V72,
Shell Mixing V16 and Parallel YM V5 have the same hashes as at
V50. Their terminal proof and scope statements were reread.
The latest Source Selection V35, Native Stationarity V38 and
Exact-Action Einstein Bridge V28 in Downloads are also unchanged.
Their stationary-reversal, boundary-obstruction and enlarged-profile
results retain their original finite-source hypotheses; none
estimates the Wilson exterior predictor covariance.

There is a new Perturbative ESM V31, beyond the previously
reviewed V27. Its scalar Euler calculation identifies a polynomial
removable sector, an explicit projector onto it, and a remainder
recovered from the \emph{complete} compatibility defect.
Its coupled extraction still requires the metric-marked action
row and reference terms. Thus a solved isolated parent is not a
proof that the combined residual vanishes. The reviewed V31
kernel, projector, reconstruction, coupled-extraction and
source-return proofs do not provide a nonperturbative YM measure
or mixing constant.

The usable transfer is this separation of a calculable coefficient
from the actual unremoved defect, with all return terms retained.
The supplied transfer suggestions also recommend examining
derivative structure before generic norms. Here the incidence
factors on both sides of the inverse Laplacian are retained.
They convert an infrared-sensitive inverse into a projection,
and its squared entries into a summable source kernel. No
unproved extra vanishing of a YM coefficient is assumed.

The nonduplication audit retains V48's fixed-graph native
Gaussian return and harmonic jump, V49's winding determinant
and Green bound, and V53's coherent compact-law pair estimate.
Those results are not reproved or promoted to joint-limit
theorems. The new statement is an absolute spatial row-sum
comparison of the \emph{winding-dependent source coefficients}.
The older Schur and restricted-fibre projection estimates
concern different operators and do not supply this comparison.

For primary-source context the flat connection Fourier spectrum
was checked again against Y.~Lin, S.~Wan and H.~Zhang,
\emph{Connection Laplacian on discrete tori with converging property},
Theorem 1.2,
\url{https://arxiv.org/html/2403.06105}.
Only its spectral subject and scope are used as background.
The projection and source estimates below are proved directly;
a connection-Laplacian convergence theorem is not imported as
a nonlinear Wilson remainder bound. The next scheduled companion
review and broad literature sweep are both V60.

\subsection{The actual flat winding family and its two gradient projections}

Let \(N\ge4\) be even, \(n=N^3\), and let \(E_N\) consist of
the \(m=3n\) positively oriented edges of the cubic torus.
Take
\[
 \varphi\in[-\pi,\pi]^3\setminus\{0\},\qquad
 T_dz=e^{(\varphi_d/2N)\mathbf i}z
                         e^{-(\varphi_d/2N)\mathbf i}.
 \tag{V55.1}\label{r55:transport}
\]
These are V48's matched lower/upper Wilson links, with charged
adjoint winding phases \(e^{i\varphi_d}\). They commute, and at
least one winding is nontrivial. The compact law is exactly
\[
 d\mu_{N,\gamma,\varphi}
 =Q_{N,\gamma,\varphi}^{-1}
   \exp\{-\gamma\sum_{x,d}(1-z_x\cdot T_dz_{x+\widehat d})\}
                                   \prod_xdH(z_x).
 \tag{V55.2}\label{r55:law}
\]
No nonperiodic gauge removal or replacement exterior law occurs.

On complex scalar vertex fields use
\[
 \begin{aligned}
 (B_\varphi f)_{x,d}&=f_x-e^{i\varphi_d/N}f_{x+\widehat d},
       & L_\varphi&=B_\varphi^*B_\varphi,\\
 P_\varphi&=B_\varphi L_\varphi^{-1}B_\varphi^*,
       & R&=B_0L_0^\dagger B_0^* .
 \end{aligned}
 \tag{V55.3}\label{r55:projectors}
\]
Here \(L_\varphi>0\), whereas the dagger removes the scalar
constant mode of \(L_0\). The real matrix \(R\), also regarded
as a complex matrix when needed, is the earlier \(R_0\).
Both displayed edge operators are orthogonal projections,
of ranks \(n\) and \(n-1\), respectively. Put
\[
 c_N=\frac1n\sum_{p\ne0}\lambda_0(p)^{-1}\le\frac{13}{16},
 \qquad r=R_{ee}=\frac{n-1}{3n},
 \qquad
 \varepsilon_N(\varphi)
   =\left(\frac1{N^3}+\frac{13|\varphi|^2}{8N^2}\right)^{1/2}.
 \tag{V55.4}\label{r55:epsilon}
\]
The Green bound and edge-transitive diagonal are inherited
from V49--V51. The Euclidean norm is used for \(\varphi\).

\subsection{A row-square estimate that keeps the harmonic mode}

\begin{lemma}[Elementary rank-one projection displacement]
\label{r55:rankone}
For nonzero \(u,v\in\mathbb C^a\), write \(P_u=uu^*/|u|^2\).
Then
\[
 \|P_u-P_v\|_{\rm F}^2
 =2\left(1-\frac{|u^*v|^2}{|u|^2|v|^2}\right)
 \le \frac{2|u-v|^2}{|v|^2}.
 \tag{V55.5}\label{r55:rankone-eq}
\]
\end{lemma}
\begin{proof}
Expand the squared Frobenius norm using the two rank-one
projection traces. The expression in parentheses is
\(\|(I-P_u)v\|^2/|v|^2\).
Since \((I-P_u)u=0\), contraction of the orthogonal
projection bounds its numerator by \(|v-u|^2\).
\end{proof}

\begin{theorem}[Volume-normalized projector displacement]
\label{r55:projection}
For the full, unpinned edge projections in
\eqref{r55:projectors},
\[
 \begin{split}
 \|P_\varphi-R\|_{\rm F}^2
       &\le 1+2c_N N|\varphi|^2,\\
 \sup_{e\in E_N}\sum_{f\in E_N}
                    |(P_\varphi-R)_{ef}|^2
       &\le \frac1n+\frac{2c_N|\varphi|^2}{N^2}
        \le\varepsilon_N(\varphi)^2 .
 \end{split}
 \tag{V55.6}\label{r55:projection-eq}
\]
The unit term is the charged zero-momentum mode; it is not
discarded by a common zero-mode convention.
\end{theorem}
\begin{proof}
Use the normalized vertex Fourier basis, with
\(p\in\{-N/2,\ldots,N/2-1\}^3\) and \(k=2\pi p/N\).
At each momentum the incidence is the three-component vector
\[
 b_\varphi(p)
   =\bigl(1-e^{i(k_d+\varphi_d/N)}\bigr)_{d=1}^3,
 \qquad |b_\varphi(p)-b_0(p)|\le|\varphi|/N .
 \tag{V55.7}\label{r55:symbol}
\]
Because \(\varphi\ne0\) in the stated cube,
\(b_\varphi(p)\ne0\) for every allowed \(p\).
The Fourier block of \(P_\varphi\) is \(P_{b_\varphi(p)}\).
For \(p\ne0\), the block of \(R\) is \(P_{b_0(p)}\), with
\(|b_0(p)|^2=\lambda_0(p)\). Thus Lemma
\ref{r55:rankone} bounds that block's squared Frobenius
difference by \(2|\varphi|^2/(N^2\lambda_0(p))\).

At \(p=0\), the block of \(R\) is zero and that of
\(P_\varphi\) has rank one, so its squared Frobenius
difference is exactly one. Parseval and the definition
of \(c_N\) prove the first assertion.

Both edge matrices commute with lattice translations.
For a fixed direction \(d\), the squared row norm of
their difference is independent of \(x\). The sum over
the three directional row norms is exactly
\(n^{-1}\|P_\varphi-R\|_{\rm F}^2\).
Each directional term is nonnegative and bounded by that sum.
This proves the second assertion without assuming edge
rotation symmetry of the winding background.
\end{proof}

Small row norms do \emph{not} mean a small linear-field
operator norm or a global Gaussian likelihood. In fact
\[
                  \|P_\varphi-R\|_{\ell^2\to\ell^2}=1 .
 \tag{V55.8}\label{r55:op-jump}
\]
Indeed, a unit vector in the \(p=0\) range of \(P_\varphi\)
is killed by \(R\), and a difference of orthogonal projections
has operator norm at most one, as follows from
\(-I\preceq P_\varphi-R\preceq I\).
The centered proper complex Gaussian edge laws with covariance
\(P_\varphi\) and \(R\) are mutually singular: the latter is
supported on a complex subspace of dimension \(n-1\), which
has zero measure for a nondegenerate Gaussian on the
\(n\)-dimensional range of the former. This concerns the
two Gaussian edge laws, not the mutually absolutely continuous
finite compact spin laws. It explains why a Gaussian
likelihood comparison is not the argument below.

\subsection{Absolute spatial summability for the native quadratic coefficient}

Keep V48's neutral real Gaussian gradient \(a\) and independent
proper complex charged gradient \(w\), with
\[
 \E a_ea_f=R_{ef},\qquad
 \E w_e\overline{w_f}=2(P_\varphi)_{ef},\qquad
 \E w_ew_f=0 .
\]
The source is the actual leading native insertion
\(q_e^\varphi=(a_e^2+|w_e|^2)/2\), not the coherent insertion
evaluated in the charged law. Its mean and covariance are
the inherited Wick coefficients
\[
 \begin{aligned}
 M^\varphi_e&=\tfrac12r+(P_\varphi)_{ee},\\
 K^\varphi_{ef}&=\tfrac12R_{ef}^2+|(P_\varphi)_{ef}|^2,
        & K^0_{ef}&=\tfrac32R_{ef}^2 .
 \end{aligned}
 \tag{V55.9}\label{r55:K}
\]
They are coefficients of normalized source cumulants.

\begin{theorem}[Summable covariance comparison at fixed winding]
\label{r55:Schur}
Uniformly for the winding phases in \eqref{r55:transport},
\[
 \begin{split}
 \sup_e\sum_f|K^\varphi_{ef}-K^0_{ef}|
                         &\le2\varepsilon_N(\varphi),\\
 \|K^\varphi-K^0\|_{\ell^2\to\ell^2}
                         &\le2\varepsilon_N(\varphi),\\
 \sup_e|M^\varphi_e-\tfrac32r|
                         &\le\varepsilon_N(\varphi).
 \end{split}
 \tag{V55.10}\label{r55:Schur-eq}
\]
In particular the first two differences are \(O(N^{-1})\)
uniformly on the full fixed phase cube, without a condition
that the V54 budget \(\Gamma\) tend to zero.
For all real source profiles \(s,t\) on all the edges,
\[
 \left|\Cov(q^\varphi(s),q^\varphi(t))
                -\Cov(q^0(s),q^0(t))\right|
 \le2\varepsilon_N(\varphi)\|s\|_2\|t\|_2,
 \quad q^\varphi(s)=\sum_es_eq^\varphi_e .
 \tag{V55.11}\label{r55:source-form}
\]
There is no support-cardinality multiplier.
\end{theorem}
\begin{proof}
For complex \(A,B\),
\(\bigl||A|^2-|B|^2\bigr|\le|A-B|(|A|+|B|)\).
The neutral terms in \eqref{r55:K} cancel. Apply this
inequality entrywise and Cauchy--Schwarz in the column index:
\[
 \begin{split}
 \sum_f\bigl||P_{ef}|^2-|R_{ef}|^2\bigr|
 &\le
  \left(\sum_f|P_{ef}-R_{ef}|^2\right)^{1/2}
     \left(\sum_f(|P_{ef}|+|R_{ef}|)^2\right)^{1/2}\\
 &\le
  \left(\sum_f|P_{ef}-R_{ef}|^2\right)^{1/2}
                    \{\sqrt{P_{ee}}+\sqrt{R_{ee}}\}\\
 &\le 2\varepsilon_N(\varphi).
 \end{split}
\]
Here \(P=P_\varphi\), and
\(\sum_f|P_{ef}|^2=P_{ee}\le1\), similarly for \(R\).
Use Theorem \ref{r55:projection}.
The difference \(K^\varphi-K^0\) is real symmetric, so
the absolute row and column sums have the same bound.
The elementary Schur test gives its \(\ell^2\) norm bound.
For completeness, that test follows by applying weighted
Cauchy--Schwarz to each row of
\(\sum_f A_{ef}t_f\) with weights \(|A_{ef}|\), then
summing and using the column sums. The bilinear
bound follows immediately. A diagonal entry of
\(P_\varphi-R\) is bounded by its row \(\ell^2\) norm,
which proves the mean assertion.
\end{proof}

The sign in this proof is not a spin-correlation inequality.
It uses explicit squared entries of the declared Gaussian
projection only. No positivity of transported compact-law
covariances has been introduced.

There are useful exact checks on the harmonic contribution:
\[
 \begin{aligned}
 \sum_fK^\varphi_{ef}&=M^\varphi_e,
       & \sum_fK^0_{ef}&=\tfrac32r,\\
 \sum_e(M^\varphi_e-\tfrac32r)&=1,
       & \sum_{e,f}(K^\varphi_{ef}-K^0_{ef})&=1 .
 \end{aligned}
 \tag{V55.12}\label{r55:rank-audit}
\]
They follow from idempotence and the projection ranks.
For the normalized constant source \(s_e=m^{-1/2}\),
the variance difference is exactly \(1/m\), rather than
zero. Thus the \(p=0\) return survives at every finite
regulator, consistently with V48, while its normalized
source effect vanishes as \(N\) grows. Deleting that mode
would give incorrect finite traces despite a plausible
large-volume conclusion.

Both \(K^\varphi\) and \(K^0\) are covariance matrices, and
their nonnegative entries and row identities give
\(\|K^\varphi\|_{2\to2}\le1+r/2\le7/6\) and
\(\|K^0\|_{2\to2}=3r/2\le1/2\).
These are upper source bounds, not lower coercivity or
physical-time spectral gaps.

\subsection{What returns to the compact law, and in which order}

Define the native scaled deficits under \eqref{r55:law} by
\[
 Y_e=\gamma(1-z_x\cdot T_dz_{x+\widehat d}),\qquad
 C_{N,\gamma,\varphi}(e,f)=\Cov_{\mu_{N,\gamma,\varphi}}(Y_e,Y_f).
\]
V48 already proves that, at fixed \(N\) and noncentral
\(\varphi\), each entry satisfies
\(C_{N,\gamma,\varphi}(e,f)=K^\varphi_{ef}
+O_{N,\varphi}(\gamma^{-1/2})\).
It proves this with source insertions under the full
normalized compact law. We consume that theorem instead
of differentiating an unmarked determinant remainder.

Write \(\|A\|_{\rm row}=\sup_e\sum_f|A_{ef}|\).
The exact unresolved absolute-row defect is
\[
 {\cal D}_{N,\gamma,\varphi}
       =\|C_{N,\gamma,\varphi}-K^\varphi\|_{\rm row}.
 \tag{V55.13}\label{r55:native-defect}
\]
At fixed \(N,\varphi\), there are finitely many pairs, so
V48 gives a finite constant \(A_{N,\varphi}\) and threshold
such that
\[
 {\cal D}_{N,\gamma,\varphi}
       \le 3N^3 A_{N,\varphi}\gamma^{-1/2}.
 \tag{V55.14}\label{r55:fixed}
\]
Neither the constant nor the threshold is asserted uniform.

\begin{corollary}[A source-normalized sequential native return]
\label{r55:native}
For every finite regulator and coupling,
\[
 \sup_e\sum_f|C_{N,\gamma,\varphi}(e,f)-K^0_{ef}|
       \le {\cal D}_{N,\gamma,\varphi}
                             +2\varepsilon_N(\varphi).
 \tag{V55.15}\label{r55:decomposition}
\]
The same bound holds for the \(\ell^2\) operator norm.
For fixed \(\varphi\ne0\) in the phase cube,
\[
 \lim_{\substack{N\to\infty\\N\ {\rm even}}}
       \lim_{\gamma\to\infty}
       \|C_{N,\gamma,\varphi}-K^0\|_{\ell^2\to\ell^2}=0 .
 \tag{V55.16}\label{r55:sequential}
\]
The inner limit is a norm of matrices on that fixed
regulator; the outer limit is a limit of nonnegative numbers.
\end{corollary}
\begin{proof}
The first assertion is the triangle inequality and
\eqref{r55:Schur-eq}; symmetry gives the operator bound.
At fixed \(N,\varphi\), \eqref{r55:fixed} makes the defect
vanish. Norm continuity in finite dimension identifies
the inner norm limit with \(\|K^\varphi-K^0\|\).
Finally \(\varepsilon_N(\varphi)\to0\).
\end{proof}

This is stronger than a fixed-pair coefficient comparison
after the sequential limit: it controls unit-\(\ell^2\)
source profiles supported anywhere on the growing corridor.
It is not a joint estimate along \(\gamma=\gamma(N)\).
In particular \eqref{r55:fixed} cannot be substituted into
a logarithmic running-coupling trajectory with its
unknown \(A_{N,\varphi}\) suppressed.

\subsection{The upstream target after the comparison}

The fixed commuting winding coefficient is no longer an
uncomputed summability obstacle: it obeys the explicit
source-norm estimate \eqref{r55:Schur-eq}. The unresolved
quantity on this family is the \emph{actual nonlinear}
defect \eqref{r55:native-defect}, or a weaker source-relative
bilinear defect if absolute row control is unnecessarily
strong. A volume-uniform estimate for it must come from
the compact law, not from the Gaussian projection theorem.
V53's uniform \emph{entrywise} coherent error cannot be
summed for this purpose without paying its spatial cost.

Moreover, general thermal exteriors need not be flat,
commuting, matched, or translation invariant. The proof
of \eqref{r55:projection-eq} uses precisely those features;
it does not estimate their Wilson probability or extend
to arbitrary \(SO(4)\) connections by notation alone.
The original exterior predictor covariance, physical
source coverage and positivity requirements remain.
An exterior mixture would require the conditional
covariance and the covariance of conditional means,
as already retained in V46.

Thus the next useful step is a connection-adapted compact
source-operator bound, first on this paid winding family
or on a genuinely controlled typical exterior class.
Repeating the winding determinant, dropping the harmonic
rank, or naming the remaining defect a small error would
not advance that task. No interacting four-dimensional
continuum construction or physical Hamiltonian mass gap
has been proved here.

The diagnostic
\path{verify_ym_restart_v55_winding_source_schur.py}
checks exact rational projection and Wick identities and
56 deterministic Fourier/real-space cases. Trigonometric
checks are not interval certificates or substitutes for
proofs. V54's algebra, source preservation and typesetting
are checked separately. Both scheduled reviews are next V60.
% END CONSOLIDATED V55 APPEND


% BEGIN CONSOLIDATED V56 APPEND
\clearpage
\section{V56: Source-resolved Ward identities and compact covariance bounds}
\label{r56:start}

V55 controls a winding-dependent Gaussian coefficient, not its
nonlinear compact-law remainder. A necessary upstream problem is
already present at zero winding: V53's compact covariance error is
uniform in the pair but is not a source-operator estimate.
We now obtain such an estimate in the \emph{actual coherent compact
law}, for arbitrary signed profiles on all corridor edges.

The new proof has two parts. First, V53's entropy density and an
infrared covariance ceiling imply a volume-independent estimate
for an entire quadratic response matrix. Second, a source-dependent
smooth vector field gives an exact identity for the native
covariance. Its nonlinear returns are explicit polynomial
observables, not an unnamed difference of two covariance matrices.
Elementary moment estimates pay these returns with visible volume
costs. The resulting sufficient regime \(\gamma/n\to\infty\),
\(n=N^3\), is \emph{not} the intended logarithmic bare trajectory.
The purpose is to identify and estimate the actual source operator,
and to locate the remaining loss before attempting winding or
typical-exterior extensions.

The previous goal turn was progress: V55's source comparison and
verification are present and their source hash was rechecked.
The V55 companion review remains current; both scheduled reviews
are next due V60. Targeted searches of low-temperature classical
spin results led back to the already credited
Giuliani--Ott paper,
\url{https://arxiv.org/abs/2302.07299}.
Its integration-by-parts and moment method is relevant background.
No selected-state expansion, quantum Heisenberg theorem, or
unverified correlation-decay statement is imported here.
The analytic inputs are exactly V50's finite coherent infrared
bound, V51's moments, and V53's unconditioned entropy.
The general entropy covariance inequality is elementary;
its use below on the full root-averaged compact corridor and the
source-resolved Ward return are the programme-specific additions.
Earlier fixed-endpoint covariance estimates in f16 V36 keep
their different chart, boundary and size hypotheses.

\subsection{A compact covariance ceiling on the common gradient image}

Keep the coherent measure \(\mu=\mu_{N,\gamma}\), even \(N\ge4\),
\(n=N^3\), \(m=3n\), \(\gamma>0\), and
\[
 \begin{gathered}
 B:\mathbb R^n\longrightarrow\mathbb R^m,\qquad
 L=B^{\mathsf T}B,\qquad C=L^\dagger,\\
 R=BCB^{\mathsf T},\qquad c=C_{xx}\le13/16,\qquad
 r=R_{ee}=(n-1)/(3n).
 \end{gathered}
\]
Choose the root \(o\) uniformly, independently of \(z\sim\mu\),
with no chart conditioning. Write
\[
 \begin{aligned}
 U_e^o&=\sqrt\gamma\,\operatorname{Im}\{(z_x-z_y)z_o^{-1}\},
 & V_e^o&=\sqrt\gamma\,(z_x-z_y)\cdot z_o,\\
 \Gamma_{ef}&=\gamma\,\E_\mu
                  [(z_x-z_y)\cdot(z_u-z_v)],
 & e&=(x,y),\quad f=(u,v).
 \end{aligned}
 \tag{V56.1}\label{r56:fields}
\]
The expectation for \(U,V\) includes the independent root.

Global quaternion conjugation rotates the three components of
\(U\) simultaneously and leaves its law invariant. Thus its mean
is zero and there is a real scalar matrix \(S\) such that
\[
 \E U_e^{o,a}U_f^{o,b}=\delta_{ab}S_{ef},\qquad
 H_{ef}=\E V_e^oV_f^o,\qquad
 \Gamma=3S+H.
 \tag{V56.2}\label{r56:decomposition}
\]
This is an exact decomposition of Euclidean inner products by
right multiplication by \(z_o^{-1}\), not a Gaussian Wick rule.
The matrices \(S,H,\Gamma\) are positive semidefinite and have
range in \(\operatorname{ran}B\). Root averaging gives translation
invariance.

\begin{lemma}[A dimension-correct compact covariance bound]
\label{r56:ceiling}
On the full edge space,
\[
 0\preceq3S\preceq\Gamma\preceq4R,\qquad
 0\preceq H\preceq4R,\qquad
 H_{ee}\le h_\gamma:=\frac{512cr}{\gamma}.
 \tag{V56.3}\label{r56:ceiling-eq}
\]
\end{lemma}
\begin{proof}
For each ambient spin component, V50's infrared bound applied
to the zero-sum site source \(B^{\mathsf T}t\) gives
\[
 \gamma\,\E_\mu\left|\sum_et_e(z_x-z_y)\right|^2
                         \le4\,t^{\mathsf T}Rt .
\]
The factor four is the number of ambient components.
The positivity and decomposition in \eqref{r56:decomposition}
then give the first two matrix inequalities.
Finally \( (V_e^o)^2=2\Delta_e^o\), and V53 retains the
full-law bound \(\E\Delta_e^o\le256cr/\gamma\).
This proves the diagonal estimate without a common chart event.
\end{proof}

\subsection{Entropy density controls a whole quadratic response matrix}

Let \(\overline P^*\) be the law of \(U^o\), and let \(G\) be
standard Gaussian on
\(\mathcal H=\operatorname{ran}B\otimes\mathbb R^3\).
V53 gives \(D(\overline P^*\Vert G)\le n\eta_\gamma\), where
\(\eta_\gamma\) is exactly \eqref{r53:budget} and is
\(O(\gamma^{-1})\) uniformly in \(N\).
In particular no global small-entropy condition is assumed.

\begin{theorem}[Compact spectral covariance error per volume]
\label{r56:entropy-covariance}
The scalar gradient covariance satisfies
\[
 \|S-R\|_{\rm F}^2\le\frac{64}{27}n\eta_\gamma,\qquad
 \sup_e\sum_f|S_{ef}-R_{ef}|^2
                          \le\frac{64}{27}\eta_\gamma .
 \tag{V56.4}\label{r56:entropy-covariance-eq}
\]
\end{theorem}
\begin{proof}
On \(\mathcal H\), the covariance of \(\overline P^*\) is
\(\Sigma=S\otimes I_3\), with the restriction to the gradient
image understood. It is positive definite: a zero covariance
direction would concentrate the mean-zero law on a Gaussian-null
hyperplane, contradicting finite entropy.
It has eigenvalues at most \(4/3\) by
Lemma \ref{r56:ceiling}.

For a centered probability with covariance \(\Sigma\) and finite
entropy relative to standard Gaussian, comparison with the
centered Gaussian of covariance \(\Sigma\) gives
\[
 D(P\Vert G)\ge
           \tfrac12\{\Tr\Sigma-\dim\mathcal H-\log\det\Sigma\}.
 \tag{V56.5}\label{r56:Gaussian-entropy}
\]
Indeed the log density ratio of those two Gaussians is quadratic;
its expectation agrees under \(P\) and under its covariance-matched
Gaussian. Subtract it from \(D(P\Vert G)\), and the remaining
relative entropy is nonnegative. All quadratic expectations
are finite here.

For \(0<t\le4/3\), Taylor's theorem with
\((t-1-\log t)''=t^{-2}\ge9/16\) along the segment to one gives
\(t-1-\log t\ge9(t-1)^2/32\).
Every eigenvalue of \(S\) on \(\operatorname{ran}B\) occurs
three times in \(\Sigma\). Consequently
\[
 n\eta_\gamma\ge\frac{27}{64}\|S-R\|_{\rm F}^2.
\]
Translation invariance makes each directional squared row norm
independent of the base vertex. The sum of the three such
norms is \(n^{-1}\|S-R\|_{\rm F}^2\); each is nonnegative.
This proves the second assertion, without taking the operator
norm of \(S-R\) to be its entropy density.
\end{proof}

Define the \emph{actual compact quadratic response} and its
Gaussian coefficient by the entrywise products
\[
 {\cal Q}_{ef}=R_{ef}\Gamma_{ef},\qquad
 {\cal Q}^0_{ef}=3R_{ef}^2=2K^0_{ef},\qquad
 d_\gamma=\frac83\sqrt{\eta_\gamma}+h_\gamma .
 \tag{V56.6}\label{r56:quadratic-data}
\]
The matrix \({\cal Q}\) uses compact-law second moments.

\begin{corollary}[A volume-independent response-operator estimate]
\label{r56:quadratic}
\[
 0\preceq{\cal Q}\preceq4(R\circ R),\qquad
 \|{\cal Q}\|_{2\to2}\le4r,\qquad
 \|{\cal Q}-2K^0\|_{2\to2}\le d_\gamma .
 \tag{V56.7}\label{r56:quadratic-eq}
\]
Here \(d_\gamma=O(\gamma^{-1/2})\) uniformly in \(N\).
\end{corollary}
\begin{proof}
The Schur product with the positive matrix \(R\) preserves
positive-semidefinite order. Since \(R\) is a projection,
\(\sum_fR_{ef}^2=r\); the nonnegative symmetric matrix
\(R\circ R\) has operator norm \(r\).
These facts and \eqref{r56:ceiling-eq} prove the first two
assertions.

For the remaining assertion decompose
\[
 {\cal Q}-2K^0=3R\circ(S-R)+R\circ H.
\]
Cauchy--Schwarz in each row and
\eqref{r56:entropy-covariance-eq} give
\[
 \sup_e\sum_f|3R_{ef}(S_{ef}-R_{ef})|
       \le\frac8{\sqrt3}\sqrt{r\eta_\gamma}
       \le\frac83\sqrt{\eta_\gamma}.
\]
Symmetry gives the same operator bound.
For the other term, decompose \(H=\sum_jh_jh_j^{\mathsf T}\).
As \(0\preceq R\preceq I\),
\[
 0\preceq R\circ H
   =\sum_j\operatorname{diag}(h_j)R\operatorname{diag}(h_j)
   \preceq\operatorname{diag}(H_{ee})
   \preceq h_\gamma I .
\]
Add the two operator estimates. No compact energy covariance
sign inequality is used.
\end{proof}

This is a genuine compact-law bound, not a second application
of V55 to a Gaussian. It also does not yet equal the native
energy covariance. The difference is calculated next.

\subsection{A smooth vector field for every signed edge source}

For \(s\in\mathbb R^m\), with arbitrary support and signs, set
\[
 \begin{aligned}
 L_s&=B^{\mathsf T}\operatorname{diag}(s)B,&
 A_s&=CL_s,&
 F_s&=\sum_es_eX_e=\frac\gamma2 z^{\mathsf T}L_sz,\\
 q_{s,x}&=z_x\cdot(A_sz)_x,&
 h_{t,x}&=z_x\cdot(L_tz)_x,&
 h_x&=h_{\mathbf1,x}.
 \end{aligned}
 \tag{V56.8}\label{r56:source-data}
\]
Contraction over the four spin components is understood.
The generally nonsymmetric matrix \(A_s\) has both row and
column sums zero, and \(A_s^{\mathsf T}L=L_s\).
With \(P_x=I-z_xz_x^{\mathsf T}\), define
\[
 ({\cal V}_s)_x=P_x(A_sz)_x,\qquad
 \rho_s=\sum_xq_{s,x}(\gamma h_x-3).
 \tag{V56.9}\label{r56:vector}
\]
These are globally smooth polynomial expressions on the compact
product of spheres, including zero-magnetization configurations.

\begin{theorem}[Source-resolved native Ward identity]
\label{r56:Ward}
For every smooth real \(F\),
\[
 \begin{aligned}
 \operatorname{div}{\cal V}_s
     &=3\Tr A_s-3\sum_xq_{s,x},\\
 {\cal V}_s(\gamma S_0)
     &=2F_s-\gamma\sum_xq_{s,x}h_x,\\
 2\E F_s&=3\Tr A_s+\E\rho_s,\\
 2\Cov(F_s,F)&=\E{\cal V}_sF+\Cov(\rho_s,F).
 \end{aligned}
 \tag{V56.10}\label{r56:Ward-eq}
\]
All expectations use the same normalized compact law.
\end{theorem}
\begin{proof}
The tangential trace at site \(x\) of the derivative of
\(P_x(A_sz)_x\) is \(3(A_s)_{xx}-3q_{s,x}\).
This includes the derivative of the \(x\)-entry of \(A_sz\);
treating that entry as an external fixed vector would omit
the trace term.

The intrinsic gradient of \(S_0\) at \(x\) is \(P_x(Lz)_x\).
Therefore
\[
 {\cal V}_sS_0
 =\langle A_sz,Lz\rangle-\sum_xq_{s,x}h_x
 =z^{\mathsf T}L_sz-\sum_xq_{s,x}h_x.
\]
Here \(CL\) is projection off constants, and \(L_s\) kills
constants, so \(A_s^{\mathsf T}L=L_s\) even for signed \(s\).
Integrating the divergence of
\(F{\cal V}_s e^{-\gamma S_0}\) on the boundaryless compact
product gives
\[
 \E\{{\cal V}_sF+
             (3\Tr A_s-2F_s+\rho_s)F\}=0.
\]
First take \(F=1\), then subtract that identity times
\(\E F\). This proves the last two assertions without
chart-boundary integration or a tilted-measure assumption.
\end{proof}

At \(s=\mathbf1\), \(A_s=I-\mathbf1\mathbf1^{\mathsf T}/n\)
and \(({\cal V}_s)_x=-P_xM\). The scalar \(\rho_s\) is
exactly V51's \(T-3nD\). Thus its collective identity is
a specialization, not a new proof being counted twice.
The nonconstant and sign-changing source directions are new here.

\subsection{The full native covariance and its nonlinear return}

Let \({\cal C}_{ef}=\Cov_\mu(X_e,X_f)\), and define the
bilinear quantities
\[
 {\cal N}(s,t)=\gamma\,\E\sum_xq_{s,x}h_{t,x},\qquad
 {\cal B}(s,t)=\Cov(\rho_s,F_t)-{\cal N}(s,t).
 \tag{V56.11}\label{r56:returns}
\]
The first is a mean of an explicit fourth-degree spin expression;
the second retains its connected return under the actual law.

\begin{proposition}[An exact source-operator decomposition]
\label{r56:native}
For every \(s,t\),
\[
 2s^{\mathsf T}{\cal C}t
       =s^{\mathsf T}{\cal Q}t+{\cal B}(s,t).
 \tag{V56.12}\label{r56:native-eq}
\]
In particular \({\cal B}\) is symmetric, although its two
terms in \eqref{r56:returns} need not separately be symmetric.
\end{proposition}
\begin{proof}
The derivative of \(F_t\) gives
\[
 \begin{split}
 \E{\cal V}_sF_t
 &=\gamma\,\E\langle CL_sz,L_tz\rangle
                                      -{\cal N}(s,t)\\
 &=\sum_{e,f}s_et_fR_{ef}\Gamma_{ef}-{\cal N}(s,t).
 \end{split}
\]
Apply \eqref{r56:Ward-eq} with \(F=F_t\).
Symmetry follows from the symmetric matrices
\({\cal C}\) and \({\cal Q}\), not by discarding the
antisymmetric pieces of the two nonlinear terms.
\end{proof}

The matrix \({\cal B}\) is the remaining native source target.
It is more explicit than V55's covariance difference: its inputs
are the source-dependent spherical score and its local energy
return. Bounding \(\rho_s\) only for the constant source cannot
control this full bank.

To quantify what has actually been paid, define the finite
source norms
\[
 \begin{aligned}
 a&=\sup_{\|s\|_2=1}\|\rho_s-\E\rho_s\|_{L^2(\mu)},\\
 b&=\sup_{\|s\|_2=\|t\|_2=1}|{\cal N}(s,t)|,\qquad
 v=\|{\cal C}\|_{2\to2}.
 \end{aligned}
 \tag{V56.13}\label{r56:ab}
\]
They are finite at every regulator, since all observables are
smooth on a compact space.

\begin{theorem}[A closed bound from the actual nonlinear source bank]
\label{r56:bootstrap}
Put \(W=\sqrt{2r+b/2}+a/2\). Then
\[
 v\le W^2,\qquad
 \|{\cal C}-K^0\|_{2\to2}
              \le\frac12\{d_\gamma+b+aW\}.
 \tag{V56.14}\label{r56:bootstrap-eq}
\]
\end{theorem}
\begin{proof}
The compact covariance is positive semidefinite. Take a unit
maximal-eigenvalue source in \eqref{r56:native-eq}.
Cauchy--Schwarz, \eqref{r56:quadratic-eq}, and
\eqref{r56:ab} give
\(2v\le4r+b+a\sqrt v\).
The positive root of this quadratic inequality is at most
\(\sqrt{2r+b/2}+a/2\), proving the first assertion.
For any unit sources,
\[
 |\Cov(\rho_s,F_t)|\le a\sqrt v\le aW .
\]
Thus the operator norm of \({\cal B}\) is at most \(aW+b\).
Subtract \(2K^0\) in \eqref{r56:native-eq} and use
Corollary \ref{r56:quadratic}.
\end{proof}

This is not a declaration that \(a,b\) are small uniformly.
Their following explicit bounds still have volume costs.

\subsection{Paid finite-volume bounds for every signed source}

Recall \(A_p=(4p!2^p)^{1/p}\) from V51, in particular
\(A_2^2=32\). To avoid confusing the scalar \(A_p\) with
the source matrix \(A_s\), its index below is always a
positive integer. Define
\[
 \begin{aligned}
 \overline a&=
       \frac{8A_4c^{3/2}n}{\gamma}(2A_4+3),\\
 \overline b&=
       \frac{512\sqrt3\,c^{3/2}r\sqrt n}{\gamma},\qquad
 \overline W=\sqrt{2r+\overline b/2}+\overline a/2 .
 \end{aligned}
 \tag{V56.15}\label{r56:constants}
\]

\begin{theorem}[Actual compact covariance in an explicit growing-volume regime]
\label{r56:finite}
One has \(a\le\overline a\), \(b\le\overline b\), and hence
\[
 \|{\cal C}-K^0\|_{2\to2}
 \le\frac12\{d_\gamma+\overline b
                            +\overline a\,\overline W\}.
 \tag{V56.16}\label{r56:finite-eq}
\]
In particular this difference tends to zero for every sequence
with \(\gamma\to\infty\) and \(n/\gamma\to0\).
The comparison is for all real unit-\(\ell^2\) source profiles,
not only a fixed set of plaquette pairs.
\end{theorem}
\begin{proof}
Write \(J=CB^{\mathsf T}\), so \(A_s=J\operatorname{diag}(s)B\).
The site resistance is
\(\mathcal R_{xy}=(e_x-e_y)^{\mathsf T}C(e_x-e_y)\le4c\),
distinct from the edge matrix \(R\).
For \(e=(u,v)\), the spherical constraint gives exactly
\[
 z_x\cdot(z_u-z_v)
      =\tfrac12\{|z_x-z_v|^2-|z_x-z_u|^2\}.
\]
V51's two-point moment bound and
\(\mathcal R_{xu},\mathcal R_{xv}\le4c\)
therefore imply
\[
 \|q_{s,x}\|_p
       \le\frac{8A_pc}{\gamma}
                         \sum_e|J_{xe}|\,|s_e|.
 \tag{V56.17}\label{r56:q-bound}
\]
This difference identity is essential: bounding
\(z_x\cdot(z_u-z_v)\) just by \(|z_u-z_v|\) loses one
half-power of the coupling.

Now
\[
 \|J\|_{\rm F}^2=\Tr(BC^2B^{\mathsf T})=\Tr C=nc,
 \qquad \sum_xJ_{xe}^2=c/3.
\]
The second equality uses edge transitivity for the diagonal
of \(BC^2B^{\mathsf T}\). By Cauchy--Schwarz in \(x\), then
in \(e\),
\[
 \sum_x\|q_{s,x}\|_p
   \le\frac{8A_pc}{\gamma}\sqrt{nc/3}\,\|s\|_1
   \le\frac{8A_pc^{3/2}n}{\gamma}\|s\|_2 .
 \tag{V56.18}\label{r56:q-sum}
\]
The full local energy is \(\gamma h_x=\sum_{e\ni x}X_e\);
there are six incident edges. Thus
\(\|\gamma h_x-3\|_4\le6A_4r+3\le2A_4+3\).
Hölder at each site and \eqref{r56:q-sum} with \(p=4\)
give \(\|\rho_s\|_2\le\overline a\|s\|_2\).
Centering reduces the \(L^2\) norm, proving the bound for \(a\).

For the bilinear return, let \(M_{xe}=\mathbf1_{\{e\ni x\}}\)
be the unsigned incidence. It has row sum six and column
sum two, hence \(\|M\|_{2\to2}\le\sqrt{12}\).
Also
\[
 \|\gamma h_{t,x}\|_2
                  \le A_2r\sum_{e\ni x}|t_e|.
\]
Use this, \eqref{r56:q-bound} at \(p=2\), and
Cauchy--Schwarz in probability. The matrix \(|J|\), with
entrywise absolute values, has norm at most its Frobenius
norm \(\sqrt{nc}\). Therefore
\[
 \begin{split}
 |{\cal N}(s,t)|
 &\le\frac{8A_2^2cr}{\gamma}
                \langle |J|\,|s|,M|t|\rangle\\
 &\le\frac{8\cdot32\,cr}{\gamma}
             \sqrt{nc}\sqrt{12}\,\|s\|_2\|t\|_2
   =\overline b\,\|s\|_2\|t\|_2 .
 \end{split}
\]
Substitute the proved bounds into
Theorem \ref{r56:bootstrap}; its right side is monotone in
\(a,b\ge0\). As \(c\le13/16\), \(r<1/3\), and
\(\eta_\gamma=O(\gamma^{-1})\) uniformly, the claimed
regime makes all displayed errors tend to zero.
\end{proof}

For example, the operator error is bounded by a universal
constant times
\[
 \gamma^{-1/2}+\gamma^{-1}
 +\frac{\sqrt n}{\gamma}
 +\frac n\gamma\sqrt{1+\frac{\sqrt n}{\gamma}}
 +\frac{n^2}{\gamma^2}
\]
for sufficiently large \(\gamma\).
This estimate is obtained from the actual score and moment
bounds, not from summing V53's uniform pairwise error over
all pairs. It does not assert an absolute row-sum bound
for \({\cal C}-K^0\), which is stronger than the stated
operator comparison.

\subsection{Where the remaining volume loss occurs}

The expensive step is visible in \eqref{r56:q-sum} and
the ensuing Minkowski bound on \(\rho_s\). It takes absolute
values before centering and before exploiting cancellation
between different sites and source directions. The bilinear
mean return already has the smaller \(\sqrt n/\gamma\)
cost. A useful next target is a volume-independent bound
on the \emph{centered} polynomial source bank \(\rho_s\),
or a direct bound on the combined return \({\cal B}\)
which preserves cancellations between its two terms.
That is not supplied by entropy density alone.

The compact quadratic response is uniformly controlled, but
the covariance also contains
\(\Cov(\rho_s,F_t)-{\cal N}(s,t)\).
The exact Ward identity retains this term for every signed
source. No correlation positivity, independence, or
differentiation of an unmarked error bound is assumed.

The coherent infrared and entropy inputs are not established
for winding or noncommuting exteriors. Combining V55's Gaussian
coefficient with this different compact law does not establish
a winding remainder. Original exterior probabilities and predictor
covariances, physical source coverage, physical-time contraction,
and the interacting four-dimensional continuum remain unresolved.
The full Yang--Mills mass gap is not proved.

The diagnostic
\path{verify_ym_restart_v56_source_ward.py}
checks rational Ward identities, normalized Haar series,
Green, entropy, Schur and source-bound fixtures. These finite
checks do not prove analytic or continuum claims. V55's
diagnostics, source preservation and visual typesetting are
checked separately.
% END CONSOLIDATED V56 APPEND


% BEGIN CONSOLIDATED V57 APPEND
\clearpage
\section{V57: Heat-balanced Ward fields and a smaller volume loss}
\label{r57:start}

V56 paid the actual coherent compact quadratic response uniformly in
volume, but its nonlinear score estimate required \(\gamma\gg n=N^3\).
We now go upstream of that estimate: change the globally smooth Ward
field, not the compact measure. A symmetric inverse of the two-sided
Laplacian distributes the source over two heat gradients. Its absolute
kernel costs only \(1+\log N\), rather than the previous one-sided
Green-kernel bound. The remaining site sum still costs \(\sqrt n\);
this note does not claim to remove that final volume dependence.

All notation and probability laws are those of V56: the even cubic
torus has \(N\ge4\), \(n=N^3\), \(m=3n\), \(L=B^{\mathsf T}B\),
\(C=L^\dagger\), \(R=BCB^{\mathsf T}\), \(c=C_{xx}\le13/16\),
and \(r=R_{ee}=(n-1)/(3n)\). The spins \(z_x\in S^3\) have the
unconditioned normalized law proportional to
\(\exp(-\gamma S_0)\prod_xdH(z_x)\), and
\(X_e=\gamma(1-z_x\cdot z_y)\).
The covariance under study is still
\({\cal C}_{ef}=\Cov(X_e,X_f)\), with \(K^0_{ef}=3R_{ef}^2/2\).
No boundary law, chart restriction, or source family is replaced.

\subsection{A two-sided source inverse with no constant-mode singularity}

Let \(\Pi=I-\mathbf1\mathbf1^{\mathsf T}/n\),
\(P_\tau=e^{-\tau L}\), and \(L_s=B^{\mathsf T}\operatorname{diag}(s)B\).
For every real signed source \(s\), define
\[
 \widehat A_s=2\int_0^\infty P_\tau L_sP_\tau\,d\tau .
 \tag{V57.1}\label{r57:A}
\]
The parameter \(\tau\) is only a spatial heat-kernel proof parameter,
not physical Euclidean time.

\begin{lemma}[Exact balanced inverse]
\label{r57:inverse}
The integral converges in matrix norm, and
\[
 \widehat A_s^{\mathsf T}=\widehat A_s,\qquad
 \widehat A_s\mathbf1=0,\qquad
 L\widehat A_s+\widehat A_sL=2L_s,\qquad
 \Tr\widehat A_s=\Tr(CL_s).
 \tag{V57.2}\label{r57:inverse-eq}
\]
It is the unique solution of the displayed matrix equation with
both row and column sums zero. At \(s=\mathbf1\), it equals \(\Pi\).
\end{lemma}
\begin{proof}
Since \(L_s=\Pi L_s\Pi\), the integrand decays exponentially at
every fixed \(N\). Differentiating it and integrating its endpoints
gives the matrix equation. Symmetry and the zero sums hold
under the integral. On \(\operatorname{ran}\Pi\), diagonalize \(L\):
the equation reads
\((\lambda_i+\lambda_j)(\widehat A_s)_{ij}=2(L_s)_{ij}\),
with strictly positive denominators. This proves uniqueness.
Multiply the matrix equation by \(C\), take traces, and use
\(CL=LC=\Pi\) to obtain the trace identity. For the constant
source, \(2\int_0^\infty P_\tau LP_\tau\,d\tau=\Pi\).
\end{proof}

The matrix integral itself is standard Lyapunov-equation algebra;
see E.~Jarlebring,
\href{https://people.kth.se/~eliasj/NLA/matrixeqs.pdf}
{\emph{Numerical methods for Lyapunov equations}}, equation (4.16).
The new use here is its exact compact Ward field and its source norm.
This is distinct from V11's configuration-dependent heat window for
a transfer message and from the archived f16 V42 spectral-band
Sylvester equation. No theorem from either different application
is being counted again.

\subsection{An elementary logarithmic heat-gradient budget}

Let \(p_\tau\) be the heat kernel on the one-dimensional cycle of
length \(N\), with generator
\(f(j+1)+f(j-1)-2f(j)\). Put
\[
 g_N(\tau)=\sum_{j\in\mathbb Z/N\mathbb Z}
                      |p_\tau(j+1)-p_\tau(j)|,\qquad
 G_N=5+2\log N .
 \tag{V57.3}\label{r57:G}
\]

\begin{lemma}[The two-gradient time integral]
\label{r57:heat}
For \(\tau>0\),
\[
 g_N(\tau)\le\min\{2,\tau^{-1/2}\},\qquad
 \int_0^\infty g_N(\tau)^2\,d\tau\le G_N .
 \tag{V57.4}\label{r57:heat-eq}
\]
\end{lemma}
\begin{proof}
The heat kernel is the law modulo \(N\) of \(U-V\), where
\(U,V\) are independent Poisson variables of mean \(\tau\).
If \(a(k)=\Pr(U=k)\), then
\(a(k-1)=k\,a(k)/\tau\), with \(a(-1)=0\). Consequently the
\(\ell^1\) distance between \(a\) and its unit translate is
\[
 \sum_{k\ge0}a(k)|k/\tau-1|
 \le\{\E(U/\tau-1)^2\}^{1/2}=\tau^{-1/2}.
\]
Convolution and reduction modulo \(N\) contract this distance.
The bound by two holds for any two probability vectors.

For the long-time cutoff, let \(\pi_N\) be uniform on the cycle.
Writing \(j=\min(k,N-k)\), its heat eigenvalues satisfy
\[
 \lambda_k=4\sin^2(\pi k/N)\ge16j^2/N^2 .
\]
Parseval and Cauchy--Schwarz give
\[
 \|p_u-\pi_N\|_1^2
 \le\sum_{k\ne0}e^{-2u\lambda_k}
 \le2\sum_{j\ge1}e^{-32j^2u/N^2}
 \le4e^{-32u/N^2}\quad(u\ge N^2/2).
\]
For the last inequality use \(j^2\ge1+3(j-1)\) and sum
the geometric series; its denominator is at least \(1-e^{-48}>1/2\).
The semigroup identity and zero sum of a discrete derivative imply
\[
 g_N(\tau)\le g_N(\tau/2)\|p_{\tau/2}-\pi_N\|_1
 \le2\sqrt{2/\tau}\,e^{-8\tau/N^2}
 \quad(\tau\ge N^2).
\]
Thus the integrals on \((0,1)\), \((1,N^2)\), and
\((N^2,\infty)\) are at most \(4\), \(2\log N\), and
\[
 8\int_{N^2}^\infty \tau^{-1}e^{-16\tau/N^2}\,d\tau
 \le\tfrac12e^{-16}<1,
\]
respectively. This proves the claimed explicit budget.
\end{proof}

For a matrix \(A\), write
\(a_x(A)=\sum_y|A_{xy}|\). These are ordinary site row sums;
no probability measure enters their definition.

\begin{proposition}[Absolute kernel norm for all signed sources]
\label{r57:kernel}
\[
 \left(\sum_xa_x(\widehat A_s)^2\right)^{1/2}
                  \le2\sqrt3\,G_N\|s\|_2 .
 \tag{V57.5}\label{r57:kernel-eq}
\]
\end{proposition}
\begin{proof}
For an oriented edge \(e\), let \(b_e=B^{\mathsf T}e_e\) and
\(k_\tau(x,e)=(P_\tau b_e)_x\).
The cubic heat kernel factors into three one-dimensional kernels.
Its nonnegative, mass-one factors in the other two coordinates give
\[
 \sum_x|k_\tau(x,e)|=g_N(\tau),\qquad
 \sum_e|k_\tau(x,e)|=3g_N(\tau).
\]
The matrix \(T_\tau(x,e)=|k_\tau(x,e)|\) therefore has
\(\ell^2\) operator norm at most \(\sqrt3\,g_N(\tau)\).
Expanding \(L_s=\sum_es_eb_eb_e^{\mathsf T}\) in
\eqref{r57:A}, and taking absolute values only after this
two-gradient factorization, gives
\[
 a_x(\widehat A_s)
 \le2\int_0^\infty g_N(\tau)
                         (T_\tau|s|)_x\,d\tau .
\]
Minkowski in site \(\ell^2\) and Lemma \ref{r57:heat} prove
\eqref{r57:kernel-eq}. No positivity of \(s\) was assumed.
\end{proof}

\subsection{The native Ward identity survives the change of field}

Use the V56 notation \(F_s=\sum_es_eX_e\),
\(h_{t,x}=z_x\cdot(L_tz)_x\), \(h_x=h_{\mathbf1,x}\).
Define the new, explicitly distinguished source score by
\[
 \begin{aligned}
 \widehat q_{s,x}&=z_x\cdot(\widehat A_sz)_x,&
 (\widehat{\cal V}_s)_x&=(I-z_xz_x^{\mathsf T})(\widehat A_sz)_x,\\
 \widehat\rho_s&=\sum_x\widehat q_{s,x}(\gamma h_x-3),&
 \widehat{\cal N}(s,t)&=\gamma\E\sum_x\widehat q_{s,x}h_{t,x}.
 \end{aligned}
 \tag{V57.6}\label{r57:fields}
\]
In general these hatted quantities differ from the V56 quantities.

\begin{proposition}[Balanced compact source identity]
\label{r57:Ward}
For every smooth real \(F\),
\[
 \begin{aligned}
 \operatorname{div}\widehat{\cal V}_s
              &=3\Tr\widehat A_s-3\sum_x\widehat q_{s,x},\\
 \widehat{\cal V}_s(\gamma S_0)
              &=2F_s-\gamma\sum_x\widehat q_{s,x}h_x,\\
 2\E F_s&=3\Tr(CL_s)+\E\widehat\rho_s,\\
 2\Cov(F_s,F)&=\E\widehat{\cal V}_sF+\Cov(\widehat\rho_s,F).
 \end{aligned}
 \tag{V57.7}\label{r57:Ward-eq}
\]
\end{proposition}
\begin{proof}
The spherical divergence calculation from V56 applies to any
constant site matrix: its contribution at \(x\) is
\(3(\widehat A_s)_{xx}-3\widehat q_{s,x}\).
For the energy drift, the only changed step is
\[
 \langle\widehat A_sz,Lz\rangle
 =\tfrac12 z^{\mathsf T}(\widehat A_sL+L\widehat A_s)z
 =z^{\mathsf T}L_sz .
\]
Integrate the divergence of
\(F\widehat{\cal V}_s e^{-\gamma S_0}\) on the compact product,
first with \(F=1\), then subtract the mean identity.
The trace is supplied by \eqref{r57:inverse-eq}.
\end{proof}

In particular, uniform \(s\) still gives exactly V51's field and
score. The gain below concerns the entire source bank and its bounds;
it is not a new scalar collective Ward identity.

\subsection{The changed quadratic response is also controlled uniformly}

Changing the field changes its quadratic response, so V56's
\({\cal Q}\) cannot simply be reused unchanged. Put
\[
 \begin{aligned}
 {\sf H}_\tau&=BP_\tau B^{\mathsf T},&
 R_\tau&=BCP_\tau B^{\mathsf T},\\
 E_\tau&=e^{-\tau BB^{\mathsf T}},&
 \Gamma_\tau&=E_\tau\Gamma,\qquad
 \widehat{\cal Q}=2\int_0^\infty
                          {\sf H}_\tau\circ\Gamma_\tau\,d\tau ,
 \end{aligned}
 \tag{V57.8}\label{r57:Q}
\]
where \(\Gamma_{ef}=\gamma\E(Bz)_e\cdot(Bz)_f\) is the actual
compact second moment from V56 and \(\circ\) is entrywise product.
The new symbol \({\sf H}_\tau\) is a deterministic heat-gradient
kernel, not V56's longitudinal Gram matrix \(H\).

\begin{lemma}[Commutation is paid by the coherent law]
\label{r57:commute}
The matrix \(\Gamma\) commutes with \(BB^{\mathsf T}\). Consequently
\[
 0\preceq\Gamma_\tau\preceq4R_\tau,\qquad
 \gamma\E\langle\widehat A_sz,L_tz\rangle
                         =s^{\mathsf T}\widehat{\cal Q}t .
 \tag{V57.9}\label{r57:commute-eq}
\]
\end{lemma}
\begin{proof}
The site matrix \(K_{xy}=\E z_x\cdot z_y\) is positive semidefinite
and translation invariant under the coherent periodic law. It is a
convolution matrix, hence commutes with \(L\). Since
\(\Gamma=\gamma BKB^{\mathsf T}\), this proves the asserted
edge commutation. The inherited \(0\preceq\Gamma\preceq4R\)
can now be conjugated by \(E_{\tau/2}\); it gives the first assertion
because \(E_{\tau/2}\Gamma E_{\tau/2}=E_\tau\Gamma\).

For the second, expand the two source Laplacians in
\(\gamma\Tr(\widehat A_sL_tK)\). The integrand for \(e,f\) is
\[
 2\gamma s_et_f
       (b_e^{\mathsf T}P_\tau b_f)
       (b_f^{\mathsf T}KP_\tau b_e)
 =2s_et_f({\sf H}_\tau)_{ef}(\Gamma_\tau)_{ef}.
\]
The last equality uses \(KP_\tau=P_\tau K\).
All integrals converge at fixed \(N\). This commutation argument
has not been established for an arbitrary fixed exterior.
\end{proof}

Let \(h_\gamma=512cr/\gamma\) and \(\eta_\gamma\) be the
already proved V53--V56 quantities. Define
\[
 \chi_N=\sqrt{1+r/2+c/12},\qquad
 \widehat d_\gamma=
 2\chi_N\left(\frac8{\sqrt3}\sqrt{\eta_\gamma}
                               +2\sqrt{h_\gamma}\right).
 \tag{V57.10}\label{r57:d}
\]

\begin{theorem}[Uniform control of the balanced quadratic response]
\label{r57:response}
\[
 0\preceq\widehat{\cal Q}\preceq4(R\circ R),\qquad
 \|\widehat{\cal Q}\|_{2\to2}\le4r,\qquad
 \|\widehat{\cal Q}-2K^0\|_{2\to2}\le\widehat d_\gamma .
 \tag{V57.11}\label{r57:response-eq}
\]
In particular \(\widehat d_\gamma=O(\gamma^{-1/2})\)
uniformly in \(N\).
\end{theorem}
\begin{proof}
Both matrices in the integrand of \eqref{r57:Q} are positive
semidefinite. Schur multiplication preserves order. Entrywise,
\(-\partial_\tau R_\tau={\sf H}_\tau\), \(R_0=R\),
and \(R_\infty=0\), so
\[
 2\int_0^\infty{\sf H}_\tau\circ R_\tau\,d\tau=R\circ R .
\]
This proves the order and norm bounds and identifies the Gaussian
value \(3R\circ R=2K^0\) without replacing the compact law.

For the error, set \(\Delta=\Gamma-3R\).
V56 gives \(\Gamma=3S+H\), \(0\preceq H\preceq4R\),
\(H_{ee}\le h_\gamma\), and
\(\sum_f|(S-R)_{ef}|^2\le64\eta_\gamma/27\).
As \(H^2\preceq4H\), every row satisfies
\[
 \|\Delta_{e,\cdot}\|_2
       \le\frac8{\sqrt3}\sqrt{\eta_\gamma}+2\sqrt{h_\gamma}
       =:\delta_\gamma .
\]
The error commutes with \(E_\tau\) by Lemma \ref{r57:commute}.
Since \(\|E_\tau\|_{2\to2}\le1\), the same row bound holds
for \(\Delta E_\tau\).

The needed integrability of the other row costs no \(\log N\).
Indeed
\(\|({\sf H}_\tau)_{e,\cdot}\|_2^2
 =(BP_{2\tau}LB^{\mathsf T})_{ee}\).
Spectral integration on \(\operatorname{ran}\Pi\) yields exactly
\[
 \begin{aligned}
 \int_0^\infty\|({\sf H}_\tau)_{e,\cdot}\|_2^2\,d\tau&=1,\\
 \int_0^\infty\tau\|({\sf H}_\tau)_{e,\cdot}\|_2^2\,d\tau&=r/4,\\
 \int_0^\infty\tau^2\|({\sf H}_\tau)_{e,\cdot}\|_2^2\,d\tau&=c/12.
 \end{aligned}
 \tag{V57.12}\label{r57:heat-moments}
\]
Here \((BB^{\mathsf T})_{ee}=2\) and
\((BC^2B^{\mathsf T})_{ee}=c/3\).
Cauchy--Schwarz in \(\tau\), with weights \((1+\tau)^{\pm2}\),
therefore gives
\[
 \int_0^\infty\|({\sf H}_\tau)_{e,\cdot}\|_2\,d\tau
 \le\sqrt{1+r/2+c/12}=\chi_N .
\]
Finally,
\[
 \sup_e\sum_f
       |(\widehat{\cal Q}-3R\circ R)_{ef}|
 \le2\sup_e\int_0^\infty
   \|({\sf H}_\tau)_{e,\cdot}\|_2
   \|(\Delta E_\tau)_{e,\cdot}\|_2\,d\tau
 \le2\chi_N\delta_\gamma .
\]
Symmetry converts this row bound to the stated operator bound.
The inherited \(\eta_\gamma=O(\gamma^{-1})\) finishes the proof.
\end{proof}

\subsection{A proved improvement for the actual whole source bank}

Define the centered source norm and the bilinear mean-return norm
for the new field, not for the old one, by
\[
 \widehat a=\sup_{\|s\|_2=1}
       \|\widehat\rho_s-\E\widehat\rho_s\|_2,\qquad
 \widehat b=\sup_{\|s\|_2=\|t\|_2=1}
                         |\widehat{\cal N}(s,t)|.
 \tag{V57.13}\label{r57:ab}
\]
With \(A_p=(4p!2^p)^{1/p}\) as in V51, put
\[
 \begin{aligned}
 a_*&=\frac{8\sqrt3\,A_4c\,G_N\sqrt n}{\gamma}(2A_4+3),\\
 b_*&=\frac{1536cr\,G_N}{\gamma},\qquad
 W_*=\sqrt{2r+b_*/2}+a_*/2 .
 \end{aligned}
 \tag{V57.14}\label{r57:budget}
\]

\begin{theorem}[Heat-balanced compact covariance bound]
\label{r57:main}
For every even \(N\ge4\) and \(\gamma>0\),
\[
 \begin{aligned}
 \widehat a&\le a_*,\qquad \widehat b\le b_*,\\
 \|{\cal C}\|_{2\to2}&\le W_*^2,\\
 \|{\cal C}-K^0\|_{2\to2}
              &\le\tfrac12\{\widehat d_\gamma+b_*+a_*W_*\}.
 \end{aligned}
 \tag{V57.15}\label{r57:main-eq}
\]
In particular the last operator difference tends to zero whenever
\[
 \gamma\longrightarrow\infty,\qquad
 \frac{N^{3/2}(1+\log N)}{\gamma}\longrightarrow0 .
 \tag{V57.16}\label{r57:regime}
\]
The estimate controls all real unit-\(\ell^2\) edge sources.
\end{theorem}
\begin{proof}
The zero row sums and \(|z_x|=1\) give the exact identity
\[
 \widehat q_{s,x}
       =-\tfrac12\sum_y(\widehat A_s)_{xy}|z_x-z_y|^2.
\]
The inherited resistance moment
\(\|(\gamma/2)|z_x-z_y|^2\|_p\le A_p{\cal R}_{xy}\),
with \({\cal R}_{xy}\le4c\), and \eqref{r57:kernel-eq} imply
\[
 \left(\sum_x\|\widehat q_{s,x}\|_p^2\right)^{1/2}
       \le\frac{8\sqrt3\,A_pcG_N}{\gamma}\|s\|_2 .
 \tag{V57.17}\label{r57:q}
\]
The mixed norm means the \(\ell^2\) norm of the individual
probability \(L^p\) norms.

Since \(\gamma h_x\) is the sum of the six incident \(X_e\),
\(\|\gamma h_x-3\|_4\le6A_4r+3\le2A_4+3\).
Hölder at each site followed by the site Cauchy--Schwarz inequality
gives
\[
 \|\widehat\rho_s\|_2
 \le(2A_4+3)\sqrt n
             \left(\sum_x\|\widehat q_{s,x}\|_4^2\right)^{1/2}
 \le a_*\|s\|_2 .
\]
Centering decreases the \(L^2\) norm.

For the other return, let \(M_{xe}=\mathbf1_{\{e\ni x\}}\).
Its norm is at most \(\sqrt{12}\), and
\(\|\gamma h_{t,x}\|_2\le A_2r(M|t|)_x\).
Use \eqref{r57:q} at \(p=2\), Cauchy--Schwarz in probability
and in the site index, and \(A_2^2=32\). This gives
\[
 |\widehat{\cal N}(s,t)|
 \le\frac{8\sqrt3\,A_2^2crG_N}{\gamma}\sqrt{12}
                 \|s\|_2\|t\|_2
 =b_*\|s\|_2\|t\|_2 .
\]

For completeness, the balanced Ward identity gives the exact
native decomposition
\[
 2s^{\mathsf T}{\cal C}t
 =s^{\mathsf T}\widehat{\cal Q}t+
   \Cov(\widehat\rho_s,F_t)-\widehat{\cal N}(s,t).
 \tag{V57.18}\label{r57:decomposition}
\]
The covariance and mean return here are kept together.
For \(v=\|{\cal C}\|_{2\to2}\), testing its maximal-eigenvalue
unit source yields \(2v\le4r+b_*+a_*\sqrt v\).
The same elementary quadratic argument as V56 gives \(v\le W_*^2\).
For arbitrary unit \(s,t\), the absolute connected return is at most
\(a_*\sqrt v\). Subtract \(2K^0\) in \eqref{r57:decomposition},
use Theorem \ref{r57:response}, and divide by two.
All constants except \(G_N,\sqrt n,\gamma\) are uniformly bounded.
Condition \eqref{r57:regime} therefore makes the three error
terms vanish.
\end{proof}

Thus V56's sufficient \(\gamma\gg N^3\) regime has been replaced
by \eqref{r57:regime}. The old result remains correct and retained;
the new estimate is not obtained by changing the normalization of
\(s\), discarding its signs, or assuming covariance positivity.
The conclusion is an operator bound, not an absolute row-sum
estimate for the full native covariance.

\subsection{Remaining obstruction and direction of the next attack}

The heat factorization removes a substantial kernel-volume loss,
but the displayed \(\sqrt n\) came from a still-uncentered site sum.
A volume-independent centered bound on \(\widehat\rho_s\), or a
bound directly on the combined return in \eqref{r57:decomposition},
has not been proved. In particular the new regime is still much
stronger than a logarithmically growing bare coupling. It is not
a physical continuum trajectory or a mass-gap theorem.

The ESM transfer suggestions remain useful as proof architecture:
change the inverse only together with its marked response and
all nonlinear returns. Here that obligation is met by
\eqref{r57:Q}--\eqref{r57:decomposition}, not by importing an ESM
existence theorem. V55 already reviewed the newer ESM V31 beyond
the supplied V23; no changed companion input was needed in this
appendix. The next five-note companion review and ten-note broad
literature sweep are both due at V60.

The remaining upstream question is connected cancellation after
the two heat gradients have been retained. Arbitrary exteriors
also remove the translation commutation used in
Lemma \ref{r57:commute}; this must be paid rather than presumed.
Winding remainders, typical-exterior probabilities and predictor
covariances, physical source coverage, physical-time contraction,
and the interacting four-dimensional construction remain open.
The full Yang--Mills mass gap is unproved.

The accompanying finite diagnostic checks the Lyapunov inverse,
its native spherical drift, the changed quadratic response, and
heat-kernel constants. Such checks do not prove the analytic
theorems above or the continuum goal. Predecessor preservation and
visual typesetting are checked separately.
% END CONSOLIDATED V57 APPEND


% BEGIN CONSOLIDATED V58 APPEND
\clearpage
\section{V58: A volume-uniform leading centered-score coefficient}
\label{r58:start}

V57's balanced Ward field improved the actual compact covariance
estimate, but its score bound still costs \(\sqrt n(1+\log N)/\gamma\).
This note investigates the connected cancellation before attempting
another compact remainder estimate. We prove that the leading
fluctuation coefficient of the \emph{whole signed-source score bank}
has neither of those volume losses. We also identify it as the
fixed-volume low-temperature limit of the actual compact score.
The order of limits is explicit: a uniform coefficient is not a
uniform remainder theorem.

Throughout, \(N\ge4\) is even, \(n=N^3\), and all graph, compact-law,
and source notation is as in V57. In particular \(c=C_{xx}\le13/16\),
\(r=R_{ee}=(n-1)/(3n)\), and
\[
 A=A_s=\widehat A_s,\qquad LA+AL=2L_s,\qquad
 A=A^{\mathsf T}=\Pi A\Pi .
 \tag{V58.1}\label{r58:A}
\]
Only within this appendix is the shorter \(A_s\) used for V57's
balanced matrix; it is not V56's generally nonsymmetric \(CL_s\).
We retain the full, unconditioned compact quantities
\(\widehat\rho_s,\widehat{\cal N},\widehat{\cal B}\), with
\[
 \widehat{\cal B}(s,t)
       =\Cov_\mu(\widehat\rho_s,F_t)-\widehat{\cal N}(s,t).
\]
No source signs or source supports are restricted.

\subsection{Two deterministic estimates before absolute kernel sums}

Let \(\circ\) denote entrywise product and set \(M=C\circ C\).
All operator orders in the following statement are finite-dimensional.

\begin{lemma}[Projected bubble and weighted source bounds]
\label{r58:matrix}
\[
 \begin{aligned}
 0\preceq\Pi M\Pi&\preceq4cC,\\
 \|A_s\|_{\rm F}^2&\le r\|s\|_2^2,\\
 \Tr(A_sCA_s)&\le\frac c3\|s\|_2^2 .
 \end{aligned}
 \tag{V58.2}\label{r58:matrix-eq}
\]
The projections in the first line cannot be omitted.
\end{lemma}
\begin{proof}
The torus symbol is
\(\lambda(p)=\sum_{j=1}^3|1-e^{ip_j}|^2\).
The Fourier eigenvalue of \(M\) at \(k\) is
\[
 \widehat M(k)=\frac1n
       \sum_{\substack{p\ne0\\k-p\ne0}}
                    \frac1{\lambda(p)\lambda(k-p)} .
\]
For every \(p,q\), the triangle inequality squared gives
\(\lambda(p+q)\le2\lambda(p)+2\lambda(q)\).
Hence \(\lambda(k)\widehat M(k)\le4c\) for \(k\ne0\).
The Schur product theorem gives \(M\succeq0\); simultaneous Fourier
diagonalization proves the first assertion on the mean-zero space.
At \(k=0\), \(\widehat M(0)=n^{-1}\Tr C^2>0\), whereas
\(\widehat C(0)=0\), explaining the projections.

For the other bounds, diagonalize \(L\) on \(\operatorname{ran}\Pi\).
Writing \(\ell_{ij}=(L_s)_{ij}\) in this basis,
\(A_{ij}=2\ell_{ij}/(\lambda_i+\lambda_j)\).
Since \(4\lambda_i\lambda_j\le(\lambda_i+\lambda_j)^2\),
\[
 \|A\|_{\rm F}^2
 \le\Tr(CL_sCL_s)
 =s^{\mathsf T}(R\circ R)s\le r\|s\|_2^2 .
\]
Symmetrizing the two indices in the weighted trace gives
\[
 \begin{aligned}
 \Tr(ACA)
 &=\sum_{i,j}\frac{2|\ell_{ij}|^2}
                         {\lambda_i\lambda_j(\lambda_i+\lambda_j)}\\
 &\le\sum_{i,j}\frac{(\lambda_i+\lambda_j)|\ell_{ij}|^2}
                              {2\lambda_i^2\lambda_j^2}
 =\Tr(CL_sCL_sC).
 \end{aligned}
\]
Put \(T=BC^2B^{\mathsf T}\). The last trace equals
\(s^{\mathsf T}(R\circ T)s\).
Because \(0\preceq R\preceq I\) and \(T\succeq0\),
\(0\preceq R\circ T\preceq I\circ T\).
Edge transitivity and \(\Tr T=\Tr C=nc\) give
\(T_{ee}=c/3\), proving the last bound.
\end{proof}

Unlike the absolute-kernel estimate in V57, these estimates retain
both source signs and the mean-zero cancellation. They do not
require an infrared estimate for a source-tilted compact law.

\subsection{The precise Gaussian coefficient and its local matrices}

Let \(u^1,u^2,u^3\) be independent centered Gaussian site fields
with covariance \(C\), supported on \(\operatorname{ran}\Pi\).
Write \(\E_G\) for this Gaussian expectation. Define
\[
 \begin{aligned}
 Q_{s,x}(u)&=u_x\cdot(A_su)_x-\tfrac12(A_s|u|^2)_x,\\
 H_{t,x}(u)&=\tfrac12\sum_{e\ni x}t_e|(Bu)_e|^2,\qquad H_x=H_{\mathbf1,x},\\
 f_t(u)&=\tfrac12\sum_et_e|(Bu)_e|^2,\qquad
 {\cal R}_s(u)=\sum_xQ_{s,x}(u)(H_x(u)-3).
 \end{aligned}
 \tag{V58.3}\label{r58:polynomials}
\]
The polynomial \({\cal R}_s\) is not assumed centered.

Work on the one-color Hilbert space \({\cal H}=\operatorname{ran}\Pi\).
Let
\[
 \begin{aligned}
 g_x&=C^{1/2}e_x,&w_x&=C^{1/2}A_se_x,&d_e&=C^{1/2}b_e,\\
 D_x&=\tfrac12(g_xw_x^{\mathsf T}+w_xg_x^{\mathsf T})
              -\tfrac12\sum_y(A_s)_{xy}g_yg_y^{\mathsf T},\\
 J_x&=\tfrac12\sum_{e\ni x}d_ed_e^{\mathsf T},&
 J_{t,x}&=\tfrac12\sum_{e\ni x}t_ed_ed_e^{\mathsf T}.
 \end{aligned}
 \tag{V58.4}\label{r58:local}
\]
If \(\xi^a\) are independent standard Gaussians on \({\cal H}\),
then \(Q_{s,x}=\sum_a(\xi^a)^{\mathsf T}D_x\xi^a\) and
\(H_{t,x}=\sum_a(\xi^a)^{\mathsf T}J_{t,x}\xi^a\).
Put \(v_x=\Tr D_x\) and
\(\|D\|_{\ell^2{\rm F}}=(\sum_x\|D_x\|_{\rm F}^2)^{1/2}\).

\begin{proposition}[Uniform local tensor bank]
\label{r58:local-bound}
\[
 \begin{aligned}
 \sum_xJ_x&=I_{\cal H},&
 \Tr J_x&=3r,&
 \left\|(\Tr(J_xJ_y))_{xy}\right\|_{2\to2}&\le3r,\\
 \|v\|_2&\le\frac c{\sqrt3}\|s\|_2,&
 \|D\|_{\ell^2{\rm F}}&\le\frac{2c}{\sqrt3}\|s\|_2,&
 \sum_xD_x&=C^{1/2}A_sC^{1/2}.
 \end{aligned}
 \tag{V58.5}\label{r58:local-bound-eq}
\]
\end{proposition}
\begin{proof}
The first equality follows from \(\sum_x\frac12
\sum_{e\ni x}b_eb_e^{\mathsf T}=L\).
Also \(|d_e|^2=r\), so \(\Tr J_x=3r\).
The nonnegative symmetric Gram matrix
\(G_{xy}=\Tr(J_xJ_y)\) has row sum
\(\Tr(J_x\sum_yJ_y)=3r\), giving its norm bound.

Since \(|g_y|^2=c\) and \(A_s\mathbf1=0\),
\(v_x=g_x\cdot w_x\). Therefore
\(\sum_xv_x^2\le c\sum_x|w_x|^2=c\Tr(A_sCA_s)\).
The first part of \(D_x\) has Frobenius norm at most
\(\sqrt c\,|w_x|\). For the other part, exactly
\[
 \sum_x\left\|\sum_y(A_s)_{xy}g_yg_y^{\mathsf T}\right\|_{\rm F}^2
       =\Tr(A_sMA_s)\le4c\Tr(A_sCA_s).
\]
The last inequality uses \(A_s=\Pi A_s\Pi\) before applying the
projected bubble bound. The triangle inequality in
\(\ell^2({\rm F})\) now gives
\(\|D\|_{\ell^2{\rm F}}\le2\sqrt{c\Tr(A_sCA_s)}\).
Lemma \ref{r58:matrix} proves both norm estimates.
Summing \eqref{r58:local}, the diagonal part cancels by the
column sums of \(A_s\), proving the last identity.
\end{proof}

\subsection{Centering and Wick contraction remove the volume loss}

For a symmetric matrix \(K\) on \({\cal H}\), set
\[
 q(K)=\sum_{a=1}^3
             \{(\xi^a)^{\mathsf T}K\xi^a-\Tr K\}.
\]
The notation \(:Q_{s,x}H_x:_4\) means the fourth homogeneous
Gaussian chaos of the product, not merely the product with its
mean removed. Define
\[
 Z_s=3\sum_xv_xJ_x-\frac3n\sum_xD_x
                       +4\operatorname{Sym}\sum_xD_xJ_x,
 \qquad \operatorname{Sym}T=(T+T^{\mathsf T})/2.
 \tag{V58.6}\label{r58:Z}
\]

\begin{lemma}[Complete centered coefficient]
\label{r58:chaos}
One has the orthogonal decomposition
\[
 {\cal R}_s-\E_G{\cal R}_s
       =q(Z_s)+{\cal R}^{[4]}_s,\qquad
 {\cal R}^{[4]}_s=\sum_x:Q_{s,x}H_x:_4 .
 \tag{V58.7}\label{r58:chaos-eq}
\]
Furthermore
\[
 \begin{aligned}
 \|Z_s\|_{\rm F}
   &\le c\left(11\sqrt r+\frac{2\sqrt3}{\sqrt n}\right)\|s\|_2,\\
 \E_G({\cal R}^{[4]}_s)^2&\le864rc^2\|s\|_2^2 .
 \end{aligned}
 \tag{V58.8}\label{r58:chaos-bounds}
\]
\end{lemma}
\begin{proof}
For two quadratic forms summed over three independent colors,
their product has second-chaos matrix
\[
 3(\Tr D)J+3(\Tr J)D+4\operatorname{Sym}(DJ).
\]
This follows by making one Wick contraction among the four
Gaussian factors: the within-form contractions produce the two
trace terms, and the four cross contractions produce the last.
Subtracting \(3Q_{s,x}\) and summing gives \eqref{r58:Z}, because
\(3\Tr J_x-3=9r-3=-3/n\). The remaining nonconstant part is
the fourth chaos, orthogonal to the second.

For the first term of \(Z_s\), the Gram bound gives
\[
 \left\|\sum_xv_xJ_x\right\|_{\rm F}\le\sqrt{3r}\|v\|_2 .
\]
The middle term is bounded by \(3\|D\|_{\ell^2{\rm F}}/\sqrt n\).
For the last, factor the product as a block row and a block column:
\[
 \sum_xD_xJ_x
       =[D_xJ_x^{1/2}]_x[J_x^{1/2}]_x^{\mathsf T}.
\]
The block column has norm one since \(\sum_xJ_x=I\). Also
\(0\preceq J_x\preceq3rI\), so
\[
 \left\|\sum_xD_xJ_x\right\|_{\rm F}^2
       \le\sum_x\Tr(J_xD_x^2)
       \le3r\sum_x\|D_x\|_{\rm F}^2 .
\]
Symmetrization contracts the Frobenius norm.
Substitute \eqref{r58:local-bound-eq}: the three contributions
are \(3c\sqrt r\), \(2\sqrt3c/\sqrt n\), and \(8c\sqrt r\),
times \(\|s\|_2\).

For the fourth chaos, put \(\overline D_x=I_3\otimes D_x\) and
\(\overline J_x=I_3\otimes J_x\) on the full color space.
The unsymmetrized fourth tensor is
\(T_s=\sum_x\overline D_x\otimes\overline J_x\).
Wick's rule yields
\(\E_G({\cal R}^{[4]}_s)^2=4!\|\operatorname{Sym}_4T_s\|^2
 \le24\|T_s\|^2\).
Indeed the two normal-ordered fourth polynomials can contract
only across the two copies; the \(4!\) permutations are precisely
this symmetric-tensor inner product. Thus no infinite-dimensional
Gaussian theorem is needed here. Now
\[
 \|T_s\|^2
   =9\sum_{x,y}\Tr(D_xD_y)\Tr(J_xJ_y)
   \le27r\sum_x\|D_x\|_{\rm F}^2
   \le36rc^2\|s\|_2^2 .
\]
Multiply by \(24\). This retains cross-site contraction before
taking norms; replacing it by the sum of individual \(L^2\)
norms would recreate the unwanted volume loss.
\end{proof}

\begin{theorem}[Volume-uniform leading centered-score bank]
\label{r58:score}
For every real signed source \(s\),
\[
 \|{\cal R}_s-\E_G{\cal R}_s\|_{L^2(G)}
                         \le24c\|s\|_2\le\frac{39}{2}\|s\|_2 .
 \tag{V58.9}\label{r58:score-eq}
\]
\end{theorem}
\begin{proof}
The two chaoses are orthogonal, and
\(\E_Gq(Z_s)^2=6\|Z_s\|_{\rm F}^2\).
The squared norm is therefore bounded by
\[
 c^2\left\{6\left(11\sqrt r+
                    \frac{2\sqrt3}{\sqrt n}\right)^2+864r\right\}\|s\|_2^2.
\]
Using \(r\le1/3\), \(n\ge64\), the braces are at most
\(4513/8<576=24^2\).
The last bound follows from \(c\le13/16\).
\end{proof}

\subsection{The mean return and the combined nonlinear coefficient}

Set
\[
 \begin{aligned}
 {\cal N}_1(s,t)&=\E_G\sum_xQ_{s,x}H_{t,x},\\
 {\cal B}_1(s,t)&=\Cov_G({\cal R}_s,f_t)-{\cal N}_1(s,t),&
 M_t&=C^{1/2}L_tC^{1/2}.
 \end{aligned}
 \tag{V58.10}\label{r58:returns}
\]

\begin{proposition}[Uniform leading return bounds]
\label{r58:return}
\[
 \begin{aligned}
 |{\cal N}_1(s,t)|&\le21cr\|s\|_2\|t\|_2,\\
 |{\cal B}_1(s,t)|
 &\le c\left(54r+6\sqrt{\frac{3r}{n}}\right)\|s\|_2\|t\|_2
 \le\frac{75c}{4}\|s\|_2\|t\|_2 .
 \end{aligned}
 \tag{V58.11}\label{r58:return-eq}
\]
\end{proposition}
\begin{proof}
Write \(w_{t,x}=\Tr J_{t,x}=(r/2)\sum_{e\ni x}t_e\).
The unsigned incidence has norm \(\sqrt{12}\), so
\(\|w_t\|_2\le\sqrt3r\|t\|_2\).
Also, by Cauchy--Schwarz over the six incident edges and then
counting each edge at its two ends,
\[
 \left(\sum_x\|J_{t,x}\|_{\rm F}^2\right)^{1/2}
                          \le\sqrt3r\|t\|_2 .
\]
Here \(\|d_ed_e^{\mathsf T}\|_{\rm F}=r\).
The Gaussian quadratic product formula gives exactly
\[
 {\cal N}_1(s,t)=9\sum_xv_xw_{t,x}
                         +6\sum_x\Tr(D_xJ_{t,x}).
 \tag{V58.12}\label{r58:Nexact}
\]
The two terms are bounded by \(9cr\) and \(12cr\), respectively,
times the source norms.

Only the second chaos of \({\cal R}_s\) pairs with \(f_t\);
thus \(\Cov_G({\cal R}_s,f_t)=3\Tr(Z_sM_t)\).
Moreover
\(\|M_t\|_{\rm F}^2=t^{\mathsf T}(R\circ R)t\le r\|t\|_2^2\).
Use \eqref{r58:chaos-bounds}, subtract \({\cal N}_1\), and
apply its proved bound. This yields \(54cr+6c\sqrt{3r/n}\).
Finally \(r\le1/3\) and \(n\ge64\) give \(18c+3c/4=75c/4\).
\end{proof}

\subsection{Identification with the actual compact law at fixed volume}

This subsection fixes \(N\); its convergence rates and domination
constants are not declared uniform in \(N\).

\begin{theorem}[The coefficient is the native fixed-volume limit]
\label{r58:compact-limit}
At every fixed even \(N\ge4\), as \(\gamma\to\infty\), the
joint moments of every finite source family satisfy
\[
 \begin{aligned}
 (\gamma\widehat\rho_s,F_t)&\longrightarrow({\cal R}_s,f_t),\\
 \gamma\widehat{\cal N}(s,t)&\longrightarrow{\cal N}_1(s,t),\\
 \gamma\widehat{\cal B}(s,t)&\longrightarrow{\cal B}_1(s,t).
 \end{aligned}
 \tag{V58.13}\label{r58:compact-limit-eq}
\]
In the first line, convergence means convergence of all polynomial
moments of the displayed random observables, not an assertion that
their two underlying probability spaces are equal.
Consequently, for the actual norms of V57,
\[
 \begin{aligned}
 \lim_{\gamma\to\infty}\gamma\widehat a
   &=\sup_{\|s\|_2=1}\|{\cal R}_s-\E_G{\cal R}_s\|_2\le24c,\\
 \lim_{\gamma\to\infty}\gamma\widehat b
   &=\|{\cal N}_1\|_{2\to2}\le21cr,\\
 \lim_{\gamma\to\infty}\gamma\|\widehat{\cal B}\|_{2\to2}
   &=\|{\cal B}_1\|_{2\to2}\le75c/4 .
 \end{aligned}
 \tag{V58.14}\label{r58:compact-norms}
\]
For bilinear forms, \(\|\cdot\|_{2\to2}\) denotes the supremum
over two unit real source vectors.
\end{theorem}
\begin{proof}
The zero set of \(S_0=\frac12\sum_e|z_x-z_y|^2\) consists exactly
of constant spins, since the graph is connected. It is a compact
copy of \(S^3\). Near that set, use the direction of the nonzero
magnetization as the collective orientation and write, in a local
orthonormal frame,
\[
 z_x=\left(\sqrt{1-|v_x|^2},v_x\right),\qquad
                  \sum_xv_x=0 .
\]
These are smooth coordinates in a fixed sufficiently small
neighborhood of the constant configurations: the longitudinal
magnetization is positive, and the transverse mean is zero.
Together with a local orientation chart their differential at
\(v=0\) separates common rotations from mean-zero variations, so
it is invertible. The product-Haar density in these coordinates
is smooth and positive at \(v=0\). Its leading factor cancels
between numerator and normalizer; all observables here are
invariant under the collective orientation.

Set \(v=u/\sqrt\gamma\). Taylor expansion at fixed \(u\) gives
\[
 \begin{aligned}
 \gamma S_0&=\tfrac12\sum_e|(Bu)_e|^2+O_N(\gamma^{-1}|u|^4),\\
 \gamma\widehat q_{s,x}
   &=Q_{s,x}(u)+O_{N,s}(\gamma^{-1}|u|^4),\\
 \gamma h_{t,x}&=H_{t,x}(u)+O_{N,t}(\gamma^{-1}|u|^4),\\
 F_t&=f_t(u)+O_{N,t}(\gamma^{-1}|u|^4).
 \end{aligned}
 \tag{V58.15}\label{r58:chart-expansion}
\]
For the second line use \(A_s\mathbf1=0\):
the transverse term is \(u_x\cdot(A_su)_x\) and the
longitudinal term is \(-\frac12(A_s|u|^2)_x\).
For the third line use the exact spherical identity
\(z_x\cdot(z_x-z_y)=|z_x-z_y|^2/2\).
It follows pointwise that \(\gamma\widehat\rho_s\to{\cal R}_s\)
and \(\gamma^2\sum_x\widehat q_{s,x}h_{t,x}\to
\sum_xQ_{s,x}H_{t,x}\).

Here is the required moment justification. Inside a sufficiently
small fixed coordinate neighborhood, the Jacobian is bounded,
and the exact transverse contribution gives
\(\gamma S_0\ge\frac12 u^{\mathsf T}Lu\).
On the mean-zero subspace this is at least
\(\frac12\lambda_{\min}^+(L)|u|^2\), with a positive constant
at this fixed \(N\). The zero-row spherical identity bounds
\(|\gamma\widehat q_{s,x}|\) by a constant times \(|u|^2\);
the analogous energy identity bounds \(|\gamma h_{t,x}|\) by
a constant times \(|u|^2\). Products of the observables are
therefore dominated by fixed polynomials times this integrable
Gaussian. The local normalizer, after the common power of
\(\gamma\) is removed, tends to a strictly positive Gaussian
integral.

Outside a smaller neighborhood of the ground-state set,
compactness gives \(S_0\ge\varepsilon_N>0\). The scores,
energies and all fixed products grow at most polynomially in
\(\gamma\), so their contribution divided by the local
normalizer tends to zero exponentially. A finite collection of
orientation charts covers the ground-state manifold.
Dominated convergence proves all claimed moments, with limiting
covariance \(C=L^\dagger\) in each of three colors.

The first moment statement gives convergence of the finite
source covariance matrices, and the second gives convergence of
the mean-return matrices. At fixed \(N\) the source dimension
\(m\) is finite, so entrywise convergence implies operator-norm
convergence. Subtraction gives the third assertion and all
norm limits. In particular \({\cal B}_1\) is symmetric, because
the exact compact \(\widehat{\cal B}\) is symmetric.
\end{proof}

The theorem makes the connection to the compact model explicit.
It does \emph{not} interchange \(\sup_N\) and
\(\lim_{\gamma\to\infty}\). For example, it proves
\(\sup_N\lim_{\gamma\to\infty}\gamma\widehat a\le39/2\),
not \(\limsup_{\gamma\to\infty}\sup_N\gamma\widehat a\le39/2\).

\subsection{An exact ambient realization of the remaining moment problem}

The fixed-volume Taylor proof identifies the coefficient, but no Taylor
remainder is needed merely to express the score under its actual law.
Extend the dot-product polynomials in \eqref{r58:polynomials} to
four-component site vectors, retaining the constant \(3\) in
\({\cal R}_s\). For every compact configuration define
\[
 M_z=\frac1n\sum_xz_x,\qquad
 y_x=\sqrt\gamma(z_x-M_z),\qquad \nu_\gamma=\operatorname{Law}_\mu(y).
 \tag{V58.16}\label{r58:ambient}
\]
This definition includes \(M_z=0\).

\begin{proposition}[Exact polynomial source realization]
\label{r58:exact-ambient}
Pointwise on the full compact configuration space,
\[
 \gamma\widehat q_{s,x}=Q_{s,x}(y),\qquad
 \gamma h_{t,x}=H_{t,x}(y),\qquad
 F_t=f_t(y),\qquad \gamma\widehat\rho_s={\cal R}_s(y).
 \tag{V58.17}\label{r58:exact-ambient-eq}
\]
Consequently
\[
 \begin{aligned}
 \gamma^2\Cov_\mu(\widehat\rho_s,\widehat\rho_t)
       &=\Cov_{\nu_\gamma}({\cal R}_s,{\cal R}_t),\\
 \gamma\widehat{\cal N}(s,t)
       &=\E_{\nu_\gamma}\sum_xQ_{s,x}H_{t,x},\\
 \gamma\widehat{\cal B}(s,t)
       &=\Cov_{\nu_\gamma}({\cal R}_s,f_t)
                              -\E_{\nu_\gamma}\sum_xQ_{s,x}H_{t,x}.
 \end{aligned}
 \tag{V58.18}\label{r58:moment-targets}
\]
\end{proposition}
\begin{proof}
Put \(v_x=z_x-M_z\). The unit-sphere constraint says
\[
 M_z\cdot v_x=\tfrac12(1-|M_z|^2-|v_x|^2).
\]
Since \(A_s\mathbf1=0\),
\[
 z_x\cdot(A_sz)_x
   =v_x\cdot(A_sv)_x-\tfrac12(A_s|v|^2)_x .
\]
Scale by \(\gamma\). For the other identities use
\(Bz=Bv\) and
\(z_x\cdot(z_x-z_u)=|z_x-z_u|^2/2\).
The definitions of the three normalized returns then give
\eqref{r58:moment-targets}. No small-field event or conditioning
was introduced.
\end{proof}

For comparison on the same ambient space, let \(\nu_0\) be the
law of \(O(0,u_x)\), where \(O\) is an independent Haar rotation
in \(SO(4)\) and \(u\) is the three-component Gaussian field above.
All the displayed polynomials depend only on dot products, so
their \(\nu_0\) expectations are exactly the Gaussian expectations
already evaluated. The actual \(\nu_\gamma\) is not declared
Gaussian; in particular replacing it by a four-independent-color
Gaussian would change the contraction \(9r-3=-3/n\) used in
Lemma \ref{r58:chaos}.

Thus the uniform question concerns the indicated fourth-, sixth-,
and eighth-degree connected moment contractions under
\(\nu_\gamma\) versus \(\nu_0\), in the full signed-source norms.
The identities do not supply smallness of those differences.
They remove a coordinate-expansion remainder from the statement
of the target, not the non-Gaussian probability-law remainder.

\subsection{What this pays, and what must now be attacked}

The leading centered coefficient, its mean return, and its
combined native return are bounded for the entire source family
uniformly in volume. The proof uses the projected bubble,
weighted balanced inverse and positive local gradient partition,
followed by second/fourth-chaos separation. These are new
coefficient estimates; V6's fixed-cube tangent-score limit and
the f5 V24 abstract cumulant-Wick pairing are different results
and are not counted again. No Ising, \(O(2)\), or unverified
\(O(4)\) correlation inequality has been imported.

The next unresolved object is the \emph{uniform compact moment remainder}
in \eqref{r58:moment-targets} after this coefficient has been subtracted.
The fixed-volume
proof above uses both \(\lambda_{\min}^+(L)\) and
\(\varepsilon_N\); it provides no uniform modulus along a
growing-volume bare-coupling trajectory. V57 remains the paid
nonperturbative bound until such a remainder is established.
Neither entropy density alone nor the present Gaussian
contractions justify making that replacement.

The useful upstream structure is now explicit: preserve the
two-sided source inverse and the projected bubble when estimating
connected marked remainders, rather than restoring absolute site
sums. Winding and general noncommuting exteriors, original
exterior probabilities and predictor covariances, physical source
coverage, physical-time contraction, and the interacting
four-dimensional continuum construction remain unresolved.
The full Yang--Mills mass gap is unproved.

The diagnostic for this appendix checks finite graph matrix
identities and Gaussian Wick coefficients; numerical checks are
distinguished from exact arithmetic. Neither type of finite check
proves the analytic or continuum statements. The next scheduled
companion review and broad primary-literature sweep are at V60.
% END CONSOLIDATED V58 APPEND


% BEGIN CONSOLIDATED V59 APPEND
\clearpage
\section{V59: An exact compact distance-polynomial recursion}
\label{r59:section}

\subsection{Purpose, transfer audit, and the unchanged law}

V58 bounded the entire leading Gaussian source bank, but identified it
with the compact law only at fixed volume. We now go upstream of that
identification. A Green-weighted rotational integration by parts gives an
\emph{exact} compact generator equal to a three-color Gaussian polynomial
generator plus one explicit \(1/\gamma\) correction. Its polynomial inverse
turns moment differences into actual compact expectations, with no chart
cutoff, selected phase, or unproved inverse of a native generator.
The new identity does not by itself bound those expectations in the full
source norm uniformly in volume.

The targeted literature input is Giuliani--Ott,
\href{https://arxiv.org/abs/2302.07299}{\emph{Low temperature asymptotic
expansion for classical \(O(N)\) vector models}},
Prob.\ Math.\ Phys.\ 6 (2025), 439--477.
Their Theorem 1.3 and integration-by-parts method motivate working on the
moment remainder, not recalculating V58's Wick coefficient. Their selected
infinite-volume state and fixed-multipoint error norm are not the present
finite periodic law and signed-source norm. No such transfer is assumed.
The required identities below are proved directly.

This also respects the supplied ESM transfer memo: keep the native marked
response and every return term, without importing a perturbative continuum
claim. For non-duplication, f16.1/f16.2 V17 already contain polynomial
Gaussian divergence solvers for a pinned massive history, and f15 V64
contains linear/cubic Haar Stein bounds for framed curvature. We credit
that architecture. The new object here is the exact global
distance-polynomial generator for the unpinned massless coherent law;
neither a general Stein solver nor a new Gaussian coefficient is being
claimed as the advance.

Keep V57--V58's even cubic torus, \(N\ge4\), \(n=N^3\), \(m=3n\),
\(L=B^{\mathsf T}B\), \(C=L^\dagger\), and \(c=C_{xx}\le13/16\).
The law throughout is exactly
\[
 d\mu_\gamma(z)=Z_\gamma^{-1}
       \exp\{-\gamma S_0(z)\}\prod_xd\sigma(z_x),\qquad
 S_0=\sum_{e=(x,y)}(1-z_x\cdot z_y),\quad z_x\in S^3 .
\]
The balanced \(A_s\) of V58 is unchanged. All polynomials below are finite;
no infinite-volume generator is constructed.

\subsection{All-pair distances and a global rotational generator}

Let \(\mathcal I=\{(i,j):i<j\}\) be the set of all unordered site pairs,
with this fixed orientation. Put
\[
 d_{ij}=\gamma(1-z_i\cdot z_j),\qquad d_{ji}=d_{ij},\quad d_{ii}=0,
 \qquad b_{ij}=e_i-e_j .
 \tag{V59.1}\label{r59:distances}
\]
The same letters \(d_{ij}\) will be indeterminates in the formal polynomial
ring \(\mathbb R[d_I:I\in\mathcal I]\). We evaluate that ring on compact
configurations only after differentiating. For \(I=(i,j)\), \(J=(k,l)\), set
\[
 \begin{aligned}
 \kappa_{IJ}&=b_I^{\mathsf T}Cb_J,& a_I&=\kappa_{II},\\
 T_{IJ}(d)&=d_{il}+d_{jk}-d_{ik}-d_{jl},&
 W_{IJ}(d)&=d_{ik}d_{jl}-d_{il}d_{jk}.
 \end{aligned}
 \tag{V59.2}\label{r59:kernels}
\]
Every nearest-neighbor edge is given the orientation of its pair.
Changing an edge orientation in intermediate incidence formulas changes
both relevant signs and does not change the resulting operator.

For ambient indices \(1\le a<b\le4\), define the sphere-tangent rotation
\[
 \mathcal R_x^{ab}=z_x^a\partial_{z_x^b}-z_x^b\partial_{z_x^a}.
\]
Write
\[
 \begin{aligned}
 \Delta_C&=\sum_{x,y}C_{xy}\sum_{a<b}\mathcal R_x^{ab}\mathcal R_y^{ab},\\
 \Gamma_C(F,G)&=\sum_{x,y}C_{xy}\sum_{a<b}
                   (\mathcal R_x^{ab}F)(\mathcal R_y^{ab}G),\\
 \mathcal A_\gamma F&=\gamma^{-1}
       \{\Delta_CF-\Gamma_C(\gamma S_0,F)\}.
 \end{aligned}
 \tag{V59.3}\label{r59:compact-generator}
\]
The positive semidefiniteness of \(C\) makes
\(\Gamma_C(F,F)\ge0\). No ellipticity or spectral gap for
\(\mathcal A_\gamma\) is asserted.

\begin{lemma}[Exact compact stationarity]
\label{r59:stationary}
For every smooth function \(F\) on \((S^3)^n\),
\(\E_{\mu_\gamma}\mathcal A_\gamma F=0\).
In particular this holds for every polynomial in the distances
\eqref{r59:distances}, at every \(\gamma>0\).
\end{lemma}
\begin{proof}
Each \(\mathcal R_x^{ab}\) is divergence-free for product Haar measure.
Integrating the first rotation in each term of \(\Delta_CF\) by parts
against \(e^{-\gamma S_0}\) gives
\(\E\Delta_CF=\E\Gamma_C(\gamma S_0,F)\).
There is no boundary. All functions and derivatives at this fixed finite
graph and coupling are integrable by compactness. The semidefinite,
rather than invertible, nature of \(C\) causes no difficulty in this
finite sum identity.
\end{proof}

\begin{theorem}[Exact two-term polynomial generator]
\label{r59:generator}
Define two operators on the formal distance-polynomial ring by
\[
 \begin{aligned}
 \mathcal A_0P={}&
   \sum_I(3a_I-2d_I)\partial_IP
       +\sum_{I,J}\kappa_{IJ}T_{IJ}\partial_I\partial_JP,\\
 \mathcal A_1P={}&
   \sum_{I,J}\kappa_{IJ}W_{IJ}\partial_I\partial_JP\\
   &-\sum_I\left(3a_Id_I+
                \sum_{e\in E}\kappa_{Ie}W_{Ie}\right)\partial_IP .
 \end{aligned}
 \tag{V59.4}\label{r59:A01}
\]
For every such \(P\), after evaluation on compact configurations,
\[
       \mathcal A_\gamma P=\mathcal A_0P+\gamma^{-1}\mathcal A_1P .
 \tag{V59.5}\label{r59:exact-generator}
\]
There are no further powers of \(1/\gamma\) in this generator identity.
\end{theorem}
\begin{proof}
The sum of the six squared rotations at one site acts on each coordinate
of \(z_x\) as multiplication by \(-3\). The cross-site rotation sum acting
on \(z_i\cdot z_j\) is \(3z_i\cdot z_j\). Consequently
\[
 \gamma^{-1}\Delta_Cd_I=3a_I(1-d_I/\gamma).
 \tag{V59.6}\label{r59:lap-distance}
\]
At the two ends of \(I\), the rotational derivatives of \(d_I\) are
opposite, with magnitude tensor \(\gamma z_i\wedge z_j\). The identity
\[
 \langle z_i\wedge z_j,z_k\wedge z_l\rangle
       =(z_i\cdot z_k)(z_j\cdot z_l)
                        -(z_i\cdot z_l)(z_j\cdot z_k)
\]
therefore gives exactly
\[
 \gamma^{-1}\Gamma_C(d_I,d_J)
       =\kappa_{IJ}\{T_{IJ}(d)+\gamma^{-1}W_{IJ}(d)\}.
 \tag{V59.7}\label{r59:diffusion}
\]
The wedge inner product here sums over \(a<b\), with no extra factor two.

To evaluate the drift, let \(D=(d_{ij})_{ij}\) be any symmetric matrix
with zero diagonal. Then \(T_{Ie}=-b_I^{\mathsf T}Db_e\) and
\[
 \sum_{e\in E}\kappa_{Ie}T_{Ie}
    =-b_I^{\mathsf T}DLCb_I
    =-b_I^{\mathsf T}Db_I=2d_I .
 \tag{V59.8}\label{r59:drift-collapse}
\]
We used \(LC=\Pi\) and \(\Pi b_I=b_I\), not a mass for the constant mode.
Since \(\gamma S_0=\sum_ed_e\), subtracting its carré-du-champ drift
from \eqref{r59:lap-distance} gives the first-derivative coefficients
of \eqref{r59:A01}. The ordinary second-order chain rule, with ordered
sum over \(I,J\), supplies \eqref{r59:diffusion} as the coefficient of
\(\partial_I\partial_JP\). This proves the identity for every polynomial.
\end{proof}

\subsection{A formal inverse that must not quotient by Gaussian rank}

Let \(\mathfrak G P\) mean expectation of \(P\) evaluated at
\[
           d^G_{ij}=\tfrac12|u_i-u_j|^2,
 \qquad u^1,u^2,u^3\ \hbox{independent }N(0,C).
 \tag{V59.9}\label{r59:Gaussian}
\]
The Gaussian operator
\(\sum_{x,y}C_{xy}\nabla_{u_x}\cdot\nabla_{u_y}
 -\sum_xu_x\cdot\nabla_{u_x}\) acts on \(P(d^G)\) as
\((\mathcal A_0P)(d^G)\): its action on \(d^G_I\) is
\(3a_I-2d^G_I\), and its diffusion on two distances is
\(\kappa_{IJ}T_{IJ}(d^G)\). Gaussian integration by parts gives
\(\mathfrak G\mathcal A_0P=0\).

This verifies the number \(3\) directly. It is not replaced by four
Gaussian colors merely because the compact spins lie in \(\mathbb R^4\).
Likewise the formal ring is \emph{not} quotiented by rank-three Gram
relations: such identities hold on the reference support but need not
hold for compact distance data. The following construction is an identity
in the full formal ring.

Write \(\mathcal A_0=-2\mathcal E+\mathcal D\), where
\(\mathcal E=\sum_Id_I\partial_I\) counts distance degree and
\[
 \mathcal D=3\sum_Ia_I\partial_I+
                 \sum_{I,J}\kappa_{IJ}T_{IJ}\partial_I\partial_J .
\]
It lowers homogeneous distance degree by one. If \(P=\sum_{j=0}^kP_j\),
define homogeneous polynomials recursively by
\[
 S_{k+1}=0,\qquad
 S_j=\frac{P_j+\mathcal D S_{j+1}}{2j}\quad(j=k,\ldots,1),\qquad
 \mathcal SP=\sum_{j=1}^kS_j-\mathfrak G\!\left(\sum_{j=1}^kS_j\right).
 \tag{V59.10}\label{r59:inverse}
\]
For constant \(P\), set \(\mathcal SP=0\).

\begin{proposition}[Finite polynomial inverse]
\label{r59:inverse-prop}
One has, as formal polynomial identities,
\[
       -\mathcal A_0\mathcal SP=P-\mathfrak GP,\qquad
       \mathfrak G\mathcal SP=0,\qquad \deg\mathcal SP\le\deg P .
 \tag{V59.11}\label{r59:inverse-eq}
\]
Moreover this is the unique solution with zero Gaussian expectation.
\end{proposition}
\begin{proof}
The recursion makes each positive-degree component on the left equal
to the corresponding component of \(P\). Their difference is constant.
Applying \(\mathfrak G\), using its stationarity for \(\mathcal A_0\),
fixes that constant as \(-\mathfrak GP\). The displayed normalization is
part of the definition. If a nonconstant polynomial of top degree \(j\)
were in the kernel of \(\mathcal A_0\), its top component would be
multiplied by \(-2j\), a contradiction. Thus the kernel consists only of
constants, proving uniqueness without an inverse on a native function space.
\end{proof}

\subsection{Exact native returns and explicit finite-volume remainder bounds}

Define the entirely explicit linear polynomial operator
\[
                     \mathcal TP=\mathcal A_1\mathcal SP .
 \tag{V59.12}\label{r59:return-operator}
\]
It increases distance degree by at most one.

\begin{theorem}[An exact compact weak expansion with a native remainder]
\label{r59:expansion}
For every distance polynomial \(P\) and every \(\gamma>0\),
\[
 \E_{\mu_\gamma}P-\mathfrak GP
                   =\gamma^{-1}\E_{\mu_\gamma}\mathcal TP .
 \tag{V59.13}\label{r59:one-step}
\]
For every integer \(J\ge1\), consequently,
\[
 \E_{\mu_\gamma}P
     =\sum_{j=0}^{J-1}\gamma^{-j}\mathfrak G(\mathcal T^jP)
                   +\gamma^{-J}\E_{\mu_\gamma}\mathcal T^JP .
 \tag{V59.14}\label{r59:iteration}
\]
The remainder is an expectation under the original complete compact law.
\end{theorem}
\begin{proof}
Apply Lemma \ref{r59:stationary} to \(\mathcal SP\), and use
\eqref{r59:exact-generator} and \eqref{r59:inverse-eq}.
This proves \eqref{r59:one-step} with its displayed positive sign.
Apply the same identity to \(\mathcal TP,\mathcal T^2P,\ldots\)
and substitute finitely many times. Constants are annihilated by
\(\mathcal T\), so no separate normalization series is omitted.
\end{proof}

Here is one fully paid but deliberately coarse quantitative norm.
For a collected polynomial \(P=\sum_\alpha p_\alpha d^\alpha\), set
\[
 \|P\|_{\rm coef}=\sum_\alpha|p_\alpha|,\qquad
 \|P\|_{\rm mom}=\sum_\alpha|p_\alpha|M_{|\alpha|},\qquad
 M_0=1,\quad M_j=4j!(8c)^j\quad(j\ge1).
 \tag{V59.15}\label{r59:norms}
\]

\begin{proposition}[A rigorous remainder norm and the location of volume loss]
\label{r59:norm-prop}
For every \(P\),
\[
       |\E_{\mu_\gamma}P|\le\|P\|_{\rm mom},\qquad
 \left|\E_{\mu_\gamma}P-
       \sum_{j=0}^{J-1}\gamma^{-j}\mathfrak G(\mathcal T^jP)\right|
                  \le\gamma^{-J}\|\mathcal T^JP\|_{\rm mom}.
 \tag{V59.16}\label{r59:remainder-bound}
\]
These inequalities are valid without a lower bound on \(\gamma\).
The coefficient norm of \(\mathcal S\) on degree at most \(k\) is
bounded by a constant depending only on \(k,c\), not the number of sites.
For a homogeneous polynomial \(P_j\), the following coarse bound holds:
\[
 \|\mathcal A_1P_j\|_{\rm coef}
 \le\{8cj(j-1)+12cj+4j\sqrt{mc}\}\|P_j\|_{\rm coef}.
 \tag{V59.17}\label{r59:A1-coef}
\]
Thus the displayed absolute-coefficient estimate for the return operator
still incurs a growing-volume loss.
\end{proposition}
\begin{proof}
The inherited coherent infrared moment bound gives
\(\|d_{ij}\|_{L^j}\le A_j a_{ij}\), where
\(A_j^j=4j!2^j\) and \(a_{ij}\le4c\). Hölder, allowing repeated pairs,
bounds the expectation of each degree-\(j\) distance monomial by \(M_j\).
Sum absolute coefficients. The second claim in \eqref{r59:remainder-bound}
then follows from the exact native remainder in \eqref{r59:iteration}.

For completeness the same \(M_j\) bounds the corresponding Gaussian
monomials: \(d^G_{ij}\) has law \((a_{ij}/2)\chi^2_3\), whose \(j\)-th
moment is at most \(4j!2^ja_{ij}^j\). Hölder again suffices.
In particular, on degree at most \(k\), subtracting a Gaussian expectation
costs at most a factor \(1+\max_{j\le k}M_j\) in coefficient norm.

Positive semidefiniteness gives
\(|\kappa_{IJ}|\le\sqrt{a_Ia_J}\le4c\).
Also \(\|T_{IJ}\|_{\rm coef}\le4\).
For a degree-\(j\) homogeneous polynomial,
the sums of coefficient norms of first and ordered second derivatives
are at most \(j\|P_j\|_{\rm coef}\) and
\(j(j-1)\|P_j\|_{\rm coef}\), respectively. Hence
\[
       \|\mathcal DP_j\|_{\rm coef}
             \le4cj(4j-1)\|P_j\|_{\rm coef}.
 \tag{V59.18}\label{r59:D-coef}
\]
The finite downward recursion \eqref{r59:inverse}, followed by its constant
normalization, proves the asserted dimension-independent bound for
\(\mathcal S\).

Finally \(\|W_{IJ}\|_{\rm coef}\le2\), and
\[
 \sum_{e\in E}\kappa_{Ie}^2
       =b_I^{\mathsf T}CLCb_I=a_I,\qquad
 \sum_{e\in E}|\kappa_{Ie}|\le\sqrt{ma_I}\le2\sqrt{mc}.
 \tag{V59.19}\label{r59:dipole}
\]
The three terms in \(\mathcal A_1\) are bounded respectively by
\(8cj(j-1)\), \(12cj\), and \(4j\sqrt{mc}\), times the coefficient norm.
This proves \eqref{r59:A1-coef}. At every fixed finite graph,
\(\|\mathcal T^JP\|_{\rm mom}\) is finite and independent of \(\gamma\);
thus \eqref{r59:iteration} is also a genuine fixed-volume asymptotic
expansion to every finite order. Nothing in this proof makes that
coefficient norm uniformly bounded for the growing full source bank.
\end{proof}

The Gaussian polynomial inversion is therefore not where a site count is
lost. The explicit force return
\(\sum_e\kappa_{Ie}W_{Ie}\), and subsequently the source contraction,
are the places where taking absolute values can destroy the needed
cancellation. Bound \eqref{r59:A1-coef} is a valid fallback, not a claimed
improvement of V57's nonperturbative source bound.

\subsection{Application to the complete V58 source bank, with centering retained}

The source polynomials can be written directly in this ring:
\[
 \begin{aligned}
 Q_{s,x}(d)&=-\sum_y(A_s)_{xy}d_{xy},&
 H_{t,x}(d)&=\sum_{e\ni x}t_ed_e,\\
 f_t(d)&=\sum_et_ed_e,&
 \mathcal R_s(d)&=\sum_xQ_{s,x}(d)(H_{\mathbf1,x}(d)-3),\\
 \mathcal N_{s,t}(d)&=\sum_xQ_{s,x}(d)H_{t,x}(d).
 \end{aligned}
 \tag{V59.20}\label{r59:source-bank}
\]
They are exactly V58's polynomials: use \(A_s\mathbf1=0\) and
\(d_{xy}=\frac12|y_x-y_y|^2\). In particular
\(\gamma\widehat\rho_s=\mathcal R_s\) pointwise under the compact law.
Their respective distance degrees are \(1,1,1,2,2\).

Put \(r_s=\mathcal R_s-\mathfrak G\mathcal R_s\) and
\(g_t=f_t-\mathfrak Gf_t\), as \emph{formal polynomials}. In the following
identities every undecorated expectation is \(\E_{\mu_\gamma}\).

\begin{corollary}[Exact contracted moment returns]
\label{r59:source-returns}
One has
\[
 \begin{aligned}
 &\gamma^2\Cov(\widehat\rho_s,\widehat\rho_t)
             -\Cov_G(\mathcal R_s,\mathcal R_t)\\
 &\hspace{8mm}
   =\gamma^{-1}\E\mathcal T(r_sr_t)
       -\gamma^{-2}(\E\mathcal T\mathcal R_s)
                         (\E\mathcal T\mathcal R_t),\\
 &\gamma\widehat{\mathcal N}(s,t)-\mathcal N_1(s,t)
       =\gamma^{-1}\E\mathcal T\mathcal N_{s,t},\\
 &\gamma\widehat{\mathcal B}(s,t)-\mathcal B_1(s,t)\\
 &\hspace{8mm}
   =\gamma^{-1}\E\mathcal T(r_sg_t-\mathcal N_{s,t})
       -\gamma^{-2}(\E\mathcal T\mathcal R_s)(\E\mathcal Tf_t).
 \end{aligned}
 \tag{V59.21}\label{r59:contracted}
\]
The largest distance degrees of the first-order returned polynomials in
these three lines are \(5,3,4\), respectively.
\end{corollary}
\begin{proof}
For any two polynomials \(P,Q\), centered at their Gaussian expectations,
\[
 \Cov_\mu(P,Q)=\E_\mu[(P-\mathfrak GP)(Q-\mathfrak GQ)]
       -(\E_\mu P-\mathfrak GP)(\E_\mu Q-\mathfrak GQ).
\]
Apply \eqref{r59:one-step} to the product and to the two individual
polynomials. This proves the first line using V58.18.
Apply the same identity to \(\mathcal N_{s,t}\) for the second.
For the third use the first product argument with
\((P,Q)=(\mathcal R_s,f_t)\), then subtract the complete mean return.
The degree claims follow from \(\deg\mathcal TP\le\deg P+1\).
In particular the last mean-product term cannot be discarded when
estimating the mixed return.
\end{proof}

For a precise remaining criterion, define
\[
 \begin{aligned}
 \tau_R&=\sup_{\|s\|_2=1}|\E\mathcal T(r_s^2)|,&
 \tau_N&=\sup_{\|s\|_2=\|t\|_2=1}|\E\mathcal T\mathcal N_{s,t}|,\\
 \tau_B&=\sup_{\|s\|_2=\|t\|_2=1}
                    |\E\mathcal T(r_sg_t-\mathcal N_{s,t})|,\\
 \alpha&=\sup_{\|s\|_2=1}|\E\mathcal T\mathcal R_s|,&
 \beta&=\sup_{\|t\|_2=1}|\E\mathcal Tf_t|.
 \end{aligned}
 \tag{V59.22}\label{r59:targets}
\]
These quantities are finite at each finite graph; no uniform bound is
part of their definition. The paid V58 coefficients and
\eqref{r59:contracted} imply
\[
 \begin{aligned}
 \gamma\widehat a&\le
                \sqrt{(24c)^2+\tau_R/\gamma},\\
 \gamma\widehat b&\le21cr+\tau_N/\gamma,\\
 \gamma\|\widehat{\mathcal B}\|_{2\to2}
             &\le75c/4+\tau_B/\gamma+\alpha\beta/\gamma^2 .
 \end{aligned}
 \tag{V59.23}\label{r59:conditional-norms}
\]
For the first inequality the subtracted squared mean in the diagonal
case of \eqref{r59:contracted} is nonnegative, so it can be dropped in an
upper variance bound. This observation does not apply to the mixed
product in the last line.

\subsection{Status and the next upstream estimate}

The new result is an exact compact-law polynomial recursion and its
explicit native remainder, including the complete source-centering
corrections. It replaces a qualitative fixed-volume Laplace comparison
by a computable finite-order identity with a rigorous remainder norm.
It does not replace the uniform moment problem by an assumption that
its answer is small.

The next useful calculation is to contract the determinant force return
\(W_{Ie}\) against V58's balanced source tensors \emph{before} using
\eqref{r59:dipole}'s absolute sum. The projected bubble and positive
gradient partition remain relevant there. The bounds in
\eqref{r59:targets}, for the actual compact law and full signed sources,
remain open uniformly in volume. V57 therefore remains the paid
nonperturbative estimate; the displayed polynomial-coefficient fallback
is not a new physical scaling regime.

No extension to typical noncommuting exteriors, original exterior
probabilities, predictor covariances, physical source coverage or
physical-time contraction has been proved. The interacting
four-dimensional continuum construction and the full Yang--Mills mass
gap remain unproved. The next companion review and broad
primary-literature sweep are both due at V60.

The V59 diagnostic independently checks the rotation/diffusion identities
on rational compact configurations, including zero magnetization, and
the formal inverse against exact three-color Wick moments on a triangle.
It also checks all three source-return polynomial types and a
single-bond coefficient sequence. A five-site exact fixture confirms
that a rank-four distance Gram determinant can be nonzero on compact
configurations although its three-color Gaussian expectation vanishes;
the formal inverse retains this polynomial. Independent single-bond
quadrature checks the native remainder identity at two resolutions.
The quadrature is floating-point, not an exact or analytic proof.
These finite checks detect algebraic mistakes; they do not prove a
volume-uniform theorem.
% END CONSOLIDATED V59 APPEND


% BEGIN CONSOLIDATED V60 APPEND
\clearpage
\section{V60: Connected returns and a log-free compact bound}
\label{r60:section}

\subsection{The scheduled reviews and the resulting adjustment}

V59's compact polynomial generator is retained exactly. The new work
first reorganizes its remainder as a \emph{native connected covariance},
then identifies an energy-source drift with an exact cycle-space
representation. This yields a genuine improved nonperturbative estimate:
the coherent compact energy covariance approaches its Gaussian coefficient
along \(\gamma/N^{3/2}\to\infty\), without V57's additional logarithm.
Neither uniformity at fixed large \(\gamma\) over all volumes nor a
physical Yang--Mills mass gap follows.

The compulsory companion checkpoint re-inventoried the active files and
verified unchanged hashes for NS V26, SM--Einstein V72, Shell Mixing V16,
and Parallel YM V5, and for the downloaded Source Selection V35,
Stationarity V38, Exact-Action Bridge V28, and Perturbative ESM V31.
Their terminal proof/scope passages and the supplied ESM transfer memo
were reread. This is not a new certification of their entire histories.
The ESM lesson remains useful: retain the combined native return before
estimating individual pieces. The SM many-copy limit and the shell
programme's flat-holonomy coefficient do not provide the missing
nonperturbative YM source estimate. The NS rank-one derivative
refinement and the bridge's finite coefficient charts do not change that.

The scheduled broad primary-literature sweep covered low-temperature
spin expansions, matrix Stein methods, multiscale functional inequalities,
nonlinear-sigma background fields, random-loop/forest representations,
lattice area laws, stochastic Yang--Mills, and physical spectral moments.
The relevant decisions are:
\begin{itemize}
\item Giuliani--Ott's
\href{https://arxiv.org/abs/2302.07299}{low-temperature expansion}
still motivates integration by parts, but its selected-state and
fixed-multipoint conclusions do not supply our full source norm.
\item Bailly--Gaunt--Ouimet--Richards--von Sachs,
\href{https://arxiv.org/html/2606.04859v1}{\emph{Stein's method for the
Wishart distribution}} (2026), supplies useful generator context.
Its Theorem 3.6 assumes a positive-definite scale and
\(\alpha>3d-3\). Our site Gram reference has only three Gaussian colors
and is singular when \(d=n-1>3\); that theorem is not imported.
V59's full formal polynomial ring is retained instead.
\item Bauerschmidt--Bodineau--Dagallier's
\href{https://arxiv.org/html/2307.07619}{Polchinski-flow framework},
Theorems 3.6 and 3.20, still requires actual renormalized-Hessian
control, and in the curved version a further continuity condition.
An averaged diffusion estimate below does not verify those hypotheses.
\item The
\href{https://arxiv.org/abs/2509.04688}{dynamical lattice area-law work}
and \href{https://arxiv.org/abs/2505.16585}{expanded area-law regimes}
are not used as weak-coupling continuum theorems.
Chevyrev--Shen's
\href{https://arxiv.org/abs/2503.03060}{stochastic three-dimensional
Yang--Mills--Higgs renormalization result} does not supply a
four-dimensional invariant measure and gap.
\item The
\href{https://arxiv.org/abs/2508.01377}{matrix spectral-moment method}
requires physical Euclidean correlators as input; the auxiliary spatial
Dirichlet form below is not such a correlator. Random-forest and
nonlinear-sigma background-field results were screened as method
comparisons, not transferred across their different measures.
\end{itemize}
Only the indicated statements, hypotheses and relevant primary passages
were reviewed; no claim of a full proof audit of every screened paper is
made. The detailed source and companion ledgers accompany this version
outside \path{latex_notes}. The next companion and broad-literature
checkpoints are V65 and V70, respectively.

\subsection{Native centering before the return is estimated}

Use V59's \(\mu_\gamma,\mathcal A_0,\mathcal A_1,\mathcal S,\mathcal T\),
with \(\mathcal T=\mathcal A_1\mathcal S\), and Gaussian functional
\(\mathfrak G\). For formal distance polynomials define
\[
 \begin{aligned}
 \Gamma_0(P,Q)&=\sum_{I,J}\kappa_{IJ}T_{IJ}
                                  (\partial_IP)(\partial_JQ),\\
 \Gamma_1(P,Q)&=\sum_{I,J}\kappa_{IJ}W_{IJ}
                                  (\partial_IP)(\partial_JQ),\\
 \Gamma_\gamma(P,Q)&=\Gamma_0(P,Q)+\gamma^{-1}\Gamma_1(P,Q).
 \end{aligned}
 \tag{V60.1}\label{r60:gamma}
\]
On compact configurations \(\Gamma_\gamma=\gamma^{-1}\Gamma_C\);
in particular it is the carré du champ of the actual reversible
generator \(\mathcal A_\gamma\).

\begin{theorem}[A connected compact Stein identity]
\label{r60:connected}
For every pair of distance polynomials,
\[
 \Cov_\mu(P,Q)=\E_\mu\Gamma_\gamma(\mathcal SP,Q)
                         +\gamma^{-1}\Cov_\mu(\mathcal TP,Q).
 \tag{V60.2}\label{r60:connected-eq}
\]
Consequently
\[
 \begin{aligned}
 \Cov_\mu(P,Q)-\Cov_G(P,Q)
  =\gamma^{-1}\big\{&
       \E_\mu\mathcal T\Gamma_0(\mathcal SP,Q)
       +\E_\mu\Gamma_1(\mathcal SP,Q)\\
       &+\Cov_\mu(\mathcal TP,Q)\big\}.
 \end{aligned}
 \tag{V60.3}\label{r60:connected-difference}
\]
These are exact full-law identities. The last covariance is centered
under \(\mu_\gamma\), not under the reference Gaussian.
\end{theorem}
\begin{proof}
Reversible integration by parts on the compact product gives
\[
       -\E_\mu[Q\mathcal A_\gamma F]
                          =\E_\mu\Gamma_\gamma(F,Q).
\]
It follows directly by applying the divergence-free rotations of V59
to \(Qe^{-\gamma S_0}\); it does not require an inverse or a gap for
\(\mathcal A_\gamma\). V59's formal inverse gives
\[
 P-\E_\mu P=-\mathcal A_\gamma\mathcal SP
                    +\gamma^{-1}(\mathcal TP-\E_\mu\mathcal TP).
\]
Multiply by \(Q-\E_\mu Q\) and integrate to obtain
\eqref{r60:connected-eq}. Gaussian integration by parts gives
\(\Cov_G(P,Q)=\mathfrak G\Gamma_0(\mathcal SP,Q)\).
Apply V59.13 to the polynomial \(\Gamma_0(\mathcal SP,Q)\) and
expand \(\Gamma_\gamma\) to prove \eqref{r60:connected-difference}.
All operations take place in V59's full formal ring before evaluation.
\end{proof}

If \(P,Q\) have distance degrees \(k,l\), the two expectation integrands
in the braces have degree at most \(k+l\), and the remaining covariance
pairs degrees at most \(k+1\) and \(l\). One may average the formula with
its \(P,Q\)-interchange. Neither form closes the hierarchy automatically.
For the V58 score bank, substituting \(P=\mathcal R_s,Q=\mathcal R_t\)
or \(Q=f_t\) retains exactly the same physical source family.
For the mixed return, subtract
\(\gamma^{-1}\E_\mu\mathcal T\mathcal N_{s,t}\) as well.
No product of separately bounded mean shifts is then needed.

\begin{proposition}[Why a raw Gaussian-centered remainder can be extensive]
\label{r60:mean-shift}
There are connected compact \(O(4)\) graph models for which the variance
of a unit-normalized energy sum is bounded independently of graph size,
while the raw first return of its Gaussian-centered square grows linearly
in the number of edges. Thus such a raw-return bound is not a
generally necessary criterion for bounded connected covariance.
\end{proposition}
\begin{proof}
Take a star tree with \(M\) leaves and the same compact nearest-neighbor
action. Conditional on the central spin the leaves are independent, and
the distribution of each bond energy
\(d=\gamma(1-z_0\cdot z_i)\) is independent of that central direction.
The \(M\) energies are therefore independent also without conditioning.
For \(F_M=M^{-1/2}\sum_{i=1}^Md_i\),
\(\Var_\mu F_M=\Var_\mu d\). On a tree the edge gradient projector is
the identity, so the three-color Gaussian bond energies are independent
copies of \(\frac12\chi^2_3\), with mean \(3/2\) and variance \(3/2\).

V59's generator algebra and one-step identity hold on this connected
graph as well: their proof only uses \(LCb_I=b_I\). Hence
\[
 \E_\mu\mathcal T(F_M-\mathfrak GF_M)^2
   =\gamma\{\Var_\mu d-3/2+M(\E_\mu d-3/2)^2\}.
 \tag{V60.4}\label{r60:raw-example}
\]
At \(\gamma=1/2\), \(0\le d\le1\), so the right side is at least
\(M/8-3/4\), whereas \(\Var_\mu F_M\le1/4\).
The same phenomenon is possible at fixed sufficiently large coupling.
Indeed the one-bond density of \(d\), after cancelling normalization
factors, is
\[
       e^{-d}d^{1/2}(2-d/\gamma)^{1/2}\mathbf1_{[0,2\gamma]}(d).
\]
For \(\gamma\ge1\) its normalizer is bounded below by the integral
over \([0,1]\) of \(e^{-d}d^{1/2}\), and every moment has an integrable
gamma-density majorant. Dominated convergence gives
\(\E d\to3/2\), \(\E d^2\to15/4\).
The exact one-bond identity
\[
       \E d-3/2=(2\gamma)^{-1}\E(d^2-3d)
\]
then gives \(\gamma(\E d-3/2)\to-3/8\).
Thus the mean shift is nonzero for every sufficiently large fixed
\(\gamma\), and the last term of \eqref{r60:raw-example} again grows
linearly with \(M\). This is a diagnostic for the criterion, not a
counterexample to a uniform covariance conjecture on the cubic torus.
\end{proof}

V59.21--V59.23 remain valid as written. The adjustment is to avoid
requiring their individually separated raw terms to be uniformly small
when a connected identity can preserve their cancellation.

\subsection{The actual energy-source diffusion matrix}

Return to the even cubic torus. Orient each edge \(e=(i,j)\) as in V59 and
put
\[
 R=BCB^{\mathsf T},\qquad P_{\rm cyc}=I-R,\qquad
 v_e=\sqrt\gamma(z_j-z_i),\qquad
 w_e=\sqrt\gamma\,z_i\wedge z_j .
 \tag{V60.5}\label{r60:currents}
\]
The bivector has six Euclidean components indexed by \(a<b\).
Let
\[
 \mathcal M_{\gamma,ef}
       =R_{ef}\E_\mu(w_e\cdot w_f),\qquad
 \mathcal Q_{ef}=R_{ef}\E_\mu(v_e\cdot v_f).
 \tag{V60.6}\label{r60:diffusion-matrix}
\]
The second matrix is precisely V56's \(\mathcal Q\), since reversing
both gradient signs leaves inner products unchanged. The first is new:
for \(f_s=\sum_es_ed_e\),
\(\E\Gamma_\gamma(f_s,f_t)=s^{\mathsf T}\mathcal M_\gamma t\).
These matrices describe an auxiliary spatial generator, not physical
Euclidean-time propagation.

\begin{theorem}[Volume-uniform compact diffusion comparison]
\label{r60:metric}
Set \(h_\gamma=512cr/\gamma\), with \(r=R_{ee}=(n-1)/(3n)\).
Then
\[
 \begin{aligned}
 0\preceq\mathcal M_\gamma&\preceq4rI,\\
 \|\mathcal M_\gamma-\mathcal Q\|_{2\to2}
                        &\le4\sqrt{rh_\gamma}+h_\gamma,\\
 \|\mathcal M_\gamma-2K^0\|_{2\to2}
                        &\le\varepsilon_\gamma
             :=4\sqrt{rh_\gamma}+2h_\gamma
                                      +\tfrac83\sqrt{\eta_\gamma}.
 \end{aligned}
 \tag{V60.7}\label{r60:metric-eq}
\]
Here \(\eta_\gamma\) is V53's already proved entropy budget.
In particular \(\varepsilon_\gamma=O(\gamma^{-1/2})\) uniformly in
the torus size.
\end{theorem}
\begin{proof}
Fix any site \(o\), without conditioning its spin or changing the law,
and let \(u_x=\sqrt\gamma(z_x-z_o)\). Exactly,
\[
 w_e=a_e+\epsilon_e,\qquad
 a_e=z_o\wedge v_e,\qquad
 \epsilon_e=\gamma^{-1/2}u_i\wedge v_e .
 \tag{V60.8}\label{r60:root-decomposition}
\]
The infrared second-moment bound gives \(\E|v_e|^2\le4r\).
Also \(|w_e|^2=2d_e-d_e^2/\gamma\le2d_e=|v_e|^2\),
and \(|a_e|\le|v_e|\).
For any edge-vector field \(b\), positivity and \(R\preceq I\) give
\[
 0\le\E\langle sb,R(sb)\rangle
             \le\sum_es_e^2\E|b_e|^2 ,
 \tag{V60.9}\label{r60:weighted-field}
\]
where \((sb)_e=s_eb_e\). This proves the first bound and the analogous
\(4rI\) bounds for the forms built from \(v,a\).

Using the inherited moments \(A_p^p=4p!2^p\), with \(A_2^2=32\),
\[
 \|u_x\|_{L^4}^2\le8cA_2,\qquad
 \|v_e\|_{L^4}^2\le2rA_2 .
\]
Hölder therefore gives \(\E|\epsilon_e|^2\le512cr/\gamma=h_\gamma\).
For \(\ell_e=z_o\cdot v_e\), the sphere constraint gives
\[
 \ell_e=-\frac{|u_j|^2-|u_i|^2}{2\sqrt\gamma},\qquad
 |\ell_e|^2\le
       \frac{|v_e|^2(|u_i|^2+|u_j|^2)}{2\gamma}.
\]
The same moment estimate yields \(\E\ell_e^2\le h_\gamma\).
These are bounds under the original law for the fixed root, including
configurations far from any chosen chart.

Since \(|z_o|=1\),
\(v_e\cdot v_f=a_e\cdot a_f+\ell_e\ell_f\).
The difference between the \(v\)- and \(a\)-forms is consequently
positive semidefinite with norm at most \(h_\gamma\), by
\eqref{r60:weighted-field}. For the \(w\)- and \(a\)-forms, write
\[
 \langle sw,R(tw)\rangle-\langle sa,R(ta)\rangle
       =\langle s\epsilon,R(tw)\rangle
                              +\langle sa,R(t\epsilon)\rangle .
\]
Apply Cauchy--Schwarz in the actual probability space and in the
positive \(R\)-form. Each term is at most
\(\sqrt{4rh_\gamma}\|s\|_2\|t\|_2\) in absolute expectation.
This gives \(4\sqrt{rh_\gamma}+h_\gamma\).
Finally V56.7 already proves
\(\|\mathcal Q-2K^0\|\le(8/3)\sqrt{\eta_\gamma}+h_\gamma\).
Add the two bounds. The factor three in \(2K^0=3R\circ R\) comes from
that paid comparison, not from replacing the four-dimensional compact
field by four independent Gaussian colors.
\end{proof}

\subsection{An exact quartic drift with the entire cycle space retained}

Define the source-linear polynomial
\(\mathcal H_s=\mathcal A_1 f_s\), of distance degree at most two.
It is not V58's score \(\mathcal R_s\).

\begin{proposition}[Rooted Hodge representation of the native drift]
\label{r60:hodge}
Pointwise on the full compact space,
\[
 \begin{aligned}
 \mathcal A_\gamma f_s
       &=3r\sum_es_e-2f_s+\gamma^{-1}\mathcal H_s,\\
 \mathcal H_s
       &=\sum_es_e\{d_e^2-3rd_e
                      +\gamma\,w_e\cdot(P_{\rm cyc}w)_e\}.
 \end{aligned}
 \tag{V60.10}\label{r60:drift}
\]
For the same arbitrary root as above put
\[
 q_e=u_i\wedge u_j,\qquad
 k_e=\tfrac12\{|u_j|^2u_i-|u_i|^2u_j\},\quad e=(i,j).
 \tag{V60.11}\label{r60:quartic-fields}
\]
Then the equivalent expression, with no explicit coupling factor, is
\[
 \mathcal H_s=
       \sum_es_e(d_e^2-3rd_e)
       +\langle sq,P_{\rm cyc}q\rangle
       +\langle sv,P_{\rm cyc}k\rangle .
 \tag{V60.12}\label{r60:quartic-drift}
\]
For the constant source it simplifies to
\[
 \mathcal H_{\mathbf1}
       =\sum_e(d_e^2-3rd_e)+\|P_{\rm cyc}q\|_2^2 .
 \tag{V60.13}\label{r60:uniform-drift}
\]
\end{proposition}
\begin{proof}
V59.6 gives the Laplacian contribution
\(3r\sum_es_e-(3r/\gamma)f_s\).
The energy drift is
\(\Gamma_\gamma(f_{\mathbf1},f_s)=\sum_es_ew_e\cdot(Rw)_e\).
Subtract it and use \(|w_e|^2=2d_e-d_e^2/\gamma\) to obtain
\eqref{r60:drift}; comparison with V59.5 identifies
\(\mathcal H_s=\mathcal A_1f_s\).

Equation \eqref{r60:root-decomposition} also reads
\(w=z_o\wedge v+\gamma^{-1/2}q\).
Since \(v\) is an exact graph gradient, \(Rv=v\) and
\[
              P_{\rm cyc}w=\gamma^{-1/2}P_{\rm cyc}q .
\]
The remaining apparent factor \(\sqrt\gamma\) is removed by the exact
constraint \(z_o\cdot u_x=-|u_x|^2/(2\sqrt\gamma)\).
For \(e=(i,j)\), \(f=(a,b)\), it implies
\[
 \sqrt\gamma\,(z_o\wedge v_e)\cdot(u_a\wedge u_b)
   =\tfrac12 v_e\cdot(|u_b|^2u_a-|u_a|^2u_b)=v_e\cdot k_f.
\]
Sum with \((P_{\rm cyc})_{ef}\) to get \eqref{r60:quartic-drift}.
For \(s=\mathbf1\), \(\langle v,P_{\rm cyc}k\rangle=0\), and the
other quadratic form is \(\|P_{\rm cyc}q\|_2^2\).
\end{proof}

Here \(P_{\rm cyc}\) is the full orthogonal complement of exact gradients.
On the torus it includes the harmonic edge modes; none has been erased.
On a tree it is zero, consistently reducing the drift to the independent
one-bond term. The cycle-space statement is a graph identity, not an
assertion that these auxiliary currents are the physical four-dimensional
Yang--Mills curvature.

\subsection{A paid centered drift bound without the heat-kernel logarithm}

\begin{theorem}[A log-free finite-volume drift estimate]
\label{r60:drift-bound}
For every real signed source and every \(\gamma>0\),
\[
 \|\mathcal H_s-\E_\mu\mathcal H_s\|_{L^2(\mu)}
        \le K_H\sqrt m\,\|s\|_2\le544\sqrt m\,\|s\|_2,
 \tag{V60.14}\label{r60:H-bound}
\]
where
\[
 K_H=r^2(A_4^2+3A_2)
          +16crA_4^2+24crA_6\sqrt{A_2A_6}.
 \tag{V60.15}\label{r60:H-constant}
\]
This is an actual compact-law estimate, not a Gaussian coefficient bound.
\end{theorem}
\begin{proof}
For even \(p\ge2\), the same infrared moments give the uniform site and
edge bounds
\[
 U_p:=\sup_x\|u_x\|_{L^p}\le\sqrt{8cA_{p/2}},\qquad
 V_p:=\sup_e\|v_e\|_{L^p}\le\sqrt{2rA_{p/2}}.
\]
Since \(q_e=u_i\wedge v_e\),
\(\sup_e\|q_e\|_{L^4}\le U_8V_8=4\sqrt{cr}\,A_4=:Q_4\).
For the cubic field use
\[
 k_e=\tfrac12\{(|u_j|^2-|u_i|^2)u_i-|u_i|^2v_e\}.
\]
The difference of squares is bounded by
\((|u_i|+|u_j|)|v_e|\). Hölder with three \(L^{12}\) factors gives
\[
       \sup_e\|k_e\|_{L^4}
           \le\tfrac32 U_{12}^2V_{12}
           =12cA_6\sqrt{2rA_6}=:K_4 .
\]

For an arbitrary vector-valued edge field \(b\) with
\(\sup_e\|b_e\|_{L^4}\le B_4\), Minkowski applied to the squared
Euclidean norm gives
\[
 \|sb\|_{L^4(\ell^2)}\le B_4\|s\|_2,\qquad
 \|P_{\rm cyc}b\|_{L^4(\ell^2)}
                          \le\sqrt m\,B_4 .
 \tag{V60.16}\label{r60:field-L4}
\]
The second inequality uses the pointwise contraction of the full
orthogonal projector, not an absolute row sum of its kernel.
Therefore
\[
 \begin{aligned}
 \|\langle sq,P_{\rm cyc}q\rangle\|_{L^2}
                    &\le\sqrt m\,Q_4^2\|s\|_2,\\
 \|\langle sv,P_{\rm cyc}k\rangle\|_{L^2}
                    &\le\sqrt m\,V_4K_4\|s\|_2 .
 \end{aligned}
\]
For the local term,
\(\|d_e^2-3rd_e\|_{L^2}\le r^2(A_4^2+3A_2)\);
sum \(|s_e|\le\sqrt m\|s\|_2\).
Equation \eqref{r60:quartic-drift} proves
\(\|\mathcal H_s\|_{L^2}\le K_H\sqrt m\|s\|_2\).
Centering cannot increase this norm.

Finally \(A_2<6\), \(A_4^2<40\), \(A_6<8\), and
\(\sqrt{A_2A_6}<7\), directly from \(A_p^p=4p!2^p\).
Using \(r\le1/3,c\le13/16\),
\[
        K_H<58r^2+1984cr\le4894/9<544.
\]
No independence of compact fields or source-tilted infrared inequality
has entered the proof.
\end{proof}

\subsection{The improved native covariance regime and its precise boundary}

Let \(\mathcal C_\gamma(s,t)=\Cov_\mu(f_s,f_t)\), and define
\[
 a_H=\sup_{\|s\|_2=1}\|\mathcal H_s-\E\mathcal H_s\|_2,\qquad
 W_H=\frac{a_H}{4\gamma}
                  +\sqrt{2r+\left(\frac{a_H}{4\gamma}\right)^2}.
 \tag{V60.17}\label{r60:bootstrap-data}
\]

\begin{theorem}[Connected drift bootstrap for every signed source]
\label{r60:bootstrap}
The exact response identity and its consequences are
\[
 \begin{aligned}
 2\mathcal C_\gamma(s,t)
      &=s^{\mathsf T}\mathcal M_\gamma t
                           +\gamma^{-1}\Cov_\mu(\mathcal H_s,f_t),\\
 \|\mathcal C_\gamma\|_{2\to2}&\le W_H^2,\\
 \|\mathcal C_\gamma-K^0\|_{2\to2}
      &\le\tfrac12\varepsilon_\gamma+\frac{a_H}{2\gamma}W_H .
 \end{aligned}
 \tag{V60.18}\label{r60:bootstrap-eq}
\]
In particular all these inequalities hold on replacing \(a_H\) by
\(544\sqrt m\) in their monotone right-hand sides.
Along any sequence with \(\gamma/\sqrt m\to\infty\),
\[
              \|\mathcal C_\gamma-K^0\|_{2\to2}\longrightarrow0 .
 \tag{V60.19}\label{r60:regime}
\]
Since \(m=3N^3\), the paid sufficient regime is
\(\gamma\gg N^{3/2}\), improving V57's
\(\gamma\gg N^{3/2}(1+\log N)\).
\end{theorem}
\begin{proof}
For a linear energy polynomial,
\(\mathcal Sf_s=\frac12(f_s-\mathfrak Gf_s)\) and
\(\mathcal Tf_s=\frac12\mathcal H_s\).
Insert this into \eqref{r60:connected-eq} to get the first identity.
It also proves that the bilinear form
\(\Cov_\mu(\mathcal H_s,f_t)\) is symmetric; symmetry need not be assumed.

Let \(x=\sqrt{\|\mathcal C_\gamma\|_{2\to2}}\), and choose a unit
maximal-variance source in this finite source space. By
\eqref{r60:metric-eq} and Cauchy--Schwarz,
\[
                  2x^2\le4r+(a_H/\gamma)x .
\]
The positive root of this quadratic is \(W_H\).
For arbitrary unit sources the covariance return has magnitude at most
\(a_Hx\). Subtract \(2K^0\) from the exact identity and use the
diffusion comparison to obtain the last bound.
Its right side is monotone in \(a_H\), so \eqref{r60:H-bound} supplies
the stated fully explicit bound. Under the displayed scaling,
\(a_H/\gamma\le544\sqrt m/\gamma\to0\), while
\(\varepsilon_\gamma\to0\) uniformly in volume.
\end{proof}

The proof uses the Hodge projection before taking norms. It does not
count V56's entropy-to-quadratic comparison or V59's Gaussian inverse
again as new results. The newly paid parts are the compact rotational
diffusion comparison, the exact rooted quartic drift, and its log-free
finite-volume bound.

The remaining upstream target is now particularly concrete:
control the \emph{native centered} bank
\(\mathcal H_s-\E\mathcal H_s\), or its connected pairing with \(f_t\),
without the \(\sqrt m\) loss in \eqref{r60:field-L4}.
The full cycle projector and both quartic terms of
\eqref{r60:quartic-drift} must be retained. The proof has not bounded
this bank uniformly in volume; the star-tree example is a warning
against switching back to separately bounded Gaussian mean shifts.

All estimates still concern the coherent auxiliary compact law.
General noncommuting exteriors, their original probabilities and
predictor covariances, physical source coverage, physical-time
contraction and the interacting four-dimensional continuum construction
remain unresolved. The full Yang--Mills mass gap is unproved.

The diagnostic checks the Hodge and diffusion identities exactly on
rational compact fixtures with every root, including zero magnetization;
checks formal product rules and Gaussian pairings; and separately tests
the full connected identities by one-bond numerical quadrature at two
resolutions. Exact arithmetic, floating quadrature and analytic uniform
proofs are distinct forms of evidence.
% END CONSOLIDATED V60 APPEND


% BEGIN CONSOLIDATED V61 APPEND
\clearpage
\section{V61: Exact current Ward decomposition and connected cycle returns}
\label{r61:section}

\subsection{Context, the upstream adjustment, and retained scope}

The full interacting four-dimensional Yang--Mills continuum construction
and physical Hamiltonian mass gap remain unproved. This iteration attacks
one concrete remaining estimate in the coherent auxiliary compact law.
V60 retained native centering and the entire graph cycle space, but its
diffusion comparison still used the entropy comparison of V53/V56, and
its centered drift estimate lost \(\sqrt m\). We now use an exact
rotational Ward identity before any comparison with a Gaussian law.

There are two new conclusions. First, the actual spin-current covariance
splits exactly into a known gradient block and a positive cycle block.
Its Schur contraction improves V60's volume-uniform
\(O(\gamma^{-1/2})\) diffusion error to \(O(\gamma^{-1})\), without the
entropy budget. Second, an energy insertion in the same Ward identity
eliminates the mixed gradient--cycle return from the \emph{connected}
drift formula, replacing it by an explicit response operator and a
cycle-energy covariance. This is an exact upstream reduction, not a
proof that the remaining covariance is volume-uniform.

The earlier ESM transfer lesson is used only as proof architecture:
retain every native return when changing representation. In particular
V60's \(k\)-term is not discarded pointwise. Its contribution to the
connected pairing is accounted for by the inserted Ward identity below.
The companion and broad-literature checkpoints completed at V60 remain
current; the next are V65 and V70. No new external theorem is needed
for the identities proved here.

Keep the even cubic torus \(N\ge4\), \(n=N^3\), \(m=3n\), and
\[
 \mu_\gamma(dz)=Z_\gamma^{-1}
       e^{-\gamma\sum_{e=(i,j)}(1-z_i\cdot z_j)}
                     \prod_x d\sigma_{S^3}(z_x),\qquad \gamma>0 .
\]
The incidence convention is \((B\phi)_e=\phi_i-\phi_j\).
Retain \(C=L^\dagger\), \(R=BCB^{\mathsf T}\),
\(P=I-R=P_{\rm cyc}\), \(c=C_{xx}\le13/16\), and
\(r=R_{ee}=(n-1)/(3n)\). The letter \(P\) in this appendix denotes
the \emph{full} cycle projector, including the torus harmonic modes.
Write
\[
 \begin{aligned}
 \sigma_e&=z_i\cdot z_j,&d_e&=\gamma(1-\sigma_e),&
 f_s&=\sum_e s_ed_e,\\
 j_e^{ab}&=z_i^az_j^b-z_i^bz_j^a,&
 w_e^{ab}&=\sqrt\gamma\,j_e^{ab},&
 \kappa_\gamma&=\E_\mu\sigma_e .
 \end{aligned}
 \tag{V61.1}\label{r61:data}
\]
Spatial symmetry makes \(\kappa_\gamma\) independent of the edge.
It is not V59's dipole coefficient \(\kappa_{IJ}\).
The total current Gram matrix is
\[
       \mathcal W_{\gamma,ef}=\E_\mu(w_e\cdot w_f),
       \qquad
       \mathcal M_\gamma=R\circ\mathcal W_\gamma ,
 \tag{V61.2}\label{r61:current-matrix}
\]
where the dot sums the six orthonormal bivector components \(a<b\),
and \(\circ\) denotes entrywise multiplication.

\subsection{The exact gradient block of the actual current covariance}

\begin{theorem}[Current Ward decomposition]
\label{r61:ward}
For every \(\gamma>0\),
\[
 \begin{aligned}
 \mathcal W_\gamma R=R\mathcal W_\gamma&=3\kappa_\gamma R,\\
 \mathcal W_\gamma&=3\kappa_\gamma R+\mathcal S_\gamma,\\
 \mathcal S_\gamma&=P\mathcal W_\gamma P
       =\E_\mu\big[(Pw)_e\cdot(Pw)_f\big]_{e,f}\succeq0 .
 \end{aligned}
 \tag{V61.3}\label{r61:ward-eq}
\]
Both mixed gradient--cycle blocks vanish. Moreover
\[
       0\le\kappa_\gamma\le1,\qquad
       1-\kappa_\gamma=\gamma^{-1}\E d_e\le2r/\gamma .
 \tag{V61.4}\label{r61:kappa}
\]
\end{theorem}

\begin{proof}
For a real site profile \(\phi\) and fixed \(a<b\), use the
divergence-free product-sphere rotation
\[
 \mathcal R_\phi^{ab}
       =\sum_x\phi_x(z_x^a\partial_{z_x^b}
                                -z_x^b\partial_{z_x^a}),
       \qquad g=B\phi .
\]
If \(S_0=\sum_e(1-\sigma_e)\) and
\(p_e^{ab}=z_i^az_j^a+z_i^bz_j^b\), direct differentiation gives
\[
 \mathcal R_\phi^{ab}S_0=-\sum_e g_ej_e^{ab},
 \qquad
 \mathcal R_\phi^{ab}j_f^{ab}=-g_fp_f^{ab}.
 \tag{V61.5}\label{r61:rotation}
\]
Integrating the derivative of \(j_f^{ab}e^{-\gamma S_0}\) against
product Haar measure, which has no boundary, yields
\[
       \gamma\,\E\left[j_f^{ab}\sum_e g_ej_e^{ab}\right]
                           =g_f\E p_f^{ab}.
 \tag{V61.6}\label{r61:component-ward}
\]
Each coordinate occurs in exactly three pairs \(a<b\), so
\(\sum_{a<b}p_f^{ab}=3\sigma_f\), pointwise.
Summing \eqref{r61:component-ward} over the six components proves
\(\mathcal W_\gamma g=3\kappa_\gamma g\).
The image of \(B\) equals the image of \(R\); symmetry of
\(\mathcal W_\gamma\) then gives both identities with \(R\).
Expanding \(I=R+P\) proves the decomposition and the vanishing mixed
blocks. Positivity follows directly from the Gram representation.

There is a nonzero gradient on this connected graph. Its quadratic
form under the positive matrix \(\mathcal W_\gamma\) implies
\(\kappa_\gamma\ge0\). The bound \(\sigma_e\le1\) gives the upper
bound. Finally V50's actual infrared second-moment bound, already
used in V60, gives
\(\E d_e=\frac12\E|\sqrt\gamma(z_j-z_i)|^2\le2r\).
No entropy estimate or Gaussian replacement was used.
\end{proof}

The factor three is thus fixed by the exact compact rotation algebra.
For completeness, full \(O(4)\) invariance also implies
\(\E w_e^{ab}=0\) and
\[
 \E(w_e^{ab}w_f^{cd})
       =\mathbf1_{\{(a,b)=(c,d)\}}\,J_{\gamma,ef},
 \qquad J_\gamma=\mathcal W_\gamma/6 .
 \tag{V61.7}\label{r61:color}
\]
Indeed a coordinate reflection kills a product with distinct component
pairs, and coordinate permutations equate the diagonal-component
expectations. Hence \(J_\gamma R=(\kappa_\gamma/2)R\).
This argument uses reflections as well as rotations; it does not
incorrectly assume irreducibility of the bivector representation
under \(SO(4)\), nor does it introduce four Gaussian colors.

\subsection{An entropy-free, volume-uniform diffusion error of order
\texorpdfstring{\(1/\gamma\)}{1/gamma}}

\begin{theorem}[Positive cycle remainder and sharp order]
\label{r61:metric}
Put \(h_\gamma=512cr/\gamma\) and
\(\mathcal E_\gamma=R\circ\mathcal S_\gamma\). Then
\[
 \begin{aligned}
 0\preceq\mathcal E_\gamma&\preceq h_\gamma I,&
 \mathcal E_\gamma\mathbf1&=0,\\
 \mathcal M_\gamma
       &=3\kappa_\gamma(R\circ R)+\mathcal E_\gamma,&
 \mathcal M_\gamma\mathbf1&=3\kappa_\gamma r\mathbf1,\\
 \|\mathcal M_\gamma-2K^0\|_{2\to2}
       &\le h_\gamma,&2K^0&=3R\circ R .
 \end{aligned}
 \tag{V61.8}\label{r61:metric-eq}
\]
All statements hold uniformly in the torus size. In particular
\[
       0\preceq\mathcal M_\gamma\preceq b_\gamma I,
       \qquad b_\gamma=\min\{4r,\,3r+h_\gamma\}.
 \tag{V61.9}\label{r61:metric-size}
\]
\end{theorem}

\begin{proof}
Fix a root \(o\), without conditioning on its spin. Use V60's exact
fields
\[
       u_x=\sqrt\gamma(z_x-z_o),\qquad
       q_e=u_i\wedge u_j=u_i\wedge(u_j-u_i).
\]
The pointwise identity \(Pw=\gamma^{-1/2}Pq\) gives
\[
 \Tr\mathcal S_\gamma
       =\gamma^{-1}\E\|Pq\|_2^2
       \le\gamma^{-1}\sum_e\E|q_e|^2
       \le mh_\gamma .
 \tag{V61.10}\label{r61:cycle-trace}
\]
For the last inequality, the inherited actual \(L^4\) moments give
\[
 \E|q_e|^2
   \le\|u_i\|_{L^4}^2\|u_j-u_i\|_{L^4}^2
   \le(8cA_2)(2rA_2)=512cr,\qquad A_2^2=32 .
\]
The current matrix and \(P\) are equivariant under spatial translations
and cubic axis symmetries, with a sign change whenever an edge
orientation is reversed. These symmetries are transitive on unoriented
edges, so all diagonal entries of \(\mathcal S_\gamma\) are equal.
Thus each is at most \(h_\gamma\). The fixed-root representation used
to prove the trace bound need not itself be translation invariant.

The Schur product theorem, applied to
\(0\preceq R\preceq I\) and \(\mathcal S_\gamma\succeq0\), now yields
\[
       0\preceq R\circ\mathcal S_\gamma
               \preceq I\circ\mathcal S_\gamma
               \preceq h_\gamma I .
\]
Also
\((\mathcal E_\gamma\mathbf1)_e=(R\mathcal S_\gamma)_{ee}=0\).
Since \(R^2=R\),
\(\sum_fR_{ef}^2=r\); hence \(R\circ R\) has constant row sum \(r\)
and operator norm \(r\). Equation \eqref{r61:ward-eq} proves the
decomposition of \(\mathcal M_\gamma\) and its constant-source identity.

In the sense of symmetric matrix order,
\[
 -3(1-\kappa_\gamma)(R\circ R)
       \preceq\mathcal M_\gamma-3R\circ R
       \preceq h_\gamma I .
 \tag{V61.11}\label{r61:one-sided}
\]
The norm is at most
\(\max\{3r(1-\kappa_\gamma),h_\gamma\}\), not the sum of those errors.
Now \(3r(1-\kappa_\gamma)\le6r^2/\gamma\).
Positivity of \(C\) and \(C_{xx}=c\) imply
\(r=(\delta_i-\delta_j)^{\mathsf T}C(\delta_i-\delta_j)\le4c\).
Consequently \(6r^2\le24cr\le512cr\), giving the stated \(h_\gamma\)
bound. The new upper bound \(3r+h_\gamma\) follows from the same
positive decomposition; combine it with V60's \(4r\) bound.
\end{proof}

This is a genuine strengthening of the \emph{paid} metric estimate.
The Gaussian coefficient \(K^0\) is retained, but no comparison of
laws with that Gaussian is needed in this proof. The constant-source
identity is exact at every coupling; the cycle correction annihilates
that source but need not vanish on nonconstant signed sources.
None of these matrices is a physical transfer operator.

\begin{corollary}[Improved paid connected-covariance error]
\label{r61:paid}
Keep V60's actual centered drift norm \(a_H\le544\sqrt m\), and set
\[
       U_H=\frac{a_H}{4\gamma}
             +\sqrt{\frac{b_\gamma}{2}
                          +\left(\frac{a_H}{4\gamma}\right)^2}.
\]
Then
\[
 \|\mathcal C_\gamma\|_{2\to2}\le U_H^2,\qquad
 \|\mathcal C_\gamma-K^0\|_{2\to2}
           \le\frac{h_\gamma}{2}+\frac{a_H}{2\gamma}U_H .
 \tag{V61.12}\label{r61:paid-eq}
\]
Both right sides remain valid on replacing \(a_H\) by \(544\sqrt m\).
\end{corollary}
\begin{proof}
V60's exact identity remains
\(2\mathcal C_\gamma=\mathcal M_\gamma+
 \gamma^{-1}[\Cov(\mathcal H_s,f_t)]_{s,t}\).
For \(x=\sqrt{\|\mathcal C_\gamma\|}\), its maximal-variance unit
source satisfies \(2x^2\le b_\gamma+(a_H/\gamma)x\).
The positive quadratic root is \(U_H\). Subtract \(2K^0\), apply
\eqref{r61:metric-eq}, and use Cauchy--Schwarz for the connected return.
Monotonicity in \(a_H\) proves the final assertion.
\end{proof}

The sufficient regime is still \(\gamma/\sqrt m\to\infty\).
The improved metric error does \emph{not} by itself remove the
\(\sqrt m\) in the centered drift estimate.

\subsection{A source insertion accounts exactly for the mixed cycle return}

For a real signed source \(t\), define the symmetric edge matrix
\[
 \mathcal Z_t
       =\big[\Cov_\mu(w_e\cdot w_f,f_t)\big]_{e,f},\qquad
 D_t=\operatorname{diag}(t_e).
\]
The column \(\mathcal C_\gamma t\) has entries \(\Cov(d_e,f_t)\).
All expectations and covariances in this subsection are under the
original \emph{untilted} compact law.

\begin{theorem}[Inserted current Ward identity]
\label{r61:inserted}
The exact matrix identity is
\[
 \mathcal Z_t R
       =\mathcal W_\gamma D_tR
                   -\frac3\gamma
                       \operatorname{diag}(\mathcal C_\gamma t)R .
 \tag{V61.13}\label{r61:inserted-eq}
\]
In particular
\[
 \begin{aligned}
 R\mathcal Z_tR
       &=3\kappa_\gamma RD_tR
          -\frac3\gamma R\operatorname{diag}(\mathcal C_\gamma t)R,\\
 P\mathcal Z_tR
       &=\mathcal S_\gamma D_tR
          -\frac3\gamma P\operatorname{diag}(\mathcal C_\gamma t)R,\\
 P\mathcal Z_tP
       &=\big[\Cov((Pw)_e\cdot(Pw)_f,f_t)\big]_{e,f}.
 \end{aligned}
 \tag{V61.14}\label{r61:inserted-blocks}
\]
Thus the two blocks having a gradient leg are explicit. The remaining
cycle--cycle covariance has not been bounded here.
\end{theorem}

\begin{proof}
In the same rotation direction as above,
\[
       \mathcal R_\phi^{ab}f_t
                    =-\sqrt\gamma\sum_e t_eg_ew_e^{ab}.
\]
Apply compact integration by parts to \(f_tw_f^{ab}\).
Using \eqref{r61:rotation} and dividing by \(\sqrt\gamma\) gives
\[
 \E\left[f_tw_f^{ab}\sum_e g_ew_e^{ab}\right]
       =\E\left[w_f^{ab}\sum_e t_eg_ew_e^{ab}\right]
                                  +g_f\E(f_tp_f^{ab}).
\]
Sum over \(a<b\) and subtract
\(\E f_t\,(\mathcal W_\gamma g)_f=3\kappa_\gamma g_f\E f_t\).
Since \(\Cov(\sigma_f,f_t)=-\gamma^{-1}\Cov(d_f,f_t)\), the result is
\[
       (\mathcal Z_tg)_f
           =(\mathcal W_\gamma D_tg)_f
                         -\frac3\gamma g_f(\mathcal C_\gamma t)_f .
\]
This holds for every \(g\) in the image of \(R\), proving
\eqref{r61:inserted-eq}. Left multiplication by \(R\) or \(P\),
followed by \eqref{r61:ward-eq}, proves the first two block identities.
The last follows by moving the deterministic projectors inside
the covariance. No source-tilted infrared estimate has been asserted.
\end{proof}

\subsection{A cycle-energy bank and an exact coercive response equation}

Define
\[
 \begin{aligned}
 \mathcal J_s
     &=\sum_es_e\{d_e^2-3rd_e+\gamma|(Pw)_e|^2\},\\
 \mathcal B_J(s,t)&=\Cov_\mu(\mathcal J_s,f_t),\\
 D&=rI-R\circ R=P\circ R,\qquad
 A_\gamma=2I+(3/\gamma)D .
 \end{aligned}
 \tag{V61.15}\label{r61:cycle-bank}
\]
Here \(A_\gamma\) is a finite edge-source matrix, not V59's
differential generator \(\mathcal A_\gamma\).
Both \(R\) and \(P\) are positive, so
\[
       0\preceq D\preceq rI,\qquad
       A_\gamma\succeq2I,\qquad
       \|A_\gamma^{-1}\|_{2\to2}\le\tfrac12 .
 \tag{V61.16}\label{r61:coercive}
\]
For the upper bound on \(D\), use \(R\circ R\succeq0\).
No native generator gap is invoked.

\begin{theorem}[Exact reduction of the connected drift bank]
\label{r61:cycle-response}
As bilinear forms on all signed edge sources,
\[
 \begin{aligned}
 [\Cov(\mathcal H_s,f_t)]_{s,t}
       &=\mathcal B_J+\gamma\mathcal E_\gamma
                                  -3D\mathcal C_\gamma,\\
 A_\gamma\mathcal C_\gamma
       &=3\kappa_\gamma(R\circ R)+2\mathcal E_\gamma
                                  +\gamma^{-1}\mathcal B_J .
 \end{aligned}
 \tag{V61.17}\label{r61:cycle-response-eq}
\]
The term \(-3D\mathcal C_\gamma\) and the extra
\(\mathcal E_\gamma\) are part of the identity, not optional errors.
\end{theorem}

\begin{proof}
Decompose \(w=Rw+Pw\) in V60's exact current drift:
\[
 \mathcal H_s-\mathcal J_s
          =\gamma\sum_es_e(Rw)_e\cdot(Pw)_e .
\]
For a signed source \(s\), let \(D_s=\operatorname{diag}(s)\).
The connected pairing with \(f_t\), divided by \(\gamma\), is
\(\Tr(D_sP\mathcal Z_tR)\). Insert the second line of
\eqref{r61:inserted-blocks}. Since all matrices in the entrywise
products are symmetric,
\[
 \begin{aligned}
 \Tr(D_s\mathcal S_\gamma D_tR)
              &=s^{\mathsf T}\mathcal E_\gamma t,\\
 \Tr\!\left(D_sP\operatorname{diag}(\mathcal C_\gamma t)R\right)
              &=s^{\mathsf T}(P\circ R)\mathcal C_\gamma t
                =s^{\mathsf T}D\mathcal C_\gamma t .
 \end{aligned}
\]
This proves the first identity. Substitute it into V60.18 and
use \(\mathcal M_\gamma=3\kappa_\gamma(R\circ R)+\mathcal E_\gamma\)
to prove the second.
\end{proof}

This reduction eliminates the mixed term only \emph{after} taking its
native connected pairing. The rooted quartic identity V60.12, including
its \(k\)-term, remains intact. Indeed
\[
       \gamma|(Pw)_e|^2=|(Pq)_e|^2
\]
for every fixed root. Thus the new bank consists of the local scalar
energy polynomial and the square of the full projected bivector current.
It is a different observable from \(\mathcal H_s\) for nonconstant \(s\).
For \(s=\mathbf1\), \(\mathcal J_{\mathbf1}=\mathcal H_{\mathbf1}\)
pointwise, consistently with
\(\mathcal E_\gamma\mathbf1=D\mathbf1=0\).

One must not assume that \(D\) and \(\mathcal C_\gamma\) commute.
The first line of \eqref{r61:cycle-response-eq} and V60's symmetry
of the old connected return give the useful audit identity
\[
       \mathcal B_J-\mathcal B_J^{\mathsf T}
             =3(D\mathcal C_\gamma-\mathcal C_\gamma D).
 \tag{V61.18}\label{r61:skew-audit}
\]
In particular \(\mathcal B_J\) should not silently be replaced by its
symmetric part.

\begin{corollary}[A sufficient remaining connected estimate]
\label{r61:criterion}
At every finite volume let \(L_J=\|\mathcal B_J\|_{2\to2}\).
Then
\[
 \|\mathcal C_\gamma-K^0\|_{2\to2}
       \le\frac1\gamma
            \left(512cr+\frac{57}{16}r^2+\frac{L_J}{2}\right).
 \tag{V61.19}\label{r61:criterion-eq}
\]
Consequently a bound on \(L_J\) uniform in volume and all sufficiently
large \(\gamma\) would prove a volume-uniform \(O(\gamma^{-1})\)
coherent covariance comparison. Such a bound is \emph{not} established
in this appendix.
\end{corollary}

\begin{proof}
Subtract \(A_\gamma K^0\) from \eqref{r61:cycle-response-eq}:
\[
 A_\gamma(\mathcal C_\gamma-K^0)
       =3(\kappa_\gamma-1)(R\circ R)+2\mathcal E_\gamma
               -(3/\gamma)DK^0+\gamma^{-1}\mathcal B_J .
\]
If \(T=R\circ R\), then \(0\preceq T\preceq rI\),
\(D=rI-T\), and \(K^0=3T/2\). Functional calculus gives
\[
       \|DK^0\|
          \le\tfrac32\max_{0\le\lambda\le r}\lambda(r-\lambda)
          =3r^2/8 .
\]
Apply \eqref{r61:coercive}, \(\|\mathcal E_\gamma\|\le h_\gamma\),
and \eqref{r61:kappa}. The resulting bound is
\[
       h_\gamma+\frac{3r(1-\kappa_\gamma)}2
                      +\frac{9r^2}{16\gamma}+\frac{L_J}{2\gamma}
       \le\frac{512cr+57r^2/16+L_J/2}{\gamma}.
\]
Only the displayed sufficient hypothesis would make the last constant
uniform; finiteness at a fixed graph is automatic from compactness.
\end{proof}

\subsection{What remains unpaid and what has been independently checked}

The new metric theorem is unconditional within the inherited coherent
model. The cycle response equation is also exact, but its final
connected bank is open. The trace/diagonal estimate
\eqref{r61:cycle-trace} controls an uninserted second moment of the
current. It is not a bound on
\(\Cov(\gamma|(Pw)_e|^2,f_t)\), and differentiating the unperturbed
bound with respect to arbitrary signed couplings is not justified.
Such a differentiation would require new uniform control of the
source-perturbed law. Likewise, marginal fourth moments do not by
themselves bound the operator norm of all connected cross-edge
correlations.

The next upstream target is therefore the \emph{combined} connected
bank
\[
       \Cov_\mu\left(
            \sum_es_e\{d_e^2-3rd_e+|(Pq)_e|^2\},\,f_t\right).
\]
One may seek a bound directly, or use a centered \(L^2\) estimate for
this combined bank followed by the coercive response equation.
Splitting its terms is allowed only if the resulting estimates retain
the needed volume uniformity. No independence of projected currents,
Gaussian Wick domination, or physical mass has been inferred from the
positive matrix decomposition.

The independent diagnostic uses exact rational rotations on a
four-site complete graph and exact product-Haar moments on a triangle.
For the latter, fix the first spin by invariance and write
\[
 \sigma_{01}=x,\quad \sigma_{02}=y,\quad
 \sigma_{12}=xy+\sqrt{1-x^2}\sqrt{1-y^2}\,u .
\]
Here \(x,y\) are independent first coordinates of uniform \(S^3\)
spins and \(u\) is independent uniform on \([-1,1]\).
The exact moments
\[
       \E x^{2a}=\frac{(2a)!}{4^a a!(a+1)!},
       \qquad \E x^{2a+1}=0,\qquad
       \E u^{2a}=\frac1{2a+1}
\]
allow finite Taylor coefficients of the normalized compact law to be
checked with fractions, without a Gaussian approximation. The script
checks the current decomposition through degree six and the inserted
and reduced response identities through degree five. Separate
three-dimensional deterministic quadrature checks the full finite
couplings \(0.2,1,3,10\) at orders \(24,40,64\). It verifies the old
and new responses, their difference, and signed-source insertions.
These triangle and complete-graph fixtures check algebra; they do not
test or replace the cubic infrared uniform estimates. The numerical
integrals are floating-point diagnostics, not exact proofs.

General noncommuting exteriors, their original probabilities and
predictor covariances, physical source coverage, physical-time
contraction, and the interacting four-dimensional continuum limit
remain unresolved. The full Yang--Mills mass gap is unproved.
% END CONSOLIDATED V61 APPEND


% BEGIN CONSOLIDATED V62 APPEND
\clearpage
\section{V62: Finite current sources and compensated loop returns}
\label{r62:section}

\subsection{The upstream target and the nonduplication audit}

V61 left one explicit actual-law obligation:
\(\mathcal B_J(s,t)=\Cov_\mu(\mathcal J_s,f_t)\), with
\[
       \mathcal J_s
          =\sum_es_e\{d_e^2-3rd_e+\gamma|(Pw)_e|^2\}.
\]
This appendix transfers that entire signed-source covariance to a
\emph{compensated observable in the positive, normalized wire law},
up to an operator bounded by \(320\), uniformly in the cubic volume
and every \(\gamma>0\). The compensation includes a projected
edge-count subtraction and the factorial contact of the local energy
square. The remaining wire covariance is not proved bounded.
The full physical Yang--Mills mass gap is still unproved.

There is also a new finite-source Laplace estimate for all six spin
currents, valid even for arbitrary positive scalar edge couplings.
It does not supply quadratic-current concentration by itself.
An explicit radial-mixture example below demonstrates that logical
distinction, without being a counterexample to the compact spin law.

The nonduplication check matters. Archive \(C\), V61--V62, already
proved orthogonal-transport alignment domination and the marked-strand
twist Hessian; the latter is its equation
\path{f16v62:exact-Q}. The present restart's V42--V45 already
retains the Haar-pairing law and marked-count differentiation.
Those are inputs, not new results. The standard spin/wire
representation was checked again against Benassi--Ueltschi,
\href{https://arxiv.org/html/1807.06564v4}
{\emph{Loop correlations in random wire models}}, Section 6,
Proposition 6.1. Its special \(O(2)\) long-loop theorems are not
imported into the \(O(4)\) model. No broad review is claimed at V62;
the next companion and broad-literature checkpoints remain V65 and V70.

The new work is the nonorthogonal matrix extension needed for finite
current sources, the exact source derivative of the old twist Hessian,
and its compensated, volume-uniform-contact transfer to V61's
\(\mathcal B_J\). Large loop lengths from V45 are not estimated anew.
The loop observable here is a \emph{net oriented circulation}, for
which forward and backward traversals cancel.

\subsection{Finite sources for the actual spin currents}

For this subsection let \(G\) be any finite loopless graph, give each
edge \(e=(i,j)\) an orientation, and let \(\lambda_e>0\).
Use the scalar ferromagnetic law
\[
       d\mu_\lambda
           =Z(\lambda)^{-1}
                    \exp\!\left(\sum_e\lambda_e z_i\cdot z_j\right)
                       \prod_xd\sigma_{S^3}(z_x).
\]
There are no one-spin fields and no imposed nontrivial exterior
transports in this law. Define
\(w_e^\lambda=\sqrt{\lambda_e}\,z_i\wedge z_j\), with its six
orthonormal components.

\begin{lemma}[The needed matrix extension of alignment domination]
\label{r62:matrix}
For arbitrary real \(4\)-by-\(4\) matrices \(M_e\), set
\[
       Z(M)=\int\exp\!\left(\sum_{e=(i,j)}z_i^{\mathsf T}M_ez_j\right)dH .
\]
Then
\[
       0<Z(M)\le Z((\|M_e\|_{\rm op}I)_e).
 \tag{V62.1}\label{r62:matrix-eq}
\]
A reverse traversal uses \(M_e^{\mathsf T}\), not \(M_e^{-1}\).
\end{lemma}
\begin{proof}
Use the inherited sphere pairing tensor
\(u(2k)=1/\{2^k(k+1)!\}\). After expanding the bond exponentials,
every pairing component is a closed index line. Its contraction is
the trace of the ordered product of the encountered \(M_e\)'s or
their transposes, and therefore has absolute value at most
\(4\prod_{\text{occurrences }e}\|M_e\|_{\rm op}\).
All factorial and vertex-pairing coefficients are nonnegative.
The product of these bounds is the corresponding diagram weight
with scalar matrices \(\|M_e\|_{\rm op}I\). Summing proves the
inequality. Absolute convergence follows either from that finite
dominating compact integral or from the absolute component expansion
of the original finite-dimensional exponential.
Strict positivity follows from the real exponential integrand.
This extends, rather than re-proves as a new discovery, the archive's
orthogonal-matrix comparison.
\end{proof}

For \(h\in\Lambda^2\mathbb R^4\), write \(A(h)\) for the skew matrix
with entries \(A_{ab}=h^{ab}\), \(A_{ba}=-h^{ab}\), \(a<b\).
Thus \(z_i^{\mathsf T}A(h)z_j=h\cdot(z_i\wedge z_j)\).

\begin{theorem}[A volume- and coupling-uniform current Laplace bound]
\label{r62:mgf}
For every real six-component edge source \(h=(h_e)_e\),
\[
 \begin{aligned}
 \log\E_{\mu_\lambda}\exp\!\left(\sum_eh_e\cdot w_e^\lambda\right)
 &\le\sum_e
    \left\{\sqrt{\lambda_e^2+
                  \lambda_e\|A(h_e)\|_{\rm op}^2}-\lambda_e\right\}\\
 &\le\tfrac12\sum_e|h_e|^2 .
 \end{aligned}
 \tag{V62.2}\label{r62:mgf-eq}
\]
All sources and all positive scalar couplings are allowed; no
small-source or low-temperature restriction enters this estimate.
\end{theorem}
\begin{proof}
The numerator is exactly \(Z(M)\) with
\(M_e=\lambda_e I+\sqrt{\lambda_e}A(h_e)\).
Since \(A\) is real skew,
\[
       M_e^{\mathsf T}M_e
          =\lambda_e^2 I-\lambda_e A(h_e)^2,\qquad
       \rho_e:=\|M_e\|
          =\sqrt{\lambda_e^2+
                         \lambda_e\|A(h_e)\|^2}\ge\lambda_e .
\]
By \eqref{r62:matrix-eq}, the Laplace transform is at most
\(Z(\rho)/Z(\lambda)\). Pointwise \(z_i\cdot z_j\le1\), so
\[
       \frac{Z(\rho)}{Z(\lambda)}
           =\E_{\mu_\lambda}
                    e^{\sum_e(\rho_e-\lambda_e)z_i\cdot z_j}
           \le e^{\sum_e(\rho_e-\lambda_e)} .
\]
The real skew normal form consists of two \(2\)-plane blocks with
speeds \(a,b\); hence
\(\|A(h)\|^2=\max(a^2,b^2)\le a^2+b^2=|h|^2\).
Finally
\(\sqrt{\lambda^2+\lambda q}-\lambda\le q/2\) for \(q\ge0\).
These facts prove both inequalities.
\end{proof}

The currents have mean zero by full \(O(4)\) invariance. For a fixed
source \(h\), Chernoff's argument applied to both signs gives
\[
 \Pr\left\{\left|\sum_eh_e\cdot w_e^\lambda\right|\ge u\right\}
          \le2\exp\{-u^2/(2\|h\|_2^2)\},\qquad u>0 ,
 \tag{V62.3}\label{r62:tail}
\]
with the zero-source case understood separately. This is a bound
for linear current statistics, not a Wick theorem for quadratic
statistics. The denominator in the proof is the original scalar
law's normalizer. Replacing it by the normalizer of a frustrated
connection is not an authorized extension of the conclusion.

\subsection{An inherited Hessian consequence in the new source norm}

Let \(\kappa_e(\lambda)=\E_{\mu_\lambda}(z_i\cdot z_j)\).
The old nonnegative twist Hessian gives the useful all-source bound
\[
 \Var_{\mu_\lambda}\!\left(\sum_eh_e\cdot w_e^\lambda\right)
                       \le\tfrac12\sum_e\kappa_e(\lambda)|h_e|^2 .
 \tag{V62.4}\label{r62:cap}
\]
Here is the precise specialization, so that it is not confused with
differentiating \eqref{r62:mgf-eq} twice at a nonzero source.
Use transports \(T_e(\epsilon)=
\exp\{\epsilon A(h_e)/\sqrt{\lambda_e}\}\).
Alignment domination implies that
\(-\log[Z(\lambda,T(\epsilon))/Z(\lambda,I)]\) has a minimum at zero.
Its second derivative is
\[
       -\sum_e\E(z_i^{\mathsf T}A(h_e)^2z_j)
                  -\Var\!\left(\sum_eh_e\cdot w_e^\lambda\right).
\]
Scalar couplings preserve \(O(4)\), so
\(\E z_i^az_j^b=(\kappa_e/4)\delta_{ab}\), and the first term is
\(\frac12\sum_e\kappa_e|h_e|^2\). This proves \eqref{r62:cap}.
The inherited count identity
\(\E_{\rm w}K_e=\lambda_e\kappa_e\ge0\) also shows
\(0\le\kappa_e\le1\).
The Hessian input is old; its consequence for V61's complete
current-source matrix is what is being used here.

Return now to V61's homogeneous cubic law, with all notation
\((R,P,r,c,\kappa_\gamma,\mathcal W_\gamma,\mathcal S_\gamma)\)
unchanged. Color isotropy and \eqref{r62:cap} give
\[
 \begin{aligned}
 0\preceq\mathcal W_\gamma&\preceq3\kappa_\gamma I,&
 0\preceq\mathcal S_\gamma&\preceq3\kappa_\gamma P,\\
 0\preceq\mathcal M_\gamma&\preceq3\kappa_\gamma rI,&
 \|\mathcal M_\gamma\|_{2\to2}&=3\kappa_\gamma r .
 \end{aligned}
 \tag{V62.5}\label{r62:sharp-metric}
\]
The last equality uses V61's exact constant-source eigenvector.
Indeed Schur multiplication of the first matrix bound by \(R\)
gives the upper bound, and
\(\mathcal M_\gamma\mathbf1=3\kappa_\gamma r\mathbf1\) attains it.
V61's error estimate \(h_\gamma=512cr/\gamma\) remains valid.
Its covariance bootstrap may use
\(b_\gamma=3\kappa_\gamma r\le3r\) instead of V61.9.
The proved sufficient scaling regime is still
\(\gamma/\sqrt m\to\infty\); no volume-uniform centered drift norm
has followed.

\subsection{Net loop circulations and the complete energy insertion}

Use the inherited positive wire law at positive couplings \(\lambda\).
A diagram \(\omega\) has labelled bond occurrences \(K_e\), even
degrees \(N_x=\sum_{e\ni x}K_e\), and a pairing at each vertex.
Its probability is proportional to
\[
       \prod_e\frac{\lambda_e^{K_e}}{K_e!}
       \prod_xu(N_x)\,4^{\#\{\text{closed components}\}} .
 \tag{V62.6}\label{r62:wire-law}
\]
The normalizer is exactly \(Z(\lambda)\); no reference probability
has been substituted. Denote wire expectation by \(\E_{\rm w}\).
Orient each closed component \(C\) arbitrarily. Let \(N_e^+(C)\)
and \(N_e^-(C)\) count its forward and backward traversals of \(e\).
Set
\[
 \begin{aligned}
 \ell_e(C)&=N_e^+(C)-N_e^-(C),\\
 \Lambda(\omega)&=\sum_C\ell(C)\ell(C)^{\mathsf T},\\
 \overline\Lambda&=\E_{\rm w}\Lambda,\qquad
 M_t=\sum_et_eK_e .
 \end{aligned}
 \tag{V62.7}\label{r62:loop-matrix}
\]
Reversing a component's orientation leaves \(\Lambda\) unchanged.
Each \(\ell(C)\) is a graph cycle, so
\(P\Lambda P=\Lambda\), including all winding components of the torus.

For a single coordinate-plane generator \(G\) with
\(\|G\|_{\rm F}^2=2\), the inherited twist Hessian with
\(A_e=a_eG/\sqrt\gamma\) reads
\[
       \mathcal W_\gamma
                   =3\kappa_\gamma I-\frac3\gamma\overline\Lambda .
 \tag{V62.8}\label{r62:loop-current}
\]
To verify the factors, each closed component contributes
\((2\gamma)^{-1}(\ell(C)\cdot a)^2\) to that one-plane Hessian,
whereas the spin expression is
\((\kappa_\gamma/2)\|a\|^2
-\Var(\sum_ea_ew_e^{ab})\).
Color isotropy multiplies the current covariance by six.
Equation \eqref{r62:loop-current} is a specialization of the old
Hessian, not a newly discovered representation.

\begin{theorem}[Normalized energy derivative of the loop Hessian]
\label{r62:loop-inserted}
For every signed source \(t\), the full edge matrix of V61 satisfies
\[
 \begin{aligned}
 \mathcal Z_t={}&D_t\mathcal W_\gamma+\mathcal W_\gamma D_t
          -3\kappa_\gamma D_t
          -\frac3\gamma\operatorname{diag}(\mathcal C_\gamma t)\\
          &+\frac3\gamma
                      \Cov_{\rm w}(\Lambda,M_t).
 \end{aligned}
 \tag{V62.9}\label{r62:loop-inserted-eq}
\]
The matrix covariance on the last line is entrywise.
All five terms are required.
\end{theorem}

\begin{proof}
Keep the observable scale \(\gamma\) and the projectors fixed.
The spin tilt \(e^{\tau f_t}\) changes the scalar couplings to
\(\lambda_e(\tau)=\gamma b_e(\tau)\), with \(b_e=1-\tau t_e>0\)
for \(\tau\) in a neighborhood of zero; its scalar prefactor
cancels on normalization.
Let \(\mathcal W(\tau)\) be the Gram matrix of the
\emph{fixed-scale} currents \(\sqrt\gamma\,z_i\wedge z_j\)
under this tilted law, and let \(\kappa_e(\tau)=\E_\tau\sigma_e\).

Use the same test twists
\(\exp(\epsilon a_eG/\sqrt\gamma)\), not twists rescaled by
\(\lambda_e(\tau)\). The spin and wire Hessians give exactly
\[
 D_b\mathcal W(\tau)D_b
       =3\operatorname{diag}(b_e\kappa_e(\tau))
                             -\frac3\gamma\overline\Lambda(\tau).
 \tag{V62.10}\label{r62:tilted-loop}
\]
Now
\[
 \mathcal W'(0)=\mathcal Z_t,\qquad
 \kappa_e'(0)=-\gamma^{-1}(\mathcal C_\gamma t)_e,\qquad
 \overline\Lambda'(0)=-\Cov_{\rm w}(\Lambda,M_t).
\]
The last identity differentiates the \emph{normalized} positive
weights: their logarithmic derivative at zero is \(-M_t\).
Differentiate \eqref{r62:tilted-loop} and move the two
\(D_b'(0)=-D_t\) terms to the right to obtain the claim.

Every differentiation is legitimate at each finite graph.
If \(M=\sum_eK_e\), then
\(\E_{\rm w}e^{uM}=Z(e^u\lambda)/Z(\lambda)<\infty\).
The entries of \(\Lambda\) are bounded in magnitude by \(M^2\).
A small coupling neighborhood and this exponential moment dominate
all the marked derivatives used above.
This justification is not a uniform higher-cumulant bound.
\end{proof}

Right multiplication by \(R\) annihilates the wire covariance, since
\(\Lambda R=0\). The remaining terms reduce exactly to V61.13.
Thus V61's gradient-leg Ward identity is respected; the additional
information in \eqref{r62:loop-inserted-eq} is its cycle--cycle
content.

\subsection{The projected count compensation and its bounded contacts}

Let
\[
       Q=P\circ P,\qquad
       \Lambda_s=\Tr(D_s\Lambda),\qquad
       \mathcal V_s^{\rm w}=\Lambda_s-M_{Qs}.
 \tag{V62.11}\label{r62:projected-compensation}
\]
The matrix \(Q\) is not the projector \(P\), and not V56's
diffusion matrix \(\mathcal Q\).
The exact count-covariance identity at homogeneous coupling is
\[
       \Gamma_{\rm w}:=
           [\Cov_{\rm w}(K_e,K_f)]_{e,f}
                 =\mathcal C_\gamma+\gamma\kappa_\gamma I .
 \tag{V62.12}\label{r62:count-cov}
\]
For example differentiate
\(\E_{\rm w}K_e=\lambda_e\E_\lambda\sigma_e\)
along \(\lambda_e=\gamma(1-\tau t_e)\).
The normalized wire derivative is
\(-(\Gamma_{\rm w}t)_e\); the spin derivative is
\(-\gamma\kappa_\gamma t_e-(\mathcal C_\gamma t)_e\).
This proves \eqref{r62:count-cov}, including its diagonal factorial
contact.

\begin{proposition}[An exact compensated cycle-energy insertion]
\label{r62:cycle-compensation}
Let \(Y_s=\gamma\sum_es_e|(Pw)_e|^2\). Then
\[
       \Cov_\mu(Y_s,f_t)
          =2\gamma\,s^{\mathsf T}(P\circ\mathcal S_\gamma)t
                         +3\Cov_{\rm w}(\mathcal V_s^{\rm w},M_t).
 \tag{V62.13}\label{r62:cycle-compensation-eq}
\]
The first term is an unconditionally bounded source operator:
\[
       \|2\gamma(P\circ\mathcal S_\gamma)\|_{2\to2}
                                \le1024cr .
 \tag{V62.14}\label{r62:cycle-contact}
\]
\end{proposition}

\begin{proof}
The left side is \(\gamma\Tr(D_sP\mathcal Z_tP)\).
Insert \eqref{r62:loop-inserted-eq}, and use
\(P\mathcal W_\gamma=\mathcal W_\gamma P=\mathcal S_\gamma\).
The result is
\[
 \begin{aligned}
 2\gamma\,s^{\mathsf T}(P\circ\mathcal S_\gamma)t
 &-3\gamma\kappa_\gamma s^{\mathsf T}Qt
     -3s^{\mathsf T}Q\mathcal C_\gamma t\\
 &+3\Cov_{\rm w}(\Lambda_s,M_t).
 \end{aligned}
\]
The last three terms combine to the compensated covariance in
\eqref{r62:cycle-compensation-eq}, by \eqref{r62:count-cov}.
No contact of order \(\gamma\) has been bounded separately.
Finally \(0\preceq P\preceq I\) and
\(\mathcal S_\gamma\succeq0\) imply
\(0\preceq P\circ\mathcal S_\gamma
 \preceq I\circ\mathcal S_\gamma\preceq h_\gamma I\);
V61's \(h_\gamma=512cr/\gamma\) gives \eqref{r62:cycle-contact}.
\end{proof}

To transfer the remaining local energy square as well, define the
single combined wire observable
\[
 \mathcal U_s^{\rm w}
    =3\Lambda_s
      -\sum_es_e\{K_e(K_e-1)+(3r-2\gamma)K_e\}
      -3M_{Qs},\qquad
 \mathcal B_{\rm w}(s,t)=\Cov_{\rm w}(\mathcal U_s^{\rm w},M_t).
 \tag{V62.15}\label{r62:full-compensation}
\]
It is a signed observable under a positive measure. Its three pieces
are not claimed separately small. A source-dependent constant would
not affect this covariance and has been omitted.

\begin{theorem}[The full native return transfers with a uniform bounded error]
\label{r62:transfer}
Put \(\mu_{1,\gamma}=\E_\mu d_e\), \(\mu_{2,\gamma}=\E_\mu d_e^2\), and
\[
       \zeta_\gamma
             =2\mu_{2,\gamma}-2\gamma\mu_{1,\gamma}
                                  +3r\gamma\kappa_\gamma .
\]
Then, exactly,
\[
       \mathcal B_J
           =\mathcal B_{\rm w}
                    +2\gamma(P\circ\mathcal S_\gamma)+\zeta_\gamma I ,
 \tag{V62.16}\label{r62:transfer-eq}
\]
and, for every even \(N\ge4\) and \(\gamma>0\),
\[
       |\zeta_\gamma|\le382/9,\qquad
       \|\mathcal B_J-\mathcal B_{\rm w}\|_{2\to2}
                    \le1024cr+382/9\le2878/9<320 .
 \tag{V62.17}\label{r62:transfer-bound}
\]
Thus volume- and coupling-uniform boundedness of the complete
\(\mathcal B_J\) and of the complete \(\mathcal B_{\rm w}\) are
equivalent, up to this absolute additive constant.
Neither boundedness assertion is proved here.
\end{theorem}

\begin{proof}
At a single edge write \(K=K_e\).
The inherited factorial moment identities are
\(\E_{\rm w}K=\gamma\E\sigma_e\) and
\(\E_{\rm w}K(K-1)=\gamma^2\E\sigma_e^2\).
For the tilted couplings in the previous proof,
\[
 \E_\tau d_e^2
       =\gamma^2
        -2\gamma^2\frac{\E_{{\rm w},\tau}K}{\lambda_e(\tau)}
        +\gamma^2
                  \frac{\E_{{\rm w},\tau}K(K-1)}{\lambda_e(\tau)^2}.
 \tag{V62.18}\label{r62:local-factorial}
\]
The observable \(d_e=\gamma(1-\sigma_e)\) still uses the fixed
base scale, not \(\lambda_e(\tau)\).
Differentiate this equality at zero. Normalizer differentiation
contributes minus covariance with \(M_t\); differentiating the
two explicit denominators contributes the diagonal contact.
With \(A_e=K_e(K_e-1)-2\gamma K_e\), this gives
\[
       \Cov_\mu(d_e^2,f_t)
           =-\Cov_{\rm w}(A_e,M_t)
                      +2t_e(\mu_{2,\gamma}-\gamma\mu_{1,\gamma}).
 \tag{V62.19}\label{r62:local-contact}
\]
Subtract \(3r\Cov_\mu(d_e,f_t)\), and use
\eqref{r62:count-cov} to rewrite this subtraction as
\(-3r\Cov_{\rm w}(K_e,M_t)+3r\gamma\kappa_\gamma t_e\).
Sum against \(s_e\) and add \eqref{r62:cycle-compensation-eq}.
The combined wire observable is precisely
\eqref{r62:full-compensation}, proving \eqref{r62:transfer-eq}.

For the numerical bound, only already paid native estimates are used:
\[
       0\le\mu_{1,\gamma}\le2r,\qquad
       0\le\mu_{2,\gamma}\le32r^2,\qquad
       |\mu_{1,\gamma}-3r/2|\le\frac{52}{3\gamma}.
\]
The first two are V50/V51 infrared moments; the last is the
edge-transitive consequence of \eqref{r51:mean-eq}.
Since \(\kappa_\gamma=1-\mu_{1,\gamma}/\gamma\),
\[
 \begin{aligned}
 |\zeta_\gamma|
   &=|2\mu_{2,\gamma}-3r\mu_{1,\gamma}
                              +\gamma(3r-2\mu_{1,\gamma})|\\
   &\le70r^2+104/3\le382/9 .
 \end{aligned}
\]
Now use \eqref{r62:cycle-contact}, \(r\le1/3\), \(c\le13/16\):
\[
       1024cr+382/9\le832/3+382/9=2878/9<320 .
\]
The two-sided norm comparison follows by the triangle inequality;
no smallness of \(\mathcal B_{\rm w}\) has been assumed in this proof.
\end{proof}

The correction in \eqref{r62:transfer-eq} is symmetric, so the
skew audit remains
\[
       \mathcal B_{\rm w}-\mathcal B_{\rm w}^{\mathsf T}
              =3(D\mathcal C_\gamma-\mathcal C_\gamma D),
       \qquad D=P\circ R .
 \tag{V62.20}\label{r62:skew}
\]
In particular, positivity of the wire weights does not make this
cross-covariance matrix positive or symmetric. The direct sufficient
criterion now reads
\[
 \|\mathcal C_\gamma-K^0\|_{2\to2}
   \le\frac1\gamma\left\{
       512cr+\frac{57}{16}r^2
                       +\frac{\|\mathcal B_{\rm w}\|_{2\to2}+320}{2}
                         \right\},
 \tag{V62.21}\label{r62:wire-criterion}
\]
by V61.19. This is a quantitative exact-law change of the remaining
target, not a conditional estimate advertised as a mass-gap proof.

\subsection{Why linear-current control does not complete the argument}

\begin{proposition}[A diagnostic radial mixture, not a spin-law counterexample]
\label{r62:mixture}
Uniform linear-current Laplace bounds, the gradient/cycle covariance
split, and an \(O(\gamma^{-1})\) cycle covariance do not alone imply
a volume-uniform variance for normalized cycle energy.
\end{proposition}
\begin{proof}
Use the same deterministic cubic projectors and fix \(\gamma\ge8\).
Set \(\kappa_*=1-3r/(2\gamma)\).
For \(\alpha=1,\ldots,6\), let \(g^\alpha,h^\alpha\) be independent
standard Gaussian vectors in \(\mathbb R^m\). Let \(\xi\), independent
of them, equal \(0\) or \(\sqrt2\), each with probability \(1/2\).
Define the abstract current field
\[
       \widetilde w^\alpha
              =\sqrt{\kappa_*/2}\,Rg^\alpha
                              +\frac{\xi}{\sqrt\gamma}Ph^\alpha .
\]
Conditional on \(\xi\), its covariance in each component is bounded
by \((\kappa_*/2)I\), since \(2/\gamma\le\kappa_*/2\).
Consequently every real six-component edge source obeys
\[
       \E e^{\sum_\alpha h_\alpha\cdot\widetilde w^\alpha}
                  \le e^{(\kappa_*/4)\sum_\alpha\|h_\alpha\|^2},
\]
which is stronger than \eqref{r62:mgf-eq}'s final bound.
Its total current matrix is
\[
       \widetilde{\mathcal W}
                =3\kappa_*R+(6/\gamma)P.
\]
It has the exact stated gradient block, vanishing mixed blocks,
translation-invariant cycle diagonal, and even a volume-uniform
cycle \emph{operator} norm \(6/\gamma\).

Nevertheless, for the unit source \(s=\mathbf1/\sqrt m\),
\[
       \widetilde Y_s
             =\frac{\gamma}{\sqrt m}
                         \sum_\alpha\|P\widetilde w^\alpha\|^2
             =\frac{\xi^2}{\sqrt m}X,\qquad
       X\sim\chi^2_\nu,\quad \nu=6(m-n+1),
\]
with \(X\) independent of \(\xi\).
Since \(\E\xi^2=1\), \(\E\xi^4=2\), and
\(\Var X=2\nu\), direct calculation gives
\[
                  \Var\widetilde Y_s
                          =\frac{\nu^2+4\nu}{m}.
 \tag{V62.22}\label{r62:mixture-variance}
\]
For \(m=3n\), this grows linearly in \(m\).
\end{proof}

These unbounded Gaussian-mixture currents are not currents of
compact unit spins. They are not claimed to satisfy the inserted
spin Ward hierarchy or the actual joint wire identities.
The example only rules out an inference from the listed linear and
second-moment estimates alone. It gives no counterexample to the
desired compact covariance estimate, much less to a physical mass gap.
The actual compensated wire structure in
\eqref{r62:full-compensation} is additional information, which the
next estimate must use rather than replace by generic sub-Gaussianity.

\subsection{Verification and next direction}

An independent diagnostic enumerates all \(49\,505\) labelled
pairing diagrams on a triangle with at most eight bond occurrences,
across \(45\) even-degree count profiles. It computes net
circulations by tracing each component; backtracking is canceled
before squaring. A separate product-Haar spin calculation supplies
exact rational Taylor coefficients. They agree on the partition
function, count covariance, loop-current Hessian, complete inserted
current identity, projected compensation and full wire transfer.
The mixed identities with one inverse coupling are tested only
through the coefficient order supported by the enumeration.
This is not an infinite-series numerical certification or a
thermodynamic argument.

Separate deterministic single-bond quadrature at two resolutions
checks the finite-source matrix-norm bound, including unequal skew
plane speeds. The radial-mixture variance and the constant
\(2878/9<320\) are checked exactly. Numerical diagnostics and
finite rational tests remain distinct from the analytic proofs.

The next target is the normalized connected covariance
\(\Cov_{\rm w}(\mathcal U_s^{\rm w},M_t)\), with all of its
compensations retained. A bound on raw loop lengths, raw
\(\Lambda_s\), or a separate raw count square could lose the
cancellations just proved; no such stronger requirement is imposed.
The finite-source current estimate is available under arbitrary
positive scalar couplings, but it is not by itself a control of
this mixed higher connected response.

Generic noncommuting exteriors, their original probabilities and
predictor covariances, physical source coverage, physical-time
contraction, and the interacting four-dimensional continuum
construction remain unresolved. The full Yang--Mills mass gap
remains unproved.
% END CONSOLIDATED V62 APPEND


% BEGIN CONSOLIDATED V63 APPEND
\clearpage
\section{V63: whole-loop Palm identities and compensated count noise}
\label{r63:section}

\subsection{Context, nonduplication correction and literature safety}
V62 moved the native cycle-return bank to a single connected
covariance under the actual normalized positive wire law. Its
observable is signed; positivity of that law does not make its
cross-covariance positive. We now go upstream of a proposed
concentration argument: even on one bond the centered wire
observable has variance of order \(\gamma^2\), while the connected
covariance of interest stays bounded. We prove this distinction
first. We then derive a whole-component insertion identity,
with all repeated-edge and repeated-vertex factors retained,
and an exact reversible Stein identity for the full bank.
No replica-switching lemma is used.

There is a nonduplication correction to the preceding appendix.
Lemma~\ref{r62:matrix}, presented there as a nonorthogonal extension,
is already contained in Lemma~\ref{r54:matrix-domination} and
\eqref{r54:matrix-eq}, after scalar rescaling. Thus V62's
finite-current Laplace estimate is an application of existing
matrix domination, not a new matrix-domination theorem. The
compensated inserted-current transfer
\eqref{r62:transfer-eq} is unaffected. We preserve the prior text
and record this correction explicitly rather than silently
rewriting its novelty claim.

A targeted primary-source check on 20 September 2026 also found
that Benjamin Lees's \emph{Griffiths inequalities for the
\(O(N)\)-spin model}, \href{https://arxiv.org/abs/2205.12928}
{arXiv:2205.12928}, was withdrawn in version 3 on 31 March 2026:
the author reports that the claimed bijection in the proof is
invalid. We therefore do not import that paper's switching
identity or claimed inequalities. This is not a proof that the
inequalities themselves are false. The current notes' earlier
sign cautions remain in force; the text audit found no dependency
on that withdrawn proof. Giuliani--Ott's
\href{https://arxiv.org/abs/2302.07299}
{low-temperature expansion paper} now has the journal reference
\emph{Probability and Mathematical Physics} 6 (2025), 439--477.
No additional theorem from it is needed below. The next scheduled
companion review is V65 and the next broad literature sweep V70.

\subsection{The exact one-bond count law}
Write \(\Phi(\gamma)=\E_0 e^{\gamma z\cdot z'}\), where \(z,z'\)
are independent uniform points of \(S^3\). The inherited Haar
moment formula gives
\begin{equation*}
 \Phi(\gamma)=\sum_{k\ge0}
 \frac{(\gamma^2/4)^k}{k!(k+1)!},\qquad
 \mathbb P_{\rm w}(K=2k)
 =\frac{(\gamma^2/4)^k}{k!(k+1)!\Phi(\gamma)}.
 \tag{V63.1}\label{r63:one-bond-law}
\end{equation*}
The count \(K\) here is the total edge multiplicity, not the
number of loop components. On a tree every closed walk has
zero net traversal through every edge. Consequently
\(R=I\), \(P=Q=0\), \(r=1\), and the one-bond version of
V62's full compensated observable is
\begin{equation*}
 U(K)=-K^2+(2\gamma-2)K.
 \tag{V63.2}\label{r63:one-bond-U}
\end{equation*}
This tree specialization does not substitute \(r=1\) into any
cubic-torus estimate requiring \(r\le1/3\).

\begin{proposition}[Exact count generator and moment identities]
\label{r63:count-generator}
The even-count law \eqref{r63:one-bond-law} is reversible for
\[
 L_{\rm b}F(K)=\gamma^2\{F(K+2)-F(K)\}
    +K(K+2)\{F(K-2)-F(K)\},
\]
where the death term vanishes at \(K=0\).
Put \(q=\E_{\rm w}K\). Then
\begin{equation*}
 \begin{aligned}
 \E_{\rm w}K^2&=\gamma^2-2q,&
 \E_{\rm w}K^3&=(\gamma^2+4)q,\\
 \E_{\rm w}K^4&=\gamma^4+4\gamma^2-8q,
 \end{aligned}
 \tag{V63.3}\label{r63:raw-moments}
\end{equation*}
and
\begin{equation*}
 \Cov_{\rm w}(U,K)=2\gamma\Var_{\rm w}K-2\gamma^2,
 \tag{V63.4}\label{r63:count-cov}
\end{equation*}
\begin{equation*}
 \Var_{\rm w}U=
 4\gamma^2\{\gamma^2-2\gamma+2-q-q^2\}.
 \tag{V63.5}\label{r63:count-var}
\end{equation*}
\end{proposition}
\begin{proof}
The ratio of successive probabilities is
\(\mathbb P(K+2)/\mathbb P(K)=
\gamma^2/[(K+2)(K+4)]\), which is precisely detailed balance.
The constant birth rate prevents explosion; deaths only reduce
the count. The factorial tails justify stationary identities for
all polynomials. Applying \(\E L_{\rm b}F=0\) successively to
\(F(K)=K,K^2,K^3\) gives \eqref{r63:raw-moments}.
In particular,
\[
 L_{\rm b}K=2\{\gamma^2-K(K+2)\},\qquad
 U=\tfrac12L_{\rm b}K+2\gamma K-\gamma^2.
\]
The expected reversible quadratic form on \(K\) is
\[
 -\E_{\rm w}K L_{\rm b}K
 =\tfrac12\E_{\rm w}
   \{4\gamma^2+4K(K+2)\}=4\gamma^2.
\]
This proves \eqref{r63:count-cov}. Expanding \(U^2\),
using \eqref{r63:raw-moments}, and subtracting
\((\E U)^2=(2\gamma q-\gamma^2)^2\), gives
\eqref{r63:count-var}.
\end{proof}

\begin{proposition}[Large count variance and a bounded connected bank]
\label{r63:one-bond-calibration}
For all \(\gamma>0\),
\[
 |\Cov_{\rm w}(U,K)|\le15.
\]
Nevertheless, as \(\gamma\to\infty\),
\begin{equation*}
 \gamma^{-2}\Var_{\rm w}U\longrightarrow2,\qquad
 \Cov_{\rm w}(U,K)\longrightarrow-\tfrac34.
 \tag{V63.6}\label{r63:count-limits}
\end{equation*}
Thus a coupling-uniform centered \(L^2\) bound on \(U\) is an
unnecessarily strong criterion for this connected covariance.
\end{proposition}
\begin{proof}
Under the compact one-bond spin law let
\(d=\gamma(1-z\cdot z')\) and \(\mu_j=\E d^j\).
Its density on \([0,2\gamma]\) is proportional to
\[
 e^{-d}d^{1/2}(2-d/\gamma)^{1/2}.
\]
Relative to a gamma random variable of shape \(3/2\) and rate
one, this is a nonnegative decreasing tilt
\(h_\gamma(d)=(2-d/\gamma)_+^{1/2}\).
For increasing \(f\), an independent-copy identity gives
\(\Cov(f(D),h_\gamma(D))\le0\). Hence
\[
 \mu_j\le(3/2)_j\quad(j\ge1),
 \qquad \mu_j\longrightarrow(3/2)_j
 \quad(\gamma\to\infty),
\]
where the limit follows by dominated convergence in the
normalized density. In particular
\(\mu_1\le3/2,\ \mu_2\le15/4,\ \mu_3\le105/8\).

For completeness, integration by parts of the one-bond compact
density gives the exact Ward identity
\begin{equation*}
 2\mu_1=3+\frac{\mu_2-3\mu_1}{\gamma}.
 \tag{V63.7}\label{r63:bond-ward}
\end{equation*}
Indeed the integral of the derivative of
\(e^{\gamma\sigma}(1-\sigma^2)^{3/2}\) over \([-1,1]\)
vanishes, so \(\gamma\E(1-\sigma^2)=3\E\sigma\).
This also shows
\[
 \gamma(\mu_1-\tfrac32)\longrightarrow-\tfrac38,
 \qquad
 q=\gamma-\mu_1
   =\gamma-\tfrac32+\frac{3}{8\gamma}
      +o(\gamma^{-1}).
\]
Substitution into \eqref{r63:count-var} gives its claimed limit.

The local factorial differentiation in
\eqref{r62:local-contact}, or direct differentiation of
\eqref{r63:one-bond-law}, gives the one-bond transfer
\[
 \Cov_{\rm w}(U,K)
  =\Cov(d^2-3d,d)-\mu_2.
\]
Here the contact
\(2\mu_2-2\gamma\mu_1+3\gamma\E\sigma\)
reduces exactly to \(\mu_2\) by
\eqref{r63:bond-ward}; it has not been dropped.
Since
\(0\le\Cov(d^2,d)\le\mu_3\) and
\(0\le\Var(d)\le\mu_2\), the right side is between
\(-4\mu_2\ge-15\) and \(\mu_3\le105/8\).
Its gamma-moment limit is
\[
 \frac{105}{8}-\frac{15}{4}\frac32
 -3\left(\frac{15}{4}-\frac94\right)-\frac{15}{4}
 =-\frac34.
\]
\end{proof}

\begin{corollary}[Uniform tree calibration]
On any finite tree with homogeneous positive coupling \(\gamma\),
the matrix
\(\bigl(\Cov_{\rm w}(U(K_e),K_f)\bigr)_{e,f}\)
has operator norm at most \(15\), independently of the tree's
size and of \(\gamma\).
\end{corollary}
\begin{proof}
Leaf integration gives the already-used tree factorization
\(Z(\lambda)=\prod_e\Phi(\lambda_e)\). Applying it to the count
generating function
\(Z((\lambda_e e^{h_e})_e)/Z(\lambda)\) shows that the edge
counts are independent with laws \eqref{r63:one-bond-law}.
The covariance matrix is diagonal; apply
Proposition~\ref{r63:one-bond-calibration} to each entry.
\end{proof}
We count neither leaf factorization nor tree energy independence
as a new result. The new calibration concerns the full compensated
wire count bank and its centered noise. Trees do not contain the
nonzero-net-traversal loops that obstruct the cubic problem.

\subsection{A whole-component Palm identity for the actual law}
Let \(G=(V,E)\) be a finite simple loopless graph with a reference
orientation and couplings \(\lambda_e>0\). Use the exact law
\eqref{r62:wire-law}: a labelled pairing diagram \(\omega\) has
edge counts \(K_e\), vertex degrees \(N_x=\sum_{e\ni x}K_e\),
and weight
\[
 \frac1{Z(\lambda)}
 \prod_e\frac{\lambda_e^{K_e}}{K_e!}
 \prod_x u(N_x)\,4^{\#\text{components}},
 \qquad u(2j)=\frac1{2^j(j+1)!}.
\]
Only even vertex degrees occur. In what follows observables are
invariant under independent relabellings of the occurrences of
each edge. Equivalently we use the quotient law on multisets of
closed-walk components, where a component is unrooted and its
two orientations are identified. Distinct components with the
same walk shape retain their multiplicity.

For \(L\ge2\), let \(\mathcal W_L\) be the set of ordered
vertex words
\(\eta=(x_0,\ldots,x_{L-1},x_L=x_0)\) following graph edges.
This is a set of rooted, oriented vertex words, not primitive
orbits modulo cyclic shift. All repeated edges, repeated
vertices, periodic words, backtracking and, on a torus, winding
are retained. Define
\[
 k_e(\eta)=\#\{\hbox{traversals of }e\},\quad
 n_x(\eta)=\#\{0\le j<L:x_j=x\},
\]
and let \(\ell_e(\eta)\) be the signed net traversal relative
to the reference orientation. Inserting this word as a separate
component is denoted \(\omega\oplus\eta\); it adds
\(k_e\) to \(K_e\), \(2n_x\) to \(N_x\), and
\(\ell\ell^{\mathsf T}\) to \(\Lambda\).

\begin{theorem}[Whole-component insertion with all revisit factors]
\label{r63:palm}
Set
\begin{equation*}
 \begin{aligned}
 \rho_\eta(N)
 &=\frac{2}{L}\prod_e\lambda_e^{k_e(\eta)}
       \prod_x\frac{u(N_x+2n_x(\eta))}{u(N_x)}
       \\
 &=\frac{2}{L}\prod_e\lambda_e^{k_e(\eta)}
       \prod_x\prod_{j=0}^{n_x(\eta)-1}
             (N_x+4+2j)^{-1}.
 \end{aligned}
 \tag{V63.8}\label{r63:rho}
\end{equation*}
For any nonnegative relabelling-invariant function \(H\) of a
remainder diagram and a component, invariant under rerooting
and reversal of that component,
\begin{equation*}
 \E_{\rm w}\sum_{C\in\omega}
       H(\omega\setminus C,C)
 =
 \E_{\rm w}\sum_{L\ge2}\sum_{\eta\in\mathcal W_L}
       \rho_\eta(N)\,H(\omega,\eta).
 \tag{V63.9}\label{r63:palm-eq}
\end{equation*}
The same identity holds for signed \(H\) whenever its absolute
version is integrable. Polynomial marks in counts, traversals
and word length satisfy that requirement.
\end{theorem}
\begin{proof}
Work first with finitely supported marks and labelled
occurrences. On the left mark a component of length \(L\),
one of its \(L\) starting occurrences, and one of the two
directions along it. Divide this marked count by \(2L\).
After deleting that component, put the remaining labels on
each edge in canonical increasing order.

Fix a remainder diagram, a rooted directed word \(\eta\), and
its ordered traversal positions. If the remainder has \(K_e\)
occurrences of edge \(e\), assigning distinct labels from the
enlarged set to the new ordered positions can be done in
\((K_e+k_e)!/K_e!\) ways. The remaining old occurrences inherit
the unused labels in their original order. The vertex pairings
on the inserted component are now uniquely specified by the
word and its assigned occurrences. Conversely the rooted,
directed, labelled marked component recovers this assignment
and the remainder, including when it revisits an edge or vertex.
Thus this is a bijection of the marked labelled objects.

The enlarged exponential-series factorials cancel these label
assignment factors. The added component supplies one factor
\(4\), while the changed vertex weights supply the ratios of
\(u\)'s. Division by \(2L\) therefore gives precisely
\eqref{r63:rho}. Rooted vertex words that coincide under a
periodic shift are not counted again as distinct words;
the corresponding marked occurrence assignments are already
counted in the label assignments above. No additional division
by a primitive period is permissible.

Summing the bijection and dividing by the same positive
\(Z(\lambda)\) proves the finite-mark identity. Tonelli's
theorem extends it to nonnegative marks. For integrability,
\[
 0\le\rho_\eta(N)\le\rho_\eta(0),\qquad
 \sum_{L,\eta}\rho_\eta(0)\le Z(\lambda)-1
       \le e^{\sum_e\lambda_e}-1.
\]
The first sum is exactly the single-component subseries of
the unnormalized positive partition expansion. Multiplying
each term by \(e^{aL}\) is equivalent to replacing \(\lambda\)
by \(e^a\lambda\), so all length moments of this sum are finite.
All count moments are also finite, by the exact count generating
function \(Z(e^a\lambda)/Z(\lambda)\). Polynomial marks are
therefore absolutely integrable; apply the identity to their
positive and negative parts. None of these finite-graph
exponential bounds is claimed uniform in graph size.
\end{proof}

For a single bond the two length-two vertex words have combined
intensity
\[
 \frac{2\lambda_e^2}{(N_i+4)(N_j+4)}.
\]
At the empty diagram this is \(\lambda_e^2/8\), agreeing with
\eqref{r63:one-bond-law}. For a double traversal of an oriented
triangle, length six and two visits at each vertex, there are
six distinct rooted oriented vertex words with nonzero net
traversal of absolute value two on each edge. Their total
empty-diagram activity is \(\gamma^6/6912\). These examples
test the root/orientation factor and the periodic-word convention.

\subsection{An exact reversible Stein identity}
On the relabelling quotient define the following auxiliary
whole-loop birth-and-death generator:
\begin{equation*}
 \begin{aligned}
 \mathcal L_{\rm w}F(\omega)
 &=\sum_{L,\eta}\rho_\eta(N)
       \{F(\omega\oplus\eta)-F(\omega)\}\\
 &\quad+\sum_{C\in\omega}
       \{F(\omega\setminus C)-F(\omega)\}.
 \end{aligned}
 \tag{V63.10}\label{r63:generator}
\end{equation*}
The sum in the first line runs over the rooted words just
defined. Several such words can yield the same quotient state;
their rates are added. Each existing component dies at rate one,
with multiplicities. This generator is not \(L_{\rm b}\), not
the spin generator used earlier, and not a physical Hamiltonian.

\begin{proposition}[Reversibility without a spectral-gap claim]
\label{r63:stein}
At every fixed finite graph and positive coupling profile,
\(\mathcal L_{\rm w}\) defines a nonexplosive irreducible
reversible process on the quotient diagrams with invariant
law \(\mathbb P_{\rm w}\). For polynomial observables,
\begin{equation*}
 -\E_{\rm w}F\mathcal L_{\rm w}G
 =\E_{\rm w}\sum_{L,\eta}
      \rho_\eta(N)\,\Delta_\eta F\,\Delta_\eta G,
 \qquad
 \Delta_\eta F=F(\omega\oplus\eta)-F(\omega).
 \tag{V63.11}\label{r63:dirichlet}
\end{equation*}
For \(M_t=\sum_e t_eK_e\), write
\begin{equation*}
 A_t(N)=\sum_{L,\eta}\rho_\eta(N)\,k_t(\eta),
 \qquad k_t(\eta)=\sum_e t_e k_e(\eta).
 \tag{V63.12}\label{r63:A}
\end{equation*}
Then \(\mathcal L_{\rm w}M_t=A_t(N)-M_t\), and
\begin{equation*}
 \Cov_{\rm w}(F,M_t)
 =\Cov_{\rm w}(F,A_t(N))
   +\E_{\rm w}\sum_{L,\eta}
      \rho_\eta(N)\,\Delta_\eta F\,k_t(\eta).
 \tag{V63.13}\label{r63:stein-eq}
\end{equation*}
\end{proposition}
\begin{proof}
The total birth rate is bounded uniformly in the state by
\(\sum\rho_\eta(0)<\infty\), with graph and couplings fixed.
Deaths only decrease the component count. In finite time the
number of births is dominated by a finite-rate Poisson process,
and the number of deaths is at most the initial number of
components plus the number of births. This proves nonexplosion.

Any finite diagram can reach the empty diagram by deaths.
Starting at the empty diagram, each prescribed component
has strictly positive birth intensity, so every finite quotient
diagram is reachable. The Palm identity with arbitrary
nonnegative functions of the pre- and post-birth states is
detailed balance between these births and component deaths.
The resulting stationary law is the normalized law already
specified, not a newly chosen reference ensemble.

The reversible Dirichlet form is one half the sum of birth
and death squared increments. By \eqref{r63:palm-eq}, the
expected death product equals the expected birth product.
This proves \eqref{r63:dirichlet}, with no missing factor two.
Deleting all components counts each occurrence exactly once,
so the death part of \(\mathcal L_{\rm w}M_t\) is \(-M_t\);
the birth part is \(A_t\). Stationarity gives \(\E A_t=\E M_t\).
Insert \(G=M_t\) in \eqref{r63:dirichlet} and subtract means
to obtain \eqref{r63:stein-eq}. Polynomial integrability follows
from Theorem~\ref{r63:palm}. Irreducibility and reversibility
alone provide no positive uniform spectral-gap lower bound.
\end{proof}

\subsection{Applying the identity to the full compensated bank}
Return now to the homogeneous cubic torus, with the full cycle
projection \(P=I-R\), including harmonic cycles, and
\(Q=P\circ P\). Retain V62's exact observable
\[
 \mathcal U_s^{\rm w}
 =3\Lambda_s-\sum_e s_e\{K_e(K_e-1)+(3r-2\gamma)K_e\}
       -3M_{Qs},
 \qquad
 \Lambda_s=\sum_e s_e\Lambda_{ee}.
\]
Put \(a_s(\eta)=\sum_e s_e\ell_e(\eta)^2\). Direct insertion
gives the following full increment, including its count contact:
\begin{equation*}
 \begin{aligned}
 \Delta_\eta\mathcal U_s^{\rm w}
 &=3a_s(\eta)-3k_{Qs}(\eta)\\
 &\quad-\sum_e s_e
 \{2K_e k_e(\eta)+k_e(\eta)(k_e(\eta)-1)
                     +(3r-2\gamma)k_e(\eta)\}.
 \end{aligned}
 \tag{V63.14}\label{r63:increment}
\end{equation*}
Net traversals must be computed before squaring; the term
\(a_s\) is not the unsigned length of a component. In particular
winding components remain present.

Combining \eqref{r63:stein-eq} with \eqref{r63:increment} yields
the exact upstream form of the remaining bank:
\begin{equation*}
 \begin{aligned}
 \mathcal B_{\rm w}(s,t)
 &=\Cov_{\rm w}\bigl(\mathcal U_s^{\rm w},A_t(N)\bigr)\\
 &\quad+\E_{\rm w}\sum_{L,\eta}
    \rho_\eta(N)\,
       \Delta_\eta\mathcal U_s^{\rm w}\,k_t(\eta).
 \end{aligned}
 \tag{V63.15}\label{r63:bank-stein}
\end{equation*}
The new source leg \(A_t\) depends only on the vertex degree
vector \(N=(N_x)_x\), not the edge pairing. Therefore
\[
 \Cov_{\rm w}(\mathcal U_s^{\rm w},A_t(N))
 =\Cov_{\rm w}\bigl(\E_{\rm w}[\mathcal U_s^{\rm w}\mid N],
                         A_t(N)\bigr).
\]
This conditional identity does not evaluate the conditional
mean or assert any independence. The insertion increment
depends only on the edge counts and the inserted word,
not on the old pairings. Thus \eqref{r63:bank-stein} separates
a degree-conditioned source response from an explicit marked
loop return. Neither term is known to be small separately;
their sum is the quantity that must remain compensated.

\begin{remark}[Why a naive free-walk bound can fail]
Dropping the shifts \(2j\) in the denominators of
\eqref{r63:rho} gives a termwise upper bound, but not necessarily
a finite one after summing all words. At the empty diagram on
one bond, the two alternating words of length \(2k\) have
combined exact activity
\[
 \frac{2}{k}\,
 \frac{(\gamma^2/4)^k}{((k+1)!)^2},
 \qquad k\ge1.
\]
Replacing all revisit denominators by \(4\) changes this to
\((2/k)(\gamma/4)^{2k}\), whose sum diverges for
\(\gamma\ge4\). The original factorial series is finite for
every finite \(\gamma\). The empty diagram has positive
probability under the actual normalized law. Consequently
one cannot replace the full revisit structure by a globally
convergent geometric resolvent without another estimate
that pays for low-degree states. This is a failure of that
upper bound, not a divergence of the physical wire law.
\end{remark}

\subsection{Verification and the precisely remaining task}
A deterministic diagnostic independently generates rooted vertex
words and enumerates all labelled pairing diagrams on a triangle
with at most eight edge occurrences. It checks 49,505 diagrams
over 45 count profiles. Exact rational coefficients agree for
the count Palm identity (27 coefficients), loop-matrix Palm
identity (81), full compensated-bank insertion (81), and the
normalized Stein identity (81). Length-six periodic triangle
words test that no extra primitive-period factor is inserted.
The calculation is a finite independent check of the analytic
bijection, not an infinite-volume proof.

The one-bond diagnostic checks 84 exact detailed-balance
relations and 57 raw-moment series coefficients. Independent
deterministic compact quadrature at orders 128 and 256 agrees
with the count sums at six couplings from \(1/2\) to \(256\);
the largest observed difference is below \(8\cdot10^{-11}\).
These floating checks support formula auditing only; the
uniform tree bound and limits above follow from the analytic
proofs, not the numerics. V62's exact-law diagnostics are retained
as regression checks.

The actual cubic obligation remains
\[
 \sup_{\substack{N\ge4\ {\rm even}\\\gamma\ge\gamma_0}}
 \ \sup_{\|s\|_2=\|t\|_2=1}
 |\mathcal B_{\rm w}(s,t)|<\infty
\]
for some fixed \(\gamma_0\), or a comparably effective native
estimate. V62 already pays
\(\|\mathcal B_J-\mathcal B_{\rm w}\|<320\);
we do not count that transfer again. The new work supplies an
exact, normalized whole-loop insertion mechanism and rules out
a needlessly strong raw count-concentration target. It does
not supply the uniform bound on the displayed combined Stein
return. A promising next step must control the degree response
and the marked increment together, retaining repeated-visit
suppression, not merely establish an auxiliary Markov-chain
spectral gap at each fixed volume.

Generic exterior normalizers, original exterior probabilities,
physical source coverage and physical-time contraction remain
separate unpaid obligations. An auxiliary reversible loop
process is not the Yang--Mills Hamiltonian. The interacting
four-dimensional continuum construction and its physical
mass gap remain unproved.
% END CONSOLIDATED V63 APPEND



% BEGIN CONSOLIDATED V64 APPEND
\clearpage
\section{V64: native degree tails and marked loop exceptions}
\label{r64:section}

\subsection{The upstream obligation and what is actually new}
V63 showed that a free-walk replacement of the exact insertion
intensity can diverge at low-degree diagrams of positive
probability. This appendix pays a quantitative part of that
exceptional set, under the actual normalized wire law.
We prove a joint lower-tail estimate for wire vertex degrees,
uniformly in the cubic volume, and exponential size tails for
their low-degree clusters. We then use the exact Palm identity
to control summed loop-insertion activity, not just unweighted
exception probabilities. Two ranges become controlled:
bounded local visits into a low-degree remainder, and extremely
large local visits. An exact local pairing calculation then
shows that the intermediate-visit activity does not vanish.
This rules out discarding the entire exceptional insertion
return merely because its underlying degree event is rare.
The full compensated connected response remains open.

The one-bond count law and gamma-moment estimate are inherited
from V63. The compact-spin quadratic exponential estimate is
exactly V50.11, not a new infrared inequality or Gaussian
replacement. Leaf integration on trees is not repeated.
Archive f9's star-degree factorials concern a forest expansion;
f15 V100's low-degree tail concerns Hermite occupation of a
Gaussian state. Neither supplies the native vertex-degree
estimate proved here. V63's nonduplication correction and its
exclusion of the withdrawn Lees switching argument remain
in force.

A targeted literature check considered G\"obel, Jenssen,
Michelen, Pappik, Perkins and Schiller,
\href{https://arxiv.org/html/2606.28009v1}
{\emph{Uniqueness, analyticity and mixing for Gibbs point processes
via spectral gaps}}, arXiv:2606.28009, introduction and definitions.
Its pair-potential spectral criterion is not automatically
available for the many-body vertex weights \(u(N_x)\) of our
loop law. We import no spectral-gap or mixing theorem from it.
The proofs below instead retain the known finite-volume spin
normalizer and V63's self-contained component insertion.
The scheduled companion review is next, at V65; the next broad
literature sweep is V70.

\subsection{An exact positive mixture for independent stars}
Use the positive wire law on a finite simple loopless graph
with couplings \(\lambda_e>0\). Its spin partition function is
\[
 Z(\lambda)=\int\exp\!\left(\sum_{e=\{x,y\}}\lambda_e z_x\cdot z_y\right)
                       \prod_xdH(z_x),\qquad z_x\in S^3.
\]
As in V63, \(K_e\) is the edge count and
\(N_x=\sum_{e\ni x}K_e\), which is always even. Define
\[
 \Phi(h)=\sum_{j\ge0}\frac{(h^2/4)^j}{j!(j+1)!},\qquad
 \pi_h(2j)=\frac{(h^2/4)^j}{j!(j+1)!\Phi(h)}.
 \tag{V64.1}\label{r64:bessel}
\]
At \(h=0\), \(\pi_0\) is concentrated at zero.

\begin{proposition}[Independent-star mixture, with its normalizer]
\label{r64:mixture}
Let \(A\subset V\) be independent in the graph: no edge has
both endpoints in \(A\). For the remaining spins put
\[
 H_x(z)=\sum_{y\sim x}\lambda_{xy}z_y,\qquad h_x=|H_x|,
 \qquad x\in A.
\]
The actual marginal spin law on \(V\setminus A\) is
\[
 d\nu_A(z)=\frac1{Z(\lambda)}
 \exp\!\left(\sum_{\{x,y\}\subset V\setminus A}
                        \lambda_{xy}z_x\cdot z_y\right)
 \prod_{x\in A}\Phi(h_x)\prod_{y\notin A}dH(z_y).
 \tag{V64.2}\label{r64:outside}
\]
For any nonnegative function \(F\) of the degree vector on \(A\),
\[
 \E_{\rm w}F((N_x)_{x\in A})
 =\int\sum_{(n_x)\in(2\mathbb N_0)^A}
       F((n_x))\prod_{x\in A}\pi_{h_x}(n_x)\,d\nu_A.
 \tag{V64.3}\label{r64:mixture-eq}
\]
Thus these degrees admit a positive mixture of independent
one-bond count laws. Independence is conditional in this
augmentation, not unconditional in the wire law.
\end{proposition}
\begin{proof}
Integrate the spins in \(A\). Their conditional spin integrals
factor because there are no internal edges, and each integral
equals \(\Phi(h_x)\). This proves \eqref{r64:outside}, with
the original \(Z(\lambda)\).

For real \(q_x\ge0\), multiplying every edge incident to \(x\in A\)
by \(q_x\) produces the exact degree generating function
\(\E_{\rm w}\prod_{x\in A}q_x^{N_x}\). No edge is multiplied
by two different selected vertices. Integrating the same
spins gives
\[
 \E_{\rm w}\prod_{x\in A}q_x^{N_x}
   =\int\prod_{x\in A}\frac{\Phi(q_xh_x)}{\Phi(h_x)}\,d\nu_A.
\]
Expanding each positive entire series identifies the
probability coefficients and proves \eqref{r64:mixture-eq}.
Tonelli extends it to arbitrary nonnegative \(F\).
This constructs an augmentation for the specified independent
set; it does not assert a single conditionally independent
degree family on all adjacent vertices.
\end{proof}

In particular, under a single-site augmentation the degree
parameter \(h_x\) can be small, despite large bare couplings.
We must estimate that event rather than replace \(h_x\) by
its mean. Conversely, on the homogeneous cubic graph
\(h_x\le6\gamma\) pointwise, which will pay an upper tail below.

\subsection{Uniform joint lower tails, without a global good event}
From now on use the homogeneous even cubic torus of V50:
side \(N\ge4\) even, \(n=N^3\), \(m=3n\), coupling \(\gamma>0\).
The letter \(N_x\) still denotes a wire degree, not the side
length. Let
\[
 a_\gamma=2^{3/2}e^{-2(1-\log2)\gamma},\qquad
 b_\gamma=16e^{-\gamma},\qquad
 p_\gamma=\min\{1,a_\gamma+b_\gamma\},\qquad
 \delta_\gamma=\sqrt{p_\gamma}.
 \tag{V64.4}\label{r64:constants}
\]

\begin{lemma}[Two paid lower-tail ingredients]
\label{r64:ingredients}
For \(0<q\le1\) and \(h\ge0\),
\[
 \frac{\Phi(qh)}{\Phi(h)}
 \le q^{-3/2}e^{-(1-q)h}.
 \tag{V64.5}\label{r64:ratio}
\]
For every independent vertex set \(J\) on the cubic torus,
\[
 \mu\{h_x<4\gamma\text{ for every }x\in J\}
 \le b_\gamma^{|J|},
 \qquad h_x=\gamma\left|\sum_{y\sim x}z_y\right|.
 \tag{V64.6}\label{r64:weak-fields}
\]
Consequently, when \(h\ge4\gamma\),
\(\pi_h\{N_x\le2\gamma\}\le a_\gamma\).
\end{lemma}
\begin{proof}
V63's one-bond bound \(\mu_1(h)\le3/2\) implies
\[
 \frac{d}{dh}\log\Phi(h)
   =1-\frac{\mu_1(h)}h\ge1-\frac{3}{2h}
 \quad(h>0).
\]
Integrate from \(qh\) to \(h\) to get \eqref{r64:ratio};
the case \(h=0\) follows by continuity. With \(q=1/2\),
Markov's inequality gives
\[
 \pi_h(N_x\le2\gamma)
 \le2^{2\gamma}\frac{\Phi(h/2)}{\Phi(h)}
 \le2^{3/2}e^{2\gamma\log2-h/2}
 \le a_\gamma
 \quad(h\ge4\gamma).
\]

For the second assertion write \(d_e=\gamma(1-z_x\cdot z_y)\)
and let \(E_J\) be the union of the edge stars of \(J\).
These stars have disjoint edge sets, so \(|E_J|=6|J|\).
Pointwise,
\[
 \sum_{e\ni x}d_e
  =6\gamma-z_x\cdot H_x\ge6\gamma-h_x.
\]
The event on the left of \eqref{r64:weak-fields} therefore
implies \(\sum_{e\in E_J}d_e>2\gamma|J|\).

Apply the inherited V50.11 with marks \(1/2\) on \(E_J\).
Its matrix is \(K=\frac12D_JRD_J\) on edge space, where
\(D_J\) selects \(E_J\). Thus \(0\preceq K\preceq I/2\) and
\(\Tr K=3r|J|\le|J|\), since \(r=R_{ee}\le1/3\).
For \(0\le u\le1/2\), convexity gives
\(-\log(1-u)\le2u\log2\). Hence
\[
 \E_\mu\exp\!\left(\tfrac12\sum_{e\in E_J}d_e\right)
 \le\det(I-K)^{-2}\le
        e^{4(\log2)\Tr K}\le16^{|J|}.
\]
Markov's inequality proves \eqref{r64:weak-fields}.
All expectations here are under the actual unconditioned
compact spin law. No simultaneous collective-chart event
has been imposed.
\end{proof}

\begin{theorem}[Joint native low-degree bound]
\label{r64:joint-tail}
For every independent set \(A\),
\[
 \mathbb P_{\rm w}(N_x\le2\gamma\text{ for every }x\in A)
      \le p_\gamma^{|A|}.
 \tag{V64.7}\label{r64:independent-tail}
\]
For every vertex set \(S\), with no separation assumption,
\[
 \mathbb P_{\rm w}(N_x\le2\gamma\text{ for every }x\in S)
      \le\delta_\gamma^{|S|}.
 \tag{V64.8}\label{r64:joint-tail-eq}
\]
The constants do not depend on volume or on the positions of
the selected vertices.
\end{theorem}
\begin{proof}
In the augmentation of Proposition~\ref{r64:mixture}, let
\(B_x=\{h_x<4\gamma\}\). The conditional probability that
\(N_x\le2\gamma\) is at most \(\mathbf1_{B_x}+a_\gamma\).
Conditional independence and expansion of this product give
\[
 \begin{aligned}
 \mathbb P_{\rm w}(N_x\le2\gamma,\ x\in A)
 &\le\sum_{J\subset A}a_\gamma^{|A\setminus J|}
                      \nu_A\Bigl(\bigcap_{x\in J}B_x\Bigr)\\
 &\le\sum_{J\subset A}a_\gamma^{|A\setminus J|}b_\gamma^{|J|}
   =(a_\gamma+b_\gamma)^{|A|}.
 \end{aligned}
\]
The fields \(h_x\), \(x\in A\), are functions of the remaining
spins, so their probability under \(\nu_A\) is their probability
under the full compact law. Lemma~\ref{r64:ingredients} therefore
applies without changing the normalizer. If \(a_\gamma+b_\gamma>1\)
use the trivial bound one, proving \eqref{r64:independent-tail}.

The even torus is bipartite. At least half of any \(S\) belongs
to one parity class, and that subset is independent. The event
for all of \(S\) is contained in the event for this subset.
Since \(0\le p_\gamma\le1\), this yields
\(p_\gamma^{\lceil|S|/2\rceil}\le\delta_\gamma^{|S|}\).
\end{proof}

This is an all-set intersection bound, not an assertion of
stochastic domination by independent Bernoulli variables.
The following union bound uses only the stated intersections.

\begin{corollary}[Exponentially small low-degree clusters]
\label{r64:clusters}
Call a vertex bad when \(N_x\le2\gamma\), and let
\(\mathcal C_x\) be its nearest-neighbor bad cluster, empty
when \(x\) is not bad. Then for every integer \(k\ge1\),
\[
 \mathbb P_{\rm w}(|\mathcal C_x|\ge k)
 \le\delta_\gamma(36\delta_\gamma)^{k-1}.
 \tag{V64.9}\label{r64:cluster-tail}
\]
In particular, for every \(\gamma\ge32\),
\[
 \delta_\gamma<\frac1{144},\qquad
 \mathbb P_{\rm w}(|\mathcal C_x|\ge k)
 <\frac1{144}\,4^{-(k-1)}.
 \tag{V64.10}\label{r64:cluster-threshold}
\]
\end{corollary}
\begin{proof}
A connected set of \(k\) vertices containing \(x\) admits a
canonical spanning-tree depth-first walk of length \(2(k-1)\).
Its visited set recovers the original set. There are at most
\(6^{2(k-1)}=36^{k-1}\) such walks, including those that wind
around the torus. A cluster of size at least \(k\) contains
a connected subset of size \(k\) through \(x\). Apply
\eqref{r64:joint-tail-eq} to every such subset and sum.

For the explicit threshold, \(\log2<3/4\) and \(2^{3/2}<3\)
give
\(p_\gamma\le3e^{-\gamma/2}+16e^{-\gamma}\).
The finite exponential series shows \(e^4>32\).
For \(\gamma\ge32\), therefore,
\[
 p_\gamma<
 \frac3{2^{20}}+\frac{16}{2^{40}}
 =\frac{196609}{68719476736}
 <\frac1{144^2}.
\]
Finally \(36/144=1/4\).
\end{proof}

There is no conclusion that the whole volume has no bad
vertex. The expected number may still be extensive at fixed
\(\gamma\). What is now paid is a local, volume-uniform
cluster tail under the native normalized wire measure.

\subsection{Low-degree loop insertions, with their activity included}
Retain V63's rooted-word intensity \(\rho_\eta(N)\), visit
count \(n_x(\eta)\), and full component Palm identity. Its
notation includes repeated vertices and winding words.
For a nonempty set \(S\), define the insertion event
\[
 \mathcal E_S(\omega,\eta)=
 \{N_x(\omega)\le\gamma,\quad
       1\le n_x(\eta)\le\gamma/2\text{ for every }x\in S\}.
\]
The remainder degree threshold is deliberately smaller than
the \(2\gamma\) threshold in Theorem~\ref{r64:joint-tail}.

\begin{theorem}[A summed, marked low-degree insertion budget]
\label{r64:marked-low}
For every nonempty vertex set \(S\),
\[
 \E_{\rm w}\sum_{L,\eta}\rho_\eta(N)
                       \mathbf1_{\mathcal E_S(\omega,\eta)}
 \le\gamma\,\delta_\gamma^{|S|}.
 \tag{V64.11}\label{r64:low-activity}
\]
For integers \(q_x\ge1\), \(x\in S\), and
\(q_*=\sum_{x\in S}q_x\),
\[
 \E_{\rm w}\sum_{L,\eta}\rho_\eta(N)
       \mathbf1_{\mathcal E_S(\omega,\eta)}
                 \prod_{x\in S}n_x(\eta)^{q_x}
 \le\gamma^{q_*}\delta_\gamma^{|S|}.
 \tag{V64.12}\label{r64:marked-low-eq}
\]
There is no restriction on the total word length or on visits
outside \(S\).
\end{theorem}
\begin{proof}
Apply V63's Palm identity with the event and marks in the
remainder. An inserted word becomes an existing component \(C\)
of the full diagram. If its contribution survives the event,
then for every \(x\in S\)
\[
 N_x(\omega)
   =N_x(\omega\setminus C)+2n_x(C)
   \le\gamma+2(\gamma/2)=2\gamma.
\]
Thus the whole sum of surviving components vanishes unless
every vertex of \(S\) is bad in the full diagram.
For a fixed anchor \(x_0\in S\), every surviving component
visits \(x_0\), so their number is bounded by
\(\sum_C n_{x_0}(C)=N_{x_0}/2\le\gamma\).
Taking the expectation and using
\eqref{r64:joint-tail-eq} proves \eqref{r64:low-activity}.

For the marked version use nonnegativity and the elementary
integer-power expansion:
\[
 \sum_C\prod_{x\in S}n_x(C)^{q_x}
 \le\prod_{x\in S}\left(\sum_C n_x(C)\right)^{q_x}
 =\prod_{x\in S}(N_x/2)^{q_x}.
\]
On the same bad event the last product is at most
\(\gamma^{q_*}\). Palm and
\eqref{r64:joint-tail-eq} prove \eqref{r64:marked-low-eq}.
Tonelli applies to all these nonnegative marks; no prior
termwise upper bound on the total birth intensity is used.
\end{proof}

This distinction matters: a small probability of a low-degree
remainder alone would not bound its possibly large insertion
intensity. The Palm removal step pays that intensity in
\eqref{r64:low-activity}--\eqref{r64:marked-low-eq}.
The bound also controls arbitrary long winding words in the
specified local-visit range; it does not classify every long
word as rare.

\subsection{Very large local visits have a separate native bound}
For integers \(q\ge1\), define the Touchard polynomial by
\[
 T_q(v)=\sum_{j=1}^q
       \left\{\begin{matrix}q\\j\end{matrix}\right\}v^j,
 \qquad
 a^q=\sum_{j=1}^q
       \left\{\begin{matrix}q\\j\end{matrix}\right\}(a)_j,
\]
where \((a)_j=a(a-1)\cdots(a-j+1)\) and the braces denote
the nonnegative Stirling coefficients.

\begin{proposition}[Marked native upper tail and high visits]
\label{r64:upper}
For every vertex, \(q\ge1\), and \(r_0\ge1\),
\[
 \E_{\rm w}N_x^q r_0^{N_x}
 \le T_q(6\gamma r_0)\,e^{6\gamma(r_0-1)}.
 \tag{V64.13}\label{r64:marked-mgf}
\]
Consequently, for \(v>6\), with
\(I(v)=v\log(v/6)-v+6>0\),
\[
 \E_{\rm w}(N_x/2)^q\mathbf1_{\{N_x\ge v\gamma\}}
 \le2^{-q}T_q(v\gamma)e^{-\gamma I(v)}.
 \tag{V64.14}\label{r64:upper-tail}
\]
The total marked activity of highly recurrent components satisfies
\[
 \E_{\rm w}\sum_{L,\eta}\rho_\eta(N)\,
       n_x(\eta)^q\mathbf1_{\{n_x(\eta)\ge4\gamma\}}
 \le2^{-q}T_q(8\gamma)
                e^{-(8\log(4/3)-2)\gamma}.
 \tag{V64.15}\label{r64:high-activity}
\]
All bounds are independent of the torus size.
\end{proposition}
\begin{proof}
Use the single-site mixture, for which \(0\le h_x\le6\gamma\).
The normalized Haar integral representation of \(\Phi\) gives
\(0\le\Phi^{(j)}(t)\le\Phi(t)\) for \(t\ge0\): nonnegativity
follows from its power series, and the upper bound from
\(\sigma^j\le1\) for \(-1\le\sigma\le1\).
Also \((\log\Phi)'(t)=\E_t\sigma\le1\), so
\(\Phi(r_0h)/\Phi(h)\le e^{(r_0-1)h}\) for \(r_0\ge1\).
For the one-bond count law,
\[
 \E_{\pi_h}(N_x)_j r_0^{N_x}
 =(r_0h)^j\frac{\Phi^{(j)}(r_0h)}{\Phi(h)}
 \le(r_0h)^j e^{(r_0-1)h}.
\]
Expand \(N_x^q\) in its falling factorials and use \(h\le6\gamma\).
This proves \eqref{r64:marked-mgf} after mixing. It is not
obtained by differentiating an inequality between generating
functions; the differentiated function is the exact one.

On \(N_x\ge v\gamma\), multiply by
\(r_0^{N_x-v\gamma}\ge1\), and choose \(r_0=v/6\).
This proves \eqref{r64:upper-tail}. Finally Palm converts
the left side of \eqref{r64:high-activity} to
\[
 \E_{\rm w}\sum_C n_x(C)^q
                       \mathbf1_{\{n_x(C)\ge4\gamma\}}.
\]
A surviving component forces \(N_x\ge8\gamma\), and
\(\sum_C n_x(C)^q\le(N_x/2)^q\). Apply
\eqref{r64:upper-tail} with \(v=8\).
\end{proof}

The Poisson-shaped right side of
\eqref{r64:marked-mgf} is a proved moment bound, not a
replacement of the wire law by independent Poisson counts.
The exponent in \eqref{r64:high-activity} is positive since
\(\log(4/3)>1/4\).

\subsection{Exact local resummation: rare remainders have nonvanishing activity}
The preceding estimates must not be strengthened into an
unproved small bound for every low-degree insertion. In fact
that strengthening is false for the actual cubic wire law.
We now sum over the local pairing before estimating.

The local matching-to-Ewens mechanism is standard, rather
than a new global loop-distribution theorem. A related
random-path formulation appears in Forien, Quattropani,
Quitmann and Taggi,
\href{https://www.wias-berlin.de/preprint/3029/wias_preprints_3029.pdf}
{\emph{Coexistence, enhancements and short loops in random walk
loop soups}}, WIAS preprint 3029 (2023), Section 4.1.2,
Lemma 4.4. We give the one-vertex counting proof needed here,
with our fugacity \(4\), and use it for the native Palm
return. No global Poisson--Dirichlet limit is imported.

\begin{lemma}[Local visit partition, conditioned on all other pairings]
\label{r64:local-ewens}
Fix all edge counts and all pairings except that at \(x\),
and write \(N_x=2k\). Let \(c_l(x)\) be the number of closed
components that visit \(x\) exactly \(l\) times. Then
\[
 \mathbb P_{\rm w}\bigl((c_l(x))=(c_l)\mid
                K,(\pi_y)_{y\ne x}\bigr)
 =\frac1{k+1}\prod_{l\ge1}\frac{(2/l)^{c_l}}{c_l!},
 \qquad \sum_{l\ge1}l c_l=k.
 \tag{V64.16}\label{r64:ewens-law}
\]
In particular, for \(1\le l\le k\),
\[
 \E_{\rm w}[c_l(x)\mid K,(\pi_y)_{y\ne x}]
   =\frac{2(k-l+1)}{l(k+1)}.
 \tag{V64.17}\label{r64:ewens-mean}
\]
These formulas also hold conditioned only on \(N_x=2k\).
\end{lemma}
\begin{proof}
Cut the pairing at \(x\). The rest of the diagram consists
of some closed components avoiding \(x\), and \(k\) paths
pairing the \(2k\) labelled ports at \(x\). The closed components
give a constant factor. Regard the \(k\) outside paths as
labelled objects with two ports each. The remaining vertex
pairing glues them into unoriented cycles.

For a specified set of \(l\) objects, the number of port
matchings making one cycle is \(2^{l-1}(l-1)!\), including
\(l=1\) and \(l=2\). For \(l\ge3\), choose a cyclic ordering
and the port directions, identifying the two overall
orientations; the same formula can be checked directly at
\(l=1,2\). Partitioning the \(k\) objects into groups therefore
gives
\[
 k!\prod_{l\ge1}\frac{(2^{l-1}/l)^{c_l}}{c_l!}
\]
matchings with profile \((c_l)\). Each resulting component
contributes fugacity \(4\). Its weighted count becomes
\[
 2^k k!\prod_{l\ge1}\frac{(2/l)^{c_l}}{c_l!}.
\]
The sum of the remaining products over profiles of total
size \(k\) is the coefficient of \(z^k\) in
\[
 \exp\!\left(\sum_{l\ge1}\frac{2z^l}{l}\right)
       =(1-z)^{-2},
\]
namely \(k+1\). All fixed-count vertex weights and edge
factorials cancel in this conditioning. Division proves
\eqref{r64:ewens-law}; its pairing normalizer is
\(2^k(k+1)!\).

Marking one size-\(l\) component leaves a profile of total
size \(k-l\), whose coefficient sum is \(k-l+1\).
This proves \eqref{r64:ewens-mean}. The result depends on
the exterior only through \(k\), so conditional expectation
gives the final assertion.
\end{proof}

For integers \(q\ge0\), define the low-remainder activity
moment
\[
 \mathcal I_{x,q}(\gamma)=
 \E_{\rm w}\sum_{L,\eta}\rho_\eta(N)\,
       n_x(\eta)^q\,
       \mathbf1_{\{N_x\le\gamma,\ n_x(\eta)\ge1\}} .
 \tag{V64.18}\label{r64:I}
\]
Here \(n_x^0=1\) on the indicated event. This is an averaged
birth intensity, not just the probability \(N_x\le\gamma\).

\begin{theorem}[A native nonvanishing exceptional return]
\label{r64:nonvanishing}
Put \(a(k,\gamma)=\max\{1,k-\lfloor\gamma/2\rfloor\}\) and
interpret an empty sum as zero. Under the actual wire law,
\[
 \mathcal I_{x,q}(\gamma)
   =\E_{\rm w}F_{\gamma,q}(N_x/2),\qquad
 F_{\gamma,q}(k)=
   \sum_{l=a(k,\gamma)}^k
          l^q\frac{2(k-l+1)}{l(k+1)}.
 \tag{V64.19}\label{r64:resummed}
\]
For every fixed \(q\ge0\), uniformly over all even side lengths
\(N\ge4\),
\[
 \gamma^{-q}\mathcal I_{x,q}(\gamma)
 \longrightarrow
 c_q:=2\,3^q\int_{5/6}^1 t^{q-1}(1-t)\,dt
 \quad(\gamma\to\infty).
 \tag{V64.20}\label{r64:activity-limit}
\]
In particular,
\[
 c_0=2\log(6/5)-\frac13>0,\qquad
 c_1=\frac1{12},\qquad c_2=\frac29.
 \tag{V64.21}\label{r64:activity-constants}
\]
Also, for \(\gamma\ge32\),
\(\mathcal I_{x,0}(\gamma)\le1+\gamma p_\gamma<2\).
\end{theorem}
\begin{proof}
Palm turns \eqref{r64:I} into the expectation of
\[
 \sum_C n_x(C)^q
    \mathbf1_{\{n_x(C)\ge1,\ N_x-2n_x(C)\le\gamma\}}.
\]
If \(N_x=2k\), the surviving components have
\(l=n_x(C)\ge k-\lfloor\gamma/2\rfloor\).
Condition on all other pairings and apply
\eqref{r64:ewens-mean}; this proves \eqref{r64:resummed}.

We next pay concentration and integrability without assuming
any continuum limit. Let \(D_x=\sum_{e\ni x}d_e\) in the
compact spin law, and \(\kappa=\E_\mu(z_i\cdot z_j)\) for a
nearest-neighbor edge. V62's exact count covariance identity
and the first count derivative give
\[
 \E_{\rm w}N_x=6\gamma\kappa=6\gamma-\E_\mu D_x,\qquad
 \Var_{\rm w}N_x=\Var_\mu D_x+6\gamma\kappa.
\]
Consequently
\[
 \E_{\rm w}(N_x-6\gamma)^2
    =\E_\mu D_x^2+6\gamma\kappa
    \le128+6\gamma.
 \tag{V64.22}\label{r64:degree-concentration}
\]
Here V51 gives \(\|d_e\|_2\le\sqrt{32}\,r\);
Minkowski and \(r\le1/3\) imply
\(\E D_x^2\le(6\sqrt{32}\,r)^2\le128\), and
\(\kappa\le1\). Thus \(k/\gamma=N_x/(2\gamma)\)
converges to \(3\) in \(L^2\), uniformly in the volume.

For \(k>\gamma\), every qualifying component has more than
\(k/2\) visits. There can be at most one, so
\(F_{\gamma,0}(k)\le1\) and
\(F_{\gamma,q}(k)\le k^q\). For \(k\le\gamma\),
the crude bound \(F_{\gamma,q}(k)\le k^{q+1}\) suffices.
Its contribution after division by \(\gamma^q\) is at most
\(\gamma\mathbb P(N_x\le2\gamma)\le\gamma p_\gamma\to0\).
This already proves
\(\mathcal I_{x,0}\le1+\gamma p_\gamma\);
for \(\gamma\ge32\), the bound in the proof of
\eqref{r64:cluster-threshold} and monotonicity of
\(\gamma e^{-\gamma/2}\), \(\gamma e^{-\gamma}\) show
\(\gamma p_\gamma<1\).

For \(q\ge1\), \eqref{r64:marked-mgf} at \(r_0=1\) and
power \(2q\) gives
\[
 \sup_{\gamma\ge1,\,N}
      \E_{\rm w}(k/\gamma)^{2q}<\infty.
\]
Together with \eqref{r64:degree-concentration}, this pays
the contribution of events where \(k/\gamma\) is outside
any fixed neighborhood of \(3\), by Cauchy--Schwarz.
For \(q=0\), use the bound one instead.
On a compact neighborhood of \(3\), the sum
\(\gamma^{-q}F_{\gamma,q}(k)\) is uniformly a Riemann sum:
its lower endpoint \(a(k,\gamma)/k\) tends to
\(1-1/(2u)\) when \(k/\gamma\to u\), and its limiting
value is
\[
 2u^q\int_{1-1/(2u)}^1 t^{q-1}(1-t)\,dt.
\]
This is continuous near \(u=3\); integer rounding contributes
an error tending to zero there. Concentration, the paid
uniform integrability, and the vanishing \(k\le\gamma\)
contribution prove \eqref{r64:activity-limit}.
Evaluating the three elementary integrals gives
\eqref{r64:activity-constants}.
\end{proof}

The two earlier tail bounds now identify where this mass lies.
For \(q\ge1\), the range \(n_x\le\gamma/2\) contributes at
most \(\gamma^q\delta_\gamma\); for \(q=0\) its contribution
is at most \(\gamma\delta_\gamma\).
The range \(n_x\ge4\gamma\) is negligible after division
by \(\gamma^q\), by \eqref{r64:high-activity}; for \(q=0\)
bound its intensity by its first visit moment.
Thus the intermediate range
\(\gamma/2<n_x(\eta)<4\gamma\), with \(N_x\le\gamma\)
in the remainder, has the same scaled limits \(c_q\).

This is an actual-law obstruction to treating the entire
low-degree insertion return as vanishing. It is not a
counterexample to boundedness of the compensated bank:
\(\mathcal I_{x,q}\) is a positive uncentered activity, whereas
\(\mathcal B_{\rm w}\) is a signed connected covariance.
Its nonzero constant \(c_0\), and the growing marked costs
for \(q=1,2\), show exactly why that distinction matters.

\subsection{Verification and the remaining combined return}
An exact rational diagnostic on a four-cycle enumerates
66,629 labelled pairing diagrams with at most eight occurrences.
For each of 85 admissible edge-count profiles, it compares
the wire partition coefficient with an independent integration
of the two opposite spins. This checks the full anisotropic
star coefficients and their degree generating variables.
The opposite degree covariance has leading coefficient
\(\gamma^4/16\), confirming that conditional independence
must not be promoted to unconditional independence.

The diagnostic also checks 27 coefficient identities for
low-remainder-degree and high-visit Palm marks, retaining the
before/after degree shift \(2n_x(C)\). It tests the elementary
marked-sum inequalities diagram by diagram, the explicit
\(\gamma=32\) threshold in rational arithmetic, and the
one-bond generating-function and upper-tail bounds by
deterministic sums. Two-resolution four-cycle quadrature
provides a separate floating check. These finite checks audit
the formulas; the uniform estimates follow from the proofs,
not extrapolation from the four-cycle.

The local pairing law is checked independently on all 146,600
matchings with at most seven outside paths, covering 45
weighted partition profiles. Its conditional first component
moments are also checked on every admissible four-cycle count
profile. Exact integral constants give \(c_1=1/12\) and
\(c_2=2/9\); finite Riemann sums check the resummation.
The uniform limit in \eqref{r64:activity-limit} is proved
analytically using \eqref{r64:degree-concentration}, not
inferred from those finite sums.

For an insertion through a fixed low-degree remainder
\(N_x\le\gamma\), V64 controls the local-visit ranges
\[
 1\le n_x(\eta)\le\gamma/2
 \qquad\hbox{and}\qquad n_x(\eta)\ge4\gamma.
\]
The intermediate range \(\gamma/2<n_x(\eta)<4\gamma\)
has now been resummed locally: its unweighted mean tends
to \(c_0>0\), and its first two visit moments have positive
limits after scaling by \(\gamma\) and \(\gamma^2\).
It must not be discarded as another vanishing tail.
None of these results bounds the full compensated V63
source response, either there or on ordinary-degree configurations.
The remaining bank is still the compensated sum
\[
 \mathcal B_{\rm w}(s,t)
 =\Cov_{\rm w}(\mathcal U_s^{\rm w},A_t(N))
 +\E_{\rm w}\sum_{L,\eta}\rho_\eta(N)
        \Delta_\eta\mathcal U_s^{\rm w}\,k_t(\eta).
\]
The \(Q=P\circ P\) count term in \(\mathcal U_s^{\rm w}\)
has nonlocal source support. A local visit estimate cannot
silently remove it, replace the connected covariance by a
positive moment, or infer a signed-source operator norm
from one-site tails.

In particular, on the complement of the bad set the crude
free-walk row sum is still bounded only by roughly
\(6\gamma/(2\gamma+4)\), not by a number below one uniformly
at large coupling. V63's exact revisit factors are still
needed. This appendix therefore pays explicit native
exceptional activity budgets, not a convergent unrestricted
free-walk resolvent or the desired full covariance bound.

The next iteration should integrate the local pairing law
into an estimate of the combined compensated degree response,
retaining rather than discarding the intermediate-visit
return just quantified. The local partition law does not
determine a component's spatial route or net circulation.
Those data still matter for \(\Lambda\) and the nonlocal
\(Q\) contact. Companion notes must
be reviewed at V65 before choosing that next direction.
Generic transported exteriors, their physical probabilities,
physical-time contraction and interacting four-dimensional
continuum construction remain separate obligations.
The full Yang--Mills mass gap remains unproved.
% END CONSOLIDATED V64 APPEND



% BEGIN CONSOLIDATED V65 APPEND
\clearpage
\section{V65: signed local re-pairing and the native current return}
\label{r65:context}

The target remains the full interacting four-dimensional Yang--Mills
construction and its physical Hamiltonian mass gap. Neither has been
proved. The calculation below concerns the actual positive wire
representation of the coherent three-dimensional conditional law,
not a replacement Gaussian ensemble or a physical-time dynamics.
V64 proved that rare low-degree remainders carry a nonvanishing
loop-birth activity. Their intermediate-visit contribution therefore
cannot be discarded. We now retain the signed spatial circulation
when summing the local pairing law.

The new results are an exact conditional cancellation of cross-path
circulations, its conditional covariance, a matrix-valued Palm
resummation, and a uniform-volume asymptotic for the incident-edge
current return. The last result improves a raw visit-square estimate
by one power of coupling, but proves that the current return still
has a nonzero order-\(\gamma\) term. A final identity places that
term inside the full count/factorial/contact compensation. It
does not estimate that connected compensation by its positive parts.

\subsection{Scheduled companion checkpoint and nonduplication}

The five-version checkpoint compared fresh sizes and SHA--256 hashes
with the V60 ledger, and reread the terminal theorem and scope
passages. The repository companions remain NS Stability V26,
SM--Einstein Continuum V72, Shell Mixing V16 and Parallel YM V5.
The latest reviewed downloads remain Source Selection V35, Native
Stationarity V38, Exact-Action Bridge V28 and Perturbative ESM V31.
All eight hashes are unchanged. The ESM transfer memorandum was
reread in full. These are research inputs, not instructions or
substitutes for the native YM estimates.

The useful ESM transfer is the discipline of computing an exact
low-order cancellation and all its returns before taking norms.
Its coupled action/consistency hypothesis remains unpaid even in
the terminal ESM result. No BV identity, fermionic construction or
perturbative continuum claim is imported into the wire law.
The other companions likewise supply no new uniform native YM
source-response estimate: the NS derivative gains, SM fixed-time
many-copy limit, finite Einstein initialization results and flat
holonomy shell coefficient keep their stated scopes.

V64 already paid the local Ewens visit partition and unmarked Palm
activity limits. V63 already paid the whole-component Palm identity.
V61--V62 already paid the current Ward identity and transfer to the
compensated wire bank. Those are inputs, not results claimed anew.
The retained signed outside-path geometry, its marked current
asymptotic, and the same-edge return formulation below are the
new development relative to those appendices.
The standard random-wire context is
\href{https://arxiv.org/abs/1807.06564}{Benassi--Ueltschi,
\emph{Loop correlations in random wire models}}.
Its special XY result is not an \(O(4)\) global loop-distribution
theorem. The new conditional calculations are proved here directly.
The related WIAS local-Ewens source cited in V64 was unavailable
on this checkpoint's retrieval attempt; no additional theorem
is attributed to an unread proof. The next scheduled companion
and broad-literature checkpoints are both V70.

\subsection{The exterior paths retain their actual signed circulation}

Keep the even cubic torus with \(N\ge4\), the original wire law
\eqref{r62:wire-law}, and loop fugacity four. Anisotropic positive
edge activities \(\lambda_e\) are allowed in the next two lemmas.
Fix a vertex \(x\) and let
\[
 \mathcal F_x=\sigma\bigl(K,(\pi_y)_{y\ne x}\bigr),
 \qquad N_x=2k.
\]
This conditioning fixes all counts and every pairing except that
at \(x\). Cutting the latter produces \(k\) labelled outside paths
between the \(2k\) ports at \(x\), and closed components avoiding
\(x\). Label the paths \(a=1,\ldots,k\). Choose either orientation
of each and let \(v_a\in\mathbb R^E\) be its signed edge-count
vector. A path begins and ends at the same graph vertex \(x\);
therefore \(B^{\mathsf T}v_a=0\), including paths with winding.
Define
\[
 T_x=\sum_{a=1}^k v_av_a^{\mathsf T},\qquad
 \Lambda_{\neg x}=\sum_{C:\,x\notin C}\ell(C)\ell(C)^{\mathsf T}.
 \tag{V65.1}\label{r65:outside}
\]
Both are \(\mathcal F_x\)-measurable and unchanged by reversing
the chosen path orientations. The outside paths may be long,
repeat edges, intersect spatially and wind around the torus.
They have no interior pairing at \(x\); this last fact, not
spatial simplicity, will matter for the incident edges.

\begin{lemma}[Relative orientation signs at one vertex]
\label{r65:signs}
Conditionally on \(\mathcal F_x\) and the partition of the \(k\)
outside paths into components, the circulation of a component
with block \(A\) has the form
\(\ell_A=\sum_{a\in A}\varepsilon_a v_a\).
Its relative signs are uniform and independent: one may fix
the sign of the least labelled member to \(+1\), and the other
\(|A|-1\) signs are independent fair signs. The sign choices
for different blocks are independent. In particular, every
nonconstant even sign monomial has mean zero.
\end{lemma}
\begin{proof}
For a specified block of \(l\) paths, orient the resulting cycle
so that its least labelled path is traversed in its chosen
direction, and start there. The other paths can appear in any
of \((l-1)!\) orders, with any of \(2^{l-1}\) relative signs.
Each order/sign choice determines exactly one port pairing.
This parametrization is still correct for \(l=1,2\).
All have the same component fugacity four and the same fixed
count and vertex factors. Thus the sign choices are uniform
and independent of the ordering. The product over blocks
proves independence between them. Fixing the overall orientation
does not affect an even monomial. A nonconstant even monomial
contains an unfixed sign to an odd power, proving the last claim.
\end{proof}

\begin{theorem}[Exact re-pairing mean and covariance]
\label{r65:conditional}
For \(\Lambda=\sum_C\ell(C)\ell(C)^{\mathsf T}\),
\[
 \E_{\rm w}[\Lambda\mid\mathcal F_x]
       =\Lambda_{\neg x}+T_x.
 \tag{V65.2}\label{r65:mean}
\]
For deterministic real symmetric edge matrices \(A,D\), put
\(\Lambda_A=\Tr(A\Lambda)\). Then
\[
 \Cov_{\rm w}(\Lambda_A,\Lambda_D\mid\mathcal F_x)
   =\frac43\sum_{1\le a<b\le k}
                (v_a^{\mathsf T}Av_b)(v_a^{\mathsf T}Dv_b).
 \tag{V65.3}\label{r65:conditional-cov}
\]
These are finite-graph identities; no global pairing
independence or concentration assertion is part of them.
\end{theorem}
\begin{proof}
For every fixed block partition, the diagonal path contributions
to \(\Lambda\) sum to \(T_x\). The sign lemma annihilates all
cross-path contributions, proving \eqref{r65:mean} already
conditionally on that partition.

The fluctuating part of \(\Lambda_A\), for a fixed partition, is
\[
 2\sum_{a<b}\mathbf1_{\{a,b\ {\rm in\ the\ same\ block}\}}
               \varepsilon_a\varepsilon_b\,v_a^{\mathsf T}Av_b.
\]
Products belonging to different unordered pairs have mean zero
by the sign lemma, including pairs with one common member.
The product belonging to the same pair has sign square one.
The conditional mean of \(\Lambda_A\) does not depend on the
partition, so there is no covariance of partition-dependent means.

It remains to average the same-block indicator. Exchangeability
and \eqref{r64:ewens-mean} give, for \(k\ge2\),
\[
 \mathbb P(a,b\hbox{ in one block}\mid\mathcal F_x)
 =\frac{\sum_{l=1}^k\binom l2\,2(k-l+1)/(l(k+1))}
        {\binom k2}
 =\frac13.
\]
The numerator is \(k(k-1)/6\), by summing
\((l-1)(k-l+1)/(k+1)\). Substitution proves
\eqref{r65:conditional-cov}. If \(k=0,1\), its sum is empty
and the conditional circulation quadratic is constant.
\end{proof}

For a signed edge source \(s\), take \(A=\operatorname{diag}s\).
The formula does not allow absolute values to be taken before
the sign average: doing so would retain cross-path terms that
are exactly zero. Conversely the surviving covariance in
\eqref{r65:conditional-cov} need not be small. Its exterior
vectors still have their original spatial geometry.

\subsection{Marked current Palm resummation, including the intermediate visits}

Let \(g\) be any real function on \(\{1,\ldots,k\}\) and set
\[
 h_g(k)=\frac{2}{k(k+1)}\sum_{l=1}^k(k-l+1)g(l)
 \quad(k\ge1),\qquad h_g(0)=0.
\]
The same sign calculation proves
\[
 \E_{\rm w}\!\left[
   \sum_{C:\,n_x(C)>0}g(n_x(C))\ell(C)\ell(C)^{\mathsf T}
                        \,\middle|\,\mathcal F_x\right]
        =h_g(k)T_x.
 \tag{V65.4}\label{r65:marked-mean}
\]
Indeed a fixed path belongs to a block of size \(l\) with
probability \(l\,\E[c_l\mid\mathcal F_x]/k
=2(k-l+1)/(k(k+1))\). Each cross-path term still has zero
sign average, even after this size mark. This proves
\eqref{r65:marked-mean}, not just its trace.

For a fixed integer \(q\ge0\), define the full edge matrix
\[
 \mathsf J_{x,q}(\gamma)=
 \E_{\rm w}\sum_{L,\eta}\rho_\eta(N)\,
       n_x(\eta)^q\,
       \mathbf1_{\{N_x\le\gamma,\ n_x(\eta)\ge1\}}\,
       \ell(\eta)\ell(\eta)^{\mathsf T}.
 \tag{V65.5}\label{r65:J}
\]
The degree in the indicator is that of the remainder, before
insertion. The word sum has exactly the rooted-oriented
normalization of \eqref{r63:palm-eq}. No primitive-period
factor is added. All entries are integrable on a finite graph:
the marked component sum is bounded by a polynomial in the
total count, whose moments V63 paid.

\begin{proposition}[Exact matrix return and a global upper bound]
\label{r65:matrix-return}
Let \(a(k,\gamma)=\max\{1,k-\lfloor\gamma/2\rfloor\}\), and put
\[
 h_{\gamma,q}(k)=
  \frac{2}{k(k+1)}
       \sum_{l=a(k,\gamma)}^k l^q(k-l+1)\quad(k\ge1),
 \qquad h_{\gamma,q}(0)=0.
 \tag{V65.6}\label{r65:h}
\]
For the actual homogeneous cubic law,
\[
 \mathsf J_{x,q}=\E_{\rm w}[h_{\gamma,q}(N_x/2)T_x].
 \tag{V65.7}\label{r65:J-exact}
\]
For \(q=0\), writing \(b=\min\{k,\lfloor\gamma/2\rfloor+1\}\),
\[
 h_{\gamma,0}(k)=\frac{b(b+1)}{k(k+1)}
 \quad(k\ge1),\qquad
 0\preceq\mathsf J_{x,0}\preceq\gamma\kappa_\gamma P.
 \tag{V65.8}\label{r65:J-bound}
\]
Here \(P=I-R\) is the full cycle projector, including its
harmonic sector.
\end{proposition}
\begin{proof}
Apply the whole-component Palm identity. The left side of
\eqref{r65:J} becomes the mean of
\[
 \sum_C n_x(C)^q\,
   \mathbf1_{\{n_x(C)\ge1,\ N_x-2n_x(C)\le\gamma\}}\,
          \ell(C)\ell(C)^{\mathsf T}.
\]
Condition on \(\mathcal F_x\). The selected sizes are precisely
\(l\ge a(k,\gamma)\). Formula \eqref{r65:marked-mean} proves
\eqref{r65:J-exact}. Summing \(1+\cdots+b\) gives
the explicit \(h_{\gamma,0}\); in particular it lies in \([0,1]\).
Since \(T_x\) is positive,
\[
 0\preceq\mathsf J_{x,0}\preceq\E T_x
 \preceq\E\Lambda
 =\gamma\kappa_\gamma P-\frac{\gamma}{3}\mathcal S_\gamma
 \preceq\gamma\kappa_\gamma P.
\]
The middle inequality follows by averaging \eqref{r65:mean},
since \(\Lambda_{\neg x}\) is positive. The equality is V62's
current identity combined with
\(\mathcal W_\gamma=3\kappa_\gamma R+\mathcal S_\gamma\)
from V61. All matrices retain the original normalization.
\end{proof}

The global bound is an uncentered matrix activity bound,
not the signed count-response operator norm. The factor
\(\gamma\) has not disappeared. The next result determines
that factor's nonzero coefficient on the incident edges.

\subsection{A paid volume-uniform current-activity asymptotic}

Let \(E_x\) denote the six incident edges, and let
\(\iota_x:\mathbb R^{E_x}\longrightarrow\mathbb R^E\)
be extension by zero in the fixed edge orientation. Write
\(A^{(x)}=\iota_x^{\mathsf T}A\iota_x\) for compression.
Every outside path has just two ports at \(x\). Its restricted
circulation has two entries of magnitude one, or is zero
when the two ports lie on the same edge. Consequently
\[
 T_x^{(x)}\succeq0,\qquad
 \Tr T_x^{(x)}\le2k=N_x,\qquad
 \Lambda_{\neg x}^{(x)}=0.
 \tag{V65.9}\label{r65:local-domination}
\]
No analogous global length bound by \(2k\) is being asserted.

\begin{theorem}[Nonvanishing incident-edge current return]
\label{r65:current-limit}
For every fixed integer \(q\ge0\), uniformly over even \(N\ge4\)
and vertices \(x\),
\[
 \left\|\gamma^{-(q+1)}\mathsf J_{x,q}^{(x)}
                  -\alpha_q P^{(x)}\right\|_{2\to2}
        \longrightarrow0,\qquad
 \alpha_q=2\,3^q\int_{5/6}^1t^q(1-t)\,dt
          =\frac{c_{q+1}}3.
 \tag{V65.10}\label{r65:current-limit-eq}
\]
In particular
\[
 \alpha_0=\frac1{36},\qquad
 \alpha_1=\frac2{27},\qquad
 \alpha_2=\frac{19}{96},\qquad
 \frac{(\mathsf J_{x,0})_{ee}}{\gamma}
             -\frac{1-r}{36}\longrightarrow0
       \quad(e\ni x).
 \tag{V65.11}\label{r65:current-constants}
\]
Here \(r=(N^3-1)/(3N^3)\) is allowed to depend on volume
throughout the uniform assertion.
\end{theorem}
\begin{proof}
Set \(u=k/\gamma\). V64 gives, uniformly in the volume,
\[
 \E(u-3)^2\le\frac{128+6\gamma}{4\gamma^2},
 \qquad
 \sup_{\gamma\ge1,N}\E u^p<\infty
          \quad\hbox{for each fixed integer }p.
 \tag{V65.12}\label{r65:u-bounds}
\]
The first is \eqref{r64:degree-concentration}; the second
follows from \eqref{r64:marked-mgf} and its Touchard
polynomial at \(r_0=1\). These are native-law bounds.

On a compact neighborhood of \(u=3\), the finite sum
\(\gamma^{-q}h_{\gamma,q}(k)\) is uniformly a Riemann sum
with limit
\[
 H_q(u)=2u^q\int_{1-1/(2u)}^1t^q(1-t)\,dt.
\]
The endpoint rounding and \(k+1\) denominator produce
errors tending uniformly to zero there. Moreover
\(0\le h_{\gamma,q}(k)\le k^q\) for \(k\ge1\):
it is the selected \(q\)-th size moment of the block of a
fixed path. For \(q=0\), use \(h_{\gamma,0}\le1\),
including \(k=0\).

Continuity of \(H_q\) at 3, the first bound in
\eqref{r65:u-bounds}, and Cauchy--Schwarz with moment
order \(2q+2\) give
\[
 \sup_{N,x}
 \E\!\left[u\,
   \left|\gamma^{-q}h_{\gamma,q}(k)-\alpha_q\right|\right]
       \longrightarrow0.
 \tag{V65.13}\label{r65:weighted-limit}
\]
For completeness, on \(|u-3|\le\varepsilon\) use the
uniform Riemann-sum error and the modulus of continuity
of \(H_q\), times the bounded first moment of \(u\).
On its complement bound the integrand by
\(u^{q+1}+\alpha_q u\), use Cauchy--Schwarz, and then
the concentration bound. Let \(\gamma\) tend to infinity
before sending \(\varepsilon\) to zero. This also treats
small \(k\); no division by a random vanishing degree
or unproved exceptional-event independence is used.

By \eqref{r65:local-domination} and \eqref{r65:weighted-limit},
\[
 \left\|\gamma^{-(q+1)}\mathsf J_{x,q}^{(x)}
       -\alpha_q\frac{\E T_x^{(x)}}{\gamma}\right\|
 \le2\,\E\!\left[u\,
       |\gamma^{-q}h_{\gamma,q}(k)-\alpha_q|\right]\to0.
\]
Compression of \eqref{r65:mean} and the current identity give
\[
 \frac{\E T_x^{(x)}}{\gamma}
    =\kappa_\gamma P^{(x)}-\frac13\mathcal S_\gamma^{(x)}.
\]
The inherited moment bound
\(\E d_e\le\sqrt{32}\,r\) implies
\(|1-\kappa_\gamma|\le\sqrt{32}\,r/\gamma\).
V61 gives positivity and
\((\mathcal S_\gamma)_{ee}\le512cr/\gamma\);
therefore the six-dimensional compression satisfies
\(\|\mathcal S_\gamma^{(x)}\|
 \le6(512cr)/\gamma\), where \(c\le13/16\) and \(r\le1/3\).
Thus \(\E T_x^{(x)}/\gamma-P^{(x)}\to0\) uniformly.
This proves \eqref{r65:current-limit-eq}.
Evaluation of the elementary integrals, and \(P_{ee}=1-r\),
proves \eqref{r65:current-constants}.
\end{proof}

The theorem uses squared \emph{net} currents, not squared
unsigned traversal lengths. For an incident edge,
\(|\ell_e(\eta)|\le2n_x(\eta)\), so the raw visit estimate
would allow order \(\gamma^{q+2}\) by V64.
The exact sign average lowers this to order \(\gamma^{q+1}\)
and computes its leading matrix. It does not make it vanish:
the coefficient for \(q=0\) is at least \(1/54\) on each
incident-edge diagonal.

The same leading limit remains if the inserted local visit
range is restricted to
\(\gamma/2<n_x(\eta)<4\gamma\).
Indeed the omitted small-visit diagonal is bounded by four
times V64's low-remainder mark of order \(q+2\), at most
\(4\gamma^{q+2}\delta_\gamma\). The large-visit diagonal
is bounded by four times its high-visit mark of order \(q+2\),
which has an exponentially decaying bound with polynomial
prefactor. After division by \(\gamma^{q+1}\), both vanish.
For the compressed positive matrices, sum the six diagonal
bounds to control the operator norm. This identifies the
same intermediate range as V64, now for the actual net
circulation rather than the unmarked component count.

\subsection{The exact same-edge return and the full compensated bank}

The preceding calculation is useful only if its surviving
order-\(\gamma\) contribution is retained in the count response.
For each edge \(e\), choose once and for all an endpoint \(x_e\).
After cutting the pairing at \(x_e\), let \(H_e\) count outside
paths whose \emph{two} ports both lie on occurrences of \(e\).
This is a nonnegative integer observable of
\(\mathcal F_{x_e}\), not a Hamiltonian.
It satisfies \(0\le2H_e\le K_e\).

\begin{proposition}[Same-edge return identity]
\label{r65:contact}
For \(e\ni x\), define \(H_{x,e}\) by the same rule. Then
\[
 (T_x)_{ee}=K_e-2H_{x,e},\qquad
 \E_{\rm w}[\Lambda_{ee}\mid\mathcal F_x]=K_e-2H_{x,e}.
 \tag{V65.14}\label{r65:return}
\]
If \(G(K)\) is an integrable count observable with the
displayed products integrable, then
\[
 \Cov_{\rm w}(\Lambda_{ee},G(K))
       =\Cov_{\rm w}(K_e-2H_e,G(K)).
 \tag{V65.15}\label{r65:return-cov}
\]
In particular, writing
\[
 \Gamma_{ef}=\Cov_{\rm w}(K_e,K_f),\quad
 F_{ef}=\Cov_{\rm w}(K_e(K_e-1),K_f),\quad
 C^{\rm ret}_{ef}=\Cov_{\rm w}(H_e,K_f),
\]
the entire V62 compensated bank is
\[
 \boxed{\quad
 \mathcal B_{\rm w}
       =(2\gamma I+3D)\Gamma-F-6C^{\rm ret},
 \qquad D=rI-R\circ R .
 \quad}
 \tag{V65.16}\label{r65:full-bank}
\]
No count contact or harmonic cycle term is omitted.
\end{proposition}
\begin{proof}
An outside path with one port on \(e\) has restricted
circulation \(+1\) or \(-1\) on that edge. A path with
two such ports traverses one occurrence away from \(x\)
and one toward \(x\); its net circulation on \(e\) is zero.
There are \(K_e\) ports on \(e\), of which \(2H_{x,e}\)
belong to the second class. This proves the first equality.
No component avoiding \(x\) can contain an incident edge;
use \eqref{r65:mean} for the second.
Because \(G(K)\) is measurable in every \(\mathcal F_x\),
conditional expectation proves \eqref{r65:return-cov},
including its centering.

Recall the complete inherited bank, with \(Q=P\circ P\):
\[
 \mathcal U_s^{\rm w}
 =3\Lambda_s-\sum_es_e
       \{K_e(K_e-1)+(3r-2\gamma)K_e\}-3M_{Qs},
 \qquad
 \mathcal B_{\rm w}(s,t)=\Cov_{\rm w}(\mathcal U_s^{\rm w},M_t).
\]
Replace each covariance involving \(\Lambda_{ee}\) using
\eqref{r65:return-cov}. The resulting matrix is
\[
 (2\gamma+3-3r)\Gamma-F-6C^{\rm ret}-3Q\Gamma .
\]
Since \(Q=(1-2r)I+R\circ R\),
\((3-3r)I-3Q=3(rI-R\circ R)=3D\).
This proves \eqref{r65:full-bank}. In particular \(D\)
multiplies \(\Gamma\) on the left; neither matrix is
assumed to commute with the other.
\end{proof}

As a useful normalization check, taking means in
\eqref{r65:return} and using \(\E K_e=\gamma\kappa_\gamma\)
gives the paid uniform local mean
\[
 \E H_e=\frac{\gamma\kappa_\gamma r}{2}
             +\frac{\gamma}{6}(\mathcal S_\gamma)_{ee},
 \qquad
 0\le\E H_e-\frac{\gamma\kappa_\gamma r}{2}
                    \le\frac{256}{3}cr.
 \tag{V65.17}\label{r65:return-mean}
\]
This is not a derivative bound for \(C^{\rm ret}\).
It holds for either endpoint choice; the random return
counts for the two choices need not agree pointwise.
Formula \eqref{r65:return-cov} proves equality of their
count-connected pairings. The conditional re-pairing
law is unchanged by count-source tilts, since the
activity factors depend only on \(K\); hence this
covariance identity has no missing derivative of a
coupling-dependent pairing kernel.

\subsection{What local averaging does not pay, and the next upstream target}

Let \(\mathsf T_xF=\E[F\mid\mathcal F_x]\).
The residual \(\Lambda_{ee}-\mathsf T_{x_e}\Lambda_{ee}\)
is orthogonal to every square-integrable count observable.
This removes a genuine noise component from the source
response exactly. But a dynamics that only resamples these
vertex pairings cannot estimate the remaining count response.
For its auxiliary generator
\(\mathsf G=\sum_x(\mathsf T_x-I)\), every \(f(K)\) satisfies
\(\mathsf Gf(K)=0\). Thus, for any nonconstant such \(f\)
of finite variance, its Dirichlet form is zero although
its variance is positive. For example the edge count
has nonzero variance at positive activities. There is
no full-law Poincare inequality with this Dirichlet form.
This elementary kernel obstruction does not concern the
physical Yang--Mills Hamiltonian.

The remaining target can now be attacked in terms of the
actual outside-return observable \(H_e\), rather than by
bounding every global loop-circulation square separately:
\[
 \sup_{\substack{N\ge4\ {\rm even}\\\gamma\ge\gamma_0}}
 \bigl\|(2\gamma I+3D)\Gamma-F-6C^{\rm ret}\bigr\|_{2\to2}
       <\infty .
 \tag{V65.18}\label{r65:remaining}
\]
This is an exact reformulation of the inherited sufficient
wire criterion, not a newly proved estimate. In particular
the individually large factorial/count terms must stay
with the return covariance. The V63 combined Stein identity
\eqref{r63:bank-stein} remains available for changing counts;
the re-pairing kernel alone cannot do that job.
An upstream next step is to compute a count-changing
insertion/deletion response of these same-edge outside
returns with its factorial and \(Q\) contacts retained.
A bound on the separate means in \eqref{r65:return-mean},
or on positive activity alone, is insufficient.

\paragraph{Executed checks and scope.}
The exact rational verifier enumerates 146,600 local
matchings through seven outside paths and 1,156 conditioned
set partitions. It checks the conditional mean already
within each partition, all nine pairs of three signed
quadratic sources, the factor \(4/3\) in the covariance,
and five size-marked first moments. A separate complete
four-cycle enumeration checks 66,629 native pairing
diagrams in 85 anisotropic count profiles. It verifies
the same-edge return identity with all 16 edge-count
marks per profile and a low-remainder marked-current
Palm identity, keeping the actual before/after degree
shift. The \(Q\)-to-\(D\) algebra and the integral
constants are checked over the rationals. Finite Riemann
sums are diagnostics, not proofs of uniform limits.

These checks supplement the written combinatorial and
analytic proofs; they are not proof-assistant formalization
or independent peer review. Every V64 source byte before
its final terminator is retained apart from the title
version. Other source files are hash-checked unchanged.
Compilation and visual checks are performed outside
\path{latex_notes}; only the current active \texttt{.tex}
replaces its predecessor there.

V65 pays exact signed re-pairing cancellation and the
native local current-return scale. It does not prove
\eqref{r65:remaining}, arbitrary transported-exterior
control, physical source coverage, physical-time
contraction, or an interacting four-dimensional
continuum limit. The full Yang--Mills mass gap remains
unproved.
% END CONSOLIDATED V65 APPEND



% BEGIN CONSOLIDATED V66 APPEND
\clearpage
\section{V66: whole-loop rate instability and return compensation}
\label{r66:context}

V65 retained the signed outside-path circulation and rewrote
the complete wire response in terms of same-edge return counts.
The present step tests the proposed next move: replacing those
counts by their whole-loop birth intensities. That replacement
has an exact Palm formula, but its raw \(L^2\) cost is
catastrophic under the \emph{actual cubic wire law}. We prove
an exponential-in-volume lower bound, and the exact
low-temperature exponential rate of every fixed \(p>1\)
moment. The connected count products nevertheless retain
polynomial, volume-independent entrywise bounds after the
Palm contact is kept. This distinguishes a failed estimate
from a failed Yang--Mills programme.

The full four-dimensional interacting continuum construction
and physical Hamiltonian mass gap remain open. No spectral
conclusion about that Hamiltonian follows from the auxiliary
birth rates below. The new obstruction is not an arbitrary
counterexample or a tree substitute: it uses the same
homogeneous cubic law as V50--V65.

\subsection{A count-changing formula, with its full contact}

Choose the endpoint \(x_e\) of each edge as in V65.
For a closed word \(\eta\), cut it at each visit to \(x_e\).
Let \(h_e(\eta)\) be the number of resulting excursions
whose departure and arrival edges are both \(e\).
This definition includes a two-step backtrack and is
invariant under rerooting and reversal. If there is no
visit to \(x_e\), set \(h_e=0\). The inherited outside-return
observable then satisfies the pathwise identities
\[
 H_e(\omega)=\sum_{C\in\omega}h_e(C),\qquad
 H_e(\omega\oplus\eta)-H_e(\omega)=h_e(\eta),\qquad
 0\le2h_e(\eta)\le k_e(\eta).
 \tag{V66.1}\label{r66:additive}
\]
Indeed cutting at \(x_e\) acts separately on each component.
Each counted excursion uses two distinct occurrences of
\(e\). Inserting a separate component does not change
any old pairing or old excursion. These elementary
identities are the starting algebra, not the main new bound.

Use exactly \(\rho_\eta(N)\) from \eqref{r63:rho}, and define
\[
 a_e(N)=\sum_{L,\eta}\rho_\eta(N)h_e(\eta),\qquad
 A_e(N)=\sum_{L,\eta}\rho_\eta(N)k_e(\eta).
 \tag{V66.2}\label{r66:rates}
\]
The second is V63's count birth intensity. Both depend on
the remainder degrees, not its pairings. At every fixed
finite graph and coupling they are finite, nonnegative,
and bounded by their values at the zero degree vector.
Also \(2a_e\le A_e\). The degree-zero argument in \(a_e(0)\)
does not mean zero coupling.

\begin{proposition}[Return Palm identity without omitted count contact]
\label{r66:palm}
Under the actual wire law,
\[
 \E a_e=\E H_e,\qquad \E A_e=\E K_e,
 \qquad
 C^{\rm ret}_{ef}
   =\Cov_{\rm w}(a_e(N),K_f)+\Pi_{ef},
 \tag{V66.3}\label{r66:palm-eq}
\]
where
\[
 \Pi_{ef}
 =\E_{\rm w}\sum_{L,\eta}\rho_\eta(N)h_e(\eta)k_f(\eta)
 =\E_{\rm w}\sum_C h_e(C)k_f(C)\ge0.
 \tag{V66.4}\label{r66:contact}
\]
Here \(C^{\rm ret}_{ef}=\Cov_{\rm w}(H_e,K_f)\)
is precisely the V65 return covariance, not a new proxy.
\end{proposition}
\begin{proof}
Apply \eqref{r63:palm-eq} to the additive component mark
\(h_e\), and then to \(k_e\), for the means. Apply it to
\(h_e(C)K_f(\omega)\). After deleting \(C\), the latter
count is \(K_f(\omega\setminus C)+k_f(C)\). Thus
\[
 \E H_eK_f=\E a_e(N)K_f+
             \E\sum_{L,\eta}\rho_\eta(N)h_e(\eta)k_f(\eta).
\]
Subtract the product of the already identified means.
Applying Palm to the component product alone gives the
second form of \(\Pi\). Nonnegative marks justify Tonelli;
finite polynomial moments were paid in V63.
\end{proof}

Consequently the full bank is
\[
 \mathcal B_{\rm w}
  =(2\gamma I+3D)\Gamma-F
       -6[\Cov_{\rm w}(a_e(N),K_f)]_{e,f}-6\Pi,
 \tag{V66.5}\label{r66:combined}
\]
with the same \(D,\Gamma,F\) as \eqref{r65:full-bank}.
The positive contact \(\Pi\) cannot be deleted because
the first covariance happens to have a useful sign.
Nor will taking an \(L^2\) norm of \(a_e\) pay this formula:
the next theorem rules out that route quantitatively.

\subsection{Rooted Euler words: the classical count and our normalization}

The counting input is the classical last-exit-tree argument
of T.~van Aardenne-Ehrenfest and N.~G.~de Bruijn,
\href{https://pure.tue.nl/ws/files/3311129/597493.pdf}
{\emph{Circuits and trees in oriented linear graphs}},
Simon Stevin \textbf{28} (1951), 203--217, Section~6,
Theorems~5a, 5b and 6. The original definitions and
last-exit proof on pages 212--214 were inspected.
We count rooted \emph{vertex words}, rather than labelled
circuits modulo cyclic rotation, so we supply the
normalization explicitly. This classical enumeration
is not claimed as a new theorem.

Consider a connected simple \(d\)-regular undirected graph
on \(n\) vertices, with both directions of each edge
available as arcs. Let \(\tau\) be the number of directed
spanning trees pointing to a specified root \(x\) in
this symmetric directed graph. This equals its
undirected spanning-tree count and is independent of
the root. For an integer \(q\ge1\), prescribe \(q\)
traversals of every directed arc. Put \(L=dnq\).

\begin{lemma}[All rooted words, including the periodic ones]
\label{r66:euler}
The number \(W_x(q)\) of rooted vertex words starting at
\(x\), with these directed traversal counts, is
\[
 W_x(q)=\tau d^{\,1-n}
             \left(\frac{(dq)!}{(q!)^d}\right)^n .
 \tag{V66.6}\label{r66:W}
\]
If \(e=\{x,y\}\) and \(q\ge2\), the number
\(W_{x,e}^{\rm bt}(q)\) of such words starting with
\(x,y,x\) satisfies
\[
 W_{x,e}^{\rm bt}(q)
       \ge\frac{2^{-(n-1)}}{d^2}W_x(q).
 \tag{V66.7}\label{r66:prefix}
\]
Every word in this latter class has \(h_e\ge1\), when
the selected endpoint of \(e\) is \(x\).
\end{lemma}
\begin{proof}
For balanced directed multiplicities \(m_{uv}\), let
\(b_v=\sum_um_{vu}\) and let \(t_x(m)\) count directed
spanning trees to \(x\), with parallel arcs distinguished.
The last outgoing arc at each nonroot vertex forms such
a tree in an Euler word. Conversely reserve a chosen
tree arc as that last exit and order the other outgoing
arcs arbitrarily. Starting at \(x\) and taking successive
unused arcs exhausts them all: it can stop only at \(x\);
an unused arc would, along the last-exit tree, force an
unused exit at \(x\). The number of labelled orders is
\(t_x(m)b_x\prod_v(b_v-1)!\).
Each vertex word has exactly \(\prod_{uv}m_{uv}!\)
assignments of distinct arc labels. Therefore
\[
 W_x(m)=
 \frac{t_x(m)b_x\prod_v(b_v-1)!}{\prod_{uv}m_{uv}!}.
 \tag{V66.8}\label{r66:general-W}
\]
For \(m_{uv}=q\), one has
\(t_x(m)=q^{n-1}\tau\) and \(b_v=dq\).
Simplification proves \eqref{r66:W}.

Remove one traversal in each direction of \(e\).
The remaining balanced multiplicities \(m'\) have
\(b'_x=b'_y=dq-1\), with other degrees unchanged.
Prepending \(x,y,x\) is a bijection from their words
rooted at \(x\) to the prescribed-prefix class.
Every original spanning tree has at least
\((q-1)^{n-1}\) available arc assignments in \(m'\).
Thus \(t_x(m')\ge(q-1)^{n-1}\tau\), and direct division
of \eqref{r66:general-W} gives
\[
 \frac{W_{x,e}^{\rm bt}(q)}{W_x(q)}
  =\frac{t_x(m')}{q^{n-1}\tau}
                      \frac{q^2}{dq(dq-1)}
  \ge\frac{2^{-(n-1)}}{d^2}.
\]
The first two steps form an excursion from \(x\) with
both ports on \(e\), proving the last assertion.
\end{proof}

No additional period divisor belongs in these formulas.
A periodic vertex sequence is counted once as that
rooted sequence. Its assignments of labels to repeated
arc positions were already included in
\(\prod m_{uv}!\). The rate \(2/L\) in \(\rho_\eta\)
will be applied exactly once, as in V63.

\subsection{A single-component sector with the full exponential rate}

Now specialize to the actual cubic torus:
\[
 n=N^3,\qquad d=6,\qquad m=|E|=3n,\qquad
 q=\lfloor\gamma/2\rfloor,\qquad \gamma\ge4.
\]
Denote its spanning-tree count by \(\tau_N\ge1\).
To make the compact normalization unambiguous, write
\[
 Z_+(\gamma)=\int\exp\!\left(\gamma\sum_{e=\{i,j\}}
                       z_i\cdot z_j\right)\prod_i d\sigma(z_i).
\]
Here \(d\sigma\) is normalized Haar measure on \(S^3\).
This is the positive wire partition function; the
constant \(e^{-m\gamma}\) in the coherent energy form
has canceled. In particular,
\[
 \mathbb P_{\rm w}(\varnothing)=Z_+(\gamma)^{-1},
 \qquad Z_+(\gamma)\le e^{m\gamma}.
 \tag{V66.9}\label{r66:empty}
\]
Set
\[
 \mathfrak b_\gamma
     =\frac{(\gamma/2)^{6q}}{(6q+1)(q!)^6},
 \qquad
 Z_1(\gamma)=\sum_{L,\eta}\rho_\eta(0).
 \tag{V66.10}\label{r66:b}
\]
The second quantity is exactly the unnormalized
single-component sector, by the inherited Palm counting.
It is not the full partition function.

\begin{proposition}[Native empty-remainder lower bounds]
\label{r66:empty-bounds}
For every edge and either selected endpoint,
\[
 \begin{aligned}
 Z_1(\gamma)&\ge
       \frac{2\tau_N}{6^nq}\mathfrak b_\gamma^n,\\
 A_e(0)&\ge4\tau_N
                  \left(\frac{\mathfrak b_\gamma}{6}\right)^n,\\
 a_e(0)&\ge\frac{\tau_N}{9nq}
                  \left(\frac{\mathfrak b_\gamma}{12}\right)^n .
 \end{aligned}
 \tag{V66.11}\label{r66:lower-rates}
\]
Also, at every degree vector of an actual diagram,
\[
 \begin{gathered}
 0\le a_e(N)\le a_e(0)\le\tfrac12\gamma Z_+(\gamma),\\
 0\le A_e(N)\le A_e(0)\le\gamma Z_+(\gamma),\\
 \E a_e=\E H_e\le\gamma/2,\qquad
 \E A_e=\E K_e\le\gamma .
 \end{gathered}
 \tag{V66.12}\label{r66:upper-rates}
\]
\end{proposition}
\begin{proof}
For the words of Lemma~\ref{r66:euler},
\(L=6nq\), each vertex has \(6q\) visits and each
undirected edge has \(2q\) traversals. At the empty
remainder, their common rate is
\[
 \frac{2}{L}\frac{(\gamma/2)^L}{((6q+1)!)^n}.
\]
Multiply by the \(nW_x(q)\) words rooted at all vertices
and use \eqref{r66:W}. This gives the first lower bound.
Multiplication by \(k_e=2q\) gives the second.
For the third, use only words rooted at \(x_e\) with the
prescribed backtrack prefix. Their fraction is at least
\(2^{-(n-1)}/36\), and their mark \(h_e\) is at least one.
Multiplying this fraction by the first bound divided by
\(n\) gives exactly the third line of
\eqref{r66:lower-rates}. No negative contributions were
discarded; all these series are positive.

Every factor in \(\rho_\eta(N)\) decreases as a remainder
degree increases. The sums at zero are unnormalized
one-component count or return moments. Enlarging to all
diagrams and using \(2H_e\le K_e\) yields
\[
 A_e(0)\le Z_+\E K_e,\qquad
 a_e(0)\le\tfrac12Z_+\E K_e.
\]
The exact count mean is \(\E K_e=\gamma\kappa_\gamma\)
with \(0\le\kappa_\gamma\le1\).
The Palm mean identities complete \eqref{r66:upper-rates}.
\end{proof}

We will need only an elementary factorial estimate,
not a fitted saddle-point approximation. Increasingness
of \(\log t\) gives
\[
 \log(q!)\le\int_1^q\log t\,dt+\log q
      =q\log q-q+1+\log q.
\]
Since \(2q\le\gamma<2q+2\), it follows that
\[
 \mathfrak b_\gamma
 \ge\frac{e^{6q-6}}{7q^7}
 \ge\frac{e^{3\gamma-12}}{7q^7}.
 \tag{V66.13}\label{r66:b-bound}
\]
The polynomial power here is deliberately loose; the
exponential rate is exact. For each fixed cubic torus,
\eqref{r66:lower-rates}, \eqref{r66:empty} and
\(Z_1\le Z_+\) already prove
\[
 \lim_{\gamma\to\infty}\frac{\log Z_1(\gamma)}{\gamma}
 =\lim_{\gamma\to\infty}\frac{\log Z_+(\gamma)}{\gamma}
 =3n.
 \tag{V66.14}\label{r66:sector-rate}
\]
Thus the one-component sector has the full fixed-volume
exponential free-energy rate. This is not an assertion
that one component has probability near one, or a
Poisson--Dirichlet theorem.

\subsection{The raw return intensities have exponentially large moments}

\begin{theorem}[An actual-law obstruction to raw-rate \(L^2\) estimates]
\label{r66:instability}
For every fixed even \(N\ge4\), every edge \(e\), and
every real \(p>1\),
\[
 \lim_{\gamma\to\infty}
   \frac1\gamma\log\E_{\rm w}a_e(N_\cdot)^p
 =\lim_{\gamma\to\infty}
   \frac1\gamma\log\E_{\rm w}A_e(N_\cdot)^p
 =(p-1)3n .
 \tag{V66.15}\label{r66:Lp-rate}
\]
Here \(N_\cdot\) is the random degree vector; the side
length \(N\) is fixed in these two limits. For \(p=2\)
the same rate \(3n\) holds for the variance of either
intensity. Quantitatively, for all even \(N\ge4\)
and all \(\gamma\ge64\),
\[
 \E_{\rm w}a_e(N_\cdot)^2
       \ge\frac{4e^{n\gamma}}{81n^2\gamma^2},\qquad
 \Var_{\rm w}(a_e(N_\cdot))
       \ge\frac{4e^{n\gamma}}{81n^2\gamma^2}
                      -\frac{\gamma^2}{4}.
 \tag{V66.16}\label{r66:volume-growth}
\]
All moments remain finite at each fixed regulator.
\end{theorem}
\begin{proof}
The empty diagram alone gives
\[
 \E a_e^p\ge Z_+^{-1}a_e(0)^p
                  \ge e^{-3n\gamma}a_e(0)^p .
\]
The last line of \eqref{r66:lower-rates} and
\eqref{r66:b-bound} imply, with \(n\) fixed,
\[
 \log a_e(0)\ge
      3n\gamma-(7n+1)\log\gamma-O_N(1).
\]
This proves the lower limit in \eqref{r66:Lp-rate}.
The upper bound follows from positivity and monotonicity:
\[
 \E a_e^p\le a_e(0)^{p-1}\E a_e
       \le(\gamma/2)^p e^{(p-1)3n\gamma}.
\]
The same argument for \(A_e\) uses the second line of
\eqref{r66:lower-rates} and
\(\E A_e^p\le\gamma^p e^{(p-1)3n\gamma}\).
The means are bounded by polynomials, so subtracting
their squares does not change the positive exponential
rate at \(p=2\).

For the explicit volume bound put
\[
 \mathfrak r_\gamma
       =e^{-3\gamma}\mathfrak b_\gamma^2/144.
\]
The same empty-event calculation, without a fixed-volume
limit, gives
\[
 \E a_e^2\ge
   \frac{\tau_N^2}{81n^2q^2}\mathfrak r_\gamma^n.
 \tag{V66.17}\label{r66:exact-volume-bound}
\]
By \(q\le\gamma/2\) and \eqref{r66:b-bound},
\[
 \mathfrak r_\gamma\ge
       \frac{1024}{441}
               \frac{e^{3\gamma-24}}{\gamma^{14}}.
\]
For \(\gamma\ge64\), one has
\(\log\gamma\le\gamma/12\).
Indeed \(\log2<3/4\), since
\(\sum_{j=0}^3(3/4)^j/j!>2\); hence
\(\log64<9/2<64/12\), and
\(\gamma/12-\log\gamma\) is increasing there.
Consequently
\[
 \log\mathfrak r_\gamma
 \ge3\gamma-24-14\log\gamma
 \ge11\gamma/6-24\ge\gamma.
\]
Use \(\tau_N\ge1\) and \(q\le\gamma/2\) in
\eqref{r66:exact-volume-bound}. This proves the first
bound in \eqref{r66:volume-growth}; the mean bound in
\eqref{r66:upper-rates} proves the second. Finiteness
follows from the finite state-independent upper rates
in that same display.
\end{proof}

In particular, even for one unit coordinate source, the
centered return intensity has no volume-uniform \(L^2\)
bound at fixed \(\gamma\ge64\). No polynomial in \(\gamma\)
can bound its \(L^2\) norm uniformly at fixed volume
either. This rules out a uniform Cauchy--Schwarz estimate
based on the raw intensity norm in \eqref{r66:combined}.
It does \emph{not} rule out a uniform bound for the
combined covariance itself. The empty state is genuinely
rare; the calculation shows why its birth rate compensates
that rarity in first moments and overwhelms it in higher ones.

\subsection{The count-connected products are nevertheless polynomial}

The exponentially large rate moments must not be confused
with the products occurring after the exact Palm contact.
Put
\[
 B_\gamma^{\rm loc}=\gamma^2+\gamma+\frac{32}{9}.
\]
The inherited count covariance and local energy moment give
\[
 \E K_e^2
   =(\gamma\kappa_\gamma)^2+\gamma\kappa_\gamma
                         +\Var_\mu(d_e)
       \le B_\gamma^{\rm loc}.
 \tag{V66.18}\label{r66:count-moment}
\]
This is uniform over all even \(N\ge4\), since
\(\E_\mu d_e^2\le32r^2\) and \(r\le1/3\).

\begin{proposition}[Polynomial mixed products with the return contact retained]
\label{r66:products}
For every pair of edges, at every positive coupling,
\[
 \begin{gathered}
 0\le\E H_eK_f\le\tfrac12 B_\gamma^{\rm loc},\qquad
 0\le\E a_e(N_\cdot)K_f\le\tfrac12 B_\gamma^{\rm loc},\\
 0\le\Pi_{ef}\le\tfrac12 B_\gamma^{\rm loc},\qquad
 |\Cov_{\rm w}(a_e(N_\cdot),K_f)|
                         \le\tfrac12 B_\gamma^{\rm loc}.
 \end{gathered}
 \tag{V66.19}\label{r66:product-bounds}
\]
These are volume-independent entrywise bounds, not an
operator-norm estimate for arbitrary signed source banks.
\end{proposition}
\begin{proof}
Pathwise \(H_e\le K_e/2\), so Cauchy--Schwarz applied
to the \emph{counts}, using \eqref{r66:count-moment},
gives the first inequality. Also
\[
 \sum_C h_e(C)k_f(C)
       \le\tfrac12\sum_C k_e(C)k_f(C)
       \le\tfrac12K_eK_f.
\]
This proves the contact bound by
\eqref{r66:contact}. The uncentered Palm formula says
\(\E a_eK_f=\E H_eK_f-\Pi_{ef}\); both terms on its
left and \(\Pi\) are nonnegative. This proves the
second bound without taking an \(L^2\) norm of \(a_e\).
Finally \((\E a_e)(\E K_f)\le\gamma^2/2\).
The absolute difference of two nonnegative quantities
is at most their maximum, proving the covariance bound.
\end{proof}

The coexistence of \eqref{r66:volume-growth} and
\eqref{r66:product-bounds} is exact: the raw rate is large
on configurations whose count marks are correspondingly
small. Treating its two factors by separate \(L^2\)
norms loses that relation. The proposition avoids an
exponential loss but still permits order \(\gamma^2\)
per entry. Summing these bounds over source supports
would again lose volume factors. Therefore it does
not prove \eqref{r65:remaining}.

\subsection{Verification, nonduplication, and the upstream change of direction}

An exact dynamic programme counts rooted vertex words
and prescribed-backtrack words independently of the
last-exit formula on a triangle, four-cycle, six-cycle
and complete four-vertex graph, with repeated directed
traversals in the first two. All nine cases agree.
The actual \(4^3\) torus spanning-tree count is checked
both by an integer Laplacian cofactor and by the
independent product of its Fourier eigenvalues.
The four-cycle native-wire check enumerates 66,629
labelled pairing diagrams in 85 anisotropic profiles.
It compares component excursion counts with independently
cut outside paths, checks the pointwise contact bound,
and verifies all 144 mixed Palm coefficients through
total occurrence order eight. The elementary logarithm
threshold is checked rationally; finite logarithmic
samples only audit the formulas. The general bounds
and limits follow from the written proofs.

V63's Palm theorem and V65's return covariance are
inherited. The Euler enumeration is classical and
credited. The new result is their quantitative use to
prove instability of the \emph{actual native} return
birth intensities, together with the contact-preserving
polynomial mixed-product estimate. Neither an isolated
toy process nor an assumed asymptotic loop distribution
is used. This is not a repetition of V64's nonvanishing
mean activity or V65's order-\(\gamma\) net-current return:
all \(p>1\) raw-rate moments have the exponential rate
\eqref{r66:Lp-rate}, and the bound
\eqref{r66:volume-growth} holds at fixed large coupling
while the cubic volume grows.

The next upstream step must therefore avoid proving
uniform \(L^2\) bounds for these whole-component rates.
One option to investigate is an explicitly normalized
local path or port surgery that changes counts without
exposing an arbitrarily large whole-loop birth clock.
Such a mechanism would require its own detailed balance,
all return terms and a quantitative estimate; it is
not supplied by the pairing-only dynamics, whose count
kernel V65 already identified. Alternatively one must
estimate the original combined connected product before
splitting it into raw rates. The factorial term \(F\),
the full \(D\Gamma\) contact and the return \(\Pi\) in
\eqref{r66:combined} all remain present.

Prior mathematical bytes are preserved apart from the
version title, and other sources are hash-checked.
Build and verification artifacts remain outside
\path{latex_notes}. Companion and broad-literature
reviews are next scheduled at V70. No generic-exterior
control, physical-time contraction, interacting
four-dimensional construction or Yang--Mills mass-gap
proof is claimed.
% END CONSOLIDATED V66 APPEND



% BEGIN CONSOLIDATED V67 APPEND
\clearpage
\section{V67: root-pinned occupation and bounded count rates}
\label{r67:context}

V66 ruled out a uniform \(L^2\) estimate for the raw whole-loop
birth intensities under the actual cubic law. We now go upstream
to a local two-link identity and integrate two genuine sources
of auxiliary noise: the loop colours and the pairing at one
endpoint. The resulting positive rate deficit is a regularized
occupation of loops avoiding that endpoint. A one-vertex
transverse deformation, controlled by the inherited native
infrared bound, gives this occupation a volume-uniform
exponential moment. The deficit multiplied by coupling has
the same kind of bound.

This is progress on the native source-response problem, not
a construction of the four-dimensional continuum theory or
a proof of its physical mass gap. The new count dynamics
preserves parity sectors and is not the physical Hamiltonian.
The remaining signed connected operator bound is stated
explicitly at the end.

\paragraph{Inputs and nonduplication.}
Colour expansions of the sphere integral and wire
representations are standard; see Benassi--Ueltschi,
\href{https://arxiv.org/html/1807.06564v4}
{\emph{Loop correlations in random wire models}},
Section~6, Proposition~6.1 and its proof, particularly
(6.7)--(6.10). Those passages were reread with the
normalization \(2J_e=\lambda_e\). No special XY
Poisson--Dirichlet result is imported for \(O(4)\).
V42--V43 already used pinned colour expansions in finite
coherent blocks; V64 paid the one-vertex pairing law;
V50 paid the native quadratic source bound. We reuse them.
The new assertions are the positive count-source deficit,
its full-cubic exponential estimates and uniform mean
asymptotic, and the bounded count update that they justify.

\subsection{The coloured count law and a bounded local update}

Use the four coordinate colours \(a=1,\ldots,4\) of \(S^3\),
not four independent Gaussian fields. Let \(k_e^a\ge0\),
\[
 K_e=\sum_a k_e^a,\qquad N_x^a=\sum_{e\ni x}k_e^a,\qquad
 N_x=\sum_a N_x^a .
\]
Every \(N_x^a\) is even. The exact coloured-count weight is
\[
 \frac1{Z_+(\lambda)}
 \prod_{e,a}\frac{\lambda_e^{k_e^a}}{k_e^a!}
 \prod_x\left[u(N_x)\prod_{a=1}^4(N_x^a-1)!!\right],
 \qquad u(2j)=\frac1{2^j(j+1)!},
 \tag{V67.1}\label{r67:colour-law}
\]
with \((-1)!!=1\). Indeed expand each coordinate coupling
and use, for even exponents,
\[
 \int_{S^3}\prod_a z_a^{N^a}\,d\sigma
 =u(\sum_aN^a)\prod_a(N^a-1)!! .
\]
Equivalently colour every component of the uncoloured
wire diagram independently and uniformly in four colours.
Summing the occurrence labels gives the factorials in
\eqref{r67:colour-law}; summing same-colour pairings
gives its double factorials. Conditional on coloured
occurrence counts and labels, vertex pairings are
independent and uniform within each colour. This
conditional statement does not hold for the uncoloured
pairing law.

Define the four-component probability vectors
\[
 p_x^a=\frac{N_x^a+1}{N_x+4},\qquad \sum_a p_x^a=1.
\]
Adding two links of colour \(a\) to \(e=\{x,y\}\) has
weight ratio
\[
 \frac{\lambda_e^2p_x^ap_y^a}
                 {(k_e^a+1)(k_e^a+2)}.
 \tag{V67.2}\label{r67:colour-ratio}
\]
This follows directly from the two changed factorials
and vertex moments. Thus birth rates
\(b_{e,a}=\lambda_e^2p_x^ap_y^a\) and death rates
\(k_e^a(k_e^a-1)\) obey detailed balance. Their total
birth rate on a homogeneous cubic graph is at most
\(\sum_e\gamma^2\sum_a p_x^ap_y^a\le m\gamma^2\).
Births add two occurrences and deaths remove two;
this bound proves nonexplosion at fixed graph and
coupling. The full coloured law is invariant, but the
update preserves every \(k_e^a\bmod2\). No full-law
ergodicity or spectral gap follows.

\subsection{Two exact averages give a positive count-source deficit}

For an ordered choice of endpoints \(e=(x,y)\), let
\[
 J_e=\sum_{C:\,n_x(C)=0}n_y(C).
 \tag{V67.3}\label{r67:J}
\]
This is the number of visits to \(y\) in components
avoiding \(x\). It is unchanged when the pairing at
\(x\) is resampled. In particular
\(0\le J_e\le N_y/2\). Define
\[
 \chi_e=\frac32\,\frac{J_e+2}{N_y+4},\qquad
 T_e=\gamma\chi_e .
 \tag{V67.4}\label{r67:chi}
\]
The constant two in the numerator will be retained.

\begin{theorem}[Positive shifted factorial identity]
\label{r67:factorial}
For every integrable uncoloured count mark \(G(K)\)
with the indicated products integrable,
\[
 \E_{\rm w}[K_e(K_e-1)G(K)]
    =\gamma^2\E_{\rm w}[(1-\chi_e)G(K+2\delta_e)].
 \tag{V67.5}\label{r67:factorial-eq}
\]
Moreover
\[
 \frac14\le1-\chi_e<1,\qquad
 \E_{\rm w}(1-\chi_e)=\E_\mu(z_x\cdot z_y)^2 .
 \tag{V67.6}\label{r67:positivity}
\]
Either endpoint may be selected. The resulting random
\(\chi_e\)'s need not coincide, but their count-connected
pairings do.
\end{theorem}
\begin{proof}
We first use arbitrary positive edge activities
\(\lambda\). Inserting \(z_x^{a\,2}z_y^{a\,2}\) in the
coloured expansion multiplies the two vertex moments
by \(p_x^ap_y^a\). Common rotation of two fixed unit
vectors \(v,w\) gives
\[
 \int_{O(4)}\sum_a (Uv)_a^2(Uw)_a^2\,dU
           =\frac{1+2(v\cdot w)^2}{6}.
 \tag{V67.7}\label{r67:rotation}
\]
To check the coefficient, take
\(v=e_1\), \(w=t e_1+\sqrt{1-t^2}e_2\);
one row of \(U\) has moments \(1/8\) and \(1/24\)
for its fourth power and mixed square product.
Odd terms vanish. Consequently
\[
 \E_{\rm col}\mathcal R_e
       =\E_{\mu_\lambda}(z_x\cdot z_y)^2,\qquad
 \mathcal R_e=3\sum_a p_x^ap_y^a-\frac12 .
\]
This remains true with every count generating factor
\(e^{h\cdot K}\), since it just replaces
\(\lambda_f\) by \(\lambda_fe^{h_f}\). The factorial
derivative of the spin partition function therefore gives
\[
 \E_{\rm w}[(K_e)_2 e^{h\cdot K}]
       =\lambda_e^2 e^{2h_e}
                        \E_{\rm col}[\mathcal R_e e^{h\cdot K}].
 \tag{V67.8}\label{r67:generating}
\]
Absolute convergence, or coefficient comparison in the
positive edge activities, yields the corresponding
identity for every count mark, with shift \(2\delta_e\).
We next remove the sign-indefinite auxiliary mark
\(\mathcal R_e\).

Condition on the uncoloured diagram and average its
independent component colours. If
\(V_{xy}=\sum_C n_x(C)n_y(C)\), direct colour counting gives
\[
 \E_{\rm col}[\mathcal R_e\mid\omega]
       =\frac14+\frac{9V_{xy}}{(N_x+4)(N_y+4)}.
 \tag{V67.9}\label{r67:colour-average}
\]
In detail,
\(\E\sum_aN_x^aN_y^a=N_xN_y/4+3V_{xy}\);
the remaining terms in the numerator sum to
\(N_x+N_y+4\).

Now condition on \(\mathcal F_x\) from V65 and put
\(k=N_x/2\). The \(k\) outside paths carry fixed
numbers \(v_a\) of visits to \(y\), whose sum is
\(N_y/2-J_e\). A component formed from a block \(A\)
contributes \(|A|\sum_{a\in A}v_a\) to \(V_{xy}\).
V64's local Ewens law gives the mean block size of
a specified outside path as \(1+(k-1)/3=(k+2)/3\).
Thus
\[
 \E[V_{xy}\mid\mathcal F_x]
       =\frac{k+2}{3}\left(\frac{N_y}{2}-J_e\right).
\]
The formula is also valid at \(k=0\), when the
parenthesis is zero. Substitution into
\eqref{r67:colour-average} cancels \(N_x+4=2(k+2)\)
and gives exactly
\[
 \frac14+\frac{3(N_y/2-J_e)}{2(N_y+4)}
            =1-\frac{3(J_e+2)}{2(N_y+4)}.
\]
Every count mark is measurable in \(\mathcal F_x\).
Both averages therefore preserve
\eqref{r67:generating}. Specialize \(\lambda_e=\gamma\)
to prove \eqref{r67:factorial-eq}.
The range follows from \(0\le J_e\le N_y/2\);
the mean identity follows from \(G=1\) and the
factorial partition derivative. Applying the
count-mark identity with either endpoint proves the
last assertion, including centering.
\end{proof}

This calculation has removed genuine colour and
partition noise, not replaced the four-colour system
by independent tangent Gaussians. The intermediate
\(\mathcal R_e\) can be negative. The final
\(1-\chi_e\) is positive and has the exact count-source
shift required for detailed balance.

\subsection{Pinning one spin preserves the avoiding-loop observable}

Fix the root spin \(z_x=e_1\) and integrate all other
spins with the original scalar couplings. Common
rotational invariance makes this pinned partition
function equal to the unpinned \(Z_+\). In the diagram
expansion, do not pair the ports at \(x\). Outside
paths ending there have forced coordinate colour one
and weight one. Closed components avoiding \(x\)
retain fugacity four.

This pinned diagram law is exactly the marginal
obtained by summing the uncoloured pairing at \(x\).
Indeed V64's local matching normalizer is
\(2^k(k+1)!\), and its product with
\(u(2k)=1/[2^k(k+1)!]\) is one. All other factors,
including occurrence factorials and the avoiding-loop
fugacities, are unchanged. Thus expectations of
\(J_e\) and other \(\mathcal F_x\)-measurable marks
are preserved. No rare-event conditioning, additional
pinning potential or change of compact normalizer
has occurred.

\begin{theorem}[Uniform exponential occupation of loops avoiding a neighbour]
\label{r67:occupation}
Let \(N\ge4\) be even and \(\gamma>0\). With
\[
 \alpha=\frac32\log\frac{61}{60},
\]
the original cubic wire law satisfies, uniformly in
the volume and the ordered neighbouring pair,
\[
 \E_{\rm w}e^{\alpha J_e}\le e^{10/27}.
 \tag{V67.10}\label{r67:occupation-eq}
\]
\end{theorem}
\begin{proof}
Work in the exactly equivalent root-pinned law.
At \(y\) alone, multiply the three transverse
coordinates in every incident coupling by \(s\ge1\).
Denote the resulting compact integral by \(Z_s\).
The coupling to the root is unchanged because
\(z_x^\perp=0\). Colour expansion gives the exact ratio
\[
 \frac{Z_s}{Z_1}
  =\E_{\rm w}\prod_{C:\,n_x(C)=0}
                 \frac{1+3s^{2n_y(C)}}4.
 \tag{V67.11}\label{r67:pin-generator}
\]
An open path is forced to colour one and receives
no factor. Each visit of a transverse closed component
to \(y\) supplies two factors \(s\). This proves the
generator with no missing endpoint or length factor.
Weighted arithmetic--geometric mean implies
\[
 \prod_{C:\,n_x(C)=0}\frac{1+3s^{2n_y(C)}}4
       \ge s^{(3/2)J_e}.
\]

Put \(\eta=s-1\). The change in the spin exponent is
\(\gamma\eta\sum_{z\sim y,\ z\ne x}
z_y^\perp\cdot z_z^\perp\). Projection and the
two-edge path \(x,y,z\) give
\[
 \begin{aligned}
 z_y^\perp\cdot z_z^\perp
 &\le\tfrac12\{|z_y-z_x|^2+|z_z-z_x|^2\}\\
 &\le\tfrac32|z_y-z_x|^2+|z_z-z_y|^2 .
 \end{aligned}
\]
There are five terms. For \(d_f=\gamma(1-z_i\cdot z_j)\),
the exponent change is bounded above by
\[
 \eta\left(15d_e+
                2\sum_{\substack{f\ni y\\f\ne e}}d_f\right).
 \tag{V67.12}\label{r67:star-domination}
\]
The right side is invariant under common rotations,
so its pinned expectation equals its unpinned
expectation. V50's native quadratic bound applies
to these six nonnegative edge marks.

Take \(\eta=1/60\). The largest mark is \(1/4\);
their sum is \(5/12\). For
\(K_{\rm mark}=\operatorname{diag}(\sqrt t)\,
R\,\operatorname{diag}(\sqrt t)\), one has
\(\|K_{\rm mark}\|\le1/4\) and
\(\Tr K_{\rm mark}=5r/12\le5/36\).
Equation~\eqref{r50:trace} gives
\[
 \log(Z_s/Z_1)
       \le\frac{2}{1-1/4}\frac5{36}
       =\frac{10}{27}.
\]
Combine this with the generator and its pointwise
lower bound to obtain \eqref{r67:occupation-eq}.
\end{proof}

The result controls \emph{occupation at \(y\)}, not
the total lengths or diameters of the avoiding loops.
Those loops may still travel far from \(y\). The
single-root proof does not assert independence for
different roots or supply a many-source clustering
bound.

\subsection{The scaled rate deficit has a paid exponential moment}

\begin{theorem}[Uniform smallness of the averaged rate defect]
\label{r67:defect}
For all even \(N\ge4\) and \(\gamma\ge32\),
\[
 \E_{\rm w}e^{\alpha T_e}<3,\qquad
 \E_{\rm w}T_e^p\le3p!\alpha^{-p}
       \quad(p=1,2,\ldots).
 \tag{V67.13}\label{r67:defect-eq}
\]
In particular
\(\|\chi_e\|_p\le(3p!)^{1/p}/(\alpha\gamma)\)
uniformly in volume.
\end{theorem}
\begin{proof}
On \(N_y>2\gamma\), formula \eqref{r67:chi} gives
\(T_e\le3(J_e+2)/4\). Its exponential contribution
is at most \(e^{3\alpha/2+10/27}\), by
\eqref{r67:occupation-eq}.
On \(N_y\le2\gamma\), use \(T_e\le3\gamma/4\)
and V64's one-site bound
\(\mathbb P(N_y\le2\gamma)\le p_\gamma\), from
\eqref{r64:independent-tail}.
The explicit constants in \eqref{r64:constants}
and \(\log2<3/4\) imply
\[
 p_\gamma\le19e^{-\gamma/2},\qquad
 \alpha<1/40.
\]
Therefore
\[
 \E e^{\alpha T_e}
 \le e^{3\alpha/2+10/27}
              +19e^{-77\gamma/160}<3
       \quad(\gamma\ge32).
\]
For the last strict inequality,
\(3/80+10/27<1/2\) makes the first term less
than two, and \(77(32)/160>5\), together with
\(e^5>2^5>19\), makes the second less than one.
Finally \(u^p\le p!e^u\) proves the moment bound.
\end{proof}

\begin{corollary}[Uniform leading occupation mean]
\label{r67:mean}
As \(\gamma\to\infty\), uniformly over even \(N\ge4\),
\[
 \E T_e-3r\longrightarrow0,\qquad
 \E J_e-\left(2-\frac4{N^3}\right)\longrightarrow0.
 \tag{V67.14}\label{r67:mean-eq}
\]
\end{corollary}
\begin{proof}
The exact mean identity \eqref{r67:positivity} gives
\[
 \E T_e
    =\gamma\{1-\E_\mu(z_x\cdot z_y)^2\}
    =2\E_\mu d_e-\gamma^{-1}\E_\mu d_e^2 .
\]
V51's uniform mean estimate
\eqref{r51:local-mean} and
\(\E d_e^2\le32r^2\) imply
\(\E T_e=3r+O(\gamma^{-1})\), uniformly.

We next compare the occupation itself with the
regularized ratio, rather than replace a random
denominator by its mean without a bound. On
\(N_y\ge2\gamma\),
\[
 \left|T_e-\frac{J_e+2}{4}\right|
 \le\frac{J_e+2}{8\gamma}
                    |N_y-6\gamma+4|.
\]
The exponential occupation bound controls
\(\E(J_e+2)^2\) uniformly. The inherited
\eqref{r64:degree-concentration} bounds
\(\E(N_y-6\gamma)^2\) by \(128+6\gamma\).
Cauchy--Schwarz therefore gives an
\(O(\gamma^{-1/2})\) contribution, uniformly in volume.
On \(N_y<2\gamma\), use
\(J_e\le\gamma\) and \(T_e\le3\gamma/4\);
the contribution is at most
\((\gamma+1/2)p_\gamma\to0\).
Hence \(\E T_e-\E(J_e+2)/4\to0\).
Finally \(12r-2=2-4/N^3\), proving the second limit.
\end{proof}

\subsection{A bounded native count update and all factorial contacts}

Let \(\nu_\gamma\) be the exact uncoloured count marginal
and set
\[
 \bar\chi_e(K)=\E_{\rm w}[\chi_e\mid K],\qquad
 b_e(K)=\gamma^2(1-\bar\chi_e(K)).
 \tag{V67.15}\label{r67:count-rate}
\]
Taking indicator count marks in
\eqref{r67:factorial-eq} proves exact detailed balance:
\[
 \nu_\gamma(K)b_e(K)
   =\nu_\gamma(K+2\delta_e)(K_e+2)(K_e+1).
 \tag{V67.16}\label{r67:count-balance}
\]
Thus births \(K\mapsto K+2\delta_e\) with rate \(b_e\)
and deaths with rate \(K_e(K_e-1)\) preserve the
original count law and satisfy
\[
 \frac{\gamma^2}{4}\le b_e(K)\le\gamma^2,\qquad
 \E_{\nu_\gamma}
       e^{\alpha\gamma(1-b_e(K)/\gamma^2)}<3
                      \quad(\gamma\ge32).
 \tag{V67.17}\label{r67:rate-control}
\]
The last assertion is conditional Jensen applied to
\eqref{r67:defect-eq}. These rates have neither V66's
exponential whole-component clock nor a changed
invariant count measure. They are nonexplosive by
the bound \(m\gamma^2\) on total births and the
decrease of total count at deaths.

The count coefficient \(b_e(K)\) is a conditional
expectation and need not be spatially local as a
function of all uncoloured counts. In contrast the
coloured update \eqref{r67:colour-ratio} is explicitly
local. Both update types preserve parity sectors.
For the count update, any admissible edge-parity
sector is connected by removing pairs to its
zero/one representative and adding them back;
distinct sectors remain disconnected. One must not
replace the full invariant mixture by one sector,
delete winding parities, or infer a full-law gap.

For the signed source bank, let
\(\bar T=\E T_e\), independent of the selected edge.
Taking \(G=K_f\) in \eqref{r67:factorial-eq} and
subtracting means gives
\[
 F_{ef}
  =-\gamma\,\Cov_{\rm w}(T_e,K_f)
           +2\gamma^2\delta_{ef}(1-\bar T/\gamma).
 \tag{V67.18}\label{r67:F}
\]
The second term is the actual \(+2\) count-shift
contact. Therefore the full inherited bank becomes
\[
 \mathcal B_{\rm w}
  =(2\gamma I+3D)\Gamma
    +\gamma[\Cov_{\rm w}(T_e,K_f)]_{e,f}
    -2\gamma^2(1-\bar T/\gamma)I-6C^{\rm ret}.
 \tag{V67.19}\label{r67:full-bank}
\]
No factorial diagonal or nonlocal \(Q\)-origin contact
has been discarded: \(D=rI-R\circ R\) retains exactly
the V65 conversion of that term. \(C^{\rm ret}\)
still involves the original same-edge outside returns
\(H_e\), not the unstable raw rates \(a_e(N)\).

The new uniform bound is on the scaled scalar mark
\(T_e\) and its conditional mean, not on the entire
matrix in \eqref{r67:full-bank}. Cauchy--Schwarz for
each entry still loses powers of coupling, and
summing entrywise bounds loses source-support factors.
Conversely, converting the new covariance back into
spin-energy derivatives reproduces the old connected
problem; that conversion would not be an independent
estimate. The needed next step is spatial control of
the combined \(T_e,H_e\) response, or a quantitative
analysis of an adequate count-changing mechanism with
all parity sectors accounted for.

\subsection{Verification and the precise remaining obligation}

The exact rational checker compares 66,629 uncoloured
pairing diagrams with 5,463 coloured count
configurations in 85 anisotropic four-cycle profiles.
All partition coefficients agree. It checks 87,408
local colour double-link balance identities,
160 shifted factorial coefficients and their four
edge-source marks, including every diagonal shift.
The two-stage colour/pairing projection is checked
against the independently computed avoiding-loop
occupation. A separate local enumeration of 146,600
matchings checks the weighted mean block size.

The root-pinned expansion is checked independently
on 1,067 coloured configurations: both the original
partition coefficients and all 85 coefficients of
the transverse deformation at \(s=3/2\) agree with
the avoiding-loop generator. The proof uses
\(s=61/60\), not a numerical extrapolation from that
test. The Haar fourth-moment constants, the
arithmetic--geometric-mean comparison, and the
star trace and low-degree thresholds are checked
exactly. No Monte Carlo estimate, independent
Gaussian replacement or proof-assistant
formalization is asserted.

V67 pays uniform exponential control of a native
root-avoiding occupation and of the scaled positive
count-rate deficit. It supplies a bounded,
correctly normalized count update and preserves
all factorial source contacts. It does not yet
pay a volume-uniform signed operator norm for
\eqref{r67:full-bank}, generic transported-exterior
control, physical source coverage, physical-time
contraction or the interacting four-dimensional
continuum construction. The full Yang--Mills mass
gap remains unproved.

The V66 mathematical body is retained apart from the
version title; the other source files are unchanged.
Only the latest active \texttt{.tex} replaces its
predecessor in \path{latex_notes}. Build products
and checks remain elsewhere. The next companion
checkpoint and broad-literature sweep are both V70.
% END CONSOLIDATED V67 APPEND



% BEGIN CONSOLIDATED V68 APPEND
\clearpage
\section{V68: two-endpoint collapse and all-boundary count control}
\label{r68:context}

V67 paid an averaged exponential bound for the count-rate deficit.
The next dynamic estimate needs to survive conditioning on the other
counts. We therefore integrate the pairings at \emph{both} endpoints
of an edge. The remaining graph supplies a polynomial with
nonnegative coefficients. This gives an exact positive mixture of
parity-conditioned Poisson laws and uniform estimates for every
admissible exterior count configuration.

In particular, the low-count edge field is dominated by independent
Bernoulli variables with exponentially small parameter. This is
stronger than V64's independent-star intersection estimate. It also
pays products of the count-averaged rate deficits on arbitrary sets
of edges. It does not yet control their signed connected response:
the polynomial's coefficients still depend on outside paths.

\paragraph{Scope, prior work and targeted literature.}
The proof applies to the zero-field ferromagnetic \(O(4)\) count
law on any finite simple graph with positive scalar edge activities.
The cubic model here is the coherent conditional model of the
preceding notes, not an arbitrary transported Wilson exterior.
The sphere-moment/wire expansion is standard; Benassi--Ueltschi,
\href{https://arxiv.org/html/1807.06564v4}
{\emph{Loop correlations in random wire models}}, Section~6 and
Proposition~6.1, were rechecked with \(2J_e=\lambda_e\).
We derive the two-endpoint conditional reduction below.
V64's one-star mixture and V67's one-root occupation bound are not
being reproved. The new estimates are conditional on \emph{all
other edge counts}, and do not use the infrared bound.

A targeted search also screened the birth--death Stein approach
of Cloez--Delplancke,
\href{https://arxiv.org/abs/1609.08390}
{\emph{Intertwinings and Stein's magic factors for birth-death processes}}.
No spectral-gap or Stein-factor theorem from that paper is imported.
The official record of Lees's
\href{https://arxiv.org/abs/2205.12928}
{\emph{Griffiths inequalities for the \(O(N)\)-spin model}}
still reports withdrawal because the proposed map is not a bijection.
We use neither that switching assertion nor an unproved
\(O(4)\) correlation-positivity inequality.
These targeted checks do not replace the broad V70 review.

\subsection{The exact conditional polynomial}

Fix \(e=\{x,y\}\), write \(\gamma=\lambda_e>0\), and condition on
\(\widehat K=(K_f:f\ne e)\). Assume this exterior count vector has
positive probability. Let
\[
 a=\sum_{\substack{f\ni x\\f\ne e}}K_f,\qquad
 b=\sum_{\substack{f\ni y\\f\ne e}}K_f .
\]
All vertices other than \(x,y\) have even degree, and
\(a\equiv b\pmod2\). Put \(r=a\bmod2\). The allowed values of
\(k=K_e\) are exactly \(k\ge0\), \(k\equiv r\pmod2\).

Remove the occurrences of \(e\), pair the ports at vertices other
than \(x,y\), and leave the ports at \(x,y\) unpaired. Each such
outside pairing \(\pi\) has closed components and open paths with
endpoints among \(x,y\). Let \(L(\pi)\) be its number of closed
components, and \(h(\pi)\) the number of paths with one endpoint
at each of \(x,y\). Thus
\[
 0\le h(\pi)\le\min(a,b),\qquad h(\pi)\equiv r\pmod2.
\]
Paths returning to the same endpoint are retained. Define
\[
 w_h=\sum_{\pi:\,h(\pi)=h}4^{L(\pi)},\qquad
 P_{\widehat K}(t)=\sum_h w_h t^h .
 \tag{V68.1}\label{r68:polynomial}
\]
Factors from moments at other vertices and occurrence factorials
are independent of \(k\) and cancel from the conditional law.
There is at least one outside pairing, so \(P\) is not zero.

For the scalar product of two independent uniform \(S^3\) spins,
let
\[
 c_\ell=\int_{S^3\times S^3}(z_x\cdot z_y)^\ell\,d\sigma_xd\sigma_y,
 \qquad
 c_{2j}=\frac{(2j)!}{4^j j!(j+1)!},\quad c_{2j+1}=0.
 \tag{V68.2}\label{r68:moments}
\]

\begin{theorem}[Two-endpoint conditional collapse]
\label{r68:collapse}
With \(\widehat K\) fixed, the original uncoloured count law is
\[
 \nu_\gamma(K_e=k\mid\widehat K)
  =\frac{1}{Z_{\widehat K}(\gamma)}
       \frac{\gamma^k}{k!}\sum_h w_h c_{k+h},
 \qquad k\equiv r\pmod2,
 \tag{V68.3}\label{r68:conditional}
\]
where \(Z_{\widehat K}\) is the sum of the displayed numerator.
No graph-size or exterior regularity assumption is required.
\end{theorem}
\begin{proof}
Keep the occurrence labels fixed and expand each scalar product
in coordinate colours. At every vertex except \(x,y\), the
sphere-moment pairing formula contracts the indices along the
chosen outside paths and closed components. A closed component
has four possible colours. A path from \(x\) to \(y\) contributes
\(z_x\cdot z_y\); a path with both endpoints at \(x\), or both at
\(y\), contributes \(|z_x|^2=1\), or \(|z_y|^2=1\).
These last paths have not been discarded: their contractions
are exactly one.

The \(k\) occurrences of \(e\) supply
\((z_x\cdot z_y)^k\). Integrating the two remaining spins
therefore gives \(c_{k+h(\pi)}\). Summing \(\pi\) gives
\(\sum_h w_hc_{k+h}\); the edge factorial contributes
\(\gamma^k/k!\). This proves the coefficient identity before
normalization. Its parity and positivity show that every allowed
\(k\) has positive weight. Finally the compact exponential
integral, or \(c_{k+h}\le1\) and the finite polynomial, gives
absolute summability. Normalizing proves the theorem.
\end{proof}

Here \(h\) is a number of \emph{outside endpoint-to-endpoint paths},
not the return observable \(H_e\) of V65. Identifying these
different marks would lose the return term in the full bank.

\subsection{A pointwise deficit and the self-coordinate drift}

For a fixed full count vector, give \(h\) the posterior weights
\[
 \rho_k(h)=\frac{w_hc_{k+h}}{\sum_jw_jc_{k+j}},\qquad
 q_k=\sum_h\rho_k(h)\frac{k+h+1}{k+h+4}.
\]
The even-moment ratio in \eqref{r68:moments} implies
\[
 \frac{(k+2)(k+1)\nu(k+2\mid\widehat K)}
                  {\nu(k\mid\widehat K)}
       =\gamma^2 q_k .
 \tag{V68.4}\label{r68:ratio}
\]
Comparison with V67's exact count detailed balance identifies
\[
 \bar\chi_e(K)=3\sum_h\frac{\rho_k(h)}{k+h+4},
 \qquad
 \tau_e(K):=\gamma\bar\chi_e(K)
       \le \frac{3\gamma}{K_e+4}=:B_e(K_e).
 \tag{V68.5}\label{r68:pointwise}
\]
This is a pointwise bound on the \emph{count-averaged} deficit.
It is not a pointwise bound on V67's unaveraged \(T_e\).
In particular
\[
 \gamma^2\frac{k+1}{k+4}\le b_e(k,\widehat K)<\gamma^2 .
 \tag{V68.6}\label{r68:birth}
\]

\begin{proposition}[Monotone self-rate and a good-count drift margin]
\label{r68:drift}
With the exterior counts fixed, the admissible two-step increment
satisfies
\[
 0\le b_e(k+2,\widehat K)-b_e(k,\widehat K)
 \le\gamma^2\left\{
 \frac6{(k+4)(k+6)}+\frac9{4(k+4)(k+1)}\right\}.
 \tag{V68.7}\label{r68:self-increment}
\]
Consequently, if \(\gamma\ge32\) and \(k\ge\gamma/2\), then
\[
 \{(k+2)(k+1)-k(k-1)\}
   -\{b_e(k+2,\widehat K)-b_e(k,\widehat K)\}\ge\gamma .
 \tag{V68.8}\label{r68:drift-margin}
\]
\end{proposition}
\begin{proof}
Put \(v_h=(k+h+1)/(k+h+4)\) and let \(\E_k\) use \(\rho_k\).
The posterior at \(k+2\) is \(v_h\rho_k(h)/q_k\). Hence
\[
 q_{k+2}-q_k
 =\frac{\operatorname{Var}_k(v_h)}{q_k}
  +\frac{\E_k\!\left[
      v_h\,\frac6{(k+h+4)(k+h+6)}\right]}{q_k}.
\]
Both terms are nonnegative. The range of \(v_h\) has length at
most \(3/(k+4)\), so its variance is at most
\(9/[4(k+4)^2]\). Use \(q_k\ge(k+1)/(k+4)\) for the first
term and the maximum of the second fraction for the other.
This proves \eqref{r68:self-increment}.
For \(k\ge\gamma/2>0\), its right side is at most
\(33\gamma^2/(4k^2)\le33\).
The death-rate increment is \(4k+2\ge2\gamma+2\).
Their difference is at least \(2\gamma-31\ge\gamma\).
\end{proof}

This is a one-coordinate drift calculation, not a spectral gap.
Cross-coordinate increments of \(b_e\) still contain changes
of the outside polynomial. No locality of those changes has
been assumed.

\subsection{The positive parity-Poisson mixture}

For \(r\in\{0,1\}\), put
\[
 C_0(v)=\cosh v,\qquad C_1(v)=\sinh v,\qquad
 Q_{r,v}(k)=\frac{v^k}{k!\,C_r(v)}
              \mathbf1_{\{k\equiv r\pmod2\}} .
 \tag{V68.9}\label{r68:poisson}
\]
These are Poisson laws conditioned on the indicated parity.
At \(v=0\), use their limits \(\delta_0,\delta_1\).
They are reference kernels, not an assertion of independent
edge counts.

\begin{lemma}[Exact positive mixture]
\label{r68:mixture}
The conditional law \eqref{r68:conditional} can be sampled by
first choosing \(t\in(0,1)\) with density proportional to
\[
  \sqrt{1-t^2}\,P_{\widehat K}(t)\,C_r(\gamma t),
 \tag{V68.10}\label{r68:t-law}
\]
and then choosing \(K_e\) from \(Q_{r,\gamma t}\).
The distribution of \(t\) stochastically dominates the baseline
density proportional to \(\sqrt{1-t^2}\cosh(\gamma t)\).
\end{lemma}
\begin{proof}
For even \(\ell\), the coordinate density on \(S^3\) gives
\(c_\ell=(4/\pi)\int_0^1t^\ell\sqrt{1-t^2}\,dt\).
Insert this in \eqref{r68:conditional} and sum
\(\sum_{k\equiv r}(\gamma t)^k/k!=C_r(\gamma t)\).
The factor \(4/\pi\) cancels. This proves the mixture.

Relative to the baseline, the density multiplier is
\(P_{\widehat K}(t)\) for \(r=0\), and
\(P_{\widehat K}(t)\tanh(\gamma t)\) for \(r=1\).
It is nonnegative and increasing. For two increasing functions
\(f,g\) on an interval,
\[
 \Cov(f(t),g(t))
 =\frac12\E[(f(t)-f(t'))(g(t)-g(t'))]\ge0,
\]
where \(t'\) is an independent copy. Apply this with the
density multiplier and indicators of upper intervals to prove
the claimed stochastic order.
\end{proof}

The scalar \(t\) is an auxiliary variable of this positive
conditional expansion. It is not being identified with a
spin scalar product conditioned on signed monomials.

\begin{lemma}[Uniform endpoint deficiency]
\label{r68:W}
For every admissible \(\widehat K\) and \(\gamma\ge2\), the
mixture variable \(W=\gamma(1-t)\) satisfies
\[
 \E W<16,\qquad \E W^2<40,\qquad
 \E e^{W/2}<36 .
 \tag{V68.11}\label{r68:W-bounds}
\]
All expectations here concern the normalized conditional mixture.
\end{lemma}
\begin{proof}
By stochastic order it suffices to bound the baseline.
After setting \(w=\gamma(1-t)\), its density is proportional,
on \([0,\gamma]\), to
\(\sqrt{w(2\gamma-w)}\cosh(\gamma-w)\).
Its normalizer is at least
\[
 \int_0^1\sqrt{\gamma w}\,\frac{e^{\gamma-1}}2\,dw
           =\frac{\sqrt\gamma\,e^\gamma}{3e}.
\]
For \(p\ge0\) and \(0\le a<1\), its unnormalized
\(w^pe^{aw}\) integral is at most
\[
 \sqrt{2\gamma}\,e^\gamma
       \int_0^\infty w^{p+1/2}e^{-(1-a)w}\,dw .
\]
Their ratio is
\[
 \E[W^pe^{aW}]
 \le \frac{3e\sqrt2\,\Gamma(p+3/2)}
                         {(1-a)^{p+3/2}}.
 \tag{V68.12}\label{r68:W-general}
\]
The choices \((p,a)=(1,0),(2,0),(0,1/2)\) give respectively
\(9e\sqrt{2\pi}/4<16\), \(45e\sqrt{2\pi}/8<40\), and
\(6e\sqrt\pi<36\). For example the first two strict bounds
follow from \(e<11/4\) and \(\sqrt{2\pi}<18/7\).
\end{proof}

\subsection{Conditional Poisson fluctuations and a uniform low-count tail}

\begin{theorem}[Every-exterior count estimates]
\label{r68:fluctuations}
For every admissible exterior count vector and \(\gamma\ge4\),
\[
 \begin{aligned}
 |\E[K_e\mid\widehat K]-\gamma|&\le17,\\
 \gamma-17\le\operatorname{Var}(K_e\mid\widehat K)
                                      &\le\gamma+161,\\
 \|\nu_\gamma(K_e\in\cdot\mid\widehat K)-Q_{r,\gamma}\|_{\rm TV}
                                      &\le\min\{1,16/\sqrt\gamma\},\\
 \nu_\gamma(K_e\le\gamma/2\mid\widehat K)
                                      &\le q_\gamma,
 \end{aligned}
 \qquad q_\gamma:=\min\{1,38e^{-\gamma/24}\}.
 \tag{V68.13}\label{r68:conditional-bounds}
\]
\end{theorem}
\begin{proof}
Write \(m_r(v),v_r(v)\) for the mean and variance of \(Q_{r,v}\).
Differentiating its elementary normalizer gives
\[
 \begin{array}{ll}
 m_0(v)=v\tanh v,&v_0(v)=v\tanh v+v^2\operatorname{sech}^2v,\\
 m_1(v)=v\coth v,&v_1(v)=v\coth v-v^2\operatorname{csch}^2v .
 \end{array}
\]
For \(v\ge0\), including the limits at zero,
\[
 |m_r(v)-v|\le1,\qquad v-1\le v_r(v)\le v+1,\qquad
                    0\le m_r'(v)\le2 .
\]
For the variance bounds use \(\sinh v\ge v\),
\(\cosh v\ge v\) and the mean bounds. For the derivative,
\(m_0'\le1+v\operatorname{sech}^2v\le2\);
\(m_1'\le1\) follows from \(1-e^{-2v}\le2v\).
Nonnegativity also follows from \(m_r'=v_r/v\).

In the mixture, \(v=\gamma t=\gamma-W\). Thus the mean
error is at most \(1+\E W<17\). The conditional variance
decomposition gives a lower bound
\(\E v_r(v)\ge\gamma-\E W-1\), and an upper bound
\[
 \E v_r(v)+\operatorname{Var}(m_r(v))
 \le\gamma+1+\E|m_r(v)-m_r(\gamma)|^2
 \le\gamma+1+4\E W^2<\gamma+161.
\]

For total variation, the derivative of the mass function
\(Q_{r,v}\) has total variation norm
\[
 \frac{\E_{Q_{r,v}}|K-m_r(v)|}{2v}
       \le\frac{\sqrt{v+1}}{2v}.
\]
On \(v\in[\gamma/2,\gamma]\), this is at most
\(1/\sqrt\gamma\), since \(\gamma\ge4\).
Integration gives
\(\|Q_{r,v}-Q_{r,\gamma}\|_{\rm TV}
 \le(\gamma-v)/\sqrt\gamma\).
For \(v<\gamma/2\) the same upper bound is at least one,
so it remains valid. Average over the mixture and use
\(\E W<16\).

Finally separate \(t<3/4\), whose probability is at most
\(36e^{-\gamma/8}\). For \(t\ge3/4\), set
\(s=\log(3/2)\) and \(v=\gamma t\ge3\).
For either parity,
\[
 \E_{Q_{r,v}}e^{-sK}
     =\frac{C_r(ve^{-s})}{C_r(v)}
     \le2e^{-v(1-e^{-s})}.
\]
For the odd case the denominator correction is
\((1-e^{-2v})^{-1}<2\); the even case is immediate.
Exponential Markov inequality therefore gives
\[
 Q_{r,v}(K\le\gamma/2)
 \le2e^{s\gamma/2-v/3}
 \le2e^{-\gamma/24},
\]
using \(\log(3/2)<5/12\). Adding the exceptional
\(t\) probability proves the tail bound.
\end{proof}

The total-variation comparison keeps the parity \(r\).
Two exterior configurations imposing opposite parities
give disjoint supports, so their count conditionals
can have total variation one. Even within a parity sector,
the same-reference comparison only gives a per-comparison
bound \(32/\sqrt\gamma\); summing it over all edges is
not a dimension-free Dobrushin estimate.

\subsection{Bernoulli domination of actual low-count edges}

Return to the homogeneous cubic law and let
\(A_e=\mathbf1_{\{K_e\le\gamma/2\}}\).
This notation is local to this subsection; it is not
V66's whole-loop count birth intensity.

\begin{theorem}[Full-field low-count domination]
\label{r68:domination}
For \(\gamma\ge4\), the field \((A_e)\) is
stochastically dominated by independent Bernoulli variables
of parameter \(q_\gamma\). In particular, for every edge set \(S\),
\[
 \mathbb P(K_e\le\gamma/2\text{ for all }e\in S)
                                  \le q_\gamma^{|S|}.
 \tag{V68.14}\label{r68:joint-tail}
\]
If \(\gamma\ge216\) and \(\mathcal C(e)\) is the connected
low-count edge cluster containing \(e\), with adjacency
meaning a shared endpoint, then
\[
 \mathbb P(|\mathcal C(e)|\ge s)
       \le q_\gamma\,2^{-(s-1)},\qquad s\ge1.
 \tag{V68.15}\label{r68:cluster}
\]
The cluster is empty when \(e\) is not a low-count edge.
\end{theorem}
\begin{proof}
The last bound in \eqref{r68:conditional-bounds} conditions
on all other counts. Taking conditional expectations again
therefore bounds the chance of \(A_e=1\) given any set of
previously revealed other indicators by \(q_\gamma\).
Reveal edges in a fixed order. With independent uniform
random variables, realize each \(A_e\) by thresholding at
its conditional probability, and a Bernoulli variable by
thresholding the same uniform variable at \(q_\gamma\).
This constructs the asserted domination. Equivalently,
iterated conditioning directly proves \eqref{r68:joint-tail}.
No independence of the native counts is being asserted.

The line graph of the cubic torus has maximum degree ten.
For a fixed root it has at most \(10^{2(s-1)}=100^{s-1}\)
connected vertex sets of size \(s\): select a deterministic
rooted spanning tree and its depth-first traversal, a walk
of length \(2(s-1)\) that determines the visited set.
A connected cluster of size at least \(s\) contains one
such rooted set. Union bound and \eqref{r68:joint-tail}
give \(q_\gamma(100q_\gamma)^{s-1}\).
At \(\gamma\ge216\),
\(q_\gamma\le38e^{-9}<38/8000<1/200\);
the elementary bound \(e^3>20\) suffices.
This proves \eqref{r68:cluster}.
\end{proof}

Unlike the V64 star argument, this conclusion concerns
the whole edge field and is actual stochastic domination,
not merely an independent-set intersection estimate.
It still does not control the distances travelled by
outside paths through edges with large counts.

\subsection{Products of rate-deficit marks without a volume loss per mark}

\begin{theorem}[Conditional and joint positive deficit bounds]
\label{r68:marks}
For \(\gamma\ge216\), every edge satisfies
\[
 \E_{\nu_\gamma}
   \left[\left.e^{B_e(K_e)/36}\right|\widehat K\right]<2,
 \qquad
 \E_{\nu_\gamma}
   \left[\left.B_e(K_e)^p\right|\widehat K\right]
                      \le2p!\,36^p,\quad p\ge1.
 \tag{V68.16}\label{r68:conditional-marks}
\]
For every finite edge set \(S\), conditionally on the counts
outside \(S\) as well as unconditionally,
\[
 \begin{aligned}
 \E\exp\left\{\frac1{36}\sum_{e\in S}\tau_e(K)\right\}
                                  &\le2^{|S|},\\
 \E\prod_{e\in S}\tau_e(K)^{p_e}
                                  &\le\prod_{e\in S}2p_e!\,36^{p_e},
                      \qquad p_e\ge1 .
 \end{aligned}
 \tag{V68.17}\label{r68:joint-marks}
\]
In the conditional form, the expectations on the left are
understood to be conditional on \(K_{S^c}\).
\end{theorem}
\begin{proof}
On \(K_e>\gamma/2\), \(B_e\le6\). On the complementary
event, \(B_e\le3\gamma/4\). The uniform conditional tail gives
\[
 \E[e^{B_e/36}\mid\widehat K]
 \le e^{1/6}+38e^{-\gamma/48}<2 .
\]
Indeed \(e^{1/6}<6/5\), and for \(\gamma\ge216\) the
second term is at most \(38e^{-9/2}<38/80\).
Here \(e^{9/2}>80\) follows from \(e^3>20\) and
\(e^{3/2}>4\). The moment bound follows from
\(u^p\le p!e^u\).

For a product of the local \(B_e(K_e)\) functions, condition
on all other counts, remove one factor by
\eqref{r68:conditional-marks}, and iterate. The same argument
works with the counts outside \(S\) fixed.
Finally use the pointwise comparison \(0\le\tau_e\le B_e\).
This last step is essential: the nonlocal functions
\(\tau_f(K)\), \(f\ne e\), themselves cannot be pulled out
of a conditional expectation over \(K_e\).
\end{proof}

These are positive product bounds with a constant per marked
edge. They are not factorization identities, independence,
or centered covariance estimates. In particular they do not
bound the variance of an arbitrary normalized signed sum
uniformly in its support.

\subsection{What this changes in the full source problem}

Since \(\tau_e=\E[T_e\mid K]\), all count-connected pairings
in V67 are unchanged when \(T_e\) is replaced by \(\tau_e\).
Thus the exact bank is still
\[
 \mathcal B_{\rm w}
 =(2\gamma I+3D)\Gamma
   +\gamma[\Cov(\tau_e,K_f)]_{e,f}
   -2\gamma^2(1-\E\tau_e/\gamma)I-6C^{\rm ret}.
 \tag{V68.18}\label{r68:bank}
\]
The factorial \(+2\) contact and the \(D\Gamma\) term are
retained. The outside-path number \(h\) has not replaced
\(H_e\) in \(C^{\rm ret}\).

We now have an all-boundary one-edge fluctuation scale,
small low-count clusters, a self-coordinate drift margin
away from those clusters, and positive joint control of
the count-averaged deficit marks. These are concrete inputs
for a count-changing argument. The missing spatial estimate
is the response of the outside polynomial, or an equivalent
connected quantity, when another edge is changed.

Small exceptional clusters alone do not supply that estimate:
\(b_e\) is a marginal rate and its outside paths can reach
remote edges through the high-count region. Nor does a
one-coordinate drift calculation give a many-coordinate
Poincare inequality. Parity sectors, including winding
parities, remain in the original mixture. No auxiliary
dynamic estimate is being called the physical Hamiltonian
mass gap.

\subsection{Verification and preservation}

The exact rational checker independently compares full
pairing sums with the two-endpoint collapse on 85 four-cycle
profiles through eight occurrences and 133 \(K_4\) profiles
through six occurrences. It checks 69,802 full pairing
diagrams, 72,914 outside pairing configurations, all 370
available shifted coefficients, and equality with V67's
count-averaged deficit. A separate 300-case exact posterior
test checks the rate increment and its constants, including
widely separated outside-path counts. Deterministic truncated
series check twenty conditional distributions; those numerical
checks are not the proofs of the uniform estimates.

The uniform bounds and Bernoulli coupling are proved above.
No Monte Carlo estimate or proof-assistant formalization is
claimed. The V67 mathematical body is preserved apart from
the version title, and the other note files are unchanged.
Only the latest active \texttt{.tex} replaces its predecessor
in \path{latex_notes}; build products and diagnostics stay
outside that folder. The next companion review and broad
literature sweep are both V70.

The full signed spatial operator bound in \eqref{r68:bank},
generic transported-exterior control, physical source
coverage, physical-time contraction and interacting
four-dimensional continuum construction remain unproved.
The full Yang--Mills mass gap is still the objective,
not a conclusion of this appendix.
% END CONSOLIDATED V68 APPEND

% BEGIN CONSOLIDATED V69 APPEND
\clearpage
\section{V69: Four-terminal reduction, signed cross-response, and a native positivity obstruction}
\label{r69:main}

\subsection{Context, result ledger, and the changed direction}

V68 paid the all-exterior \emph{one-edge} conditional law, its
positive parity-Poisson mixture, and quantitative large-count
control. It did not pay the change of that conditional when a
second edge changes. This appendix examines that missing step.
All statements concern the same positive scalar ferromagnetic
\(O(4)\) count law; they are not statements for arbitrary
transported Wilson exteriors.

The four-terminal reduction itself is an application of the
sphere-pairing framework already used in V68. Its new use here
is an explicit, exactly evaluated two-edge conditional that
occurs on every even cubic torus of side at least four. It
has three features:
\begin{itemize}
\item every admissible count coefficient is strictly positive;
\item a finite moment certificate excludes a positive mixture
of two independent even-Poisson kernels, even with unrestricted
nonnegative intensities;
\item its cross-birth increment and conditional count covariance
are strictly negative, with evaluated large-count and large-activity
decay respectively.
\end{itemize}
Thus a positive one-edge mixture does not justify a positive
joint mixture or coordinatewise attractiveness. We retain
the signed response rather than trying to prove that it has
the wrong sign. The quantitative bounds below are for the
specified conditional, not for every exterior. The full
spatial operator estimate and physical mass gap remain open.

\subsection{Four-terminal collapse with the signs retained}

Write \(\sigma_{ij}=z_i\cdot z_j\) for unit spins in \(S^3\),
and let \(dH\) denote normalized product Haar measure. For
two disjoint selected edges \(e=12\), \(f=34\), fix all other
counts, denoted by \(\widehat K\). Remove the selected
occurrences and pair the remaining occurrence ends at
every vertex other than \(1,2,3,4\). A resulting open path
either has two distinct terminal endpoints or returns to
the same terminal. Let \(h_{ij}\) count the first type
with endpoints \(i,j\), and let \(c\) count closed components
not meeting a terminal.

\begin{lemma}[Exact conditional four-terminal kernel]
\label{r69:collapse}
For a nonzero-probability exterior \(\widehat K\), put
\[
 A_h(k,l)=\int \sigma_{12}^{k}\sigma_{34}^{l}
              \prod_{1\le i<j\le4}\sigma_{ij}^{h_{ij}}\,dH.
 \tag{V69.1}\label{r69:Ah}
\]
There are nonnegative finite weights \(w_h\), independent of
\(k,l\), for which the two selected count law is
\[
 \nu(K_e=k,K_f=l\mid\widehat K)
 \ \propto\ \frac{\gamma_e^k\gamma_f^l}{k!\,l!}
                         \sum_h w_h A_h(k,l).
 \tag{V69.2}\label{r69:conditional-general}
\]
The support has the two parities imposed by the exterior
degrees. Each admissible \(A_h(k,l)\) is positive, but its
spin integrand need not be nonnegative.
\end{lemma}
\begin{proof}
At a nonterminal vertex of degree \(d\), the coordinate
moment is the sum over pairings, with the common positive
factor \(u(d)=1/[2^{d/2}(d/2+1)!]\). These degrees are fixed
by \(\widehat K\). Contract the color indices along each
outside path. Distinct terminal endpoints give
\(\sigma_{ij}\); coincident endpoints give \(|z_i|^2=1\);
each closed component gives the color dimension four.
Take \(w_h\) to be the sum of \(4^c\) over the outside
pairings with the indicated \(h\). The product of the
nonterminal \(u(d)\)'s, and the factorials and activities
of the fixed occurrences, are common factors and cancel
in the conditional normalization.

The remaining integral is exactly \eqref{r69:Ah}. A terminal
degree that is odd makes it vanish. In an admissible parity
class all terminal degrees are even. Expanding scalar
products into colors and applying the same sphere formula
at the four terminals expresses the integral as a nonempty
sum of positive pairing weights. This proves positivity
of its value, not of its integrand. This is an occurrence
contraction, with no quotient by primitive loops: periodic
and winding components have not been discarded.
\end{proof}

\subsection{An exactly evaluated signed four-spin kernel}

Choose the cross-path multiplicities
\[
 h_{13}=h_{14}=h_{23}=h_{24}=1,\qquad h_{12}=h_{34}=0,
 \qquad M=\sigma_{13}\sigma_{14}\sigma_{23}\sigma_{24}.
 \tag{V69.3}\label{r69:crosscycle}
\]
Let
\[
 a_{kl}=\int \sigma_{12}^k\sigma_{34}^l M\,dH,
 \quad
 c_{2j}=\frac{(2j)!}{4^j j!(j+1)!},
 \quad r_k=\frac{k+1}{k+4}.
 \tag{V69.4}\label{r69:coefficients}
\]
Here and below \(k,l\) are nonnegative even integers.
The symbol \(r_k\) is the ratio \(c_{k+2}/c_k\), not the
edge resistance \(r\) of earlier appendices.

\begin{proposition}[Signed scalar kernel and positive coefficients]
\label{r69:kernel}
Under independent Haar spins, put \(s=\sigma_{12}\),
\(t=\sigma_{34}\). These two scalar products are independent,
each with density \((2/\pi)\sqrt{1-s^2}\) on \([-1,1]\), and
\[
 \E[M\mid s,t]
   =g(s,t):=\frac{4s^2+4t^2+2s^2t^2-1}{72}.
 \tag{V69.5}\label{r69:g}
\]
Consequently, with \(F(x,y)=4x+4y+2xy-1\),
\[
 a_{kl}=\frac{c_kc_l}{72}F(r_k,r_l)>0.
 \tag{V69.6}\label{r69:aformula}
\]
In particular,
\[
 a_{00}=\frac1{64},\qquad
 a_{20}=a_{02}=\frac1{128},\qquad
 a_{22}=\frac7{2304}.
 \tag{V69.7}\label{r69:small}
\]
\end{proposition}
\begin{proof}
Identify \(S^3\) with unit quaternions, using
\(p\cdot q=\operatorname{Re}(\bar p q)\). A common left
rotation sets \(z_1=1\); write \(z_2=u,z_3=w,z_4=wv\).
Haar invariance makes \(u,w,v\) independent unit Haar
quaternions. Thus \(s=\operatorname{Re}u\) and
\(t=\operatorname{Re}v\) are independent coordinate
variables. Conditional on \(s,t\), the imaginary axes
of \(u,v\) are independent uniform directions in \(S^2\).
Their dot product \(\zeta\) satisfies \(\E\zeta^2=1/3\).

The four factors of \(M\), as linear functions of \(w\),
have coefficient vectors \(1,\bar v,u,u\bar v\). The
fourth coordinate moment of uniform \(w\in S^3\) is the
sum of the three pair contractions divided by \(4\cdot6\).
The contraction sum is
\[
 t^2+s^2+(st+b\zeta)(st-b\zeta),
 \qquad b^2=(1-s^2)(1-t^2).
\]
For example the last two dot products are
\(1\cdot(u\bar v)=st+b\zeta\) and
\(\bar v\cdot u=st-b\zeta\); the other pair products
are \(t^2\) and \(s^2\). Average \(\zeta^2\) and divide
by 24 to obtain \eqref{r69:g}. The endpoints \(s,t=\pm1\)
follow by continuity, or may be omitted as null sets.

Integrate \(s^kt^l g(s,t)\) against the independent
coordinate densities. The even moments are \(c_k,c_l\),
and their two-step ratios are \(r_k,r_l\). This proves
\eqref{r69:aformula}. Since \(1/4\le r_k,r_l<1\) and
\(F\) is increasing in each coordinate on this square,
\(F(r_k,r_l)\ge F(1/4,1/4)=9/8\). Hence the coefficients
are strictly positive. Substitution gives \eqref{r69:small}.
\end{proof}

The value \(g(0,0)=-1/72\) shows why it would be invalid
to call the scalar reduction a positive spin conditional.
Nevertheless \eqref{r69:aformula} defines positive count
weights. These two facts are consistent and must not
be conflated.

\subsection{The kernel occurs in the actual cubic count law}

\begin{proposition}[A positive-probability cubic exterior]
\label{r69:embedding}
For every even torus side \(N\ge4\), the count law has a
positive-probability exterior event on which
\[
 p_{\gamma_e,\gamma_f}(k,l)
  =\frac{\gamma_e^k\gamma_f^l a_{kl}}
         {k!\,l!\,Z_2(\gamma_e,\gamma_f)},\qquad k,l\text{ even},
 \quad
 Z_2(g,h)=\sum_{k,l\ {\rm even}}\frac{g^kh^l a_{kl}}{k!\,l!}.
 \tag{V69.8}\label{r69:p}
\]
Both selected edges are nearest-neighbor edges. This
statement applies in particular to homogeneous activity.
\end{proposition}
\begin{proof}
Take terminals
\[
 1=(0,0,0),\quad2=(1,0,0),\quad
 3=(0,1,0),\quad4=(1,1,0),
\]
and select edges \(12,34\). Give count one to the edges
of the following four paths, with coordinates modulo \(N\):
\[
 \begin{aligned}
 13:\;&(0,0,0)\to(0,1,0),\\
 24:\;&(1,0,0)\to(1,1,0),\\
 14:\;&(0,0,0)\to(0,0,1)\to(1,0,1)
                       \to(1,1,1)\to(1,1,0),\\
 23:\;&(1,0,0)\to(1,0,-1)\to(0,0,-1)
                       \to(0,1,-1)\to(0,1,0).
 \end{aligned}
 \tag{V69.9}\label{r69:paths}
\]
The six internal vertices are distinct and avoid the
terminals, since \(1\) and \(-1\) are distinct modulo
\(N\ge4\). No path uses a selected edge. Fix count zero
on every other nonselected edge.

Each internal vertex has degree two. Integrating it
uses
\(\int(z_a\cdot z)(z\cdot z_b)\,dH(z)
  =(z_a\cdot z_b)/4\).
The six integrations give \(4^{-6}M\).
Thus the full moment for selected counts \(k,l\) is
\(4^{-6}a_{kl}\); the common outside activity product
and this factor cancel. Each terminal has outside degree
two, so the selected counts must be even. Formula
\eqref{r69:p} follows. Its normalizer is finite because
\(|a_{kl}|\le1\), and positive because \(a_{00}>0\).
The full finite-graph partition function is finite and
all activities are positive, so this specified exterior
has strictly positive probability.
\end{proof}

There is no assertion of a lower probability bound
uniform in \(N\) or activity. This event is sufficient
to refute an \emph{every-exterior} positivity claim,
but cannot by itself determine a typical-environment
spatial estimate.

\subsection{A finite certificate against a joint positive Poisson mixture}

Recall \(Q_{0,\lambda}(k)=\lambda^k/[k!\cosh\lambda]\)
for even \(k\), with its usual limit at zero.

\begin{theorem}[Failure of the joint independent-intensity ansatz]
\label{r69:no-mixture}
For every \(\gamma_e,\gamma_f>0\), the law
\eqref{r69:p} is not a positive probability mixture of
\(Q_{0,\lambda}\otimes Q_{0,\eta}\), even if the two
nonnegative finite intensities \(\lambda,\eta\) are
allowed to be dependent and unbounded.
\end{theorem}
\begin{proof}
First record the following finite moment calculation:
\[
 \frac1{a_{00}}\sum_{i,j=0}^3(-1)^{i+j}
             \binom3i\binom3j a_{2i,2j}
                       =-\frac{49}{1024}.
 \tag{V69.10}\label{r69:witness}
\]
For a direct proof, under the unweighted coordinate
density let \(m_p=\E(1-s^2)^p=(3/2)_p/(2)_p\).
Here \((a)_p=\prod_{j=0}^{p-1}(a+j)\).
The same beta integral gives
\[
 \frac{\E[s^2(1-s^2)^p]}{m_p}=\frac1{2(p+2)}.
\]
By \eqref{r69:g} the unnormalized double integral with
the extra factor \((1-s^2)^p(1-t^2)^p\) is
\[
 \frac{m_p^2}{72}
 F\left(\frac1{2(p+2)},\frac1{2(p+2)}\right).
\]
At \(p=3\), \(m_3=35/64\) and \(F(1/10,1/10)=-9/50\).
The value is \(-49/65536\); division by \(a_{00}=1/64\)
and binomial expansion give \eqref{r69:witness}.

Suppose the claimed mixture existed and write
\(\lambda=\gamma_e s,\eta=\gamma_f t\), so its mixing
law \(\Pi\) is on \([0,\infty)^2\). Comparing all even
count coefficients gives the positive finite measure
\[
 d\mu(s,t)=
 \frac{Z_2(\gamma_e,\gamma_f)}
      {\cosh(\gamma_e s)\cosh(\gamma_f t)}\,d\Pi(s,t),
 \qquad
 \int s^{2i}t^{2j}\,d\mu=a_{2i,2j}.
 \tag{V69.11}\label{r69:momentmeasure}
\]
Its total mass is \(a_{00}\). Moreover \(a_{2i,0}\le1\)
and \(a_{0,2j}\le1\), by their original bounded-spin
integrals. For any \(R>1\), positivity gives
\(\mu(s\ge R)\le R^{-2i}a_{2i,0}\le R^{-2i}\).
Letting \(i\to\infty\), and likewise for \(t\), proves
that \(\mu\) is supported on \([0,1]^2\).
The nonnegative polynomial
\((1-s^2)^3(1-t^2)^3\) must have nonnegative integral
there. Equation \eqref{r69:witness} says its integral
is strictly negative, a contradiction.
\end{proof}

Negative covariance alone would not prove this theorem:
dependent mixing intensities can themselves have negative
covariance. The finite nonnegative-polynomial witness is
the needed additional argument. Nor does the theorem
exclude all possible positive representations of the
count law. It excludes precisely the joint mixture of
independent even-Poisson kernels. Each single-edge
conditional still has the positive representation of V68.

\subsection{Fixed padding: what the literature does and does not supply}

The coefficients also give the exact determinant
\[
 a_{22}a_{00}-a_{20}a_{02}=-\frac1{73728},
 \qquad
 \frac{a_{22}}{a_{00}}
 -\frac{a_{20}a_{02}}{a_{00}^2}=-\frac1{18}.
 \tag{V69.12}\label{r69:padded}
\]
Thus the fixed-padding inequality
\(\langle M\sigma_{12}^2\sigma_{34}^2\rangle
\langle M\rangle\ge
\langle M\sigma_{12}^2\rangle
\langle M\sigma_{34}^2\rangle\)
is false. The padding has even degree at every vertex,
although its four edge exponents are odd.

This distinction is relevant to Abdesselam's
\emph{Non-Abelian correlation inequalities and stable
determinantal polynomials}, arXiv:2207.07603v2,
\url{https://arxiv.org/html/2207.07603}.
Equation (5) and Section 4, Problem 1 ask about this
fixed-padding positivity. Theorems 2.1 and 2.3 instead
provide large-power asymptotics and inequalities for
determinantal substitutes. They do not establish
finite-count fixed-padding positivity. The calculation
\eqref{r69:padded} is a negative instance for \(O(4)\),
with padding \eqref{r69:crosscycle} and increments
\(2e_{12},2e_{34}\).

This is not a counterexample to the unpadded GKS-2
inequality in an exponential ferromagnetic spin law:
\(M\) is signed, not a nonnegative change of spin density,
and the coefficientwise sums in that question are
different. No literature-priority claim is made.
We import the distinction between asymptotic determinant
positivity and native coefficient inequalities, not an
unproved positivity theorem.

\subsection{An evaluated cross-rate, with its sign and size}

For \eqref{r69:p}, keep V67's pair-death rate \(k(k-1)\).
The pair-birth rate of the first coordinate is
\[
 b_e(k,l)=\gamma_e^2\frac{a_{k+2,l}}{a_{kl}}
       =\gamma_e^2 r_k
                 \frac{F(r_{k+2},r_l)}{F(r_k,r_l)}.
 \tag{V69.13}\label{r69:birth}
\]
Define \(\delta_k=r_{k+2}-r_k=6/[(k+4)(k+6)]\).

\begin{proposition}[Strictly negative cross-increment]
\label{r69:rate}
For all admissible \(k,l\),
\[
 b_e(k,l+2)-b_e(k,l)
 =-\frac{18\gamma_e^2r_k\delta_k\delta_l}
             {F(r_k,r_{l+2})F(r_k,r_l)}<0.
 \tag{V69.14}\label{r69:crossrate}
\]
In particular,
\[
 |b_e(k,l+2)-b_e(k,l)|
 \le\frac{512\gamma_e^2}
              {(k+4)(k+6)(l+4)(l+6)}.
 \tag{V69.15}\label{r69:rate-bound}
\]
At homogeneous activity \(\gamma\), if \(k,l\ge\gamma/2\),
this is at most \(8192/\gamma^2\).
\end{proposition}
\begin{proof}
The bilinear polynomial obeys the exact identity
\[
 F(x',y')F(x,y)-F(x',y)F(x,y')
                         =-18(x'-x)(y'-y).
 \tag{V69.16}\label{r69:detidentity}
\]
Subtract the two ratios in \eqref{r69:birth} and apply
this identity. All denominators are positive, so it
also proves the sign. Using \(r_k\le1\),
\(F\ge9/8\), and the formulas for \(\delta_k,\delta_l\)
gives \(18\cdot36\cdot(8/9)^2=512\).
When both counts are at least \(\gamma/2\), the four
denominator factors have product at least
\(\gamma^4/16\), proving the last bound.
\end{proof}

This is a genuine count cross-response, rather than
an estimate of V66's unaveraged whole-loop clocks.
It refutes coordinatewise nondecreasing pair-birth
rates for every exterior, but it does not imply that
the magnitude is large. The opposite conclusion
holds in this example at large selected counts.
Detailed balance is not compromised: the rates still
obey the cycle identity
\[
 b_e(k,l)b_f(k+2,l)=b_f(k,l)b_e(k,l+2),
\]
because both products are
\(\gamma_e^2\gamma_f^2 a_{k+2,l+2}/a_{kl}\).

\subsection{The actual conditional count covariance}

Define the one-dimensional functions
\[
 \phi(g)=\sum_{k\ {\rm even}}\frac{g^kc_k}{k!},
 \qquad R(g)=\frac{\phi''(g)}{\phi(g)}.
 \tag{V69.17}\label{r69:phi}
\]
Equivalently, \(\phi(g)=\E\cosh(gs)\) under the unweighted
coordinate law. If \(t=|s|\), then \(R(g)=\E_g t^2\)
under the positive baseline density proportional to
\(\sqrt{1-t^2}\cosh(gt)\), \(0<t<1\).

\begin{theorem}[Negative two-edge covariance and its asymptotic size]
\label{r69:cov}
For the normalized conditional law \eqref{r69:p},
\[
 Z_2(g,h)=\frac{\phi(g)\phi(h)}{72}F(R(g),R(h)),
 \tag{V69.18}\label{r69:Z}
\]
and for \(g,h>0\),
\[
 \Cov(K_e,K_f\mid\widehat K)
 =-\frac{18ghR'(g)R'(h)}{F(R(g),R(h))^2}<0.
 \tag{V69.19}\label{r69:covformula}
\]
At \(g=h=\gamma\to\infty\), this becomes
\[
 \Cov(K_e,K_f\mid\widehat K)
                      =-\frac2{\gamma^2}+O(\gamma^{-3}).
 \tag{V69.20}\label{r69:covasy}
\]
The error concerns this fixed conditional family; it
does not claim a bound for arbitrary exterior data.
\end{theorem}
\begin{proof}
Substitute \eqref{r69:aformula} in \(Z_2\).
The two sums with \(c_{k+2}\) or \(c_{l+2}\) are
\(\phi''(g)\) or \(\phi''(h)\), giving \eqref{r69:Z}.
All series and their derivatives converge absolutely
on bounded activity sets because the spin variables
are bounded. Exponential-family differentiation gives
\(\Cov(K_e,K_f)=gh\,\partial_g\partial_h\log Z_2\).
Since \(F_{xy}F-F_xF_y=-18\), this is
\eqref{r69:covformula}.

For \(g>0\), differentiate the positive baseline to get
\[
 R'(g)=\Cov_g\bigl(t^2,t\tanh(gt)\bigr)>0.
 \tag{V69.21}\label{r69:Rprime}
\]
Both functions are strictly increasing on \((0,1)\);
the independent-copy covariance identity from V68,
and the strictly positive interval density, prove
strictness. Also \(R(g)\ge R(0)=1/4\), so
the denominator in \eqref{r69:covformula} is positive.

For completeness, the needed derivative asymptotic
is proved directly, not obtained by differentiating
an unspecified error term. Put \(W=g(1-t)\).
After removal of factors independent of \(w\), its
density is proportional to
\[
 w^{1/2}e^{-w}\sqrt{1-\frac{w}{2g}}\,
       (1+e^{-2g+2w})\mathbf1_{\{0\le w\le g\}}.
 \tag{V69.22}\label{r69:Wdensity}
\]
For every fixed nonnegative integer \(p\), its
\(w^p\)-weighted integral differs by \(O(g^{-1})\)
from the gamma integral \(\Gamma(p+3/2)\).
Indeed, on \(w\le g/2\) the square-root multiplier
differs from one by at most \(w/(2g)\), and the other
correction is at most \(e^{-g}\). On the complement
the density is at most \(2w^{1/2}e^{-w}\), and the
corresponding fixed polynomial tail is exponentially
small. The same argument for \(p=0\) controls the
normalizer away from zero. Therefore
\[
 \E_g W^p=\frac{\Gamma(p+3/2)}{\Gamma(3/2)}+O(g^{-1}),
 \qquad
 \E_g W=\frac32+O(g^{-1}),\quad
 \Var_g W=\frac32+O(g^{-1}).
 \tag{V69.23}\label{r69:Wmoments}
\]
In particular \(R(g)=\E(1-W/g)^2=1-3/g+O(g^{-2})\).
In \eqref{r69:Rprime}, replace \(t\tanh(gt)\) by \(t\).
For \(t\ge1/2\) the absolute error is at most \(2e^{-g}\).
For \(t<1/2\), the V68 baseline estimate
\(\E e^{W/2}<36\) bounds its probability by
\(36e^{-g/4}\). Since the other factor \(t^2\) is bounded,
the covariance error is exponentially small. Finally,
\[
 \Cov_g(t^2,t)=\frac{2\Var_g W}{g^2}
                       -\frac{\Cov_g(W^2,W)}{g^3}
                    =\frac3{g^2}+O(g^{-3}),
\]
using the bounded first three \(W\) moments.
Thus \(R'(g)=3/g^2+O(g^{-3})\). Substitution in
\eqref{r69:covformula}, with \(F(1,1)=9\), proves
\eqref{r69:covasy}.
\end{proof}

Unlike \eqref{r69:padded}, formula \eqref{r69:covformula}
is a covariance in a genuine positive probability law:
the \emph{count law conditioned on the specified exterior}.
It is not the unconditional spin-energy covariance.
The law of total covariance has an additional term
from the exterior-dependent conditional means; there
is no license to drop it.

\subsection{Consequences for the source bank and next obligation}

No identity in V65--V68 is replaced. In particular, with
the homogeneous activity convention and count-averaged
deficit \(\tau_e=\gamma\bar\chi_e\), the full bank remains
\[
 \mathcal B_{\rm w}=(2\gamma I+3D)\Gamma
       +\gamma[\Cov(\tau_e,K_f)]_{e,f}
       -2\gamma^2(1-\E\tau_e/\gamma)I-6C^{\rm ret}.
 \tag{V69.24}\label{r69:bank}
\]
The factorial contact, return contact, \(D\Gamma\), and
the mixture over parity and winding sectors are all
retained. The \(h_{ij}\)'s of the four-terminal reduction
are outside-path multiplicities, not replacements for
the return variable \(H_e\).

The failed route is now precise: V68's scalar positive
mixture cannot be promoted to a joint independent-intensity
mixture for all exteriors, nor can one presume that
changing another count increases a pair-birth rate.
The surviving direction is quantitative \emph{signed}
four-terminal response. The example proves that a
negative response can nevertheless be small at large
counts. It does not prove that all such responses are
small, summable over remote edges, or controlled by the
V68 low-count clusters.

The next upstream obligation is to control the
connected change of the full outside-path sum in
\eqref{r69:conditional-general}, with its normalization,
as exterior counts vary. A fixed four-spin evaluation
does not bound arbitrary terminal multiplicities or
their volume-dependent distributions. A usable estimate
must also survive the exterior-mixture contribution
and retain the return term in \eqref{r69:bank}.
This is where a signed expansion, rather than a blanket
correlation-positivity import, needs to earn its bound.

\subsection{Verification, scope, and preservation}

The exact rational checker independently compares
\eqref{r69:aformula} with 55,225 full pairing diagrams
over the nine cases \(k,l\in\{0,2,4\}\). It checks
25 quaternion contractions, the finite negative
certificate, seven beta-moment identities, 2,601
cross-rate cases and fifteen large-count bounds,
including detailed-balance cycle compatibility.
It verifies the cubic path embedding at sides
\(4,6,8,10\); the construction above proves it for
every \(N\ge4\). Seven deterministic double count-series
calculations agree with the one-dimensional covariance
formula, including unequal activities. These floating
checks are diagnostics, not proofs of the asymptotics.

The cited paper was checked specifically for fixed
padding versus asymptotic determinant claims. The
scheduled broad literature sweep and companion-program
review both remain due at V70. Existing scalar bounds
were used, not reissued as new results; the earlier
component-partition four-row algebra is a different
object from this four-spin count kernel.

The V68 mathematical body is preserved apart from the
version title. Only the active \texttt{.tex} file replaces
its predecessor in \path{latex_notes}; all other notes
are preserved and build products remain outside that
folder. These are rigorous conditional results and an
explicit obstruction to two proposed positivity
extensions. They do not construct the interacting
four-dimensional continuum theory or establish the
physical Yang--Mills mass gap.
% END CONSOLIDATED V69 APPEND

% BEGIN CONSOLIDATED V70 APPEND
\clearpage
\section{V70: Every-exterior conditional count gaps and two-edge operator mixing}
\label{r70:main}

\subsection{Checkpoint and the upstream question}

V69 ruled out a positive joint parity-Poisson representation
and attractive cross-birth rates, even for an actual cubic
exterior. Neither obstruction rules out quantitative mixing.
The question addressed here is whether the original
one- and two-coordinate count dynamics have an
\emph{every-exterior} coercivity bound. Unlike V69's
evaluated cross-cycle example, the theorems below hold
for arbitrary admissible exterior counts.

The result is a uniform one-edge auxiliary spectral gap
at least \(\gamma/8192\) for \(\gamma\ge216\), followed
by a two-edge auxiliary spectral gap at least
\(\gamma/16384\) for \(\gamma\ge4096\). We also obtain
an every-exterior maximal-correlation estimate for the
two selected counts. The constants are deliberately
coarse. They are proved for the specified count clocks,
not for the physical Yang--Mills Hamiltonian. No
many-edge gap or spatial summability is concluded.

The scheduled companion review freshly checked the
repository NS V26, SM--Einstein V72, YM Shell Mixing V16
and YM Parallel V5, together with downloaded Source
Selection V35, Native Stationarity V38, Exact Action
Einstein V28 and Perturbative ESM V31. All eight hashes
are unchanged from V65; their terminal result and scope
passages were reread. The ESM transfer memorandum was
also reread as research data, not as instructions.
No companion supplies the missing native spatial bound.
The useful ESM discipline remains to identify the
complete coupled quantity before declaring a sector
closed. Here that means retaining exterior-mean,
return, contact and parity contributions after proving
a conditional gap. No BV or perturbative-continuum
theorem is imported into the nonperturbative YM claim.

\subsection{Broad literature sweep and the selected import}

The V70 sweep covered birth--death resolvents, Gibbs
sampler operator correlations, spectral independence,
stochastic localization, low-temperature vector models,
and lattice/continuum Yang--Mills results. The following
are the relevant distinctions, with primary-source links.

\paragraph{One-dimensional and two-block tools.}
Liu--Ma, \emph{Spectral gap and convex concentration
inequalities for birth--death processes}
(\url{https://doi.org/10.1214/07-AIHP149}), studies
Poisson-resolvent and spectral-gap relations.
Cloez--Delplancke, \emph{Intertwinings and Stein's magic
factors for birth-death processes}
(\url{https://arxiv.org/abs/1609.08390}), supplies a
related weighted-gradient framework. These are method
context, not a precomputed estimate for our rates.
We instead give a direct log-concavity/Cheeger proof
including the low-count region.
The conditional-expectation operator viewpoint is
standard in Liu--Wong--Kong, \emph{Covariance structure
of the Gibbs sampler}
(\url{https://doi.org/10.1093/biomet/81.1.27}).
The total-variation-to-maximal-correlation estimate
needed here is proved below, including its
countably infinite state-space step.

\paragraph{Spatial operator estimates need not be positive row sums.}
Chen--Yang--Yin--Zhang,
\emph{Spectral Independence Beyond Total Influence
on Trees and Related Graphs}
(\url{https://arxiv.org/html/2404.04668}), Section 3,
Theorem 32, separates a positive approximate inverse
from the error in its product with the correlation
matrix. This is useful architecture for a signed
source bank. Its tree factorization, model-specific
local matrices and graph estimates are not available
for our count marginal. In particular, a covariance
conditioned on all other counts is not a principal
minor of the unconditioned covariance matrix.

\paragraph{Block factorization still consumes an interaction bound.}
Caputo--Chen--Parisi,
\emph{Factorizations of Relative Entropy Using
Stochastic Localization}
(\url{https://arxiv.org/html/2503.19419}), develops
factorization for high-temperature spin systems and
weakly interacting components. The finite-spin setup
and the intercomponent hypothesis cannot be obtained
by merely naming the present counts as blocks.
Blanca--Zhang's global-chain results
(\url{https://doi.org/10.1007/s00453-025-01353-5})
likewise require spectral independence and specified
marginal assumptions. None is inferred from two-edge
conditional estimates alone.

\paragraph{The physical target is unchanged.}
Giuliani--Ott's low-temperature vector-model expansion
(\url{https://arxiv.org/abs/2302.07299}) remains relevant
to the infrared/IBP route, with the finite-source versus
full-operator distinction already audited in V59--V60.
Cao--Nissim--Sheffield's dynamical area-law note
(\url{https://arxiv.org/abs/2509.04688}) has its own
lattice parameter regime. Chatterjee's
\(\mathrm{SU}(2)\) Yang--Mills--Higgs scaling limit
(\url{https://arxiv.org/abs/2401.10507}, v4) includes
a Higgs field and a massive Gaussian limit, not
interacting pure four-dimensional YM.
Nair's recent \emph{Towards a Proof of Mass Gap in
3d Yang-Mills Theory}
(\url{https://arxiv.org/abs/2608.10133}, v2)
was screened at its primary abstract: it describes
an outline in three dimensions, not the missing
four-dimensional construction. No result from its
proof is used here. Unverified broad mass-gap claims
surfaced by search were not accepted as closure.

Thus the selected immediate import is the operator
approach to conditional mixing, with all quantitative
inputs proved for the native rates. The approximate
inverse approach remains a distinct future spatial
lead, not an imported theorem for this model.

\subsection{The unchanged count law and inherited inputs}

Work on a finite simple graph with the positive
zero-field ferromagnetic \(O(4)\) count law of V68.
Other edge activities may be arbitrary positive numbers.
On a selected edge let its activity be \(\gamma\),
let \(r\in\{0,1\}\) be the parity imposed by the
other counts, and write \(k=r+2j\), \(j\ge0\).
For its normalized conditional mass \(p_j\), V68 gives
\[
 \frac{p_{j+1}}{p_j}
  =\frac{\gamma^2q_k}{(k+2)(k+1)},\qquad
 \frac{k+1}{k+4}\le q_k<1.
 \tag{V70.1}\label{r70:ratio}
\]
Every \(p_j\) is positive. The prescribed pair-birth
and pair-death rates are
\[
 b(k)=\gamma^2q_k,\qquad d(k)=k(k-1),
 \qquad \frac{\gamma^2}{4}\le b(k)\le\gamma^2.
 \tag{V70.2}\label{r70:rates}
\]
The bottom death rate is zero for either parity.
Detailed balance reads \(p_jb(k)=p_{j+1}d(k+2)\).

For \(\gamma\ge4\), the other inherited inputs are
\[
 \Var(K_e\mid K_{e^c})\le\gamma+161,\qquad
 \|\nu(K_e\in\cdot\mid K_{e^c})-Q_{r,\gamma}\|_{\rm TV}
                  \le\min\{1,16/\sqrt\gamma\}.
 \tag{V70.3}\label{r70:inputs}
\]
They hold for every admissible exterior count vector.
They do not describe an arbitrary transported Wilson
spin exterior. We use them without repeating their
proofs or relabelling them as new results.

\begin{proposition}[A strictly log-concave conditional count sequence]
\label{r70:logconcavity}
For every exterior, every activity and every \(j\ge0\),
\[
 \frac{p_{j+2}p_j}{p_{j+1}^2}
  \le\frac{k+2}{k+3}<1,\qquad k=r+2j.
 \tag{V70.4}\label{r70:curvature}
\]
\end{proposition}
\begin{proof}
Apply \eqref{r70:ratio} twice. Since
\(q_{k+2}/q_k\le(k+4)/(k+1)\), the left side is at most
\[
 \frac{k+4}{k+1}\,
 \frac{(k+2)(k+1)}{(k+4)(k+3)}
                  =\frac{k+2}{k+3}.
\]
The normalization cancels, so no estimate of the
exterior polynomial normalizer is needed.
\end{proof}

This log-concavity is in the parity-indexed count \(j\).
It is compatible with V69's negative cross-increment:
one-coordinate log-concavity is not multivariate
log-supermodularity or attractiveness.

\subsection{An explicit discrete Cheeger comparison}

We include the elementary one-dimensional argument to
make all constants and support assumptions visible.
It is a standard isoperimetric method, not a claim of
a new general Cheeger theorem.

\begin{lemma}[Log-concavity, median mass and a bounded-rate gap]
\label{r70:cheeger}
Let \(p=(p_j)_{j\ge0}\) be a strictly positive
log-concave probability with finite variance
\(\sigma^2>0\). Set \(c_j=\min(p_j,p_{j+1})\).
For every \(f\in L^2(p)\),
\[
 \Var_p f\le16(16\sigma+4)^2
                   \sum_{j\ge0}c_j(f_{j+1}-f_j)^2.
 \tag{V70.5}\label{r70:cheeger-poincare}
\]
An infinite right side makes the assertion automatic.
\end{lemma}
\begin{proof}
Choose an integer median \(m\), so that both
\(p([0,m])\) and \(p([m,\infty))\) are at least \(1/2\).
The interval \([\E j-2\sigma,\E j+2\sigma]\) has
probability at least \(3/4\) and at most \(4\sigma+1\)
integer points. Its intersection with each side of
the median has mass at least \(1/4\).
There are therefore indices \(a\le m\le b\) with
\(p_a,p_b\ge1/[4(4\sigma+1)]\).
Log-concavity between these indices gives
\[
 p_m\ge\frac1{16\sigma+4}.
 \tag{V70.6}\label{r70:median}
\]

Put \(F_j=\sum_{i\le j}p_i\). Log-concavity makes
both \(p_j/F_j\) and \(p_{j+1}/F_j\) nonincreasing:
expand \(F_j/p_j\) as a sum of products of backward
successive ratios, each nondecreasing as \(j\) increases;
then use the nonincreasing forward ratio for the second
assertion. The corresponding right-tail ratios are
nondecreasing, by the same argument in reverse.
Consequently the left half-line ratios
\(c_j/F_j\), \(j<m\), are bounded below by the one
at \(m-1\). If \(p_{m-1}\ge p_m/2\), this ratio
is at least \(p_m/2\). Otherwise put
\(v=p_{m-1}/p_m<1/2\). Backward log-concavity gives
\(F_{m-1}\le p_{m-1}/(1-v)\), and hence the ratio
is at least \(1-v>1/2\ge p_m/2\).
When \(m=0\) there is no left side.
The identical right-tail argument, using
\(p_{m+1}/p_m\) when it is below \(1/2\), proves
\[
 \frac{c_j}{\min\{F_j,1-F_j\}}\ge h,\qquad h:=p_m/2.
 \tag{V70.7}\label{r70:cut}
\]

This also bounds the boundary of every set \(A\)
with \(p(A)\le1/2\) by \(h\,p(A)\), not just half-lines.
Decompose \(A\) into intervals. For an interval
strictly to one side of \(m\), its boundary nearer
the median already controls its mass by
\eqref{r70:cut}. For an interval \(I\) containing
\(m\), the two boundary edges control
\(h(1-p(I))\ge h p(I)\). Missing boundary edges
at zero or infinity contribute zero. Summation
over the disjoint interval components is legitimate
by nonnegativity.

To obtain the \(L^2\) inequality, let \(a\) be a
median of \(f\) and set \(g_+=(f-a)_+\),
\(g_-=(a-f)_+\). Each positive support has mass
at most \(1/2\). Discrete coarea and the preceding
set bound give, for either \(g\),
\[
 h\sum_jp_jg_j^2
 \le\sum_jc_j|g_{j+1}^2-g_j^2|
 \le
 \left(\sum_jc_j(g_{j+1}-g_j)^2\right)^{1/2}
 \left(4\sum_jp_jg_j^2\right)^{1/2}.
 \tag{V70.8}\label{r70:coarea}
\]
For the last factor use
\((g_{j+1}+g_j)^2\le2(g_{j+1}^2+g_j^2)\)
and \(c_{j-1}+c_j\le2p_j\), with \(c_{-1}=0\).
Thus \(\E g^2\le4h^{-2}\sum c_j(\Delta g_j)^2\).
The two squared differences sum to at most
\((f_{j+1}-f_j)^2\). Also
\(\Var f\le\E(f-a)^2\). Combining these facts,
\(h=p_m/2\), and \eqref{r70:median} proves
\eqref{r70:cheeger-poincare}. Coarea and the
Cauchy--Schwarz step hold by monotone truncation
on this countable state space.
\end{proof}

\subsection{A uniform gap for the original one-edge count clock}

For the fixed exterior let
\[
 \mathcal E_e(f,f)=
    \sum_{j\ge0}p_j b(r+2j)
          \bigl(f(r+2j+2)-f(r+2j)\bigr)^2 .
 \tag{V70.9}\label{r70:one-form}
\]
This is the Dirichlet form of the pair-birth/death
generator with rates \eqref{r70:rates}.

\begin{theorem}[Every-exterior one-edge coercivity]
\label{r70:one-gap}
For every admissible exterior and \(\gamma\ge4\),
\[
 \mathcal E_e(f,f)\ge\lambda_1(\gamma)\Var_p f,
 \qquad
 \lambda_1(\gamma):=
 \frac{\gamma^2}{64(8\sqrt{\gamma+161}+4)^2}.
 \tag{V70.10}\label{r70:one-gap-formula}
\]
In particular \(\lambda_1(\gamma)\ge\gamma/8192\)
for \(\gamma\ge216\), uniformly in the graph,
exterior counts and imposed parity.
\end{theorem}
\begin{proof}
Apply Lemma \ref{r70:cheeger} to the parity index
\(j=(K_e-r)/2\), whose variance is at most
\((\gamma+161)/4\) by \eqref{r70:inputs}.
The lower birth bound implies
\[
 \mathcal E_e(f,f)\ge\frac{\gamma^2}{4}
          \sum_j c_j(f(r+2j+2)-f(r+2j))^2.
\]
This proves \eqref{r70:one-gap-formula}.
For \(\gamma\ge216\),
\(\gamma+161\le16\gamma/9\) and
\(4\le\sqrt\gamma/3\), so
\[
 8\sqrt{\gamma+161}+4\le11\sqrt\gamma,\qquad
 \lambda_1(\gamma)\ge\frac{\gamma}{7744}
                                      \ge\frac{\gamma}{8192}.
\]
\end{proof}

The jump process is nonexplosive: births have rate
at most \(\gamma^2\), while every death removes two
occurrences. There cannot be more deaths than the
initial parity-indexed count plus the number of
births. Detailed balance and the normalized conditional
law therefore define the conservative reversible
process; the displayed inequality holds on its
closed Dirichlet-form domain and is a spectral-gap
statement there. This proof covers the low-count
region. It does not assume that V68's good-count
drift bound holds down to the origin.

\subsection{Two conditional channels: total variation to operator angle}

Here is the operator step that does not require
positive cross-correlations or a joint Poisson mixture.

\begin{lemma}[A two-channel maximal-correlation bound]
\label{r70:angle}
Let \((X,Y)\) have a probability law on a product
of countable sets with positive marginals. Suppose
the conditional kernels \(A:X\to Y\), \(B:Y\to X\)
have total-variation diameters at most \(a,b\),
respectively. Their maximal correlation
\[
 \rho(X,Y):=
 \sup_{\substack{\E u(X)=\E v(Y)=0\\
                  \|u(X)\|_2=\|v(Y)\|_2=1}}
             |\E[u(X)v(Y)]|
 \tag{V70.11}\label{r70:rho}
\]
satisfies \(\rho\le\sqrt{ab}\).
Moreover every \(F\in L^2(X,Y)\) satisfies
\[
 \E\Var(F\mid X)+\E\Var(F\mid Y)
                    \ge(1-\rho)\Var F .
 \tag{V70.12}\label{r70:two-projection}
\]
The nontrivial use is when \(ab<1\).
\end{lemma}
\begin{proof}
Let \(T:L^2_0(X)\to L^2_0(Y)\) be
\(Tu=\E[u(X)\mid Y]\). Its adjoint is conditional
expectation in the other direction. Thus
\(K=T^*T\) is a positive self-adjoint Markov
operator on \(X\), with transition kernel \(AB\).
Total-variation contraction under a kernel gives
\(\delta(AB)\le\delta(A)\delta(B)\le ab\).
Equivalently, for any bounded \(u\),
\[
 \operatorname{osc}(K^n u)
                       \le(ab)^n\operatorname{osc}(u).
\]
If \(\E u=0\), stationarity implies
\(\|K^n u\|_\infty\le\operatorname{osc}(K^n u)\).
Consequently
\[
 0\le\langle u,K^n u\rangle
       \le\|u\|_1(ab)^n\operatorname{osc}(u).
\]
The spectral measure of a bounded centered \(u\)
cannot charge any interval \((ab+\varepsilon,1]\):
such mass would contradict this inequality as
\(n\to\infty\). Bounded centered functions are
dense in \(L^2_0(X)\), so \(\|K\|\le ab\).
Since \(\rho=\|T\|\), the first claim follows.

For the second, subtract \(\E F\), and put
\(u=\E[F\mid X]\), \(v=\E[F\mid Y]\).
These projections are centered. By maximal correlation,
\[
 \|u+v\|_2^2\le(1+\rho)(\|u\|_2^2+\|v\|_2^2),
 \qquad
 \langle F,u+v\rangle=\|u\|_2^2+\|v\|_2^2.
\]
Cauchy--Schwarz yields
\(\|u\|_2^2+\|v\|_2^2\le(1+\rho)\|F\|_2^2\).
Subtract this from \(2\Var F\) to obtain
\eqref{r70:two-projection}. The normalization is
two coordinate-update clocks of rate one, not a
single random scan clock of total rate one.
\end{proof}

\subsection{Every-exterior two-edge gap, including signed examples}

Select any two distinct edges \(e,f\) of the finite
simple graph and fix all other counts. The selected
two-edge graph is a forest. Its leaf parity equations
determine \(r_e,r_f\) uniquely; compatibility at
the remaining vertices follows from admissibility
of the exterior. All pairs of nonnegative counts
with these two parities have positive weights.
Thus its support is a product of two parity
lattices, even when the edges share a vertex.
Take the two selected activities to be \(\gamma\).

\begin{theorem}[Two-edge mixing with no sign assumption]
\label{r70:two-gap}
For every such exterior,
\[
 \rho(K_e,K_f\mid K_{\{e,f\}^c})
                  \le\varepsilon_\gamma
                  :=\min\{1,32/\sqrt\gamma\}.
 \tag{V70.13}\label{r70:native-rho}
\]
For \(\gamma>1024\), the sum of the two original
pair-birth/death generators has conditional gap at least
\[
 \frac{\gamma}{8192}
           \left(1-\frac{32}{\sqrt\gamma}\right).
 \tag{V70.14}\label{r70:two-gap-formula}
\]
In particular, for \(\gamma\ge4096\), with conditional
expectations under this two-edge law,
\[
 \begin{split}
 &\E\!\left[
 b_e(K)\bigl(F(K+2e)-F(K)\bigr)^2+
 b_f(K)\bigl(F(K+2f)-F(K)\bigr)^2\right]\\
 &\hspace{35mm}\ge\frac{\gamma}{16384}\Var F .
 \end{split}
 \tag{V70.15}\label{r70:two-form}
\]
The constants are independent of the graph, distance
between the edges, exterior counts and selected parities.
\end{theorem}
\begin{proof}
The conditional law of \(K_e\) for each fixed \(K_f\)
is within \(16/\sqrt\gamma\) in total variation
of the \emph{same} \(Q_{r_e,\gamma}\), by
\eqref{r70:inputs}. Hence the diameter of this
channel is at most \(\varepsilon_\gamma\).
The other channel has the same bound.
Lemma \ref{r70:angle} proves \eqref{r70:native-rho}.

Apply Theorem \ref{r70:one-gap} to \(F\) with one
selected count fixed, then integrate that estimate
over the other count, and do the same with the
roles reversed. The sum of the two Dirichlet forms
is at least
\[
 \frac{\gamma}{8192}
       \{\E\Var(F\mid K_e)+\E\Var(F\mid K_f)\}.
\]
Use \eqref{r70:two-projection} and
\eqref{r70:native-rho}. At \(\gamma\ge4096\),
\(\varepsilon_\gamma\le1/2\), giving
\eqref{r70:two-form}. The sum process is
nonexplosive by the same bounded-total-birth
argument as in the one-coordinate case.
\end{proof}

The bound on maximal correlation controls
\emph{arbitrary} centered functions of the separate
counts, not only linear covariance. It is an operator
bound for a normalized two-variable conditional law,
despite V69's negative covariance example.
It does not decay with spatial distance.
It also does not bound a function such as
\(\tau_e(K)\) by pretending that it depends only
on \(K_e\); that function generally depends on both
selected counts and on the exterior.

\begin{corollary}[A uniformly normalized two-edge Poisson solver]
\label{r70:solver}
Let \(L_{ef}\) be the sum generator in
\eqref{r70:two-form}, and \(\gamma\ge4096\).
For each centered \(G\in L^2\) of its conditional
law there is a unique centered solution
\(u=(-L_{ef})^{-1}G\) in the operator domain, with
\[
 \|u\|_2\le\frac{16384}{\gamma}\|G\|_2,\qquad
 \mathcal E_{ef}(u,u)
               \le\frac{16384}{\gamma}\|G\|_2^2.
 \tag{V70.16}\label{r70:resolvent}
\]
For \(G=\tau_e-\E[\tau_e\mid K_{\{e,f\}^c}]\),
one may take \(\|G\|_2\le72\).
\end{corollary}
\begin{proof}
The closed nonnegative self-adjoint operator
\(-L_{ef}\) has spectrum at least
\(\gamma/16384\) on centered \(L^2\).
Spectral inversion proves both inequalities, the
second using \(\mathcal E(u,u)=\langle G,u\rangle\).
V68's conditional positive mark bound, conditioned
on counts outside \(\{e,f\}\), gives
\(\E\tau_e^2\le4\cdot36^2=72^2\).
Centering cannot increase this squared norm.
\end{proof}

This is a Poisson equation for the auxiliary count
clock. It does not replace a physical-time source
equation. The source is the \emph{count-averaged}
deficit; V66's unaveraged whole-loop birth clocks
still have their excluded large second moments.

\subsection{What remains between these local bounds and the full bank}

The full homogeneous identity remains
\[
 \mathcal B_{\rm w}=(2\gamma I+3D)\Gamma
       +\gamma[\Cov(\tau_e,K_f)]_{e,f}
       -2\gamma^2(1-\E\tau_e/\gamma)I-6C^{\rm ret}.
 \tag{V70.17}\label{r70:bank}
\]
Neither the \(D\Gamma\) term nor the factorial and
return contacts has been removed. Our two-coordinate
solver does not estimate the change of that solver
or its conditional mean when the exterior varies.
For any pair of observables the total covariance
still contains
\[
 \Cov(F,G)=\E\Cov(F,G\mid\widehat K)
       +\Cov\bigl(\E[F\mid\widehat K],
                  \E[G\mid\widehat K]\bigr).
 \tag{V70.18}\label{r70:total-cov}
\]
The second term is a genuine remaining spatial object.

For \(M\) count coordinates, merely summing the
available channel diameter bound would cost
\((M-1)\varepsilon_\gamma\). At fixed large \(\gamma\)
and growing volume this is not a contraction bound.
The marginal count interactions are not restricted
to geometrically adjacent edges, so replacing \(M-1\)
by the cubic degree would be unjustified. Also,
two-variable \emph{conditionals} do not control all
two-variable \emph{marginals}; arbitrary unobserved
counts can transmit dependence.

Furthermore every pair-count move preserves each
edge's parity. A global sum of these generators has
all admissible parity sectors in its kernel on the
full count mixture. Even a future sector-uniform
many-edge estimate would require a separate argument
for between-sector covariance, including winding
parities. It would still be an auxiliary count
estimate, not a physical Hamiltonian mass gap.

The new result removes a concrete local obstruction:
neither arbitrarily bad one-coordinate conditioning
nor loss of two-coordinate coercivity at high activity
can be invoked to explain the missing spatial bound.
The next task is the exterior response of the
normalized conditional operator, or an approximate
inverse with a proved signed global error. The
literature sweep suggests the latter as a useful
alternative to absolute row sums; it does not
supply that error for us.

\subsection{Verification and preservation}

The rational checker tests 480 native conditional
ratio inequalities, including both parities and
widely separated outside-path multiplicities.
It verifies the median and full-set Cheeger bounds
on 27 rational log-concave laws, enumerating 3,091
nonempty subsets with mass at most one half.
Thirty finite conditional-kernel tests verify
the exact total-variation product contraction;
their numerical singular values and heat-bath
eigenvalues check the operator-angle normalization.
Thirty truncated one-edge generators, for activities
from 16 to 4096, check the rate normalization and
gap scale. The finite spectra and truncations are
diagnostics, not proofs of the infinite-state or
every-exterior assertions. Those follow from the
inequalities and operator arguments above.

V69's mathematical body is preserved byte-for-byte
apart from the version title, and its mathematical
regression is rerun. Other manuscripts remain
unchanged. Only the latest active \texttt{.tex}
replaces its predecessor in \path{latex_notes};
all build products and certificates are outside.
The next companion review is V75 and the next
broad literature sweep is V80, with cadence retained
through the V100 archive boundary.

The interacting four-dimensional continuum
construction and physical Yang--Mills mass gap
remain unproved. No numerical eigenvalue or
conditional auxiliary gap is promoted to that goal.
% END CONSOLIDATED V70 APPEND

% BEGIN CONSOLIDATED V71 APPEND
\clearpage
\section{V71: normalized exterior response and an obstruction to uniform spatial conditioning}
\label{r71:main}

\subsection{The question after V70 and the changed direction}

V70 proves every-exterior one- and two-edge auxiliary count
gaps. The next question is how the normalized law and its
Poisson inverse change when an exterior count changes.
We answer that question for a single pair insertion, with
explicit constants. A different result then rules out a
tempting continuation: the dependence on a remote edge
cannot decay with distance uniformly over all exteriors.

The obstruction is native, not an abstract example. Stretching
the four cross paths in V69 leaves its normalized two-edge
law unchanged, even when the selected edges become arbitrarily
far apart. We evaluate the resulting nonzero one-edge
conditional response. The same exterior is exponentially
unlikely by V68. Thus the spatial task must retain the
probability of the exterior, rather than seek an impossible
worst-exterior spatial-decay estimate. This does not prove
the required probability-weighted spatial operator bound.

\paragraph{Nonduplication and method context.}
We reuse V68's coefficient-ratio bound, V69's evaluated
four-terminal kernel, and V70's conditional gaps. None is
reproved as a new result. Generic energy-normalized inverse
arguments already appear in the f7 V77 and f8 V21 material.
The new inverse estimate here evaluates the actual count-law
change of measure, conductances, centering, and source.
As a targeted primary-source check, Aldous--Fill,
\emph{Reversible Markov Chains and Random Walks on Graphs},
Section 3.6, especially Lemma 3.29
(\url{https://www.stat.berkeley.edu/~aldous/RWG/Book_Ralph/Ch3.S6.html}),
records the standard direct Dirichlet-form comparison method.
Our countable-state comparison and normalization estimates
are proved below. This targeted check does not replace the
scheduled V80 literature sweep; the next companion review
remains V75.

\subsection{One exterior pair insertion tilts the entire conditional block}

Use the finite simple-graph, positive-activity, zero-field
ferromagnetic \(O(4)\) count law of V68. For a full count
vector \(K\), write its moment coefficient as
\[
 M(K)=\int\prod_{e=xy}(z_x\cdot z_y)^{K_e}\prod_xdH(z_x).
 \tag{V71.1}\label{r71:moment}
\]
The positive sphere-pairing expansion shows that this is
strictly positive exactly when every vertex has even
occurrence degree. The integrand itself need not be
positive. The count weight is proportional to
\(M(K)\prod_e\lambda_e^{K_e}/K_e!\).
By V68, on the admissible support,
\[
 q_f(K):=\frac{M(K+2f)}{M(K)},\qquad
 \frac{K_f+1}{K_f+4}\le q_f(K)\le1.
 \tag{V71.2}\label{r71:q}
\]

Fix any set \(A\) of free edges, and admissible exterior
counts \(\eta\) on \(A^c\). Choose \(f\notin A\), put
\(n=\eta_f\), and compare exteriors \(\eta\) and
\(\eta^+=\eta+2f\). Their conditional laws on \(K_A\)
are denoted \(\mu,\mu^+\). Their supports agree because
the added pair changes no vertex parity. Cycles in \(A\)
and all compatible internal parity sectors are allowed.
In the following expressions \(q_f(x)=q_f(x,\eta)\).
Set
\[
 a=\frac{n+1}{n+4},\quad \delta=1-a=\frac3{n+4},\quad
 m=\mu(q_f),\quad \eta_n=\frac\delta a=\frac3{n+1},\quad
 \xi_n=\frac\delta{2\sqrt a}.
 \tag{V71.3}\label{r71:parameters}
\]
The symbol \(\eta_n\) is a scalar bound, distinct from
the exterior vector \(\eta\).

\begin{theorem}[Normalized exterior tilt for an arbitrary block]
\label{r71:tilt}
The exact likelihood ratio and uniform estimates are
\[
 w(x):=\frac{d\mu^+}{d\mu}(x)=\frac{q_f(x)}m,
 \qquad \frac a m\le w\le\frac1m,
 \qquad \|w-1\|_\infty\le\eta_n,
 \tag{V71.4}\label{r71:density}
\]
\[
 \|\mu^+-\mu\|_{\rm TV}
       \le\frac{1-\sqrt a}{1+\sqrt a},\qquad
 \mu[(w-1)^2]\le\xi_n^2
             =\frac9{4(n+1)(n+4)}.
 \tag{V71.5}\label{r71:tilt-bounds}
\]
For every real \(F\in L^2(\mu)\),
\[
 \mu^+F-\mu F=\frac{\Cov_\mu(F,q_f)}m,
 \qquad |\mu^+F-\mu F|\le\xi_n\sqrt{\Var_\mu F}.
 \tag{V71.6}\label{r71:mean}
\]
All estimates are independent of the size of \(A\).
They carry no spatial-distance factor.
\end{theorem}
\begin{proof}
The exterior activity and factorial ratio
\(\lambda_f^2/[(n+2)(n+1)]\) is constant in \(x\)
and cancels upon normalization. The remaining ratio is
\(q_f(x)\), proving \eqref{r71:density} and
\(a\le m\le1\). Its oscillation is at most \(\delta\),
so \(|q_f-m|/m\le\delta/a\).

For a random variable \(q\in[a,1]\) with mean \(m\),
the chord bound for \(|q-m|\), and
\((q-a)(q-1)\le0\), give respectively
\[
 \E|q-m|\le\frac{2(1-m)(m-a)}{1-a},\qquad
 \Var(q)\le(1-m)(m-a).
\]
Maximizing the first expression after division by \(2m\)
gives \((1-\sqrt a)/(1+\sqrt a)\), at \(m=\sqrt a\).
For the second, the identity
\[
 (1-a)^2m^2-4a(1-m)(m-a)=((1+a)m-2a)^2\ge0
\]
gives \(\Var(q)/m^2\le(1-a)^2/(4a)\).
Apply these facts to \(q_f\). Finally
\(\mu(w-1)=0\), so Cauchy--Schwarz gives
\eqref{r71:mean}. No spin monomial has been used as
a positive probability density.
\end{proof}

\subsection{The rate change and the form change are different}

For \(e\in A\), write \(b_e(x)=\lambda_e^2q_e(x,\eta)\)
and \(b_e^+(x)=\lambda_e^2q_e(x,\eta^+)\).
Commuting the two coefficient insertions gives
\[
 \frac{b_e^+(x)}{b_e(x)}
       =\frac{q_f(x+2e)}{q_f(x)},\qquad
 a\le\frac{b_e^+}{b_e}\le a^{-1},\qquad
 \left|\frac{b_e^+}{b_e}-1\right|\le\eta_n.
 \tag{V71.7}\label{r71:cycle}
\]
The death rates \(K_e(K_e-1)\) do not change.
For the original count clocks define the forms
\[
 \mathcal E(u,v)=\sum_x\mu(x)\sum_{e\in A}
       b_e(x)\Delta_eu(x)\Delta_ev(x),\qquad
 \Delta_eu(x)=u(x+2e)-u(x),
 \tag{V71.8}\label{r71:form}
\]
and define \(\mathcal E^+\) with \(\mu^+,b^+\).
We work with real functions; polarization also gives
the complex version with conjugation.

\begin{proposition}[Normalized conductances and comparison]
\label{r71:conductance}
For every allowed upward move,
\[
 \frac{\mu^+(x)b_e^+(x)}{\mu(x)b_e(x)}
                  =\frac{q_f(x+2e)}m.
 \tag{V71.9}\label{r71:conductance-ratio}
\]
Consequently,
\[
 \frac a m\mathcal E(u,u)\le\mathcal E^+(u,u)
                  \le\frac1m\mathcal E(u,u),
 \qquad
 \frac a m\Var_\mu u\le\Var_{\mu^+}u
                  \le\frac1m\Var_\mu u,
 \tag{V71.10}\label{r71:comparison}
\]
\[
 |(\mathcal E^+-\mathcal E)(u,v)|
         \le\eta_n\sqrt{\mathcal E(u,u)\mathcal E(v,v)}.
 \tag{V71.11}\label{r71:form-difference}
\]
The form domains agree under the equivalent \(L^2\) norms.
If a positive Poincare gap \(\lambda\) is already known
for \(\mathcal E\), then the optimal gaps satisfy
\(a\lambda\le\lambda^+\le\lambda/a\), when
\(\lambda\) denotes the optimal original gap.
\end{proposition}
\begin{proof}
Multiply \eqref{r71:cycle} by \eqref{r71:density};
the factor \(q_f(x)\) cancels and gives
\eqref{r71:conductance-ratio}. Bounds \(a\le q_f\le1\)
prove the form inequalities. Variance is
\(\inf_c\mu[(u-c)^2]\), so the same density comparison
proves the variance inequalities. Their common denominator
\(m\) cancels in the Rayleigh quotient, giving the
claimed factor \(a\), not a lost factor \(a^2\).
The conductance difference has relative absolute value
at most \(\delta/m\le\eta_n\); Cauchy--Schwarz on
the edge sum proves \eqref{r71:form-difference}.
The equivalent norms and forms give identical completed
domains. These arguments apply directly to convergent
countable sums, or by monotone convergence to the
quadratic forms.
\end{proof}

This comparison does not produce a gap on a general
block from nothing. If the block has several internal
parity sectors, the pair-move generator still preserves
them. In the application below \(A\) has one or two
edges and V70 supplies the positive gap.

\subsection{A centered inverse estimate with the changed source retained}

For \(F\in L^2(\mu)\), denote the centered Poisson
solution by \(u=R_\mu F\), so
\[
 \mu u=0,\qquad \mathcal E(u,v)=\Cov_\mu(F,v).
\]
Similarly put \(u^+=R_{\mu^+}F^+\), normalized by
\(\mu^+u^+=0\). Constants in either source do not matter.
Let \(H=F^+-F\), \(\sigma_F^2=\Var_\mu F\) and
\(\sigma_H^2=\Var_\mu H\). Both sources are assumed
in \(L^2(\mu)\), equivalently in \(L^2(\mu^+)\).

\begin{theorem}[Exterior stability of the normalized Poisson inverse]
\label{r71:inverse}
Suppose \(\mathcal E\ge\lambda\Var_\mu\) with
\(\lambda>0\). The solutions exist uniquely with
the stated centerings. Set
\(d=u^+-\mu u^+-u\). Then
\[
 \begin{split}
 \sqrt{\mathcal E(d,d)}
 &\le\frac{a^{-1}\sigma_H+(2\eta_n+\xi_n^2)\sigma_F}
                   {a\sqrt\lambda},\\
 \|d\|_{L^2(\mu)}
 &\le\frac{a^{-1}\sigma_H+(2\eta_n+\xi_n^2)\sigma_F}
                   {a\lambda}.
 \end{split}
 \tag{V71.12}\label{r71:inverse-bound}
\]
In particular, take a one- or two-edge block of activity
\(\gamma\ge4096\) and an exterior edge with
\(n\ge\gamma/2\). For a fixed source \(F^+=F\),
\[
 \|R_{\mu^+}F-\mu(R_{\mu^+}F)-R_\mu F\|_{L^2(\mu)}
             \le\frac{2^{19}}{\gamma^2}\sigma_F.
 \tag{V71.13}\label{r71:high-count-inverse}
\]
This is an inverse comparison modulo constants, not
an identification of the two centered Hilbert spaces.
\end{theorem}
\begin{proof}
The closed self-adjoint operators have positive gaps
by assumption and Proposition \ref{r71:conductance},
so spectral inversion gives existence and uniqueness.
It also gives \(\mathcal E(u,u)\le\sigma_F^2/\lambda\).
For \(\mu\)-centered \(F,v\), direct expansion shows
\[
 \Cov_{\mu^+}(F,v)-\Cov_\mu(F,v)
   =\mu[(w-1)Fv]-\mu[(w-1)F]\mu[(w-1)v].
\]
Thus for arbitrary \(F,v\) the absolute difference
is at most \((\eta_n+\xi_n^2)\sigma_F\sigma_v\),
where \(\sigma_v^2=\Var_\mu v\). Also
\(|\Cov_{\mu^+}(H,v)|\le a^{-1}\sigma_H\sigma_v\)
by the variance comparison. Subtraction of the two
weak Poisson equations, which is legitimate on their
common form domain, gives the exact identity
\[
 \begin{split}
 \mathcal E^+(d,v)
   ={}&\Cov_{\mu^+}(H,v)
      +(\Cov_{\mu^+}-\Cov_\mu)(F,v)\\
     &-(\mathcal E^+-\mathcal E)(u,v).
 \end{split}
 \tag{V71.14}\label{r71:weak-difference}
\]
Use \eqref{r71:form-difference},
\(\sigma_v\le\mathcal E(v,v)^{1/2}/\sqrt\lambda\),
and the bound on \(\mathcal E(u,u)\). The right side
is at most
\[
 \frac{a^{-1}\sigma_H+(2\eta_n+\xi_n^2)\sigma_F}
                  {\sqrt\lambda}\sqrt{\mathcal E(v,v)}.
\]
Take \(v=d\) and use \(\mathcal E^+(d,d)\ge a\mathcal E(d,d)\).
This proves the first inequality; Poincare proves the second.

For the last assertion, V70 allows \(\lambda=\gamma/16384\).
For \(n\ge\gamma/2\) in this range, \(a\ge1/2\),
\(\eta_n\le6/\gamma\), and
\(\xi_n^2=\delta\eta_n/4\le\eta_n/4\).
Therefore \((2\eta_n+\xi_n^2)/a\le27/\gamma\le32/\gamma\).
With \(H=0\), \eqref{r71:high-count-inverse} follows.
\end{proof}

\paragraph{The actual deficit is not a fixed source.}
For \(e\in A\), its count-averaged deficit is
\(\tau_e=\lambda_e(1-q_e)\). Formula \eqref{r71:cycle}
gives the separate source change
\[
 \tau_e^+(x)-\tau_e(x)
   =-\lambda_e q_e(x)
       \left(\frac{q_f(x+2e)}{q_f(x)}-1\right),
 \qquad |\tau_e^+-\tau_e|\le\lambda_e\eta_n.
 \tag{V71.15}\label{r71:source-change}
\]
Consequently \eqref{r71:inverse-bound} applies to the
actual deficit source, but its \(\sigma_H\) term must
remain. At homogeneous activity and \(n\ge\gamma/2\),
the displayed estimate only gives \(\sigma_H\le6\).
V68 gives \(\sigma_F\le72\) for a one- or two-edge
conditional block. It would be incorrect to call
\eqref{r71:high-count-inverse} an \(O(\gamma^{-2})\)
estimate for this changing source without an additional
estimate on \(H\). The raw whole-loop source is not
being restored in place of the count-averaged one.

\subsection{A stretched native exterior with no distance attenuation}

Let \(\ell\ge1\) and let the cubic torus side \(N\)
be even with \(N\ge2\ell+2\). Choose terminals
\[
 1=(0,0,0),\quad2=(1,0,0),\quad
 3=(0,\ell,0),\quad4=(1,\ell,0),
 \qquad e=12,\quad f=34.
 \tag{V71.16}\label{r71:terminals}
\]
The minimum graph distance between their endpoint sets
is \(\ell\). Connect the cross pairs by paths as follows:
\[
 \begin{aligned}
 P_{13}:{}&(0,j,0),\quad j=0,\ldots,\ell,\\
 P_{24}:{}&(1,j,0),\quad j=0,\ldots,\ell,\\
 P_{14}:{}&(0,0,0)\to(0,0,1)
       \to\cdots\to(0,\ell,1)\to(1,\ell,1)\to(1,\ell,0),\\
 P_{23}:{}&(1,0,0)\to(1,0,-1)
       \to\cdots\to(1,\ell,-1)\to(0,\ell,-1)\to(0,\ell,0).
 \end{aligned}
 \tag{V71.17}\label{r71:paths}
\]
The ellipses mean nearest-neighbor increments of the
\(j\)-coordinate. Coordinates are taken modulo \(N\).
Let \(C_\ell\) be the exterior event giving count one
to all these path edges and zero to every other edge
outside \(\{e,f\}\).

\begin{theorem}[Distance-independent normalized conditional interaction]
\label{r71:stretch}
The event \(C_\ell\) has positive probability. Given
this event, the joint law of \(K_e=k,K_f=l\) is exactly
V69's law
\[
 \mu_{g,h}(k,l)=\frac{g^kh^l a_{kl}}{k!l!Z_2(g,h)},
 \qquad
 a_{kl}=\frac{c_kc_l}{72}F(r_k,r_l),\quad k,l\text{ even},
 \tag{V71.18}\label{r71:remote-law}
\]
where \(g=\lambda_e,h=\lambda_f\),
\(r_k=(k+1)/(k+4)\) and \(F(x,y)=4x+4y+2xy-1\).
It is independent of \(\ell\), \(N\), and all positive
activities on the fixed exterior edges.
\end{theorem}
\begin{proof}
The paths have lengths \(\ell,\ell,\ell+3,\ell+3\).
Their interiors are disjoint and avoid the terminals:
the direct paths are in plane \(z=0\), and the other
two interiors are in the distinct planes \(z=1,-1\).
There are \(4\ell+6\) distinct path edges and
\(4\ell+2\) distinct internal vertices. No selected
edge is used. Each internal vertex has occurrence
degree two, so integrating it gives a factor \(1/4\)
and contracts its two adjacent scalar products.
The full moment is therefore
\[
       M(K)=4^{-(4\ell+2)}a_{kl}.
\]
The four terminal outside degrees are two, so precisely
even \(k,l\) are admissible. The common contraction
factor, exterior activities and exterior factorials
cancel after normalization. Positivity of \(a_{00}\)
and finiteness of the full partition function give
strictly positive probability of \(C_\ell\). There
is no probability lower bound uniform in \(\ell\).
\end{proof}

We next evaluate an exterior \emph{response}, not just
the inherited negative two-edge covariance. Define
\[
 \pi_\gamma(k)=\frac{\gamma^kc_k}{k!\phi(\gamma)},\quad
 \phi(\gamma)=\sum_{k\ {\rm even}}\frac{\gamma^kc_k}{k!},\quad
 R_\gamma=\E_{\pi_\gamma}r_K.
 \tag{V71.19}\label{r71:baseline}
\]
This baseline is not the parity-Poisson law of V68.
For even \(n\ge0\), let \(p_n\) be the law of \(K_e\)
conditional on \(C_\ell\) and \(K_f=n\), at activity
\(\lambda_e=\gamma>0\). Then
\[
 p_n(k)=\pi_\gamma(k)\frac{F(r_k,r_n)}{F(R_\gamma,r_n)}.
 \tag{V71.20}\label{r71:pn}
\]

\begin{proposition}[A nonzero remote insertion response]
\label{r71:remote-response}
Put \(d_n=r_{n+2}-r_n=6/[(n+4)(n+6)]\). Then
\[
 p_{n+2}(k)-p_n(k)
  =\frac{-18d_n\pi_\gamma(k)(r_k-R_\gamma)}
    {F(R_\gamma,r_{n+2})F(R_\gamma,r_n)},
 \tag{V71.21}\label{r71:pdifference}
\]
\[
 \begin{split}
 \|p_{n+2}-p_n\|_{\rm TV}
 &=\frac{9d_n\E_{\pi_\gamma}|r_K-R_\gamma|}
         {F(R_\gamma,r_{n+2})F(R_\gamma,r_n)}>0,\\
 \E_{p_{n+2}}K-\E_{p_n}K
 &=\frac{-18d_n\Cov_{\pi_\gamma}(K,r_K)}
         {F(R_\gamma,r_{n+2})F(R_\gamma,r_n)}<0.
 \end{split}
 \tag{V71.22}\label{r71:remote-values}
\]
These quantities are independent of the distance \(\ell\).
Thus, for fixed \(\gamma>0\) and fixed even \(n\),
no bound tending to zero with distance can hold for
this conditional TV response uniformly in volume and
admissible exterior. The same statement holds for
the absolute conditional mean-count response.
\end{proposition}
\begin{proof}
The affine dependence on the first variable in \(F\)
proves \eqref{r71:pn}. Direct multiplication gives
\[
 F(x,y')F(R,y)-F(x,y)F(R,y')=-18(x-R)(y'-y).
\]
This proves \eqref{r71:pdifference}; summing its
absolute value, or multiplying by \(k\) and summing,
gives \eqref{r71:remote-values}. All sums converge
absolutely by the factorial weights. The baseline
assigns positive mass to every even count, and \(r_k\)
is strictly increasing. Hence its mean absolute
deviation is positive and
\[
 \Cov_{\pi_\gamma}(K,r_K)
 =\tfrac12\E[(K-K')(r_K-r_{K'})]>0
\]
for independent baseline copies \(K,K'\). The
denominators are positive because \(R_\gamma,r_n\ge1/4\)
and \(F\ge9/8\) there. Let \(\ell\) tend to infinity
along \(N=2\ell+2\); the positive TV value is constant,
contradicting any proposed distance-decaying uniform
bound. One can take \(n\) above \(\gamma/2\) as well;
large selected counts do not remove the exterior issue.
\end{proof}

\subsection{Why this obstruction does not rule out the needed averaged estimate}

In the homogeneous law at \(\gamma\ge216\), every
path count in \(C_\ell\) is at most \(\gamma/2\).
V68's joint low-count estimate therefore implies
\[
 0<\mathbb P(C_\ell)
       \le q_\gamma^{\,4\ell+6},\qquad
 q_\gamma=\min\{1,38e^{-\gamma/24}\}<\frac1{200}.
 \tag{V71.23}\label{r71:rare-exterior}
\]
The additional zero-count restrictions only decrease
this probability. The same upper bound holds after
intersecting with \(\{K_f=n\}\). For a bounded test
\(|F|\le1\), the insertion response on that specified
exterior is at most \(2\|p_{n+2}-p_n\|_{\rm TV}\).
Its contribution multiplied by the probability of
that event is consequently at most
\[
 2\|p_{n+2}-p_n\|_{\rm TV}\,q_\gamma^{\,4\ell+6}.
 \tag{V71.24}\label{r71:weighted-fixture}
\]
This is an evaluated bound for the specified fixture,
not a summation over all exterior configurations,
all connecting paths, or all signed sources.

The conclusion is a change of quantifiers. Local
coercivity and single-insertion stability hold for
every exterior. Spatial attenuation cannot hold
uniformly over those exteriors. A future spatial
proof must instead control exterior probabilities
and connected propagation together, or prove a
signed global operator estimate that does not
require that false uniform attenuation hypothesis.
V68 suppresses long low-count structures, but does
not yet control transmission through the typical
high-count region.

\subsection{Full-bank status and verification}

The full homogeneous identity remains
\[
 \mathcal B_{\rm w}=(2\gamma I+3D)\Gamma
       +\gamma[\Cov(\tau_e,K_f)]_{e,f}
       -2\gamma^2(1-\E\tau_e/\gamma)I-6C^{\rm ret}.
 \tag{V71.25}\label{r71:bank}
\]
The new tilt theorem does not remove exterior-mean
covariances from total covariance. Its block-size
independence is not spatial decay; summing over
exterior edges still needs a volume-uniform argument.
The new inverse bound retains source motion explicitly.
All factorial contacts, returns, \(D\Gamma\), and
parity/winding sectors remain. The two-count clock
is still auxiliary, not physical Hamiltonian time.

The accompanying checker verifies the likelihood,
conductance, covariance and coefficient-cycle identities
with exact rational arithmetic, checks the inverse
comparison on finite reversible truncations, and
independently checks the stretched cubic embeddings
and their sphere-pairing coefficients. It also verifies
the evaluated remote response on finite factorial
series. Truncated spectral calculations are diagnostics,
not proofs of the countable-state or every-exterior
theorems, which are proved above.

V70's mathematical body is preserved byte-for-byte
apart from the version title, and its mathematical
regression is rerun. Other manuscripts are unchanged.
Only the latest active \texttt{.tex} replaces its
predecessor in \path{latex_notes}; diagnostics and
build products remain outside that folder.
The previous goal turn was progress, and this one
pays a normalized response estimate and rules out
an incorrect spatial quantifier. It does not close
the full spatial bound, the interacting
four-dimensional continuum construction, or the
physical Yang--Mills mass gap.
% END CONSOLIDATED V71 APPEND

% BEGIN CONSOLIDATED V72 APPEND
\clearpage
\section[V72: Parity channels in high-count paths]{V72: retained parity channels in high-count paths and fixed-pattern probability}
\label{r72:main}

\subsection{What must propagate through a high-count region?}

V71 excludes spatial attenuation uniform over all count
exteriors. Its explicit transmitting paths have count one,
so one might try to recover scalar attenuation merely by
requiring large counts on every transmitting edge. This
appendix tests that proposed repair before using it in a
global estimate. The repair is insufficient: odd-count
paths retain a four-dimensional angular channel, even
when all their counts are large and bounded above.

We compute the complete path transfer and prove decay
of its remainder after retaining the correct leading
channel. On even-count paths the leading channel is
constant. On odd-count paths it is not. In V71's cubic
geometry, replacing every transmitting count by the
same large odd integer produces a conditional law
converging to the same interacting four-terminal law
as the path lengths grow. This is a high-count-path
obstruction, not a theorem about a typical full exterior.

We also prove a uniform maximum-atom estimate for the
actual count law, and a conditional product bound for
any specified count pattern. It supplies a probability
cost for the new examples. It does not sum over the
many possible patterns or establish the required
spatial operator bound.

\paragraph{Inherited tools and scope.}
The \(\SU(2)\) character fusion and convolution
normalization are inherited from V19; V24 already
distinguishes angular sectors from the full normalized
environment. V12's positive shifted-power approximation
kernel is a different object from the fixed-count
power kernel below. The four-terminal coefficient
\(a_{kl}\), stretched cubic geometry, and its nonzero
insertion response are inherited from V69--V71.
Here the new calculation is the fixed-count path
propagation, its retained parity channel, and its
consequence for high-count exteriors.

For the polynomial convention, NIST DLMF
\url{https://dlmf.nist.gov/18.12.E4} and
\url{https://dlmf.nist.gov/18.12.E10}
identify \(C_j^{(1)}=U_j\) by their common generating
function. All power coefficients and path estimates
used here are derived below. No new spectral theorem
for the physical Hamiltonian is imported. The next
companion review remains V75 and the next broad
literature sweep V80.

\subsection{Complete spectrum of a fixed-count bond}

Identify \(S^3\) with unit quaternions and use Haar
probability. Let \(\chi_j(t)=U_j(t)\), so
\(\chi_0=1\), \(\chi_1=2t\), and
\(2t\chi_j=\chi_{j+1}+\chi_{j-1}\), with
\(\chi_{-1}=0\). Denote by \(P_j\) the orthogonal
projection onto the degree-\(j\) spherical harmonics.
Its kernel and rank, in this normalization, are
\[
 P_j(u,v)=(j+1)\chi_j(u\cdot v),\qquad
 \operatorname{rank}P_j=(j+1)^2,\qquad
 |P_j(u,v)|\le(j+1)^2.
 \tag{V72.1}\label{r72:projector}
\]
These are V19's full matrix-coefficient sectors, not
just their one-dimensional central subspaces. Schur
orthogonality gives the projection formula; the last
bound follows from the unitary character bound
\(|\chi_j|\le j+1\), or from the reproducing kernel.

For an integer count \(k\ge0\), set
\[
 (T_k f)(u)=\int (u\cdot v)^k f(v)\,dH(v).
 \tag{V72.2}\label{r72:Tk}
\]

\begin{proposition}[Parity-resolved bond spectrum]
\label{r72:spectrum}
The finite-rank self-adjoint operator is
\[
 T_k=\sum_{\substack{0\le j\le k\\j\equiv k\ (2)}}
              \lambda_j(k)P_j,\qquad
 \lambda_j(k)=
 \frac{k!}{2^k((k-j)/2)!((k+j)/2+1)!}>0.
 \tag{V72.3}\label{r72:lambda}
\]
All other eigenvalues vanish. If \(r=k\bmod2\), put
\[
 \alpha_k=\lambda_r(k),\qquad
 \beta_k=\frac{k-r}{k+r+4}.
 \tag{V72.4}\label{r72:alpha-beta}
\]
Then \(\alpha_k=c_k\) for even \(k\), and
\(\alpha_k=c_{k+1}\) for odd \(k\), with the
moments \(c_j\) of V68. Moreover,
\[
 \alpha_k^{-1}T_k=P_r+E_k,\quad E_kP_r=P_rE_k=0,
 \quad E_k\succeq0,\quad \|E_k\|=\beta_k<1.
 \tag{V72.5}\label{r72:leading}
\]
\end{proposition}
\begin{proof}
On a diagonal quaternion with eigenvalues \(z,z^{-1}\),
\(\chi_j=z^j+z^{j-2}+\cdots+z^{-j}\).
Comparison with \((z+z^{-1})^k\) gives
\[
 t^k=2^{-k}\sum_{\substack{0\le j\le k\\j\equiv k\ (2)}}
 \left\{\binom{k}{(k-j)/2}-\binom{k}{(k-j)/2-1}\right\}\chi_j(t),
\]
where a binomial coefficient with negative lower index
is zero. The difference equals
\((j+1)k!/[((k-j)/2)!((k+j)/2+1)!]\).
Division of the character coefficient by \(j+1\)
gives \eqref{r72:lambda} with the inherited convolution
normalization. Direct factorial cancellation shows
\[
 \frac{\lambda_{j+2}(k)}{\lambda_j(k)}
                         =\frac{k-j}{k+j+4}<1.
\]
Thus the largest eigenvalue is \(\lambda_r(k)\),
and the next ratio is \(\beta_k\). When \(k=r\),
there is no next sector and \(E_k=0\), as required.
The identities for \(\alpha_k\) follow from the
explicit even moments. This proves every assertion.
\end{proof}

Although \(T_k\) is positive semidefinite for odd
\(k\), its pointwise kernel is signed and \(T_k1=0\).
It is not a Markov kernel. For even \(k\),
\(T_k/\alpha_k\) is a positive Markov kernel.
In particular,
\[
 P_0(u,v)=1,\qquad P_1(u,v)=4u\cdot v.
 \tag{V72.6}\label{r72:two-leading}
\]
Positive semidefiniteness cannot turn the second
kernel into a scalar probability transition.

\subsection{Exact contraction above the retained channel}

Consider a geometric path of \(L\ge1\) edges with
counts \(k_1,\ldots,k_L\), and with no other positive
count incident to an internal vertex. Counts greater
than one are permitted; ``path'' describes the
geometric support, not occurrence degree two.
At an admissible internal vertex \(k_i+k_{i+1}\)
is even, so all path counts have one parity \(r\).
Integrating the internal spins is exactly the
convolution \(T_{k_1}\cdots T_{k_L}\).

\begin{theorem}[Path propagation with the correct retained space]
\label{r72:path}
Let \(A_P=\prod_i\alpha_{k_i}\),
\(b_P=\prod_i\beta_{k_i}\), and \(K=\max_i k_i\).
Then
\[
 A_P^{-1}T_{k_1}\cdots T_{k_L}=P_r+E_P,
 \qquad \|E_P\|=b_P,
 \tag{V72.7}\label{r72:path-op}
\]
and its kernel satisfies
\[
 |E_P(u,v)|\le D_{K,r}b_P,\qquad
 D_{K,r}=\sum_{\substack{r+2\le j\le K\\j\equiv r\ (2)}}(j+1)^2.
 \tag{V72.8}\label{r72:path-kernel}
\]
If all \(k_i\le B\), then
\[
 b_P\le\exp\left(-\frac{(2r+4)L}{B+r+4}\right).
 \tag{V72.9}\label{r72:path-rate}
\]
The leading projection is exactly retained at every
length: it has rank one for even paths and rank four
for odd paths.
\end{theorem}
\begin{proof}
All factors are diagonal in the same orthogonal
projections \(P_j\). On \(P_r\), every normalized
factor is the identity. On the next sector the
product eigenvalue is \(b_P\), and every higher
product eigenvalue is between zero and \(b_P\).
If a factor has \(k_i=r\), all higher products
vanish and both sides give \(E_P=0\).
The kernel bound follows by summing
\eqref{r72:projector}; only degrees at most
\(\min_i k_i\) can occur, so the stated use of
\(K\) is an upper bound. Finally
\(\beta_k=1-(2r+4)/(k+r+4)\), and
\(1-x\le e^{-x}\), prove \eqref{r72:path-rate}.
\end{proof}

This is genuine contraction of the higher channels
along a fixed-count path, including paths with
counts in a bounded high-count range. It is not
contraction to constants in the odd sector.
Indeed the normalized product acts as the identity
on the four coordinate functions \(u\mapsto u_a\).
Calling these four internal channels irrelevant
because an external observable is even is also
unjustified: the next calculation closes four odd
channels into a scalar interaction between even
selected counts.

\subsection{Large odd counts preserve the remote conditional interaction}

Use exactly the terminals and four disjoint paths
of \eqref{r71:terminals}--\eqref{r71:paths}, with
lengths \(\ell,\ell,\ell+3,\ell+3\), and
even \(N\ge2\ell+2\). Replace the count one on
each path edge by a fixed odd integer \(\kappa\ge1\).
All other nonselected counts are zero as before.
Call this exterior event \(C_{\ell,\kappa}\).
Every internal geometric vertex now has occurrence
degree \(2\kappa\). Each terminal has outside
degree \(2\kappa\), so the selected counts \(k,l\)
are still even.

Let \(\alpha=\lambda_1(\kappa)\),
\(\beta=(\kappa-1)/(\kappa+5)\), and
\(D=D_{\kappa,1}\). For a path of length \(L\),
write its normalized kernel as
\[
 H_L(u,v)=\frac{T_\kappa^L(u,v)}{4\alpha^L}
               =u\cdot v+R_L(u,v),\qquad
 |R_L|\le\frac D4\beta^L.
 \tag{V72.10}\label{r72:odd-H}
\]
For \(\kappa=1\), this has zero remainder and
recovers V71 exactly. For all odd \(\kappa\), put
\[
 \epsilon_\ell=\frac D4\beta^\ell,\qquad
 e_\ell=4\epsilon_\ell(1+\epsilon_\ell)^3.
 \tag{V72.11}\label{r72:epsilon}
\]

\begin{theorem}[High-count-path limit with normalization retained]
\label{r72:odd-limit}
The event \(C_{\ell,\kappa}\) has positive probability
at any positive activities. Divide its moment
coefficient by the common factor
\(4^4\alpha^{4\ell+6}\), and denote the resulting
coefficient for selected counts \(k,l\) by
\(\widehat a^{\ell,\kappa}_{kl}\). Then, uniformly
over nonnegative even \(k,l\),
\[
 \widehat a^{\ell,\kappa}_{kl}>0,\qquad
 |\widehat a^{\ell,\kappa}_{kl}-a_{kl}|\le e_\ell,
 \qquad e_\ell\longrightarrow0.
 \tag{V72.12}\label{r72:coefficient-limit}
\]
Here \(a_{kl}\) is V69's cross-cycle coefficient.
For selected activities \(g,h>0\), let
\(p_{\ell,\kappa;g,h}\) be the normalized joint
count law on this exterior, and \(p_{g,h}\) the
law \eqref{r69:p}. Then
\[
 \|p_{\ell,\kappa;g,h}-p_{g,h}\|_{\rm TV}
 \le\min\left\{1,
   \frac{e_\ell\cosh g\cosh h}{Z_2(g,h)}\right\}
       \longrightarrow0.
 \tag{V72.13}\label{r72:joint-TV}
\]
The selected activities and \(\kappa\) are fixed
in this distance limit; no uniform joint limit in
\(\ell,\kappa,g,h\) is asserted.
\end{theorem}
\begin{proof}
Integrate the internal spins of each path separately.
The normalized four-terminal coefficient is the
integral of
\(\sigma_{12}^k\sigma_{34}^l\prod_{ij}H_{L_{ij}}(z_i,z_j)\),
where \(ij=13,14,23,24\). Each leading factor is
\(\sigma_{ij}\), with absolute value at most one,
and every remainder is at most \(\epsilon_\ell\).
Telescoping the product gives uniform error
at most \(4\epsilon_\ell(1+\epsilon_\ell)^3\),
also after multiplication by the selected monomials
and Haar integration. The unnormalized full moment
is strictly positive by the sphere-pairing expansion
and the even vertex degrees. Dividing by the positive
common factor preserves positivity. This also proves
positive probability of the specified exterior.

For clarity, if nonnegative summable weights \(w,w'\)
have totals \(Z,Z'>0\), and \(\sum|w'-w|\le b\), then
\[
 \left\|\frac{w'}{Z'}-\frac wZ\right\|_{\rm TV}
                  \le\frac bZ.
\]
Indeed, split their difference into
\((w'-w)/Z+w'(1/Z'-1/Z)\), and use
\(|Z'-Z|\le b\). Here the coefficient error gives
\(b\le e_\ell\cosh g\cosh h\), proving
\eqref{r72:joint-TV}. In particular no path
attenuation factor has been left outside the
normalizer where it could be mistaken for a
normalized connected-correlation estimate.
\end{proof}

\begin{corollary}[Remote insertion response survives large odd path counts]
\label{r72:remote}
For fixed \(\gamma>0\) and even \(n\ge0\), let
\(p_{\ell,\kappa;n}\) be the conditional law of
\(K_e\) on \(C_{\ell,\kappa}\), with \(K_f=n\)
and \(\lambda_e=\gamma\). Put
\[
 Z_n(\gamma)=\sum_{k\ {\rm even}}\frac{\gamma^k}{k!}a_{kn}>0.
\]
With \(p_n\) from \eqref{r71:pn},
\[
 \|p_{\ell,\kappa;n}-p_n\|_{\rm TV}
                 \le\frac{e_\ell\cosh\gamma}{Z_n(\gamma)}.
 \tag{V72.14}\label{r72:one-TV}
\]
Consequently, if \(v_{\gamma,n}=\|p_{n+2}-p_n\|_{\rm TV}>0\)
is V71's evaluated response, then
\[
 \begin{split}
 &\left|\|p_{\ell,\kappa;n+2}-p_{\ell,\kappa;n}\|_{\rm TV}
                          -v_{\gamma,n}\right|\\
 &\qquad\le e_\ell\cosh\gamma
          \left(\frac1{Z_n(\gamma)}+\frac1{Z_{n+2}(\gamma)}\right)
                   \longrightarrow0.
 \end{split}
 \tag{V72.15}\label{r72:response-limit}
\]
\end{corollary}
\begin{proof}
Apply the same summable-weight comparison with
\(l=n\) fixed; factors involving the activity and
factorial of \(f\) cancel. The two triangle
inequalities for total variation give
\eqref{r72:response-limit}. Positivity of
\(v_{\gamma,n}\) was proved in V71, and is not
inferred from numerical convergence here.
\end{proof}

For example, at any fixed \(\gamma\ge216\) one can
choose an odd \(\kappa\in[\gamma/2,2\gamma]\).
Every transmitting path edge then has a count
in this high-count band, yet its remote conditional
response is bounded below by \(v_{\gamma,n}/2\)
for all sufficiently large \(\ell\). One can also
fix \(n\ge\gamma/2\). Thus large counts along the
transmitting paths, and a bound above on those counts,
do not by themselves yield scalar attenuation.

The surrounding zero-count edges have \emph{not}
been removed from the exterior specification. This
is not a counterexample in which all exterior edges
are in the high-count band, nor a claim about the
typical dense count network. The precise obstruction
is to treating a high-count path as an automatically
contracting scalar conduit independently of its
parity channel and attached environment.

\subsection{The contrasting even-count path limit}

For a fixed even \(\kappa\ge2\), use the same
geometric exterior. Now \(\alpha=\lambda_0(\kappa)=c_\kappa\)
and \(\beta=\kappa/(\kappa+4)\). The path kernel
\(T_\kappa^L/\alpha^L\) equals \(1+R_L\), with
\(|R_L|\le D_{\kappa,0}\beta^L\).
The same product and normalization argument proves
that the two selected even-count laws converge in
total variation to
\[
 \pi_g\otimes\pi_h,\qquad
 \pi_g(k)=\frac{g^kc_k}{k!\phi(g)}\mathbf1_{\{k\ {
m even}\}}.
 \tag{V72.16}\label{r72:even-limit}
\]
Here the moment normalizer is \(\alpha^{4\ell+6}\),
and the limiting coefficient is \(c_kc_l\).
An explicit error is \eqref{r72:joint-TV} with
\(\epsilon_\ell=D_{\kappa,0}\beta^\ell\),
\(e_\ell=4\epsilon_\ell(1+\epsilon_\ell)^3\),
and \(Z_2(g,h)\) replaced by \(\phi(g)\phi(h)\).
For either fixed even \(n\), the conditional
one-edge limit is \(\pi_\gamma\), independent of
\(n\), so its pair-insertion TV response tends to
zero. This proof treats the same normalization as
in the odd case; the difference is the retained
channel, not a discarded denominator.

\subsection{An all-exterior atom bound and the cost of prescribed patterns}

We now return to the actual full count law, not the
fixed-path transfer. V68's positive one-edge mixture
has conditional intensity \(\lambda=\gamma t\) and
\(W=\gamma(1-t)\), with \(\E e^{W/2}<36\) uniformly
in the exterior. We use this stronger mixture input,
not the coarser total-variation error in V70.

\begin{lemma}[Uniform conditional anti-concentration]
\label{r72:atom}
For every admissible exterior in the V68 count law
and every selected activity \(\gamma\ge216\),
\[
 \sup_{k\ge0}\mathbb P(K_e=k\mid K_{e^c})
                      \le\frac4{\sqrt\gamma}<\frac13.
 \tag{V72.17}\label{r72:atom-bound}
\]
\end{lemma}
\begin{proof}
For \(X\) Poisson with mean \(\lambda>0\), Fourier
inversion and \(1-\cos\theta\ge2\theta^2/\pi^2\)
for \(|\theta|\le\pi\) give
\[
 \begin{split}
 \mathbb P(X=k)
 &\le\frac1{2\pi}\int_{-\pi}^{\pi}
                  e^{-\lambda(1-\cos\theta)}\,d\theta\\
 &\le\frac{\sqrt\pi}{2\sqrt{2\lambda}}
                         <\frac1{\sqrt{2\lambda}}.
 \end{split}
 \tag{V72.18}\label{r72:poisson-atom}
\]
The trigonometric bound follows from concavity of
\(\sin\) on \([0,\pi/2]\); the last inequality
uses \(\pi<4\). Fourier inversion follows directly
by integrating the absolutely convergent exponential
series. For \(\lambda\ge1\), either Poisson parity
has probability at least \((1-e^{-2\lambda})/2>1/3\).
Thus every atom of \(Q_{r,\lambda}\) is at most
\(3/\sqrt{2\lambda}\).

On the mixture event \(\lambda\ge\gamma/2\), this
is at most \(3/\sqrt\gamma\). Its complement has
probability at most \(36e^{-\gamma/4}\) by the
exponential moment of \(W\). Therefore the required
atom is at most \(3/\sqrt\gamma+36e^{-\gamma/4}\).
For \(\gamma\ge216\),
\(36\sqrt\gamma e^{-\gamma/4}<1\): this expression
is decreasing for \(\gamma>2\), and at 216 it is
less than \(540e^{-54}<1\). This proves the first
bound. Since \(\gamma\ge216>144\), it is below
\(1/3\). The imposed parity merely makes the
incompatible atoms zero; no independence is used.
\end{proof}

\begin{theorem}[Conditional cost of a specified count pattern]
\label{r72:pattern}
In the homogeneous law with \(\gamma\ge216\), for
any finite edge set \(S\) and specified integers
\((k_e:e\in S)\), almost surely in the exterior,
\[
 \mathbb P(K_e=k_e\text{ for all }e\in S\mid K_{S^c})
             \le\left(\frac4{\sqrt\gamma}\right)^{|S|}.
 \tag{V72.19}\label{r72:pattern-bound}
\]
The unconditional bound is the same. In particular,
for the high-count path fixture,
\[
 \mathbb P(C_{\ell,\kappa})
             \le\left(\frac4{\sqrt\gamma}\right)^{4\ell+6}.
 \tag{V72.20}\label{r72:fixture-probability}
\]
\end{theorem}
\begin{proof}
In the product of the specified-count indicators,
condition on all counts except one edge in \(S\).
Its other factors are measurable and the remaining
conditional expectation is bounded by
\eqref{r72:atom-bound}. Use the tower property
and repeat, still conditioning on \(K_{S^c}\).
This proves \eqref{r72:pattern-bound}.
For \(C_{\ell,\kappa}\), take \(S\) to be its
\(4\ell+6\) transmitting edges, all fixed to
\(\kappa\), and drop the additional zero-count
restrictions. Integrating the conditional estimate
gives \eqref{r72:fixture-probability}.
\end{proof}

This is a probability bound for a \emph{specified}
pattern. A union over \(M\) allowed values at each
edge gives only \((4M/\sqrt\gamma)^{|S|}\), capped
by one. A band of typical fluctuations contains
order \(\sqrt\gamma\) values, so this union bound
does not give a small parameter for all such
patterns. In particular, it is not a small-probability
bound for arbitrary odd-count paths. The parity of
a single edge is fixed by all the other counts;
its conditional probability of being odd can be
one. We do not claim a Bernoulli domination of
odd edges.

\subsection{Consequence for the global attack and verification}

The high-count region cannot be replaced by a
product of scalar contractions simply from the
sizes of its counts. A correct path elimination
retains \(P_0\) or \(P_1\) according to parity and
pays the higher-channel remainder. The four-terminal
limit demonstrates that the retained odd channels
can contribute to even source responses after
closing through a network. They cannot be discarded
as unobserved or non-scalar intermediate states.

The next problem is therefore the joint network of
these retained channels and the actual probabilities
of its count/parity configurations, including
branching vertices. A path with two geometric
neighbors per internal vertex is not the full cubic
network. The present remainder bound and pattern
probability do not prove that this network is weakly
coupled, spatially summable, or gapped.

The full bank remains
\[
 \mathcal B_{\rm w}=(2\gamma I+3D)\Gamma
       +\gamma[\Cov(\tau_e,K_f)]_{e,f}
       -2\gamma^2(1-\E\tau_e/\gamma)I-6C^{\rm ret}.
 \tag{V72.21}\label{r72:bank}
\]
V71's moving-source term, exterior-mean covariances,
returns, factorial contacts, and parity/winding
mixture all remain. The transfer operators in this
appendix integrate fixed-count spatial paths;
neither they nor the auxiliary count clocks are
the physical Yang--Mills time generator.

The checker verifies the polynomial spectral
expansion and eigenvalue ratios exactly, compares
spectral path convolution against direct sphere
moments, and checks four-terminal coefficients by
independent sphere-pairing sums. Finite count
truncations additionally test normalized even/odd
path limits and the explicit error bounds. Those
truncations are diagnostics; the infinite count
sums are controlled analytically by the uniform
coefficient bounds above.

V71's mathematical body is preserved byte-for-byte
apart from the version title, and its regression
passes. Only the active predecessor is replaced;
all other manuscripts are unchanged. Build products
and certificates stay outside \path{latex_notes}.
The full spatial bound, interacting four-dimensional
continuum construction, and physical Yang--Mills
mass gap remain unproved.
% END CONSOLIDATED V72 APPEND

% BEGIN CONSOLIDATED V73 APPEND
\clearpage
\section[V73: Angular labels and junction weights]{V73: a positive angular-label law, junction contractions, and the cost of a fixed cutoff}
\label{r73:main}

\subsection{Going upstream from prescribed paths to normalized networks}

V72 paid the complete fixed-count path calculation. It did not
sum count patterns or control branching. This appendix does both
algebraic operations exactly in the native zero-field \(O(4)\)
model, while distinguishing an exact representation from a
spatial mixing theorem. The new conditional variables are angular
representation labels, not independent Poisson intensities.
Given these labels, the counts on different edges are independent;
the labels themselves retain an interacting network law.

This supplies a positive, fully normalized setting for the
retained \(P_0,P_1\) channels. A four-leg junction has a genuine
scalar/non-scalar separation, but its normalization and the
probability of retaining only these channels must both be paid.
On a cycle we prove that every fixed angular cutoff loses all
its probability as the activity grows. Thus the fixed-count,
long-path limit of V72 cannot be used as a uniform truncation
of the count-summed high-activity model.

\paragraph{Scope, nonduplication, and literature.}
All statements concern a finite loopless graph with free Haar
integration at its vertices, positive edge activities, and the
native dot-product interaction. They apply in particular to the
even cubic torus of V62--V72. Parallel edges are allowed for
blocked-network examples. Arbitrary transported Wilson
exteriors, boundary spin insertions, and the physical transfer
operator are not included without an additional argument.
V19's character normalization, V62--V68's sphere/wire moments,
and V72's eigenvalues are inherited tools, not new results.
The four-leg tensor below specializes those moments; the new
use is its normalized network role and comparison with the
actual count-summed law.

Positivity of these closed character networks is classical:
see H.~Pfeiffer,
\href{https://arxiv.org/abs/gr-qc/0211106}
{\emph{Positivity of relativistic spin network evaluations}},
Proposition 2.6 and Corollary 3.5, with the \(V\otimes V^*\)
convention. Its Section 4 warns against extending positivity
to arbitrary Haar-contraction diagrams. We give the
matrix-index proof needed here and do not import a general
gauge spin-foam positivity assertion. The pairing context is
also inherited from
\href{https://arxiv.org/abs/1807.06564}
{Benassi--Ueltschi, \emph{Loop correlations in random wire models}};
no global \(O(4)\) loop-distribution theorem is assumed.
This targeted check does not replace the scheduled broad
review at V80; the next companion checkpoint remains V75.

\subsection{Closed character networks have nonnegative weights}

Write \(G=(V,E)\), orient every edge arbitrarily, and set
\[
 d_j=j+1,\qquad
 S_G(J)=\int_{(S^3)^V}\prod_{e=xy}P_{J_e}(u_x,u_y)
                         \prod_{x\in V}dH(u_x),
 \quad J\in\mathbb N_0^E.
 \tag{V73.1}\label{r73:S}
\]
Here \(P_j=d_j\chi_j\) is the full harmonic projector kernel
of \eqref{r72:projector}; it is not pointwise positive for
general \(j\).

\begin{lemma}[Closed-network positivity with its convention fixed]
\label{r73:positive}
For every label assignment,
\[
 0\le S_G(J)\le\prod_e d_{J_e}^{\,2}.
 \tag{V73.2}\label{r73:S-bound}
\]
It vanishes if \(\sum_{e\ni x}J_e\) is odd at any vertex.
The latter condition is necessary, not asserted sufficient
for arbitrary higher labels.
\end{lemma}
\begin{proof}
Identify \(u_x\) with \(g_x\in\SU(2)\). Let \(D_j\) be
the unitary irreducible matrix representation of dimension
\(d_j\). On an oriented edge \(x\to y\),
\[
 \chi_j(g_x^{-1}g_y)
       =\sum_{a,b=1}^{d_j}
           \overline{D_j(g_x)_{ab}}D_j(g_y)_{ab}.
 \tag{V73.3}\label{r73:edge-indices}
\]
At a vertex \(x\), form the tensor product of the
representations on incoming edges and their conjugates
on outgoing edges. Haar integration of that tensor
representation is the orthogonal projection onto its
invariant vectors: invariance, idempotence, and
self-adjointness follow respectively from Haar translation,
convolution, and inversion. Choose an orthonormal basis
\(I_{x,\alpha}\) of this invariant space. The integrated
matrix entries, with all row indices denoted \(a\) and
all column indices \(b\), are therefore
\(\sum_\alpha I_{x,\alpha}(a)\overline{I_{x,\alpha}(b)}\).
If the invariant space is zero, this is the zero matrix.

Expand \eqref{r73:edge-indices} on every edge and integrate
at every vertex. Each edge identifies its row index at
its two ends and independently identifies its column
index. Consequently,
\[
 \int\prod_e\chi_{J_e}(g_x^{-1}g_y)\prod_xdH(g_x)
   =\sum_{(\alpha_x)}
       \left|\sum_{(a_e)}
                \prod_x I_{x,\alpha_x}((a_e)_{e\ni x})
       \right|^2\ge0.
 \tag{V73.4}\label{r73:squares}
\]
The contractions pair a representation with its conjugate
according to the chosen orientation; no convention-changing
sign has been inserted. Multiplication by \(\prod_e d_{J_e}\)
proves positivity for \(S_G\). The absolute character bound
gives its upper bound. Finally replacing \(g_x\) by \(-g_x\)
multiplies the integrand by
\((-1)^{\sum_{e\ni x}J_e}\), proving the parity selection rule.
\end{proof}

\subsection{An exact positive lift of the native count law}

For counts \(K\in\mathbb N_0^E\), put
\[
 M_G(K)=\int\prod_{e=xy}(u_x\cdot u_y)^{K_e}\prod_xdH(u_x),
 \qquad
 Z_G(\boldsymbol\gamma)=
       \sum_K M_G(K)\prod_e\frac{\gamma_e^{K_e}}{K_e!}.
 \tag{V73.5}\label{r73:count}
\]
This is the inherited native count normalizer, equal to
the integral of \(\exp\{\sum_e\gamma_e u_x\cdot u_y\}\).
It is finite and strictly positive. Define the eigenvalue
\(\lambda_j(k)\) by \eqref{r72:lambda}, with value zero
outside its stated support, and define
\[
 a_j(\gamma)=\sum_{k\ge0}\frac{\gamma^k}{k!}\lambda_j(k)
   =\sum_{n\ge0}
       \frac{(\gamma/2)^{j+2n}}{n!(n+j+1)!}>0.
 \tag{V73.6}\label{r73:aj}
\]
These series, rather than any asymptotic special-function
formula, are the definitions used below. In standard notation
\(a_j(\gamma)=2I_{j+1}(\gamma)/\gamma\).

\begin{theorem}[Positive joint law and conditionally independent counts]
\label{r73:lift}
The following is a probability law whose count marginal is
the original native count law:
\[
 {\mathbb P}_G(K,J)=
 \frac{S_G(J)}{Z_G(\boldsymbol\gamma)}
       \prod_e\frac{\gamma_e^{K_e}}{K_e!}
                       \lambda_{J_e}(K_e).
 \tag{V73.7}\label{r73:joint}
\]
Its label marginal and conditional count law are
\[
 \nu_G(J)=\frac{S_G(J)\prod_e a_{J_e}(\gamma_e)}
                       {Z_G(\boldsymbol\gamma)},\qquad
 {\mathbb P}_G(K\mid J)=\prod_e Q_{J_e,\gamma_e}(K_e),
 \tag{V73.8}\label{r73:marginals}
\]
where, for \(n\ge0\),
\[
 Q_{j,\gamma}(j+2n)=
    \frac{(\gamma/2)^{j+2n}}
         {a_j(\gamma)n!(n+j+1)!}.
 \tag{V73.9}\label{r73:Q}
\]
The conditional assertion is for \(\nu_G\)-positive label
assignments. In particular \(J_e\le K_e\) and
\(J_e\equiv K_e\pmod2\) almost surely.
\end{theorem}
\begin{proof}
For fixed \(K\), substitute V72's finite expansion into
each power in \(M_G(K)\) and integrate:
\[
 M_G(K)=\sum_J S_G(J)\prod_e\lambda_{J_e}(K_e).
 \tag{V73.10}\label{r73:moment-expansion}
\]
This is a finite sum of nonnegative terms by
Lemma~\ref{r73:positive}. Multiplication by the factorial
activities and summation over \(K\) prove that
\eqref{r73:joint} has the claimed marginal and total one.
Tonelli's theorem now applies to these nonnegative weights.
It gives the label normalizer and the product conditional
law, including all constants. Substitution \(k=j+2n\)
gives \eqref{r73:aj} and \eqref{r73:Q}.
\end{proof}

This does not contradict V69's obstruction to a positive
joint mixture of independent parity-conditioned Poisson
laws. The laws in \eqref{r73:Q} are different; for example
their consecutive supported weight ratio is
\[
 \frac{Q_{j,\gamma}(k+2)}{Q_{j,\gamma}(k)}
           =\frac{\gamma^2}{(k-j+2)(k+j+4)}.
 \tag{V73.11}\label{r73:Q-ratio}
\]
Nor does conditional independence assert independence
of \(J_e\): \(S_G(J)\) carries the vertex constraints and
the closed-network contractions.

There is an exact covariance consequence. Set
\[
 m_j(\gamma)=\gamma\,\partial_\gamma\log a_j(\gamma),
 \qquad
 v_j(\gamma)=\gamma\,\partial_\gamma m_j(\gamma).
\]
Differentiating the absolutely convergent positive series
shows that these are the conditional mean and variance
of \(K\) under \(Q_{j,\gamma}\). Total covariance gives
\[
 \Gamma_{ef}
   =\mathbf1_{\{e=f\}}\,\E_{\nu_G}v_{J_e}(\gamma_e)
       +\Cov_{\nu_G}\bigl(m_{J_e}(\gamma_e),
                         m_{J_f}(\gamma_f)\bigr).
 \tag{V73.12}\label{r73:cov}
\]
All second moments are finite: the native factorial sum
with any fixed count polynomial is bounded absolutely by
the corresponding product exponential series, since
\(|M_G(K)|\le1\). Conditional Jensen then justifies the
displayed total covariance. The last covariance, with
its possible signed off-diagonal entries, remains to
be estimated. It is not removed by the positive lift.

\subsection{The retained parity network and its true weights}

Let \(\mathcal L=\{J_e\in\{0,1\}\text{ for every }e\}\).
Under the joint law this is exactly the event
\(J_e=K_e\bmod2\) on every edge. Write
\[
 \phi(\gamma)=a_0(\gamma),\qquad
 x(\gamma)=\frac{a_1(\gamma)}{a_0(\gamma)}
                         =\frac{\phi'(\gamma)}{\phi(\gamma)}.
 \tag{V73.13}\label{r73:x}
\]
The derivative identity follows directly from the even
moment series for \(\phi\): its derivative has coefficient
\(c_{k+1}=\lambda_1(k)\) for odd \(k\).
For an edge subset \(F\), let \(S_G(F)\) mean \(S_G(J)\)
with \(J=\mathbf1_F\).

\begin{proposition}[Normalized retained sublaw]
\label{r73:retained}
The total unnormalized weight of \(\mathcal L\) is
\[
 Z_G^{\rm ret}
   =\prod_e\phi(\gamma_e)\,
        \sum_{\substack{F\subseteq E\\ \deg_F(x)\ {\rm even}\ \forall x}}
                S_G(F)\prod_{e\in F}x(\gamma_e),
 \qquad
 {\mathbb P}_G(\mathcal L)=\frac{Z_G^{\rm ret}}{Z_G}.
 \tag{V73.14}\label{r73:Zret}
\]
If \(\deg_F(x)=2h_x\), then
\[
 S_G(F)=\prod_x\frac{2^{h_x}}{(h_x+1)!}
                  \sum_{\pi}4^{\,\ell(\pi)}.
 \tag{V73.15}\label{r73:ret-loop}
\]
The sum is over all pairings of the incident \(F\)-edges
at each vertex, and \(\ell(\pi)\) is the number of closed
loops formed by these pairings. Empty vertices have
factor one. Thus every even-degree \(F\) has positive weight.
\end{proposition}
\begin{proof}
Restrict \eqref{r73:marginals} to \(J_e=0,1\);
\(P_0=1\), and the parity rule discards odd-degree subsets.
To evaluate the others write
\(P_1(u,v)=\sum_{a=1}^4(2u_a)(2v_a)\). The uniform
sphere moment at \(2h\) legs is
\[
 \int\prod_{i=1}^{2h}(2u_{a_i})\,dH(u)
   =\frac{2^h}{(h+1)!}
         \sum_{\text{pairings }\pi}
                     \prod_{\{i,j\}\in\pi}\delta_{a_i a_j}.
 \tag{V73.16}\label{r73:tensor}
\]
For completeness, a standard Gaussian vector is its
independent radius times a uniform sphere vector.
Its \(2h\)-th radius moment in four dimensions is
\(2^h(h+1)!\); expanding its Gaussian moment into
pairings and multiplying by \(2^{2h}\) gives the formula.
The Gaussian pairing rule itself follows by
differentiating \(\exp(|t|^2/2)\).
Every resulting closed colour line can have four
values, giving \eqref{r73:ret-loop}. Formula
\eqref{r73:Zret} includes the full conditional normalizer.
\end{proof}

For a single cycle of length \(L\), the two allowed
subsets are empty and the whole cycle, with weights
\(1\) and \(4x^L\). For four parallel edges between two
vertices the retained polynomial is
\[
 1+24x^2+32x^4.
 \tag{V73.17}\label{r73:four-parallel}
\]
Indeed each of the six two-edge subsets has weight \(4\).
For all four edges, the two vertices have three pairings
each: three equal choices create two loops, and the other
six choices create one. Multiplying
\(3\cdot16+6\cdot4=72\) by \((2/3)^2\) gives \(32\).
There is no unaccounted independent choice of loops at
shared vertices.

The same retained partition function can be written
algebraically as the integral of
\(\prod_e[\phi(\gamma_e)+4\phi'(\gamma_e)u_x\cdot u_y]\).
This integrand need not be positive. We use its positive
label/loop weights, not a spin probability with this
signed density. If \(p=Z_G^{\rm ret}/Z_G\), the count
marginal decomposes as \(p\mu_{\rm ret}+(1-p)\mu_{\rm rest}\),
where the second term may be omitted when \(p=1\).
Hence
\[
 \|\mu-\mu_{\rm ret}\|_{\rm TV}\le1-p.
 \tag{V73.18}\label{r73:mixture-error}
\]
Using this discarded-weight bound to justify an
approximation requires \(1-p\) to be small;
positivity alone gives no such small error.

\subsection{A concrete four-leg junction contracts only after normalization}

View the \(h=2\) tensor in \eqref{r73:tensor} as a map
on \(4\)-by-\(4\) matrices, with two incoming and two
outgoing colour indices. For the Hilbert--Schmidt norm,
\[
 (\mathcal A X)_{ab}
   =\frac23\{\delta_{ab}\Tr X+X_{ab}+X_{ba}\}
   =4\Pi_{\rm tr}X+\frac43\Pi_{\rm st}X.
 \tag{V73.19}\label{r73:junction}
\]
Here \(\Pi_{\rm tr}X=(\Tr X)I/4\), and
\(\Pi_{\rm st}X=(X+X^{\mathsf T})/2-\Pi_{\rm tr}X\).
The remaining antisymmetric subspace is killed.
These are orthogonal subspaces of dimensions \(1,9,6\).
The formula follows by inserting the three pairings
in \eqref{r73:tensor}; it proves the decomposition and
all the eigenvalues without a numerical diagonalization.
For a chain of \(n\ge1\) such junction maps,
\[
 4^{-n}\mathcal A^n
       =\Pi_{\rm tr}+3^{-n}\Pi_{\rm st},\qquad
 \|4^{-n}\mathcal A^n-\Pi_{\rm tr}\|=3^{-n}.
 \tag{V73.20}\label{r73:junction-chain}
\]
The factors \(4^n\) are normalization factors, not
optional losses in a connected-correlation bound.

There is a pointwise-positive realization of this
particular paired geometry. Two retained odd paths
between the same two junction spins contribute
\(P_1(u,v)^2=16(u\cdot v)^2\). Its row integral is \(4\).
The normalized convolution kernel is therefore
\[
 Q(u,v)=4(u\cdot v)^2,\qquad
 Q=P_0+\tfrac13P_2,\qquad
 Q^n=P_0+3^{-n}P_2.
 \tag{V73.21}\label{r73:paired-Q}
\]
The spectral identity is also the \(k=2\) case of
V72. In a stationary chain with Haar marginal and
transition \(Q\), centered \(f,g\in L^2(H)\) consequently
satisfy
\[
 |\Cov(f(U_0),g(U_n))|
       \le3^{-n}\|f\|_2\|g\|_2.
 \tag{V73.22}\label{r73:paired-cov}
\]
This proves mixing for the specified paired chain.
It does not assert that the cubic network decomposes
into such chains, or that its boundary contractions
and higher angular sectors can be dropped.

\subsection{High activity makes the retained edge weight large}

The marginal \(t=u\cdot v\) of two Haar spins has
density \(2\sqrt{1-t^2}/\pi\) on \([-1,1]\).
Let \(\rho_\gamma\) be its exponential tilt by
\(e^{\gamma t}/\phi(\gamma)\). Then
\(x(\gamma)=\E_{\rho_\gamma}t\).

\begin{lemma}[A direct endpoint estimate]
\label{r73:endpoint}
For every \(\gamma>0\),
\[
 0<x(\gamma)<1,\qquad
 1-x(\gamma)\le\frac3{2\gamma}.
 \tag{V73.23}\label{r73:x-bound}
\]
For \(r_j(\gamma)=a_j(\gamma)/a_0(\gamma)\),
\[
 0<r_j(\gamma)\le1,\qquad
 r_j(\gamma)\ge
       1-\frac{\pi^2j(j+2)}{8\gamma}.
 \tag{V73.24}\label{r73:r-bound}
\]
\end{lemma}
\begin{proof}
Positivity of \(x\) follows by pairing \(t\) and \(-t\)
in its numerator, and strict inequality \(x<1\) follows
from the nondegenerate continuous density. Set \(y=1-t\).
Its tilted density is proportional to
\(y^{1/2}e^{-\gamma y}h(y)\) on \([0,\infty)\), where
\(h(y)=\sqrt{(2-y)_+}\) is nonincreasing.
Under the gamma law of shape \(3/2\) and rate \(\gamma\),
\[
 \Cov(Y,h(Y))
     =\tfrac12\E[(Y-Y')(h(Y)-h(Y'))]\le0.
\]
Division by \(\E h(Y)>0\) gives
\(\E_{\rho_\gamma}(1-t)\le\E Y=3/(2\gamma)\).

Character orthogonality applied to
\(e^{\gamma t}=\sum_j a_j(\gamma)P_j(t)\) gives
\[
 r_j(\gamma)=\E_{\rho_\gamma}
                    \frac{\chi_j(t)}{j+1}.
 \tag{V73.25}\label{r73:ratio-average}
\]
The expansion is absolutely uniform: at \(t=1\),
\(\sum_j (j+1)^2a_j(\gamma)=e^\gamma\), either from
the positive power expansion or its trace identity.
The bound \(r_j\le1\) follows from the character bound,
and strict positivity from \eqref{r73:aj}. If
\(t=\cos\theta\), \(0\le\theta\le\pi\), the character is
the sum of the \(j+1\) phases with frequencies
\(j,j-2,\ldots,-j\). Their average squared frequency
is \(j(j+2)/3\). Hence \(1-\cos z\le z^2/2\) implies
\[
 1-\frac{\chi_j(\cos\theta)}{j+1}
       \le\frac{j(j+2)}6\theta^2
       \le\frac{\pi^2j(j+2)}{12}(1-t).
\]
The last step uses
\(\sin(\theta/2)\ge\theta/\pi\).
Taking expectations and using the endpoint estimate
proves \eqref{r73:r-bound}.
\end{proof}

Thus \(x(\gamma)\to1\), not zero, at high activity.
Already for \(\gamma>2\), the lower bound gives \(x>1/4\),
so the retained edge polynomial \(1+4xt\) is negative
on a nonempty interval near \(t=-1\). The positivity
used in \eqref{r73:Zret} is after the closed-network
integration, not pointwise spin positivity.

\subsection{A fixed angular cutoff loses its weight on a cycle}

\begin{theorem}[Count-summed cycle and cutoff obstruction]
\label{r73:cycle}
On a cycle \(C_L\), \(L\ge3\), at homogeneous activity
\(\gamma>0\), the label law is supported on constant
assignments \(J_e=j\), with probabilities
\[
 {\mathbb P}_{C_L}(J_e=j\ \forall e)
     =\frac{(j+1)^2r_j(\gamma)^L}
             {\sum_{i\ge0}(i+1)^2r_i(\gamma)^L}.
 \tag{V73.26}\label{r73:cycle-law}
\]
For a fixed cutoff \(J_*\ge0\), let
\(p_{J_*}(\gamma,L)={\mathbb P}(\max_e J_e\le J_*)\),
and put
\[
 H(q)=\sum_{j=0}^q(j+1)^2
           =\frac{(q+1)(q+2)(2q+3)}6.
\]
For every integer \(q\ge1\) such that
\(\gamma\ge(\pi^2/4)Lq(q+2)\),
\[
 p_{J_*}(\gamma,L)
             \le\min\{1,\,2H(J_*)/H(q)\}.
 \tag{V73.27}\label{r73:cutoff-bound}
\]
In particular, for \(\gamma\ge8L\),
\[
 p_{J_*}(\gamma,L)
     \le\min\{1,\,6H(J_*)(8L/\gamma)^{3/2}\}.
 \tag{V73.28}\label{r73:cutoff-rate}
\]
Every fixed \(J_*,L\) therefore has
\(p_{J_*}(\gamma,L)\to0\) as \(\gamma\to\infty\).
\end{theorem}
\begin{proof}
Integrating a degree-two vertex convolves its two
projector kernels. Orthogonality annihilates different
labels; following the cycle forces a common label.
For that label the remaining integral is
\(\Tr P_j=(j+1)^2\). Summing the positive activities
gives \eqref{r73:cycle-law}. The denominator is finite,
since \(r_j\le1\) and
\(\sum_j(j+1)^2r_j=e^\gamma/a_0(\gamma)<\infty\).

Under the hypothesis on \(q\),
\eqref{r73:r-bound} gives
\(r_j\ge1-1/(2L)\) for \(j\le q\).
Bernoulli's inequality yields \(r_j^L\ge1/2\).
Thus the denominator in \eqref{r73:cycle-law} is
at least \(H(q)/2\), whereas its numerator summed
through \(J_*\) is at most \(H(J_*)\). This proves
\eqref{r73:cutoff-bound}.

For the explicit rate take \(q=\lfloor\sqrt{\gamma/(8L)}\rfloor
\ge1\). Since \(q(q+2)\le3q^2\) and
\(3\pi^2<32\), this \(q\) meets the preceding hypothesis.
Also \(H(q)\ge(q+1)^3/3>(\gamma/(8L))^{3/2}/3\).
Substitution proves \eqref{r73:cutoff-rate}.
\end{proof}

For the parity-leading truncation \(J_*=1\), this
probability is explicitly
\[
 \frac{1+4x(\gamma)^L}
      {\sum_{j\ge0}(j+1)^2r_j(\gamma)^L}.
 \tag{V73.29}\label{r73:lead-cycle}
\]
Although its numerator stays bounded, the denominator
retains an increasing range of angular modes. This is
an actual probability in the positive joint law, not
a union bound over individual count patterns.
Vanishing retained weight does not by itself prove
that the two count marginals have total variation
distance tending to one: different labels can produce
similar count distributions. It rules out a small
discarded-weight justification of this truncation.
The theorem gives a necessary scale for a cutoff
that is to keep a fixed positive probability when
\(\gamma/L\) grows: since \(H(J_*)\le(J_*+1)^3\),
\eqref{r73:cutoff-rate} implies
\[
 p_{J_*}\ge\eta>0
 \quad\Longrightarrow\quad
 J_*+1\ge(\eta/6)^{1/3}\sqrt{\gamma/(8L)}
 \qquad(\gamma\ge8L).
 \tag{V73.30}\label{r73:necessary-cutoff}
\]
This is a necessary bound, not a sufficient truncation
theorem or a lattice-volume-uniform estimate.

The cycle can be taken even to fit bipartite geometry.
It is an isolated native cycle, not the law of a cycle
inside a fully interacting cubic exterior. That limitation
is sufficient for the present audit: a graph-independent
claim that the parity-leading channels alone dominate
at high activity is false. V72 held counts and activities
fixed and then lengthened paths; the limit here holds
cycle length fixed and increases activity. There is no
contradiction and no unproved interchange of limits.

\subsection{What has changed in the global attack}

There is now an exact positive angular-label law, a
conditional product description of counts, and an
evaluated junction contraction. These are possible
inputs to a source-sensitive spatial estimate, not
that estimate itself. The count covariance problem
can be moved to \eqref{r73:cov}; it has not disappeared.
In particular one must control the interacting label
covariance and the conditional expectations of the
actual deficit/return sources under the same law.

The cutoff obstruction means that a fixed \(P_0,P_1\)
network alone is not a justified high-activity
approximation before suitable geometric blocking.
A next calculation should retain a scale-dependent
angular range, evaluate its junction/source weights,
or bound the full label covariance directly. It must
also pay the mass outside any proposed truncation.
Neither a small prescribed-count probability nor the
paired-chain factor \(1/3\) supplies that missing step.

The full native compensation still reads
\[
 \mathcal B_{\rm w}=(2\gamma I+3D)\Gamma
       +\gamma[\Cov(\tau_e,K_f)]_{e,f}
       -2\gamma^2(1-\E\tau_e/\gamma)I-6C^{\rm ret}.
 \tag{V73.31}\label{r73:bank}
\]
All factorial contacts, returns, exterior means,
source motion, and parity/winding mixtures remain.
The auxiliary labels and junction transfers are not
the physical Yang--Mills Hamiltonian or its time.

The accompanying exact checker compares the positive
label expansion with independent sphere moments on
finite graphs, verifies the conditional covariance
identity on finite count sums, and checks the junction
projectors and normalized powers with rational
arithmetic. Bessel-series cycle evaluations are
additional diagnostics, not proofs of infinite sums
or the stated uniform inequalities. The proofs above
provide those statements.

V72's mathematical body is preserved byte-for-byte
apart from the title, and its mathematical regression
passes. Only the active predecessor is replaced;
the other manuscripts are unchanged. Certificates
and build output stay outside \path{latex_notes}.
The previous goal turn was progress. This one replaces
a pattern-wise heuristic by a positive normalized law
and proves the scale obstruction that the next
truncation must respect. The complete spatial bound,
interacting four-dimensional continuum construction,
and physical Yang--Mills mass gap remain unproved.
% END CONSOLIDATED V73 APPEND

% BEGIN CONSOLIDATED V74 APPEND
\clearpage
\section[V74: Signed sources and angular moments]{V74: the actual deficit source, angular moment domination, and a compensated label covariance}
\label{r74:main}

\subsection{The source must be transferred, not guessed}

V73 introduced a positive joint law for native counts \(K\)
and angular labels \(J\). Conditional count independence
does not identify the conditional expectation of the
actual deficit source. Here we compute its count-summed
transfer, including the label born by pair insertion
and the same-edge derivative contact. The resulting
label source is signed, even though the original count
deficit is nonnegative.

To obtain estimates rather than another formal
representation, we prove stochastic domination of a
single angular label conditional on the entire count
vector. It yields graph-size-independent moments through
order four of the angular Casimir. Elementary differential
inequalities for the conditional count law then bound the
transferred sources uniformly in activity. Finally we
combine the count and deficit contributions before
taking norms: their off-diagonal combination becomes a
covariance of two label observables with bounded moments.
Spatial summation, the return contribution, and the full
physical Yang--Mills problem are still open.

\paragraph{Scope and inherited inputs.}
Use the finite simple loopless native zero-field \(O(4)\)
graph, positive activities \(\gamma_e\), and V73's joint
law \(\mathbb P(K,J)\), with count marginal \(\mu\)
and label marginal \(\nu\). V68's two-endpoint polynomial
and V72's exact \(\lambda_j(k)\) are inputs.
The fully compensated bank at the end is stated for
the homogeneous cubic setting in which its inherited
\(D\) and \(C^{\rm ret}\) are defined.
No angular positivity or moment estimate here is
automatically exported to a transported Wilson exterior.

NIST DLMF
\href{https://dlmf.nist.gov/10.25.E1}{10.25.1} and
\href{https://dlmf.nist.gov/10.25.E2}{10.25.2}
record the standard modified Bessel equation and series.
We use the series already defined in V73 and derive
the differential identities below coefficient by
coefficient. Bounds on Bessel ratios are a classical
subject; no external asymptotic or mixing theorem is
required by the proofs here. The scheduled companion
review is next, at V75; the broad sweep remains V80.

\subsection{Pair insertion creates a new label}

On the admissible count support let
\[
 q_e(K)=\frac{M_G(K+2e)}{M_G(K)},\qquad
 \tau_e(K)=\gamma_e(1-q_e(K)).
 \tag{V74.1}\label{r74:deficit}
\]
V68--V71 give \(0\le\tau_e\) and \(q_e\le1\).
For \(k=K_e\) and \(j\le k\) of the same parity,
factorial cancellation gives
\[
 R_j(k):=\frac{\lambda_j(k+2)}{\lambda_j(k)}
       =\frac{(k+2)(k+1)}{(k-j+2)(k+j+4)}.
 \tag{V74.2}\label{r74:ratio}
\]
The old posterior expectation is not the whole answer:
\[
 q_e(K)=\E[R_{J_e}(k)\mid K]+\mathfrak b_e(K),
 \tag{V74.3}\label{r74:birth}
\]
where the nonnegative missing term is
\[
 \mathfrak b_e(K)=
 \frac{\lambda_{k+2}(k+2)}{M_G(K)}
 \sum_{J_{e^c}} S_G(k+2,J_{e^c})
                       \prod_{f\ne e}\lambda_{J_f}(K_f).
 \tag{V74.4}\label{r74:birth-weight}
\]
Indeed expand \(M_G(K+2e)\) by \eqref{r73:moment-expansion}.
The terms with \(J_e\le k\) give the posterior expectation,
and \(J_e=k+2\) is the only remaining label. V73's
closed-network positivity proves its nonnegativity.

The term need not vanish. On a four-cycle, take selected
count zero and the other three counts equal to two.
The cycle trace gives
\[
 M_G(0,2,2,2)=\frac1{64},\qquad
 q_e(0,2,2,2)=\frac14+\frac1{36}=\frac5{18}.
 \tag{V74.5}\label{r74:birth-example}
\]
The old posterior has \(J_e=0\), so it gives only
\(R_0(0)=1/4\). The additional term is
\(9(1/12)^4/(1/64)=1/36\), from the constant
label-two cycle. This example concerns a native
isolated cycle, not an arbitrary cubic exterior.

\subsection{The exact signed source and its diagonal contact}

Keep \(a_j,m_j,v_j\) from V73 and put
\[
 \kappa_j=j(j+2),\qquad
 s_j(\gamma)=\gamma\left(1-\frac{a_j''(\gamma)}{a_j(\gamma)}\right).
 \tag{V74.6}\label{r74:s}
\]
The convergent power series for \(a_j\) satisfies
\[
 \gamma^2a_j''+3\gamma a_j'-(\gamma^2+\kappa_j)a_j=0.
 \tag{V74.7}\label{r74:ode}
\]
For a coefficient of degree \(k=j+2n\), its derivative
factor is \(k(k+2)-j(j+2)=4n(n+j+1)\);
this is exactly the ratio to the preceding series
coefficient. This proves the equation also at its
lowest degree. In particular
\[
 s_j(\gamma)=\frac{3m_j(\gamma)-\kappa_j}{\gamma}.
 \tag{V74.8}\label{r74:s-casimir}
\]

\begin{theorem}[Native source transfer with the contact retained]
\label{r74:transfer}
For a bounded function \(F(K_{e^c})\), define
\(\overline F(J)=\E[F(K_{e^c})\mid J]\) under V73's
product conditional count law. Then
\[
 \E_\mu[\tau_e F]=
       \E_\nu[s_{J_e}(\gamma_e)\overline F(J)].
 \tag{V74.9}\label{r74:outside-source}
\]
In particular \(\E_\mu\tau_e=\E_\nu s_{J_e}(\gamma_e)\).
Writing \(m_e=m_{J_e}(\gamma_e)\),
\(v_e=v_{J_e}(\gamma_e)\), and \(s_e=s_{J_e}(\gamma_e)\),
the count covariance with the deficit is
\[
 \Cov_\mu(\tau_e,K_f)
   =\Cov_\nu(s_e,m_f)
      +\mathbf1_{\{e=f\}}\,
           \E_\nu\left(\frac{3v_e}{\gamma_e}-2s_e\right).
 \tag{V74.10}\label{r74:source-cov}
\]
This is not the assertion \(s_e=\E[\tau_e\mid J]\).
\end{theorem}
\begin{proof}
Multiply \(q_e\) by the count density; its denominator
\(M_G(K)\) cancels. Expanding the remaining
\(M_G(K+2e)\) gives, on edge \(e\),
\[
 \sum_{k\ge0}\frac{\gamma_e^k}{k!}\lambda_j(k+2)
                           =a_j''(\gamma_e).
 \tag{V74.11}\label{r74:shifted-series}
\]
This includes the term \(k=j-2\) when \(j\ge2\):
the born label has not been lost. The other edge
factors remain their original \(a_{J_f}\) and their
conditional \(Q_{J_f,\gamma_f}\). Summing nonnegative
terms first and then taking bounded signed \(F\)
proves \eqref{r74:outside-source}.
For \(f\ne e\), the same calculation with \(F=K_f\)
gives \(\E(\tau_e K_f)=\E_\nu(s_e m_f)\).
All relevant moments are finite by the factorial
series, so this unbounded insertion is justified.

For \(f=e\), the marked shifted series is
\(\gamma_e a_j'''(\gamma_e)\). Consequently
\[
 \E_\mu(\tau_e K_e)
  =\E_\nu\left[
       \gamma_e m_e-\gamma_e^2
                     \frac{a_{J_e}'''(\gamma_e)}{a_{J_e}(\gamma_e)}
           \right]
  =\E_\nu\left[s_em_e+\gamma_e s_{J_e}'(\gamma_e)-s_e\right].
\]
Differentiating \eqref{r74:s-casimir} and using
\(\gamma m_j'=v_j\) shows
\(\gamma s_j'-s_j=3v_j/\gamma-2s_j\).
Subtracting the means proves \eqref{r74:source-cov}.
The moment bounds proved below also directly verify
integrability of the label expressions.
\end{proof}

The label function really can be negative. For
\(j\ge2\), its series has \(m_j(\gamma)=j+O(\gamma^2)\)
as \(\gamma\downarrow0\), so
\(s_j(\gamma)=j(1-j)/\gamma+O(\gamma)<0\) for
sufficiently small \(\gamma\). Constant label \(j\)
on a cycle has positive \(\nu\)-probability.
Since \(\tau_e\ge0\), this also disproves the proposed
conditional-expectation interpretation. V68's bounds
on \(\tau_e\) must not be transferred to \(s_e\) by
a nonexistent conditional Jensen inequality.

\subsection{Conditional angular domination by a finite reference law}

For a fixed count \(k\), define a probability on
\(j=k\bmod2,k\bmod2+2,\ldots,k\) by
\[
 \pi_k(j)=(j+1)^2\lambda_j(k).
 \tag{V74.12}\label{r74:pi}
\]
It is a probability because the kernel identity
\(t^k=\sum_j\lambda_j(k)P_j(t)\), evaluated at \(t=1\),
sums these weights to one.

\begin{theorem}[Every-count posterior domination]
\label{r74:domination}
For every admissible full count vector \(K\),
the conditional law of \(J_e\) is stochastically
dominated by \(\pi_{K_e}\). That is, every increasing
real function \(f\) on the finite label support obeys
\[
 \E[f(J_e)\mid K]\le \sum_j\pi_{K_e}(j)f(j).
 \tag{V74.13}\label{r74:stochastic}
\]
\end{theorem}
\begin{proof}
Use V68's two-endpoint collapse at \(e=xy\). After
integrating all other spins its exterior polynomial
is a positive common constant times
\(\sum_h w_h t^h\), where \(w_h\ge0\) and every \(h\)
has the parity of \(k=K_e\). Marginalizing the other
labels in V73's finite fixed-count expansion gives
\[
 \mathbb P(J_e=j\mid K)
   =\frac{\pi_k(j)b_j}{\sum_l\pi_k(l)b_l},
 \qquad b_j=\sum_h w_h\lambda_j(h).
 \tag{V74.14}\label{r74:posterior}
\]
Here character orthogonality gives
\(\int P_j(t)t^h\,dH=(j+1)^2\lambda_j(h)\);
the common exterior constant cancels.
For fixed \(h\), the allowed-parity sequence
\(\lambda_j(h)\) decreases in \(j\), by V72's ratio
\((h-j)/(h+j+4)\), and is zero above \(h\).
Thus \(b_j\) is nonincreasing on the support of \(\pi_k\).
Its normalizer is positive because \(K\) is admissible.
For independent \(X,X'\) with law \(\pi_k\),
\[
 \Cov(f(X),b_X)
   =\tfrac12\E[(f(X)-f(X'))(b_X-b_{X'})]\le0.
\]
Dividing by \(\E b_X>0\) proves the claim.
\end{proof}

This is domination of one label conditional on all
counts, not independent domination of the label network.
It is enough for the single-source moments needed next.

\subsection{Uniform angular moments and a sufficient local cutoff}

On zonal functions of \(t\in[-1,1]\), let
\[
 \mathcal C=-(1-t^2)\partial_t^2+3t\partial_t.
\]
The Gegenbauer differential equation, also obtained
by differentiating the character polynomial, gives
\(\mathcal C P_j=\kappa_jP_j\).
Repeatedly applying \(\mathcal C\) to \(t^k\) and
evaluating at \(1\) therefore computes the moments
of \(\kappa_j\) under \(\pi_k\). The monomial rule is
\[
 \mathcal C t^k=k(k+2)t^k-k(k-1)t^{k-2}.
\]
It gives, including \(k=0\),
\[
 \begin{split}
 \E_{\pi_k}\kappa&=3k,\\
 \E_{\pi_k}\kappa^2&=15k^2-6k\le15k^2,\\
 \E_{\pi_k}\kappa^4
   &=945k^4-2940k^3+3660k^2-1584k
       \le945k^4.
 \end{split}
 \tag{V74.15}\label{r74:reference-moments}
\]
For the last inequality, the subtracted quantity is
\(k(2940k^2-3660k+1584)\ge0\): the quadratic has
positive leading coefficient and negative discriminant.
All three functions of \(\kappa_j\) used in the
domination theorem are increasing in \(j\).

The native factorial identity and the positive spin
measure imply
\[
 \E_\mu(K_e)_r
   =\gamma_e^r\E_{\rm spin}(u_x\cdot u_y)^r
       \le\gamma_e^r\qquad(r\in\mathbb N).
 \tag{V74.16}\label{r74:factorial}
\]
Here \((k)_r\) is a falling factorial. The equality
follows by differentiating the compact partition
function \(r\) times, or by reindexing its count
series. Its left side is nonnegative, and the spin
integrand is at most one. This does not assert a
product Poisson law or independent counts.
In particular
\(\E K_e^2\le\gamma_e^2+\gamma_e\) and
\(\E K_e^4\le\gamma_e^4+6\gamma_e^3+7\gamma_e^2+\gamma_e\).
Together with \eqref{r74:stochastic}, this proves
\[
 \begin{split}
 \E_\nu\kappa_{J_e}&\le3\gamma_e,\\
 \E_\nu\kappa_{J_e}^{\,2}
        &\le15(\gamma_e^2+\gamma_e),\\
 \E_\nu\kappa_{J_e}^{\,4}
        &\le945(\gamma_e^4+6\gamma_e^3+7\gamma_e^2+\gamma_e).
 \end{split}
 \tag{V74.17}\label{r74:angular-moments}
\]
These are bounds for the full normalized graph law,
with no volume factor.

For any finite edge set \(A\) and integer cutoffs
\(R_e\ge0\), let \(\mathcal T=\{J_e\le R_e\text{
 for all }e\in A\}\). Markov's inequality and a union
bound give the sufficient estimate
\[
 \nu(\mathcal T^c)
      \le \sum_{e\in A}
          \frac{3\gamma_e}{(R_e+1)(R_e+3)}.
 \tag{V74.18}\label{r74:cutoff}
\]
No independence is used. For a fixed number of edges
and fixed error tolerance this supplies a cutoff of
order \(\sqrt{\gamma}\), consistent with V73's
necessary cycle scale. For an entire growing graph
the displayed union bound has an explicit edge-count
cost; it is not a volume-independent global truncation.
The total variation distance between the original
count marginal and the one obtained by conditioning
the joint law on \(\mathcal T\) is at most
\(\nu(\mathcal T^c)\), by conditioning and marginalization.

\subsection{Conditional count variance and source moments without asymptotics}

\begin{lemma}[Uniform differential bounds for the Bessel-type count law]
\label{r74:conditional-bounds}
For every \(j\ge0,\gamma>0\),
\[
 m_j>\gamma-\tfrac32,\qquad
 0<v_j<\gamma,\qquad
 m_j\le\sqrt{\gamma^2+\kappa_j+1}-1\le\gamma+j.
 \tag{V74.19}\label{r74:mv}
\]
If \(\gamma\ge1\), define
\(\delta_j=m_j-\gamma\) and \(w_j=\gamma-v_j\).
Then
\[
 |\delta_j|\le\tfrac32+\frac{\kappa_j}{2\gamma},
 \qquad
 0\le w_j\le4+\frac{2\kappa_j}{\gamma}.
 \tag{V74.20}\label{r74:dw}
\]
\end{lemma}
\begin{proof}
The series defines a nondegenerate count distribution,
so \(v_j>0\). Equation \eqref{r74:ode} implies the
Riccati identity
\[
 v_j=\gamma^2+\kappa_j-m_j^2-2m_j,\qquad
 \gamma v_j'=2\gamma^2-2(m_j+1)v_j.
 \tag{V74.21}\label{r74:riccati}
\]
Near zero, \(m_j=j+\gamma^2/[2(j+2)]+O(\gamma^4)\).
The function \(m_j-(\gamma-3/2)\) starts positive.
At any zero its derivative would be
\((\kappa_j+3/4)/\gamma>0\), by the first identity.
It therefore cannot have a first crossing to
nonpositive values. This proves the first bound.

Put \(z=v_j/\gamma=m_j'\). Its differential equation is
\[
 z'=2-\frac{2m_j+3}{\gamma}z.
\]
It starts at zero in the limit \(\gamma\downarrow0\).
At \(z=1\) its derivative is strictly negative,
since \(m_j>\gamma-3/2\). Thus \(z<1\) by the
same first-crossing argument. Positivity of \(v_j\)
and the first identity in \eqref{r74:riccati} give
the upper square-root bound on \(m_j\); the final
inequality uses
\(\sqrt{\gamma^2+(j+1)^2}\le\gamma+j+1\).

The square-root bound also gives
\(m_j-\gamma\le-1+(\kappa_j+1)/(2\gamma)
\le\kappa_j/(2\gamma)\) for \(\gamma\ge1\).
This proves the bound on \(\delta_j\).
For \(w_j=\gamma-v_j\), the differential equation is
\[
 w_j'+\frac{2(m_j+1)}{\gamma}w_j=2\delta_j+3.
 \tag{V74.22}\label{r74:w-ode}
\]
On \(\gamma\ge1\), its coefficient is at least one,
and \(w_j\ge0\). If \(\gamma\ge2\), integrate
\(w_j'+w_j\le1+(\kappa_j+1)/\gamma\)
on \([\gamma/2,\gamma]\). Using
\(w_j(\gamma/2)\le\gamma/2\) gives
\[
 w_j(\gamma)\le(\gamma/2)e^{-\gamma/2}
                        +1+\frac{2(\kappa_j+1)}{\gamma}
                  \le4+\frac{2\kappa_j}{\gamma}.
\]
For \(1\le\gamma<2\), the last bound follows directly
from \(w_j\le\gamma<2\).
\end{proof}

For the rest of this subsection suppose every
activity under consideration is at least one.
Equations \eqref{r74:angular-moments} and
\eqref{r74:dw} imply, under the actual label law,
\[
 \|\delta_{J_e}\|_2<5,\qquad
 \|\delta_{J_e}\|_4<7,\qquad
 \|w_{J_e}\|_2<15,\qquad
 \|s_e\|_2<14.
 \tag{V74.23}\label{r74:norms}
\]
Indeed
\(\|\kappa_{J_e}/\gamma_e\|_2\le\sqrt{30}\) and
\(\|\kappa_{J_e}/\gamma_e\|_4\le14175^{1/4}<11\).
Minkowski gives the first three estimates.
For the last, \eqref{r74:mv} and \(3j\le\kappa_j\)
for integers \(j\ge0\) give
\[
 |s_j|\le3+\frac{2\kappa_j}{\gamma},
 \qquad 3+2\sqrt{30}<14.
\]
These estimates were proved for the signed source
itself, not borrowed from a bound on \(\tau_e\).

V73's total covariance identity now yields the
unconditional, graph-uniform estimates
\[
 \Var_\mu K_e\le\gamma_e+25,\qquad
 |\Cov_\mu(K_e,K_f)|\le25\quad(e\ne f),\qquad
 |\Cov_\mu(\tau_e,K_f)|\le70\quad(e\ne f).
 \tag{V74.24}\label{r74:entry-bounds}
\]
The first uses \(\E v_e\le\gamma_e\) and
\(\Var_\nu m_e=\Var_\nu\delta_{J_e}<25\);
the others use the exact covariance identities and
Cauchy--Schwarz. These are entrywise bounds,
not row-sum or covariance-operator bounds.
Also \eqref{r74:s-casimir}, \(\kappa_j\ge0\), and
\(\E K_e\le\gamma_e\) show
\[
 0\le\E_\mu\tau_e=\E_\nu s_e\le3.
 \tag{V74.25}\label{r74:mean-deficit}
\]

\subsection{The leading count--deficit covariance cancels}

In the homogeneous setting put
\[
 \mathcal H_j(\gamma)=
       \delta_j-\delta_j^2+w_j.
 \tag{V74.26}\label{r74:H}
\]
Its label-law norm is uniformly bounded:
\[
 \|\mathcal H_{J_e}\|_2
   \le\|\delta_{J_e}\|_2+\|\delta_{J_e}\|_4^2
                                +\|w_{J_e}\|_2<70.
 \tag{V74.27}\label{r74:H-norm}
\]
This is useful because the Riccati identity gives
the exact algebraic cancellation
\[
 (2\gamma+3)m_j-\kappa_j
       =2\gamma^2+\mathcal H_j.
 \tag{V74.28}\label{r74:algebra-cancel}
\]
For distinct edges, \eqref{r74:source-cov} therefore
proves
\[
 2\gamma\Gamma_{ef}
        +\gamma\Cov_\mu(\tau_e,K_f)
   =\Cov_\nu(\mathcal H_{J_e},\delta_{J_f}),\qquad
 \left|2\gamma\Gamma_{ef}
        +\gamma\Cov_\mu(\tau_e,K_f)\right|\le350.
 \tag{V74.29}\label{r74:compensated-entry}
\]
The activity-independent right side is obtained
after the cancellation. Bounding the two terms
separately by \eqref{r74:entry-bounds} would instead
leave an order-\(\gamma\) estimate.

\begin{proposition}[The same cancellation in the full source bank]
\label{r74:full-bank}
Retaining the exact diagonal contact from
\eqref{r74:source-cov}, the inherited full identity
can be rewritten as
\[
 \begin{split}
 \mathcal B_{\rm w}
   ={}&3D\Gamma+
       [\Cov_\nu(\mathcal H_{J_e},\delta_{J_f})]_{e,f}\\
      &+\operatorname{diag}_e
          \{(2\gamma+3)\E_\nu v_e-2\gamma^2\}
       -6C^{\rm ret}.
 \end{split}
 \tag{V74.30}\label{r74:bank}
\]
No part of \(D\Gamma\) or \(C^{\rm ret}\) has been
discarded or estimated by this rewriting.
\end{proposition}
\begin{proof}
Insert \eqref{r73:cov} and \eqref{r74:source-cov}
into \eqref{r73:bank}. The covariance part is
\(\Cov_\nu((2\gamma+3)m_e-\kappa_{J_e},m_f)\),
which becomes the covariance displayed above.
The diagonal terms before the original scalar
contact are
\((2\gamma+3)\E v_e-2\gamma\E s_e\).
The original contact is
\(-2\gamma^2+2\gamma\E\tau_e\); since
\(\E\tau_e=\E s_e\), these two source means cancel.
This proves the formula without changing the
return or \(D\Gamma\) terms.
\end{proof}

The remaining displayed diagonal equals
\(3\gamma-(2\gamma+3)\E w_e\), and may still be
of order \(\gamma\). Its cancellation with returns
or other bank terms has not been proved here.
Nor does the entrywise bound \(350\) bound the
norm of a matrix on a growing number of edges.
The next genuinely spatial obligation is a
summable or signed operator estimate for the
covariance in \eqref{r74:bank}, together with the
retained diagonal, returns, and \(D\Gamma\).

There is a limited source-aware consequence for
the cutoffs already paid in \eqref{r74:cutoff}.
If \(q=\nu(\mathcal T^c)<1\), then any square-integrable
label observable \(X\) satisfies
\[
 |\E_\nu X-\E_\nu[X\mid\mathcal T]|
       \le\sqrt{\Var_\nu X}\sqrt{\frac{q}{1-q}}.
 \tag{V74.31}\label{r74:cutoff-means}
\]
This follows by writing the difference as
\(\Cov(X,\mathbf1_{\mathcal T})/\nu(\mathcal T)\).
For \(X=\delta_{J_e},s_e,\mathcal H_{J_e}\), the
constants \(5,14,70\) apply. These are mean-error
estimates. They do not provide a small truncation
error for the covariance of unbounded products
without further moment or tail control.

\subsection{Verification and the next unpaid step}

The exact checker audits the missing-label term,
posterior stochastic domination, angular moments,
and coefficient form of the differential equation.
It compares shifted-source and diagonal-contact
identities against finite cycle sums, and verifies
the compensated-bank algebra independently.
High-precision Bessel evaluations test the uniform
differential inequalities over a range of labels
and activities; those computations are diagnostics,
not their proofs. No Monte Carlo evidence is used.

V73's mathematical body is retained byte-for-byte
apart from its title, and its regression passes.
Only the active predecessor is replaced in
\path{latex_notes}; the other manuscripts remain
unchanged and all build products stay elsewhere.
The previous goal turn was progress. This one pays
the actual source transfer and global single-edge
angular moments, and isolates an activity-independent
compensated off-diagonal covariance. It does not
pay spatial summability, the remaining return/diagonal
compensation, transport to the physical Wilson
evolution, or the interacting four-dimensional
continuum construction and Yang--Mills mass gap.
% END CONSOLIDATED V74 APPEND

% BEGIN CONSOLIDATED V75 APPEND
\clearpage
\section[V75: Exact return cancellation and the spatial target]{V75: exact Casimir--return cancellation, the covariance target,
and source-product cutoff errors}
\label{r75:main}

\subsection{Context and the scheduled companion checkpoint}

V74 paid the signed deficit-source transfer and uniform single-edge
angular moments, but left the original return covariance in the full
bank. We now compute that return before estimating it. On every finite
native graph its count-conditional mean is determined by the angular
Casimir. A pair deletion then cancels the entire factorial/return
covariance, including its same-edge contact. On the cubic torus the
bank becomes one explicit linear function of the count covariance.
This removes a real source-identification obligation. It also reveals
that the remaining uniform estimate is equivalent to the earlier
first spatial covariance remainder; changing source representations
again, without a spatial estimate, would be circular.

The scheduled V75 checkpoint found no newer versions in the eight
companion families inspected at V70. Fresh hashes are unchanged for
NS V26, SM--Einstein V72, Shell Mixing V16, parallel YM V5,
Source Selection V35, Native Stationarity V38, Exact-Action Bridge V28,
and Perturbative ESM V31. Their terminal scope statements and the ESM
transfer memorandum were reread. The useful transfer remains exact
cancellation with all returns retained, not a transfer of a
perturbative or auxiliary gap into Yang--Mills. The next companion
checkpoint and broad literature sweep are both V80.

\paragraph{Scope and nonduplication.}
The local identities below use a finite loopless native zero-field
\(O(4)\) graph, positive scalar activities, and the same count, wire,
and angular marginals as V62--V74. The matrix estimates use the
inherited homogeneous even cubic torus. V63 already proved the
one-bond specialization \eqref{r63:count-cov}; it is not a new result
here. The new identification retains cycles, arbitrary count marks,
and the original outside-return observable. No assertion is made
for arbitrary transported Wilson exteriors.

A targeted primary-source check of
\href{https://arxiv.org/html/2207.07603v2}{Abdesselam,
\emph{Non-Abelian correlation inequalities and stable determinantal
polynomials}}, Sections 1 and 4, reinforces a relevant boundary:
its free-measure/large-power determinantal results are not an
interacting \(O(4)\) covariance comparison. We do not assume
entrywise Ginibre positivity or domination of the native covariance
by a reference resolvent. The wire representation itself is inherited
from the context of
\href{https://arxiv.org/abs/1807.06564}{Benassi--Ueltschi,
\emph{Loop correlations in random wire models}}.
All new identities used here are proved below.

\subsection{A single-bond twist fixes the Casimir normalization}

Write \(M(K)=M_G(K)\), \(k=K_e\), and \(\kappa_j=j(j+2)\).
Let \(\E_{\rm w}[\cdot\mid K]\) refer to the original labelled
wire pairings, and \(\E[\cdot\mid K]\) to V73's angular lift.
These are two conditional laws with the same count marginal;
no pointwise identification of their random variables is intended.

\begin{theorem}[Conditional current--Casimir identity]
\label{r75:casimir}
For every admissible count vector and every edge,
\[
 3\E_{\rm w}[\Lambda_{ee}\mid K]
   =\E[\kappa_{J_e}\mid K]
   =k(k+2)-k(k-1)\frac{M(K-2e)}{M(K)}.
 \tag{V75.1}\label{r75:casimir-eq}
\]
The product involving \(M(K-2e)\) is defined as zero for \(k<2\).
In particular the ratio on the right is used only when all counts
of \(K-2e\) are nonnegative.
\end{theorem}
\begin{proof}
Identify \(S^3\) with unit quaternions. Let \(I_1,I_2,I_3\)
be the real matrices of left multiplication by the three imaginary
quaternion units. They satisfy \(I_a^2=-I\). For fixed unit
vectors \(u,v\), put \(t=u\cdot v\) and
\(t_a(\theta)=u\cdot e^{\theta I_a}v\). Orthogonality of
\(v,I_1v,I_2v,I_3v\) gives
\[
 t_a''(0)=-t,\qquad
 \sum_{a=1}^3(t_a'(0))^2=1-t^2.
 \tag{V75.2}\label{r75:quaternion}
\]
Thus for every polynomial \(f\),
\[
 \left.\sum_{a=1}^3\partial_\theta^2 f(t_a(\theta))\right|_0
      =(1-t^2)f''(t)-3tf'(t)=-\mathcal C f(t),
 \tag{V75.3}\label{r75:laplace}
\]
where \(\mathcal C\) is V74's positive zonal Casimir.

Twist only the occurrences of edge \(e\) by
\(e^{\theta I_a}\), and keep every other bond matrix equal to
the identity. In the finite sphere-pairing expansion of the
resulting count moment, a closed component \(C\) contracts to
\[
 \Tr e^{\theta\ell_e(C)I_a}=4\cos(\theta\ell_e(C)).
\]
A reverse traversal uses the transpose, which here is the inverse.
All occurrences of the single nontrivial matrix commute, so the
exponent is the \emph{net} oriented traversal, not the visit count.
Dividing by the untwisted loop factor four, the whole diagram is
multiplied by \(\prod_C\cos(\theta\ell_e(C))\).
Its second derivative at zero is therefore
\(-\Lambda_{ee}\) times its untwisted weight. This holds for
each of the three \(I_a\). Summing the three derivatives of
the integrated moment gives
\(-3M(K)\E_{\rm w}[\Lambda_{ee}\mid K]\).

Alternatively, expand \(t^k=\sum_j\lambda_j(k)P_j(t)\)
on this edge, using the same expansion on the others.
Equations \eqref{r75:laplace} and
\(\mathcal CP_j=\kappa_jP_j\) show that the summed derivative
is \(-M(K)\E[\kappa_{J_e}\mid K]\).
Finally, applying \(\mathcal C\) directly to \(t^k\) yields
\(k(k+2)t^k-k(k-1)t^{k-2}\). Integration proves the last equality.
All expansions are finite at fixed \(K\), so no interchange of
infinite sums or limiting twists is needed.
\end{proof}

The isoclinic quaternion generator rotates two coordinate planes.
It is not V62's one-plane generator: its normalized loop factor
is \(\cos(\theta\ell)\), whereas a one-plane twist has factor
\((1+\cos(\theta\ell))/2\). The three-generator sum above fixes
the factor three without substituting one convention for the other.

\subsection{The original outside return is an exact pair-deletion source}

Keep the endpoint choice and the actual outside-path count \(H_e\)
from V65. Combining \eqref{r65:return-cov} at the conditional-mean
level with Theorem~\ref{r75:casimir} gives
\[
 6\E_{\rm w}[H_e\mid K]
    =3k-\E[\kappa_{J_e}\mid K]
    =k(k-1)\left\{\frac{M(K-2e)}{M(K)}-1\right\}.
 \tag{V75.4}\label{r75:conditional-return}
\]
For \(k<2\), both the return and the right side are zero.
This formula concerns the count-conditional mean; it does not set
the integer-valued return equal to a deterministic count function
pointwise.

\begin{theorem}[Marked return identity and its complete contact]
\label{r75:return-source}
For a bounded count function \(G\), and more generally whenever
the displayed absolute moments are finite,
\[
 6\E_{\rm w}[H_eG(K)]
   =\gamma_e^2\E_\mu G(K+2e)
       -\E_\mu[K_e(K_e-1)G(K)].
 \tag{V75.5}\label{r75:marked-return}
\]
Consequently,
\[
 \begin{aligned}
 6\E H_e&=\gamma_e^2-\E[K_e(K_e-1)],\\
 6\Cov_{\rm w}(H_e,G(K))
   &=\gamma_e^2\E_\mu[G(K+2e)-G(K)]
       -\Cov_\mu(K_e(K_e-1),G(K)).
 \end{aligned}
 \tag{V75.6}\label{r75:return-centered}
\]
In the notation of V65, on any such finite graph,
\[
 \boxed{\ F_{ef}+6C^{\rm ret}_{ef}=2\gamma_e^2\delta_{ef}.\ }
 \tag{V75.7}\label{r75:contact-cancel}
\]
The diagonal on the right is essential.
\end{theorem}
\begin{proof}
Multiply \eqref{r75:conditional-return} by the normalized count
weight \(Z^{-1}M(K)\prod_f\gamma_f^{K_f}/K_f!\) and by \(G(K)\).
In the first term, \(K_e(K_e-1)/K_e!=1/(K_e-2)!\), so shifting
\(K\) to \(K-2e\) gives precisely the first term of
\eqref{r75:marked-return}. There is no lost support contribution:
the positive sphere-pairing formula makes \(M(K)>0\) exactly
when all vertex degrees are even, and adding two occurrences
preserves admissibility. The \(k=0,1\) terms vanish before shifting.
The finite graph count law has all exponential moments, since
its total-count moment generating function is a ratio of finite
compact partition functions. This justifies the shift for bounded
and polynomial marks; the stated absolute-integrability condition
justifies it in the more general case.

Taking \(G=1\) gives the mean. Subtracting that mean times
\(\E G\) proves the covariance identity. For \(G(K)=K_f\),
the increment is \(2\delta_{ef}\), which proves
\eqref{r75:contact-cancel} with its sign and factor two.
\end{proof}

There is an independent angular expression for the same covariance:
\[
 6C^{\rm ret}_{ef}
      =3\Gamma_{ef}-\Cov_\nu(\kappa_{J_e},m_{J_f}).
 \tag{V75.8}\label{r75:angular-return}
\]
Indeed multiply the first equality of
\eqref{r75:conditional-return} by a count mark, then use the
conditional mean \(\E[K_f\mid J]=m_{J_f}\).
No same-edge differentiation of \(\kappa_j\) is present.
The count ODE also gives
\[
 F_{ef}=\Cov_\nu(\kappa_{J_e},m_{J_f})-3\Gamma_{ef}
                                    +2\gamma_e^2\delta_{ef},
 \tag{V75.9}\label{r75:angular-factorial}
\]
in agreement with \eqref{r75:contact-cancel}. To check the contact
directly, \(\E[K_e(K_e-1)\mid J]
=\gamma_e^2+\kappa_{J_e}-3m_{J_e}\), and its conditional covariance
with \(K_e\) is \(2\gamma_e^2-3v_{J_e}\); add the covariance
of conditional means and use V73's total covariance formula.

As a finite cyclic normalization check, for a four-cycle with every
count equal to two,
\[
 M(K)=\frac5{1152},\qquad
 \E[\kappa_{J_e}\mid K]=\frac45,\qquad
 \E_{\rm w}[\Lambda_{ee}\mid K]=\frac4{15},\qquad
 \E_{\rm w}[H_e\mid K]=\frac{13}{15}.
 \tag{V75.10}\label{r75:cycle-check}
\]
The nonzero return has not been dropped or treated as a rare event.

\subsection{The full bank is a single spatial covariance remainder}

Return to the homogeneous cubic torus, and set
\[
 S=R\circ R,\qquad D=rI-S,\qquad
 A_\gamma=2\gamma I+3D,
 \qquad T_\gamma=2\gamma^2A_\gamma^{-1}.
 \tag{V75.11}\label{r75:resolvent}
\]
Here \(S\) is a fixed geometric matrix, not the angular network
weight \(S_G\). Since \(S\succeq0\) and has row sum \(r\),
\(0\preceq D\preceq rI\), and \(A_\gamma\) is invertible.

\begin{theorem}[Exact full-bank reduction]
\label{r75:bank}
The original bank of V62/V65 satisfies
\[
 \boxed{\quad
 \mathcal B_{\rm w}=A_\gamma\Gamma-2\gamma^2I
                   =A_\gamma(\Gamma-T_\gamma).
 \quad}
 \tag{V75.12}\label{r75:bank-eq}
\]
For the ordinary Euclidean operator norm,
\[
 2\gamma\|\Gamma-T_\gamma\|
   \le\|\mathcal B_{\rm w}\|
   \le(2\gamma+3r)\|\Gamma-T_\gamma\|.
 \tag{V75.13}\label{r75:equivalence}
\]
No commutation of \(D\) and \(\Gamma\) is assumed.
\end{theorem}
\begin{proof}
Insert \eqref{r75:contact-cancel} into the inherited full identity
\eqref{r65:full-bank}. The factorial term and six times the
original return covariance sum to \(2\gamma^2I\), giving
\eqref{r75:bank-eq}. The eigenvalues of the symmetric positive
matrix \(A_\gamma\) lie in \([2\gamma,2\gamma+3r]\).
Applying its maximal and minimal singular-value bounds to
\(\Gamma-T_\gamma\) proves \eqref{r75:equivalence}.
\end{proof}

For clarity, the reference itself is positive entrywise and has an
exact row sum:
\[
 T_\gamma=\frac{2\gamma^2}{2\gamma+3r}
    \sum_{n\ge0}\left(\frac{3S}{2\gamma+3r}\right)^n,
 \qquad T_\gamma\mathbf1=\gamma\mathbf1.
 \tag{V75.14}\label{r75:reference-series}
\]
This is the norm-convergent geometric series, since
\(3r/(2\gamma+3r)<1\). It proves properties of the reference,
not an ordering between \(T_\gamma\) and \(\Gamma\).

Define the first count-covariance remainder by
\[
 E_\gamma=\Gamma-\gamma I+\tfrac32D.
\]
An exact multiplication, with the order retained, gives
\[
 \mathcal B_{\rm w}=A_\gamma E_\gamma-\tfrac92D^2,
 \qquad
 T_\gamma=\gamma I-\tfrac32D
                      +\tfrac92D^2A_\gamma^{-1}.
 \tag{V75.15}\label{r75:first-remainder}
\]
Thus \(\|T_\gamma-\gamma I+\tfrac32D\|
\le9r^2/(4\gamma)\). Uniform boundedness of the full bank,
for \(\gamma\ge\gamma_0>0\) and growing even tori, is equivalent
to \(\|E_\gamma\|=O(\gamma^{-1})\) uniformly in volume.
In the old spin notation, \eqref{r62:count-cov} implies exactly
\[
 E_\gamma=\mathcal C_\gamma-K^0
        -(\mu_{1,\gamma}-3r/2)I,
 \qquad K^0=\tfrac32S.
 \tag{V75.16}\label{r75:old-target}
\]
The inherited mean bound
\(\lvert\mu_{1,\gamma}-3r/2\rvert\le52/(3\gamma)\)
therefore makes this the earlier spatial covariance target,
up to a paid diagonal error. This equivalence is \emph{not}
an estimate of its unresolved remainder.

This result changes the next attack: there is no longer a separate
unknown factorial/return cancellation to assume in this count bank.
The genuine next task is a volume-uniform bound on
\(\gamma(\Gamma-T_\gamma)\), with the graph geometry retained.
The constant source is only one vector; the exact row sum of the
reference does not establish that of the native covariance.

\subsection{A quantitative cutoff now controls covariance products}

V74's single-source cutoff error did not control unbounded products.
The full-bank reduction permits such a control using the moments
already proved there. All activities in this subsection are at
least one. Under the original label law, keep
\(\delta_e=m_{J_e}-\gamma_e\) and \(w_e=\gamma_e-v_{J_e}\).
The inherited bounds are
\(\|\delta_e\|_2<5\), \(\|\delta_e\|_4<7\), and
\(\|w_e\|_2<15\).

Let \(T\) be any label event with \(p=\nu(T)=1-q>0\).
Condition the V73 joint law on \(T\). Given \(J\), its count
law is still the original product of \(Q_{J_e,\gamma_e}\)'s.
Write \(\Gamma^T\) for this conditioned law's count covariance.

\begin{proposition}[Covariance-product error under label conditioning]
\label{r75:cutoff}
If \(q\le1/2\), then for every pair of edges
\[
 \begin{aligned}
 \left|\Cov_\nu(\delta_e,\delta_f)
             -\Cov_{\nu(\cdot\mid T)}(\delta_e,\delta_f)\right|
       &\le160\sqrt q,\\
 |\Gamma_{ef}-\Gamma^T_{ef}|&\le182\sqrt q.
 \end{aligned}
 \tag{V75.17}\label{r75:product-error}
\]
On a graph with \(m\) edges this implies
\(\|\Gamma-\Gamma^T\|\le185m\sqrt q\).
For the homogeneous cubic bank, define the \emph{approximant}
\(\mathcal B^T=A_\gamma\Gamma^T-2\gamma^2I\). Then
\[
 \|\mathcal B_{\rm w}-\mathcal B^T\|
                 \le185m(2\gamma+1)\sqrt q.
 \tag{V75.18}\label{r75:bank-cutoff}
\]
\end{proposition}
\begin{proof}
For square-integrable \(X\), the elementary identity
\(\E[X\mid T]-\E X=\Cov(X,\mathbf1_T)/p\) gives
\[
 |\E X-\E[X\mid T]|\le\sqrt{\Var X}\sqrt{q/p}.
\]
Apply this to \(X=\delta_e\delta_f\), whose \(L^2\) norm
is at most \(\|\delta_e\|_4\|\delta_f\|_4<49\).
For \(a_e=\E\delta_e\) and \(b_e=\E[\delta_e\mid T]\),
\[
 |a_e|<5,\quad |b_e|<5p^{-1/2},\quad
 |a_e-b_e|\le5\sqrt{q/p}.
\]
Subtracting the products of means therefore bounds the covariance
difference by
\[
 \{49+25(1+p^{-1/2})\}\sqrt{q/p}<160\sqrt q
 \quad (q\le1/2).
\]
The case \(q=0\) is immediate. The total covariance formula
\(\Gamma=\operatorname{diag}(\gamma_e-\E w_e)
+[\Cov(\delta_e,\delta_f)]\) holds for both laws.
The additional diagonal error is at most
\(15\sqrt{q/p}<22\sqrt q\), proving the entry bound.
The Frobenius norm of an \(m\)-by-\(m\) matrix is at most
\(m\) times its maximal absolute entry. This proves the displayed
operator bound with the harmless reserve from 182 to 185.
Finally \(\|A_\gamma\|\le2\gamma+3r\le2\gamma+1\), so
\eqref{r75:bank-cutoff} follows from the exact unconditioned
bank identity and the definition of \(\mathcal B^T\).
\end{proof}

Here \(\mathcal B^T\) is deliberately defined after the cancellation.
We have \emph{not} asserted that the original wire return identity
\eqref{r75:contact-cancel} survives an arbitrary angular cutoff.
Such a claim would require its own boundary terms and a specified
coupling of labels to wires.

For an explicit finite cutoff, take a set of edges \(A\), integers
\(R_e\ge0\), and \(T=\{J_e\le R_e\text{ for }e\in A\}\).
V74's fourth Casimir moment and Markov's inequality give
\[
 q\le\sum_{e\in A}
 \frac{945(\gamma_e^4+6\gamma_e^3+7\gamma_e^2+\gamma_e)}
      {\{(R_e+1)(R_e+3)\}^4}.
 \tag{V75.19}\label{r75:fourth-tail}
\]
For a common cutoff \(R\) on all \(m\) edges at homogeneous
\(\gamma\ge1\), put \(L_R=(R+1)(R+3)\). Then
\(q\le14175m\gamma^4/L_R^4\). In particular, for any
\(\epsilon>0\), the explicit sufficient condition
\[
 \frac{14175m\gamma^4}{L_R^4}
   \le\min\left\{\frac12,
       \frac{\epsilon^2}{185^2m^2(2\gamma+1)^2}\right\}
 \tag{V75.20}\label{r75:sufficient-cutoff}
\]
ensures \(\|\mathcal B_{\rm w}-\mathcal B^T\|\le\epsilon\).
The fourth moment was already paid; the new conclusion controls
the entire compensated matrix, not only a source mean.
For fixed tolerance this sufficient cutoff has order
\(\gamma^{3/4}m^{3/8}\), with the explicit tolerance cost in
\eqref{r75:sufficient-cutoff}; it is not asserted necessary or optimal.

When all edges are cut off, \(\Gamma^T\) is a finite label sum:
the infinite conditional counts are integrated exactly through
\(a_j,m_j,v_j\). This is an approximation certificate, not an
efficient evaluation algorithm or a uniform bound on
\(\mathcal B^T\). Its edge-count dependence remains visible.
A finite list of evaluated graphs would not prove the required
growing-volume estimate.

\subsection{Verification and the next noncircular question}

The exact checker compares 157,277 labelled pairing diagrams on
five graphs, 283 admissible count profiles, independent angular
expansions, and the original outside-path returns. It checks
1,308 current/Casimir/return entries and 5,232 marked coefficient
identities. A separate normalized triangle power-series calculation
checks the contact and full bank through degree eight. These are
coefficients of the full law, not an erroneous infinite-series
identity applied to a truncated count probability. Rational
conditioning tests verify the product-error decomposition;
noncommuting matrix fixtures check multiplication order but are
not claimed to be native Gibbs covariances.

V74's regression passes, and its mathematical body is retained
byte-for-byte except the version title. The other 46 manuscripts
remain unchanged. Only the active predecessor is replaced, and
all compilation files remain outside \texttt{latex\_notes}.
The previous goal turn was progress. This iteration pays the
native return identification, its exact covariance cancellation,
and a source-product cutoff error. It does not pay the remaining
spatial norm.

The next question is now unambiguous: can the actual interacting
count law, or equivalently the native energy covariance, be
compared to \(T_\gamma\) with error \(O(\gamma^{-1})\) uniformly
on growing cubic tori? Positivity of individual weights, a
one-edge moment bound, a reference-kernel estimate, or another
rewrite of \eqref{r75:old-target} is not such a comparison.
Any useful next argument must supply spatial control or expose
a concrete obstruction to this specified target. General Wilson
exteriors, physical-time contraction, the interacting
four-dimensional continuum construction, and the physical
Yang--Mills mass gap remain unproved.
% END CONSOLIDATED V75 APPEND

% BEGIN CONSOLIDATED V76 APPEND
\clearpage
\section{V76: angular-source control and the full Gaussian cycle-current limit}
\label{r76:main}

\subsection{Context, the nonduplication audit, and the new target}

V75 identified the original return source and reduced the complete bank
to an explicitly stated spatial covariance problem. This append goes
upstream of that problem: what is the actual joint high-activity law
of the angular Casimirs that carry its remaining covariance? We prove
that on each fixed finite connected graph it converges to the squared
norms of three independent Gaussian currents in the \emph{full cycle
space}. We also control the replacement of conditional count sources
by these Casimirs, with graph-independent entrywise error constants.
The convergence theorem itself is not uniform in graph size.

The initial rare-edge lead required an important nonduplication
correction: V68 already proves Bernoulli domination of low-count
\emph{edges}, not just the low-degree vertex-cluster bounds of V64.
Here the new statement is domination conditional on \emph{every angular
label}, with a dominating field independent of all those labels.
It consequently retains the full rare-event power for unbounded
label-source weights. We do not repackage marginal domination as a
new result. Nor do we repackage the already known leading fixed-graph
energy covariance as a new coefficient: the new identification is
of the entire joint scaled angular-label law.

The ESM transfer remains proof architecture, not a transferred YM
estimate: original-law identities, paid source moments, and a named
unpaid uniform limit remain separate. The V75 companion checkpoint
stands; the next companion review and regular broad literature sweep
are both due at V80. A targeted heat-kernel check is recorded below.
The interacting four-dimensional continuum construction and physical
Yang--Mills mass gap remain unproved.

\subsection{Native joint law and notation}

Let \(G=(V,E)\) be a finite connected simple graph, with
\(n=|V|\ge2\), \(m=|E|\), and common activity \(\gamma>0\).
Orient its edges arbitrarily. The spin space is the unit
three-sphere, identified with unit quaternions \(\SU(2)\), with
probability Haar measure \(dH\). The native zero-field weight is
\(\exp(\gamma\sum_{xy\in E}u_x\cdot u_y)\).
Use the V73 positive joint count/label lift, whose label marginal is
\[
 \nu_\gamma(J)=\frac{S_G(J)\prod_e a_{J_e}(\gamma)}{Z_G(\gamma)},
 \qquad a_j(\gamma)=\frac{2I_{j+1}(\gamma)}{\gamma},\quad
 \kappa_j=j(j+2).
 \tag{V76.1}\label{r76:law}
\]
Here \(P_j(g)=(j+1)\chi_j(g)\) and
\(S_G(J)=\int\prod_{e=xy}P_{J_e}(g_x^{-1}g_y)\prod_xdH(g_x)\).
V73 proves \(0\le S_G(J)\le\prod_e(j_e+1)^2\).
Conditional on all labels, counts are independent with
\[
 Q_{j,\gamma}(K=j+2q)=
 \frac{(\gamma/2)^{j+2q}}{q!(q+j+1)!\,a_j(\gamma)},
 \qquad q=0,1,\ldots.
 \tag{V76.2}\label{r76:Q}
\]
Write \(m_j=\E_QK=\gamma a_j'/a_j\),
\(v_j=\Var_QK=\gamma m_j'\),
\(\delta_j=m_j-\gamma\), \(w_j=\gamma-v_j\), and
\(X_e=\kappa_{J_e}/(2\gamma)\). All unqualified expectations and
norms in this append use the actual native joint law or its label
marginal, not an independently chosen reference field. In particular
\[
 \Gamma_{ef}=\Cov(K_e,K_f)
   =\gamma\delta_{ef}-\delta_{ef}\E w_{J_e}
                         +\Cov(\delta_{J_e},\delta_{J_f}).
 \tag{V76.3}\label{r76:Gamma}
\]
The Kronecker symbol \(\delta_{ef}\) is distinct from the function
\(\delta_j\).

\subsection{Rare count defects with arbitrary angular-source weights}

\begin{proposition}[Domination independent of the complete label field]
\label{r76:rare}
For \(0<s<1\), every \(j\ge0\), and every \(\gamma>0\),
\[
 \E_{Q_{j,\gamma}}s^K
 \le s^{-3/2}e^{-(1-s)\gamma}.
 \tag{V76.4}\label{r76:pgf}
\]
For \(0<\alpha<1\) set
\(p_{\gamma,\alpha,s}=
 \min\{1,s^{-3/2}e^{-\gamma(1-s+\alpha\log s)}\}\).
There is a coupling of the native law with iid Bernoulli variables
\(B_e\) of parameter \(p_{\gamma,\alpha,s}\), independent of \(J\),
such that \(A_e=\mathbf1_{\{K_e\le\alpha\gamma\}}\le B_e\)
for every edge. Therefore, for any nonnegative measurable label
weight \(W(J)\) and nonnegative increasing function \(F\) on
\(\{0,1\}^E\),
\[
 \E[W(J)F(A)]\le\E_\nu W\;\E F(B).
 \tag{V76.5}\label{r76:weighted}
\]
Extended nonnegative expectations are allowed.
\end{proposition}
\begin{proof}
The probability generating function is \(a_j(s\gamma)/a_j(\gamma)\).
Integrating V74's all-label bound \(m_j(t)>t-3/2\) gives
\[
 \log\frac{a_j(s\gamma)}{a_j(\gamma)}
 =-\int_{s\gamma}^{\gamma}\frac{m_j(t)}t\,dt
 \le-(1-s)\gamma-\tfrac32\log s.
\]
Markov's inequality applied to \(s^K\) gives
\(Q_{j,\gamma}(K\le\alpha\gamma)
 \le s^{-\alpha\gamma}\E_Qs^K\).
Take independent uniform variables \(U_e\), independent of \(J\),
and set \(A_e=\mathbf1_{\{U_e\le q_{J_e}\}}\),
\(B_e=\mathbf1_{\{U_e\le p_{\gamma,\alpha,s}\}}\), where
\(q_j=Q_{j,\gamma}(K\le\alpha\gamma)\).
Using further independent uniforms, sample each count from its
conditional bad or good part. Null parts can be assigned arbitrarily.
This reconstructs the independent conditional count laws in
\eqref{r76:Q} and proves the coupling. Monotonicity and independence
give \eqref{r76:weighted}; truncation and monotone convergence remove
boundedness restrictions on \(W\).
\end{proof}

In particular, choose \(s=\alpha=1/2\) and write
\[
 p_\gamma=\min\{1,2^{3/2}e^{-(1-\log2)\gamma/2}\}.
 \tag{V76.6}\label{r76:p}
\]
For a specified edge set \(F_0\),
\(\E[W\prod_{e\in F_0}A_e]\le p_\gamma^{|F_0|}\E W\).
On a simple cubic torus the line graph has maximum degree ten.
There are at most \(100^{k-1}\) connected edge sets of size \(k\)
containing a fixed edge: choose a deterministic spanning tree of
each such set and its depth-first walk of length \(2k-2\); distinct
sets have distinct visited sets and hence distinct walks. If the
bad cluster of \(e\) has size at least \(k\), it contains such a set.
The union bound and \eqref{r76:weighted} thus give
\[
 \E[W\mathbf1_{\{|\mathcal C_A(e)|\ge k\}}]
 \le \E W\,p_\gamma(100p_\gamma)^{k-1}.
 \tag{V76.7}\label{r76:cluster}
\]
For \(\gamma\ge64\), \(p_\gamma<3/8000<1/200\), so the last
factor is at most \(p_\gamma2^{-(k-1)}\).
Indeed \(\log2<0.7\), \(2^{3/2}<3\), and \(e^3>20\) suffice.
For example, \(\E|\delta_{J_e}\delta_{J_f}|<25\) by V74,
with no loss of half the rare-event exponent. This statement does
not make the labels independent, and
\(\Pr(\exists e:A_e=1)\le mp_\gamma\) still contains the volume.

\subsection{Uniform source replacement, with the matrix-size loss exposed}

\begin{proposition}[First Casimir approximations for the two count sources]
\label{r76:sources}
For \(\gamma\ge1\) define
\(\epsilon_j=\delta_j+3/2-\kappa_j/(2\gamma)\) and
\(\eta_j=w_j-\kappa_j/(2\gamma)\). Then
\[
 \epsilon_j=\frac{w_j-2\delta_j-\delta_j^2}{2\gamma},
 \qquad
 |\eta_j|\le\frac{16+17\kappa_j/\gamma
                           +6(\kappa_j/\gamma)^2}{\gamma}.
 \tag{V76.8}\label{r76:source-point}
\]
Under the native label law, for every edge,
\[
 \|\epsilon_{J_e}\|_2<37/\gamma,
 \qquad \|\eta_{J_e}\|_2<825/\gamma.
 \tag{V76.9}\label{r76:source-L2}
\]
Consequently there is a symmetric remainder matrix satisfying
\[
 \Gamma=\gamma I-\operatorname{diag}(\E X_e)
                    +\Cov_\nu(X)+\mathcal R_\gamma,
 \quad |(\mathcal R_\gamma)_{ef}|\le1125/\gamma,
 \quad \|\mathcal R_\gamma\|_{2\to2}\le1125m/\gamma.
 \tag{V76.10}\label{r76:source-matrix}
\]
All constants in the entry bounds are independent of the graph;
the last norm estimate is explicitly \emph{not} volume-uniform.
\end{proposition}
\begin{proof}
The Riccati identity \eqref{r74:riccati} gives the exact formula
for \(\epsilon_j\). V74's norms yield
\(2\gamma\|\epsilon\|_2\le15+2\cdot5+7^2=74\),
with strict inequality. To estimate \(\eta\), first use
\eqref{r74:w-ode} and the same Riccati identity to obtain
\[
 \eta_j'+\frac{2m_j+1}{\gamma}\eta_j
 =-\frac{2\delta_j+\delta_j^2}{\gamma}
                      -\frac{\kappa_j\delta_j}{\gamma^2}.
 \tag{V76.11}\label{r76:eta-ode}
\]
Fix \(\gamma\ge4\), put \(x=\kappa_j/\gamma\), and integrate
this linear equation on \([\gamma/2,\gamma]\). Its coefficient
is at least one there, since \(m_j(t)>t-3/2\).
Also \(|\delta_j(t)|\le3/2+x\), and hence the absolute value of
the right-hand side is at most
\((10.5+16x+6x^2)/\gamma\).
At the initial endpoint, \(|\eta_j(\gamma/2)|\le\gamma/2+x\)
because \(0<w_j(t)<t\). Its damped contribution is at most
\((2+x)/\gamma\): use
\((\gamma^2/2)e^{-\gamma/2}<2\) and
\(\gamma e^{-\gamma/2}<1\) for \(\gamma\ge4\).
Variation of constants gives
\(|\eta_j|\le(12.5+17x+6x^2)/\gamma\).
For \(1\le\gamma<4\), the direct estimate
\(|\eta_j|\le\gamma+x/2\le(16+2x)/\gamma\)
completes \eqref{r76:source-point}.

The inherited full-law moments give
\(\|\kappa_{J_e}/\gamma\|_2\le\sqrt{30}\) and
\(\|(\kappa_{J_e}/\gamma)^2\|_2\le\sqrt{14175}\).
Thus \(\gamma\|\eta\|_2\le16+17\sqrt{30}
+6\sqrt{14175}<825\).
Since \(\delta+3/2=X+\epsilon\),
\[
 \Cov(\delta_e,\delta_f)-\Cov(X_e,X_f)
   =\Cov(\epsilon_e,\delta_f)+\Cov(X_e,\epsilon_f).
\]
Here subscripts on functions abbreviate evaluation at \(J_e\).
Cauchy--Schwarz, \(\|\delta_e\|_2<5\), and
\(\|X_e\|_2\le\sqrt{30}/2<3\) bound the difference by
\(296/\gamma\). The diagonal mean correction costs at most
\(825/\gamma\). Equation \eqref{r76:Gamma} proves the entry
bound with the slightly enlarged constant 1125; the maximum
absolute row sum proves the asserted matrix bound.
\end{proof}

\subsection{An exact heat-source generating function}

Let \(C\) be the positive Laplace--Beltrami operator on the unit
three-sphere, so \(CP_j=\kappa_jP_j\). Denote its probability-Haar
heat kernel by \(h_t=\sum_{j\ge0}e^{-t\kappa_j}P_j\).
Set
\[
 b_\gamma(g)=\frac{e^{\gamma\operatorname{Re}g}}{a_0(\gamma)},
 \qquad b_{\gamma,s}=h_{s/(2\gamma)}*b_\gamma\quad(s>0),
 \qquad b_{\gamma,0}=b_\gamma.
 \tag{V76.12}\label{r76:b}
\]
Convolution uses probability Haar measure. Each is a positive
probability density. The time \(s/(2\gamma)\) is an auxiliary
Casimir-source time, not physical Euclidean time.

\begin{lemma}[Exact native Laplace transform]
\label{r76:heat-source}
For every \(\boldsymbol s\in[0,\infty)^E\), define
\[
 Z_G^\sharp(\gamma,\boldsymbol s)
  =\int\prod_{e=xy}b_{\gamma,s_e}(g_x^{-1}g_y)\prod_xdH(g_x).
\]
Then
\[
 \E_{\nu_\gamma}\exp\!\left(-\sum_es_eX_e\right)
 =\frac{Z_G^\sharp(\gamma,\boldsymbol s)}
        {Z_G^\sharp(\gamma,\boldsymbol0)}.
 \tag{V76.13}\label{r76:Laplace-exact}
\]
\end{lemma}
\begin{proof}
Convolution by \(h_t\) multiplies the \(P_j\) coefficient by
\(e^{-t\kappa_j}\). Thus the coefficient of \(P_j\) in
\(b_{\gamma,s}\) is
\((a_j/a_0)e^{-s\kappa_j/(2\gamma)}\).
The expansion and all integrations converge absolutely:
\(|P_j|\le(j+1)^2\) and
\(\sum_j(j+1)^2a_j(\gamma)=e^\gamma\).
Expanding the finite product, integrating, and cancelling the
common factor \(a_0^{-m}\) proves the formula. No label cutoff
and no replacement of the network weight have been made.
\end{proof}

\subsection{Kernel estimates with the Haar normalization proved}

\begin{lemma}[Small-scale convolution kernel]
\label{r76:local-kernel}
For each fixed \(s\ge0\) there are finite positive constants
\(C_s,c_s\) such that for \(\gamma\ge1\)
\[
 b_{\gamma,s}(g)\le C_s\gamma^{3/2}
                   e^{-c_s\gamma d(1,g)^2}.
 \tag{V76.14}\label{r76:kernel-bound}
\]
Uniformly for \(x\) in compact subsets of \(\R^3\),
\[
 \gamma^{-3/2}b_{\gamma,s}(\exp(x/\sqrt\gamma))
 \longrightarrow
 \frac{2\pi^2}{[2\pi(1+s)]^{3/2}}
              e^{-|x|^2/[2(1+s)]}.
 \tag{V76.15}\label{r76:kernel-limit}
\]
Here \(\exp x=\cos|x|+\widehat x\sin|x|\) in quaternion notation.
\end{lemma}
\begin{proof}
The probability Haar density in this chart, away from the
measure-zero cut locus, is
\((2\pi^2)^{-1}(\sin|x|/|x|)^2d^3x\), \(|x|<\pi\).
The inequalities \(1-\cos\theta\ge2\theta^2/\pi^2\) for
\(0\le\theta\le\pi\), and a lower integral bound on
\(\theta\le\gamma^{-1/2}\), show
\(a_0(\gamma)\asymp e^\gamma\gamma^{-3/2}\) for \(\gamma\ge1\).
Rescaling the same integral and using dominated convergence gives
\[
 a_0(\gamma)\sim e^\gamma\gamma^{-3/2}\sqrt{2/\pi}.
\]
These statements prove both conclusions for \(s=0\).

For completeness the heat-kernel estimate used for \(s>0\) can be
obtained directly, without importing a formal parametrix. For
\(0<\theta<\pi\), the spectral expansion and Poisson summation give
\[
 h_t(\cos\theta)=
 \frac{e^t\sqrt\pi}{4t^{3/2}\sin\theta}
 \sum_{\ell\in\mathbb Z}(\theta+2\pi\ell)
                e^{-(\theta+2\pi\ell)^2/(4t)}.
 \tag{V76.16}\label{r76:heat-images}
\]
Indeed, put \(q=j+1\) in
\(e^t\sum_{q\ge1}q e^{-tq^2}\sin(q\theta)/\sin\theta\)
and differentiate the Poisson formula for
\(\sum_{q\in\mathbb Z}e^{-tq^2}e^{iq\theta}\).
All series and differentiated series converge absolutely for
\(t>0\); the endpoint values follow by continuity. The heat
maximum principle gives positivity, and the constant eigenmode
gives integral one.

Here are details ensuring a global bound, including the apparent
denominator singularities. Write \(f_t(z)=z e^{-z^2/(4t)}\).
For \(\theta\le\pi/2\), keep \(f_t(\theta)\) separately and
pair the remaining images as
\(f_t(2\pi q+\theta)-f_t(2\pi q-\theta)\), \(q\ge1\).
The mean-value theorem removes the factor \(\theta\), while
\(\sin\theta\ge2\theta/\pi\) and
\(|f_t'(z)|\le(1+z^2/(2t))e^{-z^2/(4t)}\).
For \(\theta\ge\pi/2\), put \(\xi=\pi-\theta\) and pair
images around \((2q+1)\pi\); their differences remove \(\xi\)
in exactly the same fashion, and all arguments have absolute
value at least \(\pi/2\). Gaussian sums converge uniformly after
absorbing polynomial factors into a slightly weaker Gaussian.
For any fixed \(t_0>0\) this proves, with constants depending on
\(t_0\),
\[
 0\le h_t(g)\le C t^{-3/2}e^{-d(1,g)^2/(Ct)},\quad 0<t\le t_0.
\]
In \eqref{r76:heat-images}, take \(t=s/(2\gamma)\),
\(\theta=|x|/\sqrt\gamma\). All nonzero images are exponentially
small on compact \(x\)-sets, also at \(x=0\) by the same paired
bounds. The remaining term gives
\[
 \gamma^{-3/2}h_{s/(2\gamma)}(\exp(x/\sqrt\gamma))
 \longrightarrow\frac{2\pi^2}{(2\pi s)^{3/2}}e^{-|x|^2/(2s)}.
\]
For the convolution bound use
\(d(1,h)+d(h,g)\ge d(1,g)\), reserve half of the two Gaussian
exponents for this distance, and integrate the remaining Gaussian
in \(h\). Its integral is \(O_s(\gamma^{-3/2})\), cancelling
one of the two prefactors. This proves \eqref{r76:kernel-bound}.
For the local limit, put \(h=\exp(y/\sqrt\gamma)\) in the
convolution. On compact sets,
\(\sqrt\gamma\log(h^{-1}\exp(x/\sqrt\gamma))\to x-y\).
The preceding Gaussian bounds give an integrable dominating
function for bounded \(x\); the part outside the chart's compact
rescalings has uniformly vanishing Gaussian tails. The convolution
of Euclidean Gaussians with variances \(s\) and one has variance
\(1+s\). The Haar factor \((2\pi^2)^{-1}\) leaves precisely
the constant in \eqref{r76:kernel-limit}.
\end{proof}

The primary literature check was C.~Zhao and J.~S.~Song,
\emph{Exact Heat Kernel on a Hypersphere and Its Applications in
Kernel SVM}, Frontiers in Applied Mathematics and Statistics 4 (2018),
article 1, \href{https://doi.org/10.3389/fams.2018.00001}{doi:10.3389/fams.2018.00001}.
Its convergent hyperspherical spectral expansion provides context;
the \(S^3\) normalization, image formula, and bounds needed here
were derived above. No machine-learning application, parametrix
truncation, or uniform interacting estimate is imported.

\subsection{The joint native Casimir limit is a Gaussian cycle-current law}

Delete one root column from the oriented incidence matrix to obtain
\(B_0\in\R^{m\times(n-1)}\), and define
\[
 R=B_0(B_0^TB_0)^{-1}B_0^T,\qquad P=I-R,\qquad
 \beta=m-n+1.
 \tag{V76.17}\label{r76:projection}
\]
Let \(U\in\R^{m\times\beta}\) have orthonormal columns spanning
\(\ker B_0^T\); thus \(UU^T=P\). These definitions are independent
of the deleted root. On periodic graphs this is the \emph{full}
divergence-free space, including winding cycles, not just the span
of elementary plaquette boundaries. Empty matrices in the tree
case \(\beta=0\) have determinant one.

\begin{theorem}[Fixed-graph joint high-activity limit]
\label{r76:joint-limit}
For each fixed graph \(G\), as \(\gamma\to\infty\),
\[
 (X_e)_{e\in E}\ \Longrightarrow
 (X_e^*)_{e\in E},\qquad
 X_e^*=\tfrac12\sum_{a=1}^3(Y_e^a)^2,
 \quad Y^a=U\xi^a,
 \tag{V76.18}\label{r76:law-limit}
\]
where the \(\xi^a\) are independent standard Gaussian vectors in
\(\R^\beta\). In particular
\[
 \begin{split}
 \lim_{\gamma\to\infty}\E_\nu e^{-\sum_es_eX_e}
   &=\det\bigl(I_\beta+U^T\operatorname{diag}(s_e)U\bigr)^{-3/2},\\
 \lim_{\gamma\to\infty}\E X_e&=\tfrac32P_{ee},\\
 \lim_{\gamma\to\infty}\Cov(X_e,X_f)&=\tfrac32P_{ef}^{\,2}.
 \end{split}
 \tag{V76.19}\label{r76:limit-moments}
\]
The first assertion holds for every fixed nonnegative source vector.
No rate uniform in growing graphs is asserted.
\end{theorem}
\begin{proof}
We prove a partition-function limit with enough domination to avoid
differentiating an unproved asymptotic expansion. Put \(a_e=1+s_e\)
and \(L(\boldsymbol a)=B_0^T\operatorname{diag}(a_e^{-1})B_0\).
Global left invariance lets us fix the root spin at one. Choose a
spanning tree rooted there, oriented outwards, and use its increments
\(h_{xy}=g_x^{-1}g_y\) as independent Haar coordinates. The change
of variables preserves product Haar measure, by successively
integrating leaves. Write each tree increment as
\(h_e=\exp(y_e/\sqrt\gamma)\), \(|y_e|<\pi\sqrt\gamma\).
The rescaled Haar Jacobian is at most
\((2\pi^2)^{-1}\gamma^{-3/2}\).

For fixed tree increments \(y\), each vertex has first-order
coordinate \(x_v\) equal to the signed sum of the increments on
its unique root path; set \(x_o=0\). Multiplication of finitely
many quaternions gives
\(\sqrt\gamma\log(g_x^{-1}g_y)\to x_y-x_x\).
The linear change from the tree increments to the nonroot vertex
coordinates has determinant of absolute value one in each of the
three components. Apply Lemma~\ref{r76:local-kernel} to every edge.
After multiplication by \(\gamma^{-3\beta/2}\), the integrand
including all tree Jacobians is bounded by
\[
 C_{G,\boldsymbol s}\exp\!\left(-c_{G,\boldsymbol s}
                         \sum_{e\text{ in tree}}|y_e|^2\right).
\]
To see this, retain the Gaussian bound on each tree edge and bound
each other edge by \(C_s\gamma^{3/2}\). All powers of \(\gamma\)
cancel. Extend the integrand by zero outside its expanding chart.
Dominated convergence is now justified on \(\R^{3(n-1)}\).
Evaluating the resulting rooted Gaussian integral gives
\[
 \gamma^{-3\beta/2}Z_G^\sharp(\gamma,\boldsymbol s)
 \longrightarrow
 (2\pi^2)^\beta(2\pi)^{-3\beta/2}
 \left[\Bigl(\prod_e a_e\Bigr)\det L(\boldsymbol a)\right]^{-3/2}.
 \tag{V76.20}\label{r76:partition-limit}
\]
The determinant is positive because the graph is connected and
\(a_e>0\). Taking the ratio with \(\boldsymbol s=0\) in
\eqref{r76:Laplace-exact} produces
\[
 \left[\Bigl(\prod_e a_e\Bigr)
        \frac{\det L(\boldsymbol a)}{\det(B_0^TB_0)}\right]^{-3/2}.
 \tag{V76.21}\label{r76:det-limit}
\]
For clarity the cycle-space determinant identity can also be proved
without a combinatorial assertion. Let
\(V_0=B_0(B_0^TB_0)^{-1/2}\), so \((U,V_0)\) is orthogonal,
and put \(A=\operatorname{diag}(a_e)\). The complementary-block
Schur determinant formula gives
\[
 \det A\,\det(V_0^TA^{-1}V_0)=\det(U^TAU).
 \tag{V76.22}\label{r76:Schur}
\]
Indeed the bottom-right block of the inverse of
\((U,V_0)^TA(U,V_0)\) is the inverse of its Schur complement;
multiplying its determinant by that of the whole matrix leaves
the top-left block's determinant. This proves the identity also
when one block is empty. It converts \eqref{r76:det-limit} into
the determinant in \eqref{r76:limit-moments}.
The elementary Gaussian integral for three independent \(\xi^a\)
gives exactly that Laplace transform for \(X^*\).

For rigor in passing from transforms to laws, the inherited bound
\(\E X_e\le3/2\) gives tightness on the fixed finite-dimensional
nonnegative orthant. Every subsequential weak limit has the
displayed Laplace transform since the integrands are bounded
continuous. Laplace transforms determine probability laws there:
under \(z_e=e^{-x_e}\), their integer arguments give all polynomial
moments on \([0,1]^m\), whose polynomials are uniformly dense in
continuous functions. Thus every subsequential limit is the stated
law, proving convergence. Furthermore
\(\E X_e^4\le14175/16\) for \(\gamma\ge1\); this implies uniform
integrability of \(X_e\) and \(X_eX_f\), the latter by
\(\E(X_eX_f)^2\le\sqrt{\E X_e^4\E X_f^4}\).
Truncation therefore proves convergence of the first and second
moments without differentiating \eqref{r76:partition-limit}.
Finally a centered Gaussian vector of covariance \(P\) has
\(\Cov((Y_e)^2,(Y_f)^2)=2P_{ef}^2\); summing the three independent
components with the factor \(1/4\) proves the last two formulas.
\end{proof}

The constants in the tree domination, its dimension, and the finite
path multiplication used above depend on \(G\). This proof cannot
exchange the high-activity and growing-volume limits. In the tree
case its limit is zero; in fact integrating leaves forces every
angular label to be zero already at finite activity.

On a fixed edge-transitive cubic torus let \(r=R_{ee}\) and
\(D=rI-R\circ R\), where \(\circ\) denotes entrywise product.
Equations \eqref{r76:source-matrix} and
\eqref{r76:limit-moments} recover
\[
 \Gamma-\gamma I\longrightarrow
 \tfrac32(P\circ P-\operatorname{diag}P)=-\tfrac32D.
 \tag{V76.23}\label{r76:old-coefficient}
\]
Here \(\operatorname{diag}P\) means the diagonal matrix with
entries \(P_{ee}\). For off-diagonal entries \(P_{ef}=-R_{ef}\);
for diagonal entries \((1-r)^2-(1-r)=r^2-r\), proving the last
equality. This leading coefficient was already known from the
fixed-graph spin-wave calculation. The theorem adds its native
joint Casimir interpretation, not a new volume-uniform coefficient.

\subsection{Dense angular correlations persist after rare defects disappear}

\begin{corollary}[An exact cycle diagnostic]
\label{r76:cycle}
On a cycle of fixed length \(L\ge3\), all edge labels are equal
almost surely. As \(\gamma\to\infty\), their scaled Casimirs
converge jointly to the same variable \(\chi_3^2/(2L)\). For
any two edges,
\[
 \Cov_\nu(\delta_{J_e},\delta_{J_f})\longrightarrow\frac3{2L^2}.
 \tag{V76.24}\label{r76:cycle-cov}
\]
The joint Casimir limit, its first two moments, and this covariance
limit are unchanged if the native joint law is conditioned on
\(K_e>\gamma/2\) for every edge.
\end{corollary}
\begin{proof}
Character orthogonality at each degree-two vertex forces adjacent
labels to agree, hence all labels agree. Equivalently the exact
cycle weights are proportional to \((j+1)^2a_j(\gamma)^L\).
For a consistently oriented cycle \(P=L^{-1}\mathbf1\mathbf1^T\),
so Theorem~\ref{r76:joint-limit} gives the common chi-square limit
and its covariance \(3/(2L^2)\). Proposition~\ref{r76:sources}
transfers this covariance to \(\delta\).
The all-good event has probability at least \(1-Lp_\gamma\),
tending to one. Conditioning changes the joint law in total
variation by at most \(Lp_\gamma\). The uniform fourth-moment
bounds for both \(X_e\) and \(\delta_{J_e}\) imply uniform
integrability of their products; conditioned fourth moments are
bounded by the original ones divided by the all-good probability.
This proves the conditional moment assertions too.
\end{proof}

This diagnostic rules out discarding the dense angular covariance
on the grounds that bad counts are rare. It is an actual native
finite-cycle limit, distinct from V71--V72's worst-exterior parity
examples on high-count paths. More quantitatively on this cycle set
\(D_L=L^{-1}I-L^{-2}\mathbf1\mathbf1^T\) and
\(A_{\gamma,L}=2\gamma I+3D_L\). If one replaces the exact total
covariance by its conditional diagonal part, the omitted matrix is
\(\Var_\nu(m_J)\mathbf1\mathbf1^T\). Since
\(A_{\gamma,L}\mathbf1=2\gamma\mathbf1\),
\[
 \lim_{\gamma\to\infty}\frac1\gamma
 \bigl\|A_{\gamma,L}\Cov_\nu(m)\bigr\|_{2\to2}=\frac3L.
 \tag{V76.25}\label{r76:cycle-obstruction}
\]
It is thus not a bounded error after bank multiplication. This is
an auxiliary cycle obstruction to that approximation, not a
counterexample to the cubic spatial target; no cubic value or
bound for \(r\) has been imported onto a cycle.

\subsection{Verification, paid progress, and the next unpaid estimate}

The accompanying checker separates exact rational finite-graph
identities from numerical diagnostics. It checks the complementary
determinant against grounded Laplacians and spanning-tree
enumeration; tests conditional weighted-event inequalities on
rational latent-label fixtures; and independently evaluates positive
count series to check the source identities and tail bounds.
Spectral/image heat-kernel comparisons and exact cycle label sums
provide numerical normalization and limit diagnostics. Such floating
calculations are not interval certificates and do not replace any
proof above. No Monte Carlo estimate is used as a theorem.

V75's mathematics is preserved byte-for-byte except the version
title. The active predecessor alone is replaced; other manuscripts
are preserved and builds remain outside \texttt{latex\_notes}.
This iteration pays a native joint angular-law identification,
uniform source-replacement entries, and label-weighted rare-event
control. It does not pay a growing-volume spatial norm.

The remaining question is sharpened, not removed. Let
\(M_\gamma=\Cov_\nu(X)-\operatorname{diag}(\E X_e)\) and
\(M_*=(3/2)(P\circ P-\operatorname{diag}P)\).
We know \(M_\gamma\to M_*\) on every fixed graph and
\(\Gamma-\gamma I=M_\gamma+\mathcal R_\gamma\), with
\emph{entrywise} \(O(\gamma^{-1})\) control of \(\mathcal R\).
What V75 needs is a bound on the \emph{combined} remainder
\[
 \sup_{\text{cubic tori }G}\sup_{\gamma\ge\gamma_0}
 \gamma\,\|M_\gamma-M_*+\mathcal R_\gamma\|_{2\to2}<\infty.
 \tag{V76.26}\label{r76:unpaid}
\]
Separate uniform bounds for the two pieces would suffice but have
not been proved; assuming them would be circular. The exact source
formula \eqref{r76:Laplace-exact} suggests estimating \emph{connected}
source derivatives of the logarithmic partition-function ratio
relative to its cycle determinant, with cancellation before spatial
summation. Rare-edge weights may help localize defects, but do not
control the dense good region, as the cycle theorem demonstrates.
An extensive partition-function error or the bound
\(1125m/\gamma\) is not the required operator estimate. General
transported Wilson exteriors, physical-time contraction, continuum
existence, and the physical Yang--Mills mass gap remain open here.
% END CONSOLIDATED V76 APPEND

% BEGIN CONSOLIDATED V77 APPEND
\clearpage
\section{V77: a zero-radius obstruction and an analytic native covariance source}
\label{r77:main}

\subsection{Why the proposed complex-source route must change}

The preceding goal turn made progress: V76 identified the full
fixed-graph angular limit and proved native source estimates.
Its last suggested route was to estimate connected derivatives of
the heat-source partition function. Before applying complex-analytic
estimates, however, its domain must be established. This audit finds
a substantive obstruction: on every nonbridge edge, the actual
finite-activity Casimir has \emph{no} positive exponential moment.
The heat-source transform is smooth from its permitted side but
has Taylor radius zero at the source origin. Cauchy estimates across
that origin would be invalid, even at fixed volume.

We go upstream rather than assume an analytic continuation. A
magnetic-weight representation of V74's reference law gives a
graph-uniform tail for the actual label marginal. It pays uniform
exponential moments of the exact conditional count mean and variance
sources. A source built from those two functions, with its quadratic
diagonal term retained, is analytic and has precisely the required
covariance Hessian. This repairs the analytic tool; it does not yet
prove the required growing-volume estimate.

The character-tail technique is not new to the larger programme:
archive \(A\), f15 V91 already uses continued characters for
\emph{electric output labels of a transfer operator}. Those labels
and that spectral-window law are not the angular marginal
\(\nu_\gamma\) of V73--V76. The proof below derives the needed bound
under \(\nu_\gamma\) itself, with no transfer of a spectral-window
assumption. V68's rare-edge domination and V76's Gaussian limit
are inherited, not reproved as new results.

Throughout use V76's finite connected simple graph, positive
homogeneous activity, and positive native joint law. Set
\[
 \kappa_j=j(j+2),\quad X_e=\frac{\kappa_{J_e}}{2\gamma},\quad
 \delta_j=m_j-\gamma,\quad w_j=\gamma-v_j,
 \quad d_\gamma(x)=\sqrt{\gamma^2+x^2}-\gamma\quad(x\ge0).
 \tag{V77.1}\label{r77:notation}
\]
The scalar function \(d_\gamma\) is not the matrix \(D\) of V75.
No auxiliary source parameter below is physical Euclidean time.

\subsection{A native label tail from magnetic weights and count factorial moments}

\begin{theorem}[Uniform radial-label tail under the original marginal]
\label{r77:tail}
For every edge, every \(\gamma>0\), and \(r>0\),
\[
 \Pr_\nu\{J_e\ge r\}\le
 \min\{1,4e^{-\mathcal I_\gamma(r/2)}\}
 \le\min\{1,4e^{-d_\gamma(r)/4}\},
 \tag{V77.2}\label{r77:Jtail}
\]
where
\[
 \mathcal I_\gamma(x)=x\operatorname{arsinh}(x/\gamma)
                  -\sqrt{\gamma^2+x^2}+\gamma.
 \tag{V77.3}\label{r77:rate}
\]
Consequently, for \(0<t<1/4\),
\[
 \E_\nu e^{t d_\gamma(J_e)}
   \le1+\frac{16t}{1-4t}.
 \tag{V77.4}\label{r77:Dmgf}
\]
The constants have no graph-size factor. No independence between
different labels is asserted.
\end{theorem}
\begin{proof}
Given the full native count vector \(K\), introduce a reference
label \(\overline J\) with law \(\pi_{K_e}\) of
\eqref{r74:pi}. V74 proves that the actual conditional \(J_e\)
is stochastically dominated by this reference label. Conditional
on \(\overline J=j\), let \(M\) be uniform on
\(\{-j,-j+2,\ldots,j\}\). Character decomposition gives, for
every real \(a\),
\[
 \E[e^{aM}\mid K_e=k]
  =\sum_j\pi_k(j)\frac1{j+1}\sum_{q=0}^j e^{a(j-2q)}
  =(\cosh a)^k.
 \tag{V77.5}\label{r77:magnetic}
\]
Indeed the inner sum is \(\chi_j(\cosh a)\), and
\(\sum_j\lambda_j(k)P_j(t)=t^k\) is a polynomial identity,
valid also at \(t=\cosh a\). Thus \(M\), conditional on \(k\),
has the law of the sum of \(k\) independent symmetric signs.
This is a reference construction for one marginal, not a coupling
that replaces the full interacting angular network by independent
signs.

For any \(j\ge1\), at least half of its \(j+1\) magnetic weights
have absolute value at least \(j/2\), as direct counting of the
equally spaced weights shows. Hence, on \(\overline J\ge r>0\),
\(\Pr(|M|\ge r/2\mid\overline J)\ge1/2\). For \(a\ge0\),
the native factorial-moment bound \eqref{r74:factorial} implies
\[
 \E e^{aM}=\E(\cosh a)^{K_e}
 =\sum_{q\ge0}\frac{(\cosh a-1)^q}{q!}\E(K_e)_q
 \le e^{\gamma(\cosh a-1)}.
 \tag{V77.6}\label{r77:Skellam}
\]
Every term is nonnegative, so monotone convergence justifies the
series. Symmetry and Chernoff's inequality yield
\(\Pr(|M|\ge x)\le2e^{-\mathcal I_\gamma(x)}\);
the optimizing parameter is \(a=\operatorname{arsinh}(x/\gamma)\).
Combining this with posterior domination and the half-weight
bound proves the first inequality in \eqref{r77:Jtail}.

Both \(\mathcal I_\gamma\) and \(d_\gamma\) vanish at zero.
Their derivatives obey
\(\operatorname{arsinh}(x/\gamma)
 \ge x/\sqrt{\gamma^2+x^2}\), since the difference has
nonnegative derivative and vanishes at zero. Thus
\(\mathcal I_\gamma(x)\ge d_\gamma(x)\).
Rationalizing the square root also gives
\(d_\gamma(r/2)\ge d_\gamma(r)/4\).
As \(d_\gamma\) is increasing,
\(\Pr\{d_\gamma(J_e)>u\}\le4e^{-u/4}\) for \(u\ge0\).
Finally Tonelli's theorem applied to
\(e^{tY}=1+\int_0^Y t e^{tu}\,du\) proves
\eqref{r77:Dmgf}.
\end{proof}

\begin{corollary}[All angular moments, without assuming a Casimir mgf]
\label{r77:root-moment}
For \(\gamma\ge1\) and every edge,
\[
 \E_\nu\exp\sqrt{X_e/96}<6,\qquad
 \E_\nu X_e^p\le6\,96^p(2p)!,\quad p=0,1,\ldots.
 \tag{V77.7}\label{r77:allmoments}
\]
For every fixed finite graph, the joint law of \(X\) is determined
by all its mixed nonnegative integer moments.
\end{corollary}
\begin{proof}
For integer \(j\ge1\), \(\kappa_j\le3j^2\), and the same
inequality holds at zero. If \(Y=d_\gamma(j)\), then
\(j^2=Y^2+2\gamma Y\), so
\[
 \frac{\kappa_j}{2\gamma}
 \le\frac32Y^2+3Y\le\frac32(Y+1)^2\qquad(\gamma\ge1).
\]
Therefore \(\sqrt{X_e/96}\le(Y+1)/8\).
Equation \eqref{r77:Dmgf} at \(t=1/8\) bounds the exponential
moment by \(5e^{1/8}<6\). Expanding this nonnegative exponential
series and keeping its degree \(2p\) term proves the moment bound.

For completeness, the determinacy conclusion can be proved directly.
Give each coordinate of \(\sqrt X\) an independent symmetric random
sign. The resulting vector has all odd-coordinate moments zero;
its even-coordinate moments are the given mixed moments of \(X\).
If another nonnegative law has the same moments, the even series
for \(\cosh(c\sqrt{X_e})\), for \(0<c<1/\sqrt{96}\), converges
under that law too by the bound just proved. Since
\(e^{c|y|}\le2\cosh(cy)\), both symmetrized laws have a joint
exponential moment near zero, by Holder's inequality over the
finitely many coordinates. Their analytic moment generating
functions consequently have the same Taylor coefficients near
zero and are identical there. Uniqueness of the Fourier transform
(or analytic continuation into the common tube followed by Fourier
uniqueness) makes the symmetrized laws identical. Squaring their
coordinates recovers the same law of \(X\).
\end{proof}

This is the square-root exponential criterion discussed by
J.~Stoyanov and G.~D.~Lin, \emph{Hardy's Condition in the Moment
Problem for Probability Distributions},
\href{https://doi.org/10.1137/S0040585X9798631X}{doi:10.1137/S0040585X9798631X}.
We have supplied the native bound and the determinacy argument;
moment determinacy does not imply an analytic Casimir-source
transform.

\subsection{A genuine zero-radius obstruction at every cyclic edge}

\begin{theorem}[The finite-activity heat-source origin is not analytic]
\label{r77:obstruction}
If \(e\) is a bridge, then \(J_e=0\) almost surely. If \(e\)
belongs to a cycle, then for every \(\gamma>0\) and \(t>0\),
\[
 \E_\nu e^{t\kappa_{J_e}}=\infty.
 \tag{V77.8}\label{r77:infinite}
\]
For a cyclic edge the function
\(L_e(s)=\E_\nu e^{-sX_e}\), \(s\ge0\), has all right
derivatives at zero, but its Taylor series there has radius zero.
In particular it has no holomorphic extension to a neighborhood
of zero agreeing with this transform on the positive real axis.
The divergence in \eqref{r77:infinite} persists after conditioning
the joint law on \(K_f>\gamma/2\) for every edge \(f\).
\end{theorem}
\begin{proof}
For a bridge, rotate all spins on one component of its deletion
by a common Haar group element. Every internal bond is unchanged;
the bridge's character averages to zero unless its label is zero.
Thus \(S_G(J)=0\) if the bridge label is positive.

For a nonbridge edge choose a simple cycle \(C\ni e\) of length
\(\ell\). Give every edge of \(C\) label \(j\) and every other
edge label zero. Repeated character orthogonality around the cycle
gives \(S_G(J)=(j+1)^2\). These configurations have native weights
\[
 \nu_\gamma(J^{(j)})=
 \frac{a_0(\gamma)^{m-\ell}}{Z_G(\gamma)}
                     (j+1)^2 a_j(\gamma)^\ell,
 \qquad
 a_j(\gamma)\ge\frac{(\gamma/2)^j}{(j+1)!}.
 \tag{V77.9}\label{r77:cycle-lower}
\]
The last bound is the first term of the positive Bessel series.
The logarithm of the contribution to \(\E e^{t\kappa_{J_e}}\)
from this single label configuration is at least
\[
 t j(j+2)-\ell\log((j+1)!)+\ell j\log(\gamma/2)
                        +2\log(j+1)+c_{G,\gamma}.
\]
It tends to positive infinity, since
\(\log((j+1)!)\le(j+1)\log(j+1)\).
In particular the nonnegative summands do not tend to zero,
which proves divergence. The label law has all polynomial moments
also when \(\gamma<1\): \(J_e\le K_e\) in the joint lift and
the finite-graph count law has all exponential moments.
Dominated differentiation from \(s\ge0\) therefore gives
\(L_e^{(p)}(0+)=(-1)^p\E X_e^p\).
If its Taylor radius were positive, the series
\(\sum_p t^p\E X_e^p/p!\) would be finite for sufficiently small
\(t>0\), contrary to Tonelli and \eqref{r77:infinite}.
An analytic extension would have these same derivatives and is
therefore impossible.

Finally take \(j>\gamma/2\) in \eqref{r77:cycle-lower}.
Conditional counts on the cycle automatically exceed \(\gamma/2\),
because each is at least its label. Each off-cycle edge has label
zero and a strictly positive conditional probability
\(q_0=Q_{0,\gamma}(K>\gamma/2)\). Conditional count independence
makes the all-good probability on this configuration
\(q_0^{m-\ell}>0\), independent of \(j\). The same divergent
series survives this fixed multiplier and division by the positive
probability of the conditioning event.
\end{proof}

The obstruction is not a counterexample to V76: that theorem used
nonnegative heat sources and moment uniform integrability, not a
two-sided analytic neighborhood. It is a warning against the proposed
Cauchy route in those variables. In particular, negative sources
cannot be justified by the Gaussian limiting determinant: the
Gaussian squared-current law does have small positive exponential
moments, whereas the finite-activity law does not. A finite angular
cutoff would make the transform entire but would not supply a
cutoff-uniform analytic estimate across the origin.

\subsection{The exact count sources have better tails than their Casimir approximation}

\begin{proposition}[Exponential control of the actual two sources]
\label{r77:source-tail}
For every \(\gamma\ge1\) and \(j\ge0\),
\[
 -\tfrac32<\delta_j\le d_\gamma(j),\qquad
 0<w_j\le4+15d_\gamma(j).
 \tag{V77.10}\label{r77:source-control}
\]
Thus, for \(a,b\ge0\) with \(a+15b<1/4\),
\[
 \E_\nu e^{a|\delta_{J_e}|+b w_{J_e}}
 \le e^{3a/2+4b}
       \left(1+\frac{16(a+15b)}{1-4(a+15b)}\right).
 \tag{V77.11}\label{r77:source-mgf}
\]
The right-hand side is independent of the graph and activity.
\end{proposition}
\begin{proof}
The lower bound for \(\delta\) is V74's all-label bound.
Its upper square-root bound gives
\(\delta_j\le d_\gamma(j+1)-1\le d_\gamma(j)\), since
\(0\le d_\gamma'\le1\). For \(j=0\), V74 gives \(w_j\le4\).
If \(1\le j\le\gamma\), use \(\kappa_j\le3j^2\) and
\eqref{r74:dw} to get
\[
 w_j\le4+6j^2/\gamma
       \le4+6(1+\sqrt2)d_\gamma(j)<4+15d_\gamma(j).
\]
If \(j>\gamma\), then
\(w_j<\gamma<(1+\sqrt2)d_\gamma(j)<3d_\gamma(j)\).
These prove \eqref{r77:source-control}; applying
\eqref{r77:Dmgf} proves \eqref{r77:source-mgf}.
\end{proof}

The difference matters in the remote label tail. The quadratic
quantity \(\kappa_j/(2\gamma)\) has no positive exponential
moment, while \(\delta_j\) grows at most linearly in \(j\) and
\(w_j<\gamma\). An \(L^2\)-small replacement of one by the other,
such as V76.9, does not preserve exponential integrability. The
square-root source is viable by \eqref{r77:allmoments}, but it
would require fourth-order composite-source derivatives to recover
\(\Cov(X)\). The actual count sources give a more direct route.

\subsection{A holomorphic compensated source with the complete covariance Hessian}

For \(z\in\mathbb C^E\) define
\[
 \mathcal M_\gamma(z)=
 \E_\nu\exp\!\left\{\sum_e
       \left(z_e\delta_{J_e}-\tfrac12z_e^2w_{J_e}\right)\right\}.
 \tag{V77.12}\label{r77:source}
\]
The square \(z_e^2\), not \(|z_e|^2\), is essential for
holomorphy. Write \(F_\gamma=\log\mathcal M_\gamma\) on the
branch with \(F_\gamma(0)=0\), wherever established below.

\begin{theorem}[A graph-uniform analytic domain and exact diagonal compensation]
\label{r77:analytic}
For each fixed finite graph and activity,
\(\mathcal M_\gamma\) is entire. For all graphs and
\(\gamma\ge1\), on \(\|z\|_1\le1/64\),
\[
 |\mathcal M_\gamma(z)-1|<\tfrac13,
 \qquad |F_\gamma(z)|<\log(3/2).
 \tag{V77.13}\label{r77:zero-free}
\]
In particular the indicated branch is holomorphic throughout a
neighborhood of this closed ball. At the origin,
\[
 \partial_e F_\gamma=\E\delta_{J_e},\qquad
 \partial_e\partial_f F_\gamma
   =\Cov_\nu(\delta_{J_e},\delta_{J_f})
                   -\delta_{ef}\E w_{J_e}
   =\Gamma_{ef}-\gamma\delta_{ef}.
 \tag{V77.14}\label{r77:Hessian}
\]
This is the full native covariance remainder, not an angular
approximation with an omitted diagonal or source error matrix.
\end{theorem}
\begin{proof}
On any compact complex source set, the absolute exponential in
\eqref{r77:source} is bounded by a finite constant times
\(\exp(c\sum_eJ_e)\): use \(|\delta_j|\le j+3/2\) and
\(w_j<\gamma\). Since \(J_e\le K_e\), integrability follows
from the finite count generating function. It and all its
polynomially weighted derivatives are integrable on a slightly
larger compact set. Differentiation or the uniformly dominated
power series proves entire dependence on the finitely many sources.

For the uniform bound let \(r=\|z\|_1\), and abbreviate
\(Y_e=d_\gamma(J_e)\),
\(q=r+(15/2)r^2\), \(c=3r/2+2r^2\).
The exponent \(T\) in \eqref{r77:source} obeys
\[
 |T|\le c+\sum_e\bigl(|z_e|+(15/2)|z_e|^2\bigr)Y_e.
\]
The sum of the nonnegative coefficients is at most \(q\).
Convexity of the exponential, adding a zero dummy variable if
needed, and \eqref{r77:Dmgf} show
\[
 \E e^{|T|}\le e^c\left(1+\frac{16q}{1-4q}\right)
                    \quad(q<1/4).
 \tag{V77.15}\label{r77:joint-bound}
\]
This argument uses no label independence. Since
\(|e^T-1|\le e^{|T|}-1\), it suffices to bound this expression.
At \(r=1/64\), \(q=143/8192\) and \(c=49/2048\). Using
\(e^c\le(1-c)^{-1}\), its upper bound is
\[
 \frac{2048\cdot2477}{1999\cdot1905}<\frac43.
\]
The estimate improves for smaller \(r\), proving the first part
of \eqref{r77:zero-free}. The convergent series for
\(\log(1+(\mathcal M-1))\) proves the second and the branch
claim; strictness permits a slightly enlarged ball as well.
Finally differentiate the entire expectation twice at zero.
The first derivative of its exponent is \(\delta_{J_e}\);
the second is \(-\delta_{ef}w_{J_e}\).
The logarithm subtracts the product of the first means.
The total covariance identity \eqref{r76:Gamma} then proves
\eqref{r77:Hessian}, including its sign and its same-edge term.
\end{proof}

For real sources the label weight in \eqref{r77:source} is
positive. This does not assert that its character-expanded spin
kernel is pointwise positive, that it is a Markov heat kernel, or
that it represents a physical source in the continuum theory.
It is an exact analytic generator for the specified finite-graph
native covariance.

\subsection{The analytic reference and what can be uniform in a signed source bank}

Keep the full cycle projection \(P=UU^T\) of V76, and put
\(h(z)_e=z_e-z_e^2/2\). The corresponding limiting source is
\[
 F_*(z)=-\tfrac32\sum_ez_e
     -\tfrac32\log\det\bigl(I_\beta-U^T\operatorname{diag}(h(z))U\bigr),
 \tag{V77.16}\label{r77:reference}
\]
where the logarithm is continued from zero.

\begin{proposition}[Fixed-graph holomorphic convergence and a uniform reference bound]
\label{r77:complex-limit}
For each fixed graph, \(F_\gamma\to F_*\) locally uniformly
on \(\|z\|_1<1/64\), and every fixed source derivative converges
locally there. Moreover \(F_*\) is holomorphic on
\(\|z\|_2<1/4\), for every graph, and satisfies
\[
 |F_*(z)-DF_*(0)[z]|\le3\|z\|_2^2,
 \qquad \|z\|_2\le1/4.
 \tag{V77.17}\label{r77:reference-bank}
\]
Its origin Hessian is
\[
 D^2F_*(0)=\tfrac32\bigl(P\circ P-\operatorname{diag}P\bigr).
 \tag{V77.18}\label{r77:reference-Hessian}
\]
The first convergence assertion has no graph-uniform rate.
\end{proposition}
\begin{proof}
V76's joint convergence and its source errors give
\((\delta_{J_e},w_{J_e})_e\Longrightarrow(X_e^*-3/2,X_e^*)_e\)
for fixed \(G\). On the indicated source ball,
\eqref{r77:joint-bound} also bounds \(\E e^{2|T|}\) uniformly
in \(\gamma\), by replacing \(c,q\) with \(2c,2q\);
indeed \(2q<1/4\). This gives uniform integrability, so
expectations of the continuous complex exponentials converge.
For the limit use the Gaussian integral of V76:
\(\E e^{\sum_e h(z)_eX_e^*}
 =\det(I-U^T\operatorname{diag}(h(z))U)^{-3/2}\).
It is justified also for these complex sources by an integrable
Gaussian majorant, or by analyticity from small real sources.
Thus \(\mathcal M_\gamma\to e^{F_*}\) pointwise.
The uniform holomorphic bound \eqref{r77:zero-free} yields
equicontinuity on every compact subball, by Cauchy estimates in
coordinate polydiscs contained in the domain. Finite-dimensional
compactness and pointwise uniqueness give local uniform convergence;
Cauchy's formula then gives convergence of derivatives. The logarithm
is legitimate on the uniformly zero-free branch already proved.

For the reference bank estimate, let
\(A=U^T\operatorname{diag}(h(z))U\). If \(\|z\|_2\le r=1/4\),
then
\(\|A\|_{2\to2}\le r+r^2/2=9/32<1\) and
\(\|A\|_{\mathrm{HS}}\le(1+r/2)\|z\|_2=(9/8)\|z\|_2\).
These bounds follow by compression with \(U\), including when
\(\beta=0\). The trace series for the logarithm is therefore
holomorphic on this whole ball. After its linear part is subtracted,
\[
 F_*(z)-DF_*(0)[z]
 =-\tfrac34\sum_eP_{ee}z_e^2
                    +\tfrac32\sum_{k\ge2}\frac{\Tr A^k}{k}.
\]
For complex matrices, Schatten Holder gives
\(|\Tr A^k|\le\|A\|_{\mathrm{HS}}^2\|A\|_{2\to2}^{k-2}\).
Since \(P_{ee}\le1\), the absolute value is at most
\[
 \left(\frac34+\frac34\frac{(9/8)^2}{1-9/32}\right)\|z\|_2^2
 =\frac{381}{184}\|z\|_2^2<3\|z\|_2^2.
\]
Differentiating the trace series through degree two proves
\eqref{r77:reference-Hessian}. In particular the quadratic
compensation is exactly what produces the negative diagonal term.
\end{proof}

\subsection{The exact spatial requirement, without an invalid Cauchy disc}

On the cubic tori of V75, \(D^2F_*(0)=-(3/2)D\). Define the
centered difference, initially on the proved \(\ell^1\) ball, by
\[
 H_\gamma(z)=F_\gamma(z)-F_*(z)
                      -D(F_\gamma-F_*)(0)[z].
 \tag{V77.19}\label{r77:centered}
\]
Equations \eqref{r77:Hessian} and \eqref{r77:reference-Hessian}
give the \emph{exact} equality
\[
 D^2H_\gamma(0)=\Gamma-\gamma I+\tfrac32D.
 \tag{V77.20}\label{r77:exact-target}
\]
The source-error matrix of V76 is not discarded; the use of the
exact \(\delta,w\) has incorporated it into this source generator.

Here is a sufficient complex-source target with its missing input
explicit. Suppose constants \(r_0\in(0,1/4)\), \(C<\infty\),
and \(\gamma_0\) can be proved independent of the torus such that,
for every real \(s\) with \(\|s\|_2=1\), the germ
\(t\mapsto H_\gamma(ts)\) extends holomorphically to a
neighborhood of \(|t|\le r_0\) and
\[
 \sup_{|t|=r_0}|H_\gamma(ts)|\le C/\gamma,
                       \qquad\gamma\ge\gamma_0.
 \tag{V77.21}\label{r77:unpaid}
\]
Then Cauchy's integral formula, followed by the variational
characterization of a real symmetric matrix norm, proves
\[
 \|\Gamma-\gamma I+\tfrac32D\|_{2\to2}
                   \le\frac{2C}{r_0^2\gamma}.
 \tag{V77.22}\label{r77:criterion}
\]
Together with V75's explicit reference-resolvent remainder, this
would give the required \(\gamma\|\Gamma-T_\gamma\|\) bound.
This implication is proved; its hypothesis \eqref{r77:unpaid} is not.
It is a sufficient stronger analytic statement, not asserted
equivalent to the desired Hessian estimate.

The repaired domain alone does not supply it. For a unit
\(\ell^2\) vector \(s\), the paid \(\ell^1\) domain only
contains the disc \(|t|<1/(64\|s\|_1)\), which may have radius
\(1/(64\sqrt m)\). Even a uniform bound on that disc loses a
factor \(m\) in a second-derivative Cauchy estimate. Nor do marginal
exponential tails imply an \(\ell^2\) domain: as a purely logical
fixture, let all coordinates equal \(Y-1\) for one mean-one
exponential variable \(Y\). The centered log generator is
\(-\sum z_e-\log(1-\sum z_e)\). Its marginals have uniform
exponential moments, while the all-ones unit direction has a
singularity at \(t=m^{-1/2}\) and covariance norm \(m\).
This fixture is not the native graph law; it identifies the missing
spatial input rather than disproving it.

\subsection{Literature scope, verification, and the next step}

A targeted primary-source check considered J.~G.~Conlon and
M.~Dabkowski, \emph{Extensions of the Brascamp--Lieb Inequality
and the Dipole Gas} (2025),
\href{https://doi.org/10.1007/s10955-025-03478-x}{doi:10.1007/s10955-025-03478-x}.
Their scalar gradient action assumes a uniformly positive Hessian;
that hypothesis has not been proved for this compact interacting
angular law. No spatial covariance theorem from that paper is
imported. A square-root Poisson-semigroup approach was also screened;
we instead used the exact count sources to avoid a fourth-order
composite-source extraction. This is a targeted check, not a
replacement for the broad sweep scheduled at V80.

The checker verifies the magnetic-weight identity as an exact
rational identity, the half-weight count, the stated rational
zero-free constants, and source Hessians on finite rational
mixtures. Independent positive Bessel sums test the pointwise
source bounds. Native cycle sums check complex source convergence
and diagonal compensation, and rational projection matrices check
the Gaussian determinant's Hessian. Floating evaluations are
diagnostics, not proofs of divergence, analyticity, or a uniform
limit; those conclusions use the arguments above.

V76's regression passes and its mathematical body is preserved
byte-for-byte except the version title. Only the active predecessor
is replaced. Compilation artifacts remain outside
\texttt{latex\_notes}; companion files are unchanged. The previous
turn was progress, and this one pays a new native tail estimate,
all-order angular moment control, a concrete analytic obstruction,
and a valid compensated analytic source. It does not pay the
\(\ell^2\) spatial remainder.

The next attack should target connected response for this exact
compensated source, or a real-variable estimate of its Hessian.
It must retain dense angular correlations and prove control for
arbitrary normalized signed edge banks. Replacing \(\delta,w\)
by quadratic Casimirs inside exponentials would reintroduce the
obstruction just proved. Merely reapplying the \(\ell^1\) bound
would reproduce the volume loss. The interacting four-dimensional
continuum theory, general Wilson-exterior transport, physical-time
contraction, and the physical Yang--Mills mass gap remain unproved.
% END CONSOLIDATED V77 APPEND

% BEGIN CONSOLIDATED V78 APPEND
\clearpage
\section{V78: the complete first native covariance correction}
\label{r78:section}

\subsection{Context, transfer decision, and the exact scope of the advance}

The V75--V77 target is a volume-uniform bound for the actual compact
energy covariance, not another one-edge tail estimate. We now evaluate
its complete first inverse-activity correction and prove that this
matrix has uniformly bounded operator norm on the cubic tori. This
includes arbitrary signed edge banks. It does \emph{not} interchange
the large-activity and large-volume limits or bound the native
remainder uniformly. The full interacting continuum Yang--Mills
construction and physical Hamiltonian gap remain open.

The ESM V23 transfer suggestion is used as a test, not a theorem
about YM: combine all returns before asking which low-momentum terms
vanish. Here there is a genuine cancellation between two resistance
derivatives. However, the complete correction has a strictly positive
constant-source eigenvalue. Thus an extra derivative floor for that
whole correction is false. We retain its scalar contact instead of
discarding it as though ESM's different symmetry argument applied.
The companion audit at V75 already included the newer ESM V31;
the next regular companion and broad literature reviews are V80.

Non-duplication is important. V1 already computed the rooted scalar
coefficient \eqref{restart:graph-coefficient}. V2 already explained
coupling-differentiable Laplace expansions for a different, full-box
object. V58 bounded a centered Ward-score coefficient and its return,
not the complete coefficient of the native energy covariance.
We reuse those established starting points. The additions here are
the mean-zero resistance formula, its full conductance Hessian,
the all-volume bound for that Hessian, and the nonzero constant-bank
test. No general \(O(N)\) Griffiths inequality is used. The primary
record \url{https://arxiv.org/abs/2205.12928}, rechecked in this
iteration, remains withdrawn; that exclusion was already recorded
in V63--V64 and is not counted as a new result.

\subsection{Weighted compact integral and a differentiable coefficient}

Let \(G=(V,E)\) be a fixed finite connected simple graph, with
\(n=|V|\ge2\), \(m=|E|\), oriented incidence matrix \(B\), and
positive conductances \(\omega=(\omega_e)_{e\in E}\). Define
\[
 \begin{gathered}
 Z_\gamma(\omega)=\int_{(S^3)^V}
     \exp\{-\gamma\sum_{e=xy}\omega_e(1-u_x\cdot u_y)\}
                   \prod_x\dd H(u_x),\\
 L_\omega=B^{\mathsf T}\operatorname{diag}(\omega)B,
 \qquad C_\omega=L_\omega^\dagger,\\
 r_e(\omega)=b_e^{\mathsf T}C_\omega b_e,
 \qquad q_x(\omega)=(C_\omega)_{xx}.
 \end{gathered}
 \tag{V78.1}\label{r78:weighted}
\]
Here \(b_e=B^{\mathsf T}\mathbf e_e\), and the pseudoinverse is
the inverse on \(\mathbf1^\perp\), extended by zero on constants.
The Haar measure is normalized. Rooting one spin does not change
this integral, by common rotation invariance.

\begin{lemma}[Weighted coefficient with two differentiated remainders]
\label{r78:laplace}
On each compact subset of \((0,\infty)^E\), as \(\gamma\to\infty\),
\[
 \log Z_\gamma(\omega)
 =k_G-\frac{3(n-1)}2\log\gamma
       -\frac32\log\det L_\omega^o
       +\frac{a_G(\omega)}\gamma+O_{G,C^2}(\gamma^{-2}),
 \tag{V78.2}\label{r78:laplace-eq}
\]
where \(L_\omega^o\) is any rooted minor and \(k_G\) is independent
of \(\omega\). With \(C^o\) the inverse rooted minor extended
by a zero root row and column,
\[
 a_G(\omega)=\sum_{e=xy}\omega_e
       \left\{\frac58r_e^2+C^o_{xx}C^o_{yy}-(C^o_{xy})^2\right\}
                              -\sum_xC^o_{xx}.
 \tag{V78.3}\label{r78:rooted}
\]
The remainder and its first two conductance derivatives have the
stated order. All constants in this lemma may depend on the graph.
\end{lemma}
\begin{proof}
After rooting, the zero set of the action is the single constant
configuration. In exponential coordinates \(u_x=\exp(A_x/\sqrt\gamma)\),
the quadratic action is \(\frac12\sum_e\omega_e|A_y-A_x|^2\).
The weighted version of V1's exact quaternion expansion has quartic
contribution
\[
 -\frac1\gamma\sum_{e=xy}\omega_e
       \left\{\frac{|A_y-A_x|^4}{24}
                         +\frac{|A_x\times A_y|^2}{6}\right\}.
\]
The Haar correction is
\(-\sum_{x\ne o}|A_x|^2/(3\gamma)\).
The three independent Gaussian coordinates have covariance \(C^o\).
Their fourth moments give exactly \eqref{r78:rooted}, by the same
pairings as in V1; this does not omit the cross-product term.

For clarity, the derivative assertion is not obtained by formally
differentiating a pointwise big-O. Differentiate the original compact
integral zero, one, or two times before expanding. Each derivative
inserts \(-\gamma(1-u_x\cdot u_y)\); two derivatives insert the
corresponding product. In the rescaled chart these insertions and
their Taylor remainders are bounded by fixed polynomials in \(A\).
On a compact positive conductance set, the rooted quadratic form
has a common strictly positive least eigenvalue. Shrink the chart
so the exact rescaled action dominates a fixed positive quadratic
form there. On \(|A|\le\gamma^{1/12}\), Taylor expansion of the
action, Haar density, and these insertions through relative order
\(\gamma^{-1}\) has an integrated error at most
\(C_G\gamma^{-2}\): the pointwise remainders are bounded by this
power times fixed polynomials against a common integrable Gaussian.
The annular part of the chart is smaller than every inverse power,
by the same quadratic lower bound. Outside the chart, compactness
and the unique rooted minimum give exponential suppression, even
after the at-most-quadratic powers of \(\gamma\) from the insertions.

Thus the rescaled integral and its first two derivatives have
uniform relative two-term expansions. Their leading Gaussian
normalizer is positive and bounded away from zero on the chosen
conductance set. Taking its logarithm preserves the \(C^2\)
remainder and proves \eqref{r78:laplace-eq}.
\end{proof}

\begin{proposition}[Root-free formula with the collective-mode return]
\label{r78:rootfree}
The same coefficient is
\[
 \boxed{a_G(\omega)=\frac38\sum_e\omega_e r_e(\omega)^2
       -\frac34q(\omega)^{\mathsf T}L_\omega q(\omega)
       -\frac3n\Tr C_\omega.}
 \tag{V78.4}\label{r78:rootfree-eq}
\]
In particular the last term is present even when the diagonal of
\(C_\omega\) is constant.
\end{proposition}
\begin{proof}
Set \(p_x=C^o_{xx}\), with \(p_o=0\). Since
\(r_e=p_x+p_y-2C^o_{xy}\),
\[
 C^o_{xx}C^o_{yy}-(C^o_{xy})^2
 =-\frac14(p_x-p_y)^2+\frac12r_e(p_x+p_y)-\frac14r_e^2.
\]
Furthermore
\(\sum_{y\sim x}\omega_{xy}r_{xy}
  =2\mathbf1_{x\ne o}-(L_\omega p)_x\).
Multiply this identity by \(p_x/2\) and sum over vertices.
Substitution in \eqref{r78:rooted} yields
\(a_G=\frac38\sum_e\omega_e r_e^2-\frac34p^{\mathsf T}L_\omega p\).
The two Green functions satisfy
\[
 C^o_{xy}=(C_\omega)_{xy}-(C_\omega)_{xo}
                 -(C_\omega)_{oy}+(C_\omega)_{oo},\qquad
 p=q+q_o\mathbf1-2C_\omega\mathbf e_o.
\]
Use \(L_\omega C_\omega=I-\mathbf1\mathbf1^{\mathsf T}/n\)
and \(C_\omega L_\omega C_\omega=C_\omega\). Direct multiplication
then gives
\[
 p^{\mathsf T}L_\omega p
 =q^{\mathsf T}L_\omega q-4(q_o-n^{-1}\Tr C_\omega)+4q_o
 =q^{\mathsf T}L_\omega q+\frac4n\Tr C_\omega.
\]
This proves the formula without changing the native measure or
introducing an uncomputed collective-coordinate Jacobian.
\end{proof}

\subsection{A complete Hessian on the cubic torus}

Now take \(G=(\mathbb Z/N\mathbb Z)^3\), even \(N\ge4\), at
\(\omega=\mathbf1\). Thus \(n=N^3\), \(m=3n\), and set
\[
 \begin{aligned}
 L&=B^{\mathsf T}B,& C&=L^\dagger,& R&=BCB^{\mathsf T},
 &H&=R\circ R,\\
 S&=BC^2B^{\mathsf T},&Q_{xe}&=((Cb_e)_x)^2,
 &c&=C_{xx},&r&=R_{ee}=\frac{n-1}{3n}.
 \end{aligned}
 \tag{V78.5}\label{r78:matrices}
\]
All symbols in this display are local to V78; \(H\) is an edge
matrix, not a physical Hamiltonian. The \(n\times m\) matrix \(Q\)
is not V58's source-indexed quadratic random field. Translation
and cubic symmetry make \(c,r\) constant. The earlier Fourier
bound gives \(c\le13/16\). Also \(R^2=R\), \(0\preceq R\preceq I\),
and \(H\mathbf1=r\mathbf1\).

For the actual compact law define
\(d_e=\gamma(1-u_x\cdot u_y)\) and
\(\mathcal C_\gamma=(\Cov(d_e,d_f))_{ef}\).

\begin{theorem}[Full first energy-covariance coefficient]
\label{r78:hessian}
At each fixed \(N\), in edge-operator norm,
\[
 \mathcal C_\gamma=\frac32H+\frac1\gamma\mathcal C_1
                         +O_N(\gamma^{-2}),\qquad
 \boxed{\mathcal C_1=\frac34H^2-\frac32Q^{\mathsf T}LQ
                               -\frac6n(R\circ S).}
 \tag{V78.6}\label{r78:hessian-eq}
\]
This is the conductance Hessian of \(a_G\) at \(\mathbf1\),
not a selected subdiagram of that Hessian.
\end{theorem}
\begin{proof}
For real banks \(s,t\), put \(L_s=B^{\mathsf T}\operatorname{diag}(s)B\).
Differentiating the inverse on the fixed space \(\mathbf1^\perp\),
\[
 \partial_sC=-CL_sC,\qquad
 \partial_s\partial_t C=CL_sCL_tC+CL_tCL_sC.
 \tag{V78.7}\label{r78:inverse-jets}
\]
Consequently \(\partial_s q=-Qs\),
\(\partial_s r=-Hs\), and
\[
 \partial_s\partial_t r_e
     =2\sum_{f,h}s_ft_hR_{ef}R_{fh}R_{he}.
\]
Let \(F(\omega)=\sum_e\omega_e r_e(\omega)^2\).
Its mixed derivative is the sum of
\[
 -4r\,s^{\mathsf T}Ht,
 \qquad 2s^{\mathsf T}H^2t,
 \qquad 4r\sum_{e,f,h}s_ft_hR_{ef}R_{fh}R_{he}.
\]
The last term equals \(4r\,s^{\mathsf T}Ht\), since
\(\sum_eR_{he}R_{ef}=R_{hf}\). The first and last terms cancel
\emph{before} any norm is taken; thus \(D^2F=2H^2\).

At the base point \(q=c\mathbf1\). Since \(L\mathbf1=L_s\mathbf1=0\),
every term differentiating \(q^{\mathsf T}Lq\) twice vanishes
except \(2(\partial_s q)^{\mathsf T}L(\partial_t q)\).
This gives the \(-\frac32Q^{\mathsf T}LQ\) term. Finally
\[
 \partial_s\partial_t\Tr C
   =2\Tr(CL_sCL_tC)=2s^{\mathsf T}(R\circ S)t,
\]
which supplies the last term in \eqref{r78:hessian-eq}.

The exact identities
\(-\partial_e\log Z_\gamma=\E d_e\) and
\(\partial_e\partial_f\log Z_\gamma=\Cov(d_e,d_f)\)
now apply to the \(C^2\) expansion of Lemma \ref{r78:laplace}.
The Hessian of \(-\frac32\log\det L_\omega^o\) is
\(\frac32R_{ef}^2\); inverse differences of rooted and mean-zero
Green functions give the same edge matrix. Finite dimension at
this fixed \(N\) converts the entrywise remainder to operator norm.
No uniformity in \(N\) is asserted for that conversion.
\end{proof}

\subsection{A spatial operator bound with no volume factor}

\begin{lemma}[Squared Green-dipole bank]
\label{r78:dipoles}
The following are positive-semidefinite inequalities:
\[
 0\preceq Q^{\mathsf T}LQ\preceq4rcI,
 \qquad 0\preceq R\circ S\preceq\frac c3 I,
 \qquad 0\preceq H\preceq rI.
 \tag{V78.8}\label{r78:dipoles-eq}
\]
\end{lemma}
\begin{proof}
Write \(h_e=Cb_e\). For an edge \(f=xy\),
\[
 (BQs)_f=\sum_e s_eR_{fe}\{h_e(x)+h_e(y)\}.
\]
The orientation sign follows from using the same orientation for
\(B\) and \(R\); it does not affect the following squared estimate.
Cauchy--Schwarz and \(\sum_eR_{fe}^2=r\) give
\[
 \|BQs\|_2^2
 \le r\sum_e s_e^2\sum_{f=xy}\{h_e(x)+h_e(y)\}^2
 \le12r\sum_e s_e^2\|h_e\|_2^2.
\]
The factor \(12\) is twice the degree six: each vertex occurs
in six unordered edges. The diagonal of \(S\) is constant by
edge symmetry, and
\(\Tr S=\Tr(C^2L)=\Tr C=nc\); hence
\(\|h_e\|_2^2=S_{ee}=c/3\). This proves the first bound.

The Schur product is positive: if \(S=\sum_jv_jv_j^{\mathsf T}\),
then \(R\circ S=\sum_j\operatorname{diag}(v_j)
R\operatorname{diag}(v_j)\). Applying \(0\preceq R\preceq I\)
inside this sum gives \(0\preceq R\circ S\preceq\operatorname{diag}S\).
Finally \(H\) is positive by the same factorization, has
nonnegative entries, and has row sum \(r\). The elementary
row-sum norm bound yields \(H\preceq rI\).
\end{proof}

\begin{theorem}[Uniform bound for the complete first coefficient]
\label{r78:uniform}
For every even \(N\ge4\),
\[
 -2cI\preceq\mathcal C_1\preceq\frac34r^2I,
 \qquad \|\mathcal C_1\|_{2\to2}\le\frac{13}{8}.
 \tag{V78.9}\label{r78:uniform-eq}
\]
In particular
\[
 \sup_{N\ {
m even},\ N\ge4}\ 
 \lim_{\gamma\to\infty}
       \gamma\|\mathcal C_\gamma-\tfrac32H\|_{2\to2}
                                                        \le\frac{13}{8}.
 \tag{V78.10}\label{r78:order}
\]
This order of supremum and limit is part of the statement.
\end{theorem}
\begin{proof}
The two subtracted matrices in \eqref{r78:hessian-eq} are positive.
Their sum is at most
\((6rc+2c/n)I=2cI\), using \(3r=1-1/n\).
The positive term is at most \(\frac34r^2I\). Thus
\(\|\mathcal C_1\|\le\max\{2c,3r^2/4\}\le13/8\), since
\(c\le13/16\) and \(r\le1/3\). The fixed-volume expansion
then proves \eqref{r78:order}. Taking the supremum before the
large-activity limit is not justified by this argument.
\end{proof}

\subsection{Count contact, the wire target, and a nonvanishing constant bank}

Let \(\Gamma_\gamma=\Cov(K,K)\) be the actual count covariance
from V75, and set \(D=rI-H\). Homogeneity and edge symmetry give
\[
 a_G(\mathbf1)=\frac38mr^2-3c,\qquad
 \alpha:=\partial_ea_G(\mathbf1)=-\frac38r^2+\frac cn.
 \tag{V78.11}\label{r78:alpha}
\]
Indeed \(a_G(t\omega)=t^{-1}a_G(\omega)\), so the sum of its
edge derivatives is \(-a_G\). This also verifies the sign of
the contact in the next result.

\begin{corollary}[Complete count and wire coefficients]
\label{r78:count}
At fixed \(N\),
\[
 \begin{aligned}
 \E d_e&=\frac32r-\frac\alpha\gamma+O_N(\gamma^{-2}),\\
 \Gamma_\gamma&=\gamma I-\frac32D
                 +\frac1\gamma\mathcal J_1+O_N(\gamma^{-2}),
 &\mathcal J_1&=\mathcal C_1+\alpha I,
 &\|\mathcal J_1\|&\le\frac53.
 \end{aligned}
 \tag{V78.12}\label{r78:count-eq}
\]
For the V75 target \(T_\gamma=2\gamma^2(2\gamma I+3D)^{-1}\),
\[
 \begin{aligned}
 \lim_{\gamma\to\infty}\gamma(\Gamma_\gamma-T_\gamma)
       &=\mathcal J_1-\frac94D^2=:\mathcal W_1,
 &\|\mathcal W_1\|&\le\frac{23}{12},\\
 \lim_{\gamma\to\infty}
       \{(2\gamma I+3D)\Gamma_\gamma-2\gamma^2I\}
       &=2\mathcal W_1,
 &\|2\mathcal W_1\|&\le\frac{23}{6}.
 \end{aligned}
 \tag{V78.13}\label{r78:wire}
\]
These coefficient bounds are independent of \(N\); the limiting
statements are still fixed-volume statements.
\end{corollary}
\begin{proof}
The mean expansion follows by differentiating \eqref{r78:laplace-eq}.
Retain the exact count contact
\(\Gamma_\gamma=\gamma I-\operatorname{diag}(\E d_e)
 +\mathcal C_\gamma\), proved in V75. It gives the stated
\(\mathcal J_1\), including \(\alpha I\).
Its lower eigenvalue is at least
\(-2c-3r^2/8+c/n\ge-13/8-1/24=-5/3\).
Its upper eigenvalue is at most
\(3r^2/8+c/n\le1/24+13/1024<5/3\).
Also \(0\preceq D\preceq rI\). Expanding the finite matrix inverse
gives
\(T_\gamma=\gamma I-\frac32D+\frac9{4\gamma}D^2
 +O(\gamma^{-2})\), with a volume-uniform inverse-series error.
Thus \(\|\mathcal W_1\|\le5/3+1/4=23/12\).
Multiplication by \(2\gamma I+3D\) proves the last line;
the native remainder in this multiplication is only known at fixed \(N\).
\end{proof}

\begin{proposition}[The complete scalar constant mode does not vanish]
\label{r78:constant}
For \(\mathbf1_E\) the constant edge bank, define
\(\kappa_N=3r^2/8-c/n\). Then
\[
 \begin{aligned}
 Q^{\mathsf T}LQ\mathbf1_E&=0,
 &\mathcal C_1\mathbf1_E&=2\kappa_N\mathbf1_E,\\
 \mathcal J_1\mathbf1_E&=\kappa_N\mathbf1_E,
 &\mathcal W_1\mathbf1_E&=\kappa_N\mathbf1_E,\\
 \frac{907}{32768}&\le\kappa_N\le\frac1{24}.
 \end{aligned}
 \tag{V78.14}\label{r78:constant-eq}
\]
Thus a zero-momentum derivative floor for these complete scalar-bank
corrections is false. This says nothing about a different, separately
defined physical transverse sector.
\end{proposition}
\begin{proof}
The identity \(\sum_e(Cb_e)(Cb_e)^{\mathsf T}=CLC=C\) gives
\(Q\mathbf1_E=c\mathbf1_V\). Hence its Dirichlet-energy return
annihilates constants. Also \(H^2\mathbf1_E=r^2\mathbf1_E\).
Since \(RS=S\),
\(\sum_fR_{ef}S_{ef}=(RS)_{ee}=S_{ee}=c/3\).
Insert these identities in \eqref{r78:hessian-eq} and add the
contact \(\alpha I=-\kappa_NI\). Finally \(D\mathbf1_E=0\).
For the lower bound, \(n\ge64\) implies \(r\ge21/64\), so
\[
 \kappa_N\ge\frac38\left(\frac{21}{64}\right)^2
                       -\frac{13}{16\cdot64}
                  =\frac{907}{32768}>0.
\]
The upper bound uses \(c\ge0\) and \(r\le1/3\).
\end{proof}

There is a legitimate bookkeeping split
\(\mathcal J_1=\kappa_NI+
 (\mathcal J_1-\kappa_NI)\), whose second term annihilates
\(\mathbf1_E\). It is not permission to remove the first term
from the native observable. Nor does a row-sum-zero identity alone
prove a uniform quadratic Fourier bound near zero momentum.

\subsection{Verification and the remaining nonperturbative obligation}

The accompanying checker compares \eqref{r78:rootfree-eq} against
the independent rooted formula, at uniform and nonuniform rational
conductances and at every root of six small graphs. It differentiates
the rooted formula directly, rather than reusing the simplified
Hessian, and checks the complete Hessian on signed rational banks
for cycles, the tetrahedral graph, and the cube. Deterministic
cubic-torus matrix checks and actual character-sum cycle covariances
are separate floating diagnostics. They do not certify the
infinite-volume claim or the analytic Laplace remainder; those
claims use the proofs above. No Monte Carlo estimate is used.

The exact next remainder is now identifiable without missing contacts:
\[
 \mathcal R_{\gamma,N}
   :=\mathcal C_\gamma-\frac32H-\frac1\gamma\mathcal C_1.
 \tag{V78.15}\label{r78:remainder}
\]
At each fixed \(N\), \(\|\mathcal R_{\gamma,N}\|=O_N(\gamma^{-2})\).
Since \(\sup_N\|\mathcal C_1\|\le13/8\), the still-unpaid
volume-uniform target is equivalently
\[
 \sup_{\substack{N\ge4\ {
m even}\\\gamma\ge\gamma_0}}
           \gamma\|\mathcal R_{\gamma,N}\|_{2\to2}<\infty.
 \tag{V78.16}\label{r78:unpaid}
\]
Here \(\gamma_0\) must be independent of \(N\). Demanding a
uniform \(O(\gamma^{-2})\) remainder would be stronger than needed
and is not silently substituted for this target. A useful next
attack must control the actual compact spatial return for arbitrary
unit signed banks. It cannot infer that control from the finite
coefficient, the V77 marginal exponential tails, or the cancellation
of just \(Q^{\mathsf T}LQ\)'s constant mode.

This iteration pays the complete first coefficient of the native
covariance, with explicit volume-independent constants, and rules
out an overly strong derivative-floor shortcut. The V77 mathematical
body is preserved byte-for-byte except for its version title; only
the active predecessor is replaced. All build and diagnostic files
remain outside \texttt{latex\_notes}. The general Wilson-exterior
transport, physical-time contraction, interacting continuum
construction, and full physical Yang--Mills mass gap remain unproved.
% END CONSOLIDATED V78 APPEND

% BEGIN CONSOLIDATED V79 APPEND
\clearpage
\section{V79: the full orientation contact and a uniform spatial floor}
\label{r79:section}

\subsection{Companion review, scope, and the upstream change}

The four newly supplied companions have been reviewed at their current
terminal developments: source selection V43, exact-action Einstein bridge
V36, native stationarity V46, and perturbative ESM continuum V51. Their
new proof mechanisms are relevant, but none supplies the missing
Yang--Mills spatial estimate. The transfer decisions are as follows.
\begin{itemize}
\item Source-selection V43 proves a seven-event, duration-retaining
covariant stopping construction and injectivity of an independent
renewal-clock transform on count instruments. The latter uses
independence of the waiting times from labels and stopping decisions.
Those hypotheses are not established for the physical YM transfer;
neither its clock nor its source can be replaced by that example.
\item Exact-action bridge V36 retains constant-translation rows in a
spectral-parameter-dependent Schur reduction and finds an additional
real pair at the smaller inverse-amplitude scale. This warns against
discarding a sector merely because it is absent from a leading rapid
block. Its finite canonical pencil is not a positive YM Hamiltonian.
\item Native stationarity V46 finds two genuine neutral forcing channels,
then uses the exact constraints and a graded coordinate change to gain
an amplitude factor in the feedback. Its longer physical interval still
shrinks with amplitude. We import the requirement to retain the entire
constraint-compatible return, not that finite-dimensional lifespan.
\item Perturbative ESM V51 constructs a polynomial metric Noether section
at each fixed scalar degree, including the extra curl term for nonclosed
coefficient one-forms. Its full BV subtraction retains higher-matter
partners. Fixed-degree coefficient bounds are not convergence of an
infinite field series or an estimate for the original compact YM law.
\end{itemize}
The review checks these terminal arguments and their stated inherited
inputs; it is not an independent reproof of every theorem in the four
large files. A file-and-hash ledger records the inspected sources and
the limitations of transfer. The regular companion and broad-literature
checkpoints at V80 remain due.

V78 computed the complete first covariance coefficient and bounded its
operator norm uniformly in volume. It also found a nonzero scalar
constant-mode coefficient. Its suggested split by that scalar does
\emph{not} settle the full spatial question: there are three edge
orientations at zero momentum, not one. We now evaluate this entire
block, prove a uniform separation of its scalar and traceless
orientation responses, and prove an actual second-moment estimate
for the contact-subtracted coefficient. The last step is stronger
than a row-sum-zero identity. All count contacts and the collective-mode
return of V78 are retained.

This is a spatial theorem for the \emph{first coefficient}, not for
the nonlinear remainder. The earlier current-metric estimate in V61
and the mixed-link variance witness in archive f16.2 V85 already cover
two tempting alternative routes; they are not claimed again as new.
The entire V78 mathematical prefix is preserved.

\subsection{Centered edge Fourier coordinates and the exact zero block}

Keep all V78 notation on the cubic torus, with even \(N\ge4\),
\(n=N^3\), and \(r=(n-1)/(3n)\). An edge \((x,j)\) goes from
\(x\) to \(x+e_j\); its geometric location is \(x+e_j/2\).
We use the unitary Fourier convention
\[
 \widehat f_j(p)=n^{-1/2}\sum_x
          e^{-ip\cdot(x+e_j/2)}f_{x,j},\qquad
 p\in(2\pi/N)\mathbb Z^3\pmod{2\pi}.
 \tag{V79.1}\label{r79:fourier}
\]
For estimates choose each component in \([-\pi,\pi]\).
Symbols in these coordinates are conjugate by a diagonal unitary
matrix to the ordinary site-based edge symbols. Thus positivity and
operator norms are unchanged. An off-diagonal symbol need not be
periodic as a scalar function when a single component of \(p\) is
increased by \(2\pi\); its prescribed sign change is retained.

Put \(a_j(p)=4\sin^2(p_j/2)\) and \(\lambda(p)=\sum_j a_j(p)\).
Every sum below omits \(p=0\), or assigns the summand zero there.
Define
\[
 \begin{aligned}
 c&=\frac1n\sum_{p\ne0}\lambda(p)^{-1},&
 \beta&=\frac1n\sum_{p\ne0}\frac{a_1(p)a_2(p)}{\lambda(p)^2},\\
 \beta'&=\frac1n\sum_{p\ne0}\frac{a_1(p)a_2(p)}{\lambda(p)^3},&
 d&=r-3\beta,&
 \kappa&=\frac38r^2-\frac cn.
 \end{aligned}
 \tag{V79.2}\label{r79:sums}
\]
The prime on \(\beta'\) is a name for the second displayed sum,
not differentiation. Let \(P_{\parallel}=\mathbf1_3\mathbf1_3^T/3\)
and \(P_{\perp}=I_3-P_{\parallel}\). Here ``perpendicular'' means
traceless \emph{edge-orientation} coefficients. It is not a physical
gauge-transverse projection.

\begin{proposition}[The complete three-orientation block]
\label{r79:zero}
The first count coefficient \(\mathcal J_1\) and the wire coefficient
\(\mathcal W_1\) of V78 have zero-momentum symbols
\[
 \begin{aligned}
 \widehat{\mathcal J}_1(0)
     &=\kappa P_{\parallel}+\eta P_{\perp},\\
 \widehat{\mathcal W}_1(0)
     &=\kappa P_{\parallel}
          +\left(\eta-\frac{81}{4}\beta^2\right)P_{\perp},\\
 \eta&=\frac34d^2-\frac6n\left(\frac c3-3\beta'\right)-\kappa.
 \end{aligned}
 \tag{V79.3}\label{r79:zero-eq}
\]
In particular this is a matrix contact, not merely the scalar
\(\kappa I_3\).
\end{proposition}
\begin{proof}
Write \(u_j(p)=2\sin(p_j/2)\). The centered symbol of \(R\) is
\(u(p)u(p)^T/\lambda(p)\). Parseval on each orientation pair gives
\[
 \begin{aligned}
 \widehat H_{ij}(0)
     &=\frac1n\sum_{p\ne0}\frac{a_i(p)a_j(p)}{\lambda(p)^2},\\
 \widehat{(R\circ S)}_{ij}(0)
     &=\frac1n\sum_{p\ne0}\frac{a_i(p)a_j(p)}{\lambda(p)^3}.
 \end{aligned}
 \tag{V79.4}\label{r79:parseval-zero}
\]
Their diagonal entries are respectively \(r-2\beta\) and
\(c/3-2\beta'\), and their off-diagonal entries are
\(\beta\) and \(\beta'\). This follows by summing over \(j\),
using cubic symmetry, and retaining the missing constant vertex mode.
Thus their eigenvalues on the two projections are respectively
\((r,d)\) and \((c/3,c/3-3\beta')\).
The zero symbol of \(Q^TLQ\) is zero because \(\lambda(0)=0\).
Insert these identities into V78's complete formula
\[
 \mathcal J_1=\frac34H^2-\frac32Q^TLQ
       -\frac6n(R\circ S)-\kappa I.
 \tag{V79.5}\label{r79:J}
\]
On \(P_{\parallel}\) the answer is \(\kappa\); on
\(P_{\perp}\) it is \(\eta\). Finally \(D=rI-H\) has zero
symbol \(3\beta P_{\perp}\), and
\(\mathcal W_1=\mathcal J_1-9D^2/4\). This proves all coefficients,
including the finite-volume collective return.
\end{proof}

\begin{theorem}[Uniformly opposite orientation responses]
\label{r79:signs}
For every even \(N\ge4\),
\[
 \boxed{\beta\ge\frac1{16},\qquad
  \kappa\ge\frac{907}{32768},\qquad
  \eta\le-\frac{421}{32768},\qquad
  \kappa-\eta\ge\frac{83}{2048}.}
 \tag{V79.6}\label{r79:signs-eq}
\]
Consequently
\[
 \inf_{b\in\R}\|
          \widehat{\mathcal J}_1(0)-bI_3\|
       =\frac{\kappa-\eta}{2}\ge\frac{83}{4096}.
 \tag{V79.7}\label{r79:scalar-obstruction}
\]
No scalar subtraction gives a vanishing full zero-momentum block.
The count and wire zero blocks have one positive and two negative
eigenvalues.
\end{theorem}
\begin{proof}
For the uniform frequency measure, independence of the three coordinate
frequencies and \(N\ge4\) give
\(\langle a_j\rangle=2\) and \(\langle a_j^2\rangle=6\).
Therefore \(\langle a_1a_2\rangle=4\) and
\(\langle a_1a_2\lambda\rangle=32\).
The tangent-line inequality for the convex function \(t^{-2}\) at
\(t=8\) is
\[
                 t^{-2}\ge\frac3{64}-\frac{t}{256},\qquad t>0.
\]
Multiplying by \(a_1a_2\) and averaging proves
\(\beta\ge12/64-32/256=1/16\). The zero frequency has zero
weight and creates no exception to this multiplication.

Positivity of \(H\) and \(R\circ S\), already established in V78,
implies \(d\ge0\) and \(c/3-3\beta'\ge0\). Hence
\(0\le d\le r-3/16\) and
\[
 \eta\le\frac34(r-3/16)^2-\frac38r^2+\frac cn
      =\frac38r^2-\frac9{32}r+\frac{27}{1024}+\frac cn.
\]
The quadratic on the right decreases for
\(21/64\le r\le1/3\). Since \(n\ge64\) and \(c\le13/16\),
its upper bound is
\[
 \frac38\left(\frac{21}{64}\right)^2
 -\frac9{32}\frac{21}{64}+\frac{27}{1024}+\frac{13}{1024}
                       =-\frac{421}{32768}.
\]
V78 gives the stated lower bound for \(\kappa\). Subtraction gives
the separation. A real symmetric matrix with just the two extreme
eigenvalues \(\kappa,\eta\) has optimal scalar approximation at
their midpoint, with error half their separation. The wire's second
eigenvalue is even smaller, proving its sign statement as well.
\end{proof}

The signed coefficient does not imply a negative covariance:
\(\Gamma_\gamma\) still has leading term \(\gamma I\).
Nor does the traceless orientation sector describe physical gluon
states. The conclusion is an obstruction to a particular scalar
bookkeeping shortcut for the native edge bank.

\subsection{A discrete second-moment estimate with the midpoint twist retained}

For a translation-invariant edge matrix \(A\), let
\(A_{ij}(\delta)\) denote its kernel at midpoint displacement
\(\delta=x+(e_i-e_j)/2\). Representatives are chosen componentwise
with \(|\delta_\ell|\le N/2\). Boundary representatives differing by
\(N\) have the same Fourier phase at the allowed frequencies. Define
\[
 M_R=\sum_{i,j,\delta}|\delta|^2|R_{ij}(\delta)|^2,
 \qquad M_S=\sum_{i,j,\delta}|\delta|^2|S_{ij}(\delta)|^2.
 \tag{V79.8}\label{r79:moments}
\]
Each sum fixes one input cell and sums the relative displacement;
there is no additional factor of the total number of cells.

\begin{lemma}[Uniform squared-current moment]
\label{r79:moment}
The exact periodic kernels satisfy
\[
 \boxed{M_R\le3\pi^2c,\qquad
        M_S\le\frac{3cN^4}{64},\qquad
 \sum_{i,j,\delta}|\delta|^2
       |R_{ij}(\delta)S_{ij}(\delta)|
                     \le\frac{3\pi cN^2}{8}.}
 \tag{V79.9}\label{r79:moment-eq}
\]
The apparently growing last moment is multiplied by \(1/n\) in the
complete coefficient, so it causes no volume loss there.
\end{lemma}
\begin{proof}
For a real vector \(v\ne0\), set
\(\Pi(v)=vv^T/|v|^2\), and set \(\Pi(0)=0\).
For nonzero \(v,w\) the squared Frobenius difference is
\(2\sin^2\theta\), where \(\theta\) is their angle.
Projecting the longer vector perpendicular to the shorter gives
\[
 \|\Pi(v)-\Pi(w)\|_F^2
 \le2|v-w|^2\left(
       \frac{\mathbf1_{v\ne0}}{|v|^2}
       +\frac{\mathbf1_{w\ne0}}{|w|^2}\right).
 \tag{V79.10}\label{r79:projection}
\]
The inequality also holds when one vector is zero, directly from
\(\|\Pi(v)\|_F^2=1\).

Set \(h=2\pi/N\). On an unfolded frequency grid,
\(|u(p+he_\ell)-u(p)|\le h\). At a boundary of a period,
\(u_\ell\) changes sign under a full \(2\pi\) shift;
we keep that sign rather than identifying every centered matrix
entry periodically. Applying \eqref{r79:projection} and averaging,
\[
 \frac1n\sum_p
   \|\Pi(u(p+he_\ell))-\Pi(u(p))\|_F^2
                         \le4h^2c.
 \tag{V79.11}\label{r79:freq-difference}
\]
The zero vector has the zero convention just specified, and the shifted
frequency sum of the nonzero reciprocals is again \(nc\).

The twisted Fourier Parseval identity gives, entry by entry,
\[
 \sum_{i,j,\delta}4\sin^2(\pi\delta_\ell/N)
              |R_{ij}(\delta)|^2
 =\frac1n\sum_p
   \|\Pi(u(p+he_\ell))-\Pi(u(p))\|_F^2.
 \tag{V79.12}\label{r79:twisted-parseval}
\]
Indeed shifting the Fourier argument multiplies a kernel coefficient
by \(e^{-ih\delta_\ell}\); a half-integer displacement accounts for
precisely the boundary sign above. This also proves the identity for
the off-diagonal orientation entries, where ordinary periodic
identification would be wrong.
For \(|\delta_\ell|\le N/2\),
\(\sin(\pi|\delta_\ell|/N)\ge2|\delta_\ell|/N\).
Multiplying \eqref{r79:twisted-parseval} by \(N^2/16\) therefore
bounds the \(\ell\)-th coordinate moment by \(\pi^2c\).
Summing over three coordinates proves the bound for \(M_R\).

The centered symbol of \(S\) is \(\Pi(u)/\lambda\). Parseval and
\(\lambda_{\min}=4\sin^2(\pi/N)\ge16/N^2\) give
\[
 \sum_{i,j,\delta}|S_{ij}(\delta)|^2
   =\frac1n\sum_{p\ne0}\lambda(p)^{-2}
                         \le\frac{cN^2}{16}.
\]
Since \(|\delta|^2\le3N^2/4\), this proves the asserted bound
for \(M_S\). Weighted Cauchy--Schwarz then gives
\(\sum|\delta|^2|RS|\le\sqrt{M_RM_S}\le3\pi cN^2/8\).
No infinite-volume Green asymptotic or unproved pointwise decay
estimate was used.
\end{proof}

\subsection{A genuine quadratic floor for the complete first coefficient}

The midpoint kernels \(R,S\) are real and even under simultaneous
spatial inversion. To check evenness directly, invert the lattice
through the origin; both oriented edge gradients change sign, so
their product does not. The squared vertex-to-edge dipole kernel
\(Q\) is similarly even about its edge midpoint. Thus the centered
symbols of \(H=R\circ R\) and \(R\circ S\) may be written as
cosine sums. Representatives on a boundary of the displacement
cube cause no sine contribution: an integer displacement \(N/2\)
has sine zero at every allowed coordinate frequency. Half-integer
displacements have no such boundary tie when \(N\) is even.

\begin{theorem}[Uniform quadratic symbol differences]
\label{r79:floor}
For every allowed principal momentum \(p\),
\[
 \begin{aligned}
 \|\widehat{\mathcal C}_1(p)-\widehat{\mathcal C}_1(0)\|
   &\le K_C|p|^2,\\
 \|\widehat{\mathcal J}_1(p)-\widehat{\mathcal J}_1(0)\|
   &\le K_C|p|^2,\\
 \|\widehat{\mathcal W}_1(p)-\widehat{\mathcal W}_1(0)\|
   &\le K_W|p|^2,
 \end{aligned}
 \tag{V79.13}\label{r79:floor-eq}
\]
where one may take
\[
 \begin{aligned}
 K_C&=\frac94r\pi^2c+\frac{c^2}{2}+\frac{9\pi c}{8N},\\
 K_W&=9r\pi^2c+\frac{c^2}{2}+\frac{9\pi c}{8N}.
 \end{aligned}
 \tag{V79.14}\label{r79:constants}
\]
In particular both are bounded independently of \(N\), using
\(r\le1/3\), \(c\le13/16\), and \(N\ge4\).
\end{theorem}
\begin{proof}
The cosine identity and \(|1-\cos t|\le t^2/2\), followed by the
entrywise absolute-sum upper bound for matrix norm, give
\[
 \begin{aligned}
 \|\widehat H(p)-\widehat H(0)\|
       &\le\frac12 M_R|p|^2
          \le\frac32\pi^2c|p|^2,\\
 \|\widehat{(R\circ S)}(p)-\widehat{(R\circ S)}(0)\|
       &\le\frac{3\pi cN^2}{16}|p|^2.
 \end{aligned}
 \tag{V79.15}\label{r79:component-floors}
\]
The first matrix has \(0\preceq\widehat H(p)\preceq rI_3\)
at every allowed momentum. Therefore
\[
 \|\widehat H(p)^2-\widehat H(0)^2\|
       \le2r\|\widehat H(p)-\widehat H(0)\|.
\]
This uses \(A^2-B^2=A(A-B)+(A-B)B\), without assuming that
the two matrices commute.

For each orientation the nonnegative dipole square sums to
\(S_{ee}=c/3\). Hence every component of the row symbol
\(\widehat Q(p)\) has magnitude at most \(c/3\). It follows that
\[
 \|\widehat{Q^TLQ}(p)\|
  =\lambda(p)\|\widehat Q(p)\|_2^2
       \le\frac{c^2}{3}|p|^2,
 \qquad \widehat{Q^TLQ}(0)=0.
 \tag{V79.16}\label{r79:Qfloor}
\]
Insert these bounds into the complete V78 coefficient. Its three
contributions have constants \(9r\pi^2c/4\), \(c^2/2\), and
\(9\pi c/(8N)\). The scalar count contact has zero symbol
difference, so the same bound holds for \(\mathcal J_1\).
Finally \(0\preceq\widehat D(p)\preceq rI_3\) and
\(\widehat D(p)-\widehat D(0)=-(\widehat H(p)-\widehat H(0))\).
The term \(-9D^2/4\) adds \(27r\pi^2c/4\), proving the
wire constant. Every estimate concerns the full orientation matrix.
\end{proof}

\subsection{Realizing the reference by a finite-range native edge operator}

A constant matrix in midpoint Fourier coordinates is not automatically
a cell-local operator on differently centered edge components. We do
not silently identify these two notions. There is a finite-range
replacement with exactly the same zero block and the same quadratic
order of error.

For a real symmetric \(3\)-by-\(3\) matrix \(Z\), define
\(\mathfrak C_Z\) by its centered symbol
\[
 (\widehat{\mathfrak C}_Z)_{ii}(p)=Z_{ii},\qquad
 (\widehat{\mathfrak C}_Z)_{ij}(p)
       =Z_{ij}\cos(p_i/2)\cos(p_j/2)\quad(i\ne j).
 \tag{V79.17}\label{r79:contact}
\]
For \(i\ne j\), this is the average of the four orientation-\(j\)
edges incident to either endpoint of the orientation-\(i\) edge,
multiplied by \(Z_{ij}\). The four midpoint displacements are
\(\pm e_i/2\pm e_j/2\). Thus the definition is a genuine real,
self-adjoint finite-range operator on the original edge lattice.
It satisfies the required Fourier boundary signs.

\begin{corollary}[A retained finite-range contact and a derivative-order residual]
\label{r79:local}
For \(A=\mathcal J_1\) or \(\mathcal W_1\), set
\(Z=\widehat A(0)\) from \eqref{r79:zero-eq}. Then
\[
 A=\mathfrak C_Z+\mathfrak E_A,\qquad
 \|\widehat{\mathfrak E}_A(p)\|
        \le\left(K_A+\frac{\|Z\|}{4}\right)|p|^2,
 \tag{V79.18}\label{r79:local-floor}
\]
where \(K_A=K_C\) for the count coefficient and \(K_W\) for
the wire coefficient. Both constants are uniform in volume;
\(\|Z\|\le5/3\) and \(23/12\), respectively, suffice.
For every real signed edge bank \(s\), this also gives the form bound
\[
 |\langle s,\mathfrak E_A s\rangle|
 \le\frac{\pi^2}{4}\left(K_A+\frac{\|Z\|}{4}\right)
         \sum_p\lambda(p)\|\widehat s(p)\|_2^2.
 \tag{V79.19}\label{r79:form}
\]
The contact is retained as a summand, not removed from the actual
observable or used to change the measure.
\end{corollary}
\begin{proof}
For principal momenta, both cosines in \eqref{r79:contact} lie in
\([0,1]\). Hence
\[
 |1-\cos(p_i/2)\cos(p_j/2)|
       \le\frac{p_i^2+p_j^2}{8}.
\]
Each row has two off-diagonal entries, each with \(|Z_{ij}|\le\|Z\|\).
The absolute row-sum bound gives
\(\|\widehat{\mathfrak C}_Z(p)-Z\|\le\|Z\||p|^2/4\).
Combine this with Theorem \ref{r79:floor}. The V78 coefficient
norm bounds apply at zero momentum. Finally
\(|p|^2\le\pi^2\lambda(p)/4\) follows from
\(\sin(|p_i|/2)\ge|p_i|/\pi\), and Fourier Parseval proves
the asserted form inequality. No positivity of \(\mathfrak E_A\)
is required.
\end{proof}

\subsection{Verification, non-duplication, and the remaining original-law estimate}

The independent checker evaluates the zero block over rational frequency
data for \(N=4,6\), where every \(a_j(p)\) is an integer. For example,
at \(N=4\) it gives exactly
\[
 c=\frac{1517}{7680},\quad \beta=\frac{2263}{28800},\quad
 \kappa=\frac{2291}{61440},\quad
 \eta=-\frac{4117921}{122880000}.
 \tag{V79.20}\label{r79:exact-check}
\]
The all-volume inequalities above use analytic proofs, not extrapolation
from these two cases. Separate floating FFT checks at \(N=4,6,8,12\)
test every allowed momentum, the centered boundary signs, both weighted
moments, the exact twisted Parseval equality, the zero-block formulas,
and the finite-range contact residual. They are deterministic diagnostics,
not interval-certified proofs. The prior V78 mathematical checker is
also replayed, and recovery of its complete source prefix is checked by
SHA-256. All diagnostic and typesetting outputs are outside
\texttt{latex\_notes}.

The new information is genuinely spatial: the full first coefficient
has a non-scalar, signed orientation contact and, after that contact is
retained explicitly, a volume-uniform derivative-order residual. The
finite-range realization prevents midpoint phases from being mistaken
for an already local subtraction. None of this repeats the V61 current
Ward estimate, the archive mixed-link witness, or V78's unweighted
operator bound.

The unchanged unpaid quantity is still the actual compact remainder
\[
 \mathcal R_{\gamma,N}
     =\mathcal C_\gamma-\frac32H-\frac1\gamma\mathcal C_1,
 \qquad
 \sup_{\substack{N\ge4\ {\rm even}\\\gamma\ge\gamma_0}}
            \gamma\|\mathcal R_{\gamma,N}\|<\infty.
 \tag{V79.21}\label{r79:unpaid}
\]
The threshold \(\gamma_0\) must be independent of \(N\). We have
\emph{not} proved this bound, nor a quadratic symbol estimate for this
nonlinear remainder. In particular the new derivative floor cannot
be applied termwise to an unproved convergent large-activity expansion.
The justified next direction is to seek an original-law counterpart of
the weighted second-moment control, while retaining its complete
orientation contact and higher connected returns. A further scalar
contact ansatz is now rigorously excluded; higher formal coefficients
alone would not pay the remaining estimate.

General Wilson exteriors, physical-time contraction, the interacting
four-dimensional continuum construction, and the full physical YM mass
gap remain unproved. The companion developments sharpen the discipline
of the reduction; they do not supply any of these missing conclusions.
% END CONSOLIDATED V79 APPEND

% BEGIN CONSOLIDATED V80 APPEND
\clearpage
\section{V80: Flat mismatched exteriors and linear energy returns}
\label{r80:section}

\subsection{Context, checkpoint, and the distinction being tested}

The goal remains the interacting four-dimensional continuum YM theory
with a positive physical Hamiltonian gap; it is not achieved. V79 pays
a spatial bound for the complete first coherent covariance coefficient,
not for the actual compact remainder \eqref{r79:unpaid}. At this
every-ten-note checkpoint we review the companions and the primary
literature, then go upstream to the exterior quantifier. A theorem
uniform over \emph{all flat pairs} of Wilson boundaries would already
have to handle different central winding holonomies. V48's matched
adjoint connection has a parallel section; the example below does not.

The result is not another coherent perturbative coefficient. We compute
the complete minimum manifold, normalizer, and leading covariance
under an \emph{actual compact conditional Wilson law}. For one fixed
flat mismatched exterior its rescaled local-energy covariance grows
like \(\gamma\), even though the variance of the summed energy stays
bounded. Thus the coherent order-one covariance estimate cannot be
extended to all flat pairs without an additional term or an exterior
averaging argument. This does not contradict the coherent estimate
still sought in V79, and does not disprove a physical YM mass gap.

The four latest supplied companions remain Source Selection V43,
Exact-Action Einstein Bridge V36, Native Stationarity V46, and
Perturbative ESM V51. Their hashes agree with the V79 review. The
folder's NS V26, SM--Einstein V72, Shell Mixing V16, and Parallel YM V5
also remain unchanged. The current review rechecks their terminal
scope, rather than claiming an independent recertification of their
entire histories. The transferable discipline is to retain constrained
zero modes, elapsed-time information and induced returns. Independent
renewal clocks, shrinking finite-dimensional lifespans, and fixed-degree
polynomial sections are not native YM spatial covariance estimates.
Shell Mixing's measured one-loop flat coefficient is not a uniform
nonlinear theorem; Parallel YM's distinction between local topology
and extensive total charge remains important for exterior averaging.

\subsection{The primary-literature checkpoint: usable ideas and unpaid hypotheses}

This is a targeted broad search across low-temperature spin expansions,
gradient fields, functional inequalities, spectral independence and
lattice/continuum gauge constructions. It is not a claim to have
exhausted the literature. Relevant primary-source checks are as follows.
\begin{itemize}
\item Giuliani and Ott's
\href{https://arxiv.org/html/2302.07299v1}{low-temperature expansion paper}
retains its fixed-coordinate scope: the introduction explicitly
separates the expansion from large-distance asymptotics. The earlier
V59--V60 audit therefore still applies. The Balaban--O'Carroll
\href{https://doi.org/10.1007/s002200050510}{1999 primary abstract}
concerns one- and two-point asymptotics for its low-temperature
\(N\)-vector setting. That abstract is not a hard-spin,
all-periodic-volume operator-remainder theorem; the full text was not
available in this checkpoint.
\item Buchholz--Cotar--Schweiger,
\emph{Gradient Gibbs measures with non-convex potentials and the
universality class of the Gaussian Free Field},
\href{https://arxiv.org/html/2608.14526v1}{arXiv:2608.14526v1},
is a new useful lead. Its Definition 1.2 and Theorem 1.3 require
\(V'(t)\ge\alpha t\) for \(t\ge0\), with \(\alpha>0\), in the
stated even scalar setting. The Poincare part of Theorem 1.10(a)
requires oddness in at least one coordinate; part (b) requires an odd
observable. Lemma 3.16 similarly uses vanishing conditional means of
odd observables. These are substantive hypotheses, not cosmetic ones.
The compact bond potential \(1-\cos t\) fails the stated monotonicity
test at \(t=\pi\); native energy banks are even. Thus neither this
mixture representation nor its odd-observable inequality is imported
as the missing energy estimate.
\item Bodineau and Graham's
\href{https://arxiv.org/abs/1006.5622}{conservative-dynamics paper}
treats convex Ginzburg--Landau interactions. The approximate-inverse
spectral-independence results for the models in Chen--Yang--Yin--Zhang,
\href{https://arxiv.org/html/2404.04668}{arXiv:2404.04668},
also suggest useful representations, not an identified inequality for our
compact cyclic \(O(4)\) law. Caputo--Chen--Parisi,
\href{https://arxiv.org/html/2503.19419}{arXiv:2503.19419},
Theorem 2.8 retains a quantitative weak-interaction condition
\(R\lambda(\Gamma)<1\). No such condition is established here.
\item Cao--Nissim--Sheffield's lattice area-law developments,
\href{https://arxiv.org/abs/2505.16585}{arXiv:2505.16585} and
\href{https://arxiv.org/abs/2509.04688}{arXiv:2509.04688},
do not by themselves supply the weak-coupling interacting
four-dimensional continuum limit required here. The first abstract's
truncated-model/merger-term method is an architectural lead only;
its theorem is not asserted for our conditional model.
Chatterjee's \href{https://arxiv.org/abs/2401.10507}{arXiv:2401.10507}
has a Higgs/massive-Gaussian scaling target. The local-in-time
three-dimensional stochastic construction in
\href{https://arxiv.org/abs/2201.03487}{arXiv:2201.03487}
has a different dimension and conclusion. Neither is a pure interacting
four-dimensional YM gap theorem.
\item A negative bibliographic check matters: Lees's
\href{https://arxiv.org/abs/2205.12928}{arXiv:2205.12928}
was withdrawn on 31 March 2026, with a defect in the proposed bijection
reported on the primary record. Its advertised general Griffiths
inequalities are not used.
\end{itemize}

The following bookkeeping observation explains the odd/even issue
without importing any external theorem. For any normalized positive
mixture \(\mu=\int\mu_\kappa\,\rho(d\kappa)\), conditional variance gives
\[
 \Var_\mu(F)
 =\int\Var_{\mu_\kappa}(F)\,\rho(d\kappa)
       +\Var_\rho\!\left(\mu_\kappa(F)\right).
 \tag{V80.1}\label{r80:mixture}
\]
If all components are symmetric and \(F\) is odd, the last term
vanishes. Centering an even \(F\) only under \(\mu\) does not make its
component means equal. A conditional convexity argument for energies
would therefore still owe this second term, as well as a valid native
representation. This familiar identity is a scope audit, not a new
theorem or a closure of the original count covariance return.

\subsection{A literal flat Wilson pair with central winding mismatch}

Fix even \(N\ge4\), \(\Lambda_N=(\mathbb Z/N\mathbb Z)^3\),
\(n=N^3\), and the \(3n\) positively oriented nearest-neighbour edges
\(e=(x,x+\widehat e_j)\). Choose lower and upper spatial Wilson links
\[
 U_e=1,\qquad V_e=\varepsilon_e1,\qquad
 \varepsilon_{x,j}=\begin{cases}
 -1,&j=1,\ x_1=N-1,\\
 +1,&\text{otherwise}.
 \end{cases}
 \tag{V80.2}\label{r80:exterior}
\]
Every elementary spatial plaquette has product \(+1\) on either
boundary: a plaquette crossing the cut contains either zero or two
negative direction-one edges. Their winding holonomies in direction
one are respectively \(+1\) and \(-1\); in the other directions they
are \(+1\). Thus both connections are flat, but their fundamental
central holonomies cannot be identified by periodic vertex gauges.
No twisted boundary condition or altered magnetic plaquette action is
being imposed on the four-dimensional lattice. These are literal
allowed values of its two spatial slices.

For temporal links \(g_x\in\SU(2)\), write \(z_x\in S^3\) for their
real quaternion coordinates. The original incident Wilson trace is
\[
 q_e=\tfrac12\operatorname{Re}\Tr
       (U_e g_{x+\widehat e_j}V_e^{-1}g_x^{-1})
      =\varepsilon_e z_x\cdot z_{x+\widehat e_j}.
 \tag{V80.3}\label{r80:trace}
\]
Define, with every Haar factor of mass one,
\[
 \begin{gathered}
 S_\varepsilon(z)=\sum_e(1-q_e),\qquad
 Q_{\gamma,N}^{\varepsilon}=\int e^{-\gamma S_\varepsilon}\,dH^{\otimes n},
 \\
 d\mu_{\gamma,N}^{\varepsilon}
  =(Q_{\gamma,N}^{\varepsilon})^{-1}
       e^{-\gamma S_\varepsilon}\,dH^{\otimes n},\qquad
 d_e=\gamma(1-q_e),\qquad
 C_{\gamma,N}^{\varepsilon}=\Cov_{\mu_{\gamma,N}^{\varepsilon}}(d).
 \end{gathered}
 \tag{V80.4}\label{r80:law}
\]
This density is strictly positive. Its signed bonds do \emph{not}
justify importing the untwisted positive count/angular lift.
All arguments below work directly with \eqref{r80:law}.

Put \(\theta=\pi/N\), \(c_\theta=\cos\theta>0\),
\(s_\theta=\sin\theta>0\). These symbols are unrelated to V79's
Green diagonal \(c\). Let \(L_\varepsilon\) be the real signed
Laplacian defined by
\[
 a^{\mathsf T}L_\varepsilon a
       =\sum_{e=xy}(a_x-\varepsilon_e a_y)^2.
 \tag{V80.5}\label{r80:signed-laplacian}
\]

\begin{lemma}[Complete positive-energy minimum]
\label{r80:minima}
The minimum action and its complete minimum manifold are
\[
 \begin{gathered}
 E_N=n(1-c_\theta),\\
 \mathcal M_N=\left\{z_x=\cos(\theta x_1)v+\sin(\theta x_1)w:
        |v|=|w|=1,\ v\cdot w=0\right\}.
 \end{gathered}
 \tag{V80.6}\label{r80:ground}
\]
The manifold is an embedded copy of the five-dimensional Stiefel
manifold of orthonormal ordered pairs in \(\mathbb R^4\). At every
minimum, \(q_{x,1}=c_\theta\) and \(q_{x,2}=q_{x,3}=1\).
\end{lemma}
\begin{proof}
The signed scalar Laplacian is the ordinary difference operator with
antiperiodic continuation across the direction-one cut. Its complete
complex Fourier basis has momenta
\[
 k_1=(2\ell_1+1)\pi/N,\quad k_j=2\pi\ell_j/N\ (j=2,3),
 \qquad \ell_j=0,\ldots,N-1,
\]
and eigenvalues \(\sum_j2(1-\cos k_j)\). The smallest is
\(\lambda_0=2(1-c_\theta)\). Its real eigenspace is exactly the
span of \(\cos(\theta x_1)\) and \(\sin(\theta x_1)\), constant
in the other two coordinates. Since \(\sum_x|z_x|^2=n\),
\[
 S_\varepsilon(z)
    =\tfrac12\sum_{a=0}^3 (z^a)^{\mathsf T}L_\varepsilon z^a
    \ge\tfrac12 n\lambda_0=E_N.
\]
Equality forces every coordinate field into this lowest eigenspace.
Write the resulting spin field as in \eqref{r80:ground}. The unit
constraints at \(x_1=0\) and \(x_1=N/2\) force \(|v|=|w|=1\);
the constraint at \(x_1=1\) then forces \(v\cdot w=0\), because
\(\sin\theta\cos\theta\ne0\). Conversely these three constraints
give unit spins and attain equality at every site. Evaluation at
\(0,N/2\) recovers \((v,w)\), so the parametrization is an embedding.
The direct edge dot products, including the negative wrap edge, give
the last assertion.
\end{proof}

\subsection{The exact normal Hessian and the previously absent linear source}

Let \(B_j\) be the ordinary \emph{periodic} direction-\(j\) difference
matrix, \((B_j\phi)_x=\phi_{x+\widehat e_j}-\phi_x\), and put
\[
 L_{c_\theta}=c_\theta B_1^{\mathsf T}B_1+
                    B_2^{\mathsf T}B_2+B_3^{\mathsf T}B_3,
 \qquad K_\theta=L_\varepsilon-\lambda_0 I.
 \tag{V80.7}\label{r80:normal-operators}
\]
The dagger below denotes the Moore--Penrose inverse, with the constant
phase mode removed. Write \(\iota_1:\mathbb R^n\to\mathbb R^{3n}\)
for the isometry placing a field on the direction-one edges.

\begin{lemma}[Five collective modes and the normal derivative]
\label{r80:hessian}
At every minimum the action Hessian on the product-sphere tangent
space is orthogonally equivalent to
\[
                L_{c_\theta}\oplus K_\theta\oplus K_\theta.
 \tag{V80.8}\label{r80:hessian-eq}
\]
Its kernel has dimension \(1+2+2=5\) and is exactly
\(T\mathcal M_N\); hence it is positive on the normal bundle.
In these coordinates the derivative of the full edge trace vector is
\[
                 Dq(\phi,\psi_2,\psi_3)
                       =-s_\theta\iota_1 B_1\phi.
 \tag{V80.9}\label{r80:dq}
\]
\end{lemma}
\begin{proof}
By a global orthogonal quaternion rotation choose \(v=e_0,w=e_1\)
at the chosen minimum. Define the unit phase vector
\(t_x=-\sin(\theta x_1)e_0+\cos(\theta x_1)e_1\).
Every tangent vector is uniquely
\(h_x=\phi_x t_x+\psi_{2,x}e_2+\psi_{3,x}e_3\), an isometry of
the tangent space with three copies of \(\mathbb R^n\).

On the product sphere the constrained second variation is
\[
 \operatorname{Hess}S_\varepsilon[h,h]
             =\sum_{a=0}^3(h^a)^{\mathsf T}
                         (L_\varepsilon-\lambda_0 I)h^a.
\]
Indeed, \(L_\varepsilon z^a=\lambda_0z^a\), and the second
derivative of the unit-sphere constraint contributes
\(-\lambda_0\sum_x|h_x|^2\). The two constant transverse frames
therefore give \(K_\theta\). The phase frame has signed adjacent
dot product \(c_\theta\) in direction one and \(1\) in the other
directions. Its diagonal is \(6-\lambda_0=2c_\theta+4\), so its
block is exactly \(L_{c_\theta}\). Mixed blocks vanish.

The positive weighted periodic Laplacian has only one zero mode.
The preceding Fourier proof shows that \(K_\theta\) is nonnegative
with exactly two zero modes. A smooth manifold of minima contributes
its five tangent modes, so these exhaust the kernel. Directly
differentiating \(\varepsilon_e z_x\cdot z_y\) gives
\(s_\theta\phi_x-s_\theta\phi_y\) on a direction-one edge,
including the wrap edge, and zero on the other edges. Transverse
variations have zero first derivative. This proves \eqref{r80:dq};
it also shows that \(Dq\) annihilates all collective modes.
\end{proof}

\begin{theorem}[Native fixed-volume linear covariance return]
\label{r80:covariance}
For each fixed even \(N\ge4\), under the exact compact law
\eqref{r80:law},
\[
 \left\|\gamma^{-1}C_{\gamma,N}^{\varepsilon}-A_N\right\|
         =O_N(\gamma^{-1/2}),\qquad
 A_N=s_\theta^2\iota_1B_1L_{c_\theta}^{\dagger}
                       B_1^{\mathsf T}\iota_1^{\mathsf T}.
 \tag{V80.10}\label{r80:cov-limit}
\]
Here and below matrix norms are Euclidean operator norms. Moreover,
\[
 \|A_N\|=\frac{s_\theta^2}{c_\theta}>0,\qquad
 \operatorname{rank}A_N=n-N^2,\qquad A_N\boldsymbol1=0.
 \tag{V80.11}\label{r80:cov-size}
\]
In particular
\(\|C_{\gamma,N}^{\varepsilon}\|/
\gamma\longrightarrow s_\theta^2/c_\theta\).
\end{theorem}
\begin{proof}
We give the compact-measure justification, not merely the formal
Hessian calculation. Lemmas \ref{r80:minima}--\ref{r80:hessian}
identify the entire compact minimum manifold and a positive normal
Hessian \(H_u\) at each \(u\in\mathcal M_N\). At fixed \(N\) its
positive eigenvalues have a common positive lower bound. In normal
tubular coordinates \((u,x)\), uniformly in \(u\),
\[
 \begin{split}
 S_\varepsilon(u,x)&=E_N+\tfrac12x^{\mathsf T}H_ux+O_N(|x|^3),\\
 q(u,x)&=q_*+Dq_u x+O_N(|x|^2),\\
 dH^{\otimes n}&=(2\pi^2)^{-n}(1+O_N(|x|))\,
                     dx\,d\operatorname{vol}_{\mathcal M_N}(u).
 \end{split}
 \tag{V80.12}\label{r80:tube}
\]
The constant vector \(q_*\) is given in Lemma \ref{r80:minima}.
Decrease the tube radius so that the action above its minimum is
bounded above and below by positive multiples of \(|x|^2\).
On the compact complement it exceeds \(E_N\) by a fixed positive
amount. Thus off-tube contributions, even multiplied by a fixed
power of \(\gamma\), are exponentially small relative to the
normalizing integral.

Set \(x=\gamma^{-1/2}y\). The mean-value formula for the exponential
and the two positive quadratic lower bounds dominate its error by
\(C_N\gamma^{-1/2}|y|^3e^{-c_N|y|^2}\). The volume factor and the
Taylor errors of \(\sqrt\gamma(q-q_*)\), including its second
moment, add only fixed polynomial factors. Consequently the first
two moments of this vector differ by \(O_N(\gamma^{-1/2})\) from
those of \(Dq_uY_u\), where \(Y_u\) is the centered normal Gaussian
with covariance \(H_u^{-1}\) on the normal fibre. The Gaussian mean
is zero. Its covariance is independent of \(u\): the global
\(O(4)\) action is transitive on the ordered orthonormal pairs and
leaves every \(q_e\) invariant. Therefore averaging over the whole
minimum manifold introduces no omitted random mean or chosen phase.

Only the phase block in \eqref{r80:dq} contributes. On its normal
space its inverse is \(L_{c_\theta}^{\dagger}\); hence the covariance
is exactly \(A_N\). Since
\(\gamma^{-1}\Cov(d)=\gamma\Cov(q)
=\Cov(\sqrt\gamma(q-q_*))\), this proves \eqref{r80:cov-limit}.
Entrywise bounds imply the stated matrix bound at fixed finite \(N\).

For the ordinary periodic Fourier modes put
\(a_j(p)=4\sin^2(p_j/2)\). The only nonzero orientation block of
\(A_N\) has scalar symbol
\[
  s_\theta^2\frac{a_1(p)}{c_\theta a_1(p)+a_2(p)+a_3(p)},
  \qquad p\ne0,
 \tag{V80.13}\label{r80:symbol}
\]
with value zero at \(p=0\). Its maximum is
\(s_\theta^2/c_\theta\), attained whenever
\(p_1\ne0,p_2=p_3=0\). It is positive exactly when \(p_1\ne0\),
giving rank \((N-1)N^2\). The constant vector is annihilated by
\(B_1^{\mathsf T}\). This proves every assertion.
\end{proof}

\subsection{Normalizer and why the total-energy test misses the return}

Write \(\det''K_\theta\) for the product of its \(n-2\) positive
eigenvalues, and set
\[
 D_N=3n-5,\qquad
 \Delta_N=(\det' L_{c_\theta})(\det''K_\theta)^2,\qquad
 b_N=\frac{(2\pi)^{D_N/2}\operatorname{vol}(\mathcal M_N)}
                 {(2\pi^2)^n\sqrt{\Delta_N}}>0.
 \tag{V80.14}\label{r80:prefactor}
\]
All volumes use the metric induced from the product of unit round
spheres. More explicitly, the parametrization in \eqref{r80:ground}
multiplies the usual ambient Stiefel metric
\(|\delta v|^2+|\delta w|^2\) by \(n/2\). This follows from
\(\sum_x\cos^2(\theta x_1)=\sum_x\sin^2(\theta x_1)=n/2\) and
\(\sum_x\cos(\theta x_1)\sin(\theta x_1)=0\). Thus its volume is
\((n/2)^{5/2}\operatorname{vol}(V_2(\mathbb R^4))\) in that metric;
no collective-mode factor is silently normalized away.

\begin{proposition}[Suppression factor and summed-energy fluctuations]
\label{r80:normalization}
For each fixed \(N\),
\[
 Q_{\gamma,N}^{\varepsilon}
     =b_N e^{-\gamma E_N}\gamma^{-D_N/2}
                       (1+O_N(\gamma^{-1/2})).
 \tag{V80.15}\label{r80:normalizer}
\]
If \(Q_{\gamma,N}^{0}\) is the coherent integral with all
\(\varepsilon_e=1\), and \(b_N^0>0\) is its leading coefficient,
then
\[
 \frac{Q_{\gamma,N}^{\varepsilon}}{Q_{\gamma,N}^{0}}
   =\frac{b_N}{b_N^0}\,\gamma e^{-\gamma E_N}
                         (1+O_N(\gamma^{-1/2})).
 \tag{V80.16}\label{r80:ratio}
\]
Under the mismatched law itself,
\[
 \begin{split}
 \E\sum_e d_e&=\gamma E_N+\frac{D_N}{2}+O_N(\gamma^{-1/2}),\\
 \Var\!\left(\sum_e d_e\right)&=\frac{D_N}{2}+O_N(\gamma^{-1/2}).
 \end{split}
 \tag{V80.17}\label{r80:total-energy}
\]
\end{proposition}
\begin{proof}
Integrating the zeroth-order normal Gaussian in \eqref{r80:tube}
gives \eqref{r80:normalizer}; the determinant is independent of the
minimum by \eqref{r80:hessian-eq}. The coherent minimum has dimension
three and normal dimension \(3n-3=D_N+2\), with zero minimum action,
as already proved in V48. Dividing the two correctly measured
normalizers leaves the factor \(\gamma\) in \eqref{r80:ratio}.

For the last assertions insert \(\gamma(S_\varepsilon-E_N)\) and
its square into the same tubular integral. Their limits are the
first two moments of \(Y^{\mathsf T}H_uY/2\), which is half a
chi-square variable with \(D_N\) degrees of freedom. Its mean and
variance are both \(D_N/2\). The polynomial-weighted domination in
the preceding proof gives the stated errors. No derivative of an
uncontrolled asymptotic remainder is taken.
\end{proof}

For clarity, the divergence in Theorem \ref{r80:covariance} is not
seen by any spatially constant orientation bank. Formula
\eqref{r80:symbol} vanishes at zero momentum. On a nonzero
direction-one Fourier mode, however, its coefficient is
\(s_\theta^2/c_\theta\). Real sine or cosine banks yield the same
Rayleigh quotient. Thus a global energy fluctuation test cannot replace
the signed spatial operator estimate. Deterministically subtracting
\(\gamma(1-q_*)\), or subtracting the exact mean, leaves the covariance
unchanged. Subtracting the random linear phase term would be a different
observable and its return would have to be restored.

\subsection{Non-duplication, verification, and revised proof obligations}

V46's endpoint perturbations shrink at the thermal scale and produce
order-one noncentral Gaussian energy terms. V48--V49's matched flat
holonomies retain a parallel section and zero minimum action. Here
the central mismatch is fixed, the minimum is a five-dimensional
positive-energy helical manifold, and the leading local-energy
covariance is order \(\gamma\). The hypothesis \(r\ge1\) in V48 is
absent: the relative winding transport is \(-I\) on \(\mathbb R^4\),
so it has no nonzero parallel section. Hence that theorem was not
already a proof of this result. Archive and current-source searches
for the antiperiodic/Stiefel calculation found no earlier version.
The normal-tube method itself is inherited, not claimed as new.

An independent diagnostic assembles the tangent Hessian and edge
derivatives directly from the literal quaternion trace rather than
from \eqref{r80:hessian-eq}. Exact symbolic arithmetic on the four-site
signed ring gives five zero modes and the nonzero source-covariance
eigenvalue \(1/\sqrt2\). Separate floating checks on the cubic
\(N=4,6,8\) tori verify both flat boundaries, the two Hessian blocks,
their nullities, covariance ranks \(48,180,448\), and operator norms
approximately \(0.707107,0.288675,0.158513\). Geodesic finite
differences of the original compact action check the second variation.
These computations are diagnostics, not the proofs of the
all-\(N\) statements above. V79's mathematical checker is replayed,
and its complete source is recovered byte-for-byte by reversing only
the title and this append. All builds and diagnostics remain outside
\texttt{latex\_notes}; the active main source is V80.

The direction adjustment is precise:
\begin{enumerate}
\item Continue seeking the original \emph{coherent} spatial remainder
\eqref{r79:unpaid}; this note neither proves nor refutes it. In
particular V61's bounded current metric and V79's bounded coefficient
must not be relabelled as a nonlinear remainder estimate.
\item Do not impose its order-one energy-covariance scale uniformly
on all individually flat Wilson pairs. Theorem \ref{r80:covariance}
rules that out already at a fixed finite \(N\). General exteriors
require a transported-minimum expansion retaining linear source
returns, or an estimate integrated against the actual exterior law.
\item Keep the suppression \eqref{r80:ratio} together with the
large conditional covariance. A ratio of kernel values at two exact
boundary pairs is \emph{not} their marginal probability ratio: the
exact configurations have Haar measure zero, and neighbourhood volume,
spatial weights, and endpoint correlations remain to be controlled.
No exterior-averaged smallness is proved by this calculation.
\item All large-\(\gamma\) limits here are at fixed \(N\). Although
\(s_\theta^2/c_\theta\sim\pi^2/N^2\), the remainders and prefactors
are not uniform in \(N\). Neither this expression nor
\(E_N\sim\pi^2N/2\) licenses a simultaneous continuum scaling limit.
\end{enumerate}

The companion and broad-literature checkpoints are completed at V80;
the next regular reviews are V85 and V90 respectively. No external
theorem is imported with weakened hypotheses. The full continuum
construction, native spatial closure, physical-time contraction and
physical YM mass gap remain open.
% END CONSOLIDATED V80 APPEND

% BEGIN CONSOLIDATED V81 APPEND
\clearpage
\section{V81: Exterior-averaged energy bounds on Wilson time strips}
\label{r81:section}

\subsection{What is advanced, and the regime that is not being changed silently}

V80 rules out an order-one rescaled energy covariance bound at every
flat Wilson boundary pair. It does not rule out averaging that covariance
under the actual boundary distribution. This installment proves such an
average, together with the full predictor covariance and a centered
signed-source estimate, in an explicit anisotropic regime. The electric
couplings can be arbitrarily large; the \emph{total magnetic interaction
along each spatial-link history} must be small. The time interval is open
with integrated Haar endpoints, not a periodic-time or vacuum path law.

The result is uniform in spatial volume, electric couplings, and the
number of time steps while the declared magnetic budget is fixed. It is
not the isotropic weak-coupling estimate, the coherent all-volume remainder
of V79, or a physical mass-gap theorem. Its point is to obtain a genuine
original-law exterior average, rather than infer one from V80's ratio of
two fixed kernel values. Temporal gauge is used only as an exact
change of variables in the full integral; the original conditional
expectation is retained when the variance is split.

For nonduplication, archive f7 V91 already identifies link histories as
the independent electric reference factors and gives a conditional
polymer construction. Archive f14 V52 already proves a fine-lattice
Dobrushin theorem at small \emph{isotropic} Wilson coupling. Neither
supplies the following explicit rescaled-energy operator bound uniformly
at large electric coupling, with its entire centered real source domain
and its original temporal-link conditional projection. The product-history
representation, normalized-tilt estimate, and Dobrushin method themselves
are not new. We prove their needed specializations below.

A targeted primary-source check of L.~Wu,
\href{https://arxiv.org/pdf/math/0611635}{\emph{Poincare and transportation
inequalities for Gibbs measures under the Dobrushin uniqueness condition}},
Section 2 and Theorem 2.1, confirms the finite-product conditional-variance
inequality for general measurable spin spaces with the discrete metric.
The finite-block proof used here is included to fix the update-rate
normalization and avoid confusing that auxiliary gap with physical time.
The V80 broad-literature and companion reviews remain current; the next
regular checkpoints are V90 and V85 respectively.

\subsection{The exact anisotropic open-strip Wilson law}

Let \(\Lambda=(\mathbb Z/N\mathbb Z)^3\), \(N\ge3\), with positive
spatial edge set \(E\) and unoriented elementary spatial plaquettes
\(\mathcal P\), one chosen orientation per plaquette. Every spatial
edge belongs to four such plaquettes, each containing four distinct
edges. Use time slices \(t=0,\ldots,T\), \(T\ge1\), spatial links
\(U_{e,t}\in\SU(2)\), and temporal links \(g_{x,t}\), \(0\le t<T\).
Set \(\epsilon(q)=1-\tfrac12\operatorname{Re}\Tr q\in[0,2]\).
For \(e=(x,y)\), put
\[
 P_{e,t}=U_{e,t}g_{y,t}U_{e,t+1}^{-1}g_{x,t}^{-1},\qquad
 D_{e,t}=\gamma_{e,t}\epsilon(P_{e,t}),\qquad \gamma_{e,t}>0.
 \tag{V81.1}\label{r81:plaquette}
\]
For arbitrary \(b_{p,t}\ge0\), define
\[
 \begin{split}
 S(U,g)&=\sum_{e,t<T}\gamma_{e,t}\epsilon(P_{e,t})
             +\sum_{p,t\le T}b_{p,t}\epsilon(U_{p,t}),\\
 d\mu&=Z^{-1}e^{-S(U,g)}
           \prod_{e,t\le T}dH(U_{e,t})\prod_{x,t<T}dH(g_{x,t}).
 \end{split}
 \tag{V81.2}\label{r81:law}
\]
All Haar factors have mass one. There is no endpoint insertion,
identified temporal seam, selected holonomy, or Perron eigenfunction
in this definition. The magnetic part of \eqref{r81:law} is retained
in every final expectation.

The local time-integrated magnetic budget and the resulting constant are
\[
 a_e=\sum_{t=0}^T\sum_{p\ni e}b_{p,t},\qquad
 A=\max_e a_e<\frac13,\qquad
 K_A=\frac{3e^{2A}}{2(1-3A)}.
 \tag{V81.3}\label{r81:budget}
\]
If \(b_{p,t}=b_t\), write \(B=\sum_{t=0}^T b_t\). Then
\(A=4B\), so the sufficient condition and constant become
\[
                 B<\frac1{12},\qquad
                 K_A=\frac{3e^{8B}}{2(1-12B)}.
 \tag{V81.4}\label{r81:hom-budget}
\]

\begin{lemma}[Exact temporal gauge and independent electric histories]
\label{r81:gauge}
For every gauge-invariant observable, its integral under \eqref{r81:law}
equals its temporally gauge-fixed integral on spatial-link histories
\(X_e=(V_{e,0},\ldots,V_{e,T})\), with probability density
\[
 d\widehat\mu(X)=\widehat Z^{-1}e^{-\mathcal V(X)}
                              \prod_{e\in E}d\pi_e(X_e),\qquad
 \mathcal V(X)=\sum_{p,t}b_{p,t}\epsilon(V_{p,t}).
 \tag{V81.5}\label{r81:paths}
\]
In \(\pi_e\), the variables \(V_{e,0}\) and
\(Q_{e,t}=V_{e,t}V_{e,t+1}^{-1}\) are independent, with respective
laws Haar and
\[
 d\nu_{\gamma_{e,t}}(q)=z_{\gamma_{e,t}}^{-1}
                         e^{-\gamma_{e,t}\epsilon(q)}dH(q).
 \tag{V81.6}\label{r81:increment}
\]
The full electric energy vector \(D\) becomes
\(D_{e,t}=\gamma_{e,t}\epsilon(Q_{e,t})\).
\end{lemma}
\begin{proof}
For fixed temporal links define \(G_{x,0}=1\) recursively by
\(G_{x,t+1}=G_{x,t}g_{x,t}\). Set
\(V_{e,t}=G_{x,t}U_{e,t}G_{y,t}^{-1}\). Every temporal link
then becomes the identity, and
\[
 V_{e,t}V_{e,t+1}^{-1}=G_{x,t}P_{e,t}G_{x,t}^{-1}.
\]
Spatial plaquette traces are also unchanged. For each fixed \(g\)
this is a product of left/right Haar translations of the spatial
links, hence has unit Jacobian in the normalized Haar integral.
The resulting integrand no longer depends on \(g\), whose integral
is one. The open time interval is essential: no temporal closing
constraint or residual seam has been omitted.

For each edge the triangular change
\((V_0,\ldots,V_T)\mapsto(V_0,Q_0,\ldots,Q_{T-1})\)
preserves product Haar. Its inverse is
\(V_{t+1}=Q_t^{-1}V_t\). Factoring the electric normalizers
gives \(\pi_e\) and \eqref{r81:paths}. This is a representation
of the same integral, not a reset of the magnetic law.
\end{proof}

\subsection{The electric radial variance is uniform in its coupling}

\begin{lemma}[A scalar bound with its Haar factor retained]
\label{r81:radial}
For \(q\sim\nu_\gamma\), \(\gamma>0\), let
\(Y=\gamma\epsilon(q)\). Then
\[
                 \Var(Y)\le\E Y\le\frac32.
 \tag{V81.7}\label{r81:radial-bound}
\]
The constant \(3/2\) cannot be decreased uniformly over \(\gamma\).
\end{lemma}
\begin{proof}
Normalized \(S^3\) Haar measure gives the action variable
\(x=1-\operatorname{Re}q\) density proportional to
\(x^{1/2}(2-x)^{1/2}\) on \((0,2)\). Thus the density \(p\) of
\(Y\), on \((0,L)\), \(L=2\gamma\), is proportional to
\[
 y^{\alpha-1}e^{-y}(L-y)^{1/2},\qquad \alpha=\frac32.
\]
Define \(F(y)=y/[2(L-y)]\), an increasing nonnegative function.
Direct differentiation gives
\((yp(y))'=(\alpha-y-F(y))p(y)\).
The boundary terms \(yp(y)\) and \(y^2p(y)\) vanish at both
endpoints. Also \(Fp\) and \(yFp\) are integrable, including
their inverse-square-root singularity at \(L\). Integrating the
derivative, and then its product with \(y\), yields
\[
 m+\E F=\alpha,\qquad
 \E Y^2+\E(YF)=(\alpha+1)m,\qquad m=\E Y.
\]
Consequently
\[
               \Var(Y)=m-\Cov(Y,F(Y))\le m\le\alpha.
\]
The covariance is nonnegative: with an independent copy \(Y'\),
it is half the expectation of
\((Y-Y')(F(Y)-F(Y'))\ge0\). All terms are integrable.
Finally, dividing the radial density by its harmless factor
\(L^{1/2}\), the remaining multiplier
\((1-y/L)^{1/2}\mathbf1_{y<L}\) tends to one and is bounded by
one. Dominated convergence for the first two gamma-weighted moments
gives the \(\operatorname{Gamma}(3/2,1)\) limit and
\(\Var(Y)\to3/2\).
\end{proof}

In particular, under the product-history reference,
\[
 \Var_{\pi_e}\!\left(\sum_{t<T}s_{e,t}D_{e,t}\right)
                         \le\frac32\sum_{t<T}s_{e,t}^2
 \tag{V81.8}\label{r81:reference-bank}
\]
for every real bank, with no factor proportional to the history length.
This uses the independent \emph{increments}; the link values at
different times are not declared independent.

\subsection{Magnetic path interactions and the finite-block variance inequality}

We record the normalized-tilt fact used in the influence calculation.
If two probabilities have log density-ratio oscillation at most \(D\),
their total-variation distance, normalized as \(\sup_B|\mu(B)-\nu(B)|\),
is at most \(\tanh(D/4)\). Indeed their likelihood ratio has mean one
and range in \([a,b]\), \(b/a\le e^D\). The chord bound for
\(|r-1|\) gives distance at most
\((1-a)(b-1)/(b-a)\); maximizing with fixed \(b/a=k\) gives
\((\sqrt k-1)/(\sqrt k+1)\). This also proves the assertion by
essential bounds and limits. It is the familiar normalized-tilt
inequality already used in the archives.

\begin{lemma}[History influence and conditional density]
\label{r81:influence}
The conditional laws of \(\widehat\mu\), with one whole spatial-link
history as one block, have interdependence coefficients bounded by
\[
 c_{ef}\le\tanh w_{ef}\le w_{ef},\qquad
 w_{ef}=\sum_t\sum_{p\ni e,f}b_{p,t}\quad(e\ne f),\qquad
 \sup_e\sum_{f\ne e}c_{ef}\le3A.
 \tag{V81.9}\label{r81:influence-eq}
\]
Moreover, for every exterior and every square-integrable function of
one history,
\[
 \Var_{\widehat\mu_e(\cdot\mid X_{-e})}(f)
                     \le e^{2a_e}\Var_{\pi_e}(f).
 \tag{V81.10}\label{r81:conditional-var}
\]
Both assertions are independent of all electric couplings.
\end{lemma}
\begin{proof}
The part of \(\mathcal V\) involving edge \(e\) lies in
\([0,2a_e]\), so the conditional density relative to \(\pi_e\)
is at most \(e^{2a_e}\). In the conditional variance use the
trial center \(\pi_e(f)\) and this density bound to prove
\eqref{r81:conditional-var}.

Changing only history \(f\) changes only faces containing both
\(e\) and \(f\). The log-weight difference has oscillation at most
\(4w_{ef}\) as a function of history \(e\), since each original
plaquette energy is in \([0,2]\). The normalized-tilt inequality
gives the first two bounds. Each face containing \(e\) contains
exactly three other edges, so
\(\sum_{f\ne e}w_{ef}=3a_e\). This proves the row estimate.
No global density-ratio bound of size exponential in spatial volume
has been used.
\end{proof}

\begin{lemma}[Finite-block Poincare inequality, with rate convention]
\label{r81:poincare}
Let a probability on a finite product of compact spaces have strictly
positive continuous density relative to a product probability. If its
total-variation interdependence matrix has row sums at most
\(c<1\), then
\[
 (1-c)\Var(f)\le\sum_i\E\Var(f\mid X_{-i})
                  \qquad(f\in L^2).
 \tag{V81.11}\label{r81:poincare-eq}
\]
\end{lemma}
\begin{proof}
Update each block at rate one from its exact conditional distribution.
The generator is \(\mathcal L=\sum_i(E_i-I)\), where \(E_i\) is
conditional expectation given all other blocks. Each \(E_i\) is an
orthogonal projection in \(L^2\) of the Gibbs law. Thus the semigroup
is self-adjoint and its nonnegative generator \(-\mathcal L\) has
Dirichlet form \(\sum_i\E\Var(f\mid X_{-i})\).

Couple two copies by common update clocks and maximal couplings of
their conditional laws. Such couplings are measurable here: use the
common part of their densities and the two residual densities.
By telescoping differing exterior coordinates, the probability of a
disagreement after updating block \(i\) is at most
\(\sum_jc_{ij}\mathbf1_{X_j\ne Y_j}\).
Writing \(u_i(t)=\mathbb P(X_i(t)\ne Y_i(t))\) gives
\[
 D^+\max_i u_i(t)\le-(1-c)\max_i u_i(t),
 \qquad \max_i u_i(t)\le e^{-(1-c)t}.
\]
The differential inequality can equivalently be read from the finite
linear comparison system for the \(u_i\); no differentiability at a
tie is required. Couple an arbitrary initial state to a stationary
copy. If there are \(M\) blocks, the union bound gives
\(\|P_tf\|_\infty\le2M\|f\|_\infty e^{-(1-c)t}\) for bounded
centered \(f\).

The positive spectral measure of \(-\mathcal L\) at this \(f\)
satisfies
\(\langle f,P_tf\rangle\le2M\|f\|_\infty^2e^{-(1-c)t}\).
Any mass on \([0,1-c-\delta]\), \(\delta>0\), would contradict
this as \(t\to\infty\). The spectral support on the centered
subspace is therefore in \([1-c,\infty)\), first for bounded
functions and then by their density in \(L^2\). The spectral theorem
and the displayed Dirichlet form give \eqref{r81:poincare-eq}.
The finite prefactor \(M\) is not part of the final variance constant.
\end{proof}

This is a variance inequality for an auxiliary sampler that replaces
whole histories. Its update time is not Wilson transfer time or a
physical Euclidean time, so its gap is not a YM mass gap.

\subsection{Uniform full-law covariance and the original exterior average}

\begin{theorem}[All signed electric energy banks]
\label{r81:main}
Under \eqref{r81:law} and \eqref{r81:budget}, for every deterministic
real bank \(s=(s_{e,t})\),
\[
 \Var_\mu\!\left(\sum_{e,t<T}s_{e,t}D_{e,t}\right)
                         \le K_A\sum_{e,t<T}s_{e,t}^2.
 \tag{V81.12}\label{r81:covariance}
\]
Equivalently \(0\preceq\Cov_\mu(D)\preceq K_A I\), with constants
uniform in \(N,T,\gamma_{e,t}\) subject only to \eqref{r81:budget}.
\end{theorem}
\begin{proof}
The observable is gauge invariant, so Lemma \ref{r81:gauge} gives
its exact variance under \(\widehat\mu\). In the history variables
it is a sum \(\sum_e f_e(X_e)\), with
\(f_e=\sum_t s_{e,t}\gamma_{e,t}\epsilon(Q_{e,t})\).
Lemmas \ref{r81:influence} and \ref{r81:poincare} give
\[
 \begin{split}
 \Var_{\widehat\mu}\Big(\sum_e f_e\Big)
 &\le\frac1{1-3A}\sum_e
       \E\Var_{\widehat\mu_e(\cdot\mid X_{-e})}(f_e)\\
 &\le\frac{e^{2A}}{1-3A}\sum_e\Var_{\pi_e}(f_e)
 \le K_A\sum_{e,t}s_{e,t}^2.
 \end{split}
\]
The last step is \eqref{r81:reference-bank}. Each use of the
conditional law includes its own normalizer.
\end{proof}

Let \(\mathcal Y=\sigma(U_{e,t}:e\in E,0\le t\le T)\) be the
\emph{original}, unfixed spatial exterior sigma algebra, and define
\[
 H=\E_\mu[D\mid\mathcal Y],\qquad
 C(\mathcal Y)=\Cov_\mu(D\mid\mathcal Y).
 \tag{V81.13}\label{r81:original-projection}
\]

\begin{corollary}[Measured conditional and predictor covariance]
\label{r81:averaged}
Under the same hypotheses,
\[
 \E_\mu C(\mathcal Y)+\Cov_\mu(H)=\Cov_\mu(D)\preceq K_A I.
 \tag{V81.14}\label{r81:tower}
\]
Both summands on the left are positive semidefinite and individually
bounded by \(K_A I\). Conditional on \(\mathcal Y\), the temporal
links in different slabs are independent, so \(C(\mathcal Y)\) is
block diagonal in the slab index. The bound nevertheless controls
the entire cross-slab predictor covariance as well.
\end{corollary}
\begin{proof}
Given all original spatial links, the magnetic action is constant
and the temporal action is a sum of terms, one for each slab,
in disjoint families \(g_{x,t}\). The conditional law therefore
factorizes over slabs, and each factor is exactly the compact Wilson
corridor integral with its two actual boundary slices. The matrix
conditional-variance identity, applied before any gauge change,
gives the equality in \eqref{r81:tower}. Positivity and
Theorem \ref{r81:main} give the remaining claims.
\end{proof}

This does \emph{not} identify \(\mathcal Y\) with the sigma algebra
generated by the gauge-fixed histories: those histories depend also on
the temporal links being integrated out. Temporal gauge was used to
bound the full gauge-invariant variance, not to replace the original
conditional projection. Nor does \eqref{r81:tower} bound
\(C(\mathcal Y)\) pointwise. V80's large conditional covariance is
consistent with this actual-law average.

\subsection{A centered signed-source domain without a volume cost}

Define
\(\Psi(s)=\log\E_\mu\exp(s\cdot D)\). At every finite regulator
all components of \(D\) are bounded, so real source derivatives exist
of all orders. The following estimate is uniform on its stated domain.

\begin{theorem}[Real source Hessian and gamma-type centered budget]
\label{r81:source}
For all real banks with \(s_{e,t}<1\) at every coordinate,
\[
 \nabla^2\Psi(s)\preceq
        K_A\operatorname{diag}\big((1-s_{e,t})^{-2}\big),
 \tag{V81.15}\label{r81:source-hessian}
\]
and, with \(h(u)=-\log(1-u)-u\),
\[
 \log\E_\mu e^{s\cdot(D-\E D)}
                    \le K_A\sum_{e,t}h(s_{e,t}).
 \tag{V81.16}\label{r81:centered-source}
\]
In particular, if \(\|s\|_\infty\le\delta<1\), the right side
is at most \(K_A\|s\|_2^2/[2(1-\delta)]\).
The same centered exponential bound holds with \(D\) replaced by
its original exterior predictor \(H\).
\end{theorem}
\begin{proof}
The source tilt changes exactly the electric activities to
\(\gamma'_{e,t}=(1-s_{e,t})\gamma_{e,t}>0\). Magnetic terms and
the budget \(A\) are unchanged. Apply the proof of
Theorem \ref{r81:main} in this tilted law. The observed variable
is still \(D_{e,t}\), not its newly rescaled version, so its
one-increment variance bound is
\(\tfrac32(\gamma_{e,t}/\gamma'_{e,t})^2
=\tfrac32(1-s_{e,t})^{-2}\). The proof for every signed test bank
gives \eqref{r81:source-hessian}, because the source Hessian is
the tilted covariance of these original variables.

Integrate this quadratic-form bound along \(t\mapsto ts\),
using Taylor's integral remainder. Coordinatewise,
\[
 \int_0^1\frac{(1-t)u^2}{(1-tu)^2}\,dt
                          =-\log(1-u)-u=h(u),\qquad u<1.
\]
This proves \eqref{r81:centered-source}. For \(|u|<1\), its
power series gives \(h(u)\le u^2/[2(1-|u|)]\), yielding the
claimed \(\ell^2\) estimate. Finally conditional Jensen gives
\(\E e^{s\cdot(H-\E H)}\le\E e^{s\cdot(D-\E D)}\).
No independence is assigned to the predictor coordinates.
\end{proof}

For example, for a real bank \(v\), let
\(F=v\cdot(D-\E D)\), \(V=K_A\|v\|_2^2\), and
\(c=\|v\|_\infty\). The preceding bound gives
\[
 \log\E e^{\lambda F}\le\frac{V\lambda^2}{2(1-c|\lambda|)}
             \quad(c|\lambda|<1),\qquad
 \mu\{|F|\ge\sqrt{2Vx}+cx\}\le2e^{-x}\quad(x>0).
 \tag{V81.17}\label{r81:tail}
\]
For the upper tail choose
\(\lambda=\sqrt{2x/V}/(1+c\sqrt{2x/V})\) when \(V>0\);
the lower tail uses \(-F\). The zero bank is immediate. The same
tail holds for the centered predictor. These are bounds at the actual
means with arbitrary support, not V47's noncentered positive
\(\ell^1\) budget. A complex zero-free domain is not asserted.

\subsection{What the bound pays, and what remains upstream}

The new paid statement is an original open-strip Wilson
\emph{exterior-averaged} covariance estimate, including its predictor
return, uniform in arbitrarily large electric couplings. The signed
source Hessian is controlled without expanding the magnetic
normalizer into an extensive absolute error. The scalar factor
\(e^{2A}\) pays one history's conditional distortion, not the whole
lattice density ratio. In particular the result does not condition on
an all-good boundary event or delete V80's frustrated sectors.

Its remaining restrictions must be kept:
\begin{enumerate}
\item For homogeneous isotropic coupling \(\gamma\),
\(A=4(T+1)\gamma\). Thus \eqref{r81:budget} excludes the large
isotropic couplings relevant to the desired four-dimensional
continuum. The theorem establishes no entry into that regime.
\item If a chosen time discretization has magnetic coefficients
\(b_t\) of order its time step, then \(B\) is their time-integrated
budget, not the number of slices alone. Refining time while keeping
that budget small is allowed. This algebraic observation does not
construct a physical continuum trajectory or a spatial-refinement
uniform time window.
\item The endpoints are integrated Haar in the stated finite strip.
Adding the physical vacuum's Perron weights, prescribing endpoints,
or closing time changes the path factors or their interactions.
It needs its own conditional-moment and influence proof. No vacuum
or infinite-time conclusion follows by renaming these endpoints.
\item Averaging over an internal slice before multiplying adjacent
strip kernels does not make it independent. Iterating the theorem
as if every strip had fresh Haar endpoints would change the law.
The needed extension is boundary-weight stability or an exact
sewing estimate retaining the intermediate-slice dependence.
\item The unaveraged coherent estimate \eqref{r79:unpaid}, and the
original-law connected Ward return left in V61, remain separate
unpaid questions. This result supplies no bound on their exceptional
fixed exterior merely by the conditional-variance identity.
\end{enumerate}

The diagnostic checks the four-face/twelve-neighbour incidence on
\(N=3,4,6\), and the full noncommuting temporal-gauge identity on
a three-step \(N=3\) strip, including every spatial plaquette and
the triangular inverse. Independent high-precision Bessel moments
and radial integration check \eqref{r81:radial-bound} and its
integration-by-parts identities at sampled couplings. A fully
enumerated center-valued square checks the original conditional
variance decomposition against the gauge-fixed full law and the
finite history heat-bath spectrum. That discrete analogue is
\emph{not} Haar \(\SU(2)\), and is not used to prove its constants
or extrapolate a mass gap. The analytic proofs above supply the
stated uniformities. V80's mathematical checker is also replayed.

All prior mathematical content is preserved. Diagnostics, compilation
and render checks are outside \texttt{latex\_notes}; V81 is the sole
active main version. The programme has gained a rigorously measured
averaging regime, not the full YM mass gap. The physical-time,
vacuum-boundary, spatial-refinement and interacting-continuum
obligations remain open.
% END CONSOLIDATED V81 APPEND

% BEGIN CONSOLIDATED V82 APPEND
\clearpage
\section{V82: Sewn vacuum energy sources and temporal defect patterns}
\label{r82:section}

\subsection{Upstream choice, nonduplication, and the actual advance}

V81 obtains a centered signed-source bound on an open Haar-ended strip,
but only under a small integrated magnetic budget. Instead of assuming
that its endpoints may be replaced by vacuum weights, we return to the
native transfer operator. An archive check is decisive here: f15 V88
already proves the one-step kinetic-source order and relative Chernoff
bound, and f15 V93 already differentiates it to bound the vacuum mean
increment. Those are inherited results, not new V82 claims. Likewise the
connected-set counting argument below is elementary and was already
used for different bad patches in f15 V20.

What is added is an exact \emph{arbitrary-time vacuum sewing} theorem
for these sources, and a relative \emph{path-operator} bound for specified
temporal defect patterns. It applies at every positive electric coupling
and every finite magnetic coupling, without a small time-integrated
magnetic budget or a denominator involving endpoint overlaps. Its
corollaries give joint upper tails and rootwise exponential cluster
bounds for large temporal plaquette energies in the actual vacuum law.
It does not supply V81's centered signed covariance bound in this larger
regime. The distinction is quantified, not hidden in a change of law.

The current companion audit remains V80 (next V85), and the regular
literature sweep remains V80 (next V90). A targeted check of
\href{https://arxiv.org/abs/1810.09756}{N. Garofalo, \emph{Two classical
properties of the Bessel quotient} (2018)}, and the
\href{https://dlmf.nist.gov/10.32.E2}{DLMF integral representation
10.32.2}, confirms the classical monotonicity used in the archived
kinetic order. We include its short proof only to make the new sewing
argument self-contained. No external physical gap conclusion is imported.

\subsection{Exact transfer, normalization, and the inherited one-step input}

Fix a finite cubic spatial torus of side at least three, with positive
edges $E$, and put $\mathcal X=\SU(2)^E$ with normalized product Haar
measure $dm$. Let $\gamma_e>0$ be time-independent electric activities.
For $q\in\SU(2)$ set
\[
 \epsilon(q)=1-\tfrac12\operatorname{Re}\Tr q,\qquad
 c_0(x)=\int e^{-x\epsilon(q)}dH(q)=2e^{-x}I_1(x)/x.
\]
Let $C_x$ be convolution by $e^{-x\epsilon}/c_0(x)$, and
$C_{\boldsymbol\gamma}=\bigotimes_e C_{\gamma_e}$.
For any fixed finite nonnegative magnetic activities $b_p$, retain
\begin{equation}
 V(U)=\sum_{p\ {
m spatial}}b_p\epsilon(U_p),\quad
 b(U)=e^{-V(U)/2},\quad
 T=M_b C_{\boldsymbol\gamma}M_b,\quad P=T/\Lambda.
 \tag{V82.1}\label{r82:transfer}
\end{equation}
Here $T\Omega=\Lambda\Omega$, $\Omega>0$, $\|\Omega\|_2=1$.
The strictly positive continuous kernel and compactness give the simple
positive top eigenfunction. The operator is positive semidefinite by
its character multipliers. Thus $\Lambda=\|T\|>0$ and $\|P\|=1$.
These are the inherited finite-regulator transfer facts, not a
regulator-uniform lower gap. Throughout this section $P$ is the
normalized selfadjoint transfer on Haar space, not its Markov conjugate.

All these operators commute with the simultaneous vertex gauge action;
uniqueness makes $\Omega$ gauge invariant. The physical transfer is the
restriction to invariant functions. A temporal-gauge representation of
an $n$-step vacuum path therefore has probability density
\begin{equation}
 d\mathbb P_n(U_0,\ldots,U_n)=
 \Omega(U_0)\prod_{t=0}^{n-1}P(U_t,U_{t+1})\Omega(U_n)
 \prod_{t=0}^n dm(U_t).
 \tag{V82.2}\label{r82:path}
\end{equation}
It has mass one, consistent endpoint marginals, and stationary extensions
to all integer times, since $P\Omega=\Omega$. Every magnetic interior
slice receives its full weight; the two end slices receive the displayed
half weights together with the actual Perron factors.

For the original unfixed augmented path, integrating the temporal links
and changing the later spatial slices by successive vertex gauge
transformations gives \eqref{r82:path} for gauge-invariant observables.
Haar measure is preserved at each step, the magnetic action is unchanged,
and gauge invariance of $\Omega$ removes the terminal gauge rotation.
In particular the original temporal plaquette energies become exactly
\begin{equation}
 D_{e,t}=\gamma_e\epsilon(U_{e,t}U_{e,t+1}^{-1}).
 \tag{V82.3}\label{r82:energy}
\end{equation}
This is a full-path change of variables, not an identification of the
original exterior sigma-algebra with gauge-fixed slices.

For $0\le s_e<1$, let $T_s$ be $T$ with kernel multiplied by
$\exp(\sum_e s_e\gamma_e\epsilon(U_eV_e^{-1}))$. Define
\begin{equation}
 R(s)=\prod_e\frac{c_0((1-s_e)\gamma_e)}{c_0(\gamma_e)}.
 \qquad 0\preceq T_s\preceq R(s)T,
 \qquad R(s)\le\prod_e(1-s_e)^{-3/2}.
 \tag{V82.4}\label{r82:one-step}
\end{equation}
The order is operator order, not pointwise kernel order. This is f15
V88's inherited estimate, with possibly different electric activities.

The normalized degree-$j$ convolution multiplier is
$I_{j+1}(x)/I_1(x)$. For completeness, write
\[
 Z_\nu(x)=\int_{-1}^1e^{xu}(1-u^2)^{\nu-1/2}du.
\]
The Bessel integral and recurrence give
$I_{\nu+1}/I_\nu=Z_\nu'/Z_\nu$ for $\nu\ge1$; differentiation
is the strictly positive variance of $u$ under that tilted density.
The quotient is positive by symmetry. Products of these quotients
therefore show that $I_{j+1}/I_1$ increases with $x$. Hence
$0\preceq C_{\alpha x}\preceq C_x$ for $0<\alpha\le1$.
Tensoring these commuting positive inequalities and taking the
congruence by $M_b$ proves the operator part of \eqref{r82:one-step}.
Finally
\[
 x^{3/2}c_0(x)=\frac2\pi\int_0^{2x}
 e^{-y}y^{1/2}(2-y/x)^{1/2}\,dy
\]
is increasing in $x$, proving its last bound. Nothing in this argument
requires small $V$, a bound on $\Omega_{\max}/\Omega_{\min}$, or an
approximation to $\Lambda$.

\subsection{Sewing sources at any finite set of times}

\begin{theorem}[Native vacuum positive sources at arbitrary times]
\label{r82:sources}
For every finite nonnegative spacetime source bank $0\le s_{e,t}<1$,
in the stationary vacuum law just specified,
\begin{equation}
 \mathbb E_{\rm vac}\exp\Big(\sum_{e,t}s_{e,t}D_{e,t}\Big)
 \le\prod_t R(s_{\cdot,t})
 \le\prod_{e,t}(1-s_{e,t})^{-3/2}.
 \tag{V82.5}\label{r82:mgf}
\end{equation}
The constants are independent of the time span, spatial volume, and
all magnetic and positive electric activities. The time-independent
transfer and its own vacuum are held fixed while the insertions vary.
\end{theorem}
\begin{proof}
Enclose the support in steps $0,\ldots,n-1$, allowing zero sources in
gaps. Exact integration of \eqref{r82:path} gives the matrix element
\[
 \Lambda^{-n}\langle\Omega,T_{s_{\cdot,0}}T_{s_{\cdot,1}}
 \cdots T_{s_{\cdot,n-1}}\Omega\rangle.
\]
It is nonnegative because it is the displayed path integral. By
\eqref{r82:one-step}, each factor has norm at most
$R(s_{\cdot,t})\Lambda$. Apply Cauchy--Schwarz and norm
submultiplicativity. No order between noncommuting products is claimed;
one cannot simply multiply their Loewner inequalities. Zero-source
steps cost exactly $\|P\|=1$, even across arbitrarily long gaps.
Stationarity shifts the chosen support to any integer times.
\end{proof}

This proof retains every sewn spatial slice. It neither replaces the
vacuum by a product endpoint law nor assumes temporal independence.
It also does not differentiate a source-dependent Perron vector.

Let $m_0(x)=\mathbb E_{\nu_x}[x\epsilon]\le3/2$ be V81's
single-increment mean, and $h(s)=-\log(1-s)-s$. Its proved variance
bound $\operatorname{Var}_{\nu_x}(x\epsilon)\le3/2$ gives
\begin{equation}
 \log\frac{c_0((1-s)x)}{c_0(x)}-s m_0(x)
 \le\tfrac32h(s),\qquad 0\le s<1.
 \tag{V82.6}\label{r82:scalar-centered}
\end{equation}
Indeed the second derivative of the logarithmic ratio is
$\operatorname{Var}_{\nu_{(1-s)x}}((1-s)x\epsilon)/(1-s)^2$;
integrate its upper bound twice with the value and derivative at zero.
Consequently
\begin{equation}
 \log\mathbb E_{\rm vac}e^{\sum s_{e,t}(D_{e,t}-m_0(\gamma_e))}
 \le\tfrac32\sum_{e,t}h(s_{e,t}).
 \tag{V82.7}\label{r82:free-center}
\end{equation}
For nonnegative finitely supported $v$, $c=\|v\|_\infty>0$, this implies
\begin{equation}
 \mathbb P_{\rm vac}\left\{
 \sum v_{e,t}(D_{e,t}-m_0(\gamma_e))
 \ge\sqrt{3\|v\|_2^2x}+cx\right\}\le e^{-x},\quad x>0.
 \tag{V82.8}\label{r82:upper-tail}
\end{equation}
To check constants, $h(u)\le u^2/[2(1-u)]$ on $[0,1)$;
the log moment generating function is at most
$V\lambda^2/[2(1-c\lambda)]$ with $V=3\|v\|_2^2/2$.
In Chernoff's inequality take
$\lambda=\sqrt{2x/V}/(1+c\sqrt{2x/V})$.
The zero bank is trivial. These are upper tails about a known
\emph{free-increment center}, not about the actual vacuum mean.

Conditional Jensen also passes \eqref{r82:mgf} and
\eqref{r82:free-center} to $H_{e,t}=\mathbb E[D_{e,t}\mid\mathcal Y]$
for any one common sigma-algebra $\mathcal Y$ in the original augmented
law, including all original spatial slices. This claim is unconditional
after averaging $\mathcal Y$; it gives no uniform conditional bound.

\subsection{Operator bounds for a prescribed temporal defect pattern}

For $\nu>0$ and $L\ge0$ set
\begin{equation}
 q_\nu(L)=
 \begin{cases}1,&L\le\nu,\\
 e^{\nu-L}(L/\nu)^\nu,&L>\nu.
 \end{cases}
 \tag{V82.9}\label{r82:q}
\end{equation}
At each marked time $t$, choose pairwise disjoint nonempty edge sets
$J_{a,t}$ and thresholds $L_{a,t}\ge0$. Put
\[
 \mathcal B_{a,t}=
 \left\{\sum_{e\in J_{a,t}}D_{e,t}\ge L_{a,t}\right\},
 \qquad \nu_{a,t}=3|J_{a,t}|/2.
\]
Sets at different times may use exactly the same spatial links.

\begin{theorem}[Relative path-mask bound]
\label{r82:mask}
Let $0\le a(U_0,\ldots,U_n)\le1$ be measurable and vanish outside
$\bigcap_{a,t}\mathcal B_{a,t}$. Integrate the internal slices of the
kernel $a\prod_{t=0}^{n-1}P(U_t,U_{t+1})$ to obtain an endpoint
operator $K_a$. Then
\begin{equation}
 \|K_a\|_{L^2(dm)\to L^2(dm)}
 \le\prod_{a,t}q_{\nu_{a,t}}(L_{a,t}).
 \qquad
 \mathbb P_{\rm vac}\Big(\bigcap_{a,t}\mathcal B_{a,t}\Big)
 \le\prod_{a,t}q_{\nu_{a,t}}(L_{a,t}).
 \tag{V82.10}\label{r82:pattern}
\end{equation}
No selfadjointness or positive-semidefiniteness of $K_a$ is assumed.
For gauge-invariant masks the bound restricts to the physical space.
\end{theorem}
\begin{proof}
At a fixed step let $A_t$ have kernel $P$ times the product of the
relevant event indicators. For numbers $0\le\theta_{a,t}<1$,
pointwise kernel domination gives, for every $f$,
\[
 |A_tf|\le
 e^{-\sum_a\theta_{a,t}L_{a,t}}\Lambda^{-1}T_{s_t}|f|,
 \quad (s_t)_e=\theta_{a,t}\ (e\in J_{a,t}),
\]
with zero elsewhere. Disjointness makes this source well-defined.
Use \eqref{r82:one-step} to bound the norm by
$\prod_a e^{-\theta_{a,t}L_{a,t}}(1-\theta_{a,t})^{-\nu_{a,t}}$.
Separate scalar minimizations give $\prod_a q_{\nu_{a,t}}(L_{a,t})$:
the minimizer is zero for $L\le\nu$ and $1-\nu/L$ otherwise.
Unmarked steps are $P$ and cost one. Since $a$ is bounded by the
product of these step indicators, the endpoint kernel of $K_a$ is
pointwise dominated by that of $A_0\cdots A_{n-1}$, all kernels being
nonnegative. Thus $|K_af|\le A_0\cdots A_{n-1}|f|$, proving the
operator estimate by submultiplicativity. Finally sandwich the product
by the unit vector $\Omega$, as in \eqref{r82:path}.
\end{proof}

Thus a specified defect pattern is small even as an operator inserted
between arbitrary normalized endpoint vectors. It is not merely small
under the vacuum distribution. This does not estimate a resolvent or
the cost of summing all possible defect patterns.

\subsection{A native large-increment cluster estimate, without independence}

Call the temporal plaquette index $z=(x,j,t)$ bad if $D_{j,x,t}\ge L$,
and write $I_z$ for its indicator and $q=q_{3/2}(L)$. Distinct indices
have disjoint single-edge supports whenever their times coincide.
The preceding theorem therefore gives, for every finite set $A$,
\begin{equation}
 \mathbb E_{\rm vac}\prod_{z\in A}I_z\le q^{|A|},\qquad
 \mathbb E_{\rm vac}\prod_{z\in F}(1+u_zI_z)
 \le\prod_{z\in F}(1+q u_z),\quad u_z\ge0.
 \tag{V82.11}\label{r82:joint}
\end{equation}
The second assertion follows by the finite subset expansion. In
particular the number of bad indices in $F$ obeys the resulting
binomial Chernoff bound. Neither assertion is full stochastic
domination by independent indicators, nor a bound after arbitrary
conditioning; f15 V20 already records that distinction.

Join indices with the same $j$ and nearest-neighbor base points in
$(\mathbb Z/N\mathbb Z)^3\times\mathbb Z$, and join distinct orientations
at the same base point. The degree is at most $\Delta=10$.
Let $\mathcal C(z)$ be the bad component containing $z$, empty when
$z$ is good. For every $k\ge1$,
\begin{equation}
 \mathbb P_{\rm vac}\{|\mathcal C(z)|\ge k\}
 \le q(\Delta^2q)^{k-1}.
 \tag{V82.12}\label{r82:clusters}
\end{equation}
Indeed an occupied component of size at least $k$ contains a connected
$k$-set rooted at $z$. Order neighbors deterministically. The canonical
depth-first walk of a chosen spanning tree has $2(k-1)$ steps and
recovers its visited set. Hence there are at most
$\Delta^{2(k-1)}$ such sets. Apply \eqref{r82:joint} and a union bound.
For $L=16$, $100q_{3/2}(16)<0.0018<1$, so the estimate is genuinely
exponential, uniformly in volume and all couplings. At small coupling
some bad events are simply impossible because $D\le2\gamma_e$.
There is no volume factor for a specified root, and no assertion that
all roots are good simultaneously. A cluster's graph distance here is
not the correlation length of an arbitrary physical observable.

\subsection{Exactly what still blocks a centered or physical conclusion}

Let $\mu_{e,t}=\mathbb E_{\rm vac}D_{e,t}$.
Right differentiation of \eqref{r82:mgf} at zero gives the already known
$\mu_{e,t}\le m_0(\gamma_e)$. Define the nonnegative deficit
$\delta_{e,t}=m_0(\gamma_e)-\mu_{e,t}$. The actual conclusion for
centered positive sources is only
\begin{equation}
 \log\mathbb E_{\rm vac}e^{\sum s_{e,t}(D_{e,t}-\mu_{e,t})}
 \le\sum s_{e,t}\delta_{e,t}+\tfrac32\sum h(s_{e,t}).
 \tag{V82.13}\label{r82:deficit}
\end{equation}
An upper bound with this nonzero linear slack cannot be differentiated
twice at zero to produce a covariance bound. For example, for one common
Bernoulli variable $B$ of parameter $1/2$, let $X_i=B$ for $1\le i\le m$.
For every $s_i\in[0,1)$,
\[
 \mathbb E e^{\sum_i s_iX_i}\le e^{\sum_i s_i}
 \le\prod_i(1-s_i)^{-3/2},\qquad
 \|\operatorname{Cov}(X)\|=m/4.
\]
This is a counterexample to an inference from the moment bound alone,
not a Yang--Mills counterexample. Nor does the proof permit negative
sources: the normalized kinetic comparison reverses when an activity
is increased, and gives no corresponding upper form bound.

The defects controlled here concern \emph{temporal increments}, not
all spatial magnetic roughness or low-staple sectors. Moreover even
the all-tame region retains a concrete global obstruction. On V80's
flat relative-center-twist boundary, every ground configuration has
\[
 D_{e}=\gamma(1-\cos(\pi/N))\quad(e\parallel1),\qquad
 D_e=0\quad(e\parallel2,3).
\]
For each fixed $L>0$ these are all below $L$ when
$N>\pi\sqrt{\gamma/(2L)}$, using $1-\cos u\le u^2/2$.
The endpoints still have different winding holonomy. This exhibits
an all-tame configuration, not a lower bound on its vacuum probability;
it explains why the local cluster estimate alone cannot eliminate a
distributed soft mode or repair V80's conditional obstruction.

The next useful estimate must therefore concern \emph{connected signed
sources inside the retained tame law}, with its original normalization
and coherent long-range returns still present, or the V61 Ward return
in the coherent compact model. Merely repeating rare-event estimates
cannot pay that estimate. No physical gap, all-volume covariance bound,
or interacting four-dimensional continuum construction is proved here.

The diagnostic checks classical SU(2) multiplier/source inequalities at
sampled parameters and the sewing and mask algebra in a fully specified
finite center-valued transfer model, including noncommuting insertions
and nontrivial magnetic weights. The finite model is not SU(2) Haar
quadrature and proves none of the stated continuum uniformities; the
proofs above do. It also checks the tail constants, rooted-set encoding,
and the covariance-inference counterexample. V81's diagnostic is rerun.
All previous mathematical content is preserved. Only V82 remains as
the active main source, and build products stay outside
\texttt{latex\_notes}.
% END CONSOLIDATED V82 APPEND

% BEGIN CONSOLIDATED V83 APPEND
\clearpage
\section{V83: Native increment entropy and centered signed banks}
\label{r83:section}

\subsection{The estimate being paid, and the parameter that must be small}

V82 gives all-time positive-source bounds in the actual vacuum, but its
free-increment center differs from the true mean. This section pays a
volume-uniform bound for that deficit and a stronger finite-bank
relative-entropy estimate. The latter gives a genuinely centered,
signed-source covariance comparison, with its bank-size dependence
explicit. No restriction on the time separation of the selected
increments is required. The useful small parameter is the ratio of
magnetic to electric coupling, not a small integrated magnetic budget.
It is an anisotropic result, not a solution of the isotropic
weak-coupling continuum problem.

The native transfer, its vacuum, V81's scalar moment bounds, and V82's
exact path law are retained. Archive f16.2 V51 already compares Wilson
and heat kinetic deficits uniformly over representations; we do not
count heat comparison as new. The elementary weaker lower bound below
is included because it supplies an explicit variational constant at
every positive coupling. Generic entropy chain rules, Pinsker's
inequality, and fractional-packing inequalities are classical. For
the latter, see M. Madiman and P. Tetali,
\href{https://arxiv.org/pdf/0901.0044}{\emph{Information Inequalities for
Joint Distributions, with Interpretations and Applications}}, Section IX.
We prove the needed product-reference case directly. The new work is
the native variational cost, its exact path-entropy accounting, the
space-time averaging without a volume loss, and the resulting actual
energy-bank estimate. The regular companion and literature checkpoints
remain V85 and V90.

\subsection{Homogeneous native transfer and an explicit intensive cost}

Use V82's cubic torus with $m=3N^3$ positive spatial edges, $N\ge3$.
All electric activities are $\gamma>0$ and all spatial magnetic
activities are $\beta\ge0$. Write
\begin{equation}
 V(U)=\beta\sum_{p\ {\rm spatial}}\epsilon(U_p),\quad
 B(U)=e^{-V(U)/2},\quad
 T=M_B C_\gamma^{\otimes m}M_B,\quad T\Omega=\Lambda\Omega,
 \quad \|\Omega\|_2=1.
 \tag{V83.1}\label{r83:transfer}
\end{equation}
As before $\Omega$ is positive and gauge invariant, and
$\Lambda=\|T\|$ also on the full Haar space. Define
\begin{equation}
 F=-\log\Lambda,\qquad \rho=\beta/\gamma,
 \qquad f(\rho)=16\rho+4\sqrt{6\rho}.
 \tag{V83.2}\label{r83:f}
\end{equation}
All logarithms and relative entropies below use natural logarithms.

\begin{lemma}[A full-spectrum kinetic lower comparison]
\label{r83:kinetic}
For the positive round-sphere Laplacian $\mathcal L$ on $\SU(2)=S^3$,
whose degree-$j$ eigenvalue is $j(j+2)$,
\begin{equation}
 C_\gamma\succeq e^{-4\mathcal L/(3\gamma)},\qquad
 C_\gamma^{\otimes m}\succeq
 e^{-4\sum_e\mathcal L_e/(3\gamma)}.
 \tag{V83.3}\label{r83:heat-lower}
\end{equation}
\end{lemma}
\begin{proof}
Put $r_\nu(x)=I_{\nu+1}(x)/I_\nu(x)$. For integers $\nu\ge1$,
V82's tilted integral gives $0<r_\nu<1$. The Bessel recurrence gives
\[
 r_\nu(x)=\frac{x}{2\nu+2+x r_{\nu+1}(x)}
 \ge\frac{x}{x+2\nu+2}.
\]
Thus $-\log r_\nu(x)\le(2\nu+2)/x$. Summing for $\nu=1,\ldots,j$
yields
\[
 -\log\frac{I_{j+1}(x)}{I_1(x)}
 \le\frac{j(j+3)}x\le\frac{4j(j+2)}{3x},\qquad j\ge1.
\]
The trivial representation is exact. Peter--Weyl diagonalization and
products of the scalar inequalities prove both assertions, without a
representation cutoff or an accumulated additive link error.
\end{proof}

\begin{theorem}[Native intensive free-energy bound]
\label{r83:pressure}
At every finite spatial volume and the stated couplings,
\begin{equation}
 0\le F/m\le f(\beta/\gamma).
 \tag{V83.4}\label{r83:pressure-eq}
\end{equation}
\end{theorem}
\begin{proof}
$B\le1$ and $C_\gamma$ is a contraction, so $\Lambda\le1$.
For $\beta>0$ choose $\alpha>16\beta$ and the unit product vector
\[
 \phi_\alpha(U)=\prod_e
 c_0(\alpha)^{-1/2}e^{-\alpha\epsilon(U_e)/2},\qquad
 \psi=B^{-1}\phi_\alpha/\|B^{-1}\phi_\alpha\|_2.
\]
This trial vector need not be gauge invariant: the full top eigenvalue
is the same physical $\Lambda$, so its Rayleigh quotient remains a
valid lower bound. No trial wavefunction is substituted for $\Omega$.
The exact quotient is
\[
 \langle\psi,T\psi\rangle
 =\frac{\langle\phi_\alpha,C_\gamma^{\otimes m}\phi_\alpha\rangle}
 {\mathbb E_{\nu_\alpha^{\otimes m}}e^{V}}.
\]
The chord triangle inequality for four group factors and
Cauchy--Schwarz give $\epsilon(U_p)\le4\sum_{e\in p}\epsilon(U_e)$.
Every edge belongs to four spatial plaquettes; therefore
$V\le16\beta\sum_e\epsilon(U_e)$. The scalar ratio in V82 yields
\[
 \mathbb E_{\nu_\alpha^{\otimes m}}e^V
 \le\left(\frac{c_0(\alpha-16\beta)}{c_0(\alpha)}\right)^m
 \le\left(\frac{\alpha}{\alpha-16\beta}\right)^{3m/2}.
\]
On one sphere $|\nabla\epsilon|^2=\sin^2\theta\le2\epsilon$.
By V81's mean bound,
\[
 \left\langle\phi_\alpha,\sum_e\mathcal L_e\phi_\alpha\right\rangle
 \le\frac{m\alpha^2}{2}\mathbb E_{\nu_\alpha}\epsilon
 \le\frac{3m\alpha}{4}.
\]
The spectral Jensen inequality $\langle\phi,e^{-tH}\phi\rangle
\ge e^{-t\langle\phi,H\phi\rangle}$ for a unit form-domain vector,
together with \eqref{r83:heat-lower}, now proves
\[
 \frac Fm\le\frac\alpha\gamma+
 \frac32\log\frac\alpha{\alpha-16\beta}.
\]
Choose $\alpha=16\beta+\sqrt{24\beta\gamma}$ and use
$\log(1+u)\le u$. The bound becomes
$16\beta/\gamma+2\sqrt{24\beta/\gamma}=f(\rho)$.
When $\beta=0$, $\Omega=1$, $\Lambda=1$, and the conclusion is exact.
\end{proof}

\subsection{The vacuum mean deficit is paid, not differentiated from an upper bound}

The single free-increment energy mean is $m_0(\gamma)\le3/2$.
Translation and coordinate-permutation symmetry make
$\mu=\mathbb E_{\rm vac}D_{e,t}$ independent of $e,t$. Put
$\delta=m_0(\gamma)-\mu\ge0$, as proved in V82.

\begin{lemma}[A Bessel elasticity estimate]
\label{r83:elasticity}
For $k_j(x)=I_{j+1}(x)/I_1(x)$,
\begin{equation}
 0\le x\,\partial_x\log k_j(x)\le-2\log k_j(x),
 \qquad x>0,\quad j\ge0.
 \tag{V83.5}\label{r83:elasticity-eq}
\end{equation}
\end{lemma}
\begin{proof}
It suffices to prove $xr_\nu'/r_\nu\le2(1-r_\nu)$ for integers
$\nu\ge1$, since $1-r\le-\log r$ and $k_j=\prod_{\nu=1}^j r_\nu$.
Set $a=2\nu+1\ge3$ and $r=r_\nu(x)$. Its Riccati equation is
$r'=1-r^2-ar/x$. Define
\[
 G(x)=xr'-2r(1-r)=x(1-r^2)-(a+2)r+2r^2.
\]
The convergent small-$x$ Bessel series gives
$r=x/(a+1)+O(x^3)$, so $G<0$ initially. At any zero of $G$,
$x=r(a+2-2r)/(1-r^2)$ and direct differentiation gives
\[
 G'=\frac{1-r}{a+2-2r}
 \bigl[-(a+2)+(a+8)r-(2a+6)r^2\bigr]<0.
\]
Indeed at $a=3$ the bracket is $-5+11r-12r^2<0$, and its
derivative with respect to $a$ is $-1+r-2r^2<0$.
Since $0<r<1$, a first crossing of $G$ from below is impossible.
This proves the upper bound. The lower bound follows from the
inherited strict monotonicity of each $r_\nu$. The case $j=0$ is zero.
\end{proof}

\begin{proposition}[Explicit true-mean centering cost]
\label{r83:mean}
In the unchanged native vacuum,
\begin{equation}
 0\le\delta\le\min\{3/2,\,2F/m\}
 \le\min\{3/2,\,2f(\rho)\}.
 \tag{V83.6}\label{r83:mean-eq}
\end{equation}
\end{proposition}
\begin{proof}
Put $r=B\Omega$. Diagonalize the product convolution, with eigenvalues
$c_\xi=\prod_e k_{j_e}(\gamma)>0$, and define probabilities
\[
 p_\xi=|r_\xi|^2/\|r\|_2^2,\quad
 Z=\Lambda/\|r\|_2^2=\sum_\xi p_\xi c_\xi,\quad
 w_\xi=p_\xi c_\xi/Z,\quad \ell_\xi=-\log c_\xi.
\]
Multiplicities are included in $\xi$; zero coefficients are omitted.
Differentiating only the kinetic kernel, with $r$ held fixed, gives
\[
 m\delta=\Lambda^{-1}\langle r,\gamma\partial_\gamma
 (C_\gamma^{\otimes m})r\rangle
 =\sum_\xi w_\xi\gamma\partial_\gamma\log c_\xi
 \le2\sum_\xi w_\xi\ell_\xi.
\]
The first equality follows directly from its kernel logarithmic
derivative $m m_0(\gamma)-\sum_eD_e$; no derivative of a bounding
inequality or source-dependent vacuum is used. Relative entropy gives
\[
 0\le D(w\|p)=-\sum_\xi w_\xi\ell_\xi-\log Z,
 \qquad \sum_\xi w_\xi\ell_\xi\le\log(\|r\|_2^2/\Lambda)\le F.
\]
The sums are finite in value since $u|\log u|$ is bounded on $[0,1]$;
the derivative sum is then controlled by \eqref{r83:elasticity-eq}.
Also $\|r\|_2\le1$. Divide by $m$, use symmetry, and retain
$\delta\le m_0\le3/2$.
\end{proof}

For example V82.13 now has the explicit centering error
$2f(\rho)\|s\|_1$ for a positive source bank. This still has linear
slack at the origin. A centered covariance theorem requires the
additional entropy argument that follows.

\subsection{Exact path relative entropy, then spatial and temporal dilution}

Let $\mathbb P_n$ be V82's native $n$-step vacuum path law. Let
$\mathbb R_n$ start from the \emph{same} slice law $\Omega^2dm$ and
then use the product electric transition $C_\gamma^{\otimes m}$ at
every step, without magnetic factors. This is a comparison measure
only; all desired expectations remain under $\mathbb P_n$. Its
increments $Q_{e,t}=U_{e,t}U_{e,t+1}^{-1}$ are independent with law
$\nu_\gamma$, independently of its initial slice. Define
\begin{equation}
 h_{\rm kin}=\frac{F-\mathbb E_{\Omega^2}V}{m}.
 \tag{V83.7}\label{r83:entropy-rate}
\end{equation}

\begin{theorem}[Native finite-bank increment entropy]
\label{r83:entropy}
One has $0\le h_{\rm kin}\le f(\rho)$. For any finite set $A$ of
distinct space-time increment indices, $k=|A|$, their joint native
temporal-gauge law $\mathsf P_A$ satisfies
\begin{equation}
 D(\mathsf P_A\|\nu_\gamma^{\otimes k})
 \le k h_{\rm kin}\le k f(\rho).
 \tag{V83.8}\label{r83:entropy-eq}
\end{equation}
The same bound holds after pushing forward to the original
gauge-invariant temporal plaquette energies. There is no dependence
on their separation or total time span.
\end{theorem}
\begin{proof}
The exact path density ratio is
\[
 \frac{d\mathbb P_n}{d\mathbb R_n}
 =\frac{\Omega(U_n)}{\Omega(U_0)}\Lambda^{-n}
 \exp\left[-\tfrac12V(U_0)-\sum_{t=1}^{n-1}V(U_t)
                  -\tfrac12V(U_n)\right].
\]
At each fixed regulator all functions are bounded, $\Omega$ is
bounded away from zero, and the logarithm is integrable.
Stationarity cancels the two logarithmic vacuum terms and gives
\begin{equation}
 D(\mathbb P_n\|\mathbb R_n)
 =n(F-\mathbb E_{\Omega^2}V)=nm h_{\rm kin}.
 \tag{V83.9}\label{r83:path-entropy}
\end{equation}
Nonnegativity and \eqref{r83:pressure-eq} prove the first assertion.
Data processing bounds the entropy of all $mn$ increments against
their product reference by the same number.

We spell out the classical fractional-packing step for a product
reference $Q=\bigotimes_{i=1}^M Q_i$. If a finite multiset of coordinate
subsets $\mathcal A$ covers each coordinate at most $r$ times, then
\begin{equation}
 \sum_{S\in\mathcal A}D(P_S\|Q_S)\le r D(P\|Q).
 \tag{V83.10}\label{r83:packing}
\end{equation}
Order the coordinates. The chain rule writes the full entropy as
$\sum_i d_i$, where
$d_i=\mathbb E D(P(X_i\mid X_{<i})\|Q_i)\ge0$.
The corresponding term for a marginal $S$ is no larger than $d_i$:
its conditional density is a mixture of the more-conditioned
densities, and relative entropy is convex in its first argument.
Sum over $i\in S$, then over $S$. This proof applies on the present
compact Borel spaces; alternatively finite measurable partitions and
monotone refinement give the same statement.

Translate $A$ so its times lie in $0,\ldots,w$. For $n>w$, take all
its time translates by $0,\ldots,n-w-1$ and all spatial transforms
in the group $G$ generated by torus translations and the six
coordinate permutations. This group has size $6N^3=2m$ and acts
transitively on positive spatial edges, with two elements taking
any specified edge to another. It preserves the magnetic action,
the product kinetic law, and $\Omega$ by uniqueness. Only coordinate
permutations are used, not edge reversals depending on link values.
Thus every transformed bank has entropy $D(\mathsf P_A\|\nu^k)$.
Each coordinate in the $n$-step block belongs to at most $2k$ of
these banks, counted with multiplicity: for each original bank
index its time shift is determined and there are two possible
spatial transforms. Apply \eqref{r83:packing} to all block increments:
\[
 2m(n-w)D(\mathsf P_A\|\nu_\gamma^k)
 \le 2k\,nm h_{\rm kin}.
\]
Let $n\to\infty$ at this same finite spatial regulator. Consistency
of the stationary path law keeps the left marginal fixed and gives
\eqref{r83:entropy-eq}. This step proves a bound independent of $w$;
it is not an exchange of infinite-volume and vacuum limits.
Finally use data processing under $Q\mapsto\gamma\epsilon(Q)$ and
V82's exact identification of these energies in the original law.
\end{proof}

\subsection{A centered covariance estimate for every signed finite bank}

Let $\mathsf L_A$ be the native law of the energy vector $D_A$ and
let $\mathsf L_A^0$ be the product of its $k$ single-increment free
laws. Set
\begin{equation}
 \eta=\min\{1,\sqrt{k f(\rho)/2}\},\qquad
 v_0(\gamma)=\operatorname{Var}_{\nu_\gamma}(\gamma\epsilon)\le3/2.
 \tag{V83.11}\label{r83:eta}
\end{equation}
Pinsker's inequality, with total variation defined as
$\sup_B|P(B)-Q(B)|$, gives
$\|\mathsf L_A-\mathsf L_A^0\|_{\rm TV}\le\eta$.
To recall its elementary normalization, apply data processing to
an event where the density difference is positive; the resulting
Bernoulli relative entropy is at least twice the squared difference
of its parameters, by its second derivative in the first parameter.
Approximation handles a supremum or endpoint parameter.

\begin{theorem}[Native centered signed-source comparison]
\label{r83:covariance}
For every real bank coefficient vector $v$ and $\eta>0$,
\begin{equation}
 \left|\operatorname{Var}_{\rm vac}(v\cdot D_A)
                   -v_0(\gamma)\|v\|_2^2\right|
 \le160\|v\|_1^2\eta\log^2(e/\eta).
 \tag{V83.12}\label{r83:covariance-eq}
\end{equation}
In particular
\begin{equation}
 \|\operatorname{Cov}_{\rm vac}(D_A)-v_0(\gamma)I_k\|_{2\to2}
 \le160k\eta\log^2(e/\eta).
 \tag{V83.13}\label{r83:covariance-op}
\end{equation}
At $\beta=0$ the comparison is exact and the right sides are defined
as zero. The variances use the \emph{actual means}.
\end{theorem}
\begin{proof}
For nonzero $v$ normalize $X=v\cdot D_A/\|v\|_1$.
Under either measure, $|X|\le\sum_i a_iD_i$ with
$a_i=|v_i|/\|v\|_1$, $\sum_i a_i=1$. V82's one-coordinate bounds
and convexity of the exponential imply
\[
 \mathbb E e^{|X|/2}\le\sum_i a_i\mathbb E e^{D_i/2}
 \le M:=2^{3/2}<3,\qquad \mathbb E|X|\le3/2.
\]
These estimates hold also for the product reference; no native
independence is used. For $R>0$, truncate $X$ to $[-R,R]$ and
$X^2$ to $[0,R^2]$. The total-variation bound and integration of
$\mathbb P(|X|>t)\le M e^{-t/2}$ give
\[
 \begin{split}
 |\mathbb E X-\mathbb E_0X|&\le2R\eta+4M e^{-R/2},\\
 |\mathbb E X^2-\mathbb E_0X^2|
   &\le R^2\eta+M(8R+16)e^{-R/2}.
 \end{split}
\]
Since the sum of the absolute means is at most three, subtraction
of their squares proves
\[
 |\operatorname{Var}X-\operatorname{Var}_0X|
 \le\eta(R^2+6R)+M(8R+28)e^{-R/2}.
\]
Choose $R=2\log(e/\eta)\ge2$. Then $e^{-R/2}\le\eta$, and
\[
 R^2+30R+84\le37R^2=148\log^2(e/\eta)
 \le160\log^2(e/\eta).
\]
Multiply by $\|v\|_1^2$ and use the product reference variance.
For the operator bound use $\|v\|_1^2\le k\|v\|_2^2$ and the
variational characterization of a real symmetric matrix norm.
If $\beta=0$, $\Omega=1$ and the whole increment field is exactly
product, not merely entropy-close. The zero source is trivial.
\end{proof}

For fixed $k$, \eqref{r83:entropy-eq} and
\eqref{r83:covariance-op} tend to zero as $\beta/\gamma\to0$,
uniformly over spatial volume and all placements of the bank. If
also $\gamma\to\infty$, V81's exact radial limit identifies the
limiting energies as independent Gamma$(3/2,1)$ variables and
$v_0(\gamma)\to3/2$. This is a statement about local ultraviolet
increments, not Gaussianity or triviality of the full gauge theory.
A growing bank is admitted only when the displayed error tends to
zero; for example it suffices that
\[
 k^3 f(\rho)\log^4\!\left(\frac{e}{k f(\rho)}\right)\longrightarrow0
 \quad\hbox{and}\quad k f(\rho)\longrightarrow0.
\]
The size restriction is part of the theorem, not an all-bank bound.

For any one common original-law exterior sigma-algebra $\mathcal Y$,
the exact tower decomposition consequently gives
\[
 \mathbb E\operatorname{Cov}(D_A\mid\mathcal Y)
 +\operatorname{Cov}(\mathbb E[D_A\mid\mathcal Y])
 \preceq\left[v_0(\gamma)+160k\eta\log^2(e/\eta)\right]I_k.
\]
Both terms are positive semidefinite. Neither is separately asserted
close to the free covariance; the conclusion is averaged, not a
uniform-exterior estimate.

\subsection{Physical scaling, the unpaid bank limit, and verification}

The result genuinely removes the small integrated magnetic budget
from a \emph{specified-bank} vacuum covariance problem, while
retaining its exact law. It does not remove all costs. Under the
anisotropic convention
$\gamma=c\,a_s/(g^2a_t)$, $\beta=c\,a_t/(g^2a_s)$ with the same
positive constant $c$, one has $\rho=(a_t/a_s)^2$. Thus
$f(\rho)=16(a_t/a_s)^2+4\sqrt6\,a_t/a_s$ tends to zero under time
refinement at fixed spatial cutoff. This algebra does not identify
or construct a full spatial continuum trajectory.

In particular a fixed physical region contains a number of elementary
increments growing with the spatial cutoff and inverse time step.
The factors $k$ in \eqref{r83:entropy-eq} and
\eqref{r83:covariance-op} cannot be discarded for that bank. At
isotropic coupling $\beta=\gamma$, the present small-ratio estimate
does not tend to zero. Even small entrywise or fixed-bank correlations
need not give a summable connected kernel on an unrestricted bank.
Likewise a small amplitude independent of time separation is not
exponential decay in physical time. The source normalization, coherent
long-range returns, and spatial continuum construction remain unpaid.
The next upstream task is an estimate with source cancellation that
survives the growing physical bank, not another assertion that a fixed
bank is small. The full YM mass gap remains open.

The diagnostic independently checks the Bessel bounds and Riccati
barrier algebra, the chord/plaquette trial-state budget, and exact
finite-state path entropy before and after taking increments. It
checks fractional packing on correlated finite laws, space-time
coverage counts, and centered covariance comparisons. The discrete
transfer fixtures are not Haar $\SU(2)$ and the floating tests are not
uniform proofs; those are the analytic arguments above. V82's
mathematical checks are rerun. All earlier mathematical content is
preserved, only V83 is kept as the active main version, and all build
products remain outside \texttt{latex\_notes}.
% END CONSOLIDATED V83 APPEND

% BEGIN CONSOLIDATED V84 APPEND
\clearpage
\section{V84: Local kinetic defects and physically scaled positive-time sources}
\label{r84:section}

\subsection{Direction, inherited facts, and the new question}

V83 paid the native kinetic mean deficit and arbitrary fixed-bank
entropy, but its centered covariance estimate deteriorates with bank
size. A fixed physical time interval contains increasingly many
increments as the time step decreases. Repeating fixed-bank convergence
therefore does not close the next task. We instead retain the physical
transfer and remove only the magnetic plaquettes touching a specified
link, as a comparison inside a proof. This yields an operator form
bound for the actual energy insertion. Positive-time propagation then
pays a second power of the time step, without a spatial-volume factor.
The resulting estimate controls sources at separated times and their
time smearings; it does not control coincident-time contact terms,
spatial refinement, or exponential decay.

The four newly supplied companion versions, SM source V43, Einstein
bridge V36, native stationarity V46, and perturbative ESM V51, have the
same fingerprints as at their preceding review. Their useful discipline
is to retain the actual operator, centering, and scaling before passing
to a limiting claim. None supplies the missing YM infrared estimate.
The regular companion checkpoint remains V85 and the broad literature
checkpoint V90. This appendix uses the Bessel elasticity proved in
V83, not a new assertion of Wilson transfer positivity. The latter and
the physical insertion identity are inherited from the earlier transfer
construction and V39. Archive checks distinguish the local comparison
below from the previous Perron--Harnack and conditional-fibre results.

\subsection{A one-link kinetic form bound, with the native vacuum retained}

On a finite spatial cubic torus of side $N\ge3$, let $E$ be the positive
links and $\mathcal P$ the spatial plaquettes. We allow nonuniform
electric couplings $\gamma_e>0$ and magnetic couplings $\beta_p\ge0$.
With normalized Haar measure and $\epsilon(q)=1-\tfrac12\operatorname{Re}
\operatorname{Tr}q\in[0,2]$, set
\[
 C=\bigotimes_{e\in E}C_{\gamma_e},\qquad
 V(U)=\sum_{p\in\mathcal P}\beta_p\epsilon(U_p),\qquad
 B=M_{e^{-V/2}},\qquad T=BCB.
\]
The operator $C_\gamma$ is convolution with
$e^{-\gamma\epsilon(q)}\,dq/c_0(\gamma)$, where
$c_0(\gamma)=2e^{-\gamma}I_1(\gamma)/\gamma$. Write
\[
 T\Omega=\Lambda\Omega,\quad \Omega>0,\quad \|\Omega\|=1,
 \qquad P=T/\Lambda,\quad Q=I-|\Omega\rangle\langle\Omega|.
 \tag{V84.1}\label{r84:setup}
\]
Here $P$ is the self-adjoint physical transfer, not a Markov conjugate.
On the full Haar space $0<P\le I$; the unique positive top vector is
gauge invariant, so all statements restrict to the Gauss-law space.
All comparisons below are first made on the full tensor product, where
one-link factorization is literal, and then restricted. No tensor
factorization of the constrained physical Hilbert space is assumed.

For the temporal plaquette energy $D_{e,t}=\gamma_e\epsilon
(U_{e,t}U_{e,t+1}^{-1})$ in temporal gauge, let $T_{D_e}$ be the kernel
$T(U,U')$ multiplied by $D_e(U,U')$. These are the original
gauge-invariant temporal plaquettes after gauge fixing, not substituted
link observables. Define
\[
 \begin{split}
 m_{0,e}&=-\gamma_e\partial_{\gamma_e}\log c_0(\gamma_e),\qquad
 \mu_e=\Lambda^{-1}\langle\Omega,T_{D_e}\Omega\rangle,\\
 A_e&=\Lambda^{-1}(m_{0,e}T-T_{D_e}),\qquad
 \delta_e=m_{0,e}-\mu_e=\langle\Omega,A_e\Omega\rangle,\\
 v_e&=2\sum_{p\ni e}\beta_p,\qquad h_e=e^{v_e}-1.
 \end{split}
 \tag{V84.2}\label{r84:defect}
\]
The derivative in the definition of the kinetic insertion holds every
other electric coupling and every magnetic coupling fixed. It does
not differentiate $\Lambda$ or $\Omega$.

\begin{theorem}[Local defect form domination]
\label{r84:local}
At every finite regulator and for all the couplings just specified,
\[
 0\preceq A_e\preceq2\bigl[h_eI+(I-P)\bigr],
 \qquad 0\le\delta_e\le\min\{m_{0,e},2h_e\}.
 \tag{V84.3}\label{r84:form}
\]
In particular $m_{0,e}\le3/2$. For homogeneous couplings, the four
spatial plaquettes incident on $e$ give $v_e=8\beta$. In that case
the previous bound can be retained as well:
\[
 \delta_e\le\min\{3/2,\,2(e^{8\beta}-1),\,2f(\beta/\gamma)\},
 \quad f(\rho)=16\rho+4\sqrt{6\rho}.
 \tag{V84.4}\label{r84:mean}
\]
There is no translation-invariance hypothesis in \eqref{r84:form}.
\end{theorem}
\begin{proof}
On degree-$j$ Peter--Weyl modes of one $\SU(2)$ link, $C_\gamma$
has eigenvalue $k_j(\gamma)=I_{j+1}(\gamma)/I_1(\gamma)$, with
$k_0=1$ and $0<k_j<1$ for $j>0$. By \eqref{r83:elasticity-eq},
\[
 0\le\gamma k'_j(\gamma)
 \le-2k_j(\gamma)\log k_j(\gamma)
 \le2(1-k_j(\gamma)).
 \tag{V84.5}\label{r84:bessel-form}
\]
The last inequality is $-x\log x\le1-x$ on $(0,1]$.
Thus the bounded self-adjoint operator
$K_e=\gamma_e\partial_{\gamma_e}C_{\gamma_e}$ satisfies
$0\preceq K_e\preceq2(I-C_{\gamma_e})$. Differentiation of the
normalized one-link kernel gives exactly $K_e$ equal to convolution
with $(m_{0,e}-D_e)e^{-\gamma_e\epsilon}/c_0(\gamma_e)$.
This kernel need not be pointwise positive; its operator positivity is
the spectral fact just proved.

Split $V=V_e+V_{\bar e}$, where $V_e$ sums precisely the plaquettes
containing $e$. Then $0\le V_e\le v_e$, while $V_{\bar e}$ does not
depend on $U_e$. Put $M_e=M_{e^{-V_e/2}}$ and
\[
 R_e=M_{e^{-V_{\bar e}/2}}
       \left(\bigotimes_{f\ne e}C_{\gamma_f}\right)
       M_{e^{-V_{\bar e}/2}},\qquad
 T^{(e)}=C_{\gamma_e}\otimes R_e.
\]
The multiplications in $R_e$ act on the other links only.
$R_e$ is positive, $\Lambda_e:=\|R_e\|=\|T^{(e)}\|$, and
$T=M_eT^{(e)}M_e$. Since $T^{(e)}=M_e^{-1}TM_e^{-1}$,
\[
       \Lambda_e/\Lambda\le\|M_e^{-1}\|^2\le e^{v_e}.
 \tag{V84.6}\label{r84:local-ratio}
\]
This is a ratio for deleting finitely many local plaquettes, not a
global vacuum-density comparison. The exact insertion is
$A_e=\Lambda^{-1}M_e(K_e\otimes R_e)M_e$. Consequently,
\[
 \begin{split}
 0\preceq A_e
 &\preceq \frac{2}{\Lambda}
    M_e\bigl[(I-C_{\gamma_e})\otimes R_e\bigr]M_e\\
 &\preceq 2\bigl[(\Lambda_e/\Lambda)M_e^2-P\bigr]
 \preceq2\bigl[e^{v_e}I-P\bigr].
 \end{split}
\]
The second inequality uses $I\otimes R_e\preceq\Lambda_e I$;
the last uses $M_e^2\preceq I$. These are congruences and sums of
operator inequalities, not products of noncommuting orders. Taking
the vacuum expectation proves $\delta_e\le2h_e$. Since the native
kernel energy is nonnegative, $\mu_e\ge0$, giving
$\delta_e\le m_{0,e}\le3/2$. The homogeneous pressure bound is
\eqref{r83:mean-eq}.
\end{proof}

\subsection{Source vectors, a spectral-window bound, and propagation}

Define the exactly centered insertion vector
\[
 \psi_e=Q\Lambda^{-1}T_{D_e}\Omega=-QA_e\Omega.
 \tag{V84.7}\label{r84:vector}
\]
For a nonnegative bounded operator $A$ and a unit vector $\Omega$,
Cauchy--Schwarz for $A^{1/2}$ gives the rank-one order
$|A\Omega\rangle\langle A\Omega|preceq
\langle\Omega,A\Omega\rangle A$. Hence \eqref{r84:form} implies
\[
 |\psi_e\rangle\langle\psi_e|
 \preceq\delta_e QA_eQ
 \preceq2\delta_e Q\bigl[h_eI+(I-P)\bigr]Q.
 \tag{V84.8}\label{r84:rank-one}
\]
This step is stronger than a bound on the unsmoothed source norm:
it retains the transfer Dirichlet form on the right.

\begin{theorem}[Positive-time and low-energy source bounds]
\label{r84:propagation}
For every real $q\ge0$, define
\[
 b_q=\sup_{0\le x\le1}x^q(1-x)
 =\begin{cases}1,&q=0,\\
 q^q/(q+1)^{q+1},&q>0.
 \end{cases}
 \qquad b_q\le\frac{1}{q+1}.
\]
Fractional powers are defined by spectral calculus. Then
\[
 \|P^{q/2}\psi_e\|^2\le2\delta_e(h_e+b_q).
 \tag{V84.9}\label{r84:smooth}
\]
For a specified time step $a>0$ let $H_a=-a^{-1}\log P$, the actual
finite-regulator transfer Hamiltonian. Its spectral projector satisfies
\[
 \left\|\mathbf1_{[0,E]}(H_a)\frac{\psi_e}{a}\right\|^2
 \le \frac{2\delta_e}{a^2}(h_e+1-e^{-aE})
 \le2\frac{\delta_e}{a}\left(\frac{h_e}{a}+E\right),
 \quad E\ge0.
 \tag{V84.10}\label{r84:spectral}
\]
The vacuum is already removed by $Q$. These are upper bounds on
spectral weight, not the exclusion of an interval of positive energies.
\end{theorem}
\begin{proof}
For any bounded real Borel function $g$, multiply
\eqref{r84:rank-one} on both sides by $g(P)$ and take operator norms.
The norm of the left rank-one operator is $\|g(P)\psi_e\|^2$.
The right side is bounded by
$2\delta_e\sup_{0\le x\le1}|g(x)|^2(h_e+1-x)$.
For $g(x)=x^{q/2}$ this is at most
$2\delta_e(h_e+b_q)$. For the spectral projector it is at most
$2\delta_e(h_e+1-e^{-aE})$; divide by $a^2$ and use
$1-e^{-aE}\le aE$. Elementary maximization gives $b_q$.
Strict positivity of the finite-regulator transfer defines $H_a$ by
spectral calculus even though its spectrum is unbounded above.
\end{proof}

For distinct temporal steps separated by an integer $k\ge1$, exact
vacuum sewing gives
\[
 \operatorname{Cov}_{\rm vac}(D_{e,0},D_{f,k})
       =\langle\psi_e,P^{k-1}\psi_f\rangle.
 \tag{V84.11}\label{r84:cov-sew}
\]
Indeed the uncentered expectation is
$\Lambda^{-2}\langle\Omega,T_{D_e}P^{k-1}T_{D_f}\Omega\rangle$;
the vacuum component contributes exactly $\mu_e\mu_f$.
No independence is used, including for adjacent steps $k=1$.
Cauchy--Schwarz and \eqref{r84:smooth} therefore prove
\[
 \left|\operatorname{Cov}_{\rm vac}(D_{e,0},D_{f,k})\right|
 \le2\sqrt{\delta_e\delta_f}
          \sqrt{(h_e+b_{k-1})(h_f+b_{k-1})}.
 \tag{V84.12}\label{r84:cross}
\]
For $e=f$ the covariance is nonnegative. In fact its sequence in
$k-1$ is a moment sequence of a positive measure on $[0,1]$;
all alternating finite differences have the appropriate nonnegative
sign. This concerns one identical energy insertion at separated times,
not arbitrary distinct-link covariance signs. In the homogeneous case,
with $h=e^{8\beta}-1$, a convenient common bound is
\[
 \left|\operatorname{Cov}_{\rm vac}(D_{e,0},D_{f,k})\right|
       \le4h\left(h+\frac{1}{k}\right),\qquad k\ge1.
 \tag{V84.13}\label{r84:hom-cov}
\]

\subsection{Physical-time normalization and growing time banks}

Fix the spatial cutoff and take $\beta=ba$, with $b\ge0$ fixed,
$0<a\le a_*$, and arbitrary $\gamma>0$ at each step size. The
anisotropic Wilson choice is $b=c/(g^2a_s)$ and
$\gamma=c a_s/(g^2a)$; no conclusion here requires uniformity as
$a_s\downarrow0$. Put
\[
 L_*=8b e^{8ba_*},\qquad
 \frac{e^{8ba}-1}{a}\le L_*,\qquad
 \mathcal E_{e,t}^{(a)}=\frac{D_{e,t}-\mu_e}{a}.
 \tag{V84.14}\label{r84:scaling}
\]
This source normalization is the one obtained by dividing the
dimensionless centered temporal-action insertion by the time step.
It is not yet an identification with a renormalized continuum
energy-density field; that identification still needs its matching
and spatial renormalization.

\begin{corollary}[Volume-uniform separated-time source amplitudes]
\label{r84:physical}
At physical separation $\tau=ka>0$, for every pair of spatial links,
\[
 \left|\operatorname{Cov}_{\rm vac}
       (\mathcal E_{e,0}^{(a)},\mathcal E_{f,k}^{(a)})\right|
 \le4L_*\left(L_*+\frac{1}{\tau}\right).
 \tag{V84.15}\label{r84:physical-cov}
\]
Uniformly in spatial volume, the corresponding spectral weight obeys
\[
 \left\|\mathbf1_{[0,E]}(H_a)\frac{\psi_e}{a}\right\|^2
       \le4L_*(L_*+E).
 \tag{V84.16}\label{r84:physical-spectral}
\]
Let $f_{e,t},g_{e,t}$ be finitely supported real arrays, with all
steps in the first time support preceding all steps in the second by
at least $\tau_*/a$, where $\tau_*>0$. Define
\[
 \mathcal O_a(f)=\sum_{e,t}a f_{e,t}\mathcal E_{e,t}^{(a)},
 \qquad \|f\|_{1,a}=\sum_{e,t}a|f_{e,t}|.
\]
Then
\[
 |\operatorname{Cov}_{\rm vac}(\mathcal O_a(f),\mathcal O_a(g))|
 \le4L_*\left(L_*+\frac{1}{\tau_*}\right)
                    \|f\|_{1,a}\|g\|_{1,a}.
 \tag{V84.17}\label{r84:time-bank}
\]
In particular the number of increments in a fixed physical-time
smearing may grow like $a^{-1}$ without an additional cardinality
factor, provided these discrete $L^1$ norms stay bounded and the time
supports remain separated.
\end{corollary}
\begin{proof}
Divide \eqref{r84:hom-cov} by $a^2$ and use $ka=\tau$ and
$h/a\le L_*$. For \eqref{r84:physical-spectral} use
$\delta_e\le2h$ in \eqref{r84:spectral}. Expand the covariance of
the finite sums, apply \eqref{r84:physical-cov} termwise, and sum
absolute weights. No interchange of an infinite sum and a limiting
law is involved. The assertion is a uniform bound, not a claim that
the time-continuum law or these smearings have already been constructed.
\end{proof}

This pays a specific growing-time-bank problem that V83's entropy
estimate did not pay. It remains an $L^1$ spatial bound: for a spatial
bank of $s$ equal-sized signed coefficients, replacing the $L^1$ norm
by its $L^2$ norm can cost $\sqrt{s}$. There is no unrestricted spatial
covariance-operator bound in \eqref{r84:time-bank}.

There is also an operator statement on the whole finite-energy sector,
not only on the vacuum vector. Let $\Pi_E=\mathbf1_{[0,E]}(H_a)$ and
let the centered one-step kernel operator be
$J_e^{(a)}=a^{-1}(\Lambda^{-1}T_{D_e}-\mu_eP)
=a^{-1}(\delta_eP-A_e)$. Since $I-P\preceq aH_a$, compression of
\eqref{r84:form} gives
\[
 0\preceq\Pi_E\frac{A_e}{a}\Pi_E
       \preceq2(L_*+E)\Pi_E,\qquad
 \|\Pi_EJ_e^{(a)}\Pi_E\|\le2(L_*+E).
 \tag{V84.18}\label{r84:compressed}
\]
For the second assertion, its upper quadratic-form bound is
$\delta_e/a\le2L_*$ and its lower one is $-2(L_*+E)$; take the
larger absolute endpoint. The spectral projection preserves the
physical Gauss-law space. This is a uniform bound on a specified
energy window; a common continuum Hilbert space and convergence of
these compressed operators have not been constructed here.

\subsection{Contact terms and an explicit infrared nonimplication}

The sewing formula \eqref{r84:cov-sew} cannot be used at $k=0$ by
inserting $P^{-1}$. At coincident times the correct second moment
involves $T_{D_eD_f}$, not $T_{D_e}T_{D_f}$. Already at $\beta=0$,
$\Omega=1$, $\delta_e=0$, and $\psi_e=0$, so all distinct-time energy
covariances vanish. Nevertheless
\[
 \operatorname{Var}_{\rm vac}(\mathcal E_{e,t}^{(a)})
       =\frac{v_0(\gamma)}{a^2},\qquad
 v_0(\gamma)=\operatorname{Var}_{\nu_\gamma}
                         (\gamma\epsilon)>0.
 \tag{V84.19}\label{r84:contact}
\]
If $\gamma\to\infty$, $v_0(\gamma)\to3/2$ by V81. For a fixed
nonzero continuous compactly supported time test $f$, the same
free increment law gives
$\operatorname{Var}(\mathcal O_a(f))=
v_0(\gamma)\sum_t f(ta)^2\sim(3/(2a))\int f^2$.
Thus the separated-support bound does not assert finite variance of
an unrenormalized full time smearing. Contact subtraction must be
specified and justified, not silently inferred from propagation.

Nor does \eqref{r84:rank-one} imply a gap. To see this at the level
of the proved operator certificate, take an abstract two-dimensional
Hilbert space and, for $a,b,\varepsilon>0$, set
\[
 \Omega=\binom10,\quad
 P=\begin{pmatrix}1&0\\0&e^{-a\varepsilon}\end{pmatrix},\quad
 h=ab,\quad A=ab\begin{pmatrix}1&1\\1&1\end{pmatrix}.
 \tag{V84.20}\label{r84:counterexample}
\]
Here $\delta=ab$, $\psi=-QA\Omega=-ab\binom01$, and
\[
 2\bigl[hI+(I-P)\bigr]-A
 =\begin{pmatrix}ab&-ab\\-ab&ab+2(1-e^{-a\varepsilon})\end{pmatrix}
 \succeq0.
\]
Indeed its first principal minor is positive and its determinant is
$2ab(1-e^{-a\varepsilon})>0$. Thus the entire local form certificate
holds, while the normalized source has spectral weight $b^2$ at the
arbitrarily small energy $\varepsilon$. The example is not claimed
to be a Wilson gauge theory; it disproves the implication from this
operator certificate alone to a regulator-uniform gap. The constant
term $4L_*^2$ at large $\tau$ in \eqref{r84:physical-cov} is consistent
with exactly this limitation.

\subsection{Verification and next unpaid step}

The new mechanism is local magnetic deletion, the positive kinetic
defect, and propagation of its retained Dirichlet form. It proves
volume-uniform physically scaled source amplitudes at positive time
and a finite-energy spectral-weight bound. It neither changes the
native vacuum nor assumes the unknown gap. It also does not solve the
coherent $B_J$ estimate, distributed holonomy returns, spatial
renormalization, or the interacting four-dimensional continuum problem.
The full YM mass gap remains unproved.

The diagnostic checks \eqref{r84:bessel-form} at high precision for
multiple degrees and couplings, verifies local deletion and insertion
algebra in finite tensor-transfer fixtures, and tests the propagation,
spectral-window, and two-level nonimplication formulas. The finite
fixtures use exact $\SU(2)$ Bessel multipliers but only a two-mode
auxiliary link space; they are explicitly not Haar $\SU(2)$ lattice
theories. Their role is to detect normalization or operator-algebra
mistakes, not to certify the analytic uniform bounds. V83's independent
mathematical diagnostics are rerun. All prior mathematical text is
preserved and all build products remain outside \texttt{latex\_notes}.

The next upstream target is to replace the nondecaying local allowance
$h_eI$ in the low-energy source sector by a spatially controlled
connected estimate, or return to the exact coherent $B_J$ identity
where that missing cancellation is visible. Merely improving the
constant $L_*$ would not address this infrared obstruction. At V85
the parallel-program checkpoint must also be performed.
% END CONSOLIDATED V84 APPEND

% BEGIN CONSOLIDATED V85 APPEND
\clearpage
\section{V85: Conserved local energy and a volume-uniform spatial sum rule}
\label{r85:section}

\subsection{Companion checkpoint and the change of coordinate}

V84 controls individual kinetic insertions after physical-time propagation,
but its scalar local allowance does not decay and its spatial source bound
is of $\ell^1$ type. We now retain the magnetic part in the local energy.
The spatially constant sum is then the actual Hamiltonian, which annihilates
the centered vacuum exactly. The resulting double-commutator calculation
gives a nonperturbative spatial-gradient bound for the \emph{complete}
energy-weighted source Gram matrix, not only for a fixed bank. Its spectral
energy weight is essential: an unweighted covariance bound or a mass gap
does not follow by dropping that weight.

The scheduled V85 companion review checked the current inventory and
fingerprints and reread the terminal mathematical scope of NS V26,
SM--Einstein V72, shell mixing V16, and parallel YM V5. These files, and
the four latest supplied attachments (SM source V43, Einstein bridge V36,
native stationarity V46, and perturbative ESM V51), are unchanged. NS's
rank-one Oseen calculations and SM's fixed-time, fixed-mode many-copy
transport do not control the native YM infrared sector. The shell-mixing
calculation retains all contributions to its measured one-loop coefficient
but supplies no nonperturbative remainder. The parallel YM warning against
volume-uniform tails for extensive total charge remains applicable. No
auxiliary limit, selected trajectory, or perturbative coefficient is
imported as a YM closure theorem. The next companion and broad-literature
checkpoints are both V90.

The finite differential gauge Hamiltonian and its ground-state representation
are already in f1 V6; calibrated fixed-spatial-volume temporal Wilson
identification is in f1 V10--V11. Those results are not reproved as new.
The canonical Hamiltonian background is classical
\href{https://doi.org/10.1103/PhysRevD.11.395}{Kogut--Susskind (1975)}.
The energy-weighted double-commutator identity is likewise a standard
sum-rule tool; see, for example,
\href{https://arxiv.org/abs/1710.03187}{Lu--Johnson (2018)}.
We prove the identity and the needed compact-group estimates here. The new
object in this continuation is the native, volume-uniform conserved-energy
bank estimate below, not a claim of priority for the general sum-rule method.

\subsection{The Hamiltonian and a local kinetic estimate}

Let $N\ge3$ and let $E,\mathcal P$ denote the positive links and spatial
plaquettes of the cubic torus. On product probability Haar measure on
$\SU(2)^E$, use the unit-round quaternion metric and $L_e=-\Delta_e$,
whose one-link eigenvalues are $j(j+2)$. Fix $\kappa>0$ and $b\ge0$ and set
\[
 V_p=b\epsilon(U_p),\qquad
 \mathcal H=\kappa\sum_{e\in E}L_e+\sum_{p\in\mathcal P}V_p,
 \qquad H=\mathcal H-E_0,
 \quad \mathcal H\Omega=E_0\Omega,\quad\|\Omega\|=1.
 \tag{V85.1}\label{r85:H}
\]
This is the differential Hamiltonian after the fixed-spatial-regulator
temporal limit, not the finite-step operator $-a^{-1}\log P$ of V84.
In round-metric normalization the calibrated Wilson schedule can be written
$a=-\log k_1(\gamma)/(3\kappa)$, $\beta=ba$; asymptotically
$\gamma\sim(2\kappa a)^{-1}$. The inherited identification is at each
fixed spatial graph. No uniform interchange with spatial refinement is
asserted here.

The bounded smooth real potential is a perturbation of the compact elliptic
product Laplacian. The closed operator has compact resolvent. Its heat
semigroup is positivity improving: the positive free heat kernel and the
bounded-potential Feynman--Kac weight give this directly. Thus the ground
state is unique and strictly positive; elliptic regularity makes it smooth.
Gauge invariance forces $\Omega$ to be physical. We may prove estimates on
the full Haar space and restrict to Gauss law, without factorizing the
constrained physical Hilbert space. All commutators below are evaluated on
the smooth core and on $\Omega$, so their displayed expectations have no
unbounded-domain ambiguity.

\begin{lemma}[Local kinetic energy without a volume factor]
\label{r85:local}
As quadratic forms,
\[
 \kappa L_e\preceq H+8bI,
 \qquad \kappa\int|\nabla_e\Omega|^2\le8b.
 \tag{V85.2}\label{r85:local-eq}
\]
\end{lemma}
\begin{proof}
Remove $\kappa L_e$ and the four $V_p$ containing $e$ from $\mathcal H$.
The remaining operator $\mathcal H_{\bar e}$ acts on the other links only;
write $E_{\bar e}$ for its bottom energy. Since $0\le\sum_{p\ni e}V_p\le8b$,
\[
 E_{\bar e}\le E_0\le E_{\bar e}+8b.
\]
For the upper bound use the ground state of $\mathcal H_{\bar e}$ tensored
with the constant function of $U_e$ as a variational vector on the full
space. Its bottom energy equals the physical bottom by uniqueness and gauge
invariance. Also $\mathcal H\succeq\kappa L_e+E_{\bar e}I$, proving the form
bound. Take its expectation in $\Omega$ and integrate by parts.

The same deletion also pays a local second kinetic moment. Put
$A=\kappa L_e+(\mathcal H_{\bar e}-E_{\bar e})$. Its two nonnegative
summands strongly commute because they act on distinct tensor factors.
Thus $\|\kappa L_eu\|\le\|Au\|$ on their common operator domain.
The ground equation gives
$A\Omega=(E_0-E_{\bar e}-\sum_{p\ni e}V_p)\Omega$, whose norm is at most
$8b$. Consequently $\|\kappa L_e\Omega\|\le8b$, uniformly in volume.
\end{proof}

\subsection{Compact plaquette derivatives and the double commutator}

For $e,f\in p$, quaternion multilinearity gives
\[
 L_eV_p=3(V_p-b),\qquad L_fL_eV_p=9(V_p-b),
 \quad |\nabla_eV_p|\le b,\quad
 \|\operatorname{Hess}_{f,e}V_p\|\le b.
 \tag{V85.3}\label{r85:jets}
\]
Here the Hessian is intrinsic when $e=f$ and mixed between the two product
factors otherwise. Holding the other links fixed, $V_p$ is $b$ minus a
height function of amplitude $b$ on the round $S^3$. A height function has
Laplacian eigenvalue three and Hessian equal to minus itself times the
metric, proving the same-link statements. For distinct links, differentiation
in two unit tangent directions replaces two unit quaternion factors in the
plaquette word by unit tangent quaternions. The absolute real part of their
product is at most one. Inversion is quaternion conjugation and is a linear
isometry, so the argument includes both orientations. This proves the
mixed-Hessian bound. If a differentiated link is not in $p$, the derivative
vanishes.

\begin{lemma}[An integrated second-commutator identity]
\label{r85:double}
For smooth real $u,V$ on the product group,
\[
 \begin{split}
 \langle u,[L_f,[L_e,V]]u\rangle
 &=\int u^2 L_fL_eV
   -4\int\operatorname{Hess}_{f,e}V(\nabla_f u,\nabla_e u),\\
 [W,[L_e,V]]&=2M_{\langle\nabla_eW,\nabla_eV\rangle}
 \end{split}
 \tag{V85.4}\label{r85:double-eq}
\]
for real multiplication $W$. Consequently, in the actual ground state,
\[
 \begin{split}
 |\langle[L_f,[L_e,V_p]]\rangle_\Omega|
       &\le 9b+32b^2/\kappa\quad(e,f\in p),\\
 |\langle[V_q,[L_e,V_p]]\rangle_\Omega|
       &\le2b^2\quad(e\in p\cap q).
 \end{split}
 \tag{V85.5}\label{r85:double-bounds}
\]
The corresponding commutators vanish outside the indicated incidences.
\end{lemma}
\begin{proof}
The first commutator is $[L_e,V]=-(\Delta_eV)-2\nabla_eV\cdot\nabla_e$.
This immediately gives the second identity. If $e\ne f$, expand the next
commutator and integrate once in each independent product factor; the
third-derivative terms cancel, leaving the first line of
\eqref{r85:double-eq}. For $e=f$ the same calculation is intrinsic.
Explicitly, with $L=-\Delta$, the second-commutator differential expression is
\[
 \Delta^2V+4\nabla\Delta V\cdot\nabla
 +4\operatorname{Hess}V:\operatorname{Hess}
 +4\operatorname{Ric}(\nabla V,\nabla).
\]
Integration of the Hessian term uses
$\operatorname{div}\operatorname{Hess}V=\nabla\Delta V+
\operatorname{Ric}(\nabla V,\cdot)$. These two first-order terms cancel
the two displayed first-order contributions in the expectation, leaving
$\int u^2\Delta^2V-4\int\operatorname{Hess}V(\nabla u,\nabla u)$.
Thus no curvature term has been omitted.

Apply \eqref{r85:jets} to $u=\Omega$ and use
\eqref{r85:local-eq}:
\[
 \int|\nabla_f\Omega||\nabla_e\Omega|
 \le\|\nabla_f\Omega\|_2\|\nabla_e\Omega\|_2\le8b/\kappa.
\]
This proves the first bound. The second follows from the gradient bound
in \eqref{r85:jets}. The formulas also prove the support assertions.
\end{proof}

\subsection{A complete signed-bank estimate with conservation retained}

Partition the magnetic energy equally among the four links of each face:
\[
 h_e=\kappa L_e+\frac14\sum_{p\ni e}V_p,\qquad
 \sum_e h_e=\mathcal H.
\]
In particular $\|h_e\Omega\|\le10b$ by the second-moment conclusion
above and $\|\tfrac14\sum_{p\ni e}V_p\|\le2b$. This is a local
Hamiltonian-source bound, not a bound on V84's raw coincident temporal
plaquette measurement or its contact term.
For real weights $w=(w_e)$, put
\[
 \begin{split}
 \bar w_p&=\tfrac14\sum_{e\in p}w_e,\qquad
 F_w=\sum_e w_eh_e
       =\kappa\sum_e w_eL_e+\sum_p\bar w_pV_p,\\
 \Phi_w&=(F_w-\langle F_w\rangle_\Omega)\Omega,\qquad
 S(w)=\sum_p\sum_{e\in p}|w_e-\bar w_p|^2
       =\langle w,\mathsf L_\square w\rangle.
 \end{split}
 \tag{V85.6}\label{r85:bank}
\]
The real symmetric matrix $\mathsf L_\square$ is defined by this quadratic
form. Each plaquette contributes its four-coordinate projection off the
constant vector; each link lies in four plaquettes. Hence
\[
 0\preceq\mathsf L_\square\preceq4I,
 \qquad \mathsf L_\square\mathbf1=0,
 \qquad \Phi_{\mathbf1}=0.
 \tag{V85.7}\label{r85:incidence}
\]
The last identity uses the full Hamiltonian, not separate conservation of
its electric and magnetic parts.

\begin{theorem}[Volume-uniform energy-weighted spatial Gram]
\label{r85:main}
Set $C_{\kappa,b}=36\kappa^2b+132\kappa b^2$. For every spatial volume
and every real or complex bank $w$,
\[
 \boxed{\quad
 0\le\langle\Phi_w,H\Phi_w\rangle
       \le C_{\kappa,b}S(w)\le4C_{\kappa,b}\|w\|_2^2.
 \quad}
 \tag{V85.8}\label{r85:main-eq}
\]
Equivalently the complete matrix
$\mathsf M_{ef}=\langle\Phi_{\mathbf1_e},H\Phi_{\mathbf1_f}\rangle$
satisfies $0\preceq\mathsf M\preceq C_{\kappa,b}\mathsf L_\square$.
It is exactly finite range: $\mathsf M_{ef}=0$ when $e,f$ have distance
greater than two in the link graph whose adjacency is sharing a plaquette.
It also has zero row sums. No sign is asserted for individual off-diagonal
entries, and this is not an unweighted covariance estimate.
\end{theorem}
\begin{proof}
For real $w$, the eigenvector equation gives the exact sum rule
\[
 \langle\Phi_w,H\Phi_w\rangle
       =\tfrac12\langle[F_w,[\mathcal H,F_w]]\rangle_\Omega.
 \tag{V85.9}\label{r85:sum-rule}
\]
Indeed expansion of the double commutator gives twice
$\langle F_w\Omega,(\mathcal H-E_0)F_w\Omega\rangle$.
The kinetic terms commute and the magnetic terms commute, so, writing
$d_{ep}=\bar w_p-w_e$,
\[
 [\mathcal H,F_w]=\kappa\sum_p\sum_{e\in p}d_{ep}[L_e,V_p].
 \tag{V85.10}\label{r85:current}
\]
For each separate summand, stationarity gives
$\langle[\mathcal H,[L_e,V_p]]\rangle_\Omega=0$. Thus one can subtract
$w_e\mathcal H$ from the \emph{outer} $F_w$ inside that summand without
changing its expectation. This is the local cancellation that would be
lost by estimating the two energies separately. Equations
\eqref{r85:sum-rule}--\eqref{r85:current} give
\[
 \begin{split}
 \langle\Phi_w,H\Phi_w\rangle
 ={}&\frac{\kappa^2}{2}\sum_{p,e\in p}d_{ep}
       \sum_{f\in p}(w_f-w_e)
                   \langle[L_f,[L_e,V_p]]\rangle_\Omega\\
 &+\frac\kappa2\sum_{p,e\in p}d_{ep}
       \sum_{q\ni e}(\bar w_q-w_e)
                   \langle[V_q,[L_e,V_p]]\rangle_\Omega.
 \end{split}
 \tag{V85.11}\label{r85:expanded}
\]
The notation $\sum_{p,e\in p}$ means $\sum_p\sum_{e\in p}$.
For each plaquette, $w_f-w_e=d_{ep}-d_{fp}$, whence
\[
 \sum_{p,e\in p}|d_{ep}|\sum_{f\in p}|w_f-w_e|
 \le\sum_p\left[4\sum_{e\in p}|d_{ep}|^2+
                         \left(\sum_{e\in p}|d_{ep}|\right)^2\right]
 \le8S(w).
 \tag{V85.12}\label{r85:first-count}
\]
Similarly, four plaquettes meet a link, so
\[
 \sum_{p,e\in p}|d_{ep}|\sum_{q\ni e}|\bar w_q-w_e|
 =\sum_e\left(\sum_{p\ni e}|d_{ep}|\right)^2\le4S(w).
 \tag{V85.13}\label{r85:second-count}
\]
Using \eqref{r85:double-bounds} in \eqref{r85:expanded} yields
\[
 \langle\Phi_w,H\Phi_w\rangle
 \le\left[4\kappa^2(9b+32b^2/\kappa)+4\kappa b^2\right]S(w),
\]
which is exactly the stated constant. Positivity follows from $H\ge0$.
All coefficients, $\Omega$, and differential operators are real; for
$w=u+iv$ the quadratic form is the sum of the forms for $u,v$. This
extends the estimate to complex weights with the same constant.

For the support assertion polarize \eqref{r85:sum-rule}:
$\mathsf M_{gf}=\tfrac12\langle[h_g,[\mathcal H,h_f]]\rangle_\Omega$.
The reality just noted removes the imaginary part. The inner commutator
is a sum of $[L_e,V_p]$ with $f\in p$ and $e\in p$.
An outer $L_g$ contributes only if $g\in p$; an outer $V_q$ from $h_g$
only if $g\in q$ and $e\in q$. These are paths of at most two plaquette
adjacencies. Finally $\Phi_{\mathbf1}=0$ proves the zero row sums.
\end{proof}

\subsection{The long-wavelength channel and useful spectral consequences}

Let $n=N^3$ and choose a nonzero torus momentum
$p\in(2\pi/N)\mathbb Z^3$, represented in $[-\pi,\pi]^3$. Define the
unit-norm, equal-orientation midpoint bank
\[
 w_{(x,j)}=(3n)^{-1/2}e^{ip\cdot(x+\widehat e_j/2)}.
\]
For the four edge midpoints of an $ij$ face, its phase sum is a common
phase times $2[\cos(p_i/2)+\cos(p_j/2)]$. Therefore
\[
 S(w)=\frac13\sum_{i<j}
       \left[4-\bigl(\cos(p_i/2)+\cos(p_j/2)\bigr)^2\right]
       \le\frac{|p|^2}{3}.
 \tag{V85.14}\label{r85:momentum}
\]
Indeed the cosines are nonnegative on the chosen momentum representatives,
and each summand is at most $(p_i^2+p_j^2)/2$. Thus
$\langle\Phi_w,H\Phi_w\rangle\le C_{\kappa,b}|p|^2/3$, with the
full interacting ground state retained. At $p=0$ the source vector is
exactly zero. Midpoint phases are part of this convention and cannot
be dropped at nonzero momentum.

Let $\mu_w$ be the spectral measure of $H$ in $\Phi_w$. The exact result
is a first-moment bound
\[
 \int_0^\infty\lambda\,d\mu_w(\lambda)\le C_{\kappa,b}S(w).
 \tag{V85.15}\label{r85:moment}
\]
For example, for every $r>0$,
$\mu_w([r,\infty))\le C_{\kappa,b}S(w)/r$. An all-bank bound with a
time cancellation is also available. For $s>0$ set
$\Psi_{w,s}=(I-e^{-sH})\Phi_w$. Spectral calculus gives
\[
 \begin{split}
 \|\Psi_{w,s}\|^2&\le C_{\kappa,b}sS(w),\\
 |\langle\Psi_{u,s},e^{-tH}\Psi_{w,s}\rangle|
 &\le C_{\kappa,b}\sqrt{S(u)S(w)}
                    \min\{s,s^2/(et)\},\qquad t>0.
 \end{split}
 \tag{V85.16}\label{r85:filtered}
\]
To prove these bounds use $1-e^{-s\lambda}\le\min\{s\lambda,1\}$,
so $(1-e^{-s\lambda})^2/\lambda\le s$, and also
$e^{-t\lambda}(1-e^{-s\lambda})^2/\lambda
\le s^2\lambda e^{-t\lambda}\le s^2/(et)$. Apply
\eqref{r85:moment}, then Cauchy--Schwarz for the cross term.
This is an $\ell^2$-bank estimate for an explicitly time-differenced
physical source. Its algebraic decay is not exponential clustering of
the unfiltered energy density, and it is not a contact-subtraction rule
for V84's coincident temporal plaquettes.

\subsection{Why the remaining low-energy estimate cannot be skipped}

Neither a nonzero first spectral moment nor its spatial-gradient upper
bound excludes weight at arbitrarily small positive energy. This is
visible even while retaining the entire finite-range Gram structure.
On the same edge coefficient space, $\ker\mathsf L_\square$ consists
exactly of constants: $S(w)=0$ forces equality on every face, and the
face-sharing link graph of the cubic torus is connected. Consider the
abstract Hilbert space and positive Hamiltonian
\[
 \mathcal K_N=\mathbb C\Omega\oplus\mathbf1^\perp,\qquad
 H_N=0\oplus\mathsf L_\square|_{\mathbf1^\perp},\qquad
 \Phi_w=0\oplus Q_\mathbf1w.
 \tag{V85.17}\label{r85:soft-example}
\]
Here $Q_\mathbf1$ is orthogonal projection off constants. The centered
source has $\Phi_{\mathbf1}=0$, its unweighted Gram is $Q_\mathbf1$,
and its energy-weighted Gram is exactly the finite-range
$\mathsf L_\square$. Yet the nonzero momentum test in
\eqref{r85:momentum} gives
\[
 0<\operatorname{gap}(H_N)\le\frac{4\pi^2}{3N^2}\longrightarrow0.
 \tag{V85.18}\label{r85:soft-gap}
\]
If desired these source vectors are realized by self-adjoint bounded
operators $|\Phi_w\rangle\langle\Omega|+
|\Omega\rangle\langle\Phi_w|$ for real $w$. This is not a Wilson
counterexample. It proves only that the newly paid Gram bound, conservation,
and finite range do not alone imply a mass gap.

In the actual gauge Hamiltonian let $\mathsf G_0$ be the unweighted Gram
of the same energy bank. A positive physical gap $m$ would imply the
\emph{opposite-direction} relative bound
\[
              m\mathsf G_0\preceq\mathsf M
                     \preceq C_{\kappa,b}\mathsf L_\square.
 \tag{V85.19}\label{r85:unpaid}
\]
The left inequality has not been proved. The variational ratio
$\langle\Phi_w,H\Phi_w\rangle/\|\Phi_w\|^2$ gives an \emph{upper}
bound on the lowest energy visible to that vector, not a lower mass bound.
Moreover this scalar-energy bank is not asserted to be cyclic for all
physical sectors. Even a lower inequality on this bank alone would not
settle the full physical spectrum without a dense-core argument.

The constants here are volume-uniform at fixed spatial cutoff, not
ultraviolet-uniform. For the anisotropic convention of V83--V84,
$\kappa=g^2/(2ca_s)$ and $b=c/(g^2a_s)$, so
\[
 C_{\kappa,b}=a_s^{-3}\left(9g^2/c+66c/g^2\right).
 \tag{V85.20}\label{r85:UV}
\]
Energy-density normalization, renormalization, and spatial refinement
must therefore still be handled. This appendix advances spatial source
control by keeping an exact conservation cancellation, but it does not
prove the interacting four-dimensional continuum construction or YM gap.

\subsection{Verification and the next upstream obligation}

The checker integrates the same-link and mixed-link second-commutator
identities exactly against polynomial tests on $S^3$ and $S^3\times S^3$,
including the intrinsic curvature cancellation. Finite qubit Hamiltonians
independently check the energy partition, stationary subtraction,
double-commutator Gram, and zero row sums; they are not Haar $\SU(2)$
models. Cubic incidence matrices check all face multiplicities, the two
counting inequalities, midpoint momentum symbols, and the explicit soft
model. V84's mathematical checks are rerun. Finite tests supplement, not
replace, the analytic proofs of the uniform constants. All prior text is
preserved apart from the version title, and only the latest main \texttt{.tex}
is retained in \texttt{latex\_notes}; builds and diagnostics are elsewhere.

The next task is the low-energy, unweighted connected source law, not
another first-moment estimate. It may be attacked through the native
electric--magnetic cross Gram and its conservation identity, while retaining
the coherent $B_J$ remainder as an independent unresolved route. A further
upper moment or another time filter would not pay the missing direction
in \eqref{r85:unpaid}. The full mass-gap goal remains active and unproved.
% END CONSOLIDATED V85 APPEND

% BEGIN CONSOLIDATED V86 APPEND
\clearpage
\section{V86: Unweighted block-energy fluctuations from a bounded cut}
\label{r86:section}

\subsection{Direction, scope, and nonduplication}

V85's complete spatial sum rule controls the first spectral moment of
the energy source, not its unweighted covariance. We now bound that
unweighted covariance for block sources, without a gap hypothesis and
without discarding low energies. The gain is geometric: the standard
deviation costs the number of crossing plaquettes, not the block volume.
This is a boundary-squared \emph{variance} bound, not an area-law variance
bound, and not a uniform bound for arbitrary unit-$\ell^2$ banks.

The single-link deletion argument in V85 is inherited. The new use is a
two-sided cut into arbitrary link regions, keeping the internal magnetic
energies on both sides and then restoring the correct fractions of every
crossing plaquette. A level-set argument extends the resulting unweighted
bound to spatial weights of bounded discrete variation. This does not
repeat V85's energy-weighted commutator estimate or V84's temporal filter.

Arad--Kuwahara--Landau,
\href{https://arxiv.org/abs/1406.3898}{\emph{Connecting global and local
energy distributions in quantum spin models on a lattice}}, Section 2.1,
equation (3), and Theorem 2.2, provide related local-energy distribution
bounds under bounded local interaction norms. The unbounded link
Laplacians here do not satisfy that norm assumption. We therefore do not
import their exponential tail estimate: the weaker cut estimate needed
below is proved directly for the differential gauge Hamiltonian.
The companion review remains V85; the next regular companion review and
broad literature sweep are both due at V90.

\subsection{A cut lemma allowing unbounded internal Hamiltonians}

\begin{lemma}[Unweighted second moment across a bounded positive cut]
\label{r86:abstract-cut}
Let $A$ and $B$ be nonnegative strongly commuting self-adjoint operators
on a Hilbert space, and let $V$ be bounded self-adjoint with
$0\preceq V\preceq M I$. Assume $S=A+B$ is self-adjoint, has a normalized
zero vector, and $S+V$ has normalized ground state $\Omega$ with ground
energy $\delta$. Then
\[
 0\le\delta\le M,\qquad
 \|A\Omega\|,\ \|B\Omega\|\le M.
 \tag{V86.1}\label{r86:cut-second}
\]
If $0\preceq W_A\preceq M_A I$, $0\preceq W_B\preceq M_B I$,
and $W_A+W_B=V$, set $F_A=A+W_A$, $F_B=B+W_B$, and
$Q=I-|\Omega\rangle\langle\Omega|$. Then
\[
 \|QF_A\Omega\|=\|QF_B\Omega\|
 \le M+\tfrac12\min(M_A,M_B).
 \tag{V86.2}\label{r86:split-cut}
\]
No commutation of $V$ or $W_A$ with $A$ or $B$ is assumed.
\end{lemma}

\begin{proof}
Nonnegativity and a zero-vector trial for $S$ give $0\le\delta\le M$.
Bounded perturbation leaves the domain of $S$ unchanged, and
\[
              S\Omega=(\delta-V)\Omega.
\]
The spectral theorem for $V$ gives
$\|S\Omega\|\le\max(\delta,M-\delta)\le M$.
In the joint spectral calculus of $A,B$, the scalar inequalities
$a^2\le(a+b)^2$ and $b^2\le(a+b)^2$ for $a,b\ge0$ imply both
the required domain inclusions and
$\|A\Omega\|,\|B\Omega\|\le\|S\Omega\|$.
This is not the invalid assertion that squaring arbitrary noncommuting
positive operators preserves their order.

For any normalized vector, the variance of an operator with spectrum in
$[0,M_A]$ is at most $M_A^2/4$: if its mean is $u$, then
$\langle W_A^2\rangle-u^2\le M_Au-u^2\le M_A^2/4$.
Consequently
\[
 \|QF_A\Omega\|\le\|QA\Omega\|+\|QW_A\Omega\|
                    \le M+M_A/2.
\]
The same argument applies to $B$. Finally
$(F_A+F_B)\Omega=\delta\Omega$, so $QF_A\Omega=-QF_B\Omega$.
Taking the smaller bound proves \eqref{r86:split-cut}.
\end{proof}

\subsection{Native gauge-invariant block energy}

Use exactly the finite spatial Hamiltonian of V85, on a cubic torus
$N\ge3$, with $m=3N^3$ links:
\[
 H_{\rm raw}=\kappa\sum_e L_e+\sum_p V_p,\qquad
 L_e=-\Delta_{S^3,e},\quad
 V_p=b\epsilon(U_p),\quad 0\le V_p\le2b,
 \qquad \kappa>0,\quad b\ge0.
 \tag{V86.3}\label{r86:H}
\]
Let $E_0$ and $\Omega$ be its bottom and normalized positive ground state;
write $H=H_{\rm raw}-E_0$. For a set $D$ of links, define
\[
 \partial_\square D=\{p:0<|p\cap D|<4\},\quad
 k_D=|\partial_\square D|,\qquad
 H_D^\circ=\kappa\sum_{e\in D}L_e+\sum_{p\subset D}V_p.
 \tag{V86.4}\label{r86:cut-def}
\]
The bottom $E_D$ of $H_D^\circ$ refers to its own link-factor Hilbert
space. For the empty region use $E_D=0$. With
\[
 A=H_D^\circ-E_D,\quad B=H_{D^c}^\circ-E_{D^c},\quad
 V_\partial=\sum_{p\in\partial_\square D}V_p,
 \tag{V86.5}\label{r86:AB}
\]
$A,B$ are nonnegative and strongly commute on separate factors. Their
sum has a product zero vector and
$H_{\rm raw}=E_D+E_{D^c}+A+B+V_\partial$. Hence
\[
 0\le E_0-E_D-E_{D^c}\le 2b k_D,\qquad
 \|(H_D^\circ-E_D)\Omega\|\le2b k_D.
 \tag{V86.6}\label{r86:internal}
\]
These are statements on the full product Haar space. Each cut Hamiltonian
is gauge invariant, and the unique positive full ground state is gauge
invariant. Thus the resulting source vectors and estimates restrict to
the physical space. A tensor product decomposition of the Gauss-law
subspace is neither asserted nor used.

Retain V85's energy partition
$h_e=\kappa L_e+\frac14\sum_{p\ni e}V_p$ and set
\[
 F_D=\sum_{e\in D}h_e
     =H_D^\circ+\sum_{p\in\partial_\square D}r_p V_p,
 \qquad r_p=|p\cap D|/4.
 \tag{V86.7}\label{r86:FD}
\]
The coefficients $r_p$ on crossing plaquettes belong to
$\{1/4,1/2,3/4\}$. Put $R_D=\sum_{p\in\partial_\square D}r_p$.

\begin{theorem}[Unfiltered cut covariance]
\label{r86:block}
For every link set $D$, uniformly in $N$ and $\kappa>0$,
\[
 \begin{split}
 \Phi_D&=(F_D-\langle F_D\rangle)\Omega,\\
 \|\Phi_D\|&\le 2b k_D+b\min(R_D,k_D-R_D)
                 \le\tfrac52 b k_D,\\
 \Var_\Omega(F_D)&\le\tfrac{25}{4}b^2 k_D^2.
 \end{split}
 \tag{V86.8}\label{r86:block-bound}
\]
Here the variance is the actual quantum ground-state variance, not a
conditional or a reference-field variance. In particular the bound
includes all low-energy spectral weight:
\[
 \|\mathbf1_{(0,\eta]}(H)\Phi_D\|^2
 \le\|\Phi_D\|^2\le\tfrac{25}{4}b^2k_D^2
 \qquad (\eta>0).
 \tag{V86.9}\label{r86:unweighted-window}
\]
\end{theorem}

\begin{proof}
Apply Lemma~\ref{r86:abstract-cut} with $M=2bk_D$,
$W_A=\sum r_pV_p$, $W_B=\sum(1-r_p)V_p$,
$M_A=2bR_D$, and $M_B=2b(k_D-R_D)$. Centering removes the scalar $E_D$.
The last inequality uses $\min(R_D,k_D-R_D)\le k_D/2$.
The spectral assertion is the contraction property of an orthogonal
projection; no inverse power of $H$ or positive lower spectral bound is
being inserted. If $k_D=0$, the lemma gives $\Phi_D=0$ exactly.
\end{proof}

The right side of \eqref{r86:unweighted-window} does not tend to zero
with $\eta$. Thus it controls total low-energy weight but does not
exclude any interval above zero. For two blocks and any $t\ge0$, contraction
of $e^{-tH}$ gives only
\[
 |\langle\Phi_D,e^{-tH}\Phi_{D'}\rangle|
       \le\tfrac{25}{4}b^2 k_D k_{D'},
 \tag{V86.10}\label{r86:cross}
\]
with no distance decay claimed.

\subsection{Level sets and signed spatial weights}

For real link weights $w$, define the plaquette total variation
\[
 \mathcal T_\square(w)
   =\sum_p\left(\max_{e\in p}w_e-\min_{e\in p}w_e\right).
 \tag{V86.11}\label{r86:TV}
\]
As in V85, $F_w=\sum_e w_eh_e$ and $\Phi_w=QF_w\Omega$.

\begin{corollary}[Unweighted spatial variation bound]
\label{r86:coarea}
For arbitrary real weights, without a restriction on their support size,
\[
       \|\Phi_w\|\le\tfrac52 b\mathcal T_\square(w).
 \tag{V86.12}\label{r86:TV-bound}
\]
For complex $w=u+iv$ with real $u,v$,
\[
 \|\Phi_w\|^2\le\tfrac{25}{4}b^2
       \bigl(\mathcal T_\square(u)^2+\mathcal T_\square(v)^2\bigr).
 \tag{V86.13}\label{r86:complex}
\]
\end{corollary}

\begin{proof}
Let $a=\min_e w_e$, $c=\max_e w_e$ and $D_t=\{e:w_e>t\}$.
The finite-dimensional layer-cake identity and total-energy conservation
give
\[
 w=a\mathbf1+\int_a^c\mathbf1_{D_t}\,dt,
 \qquad \Phi_w=\int_a^c\Phi_{D_t}\,dt,
 \qquad \Phi_{\mathbf1}=0.
\]
These integrals are finite sums of constant vectors over intervals.
For an individual plaquette, the set of thresholds at which it crosses
$D_t$ has length $\max_{e\in p}w_e-\min_{e\in p}w_e$. Therefore
\[
       \int_a^c k_{D_t}\,dt=\mathcal T_\square(w).
\]
Integrate the norm bound \eqref{r86:block-bound} and use the triangle
inequality. For complex weights, the real ground state and real
differential operators make $\Phi_u,\Phi_v$ real vectors. Their scalar
product is real, so
$\|\Phi_u+i\Phi_v\|^2=\|\Phi_u\|^2+\|\Phi_v\|^2$.
This proves \eqref{r86:complex} without treating complex level sets.
\end{proof}

\subsection{Cubic block averages and the remaining bank-size loss}

Let $1\le\ell\le N-1$ and take the nonwrapping coordinate cube
$C_\ell=\{0,\ldots,\ell-1\}^3$ in the torus. Let $D_\ell$ consist of
the three positively oriented links \emph{anchored} at each vertex of
$C_\ell$; thus $|D_\ell|=3\ell^3$. This convention includes some links
exiting the vertex cube and removes any ambiguity about the boundary.
For a plaquette in the $ij$ plane anchored at $x$, its four link anchors
are $x,x+e_i,x+e_j$, with $x$ counted twice. Consequently
\[
                   k_{D_\ell}=12\ell^2-3\ell\le12\ell^2.
 \tag{V86.14}\label{r86:cube-cut}
\]
Indeed, for each of the three coordinate pairs the union of the three
translated cubes has size $\ell^3+2\ell^2$, while their intersection has
size $\ell(\ell-1)^2$. The difference is $4\ell^2-\ell$.
These counts also hold for $\ell=N-1$; the one added cyclic boundary
layer remains disjoint from the original cube in its shifted coordinate.

Translation and cubic symmetry of the unique ground state give
$\langle h_e\rangle=E_0/(3N^3)$. Theorem~\ref{r86:block} yields
\[
 \left\|\left(\frac{F_{D_\ell}}{3\ell^3}
                 -\frac{E_0}{3N^3}\right)\Omega\right\|^2
                 \le\frac{100b^2}{\ell^2}.
 \tag{V86.15}\label{r86:average}
\]
Thus at fixed spatial cutoff and coupling $b$, energy per link averaged
over a large block has vanishing variance uniformly in the surrounding
torus. This is stronger than applying V85's local vector bound separately
and paying the whole block volume. By the spectral theorem for the
self-adjoint observable $F_{D_\ell}/(3\ell^3)$, the corresponding
measurement probability of a deviation exceeding $s>0$ is at most
$100b^2/(s^2\ell^2)$. This measurement distribution is not a classical
joint law of all noncommuting local energies.

For the unit-$\ell^2$ bank
$w=\mathbf1_{D_\ell}/\sqrt{3\ell^3}$ the \emph{same} estimate is only
\[
                       \|\Phi_w\|^2\le300b^2\ell.
 \tag{V86.16}\label{r86:bank-loss}
\]
The linear loss has not disappeared. In particular,
\eqref{r86:average} must not be advertised as a volume-uniform bound on
the operator norm of the complete unweighted Gram matrix. Nor is it
exponential clustering, a bound on the coherent $B_J$ remainder, or a
renormalized continuum energy-density estimate. The spatial refinement
$b=c/(g^2a_s)$ remains to be handled separately.

\subsection{Compatibility with arbitrarily soft spectra}

The cut bound still cannot pay V85's missing inequality
$m\mathsf G_0\preceq\mathsf M$. This can be checked while retaining a
stronger unweighted bank estimate than the one just proved. Use V85's
plaquette incidence matrix $\mathsf L_\square$, with
$0\preceq\mathsf L_\square\preceq4I$, kernel the constant vectors, and
nonzero spectral bottom at most $4\pi^2/(3N^2)$. For fixed $\kappa,b>0$
choose
\[
 \alpha=\min\{b,\tfrac12\sqrt{C_{\kappa,b}}\}>0,\qquad
 \mathcal K_N=\mathbb C\Omega\oplus\mathbf1^\perp,
 \quad H_N=0\oplus\mathsf L_\square|_{\mathbf1^\perp},
 \quad \Phi_w=0\oplus\alpha\mathsf L_\square^{1/2}w.
 \tag{V86.17}\label{r86:soft}
\]
Here $C_{\kappa,b}$ is V85's explicitly paid constant. Then
\[
 \mathsf G_0=\alpha^2\mathsf L_\square\preceq4b^2I,
 \qquad
 \mathsf M=\alpha^2\mathsf L_\square^2
                  \preceq C_{\kappa,b}\mathsf L_\square.
 \tag{V86.18}\label{r86:soft-Grams}
\]
Both are finite range and annihilate constants. Moreover
\[
 \|\Phi_D\|^2
   =\alpha^2\sum_{p\in\partial_\square D}
           \frac{|p\cap D|(4-|p\cap D|)}4
   \le b^2 k_D\le b^2 k_D^2
 \tag{V86.19}\label{r86:soft-cuts}
\]
when $k_D\ge1$, and is zero otherwise. The level-set argument gives the
stronger $\|\Phi_w\|\le b\mathcal T_\square(w)$ for real $w$.
Thus this abstract model satisfies all the new cut/variation estimates,
has a uniform complete unweighted Gram bound, and its source vectors span
the full excited space at every finite $N$. Nevertheless its gap tends
to zero. For an eigenvector of $\mathsf L_\square$ of eigenvalue $\lambda$,
the ratio of its two Gram forms is exactly $\lambda$.
As in V85, bounded self-adjoint source operators can realize these vectors;
no local tensor-product realization or Wilson identification is claimed.
This is a nonimplication test, not a counterexample to Yang--Mills.

\subsection{Verification and what to attack next}

The checker verifies the cut lemma on finite tensor Hamiltonians with
noncommuting positive boundary interactions, including splits for which
neither side's internal Hamiltonian commutes with the cut. It separately
checks the exact plaquette fractions in a bounded algebra fixture,
all link subsets in that fixture, cubic crossing counts, layer-cake
identities for signed weights, and the strengthened soft-spectrum example.
These finite tests are not substitutes for the above operator proofs and
are not simulations of Haar $\SU(2)$. V85's mathematical regression is
also rerun. The prior source is preserved byte-for-byte except its version
title; build and diagnostic files stay outside \texttt{latex\_notes}.

The next useful gain would be a connected \emph{quadratic} boundary
estimate, instead of squaring a sum of boundary interaction norms. It
would require cancellation or decorrelation under the actual ground state;
boundedness of each plaquette alone does not supply it. Even such a gain
would still have to be followed by a genuine lower spectral estimate and
control of all physical sectors and the continuum limit. Reapplying
Cauchy--Schwarz to \eqref{r86:block-bound}, or changing the spatial
normalization, does not meet those obligations. The full interacting
four-dimensional YM mass-gap goal remains active and unproved.
% END CONSOLIDATED V86 APPEND

% BEGIN CONSOLIDATED V87 APPEND
\clearpage
\section{V87: A complete cut-source bank and its first neutral return}
\label{r87:section}

\subsection{What is inherited and what is gained}

V86 controls unweighted block-energy fluctuations by a squared boundary
size in the fully coupled vacuum. To improve that loss, this continuation
goes upstream to the operator which restores the crossing plaquettes.
The two interiors remain fully interacting; only their crossing
plaquettes are removed at the reference point.

Charged-sector domination is \emph{not} a new discovery here: f15 V82--V83
already decomposes the cut and proves native transfer bounds relative to
the regional Perron eigenvalues. The ground-state representation is in
f1 V6. We give its short differential-Hamiltonian consequence to specify
the normalization used below. The new estimates concern the
\emph{complete crossing-plaquette source operator}, including arbitrary
neutral input vectors, its charge-pattern multiplicities, its inverse-energy
Gram, and the first response of the actual block-energy fluctuation.
These are not inferred merely from a charged-sector eigenvalue.

There are two different conclusions. At the decoupled cut, an exact
orthogonality argument pays a whole signed source bank with no boundary
cardinality loss in its operator norm. At arbitrary restoring coupling,
integration by parts pays a quadratic binding-energy bound, but not a
variance or a mass gap. A neutral-return calculation below identifies why
the first conclusion cannot simply be extended to the physical endpoint.
The next companion and broad-literature checkpoints remain V90.

\subsection{The cut Hamiltonian and the enlarged symmetry}

Keep the cubic torus $N\ge3$, unit-round $\SU(2)$ metric and $\kappa>0$.
For a link partition $D,D^c$, write $\mathcal X=\partial_\square D$ for
the crossing plaquettes and $k=|\mathcal X|$. Put
\[
 K_0=H_D^\circ+H_{D^c}^\circ-E_D-E_{D^c}=A+B,\qquad
 \Omega_0=\Omega_D\otimes\Omega_{D^c}.
 \tag{V87.1}\label{r87:cut-H}
\]
The internal plaquette couplings may be arbitrary fixed nonnegative
numbers; homogeneous coupling $b$ is the case needed for V86.
Here $A,B\ge0$ act on separate full Haar factors and strongly commute.
Their positive smooth normalized ground states give the unique ground
state $\Omega_0$ of $K_0$. No factorization of the physical Gauss-law
space is used.

At a vertex $x$, let $d_R(x)$ be the number of incident links in
$R=D,D^c$. When $d_R(x)>0$, partial gauge transformation at $(x,R)$
acts only on those links. Let $G^R_{x,a}$, $a=1,2,3$, be its
skew-adjoint infinitesimal generators, normalized so that the fundamental
Casimir is $3$. Define
\[
 C_x^R=-\sum_{a=1}^3(G^R_{x,a})^2,\qquad
 \mathcal C=\sum_{R=D,D^c}\sum_{x:d_R(x)>0}
                         \frac{C_x^R}{d_R(x)}.
 \tag{V87.2}\label{r87:C}
\]
All these gauge actions commute with $K_0$. The two partial actions at a
vertex act on disjoint link sets; actions at different vertices commute
through left/right multiplication. Denote the orthogonal projection onto
their joint invariant space by $P$. This neutral space contains
arbitrarily many interior loop excitations, not only $\Omega_0$.

For every smooth $u$, Cauchy--Schwarz on a star and then counting each
link's two endpoints give the pointwise inequality
\[
 \Gamma_{\mathcal C}(u):=
 \sum_{R,x}\frac{\sum_a|G^R_{x,a}u|^2}{d_R(x)}
          \le 2\sum_e|\nabla_e u|^2.
 \tag{V87.3}\label{r87:star}
\]
Since every partial generator annihilates $\Omega_0$, the inherited
ground-state representation, applied to $f=\Omega_0u$, gives
\[
 \langle f,K_0f\rangle=\kappa\int\Omega_0^2\sum_e|\nabla_e u|^2,
 \qquad
 \langle f,\mathcal Cf\rangle=\int\Omega_0^2\Gamma_{\mathcal C}(u),
 \qquad K_0\succeq\frac{\kappa}{2}\mathcal C.
 \tag{V87.4}\label{r87:charge-form}
\]
The ground-state identity follows by integrating the product Laplacian
in $K_0(\Omega_0u)$ and using $K_0\Omega_0=0$; it holds for arbitrary
smooth complex $u$, before Gauss restriction. Smoothness, strict
positivity on the compact configuration space and form closure justify
the displayed form inequality on its full domain.

\subsection{Crossing-plaquette signatures and their small multiplicity}

Write $w_p=\frac12\operatorname{ReTr}U_p$ and $V_p=b(1-w_p)$ for
$p\in\mathcal X$, with $b\ge0$. Define $\sigma(p)$ to be the set of
vertices of $p$ at which exactly one of its two incident boundary links
belongs to $D$. It has either two or four vertices. At each such vertex,
$w_p$ is a fundamental matrix coefficient under \emph{each} of the two
partial gauge actions; at every other vertex it is invariant. Thus
\[
 C_x^R(w_pf)=
 \begin{cases}3w_pf,&x\in\sigma(p),\\0,&x\notin\sigma(p),\end{cases}
 \quad f\in P\mathcal H,\qquad
 \mathcal C(w_pf)=\lambda_p w_pf,
 \tag{V87.5}\label{r87:signature}
\]
initially on smooth invariant $f$, where
\[
 \lambda_p=3\sum_{x\in\sigma(p)}
       \left(\frac1{d_D(x)}+\frac1{d_{D^c}(x)}\right)\ge4.
 \tag{V87.6}\label{r87:lambda}
\]
Indeed $d_D+d_{D^c}=6$ and both degrees are positive at a transition;
their reciprocal sum is at least $2/3$. There are at least two transitions.
Equation \eqref{r87:signature} can also be seen by inserting a gauge
matrix into the single affected factor of the plaquette word. When both
adjacent links are on the same side, the two inserted matrices cancel.

The number of plaquettes with any one signature is at most four.
For four transition vertices the vertex set specifies the plaquette.
For two opposite transition vertices the diagonal pair specifies it.
For two adjacent transition vertices, every possible plaquette contains
their common link, and a cubic link belongs to four plaquettes.
These elementary incidence statements hold for $N\ge3$.
Moreover
\[
 r_p:=|p\cap D|/4
 \quad\hbox{is constant on every signature class, and}\quad
 r_p\in\{1/4,1/2,3/4\}.
 \tag{V87.7}\label{r87:rclass}
\]
For adjacent transitions the common link's membership in $D$ decides
between $1/4$ and $3/4$ for the whole class. Opposite transitions or four
transitions give $r_p=1/2$.

\begin{lemma}[Exact partial-gauge averages]
\label{r87:averages}
On the neutral space,
\[
 Pw_pP=0,\qquad Pw_p^2P=\tfrac14P,\qquad
 Pw_pw_qP=0\quad\hbox{if }\sigma(p)\ne\sigma(q).
 \tag{V87.8}\label{r87:Haar}
\]
These are operator identities, independent of the interior vacuum law.
\end{lemma}

\begin{proof}
Average over one partial gauge group at a transition vertex.
The affected holonomy is a left/right translate of a Haar $\SU(2)$
matrix or its inverse. Its normalized real trace has mean zero and
second moment $1/4$, the second moment of a coordinate of the uniform
unit three-sphere. This gives the first two identities after full gauge
averaging. If signatures differ, at a vertex in their symmetric difference
the partial central transformation $-I$ changes the sign of exactly one
factor in $w_pw_q$ while fixing every neutral input. Averaging gives zero.
\end{proof}

\subsection{A complete operator-valued source bank}

Let $\mathcal S$ be the direct sum of joint gauge-isotypic sectors with
fundamental labels on both sides at the vertices of one of the signatures
above, and trivial labels elsewhere. It is a reducing space for $K_0$,
orthogonal to $P\mathcal H$, and \eqref{r87:charge-form} implies
\[
                         K_0|_{\mathcal S}\succeq2\kappa I.
 \tag{V87.9}\label{r87:source-gap}
\]
This is a cut-source spectral bound. It says nothing about the bottom
positive spectrum of $K_0$ in its neutral sector.

\begin{theorem}[No boundary-size loss in the complete source operator]
\label{r87:bank}
Define
\[
 \mathcal W:\bigoplus_{p\in\mathcal X}P\mathcal H\longrightarrow\mathcal S,
       \qquad \mathcal W(\xi_p)_p=-b\sum_{p\in\mathcal X}w_p\xi_p.
\]
Then, for $t\ge0$ and real $z<2\kappa$,
\[
 \begin{split}
 \mathcal W^*\mathcal W&\preceq b^2 I,\\
 0\preceq\mathcal W^*e^{-tK_0}\mathcal W
                         &\preceq b^2e^{-2\kappa t}I,\\
 0\preceq\mathcal W^*(K_0|_{\mathcal S}-z)^{-1}\mathcal W
                         &\preceq\frac{b^2}{2\kappa-z}I.
 \end{split}
 \tag{V87.10}\label{r87:operator-bank}
\]
Every displayed operator matrix has zero $(p,q)$ entry when
$\sigma(p)\ne\sigma(q)$.
\end{theorem}

\begin{proof}
Different signature classes are orthogonal, even for different neutral
input vectors. Inside a class $\mathcal A$ of size at most four,
Lemma~\ref{r87:averages} gives
\[
 \left\|\sum_{p\in\mathcal A}w_p\xi_p\right\|^2
 \le|\mathcal A|\sum_{p\in\mathcal A}\|w_p\xi_p\|^2
 =\frac{|\mathcal A|}{4}\sum_{p\in\mathcal A}\|\xi_p\|^2
 \le\sum_{p\in\mathcal A}\|\xi_p\|^2.
\]
Sum over classes. The spectral theorem on $\mathcal S$ gives the other
inequalities. $K_0$ and its functions preserve all partial gauge labels,
which also proves the matrix zero statements.
\end{proof}

In particular, for scalar weights $s_p$ and
$W_s=-b\sum_p s_pw_p$, the map $W_sP$ has norm at most $b\|s\|_2$.
Taking every input proportional to $\Omega_0$ proves a complete unweighted
vacuum Gram bound, with diagonal $b^2/4$, and an inverse-energy Gram bound
$b^2/(2\kappa)$. Taking \emph{arbitrary} neutral inputs instead gives
\[
 0\preceq PW_s(K_0|_{\mathcal S}-z)^{-1}W_sP
       \preceq\frac{b^2\|s\|_2^2}{2\kappa-z}P
 \quad(s\hbox{ real}).
 \tag{V87.11}\label{r87:return-bound}
\]
This is the unperturbed second-order return operator, not the fully
dressed Schur complement. Finite range in its plaquette labels does not
prove that each entry is local in all the interior link variables:
the interior Hamiltonian remains inside its resolvent.

\subsection{The actual first response of block-energy fluctuations}

For real $s$, restore the cut along the family
\[
 \mathcal H(t)=H_D^\circ+H_{D^c}^\circ
                     +t\sum_{p\in\mathcal X}s_pV_p,
 \qquad
 F_D(t)=H_D^\circ+t\sum_{p\in\mathcal X}r_ps_pV_p.
 \tag{V87.12}\label{r87:family}
\]
Let $E(t),\Omega(t)$ be its bottom and normalized positive ground vector,
and $\Phi_D(t)=(F_D(t)-\langle F_D(t)\rangle_t)\Omega(t)$.
For fixed regulator these objects are analytic near zero: the
perturbation is bounded, the elliptic operator has compact resolvent, and
its ground state is simple. No uniform analytic radius is asserted.
At zero $\Phi_D(0)=0$.

\begin{theorem}[Boundary-linear quadratic fluctuation coefficient]
\label{r87:response}
With derivatives evaluated at $t=0$,
\[
 \begin{aligned}
 E'(0)&=b\sum_p s_p,&
 0\le -E''(0)&\le\frac{b^2}{\kappa}\|s\|_2^2,\\
 \|\Omega'(0)\|&\le\frac{b}{2\kappa}\|s\|_2,&
 \|\Phi_D'(0)\|&\le\frac{3b}{4}\|s\|_2.
 \end{aligned}
 \tag{V87.13}\label{r87:response-bounds}
\]
Consequently, for the actual restoring path $s_p=1$,
\[
 \lim_{t\to0}\frac{\Var_{\Omega(t)}F_D(t)}{t^2}
                      \le\frac9{16}b^2k.
 \tag{V87.14}\label{r87:quadratic}
\]
\end{theorem}

\begin{proof}
Put $Q_0=I-|\Omega_0\rangle\langle\Omega_0|$ and
$R=Q_0(K_0|_{\Omega_0^\perp})^{-1}Q_0$, at the fixed regulator.
Let $\phi_s=W_s\Omega_0$. It is in $\mathcal S$.
Differentiating the normalized eigenvalue equation gives
\[
 \Omega'(0)=-R\phi_s,\qquad
 E''(0)=-2\langle\phi_s,R\phi_s\rangle.
\]
The scalar first derivative follows from \eqref{r87:Haar}; the first
three estimates follow from Theorem~\ref{r87:bank}.

The derivative of the centered regional source is
\[
             \Phi_D'(0)=-AR\phi_s+\phi_{rs}.
 \tag{V87.15}\label{r87:regional-derivative}
\]
Here the derivative of its mean is $b\sum r_ps_p$, since
$H_D^\circ\Omega_0=E_D\Omega_0$ and
$\langle\Omega_0,\Omega'(0)\rangle=0$.
In the joint spectral calculus of $A,B$, $AR$ has multiplier
$a/(a+b)$ off the joint vacuum, hence $0\preceq AR\preceq I$.
It preserves all signature classes. By \eqref{r87:rclass}, on each class
the right side of \eqref{r87:regional-derivative} is
$(r_{\mathcal A}I-AR)\phi_{s,\mathcal A}$, of norm at most
$\max(r_{\mathcal A},1-r_{\mathcal A})\|\phi_{s,\mathcal A}\|
\le(3/4)\|\phi_{s,\mathcal A}\|$.
Sum squared norms of the orthogonal classes and use
$\|\phi_s\|\le b\|s\|_2$. Finally
$\Phi_D(t)=t\Phi_D'(0)+O(t^2)$ in norm at the fixed regulator.
This proves \eqref{r87:quadratic}, not a uniform bound on its remainder.
\end{proof}

Thus a connected, boundary-linear variance coefficient is paid without a
neutral interior gap. It is a coefficient at the decoupled cut, not
V86's endpoint variance under $\Omega(1)$.

\subsection{A nonperturbative quadratic binding-energy estimate}

The partial gauge directions also give a statement valid at the fully
restored endpoint, although for energy rather than covariance.
For real $s$ let
$a_s(U)=\sum_p s_pw_p(U)/\lambda_p$.
At a transition vertex each fundamental height has
$\sum_a|G^R_{x,a}w_p|^2=1-w_p^2\le1$; otherwise this is zero.
At most twelve cubic plaquettes contain any fixed vertex.
Cauchy--Schwarz in each derivative group therefore gives
\[
 \begin{split}
 \Gamma_{\mathcal C}(a_s)
 &\le12\sum_p\frac{s_p^2}{\lambda_p^2}
                      \Gamma_{\mathcal C}(w_p)\\
 &\le4\sum_p\frac{s_p^2}{\lambda_p}\le\|s\|_2^2,
 \qquad \mathcal Ca_s=\sum_p s_pw_p.
 \end{split}
 \tag{V87.16}\label{r87:carre}
\]
Integration by parts against product Haar measure, followed by
Cauchy--Schwarz, yields for every smooth complex $f$
\[
 |\langle f,W_sf\rangle|
 \le2b\|s\|_2\|f\|\sqrt{\langle f,\mathcal Cf\rangle}
 \le2\sqrt{2/\kappa}\,b\|s\|_2\|f\|
                       \sqrt{\langle f,K_0f\rangle}.
 \tag{V87.17}\label{r87:relative}
\]
Indeed $G|f|^2=2\operatorname{Re}(\bar fGf)$, and
\eqref{r87:carre} bounds the other gradient factor pointwise.
Form closure extends this inequality beyond smooth $f$.

For normalized $f$, write $h=\langle f,K_0f\rangle\ge0$.
Completing the square in
$h-2\sqrt{2/\kappa}\,b|t|\|s\|_2\sqrt h$ gives
\[
 -\frac{2b^2t^2}{\kappa}\|s\|_2^2
 \le E(t)-E_D-E_{D^c}-tb\sum_p s_p\le0
 \qquad(t\in\mathbb R).
 \tag{V87.18}\label{r87:binding}
\]
The upper bound is the trial $\Omega_0$. The same argument for
$f=\Omega(t)$ gives
\[
 \langle K_0\rangle_t\le
             \frac{8b^2t^2}{\kappa}\|s\|_2^2,\qquad
 \langle\mathcal C\rangle_t\le
             \frac{16b^2t^2}{\kappa^2}\|s\|_2^2.
 \tag{V87.19}\label{r87:charge-moment}
\]
For $s_p=1$ and $t\ge0$, positivity of the restoring plaquette potential
adds $E(t)\ge E_D+E_{D^c}$, while
$E(t)\le E_D+E_{D^c}+tbk$ improves the crude cut trial in V86.
The quadratic bound has no power $k^2$ in its binding-energy debit.
It does not bound the curvature of $E$ at every $t$, and must not be
differentiated as though it did. Its constants also deteriorate when
$b/\kappa$ grows; it is not an ultraviolet-uniform estimate.

\subsection{The neutral return that survives these bounds}

Already at multiplication level the first return is not confined to the
vacuum. Take $D$ to be one link $e$ and $p\ne q$ two incident plaquettes.
Orient their traces as $w_p=\frac12\operatorname{ReTr}(U_eS_p)$ and
$w_q=\frac12\operatorname{ReTr}(U_eS_q)$, where $S_p,S_q$ are the three-link
complementary staples. A Haar average of $U_e$ gives exactly
\[
       P(w_pw_q)=\frac14\,w_{pq},\qquad
       w_{pq}=\frac12\operatorname{ReTr}(S_pS_q^{-1}).
 \tag{V87.20}\label{r87:neutral-loop}
\]
The loop on the right lies entirely in $D^c$, so it is already neutral
under all the partial gauge transformations. It is nonconstant.
Since $\Omega_0$ is strictly positive, its centered multiplication on
$\Omega_0$ is a nonzero neutral excited vector for arbitrary finite
interior couplings. In contrast $P(w_p^2)=1/4$ is constant.

With all internal plaquette couplings zero, the statement remains true
with the charged resolvent inserted. Now $\Omega_0=1$ and
$K_0w_q=12\kappa w_q$, one fundamental Casimir $3$ on each of four links.
Thus
\[
        Pw_pRw_q\Omega_0=\frac1{48\kappa}w_{pq},
 \tag{V87.21}\label{r87:resolved-loop}
\]
a nonzero nonconstant neutral vector. This is an exact compact-group
example, not a Gaussian or qubit approximation. It does not assert the
same coefficient in the interacting interiors.

More generally, if $s$ is fixed and the scalar $b\sum s_p$ has been
removed from the Hamiltonian, differentiating once more gives
\[
 Q_0\Omega''(0)=2R\,W_sR\,W_s\Omega_0.
 \tag{V87.22}\label{r87:second}
\]
The first resolvent from the right acts on the paid charged sector.
After the next multiplication, the left resolvent can act on a neutral
excitation such as \eqref{r87:neutral-loop}. Its spectral denominator
is not bounded by \eqref{r87:source-gap}.
The selection rule still kills unequal signatures, so this is a
same-signature, nearby-plaquette return, not a sum over arbitrary distant
pairs. But locality of those labels does not pay the neutral resolvent.

This identifies the next substantive task: control the connected neutral
return in the interacting interiors, or an equivalent relative form
estimate after that return. Neither the cut-source gap nor the global
binding-energy bound supplies it. The restored Hamiltonian at $t=1$
breaks the independent partial gauge symmetries on which the exact
orthogonality depended. Accordingly the full physical mass gap and the
interacting four-dimensional continuum limit remain unproved.

\subsection{Verification and preservation}

The diagnostics enumerate crossing signatures on several cubic tori and
cuts, including one-link cuts realizing multiplicity four, and verify
the class-constant fractions, charge lower bound and star-incidence
constants. Exact rational quaternion Haar calculations check the
fundamental means, second moments and staple contraction.
Polynomial Casimir calculations check the free charged denominator.
Finite commuting-interior algebra fixtures independently check the
regional derivative and its $3/4$ bound; they are explicitly not Haar
gauge theories. These tests supplement the analytic operator proofs,
not certify them by finite sampling. V86's mathematical checks are
rerun. The entire predecessor is preserved apart from its version title,
and only the latest active main source is kept in \texttt{latex\_notes}.
% END CONSOLIDATED V87 APPEND

% BEGIN CONSOLIDATED V88 APPEND
\clearpage
\section{V88: A uniformly gapped magnetic rotor and the cost of adapting its fibre}
\label{r88:section}

\subsection{Why change the reference projection?}

V87 pays the charged cut-source resolvent, but its second return creates
nonconstant neutral staple-pair loops. Repeating the charged estimate
does not pay the neutral denominator. This continuation therefore absorbs
the four restoring plaquettes around one link into its reference
Hamiltonian, with the exterior left as a variable. We prove a uniform
gap for this \emph{magnetic one-link differential Hamiltonian}, and then
keep the cost of its exterior-dependent ground projection exactly.

This is distinct from the auxiliary root-field rotor law in V46, the
free one-link heat-channel estimate in f6 V72, and the charged transfer
decompositions in f15 V82--V83. None of those results is claimed anew.
The one-link ground state below is \emph{not} identified with a
conditional distribution of the actual full vacuum. Its usefulness is
an exact operator decomposition, not that identification.
The companion audit at V85 remains in force; the next companion and
broad-literature checkpoints are V90.

\subsection{A direct one-dimensional gap lemma}

We first record the elementary factorization estimate used in the rotor.
The constant is deliberately nonsharp.

\begin{lemma}[Monotone Dirichlet potential]\label{r88:monotone-lemma}
Let $T=-\kappa\partial_\theta^2+V(\theta)$ on $(0,L)$ with Dirichlet
boundary conditions, $\kappa>0$, and real $V\in C^1([0,L])$ satisfying
$V'\ge0$. If $\mu_0<\mu_1$ are its two lowest eigenvalues, then
\[
 \mu_1-\mu_0\ge\frac{\kappa\pi^2}{L^2}.
 \tag{V88.1}\label{r88:interval-gap}
\]
\end{lemma}

\begin{proof}
Let $u>0$ be the ground eigenfunction on $(0,L)$ and $p=u'/u$.
The equation and a differentiation give
\[
 p'=\frac{V-\mu_0}{\kappa}-p^2,\qquad
 (u^2p')'=\frac{u^2V'}{\kappa}\ge0.
 \tag{V88.2}\label{r88:riccati}
\]
At the right endpoint, $u^2p'=uu''-(u')^2$ tends to
$-(u'(L))^2<0$. The derivative is nonzero by uniqueness for the
second-order ODE with specified value and derivative. Monotonicity in
\eqref{r88:riccati} implies $p'<0$ throughout the open interval.

Set $D=\partial_\theta-p$, with formal adjoint
$D^*=-\partial_\theta-p$. Then
\[
 \frac{T-\mu_0}{\kappa}=D^*D,\qquad
 DD^*=-\partial_\theta^2+p^2-p'.
 \tag{V88.3}\label{r88:factor}
\]
For any Dirichlet eigenfunction $v$ with eigenvalue $\mu>\mu_0$,
$w=Dv$ is nonzero and satisfies
$DD^*w=(\mu-\mu_0)w/\kappa$ in the interior. Near $0$, the ODE gives
$u=c\theta+O(\theta^3)$ and $v=d\theta+O(\theta^3)$, whence
$w=O(\theta^2)$ and $w'=O(\theta)$; the corresponding assertions hold
near $L$. Consequently $w\in H^1_0(0,L)$, all terms below are integrable,
and integration on $[\epsilon,L-\epsilon]$ followed by
$\epsilon\downarrow0$ has no boundary remainder. Thus
\[
 \frac{\mu-\mu_0}{\kappa}\|w\|^2
 =\int_0^L\bigl(|w'|^2+(p^2-p')|w|^2\bigr)
 \ge\int_0^L|w'|^2
 \ge\frac{\pi^2}{L^2}\|w\|^2.
\]
This proves the lemma without assuming convexity of $V$.
\end{proof}

For literature orientation, one-dimensional factorization is classical;
the abstract of D.~J.~Fern\'andez C., \emph{Supersymmetric Quantum
Mechanics}, \href{https://arxiv.org/abs/0910.0192}{arXiv:0910.0192},
describes this general method. The quantitative assertion just needed
has been proved above, not imported from that abstract. The convex-potential
fundamental-gap theorem of Andrews--Clutterbuck,
\href{https://arxiv.org/abs/1006.1686}{arXiv:1006.1686}, is not being
applied: the potential $-h\cos\theta$ below is not convex on $(0,\pi)$.
This is a targeted scope check, not the scheduled V90 literature sweep.

\subsection{All angular sectors of the magnetic SU(2) rotor}

Identify $\SU(2)$ with the unit round $S^3$, with normalized Haar
measure, coordinate $q=(q_0,q_1,q_2,q_3)$, and
$L=-\Delta_{S^3}$. Its fundamental coordinate eigenvalue is $3$.
For $f\in\mathbb R^4$ define
\[
 h_f=\kappa L-f\cdot q,\qquad
 e(f)=\inf\operatorname{spec}h_f=e_\kappa(|f|).
 \tag{V88.4}\label{r88:rotor}
\]
The compact elliptic operator has a simple normalized strictly positive
real ground function $\psi_f$: the heat kernel is strictly positive,
and multiplication by a bounded real potential preserves positivity
improvement, for example by the Feynman--Kac formula.

\begin{theorem}[Uniform full-rotor gap]\label{r88:rotor-theorem}
With $p_f=|\psi_f\rangle\langle\psi_f|$, for every real $f$,
\[
 h_f-e(f)\succeq\kappa(I-p_f).
 \tag{V88.5}\label{r88:rotor-gap}
\]
This holds on all of $L^2(S^3)$, not only on class functions, and includes
$f=0$ without a limiting convention.
\end{theorem}

\begin{proof}
An orthogonal rotation of $S^3$ sends $f$ to $h(1,0,0,0)$, $h=|f|$.
Uniqueness and positivity make the ground function invariant under
rotations of the last three coordinates. Write
$q=(\cos\theta,\sin\theta\,\omega)$, $\omega\in S^2$, and expand
in spherical harmonics of angular degree $\ell$. Multiplying the radial
function by $\sin\theta$ converts the radial measure to Lebesgue
measure. The resulting Friedrichs operator is
\[
 T_\ell=-\kappa\partial_\theta^2-\kappa-h\cos\theta
       +\frac{\kappa\ell(\ell+1)}{\sin^2\theta},
 \qquad 0<\theta<\pi.
 \tag{V88.6}\label{r88:radial}
\]
For $\ell=0$ this has the usual Dirichlet domain. Its potential has
derivative $h\sin\theta\ge0$, so Lemma~\ref{r88:monotone-lemma}
pays a radial excitation gap at least $\kappa$. For $\ell\ge1$, the
centrifugal form domain is contained in the $\ell=0$ form domain, and
$\sin^{-2}\theta\ge1$ gives
\[
 \inf\operatorname{spec}T_\ell
 \ge e_\kappa(h)+\kappa\ell(\ell+1)
 \ge e_\kappa(h)+2\kappa.
\]
The full harmonic decomposition proves \eqref{r88:rotor-gap}.
At $h=0$ the actual gap is $3\kappa$; no sharp constant at nonzero
$h$ is asserted here.
\end{proof}

\subsection{A nonsingular projector even at cancelling staples}

Because $h_f$ is a common-domain operator with bounded affine
parameter dependence and an isolated simple eigenvalue, its local
Riesz projection is real analytic in $f$. The positive real normalized
choice of $\psi_f$ patches these local choices globally. In particular
there is no division by $|f|$ and no singular choice of axis at $f=0$.
Let $Q_f=I-p_f$ and
$R_f=Q_f(h_f-e(f))^{-1}Q_f$, with zero action on $p_f$.
For a real direction $v$, differentiating the eigenvalue equation and
normalization gives
\[
 \begin{aligned}
  \partial_v\psi_f&=R_f(v\cdot q)\psi_f,
       &\langle\psi_f,\partial_v\psi_f\rangle&=0,\\
  De(f)[v]&=-\langle\psi_f,(v\cdot q)\psi_f\rangle,
       &\|R_f\|&\le\kappa^{-1}.
 \end{aligned}
 \tag{V88.7}\label{r88:derivative}
\]
Here and below the fibre scalar product is in Haar $L^2(S^3)$.
Since the variance of a variable in $[-|v|,|v|]$ is at most $|v|^2$,
\[
 \begin{aligned}
  \|\partial_v\psi_f\|&\le |v|/\kappa,
  &\|\partial_vp_f\|&=\|\partial_v\psi_f\|\le |v|/\kappa,\\
  |De(f)[v]|&\le |v|,
  &0\le-D^2e(f)[v,v]&\le 2|v|^2/\kappa.
 \end{aligned}
 \tag{V88.8}\label{r88:smooth-bounds}
\]
For the projector norm, differentiate $p_f$ to obtain the two
off-diagonal rank-one terms between $\psi_f$ and
$\partial_v\psi_f\perp\psi_f$; their sum has norm
$\|\partial_v\psi_f\|$, not twice that value. For the last bound,
differentiate the middle expression in \eqref{r88:derivative}:
\[
 -D^2e(f)[v,v]
 =2\langle Q_f(v\cdot q)\psi_f,
                R_f Q_f(v\cdot q)\psi_f\rangle.
\]
The estimates are uniform for arbitrarily large or cancelling fields.

\subsection{An exact magnetic fibre inside the native differential Hamiltonian}

Fix a positive link $e$ of the cubic torus $N\ge3$; let $y$ denote all
other links. Keep $\kappa>0$ and arbitrary nonnegative plaquette
couplings $b_p$. Orient each of the four incident plaquette traces as
\[
 w_p=\tfrac12\operatorname{Re}\Tr(U_e S_p(y)),\qquad
 B_e=\sum_{p\ni e}b_p,\qquad
 f_e(y)=\sum_{p\ni e}b_p\,\overline{S_p(y)}.
 \tag{V88.9}\label{r88:staple-field}
\]
Here a unit quaternion and its coordinate vector are identified, and
the bar is quaternion conjugation. The choice follows from
$\operatorname{Re}(qS)=q\cdot\overline S$; reversing a plaquette
orientation does not change its real trace. The exact one-link operator
and the full, uncentered differential Hamiltonian are
\[
 \begin{aligned}
  h_e(y)&=\kappa L_e+B_e-f_e(y)\cdot q,&
  \lambda_e(y)&=B_e+e(f_e(y)),\\
  \mathcal H&=\kappa\sum_{j\ne e}L_j+V_{\not e}(y)+h_e(y),&
  V_{\not e}&=\sum_{p\not\ni e}b_p(1-w_p).
 \end{aligned}
 \tag{V88.10}\label{r88:full-H}
\]
This is the same fixed-spatial-regulator Hamiltonian as V85--V87,
with inhomogeneous couplings allowed. It is not a new identification of
the joint four-dimensional continuum limit.

Define the direct-integral projections
$\Pi_e(y)=p_{f_e(y)}$, $\mathcal Q_e=I-\Pi_e$, the exterior operator,
and an isometry from the exterior Haar space by
\[
 \begin{aligned}
  H_{\rm lo}&=\kappa\sum_{j\ne e}L_j+V_{\not e}+\lambda_e,\\
  (J_e\alpha)(q,y)&=\psi_{f_e(y)}(q)\alpha(y),
       &J_e^*J_e&=I,\quad J_eJ_e^*=\Pi_e.
 \end{aligned}
 \tag{V88.11}\label{r88:exterior}
\]
The lower operator is a scalar elliptic operator on the exterior, not
a ground-state projection and not a claimed gapped effective theory.

\begin{proposition}[Lower form and exact ground-fibre compression]
\label{r88:compression-prop}
As closed quadratic forms on the full Haar space,
\[
 \mathcal H\succeq I_e\otimes H_{\rm lo}+\kappa\mathcal Q_e.
 \tag{V88.12}\label{r88:form-lower}
\]
The exact quadratic-form compression under $J_e$ is
\[
 \begin{aligned}
 J_e^*\mathcal HJ_e&=H_{\rm lo}+V_{\rm BH},\\
 V_{\rm BH}(y)&=\kappa\sum_{j\ne e}\sum_{a=1}^3
       \|\partial_{j,a}\psi_{f_e(y)}\|^2,
       &0\le V_{\rm BH}&\le
       M_e:=\frac{9}{\kappa}\sum_{p\ni e}b_p^2.
 \end{aligned}
 \tag{V88.13}\label{r88:BH}
\]
The derivatives $\partial_{j,a}$ are an orthonormal invariant frame
for the unit-round link metric. In the homogeneous case $M_e=36b^2/\kappa$.
\end{proposition}

\begin{proof}
Apply \eqref{r88:rotor-gap} in each fibre and add the unchanged exterior
kinetic and potential forms. This proves \eqref{r88:form-lower}.
It uses neither commutation of $\mathcal Q_e$ with the exterior kinetic
energy nor squaring a noncommuting operator inequality.

The positive real phase and normalization give
$\langle\psi_{f_e},\partial_{j,a}\psi_{f_e}\rangle=0$. Thus the
Berry connection vanishes and, for smooth complex $\alpha$,
\[
 \|\partial_{j,a}(\psi_{f_e}\alpha)\|^2
 =|\partial_{j,a}\alpha|^2+
                    |\alpha|^2\|\partial_{j,a}\psi_{f_e}\|^2.
 \tag{V88.14}\label{r88:kinetic-product}
\]
Integrate this identity over the exterior and add the fibre energy
$\lambda_e|\alpha|^2$ to prove the compression formula. This is
the Born--Huang scalar term, derived here directly; dropping it would
not be an exact compression.

There are twelve distinct exterior links in the four three-link staples
when $N\ge3$. Each belongs to just one of these staples. An invariant
unit derivative of a unit quaternion factor has Euclidean norm one;
left/right multiplication and conjugation are isometries. Therefore,
for each $j$ in staple $p$ and each $a$,
$|\partial_{j,a}f_e|=b_p$, and exactly
\[
 \sum_{j\ne e,a}|\partial_{j,a}f_e|^2
       =9\sum_{p\ni e}b_p^2.
 \tag{V88.15}\label{r88:field-gradient}
\]
The chain rule and \eqref{r88:smooth-bounds} prove the bound in
\eqref{r88:BH}. Smoothness on the compact exterior and the derivative
bound extend the identities from smooth functions to the form domains.
\end{proof}

All this was done before the Gauss restriction. Under a lattice gauge
transformation, the transformed one-link operator is conjugated by the
unitary action $U_e\mapsto g_sU_eg_t^{-1}$, together with the transformed
exterior. Uniqueness of the positive normalized ground function makes
$J_e$ equivariant. Consequently gauge-invariant exterior amplitudes
map to physical full states; $\lambda_e$ and $V_{\rm BH}$ are gauge
invariant. The bounds restrict to those spaces without factoring the
physical Hilbert space into link factors.

\subsection{What is paid in the coupled vacuum, and what still mixes}

Let $E_{\rm lo}$ and $E_0$ be the lowest eigenvalues of $H_{\rm lo}$
and $\mathcal H$, and $\alpha_{\rm lo}$, $\Omega$ their positive
normalized ground functions. Lower form comparison and the variational
trial $J_e\alpha_{\rm lo}$ give
\[
 \begin{aligned}
 0\le E_0-E_{\rm lo}
   &\le \langle\alpha_{\rm lo},V_{\rm BH}\alpha_{\rm lo}\rangle
    \le M_e,\\
 \kappa\|\mathcal Q_e\Omega\|^2
   &\le E_0-E_{\rm lo}\le M_e.
 \end{aligned}
 \tag{V88.16}\label{r88:vacuum-cost}
\]
For the second line, $I_e\otimes H_{\rm lo}\succeq E_{\rm lo}I$
in \eqref{r88:form-lower}. In particular
$\|\mathcal Q_e\Omega\|^2\le\min\{1,M_e/\kappa\}$.
The expectation in the first line is in the scalar lower-operator
vacuum, not an asserted conditional law of $\Omega$.

The same argument pays a genuine complementary-block denominator in
a restricted coupling range. On the closed form domain in
$\operatorname{Ran}\mathcal Q_e$,
\[
 \mathcal Q_e(\mathcal H-E_0)\mathcal Q_e
       \succeq(\kappa-M_e)\mathcal Q_e.
 \tag{V88.17}\label{r88:fast-compression}
\]
Its Friedrichs inverse therefore has norm at most
$(\kappa-M_e)^{-1}$ if $M_e<\kappa$. For homogeneous coupling this
condition is $6b<\kappa$. Multiplication by $\mathcal Q_e(y)$ preserves
the form domain by the derivative bound, so the restricted form is
closed and densely defined in its range. Equation
\eqref{r88:fast-compression} concerns this \emph{compressed block},
not a reducing spectral subspace of the full Hamiltonian or the full gap.

For clarity the mixing that has not vanished can also be written exactly.
Let $\chi\in\operatorname{Ran}\mathcal Q_e$ be in the form domain,
and let $\mathcal S_e$ denote the twelve exterior staple links. Put
\[
 \|\nabla_{\mathcal S_e}\alpha\|^2
 =\sum_{j\in\mathcal S_e,a}\|\partial_{j,a}\alpha\|^2
\]
and use the analogous full-Haar norm for $\chi$. The off-diagonal
sesquilinear form $\mathfrak c_e(\alpha,\chi)
=\mathfrak h(J_e\alpha,\chi)$ has only exterior kinetic terms, since
$h_e\psi_{f_e}=\lambda_e\psi_{f_e}$ and $J_e^*\chi=0$. Differentiating
the latter identity gives, with fibre scalar products inside the integral,
\[
 \begin{aligned}
 \mathfrak c_e(\alpha,\chi)=\kappa\sum_{j\in\mathcal S_e,a}\int
 \bigl[&-(\partial_{j,a}\alpha)^*
          \langle\partial_{j,a}\psi_{f_e},\chi\rangle\\
       &+\alpha^*\langle\partial_{j,a}\psi_{f_e},
                                      \partial_{j,a}\chi\rangle\bigr]\,dy.
 \end{aligned}
 \tag{V88.18}\label{r88:offdiag}
\]
Terms outside $\mathcal S_e$ vanish. Cauchy--Schwarz, the chain rule,
and \eqref{r88:field-gradient} yield
\[
 |\mathfrak c_e(\alpha,\chi)|
 \le 3\Bigl(\sum_{p\ni e}b_p^2\Bigr)^{1/2}
 \bigl(\|\nabla_{\mathcal S_e}\alpha\|\,\|\chi\|
       +\|\alpha\|\,\|\nabla_{\mathcal S_e}\chi\|\bigr).
 \tag{V88.19}\label{r88:offdiag-bound}
\]
This explicit, volume-independent relative-form estimate does not
make its coefficient small at large $b/\kappa$. The decomposition has
not decoupled the exterior kinetic energy.

\subsection{The retained potential contains the neutral loops}

The dependence of $\lambda_e$ on the exterior does not erase V87's
neutral return. At small field it displays it explicitly. The even
analytic branch at zero has
\[
 e_\kappa(h)=-\frac{h^2}{12\kappa}+\mathcal R_\kappa(h),
 \qquad
 |\mathcal R_\kappa(h)|\le\frac{4h^4}{3\kappa^3}
 \quad(0\le h\le\kappa/2).
 \tag{V88.20}\label{r88:small-field}
\]
Indeed $q_0$ has Haar mean zero, second moment $1/4$ and free energy
$3\kappa$, so second-order perturbation gives $-1/(12\kappa)$.
The symmetry $q\mapsto-q$ makes the branch even. To justify the stated
remainder independently of formal perturbation, use the Riesz energy
circle $|\zeta|=3\kappa/2$. For complex parameter $z$ with
$|z|<3\kappa/2$, the resolvent of $\kappa L-zq_0$ exists on this
circle by the Neumann criterion; the interior Riesz subspace has rank
one. On $|z|=\kappa$, the corresponding eigenvalue lies inside the
circle and, by the resolvent-distance bound for the normal operator
$\kappa L$, within $|z|$ of its spectrum. All nonzero free eigenvalues
are at least $3\kappa$ away from zero, so this forces $|e(z)|\le\kappa$.
Cauchy's coefficient estimate and evenness bound the tail by
\[
 \kappa\sum_{n\ge2}(h/\kappa)^{2n}
   =\frac{h^4/\kappa^3}{1-h^2/\kappa^2}
   \le\frac{4h^4}{3\kappa^3}.
\]

Write $w_{pq}(y)=\tfrac12\operatorname{Re}\Tr(S_p(y)S_q(y)^{-1})$.
The exact quaternion identity
\[
 |f_e(y)|^2=\sum_{p\ni e}b_p^2
             +2\sum_{\substack{p<q\\p,q\ni e}}b_pb_qw_{pq}(y)
 \tag{V88.21}\label{r88:pair-field}
\]
and \eqref{r88:small-field} give, whenever $B_e\le\kappa/2$,
\[
 \lambda_e(y)=B_e-\frac{\sum_{p\ni e}b_p^2}{12\kappa}
   -\frac1{6\kappa}\sum_{\substack{p<q\\p,q\ni e}}b_pb_qw_{pq}(y)
   +\mathcal R_\kappa(|f_e(y)|).
 \tag{V88.22}\label{r88:retained-loops}
\]
The $w_{pq}$ are precisely the nonconstant neutral staple-pair loops
encountered in V87; here they remain in the exterior scalar potential.
For large field, the exact definition \eqref{r88:exterior}, not this
Taylor expansion, is used.

There is an important denominator distinction. In the frozen fibre,
the free inverse on $q_0$ is $(3\kappa)^{-1}$. In V87's free
\emph{whole cut}, $w_q$ excites four link Laplacians and its inverse
is $(12\kappa)^{-1}$, giving
$P w_pR w_q\Omega_0=w_{pq}/(48\kappa)$ there.
Thus \eqref{r88:retained-loops} is not that resolved return with a
renamed coefficient. The exterior kinetic terms, including
\eqref{r88:BH} and \eqref{r88:offdiag}, distinguish the two problems.
In particular the scalar compression is not an exact Feshbach operator
obtained by discarding its complementary block.

\subsection{Status, next obligation, and verification}

The new paid statements are a uniform full magnetic-rotor gap, a
globally regular adapted projector, an exact compression retaining its
derivative cost, and a complementary-block inverse in the explicit
range $M_e<\kappa$. All are proved without a bulk-gap assumption.
They do not prove neutral exterior coercivity. In the physical
weak-coupling scaling, $b/\kappa$ grows, while the displayed derivative
bound grows like $b^2/\kappa$; the sufficient inverse criterion is
therefore not a cofinal continuum estimate.

The next upstream obligation is to improve the actual weighted
derivative and off-diagonal return estimates, or prove coercivity in
the retained neutral exterior. Simply summing the uniform fibre gaps
would omit their noncommuting, exterior-dependent projections and
would not accomplish either task. The full interacting four-dimensional
continuum construction and physical Yang--Mills mass gap remain unproved.

Exact diagnostics check the radial conjugation and factorization
identities, the quaternion staple-pair identity, the twelve-link
incidence count and the derivative coefficient in
\eqref{r88:field-gradient}. Finite Gegenbauer truncations check the
radial and first angular spectra at several field strengths, with two
truncation sizes compared; these are consistency tests, not uniform
spectral proofs. A separate moving two-component real-vector fixture
checks the kinetic product rule and its nonzero scalar derivative cost;
it is not a Yang--Mills model. V87's mathematical diagnostics are rerun.
The predecessor body is preserved byte for byte except for its title;
only the latest active main source remains in \texttt{latex\_notes},
and all build products are kept outside that folder.
% END CONSOLIDATED V88 APPEND

% BEGIN CONSOLIDATED V89 APPEND
\clearpage
\section{V89: Exact fibre geometry, tempered compression, and a native weighted replacement}
\label{r89:section}

\subsection{The new question and the inherited inputs}

V88 proved a uniform magnetic one-link gap but retained a derivative
cost bounded by $36b^2/\kappa$. We first determine whether that poor
weak-coupling order is only an artifact of the bound. It is not:
at a realizable cancelling exterior the exact ground-fibre derivative
cost equals $b^2/\kappa$. The obstruction is a large response to changing
staples, not closure of the frozen-link spectral gap.

We then use a smoother, explicitly normalized \emph{trial} fibre.
This is not substituted for the native conditional vacuum, and its
energy error and derivative cost are both retained. Two different
estimates result: an all-exterior operator compression bound, and a
stronger bound averaged against the \emph{actual} vacuum exterior
marginal. The latter produces a gauge-invariant trial state with that
same marginal and a volume-independent replacement-energy cost.
It does not prove a norm approximation or a mass gap.

The spherical Gamma domination is inherited from f15 V38, including
its all-field proof, and the frustration split from f14 V95--V96 and
f15 V38. They are not new results here. The archived Gamma law is used
only for an explicitly chosen trial density; the differential vacuum
conditional is not asserted to have that density. The full-space
ground-state variational principle followed by gauge invariance was
already used, for example, in f16.2 V51. The new quantitative statements
below concern the differential Hamiltonian and its moving fibres.

A targeted literature check of S.~Wiggins, \emph{Born--Oppenheimer,
Born--Huang, and exact factorization},
\href{https://arxiv.org/html/2608.08668v1}{arXiv:2608.08668v1},
Sections 2.4 and 4--5, supports the useful distinction between a clamped
eigenstate, its derivative correction, and a state-dependent exact
conditional factor. We use that distinction, not a molecular
mass-ratio expansion or a quantitative spectral theorem from that paper.
All new estimates are proved below. The scheduled companion review
and broad literature sweep are both due at V90.

\subsection{The exact ground-fibre metric and its cancellation cost}

Keep V88's unit-round normalization. For $h>0$, let
$\psi_h=\psi_{h e_0}$, $e(h)=e_\kappa(h)$, and define
\[
 m(h)=\langle\psi_h,q_0\psi_h\rangle=-e'(h),\qquad
 T(h)=\langle\psi_h,L\psi_h\rangle
       =\frac{e(h)+h m(h)}{\kappa}.
 \tag{V89.1}\label{r89:rotor-observables}
\]
The real positive phase makes the ground-fibre quantum metric simply
\[
 \mathsf G_f(v,w)=\langle\partial_v\psi_f,\partial_w\psi_f\rangle.
 \tag{V89.2}\label{r89:metric-definition}
\]
For $f=h n$, $|n|=1$, its matrix has the exact decomposition
\[
 \mathsf G_f=g_\parallel(h)nn^{\mathsf T}
             +g_\perp(h)(I-nn^{\mathsf T}),\qquad
 g_\parallel(h)=\|\partial_h\psi_h\|^2,\qquad
 g_\perp(h)=\frac{T(h)}{3h^2}.
 \tag{V89.3}\label{r89:metric-split}
\]
Indeed the stabilizer of $n$ makes the transverse metric scalar and
kills mixed terms. At $n=e_0$, write $\psi_h(q)=u_h(q_0)$.
A transverse field derivative is $q_i u_h'(q_0)/h$. Summing its
squared norm over $i=1,2,3$ gives
$h^{-2}\int(1-q_0^2)|u_h'(q_0)|^2=T(h)/h^2$.
Transverse symmetry divides this by three, proving the last formula.

The apparent singularity at $h=0$ is removable. V88's derivative
formula and $Lq_i=3q_i$ give
\[
 \left.\partial_{f_i}\psi_f\right|_{f=0}=\frac{q_i}{3\kappa},
 \qquad \mathsf G_0=\frac{I_4}{36\kappa^2},
 \tag{V89.4}\label{r89:metric-zero}
\]
because $\psi_0=1$ and $\int q_iq_j=\delta_{ij}/4$.

For a fixed link $e$, write $a_p=\overline{S_p}\in S^3$,
$f=\sum_{p\ni e}b_pa_p$, $B=\sum_{p\ni e}b_p$ and
$S_2=\sum_{p\ni e}b_p^2$. The full tensor strengthening of
\eqref{r88:field-gradient} is
\[
 \mathsf A_e:=\sum_{j\ne e,a}
       (\partial_{j,a}f)(\partial_{j,a}f)^{\mathsf T}
       =3\sum_{p\ni e}b_p^2(I-a_pa_p^{\mathsf T}),
 \qquad \Tr\mathsf A_e=9S_2.
 \tag{V89.5}\label{r89:field-tensor}
\]
For each of a staple's three exterior links, its three invariant
derivatives form an orthonormal tangent frame at the unit quaternion
$a_p$, possibly after an orthogonal change of frame. Their outer
products therefore sum to $I-a_pa_p^{\mathsf T}$. The twelve exterior
links are distinct for $N\ge3$, so contributions add without cross
terms. This proves the tensor identity, not just its trace.

Consequently the exact V88 Born--Huang term is
\[
 \begin{aligned}
 V_{\rm BH}&=\kappa\Tr(\mathsf G_f\mathsf A_e)\\
 &=3\kappa\sum_{p\ni e}b_p^2
       \{g_\parallel(h)(1-c_p^2)+g_\perp(h)(2+c_p^2)\},
       \qquad c_p=n\cdot a_p,
 \end{aligned}
 \tag{V89.6}\label{r89:exact-BH}
\]
when $h>0$. At a cancelling field, \eqref{r89:metric-zero} instead gives
the nonsingular exact value
\[
       f=0\quad\Longrightarrow\quad
       V_{\rm BH}=\frac{S_2}{4\kappa}.
 \tag{V89.7}\label{r89:cancellation-cost}
\]
At homogeneous coupling take two staples equal to $1$ and two equal to
$-1$. Their twelve distinct exterior links allow this assignment: set
all their factors to $1$, except one chosen factor in each desired
negative staple. Thus $f=0$, $S_2=4b^2$ and $V_{\rm BH}=b^2/\kappa$
on an actual finite-lattice exterior. Continuity also supplies nearby
exteriors with comparable cost. No uniform estimate
$V_{\rm BH}\le C\sqrt{\kappa b}$ with fixed $C$ can therefore hold
as $b/\kappa\to\infty$. This is not a probability lower bound for
that exterior in the native vacuum.

\subsection{The inherited thermal input, in trial-state normalization}

Let $\nu_a$ be the normalized Haar tilt $e^{a q_0}\,dq$, $a\ge0$,
and put $\mu(a)=\mathbb E_{\nu_a}q_0$. Here $\mu$ refers to this
explicit thermal law, whereas $m$ in \eqref{r89:rotor-observables}
refers to the magnetic Hamiltonian ground state. They are not identified.
The f15 V38 Gamma domination gives, for $a>0$,
\[
 0\le\mu(a)<1,\qquad
 a(1-\mu(a))\le\frac32,\qquad
 a^2\mathbb E_{\nu_a}(1-q_0)^2\le\frac{15}{4}.
 \tag{V89.8}\label{r89:thermal-moments}
\]
Here $s=a(1-q_0)$ has density proportional to
$e^{-s}s^{1/2}\sqrt{(2-s/a)_+}$. Its decreasing final factor is the
archived monotone tilt of a rate-one Gamma variable of shape $3/2$.
Thus no large-field asymptotic or Bessel-ratio approximation is used.

Integration of
$\operatorname{div}(e^{a q_0}\nabla q_0)$ and $\Delta q_0=-3q_0$
give an additional exact identity
\[
 \mathbb E_{\nu_a}(1-q_0^2)=\frac{3\mu(a)}a.
 \tag{V89.9}\label{r89:thermal-ibp}
\]
In a frame with axis $e_0$, the covariance matrix is diagonal, with
longitudinal entry $\operatorname{Var}_{\nu_a}q_0$ and each transverse
entry $\mu(a)/a$. Hence
\[
 \begin{aligned}
 0\preceq\operatorname{Cov}_{\nu_a}(q)&\preceq I_4
                       &&(a\ge0),\\
 \|\operatorname{Cov}_{\nu_a}(q)\|
       &\le\max\{15/(4a^2),1/a\}
        \le15/(8a) &&(a\ge2).
 \end{aligned}
 \tag{V89.10}\label{r89:thermal-covariance}
\]
The first bound follows from $|v\cdot q|\le|v|$; the second from
\eqref{r89:thermal-moments} and \eqref{r89:thermal-ibp}. These are
bounds for a declared trial law, not for the unknown vacuum conditional.

One useful consequence for the \emph{exact} rotor follows by testing
it with $Z_a^{-1/2}e^{a q_0/2}$. Its kinetic expectation is
$3\kappa a\mu(a)/4$, so
\[
 0\le e(h)+h
    \le\min\{h,\,3\kappa a/4+3h/(2a)\}.
\]
The first upper bound uses the constant trial. Choosing
$a=\sqrt{2h/\kappa}$ and using $\kappa T(h)\le e(h)+h$ gives
\[
 g_\perp(h)\le
 \min\left\{\frac1{\kappa^2},\frac1{3\kappa h},
                   \frac1{\sqrt{2\kappa}\,h^{3/2}}\right\}.
 \tag{V89.11}\label{r89:transverse-bound}
\]
The first entry is V88's derivative bound. The large-field improvement
in the transverse directions coexists with, and does not remove,
the cancellation cost \eqref{r89:cancellation-cost}.

\subsection{A tempered fibre with its full compression error paid}

Fix an energy parameter $\tau>0$, independent of the exterior, and set
\[
 \chi_{\tau,f}(q)=
 \frac{\exp(f\cdot q/(2\tau))}
      {\left(\int_{S^3}\exp(f\cdot v/\tau)\,dv\right)^{1/2}},
 \qquad (J_\tau\alpha)(q,y)=\chi_{\tau,f(y)}(q)\alpha(y).
 \tag{V89.12}\label{r89:tempered-fibre}
\]
This positive normalized family is globally smooth, including $f=0$.
It is gauge equivariant by the same left/right quaternion covariance
as the exact fibre. No new Hamiltonian or modified gauge constraint is
introduced: $J_\tau$ only specifies a trial subspace in the original one.
Let
\[
 \delta_\tau(h)=
   \langle\chi_{\tau,h e_0},(\kappa L-hq_0)\chi_{\tau,h e_0}\rangle-e(h)
       \ge0,
 \qquad
 W_\tau(y)=\kappa\sum_{j\ne e,a}\|\partial_{j,a}\chi_{\tau,f(y)}\|^2.
 \tag{V89.13}\label{r89:tempered-errors}
\]

\begin{theorem}[All-exterior tempered compression]\label{r89:uniform-theorem}
For the full native differential Hamiltonian and $H_{\rm lo}$ of
V88, the exact quadratic-form compression is
\[
 J_\tau^*\mathcal HJ_\tau=H_{\rm lo}+\delta_\tau(|f|)+W_\tau.
 \tag{V89.14}\label{r89:tempered-compression}
\]
At every exterior, for arbitrary nonnegative incident couplings,
\[
 \begin{aligned}
 0\le\delta_\tau(|f|)&\le\frac{3\kappa B}{4\tau}+\frac{3\tau}{2},\\
 0\le W_\tau&\le\frac{9\kappa S_2}{4\tau^2},\\
 0\le J_\tau^*\mathcal HJ_\tau-H_{\rm lo}
       &\preceq C_\tau I,\qquad
 C_\tau=\frac{3\kappa B}{4\tau}+\frac{3\tau}{2}
                         +\frac{9\kappa S_2}{4\tau^2}.
 \end{aligned}
 \tag{V89.15}\label{r89:uniform-errors}
\]
Consequently, uniformly in volume and the other plaquette couplings,
\[
 0\le E_0-E_{\rm lo}\le
           \min\{9S_2/\kappa,\,C_\tau\}.
 \tag{V89.16}\label{r89:energy-shift}
\]
\end{theorem}

\begin{proof}
For $h=|f|>0$ let $a=h/\tau$. The fibre probability
$\chi_{\tau,f}^2\,dq$ is a rotated $\nu_a$, by definition. Direct
differentiation and \eqref{r89:thermal-ibp} give the kinetic expectation
$3\kappa a\mu(a)/4$. Since $e(h)\ge-h$,
\[
 \delta_\tau(h)
 \le\frac{3\kappa h}{4\tau}\mu(a)+h(1-\mu(a))
 \le\frac{3\kappa B}{4\tau}+\frac{3\tau}{2}.
\]
At $h=0$ the actual defect is zero and the inequality still holds.
For each field direction $v$,
\[
 \partial_v\chi_{\tau,f}
   =\frac{v\cdot(q-\mathbb E_{\nu_{f/\tau}}q)}{2\tau}\,
           \chi_{\tau,f},\qquad
 \mathsf G^{(\tau)}_f
   =\frac{\operatorname{Cov}_{\nu_{f/\tau}}(q)}{4\tau^2}.
 \tag{V89.17}\label{r89:tempered-metric}
\]
Here $\nu_{f/\tau}$ denotes the vector tilt, not a native conditional.
The positive real normalization makes the Berry term zero.
Equation \eqref{r89:field-tensor} and the first covariance bound give
$W_\tau=\kappa\Tr(\mathsf G^{(\tau)}_f\mathsf A_e)
\le9\kappa S_2/(4\tau^2)$.
The kinetic product rule, exactly as a product rule and not a claim
that $\chi$ is an eigenfunction, proves \eqref{r89:tempered-compression}.
Use $J_\tau\alpha_{\rm lo}$ as a normalized variational trial to get
the upper energy bound; the lower bound and the alternative
$9S_2/\kappa$ are V88. Smoothness and compactness justify passage to
the closed forms, and gauge equivariance permits the physical restriction.
\end{proof}

For $S_2>0$, a useful explicit choice is
\[
 \tau_0=(3\kappa S_2)^{1/3},\qquad
 C_{\tau_0}=\frac94(3\kappa S_2)^{1/3}
             +\frac{3\kappa B}{4(3\kappa S_2)^{1/3}}.
 \tag{V89.18}\label{r89:cubic-scale}
\]
It minimizes the last two terms of $C_\tau$, not necessarily their
sum with its first term. In the homogeneous case $B=4b$, $S_2=4b^2$,
this is
$O(\kappa^{1/3}b^{2/3}+\kappa^{2/3}b^{1/3})$ as $b/\kappa$ grows.
Thus the uniform compression error is $o(b)$ in that ratio, with
no excluded exterior. If $S_2=0$, the exact constant fibre instead
gives zero error. The softened fibre has not been assigned V88's
exact eigenprojection or its complementary-block gap.

\subsection{A native energy-density bound and the probability of cancellation}

For this subsection onward the entire cubic torus has homogeneous
coupling $b>0$. Let $n_\ell=3N^3$ be its number of positive links;
the number of plaquettes is also $n_\ell$. Uniqueness of the positive
full ground function makes it gauge invariant and invariant under
lattice translations and cubic axis symmetries. Hence every link has
the same kinetic expectation and every plaquette the same magnetic
expectation. In particular a full-Haar ground-energy trial is legitimate,
although a non-invariant trial alone would not bound an excited
physical energy.

\begin{proposition}[Volume-uniform native local energy]\label{r89:energy-density-prop}
The actual ground energy satisfies
\[
 0\le\frac{E_0}{n_\ell}
       \le\min\{b,\,3\sqrt{2\kappa b}\}.
 \tag{V89.19}\label{r89:energy-density}
\]
Consequently, in the actual vacuum $\Omega$,
\[
 \kappa\|\nabla_e\Omega\|^2\le3\sqrt{2\kappa b},\qquad
 b\langle1-w_p\rangle_\Omega\le3\sqrt{2\kappa b}.
 \tag{V89.20}\label{r89:local-native-energy}
\]
\end{proposition}

\begin{proof}
Take the normalized product wavefunction
$\prod_e Z_a^{-1/2}e^{a q_{e,0}/2}$, $a>0$, in full Haar space.
Each link contributes kinetic energy $3\kappa a\mu(a)/4$.
The mean quaternion of a link under its squared trial density is
$(\mu(a),0,0,0)$, and the same is true for its inverse. The four
links of each plaquette are distinct and independent in this trial,
so multilinearity of quaternion multiplication gives
$\langle w_p\rangle_{\rm trial}=\mu(a)^4$. Thus
\[
 \frac{E_0}{n_\ell}
 \le\frac{3\kappa a\mu(a)}4+b(1-\mu(a)^4)
 \le\frac{3\kappa a}4+\frac{6b}{a}.
\]
Here $1-\mu^4\le4(1-\mu)$ and the inherited first-moment bound
were used. Choosing $a=\sqrt{8b/\kappa}$ gives
$3\sqrt{2\kappa b}$. The constant full-Haar trial gives $b$.
All actual kinetic and magnetic contributions are nonnegative;
the stated symmetries then give \eqref{r89:local-native-energy}.
No trial expectation has been equated to an actual vacuum expectation.
\end{proof}

The exterior-measurable star frustration is
\[
 \mathcal F_e(y)=4b-|f_e(y)|,\qquad
 0\le\mathcal F_e\le\sum_{p\ni e}b(1-w_p(q,y)).
 \tag{V89.21}\label{r89:frustration}
\]
This is the inherited minimum-over-one-link inequality, applied to
the four spatial plaquettes. It is pointwise for every $q,y$, not
a statement about a Gibbs conditional of $\Omega$. Therefore the
actual exterior marginal satisfies
\[
 \begin{aligned}
 \langle\mathcal F_e\rangle_\Omega
       &\le\min\{4b,12\sqrt{2\kappa b}\},\\
 \mathbb P_\Omega\{|f_e|\le4b(1-\eta)\}
       &\le\min\left\{1,\frac{3\sqrt2}{\eta}
                                  \sqrt{\frac\kappa b}\right\},
               \qquad 0<\eta\le1.
 \end{aligned}
 \tag{V89.22}\label{r89:native-bad-field}
\]
The second line is Markov's inequality applied to the first.
This is a one-star first-moment estimate. It supplies neither joint
rare-region tails nor a many-source operator norm.

\subsection{A paid replacement at the actual exterior marginal}

Define the marginal amplitude and actual conditional factor by
\[
 \rho(y)=\left(\int|\Omega(q,y)|^2\,dq\right)^{1/2},\qquad
 \zeta(q,y)=\Omega(q,y)/\rho(y).
 \tag{V89.23}\label{r89:actual-factor}
\]
At fixed finite regulator both are smooth and strictly positive, and
$\|\zeta(\cdot,y)\|=1$. The probability $\rho(y)^2dy$ is retained
exactly below. In particular it is not the ground-state density of
$H_{\rm lo}$ and is not replaced by a product exterior law.

\begin{theorem}[Native weighted tempered replacement]\label{r89:native-theorem}
Set $\tau_*=\sqrt{\kappa b}$ and
$\Phi_e=J_{\tau_*}\rho$. It is a normalized physical trial state
with exactly the same exterior marginal as $\Omega$, and
\[
 0\le\langle\Phi_e,(\mathcal H-E_0)\Phi_e\rangle
                   \le90\sqrt{\kappa b}.
 \tag{V89.24}\label{r89:native-replacement}
\]
The constant is independent of $N\ge3$, $\kappa,b>0$, and the
chosen link. The result makes no identification
$\zeta=\chi_{\tau_*,f}$ or $\zeta=\psi_f$.
\end{theorem}

\begin{proof}
The normalized real product rule gives the exact factorization identity
\[
 E_0=\mathfrak h_{\rm lo}[\rho]+\int\rho^2
 \left\{\langle\zeta,(h_e-\lambda_e)\zeta\rangle
         +\kappa\sum_{j\ne e,a}\|\partial_{j,a}\zeta\|^2\right\}dy
       \ge\mathfrak h_{\rm lo}[\rho].
 \tag{V89.25}\label{r89:actual-energy-split}
\]
Both terms in braces are nonnegative. In the derivative term the link
coordinate $q$ is held fixed, exactly as in the original kinetic form.
Subtract this identity from the tempered compression to obtain
\[
 \langle\Phi_e,\mathcal H\Phi_e\rangle-E_0
       \le\int\rho^2\{\delta_{\tau_*}(|f|)+W_{\tau_*}\}\,dy.
 \tag{V89.26}\label{r89:weighted-cost}
\]
The reverse lower bound zero is the variational principle. Normalization
and the unchanged marginal follow from the fibre normalization;
gauge invariance follows from equivariance and the gauge-invariant
native marginal amplitude.

The defect bound in \eqref{r89:uniform-errors}, with $B=4b$, gives
$\delta_{\tau_*}\le9\tau_*/2$. The global derivative bound is
$W_{\tau_*}\le9b$. If $b\le\kappa$, this already gives a total
weighted cost at most $27\tau_*/2<90\tau_*$.

Suppose $b\ge\kappa$. Split the exterior into
$\mathcal G=\{|f|\ge2b\}$ and its complement. On $\mathcal G$,
$a=|f|/\tau_*\ge2$, so \eqref{r89:thermal-covariance},
\eqref{r89:tempered-metric}, and $\Tr\mathsf A_e=36b^2$ give
\[
 W_{\tau_*}
 \le9b\frac{15}{8a}
 \le\frac{135}{16}\tau_*.
\]
On $\mathcal G^c$, use the global bound and
\eqref{r89:native-bad-field} at $\eta=1/2$:
\[
 \int_{\mathcal G^c}\rho^2 W_{\tau_*}\,dy
       \le9b\,\mathbb P_\Omega(\mathcal G^c)
       \le54\sqrt2\,\tau_*.
\]
The cost in \eqref{r89:weighted-cost} is therefore at most
$(207/16+54\sqrt2)\tau_*<90\tau_*$.
No conditional independence or approximation of $\rho^2$ was used.
\end{proof}

This bound also implies only a high-energy tail estimate,
\[
 \|\mathbf1_{[R,\infty)}(\mathcal H-E_0)\Phi_e\|^2
          \le\frac{90\sqrt{\kappa b}}R,\qquad R>0,
 \tag{V89.27}\label{r89:high-energy-only}
\]
by the spectral theorem. It does not bound the weight arbitrarily
close to zero above the vacuum. In particular a small replacement
energy relative to the magnetic scale $b$ must not be turned into a
norm estimate using an unproved full spectral gap.

\subsection{Status and the remaining neutral problem}

The exact fibre's uniform derivative obstruction has been located in
the actual compact geometry, not inferred from a qubit countermodel.
The tempered compression avoids that worst-case loss while paying its
fibre-energy error. Its native weighted version additionally retains
the actual exterior marginal and costs $O(\sqrt{\kappa b})$ per
single-link replacement. These results are compatible with arbitrarily
soft retained neutral excitations. In the physical normalization
$\kappa=g^2/(2c a_s)$, $b=c/(g^2a_s)$ of V85, even
$\sqrt{\kappa b}=1/(\sqrt2 a_s)$ is an ultraviolet energy scale,
not a finite physical mass lower bound.

The next obligation is a neutral, growing-region coercivity or
connected-return estimate that survives repeated replacements while
retaining their exterior correlations and physical normalization.
Single-star probability bounds do not permit independent resampling
of overlapping stars or summing spectral gaps of moving fibres.
The continuum construction and the full Yang--Mills mass gap remain
unproved. The V90 literature and companion checkpoint should target
this interacting retained problem, rather than re-derive the local rotor.

Exact diagnostics check the full quaternion derivative tensor, its
cancelling-field value, the thermal product plaquette mean, and the
constants in the uniform and native weighted budgets. Independent
finite angular-sector calculations compare the rotation metric with
the transverse resolvent metric; thermal quadrature checks the declared
trial-law formulas. They do not sample or certify the actual many-link
vacuum, and numerical spectra are not proofs of uniform bounds.
V88's mathematical checks are rerun. The full predecessor body is
preserved except for its title, and all build products remain outside
the \texttt{latex\_notes} folder.
% END CONSOLIDATED V89 APPEND

% BEGIN CONSOLIDATED V90 APPEND
\clearpage
\section{V90: Consistent trial law and exact neutral penalty}
\label{r90:section}

\subsection{Checkpoint, provenance, and a change of scale}

V89 supplied a controlled replacement of one link at the actual vacuum's
exterior marginal. It did not justify repeatedly replacing overlapping
links, identify the trial conditional with the native conditional, or
bound low-energy neutral excitations. This append moves upstream: assemble
the tempered conditionals into their unique compatible full-system law,
then compute exactly which Hamiltonian has that law as its vacuum and
which terms distinguish it from the native Wilson Hamiltonian.
The outcome is an explicit nonnegative neutral penalty and a two-sided,
volume-uniform bound on the native \emph{ground-energy density}. The
excitation gap is not that density.

The scheduled V90 companion review found no changes in NS V26,
SM--Einstein V72, Shell Mixing V16, parallel YM V5, or the four supplied
attachments SM Source V43, Einstein Bridge V36, Native Stationarity V46,
and Perturbative ESM V51. Their terminal scope statements were reread and
their fingerprints checked against V85. This is not a new independent
recertification of their complete histories. In particular, fixed-degree
ESM identities are not analytic control of all fluctuation orders;
fixed-time or selected-history stability is not a uniform physical-time
estimate; the assembled flat-holonomy coefficient does not control all
noncommuting backgrounds. None supplies the missing YM infrared bound.
The review record is \path{artifacts/ym_restart_v90_companion_review.json}.

Nonduplication matters here. The general ground-state transform and
Perron/local-law mismatch already occur in archive f1 V6 and V21--V23.
Archive f15 V38 supplies the one-star frustration split and thermal
moment estimates. V87 identifies neutral two-staple returns; V88--V89
give the differential rotor, moving fibres, and their derivative costs.
We reuse these ideas, not claim them anew. The new calculation is the
full three-dimensional Wilson force-square identity, its tuned
positive decomposition of the native differential Hamiltonian, and an
explicit all-temperature obstruction to identifying its vacuum with
the compatible tempered law.

\subsection{Broad literature sweep: what is and is not imported}

The V90 sweep covered continuum construction and RG, gauge-orbit
Hamiltonians, finite-size gap certificates, stochastic quantization,
entropy factorization, and Wilson-loop area laws. The following are
scope decisions, not an assertion that every cited paper has been
fully independently verified.

\begin{itemize}
\item Nair, \emph{Towards a Proof of Mass Gap in 3d Yang--Mills Theory},
\href{https://arxiv.org/html/2608.10133v2}{arXiv:2608.10133v2},
uses two spatial dimensions and a gauge-orbit volume element. Equations
(35)--(40) concern a kinetic operator. The passage following (40) adds
a nonnegative magnetic operator. Absolute spectral order alone does
not compare gaps above different vacua; the elementary check below
is mandatory. No three-dimensional or kinetic bound is imported into
our four-dimensional target.

\item Faizal--Shabir,
\emph{Reflection-Positive Construction of a Four-Dimensional
SU(N) Yang--Mills Theory with Mass Gap and Confinement},
\href{https://arxiv.org/pdf/2606.19362v1}{arXiv:2606.19362v1},
was screened at its abstract and Section 10, especially Definition
10.1, Lemma 10.2, and the entry argument. Equation (10.8) uses only
finitely many derivatives at the origin, whereas (10.14) claims
stability under arbitrary fluctuation shifts. Such a seminorm is
not shift stable on general activities. Moreover the displayed
$g\|K\|+\|K\|^2$ remainder in (10.41) is not uniformly dominated by
$g^3$ merely by fixing a small, $g$-independent polydisc.
The explicit tests below rule out these inferences without extra
hypotheses. We do not import the claimed continuum conclusion.
This is a targeted audit, not a review of all 593 PDF pages.

\item Rai and collaborators,
\emph{A Hierarchy of Spectral Gap Certificates for Frustration-Free
Spin Systems},
\href{https://arxiv.org/html/2411.03680v3}{arXiv:2411.03680v3},
Section 2, assumes finite-dimensional local spaces and a common
zero-energy ground state for the positive local terms.
Its local positive-decomposition strategy is a useful candidate
after these hypotheses are established, not a ready certificate for
unbounded Wilson rotors with a nonzero vacuum-energy subtraction.
The present append computes the positive parent and the missing
penalty explicitly.

\item Sevostyanov's
\href{https://arxiv.org/abs/2605.14616v2}{arXiv:2605.14616v2}
constructs a regularity-structure model for three-dimensional
Yang--Mills Langevin dynamics, according to its abstract; that is the
scope screened here. It is not used as a four-dimensional physical
Hamiltonian gap or as a proof that our spatial Gibbs trial is the
native vacuum.

\item The abstracts and versions of Caputo--Chen--Parisi,
\href{https://arxiv.org/abs/2503.19419v1}{arXiv:2503.19419v1},
and Cao--Nissim--Sheffield,
\href{https://arxiv.org/abs/2509.04688v2}{arXiv:2509.04688v2},
were rechecked. The former's weak-interaction factorization conditions
remain to be verified for a proposed retained law; the latter's
lattice area-law regime does not supply a weak-coupling continuum
transfer theorem. Their previously recorded scope restrictions remain.
\end{itemize}

\begin{proposition}[Two elementary import checks]
\label{r90:import-checks}
Neither of the following implications holds without additional input:
positive operator order preserves a gap above the lowest eigenvalue;
finite jets at one point control Gaussian fluctuation averaging.
\end{proposition}
\begin{proof}
For $0<\varepsilon<1$ take
\[
 T=\operatorname{diag}(0,1),\qquad
 W=\operatorname{diag}(1-\varepsilon,0)\ge0.
\]
Although $T+W\ge T$ and both ground states are unique, their gaps are
$\varepsilon$ and $1$, respectively. Adding a positive term can thus
close the centered gap without contradicting eigenvalue monotonicity.
This fixture is a test of an implication, not a Wilson counterexample.

For an integer $p\ge0$, set
$J_p(F)=\sum_{j=0}^p|F^{(j)}(0)|/j!$ and $F(x)=x^{2p+2}$.
Then $J_p(F)=0$. If $Z$ is a nondegenerate centered real Gaussian,
$(\mathbb E F(\,\cdot+Z))(0)=\mathbb E Z^{2p+2}>0$.
Hence there is no finite constant $C$ with
$J_p(\mathbb E F(\,\cdot+Z))\le C J_p(F)$ for all such polynomials.
Nor is the kernel of $J_p$ preserved by this averaging, so passing
to its quotient does not repair the problem. These are own elementary
checks of the proposed inference, not imported theorems.
\end{proof}

For completeness, if $\delta,c,\beta>0$ are fixed, then
$cg\delta/g^3=c\delta/g^2\to\infty$ as $g\downarrow0$.
A bound by $c(g^5+g\|K\|+\|K\|^2)$ is therefore insufficient to
deduce a negative cubic flow uniformly over $\|K\|\le\delta$.
A suitably small $g$-dependent condition such as
$\|K\|\le a g^2$, with its own preserved-domain proof, is a
different hypothesis. Likewise, a norm controlling actual shifted
activities and their integrable tails is different from $J_p$.
These tests reinforce the fixed-degree/analytic distinction in the
ESM attachment rather than promoting that attachment to an RG closure.

\subsection{The compatible full-system trial law}

Use the homogeneous cubic torus of V85--V89, of side $N\ge3$, with
$n=3N^3$ links and $n$ plaquettes, round unit-radius $\SU(2)=S^3$,
and normalized product Haar measure $dU$. Let
\[
 L=\sum_e L_e=-\Delta,\qquad
 V=b\sum_p(1-w_p),\qquad
 \mathcal H=\kappa L+V,\qquad b,\kappa>0.
 \tag{V90.1}\label{r90:H}
\]
All statements below concern this specified finite spatial regulator.
For each fixed energy parameter $\tau>0$ define
\[
 \Psi_\tau=Z_\tau^{-1/2}e^{-V/(2\tau)},\quad
 Z_\tau=\int e^{-V/\tau}\,dU,\quad d\mu_\tau=\Psi_\tau^2dU.
 \tag{V90.2}\label{r90:trial}
\]
This is a declared \emph{spatial Gibbs trial law}. It is not assigned
the physical vacuum interpretation.

\begin{proposition}[Compatibility and uniqueness]
\label{r90:compatible}
For every link $e$ the normalized conditional amplitude of
$\Psi_\tau$ at fixed exterior $y$ is exactly V89's
$\chi_{\tau,f_e(y)}$. Conversely, a strictly positive smooth
normalized amplitude with all these single-link conditionals equals
$\Psi_\tau$.
\end{proposition}
\begin{proof}
The potential terms containing $e$ are
$E_e(q,y)=4b-f_e(y)\cdot q$; all other terms are independent of $q$.
Conditional normalization of $e^{-V/(2\tau)}$ therefore leaves
$Z(f_e/\tau)^{-1/2}e^{f_e\cdot q/(2\tau)}$, precisely the
V89 amplitude. If another positive amplitude $\Phi$ has the same
normalized conditional for $e$, the ratio $\Phi/\Psi_\tau$ is
independent of that coordinate. Applying this statement to every
link makes the ratio constant on the connected product manifold.
Normalization and positivity fix it to one.
\end{proof}

Compatibility does not prove that sequential replacements starting
from the native vacuum converge with a useful rate, or keep the
exterior marginal in the V89 weighted class. It only identifies their
common compatible law if all conditional equations are imposed.

\subsection{The full Wilson force and its retained neutral part}

Write the four incident staples as unit quaternions $a_{ep}(y)$ in
the common link frame, so $w_p=a_{ep}\cdot q_e$. Set
\[
 \begin{gathered}
 f_e=b\sum_{p\ni e}a_{ep},\qquad s_e=|f_e|,\qquad
 E_e=\sum_{p\ni e}b(1-w_p)=4b-f_e\cdot q_e,\\
 D_e^{\,0}=16b^2-s_e^2,\qquad F_e=4b-s_e.
 \end{gathered}
 \tag{V90.3}\label{r90:star}
\]
The superscript distinguishes the scalar $D_e^{\,0}$ from the
first-order operator introduced below.

\begin{lemma}[Exact force-square decomposition]
\label{r90:force}
Pointwise on the full compact configuration space,
\[
 \begin{split}
 LV&=12(V-bn),\\
 |\nabla_e V|^2&=s_e^2-(4b-E_e)^2
                  =8bE_e-E_e^2-D_e^{\,0},\\
 |\nabla V|^2&=32bV-\sum_e\bigl(E_e^2+D_e^{\,0}\bigr).
 \end{split}
 \tag{V90.4}\label{r90:force-id}
\]
Moreover $D_e^{\,0}$ is an exterior-only, gauge-invariant,
nonnegative two-staple expression:
\[
 D_e^{\,0}
  =b^2\sum_{\substack{p<q\\p,q\ni e}}|a_{ep}-a_{eq}|^2
  =F_e(8b-F_e).
 \tag{V90.5}\label{r90:neutral}
\]
\end{lemma}
\begin{proof}
A plaquette is linear in each of its four separate quaternion
coordinates. Since $L_eq_e^i=3q_e^i$, summing its four link
Laplacians gives $Lw_p=12w_p$. This proves the first identity.
The tangential gradient on $S^3$ of $-f_e\cdot q_e$ is
$-f_e+(f_e\cdot q_e)q_e$, of squared norm
$s_e^2-(f_e\cdot q_e)^2$. Substitute $f_e\cdot q_e=4b-E_e$.
Every plaquette enters four stars, hence $\sum_e E_e=4V$,
which proves the last line.

For four unit vectors, expansion of the six squared differences
gives $\sum_{p<q}|a_p-a_q|^2=16-|\sum_p a_p|^2$.
This gives the first expression for $D_e^{\,0}$; factoring
$(4b-s_e)(4b+s_e)$ gives the second. Gauge transformations rotate
all four vectors by the same orthogonal left-right quaternion action,
so their pairwise differences and inner products are invariant.
Equivalently $a_{ep}\cdot a_{eq}$ is the normalized trace of a
closed two-staple loop. The term survives elimination of the
central link; it is not a charged fibre excitation.
\end{proof}

\subsection{A positive parent and an exact native comparison}

Define on smooth scalar functions
\[
 \mathcal D_{\tau,e}=\nabla_e+\frac{\nabla_e V}{2\tau},
 \qquad
 K_\tau=\kappa\sum_e\mathcal D_{\tau,e}^{\,*}
                           \mathcal D_{\tau,e}.
 \tag{V90.6}\label{r90:parent}
\]
The adjoint uses product Haar measure and the round tangent metric.
The sum of squares defines a closed nonnegative form on $H^1$;
bounded smooth coefficients on the finite compact manifold give the
usual self-adjoint elliptic realization. Its form kernel is
$\operatorname{span}\{\Psi_\tau\}$: the equation
$\mathcal D_\tau u=0$ is
$\nabla(u e^{V/(2\tau)})=0$. Each local term annihilates this same
positive state. Gauge invariance of $V$ and of the metric implies
that $K_\tau$ preserves the physical subspace. Each term involves
only a link and its incident plaquette links, not the entire volume.

\begin{theorem}[Native Hamiltonian equals parent plus neutral penalty]
\label{r90:decomposition}
As self-adjoint operators, or equivalently closed forms,
\[
 \begin{split}
 \mathcal H&=K_\tau+c_\tau n+A_\tau V+R_\tau,\\
 c_\tau&=\frac{6\kappa b}{\tau},\qquad
 A_\tau=1-\frac{6\kappa}{\tau}
                       -\frac{8\kappa b}{\tau^2},\\
 R_\tau&=\frac{\kappa}{4\tau^2}
                 \sum_e\bigl(E_e^2+D_e^{\,0}\bigr)\ge0.
 \end{split}
 \tag{V90.7}\label{r90:decomposition-id}
\]
In particular, for
\[
 \tau_{\rm P}=3\kappa+\sqrt{9\kappa^2+8\kappa b},
 \qquad c_{\rm P}=\frac34
                 \bigl(\sqrt{9\kappa^2+8\kappa b}-3\kappa\bigr),
 \tag{V90.8}\label{r90:tuning}
\]
one has the exact positive decomposition
\[
             \mathcal H= c_{\rm P}n+K_{\tau_{\rm P}}
                                               +R_{\tau_{\rm P}}.
 \tag{V90.9}\label{r90:positive}
\]
Here $\tau_{\rm P}$ is a new constant choice of trial parameter;
it is not V89's choice $\sqrt{\kappa b}$.
\end{theorem}
\begin{proof}
On a boundaryless compact manifold, expanding the square gives
\[
 K_\tau=\kappa L+\frac{\kappa LV}{2\tau}
                           +\frac{\kappa|\nabla V|^2}{4\tau^2}.
\]
Indeed the first-order cross terms cancel, leaving
$-\operatorname{div}(\nabla V/(2\tau))=LV/(2\tau)$.
Subtract this expression from $\mathcal H$ and insert
\eqref{r90:force-id}. The resulting coefficient of $V$ is $A_\tau$,
the constant is $c_\tau n$, and the remaining multiplication
operator is $R_\tau$. All coefficients are bounded at fixed
regulator, so equality on smooth functions extends to the stated
realizations. The positive root of
$\tau^2-6\kappa\tau-8\kappa b=0$ is $\tau_{\rm P}$.
Rationalizing $6\kappa b/\tau_{\rm P}$ gives $c_{\rm P}$.
The same identities hold after physical restriction.
\end{proof}

\begin{corollary}[Two-sided native vacuum-energy density]
\label{r90:energy}
Let $E_0$ be the lowest native eigenvalue. Uniformly in $N\ge3$,
\[
 c_{\rm P}\ \le\ \frac{E_0}{n}\
                    \le\ \min\{b,3\sqrt{2\kappa b}\}.
 \tag{V90.10}\label{r90:energy-bounds}
\]
At each finite regulator the lower inequality is strict.
For $b/\kappa\to\infty$ its lower endpoint is
\[
 c_{\rm P}=\frac3{\sqrt2}\sqrt{\kappa b}
                      -\frac94\kappa
                      +O(\kappa^{3/2}b^{-1/2}).
 \tag{V90.11}\label{r90:energy-asymptotic}
\]
\end{corollary}
\begin{proof}
The lower bound follows from \eqref{r90:positive}.
The upper bounds are exactly V89's constant and product trial
estimates; they are not rederived. The strictly positive smooth
native ground $\Omega$ has positive integral against the continuous
nonnegative function $R_{\tau_{\rm P}}$, which is nonzero on an
open set, as also shown by the example below. Thus
$E_0-c_{\rm P}n=\langle\Omega,K_{\tau_{\rm P}}\Omega\rangle
+\int R_{\tau_{\rm P}}\Omega^2dU>0$.
The expansion is that of the square root in \eqref{r90:tuning}.
Full and physical lowest eigenvalues agree because uniqueness and
positivity make the native ground gauge invariant, as used in V89.
\end{proof}

This lower bound controls the zero-point energy, not its separation
from the next level. In particular the two positive summands in
\eqref{r90:positive} have no common zero vector. The full operator
after this constant subtraction is \emph{not} frustration free.

\subsection{An explicit obstruction to an exact tempered vacuum}

\begin{proposition}[No value of the scalar temperature is exact]
\label{r90:not-native}
For every $b,\kappa,\tau>0$ and $N\ge3$, $\Psi_\tau$ is not an
eigenfunction of $\mathcal H$. Consequently the native positive
vacuum cannot have all of V89's tempered conditionals at one fixed
$\tau$, even though those conditionals are mutually compatible.
\end{proposition}
\begin{proof}
Keep all links at the identity except one link, set to
$q(\theta)=\cos\theta+\mathbf i\sin\theta$.
Precisely four plaquettes have trace $\cos\theta$, hence
$V=4b(1-\cos\theta)$. The selected link has force squared
$16b^2\sin^2\theta$. The twelve distinct exterior links of
these four plaquettes each have force squared $b^2\sin^2\theta$:
each belongs to only one of these four plaquettes, while all
other incident plaquettes are at the identity and have zero
first derivative. Thus
\[
                |\nabla V|^2=28b^2\sin^2\theta .
 \tag{V90.12}\label{r90:one-link-curve}
\]
Here $N\ge3$ ensures that the exterior links are distinct.
Since $K_\tau\Psi_\tau=0$, with $x=\cos\theta$ one obtains
\[
 \frac{\mathcal H\Psi_\tau}{\Psi_\tau}
   =c_\tau n+4b\left(1-\frac{6\kappa}{\tau}\right)(1-x)
                     -\frac{7\kappa b^2}{\tau^2}(1-x^2).
 \tag{V90.13}\label{r90:local-energy-curve}
\]
Its coefficient of $x^2$ is $7\kappa b^2/\tau^2>0$.
It cannot be constant on $[-1,1]$, whereas an eigenfunction
would have constant local energy everywhere. Smoothness and
positivity exclude an almost-everywhere loophole. At the tuned
parameter the nonconstant penalty along this curve is explicitly
\[
 R_{\tau_{\rm P}}
     =\frac{\kappa b^2}{\tau_{\rm P}^2}(1-x)(25-7x);
 \tag{V90.14}\label{r90:penalty-curve}
\]
it vanishes at $x=1$ and is strictly positive at $x=-1$.
Finally apply Proposition~\ref{r90:compatible}.
\end{proof}

\subsection{The native density correction is retained, not discarded}

For $\tau=\tau_{\rm P}$ abbreviate $\mu=\mu_\tau$, $\Psi=\Psi_\tau$,
$R=R_\tau$, and $\delta_N=E_0-c_{\rm P}n>0$.
Multiplication by $\Psi$ is a unitary map from $L^2(\mu)$ to
$L^2(dU)$. Direct application of the product rule gives
\[
 \Psi^{-1}K_\tau\Psi=\mathcal L_\tau,\qquad
 \mathcal L_\tau=\kappa\left(
             -\Delta+\frac1\tau\nabla V\cdot\nabla\right),\qquad
 q_{\mathcal L_\tau}[f]=\kappa\int|\nabla f|^2d\mu.
 \tag{V90.15}\label{r90:weighted}
\]
The minus sign in front of $\Delta$ makes this a nonnegative
operator; the corresponding Markov generator is $-\mathcal L_\tau$.
It is an auxiliary spatial dynamics, not physical time.
For the positive normalized native ground define $r=\Omega/\Psi$.
The exact correction equation and energy identity are
\[
 \begin{gathered}
 (\mathcal L_\tau+R)r=\delta_N r,\qquad
                  \int r^2d\mu=1,\\
 \delta_N=\kappa\int|\nabla r|^2d\mu+\int Rr^2d\mu,\qquad
                  \Omega^2dU=r^2d\mu.
 \end{gathered}
 \tag{V90.16}\label{r90:correction}
\]
They follow by unitary conjugation of \eqref{r90:positive} and
integration by parts. In relative-Fisher-information notation this
also reads
\[
 \frac{\kappa}{4}I(\Omega^2dU\,|\,\mu)
                       +\mathbb E_\Omega R
          =\delta_N
          \le n\bigl(\min\{b,3\sqrt{2\kappa b}\}-c_{\rm P}\bigr),
 \tag{V90.17}\label{r90:fisher}
\]
where $I(\nu|\mu)=\int|\nabla\log(d\nu/d\mu)|^2d\nu$.
Indeed $d\nu/d\mu=r^2$, so $I=4\int|\nabla r|^2d\mu$.
This is an extensive upper bound, not uniform density-ratio or
total-variation control.

For clarity the already standard exact native ground transform now
has a fully specified residual equation:
\[
 \left\langle\Omega f,(\mathcal H-E_0)\Omega f\right\rangle
       =\kappa\int|\nabla f|^2r^2d\mu .
 \tag{V90.18}\label{r90:actual-gap-form}
\]
To check it, expand $\nabla(rf)$ in the form of
$\mathcal L_\tau+R-\delta_N$ and test
\eqref{r90:correction} against $r|f|^2$; the cross term and
the terms proportional to $|f|^2$ cancel.
Thus a physical gap lower bound $m$ is equivalent, at this regulator,
to the Poincare inequality
\[
 m\,\Var_{r^2\mu}(f)
           \le\kappa\int|\nabla f|^2r^2d\mu
 \quad\hbox{for all gauge-invariant }f.
 \tag{V90.19}\label{r90:target}
\]
No Poincare estimate under $\mu$ has been substituted for this one.
In particular
\[
 \mathcal H-E_0=K_\tau+R-\delta_N I
 \tag{V90.20}\label{r90:centered}
\]
contains a positive penalty \emph{and} a nonzero ground-energy
subtraction. Dropping the former while forgetting the latter is
exactly the spectral-order error detected earlier.

\subsection{Status, verification, and next upstream obligation}

V90 makes the global parent comparison explicit, retains rather than
discards the neutral two-staple geometry, gives a nonzero lower
vacuum-energy density of the correct $\sqrt{\kappa b}$ order, and
rules out exact scalar-tempered native conditionals for every
temperature. The new results hold for the original finite-volume
Hamiltonian, not a claimed continuum replacement. Under V85's
physical normalization $\sqrt{\kappa b}=1/(\sqrt2a_s)$, the density
scale is ultraviolet. No finite physical mass follows from it.

The next problem is quantitative control of the interacting correction
$r$ or of the centered operator
$\mathcal L_{\tau_{\rm P}}+R-\delta_N$, on growing regions, with
the physical scale retained. A local certificate for the parent alone
is insufficient. The explicit $E_e^2+D_e^{\,0}$ structure offers a
concrete target for a connected correction or a native low-energy
certificate; global comparison by the oscillation of an extensive
potential is not a uniform solution.
The four-dimensional continuum construction and the full physical
Yang--Mills mass gap remain unproved.

The diagnostic \path{scripts/verify_ym_restart_v90_parent_penalty.py}
checks the exact tuning algebra, rational quaternion force identities
on complete periodic lattices, the single-link curve with all exterior
derivatives retained, and the two import counterchecks. These finite
algebra checks are not uniform spectral proofs.
V89's mathematical diagnostics are rerun. The V89 body is preserved
byte for byte except its title; builds and review ledgers are outside
\texttt{latex\_notes}. The next companion checkpoint is V95 and the
next broad literature checkpoint is V100, followed by the requested
archive/restart.
% END CONSOLIDATED V90 APPEND

% BEGIN CONSOLIDATED V91 APPEND
\clearpage
\section{V91: Native conditional coercivity and heat comparison}
\label{r91:section}

\subsection{The upstream question and what is new}

V90's positive trial parent retains an explicit neutral penalty, but
its vacuum differs from the native one. Its gap cannot be transferred
by dropping that penalty and its new ground-energy subtraction.
This append instead obtains a coercive comparison which annihilates
the \emph{actual} native vacuum. The missing ingredient is an
all-exterior bound on a native single-link conditional, not on a
frozen magnetic rotor or a spatial Gibbs trial.

We prove that ingredient by isolating a link's continuous-time heat
kernel inside the native semigroup. Optimizing the comparison time
gives a volume-independent logarithmic oscillation of order
$\sqrt{b/\kappa}$ at large $b/\kappa$. It follows that the native
centered Hamiltonian dominates a sum of bounded projections, all
with the correct vacuum in their kernels. These projections are
not asserted to have finite spatial support: controlling their
long-range conditional dependence is the remaining obstruction.

Archive f1 V13--V14 already proves finite-step Perron--Harnack
inequalities and conditional density comparisons. Archive f16.1 V64
records a crude finite-step bound and an unproved improved small-move
bound. Archive f5 V55 supplies a qualitative antipodal heat estimate;
f7 V12 and f12 V60 supply density/Poincare and diffusion/heat-bath
comparisons. Those general principles are credited, not rediscovered.
The new input here is their quantitative application to the actual
differential Wilson vacuum using its continuous semigroup, including
an explicit heat-kernel oscillation constant and the optimized
coupling dependence. This does not prove the differently normalized
finite-step small-move assertion in f16.1 V64.

For heat-kernel context the targeted primary source is
Nowak--Sj\"ogren--Szarek,
\emph{Sharp estimates of the spherical heat kernel},
\href{https://arxiv.org/pdf/1802.09385v2}{arXiv:1802.09385v2},
Sections 1--2. Its monotonicity statement and sphere recurrence are
consistent with the direct three-sphere calculation below; no
unspecified Gaussian constants or external gap theorem are imported.
The V90 broad review remains the latest checkpoint, with companion
review next at V95 and broad review next at V100.

\subsection{A completely specified one-link heat comparison}

Let $p_s(q,q')$ be the kernel of $e^{-sL_e}$ relative to normalized
Haar measure on the unit round $S^3$, with $L_e=-\Delta_{S^3}$.
By invariance it is a function $p_s(\theta)$ of distance
$0\le\theta\le\pi$. Put
\[
 M_s=p_s(0),\qquad m_s=p_s(\pi),\qquad
 a(u)=\inf_{s>0}\left\{us+\log\frac{M_s}{m_s}\right\},
 \quad u\ge0 .
 \tag{V91.1}\label{r91:a}
\]
All logarithms are natural. These quantities refer to an ordinary
one-link heat kernel, not a many-link vacuum computation.

\begin{lemma}[Explicit heat-kernel oscillation]
\label{r91:heat}
For all $s>0$ and $q,q'\in S^3$,
$0<m_s\le p_s(q,q')\le M_s$. Moreover
\[
 \log\frac{M_s}{m_s}\le
 \begin{cases}
       \pi^2/(4s),&0<s\le1,\\
       14e^{-3s},&s\ge1.
 \end{cases}
 \tag{V91.2}\label{r91:heat-ratio}
\]
Consequently $a(0)=0$, and for $u>0$,
\[
 a(u)\le
 \begin{cases}
       u+\pi^2/4,&u\le\pi^2/4,\\
       \pi\sqrt u,&u\ge\pi^2/4 .
 \end{cases}
 \tag{V91.3}\label{r91:a-large}
\]
There is also the small-$u$ bound
\[
 a(u)\le\frac u3\left(1+\log\frac{42}{u}\right),
             \qquad 0<u\le42e^{-3}.
 \tag{V91.4}\label{r91:a-small}
\]
\end{lemma}
\begin{proof}
The spectral expansion, in the metric normalization just stated, is
\[
 p_s(\theta)=\sum_{\ell=0}^\infty
    (\ell+1)e^{-s\ell(\ell+2)}
                 \frac{\sin((\ell+1)\theta)}{\sin\theta}.
 \tag{V91.5}\label{r91:heat-series}
\]
The apparent endpoint singularities are removable. Uniform convergence
with derivatives for $s$ in a compact subset of $(0,\infty)$ follows
from the quadratic eigenvalues. Positivity is the usual positivity
of the heat semigroup on the connected compact sphere.

For an elementary monotonicity argument, start the heat equation at
the smooth radial probability density proportional to $e^{j\cos\theta}$.
Its derivative is nonpositive. The radial equation is
$v_s=v_{\theta\theta}+2\cot\theta\,v_\theta$; hence
$w=v_\theta$ satisfies
$w_s=w_{\theta\theta}+2\cot\theta\,w_\theta-2\csc^2\theta\,w$,
with $w=0$ at both poles. The maximum principle excludes a positive
interior maximum, so $w\le0$. The singular coefficients at the poles
cause no positive-maximum exception since smooth radial solutions have
zero derivative there. Letting $j\to\infty$ gives the heat kernel and
its monotonicity. Thus its extrema are $M_s,m_s$.

At the two poles, \eqref{r91:heat-series} gives
\[
 M_s=e^s\sum_{k=1}^\infty k^2e^{-sk^2},\qquad
 m_s=e^s\sum_{k=1}^\infty(-1)^{k-1}k^2e^{-sk^2}.
 \tag{V91.6}\label{r91:poles}
\]
For $0<s\le1$, comparison on intervals $[k-1,k]$ gives
\[
 M_s\le e^s\int_0^\infty(x+1)^2e^{-sx^2}dx
      \le e^s s^{-3/2}\left(\frac{3\sqrt\pi}{4}+1\right).
 \tag{V91.7}\label{r91:upper}
\]
Poisson summation for the alternating Gaussian, followed by
differentiation, gives the exact positive representation in this range
\[
 m_s=\frac{e^s\sqrt\pi}{2s^{3/2}}
       \sum_{j\in\mathbb Z}
       \left(\frac{\pi^2(j+1/2)^2}{s}-\frac12\right)
                    e^{-\pi^2(j+1/2)^2/s}.
 \tag{V91.8}\label{r91:poisson}
\]
Indeed the alternating theta sum is
$\sum_{k\in\mathbb Z}(-1)^ke^{-sk^2}
=\sqrt{\pi/s}\sum_{j\in\mathbb Z}e^{-\pi^2(j+1/2)^2/s}$,
and its derivative times $e^s/2$ is $m_s$.
All differentiated series converge locally uniformly for $s>0$.
Every term in \eqref{r91:poisson} is positive when $s\le1$.
Keeping $j=0,-1$ and putting $z=\pi^2/(4s)$ yields
\[
 m_s\ge e^s\sqrt\pi\,s^{-3/2}(z-\tfrac12)e^{-z}.
 \tag{V91.9}\label{r91:lower}
\]
Since $z\ge\pi^2/4$ and
$3/4+1/\sqrt\pi<\pi^2/4-1/2$, division of
\eqref{r91:upper} by \eqref{r91:lower} proves the first branch
of \eqref{r91:heat-ratio}.

For $s\ge1$, the character bound
$|\sin((\ell+1)\theta)/\sin\theta|\le\ell+1$ implies
\[
 |p_s(\theta)-1|\le
 d_s:=\sum_{k=2}^\infty k^2e^{-s(k^2-1)}
                  \le5e^{-3s}\le\tfrac14 .
 \tag{V91.10}\label{r91:large-time}
\]
For the middle inequality, use $k^2-4\ge3(k-2)$ for $k\ge3$
and $e^{-3}<1/20$. With $z_0=1/20$,
$\sum_{j\ge1}(j+2)^2z_0^j
=z_0(9-11z_0+4z_0^2)/(1-z_0)^3<1$.
Then
$\log(M_s/m_s)\le2d_s/(1-d_s)\le(40/3)e^{-3s}<14e^{-3s}$.
Taking $s=1$ or $s=\pi/(2\sqrt u)\le1$ proves
\eqref{r91:a-large}. Taking $s=\frac13\log(42/u)\ge1$
proves \eqref{r91:a-small}; letting $s\to\infty$ proves $a(0)=0$.
\end{proof}

\subsection{The native conditional oscillation, without a gap assumption}

Use the actual homogeneous Hamiltonian \eqref{r90:H}, with normalized
positive ground $\Omega$ and energy $E_0$. For a link $e$ let
$y$ denote all exterior links. Delete only the magnetic plaquettes
touching $e$ in a comparison operator:
\[
 \mathcal H=\overline{\mathcal H}_e+W_e,\qquad
 \overline{\mathcal H}_e=\kappa L_e+\mathcal H_{\mathrm{out},e},
 \qquad 0\le W_e=\sum_{p\ni e}b(1-w_p)\le8b.
 \tag{V91.11}\label{r91:deletion}
\]
Here $\mathcal H_{\mathrm{out},e}$ acts only on the exterior. This
deletion is inside a proof; the native operator and vacuum are unchanged.

\begin{theorem}[All-exterior native Harnack bound]
\label{r91:harnack}
Put $u=8b/\kappa$. For every spatial volume $N\ge3$, every exterior
$y$, and every $q,q'\in S^3$,
\[
 \left|\log\frac{\Omega(q,y)}{\Omega(q',y)}\right|\le a(u).
 \tag{V91.12}\label{r91:harnack-bound}
\]
In particular, when $b/\kappa\ge\pi^2/32$ the right side is at most
$2\pi\sqrt{2b/\kappa}$. No many-link gap, trial/native density
identification, or exterior probability exception is assumed.
\end{theorem}
\begin{proof}
Feynman--Kac with the bounded multiplier $0\le W_e\le8b$ gives,
for every nonnegative $F$ and $t>0$, the positivity-order inequalities
\[
 e^{-8bt}e^{-t\overline{\mathcal H}_e}F
      \le e^{-t\mathcal H}F
      \le e^{-t\overline{\mathcal H}_e}F.
 \tag{V91.13}\label{r91:semigroup-order}
\]
These are pointwise inequalities on the positive cone, not an
assertion of operator monotonicity of the exponential in the
quadratic-form order. They follow by bounding the path multiplier
$e^{-\int_0^tW_e(U_s)ds}$ between $e^{-8bt}$ and one.

Take $F=\Omega$. The two summands of
$\overline{\mathcal H}_e$ act on separate factors. Define
\[
 F_t(\,\cdot,y)=(e^{-t\mathcal H_{\mathrm{out},e}}\Omega)
                                                   (\,\cdot,y)>0 .
\]
The eigenfunction equation gives
\[
 e^{-8bt}\int p_{\kappa t}(q,x)F_t(x,y)dx
 \le e^{-tE_0}\Omega(q,y)
 \le\int p_{\kappa t}(q,x)F_t(x,y)dx .
 \tag{V91.14}\label{r91:heat-sandwich}
\]
Use the upper inequality at $q$, the lower at $q'$, and
$m_{\kappa t}\le p_{\kappa t}\le M_{\kappa t}$.
The positive integral of $F_t(x,y)$ cancels, giving the ratio bound
$e^{8bt}M_{\kappa t}/m_{\kappa t}$.
Interchange $q,q'$ and then infimize over $s=\kappa t>0$.
Lemma~\ref{r91:heat} gives the stated explicit specialization.
All kernels and $\Omega$ are smooth at this finite regulator, so
the bound holds everywhere, not only almost everywhere.
\end{proof}

With nonnegative nonuniform magnetic coefficients $b_p$ and positive
electric coefficient $\kappa_e$, the same proof uses
$u_e=2\sum_{p\ni e}b_p/\kappa_e$ and leaves all exterior couplings
arbitrary. Only the incident potential budget enters.

\subsection{Native conditional coercivity and an exact centered comparison}

Define the actual marginal, conditional amplitude, and conditional law
\[
 \rho_e(y)^2=\int\Omega(q,y)^2dq,\quad
 \zeta_e(q,y)=\frac{\Omega(q,y)}{\rho_e(y)},\quad
 \nu_e^y(dq)=\zeta_e(q,y)^2dq,\quad d\nu=\Omega^2dU.
 \tag{V91.15}\label{r91:conditional}
\]
Let $\mathbb E_e$ be conditional expectation under $\nu$ given $y$.
It is an orthogonal projection in $L^2(\nu)$.

\begin{theorem}[Conditional floor on the actual ground law]
\label{r91:conditional-floor}
Set
\[
                   \eta_{\kappa,b}=3\kappa e^{-2a(8b/\kappa)}>0.
 \tag{V91.16}\label{r91:eta}
\]
For every exterior and smooth $f$,
\[
 \eta_{\kappa,b}\Var_{\nu_e^y}(f)
                  \le\kappa\int|\nabla_e f|^2d\nu_e^y .
 \tag{V91.17}\label{r91:conditional-poincare}
\]
On the full Haar Hilbert space define
\[
 (P_e^{\rm nat}v)(q,y)=
       \zeta_e(q,y)\int\zeta_e(x,y)v(x,y)dx,\qquad
 G_N=\sum_e(I-P_e^{\rm nat}).
 \tag{V91.18}\label{r91:projectors}
\]
Then $G_N$ is a bounded positive operator with kernel
$\operatorname{span}\{\Omega\}$, and the following closed-form
inequality is valid on both the full and physical spaces:
\[
             \mathcal H-E_0\ \ge\ \eta_{\kappa,b}G_N .
 \tag{V91.19}\label{r91:native-comparison}
\]
There is no unaccounted additive energy shift in this comparison.
\end{theorem}
\begin{proof}
The maximum/minimum density ratio of $\nu_e^y$ is at most
$e^{2a(u)}$, by \eqref{r91:harnack-bound}; the conditional
normalizer cancels in this ratio. Haar measure on the unit $S^3$
has Poincare gap $3$, from the eigenvalues $\ell(\ell+2)$.
If $p_-\le d\nu_e^y/dq\le p_+$, then
\[
 \Var_{\nu_e^y}(f)
 \le p_+\Var_{dq}(f)
 \le\frac{p_+}{3}\int|\nabla_e f|^2dq
 \le\frac{e^{2a(u)}}3\int|\nabla_e f|^2d\nu_e^y .
\]
This proves the conditional claim using the inherited elementary
density comparison, now with a native, volume-independent input.

Integrate in $y$ and sum in $e$. The exact native ground transform
of V90 gives
\[
 \begin{split}
 \langle\Omega f,(\mathcal H-E_0)\Omega f\rangle
 &=\kappa\sum_e\int|\nabla_e f|^2d\nu\\
 &\ge\eta_{\kappa,b}\sum_e
                  \|(I-\mathbb E_e)f\|_{L^2(\nu)}^2.
 \end{split}
 \tag{V91.20}\label{r91:sum}
\]
Multiplication by $\Omega$ conjugates $\mathbb E_e$ to
$P_e^{\rm nat}$, as follows directly by inserting
$\Omega=\rho_e\zeta_e$. Hence each $P_e^{\rm nat}$ is an
orthogonal projection and annihilation of the right side is
equivalent to $f=\mathbb E_e f$ for every $e$.
Such an $f$ is independent of every coordinate and therefore
constant: $\nu$ has a strictly positive smooth density, so the
coordinate measurability statements also hold under product Haar.
This identifies the common kernel. Gauge covariance of $\Omega$
and invariance of the conditional Haar integral show that every
$P_e^{\rm nat}$ commutes with the gauge action. Smooth
approximation extends \eqref{r91:sum} to closed forms.
\end{proof}

This comparison differs essentially from V88's frozen-rotor one.
Here $P_e^{\rm nat}\Omega=\Omega$ exactly, but the paid local
coercivity is exponentially small in $\sqrt{b/\kappa}$.
The projection is not the V88 magnetic eigenprojection and is not
V89's tempered trial projection.

Writing $h_N$ for the gap of $G_N$ on the physical
$\Omega^\perp$ subspace, \eqref{r91:native-comparison} implies
\[
                 \operatorname{gap}_{\rm phys}(\mathcal H)
                         \ge\eta_{\kappa,b}h_N.
 \tag{V91.21}\label{r91:gap-reduction}
\]
No positive lower bound on $h_N$ uniform in $N$ is asserted.
Although $P_e^{\rm nat}$ integrates just one link, its coefficients
$\zeta_e(q,y)$ may depend on \emph{all} exterior links.
Boundedness and a common vacuum do not make $G_N$ a finite-range
spin Hamiltonian. The finite-size certificates screened in V90
therefore still require a locality or conditional-mixing estimate.

\subsection{Blocks and the retained V90 density defect}

For a nonempty block $B$ of $k$ links, changing one coordinate at a
time in \eqref{r91:harnack-bound} gives
\[
       \operatorname{osc}_{U_B}\log\Omega(U_B,y)\le k a(u).
\]
The same density comparison, now using the product-Haar gap $3$,
therefore proves the all-exterior block bound
\[
 3\kappa e^{-2ka(u)}
       \Var_{\nu(\cdot\mid U_{B^c})}(f)
       \le\kappa\int\sum_{e\in B}|\nabla_e f|^2
                                      d\nu(\cdot\mid U_{B^c}).
 \tag{V91.22}\label{r91:block}
\]
For a family of blocks of size at most $k$, with each link in at
most $d$ blocks, summing gives a centered comparison with their
native conditional projections and coefficient
$3\kappa e^{-2ka(u)}/d$. The cost grows with block cardinality.
It is not a volume-independent large-block mixing theorem.

The actual correction $r=\Omega/\Psi_{\tau_{\rm P}}$ of V90 is
also now controlled on each coordinate in a limited but precise sense:
\[
 \operatorname{osc}_{q_e}\log r(q_e,y)
                \le a(8b/\kappa)+\frac{4b}{\tau_{\rm P}} .
 \tag{V91.23}\label{r91:local-r}
\]
Indeed the corresponding oscillation of $V$ is at most $8b$,
so that of $\log\Psi_{\tau_{\rm P}}$ is at most $4b/\tau_{\rm P}$.
Subtract the two logarithms and use the triangle inequality.
This pays a local density window; it does not bound a mixed
exterior derivative, a spatial influence matrix, or the global
oscillation of $r$ independently of volume.

\subsection{Physical scaling and the exact remaining task}

For $b/\kappa\ge\pi^2/32$, the explicit lower bound is
\[
 \eta_{\kappa,b}\ge3\kappa
                   e^{-4\pi\sqrt{2b/\kappa}}.
 \tag{V91.24}\label{r91:eta-explicit}
\]
With the normalization $\kappa=g^2/(2c a_s)$,
$b=c/(g^2a_s)$ from V85, this becomes
\[
 \eta_{\kappa,b}\ge\frac{3g^2}{2c a_s}e^{-8\pi c/g^2}.
 \tag{V91.25}\label{r91:physical}
\]
An exponential in $1/g^2$ is not, by itself, dimensional
transmutation or a positive continuum mass: neither a uniform
$h_N$ nor the necessary positive cofinal value of the complete
lower bound has been proved.

The new lower comparison uses the original vacuum, all exteriors,
and the correct ground subtraction. The next unresolved estimate
is spatial control of the native conditional projectors, or an
alternative direct neutral coercivity argument. Replacing
$\zeta_e$ by a local trial conditional would lose the exact kernel
property which makes \eqref{r91:native-comparison} useful.
The interacting four-dimensional construction and the full YM
mass gap remain unproved.

The diagnostic \path{scripts/verify_ym_restart_v91_native_conditionals.py}
checks exact constants and optimization, both heat-kernel pole
representations, and positive-semigroup/conditional-projection
comparisons in a separate finite fixture. That fixture is not a
Wilson lattice and its numerical gap is not a YM estimate.
The analytic proof, not a grid of heat values, supplies uniformity.
V90's algebra checks are rerun. The preceding body is preserved
byte for byte except its title; only the latest active source is
retained, with builds and ledgers outside \texttt{latex\_notes}.
% END CONSOLIDATED V91 APPEND

% BEGIN CONSOLIDATED V92 APPEND
\clearpage
\section{V92: Native score response and summed conditional sensitivity}
\label{r92:section}

\subsection{The upstream estimate, and the archive audit}

V91 constructs conditional projections which preserve the actual
vacuum, but does not control their spatial dependence. We now pay
one such control: a volume-independent, native-law average of the
\emph{sum over all exterior links} of their squared derivatives.
The input is an exact inverse-energy force response and a mixed
log-ground derivative identity. No trial law, deleted magnetic
interaction, large-volume limit, or assumed neutral gap enters.

The distinction from earlier work matters. Archive f1 V17 gives
the first Fisher-energy identity, not the following second-derivative
row. Archive f7 V11 gives the general moving-ground Poisson equation;
f6 V92 gives conditional-score differentiation and warns that its
average is not an essential supremum. Archive f12 V62 and f16.1 V13
already give conditional-angle/block-gap assembly. V85 controls an
energy-weighted \emph{physical local-energy} bank, and V87 an
inverse-energy bank at a cut. We reuse that algebraic architecture,
but evaluate the new force bank in the fully coupled native vacuum.
The conditional-variance identity alone is not claimed as new.

For literature context, a targeted check of J.~Stubbe,
\href{https://arxiv.org/html/2605.24694v1}{\emph{Sum rules and a
second order Feynman--Hellmann theorem for abstract operators with
applications}}, arXiv:2605.24694v1, Introduction and equations
(1.15)--(1.20), confirms the standard second-order response identity.
Our link translations are isospectral; the needed specialization is
proved directly below. No general eigenvalue bound in that paper
supplies the YM lower gap. The regular companion and broad-literature
checkpoints remain V95 and V100 respectively, following V90.

\subsection{An exact force-response matrix in the original vacuum}

Keep the homogeneous cubic torus, $N\ge3$,
\[
 H=\kappa\sum_e L_e+V,\qquad V=b\sum_p(1-w_p),\qquad
 K=H-E_0,\qquad \nu=\Omega^2\,dU ,
 \tag{V92.1}\label{r92:setup}
\]
where $\Omega>0$ is the normalized smooth actual ground, $\kappa>0$,
and $b\ge0$. Let $X_{e,a}$, $a=1,2,3$, be the unit-speed
infinitesimal left-multiplication fields on link $e$. They are
skew-adjoint in Haar measure, commute with every $L_f$, and
$L_e=-\sum_a X_{e,a}^2$. Fields on different links commute.
On one link their brackets have structure constants of absolute
size $2$; the sign depends on the invariant-field convention.
Write $i=(e,a)$, and set
\[
 \ell=\log\Omega,\qquad s_i=X_i\ell,\qquad
 F_i=(X_iV)\Omega,\qquad R=(K|_{\Omega^\perp})^{-1}.
 \tag{V92.2}\label{r92:sources}
\]
At each finite regulator $R$ exists by compact ellipticity and
simplicity. Its full operator norm is not assumed uniform.
All displayed differentiations act on smooth functions on a compact
manifold; in particular the sources and the differentiated ground
belong to the required operator domains.

\begin{theorem}[Native inverse-energy response and a complete bank bound]
\label{r92:force}
The scores and forces have zero means in the appropriate sense,
$\E_\nu s_i=0$ and $\langle\Omega,F_i\rangle=0$, and
\[
 K(X_i\Omega)=-F_i,\qquad RF_i=-X_i\Omega.
 \tag{V92.3}\label{r92:poisson}
\]
The real symmetric positive semidefinite matrix
\[
 M_{ij}:=\langle F_i,RF_j\rangle
   =\kappa\sum_{f,c}\E_\nu[(X_{f,c}s_i)(X_{f,c}s_j)]
   =\frac12\E_\nu[X_iX_jV]
 \tag{V92.4}\label{r92:M}
\]
satisfies, on the entire $3|E|$-dimensional coefficient space,
\[
                         0\preceq M\preceq8b I .
 \tag{V92.5}\label{r92:bank}
\]
For distinct links sharing no plaquette, its corresponding block
vanishes exactly.
\end{theorem}
\begin{proof}
Haar integration gives $\langle\Omega,X_i\Omega\rangle=0$.
Commuting $X_i$ through the ground equation proves the first identity
in \eqref{r92:poisson}. Pairing it with $\Omega$ proves the force
centering; inversion on $\Omega^\perp$ proves the second identity.
Consequently
\[
 \langle F_i,RF_j\rangle
 =\langle X_i\Omega,KX_j\Omega\rangle
 =-\E_\nu[s_iX_jV].
\]
The exact ground transform gives the middle expression in
\eqref{r92:M}. Finally,
$\int X_i(\Omega^2X_jV)\,dU=0$ proves the last expression.
This also proves its symmetry when $i,j$ are noncommuting fields
on the same link. Positivity follows from $R\ge0$.

For real coefficients $t_{e,a}$, put
$X_t=\sum_{e,a}t_{e,a}X_{e,a}$ and
$|t_e|^2=\sum_a t_{e,a}^2$.
Differentiate a single literal plaquette word twice along the flow
of $X_t$. A link or its inverse has first matrix derivative norm
$|t_e|$ and second derivative norm $|t_e|^2$. Unitarity of all other
factors and $|\operatorname{Tr}A|/2\le\|A\|$ give
\[
 |X_t^2 w_p|\le\left(\sum_{e\in p}|t_e|\right)^2
                 \le4\sum_{e\in p}|t_e|^2 .
\]
Every link belongs to four plaquettes. Thus
$|X_t^2V|\le16b\sum_e|t_e|^2$. Equation \eqref{r92:M} proves
\eqref{r92:bank}; real symmetry extends it to complex coefficients.
If no plaquette contains both differentiated links,
$X_iX_jV=0$ pointwise, proving the support assertion.
\end{proof}

This finite-range \emph{response matrix} is not the Hessian of
$\ell$, nor the covariance of arbitrary physical sources. In particular
\eqref{r92:M} is not permission to remove its inverse-energy weight.
The isospectral flow $e^{tX_i}He^{-tX_i}$ is another way to see why
its zero ground-energy curvature leaves precisely this identity.

\subsection{The native mixed-score row, with its same-link cost retained}

Define the actual Fisher quantity
$I_e=\sum_a\E_\nu|s_{e,a}|^2$ and, for $f\ne e$,
\[
 A_{ef}=\sum_{a,c}\E_\nu|X_{f,c}X_{e,a}\ell|^2,\qquad A_{ee}=0 .
 \tag{V92.6}\label{r92:A}
\]
These are squared mixed derivatives, not squared entries of their
expectation. In particular a vanishing mean derivative is not enough
to make $A_{ef}$ zero.

\begin{theorem}[A summed second-derivative budget without a neutral gap]
\label{r92:mixed}
The matrix $A$ is entrywise nonnegative and symmetric. For every link,
\[
 \begin{split}
 \sum_{f,a,c}\E_\nu|X_{f,c}X_{e,a}\ell|^2
     &=\frac{3b}{2\kappa}\sum_{p\ni e}\E_\nu w_p,\\
 \sum_{f\ne e}A_{ef}+2I_e
     &\le\frac{3b}{2\kappa}\sum_{p\ni e}\E_\nu w_p
       \le\frac{6b}{\kappa}.
 \end{split}
 \tag{V92.7}\label{r92:mixed-row}
\]
Consequently $\|A\|_{\ell^2\to\ell^2}\le6b/\kappa$.
For $b>0$ the incident expected Wilson-loop sum in this formula is
nonnegative. No pointwise positivity of an individual $w_p$ is used.
\end{theorem}
\begin{proof}
Sum the diagonal entries of \eqref{r92:M} over the three fields on
$e$. A plaquette containing $e$ is a fundamental matrix coefficient
in that link, so $\sum_a X_{e,a}^2 w_p=-3w_p$. Hence
$\sum_a X_{e,a}^2V=3b\sum_{p\ni e}w_p$, proving the first line.

For fixed configuration set $D_{ca}=X_{e,c}X_{e,a}\ell$.
Its antisymmetric part has entries
$\pm\sum_d\epsilon_{cad}s_{e,d}$. Its squared Frobenius norm is
$2\sum_d|s_{e,d}|^2$. Symmetric and antisymmetric matrices are
orthogonal, so $\sum_{c,a}|D_{ca}|^2\ge2\sum_d|s_{e,d}|^2$.
Subtract this paid same-link cost from the first line.
The upper bound uses only $w_p\le1$ and the incidence count.
For $e\ne f$, commutation of the two link fields gives $A_{ef}=A_{fe}$.
The row bound and symmetry give the operator bound by the Schur test.
The nonnegative left side of the first line gives the sign assertion.
\end{proof}

This is a gain over bounding each $A_{ef}$ separately and then summing
over $f$: the right side has no factor $|E|$. It does not yet give a
uniform tail as the spatial distance between $e$ and $f$ increases.
The elementary first Fisher-energy identity, already in f1 V17,
is not being relabelled as this second-derivative estimate.

\subsection{What charged symmetry does, and what it does not do}

The following application of the inherited gauge-star estimate records
why the force response does not secretly invert a neutral small gap.
For a positive link $e=(x,y)$, both $X_{e,a}\Omega$ and $F_{e,a}$
transform as components of spin one at $x$, and are invariant at all
other vertices. Indeed left multiplication commutes with the right
action at $y$; covariance at $x$ rotates its Lie-algebra index.
Thus the local gauge Casimir is $8$ at $x$ and zero elsewhere.
Components based at different vertices lie in orthogonal isotypic
spaces.

Let $C_x$ be this Casimir. A single degree-six gauge star obeys
\[
 \sum_a|G_{x,a}u|^2\le6\sum_{e\ni x}|\nabla_e u|^2 .
\]
Since $G_{x,a}\Omega=0$, the ground transform proves
$K\succeq(\kappa/6)C_x$ on the full Haar space, just as in V87's
regional calculation. Therefore on each of the present spin-one
sectors,
\[
 K\ge\frac{4\kappa}{3},\qquad
 \|R\|\le\frac{3}{4\kappa}.
 \tag{V92.8}\label{r92:charged}
\]
Their orthogonal direct sum has the same bound. In particular,
for every real or complex deterministic coefficient bank $t$,
\[
 \E_\nu\left|\sum_i t_i s_i\right|^2
 \le\frac{3}{4\kappa}\,t^*Mt
 \le\frac{6b}{\kappa}\sum_i|t_i|^2 .
 \tag{V92.9}\label{r92:score-bank}
\]
The first inequality applies \eqref{r92:charged} to
$\sum_i t_iX_i\Omega$, and the second is \eqref{r92:bank}.
These sectors are gauge charged, not the physical singlet sector.
Contracting colours or differentiating a second time can return to
neutral representations; \eqref{r92:charged} cannot then be reused
without a new sector argument. Gauge-star coercivity itself was
already paid in V87 and f15 V83, not discovered here.

\subsection{Actual conditional projectors: a two-sided averaged sensitivity row}

Use V91's actual conditional normalization and its proved floor:
\[
 \rho_e(y)^2=\int\Omega(q,y)^2\,dq,\qquad
 \zeta_e(q,y)=\frac{\Omega(q,y)}{\rho_e(y)},\qquad
 \eta=3\kappa e^{-2a(8b/\kappa)}>0 .
 \tag{V92.10}\label{r92:conditional}
\]
Write $\bar\nu_e(dy)=\rho_e(y)^2dy$ and
$\nu_e^y(dq)=\zeta_e(q,y)^2dq$.
For $f\ne e$ define
\[
 J_{ef}:=\sum_c\int
       \|X_{f,c}\zeta_e(\cdot,y)\|_{L^2(dq)}^2\,\bar\nu_e(dy),
       \qquad J_{ee}=0 .
 \tag{V92.11}\label{r92:J}
\]
The density and derivatives are the original vacuum's; no frozen rotor
eigenfunction or declared spatial Gibbs amplitude replaces $\zeta_e$.

\begin{theorem}[Summed native conditional sensitivity]
\label{r92:sensitivity}
For $f\ne e$,
\[
 X_{f,c}\zeta_e
   =\zeta_e\bigl(s_{f,c}-\E_{\nu_e^y}s_{f,c}\bigr),\qquad
 J_{ef}=\sum_c\E_{\bar\nu_e}\operatorname{Var}_{\nu_e^y}(s_{f,c})
           \le\frac{\kappa}{\eta}A_{ef}.
 \tag{V92.12}\label{r92:sensitivity-eq}
\]
In particular,
\[
 \sup_e\sum_fJ_{ef}\le\frac{6b}{\eta},\qquad
 \sup_f\sum_eJ_{ef}\le\frac{6b}{\eta},\qquad
 \|J\|_{\ell^2\to\ell^2}\le\frac{6b}{\eta}.
 \tag{V92.13}\label{r92:J-row}
\]
The matrix $J$ need not be symmetric.
\end{theorem}
\begin{proof}
Differentiate $\rho_e^2=\int\Omega^2dq$ under the compact integral.
It gives $X_{f,c}\log\rho_e=\E_{\nu_e^y}s_{f,c}$, and therefore the
first identity. Its squared Haar norm is the stated conditional variance.
V91 proves for every exterior the inequality
$\eta\operatorname{Var}_{\nu_e^y}v\le
\kappa\int|\nabla_e v|^2d\nu_e^y$.
Apply it to $v=s_{f,c}$, sum over $c$, and integrate the exterior
against its actual marginal. Commutation of distinct-link derivatives
identifies the resulting mixed derivative sum with $A_{ef}$.
The symmetric majorant $(\kappa/\eta)A$ has both row and column sums
at most $6b/\eta$. This proves all three bounds.
\end{proof}

These quantities have an exact operator interpretation. In each
$e$-fibre write $P_e^{\rm nat}=|\zeta_e\rangle\langle\zeta_e|$.
For an exterior field $X=X_{f,c}$, normalization and reality give
$\langle\zeta_e,X\zeta_e\rangle=0$. Hence
\[
 XP_e^{\rm nat}
   =|X\zeta_e\rangle\langle\zeta_e|
        +|\zeta_e\rangle\langle X\zeta_e|,\qquad
 \|XP_e^{\rm nat}\|_{\rm op}^2
   =\tfrac12\|XP_e^{\rm nat}\|_{\rm HS}^2
   =\|X\zeta_e\|_2^2 .
 \tag{V92.14}\label{r92:projection-derivative}
\]
To check this, restrict to the span of $\zeta_e$ and $X\zeta_e$;
the two nonzero eigenvalues are $\pm\|X\zeta_e\|_2$.
Thus \eqref{r92:J-row} really bounds a complete exterior derivative
sum for each native rank-one fibre projection, averaged in its
native marginal.

\subsection{Why this is not yet spatial mixing or a gap}

There are two precise losses which must not be hidden.
First, let
\[
 j_e(y)=\sum_{f\ne e,c}\|X_{f,c}\zeta_e(\cdot,y)\|_2^2 .
 \tag{V92.15}\label{r92:je}
\]
We have $\int j_e\,d\bar\nu_e\le6b/\eta$, not a uniform essential
supremum. On smooth sections,
$[X_{f,c},P_e^{\rm nat}]=X_{f,c}P_e^{\rm nat}$.
For the specific vector
$u(q,y)=\rho_e(y)h(y)\zeta_e(q,y)$ one has exactly
\[
 \sum_{f\ne e,c}\|[X_{f,c},P_e^{\rm nat}]u\|_2^2
      =\int |h(y)|^2j_e(y)\,d\bar\nu_e(y),\qquad
 \|u\|_2^2=\int|h|^2\,d\bar\nu_e .
 \tag{V92.16}\label{r92:multiplier}
\]
Thus a bound by $C\|u\|_2^2$ for all such vectors would require
$j_e\le C$ almost everywhere, by testing measurable level sets and
smooth approximation. This is the old f6 V92 multiplier obstruction,
now attached to the actual coefficients just bounded. Averaging is
genuine progress but does not remove that obstruction.

Second, $J$ contains \emph{squared} derivative norms. It is not the
matrix of conditional angles used by the archived block criterion.
For example the abstract nonnegative matrix
$J_{ef}=C/(n-1)$ for $e\ne f$, with zero diagonal, has both sums $C$,
but its entrywise square-root matrix has operator norm
$\sqrt{C(n-1)}$. This follows by applying that constant-row matrix to
the all-ones vector and inspecting its other eigenvalues.
This example is not a YM counterexample; it prevents an invalid
matrix-norm inference. No decay with spatial separation is included
in \eqref{r92:J-row}.

The next model-specific target is a distance-weighted tail or a
form-bound upgrade for these \emph{actual} mixed-score coefficients.
The force-response support in \eqref{r92:M} alone does not supply it:
it concerns the signed Gram of complete score gradients, whereas
$A_{ef}$ resolves nonnegative pieces of a single score's gradient.
Neither cancellations in that Gram nor charged-sector coercivity
establish a neutral-sector lower bound. V91 still reads
$\operatorname{gap}_{\rm phys}(H)\ge\eta h_N$, with a uniform
positive $h_N$ unproved. The interacting four-dimensional continuum
construction and physical mass gap remain unproved.

\subsection{Verification and a genuinely Wilson, but isolated, check}

The diagnostic
\path{scripts/verify_ym_restart_v92_score_response.py}
checks the quaternion invariant-field brackets, the fundamental
Casimir, and the same-link antisymmetric cost exactly. It also checks
the derivative identity in a separate Wilson model with a single
oriented loop of $m$ links and potential $b(1-w)$.
The true positive ground there is $\Omega(U)=\psi(U_1\cdots U_m)$,
where $\psi$ is the positive class ground of
$m\kappa L+b(1-q_0)$ on $\SU(2)$.
Indeed the pullback is an isometry by Haar invariance, every link
Laplacian acts as $L$, and the positive pulled-back eigenfunction
must be the full ground. Every one-link marginal normalizer is
$\rho_e=1$.

Writing $x=q_0$, $r=\psi'/\psi$, direct quaternion differentiation gives,
for every ordered link pair, the squared score-derivative block
\[
 \mathcal B(x)=(r')^2(1-x^2)^2-2rr'x(1-x^2)+r^2(2+x^2).
 \tag{V92.17}\label{r92:loop}
\]
Rotations of the Lie-algebra frames do not alter its Frobenius norm.
Thus the specialized identity is
$m\kappa\E_{\psi^2}\mathcal B=(3b/2)\E_{\psi^2}x$, and for $f\ne e$,
$J_{ef}=\E_{\psi^2}[r^2(1-x^2)]$.
The script verifies the polynomial in \eqref{r92:loop} symbolically,
then checks these quantities using a character truncation for
$m=4$ and four magnetic couplings. These floating spectral integrals
are diagnostics, not certified bounds on an infinite-volume cubic
model and not the proof of \eqref{r92:mixed-row}.
The analytic arguments above supply the volume uniformity.

V91's mathematical checks are replayed. Its entire source body is
retained byte for byte except the version title, and the prior source
is recoverable exactly by removing this append and reversing that title.
The latest active file is V92; build products and ledgers remain
outside \texttt{latex\_notes}.
% END CONSOLIDATED V92 APPEND

% BEGIN CONSOLIDATED V93 APPEND
\clearpage
\section{V93: Weighted native scores and conditional Schur control}
\label{r93:section}

\subsection{What changes beyond the averaged estimate}

V92 bounds the native conditional-projector sensitivity averaged against
its actual exterior law. Its multiplier warning prevents replacing that
average by a bound against arbitrary test functions. Here we prove an
upgrade with an explicit derivative cost:
\[
 \int j_e|h|^2\,d\bar\nu_e
 \le c_0\|h\|_{L^2(\bar\nu_e)}^2+c_1\bar{\mathcal E}_e(h).
 \tag{V93.1}\label{r93:target}
\]
The constants are independent of spatial volume. This is a form bound,
not a zero-derivative multiplier bound. We then transport it to the
complete first-order commutator bank and to the actual native
Hamiltonian's off-diagonal conditional block. No vacuum mismatch or
additive ground-energy subtraction is introduced.

The new input is a pointwise single-link logarithmic-gradient bound,
followed by a weighted score identity. The maximum-principle and
Caccioppoli methods are classical; their use here retains the compact
group frame, the original ground, and the complete exterior derivative
sum. The archive audit rechecked f6 V92's conditional-gradient and
multiplier obstruction, V92's score equation, and V88's different
frozen-fibre Schur block. Searches for a native pointwise score,
Caccioppoli row, and conditional off-diagonal bound found no preceding
version of the estimates below. No external gradient theorem is
substituted for their proof. The regular companion and broad-literature
checkpoints remain V95 and V100, following V90.

\subsection{A pointwise bound with no global dimension factor}

Keep $H,K,\Omega,\nu,X_{e,a},\ell$ and $s_{e,a}$ from V92, with
$N\ge3$, $\kappa>0$ and $b\ge0$. Set
\[
 \mathcal L=\Omega^{-1}K\Omega
    =-\kappa\sum_{f,c}(X_{f,c}^2+2s_{f,c}X_{f,c}),\qquad
 Q_e=\sum_{f,c,a}|X_{f,c}s_{e,a}|^2 .
 \tag{V93.2}\label{r93:L}
\]
In this subsection only, $Q_e$ denotes this nonnegative scalar
derivative sum; the orthogonal complement projection below will be
denoted $\mathsf Q_e$. Let $\mathbf s_e=(s_{e,1},s_{e,2},s_{e,3})$
and $\mathbf F_e=(X_{e,1}V,X_{e,2}V,X_{e,3}V)$.
V92 gives $\mathcal L\mathbf s_e=-\mathbf F_e$, componentwise, and
its frame-commutator calculation gives the pointwise inequality
$Q_e\ge2|\mathbf s_e|^2$.

\begin{theorem}[Uniform single-link logarithmic gradient]
\label{r93:pointwise}
For every configuration and every link,
\[
                    |\nabla_e\log\Omega|\le S:=\frac{2b}{\kappa}.
 \tag{V93.3}\label{r93:S}
\]
More generally the same proof gives
$|\nabla_e\log\Omega|\le\|\nabla_eV\|_\infty/(2\kappa)$.
This bounds one link, not the norm of the gradient on the entire
configuration space.
\end{theorem}
\begin{proof}
The product rule with the sign convention in \eqref{r93:L} yields
\[
 \mathcal L|\mathbf s_e|^2
     =-2\mathbf s_e\cdot\mathbf F_e-2\kappa Q_e .
 \tag{V93.4}\label{r93:bochner}
\]
At a global maximum of the smooth scalar $|\mathbf s_e|^2$ on the
compact configuration space, its drift term vanishes and its negative
Laplacian is nonnegative. Hence at that point
$0\le2|\mathbf s_e|\|\mathbf F_e\|_\infty
-4\kappa|\mathbf s_e|^2$.
If the maximum is zero there is nothing to prove; otherwise divide
by $2|\mathbf s_e|$.
A plaquette trace is a unit linear quaternion functional in each of
its links, so $|\nabla_e w_p|\le1$ when $e\in p$ and is zero otherwise.
Four incident plaquettes give $\|\mathbf F_e\|_\infty\le4b$.
There is no maximum principle in a growing coordinate chart: this
argument is on the full compact manifold at each finite regulator,
with a constant depending only on the one-link frame and incidence.
\end{proof}

In particular the actual one-link oscillation admits the two simultaneous
bounds
\[
 |\ell(q,y)-\ell(q',y)|
       \le\min\{a(8b/\kappa),\,S\,d(q,q')\}.
 \tag{V93.5}\label{r93:osc}
\]
The second follows by integrating \eqref{r93:S} along a shortest
round-sphere geodesic; the first is V91. Define
\[
 \alpha=\min\{a(8b/\kappa),\,2\pi b/\kappa\},\qquad
 \widehat\eta=3\kappa e^{-2\alpha}.
 \tag{V93.6}\label{r93:eta}
\]
V91's conditional-density comparison therefore holds with
$\widehat\eta$ in place of $\eta$:
$\widehat\eta\operatorname{Var}_{\nu_e^y}v
\le\kappa\int|\nabla_ev|^2d\nu_e^y$ for every exterior.
The new floor is at least V91's and satisfies
$0<\widehat\eta\le3\kappa$. In particular it is
$3\kappa(1-O(b/\kappa))$ as $b/\kappa\to0$.
Neither this small-ratio regime nor the bound on one gradient is a
weak-coupling continuum limit.

\subsection{A weighted mixed-score estimate with its exact constants}

Let
$Q_e^{\rm out}=\sum_{f\ne e,c,a}|X_{f,c}s_{e,a}|^2$.
All integrals in the following argument use the native $\nu$.
For any smooth complex $h$ the score equation and integration by parts
give
\[
 \begin{split}
 \kappa\int |h|^2 Q_e\,d\nu
  ={}&-\int |h|^2\mathbf s_e\cdot\mathbf F_e\,d\nu\\
     &-2\kappa\operatorname{Re}\int\overline h
       \sum_{f,c,a}(X_{f,c}h)s_{e,a}(X_{f,c}s_{e,a})\,d\nu .
 \end{split}
 \tag{V93.7}\label{r93:weighted-identity}
\]
It is obtained by testing each component of
$\mathcal L\mathbf s_e=-\mathbf F_e$ against
$|h|^2s_{e,a}$, summing, and using the native Dirichlet form.

Now suppose $h=h(y)$ is independent of link $e$. Integration by parts
in that link alone gives a stronger local forcing expression:
\[
 -\int |h|^2\mathbf s_e\cdot\mathbf F_e\,d\nu
    =\frac12\int |h|^2\Delta_eV\,d\nu
    =\frac{3b}{2}\int |h|^2\sum_{p\ni e}w_p\,d\nu
    \le6b\int|h|^2d\nu .
 \tag{V93.8}\label{r93:forcing}
\]
No derivative of $h$ is lost here: all of its $e$-derivatives are zero.
This improvement is why the constant term below is proportional to
$b$, rather than the cruder $b^2/\kappa$ from $|\mathbf s_e||\mathbf F_e|$.

\begin{theorem}[Uniform exterior weighted row]
\label{r93:weighted}
For $0<\delta<1$ and every smooth exterior $h$,
\[
 \kappa\int |h|^2Q_e^{\rm out}\,d\nu
 \le\frac{6b}{1-\delta}\int|h|^2d\bar\nu_e
   +\frac{\kappa S^2}{\delta(1-\delta)}
                  \int|\nabla_yh|^2d\bar\nu_e .
 \tag{V93.9}\label{r93:Q-form}
\]
Define the closed exterior form and its nonnegative generator by
\[
 \bar{\mathcal E}_e(h)=\kappa\int|\nabla_yh|^2d\bar\nu_e,
 \qquad
 \bar{\mathcal L}_e
       =-\kappa\bigl(\Delta_y+2\nabla_y\log\rho_e\cdot\nabla_y\bigr).
 \tag{V93.10}\label{r93:marginal-form}
\]
With $j_e$ exactly as in V92, the form bound \eqref{r93:target} holds
on the full exterior form domain with
\[
 c_0(\delta)=\frac{6b}{(1-\delta)\widehat\eta},\qquad
 c_1(\delta)=\frac{S^2}{\delta(1-\delta)\widehat\eta}.
 \tag{V93.11}\label{r93:c}
\]
In particular $\delta=1/2$ gives
\[
                  c_0=\frac{12b}{\widehat\eta},\qquad
                  c_1=\frac{16b^2}{\kappa^2\widehat\eta}.
 \tag{V93.12}\label{r93:half}
\]
\end{theorem}
\begin{proof}
Only exterior derivatives of $h$ occur in the second line of
\eqref{r93:weighted-identity}. Pointwise Cauchy--Schwarz bounds its
absolute value by
$2\kappa S|h||\nabla_yh|\sqrt{Q_e^{\rm out}}$.
Young's inequality bounds this by
$\delta\kappa|h|^2Q_e^{\rm out}
+\kappa S^2|\nabla_yh|^2/\delta$.
On the left $Q_e=Q_e^{\rm out}+Q_e^{\rm same}$ with
$Q_e^{\rm same}\ge0$. Drop this same-link contribution,
apply \eqref{r93:forcing}, and absorb the $\delta$ term.
Because $h$ is exterior measurable, disintegration proves
\eqref{r93:Q-form}.

The conditional Poincare estimate in \eqref{r93:eta}, applied to
every exterior score, gives pointwise in $y$
\[
       j_e(y)\le\frac{\kappa}{\widehat\eta}
                         \E_{\nu_e^y} Q_e^{\rm out}.
 \tag{V93.13}\label{r93:conditional-Q}
\]
Multiply by $|h(y)|^2$, integrate, and use
\eqref{r93:Q-form}. This proves \eqref{r93:target} with
\eqref{r93:c}. Smooth functions are dense in the weighted $H^1$
form domain: the finite-regulator density is smooth and strictly
positive on a compact manifold. Applying the estimate to differences
of approximants extends multiplication by $\sqrt{j_e}$ continuously
to that domain. This proves the asserted closed-form statement,
without needing a uniform lower bound on the global density.
\end{proof}

Thus $M_{j_e}\preceq c_0 I+c_1\bar{\mathcal L}_e$ as forms.
For example, for every $\mu>0$,
\[
 \|\sqrt{j_e}\,(\bar{\mathcal L}_e+\mu)^{-1/2}\|^2
                      \le\max\{c_0/\mu,c_1\}.
 \tag{V93.14}\label{r93:resolvent}
\]
Indeed the scalar function $(c_0+c_1\lambda)/(\mu+\lambda)$ is
bounded by this maximum for $\lambda\ge0$; functional calculus
applies after the form bound. This is genuine operator control
with one exterior derivative paid, not a claim that
$\|\sqrt{j_e}\|_\infty$ is bounded uniformly.

\subsection{The first-order cancellation and the complete adjoint source bank}

Fix $e$ and write $P=P_e^{\rm nat}$, $\mathsf Q=I-P$.
Use the isometry from the actual marginal space
\[
 U_e:L^2(\bar\nu_e)\longrightarrow L^2(dU),\qquad
 U_eh=\Omega h,\qquad P=U_eU_e^* .
 \tag{V93.15}\label{r93:U}
\]
On the full Haar space define the ground-annihilating derivatives
\[
 D_{f,c}=X_{f,c}-s_{f,c},\qquad
 K=\kappa\sum_{f,c}D_{f,c}^*D_{f,c}.
 \tag{V93.16}\label{r93:D}
\]
The last equality is the ground-state form identity, with adjoints
taken in Haar measure. These $D$'s are not a redefinition of the
physical gauge connection.

For $f\ne e$, put
$z_{f,c}=X_{f,c}\zeta_e
=\zeta_e(s_{f,c}-\E_{\nu_e^y}s_{f,c})$.
V92 gives $\langle\zeta_e,z_{f,c}\rangle=0$. A direct fibre
calculation now gives
\[
 \begin{aligned}
 [D_{f,c},P]&=2|\zeta_e\rangle\langle z_{f,c}|=:\mathcal C_{f,c},
       &\mathcal C_{f,c}P&=0,\\
 D_{f,c}U_eh&=U_eX_{f,c}h,
       &D_{e,a}U_eh&=0 .
 \end{aligned}
 \tag{V93.17}\label{r93:triangular}
\]
For the first equality, $[X_{f,c},P]=
|z_{f,c}\rangle\langle\zeta_e|+
|\zeta_e\rangle\langle z_{f,c}|$, whereas
$[s_{f,c},P]=|z_{f,c}\rangle\langle\zeta_e|-
|\zeta_e\rangle\langle z_{f,c}|$.
Their difference proves the factor $2$ and its direction.
The second line follows from $D(\Omega h)=\Omega Xh$.
Thus $\mathsf QD_{f,c}P=0$: ordinary gradient leakage on the range
of $P$ cancels after the actual ground derivative is included.
The nonzero coupling is the opposite triangular block,
$PD_{f,c}\mathsf Q=-\mathcal C_{f,c}$.

Let $\mathcal C$ stack all the $\mathcal C_{f,c}$ with $f\ne e$.
Its adjoint acts on exterior-measurable vector banks by
\[
 \mathcal C^*((U_eh_{f,c})_{f\ne e,c})
   =2\Omega\sum_{f\ne e,c}
          (s_{f,c}-\E_{\nu_e^y}s_{f,c})h_{f,c}(y).
 \tag{V93.18}\label{r93:Cstar}
\]

\begin{theorem}[Form control of the whole returning source bank]
\label{r93:bank}
For arbitrary smooth exterior functions $h_{f,c}$,
\[
 \|\mathcal C^*(U_e\mathbf h)\|^2
 \le4c_0\sum_{f\ne e,c}\|h_{f,c}\|^2
       +4c_1\sum_{f\ne e,c}\bar{\mathcal E}_e(h_{f,c}),
 \tag{V93.19}\label{r93:bank-bound}
\]
where the norms on the right use the actual marginal.
There is no extra number-of-links factor.
\end{theorem}
\begin{proof}
The conditional covariance matrix of the complete exterior score
vector is positive semidefinite and has trace $j_e(y)$. Its quadratic
form is at most $j_e(y)\sum_{f,c}|h_{f,c}(y)|^2$.
Equation \eqref{r93:Cstar} therefore bounds the left side by
$4\sum_{f,c}\int j_e|h_{f,c}|^2d\bar\nu_e$.
Apply \eqref{r93:target} separately to each component and sum.
This uses the natural direct-sum norm of the bank, not an estimate
by the number of its components. The same form-density argument
extends the result to vector banks in the direct-sum form domain.
\end{proof}

\subsection{A native conditional Schur source, with the derivative loss exposed}

Compression has exactly the marginal ground form:
$U_e^*KU_e=\bar{\mathcal L}_e$.
For smooth $h$ let
\[
 B_eh=\mathsf QKU_eh,\qquad
 \mathcal T_{2,e}(h)=
        \sum_{f\ne e,c}\bar{\mathcal E}_e(X_{f,c}h).
 \tag{V93.20}\label{r93:B}
\]
The latter is a summed second-derivative quantity, not assumed
bounded by a multiple of $\bar{\mathcal E}_e(h)$.
The exact triangular identity gives
\[
 B_eh=-\kappa\mathcal C^*((U_eX_{f,c}h)_{f\ne e,c}),
 \qquad
 \|B_eh\|^2
 \le4\kappa c_0\bar{\mathcal E}_e(h)
                      +4\kappa^2c_1\mathcal T_{2,e}(h).
 \tag{V93.21}\label{r93:B-bound}
\]
To verify the first equality, pair $KU_eh$ with a smooth
$v\in\mathsf QL^2$. The ground form and
$PD_{f,c}\mathsf Q=-\mathcal C_{f,c}$ give
\[
 \begin{aligned}
 \langle KU_eh,v\rangle
   &=\kappa\sum_{f\ne e,c}\langle U_eX_{f,c}h,D_{f,c}v\rangle\\
   &=-\kappa\langle\mathcal C^*(U_e\nabla_yh),v\rangle .
 \end{aligned}
\]
Smoothness on the finite compact regulator identifies the resulting
$L^2$ vectors. The bound follows from \eqref{r93:bank-bound} and
$\sum\|X_{f,c}h\|^2=\bar{\mathcal E}_e(h)/\kappa$.
In particular $B_e1=0$ exactly: this is the original vacuum.

Let $K_{\mathsf Q}$ be the self-adjoint operator of the restricted
form of $K$ on $\mathsf QL^2$. The improved V91 comparison gives
$K\succeq\widehat\eta(I-P)$, so
\[
                     K_{\mathsf Q}\ge\widehat\eta I .
 \tag{V93.22}\label{r93:Q-floor}
\]
This is a paid conditional complement floor, not an assumed
many-link neutral gap. Its inverse is bounded. The actual return
and its Schur subtraction, on smooth exterior $h$, are
\[
 \mathfrak r_e(h)=\langle B_eh,K_{\mathsf Q}^{-1}B_eh\rangle,\qquad
 \mathfrak s_e(h)=\bar{\mathcal E}_e(h)-\mathfrak r_e(h)\ge0 .
 \tag{V93.23}\label{r93:Schur}
\]
The inequality follows by minimizing the nonnegative form of $K$
over vectors $U_eh+v$ with $v\in\mathsf QL^2$ in its form domain.
The minimum is attained at $v=-K_{\mathsf Q}^{-1}B_eh$.
Using \eqref{r93:B-bound} gives the additional explicit estimate
\[
 \mathfrak r_e(h)
 \le\frac{4\kappa c_0}{\widehat\eta}\bar{\mathcal E}_e(h)
       +\frac{4\kappa^2c_1}{\widehat\eta}\mathcal T_{2,e}(h).
 \tag{V93.24}\label{r93:Schur-bound}
\]
For $\delta=1/2$, these two coefficients are
$48\kappa b/\widehat\eta^2$ and $64b^2/\widehat\eta^2$.
This is an estimate of the native off-diagonal response, not of
the frozen-fibre source from V88. The complement projection
annihilates the actual ground, and no extensive energy defect occurs.

\subsection{An explicit non-closure test and the next upstream target}

The new form estimates resolve a specific average-to-weighted gap
in V92. They do not establish a strict Schur margin. In fact, for
the entire parameter family \eqref{r93:c},
\[
 \frac{4\kappa c_0(\delta)}{\widehat\eta}
       =\frac{24\kappa b}{(1-\delta)\widehat\eta^2}
       \ge\frac{8}{3(1-\delta)}\,\frac b\kappa .
 \tag{V93.25}\label{r93:failure}
\]
Here only $\widehat\eta\le3\kappa$ was used. If $b/\kappa\ge3/8$,
even the first coefficient is greater than one for every
$0<\delta<1$. The nonnegative second-derivative term cannot
improve that upper majorant. Hence optimizing this absorption
parameter cannot make \eqref{r93:Schur-bound} prove a strict
return bound $\mathfrak r_e\le\theta\bar{\mathcal E}_e$ with
$\theta<1$ in the weak-coupling regime. This is a limitation of
the \emph{estimate}, not a lower bound on the actual return or
a counterexample to a YM mass gap.

Nor may $\mathcal T_{2,e}$ be silently replaced by
$\|\bar{\mathcal L}_eh\|^2$ with a uniform coefficient: the latter
contains derivatives of the unknown native marginal drift.
The first-order cancellation in \eqref{r93:triangular} must be
kept, but it does not remove the adjoint return.
For physical input $U_eh$, $B_eh$ is physical as well because $K$
and $P$ commute with gauge transformations. Thus the contracted
return source is in the singlet sector. V92's spin-one resolvent
floor cannot replace $K_{\mathsf Q}^{-1}$ on this source.

The next useful move is upstream of the scalar replacement
$K_{\mathsf Q}^{-1}\preceq\widehat\eta^{-1}$: control that resolvent
on the actual returning source, or find a spatially resolved
cancellation which avoids its second-derivative loss. Repeating
the averaged estimate or changing $\delta$ cannot pay this step.
The fully interacting four-dimensional continuum construction,
a uniform neutral margin, and the physical YM mass gap remain
unproved.

\subsection{Verification and preservation}

The script
\path{scripts/verify_ym_restart_v93_weighted_score.py}
checks the score equation and product-rule sign symbolically,
the Young-inequality constants, and all factors in
\eqref{r93:half}--\eqref{r93:B-bound}. A separate two-$S^3$ fixture
uses the exact positive ground proportional to $e^{a q_0r_0}$,
with engineered potential $\kappa\Delta\Omega/\Omega$.
It verifies the weighted identities, conditional form bound,
the factor $2$ and direction in \eqref{r93:triangular}, and the
off-diagonal source formula under a non-product ground law.
That potential is explicitly not a Wilson action.
An isolated four-link Wilson loop also checks the pointwise
score scale using a character truncation. Neither numerical
fixture is a many-volume or continuum proof.

V92's mathematical checks are replayed. The preceding source body
is retained byte for byte except the version title, and can be
recovered exactly by reversing that title and removing this append.
Only the latest active V93 source is kept; all builds and checking
records remain outside \texttt{latex\_notes}.
% END CONSOLIDATED V93 APPEND

% BEGIN CONSOLIDATED V94 APPEND
\clearpage
% Widen only the new three-digit section-number slot in the contents.
\makeatletter
\addtocontents{toc}{\protect\makeatletter
  \protect\patchcmd{\protect\l@section}{1.5em}{2em}{}{}
  \protect\patchcmd{\protect\l@subsection}{2.3em}{2.9em}{}{}
  \protect\makeatother}
\makeatother
\section{V94: Positive native feedback and the conditional resolvent metric}
\label{r94:section}

\subsection{The upstream correction to V93's non-absorption}

V93 established the native off-diagonal source but showed that replacing
its resolvent by the conditional scalar floor cannot yield its proposed
strict subtraction estimate at large $b/\kappa$. We now retain a positive
square in the discarded restriction which is built from exactly the same
source. Completing that square gives a positive, source-dependent
lower form without a smallness assumption. The underlying inverse
identity is classical; the new application here is the actual-vacuum
triangular factorization, including the full exterior source bank.

The archived repaired Witten Woodbury channel in
\path{Renewal_Yang_Mills_Mass_Gap_Research_Notes_V135_f1.tex},
V21's corollary \emph{Exact repaired Woodbury channel}, concerns
subtraction of a curvature defect. The Gaussian posterior identity in
\path{Renewal_Yang_Mills_Mass_Gap_Research_Notes_V100_f15.tex},
V99's Gaussian bridge Hessian, concerns a different finite Gaussian form.
Neither supplies the native positive square proved below. We do not
claim a new abstract Woodbury formula. For the related variational
parallel-sum viewpoint, see Z.~Tarcsay,
\emph{On the parallel sum of positive operators, forms, and functionals},
\href{https://arxiv.org/html/1501.01922v1}{arXiv:1501.01922v1},
Introduction~(1.1) and Section~2. All operator identities needed here
are proved explicitly.

Fix the finite periodic cubic regulator $N\ge3$, $n=|E_N|$,
$\kappa>0$, and $b\ge0$. Retain V91--V93's actual positive ground
$\Omega$, $K=H-E_0$, $\nu=\Omega^2dU$, $P=P_e^{\rm nat}$,
$\mathsf Q=I-P$, $\zeta_e$, $\bar\nu_e$, and $U_eh=\Omega h$.
Write $i=(f,c)$ for the $m=3(n-1)$ exterior directions $f\ne e$.
Let $\mathcal C$ be the complete stack in \eqref{r93:triangular},
mapping $\mathsf QL^2(dU)$ into
$\mathcal G_e=\bigoplus_{i=1}^m PL^2(dU)$.
Put $g_h=(U_eX_i h)_i$; thus
$\bar{\mathcal E}_e(h)=\kappa\|g_h\|^2$.
The improved conditional floor is $\widehat\eta$ from
\eqref{r93:eta}, and $S=2b/\kappa$.

\subsection{An exact square in the actual discarded form}

On $\mathsf QL^2(dU)$ define $A_e$ by the closed fibre form
\[
 a_e[v]=\kappa\sum_{a=1}^3\|D_{e,a}v\|^2,\qquad
 A_e=\int^\oplus A_e(y)\,dy,\qquad A_e\ge\widehat\eta I .
 \tag{V94.1}\label{r94:A}
\]
The last inequality is the native conditional Poincare bound.
On each fibre, $A_e(y)$ is the restriction of
$\kappa\sum_a D_{e,a}^*D_{e,a}$ to $\zeta_e(\cdot,y)^\perp$;
it is not the frozen magnetic rotor of V88. The direct-integral form
requires only the $e$ derivatives. The restricted full form
$K_{\mathsf Q}$ has the smaller domain with all link derivatives.

\begin{theorem}[The matching positive square]
\label{r94:square-theorem}
For smooth exterior $h$ and smooth $v\in\mathsf QL^2$,
\[
 q_K[U_eh+v]
 =a_e[v]+\kappa\|g_h-\mathcal C v\|^2
       +\kappa\sum_{i=1}^m\|\mathsf QD_i v\|^2 .
 \tag{V94.2}\label{r94:square}
\]
The identity extends to the respective full form domains. In particular,
as an order of closed forms on $\mathsf QL^2$,
\[
 K_{\mathsf Q}\succeq A_e+\kappa\mathcal C^*\mathcal C .
 \tag{V94.3}\label{r94:feedback-floor}
\]
\end{theorem}
\begin{proof}
V93 gives $D_{e,a}U_eh=0$, $D_iU_eh=U_eX_i h$,
and $PD_i v=-\mathcal C_i v$ for $v\in\mathsf QL^2$.
The $P$ and $\mathsf Q$ parts of $D_i(U_eh+v)$ are therefore
$U_eX_i h-\mathcal C_i v$ and $\mathsf QD_i v$.
They are orthogonal. Sum their squared norms and the $e$ terms
in $q_K=\kappa\sum_{f,c}\|D_{f,c}\,\cdot\,\|^2$.
Set $h=0$ to obtain \eqref{r94:feedback-floor}.

At fixed $N$, all coefficients are smooth on a compact product.
The smooth positive $\Omega$ and its marginal have positive minima;
$P$ preserves the ordinary $H^1$ domain, with a finite constant not
claimed uniform in $N$. Smooth vectors in the range of $\mathsf Q$
are a form core, by projecting smooth approximants.
The operators $\mathcal C_i$ are bounded fibre multipliers at
fixed $N$. Consequently all terms pass to the full form limit.
The right side of \eqref{r94:feedback-floor} is closed on the
larger domain of $a_e$ because $\mathcal C$ is bounded.
This establishes both the inequality and its domain inclusion.
\end{proof}

The same $\mathcal C$ determines the off-diagonal source
$B_eh=-\kappa\mathcal C^*g_h$. Dropping
$\kappa\mathcal C^*\mathcal C$ before bounding this source is the
loss in V93's scalar subtraction. No negative interaction has been
discarded in \eqref{r94:square}.

\subsection{The positive source-adapted inverse weight}

Define a bounded positive operator on $\mathcal G_e$ by
\[
 \mathcal M_e=\kappa\mathcal C A_e^{-1}\mathcal C^* .
 \tag{V94.4}\label{r94:metric}
\]
Boundedness here is at each finite regulator; no uniform essential
supremum is built into this definition.

\begin{theorem}[Native Schur lower bound without return smallness]
\label{r94:short-theorem}
The actual Schur form from \eqref{r93:Schur} satisfies
\[
 \boxed{\quad
 \bar{\mathcal E}_e(h)\ \ge\ \mathfrak s_e(h)\
 \ge\ \kappa\langle g_h,(I+\mathcal M_e)^{-1}g_h\rangle .
 \quad}
 \tag{V94.5}\label{r94:short}
\]
The assertion holds on the marginal $H^1$ form domain. There is no
hypothesis $\|\mathcal M_e\|<1$, and the lower form contains only
first derivatives of $h$.
\end{theorem}
\begin{proof}
Put $L=\sqrt\kappa\,\mathcal C A_e^{-1/2}$, a bounded operator.
The positive form sum and its inverse factor as
\[
 (A_e+\kappa\mathcal C^*\mathcal C)^{-1}
 =A_e^{-1/2}(I+L^*L)^{-1}A_e^{-1/2}.
\]
This follows either by solving its coercive form equation or by
substituting $w=A_e^{1/2}v$ into that equation. In particular it does
not assume $A_e$ bounded. The bounded identity
$(I+LL^*)L=L(I+L^*L)$ gives
\[
 \kappa\mathcal C(A_e+\kappa\mathcal C^*\mathcal C)^{-1}
                    \mathcal C^*
 =L(I+L^*L)^{-1}L^*
 =\mathcal M_e(I+\mathcal M_e)^{-1}.
 \tag{V94.6}\label{r94:pushthrough}
\]
Inverse order for \eqref{r94:feedback-floor} is valid because
both operators are bounded below by $\widehat\eta I$.
For completeness, its quadratic-form proof is
\[
 \langle f,T^{-1}f\rangle
 =\sup_{v\in\mathcal D(q_T)}
       \{2\operatorname{Re}\langle f,v\rangle-q_T[v]\}.
\]
Increasing the form and reducing its domain cannot increase this
supremum. Hence the return obeys
\[
 \mathfrak r_e(h)
 \le\kappa\langle g_h,
                  \mathcal M_e(I+\mathcal M_e)^{-1}g_h\rangle .
 \tag{V94.7}\label{r94:return}
\]
Subtract from $\bar{\mathcal E}_e(h)=\kappa\|g_h\|^2$.
The upper bound in \eqref{r94:short} follows from the
nonnegativity of the exact return.

Equivalently, minimize \eqref{r94:square} after discarding its
last nonnegative term and enlarging the minimization domain to
$\mathcal D(a_e)$. The resulting minimum is
$\kappa\langle g_h,(I+\mathcal M_e)^{-1}g_h\rangle$.
The enlarged-domain minimizer need not minimize the full form;
this explains why \eqref{r94:short} is a lower bound, not an
identification of the exact effective operator.

Finally $B_eh=-\kappa\mathcal C^*g_h$ extends continuously
from marginal $H^1$ to $L^2$ at fixed $N$, so the return and
the displayed inequalities extend by smooth density.
The finite-regulator coercivity established below also shows
that the Schur form is closed on this domain.
\end{proof}

There is no additional centering error: $g_1=0$ and
$\mathfrak s_e(1)=0$ for the original vacuum.
For physical $U_eh$, both $B_eh$ and the exact complement resolvent
remain in the physical singlet sector. The proof did not use the
charged spin-one gap from V92.

\subsection{The conditional Poisson metric and its paid bounds}

Under the componentwise isometry $U_e$, $\mathcal M_e$ is
multiplication by the positive $m$-by-$m$ matrix
\[
 (\mathbf M_e(y))_{ij}
   =4\kappa\langle z_i(\cdot,y),
                   A_e(y)^{-1}z_j(\cdot,y)\rangle_{L^2(dq)},
 \qquad z_i=X_i\zeta_e .
 \tag{V94.8}\label{r94:fibre-matrix}
\]
Indeed, with $\rho_e^2=\int\Omega^2dq$,
\[
 \mathcal C_j^*(U_e h_j)=2\rho_e z_j h_j,\qquad
 \mathcal C_i A_e^{-1}\mathcal C_j^*(U_eh_j)
 =4\rho_e\zeta_e\langle z_i,A_e(y)^{-1}z_j\rangle h_j .
\]
Thus the lower bound can be written entirely in the actual marginal:
\[
 \mathfrak s_e(h)\ge
 \kappa\int (\nabla_yh)^*(I+\mathbf M_e(y))^{-1}
                            \nabla_yh\,d\bar\nu_e(y).
 \tag{V94.9}\label{r94:weighted-form}
\]
The matrix is not asserted spatially sparse. Locality in configuration
space of this comparison form is not finite interaction range in
the link labels, nor is it locality of the exact Schur form.

To specify exactly which inverse appears, set
$\xi_i=s_i-\E_{\nu_e^y}s_i$ and let
\[
 \mathcal L_e^y=-\kappa\left(\Delta_e+
                     2\sum_a s_{e,a}X_{e,a}\right)
\]
act on the mean-zero part of $L^2(\nu_e^y)$.
Multiplication by $\zeta_e$ intertwines this operator with $A_e(y)$.
For $t\in\mathbb C^m$, put $\xi_t=\sum_i t_i\xi_i$. Then
\[
 \begin{aligned}
 t^*\mathbf M_e(y)t
 &=4\kappa\langle\xi_t,(\mathcal L_e^y)^{-1}\xi_t\rangle_{\nu_e^y}\\
 &=4\kappa
 \sup_{\substack{\phi\in H^1(\nu_e^y),\ \E_{\nu_e^y}\phi=0\\
                  \int|\nabla_e\phi|^2d\nu_e^y>0}}
 \frac{|\langle\xi_t,\phi\rangle_{\nu_e^y}|^2}
             {\kappa\int|\nabla_e\phi|^2d\nu_e^y}.
 \end{aligned}
 \tag{V94.10}\label{r94:negative-sobolev}
\]
This is the conditional negative-Sobolev norm of the actual score
source. To prove the last equality, apply Cauchy--Schwarz to
$(\mathcal L_e^y)^{-1/2}\xi_t$ and
$(\mathcal L_e^y)^{1/2}\phi$. Equality is attained at
$\phi=(\mathcal L_e^y)^{-1}\xi_t$ when $\xi_t\ne0$; both sides
vanish when $\xi_t=0$. The conditional floor proves existence.
This variational description does not replace the unknown native
conditional law by a Gibbs trial law.

\begin{proposition}[Mean control and finite-regulator ellipticity]
\label{r94:bounds}
The following bounds require no new many-link spectral assumption:
\[
 \begin{aligned}
 0\le\|\mathbf M_e(y)\|
 &\le\Tr\mathbf M_e(y)\le\frac{4\kappa}{\widehat\eta}j_e(y),\\
 \int\Tr\mathbf M_e\,d\bar\nu_e
 &\le C_*:=\frac{24\kappa b}{\widehat\eta^2},\\
 \sup_y\|\mathbf M_e(y)\|
 &\le C_N:=\frac{16(n-1)b^2}{\kappa\widehat\eta}.
 \end{aligned}
 \tag{V94.11}\label{r94:metric-bounds}
\]
In particular,
\[
 \mathfrak s_e(h)\ge(1+C_N)^{-1}\bar{\mathcal E}_e(h).
 \tag{V94.12}\label{r94:finite-ellipticity}
\]
This last coefficient is volume-dependent.
\end{proposition}
\begin{proof}
Apply $A_e(y)^{-1}\preceq\widehat\eta^{-1}$ to the diagonal
entries in \eqref{r94:fibre-matrix}. Their source norms sum to $j_e$.
For its mean use \eqref{r93:conditional-Q} and V92's
$\kappa\int Q_e^{\rm out}d\nu\le6b$:
$\int j_e\,d\bar\nu_e\le6b/\widehat\eta$.
This is the V92 argument with V93's improved conditional floor.
For the pointwise bound use the conditional variance representation
and the per-link score bound:
\[
 j_e(y)=\sum_{f\ne e}
       \E_{\nu_e^y}|s_f-\E_{\nu_e^y}s_f|^2
       \le\sum_{f\ne e}\E_{\nu_e^y}|s_f|^2
       \le(n-1)S^2 .
\]
All three components of each link are already included in $|s_f|^2$;
there is no further factor three. Substitute $S=2b/\kappa$.
Spectral calculus in \eqref{r94:weighted-form} proves
\eqref{r94:finite-ellipticity}.
Together with $\mathfrak s_e\le\bar{\mathcal E}_e$, this makes
$\mathfrak s_e+\|\cdot\|^2$ equivalent to the closed marginal
$H^1$ norm at fixed $N$, completing the closedness assertion.
\end{proof}

The mean estimate also gives, for each \emph{constant} bank vector $t$,
\[
 \int t^*(I+\mathbf M_e(y))^{-1}t\,d\bar\nu_e
       \ge\frac{|t|^2}{1+C_*}.
 \tag{V94.13}\label{r94:constant-bank}
\]
Indeed matrix-weighted Cauchy--Schwarz bounds $|t|^4$ by the
product of the left side and
$\int t^*(I+\mathbf M_e)t\,d\bar\nu_e\le(1+C_*)|t|^2$.
This does not permit replacing $t$ by the arbitrary
configuration-dependent $\nabla h(y)$ in the resulting coefficient.
Its correlation with large eigenvalues of $\mathbf M_e$ is unpaid.

\subsection{A positive scalar corollary instead of failed subtraction}

\begin{corollary}[A harmonic-mean bound with the derivative cost visible]
\label{r94:harmonic}
For smooth exterior $h$ with $E=\bar{\mathcal E}_e(h)>0$, let
$T=\mathcal T_{2,e}(h)$ from \eqref{r93:B}. With any V93 choice
$0<\delta<1$, set $a=4\kappa c_0/\widehat\eta$ and
$d=4\kappa^2c_1/\widehat\eta$. Then
\[
 \boxed{\qquad
 \mathfrak s_e(h)\ge\frac{E^2}{(1+a)E+dT}.
 \qquad}
 \tag{V94.14}\label{r94:harmonic-bound}
\]
For $\delta=1/2$,
\[
 a=\frac{48\kappa b}{\widehat\eta^2},\qquad
 d=\frac{64b^2}{\widehat\eta^2}.
 \tag{V94.15}\label{r94:constants}
\]
In particular, on any specified class satisfying $T\le\Lambda E$,
one has $\mathfrak s_e\ge(1+a+d\Lambda)^{-1}E$.
No such uniform derivative bound on the full marginal space is claimed.
\end{corollary}
\begin{proof}
Cauchy--Schwarz applied to
$(I+\mathcal M_e)^{-1/2}g_h$ and
$(I+\mathcal M_e)^{1/2}g_h$ gives
\[
 E^2\le
 [\kappa\langle g_h,(I+\mathcal M_e)^{-1}g_h\rangle]\,
 [E+\kappa\langle g_h,\mathcal M_e g_h\rangle].
\]
For the second factor, \eqref{r94:metric}, the conditional floor,
and V93's complete bank estimate imply
\[
 \begin{aligned}
 \kappa\langle g_h,\mathcal M_e g_h\rangle
 &\le\frac{\kappa^2}{\widehat\eta}\|\mathcal C^*g_h\|^2\\
 &\le\frac{4\kappa c_0}{\widehat\eta}E
          +\frac{4\kappa^2c_1}{\widehat\eta}T .
 \end{aligned}
\]
Combine with \eqref{r94:short} and divide by the positive denominator.
If $E=0$, connectedness makes $h$ constant and the conclusion is
interpreted as the trivial zero lower bound, without division.
\end{proof}

The same constants which made V93's upper majorant non-absorbable
now enter a positive denominator. This does not invalidate V93's
non-absorption calculation; it uses the additional matching square
which that calculation discarded. In particular this is not merely
a new choice of its Young parameter.

\subsection{Checks, scope, and the next source-level obligation}

An exact non-Wilson check helps distinguish positivity from smallness.
On two copies of $S^3$ use the actual positive ground
$\Omega(x,y)\propto e^{a x_0y_0}$ of the engineered potential
$\kappa\Delta\Omega/\Omega$. At $y_0=0$ the $x$ conditional is Haar,
$A_e=\kappa L_{S^3}$ on constants' orthogonal complement, and
$z_i=a(X_i y_0)x_0$. Since $L_{S^3}x_0=3x_0$,
$\int x_0^2dq=1/4$, and $\sum_i|X_i y_0|^2=1$ at this point,
the only nonzero eigenvalue of $\mathbf M_e$ is $a^2/3$.
At $a=3$, $I-\mathbf M_e$ has a negative eigenvalue, whereas
$(I+\mathbf M_e)^{-1}$ has smallest eigenvalue $1/4$.
This is an exact conditional calculation, not a Wilson-gap example.

The script
\path{scripts/verify_ym_restart_v94_positive_feedback.py}
checks the constants and scalar square symbolically. Fifteen rectangular
matrix fixtures check the noncommuting inverse identity, extra positive
discarded energy, rank deficiency, and large source norms.
Twelve two-$S^3$ conditional Poisson fixtures use Galerkin dimensions
16, 32, and 48 to check the source metric and its Fisher-floor bound.
Refinement agreement is a diagnostic, not a certified truncation error
or a many-volume proof. V93's mathematical checks are replayed.

The new comparison retains the actual vacuum, the actual one-link
conditional resolvent, and its whole returning source. It supplies
neither a volume-uniform lower bound for the matrix-weighted marginal
form nor a cofinal iteration theorem preserving that bound.
V93's second-derivative quantity must not be identified silently
with $\|\bar{\mathcal L}_e h\|^2$ or bounded by $E$ on all inputs.
Nor may the mean trace estimate be promoted to a uniform multiplier
bound. The next useful target is therefore the native
Poisson metric \eqref{r94:negative-sobolev} on actual physical
gradients, or a justified multiscale treatment of its large directions.
Optimizing the old scalar subtraction is no longer the relevant step.
The interacting four-dimensional continuum construction and its
physical Yang--Mills mass gap remain unproved.

V93's preceding body is preserved byte for byte, apart from the title,
and is exactly recoverable by removing this append and reversing
that title. Only active V94 is retained in \texttt{latex\_notes};
builds and verification records are outside that folder.
The next regular companion review is V95, and the next broad
literature sweep is V100.
% END CONSOLIDATED V94 APPEND

% BEGIN CONSOLIDATED V95 APPEND
\clearpage
\section{V95: Gauge transport and the native physical shape channel}
\label{r95:section}

\subsection{Companion checkpoint and the calculation chosen}

V94 replaced a failed scalar subtraction by a positive inverse metric
on the actual conditional source. The next issue is which directions
of that metric are controlled. This note proves an all-coupling
gauge-transport bound, tests its sharpness in a native isolated Wilson
loop, and computes a nonzero physical shape channel in the original
cubic Wilson Hamiltonian. The last calculation is a fixed-regulator
strong-electric expansion, not a weak-physical-coupling theorem.

The regular V95 checkpoint freshly inventories and hashes the eight
companion sources reviewed at V90. All are unchanged. Their terminal
scope statements were reread: folder NS~V26, SM--Einstein~V72,
Shell Mixing~V16, parallel YM~V5, and the supplied Source Selection~V43,
Exact-Action Bridge~V36, Native Stationarity~V46, and Perturbative
ESM~V51. The first two retain their kernel-norm and fixed-mode/time
qualifications; the shell calculation does not supply a general
noncommuting nonlinear remainder. The supplied action/constraint notes
do not supply a native YM resolvent estimate. No companion theorem is
promoted to the missing continuum or neutral-sector bound.
The checkpoint is recorded outside the source folder in
\path{artifacts/ym_restart_v95_companion_review.json}.

The conceptual transport language is not new. W.~Li and J.~Zhao,
\emph{Wasserstein information matrix},
\href{https://arxiv.org/html/1910.11248v5}{arXiv:1910.11248v5},
Section~2.2, Definition~3 and Proposition~4, express a density-family
metric through a weighted Poisson equation. Here
$X_i\log\nu_e^y=2\xi_i$, in V94's notation, and the solution
$\Phi_i=2\kappa(\mathcal L_e^y)^{-1}\xi_i$ gives exactly
\[
 (\mathbf M_e)_{ij}
      =\int\langle\nabla_e\Phi_i,\nabla_e\Phi_j\rangle\,d\nu_e^y .
 \tag{V95.1}\label{r95:transport-identification}
\]
This follows directly from \eqref{r94:negative-sobolev} by polarization.
We use this familiar interpretation, not the cited paper's statistical
inequalities as YM estimates. The archived f6 V39--V40 coherent
Gibbs-staple response and current V36's co-moving auxiliary Poisson
transport concern different conditional laws. They do not prove
the actual-ground identities below.

\subsection{An exact, coupling-independent gauge-transport estimate}

Retain V94's native notation. Orient $e$ from $x$ to $x'$ and identify
$\SU(2)$ with unit quaternions. For vertex parameters
$\alpha_v\in\mathfrak{su}(2)\simeq\R^3$, the full infinitesimal gauge
action on a link $f:u\to v$ is
$\dot U_f=\alpha_uU_f-U_f\alpha_v$.
In the exterior left-multiplication orthonormal frames, its coefficient
bank $t=G_y\alpha$ has link components
\[
       t_f=\alpha_u-\Ad_{U_f}\alpha_v,\quad f\ne e .
 \tag{V95.2}\label{r95:gauge-bank}
\]
The missing link carries the Killing field
$Y_\alpha(q)=\alpha_xq-q\alpha_{x'}$.
Let $\Pi_y$ denote the orthogonal projection in
$L^2(\nu_e^y;TS^3)$ onto the closed range of the fibre gradient.
That range is closed by the conditional Poincare inequality.

\begin{theorem}[Native gauge pullback of the Poisson metric]
\label{r95:ward}
For every exterior configuration, every $\kappa>0$, $b\ge0$,
and every vertex parameter bank,
\[
 \boxed{\quad
 (G_y\alpha)^*\mathbf M_e(y)(G_y\alpha)
    =\|\Pi_yY_\alpha\|_{L^2(\nu_e^y)}^2
    \le\int|\alpha_x-\Ad_q\alpha_{x'}|^2d\nu_e^y(q).
 \quad}
 \tag{V95.3}\label{r95:ward-bound}
\]
The right side is at most
$2(|\alpha_x|^2+|\alpha_{x'}|^2)$.
With just one nonzero endpoint parameter it is $|\alpha|^2$.
If both endpoint parameters vanish, then
$\mathbf M_e(y)G_y\alpha=0$ exactly.
None of these constants contains volume, coupling, or
$\widehat\eta^{-1}$.
\end{theorem}
\begin{proof}
The uniqueness and positivity of the ground make $\Omega$ gauge
invariant. Its marginal is invariant under the exterior gauge action:
change the integration variable $q$ by
$q\mapsto g_xqg_{x'}^{-1}$ and use Haar invariance.
Thus the normalized conditional ground satisfies
\[
       X_t\zeta_e=-Y_\alpha\zeta_e,\qquad
       \xi_t=-Y_\alpha\log\zeta_e .
 \tag{V95.4}\label{r95:ward-source}
\]
The Killing field has zero Haar divergence. Therefore
$\operatorname{div}_{\nu_e^y}Y_\alpha
=2Y_\alpha\log\zeta_e=-2\xi_t$.
With $\nabla_e^*=-\operatorname{div}_{\nu_e^y}$ and
$Q_y=\nabla_e^*\nabla_e=\mathcal L_e^y/\kappa$, the transport
Poisson equation is
\[
       Q_y\Phi_t=2\xi_t=\nabla_e^*Y_\alpha .
\]
For every fibre $H^1$ test $\phi$ it follows that
$\langle\nabla_e\Phi_t,\nabla_e\phi\rangle
=\langle Y_\alpha,\nabla_e\phi\rangle$.
Consequently $\nabla_e\Phi_t=\Pi_yY_\alpha$.
Equation \eqref{r95:transport-identification} proves the equality
and the projection contraction proves the bound.
The norm of $Y_\alpha(q)$ is
$|\alpha_x-\Ad_q\alpha_{x'}|$. The endpoint estimates follow by
the triangle inequality and $2ab\le a^2+b^2$.
If both endpoint parameters vanish, the source in
\eqref{r95:ward-source} is zero; the positive matrix then annihilates
$G_y\alpha$. All statements use the original normalized conditional,
not a gauge-fixed substitute.
\end{proof}

These are bounds in the vertex-parameter norm, not in an assumed
uniform inverse norm for $G_y$. Gauge stabilizers need not be removed
to state them. At fixed $y$ let
$\mathcal V_y=\operatorname{Ran}G_y$ and
$\mathcal H_y=\ker G_y^*$, using the original exterior product metric.
If $U_eh=\Omega h$ is physical, then $h$ is exterior gauge invariant,
so
\[
                      \nabla_y h(y)\in\mathcal H_y .
 \tag{V95.5}\label{r95:horizontal}
\]
Indeed differentiating that invariance gives
$\langle\nabla h,G_y\alpha\rangle=0$ for every $\alpha$.
Thus \eqref{r95:ward-bound} does not bound the very horizontal
compression of $\mathbf M_e$ needed for all physical inputs.
No smooth global gauge slice or constant orbit rank is assumed.

\subsection{The constant one is sharp in an actual isolated Wilson loop}

Consider the isolated $m$-link loop already used in V92,
with positive normalized class ground $\psi_b(q_0)$ for
\[
             m\kappa L_{S^3}+b(1-q_0).
 \tag{V95.6}\label{r95:loop}
\]
Its full ground is $\psi_b(U_1\cdots U_m)$; after fixing the exterior
product to the identity, its native conditional is $\nu_b=\psi_b^2dq$.
This is the actual Hamiltonian ground, not the thermal density
$e^{-b(1-q_0)/\tau}$.
For a unit left rotation $Y(q)=E_1q$, set
$\xi=-Y\log\psi_b$ and let $\mu_b$ be its scalar Poisson-metric value.
The proof above applies verbatim and gives
\[
  \frac{\langle q_0\rangle_{\nu_b}^2}
       {1-\langle q_1^2\rangle_{\nu_b}}
       \ \le\ \mu_b\ \le\ 1,\qquad
       \lim_{b/\kappa\to\infty}\mu_b=1
       \quad(m\ \hbox{fixed}).
 \tag{V95.7}\label{r95:loop-sharp}
\]
To prove the lower bound, insert the mean-zero test $q_1$ into
the variational formula \eqref{r94:negative-sobolev}.
Weighted integration by parts gives
$\langle\xi,q_1\rangle=\frac12\langle Yq_1\rangle
=\frac12\langle q_0\rangle$, and
$|\nabla q_1|^2=1-q_1^2$.

For the limit, divide \eqref{r95:loop} by $b$.
For every $\varepsilon>0$, a smooth normalized test supported in
$q_0>1-\varepsilon$ shows
$E_0(b)/b\le\varepsilon+C_\varepsilon m\kappa/b$.
Nonnegative kinetic energy gives
$\langle1-q_0\rangle_{\nu_b}\le E_0(b)/b\to0$.
Also $q_1^2\le1-q_0^2\le2(1-q_0)$, so the lower bound tends to one.
This proves the limit without an asymptotic ground-state ansatz.
Removing one edge from the isolated cycle leaves a tree; its
exterior gauge-invariant functions are constant, by successively
gauging its edges to the identity. Thus this sharp gauge channel
is not a physical gapless excitation or a counterexample to a YM gap.

\subsection{A genuinely native Wilson coefficient on the full cubic regulator}

We now evaluate what is present beyond gauge transport.
Put $\beta=b/\kappa$ and
\[
 \kappa^{-1}H=\mathcal K+\beta W,\qquad
 \mathcal K=\sum_f L_f,\qquad
 W=\sum_p(1-w_p).
\]
Throughout this subsection $N$ is fixed and $\beta\downarrow0$.
This is the strong-electric regime, opposite to the
$\beta\to\infty$ weak-physical-coupling regime in V91.
Write the four plaquettes incident to $e$ as
\[
      w_p(q,y)=q\cdot a_p(y),\qquad
      |a_p|=1,\qquad a_\Sigma(y)=\sum_{p\ni e}a_p(y).
 \tag{V95.8}\label{r95:staples}
\]
Let $J_e(y):\R^{3(n-1)}\to\R^4$ be the actual derivative of
$a_\Sigma$ in the exterior orthonormal link frames:
$J_et=\sum_i t_iX_i a_\Sigma$.

\begin{theorem}[First native conditional metric coefficient]
\label{r95:coefficient}
For each fixed regulator, the following expansions hold uniformly
in exterior configuration, with regulator-dependent remainders:
\[
 \begin{aligned}
 \Omega_\beta
    &=1+\frac{\beta}{12}\sum_p w_p+O_N(\beta^2),\\
 \zeta_{e,\beta}(q,y)
    &=1+\frac{\beta}{12}q\cdot a_\Sigma(y)+O_N(\beta^2),\\
 \mathbf M_{e,\beta}(y)
    &=\frac{\beta^2}{432}J_e(y)^*J_e(y)+O_N(\beta^3).
 \end{aligned}
 \tag{V95.9}\label{r95:expansion}
\]
The first two expansions hold in every prescribed finite $C^r$ norm,
with constants depending also on $r$; the last is a matrix operator-norm
expansion. The leading coefficient is supported on the twelve exterior
links of the incident plaquettes and has rank at most four.
Exactly,
\[
 \begin{aligned}
 J_eJ_e^*&=3\sum_{p\ni e}(I_4-a_pa_p^*)
             =12I_4-3\sum_{p\ni e}a_pa_p^*,\\
 \Tr(J_e^*J_e)&=36,\qquad \|J_e^*J_e\|\le12 .
 \end{aligned}
 \tag{V95.10}\label{r95:gram}
\]
In particular,
\[
 \Tr\mathbf M_{e,\beta}=\frac{\beta^2}{12}+O_N(\beta^3),
 \qquad
 \|\mathbf M_{e,\beta}\|\le\frac{\beta^2}{36}+O_N(\beta^3).
 \tag{V95.11}\label{r95:leading-bounds}
\]
These uniform leading coefficients are not uniform remainder bounds.
\end{theorem}
\begin{proof}
The free ground of $\mathcal K$ is the constant $1$, is simple,
and has full-space spectral gap $3$. Multiplication by $W$ is bounded
at fixed $N$. A small complex circle about zero in the resolvent plane
and its Neumann series therefore define an analytic rank-one ground
projection for sufficiently small $|\beta|$. Normalize its eigenvector
by $\|\Omega_\beta\|_2=1$ and positivity for real $\beta$ near zero.
Analyticity in each finite Sobolev order follows iteratively from
\[
 \Omega_\beta=(\mathcal K+1)^{-1}
       (1+E_\beta-\beta W)\Omega_\beta .
\]
The inverse gains two derivatives and the smooth multiplier preserves
each order. Sobolev embedding on the fixed compact product then
gives each stated $C^r$ expansion. No radius uniform in dimension
is asserted.

Each plaquette has Haar mean zero and
$\mathcal K w_p=12w_p$, one fundamental Casimir $3$ for each of its
four distinct links. Differentiating the eigenvalue equation at zero,
with $\langle1,\Omega'_0\rangle=0$, gives
$E'_0=|\mathcal P_N|$ and
$\Omega'_0=\frac1{12}\sum_p w_p$.
In the normalized conditional, division by
$\rho_e=(\int\Omega_\beta^2dq)^{1/2}$ subtracts precisely the
plaquettes not containing $e$. The remaining incident plaquettes
are linear in $q$ with zero fibre mean. This proves the first two
lines of \eqref{r95:expansion}, including their differentiated forms.

For clarity about the moving fibre null space, temporarily define
the full fibre operator
$\widetilde A_\beta=\kappa\sum_aD_{e,a}^*D_{e,a}$ and its null
projection $P_\beta=|\zeta_{e,\beta}\rangle\langle\zeta_{e,\beta}|$.
On the fixed full fibre space the operator
$T_\beta=\kappa^{-1}\widetilde A_\beta+P_\beta$ has domain $H^2(S^3)$.
Its coefficients and $P_\beta$ differ from their $\beta=0$ values
by $O_N(\beta)$ uniformly in $y$. Since
$T_0=L_{S^3}+P_0$ is invertible and its inverse maps $L^2$ to $H^2$,
a Neumann series gives
$T_\beta^{-1}=T_0^{-1}+O_N(\beta)$ in $L^2$ operator norm.
For $z_i=X_i\zeta_{e,\beta}\perp\zeta_{e,\beta}$,
V94's reduced inverse matrix can consequently be evaluated using
$\kappa^{-1}T_\beta^{-1}$.
But
$z_i=\frac{\beta}{12}q\cdot X_i a_\Sigma+O_N(\beta^2)$,
$L_{S^3}(q\cdot v)=3q\cdot v$, and
$\int(q\cdot u)(q\cdot v)dq=u\cdot v/4$.
The coefficient in \eqref{r94:fibre-matrix} is therefore
\[
       4\kappa\left(\frac{\beta}{12}\right)^2
                     \frac1{3\kappa}\frac14
                       =\frac{\beta^2}{432}.
\]
This proves the matrix expansion with the actual conditional inverse.

On a cubic regulator $N\ge3$, the four three-edge return paths
opposite $e$ have twelve distinct edges. Each edge occurs in just
one of these four paths. Holding the other two path links fixed,
its map to $a_p\in S^3$ is a composition of quaternion multiplication
and possibly inversion, hence an isometry. Its differential times
its adjoint is the tangent projection $I_4-a_pa_p^*$.
Three links per path give the factor $3$ in \eqref{r95:gram};
different paths have disjoint derivative banks.
Taking the trace gives $3\cdot4\cdot3=36$, and discarding the
positive $3\sum_pa_pa_p^*$ gives the norm bound $12$.
Nonzero spectra of $J_e^*J_e$ and $J_eJ_e^*$ agree.
Equation \eqref{r95:leading-bounds} follows.
\end{proof}

The factor $12$ in the native first ground correction is important:
it comes from all four electric links of a plaquette. Replacing
the conditional ground by the frozen one-link magnetic rotor would
give a different coefficient. Neither the rank nor the finite-range
support just proved for the leading matrix has been proved for the
full nonlinear remainder. Nor is the full native conditional asserted
to depend only on $a_\Sigma$ beyond this first order.

\subsection{An explicit physical shape direction which survives}

Choose any two of the four incident return paths and put
$h(y)=a_1(y)\cdot a_2(y)$.
Both $a_p$ transform by the same orthogonal left-right quaternion
action at the endpoints. Hence $h$ is exterior gauge invariant and
$g=\nabla_yh$ is a genuine physical horizontal gradient.
The other two paths have disjoint derivative banks and do not
enter $g$.

\begin{proposition}[Nonzero native metric on a physical relative-staple source]
\label{r95:shape}
At every exterior configuration,
\[
 \begin{aligned}
 |g|^2&=6(1-h^2),\\
 J_eg&=3(1-h)(a_1+a_2),\\
 |J_eg|^2&=18(1-h)^2(1+h).
 \end{aligned}
 \tag{V95.12}\label{r95:shape-geometry}
\]
Consequently, when $-1<h<1$,
\[
 \boxed{\qquad
 \frac{g^*\mathbf M_{e,\beta}(y)g}{|g|^2}
      =\frac{\beta^2}{144}(1-h)+O_N(\beta^3).
 \qquad}
 \tag{V95.13}\label{r95:physical-metric}
\]
In particular the physical compression is not identically zero
for sufficiently small positive $\beta$.
\end{proposition}
\begin{proof}
The intrinsic gradients of $a_1\cdot a_2$ on the two unit spheres
are $a_2-ha_1$ and $a_1-ha_2$. Each return path supplies three
isometric copies of its derivative. Thus their squared norms add to
$3(1-h^2)+3(1-h^2)$, while applying the derivative of $a_\Sigma$
gives $3(a_2-ha_1)+3(a_1-ha_2)$.
This proves \eqref{r95:shape-geometry}. Substitution into
\eqref{r95:expansion} and division by $|g|^2$ proves
\eqref{r95:physical-metric}; the matrix-norm remainder remains
$O_N(\beta^3)$ after this division.

The four path products may be assigned independently because their
edges are disjoint. In particular configurations with $h=0$ exist;
there the leading ratio is $\beta^2/144>0$. If desired choose also
$a_3+a_4=0$. At such a configuration,
\[
               d|a_\Sigma|^2[g]=12(1-h^2)>0 .
\]
Gauge rotations preserve $|a_\Sigma|$, so this exhibits an actual
change in conditional staple magnitude, not a hidden endpoint
rotation. It is compatible with the exact horizontality
\eqref{r95:horizontal}.
\end{proof}

This calculation does not say that physical directions produce a
large metric. It says that the all-coupling Ward estimate does not
exhaust the source: the shape channel is explicitly nonzero in the
native Wilson model. A prospective proof cannot close the problem
by declaring all conditional motion to be gauge.

\subsection{Verification, scope, and the next upstream target}

The checker
\path{scripts/verify_ym_restart_v95_gauge_shape.py}
reconstructs literal oriented quaternion plaquettes on sides three
and four. Four rational fixtures verify the derivative Gram,
the twelve-edge incidence, and the relative-staple formulas exactly;
36 gauge-generator tests check the covariance and horizontality
without numerical differentiation. Separate native isolated-loop
character calculations test the angular Poisson metric, its
integration-by-parts source, and the two bounds in
\eqref{r95:loop-sharp}. Those use character cutoff 96 and angular
Galerkin degrees 12, 20, and 28. They are diagnostics without a
certified truncation error. The analytic bounds and limiting argument
are the written proofs above. V94's mathematical checks are replayed.

The new all-coupling result controls gauge transport with a sharp
constant, while the new original-Wilson coefficient identifies a
physical shape contribution and its exact local norm. What remains
is nonperturbative control of that shape contribution on actual
physical gradients. The $O_N(\beta^3)$ statement does not give a
many-volume remainder, and small $\beta$ does not reach the physical
weak-coupling regime. The next acquisition should target the
conditional shape Poisson source itself, possibly through a
source-resolved nonperturbative identity, rather than repeat the
gauge rotation calculation or extrapolate the leading coefficient.
V94's positive Schur comparison is retained, but no uniform neutral
gap, cofinal continuum construction, or full YM mass gap is proved.

V94's body is preserved byte for byte except the version title.
Only active V95 is kept in \texttt{latex\_notes}; the predecessor is
exactly recoverable by reversing the title and removing this append.
All builds and verification records are outside that folder.
The next regular companion checkpoint, broad literature sweep,
and requested archive/restart boundary are at V100.
% END CONSOLIDATED V95 APPEND

% BEGIN CONSOLIDATED V96 APPEND
\clearpage
\section{V96: Native conditional variational budgets and resolvent control}
\label{r96:section}

\subsection{The upstream change: use vacuum minimality before a conditional inverse}

V95 identified a genuine physical shape source, but its explicit
original-Wilson coefficient was restricted to small $b/\kappa$.
V93's nonperturbative weighted source estimate instead paid an
inverse conditional gap before controlling that source. We now go
upstream to the variational principle for the actual positive vacuum.
Retain its exact exterior marginal, insert an arbitrary exterior
multiplier, and vary only the normalized conditional state. Positivity
of the centered full Hamiltonian then bounds the \emph{native}
conditional derivative cost by the cost of a declared trial family.
This produces a weighted all-coupling estimate without an inverse-gap
factor in the source bound.

The factorization, smooth trial family, spherical moments, and Wilson
staple derivative tensor are inherited from V89. They are not new
claims here, and the trial conditional is never identified with the
native one. The new step is the arbitrary-multiplier variational
inequality and its use on the complete native source bank. It differs
from V89's upper bound on a single replacement state's energy.
The latter did not by itself state a weighted multiplier bound.

For provenance, exact conditional factorization and its induced
scalar/vector potentials are classical; see A.~Abedi, N.~T.~Maitra
and E.~K.~U.~Gross,
\href{https://arxiv.org/html/1006.2638v2}{arXiv:1006.2638v2},
Theorem I and equations (4)--(10). That molecular result is not a
Yang--Mills spectral theorem. We use only the underlying product-rule
architecture and prove the compact positive-ground-state statements
needed here directly. This targeted check supplements V89's review;
the scheduled broad sweep and companion checkpoint remain due at V100.

\subsection{The exact marginal equation and weighted trial identity}

Work on the same finite homogeneous cubic torus, $N\ge3$, with
$\kappa>0$, $b\ge0$, and its native differential Hamiltonian
$\mathcal H=-\kappa\Delta+V$. Write
$K=\mathcal H-E_0\succeq0$ and normalize its positive ground
function $\Omega$ in full Haar space. Fix a link $e$, write its
coordinate as $q\in S^3$ and all exterior coordinates as $y$, and
retain exactly the factorization \eqref{r89:actual-factor}:
\[
 \Omega(q,y)=\rho(y)\zeta(q,y),\qquad
 d\bar\nu_e(y)=\rho(y)^2dy,\qquad
 \int_{S^3}\zeta(q,y)^2dq=1.
 \tag{V96.1}\label{r96:factor}
\]
Let $i$ index every exterior link and invariant-frame component,
not only the twelve links in the incident plaquettes. Define
\[
 \begin{aligned}
 i_e(y)&=\sum_{a=1}^3\|X_{e,a}\zeta(\cdot,y)\|_{L^2(dq)}^2,\\
 j_e(y)&=\sum_i\|X_i\zeta(\cdot,y)\|_{L^2(dq)}^2,\\
 W_e(q,y)&=4b-f_e(y)\cdot q,\qquad
 V(q,y)=V_{\rm out}(y)+W_e(q,y),\\
 w_e(y)&=\langle\zeta,W_e\zeta\rangle_{dq},\qquad
 F_e(y)=\kappa(i_e+j_e)+w_e .
 \end{aligned}
 \tag{V96.2}\label{r96:costs}
\]
Here $j_e$ is precisely V92--V95's native conditional sensitivity,
$f_e=b\sum_{p\ni e}a_p$, and $s=|f_e|\le4b$.
In particular $w_e\ge4b-s$. The actual $\zeta$ need not depend on
$y$ only through $f_e$.

Use the native marginal Dirichlet form
\[
 \bar{\mathcal E}_e(h)=\kappa\int\sum_i|X_i h|^2d\bar\nu_e,\qquad
 \bar L_e=-\kappa\sum_i\bigl(X_i^2+
                         2(X_i\log\rho)X_i\bigr).
 \tag{V96.3}\label{r96:marginal}
\]
Its closed form domain is the weighted $H^1$ space. At fixed finite
regulator all coefficients and positive density ratios are smooth
and bounded on the compact product; no uniform lower bound for
$\rho$ as the volume grows is being assumed.

\begin{lemma}[Weighted native competitor identity]\label{r96:identity-theorem}
For any smooth real normalized conditional family $\chi(q,y)$, set
\[
 F_\chi(y)=\kappa\|\nabla_q\chi\|_{dq}^2+
          \kappa\sum_i\|X_i\chi\|_{dq}^2+
          \langle\chi,W_e\chi\rangle_{dq}.
 \tag{V96.4}\label{r96:trial-cost}
\]
The exact marginal equation and the centered trial form are
\[
 \begin{aligned}
 [-\kappa\Delta_y+V_{\rm out}+F_e]\rho&=E_0\rho,\\
 \mathfrak q_K[\rho h\chi]
   &=\bar{\mathcal E}_e(h)+
       \int |h|^2(F_\chi-F_e)\,d\bar\nu_e\ \ge0
 \end{aligned}
 \tag{V96.5}\label{r96:identity}
\]
for every smooth complex exterior $h$, and by density for its
weighted $H^1$ closure. In particular
\[
 \kappa\int(i_e+j_e)|h|^2d\bar\nu_e
 \le\bar{\mathcal E}_e(h)+
                \int|h|^2(F_\chi-w_e)\,d\bar\nu_e .
 \tag{V96.6}\label{r96:variational-budget}
\]
\end{lemma}

\begin{proof}
All conditional scalar products in this proof are in $dq$.
Real normalization gives $\langle\zeta,X_i\zeta\rangle=0$.
Differentiating that identity once more gives
$\langle\zeta,-X_i^2\zeta\rangle=\|X_i\zeta\|^2$.
Multiply $\mathcal H(\rho\zeta)=E_0\rho\zeta$ by $\zeta$ and
integrate out $q$. The terms containing
$2(X_i\rho)\langle\zeta,X_i\zeta\rangle$ vanish, the $q$ kinetic
term gives $\kappa i_e\rho$, and the remaining exterior conditional
term gives $\kappa j_e\rho$. This proves the first line of
\eqref{r96:identity}.

The same normalization identities for $\chi$ imply the exact
compression
\[
 \mathfrak q_K[\rho h\chi]
 =\kappa\int|\nabla_y(\rho h)|^2dy+
   \int(V_{\rm out}+F_\chi-E_0)\rho^2|h|^2dy .
\]
The ground-state product rule for the scalar equation just proved
replaces the expression with $F_\chi$ changed to $F_e$ by
$\bar{\mathcal E}_e(h)$. Subtracting gives the equality in the
second line of \eqref{r96:identity}. Its sign is the full-space
variational principle $K\succeq0$, with no excited-state estimate.
Substitute \eqref{r96:costs} and rearrange to obtain
\eqref{r96:variational-budget}.

For fixed smooth $\chi$, all multiplication terms are bounded at
fixed volume, so smooth density proves the stated form extension.
If $\chi$ is gauge equivariant and $h$ is exterior gauge invariant,
then $\rho h\chi$ is a physical state. Thus the trial families used
below respect the Gauss restriction as well. The proof did not
factor the physical Hilbert space into link spaces.
\end{proof}

The exterior multiplier changes the trial marginal to
$|h|^2d\bar\nu_e$; the \emph{reference} marginal $d\bar\nu_e$ has not
been replaced by an auxiliary Gibbs law. At $h=1$ the marginal is
exactly unchanged. The sign in \eqref{r96:identity} is essential:
it is a variational upper bound on native conditional cost, not
an assertion that the trial conditional solves the native equation.

\subsection{An all-coupling source budget with polynomial coupling dependence}

\begin{theorem}[Native weighted kinetic and shape budget]\label{r96:budget-theorem}
Put
\[
 C_0=\frac94\left(\frac{52}{3}\right)^{1/3},\qquad
 U_0(\kappa,b)=\min\{4b,\ C_0\kappa^{1/3}b^{2/3}\}.
 \tag{V96.7}\label{r96:U}
\]
Then $C_0^3=3159/16$, and for every exterior $h$ in the form domain,
\[
 \boxed{\;
 \kappa\int(i_e+j_e)|h|^2d\bar\nu_e
       \le\bar{\mathcal E}_e(h)+U_0(\kappa,b)\|h\|_{\bar\nu_e}^2 .
 \;}
 \tag{V96.8}\label{r96:main-budget}
\]
The constants are independent of $N\ge3$, the link, and the exterior.
There is no conditional-gap factor in this bound.
\end{theorem}

\begin{proof}
First take $\chi=1$. Its kinetic costs vanish and $F_\chi=4b$.
Writing $m_e(y)=\int q\zeta(q,y)^2dq$, the stronger signed identity
gives
\[
 \kappa\int(i_e+j_e)|h|^2d\bar\nu_e
 \le\bar{\mathcal E}_e(h)+
                    \int|h|^2 f_e\cdot m_e\,d\bar\nu_e
 \le\bar{\mathcal E}_e(h)+4b\|h\|_{\bar\nu_e}^2.
 \tag{V96.9}\label{r96:Haar}
\]
The intermediate integrand need not be nonnegative pointwise.
No pointwise bound on $i_e+j_e$ follows by dropping the exterior
Dirichlet cost.

Next use the globally smooth tempered trial
$\chi_{\tau,f_e}$ of \eqref{r89:tempered-fibre}, with $\tau>0$
constant in $y$. The inherited covariance formula
\eqref{r89:tempered-metric} and exact tensor
\eqref{r89:field-tensor} give
\[
 \kappa\sum_i\|X_i\chi_{\tau,f_e}\|^2
       \le\frac{9\kappa b^2}{\tau^2}.
 \tag{V96.10}\label{r96:trial-outside}
\]
Indeed the complete outside field-derivative trace is $36b^2$;
only the trial, not the native conditional, has twelve-link support.
For $s=|f_e|>0$ put $x=(f_e/s)\cdot q$ and $r=s/\tau$.
Since $|\nabla_q x|^2=1-x^2$, the trial $q$ kinetic term is
\[
 \kappa\|\nabla_q\chi_{\tau,f_e}\|^2
 =\frac{\kappa s^2}{4\tau^2}
       \mathbb E_{\nu_r}(1-x^2)
 \le\frac{4\kappa b^2}{\tau^2}.
 \tag{V96.11}\label{r96:trial-inside}
\]
Here $\nu_r$ is the declared spherical tilt, not a native law.
Its inherited moment bound \eqref{r89:thermal-moments}, together
with $w_e\ge4b-s$, yields
\[
 \langle\chi_{\tau,f_e},W_e\chi_{\tau,f_e}\rangle-w_e
 \le s(1-\mathbb E_{\nu_r}x)\le\frac{3\tau}{2}.
 \tag{V96.12}\label{r96:potential-defect}
\]
At $s=0$ the potential difference is zero and the smooth formulas
remain valid. Thus the possibly large frustration $4b-s$ cancels
in the comparison with the native potential expectation; no
exceptional-exterior probability estimate is needed for this bound.

Combining \eqref{r96:trial-outside}--\eqref{r96:potential-defect}
gives $F_\chi-w_e\le3\tau/2+13\kappa b^2/\tau^2$ at every exterior.
For $b>0$ this majorant has its unique minimum at
$\tau^3=52\kappa b^2/3$, with value
$C_0\kappa^{1/3}b^{2/3}$. This is optimization of the displayed
majorant, not of all possible trial states. Apply
\eqref{r96:variational-budget} and take the smaller bound with
\eqref{r96:Haar}. At $b=0$ the Haar trial already gives the result;
the native vacuum is constant. This proves the theorem.
\end{proof}

The exact thermal integration by parts in \eqref{r89:thermal-ibp}
also permits the sharper trial $q$ bound $3\kappa b/\tau$.
Thus one can replace $U_0$ in \eqref{r96:main-budget} by its minimum
with
\[
 \inf_{\tau>0}
 \left\{\frac{3\tau}{2}+\frac{3\kappa b}{\tau}
                          +\frac{9\kappa b^2}{\tau^2}\right\}.
 \tag{V96.13}\label{r96:refined-budget}
\]
This expression was already V89's all-exterior trial error; its new
use here is to bound the weighted \emph{native} cost. We retain
the simpler explicit $U_0$ in subsequent formulas.

There is a stronger mean estimate at large $b/\kappa$, obtained
without redoing V89's probability argument. Define
\[
 \bar U_0=\min\{U_0,\ 90\sqrt{\kappa b}\}.
 \qquad
 \kappa\int(i_e+j_e)d\bar\nu_e\le\bar U_0.
 \tag{V96.14}\label{r96:mean-budget}
\]
For completeness, set $h=1$ and $\tau=\sqrt{\kappa b}$ in
\eqref{r96:variational-budget}. The $q$ kinetic term and potential
defect total at most $9\tau/2$. V89's proof of
\eqref{r89:native-replacement} already bounds the mean exterior
trial term by $9b$ when $b\le\kappa$, and by
$(135/16+54\sqrt2)\tau$ otherwise. Their sums are below $90\tau$.
Combine with \eqref{r96:main-budget}. This is a consequence of
the inherited mean trial estimate, not a new rare-region theorem.
It does not replace $U_0$ by $\bar U_0$ for arbitrary multipliers.

\subsection{The complete source bank and a marginal-resolvent comparison}

Keep V93--V94's actual $\mathcal C$, $U_eh=\Omega h$,
off-diagonal source $B_eh$, conditional floor $\widehat\eta>0$,
matrix multiplication operator $M_e$, and Schur form
$\mathfrak s_e$. All are associated with the native conditional
vacuum, not the tempered trial. Dropping $i_e\ge0$ in
\eqref{r96:main-budget} gives the new admissible coefficients
\[
 \int j_e|h|^2d\bar\nu_e
 \le\frac{U_0}{\kappa}\|h\|^2+
                  \frac1\kappa\bar{\mathcal E}_e(h).
 \tag{V96.15}\label{r96:j-form}
\]
For a complete exterior bank $\mathbf h=(h_i)$, V93's exact
conditional Cauchy--Schwarz argument now yields
\[
 \begin{aligned}
 \|\mathcal C^*U_e\mathbf h\|^2
 &\le\frac{4U_0}{\kappa}\sum_i\|h_i\|^2+
                       \frac4\kappa\sum_i\bar{\mathcal E}_e(h_i),\\
 \|B_eh\|^2
 &\le4U_0\bar{\mathcal E}_e(h)+4\kappa\mathcal T_{2,e}(h),
 \qquad
 \mathcal T_{2,e}(h)=\sum_i\bar{\mathcal E}_e(X_i h).
 \end{aligned}
 \tag{V96.16}\label{r96:bank}
\]
These hold for the indicated component form domains, in particular
for smooth $h$. To check the normalization, use
$B_eh=-\kappa\mathcal C^*U_e\nabla_yh$ from \eqref{r93:B} and
$\sum_i\|X_i h\|^2=\bar{\mathcal E}_e(h)/\kappa$.
No conditional Poincare inequality was used to obtain
\eqref{r96:j-form} or \eqref{r96:bank}.
V93's bounds remain available and can be smaller in some parameter
and test-function ranges; the new pair does not dominate them
term by term for every coupling.

\begin{corollary}[Positive return with the new source budget]\label{r96:return-theorem}
For $E=\bar{\mathcal E}_e(h)>0$ and
$T=\mathcal T_{2,e}(h)<\infty$,
\[
 \mathfrak s_e(h)\ge
 \frac{E^2}
 {(1+4U_0/\widehat\eta)E+(4\kappa/\widehat\eta)T}.
 \tag{V96.17}\label{r96:harmonic}
\]
There is also a comparison defined on the first-derivative form
domain without displaying a second-derivative norm:
\[
 \boxed{\;
 \mathfrak s_e(h)\ge\frac{\kappa\widehat\eta}{4}
  \sum_i\left\langle X_i h,
    \bigl(\bar L_e+U_0+\widehat\eta/4\bigr)^{-1}X_i h
                       \right\rangle_{\bar\nu_e}.
 \;}
 \tag{V96.18}\label{r96:resolvent}
\]
Finally,
\[
 \int\Tr M_e\,d\bar\nu_e\le\frac{4\bar U_0}{\widehat\eta}.
 \tag{V96.19}\label{r96:mean-metric}
\]
\end{corollary}

\begin{proof}
For \eqref{r96:harmonic}, apply V94's positive harmonic comparison
\eqref{r94:harmonic-bound} with $c_0=U_0/\kappa$ and
$c_1=1/\kappa$. Equivalently apply Cauchy--Schwarz with weights
$I+M_e$ and $(I+M_e)^{-1}$, bound
$M_e=\kappa\mathcal C\mathsf A_e^{-1}\mathcal C^*$
using $\mathsf A_e\succeq\widehat\eta$, and use the first line of
\eqref{r96:bank}. This is the already proved positive comparison,
not subtraction of a return majorant that must be made smaller
than one.

To prove \eqref{r96:resolvent}, the same native inverse floor and
\eqref{r96:bank} give the bank quadratic-form inequality
\[
 \langle\mathbf h,M_e\mathbf h\rangle
 \le\frac{4U_0}{\widehat\eta}\|\mathbf h\|^2+
                  \frac4{\widehat\eta}
                         \sum_i\bar{\mathcal E}_e(h_i).
 \tag{V96.20}\label{r96:metric-form}
\]
At fixed volume $M_e$ is a bounded nonnegative matrix-valued
multiplication operator. On the component weighted $H^1$ domain,
\[
 I+M_e\preceq
 \frac4{\widehat\eta}
       \bigl(\bar L_e+U_0+\widehat\eta/4\bigr)^{\oplus}.
\]
The right side is the direct sum of the indicated scalar
Friedrichs operators, with strictly positive shift. Inverse order
therefore gives an inequality between their bounded inverses on
the whole bank $L^2$ space. Combine it with
$\mathfrak s_e(h)\ge
\kappa\langle\nabla_yh,(I+M_e)^{-1}\nabla_yh\rangle$ from
\eqref{r94:weighted-form}. Approximation in weighted $H^1$ proves
the stated domain extension. Finally V94's pointwise
$\Tr M_e\le4\kappa j_e/\widehat\eta$, followed by
\eqref{r96:mean-budget}, proves \eqref{r96:mean-metric}.
\end{proof}

The superscript $\oplus$ denotes independent copies in the global
left-invariant frame. It is not a Hodge Laplacian, nor is any
commutation of $\nabla_y$ with $\bar L_e$ asserted. Individual
gradient components of a physical function need not be gauge
invariant, which is why the bank form was proved in the full
exterior space first. The mean metric bound is not a pointwise
bound on its physical compression.

\subsection{What has improved, and what is still unpaid}

The weighted source estimate is now all-coupling and volume uniform,
with explicit zeroth-order energy scale
$U_0=O(\kappa^{1/3}b^{2/3})$ at large $b/\kappa$. Its proof never
passes from a score gradient to a score through an inverse
conditional gap. The native source may be nonlocal in the exterior;
the complete derivative sum is nevertheless bounded as a weighted
form. This is stronger information than its small-coupling leading
coefficient or an unweighted sensitivity estimate alone.

However the remaining resolvent comparison still contains
$\widehat\eta$, and \eqref{r96:resolvent} is not a uniform lower
bound by $\|h-\bar\nu_eh\|^2$. The native exterior drift, the
neutral low modes, and control through growing regions have not
been removed. One cannot commute the gradient through the scalar
resolvent, replace $T$ by a constant times $E$, or turn a mean
metric bound into such coercivity without another proof.
In the physical normalization of V85,
\[
 \kappa^{1/3}b^{2/3}
   =\frac{c^{1/3}}{2^{1/3}g^{2/3}a_s},
 \qquad
 \sqrt{\kappa b}=\frac1{\sqrt2\,a_s}.
\]
These are ultraviolet upper-cost scales, not a finite physical
mass lower bound. No interacting four-dimensional continuum
construction or full Yang--Mills mass gap is proved.

The next upstream obligation is source-resolved control of the
native conditional inverse or a covariant marginal-resolvent
estimate on physical gradients that survives growing volume and
physical weak coupling. Merely optimizing the constant $C_0$
again would not address that obligation. V95's exact gauge
transport bound remains available to separate known gauge
directions from the genuinely physical shape source.

Verification checks the exact constants, the literal oriented
Wilson-star trace on three finite fixtures, and the spherical
moment and integration-by-parts identities. A separate exact
nonproduct ground $\Omega\propto e^{a x_0y_0}$ on two spheres,
with engineered potential $\kappa\Delta\Omega/\Omega$, checks
the marginal equation and 64 weighted competitor identities
against a direct ground-form calculation. That fixture is
explicitly not a Wilson model. Native isolated four-link Wilson
loops at five couplings check the energy decomposition and their
own one-plaquette variational budgets by finite character
truncation; no certified truncation error or full-lattice inference
is claimed. V95's mathematical diagnostics are also replayed.
The inequalities themselves are proved above, not inferred from
these finite tests.

V95's body is preserved byte for byte except its version title.
Only the latest active main source, V96, is retained in
\texttt{latex\_notes}; its predecessor is exactly recoverable.
All diagnostic and build files remain outside that folder.
The next regular companion review, broad literature sweep and
archive/restart boundary are at V100.
% END CONSOLIDATED V96 APPEND

% BEGIN CONSOLIDATED V97 APPEND
\clearpage
\section{V97: A native conditional spectral count and the retained soft return}
\label{r97:section}

\subsection{Separate the inverse loss instead of paying it in every channel}

V96 bounded the complete native shape source without a conditional
inverse gap, but its subsequent resolvent comparison reintroduced
$\widehat\eta^{-1}$. We now separate the conditional spectrum before
making that estimate. V93's pointwise native score bound gives
a volume-uniform comparison of \emph{every} conditional eigenvalue
with the free sphere spectrum. Consequently only polynomially many
conditional modes can lie below a specified energy. The remaining
metric has a gap-independent weighted bound; the low-mode return is
retained exactly as a finite-rank correction at each exterior.

This is an all-coupling statement about the actual conditional
operator, not V95's small-coupling rank-four coefficient and not
the frozen magnetic rotor of V88. It does not assert a finite-rank
full metric or a finite-dimensional global field theory.
The min--max and Woodbury arguments are standard; their native
spectral count and source-budget combination are the new acquisition.
The targeted publisher abstract for F.~Nier,
\emph{Quantitative analysis of metastability in reversible diffusion
processes via a Witten complex approach},
\href{https://proceedings.centre-mersenne.org/articles/10.5802/jedp.8/}
{JEDP (2004), article 8},
provides relevant provenance for separating small eigenvalues through
the first-order factorization. No Morse landscape, tunnelling
asymptotic or quantitative theorem from that work is assumed for
the native YM vacuum.

\subsection{All conditional eigenvalues from the paid native score}

Keep $N\ge3$, $\kappa>0$, $b\ge0$, and the exact factorization
$\Omega=\rho(y)\zeta(q,y)$ of V96. On full one-link Haar space define
\[
 D_yv=\nabla_qv-v\nabla_q\log\zeta,\qquad
 \widetilde A_y=\kappa D_y^*D_y,\qquad
 S=\frac{2b}{\kappa}.
 \tag{V97.1}\label{r97:A}
\]
The quadratic form domain is $H^1(S^3)$. V93 proved
$|\nabla_q\log\zeta|\le S$ at every exterior.
The smooth positive normalized $\zeta$ spans the kernel of
$\widetilde A_y$: if $D_yv=0$, then $\nabla_q(v/\zeta)=0$.
Its restriction to $\zeta^\perp$ is precisely the native $A_e(y)$
in \eqref{r94:A}. There is no subtraction of a frozen-rotor energy
or change of conditional law in this definition.

List the full eigenvalues, with multiplicity, as
$0=\lambda_0(y)<\lambda_1(y)\le\cdots$, and let
$0=\mu_0<\mu_1\le\cdots$ be those of the unit-round
$L=-\Delta_{S^3}$. The latter are $\ell(\ell+2)$ with multiplicity
$(\ell+1)^2$, $\ell=0,1,\ldots$.

\begin{theorem}[Uniform native square-root spectral comparison]
\label{r97:count-theorem}
For every $k\ge0$ and every exterior,
\[
 \boxed{\quad
 \kappa(\sqrt{\mu_k}-S)_+^2
       \le\lambda_k(y)\le
                  \kappa(\sqrt{\mu_k}+S)^2 .
 \quad}
 \tag{V97.2}\label{r97:eigenvalues}
\]
For $\gamma>0$, let $\Pi_<(y)$ be the spectral projection of
$A_e(y)$ onto $(0,\gamma)$. Put
\[
 \begin{aligned}
 T_\gamma&=S+\sqrt{\gamma/\kappa},\\
 L_\gamma&=\left\lfloor\sqrt{1+T_\gamma^2}-1\right\rfloor,\\
 d_\gamma&=
       \frac{(L_\gamma+1)(L_\gamma+2)(2L_\gamma+3)}6-1.
 \end{aligned}
 \tag{V97.3}\label{r97:count}
\]
Then
\[
 \operatorname{rank}\Pi_<(y)\le d_\gamma
       =O\bigl((1+S+\sqrt{\gamma/\kappa})^3\bigr)
 \tag{V97.4}\label{r97:rank}
\]
uniformly in volume. Inclusion of a free eigenvalue at the endpoint
in \eqref{r97:count} only makes this an upper bound.
\end{theorem}

\begin{proof}
The multiplication map $v\mapsto v\nabla_q\log\zeta$ has norm at
most $S$ from scalar $L^2$ to tangent-vector $L^2$. Therefore
\[
 \|\nabla_qv\|-S\|v\|\le\|D_yv\|
                         \le\|\nabla_qv\|+S\|v\|.
 \tag{V97.5}\label{r97:triangle}
\]
For the upper min--max bound, take the span of the first $k+1$
free eigenfunctions. Every vector there satisfies
$\|\nabla_qv\|\le\sqrt{\mu_k}\|v\|$, so the upper inequality
gives the upper bound in \eqref{r97:eigenvalues}.
For the lower bound, take the orthogonal complement of the first
$k$ free eigenfunctions. If $\sqrt{\mu_k}\ge S$, every vector
there satisfies
$\|D_yv\|\ge(\sqrt{\mu_k}-S)\|v\|$. The max--min form of the
same principle gives the lower bound. Otherwise use
$\widetilde A_y\succeq0$. This argument compares singular values
of the first-order operators; it does not square an invalid
operator inequality between noncommuting terms.

If $\lambda_k(y)<\gamma$, the lower bound implies
$\mu_k<T_\gamma^2$. Count the free eigenvalues up to that
threshold and remove the one native zero mode. The sphere
multiplicity follows, for example, by subtracting the dimensions
of homogeneous polynomial spaces of degrees $\ell$ and $\ell-2$
in four variables:
$\binom{\ell+3}{3}-\binom{\ell+1}{3}=(\ell+1)^2$,
with the absent low-degree space taken as zero.
Summing the squares gives \eqref{r97:count}.
Smoothness and compactness ensure compact resolvent throughout,
so both applications of min--max and the counting argument
apply to the closed native form.
\end{proof}

In particular $\lambda_1(y)\ge
\kappa(\sqrt3-2b/\kappa)_+^2$. This is useful only in its stated
small-coupling range; it does not supply a positive bound at
arbitrarily large $b/\kappa$. At the illustrative cutoff
$\gamma=\kappa$, \eqref{r97:count} gives respectively
$d_\gamma=0,13,3310$ at $b/\kappa=0,1,10$.
The polynomial count can be large. Its content is independence
of the total number of exterior links, not small numerical rank
at physical weak coupling.

\subsection{A controlled regular metric and an exact low-mode correction}

Use the exterior left-invariant frame with $m=3(|E_N|-1)$
components. At each $y$ write $z_i=X_i\zeta$ and
\[
 C_yv=(2\langle z_i,v\rangle_{dq})_{i=1}^m,\qquad
 M_y=\kappa C_yA_e(y)^{-1}C_y^* .
 \tag{V97.6}\label{r97:C}
\]
This is V94's metric in the native marginal representation.
Let $\Pi_\ge=I-\Pi_<$ on $\zeta^\perp$, and define
\[
 \begin{aligned}
 M_{\ge,\gamma}
    &=\kappa C_yA_e(y)^{-1}\Pi_\ge C_y^*,\\
 W_\gamma&=I+M_{\ge,\gamma},\\
 \mathcal H_<(y)&=\operatorname{Ran}\Pi_<(y),\qquad
 A_< = A_e(y)|_{\mathcal H_<(y)},\\
 Z_\gamma&=\sqrt\kappa\,C_y|_{\mathcal H_<(y)} .
 \end{aligned}
 \tag{V97.7}\label{r97:split}
\]
All quantities in this subsection depend on $y$ unless indicated.
The cutoff is constant in the exterior. Its spectral projection
may jump when an eigenvalue crosses $\gamma$; we use measurable
functional calculus and never differentiate that projection.
No smooth choice of eigenvectors or constant-rank bundle is needed.

\begin{theorem}[Regular source bound and retained conditional return]
\label{r97:return-theorem}
With $U_0,\bar U_0$ from V96, every smooth exterior bank
$\mathbf u=(u_i)$ satisfies
\[
 \begin{aligned}
 \int\mathbf u^*M_{\ge,\gamma}\mathbf u\,d\bar\nu_e
  &\le\frac{4U_0}{\gamma}\|\mathbf u\|^2+
                         \frac4\gamma\sum_i\bar{\mathcal E}_e(u_i),\\
 \int\Tr M_{\ge,\gamma}\,d\bar\nu_e
  &\le\frac{4\bar U_0}{\gamma}.
 \end{aligned}
 \tag{V97.8}\label{r97:regular}
\]
The exact inverse weight is
\[
 \begin{aligned}
 (I+M_y)^{-1}&=W_\gamma^{-1}-R_{<,\gamma},\\
 R_{<,\gamma}
  &=W_\gamma^{-1}Z_\gamma
       (A_<+Z_\gamma^*W_\gamma^{-1}Z_\gamma)^{-1}
                        Z_\gamma^*W_\gamma^{-1},\\
 0\preceq R_{<,\gamma}&\preceq W_\gamma^{-1}\preceq I,\qquad
 \operatorname{rank}R_{<,\gamma}\le d_\gamma .
 \end{aligned}
 \tag{V97.9}\label{r97:woodbury}
\]
When $\mathcal H_<=\{0\}$ the correction is zero. Consequently,
for $g=\nabla_yh$ and $u=Z_\gamma^*W_\gamma^{-1}g$, the native
Schur comparison is
\[
 \begin{aligned}
 \mathfrak s_e(h)\ge\kappa\int
 \bigl[\,g^*W_\gamma^{-1}g
     -u^*(A_<+Z_\gamma^*W_\gamma^{-1}Z_\gamma)^{-1}u\,\bigr]
                                                      d\bar\nu_e .
 \end{aligned}
 \tag{V97.10}\label{r97:retained}
\]
The bracket is nonnegative pointwise. No estimate using
$\widehat\eta^{-1}$ occurs in \eqref{r97:regular}.
\end{theorem}

\begin{proof}
The regular spectral inverse is at most $\gamma^{-1}I$.
Moreover
$\|C_y^*\mathbf u(y)\|^2
\le4j_e(y)|\mathbf u(y)|^2$.
Apply V96's complete bank bound \eqref{r96:bank}, and its mean
bound \eqref{r96:mean-budget} for the trace. This proves
\eqref{r97:regular}, including its extension to the component
weighted $H^1$ domain.

The spectral split gives exactly
$I+M_y=W_\gamma+Z_\gamma A_<^{-1}Z_\gamma^*$.
At each finite regulator every positive conditional eigenvalue
is strictly positive, so the finite-dimensional $A_<$ is
invertible on its range. This uses existence, not a quantitative
lower-gap estimate. Set $Y=W_\gamma^{-1/2}Z_\gamma$.
Multiplication verifies
\[
 (I+Y A_<^{-1}Y^*)^{-1}
        =I-Y(A_<+Y^*Y)^{-1}Y^* .
\]
Conjugation by $W_\gamma^{-1/2}$ proves the first two lines of
\eqref{r97:woodbury}. Positivity of both the difference and the
correction follows from this inverse identity and inverse order.
The correction factors through $\mathcal H_<$, whose dimension
is bounded by \eqref{r97:rank}. This proves the rank claim and
the remaining order bounds. No new abstract Woodbury identity
is being claimed.

The spectral fields and their compositions are measurable;
the inverse weight and correction are bounded by $I$.
Thus the displayed nonnegative integrand defines the asserted
form on native weighted $H^1$. Substitute the identity into
V94's comparison \eqref{r94:weighted-form} to obtain
\eqref{r97:retained}. Gauge covariance of the native fibre,
its spectral projections and the exterior frame ensures that
the comparison restricts to physical exterior amplitudes.
\end{proof}

There is also the regular bank inverse comparison
\[
 W_\gamma^{-1}\succeq
   \frac{\gamma}{4}
       \bigl(\bar L_e+U_0+\gamma/4\bigr)^{-1,\oplus}.
 \tag{V97.11}\label{r97:regular-inverse}
\]
Indeed \eqref{r97:regular} gives
$W_\gamma\preceq
(4/\gamma)(\bar L_e+U_0+\gamma/4)^\oplus$
as forms on the component weighted $H^1$ domain. Inverse order
is legitimate because the right-hand Friedrichs operator has
a strictly positive shift. This is V96's argument applied only
after removing the low conditional spectrum, so $\gamma$ is a
freely chosen energy and not an exponentially small bound.
It does not discard the correction in \eqref{r97:retained}.

For example choose $\gamma=\kappa+U_0$. Then
$4U_0/\gamma<4$ and $4\bar U_0/\gamma<4$, uniformly in coupling
and volume. The regular metric has these bounded zeroth-order
and mean-trace constants, while the correction has conditional
rank $O((1+b/\kappa)^3)$ by \eqref{r97:count}.
Also $\Tr R_{<,\gamma}\le d_\gamma$ pointwise.
These statements do not require small coupling or that the
actual conditional wavefunction depend only on the four staples.

\subsection{Why the retained directions cannot be removed by counting}

In a single retained direction with $W_\gamma=1$, $A_<=a>0$
and $Z_\gamma=z$, the inverse weight is
\[
 1-\frac{|z|^2}{a+|z|^2}=\frac{a}{a+|z|^2}.
 \tag{V97.12}\label{r97:one-channel}
\]
It can approach zero even though the correction has rank one.
Thus a small rank does not prove coercivity on physical gradients.
As $y$ ranges over the exterior, even a rank-one fibre field
generally defines an infinite-rank global multiplication
operator. Nor does a small fraction of bad bank directions force
$\nabla_yh$ to have small overlap with them.
V95's exact gauge-transport estimate still controls known gauge
directions, but does not remove the physical shape source.

It is also false that bounded score alone forces an inverse gap
with only polynomial loss. Here is a self-contained check of
that logical limitation, not a YM counterexample. On $S^3$ let
$x=q_0$, $a>0$, and
$\zeta_a=Z_a^{-1/2}e^{a x^2}$, where
$Z_a=\int_{S^3}e^{2ax^2}dq$. Its score has supremum $a$,
and the positive-ground operator is
\[
 \widetilde A^{(a)}
 =\kappa\left[-\Delta+
          a(2-8x^2)+4a^2x^2(1-x^2)\right].
 \tag{V97.13}\label{r97:two-peak}
\]
This follows from
$\Delta x^2=2-8x^2$ and
$|\nabla(a x^2)|^2=4a^2x^2(1-x^2)$.
Its ground-transform form has density proportional to
$e^{2a x^2}dq$. The odd Lipschitz function
$f(x)=\max\{-1,\min\{2x,1\}\}$ has mean zero in this density.
Let $c_{\rm cap}=\int_{|x|\ge3/4}dq>0$.
After cancelling normalization in its Rayleigh quotient,
\[
 \lambda_1(\widetilde A^{(a)})
 \le
 \frac{\kappa\int(1-x^2)|f'(x)|^2e^{2a x^2}dq}
      {\int |f(x)|^2e^{2a x^2}dq}
 \le \frac{4\kappa}{c_{\rm cap}}e^{-5a/8}.
 \tag{V97.14}\label{r97:soft-example}
\]
The numerator is at most $4\kappa e^{a/2}$ since $f'$ is
supported on $|x|<1/2$; the denominator is at least
$c_{\rm cap}e^{9a/8}$. Smooth approximation gives the same
variational bound if desired. This operator has a polynomial
score bound but an exponentially soft mode. Its potential is
not the original Wilson potential, so the example neither
disproves a YM gap nor verifies any such soft mode in the
native cubic conditional law. It explains why the latter law,
not just its score norm, must enter the next step.

Conversely the native isolated $m$-link Wilson loop provides a
useful calibration with no mode in $(0,\kappa)$. If its reduced
ground solves
$(m\kappa L+b(1-x))\psi=E_{\rm loop}\psi$, then its actual
one-link conditional operator is exactly
\[
 \widetilde A_{\rm loop}
   =\kappa L+\frac{b}{m}(1-x)-\frac{E_{\rm loop}}m .
 \tag{V97.15}\label{r97:loop}
\]
It is the magnetic rotor with coupling $b/m$, centered at its
ground energy. V88's already proved rotor gap therefore gives
its conditional gap at least $\kappa$. The equality uses the
special one-cycle reduction; it is not valid for the coupled
cubic vacuum. Removing one link from that cycle also leaves no
nonconstant physical exterior amplitude, as V95 explained.

\subsection{Verification and the next genuinely unpaid object}

Exact diagnostics verify the spherical multiplicity and count,
and a noncommuting rational-matrix instance of the retained
Woodbury identity, including its rank and sign. Five native
isolated-loop spectra check the conditional normalization,
the first 99 eigenvalue comparisons, and finite low-mode counts
using radial Gegenbauer matrices and angular multiplicities.
They are consistency tests without certified truncation error,
not proofs for the full cubic lattice. The two-peak diagnostic
checks the odd Rayleigh test and score bound; its analytic proof
is \eqref{r97:soft-example}. V96's mathematical checks are replayed.

The new unconditional result isolates the entire inverse-gap
loss into a native low-mode return of bounded conditional rank;
the regular metric is controlled without that loss. What remains
is quantitative control of
$A_<+Z_\gamma^*W_\gamma^{-1}Z_\gamma$ together with the actual
physical source overlap in \eqref{r97:retained}, or a proof that
the relevant low conditional modes cannot occur in the native
Wilson law. Rank alone, a frozen-rotor gap, and a Gibbs trial
substitution supply none of these. The number $d_\gamma$ still
grows at weak coupling and is not a uniform continuum rank.
Even control of this conditional return would leave the growing
neutral-region and continuum-construction obligations.
The full interacting four-dimensional YM mass gap remains unproved.

V96's body is preserved byte for byte except the version title,
and its exact predecessor is recoverable. Only active V97 is
retained in \texttt{latex\_notes}; all build and diagnostic files
are outside it. The regular companion checkpoint, broad
literature sweep and requested archive/restart remain due at V100.
% END CONSOLIDATED V97 APPEND

% BEGIN CONSOLIDATED V98 APPEND
\clearpage
\section{V98: Smooth native harmonic compression and a coercive full high block}
\label{r98:section}

\subsection{Replace a crossing spectral cutoff by a vacuum-compatible subspace}

V97 isolated soft conditional modes, but its sharp spectral
projection need not vary smoothly with the exterior. To retain
the exterior kinetic energy without differentiating such a
projection, we now construct a different cutoff. Start from a fixed
free harmonic space and replace one of its directions by the exact
native conditional vacuum. The resulting subspace is smooth,
has fixed rank, and contains the full vacuum exactly. Its
orthogonal complement has a freely prescribed lower energy bound
for the \emph{full} native Hamiltonian, not just a frozen rotor.
Moreover its moving-isometry cost decays with the harmonic cutoff,
with volume-independent constants and without an inverse
conditional gap.

Elementary geometry of close projections is classical; see the
publisher abstract of J.~Avron, R.~Seiler and B.~Simon,
\emph{The Index of a Pair of Projections},
\href{https://www.sciencedirect.com/science/article/pii/S0022123684710317}
{J. Functional Analysis 120 (1994), 220--237}.
We do not import an index theorem or infer any YM spectral estimate
from that abstract. The explicit projection and all estimates below
are proved directly. The new point relative to the previous notes
is their compatibility with the \emph{native} vacuum and the paid
native mixed-score form. Finite fibre rank will not be confused
with finite global dimension.

\subsection{Construction and the actual complementary Hamiltonian}

Keep the finite cubic regulator $N\ge3$, $\kappa>0$, $b\ge0$,
and the Haar-space decomposition
$\Omega(q,y)=\rho(y)\zeta(q,y)$ from V96.
The positive real $\zeta$ has unit one-link Haar norm.
Let $S=2b/\kappa$, and recall V93's pointwise score bound
$|\nabla_q\log\zeta|\le S$. All fibre inner products below use
Haar measure $dq$; complex amplitudes are allowed even though
$\zeta$ and its derivatives are real.

Fix an energy $\gamma>0$. Choose an integer $L\ge0$ such that
\[
 \mu=(L+1)(L+3)\ge 4S^2+\frac{2\gamma}{\kappa}.
 \tag{V98.1}\label{r98:threshold}
\]
Let $F=F_L$ be the free $S^3$ harmonic projection onto degrees
$0,\ldots,L$ and $E_L=\operatorname{Ran}F$. Then
\[
 D_L=\dim E_L=\frac{(L+1)(L+2)(2L+3)}6,\qquad
 \lambda_L=L(L+2).
 \tag{V98.2}\label{r98:rank}
\]
These are the full sphere harmonics, not only class functions.
The first omitted eigenvalue is $\mu$.
Define, at each exterior,
\[
 a=\|F\zeta\|,\qquad p=\frac{F\zeta}{a},\qquad
 \Pi=P_\zeta+F-P_p,\qquad
 Jv=v+\langle p,v\rangle(\zeta-p)\quad(v\in E_L).
 \tag{V98.3}\label{r98:construction}
\]
Here $P_u=|u\rangle\langle u|$ for a unit vector $u$.

\begin{theorem}[Smooth native cutoff and coercive high block]
\label{r98:high-theorem}
The definitions in \eqref{r98:construction} are globally smooth
in $y$, and
\[
 \begin{gathered}
 a^2\ge1-\frac{S^2}{\mu}\ge\frac34,\qquad
 J^*J=I_{E_L},\qquad JJ^*=\Pi,\qquad Jp=\zeta,\\
 \Pi^2=\Pi=\Pi^*,\qquad \operatorname{rank}\Pi=D_L,\qquad
 \|\Pi-F\|=\sqrt{1-a^2}\le\frac{S}{\sqrt\mu}.
 \end{gathered}
 \tag{V98.4}\label{r98:projection}
\]
Let $\mathcal P=\int^\oplus\Pi(y)\,dy$ and
$\mathcal Q=I-\mathcal P$ on full Haar configuration space.
For $K=H-E_0$, its closed form restricted to
$\operatorname{Ran}\mathcal Q$ satisfies
\[
 \boxed{\quad K_{\mathcal Q}\succeq
       \Gamma_L I_{\operatorname{Ran}\mathcal Q},\qquad
 \Gamma_L=\kappa\bigl(\sqrt{\mu-S^2}-S\bigr)^2\ge\gamma .\quad}
 \tag{V98.5}\label{r98:high-floor}
\]
In particular $\|K_{\mathcal Q}^{-1}\|\le\Gamma_L^{-1}$.
The projection is gauge equivariant, so the same bound holds
on its gauge-invariant restricted space. It retains the exact
full vacuum: $\mathcal P\Omega=\Omega$.
\end{theorem}

\begin{proof}
Write $r_0=(I-F)\zeta$. The free spectral inequality and the
score bound give
\[
 \mu\|r_0\|^2\le\|\nabla_q\zeta\|^2\le S^2.
\]
Thus $a$ is uniformly separated from zero. Smoothness follows
from that of the finite-regulator positive elliptic ground
state and finite-dimensional harmonic projection.
The subspace $E_L\cap p^\perp$ is orthogonal to $\zeta$ since
$F\zeta=ap$. The map $J$ is the identity there and sends the
unit vector $p$ to the unit vector $\zeta$. This proves its
isometry, the orthogonal-sum formula for $\Pi$, and all rank
and vacuum statements. The difference $\Pi-F=P_\zeta-P_p$
acts on $\operatorname{span}\{p,\zeta\}$; its nonzero eigenvalues,
when this span has dimension two, are
$\pm\sqrt{1-|\langle p,\zeta\rangle|^2}
=\pm\sqrt{1-a^2}$. If $a=1$ the difference vanishes.
There is no division by $\|r_0\|$, so this case is smooth too.

For a fibre vector $v\perp\operatorname{Ran}\Pi$, write
$v=cp+w$, $w=(I-F)v$. Orthogonality to $\zeta$ implies
$ac+\langle r_0,w\rangle=0$. Therefore
\[
 |c|^2\le\frac{1-a^2}{a^2}\|w\|^2,\qquad
 \|w\|^2\ge a^2\|v\|^2,\qquad
 \|\nabla_qv\|^2\ge\mu a^2\|v\|^2
                       \ge(\mu-S^2)\|v\|^2.
 \tag{V98.6}\label{r98:high-kinetic}
\]
V97's first-order triangle inequality now gives
$\kappa\|D_qv\|^2\ge\Gamma_L\|v\|^2$.
Its right-hand square is legitimate because
$\sqrt{\mu-S^2}>S$ under \eqref{r98:threshold}.
If $t=\sqrt{\gamma/\kappa}$, then
\[
 3S^2+2t^2-(S+t)^2=(S-t)^2+S^2\ge0;
\]
hence $\sqrt{\mu-S^2}\ge S+t$ and $\Gamma_L\ge\gamma$.

On the full configuration space the exact ground form is
$q_K[u]=\kappa\sum_{\text{all }i}\|D_i u\|^2$.
For $u=\mathcal Q u$, apply the fibre inequality to its
$q$ derivatives and keep the remaining nonnegative terms.
This proves \eqref{r98:high-floor} for smooth vectors.
At fixed finite regulator the fibre kernel of $\Pi$ is
smooth in $(q,q',y)$ and has finite rank in the link variables.
The direct-integral projection therefore maps $H^1$ to $H^1$;
it is not a finite-rank operator on the full space.
Its complementary form domain is
$H^1\cap\operatorname{Ran}\mathcal Q$. The restriction is
closed, and is dense in $\operatorname{Ran}\mathcal Q$
because projections of smooth approximants are smooth.
The bound extends by closure and defines the stated
self-adjoint operator and inverse. No volume-uniform norm
of this auxiliary $H^1$ projection is assumed.

Finally the free Casimir projection commutes with both
endpoint group actions on the link. The native conditional
vacuum transforms covariantly under the simultaneous action
on $(q,y)$, as do $a,p,\Pi,J$. Thus $\mathcal P$ commutes
with gauge transformations and the restricted bound is valid.
This uses no gauge fixing or physical tensor factorization.
\end{proof}

For real $z<\Gamma_L$ the same argument yields
$\|(K_{\mathcal Q}-z)^{-1}\|\le(\Gamma_L-z)^{-1}$.
The required rank is
$O((1+S+\sqrt{\gamma/\kappa})^3)$, independent of volume
but growing at weak coupling. Although each fibre is finite
dimensional, $\mathcal P$ still has an exterior function space.

\subsection{Exact moving geometry with no fibre-rank factor}

Let $i$ range over all exterior left-invariant derivative
components, and write
\[
 z_i=X_i\zeta,\quad p_i=X_i p,\quad a_i=X_i a,\quad
 j_e=\sum_i\|z_i\|^2,\quad
 \mathcal A_i=J^*(X_iJ),\quad
 \mathcal B_i=(I-\Pi)(X_iJ).
 \tag{V98.7}\label{r98:geometry-def}
\]
All are fibre fields. The same fixed space $E_L$ is used at
every exterior; no eigenvector-frame choice is involved.
The connection is skew adjoint and is explicitly
\[
 \begin{aligned}
 X_iJ&=|z_i-p_i\rangle\langle p|
                     +|\zeta-p\rangle\langle p_i|,\\
 \mathcal A_i&=(1-a)
       \bigl(|p\rangle\langle p_i|-|p_i\rangle\langle p|\bigr),\\
 \mathcal B_i&=|z_i-a p_i\rangle\langle p|
                          -|p-a\zeta\rangle\langle p_i|.
 \end{aligned}
 \tag{V98.8}\label{r98:connection}
\]

\begin{lemma}[Exact derivative accounting]
\label{r98:geometry-lemma}
With Hilbert--Schmidt norms between the indicated fibre spaces,
\[
 \begin{aligned}
 \|X_iJ\|_{\rm HS}^2
     &=\|z_i\|^2+(3-4a)\|p_i\|^2,\\
 \|\mathcal B_i\|_{\rm HS}^2
     &=\|z_i\|^2+(1-2a^2)\|p_i\|^2,\\
 \|X_i\Pi\|_{\rm HS}^2&=2\|\mathcal B_i\|_{\rm HS}^2.
 \end{aligned}
 \tag{V98.9}\label{r98:HS}
\]
Consequently, with $\mathcal G_L=\sum_i\|X_iJ\|_{\rm HS}^2$,
\[
 \mathcal G_L\le j_e,\qquad
 0\preceq V_{\rm BH}^{(L)}
       :=\kappa\sum_i\mathcal B_i^*\mathcal B_i
       \preceq\kappa\mathcal G_L I_{E_L}.
 \tag{V98.10}\label{r98:BH}
\]
There is no factor $D_L$ or number of exterior links here.
\end{lemma}

\begin{proof}
Differentiate the real unit normalizations to obtain
$\langle\zeta,z_i\rangle=\langle p,p_i\rangle=0$, and
$Fz_i=a_i p+a p_i$. Substitution into the derivative of $J$
gives \eqref{r98:connection}; in particular
$\Pi z_i=a p_i$ and $(I-\Pi)p=p-a\zeta$.
The two rank-one summands of $X_iJ$ have orthogonal domain
directions $p,p_i$. Also
$\langle z_i,p_i\rangle=a\|p_i\|^2$ and
$\|\zeta-p\|^2=2(1-a)$. Expanding their squared norms gives
the first identity in \eqref{r98:HS}.
Since $X_iJ=J\mathcal A_i+\mathcal B_i$ is an orthogonal
range decomposition and
$\|\mathcal A_i\|_{\rm HS}^2=2(1-a)^2\|p_i\|^2$,
subtraction gives the second identity.
Differentiating $\Pi=JJ^*$ and using
$\mathcal A_i^*=-\mathcal A_i$ yields
$X_i\Pi=\mathcal B_iJ^*+J\mathcal B_i^*$.
These are its two orthogonal off-diagonal blocks, proving
the last identity. Since $a\ge\sqrt3/2>3/4$, the correction
in the first line of \eqref{r98:HS} is nonpositive.
Sum it, and use operator norm at most Hilbert--Schmidt norm
in each nonnegative $\mathcal B_i^*\mathcal B_i$.
\end{proof}

The vacuum's internal motion is not lost:
\[
 p_i+\mathcal A_i p=a p_i,\qquad
 \mathcal B_i p=z_i-a p_i,\qquad
 a^2\|p_i\|^2+\|\mathcal B_i p\|^2=\|z_i\|^2.
 \tag{V98.11}\label{r98:vacuum-motion}
\]
Thus the coefficient of the retained vacuum will be $p(y)$,
not a constant vector.

\subsection{A cutoff-decaying bound from the paid mixed-score form}

Write $d\bar\nu_e=\rho^2dy$ and
$\bar{\mathcal E}_e(h)=\kappa\int|\nabla_yh|^2d\bar\nu_e$.
For vector amplitudes the energy means the sum of component
energies. Retain V96's constant
\[
 U_0=\min\{4b,C_0\kappa^{1/3}b^{2/3}\},\qquad
 C_0=\frac94(52/3)^{1/3}.
 \tag{V98.12}\label{r98:U}
\]
Let $s_{e,a}=X_{e,a}\log\Omega$ and put
\[
 Q_c(y)=
 \mathbb E_{\nu_e^y}
       \sum_{i\ {\rm exterior},\,a}|X_i s_{e,a}|^2.
 \tag{V98.13}\label{r98:Qc}
\]
The two already proved inputs, used before any conditional
Poincare inequality, are
\[
 \begin{aligned}
 \kappa\int |h|^2Q_c\,d\bar\nu_e
    &\le12b\|h\|^2+4S^2\bar{\mathcal E}_e(h),\\
 \kappa\int |h|^2j_e\,d\bar\nu_e
    &\le U_0\|h\|^2+\bar{\mathcal E}_e(h).
 \end{aligned}
 \tag{V98.14}\label{r98:paid-inputs}
\]
The first is \eqref{r93:Q-form} with its parameter $1/2$;
the second is V96's complete weighted native budget after
discarding the nonnegative $i_e$ term.

\begin{theorem}[Vanishing moving-cutoff cost]
\label{r98:tail-theorem}
For every smooth scalar or $E_L$-valued exterior amplitude $h$,
\[
 \boxed{\quad
 \kappa\int |h|^2\mathcal G_L\,d\bar\nu_e
       \le\delta_L\|h\|^2+\epsilon_L\bar{\mathcal E}_e(h),
 \quad}
 \tag{V98.15}\label{r98:tail-budget}
\]
where
\[
 \delta_L=\frac{32b+(16/3)S^2U_0}{\mu},\qquad
 \epsilon_L=\frac{16S^2}{\mu}.
 \tag{V98.16}\label{r98:tail-constants}
\]
The same right side bounds
$\int\langle h,V_{\rm BH}^{(L)}h\rangle d\bar\nu_e$
and $\kappa\sum_i\|\mathcal A_i h\|_{L^2(\bar\nu_e)}^2$.
Both constants tend to zero with $L$ at fixed $\kappa,b$,
uniformly in volume and without $\widehat\eta^{-1}$.
\end{theorem}

\begin{proof}
Let $T_i=\|(I-F)z_i\|^2$. Differentiating
$a^2+\|(I-F)\zeta\|^2=1$ gives
$aa_i=-\langle(I-F)\zeta,(I-F)z_i\rangle$, so
$a_i^2\le(1-a^2)T_i/a^2$.
The orthogonal decomposition
$Fz_i=a_i p+a p_i$ and \eqref{r98:HS} give the useful
nonnegative expansion
\[
 \|X_iJ\|_{\rm HS}^2
   =a_i^2+T_i+(1-a)(3-a)\|p_i\|^2.
 \tag{V98.17}\label{r98:tail-expansion}
\]
Use $a^{-2}\le4/3$,
$(1-a)(3-a)\le2(1-a^2)\le2S^2/\mu$,
$\|p_i\|^2\le\|z_i\|^2/a^2$, and
$T_i\le\mu^{-1}\|\nabla_qz_i\|^2$. Summation yields
\[
 \mathcal G_L\le
 \frac4{3\mu}
      \left(\sum_i\|\nabla_qz_i\|^2+2S^2j_e\right).
 \tag{V98.18}\label{r98:tail-sobolev}
\]
The normalization of the conditional state gives
$z_i=\zeta\,\xi_i$ with
$\xi_i=s_i-\mathbb E_{\nu_e^y}s_i$.
The conditional mean is independent of $q$, and different
link derivatives commute. Consequently
\[
 \nabla_qz_i=\zeta\bigl(X_i s_e+s_e\xi_i\bigr),\qquad
 \sum_i\|\nabla_qz_i\|^2\le2Q_c+2S^2j_e.
\]
Substitute this into \eqref{r98:tail-sobolev}:
\[
 \mathcal G_L\le\frac8{3\mu}Q_c+
                               \frac{16S^2}{3\mu}j_e.
 \tag{V98.19}\label{r98:tail-pointwise}
\]
Multiplication by $\kappa|h|^2$, integration, and the two
paid estimates \eqref{r98:paid-inputs} give exactly
\eqref{r98:tail-constants}. For vector $h$, apply the scalar
estimate to each component and sum; this introduces no
dimension factor. The Born--Huang assertion follows from
\eqref{r98:BH}. Since
$\|\mathcal A_i\|_{\rm HS}\le\|X_iJ\|_{\rm HS}$,
the connection assertion follows as well.
The estimates extend to the component weighted $H^1$ domain
by closure (or smooth approximation and Fatou's lemma).
\end{proof}

For example, increasing the cutoff also to satisfy
$\mu\ge64S^2$ makes $\epsilon_L\le1/4$.
This is a controllable geometric error, not an estimate on
the centered retained matrix potential below.
At $b=0$, $\zeta$ is constant, $\Pi=F$, $J$ is the fixed
inclusion, and all these geometric costs vanish.

\subsection{Exact retained matrix form, including its nonconstant vacuum}

Use the original splitting
$V(q,y)=V_{\rm out}(y)+W_e(q,y)$,
$W_e=4b-f(y)\cdot q\in[0,8b]$. Define the finite matrix
\[
 H_{e,L}(y)=J_y^*(\kappa L_q+W_e(q,y))J_y,\qquad
 F_e=\kappa(i_e+j_e)+w_e
 \tag{V98.20}\label{r98:matrix-def}
\]
with $i_e=\|\nabla_q\zeta\|^2$ and
$w_e=\langle\zeta,W_e\zeta\rangle$ as in V96.
In particular
\[
 0\preceq H_{e,L}
       \preceq[\kappa(\lambda_L+S^2)+8b]I.
 \tag{V98.21}\label{r98:matrix-upper}
\]
Indeed $Jv=v_\perp+c\zeta$ with
$v_\perp\in E_L\cap p^\perp$ and
$\|v\|^2=\|v_\perp\|^2+|c|^2$.
Then
$\|\nabla_qJv\|
\le\sqrt{\lambda_L}\|v_\perp\|+S|c|$,
and Cauchy--Schwarz proves the kinetic upper bound.
Potential positivity and its supremum give the rest.

\begin{proposition}[Exact compression]
\label{r98:compression-proposition}
For smooth Haar exterior coefficients $\alpha(y)\in E_L$,
\[
 \begin{aligned}
 q_H[J\alpha]
   ={}&\kappa\sum_i\int
                   |X_i\alpha+\mathcal A_i\alpha|^2dy\\
     &+\int\langle\alpha,
        (V_{\rm out}I+H_{e,L}+V_{\rm BH}^{(L)})\alpha\rangle dy .
 \end{aligned}
 \tag{V98.22}\label{r98:H-compression}
\]
In native marginal coefficients $\alpha=\rho h$ this becomes
\[
 \begin{aligned}
 q_K[J(\rho h)]
   ={}&\kappa\sum_i\int
                    |X_i h+\mathcal A_i h|^2d\bar\nu_e\\
     &+\int\langle h,
        (H_{e,L}+V_{\rm BH}^{(L)}-F_eI)h\rangle d\bar\nu_e .
 \end{aligned}
 \tag{V98.23}\label{r98:K-compression}
\]
The exact zero-energy coefficient is $h=p(y)$.
\end{proposition}

\begin{proof}
The identity
$X_i(J\alpha)=J(X_i\alpha+\mathcal A_i\alpha)
                  +\mathcal B_i\alpha$
is an orthogonal range decomposition. Taking its squared
norm and adding the $q$ kinetic and potential terms proves
\eqref{r98:H-compression}.
For $\alpha=\rho h$, skew adjointness implies
$\operatorname{Re}\langle h,\mathcal A_i h\rangle=0$.
Thus all terms involving $X_i\rho$ reduce, by Haar integration
by parts, to $-\kappa\rho\Delta_y\rho\,|h|^2$.
Use the exact V96 marginal equation
\[
 (-\kappa\Delta_y+V_{\rm out}+F_e)\rho=E_0\rho
\]
and subtract $E_0\|J(\rho h)\|^2$ to obtain
\eqref{r98:K-compression}. Finally
$J(\rho p)=\rho\zeta=\Omega$, proving the zero-energy claim.
The displayed forms extend from the smooth core at each
finite regulator.
\end{proof}

It is essential to keep the centered matrix
$H_{e,L}+V_{\rm BH}^{(L)}-F_eI$ and the connection.
Neither positivity of $H_{e,L}$ nor smallness of
\eqref{r98:tail-budget} permits discarding $-F_eI$.
Likewise replacing $p(y)$ by a constant coefficient would
change the exact vacuum. The formula is a compression of
the full form, not yet a Schur complement.

\subsection{What has been paid, and what remains in the low return}

For a smooth retained Haar coefficient $\alpha$ the actual
zero-energy Schur form is
\[
 \mathfrak s_L(\alpha)=
    \inf_{v\in H^1\cap\operatorname{Ran}\mathcal Q}
                          q_K[J\alpha+v].
 \tag{V98.24}\label{r98:Schur}
\]
The infimum is attained uniquely. Indeed the restricted
high form is a Hilbert norm by \eqref{r98:high-floor}, and
the linear functional $v\mapsto q_K(J\alpha,v)$ is bounded
in that norm by form Cauchy--Schwarz. Riesz representation
gives its minimizer. Nonnegativity follows from $K\succeq0$,
and $\mathfrak s_L(\rho p)=0$. The high inverse in this
operation is now paid by a prescribed $\Gamma_L^{-1}$;
no conditional $\widehat\eta^{-1}$ is required.

The off-diagonal return has \emph{not} thereby been bounded.
In particular the special rank-one native triangular
identity of V94 must not be transferred to this larger
projection. For exterior $D_i=X_i-s_i$,
\[
 [D_i,\Pi]=X_i\Pi-[s_i,\Pi].
 \tag{V98.25}\label{r98:covariant-warning}
\]
The second commutator, and the one-link Hamiltonian
off-diagonal block, are not estimated merely by bounding
$X_iJ$. Nor does \eqref{r98:Schur} at zero energy identify
the entire excited spectrum without the corresponding
energy-dependent reduction.

Thus V98 pays a smooth fixed-rank native cutoff, a coercive
\emph{full} discarded block, and a cutoff-decaying derivative
budget. The next upstream object is the actual retained
neutral matrix dynamics and its coupling back through that
high block. A useful next estimate must involve these
coefficients or their physical source overlap; repeating
projection algebra or shrinking only the geometric error
would not resolve it. Growing-region coercivity, physical
weak-coupling control and construction of the interacting
four-dimensional continuum remain unpaid. The full YM
mass gap is not proved.

During this installment the user supplied four newer companion
manuscripts. Their terminal developments and relevant scope
statements were reviewed before finalization; the large inherited
source banks and all earlier proofs were not independently
recertified. They lead to the following transfer decisions.

\paragraph{Perturbative ESM V64: assemble the same action and its marks.}
Its common order-one counterterm contributes the single next-order
insertion $\tfrac12\operatorname{STr}(G D_\Phi^2 C_1)$.
The six second-mark terms and the local Hessian compatibility
test prevent separately chosen subgraph tensors from being
mistaken for derivatives of one action. Its evaluated
lower-reference scalar-kinetic response also shows why a
field-independent subtraction does not remove a marked
coefficient. This supports retaining the \emph{whole} expression
\eqref{r98:K-compression}, including its reference centering,
when differentiating a cutoff reduction. It supplies no YM
counterterm or mass estimate: the concrete example has four
massless scalars with internal gauge, Yukawa and self-interaction
couplings zero, and its analytic bounds are at fixed source
and loop order. No common critical normalization is filled in
by that companion's compatibility test.

\paragraph{Native stationarity V54: estimate the actual source response.}
Its fixed-coefficient descriptor contains a rank-two map of a
differentiated source. The prepared lift controls a prescribed
$C^3$ covector through both fast scales, rather than charging
only its amplitude by the full mass inverse. The nonlinear
original remainder is explicitly still unpaid, and growing-mode
terminal preparation is not causal evolution or vacuum selection.
The useful transfer is to calculate the \emph{actual}
off-diagonal covector
$v\mapsto q_K(J(\rho h),v)$ in \eqref{r98:Schur} before estimating
its response. Its differentiated-coordinate coefficients,
preparation conditions and finite-dimensional constants do not
transfer to the YM Hamiltonian.

\paragraph{Exact-action Einstein V45: keep exact neutral modes and scales.}
Its fixed-$N=3$ weak-rank certificate resolves two genuine
central Jordan chains. A positive intermediate pencil then
has rank four and two leading oscillatory pairs; the remaining
next-order skew coefficient and moving nonlinear source are
not set to zero. This reinforces resolving the native retained
matrix and its vacuum before claiming a uniform gap.
The companion's amplitude scale is not lattice refinement;
its shrinking selected lifespan and frozen canonical
frequencies are not physical YM spectral lower bounds.

\paragraph{SM source selection V52: retain the actual admissible object.}
Its one-query optima are proved for specified invariant
preparation--steering classes. The faithful spanning-bank
example distinguishes signed reconstruction from executable
pure preparation, while a retained compensating reference
changes exact access. These are finite quantum-source results,
not a continuum mass-gap theorem. Here $J$ and $\Pi$ are
mathematical decompositions of the unchanged native state,
not a filtering operation or permission to choose a more
convenient vacuum. They depend on the actual $\zeta$; their
existence is not an algorithm for obtaining that vacuum.
No finite-range exterior dependence has been proved for them.

Accordingly the direction is sharpened, not replaced by a
companion's auxiliary dynamics: evaluate the native retained
matrix and its complete high-return source, check their
vacuum cancellation at $h=p$, and then seek a strict
physical coercivity estimate uniform in growing regions.
The small moving-isometry budget alone estimates only one
part of that source. This additional attachment review does
not replace the scheduled V100 folder and literature audits.

Verification includes exact rational projection, connection,
Hilbert--Schmidt, vacuum-motion and kinetic-compression
identities, including the zero-discarded-tail case.
Five native isolated four-link Wilson-loop truncations test
the high-block floor and moving-cutoff budget; they are not
full cubic-lattice proofs and have no certified truncation
error. An engineered exact ground
$\Omega\propto e^{t x_0y_0}$ on two spheres independently
checks \eqref{r98:K-compression} against the full positive
ground form for complex coefficients and $h=p$. Its
potential is $\kappa\Delta\Omega/\Omega$, not the Wilson
potential. These diagnostics verify identities; the uniform
statements are justified by the proofs above.
V97's mathematical diagnostics are replayed.

The complete V97 body is preserved except its version title,
and is exactly recoverable. Only active V98 is kept in
\texttt{latex\_notes}; build and diagnostic files stay outside.
The companion review, broad literature checkpoint and
archive/restart remain due at V100.
% END CONSOLIDATED V98 APPEND

% BEGIN CONSOLIDATED V99 APPEND
\clearpage
\section{V99: Incident-star Schur reduction with a bounded native source}
\label{r99:section}

\subsection{Go upstream of the moving-cutoff derivative cost}

V98 preserves the vacuum inside a moving harmonic space. Its
geometric estimates are valid, but that particular requirement
makes its complete off-diagonal source harder to estimate.
The new calculation goes upstream: a fixed free harmonic
projection commutes with \emph{every kinetic term}. An
incident-plaquette deletion compares the actual ground energy
with the exterior ground energy before either is charged as
an extensive constant. This gives a volume-independent high
block and a bounded off-diagonal source in the original
Hamiltonian. The exact vacuum is reconstructed by its Schur
graph rather than inserted into the cutoff space.

The archived f1 V3 Casimir theorem already gave a high block,
bounded magnetic source and Schur reduction using the
\emph{total} magnetic norm. We do not claim that algebra again.
The new input is the local energy comparison and the exact
incident-star harmonic source, with constants independent of
the number of plaquettes outside the selected link. The f15
V91 electric-window theorem instead concerned the normalized
compact transfer; it is not this continuous-time Hamiltonian
estimate. Current V91's potential deletion and V96's Haar
trial are inherited ingredients used together here.
Standard Feshbach/spectral-discretization context is supplied
by the primary abstracts of G.~Dusson, I.~M.~Sigal and
B.~Stamm,
\href{https://arxiv.org/abs/2105.02058}{arXiv:2105.02058}
and \href{https://arxiv.org/abs/2008.10871}{arXiv:2008.10871}.
No numerical or YM bound is imported from those abstracts.

\subsection{Cancel the exterior ground energy before taking the cutoff}

Work with the same finite cubic $\mathrm{SU}(2)$ Hamiltonian,
$N\ge3$, $\kappa>0$, $b\ge0$, on full product Haar space.
Its positive normalized ground is $\Omega$, its ground energy
is $E_0$, and $K=H-E_0\ge0$. For one chosen link $q$ write
the exact incident-star decomposition as
\[
 H=H_{\rm out}+\kappa L_q+4b+U_e,\qquad
 U_e(q,y)=-f(y)\cdot q,\qquad |f(y)|\le4b.
 \tag{V99.1}\label{r99:split}
\]
Thus $4b+U_e=W_e\ge0$ and $\int U_e(q,y)dq=0$.
The deleted operator $H_{\rm out}$ acts only on the exterior;
its ground energy is denoted $E_{\rm out}$.

\begin{lemma}[Incident ground-energy comparison]
\label{r99:ground-lemma}
At every volume and coupling,
\[
 E_{\rm out}\le E_0\le E_{\rm out}+4b,\qquad
 D_e:=H_{\rm out}+4b-E_0\ge0.
 \tag{V99.2}\label{r99:local-energy}
\]
Consequently the exact native operator has the decomposition
\[
                  K=D_e+\kappa L_q+U_e.
 \tag{V99.3}\label{r99:positive-exterior}
\]
The same statements hold on the physical subspace.
\end{lemma}

\begin{proof}
The lower comparison follows from
$H\ge H_{\rm out}+\kappa L_q\ge E_{\rm out}$.
For the upper one, use the tensor product of the constant
normalized link state and the normalized positive ground of
$H_{\rm out}$. The electric link energy and the Haar average
of $U_e$ vanish, leaving $E_{\rm out}+4b$.
The exterior ground is gauge invariant by positivity and
uniqueness on the finite connected compact configuration
space; its extension independent of $q$ is an admissible
physical trial. The full ground has the same invariance.
Finally $D_e\ge E_{\rm out}+4b-E_0\ge0$.
This deletion is an estimate inside the proof, not a change
of the native action or its vacuum.
\end{proof}

Let $P_\ell$ now denote the fixed one-link projection onto
harmonic degree $\ell$, tensored with exterior identity, and put
\[
 F_L=\sum_{\ell=0}^L P_\ell,\quad Q_L=I-F_L,\quad
 \mu_L=(L+1)(L+3),\quad
 d_L=\kappa\mu_L-4b.
 \tag{V99.4}\label{r99:cutoff}
\]
Unlike V98's $\Pi$, $F_L$ is independent of $y$ and is not
asserted to contain $\Omega$. Its fibre rank is
$D_L=(L+1)(L+2)(2L+3)/6$; its full range still contains
all exterior functions.

\begin{theorem}[Volume-independent native high block]
\label{r99:high-theorem}
If $d_L>0$, the actual complementary block satisfies
\[
 C_L:=Q_LKQ_L\ge d_LI_{\operatorname{Ran}Q_L}.
 \tag{V99.5}\label{r99:high}
\]
Its mixed block is exactly
\[
 B_L:=Q_LKF_L=Q_LU_eF_L.
 \tag{V99.6}\label{r99:mixed}
\]
Both statements restrict to the gauge-invariant space.
For a prescribed $\gamma>0$, it suffices to take
$\kappa\mu_L\ge4b+\gamma$; the required fibre rank is
\[
 D_L=O\!\left((1+\sqrt{b/\kappa+\gamma/\kappa})^3\right),
 \tag{V99.7}\label{r99:rank}
\]
with no volume factor.
\end{theorem}

\begin{proof}
In \eqref{r99:positive-exterior}, $D_e$ is nonnegative
and commutes with the link projections; on $Q_L$ the
electric part is at least $\kappa\mu_L$, while $U_e\ge-4b$.
This proves the closed-form bound. Every term except $U_e$
commutes with $F_L$, giving \eqref{r99:mixed}.
More concretely $D_e+\kappa L_q$ is a nonnegative
self-adjoint sum on separate factors and reduces $F_L$.
The perturbation $U_e$ is bounded. Its compressed diagonal
blocks therefore define self-adjoint operators on the
corresponding unperturbed domains, and its off-diagonal
blocks are bounded. This also justifies the operator block
notation in \eqref{r99:high}.
The Casimir projections commute with both endpoint actions;
all terms preserve the full gauge-invariant space.
The rank estimate follows from the free sphere multiplicities.
\end{proof}

For example at $\gamma=\kappa$ and $b/\kappa=1,10,100$,
sufficient choices are $L=1,5,19$, with fibre ranks
$5,91,2870$. These are ranks for a different cutoff from
V98, with a Schur-reconstructed rather than exactly retained
vacuum. They are not dimensions of the whole physical theory.

\subsection{The entire mixed source lies at one harmonic boundary}

\begin{lemma}[Exact source shell and its norm]
\label{r99:shell-lemma}
Multiplication by a coordinate of $S^3$ maps degree $\ell$
only to degrees $\ell-1$ and $\ell+1$. Moreover, for every
$v\in\mathbb R^4$,
\[
 \|P_{\ell+1}(v\cdot q)P_\ell\|=\frac{|v|}{2}.
 \tag{V99.8}\label{r99:raising}
\]
Consequently the native source obeys the sharper bound
\[
 \boxed{\quad
 B_L=P_{L+1}U_eP_L,\qquad \|B_L\|\le2b.
 \quad}
 \tag{V99.9}\label{r99:boundary-source}
\]
Here $P_L$ on the right is the single top-degree shell,
not the cumulative projection $F_L$.
\end{lemma}

\begin{proof}
If $p_\ell$ is a homogeneous harmonic polynomial on
$\mathbb R^4$, then
\[
 x_jp_\ell=
 \left(x_jp_\ell-\frac{|x|^2}{2\ell+2}\partial_jp_\ell\right)
       +\frac{|x|^2}{2\ell+2}\partial_jp_\ell.
\]
The two terms restrict to harmonic degrees $\ell+1$ and
$\ell-1$: $\partial_jp_\ell$ is harmonic, and applying the
Euclidean Laplacian verifies that the first parenthesis is
harmonic. For $\ell=0$ the second term is zero.
This proves the selection rule without discarding
angular sectors.

Rotate $v$ to the $q_0=x$ axis; the free degree projections
are rotation invariant. In the angular degree $m$ on the
remaining $S^2$, $0\le m\le\ell$, the normalized radial
factor is proportional to
$(1-x^2)^{m/2}C_{\ell-m}^{m+1}(x)$.
The normalized raising coefficient for multiplication by $x$
has squared value
\[
 c_{\ell,m}^2
 =\frac{(\ell-m+1)(\ell+m+2)}
        {4(\ell+1)(\ell+2)}
 =\frac14\left(1-\frac{m(m+1)}{(\ell+1)(\ell+2)}\right).
 \tag{V99.10}\label{r99:Gegenbauer}
\]
To check the normalization directly, the Gegenbauer
recurrence has upward coefficient
$(k+1)/(2(k+\alpha))$ and downward coefficient
$(k+2\alpha-1)/(2(k+\alpha))$.
Self-adjointness of multiplication by $x$ in its
orthogonality weight makes the normalized upward square
the product of the upward coefficient at $k$ and the
downward one at $k+1$. Insert $k=\ell-m$, $\alpha=m+1$.
These classical recurrence conventions are recorded in
\href{https://dlmf.nist.gov/18.9}{NIST DLMF, Section 18.9}.
The angular factors and these radial factors exhaust degree
$\ell$, since their dimensions sum to
$\sum_{m=0}^{\ell}(2m+1)=(\ell+1)^2$.
Thus all raising singular values are at most $1/2$, and
$m=0$ attains $1/2$, proving \eqref{r99:raising}.

Only the degree-$L$ input can leave $F_L$ under $U_e$.
Apply \eqref{r99:raising} pointwise in $y$ to $f(y)$ and
use $|f(y)|\le4b$. No derivatives of its rotating direction
are taken. Integration gives the full-space source bound.
\end{proof}

\subsection{A paid Schur return and the exact vacuum graph}

Write $A_L=F_LKF_L$. For real $z<d_L$ define
\[
 \begin{aligned}
 \Sigma_L(z)&=B_L^*(C_L-z)^{-1}B_L,\\
 S_L(z)&=A_L-z-\Sigma_L(z).
 \end{aligned}
 \tag{V99.11}\label{r99:Schur}
\]
This acts on the fixed finite-fibre exterior space.
The elementary block elimination already used in the
archives now has fully paid native constants:
\[
 \begin{gathered}
 0\le\Sigma_L(z)=P_L\Sigma_L(z)P_L,\qquad
 \|\Sigma_L(z)\|\le\frac{4b^2}{d_L-z},\\
 \|(C_L-z)^{-1}B_L\|\le\frac{2b}{d_L-z},\\
 \|\Sigma_L(z)-\Sigma_L(w)\|
 \le\frac{4b^2|z-w|}{(d_L-z)(d_L-w)}
 \quad(z,w<d_L).
 \end{gathered}
 \tag{V99.12}\label{r99:return-bounds}
\]
The first two estimates follow from
\eqref{r99:high} and \eqref{r99:boundary-source};
the last is the resolvent identity. In particular these
are operator-norm bounds on the actual return, not
higher-derivative amplitude estimates.

For completeness, if $h\in\operatorname{Dom}A_L$, then
\[
 (K-z)\psi=0
 \quad\Longleftrightarrow\quad
 S_L(z)h=0,\qquad
 \psi=h-(C_L-z)^{-1}B_Lh .
 \tag{V99.13}\label{r99:lift}
\]
The two components on the right occupy orthogonal
subspaces. Solve the high block equation and substitute
into the low equation; boundedness of $B_L$ and the
resolvent justifies both steps on the operator domains.
This correspondence preserves kernel multiplicity and
gauge invariance. It does not claim that $S_L(0)$ alone
has the same nonzero eigenvalues as $K$.

The exact vacuum head $h_0=F_L\Omega$ is nonzero and obeys
\[
 \begin{gathered}
 S_L(0)\ge0,\qquad \ker S_L(0)=\mathbb C h_0,\qquad
 Q_L\Omega=-C_L^{-1}B_Lh_0,\\
 \|Q_L\Omega\|\le
                 \frac{2b}{\sqrt{d_L^2+4b^2}} .
 \end{gathered}
 \tag{V99.14}\label{r99:vacuum}
\]
Indeed $C_L>0$ precludes a zero-energy state entirely
in $Q_L$; completion of the high square proves
$S_L(0)\ge0$; and \eqref{r99:lift} transfers the native
simple ground kernel. The router bound gives
$\|Q_L\Omega\|\le(2b/d_L)\|F_L\Omega\|$.
Use $\|\Omega\|=1$ and orthogonality to get the last line.
This is a global Haar norm estimate, equivalently an
averaged conditional tail. It is not a pointwise
lower bound for $\|F_L\zeta(\cdot,y)\|$.

The norm of the zero-energy lifted graph is governed by
\[
 G_L=I+B_L^*C_L^{-2}B_L,\qquad
 I\le G_L\le\left(1+\frac{4b^2}{d_L^2}\right)I.
 \tag{V99.15}\label{r99:graph-metric}
\]
For $\psi=h-C_L^{-1}B_Lh$, one has
$\|\psi\|^2=\langle h,G_Lh\rangle$ and
$\langle\Omega,\psi\rangle=\langle h_0,G_Lh\rangle$.
Thus even on the minimizing graph the physical vacuum
must be removed in this metric, not by subtracting
the constant coefficient. General vectors also have
their high innovation; the graph identity does not
discard that component.

\subsection{A convergent operator tail recursion with its error paid}

There is additional structure beyond a bounded mixed block.
In degree coordinates the exact operator
\eqref{r99:positive-exterior} is block tridiagonal:
its diagonal at degree $\ell$ is
$D_e+\kappa\ell(\ell+2)$, and its upward block is
\[
                   U_\ell=P_{\ell+1}U_eP_\ell.
 \tag{V99.16}\label{r99:ladder}
\]
There is no same-degree part of $U_e$.
All blocks still act on the exterior Hilbert space.
For $z<d_L$, elimination inside $C_L-z$ gives
\[
 \Sigma_L(z)|_{P_L}
 =U_L^*\left(D_e+\kappa\mu_L-z
          -\Sigma_{L+1}(z)|_{P_{L+1}}\right)^{-1}U_L .
 \tag{V99.17}\label{r99:recursion}
\]
This is an operator recursion: $D_e$ is not commuted
through $U_L$, and the exterior is not frozen.

\begin{theorem}[Uniform truncation error for the native high return]
\label{r99:recursion-theorem}
For $R\ge L$, set $\Sigma_R^{[R]}(z)=0$ and evaluate
\eqref{r99:recursion} downward to define
$\Sigma_L^{[R]}(z)$. It retains high degrees
$L+1,\ldots,R$, with the exterior unchanged. Then
\[
 0\le\Sigma_L^{[R]}(z)
       \le\Sigma_L^{[R+1]}(z)\le\Sigma_L(z),
 \tag{V99.18}\label{r99:recursion-order}
\]
and
\[
 \boxed{\quad
 \|\Sigma_L(z)-\Sigma_L^{[R]}(z)\|
 \le \frac{4b^2}{d_R-z}
          \prod_{\ell=L}^{R-1}
                      \frac{4b^2}{(d_\ell-z)^2}.
 \quad}
 \tag{V99.19}\label{r99:recursion-error}
\]
The empty product is one. At fixed coupling, fixed $L$
with $d_L>z$, and fixed $z$, this tends to zero, uniformly
in lattice volume.
\end{theorem}

\begin{proof}
Split the high space first into degree $L+1$ and its
higher complement. Its diagonal and mixed blocks are
exactly those in \eqref{r99:ladder}, so block elimination
proves \eqref{r99:recursion}. Every high tail, including
its finite-degree compression, has floor $d_\ell-z$
when its first degree is $\ell+1$.
The corresponding first-shell Schur operator has the
same lower bound: minimize its positive high-tail form,
whose lower bound contains $(d_\ell-z)$ times the
first-shell norm. Thus both inverses in any comparison
step have norm at most $(d_\ell-z)^{-1}$.

Starting with a zero terminal return, positivity and
inverse order prove the inequalities in
\eqref{r99:recursion-order} by downward induction.
The resolvent identity at one step bounds the difference
by
\[
 \frac{\|U_\ell\|^2}{(d_\ell-z)^2}
       \|\Sigma_{\ell+1}(z)-\Sigma_{\ell+1}^{[R]}(z)\|.
 \]
At the terminal shell use
$\|\Sigma_R(z)\|\le4b^2/(d_R-z)$ and
$\|U_\ell\|\le2b$. Iterating gives
\eqref{r99:recursion-error}. Since
$d_\ell\sim\kappa\ell^2$, the product tends to zero
faster than any fixed inverse power of $R$.
All constants used are incident-link constants.
For $b=0$ all mixed blocks and returns vanish identically.
\end{proof}

This is a certified operator approximation to the
\emph{one-link high return}. It is not a finite matrix
approximation to the whole theory: the exact, coupled
exterior operator remains in every denominator.

\subsection{Factorial electric tails for the entire low-energy window}

The bounded source also controls a full low-energy
spectral subspace, without selecting individual eigenvectors.
Fix $E_*\ge0$, put
$\mathcal P_*=1_{[0,E_*]}(K)$, and define
\[
 T_n=\sum_{\ell\ge n}P_\ell,\qquad
 t_n=\|T_n\mathcal P_*\|,\qquad t_0\le1 .
 \tag{V99.20}\label{r99:window-def}
\]
The same notation can be restricted throughout to the
physical spectral subspace.

\begin{theorem}[Native spectral-window tail recurrence]
\label{r99:window-theorem}
Whenever $\kappa n(n+2)>4b+E_*$,
\[
 \boxed{\quad
 t_n\le
 \frac{2b}{\kappa n(n+2)-4b-E_*}
            \|P_{n-1}\mathcal P_*\|
 \le q_n t_{n-1},\quad
 q_n=\frac{2b}{\kappa n(n+2)-4b-E_*}.
 \quad}
 \tag{V99.21}\label{r99:tail-recurrence}
\]
Choose $n_0\ge1$ with
$\kappa n_0(n_0+2)\ge2(4b+E_*)$. Then for $n\ge n_0$,
\[
 t_n\le\prod_{j=n_0}^n q_j
 \le\left(\frac{4b}{\kappa}\right)^{n-n_0+1}
           \left(\frac{(n_0-1)!}{n!}\right)^2 ,
 \tag{V99.22}\label{r99:factorial}
\]
with the independent cap $t_n\le1$.
\end{theorem}

\begin{proof}
Let $Q=T_n$, $F=I-Q$, $C=QKQ$, $B=QKF$,
and let $K_*$ be $K$ restricted to $\operatorname{Ran}\mathcal P_*$.
The operator $X=Q\mathcal P_*$ solves the exact equation
\[
 CX-XK_*=-BF\mathcal P_*.
 \]
The high floor is $d=\kappa n(n+2)-4b>E_*$.
Its norm-convergent solution is
\[
 Q\mathcal P_*=
 -\int_0^\infty e^{-tC} B F\mathcal P_* e^{tK_*}\,dt.
 \tag{V99.23}\label{r99:Sylvester}
\]
Differentiate the integrand involving $X$ to verify this
identity; the terminal term vanishes because
$\|e^{-tC}\|\|e^{tK_*}\|\le e^{-t(d-E_*)}$.
At each finite regulator the spectral window consists
of smooth vectors and has bounded energy, so this
calculation is valid on the operator domains. The
norm estimate also gives its bounded-operator extension.

The exact source is $B=P_nU_eP_{n-1}$.
Use \eqref{r99:raising} in the integrand, followed by
the norm of $P_{n-1}\mathcal P_*$ and the exponential
bound. Integration gives the first inequality in
\eqref{r99:tail-recurrence}. The second follows from
$P_{n-1}\le T_{n-1}$.
This proof is for the whole spectral window; replacing
it by estimates on each eigenvector separately would
not justify the operator-norm conclusion.

For $j\ge n_0$, the denominator of $q_j$ is at least
$\kappa j(j+2)/2$. Hence
$q_j\le4b/(\kappa j^2)$. Iteration from $t_{n_0-1}\le1$
proves \eqref{r99:factorial}. If $b=0$, the vanishing
mixed source directly gives zero tails under the
stated energy separation, with the same interpretation.
\end{proof}

For a finite set of all lattice links $\mathcal E_N$,
the commuting full product cutoff
$\mathcal F_L=\prod_{e\in\mathcal E_N}F_{e,L}$ satisfies
\[
 \|(I-\mathcal F_L)\mathcal P_*\|
       \le\sqrt{|\mathcal E_N|}\,t_{L+1},\qquad
 \operatorname{rank}\mathcal F_L\le D_L^{|\mathcal E_N|}.
 \tag{V99.24}\label{r99:global-window}
\]
Indeed $I-\prod_eF_{e,L}\le\sum_e(I-F_{e,L})$;
apply it to a vector in the spectral window and use
the identical local bound at each edge. Restriction
to physical states cannot increase the error or rank.
This is a state-window approximation, not a bounded
operator-norm approximation to the unbounded Hamiltonian.
The growing global rank and explicit volume factor
have not been suppressed.

\subsection{What this upstream change closes, and what it does not}

The native \emph{full-Hamiltonian} high block and its
complete return now have volume-independent operator
bounds with no conditional inverse gap, no unknown
vacuum derivatives and no loss of amplitude derivatives.
The error of a further high-degree truncation is also
paid. This is stronger for that particular obligation
than estimating the moving source through a higher
Sobolev norm. V98's smooth vacuum-compatible construction
is not invalidated, but is no longer the preferred
route for the high-block elimination.

The remaining object is $S_L(z)$ on the actual coupled
neutral exterior, with its native vacuum graph and
physical normalization. The estimates above do not
bound its first physical excitation. A small high
return does not make $A_L$ or its centered Schur head
uniformly coercive. In particular the retained head
still permits arbitrarily complicated exterior
collective modes; an incident-shell source is not a
spatially local effective interaction after inversion
of the exterior operator. Increasing $L$ alone would
repeat an ultraviolet approximation, not resolve
that infrared question.

Nor have we proved a positive conditional gap in
V94--V97's native one-link law by this different
projection. The new argument bypasses that inverse
on the full-Hamiltonian branch. It remains a theorem
about the stated lattice Hamiltonian regulator, not
an automatic identification with a normalized
finite-time Wilson transfer generator. Spatial
refinement, the physical weak-coupling trajectory,
nontrivial continuum construction, and the required
uniform neutral-sector estimate remain open.
The full interacting four-dimensional YM mass gap
is unproved.

Exact polynomial checks verify the harmonic selection
rule and the normalized raising coefficient. Ninety-one
Gegenbauer quadratures check all angular sectors
through degree twelve. Three finite matrix-exterior
models test the noncommuting tail recursion, graph
metric and whole-window identity; these are not
Wilson lattices. Five native isolated four-link
Wilson loops check the same source and return
identities in their physical class-function sectors.
Their finite character matrices are diagnostics,
not certified infinite-spectrum bounds and not the
coupled cubic model. V98's mathematical diagnostics
are replayed. The uniform statements rely on the
written operator proofs, not these finite tests.

V98's preceding body is preserved except the version
title and is exactly recoverable. Only active V99 is
kept in \texttt{latex\_notes}; all build and diagnostic
files are outside it. The next installment, V100,
is the scheduled companion review, broad literature
sweep and archive/restart boundary. That review
should target the retained interacting neutral
operator and physical continuum bridge, not another
optimization of the already paid electric tail.
% END CONSOLIDATED V99 APPEND

% BEGIN CONSOLIDATED V100 APPEND
\clearpage
\section{V100: Native electric resonances and a local averaging step}
\label{r100:section}

\subsection{The checkpoint changes the proposed import}

V99 pays the full high return for a fixed one-link harmonic
cutoff. The next question is spatial locality and retained
neutral dynamics, not another improvement of its ultraviolet
tail. The scheduled companion audit found all eight current
sources unchanged since their last relevant reviews. NS V26,
SM--Einstein V72, Shell Mixing V16 and parallel YM V5 were
reread at their terminal scope statements. The four Downloads
reviewed in V98 were fingerprinted and their conclusions
reread. Their finite-order, selected-source and finite-cutoff
qualifications remain; none supplies a nonperturbative YM
spectral estimate.

A broad primary-source sweep covered effective Hamiltonians,
quasi-local dynamics, finite-size gap criteria, lattice
functional inequalities, and recent lower-dimensional continuum
limits. The following is a transfer audit, not certification
of all results in the cited papers.
\begin{itemize}
\item Bravyi--DiVincenzo--Loss,
\href{https://arxiv.org/abs/1105.0675}{arXiv:1105.0675},
Sections 3.1, 3.5 and 4.1--4.4: spectral separation,
additivity and connected expansions are relevant. A
one-link high-block floor does not verify separation
of its entire retained band.
\item Nachtergaele--Sims--Young,
\href{https://arxiv.org/abs/1810.02428}{arXiv:1810.02428},
Section 3.2 and Theorem 3.3: unbounded on-site generators
with bounded local interactions fit their dynamical framework.
Such locality estimates do not manufacture a vacuum gap.
\item Lemm--Lucia,
\href{https://arxiv.org/abs/2409.09685}{arXiv:2409.09685v2}:
the finite-size result concerns frustration-free interactions.
The native centered Hamiltonian has not been put in that class.
\item Shen--Zhu--Zhu,
\href{https://arxiv.org/abs/2204.12737}{arXiv:2204.12737}
and
\href{https://arxiv.org/abs/2401.13299}{arXiv:2401.13299v3}:
strong-coupling lattice functional inequalities and
Yang--Mills--Higgs extensions do not provide a weak-coupling
four-dimensional continuum estimate in our normalization.
\item The 2026 primary records
\href{https://arxiv.org/abs/2602.08591v2}{Dang--Nohra},
\href{https://arxiv.org/abs/2608.27828v2}{Chevyrev--Qu--Shen},
and
\href{https://doi.org/10.1007/s00205-025-02163-3}{Chevyrev--Shen (2026)}
concern respectively surfaces, three-dimensional abelian
scaling, and three-dimensional stochastic gauge-covariant
renormalisation. Their dimension, group and dynamical scopes
are not replaced by four-dimensional nonabelian vacuum claims.
\end{itemize}
For the last three bullets the primary abstracts/records
were used to check scope, not to import quantitative theorems.
The new estimates below have their own proofs.
General Schur algebra, resonant averaging and connected
commutator expansions are established techniques; the native
resonant matrix element and the explicit application to the
untruncated cubic regulator are the present calculations.
Archive searches did not locate these calculations in the
existing notes; no universal priority claim is intended.

\subsection{A Schur correction need not be a local interaction}

Here is an exact test of what V99's norm estimate does
\emph{not} imply. It is an auxiliary finite model, not a
counterexample to Wilson Yang--Mills.
Take a two-level local positive operator and an independent,
arbitrarily distant two-level spectator:
\[
 k=\begin{pmatrix}r^2/d&r\\r&d\end{pmatrix},\qquad
 D=\varepsilon n,\quad n=\begin{pmatrix}0&0\\0&1\end{pmatrix},
 \qquad K=k\otimes I+I\otimes D,
 \tag{V100.1}\label{r100:spectator}
\]
where $d,\varepsilon>0$, $r\ne0$. The matrix $k$ has a simple
zero eigenvalue and second eigenvalue $d+r^2/d$.
Retain its first coordinate. Then $C=d+D\ge d$,
$B=rI$, but
\[
 \Sigma(z)=r^2(d+D-z)^{-1}.
 \tag{V100.2}\label{r100:spectator-return}
\]
For the spectator flip $X$,
\[
 \|[\Sigma(z),X]\|
 =\frac{r^2\varepsilon}{(d-z)(d+\varepsilon-z)}
 \quad(z<d).
 \tag{V100.3}\label{r100:spectator-commutator}
\]
Direct diagonal multiplication proves this formula. It
does not decrease when the spectator is moved farther away.
Thus even zero physical interaction between the two sites
does not make the fixed-energy Schur return a local
observable near the eliminated site.

There is no failure of spectral additivity. On the spectator
eigenvalue $t$, the full Schur pencil is
\[
 S(z)|_t=r^2/d+t-z-\frac{r^2}{d+t-z}
                   =S_{\rm local}(z-t).
 \tag{V100.4}\label{r100:spectator-shift}
\]
It vanishes at $z=t$. At $z=0$, however, its value at
$t=\varepsilon$ is
$\varepsilon[1+r^2/(d(d+\varepsilon))]$.
Freezing the energy dependence creates an apparent
spectator correction that is absent from the exact spectrum.
The same operator-valued spectral shift follows by the
spectral theorem for any nonnegative independent spectator.
This is why quasi-local Hamiltonian theorems cannot simply
be applied to $\Sigma_L(0)$.

\subsection{Exact physical resonances across every one-link cutoff}

Return to the native cubic $\mathrm{SU}(2)$ torus, $N\ge3$.
Write the uncentered Hamiltonian as
\[
 H=T+b|\mathcal P_N|+V,\qquad
 T=\kappa\sum_e L_e,\qquad
 V=\sum_p V_p,\quad V_p=-b w_p,\quad
 w_p=\tfrac12\Tr U_p .
 \tag{V100.5}\label{r100:native}
\]
The scalar constant will be restored unchanged after any
unitary conjugation.

Choose opposite vertices of one plaquette, with its two
length-two paths $a,b$ oriented from the first vertex to the
second. Add an exterior path $c$ of length four, internally
disjoint from both. One concrete choice connects
$(0,0,0)$ to $(1,1,0)$ through
$(0,0,1),(1,0,1),(1,1,1)$.
This theta graph embeds in every stated torus.

\begin{theorem}[A nonzero native resonant crossing]
\label{r100:resonance-theorem}
For every integer $\ell\ge0$ there are normalized, orthogonal,
gauge-invariant spin-network vectors $\Psi_\ell,\Psi'_\ell$
such that
\[
 \begin{gathered}
 T\Psi_\ell=E_\ell\Psi_\ell,\qquad
 T\Psi'_\ell=E_\ell\Psi'_\ell,\qquad
 E_\ell=\kappa(4\ell^2+12\ell+18),\\
 F_{e,\ell}\Psi_\ell=\Psi_\ell,\qquad
 Q_{e,\ell}\Psi'_\ell=\Psi'_\ell,\\
 \langle\Psi'_\ell,V\Psi_\ell\rangle
          =-\frac{b}{2(\ell+1)(\ell+2)} ,
 \end{gathered}
 \tag{V100.6}\label{r100:resonance}
\]
where $e$ is either edge of path $a$ and $F_{e,\ell}$
is V99's one-link cutoff.
\end{theorem}

\begin{proof}
Put spins $(j,j+\tfrac12,\tfrac12)$ on the three paths,
where $j=\ell/2$, and use the unique normalized invariant
three-valent tensor at both endpoints. All other edges
carry the trivial representation. The primed vector
interchanges the first two spins. Bivalent vertices along
a path simply compose its representation matrices.

For explicit normalization, let $g_a,g_b,g_c$ be the path
holonomies, $d_s=2s+1$, and
$t_{m_a m_b m_c}$ the real Wigner $3j$ invariant of norm one.
The unprimed function is
\[
 \sqrt{d_jd_{j+1/2}d_{1/2}}\!
 \sum_{\boldsymbol m,\boldsymbol n}
 t_{\boldsymbol m}t_{\boldsymbol n}
 D^j_{m_a n_a}(g_a)
 D^{j+1/2}_{m_b n_b}(g_b)
 D^{1/2}_{m_c n_c}(g_c).
 \tag{V100.7}\label{r100:theta-state}
\]
Independent product Haar on the edges pushes to independent
Haar on these paths. Representation orthogonality proves
norm one; different representation labels on path $a$
prove orthogonality of the two vectors. Endpoint
invariance gives the Gauss constraint, including at the
interior path vertices.

The Casimir degrees on $a,b,c$ are respectively
$\ell,\ell+1,1$. The path lengths give
$E_\ell/\kappa=2\ell(\ell+2)+2(\ell+1)(\ell+3)+12$,
unchanged by the interchange. This proves the electric
and cutoff assertions.

We evaluate the fundamental Wilson insertion rather
than infer a nonzero recoupling coefficient.
Since $U_p=g_a g_b^{-1}$,
\[
 \Tr(g_a g_b^{-1})
   =\sum_{\alpha,\beta}
       D^{1/2}_{\alpha\beta}(g_a)
       \overline{D^{1/2}_{\alpha\beta}(g_b)} .
\]
The triple Haar integrals reduce each endpoint to the
same real contraction
\[
 \begin{split}
 A_\ell=\sum&
 t^{j,j+1/2,1/2}_{m_a m_b m_c}
 t^{j+1/2,j,1/2}_{m'_a m'_b m_c}\\[-.5ex]
 &{}\times
 C^{j+1/2,m'_a}_{j,m_a;\,1/2,\alpha}
 C^{j+1/2,m_b}_{j,m'_b;\,1/2,\alpha}.
 \end{split}
 \tag{V100.8}\label{r100:vertex}
\]
The sum is over admissible indices, so
$m'_a=m_a+\alpha$, $m'_b=m_b-\alpha$.
Haar denominators and the two network normalizations
leave the factor $(\ell+1)/(\ell+2)$.

Here is an evaluation valid for all $\ell$, not a finite
spin table. With $r=j+m_a$, use
\[
 \left|C^{j+1/2,m+\alpha}_{j,m;\,1/2,\alpha}\right|^2
       =\frac{j+2\alpha m+1}{2j+1}.
\]
The standard real $3j$ phases give a positive term when
$m_c=\alpha$ and a negative term when $m_c=-\alpha$.
The two equal-sign cases and the two opposite-sign cases
therefore give
\[
 A_\ell=
 \frac{2\sum_{r=0}^{\ell-1}(r+1)(\ell-r)
       -2\sum_{r=0}^{\ell}(r+1)(\ell-r+1)}
      {(\ell+1)^2(\ell+2)}
       =-\frac1{\ell+1}.
 \tag{V100.9}\label{r100:vertex-sum}
\]
The first sum is empty at $\ell=0$. The sums evaluate to
$\ell(\ell+1)(\ell+2)/6$ and
$(\ell+1)(\ell+2)(\ell+3)/6$.
Thus the normalized $w_p$ matrix element is
\[
 \frac12\,\frac{\ell+1}{\ell+2}\,A_\ell^2
              =\frac1{2(\ell+1)(\ell+2)}.
\]
A different plaquette omits at least one of the four
edges on which the two vectors have different degrees.
Integration over that edge makes its matrix element zero.
The scalar part also has zero matrix element.
Consequently the full $V$ has the value in
\eqref{r100:resonance}.
\end{proof}

For example the $w_p$ values at $\ell=0,1,2$ are
$1/4,1/12,1/24$. The first is independently
$\int (2a_0)(a\cdot b)(2b_0)\,da\,db=1/4$.
This uses actual physical states of the full regulator;
it is not the isolated one-loop model.

\begin{corollary}[No whole-head electric homological inverse]
\label{r100:no-inverse}
For $b>0$, no operator $X$ with the requisite matrix
elements and electric domains can solve
\[
 (Q_{e,\ell}TQ_{e,\ell})X
       -X(F_{e,\ell}TF_{e,\ell})=-Q_{e,\ell}VF_{e,\ell}
 \tag{V100.10}\label{r100:homological-failure}
\]
on the whole physical retained space.
\end{corollary}
\begin{proof}
Take its matrix element between $\Psi'_\ell$ and $\Psi_\ell$.
The left side is $(E_\ell-E_\ell)
\langle\Psi'_\ell,X\Psi_\ell\rangle=0$;
the right side is nonzero by \eqref{r100:resonance}.
\end{proof}
In particular a near-identity differentiable block rotation
about $b=0$, with a well-defined first homological equation,
cannot eliminate the entire $F_{e,\ell}$--$Q_{e,\ell}$ source.
Taking $\ell$ large enough to make V99's $d_\ell>0$
does not remove this resonance. The resonant total energies
are high, so this does not contradict V99's Sylvester
inverse on a \emph{total} low-energy window below $d_\ell$.
It identifies the extra hypothesis missing from a
whole-band Schrieffer--Wolff import.

\subsection{Keep the resonance: a bounded, strictly local averaging map}

There is a constructive alternative to dividing by zero.
The electric spectrum is contained in $\kappa\mathbb N_0$,
so its unitary dynamics has period
$\tau=2\pi/\kappa$. For a bounded local operator $A$ define
strong operator integrals
\[
 \begin{aligned}
 {\cal R}(A)&=\frac1\tau\int_0^\tau e^{itT}Ae^{-itT}\,dt,\\
 {\cal J}(A)&=\frac{i}{\tau}\int_0^\tau
             (t-\tau/2)e^{itT}Ae^{-itT}\,dt .
 \end{aligned}
 \tag{V100.11}\label{r100:averages}
\]
These are not averages under the unknown interacting
vacuum. They are explicit maps on the original operator.

\begin{lemma}[Local resonant decomposition]
\label{r100:average-lemma}
For $R_p={\cal R}(V_p)$ and $S_p={\cal J}(V_p)$,
\[
 \begin{gathered}
 R_p^*=R_p,\quad [T,R_p]=0,\quad S_p^*=-S_p,\\
 [T,S_p]=V_p-R_p,\qquad
 \|R_p\|\le b,\qquad
 \|S_p\|\le s:=\frac{\pi b}{2\kappa}.
 \end{gathered}
 \tag{V100.12}\label{r100:average-bounds}
\]
Both operators are supported on the same plaquette as
$V_p$ and commute with every gauge transformation.
The commutator identity holds on $\operatorname{Dom}T$
with a bounded right side.
\end{lemma}

\begin{proof}
The electric evolution factors over links. It changes
neither the support nor the gauge invariance of $V_p$.
Its norm is $b$, and
$\tau^{-1}\int_0^\tau |t-\tau/2|dt=\tau/4$ gives the
stated norms. Periodicity gives $[T,R_p]=0$.
Integration by parts, using
$[T,e^{itT}V_pe^{-itT}]=-i\,\partial_t(e^{itT}V_pe^{-itT})$,
gives $[T,S_p]=V_p-R_p$.

More explicitly, between electric eigenvectors with
energies $E_m,E_n$, \({\cal R}\) keeps just $E_m=E_n$,
while \({\cal J}\) gives $A_{mn}/(E_m-E_n)$ when
$E_m\ne E_n$ and zero otherwise. Each $V_p$ maps the
finite Peter--Weyl core into itself, as do these two
maps. On that core the commutator is bounded. Approximate
a vector in $\operatorname{Dom}T$ in graph norm by core
vectors and use closedness of $T$ to extend the identity.
This proves domain preservation. Strong integrals suffice;
norm continuity of the unbounded electric dynamics is
not being assumed.
\end{proof}

The native theta matrix element survives unchanged in
$R=\sum_pR_p$, since its two energies are equal.
Thus $R$ is not a scalar vacuum counterterm.
Let $P_{0,p}$ project onto constants on the four links
of $p$. Then
\[
 R_pP_{0,p}=0,\qquad
 |\langle u,Ru\rangle|
      \le\frac{4b}{3\kappa}\langle u,Tu\rangle .
 \tag{V100.13}\label{r100:resonant-relative}
\]
Indeed $R_p$ commutes with the local electric energy,
whose zero space is $P_{0,p}$, and the Haar mean of
$w_p$ is zero. Hence
$-b(I-P_{0,p})\le R_p\le b(I-P_{0,p})$.
The local electric gap $3\kappa$ bounds $I-P_{0,p}$,
and each link belongs to four plaquettes. Summation
proves the form inequality.
For $4b<3\kappa$ it follows that $T+R$ has the constant
ground and a full-space gap at least $3\kappa-4b$.
This is a bound for the \emph{resonant comparison operator},
not for $H-E_0$.

\subsection{An exact conjugation with a spatially summable remainder}

At finite volume put $S=\sum_p S_p$.
It is bounded and anti-self-adjoint; by the bounded
commutator identity $e^S$ preserves $\operatorname{Dom}T$.
Consequently
\[
 e^S(H-b|\mathcal P_N|)e^{-S}=T+R+W
 \tag{V100.14}\label{r100:normal-form}
\]
is an exact unitary equivalence.
The following estimate has no electric truncation.

\begin{theorem}[One local resonance-preserving step]
\label{r100:local-normal-theorem}
Let
\[
 q=26s=13\pi b/\kappa<1 .
 \tag{V100.15}\label{r100:q}
\]
The bounded remainder admits a connected plaquette
expansion with rooted order-$n$ norm budget at most
$bq^n$, $n\ge1$. In particular,
\[
 \|W\|\le|\mathcal P_N|\,b\,\frac{q}{1-q}.
 \tag{V100.16}\label{r100:global-remainder}
\]
More importantly, assign each term to the union $Z$
of its contributing plaquettes. Its interaction
strength at every link obeys
\[
 \sup_e\sum_{Z\ni e}\|W_Z\|
   \le b\sum_{n\ge1}3(4n+5)^3q^n<\infty ,
 \tag{V100.17}\label{r100:local-remainder}
\]
uniformly in $N$. This is $O(b^2/\kappa)$ as $b/\kappa\to0$.
The interaction tail beyond $n$ commutator steps is
bounded by the corresponding tail of that convergent series.
\end{theorem}

\begin{proof}
Since $[S,T]=-(V-R)$, the exact bounded-commutator
expansion of the conjugation gives
\[
 W=\sum_{n\ge1}
       \frac{n\,\operatorname{ad}_S^n(V)
                +\operatorname{ad}_S^n(R)}{(n+1)!},
 \qquad \operatorname{ad}_S(A)=[S,A].
 \tag{V100.18}\label{r100:series}
\]
At finite volume this follows from the usual exponential
series, with the first commutator with $T$ bounded.
All remaining terms are bounded operators.

Fix the initial plaquette $p_0$. A nested commutator
vanishes whenever its new plaquette is disjoint from
the accumulated support. Each plaquette has at most
thirteen plaquettes in its closed overlap neighbourhood:
it has four links and each link has three other incident
plaquettes. After $k$ plaquettes the next choice has at
most $13k$ possibilities. Thus at order $n$ there are at
most $13^n n!$ nonzero ordered choices for each root.
Their individual norms are bounded by
$(2s)^n b$, for either initial $V_{p_0}$ or $R_{p_0}$.
The two coefficients in \eqref{r100:series} add to
$1/n!$. The root budget is therefore $b(26s)^n=bq^n$.
This proves absolute convergence and
\eqref{r100:global-remainder}.

For a fixed link $e$ in such a union, the root is
connected to an incident plaquette of $e$ by at most
$n$ overlap steps. In the cubic lattice its base vertex
lies in a coordinate box of side at most $4n+5$ about
$e$; there are three plaquette orientations. Therefore
at most $3(4n+5)^3$ roots can contribute. This is still
an upper bound on a periodic torus, by lifting a path
to the infinite cubic lattice. Sum their rooted norm
budgets. Grouping terms with the same union can only
decrease the bound, proving \eqref{r100:local-remainder}.
Every term is gauge invariant and supported on its
assigned connected union. The same argument applied
to the tail proves the final assertion.
\end{proof}

The global norm estimate is deliberately extensive;
the useful spatial estimate is
\eqref{r100:local-remainder}. An extensive error cannot
be subtracted from a vacuum excitation gap as if it
were a volume-independent spectral perturbation.

This step is different from V99's exact one-link
Schur pencil. It retains all equal-total-electric-energy
interactions and does not attempt to separate the
overlapping entire $F_L$ and $Q_L$ spectra.
It is also not a continuum result. The convergence
window $13\pi b/\kappa<1$ is a declared strong-electric
regime, not the physical weak-coupling trajectory.
One must still control the transformed interacting
vacuum, the resonant neutral dynamics and any iterative
renormalisation before obtaining a physical gap.
Neither the comparison gap of $T+R$ nor the unitary
identity with its nonzero $W$ closes that problem.

\subsection{Verification, archive, and the next calculation}

Exact $3j$/Clebsch--Gordan contractions for eleven spins
agree with the all-$\ell$ polynomial proof; the
$\ell=0$ Haar moment is an independent check. Three
finite qutrit tests verify the sign, commutator,
periodic energy selection and bounded-conjugation
remainder formula. They are not Wilson lattices.
Cubic incidence checks at $N=3,4,5$ verify the
thirteen-neighbour count. The spectator matrix is
checked exactly. These diagnostics do not replace
the domain and uniform connected-series proofs.
V99's mathematical diagnostics are replayed.

The checkpoint closes this consolidated cycle.
V1--V99 are preserved before this append except the
version title; the source is archived as
\path{Renewal_Yang_Mills_Mass_Gap_Research_Notes_V100_f17.tex}.
A fresh active V1 summarizes the major results and
continues with the connected second-order resonant
operator. All older archives and companion files
are retained. Build files stay outside
\texttt{latex\_notes}. The new cycle's regular companion
and broad literature checkpoints are V5 and V10.

The outstanding goal is unchanged: a nontrivial,
interacting four-dimensional continuum Yang--Mills
theory with a positive physical mass gap. It remains
unproved.
% END CONSOLIDATED V100 APPEND

\end{document}
